\documentclass[11pt,a4paper]{article}
\pdfextension mapfile{+cm.map}
\pdfextension mapfile{+cmextra.map}
\pdfextension mapfile{+symbols.map}
\pdfextension mapfile{+euler.map}
\pdfextension mapfile{+latxfont.map}
\pdfextension mapfile{+rsfs.map}
\usepackage[margin=25mm,headheight=15pt]{geometry}
\usepackage{fontspec}
\setmainfont{Latin Modern Roman}
\usepackage{amsmath,amssymb,amsthm,amscd,mathtools,mathrsfs}
\usepackage{microtype,booktabs,longtable,array,xcolor,graphicx,enumitem,needspace,fancyhdr}
\usepackage[round,authoryear]{natbib}
\usepackage[unicode,colorlinks=true,linkcolor=blue!40!black,citecolor=blue!40!black,urlcolor=blue!40!black]{hyperref}
\usepackage{xurl}
\usepackage{listings}
\usepackage{tabularx}
\lstset{basicstyle=\ttfamily\small,breaklines=true,columns=fullflexible,keepspaces=true}
\usepackage{tikz}
\usepackage{tikz-cd}
\usetikzlibrary{arrows.meta,calc,patterns,positioning}
\usepackage{aliascnt}
\usepackage[nameinlink,noabbrev]{cleveref}
\allowdisplaybreaks[2]
\setcounter{MaxMatrixCols}{32}
\setlength{\emergencystretch}{3em}
\setlength{\parskip}{2pt}
\newtheorem{theorem}{Theorem}[section]
\newaliascnt{proposition}{theorem}
\newtheorem{proposition}[proposition]{Proposition}
\aliascntresetthe{proposition}
\crefname{proposition}{proposition}{propositions}
\newaliascnt{lemma}{theorem}
\newtheorem{lemma}[lemma]{Lemma}
\aliascntresetthe{lemma}
\crefname{lemma}{lemma}{lemmas}
\newaliascnt{corollary}{theorem}
\newtheorem{corollary}[corollary]{Corollary}
\aliascntresetthe{corollary}
\crefname{corollary}{corollary}{corollarys}
\newaliascnt{claim}{theorem}
\newtheorem{claim}[claim]{Claim}
\aliascntresetthe{claim}
\crefname{claim}{claim}{claims}
\newaliascnt{statement}{theorem}
\newtheorem{statement}[statement]{Statement}
\aliascntresetthe{statement}
\crefname{statement}{statement}{statements}
\theoremstyle{definition}
\newaliascnt{definition}{theorem}
\newtheorem{definition}[definition]{Definition}
\aliascntresetthe{definition}
\crefname{definition}{definition}{definitions}
\newaliascnt{remark}{theorem}
\newtheorem{remark}[remark]{Remark}
\aliascntresetthe{remark}
\crefname{remark}{remark}{remarks}
\newaliascnt{example}{theorem}
\newtheorem{example}[example]{Example}
\aliascntresetthe{example}
\crefname{example}{example}{examples}
\newaliascnt{problem}{theorem}
\newtheorem{problem}[problem]{Open problem}
\aliascntresetthe{problem}
\crefname{problem}{open problem}{open problems}
\newaliascnt{question}{theorem}
\newtheorem{question}[question]{Question}
\aliascntresetthe{question}
\crefname{question}{question}{questions}
\newaliascnt{auditrecord}{theorem}
\newtheorem{auditrecord}[auditrecord]{Verification record}
\aliascntresetthe{auditrecord}
\crefname{auditrecord}{verification record}{verification records}
\newaliascnt{test}{theorem}
\newtheorem{test}[test]{Verification}
\aliascntresetthe{test}
\crefname{test}{verification}{verifications}
\newaliascnt{conjecture}{theorem}
\newtheorem{conjecture}[conjecture]{Conjecture}
\aliascntresetthe{conjecture}
\crefname{conjecture}{conjecture}{conjectures}
\newcommand{\C}{\mathbb C}\newcommand{\R}{\mathbb R}\newcommand{\Z}{\mathbb Z}\newcommand{\Q}{\mathbb Q}\newcommand{\N}{\mathbb N}\newcommand{\F}{\mathbb F}
\newcommand{\CC}{\mathbb C}\newcommand{\RR}{\mathbb R}\newcommand{\ZZ}{\mathbb Z}\newcommand{\QQ}{\mathbb Q}\newcommand{\NN}{\mathbb N}
\newcommand{\PP}{\mathbb P}\newcommand{\Pp}{\mathbb P}\newcommand{\Pone}{\mathbb P^1}\newcommand{\HH}{\mathbb H}
\newcommand{\RP}{\mathbb{RP}}\newcommand{\CP}{\mathbb{CP}}
\newcommand{\OO}{\mathcal O}\newcommand{\cO}{\mathcal O}\newcommand{\cM}{\mathcal M}\newcommand{\cF}{\mathcal F}
\newcommand{\h}{\mathfrak h}\newcommand{\hz}{\mathfrak h_z}\newcommand{\Acon}{\mathfrak A}\newcommand{\wt}[1]{\widehat{#1}}
\newcommand{\red}{\mathrm{red}}\newcommand{\tor}{\mathrm{tor}}\newcommand{\dP}{\mathrm{dP}}
\newcommand{\gh}{\widehat\gamma}\newcommand{\uh}{\widehat u}\newcommand{\wh}{\widehat w}\newcommand{\dhat}{\widehat\delta}\renewcommand{\dh}{\widehat\delta}
\newcommand{\dd}{\mathrm d}\newcommand{\ii}{\mathrm i}\newcommand{\eps}{\varepsilon}\newcommand{\partt}{\mathrm{part}}\newcommand{\Bcirc}{B^\circ}\newcommand{\ddbar}{\partial\bar\partial}
\DeclareMathOperator{\Pic}{Pic}\DeclareMathOperator{\Alb}{Alb}\DeclareMathOperator{\Aut}{Aut}\DeclareMathOperator{\NS}{NS}
\DeclareMathOperator{\Hom}{Hom}\DeclareMathOperator{\End}{End}\DeclareMathOperator{\Ext}{Ext}\DeclareMathOperator{\Tor}{Tor}
\DeclareMathOperator{\Spec}{Spec}\DeclareMathOperator{\Proj}{Proj}\DeclareMathOperator{\Sym}{Sym}\DeclareMathOperator{\Diff}{Diff}
\DeclareMathOperator{\GL}{GL}\DeclareMathOperator{\SL}{SL}\DeclareMathOperator{\PSL}{PSL}\DeclareMathOperator{\Sp}{Sp}
\DeclareMathOperator{\im}{im}\DeclareMathOperator{\coker}{coker}\DeclareMathOperator{\rank}{rank}\DeclareMathOperator{\diag}{diag}
\DeclareMathOperator{\ord}{ord}\DeclareMathOperator{\vol}{vol}\DeclareMathOperator{\pr}{pr}\DeclareMathOperator{\id}{id}
\DeclareMathOperator{\divisor}{div}\DeclareMathOperator{\Tr}{Tr}\DeclareMathOperator{\Nm}{Nm}\DeclareMathOperator{\Sha}{Sha}
\DeclareMathOperator{\Imn}{Im}\DeclareMathOperator{\Ren}{Re}\DeclareMathOperator{\trdeg}{tr.deg}
\newcommand{\Uses}[1]{\par\smallskip\noindent\textit{Uses:} #1.\par\smallskip}
\newcommand{\Verified}{}\newcommand{\AuthorSelfCheck}{}\newcommand{\NeedsCheck}{}\newcommand{\Open}{\textbf{Open.}\ }

\newcommand{\RungZero}[1]{\par\noindent\textbf{Result and mechanism.} #1\par\smallskip}
\newcommand{\RungOne}{\par\Needspace{6\baselineskip}\noindent\textbf{Proof at the first rung.}\par\nobreak\smallskip}
\newcommand{\CurrentChecker}{GPT-6 Astra (OpenAI, ChatGPT Pro mode, effort Pro), 18 November 2026}
\newaliascnt{cited}{theorem}
\newtheorem{cited}[cited]{Cited result}
\aliascntresetthe{cited}
\crefname{cited}{cited result}{cited results}
\newaliascnt{refuted}{theorem}
\newtheorem{refuted}[refuted]{Correction}
\aliascntresetthe{refuted}
\crefname{refuted}{correction}{corrections}
\newaliascnt{computed}{theorem}
\newtheorem{computed}[computed]{Computation}
\aliascntresetthe{computed}
\crefname{computed}{computation}{computations}
\newaliascnt{openq}{theorem}
\newtheorem{openq}[openq]{Open question}
\aliascntresetthe{openq}
\crefname{openq}{open question}{open questions}

\providecommand{\rk}{\operatorname{rk}}
\providecommand{\Ann}{\operatorname{Ann}}
\providecommand{\Log}{\operatorname{Log}}
\providecommand{\vir}{\mathrm{vir}}

\usepackage{etoolbox}
\setmonofont[BoldFont=lmmonolt10-bold.otf,ItalicFont=lmmono10-italic.otf]{lmmono10-regular.otf}
\hypersetup{bookmarksnumbered=true,bookmarksdepth=2}
\setcounter{tocdepth}{2}
\makeatletter\renewcommand{\@pnumwidth}{2.5em}\renewcommand{\@tocrmarg}{3.5em}\renewcommand*{\l@subsection}{\@dottedtocline{2}{1.5em}{3.5em}}\makeatother
\newtheorem{maintheorem}{Theorem}
\renewcommand{\themaintheorem}{\Alph{maintheorem}}
\crefname{maintheorem}{theorem}{theorems}
\newcommand{\StatementRecord}[4]{\par\smallskip\begingroup\upshape\raggedright\noindent
\textbf{\texttt{[#3]} #2}\quad
\href{S6_SERIES_GRIND_checks_\CurrentPaper.md\#tag-#3}{checks}\quad
\href{S6_SERIES_GRIND_checks_\CurrentPaper.md\#uses-#3}{uses and used by}\par
\noindent\textit{uses:} #4.\par
\ifcsname SBack@#1\endcsname\noindent\csname SBack@#1\endcsname\par\fi
\endgroup\smallskip}
\newcommand{\ExternalPaperRef}[3]{#1#2 (Paper~#3)}
\newcommand{\CompendiumRef}[2]{\ifmmode\text{\upshape S6\_09, Paper~#1, label~\nolinkurl{#2}}\else S6\_09, Paper~#1, label~\nolinkurl{#2}\fi}
\newcommand{\SeriesRefText}[1]{\ifcsname SRef@#1\endcsname\csname SRef@#1\endcsname\else\textbf{UNRESOLVED}\fi}
\newcommand{\SeriesRef}[1]{\ifmmode\text{\SeriesRefText{#1}}\else\SeriesRefText{#1}\fi}
\AfterEndEnvironment{proof}{\par\begingroup\footnotesize\noindent\hyperref[sec:results]{Return to \S1.}\par\endgroup\smallskip}

\newcommand{\CurrentPaper}{I}
\expandafter\def\csname SRef@wb:pshift:eq:periods\endcsname{\hyperlink{tag-03CN}{I.\ref*{wb:pshift:eq:periods}}}
\expandafter\def\csname SRef@wb:pshift:eq:rank-one-map\endcsname{\hyperlink{tag-03CO}{I.\ref*{wb:pshift:eq:rank-one-map}}}
\expandafter\def\csname SRef@wb:pshift:eq:finite-data\endcsname{\hyperlink{tag-03C8}{I.\ref*{wb:pshift:eq:finite-data}}}
\expandafter\def\csname SRef@wb:pshift:eq:covariance\endcsname{\hyperlink{tag-03C2}{I.\ref*{wb:pshift:eq:covariance}}}
\expandafter\def\csname SRef@wb:pshift:eq:integer-shear\endcsname{\hyperlink{tag-03CF}{I.\ref*{wb:pshift:eq:integer-shear}}}
\expandafter\def\csname SRef@wb:pshift:eq:centralizer\endcsname{\hyperlink{tag-03BZ}{I.\ref*{wb:pshift:eq:centralizer}}}
\expandafter\def\csname SRef@wb:pshift:eq:period-morphism\endcsname{\hyperlink{tag-03CM}{I.\ref*{wb:pshift:eq:period-morphism}}}
\expandafter\def\csname SRef@wb:pshift:eq:norms\endcsname{\hyperlink{tag-03CK}{I.\ref*{wb:pshift:eq:norms}}}
\expandafter\def\csname SRef@wb:pshift:eq:finite-equation\endcsname{\hyperlink{tag-03C9}{I.\ref*{wb:pshift:eq:finite-equation}}}
\expandafter\def\csname SRef@wb:pshift:eq:finite-norm-obstruction\endcsname{\hyperlink{tag-03CB}{I.\ref*{wb:pshift:eq:finite-norm-obstruction}}}
\expandafter\def\csname SRef@wb:pshift:eq:finite-class\endcsname{\hyperlink{tag-03C7}{I.\ref*{wb:pshift:eq:finite-class}}}
\expandafter\def\csname SRef@wb:pshift:eq:local-corrections\endcsname{\hyperlink{tag-03CH}{I.\ref*{wb:pshift:eq:local-corrections}}}
\expandafter\def\csname SRef@wb:pshift:eq:local-matrix-identity\endcsname{\hyperlink{tag-03CI}{I.\ref*{wb:pshift:eq:local-matrix-identity}}}
\expandafter\def\csname SRef@wb:pshift:eq:finite-middle-sections\endcsname{\hyperlink{tag-03CA}{I.\ref*{wb:pshift:eq:finite-middle-sections}}}
\expandafter\def\csname SRef@wb:pshift:eq:local-affine-cocycle\endcsname{\hyperlink{tag-03CG}{I.\ref*{wb:pshift:eq:local-affine-cocycle}}}
\expandafter\def\csname SRef@wb:pshift:eq:actual-overlap-difference\endcsname{\hyperlink{tag-03BX}{I.\ref*{wb:pshift:eq:actual-overlap-difference}}}
\expandafter\def\csname SRef@wb:pshift:eq:cusp-periods\endcsname{\hyperlink{tag-03C5}{I.\ref*{wb:pshift:eq:cusp-periods}}}
\expandafter\def\csname SRef@wb:pshift:eq:cusp-actions-identical\endcsname{\hyperlink{tag-03C3}{I.\ref*{wb:pshift:eq:cusp-actions-identical}}}
\expandafter\def\csname SRef@wb:pshift:eq:cocycle-law\endcsname{\hyperlink{tag-03C0}{I.\ref*{wb:pshift:eq:cocycle-law}}}
\expandafter\def\csname SRef@wb:pshift:eq:cusp-cocycle\endcsname{\hyperlink{tag-03C4}{I.\ref*{wb:pshift:eq:cusp-cocycle}}}
\expandafter\def\csname SRef@wb:pshift:eq:cusp-translation\endcsname{\hyperlink{tag-03C6}{I.\ref*{wb:pshift:eq:cusp-translation}}}
\expandafter\def\csname SRef@wb:pshift:eq:affine-bundle-action\endcsname{\hyperlink{tag-03BY}{I.\ref*{wb:pshift:eq:affine-bundle-action}}}
\expandafter\def\csname SRef@wb:pshift:eq:vector-sheaf\endcsname{\hyperlink{tag-03CP}{I.\ref*{wb:pshift:eq:vector-sheaf}}}
\expandafter\def\csname SRef@wb:pshift:eq:no-analytic-obstruction\endcsname{\hyperlink{tag-03CJ}{I.\ref*{wb:pshift:eq:no-analytic-obstruction}}}
\expandafter\def\csname SRef@wb:pshift:eq:global-correction\endcsname{\hyperlink{tag-03CE}{I.\ref*{wb:pshift:eq:global-correction}}}
\expandafter\def\csname SRef@wb:pshift:eq:full-isomorphism\endcsname{\hyperlink{tag-03CD}{I.\ref*{wb:pshift:eq:full-isomorphism}}}
\expandafter\def\csname SRef@wb:pshift:eq:finite-overlap-identity\endcsname{\hyperlink{tag-03CC}{I.\ref*{wb:pshift:eq:finite-overlap-identity}}}
\expandafter\def\csname SRef@wb:pshift:eq:coherent-shift-action\endcsname{\hyperlink{tag-03C1}{I.\ref*{wb:pshift:eq:coherent-shift-action}}}
\expandafter\def\csname SRef@wb:pshift:eq:parameter-disc\endcsname{\hyperlink{tag-03CL}{I.\ref*{wb:pshift:eq:parameter-disc}}}
\expandafter\def\csname SRef@mapsI:quotient-eq-2\endcsname{\hyperlink{tag-02OZ}{I.\ref*{mapsI:quotient-eq-2}}}
\expandafter\def\csname SRef@mapsI:quotient-eq-3\endcsname{\hyperlink{tag-02P0}{I.\ref*{mapsI:quotient-eq-3}}}
\expandafter\def\csname SRef@mapsI:quotient-eq-4\endcsname{\hyperlink{tag-02P1}{I.\ref*{mapsI:quotient-eq-4}}}
\expandafter\def\csname SRef@mapsI:quotient-eq-5\endcsname{\hyperlink{tag-02P2}{I.\ref*{mapsI:quotient-eq-5}}}
\expandafter\def\csname SRef@mapsI:quotient-eq-7\endcsname{\hyperlink{tag-02P3}{I.\ref*{mapsI:quotient-eq-7}}}
\expandafter\def\csname SRef@app:rederive:fam-src-D.1.2\endcsname{\hyperlink{tag-000F}{I.\ref*{app:rederive:fam-src-D.1.2}}}
\expandafter\def\csname SRef@app:rederive:fam-src-D.1.6\endcsname{\hyperlink{tag-000G}{I.\ref*{app:rederive:fam-src-D.1.6}}}
\expandafter\def\csname SRef@app:rederive:fam-src-D.2.3\endcsname{\hyperlink{tag-000H}{I.\ref*{app:rederive:fam-src-D.2.3}}}
\expandafter\def\csname SRef@mapsI:real-eq-1\endcsname{\hyperlink{tag-02P4}{I.\ref*{mapsI:real-eq-1}}}
\expandafter\def\csname SRef@mapsI:real-eq-2\endcsname{\hyperlink{tag-02P4}{I.\ref*{mapsI:real-eq-2}}}
\expandafter\def\csname SRef@mapsI:real-eq-4\endcsname{\hyperlink{tag-02P5}{I.\ref*{mapsI:real-eq-4}}}
\expandafter\def\csname SRef@mapsI:real-eq-5\endcsname{\hyperlink{tag-02P6}{I.\ref*{mapsI:real-eq-5}}}
\expandafter\def\csname SRef@mapsI:real-eq-6\endcsname{\hyperlink{tag-02P6}{I.\ref*{mapsI:real-eq-6}}}
\expandafter\def\csname SRef@mapsI:real-eq-7\endcsname{\hyperlink{tag-02P7}{I.\ref*{mapsI:real-eq-7}}}
\expandafter\def\csname SRef@mapsI:real-eq-8\endcsname{\hyperlink{tag-02P8}{I.\ref*{mapsI:real-eq-8}}}
\expandafter\def\csname SRef@ari:open:norm\endcsname{\hyperlink{tag-003T}{I.\ref*{ari:open:norm}}}
\expandafter\def\csname SRef@push:src:C.3.5\endcsname{\hyperlink{tag-0322}{I.\ref*{push:src:C.3.5}}}
\expandafter\def\csname SRef@push:src:B.4.3\endcsname{\hyperlink{tag-030Z}{I.\ref*{push:src:B.4.3}}}
\expandafter\def\csname SRef@push:src:D.4.6\endcsname{\hyperlink{tag-032O}{I.\ref*{push:src:D.4.6}}}
\expandafter\def\csname SRef@inv:A1.0\endcsname{\hyperlink{tag-01RQ}{I.\ref*{inv:A1.0}}}
\expandafter\def\csname SRef@wb:lem:cusp-structure-sheaf-audit\endcsname{\hyperlink{tag-01RQ}{I.\ref*{wb:lem:cusp-structure-sheaf-audit}}}
\expandafter\def\csname SRef@wb:thm:structure-direct-images-audit\endcsname{\hyperlink{tag-01RQ}{I.\ref*{wb:thm:structure-direct-images-audit}}}
\expandafter\def\csname SRef@inv:C3.6\endcsname{\hyperlink{tag-01SJ}{I.\ref*{inv:C3.6}}}
\expandafter\def\csname SRef@push:src:E.3.1\endcsname{\hyperlink{tag-033I}{I.\ref*{push:src:E.3.1}}}
\expandafter\def\csname SRef@ari:C.5.1\endcsname{\hyperlink{tag-000U}{I.\ref*{ari:C.5.1}}}
\expandafter\def\csname SRef@ari:C.6.R4\endcsname{\hyperlink{tag-0013}{I.\ref*{ari:C.6.R4}}}
\expandafter\def\csname SRef@ari:C.4.1\endcsname{\hyperlink{tag-000P}{I.\ref*{ari:C.4.1}}}
\expandafter\def\csname SRef@ari:C.7.7\endcsname{\hyperlink{tag-001G}{I.\ref*{ari:C.7.7}}}
\expandafter\def\csname SRef@ari:C.7.8\endcsname{\hyperlink{tag-001H}{I.\ref*{ari:C.7.8}}}
\expandafter\def\csname SRef@push:src:A.4.5\endcsname{\hyperlink{tag-02ZV}{I.\ref*{push:src:A.4.5}}}
\expandafter\def\csname SRef@push:A45\endcsname{\hyperlink{tag-02ZV}{I.\ref*{push:A45}}}
\expandafter\def\csname SRef@push:src:B.3.5\endcsname{\hyperlink{tag-0309}{I.\ref*{push:src:B.3.5}}}
\expandafter\def\csname SRef@push:B35\endcsname{\hyperlink{tag-0309}{I.\ref*{push:B35}}}
\expandafter\def\csname SRef@push:src:B.4.1\endcsname{\hyperlink{tag-030D}{I.\ref*{push:src:B.4.1}}}
\expandafter\def\csname SRef@push:B41\endcsname{\hyperlink{tag-030D}{I.\ref*{push:B41}}}
\expandafter\def\csname SRef@push:src:B.2.2\endcsname{\hyperlink{tag-0303}{I.\ref*{push:src:B.2.2}}}
\expandafter\def\csname SRef@push:B22\endcsname{\hyperlink{tag-0303}{I.\ref*{push:B22}}}
\expandafter\def\csname SRef@push:src:B.5.1\endcsname{\hyperlink{tag-031A}{I.\ref*{push:src:B.5.1}}}
\expandafter\def\csname SRef@push:B51\endcsname{\hyperlink{tag-031A}{I.\ref*{push:B51}}}
\expandafter\def\csname SRef@push:src:B.2.3\endcsname{\hyperlink{tag-0304}{I.\ref*{push:src:B.2.3}}}
\expandafter\def\csname SRef@push:B23\endcsname{\hyperlink{tag-0304}{I.\ref*{push:B23}}}
\expandafter\def\csname SRef@push:src:D.3.2\endcsname{\hyperlink{tag-032I}{I.\ref*{push:src:D.3.2}}}
\expandafter\def\csname SRef@push:src:E.1.3\endcsname{\hyperlink{tag-032T}{I.\ref*{push:src:E.1.3}}}
\expandafter\def\csname SRef@fam:eq:cd-beta\endcsname{\hyperlink{tag-01P5}{I.\ref*{fam:eq:cd-beta}}}
\expandafter\def\csname SRef@ari:eq:ca-GHbounds\endcsname{\hyperlink{tag-0024}{I.\ref*{ari:eq:ca-GHbounds}}}
\expandafter\def\csname SRef@eq:P5L\endcsname{\hyperlink{tag-01OG}{I.\ref*{eq:P5L}}}
\expandafter\def\csname SRef@eq:P5floor\endcsname{\hyperlink{tag-01OI}{I.\ref*{eq:P5floor}}}
\expandafter\def\csname SRef@eq:P5normalization\endcsname{\hyperlink{tag-01OJ}{I.\ref*{eq:P5normalization}}}
\expandafter\def\csname SRef@eq:P5exponential\endcsname{\hyperlink{tag-01OH}{I.\ref*{eq:P5exponential}}}
\expandafter\def\csname SRef@eq:P5H2\endcsname{\hyperlink{tag-01OF}{I.\ref*{eq:P5H2}}}
\expandafter\def\csname SRef@ari:eq:cp-sigma\endcsname{\hyperlink{tag-003K}{I.\ref*{ari:eq:cp-sigma}}}
\expandafter\def\csname SRef@ari:eq:cp-klein\endcsname{\hyperlink{tag-003F}{I.\ref*{ari:eq:cp-klein}}}
\expandafter\def\csname SRef@ari:eq:cp-periods\endcsname{\hyperlink{tag-003G}{I.\ref*{ari:eq:cp-periods}}}
\expandafter\def\csname SRef@ari:eq:cp-v\endcsname{\hyperlink{tag-003L}{I.\ref*{ari:eq:cp-v}}}
\expandafter\def\csname SRef@ari:eq:cp-chi\endcsname{\hyperlink{tag-003E}{I.\ref*{ari:eq:cp-chi}}}
\expandafter\def\csname SRef@ari:eq:ca-derivatives\endcsname{\hyperlink{tag-002B}{I.\ref*{ari:eq:ca-derivatives}}}
\expandafter\def\csname SRef@ari:eq:cp-MTdim\endcsname{\hyperlink{tag-003D}{I.\ref*{ari:eq:cp-MTdim}}}
\expandafter\def\csname SRef@ari:eq:cp-positive-centre\endcsname{\hyperlink{tag-003H}{I.\ref*{ari:eq:cp-positive-centre}}}
\expandafter\def\csname SRef@ari:eq:cp-real-linear\endcsname{\hyperlink{tag-003I}{I.\ref*{ari:eq:cp-real-linear}}}
\expandafter\def\csname SRef@ari:eq:cp-relation\endcsname{\hyperlink{tag-003J}{I.\ref*{ari:eq:cp-relation}}}
\expandafter\def\csname SRef@ari:eq:ct-CMmult\endcsname{\hyperlink{tag-003Q}{I.\ref*{ari:eq:ct-CMmult}}}
\expandafter\def\csname SRef@ari:eq:ct-Gaussian\endcsname{\hyperlink{tag-003R}{I.\ref*{ari:eq:ct-Gaussian}}}
\expandafter\def\csname SRef@ari:eq:cd-slope\endcsname{\hyperlink{tag-0033}{I.\ref*{ari:eq:cd-slope}}}
\expandafter\def\csname SRef@ari:eq:cd-compatibility\endcsname{\hyperlink{tag-0032}{I.\ref*{ari:eq:cd-compatibility}}}
\expandafter\def\csname SRef@ari:eq:ce-ribet-r\endcsname{\hyperlink{tag-0036}{I.\ref*{ari:eq:ce-ribet-r}}}
\expandafter\def\csname SRef@ari:eq:ce-zeta\endcsname{\hyperlink{tag-0037}{I.\ref*{ari:eq:ce-zeta}}}
\expandafter\def\csname SRef@ari:eq:ca-delta-h\endcsname{\hyperlink{tag-0028}{I.\ref*{ari:eq:ca-delta-h}}}
\expandafter\def\csname SRef@ari:eq:ce-gap\endcsname{\hyperlink{tag-0035}{I.\ref*{ari:eq:ce-gap}}}
\expandafter\def\csname SRef@ari:eq:ce-explicit-witness\endcsname{\hyperlink{tag-0034}{I.\ref*{ari:eq:ce-explicit-witness}}}
\expandafter\def\csname SRef@ari:eq:cc-fixedfield\endcsname{\hyperlink{tag-002Q}{I.\ref*{ari:eq:cc-fixedfield}}}
\expandafter\def\csname SRef@ari:eq:cc-rankcriterion\endcsname{\hyperlink{tag-002X}{I.\ref*{ari:eq:cc-rankcriterion}}}
\expandafter\def\csname SRef@ari:eq:division-ratio-model\endcsname{\hyperlink{tag-003S}{I.\ref*{ari:eq:division-ratio-model}}}
\expandafter\def\csname SRef@ari:eq:ca-ladder\endcsname{\hyperlink{tag-002D}{I.\ref*{ari:eq:ca-ladder}}}
\expandafter\def\csname SRef@fam:eq:F-g\endcsname{\hyperlink{tag-01OT}{I.\ref*{fam:eq:F-g}}}
\expandafter\def\csname SRef@fam:eq:F-Gamma\endcsname{\hyperlink{tag-01OQ}{I.\ref*{fam:eq:F-Gamma}}}
\expandafter\def\csname SRef@fam:eq:F-law\endcsname{\hyperlink{tag-01OV}{I.\ref*{fam:eq:F-law}}}
\expandafter\def\csname SRef@fam:eq:F-SLpresentation\endcsname{\hyperlink{tag-01OS}{I.\ref*{fam:eq:F-SLpresentation}}}
\expandafter\def\csname SRef@fam:eq:F-H\endcsname{\hyperlink{tag-01OR}{I.\ref*{fam:eq:F-H}}}
\expandafter\def\csname SRef@fam:eq:F-qD\endcsname{\hyperlink{tag-01OW}{I.\ref*{fam:eq:F-qD}}}
\expandafter\def\csname SRef@fam:eq:cd-wp\endcsname{\hyperlink{tag-01P8}{I.\ref*{fam:eq:cd-wp}}}
\expandafter\def\csname SRef@fam:eq:F-tauaction\endcsname{\hyperlink{tag-01OX}{I.\ref*{fam:eq:F-tauaction}}}
\expandafter\def\csname SRef@fam:eq:F-j\endcsname{\hyperlink{tag-01OU}{I.\ref*{fam:eq:F-j}}}
\expandafter\def\csname SRef@fam:eq:cd-NS\endcsname{\hyperlink{tag-01P4}{I.\ref*{fam:eq:cd-NS}}}
\expandafter\def\csname SRef@fam:eq:cd-theta-sections\endcsname{\hyperlink{tag-01P7}{I.\ref*{fam:eq:cd-theta-sections}}}
\expandafter\def\csname SRef@mapsI:closing:eq:C-explicit-closing\endcsname{\hyperlink{tag-02OP}{I.\ref*{mapsI:closing:eq:C-explicit-closing}}}
\expandafter\def\csname SRef@mapsI:closing:eq:C-twelve-scalar\endcsname{\hyperlink{tag-02OQ}{I.\ref*{mapsI:closing:eq:C-twelve-scalar}}}
\expandafter\def\csname SRef@eq:d-period-normal\endcsname{\hyperlink{tag-01ON}{I.\ref*{eq:d-period-normal}}}
\expandafter\def\csname SRef@eq:d-limit-group\endcsname{\hyperlink{tag-01OM}{I.\ref*{eq:d-limit-group}}}
\expandafter\def\csname SRef@eq:d-dual-exact\endcsname{\hyperlink{tag-01OL}{I.\ref*{eq:d-dual-exact}}}
\expandafter\def\csname SRef@eq:P2H\endcsname{\hyperlink{tag-01O8}{I.\ref*{eq:P2H}}}
\expandafter\def\csname SRef@fam:eq:cd-norm\endcsname{\hyperlink{tag-01P6}{I.\ref*{fam:eq:cd-norm}}}
\expandafter\def\csname SRef@eq:P2zg\endcsname{\hyperlink{tag-01OC}{I.\ref*{eq:P2zg}}}
\expandafter\def\csname SRef@eq:P2F\endcsname{\hyperlink{tag-01O5}{I.\ref*{eq:P2F}}}
\expandafter\def\csname SRef@eq:P2FQ\endcsname{\hyperlink{tag-01O6}{I.\ref*{eq:P2FQ}}}
\expandafter\def\csname SRef@eq:P2FR\endcsname{\hyperlink{tag-01O7}{I.\ref*{eq:P2FR}}}
\expandafter\def\csname SRef@eq:P2chain\endcsname{\hyperlink{tag-01O9}{I.\ref*{eq:P2chain}}}
\expandafter\def\csname SRef@eq:P2split\endcsname{\hyperlink{tag-01OB}{I.\ref*{eq:P2split}}}
\expandafter\def\csname SRef@eq:P2ratio\endcsname{\hyperlink{tag-01OA}{I.\ref*{eq:P2ratio}}}
\expandafter\def\csname SRef@inv:eq:A-E-line\endcsname{\hyperlink{tag-01SZ}{I.\ref*{inv:eq:A-E-line}}}
\expandafter\def\csname SRef@inv:eq:A-H-images\endcsname{\hyperlink{tag-01T0}{I.\ref*{inv:eq:A-H-images}}}
\expandafter\def\csname SRef@inv:eq:A-N-images\endcsname{\hyperlink{tag-01T1}{I.\ref*{inv:eq:A-N-images}}}
\expandafter\def\csname SRef@ari:eq:cs-isogeny\endcsname{\hyperlink{tag-003O}{I.\ref*{ari:eq:cs-isogeny}}}
\expandafter\def\csname SRef@inv:eq:A-absolute-table\endcsname{\hyperlink{tag-01T2}{I.\ref*{inv:eq:A-absolute-table}}}
\expandafter\def\csname SRef@inv:eq:A-s-branch\endcsname{\hyperlink{tag-01T5}{I.\ref*{inv:eq:A-s-branch}}}
\expandafter\def\csname SRef@inv:eq:A-modular-K\endcsname{\hyperlink{tag-01T3}{I.\ref*{inv:eq:A-modular-K}}}
\expandafter\def\csname SRef@inv:eq:A-mu-value\endcsname{\hyperlink{tag-01T4}{I.\ref*{inv:eq:A-mu-value}}}
\expandafter\def\csname SRef@inv:eq:A-DI-branch\endcsname{\hyperlink{tag-01SY}{I.\ref*{inv:eq:A-DI-branch}}}
\expandafter\def\csname SRef@inv:eq:C-hodge-input\endcsname{\hyperlink{tag-01TF}{I.\ref*{inv:eq:C-hodge-input}}}
\expandafter\def\csname SRef@inv:eq:C-BC-A\endcsname{\hyperlink{tag-01T6}{I.\ref*{inv:eq:C-BC-A}}}
\expandafter\def\csname SRef@inv:eq:C-arrows\endcsname{\hyperlink{tag-01TA}{I.\ref*{inv:eq:C-arrows}}}
\expandafter\def\csname SRef@inv:eq:C-Pic-map\endcsname{\hyperlink{tag-01T8}{I.\ref*{inv:eq:C-Pic-map}}}
\expandafter\def\csname SRef@inv:eq:C-G-Green\endcsname{\hyperlink{tag-01T7}{I.\ref*{inv:eq:C-G-Green}}}
\expandafter\def\csname SRef@ari:eq:cs-division\endcsname{\hyperlink{tag-003N}{I.\ref*{ari:eq:cs-division}}}
\expandafter\def\csname SRef@inv:eq:C-factor12\endcsname{\hyperlink{tag-01TE}{I.\ref*{inv:eq:C-factor12}}}
\expandafter\def\csname SRef@inv:eq:C-collar-sum\endcsname{\hyperlink{tag-01TB}{I.\ref*{inv:eq:C-collar-sum}}}
\expandafter\def\csname SRef@inv:eq:C-delta\endcsname{\hyperlink{tag-01TC}{I.\ref*{inv:eq:C-delta}}}
\expandafter\def\csname SRef@inv:eq:C-Ricci\endcsname{\hyperlink{tag-01T9}{I.\ref*{inv:eq:C-Ricci}}}
\expandafter\def\csname SRef@inv:eq:C-eta-negative\endcsname{\hyperlink{tag-01TD}{I.\ref*{inv:eq:C-eta-negative}}}
\expandafter\def\csname SRef@ari:eq:cs-Li\endcsname{\hyperlink{tag-003M}{I.\ref*{ari:eq:cs-Li}}}
\expandafter\def\csname SRef@ari:eq:ch-global\endcsname{\hyperlink{tag-0039}{I.\ref*{ari:eq:ch-global}}}
\expandafter\def\csname SRef@ari:eq:ch-local-mult\endcsname{\hyperlink{tag-003B}{I.\ref*{ari:eq:ch-local-mult}}}
\expandafter\def\csname SRef@ari:eq:ch-local-exact\endcsname{\hyperlink{tag-003A}{I.\ref*{ari:eq:ch-local-exact}}}
\expandafter\def\csname SRef@ari:eq:ch-bound\endcsname{\hyperlink{tag-0038}{I.\ref*{ari:eq:ch-bound}}}
\expandafter\def\csname SRef@ari:eq:cs-torsionvalue\endcsname{\hyperlink{tag-003P}{I.\ref*{ari:eq:cs-torsionvalue}}}
\expandafter\def\csname SRef@ari:eq:ch-wp-bracket\endcsname{\hyperlink{tag-003C}{I.\ref*{ari:eq:ch-wp-bracket}}}
\expandafter\def\csname SRef@ari:eq:Cr-first\endcsname{\hyperlink{tag-0020}{I.\ref*{ari:eq:Cr-first}}}
\expandafter\def\csname SRef@ari:eq:Cr-second\endcsname{\hyperlink{tag-0022}{I.\ref*{ari:eq:Cr-second}}}
\expandafter\def\csname SRef@ari:eq:Cr-tails\endcsname{\hyperlink{tag-0023}{I.\ref*{ari:eq:Cr-tails}}}
\expandafter\def\csname SRef@ari:eq:Cr-rational-factor\endcsname{\hyperlink{tag-0021}{I.\ref*{ari:eq:Cr-rational-factor}}}
\expandafter\def\csname SRef@ari:eq:Cgz-central-integrality\endcsname{\hyperlink{tag-001X}{I.\ref*{ari:eq:Cgz-central-integrality}}}
\expandafter\def\csname SRef@ari:eq:Cgz-product\endcsname{\hyperlink{tag-001Z}{I.\ref*{ari:eq:Cgz-product}}}
\expandafter\def\csname SRef@ari:eq:Cgz-denominator\endcsname{\hyperlink{tag-001Y}{I.\ref*{ari:eq:Cgz-denominator}}}
\expandafter\def\csname SRef@ari:eq:ca-mahler\endcsname{\hyperlink{tag-002E}{I.\ref*{ari:eq:ca-mahler}}}
\expandafter\def\csname SRef@ari:eq:ca-wpmean\endcsname{\hyperlink{tag-002H}{I.\ref*{ari:eq:ca-wpmean}}}
\expandafter\def\csname SRef@ari:eq:ca-sup\endcsname{\hyperlink{tag-002G}{I.\ref*{ari:eq:ca-sup}}}
\expandafter\def\csname SRef@ari:eq:ca-nu\endcsname{\hyperlink{tag-002F}{I.\ref*{ari:eq:ca-nu}}}
\expandafter\def\csname SRef@ari:eq:ca-current\endcsname{\hyperlink{tag-0026}{I.\ref*{ari:eq:ca-current}}}
\expandafter\def\csname SRef@ari:eq:ca-density0\endcsname{\hyperlink{tag-0029}{I.\ref*{ari:eq:ca-density0}}}
\expandafter\def\csname SRef@ari:eq:ca-density1\endcsname{\hyperlink{tag-002A}{I.\ref*{ari:eq:ca-density1}}}
\expandafter\def\csname SRef@ari:eq:ca-cuspexp\endcsname{\hyperlink{tag-0027}{I.\ref*{ari:eq:ca-cuspexp}}}
\expandafter\def\csname SRef@ari:eq:cc-EZ\endcsname{\hyperlink{tag-002I}{I.\ref*{ari:eq:cc-EZ}}}
\expandafter\def\csname SRef@ari:eq:cc-radial\endcsname{\hyperlink{tag-002W}{I.\ref*{ari:eq:cc-radial}}}
\expandafter\def\csname SRef@ari:eq:cc-FR\endcsname{\hyperlink{tag-002J}{I.\ref*{ari:eq:cc-FR}}}
\expandafter\def\csname SRef@ari:eq:cc-cuspbound\endcsname{\hyperlink{tag-002O}{I.\ref*{ari:eq:cc-cuspbound}}}
\expandafter\def\csname SRef@ari:eq:cc-ode\endcsname{\hyperlink{tag-002T}{I.\ref*{ari:eq:cc-ode}}}
\expandafter\def\csname SRef@ari:eq:ca-exp\endcsname{\hyperlink{tag-002C}{I.\ref*{ari:eq:ca-exp}}}
\expandafter\def\csname SRef@ari:eq:cc-h\endcsname{\hyperlink{tag-002R}{I.\ref*{ari:eq:cc-h}}}
\expandafter\def\csname SRef@ari:eq:cc-Fseries\endcsname{\hyperlink{tag-002L}{I.\ref*{ari:eq:cc-Fseries}}}
\expandafter\def\csname SRef@ari:eq:cc-Ffactored\endcsname{\hyperlink{tag-002K}{I.\ref*{ari:eq:cc-Ffactored}}}
\expandafter\def\csname SRef@ari:eq:cc-norms\endcsname{\hyperlink{tag-002S}{I.\ref*{ari:eq:cc-norms}}}
\expandafter\def\csname SRef@ari:eq:cc-recurrence\endcsname{\hyperlink{tag-002Y}{I.\ref*{ari:eq:cc-recurrence}}}
\expandafter\def\csname SRef@ari:eq:cc-power\endcsname{\hyperlink{tag-002V}{I.\ref*{ari:eq:cc-power}}}
\expandafter\def\csname SRef@ari:eq:cc-constants\endcsname{\hyperlink{tag-002N}{I.\ref*{ari:eq:cc-constants}}}
\expandafter\def\csname SRef@ari:eq:cc-outer\endcsname{\hyperlink{tag-002U}{I.\ref*{ari:eq:cc-outer}}}
\expandafter\def\csname SRef@ari:eq:cc-tail60\endcsname{\hyperlink{tag-0031}{I.\ref*{ari:eq:cc-tail60}}}
\expandafter\def\csname SRef@ari:eq:cc-tail100\endcsname{\hyperlink{tag-0030}{I.\ref*{ari:eq:cc-tail100}}}
\expandafter\def\csname SRef@ari:eq:ca-H\endcsname{\hyperlink{tag-0025}{I.\ref*{ari:eq:ca-H}}}
\expandafter\def\csname SRef@ari:eq:cc-finite-table\endcsname{\hyperlink{tag-002P}{I.\ref*{ari:eq:cc-finite-table}}}
\expandafter\def\csname SRef@ari:eq:cc-scalar\endcsname{\hyperlink{tag-002Z}{I.\ref*{ari:eq:cc-scalar}}}
\expandafter\def\csname SRef@ari:eq:cc-P\endcsname{\hyperlink{tag-002M}{I.\ref*{ari:eq:cc-P}}}
\expandafter\def\csname SRef@fam:eq:ab-root-degree\endcsname{\hyperlink{tag-01P3}{I.\ref*{fam:eq:ab-root-degree}}}
\expandafter\def\csname SRef@fam:eq:ab-cusp-resolution\endcsname{\hyperlink{tag-01OY}{I.\ref*{fam:eq:ab-cusp-resolution}}}
\expandafter\def\csname SRef@fam:eq:ab-log-tangent\endcsname{\hyperlink{tag-01P0}{I.\ref*{fam:eq:ab-log-tangent}}}
\expandafter\def\csname SRef@fam:eq:ab-relative-dualizing\endcsname{\hyperlink{tag-01P2}{I.\ref*{fam:eq:ab-relative-dualizing}}}
\expandafter\def\csname SRef@fam:eq:ab-direct-images\endcsname{\hyperlink{tag-01OZ}{I.\ref*{fam:eq:ab-direct-images}}}
\expandafter\def\csname SRef@fam:eq:ab-normal-quotient\endcsname{\hyperlink{tag-01P1}{I.\ref*{fam:eq:ab-normal-quotient}}}
\expandafter\def\csname SRef@eq:P5tangent\endcsname{\hyperlink{tag-01OK}{I.\ref*{eq:P5tangent}}}
\expandafter\def\csname SRef@fig:I:orbifold\endcsname{\hyperlink{tag-01RL}{I.\ref*{fig:I:orbifold}}}
\expandafter\def\csname SRef@fig:I:parameter-disc\endcsname{\hyperlink{tag-01RM}{I.\ref*{fig:I:parameter-disc}}}
\expandafter\def\csname SRef@fig:I:monodromy-domain\endcsname{\hyperlink{tag-01RK}{I.\ref*{fig:I:monodromy-domain}}}
\expandafter\def\csname SRef@push:src:A.3.1\endcsname{\hyperlink{tag-02ZO}{I.\ref*{push:src:A.3.1}}}
\expandafter\def\csname SRef@push:A31\endcsname{\hyperlink{tag-02ZO}{I.\ref*{push:A31}}}
\expandafter\def\csname SRef@push:src:D.1.2\endcsname{\hyperlink{tag-032B}{I.\ref*{push:src:D.1.2}}}
\expandafter\def\csname SRef@push:src:D.2.1\endcsname{\hyperlink{tag-032D}{I.\ref*{push:src:D.2.1}}}
\expandafter\def\csname SRef@I:weierstrass-comparison\endcsname{\hyperlink{tag-0002}{I.\ref*{I:weierstrass-comparison}}}
\expandafter\def\csname SRef@wb:pshift:lem:centralizer\endcsname{\hyperlink{tag-03CQ}{I.\ref*{wb:pshift:lem:centralizer}}}
\expandafter\def\csname SRef@fam:src:F.1.2\endcsname{\hyperlink{tag-01RA}{I.\ref*{fam:src:F.1.2}}}
\expandafter\def\csname SRef@fam:st:F.1.2\endcsname{\hyperlink{tag-01RA}{I.\ref*{fam:st:F.1.2}}}
\expandafter\def\csname SRef@fam:src:F.2.1\endcsname{\hyperlink{tag-01RB}{I.\ref*{fam:src:F.2.1}}}
\expandafter\def\csname SRef@fam:st:F.2.1\endcsname{\hyperlink{tag-01RB}{I.\ref*{fam:st:F.2.1}}}
\expandafter\def\csname SRef@fam:src:F.2.3\endcsname{\hyperlink{tag-01RD}{I.\ref*{fam:src:F.2.3}}}
\expandafter\def\csname SRef@fam:st:F.2.3\endcsname{\hyperlink{tag-01RD}{I.\ref*{fam:st:F.2.3}}}
\expandafter\def\csname SRef@fam:src:F.3.2\endcsname{\hyperlink{tag-01RF}{I.\ref*{fam:src:F.3.2}}}
\expandafter\def\csname SRef@fam:st:F.3.2\endcsname{\hyperlink{tag-01RF}{I.\ref*{fam:st:F.3.2}}}
\expandafter\def\csname SRef@fam:src:F.3.3\endcsname{\hyperlink{tag-01RG}{I.\ref*{fam:src:F.3.3}}}
\expandafter\def\csname SRef@fam:st:F.3.3\endcsname{\hyperlink{tag-01RG}{I.\ref*{fam:st:F.3.3}}}
\expandafter\def\csname SRef@fam:src:D.1.1\endcsname{\hyperlink{tag-01QR}{I.\ref*{fam:src:D.1.1}}}
\expandafter\def\csname SRef@fam:st:D.1.1\endcsname{\hyperlink{tag-01QR}{I.\ref*{fam:st:D.1.1}}}
\expandafter\def\csname SRef@fam:src:C.3.1\endcsname{\hyperlink{tag-01QM}{I.\ref*{fam:src:C.3.1}}}
\expandafter\def\csname SRef@fam:st:C.3.1\endcsname{\hyperlink{tag-01QM}{I.\ref*{fam:st:C.3.1}}}
\expandafter\def\csname SRef@fam:src:C.3.2\endcsname{\hyperlink{tag-01QN}{I.\ref*{fam:src:C.3.2}}}
\expandafter\def\csname SRef@fam:st:C.3.2\endcsname{\hyperlink{tag-01QN}{I.\ref*{fam:st:C.3.2}}}
\expandafter\def\csname SRef@fam:src:C.3.4\endcsname{\hyperlink{tag-01QP}{I.\ref*{fam:src:C.3.4}}}
\expandafter\def\csname SRef@fam:st:C.3.4\endcsname{\hyperlink{tag-01QP}{I.\ref*{fam:st:C.3.4}}}
\expandafter\def\csname SRef@fam:src:D.1.2\endcsname{\hyperlink{tag-01QS}{I.\ref*{fam:src:D.1.2}}}
\expandafter\def\csname SRef@fam:st:D.1.2\endcsname{\hyperlink{tag-01QS}{I.\ref*{fam:st:D.1.2}}}
\expandafter\def\csname SRef@fam:src:D.1.4\endcsname{\hyperlink{tag-01QU}{I.\ref*{fam:src:D.1.4}}}
\expandafter\def\csname SRef@fam:st:D.1.4\endcsname{\hyperlink{tag-01QU}{I.\ref*{fam:st:D.1.4}}}
\expandafter\def\csname SRef@fam:src:D.1.5\endcsname{\hyperlink{tag-01QV}{I.\ref*{fam:src:D.1.5}}}
\expandafter\def\csname SRef@fam:st:D.1.5\endcsname{\hyperlink{tag-01QV}{I.\ref*{fam:st:D.1.5}}}
\expandafter\def\csname SRef@fam:src:D.1.7\endcsname{\hyperlink{tag-01QX}{I.\ref*{fam:src:D.1.7}}}
\expandafter\def\csname SRef@fam:st:D.1.7\endcsname{\hyperlink{tag-01QX}{I.\ref*{fam:st:D.1.7}}}
\expandafter\def\csname SRef@fam:src:D.2.1\endcsname{\hyperlink{tag-01QY}{I.\ref*{fam:src:D.2.1}}}
\expandafter\def\csname SRef@fam:st:D.2.1\endcsname{\hyperlink{tag-01QY}{I.\ref*{fam:st:D.2.1}}}
\expandafter\def\csname SRef@fam:src:D.2.2\endcsname{\hyperlink{tag-01QZ}{I.\ref*{fam:src:D.2.2}}}
\expandafter\def\csname SRef@fam:st:D.2.2\endcsname{\hyperlink{tag-01QZ}{I.\ref*{fam:st:D.2.2}}}
\expandafter\def\csname SRef@mapsI:src:A.3.1\endcsname{\hyperlink{tag-02PE}{I.\ref*{mapsI:src:A.3.1}}}
\expandafter\def\csname SRef@mapsI:src:A.3.4\endcsname{\hyperlink{tag-02PH}{I.\ref*{mapsI:src:A.3.4}}}
\expandafter\def\csname SRef@mapsI:real-eq-9\endcsname{\hyperlink{tag-02PH}{I.\ref*{mapsI:real-eq-9}}}
\expandafter\def\csname SRef@I:monodromy-fundamental-set\endcsname{\hyperlink{tag-0001}{I.\ref*{I:monodromy-fundamental-set}}}
\expandafter\def\csname SRef@fam:src:E.1.1\endcsname{\hyperlink{tag-01R1}{I.\ref*{fam:src:E.1.1}}}
\expandafter\def\csname SRef@fam:st:E.1.1\endcsname{\hyperlink{tag-01R1}{I.\ref*{fam:st:E.1.1}}}
\expandafter\def\csname SRef@fam:src:E.1.2\endcsname{\hyperlink{tag-01R2}{I.\ref*{fam:src:E.1.2}}}
\expandafter\def\csname SRef@fam:st:E.1.2\endcsname{\hyperlink{tag-01R2}{I.\ref*{fam:st:E.1.2}}}
\expandafter\def\csname SRef@fam:src:E.1.3\endcsname{\hyperlink{tag-01R3}{I.\ref*{fam:src:E.1.3}}}
\expandafter\def\csname SRef@fam:st:E.1.3\endcsname{\hyperlink{tag-01R3}{I.\ref*{fam:st:E.1.3}}}
\expandafter\def\csname SRef@fam:src:E.1.4\endcsname{\hyperlink{tag-01R4}{I.\ref*{fam:src:E.1.4}}}
\expandafter\def\csname SRef@fam:st:E.1.4\endcsname{\hyperlink{tag-01R4}{I.\ref*{fam:st:E.1.4}}}
\expandafter\def\csname SRef@fam:src:E.2.1\endcsname{\hyperlink{tag-01R6}{I.\ref*{fam:src:E.2.1}}}
\expandafter\def\csname SRef@fam:st:E.2.1\endcsname{\hyperlink{tag-01R6}{I.\ref*{fam:st:E.2.1}}}
\expandafter\def\csname SRef@push:src:A.1.1\endcsname{\hyperlink{tag-02ZH}{I.\ref*{push:src:A.1.1}}}
\expandafter\def\csname SRef@push:A11\endcsname{\hyperlink{tag-02ZH}{I.\ref*{push:A11}}}
\expandafter\def\csname SRef@push:src:A.4.1\endcsname{\hyperlink{tag-02ZR}{I.\ref*{push:src:A.4.1}}}
\expandafter\def\csname SRef@push:A41\endcsname{\hyperlink{tag-02ZR}{I.\ref*{push:A41}}}
\expandafter\def\csname SRef@push:src:A.4.2\endcsname{\hyperlink{tag-02ZS}{I.\ref*{push:src:A.4.2}}}
\expandafter\def\csname SRef@push:A42\endcsname{\hyperlink{tag-02ZS}{I.\ref*{push:A42}}}
\expandafter\def\csname SRef@push:src:A.4.3\endcsname{\hyperlink{tag-02ZT}{I.\ref*{push:src:A.4.3}}}
\expandafter\def\csname SRef@push:A43\endcsname{\hyperlink{tag-02ZT}{I.\ref*{push:A43}}}
\expandafter\def\csname SRef@push:src:A.4.4\endcsname{\hyperlink{tag-02ZU}{I.\ref*{push:src:A.4.4}}}
\expandafter\def\csname SRef@push:A44\endcsname{\hyperlink{tag-02ZU}{I.\ref*{push:A44}}}
\expandafter\def\csname SRef@push:src:A.4.6\endcsname{\hyperlink{tag-02ZW}{I.\ref*{push:src:A.4.6}}}
\expandafter\def\csname SRef@push:A46\endcsname{\hyperlink{tag-02ZW}{I.\ref*{push:A46}}}
\expandafter\def\csname SRef@push:src:A.4.9\endcsname{\hyperlink{tag-02ZZ}{I.\ref*{push:src:A.4.9}}}
\expandafter\def\csname SRef@push:A49\endcsname{\hyperlink{tag-02ZZ}{I.\ref*{push:A49}}}
\expandafter\def\csname SRef@push:src:A.2.1\endcsname{\hyperlink{tag-02ZJ}{I.\ref*{push:src:A.2.1}}}
\expandafter\def\csname SRef@push:A21\endcsname{\hyperlink{tag-02ZJ}{I.\ref*{push:A21}}}
\expandafter\def\csname SRef@push:src:A.2.5\endcsname{\hyperlink{tag-02ZN}{I.\ref*{push:src:A.2.5}}}
\expandafter\def\csname SRef@push:A25\endcsname{\hyperlink{tag-02ZN}{I.\ref*{push:A25}}}
\expandafter\def\csname SRef@push:src:A.3.3\endcsname{\hyperlink{tag-02ZQ}{I.\ref*{push:src:A.3.3}}}
\expandafter\def\csname SRef@push:A33\endcsname{\hyperlink{tag-02ZQ}{I.\ref*{push:A33}}}
\expandafter\def\csname SRef@push:src:B.1.1\endcsname{\hyperlink{tag-0300}{I.\ref*{push:src:B.1.1}}}
\expandafter\def\csname SRef@push:B11\endcsname{\hyperlink{tag-0300}{I.\ref*{push:B11}}}
\expandafter\def\csname SRef@push:src:B.3.6\endcsname{\hyperlink{tag-030A}{I.\ref*{push:src:B.3.6}}}
\expandafter\def\csname SRef@push:B36\endcsname{\hyperlink{tag-030A}{I.\ref*{push:B36}}}
\expandafter\def\csname SRef@push:src:B.3.7\endcsname{\hyperlink{tag-030B}{I.\ref*{push:src:B.3.7}}}
\expandafter\def\csname SRef@push:B37\endcsname{\hyperlink{tag-030B}{I.\ref*{push:B37}}}
\expandafter\def\csname SRef@push:src:B.3.8\endcsname{\hyperlink{tag-030C}{I.\ref*{push:src:B.3.8}}}
\expandafter\def\csname SRef@push:B38\endcsname{\hyperlink{tag-030C}{I.\ref*{push:B38}}}
\expandafter\def\csname SRef@push:src:B.4.4\endcsname{\hyperlink{tag-0314}{I.\ref*{push:src:B.4.4}}}
\expandafter\def\csname SRef@push:src:B.4.6\endcsname{\hyperlink{tag-0316}{I.\ref*{push:src:B.4.6}}}
\expandafter\def\csname SRef@push:B46\endcsname{\hyperlink{tag-0316}{I.\ref*{push:B46}}}
\expandafter\def\csname SRef@push:src:B.4.7\endcsname{\hyperlink{tag-0317}{I.\ref*{push:src:B.4.7}}}
\expandafter\def\csname SRef@push:src:B.4.8\endcsname{\hyperlink{tag-0318}{I.\ref*{push:src:B.4.8}}}
\expandafter\def\csname SRef@push:src:B.4.9\endcsname{\hyperlink{tag-0319}{I.\ref*{push:src:B.4.9}}}
\expandafter\def\csname SRef@push:src:B.4.10\endcsname{\hyperlink{tag-030E}{I.\ref*{push:src:B.4.10}}}
\expandafter\def\csname SRef@push:src:B.4.11\endcsname{\hyperlink{tag-030F}{I.\ref*{push:src:B.4.11}}}
\expandafter\def\csname SRef@push:src:B.4.12\endcsname{\hyperlink{tag-030G}{I.\ref*{push:src:B.4.12}}}
\expandafter\def\csname SRef@push:src:B.4.13\endcsname{\hyperlink{tag-030H}{I.\ref*{push:src:B.4.13}}}
\expandafter\def\csname SRef@push:src:B.4.14\endcsname{\hyperlink{tag-030I}{I.\ref*{push:src:B.4.14}}}
\expandafter\def\csname SRef@push:src:B.4.15\endcsname{\hyperlink{tag-030J}{I.\ref*{push:src:B.4.15}}}
\expandafter\def\csname SRef@push:B415\endcsname{\hyperlink{tag-030J}{I.\ref*{push:B415}}}
\expandafter\def\csname SRef@push:src:B.4.16\endcsname{\hyperlink{tag-030K}{I.\ref*{push:src:B.4.16}}}
\expandafter\def\csname SRef@push:B416\endcsname{\hyperlink{tag-030K}{I.\ref*{push:B416}}}
\expandafter\def\csname SRef@push:src:B.2.1\endcsname{\hyperlink{tag-0302}{I.\ref*{push:src:B.2.1}}}
\expandafter\def\csname SRef@push:B21\endcsname{\hyperlink{tag-0302}{I.\ref*{push:B21}}}
\expandafter\def\csname SRef@push:src:B.4.17\endcsname{\hyperlink{tag-030L}{I.\ref*{push:src:B.4.17}}}
\expandafter\def\csname SRef@push:src:B.4.18\endcsname{\hyperlink{tag-030M}{I.\ref*{push:src:B.4.18}}}
\expandafter\def\csname SRef@push:B418\endcsname{\hyperlink{tag-030M}{I.\ref*{push:B418}}}
\expandafter\def\csname SRef@push:src:B.4.19\endcsname{\hyperlink{tag-030N}{I.\ref*{push:src:B.4.19}}}
\expandafter\def\csname SRef@push:B419\endcsname{\hyperlink{tag-030N}{I.\ref*{push:B419}}}
\expandafter\def\csname SRef@push:src:B.4.20\endcsname{\hyperlink{tag-030P}{I.\ref*{push:src:B.4.20}}}
\expandafter\def\csname SRef@push:B420\endcsname{\hyperlink{tag-030P}{I.\ref*{push:B420}}}
\expandafter\def\csname SRef@push:src:B.4.21\endcsname{\hyperlink{tag-030Q}{I.\ref*{push:src:B.4.21}}}
\expandafter\def\csname SRef@push:src:B.4.22\endcsname{\hyperlink{tag-030R}{I.\ref*{push:src:B.4.22}}}
\expandafter\def\csname SRef@push:B422\endcsname{\hyperlink{tag-030R}{I.\ref*{push:B422}}}
\expandafter\def\csname SRef@push:src:B.4.23\endcsname{\hyperlink{tag-030S}{I.\ref*{push:src:B.4.23}}}
\expandafter\def\csname SRef@push:src:B.4.24\endcsname{\hyperlink{tag-030T}{I.\ref*{push:src:B.4.24}}}
\expandafter\def\csname SRef@push:B424\endcsname{\hyperlink{tag-030T}{I.\ref*{push:B424}}}
\expandafter\def\csname SRef@push:src:B.4.25\endcsname{\hyperlink{tag-030U}{I.\ref*{push:src:B.4.25}}}
\expandafter\def\csname SRef@push:B425\endcsname{\hyperlink{tag-030U}{I.\ref*{push:B425}}}
\expandafter\def\csname SRef@push:src:B.4.26\endcsname{\hyperlink{tag-030V}{I.\ref*{push:src:B.4.26}}}
\expandafter\def\csname SRef@push:B426\endcsname{\hyperlink{tag-030V}{I.\ref*{push:B426}}}
\expandafter\def\csname SRef@push:src:B.4.27\endcsname{\hyperlink{tag-030W}{I.\ref*{push:src:B.4.27}}}
\expandafter\def\csname SRef@push:B427\endcsname{\hyperlink{tag-030W}{I.\ref*{push:B427}}}
\expandafter\def\csname SRef@push:src:B.4.28\endcsname{\hyperlink{tag-030X}{I.\ref*{push:src:B.4.28}}}
\expandafter\def\csname SRef@push:src:B.4.29\endcsname{\hyperlink{tag-030Y}{I.\ref*{push:src:B.4.29}}}
\expandafter\def\csname SRef@push:B429\endcsname{\hyperlink{tag-030Y}{I.\ref*{push:B429}}}
\expandafter\def\csname SRef@push:src:B.4.30\endcsname{\hyperlink{tag-0310}{I.\ref*{push:src:B.4.30}}}
\expandafter\def\csname SRef@push:B430\endcsname{\hyperlink{tag-0310}{I.\ref*{push:B430}}}
\expandafter\def\csname SRef@push:src:B.4.31\endcsname{\hyperlink{tag-0311}{I.\ref*{push:src:B.4.31}}}
\expandafter\def\csname SRef@push:src:B.4.32\endcsname{\hyperlink{tag-0312}{I.\ref*{push:src:B.4.32}}}
\expandafter\def\csname SRef@push:B432\endcsname{\hyperlink{tag-0312}{I.\ref*{push:B432}}}
\expandafter\def\csname SRef@push:src:B.3.2\endcsname{\hyperlink{tag-0306}{I.\ref*{push:src:B.3.2}}}
\expandafter\def\csname SRef@push:B32\endcsname{\hyperlink{tag-0306}{I.\ref*{push:B32}}}
\expandafter\def\csname SRef@push:src:B.3.3\endcsname{\hyperlink{tag-0307}{I.\ref*{push:src:B.3.3}}}
\expandafter\def\csname SRef@push:B33\endcsname{\hyperlink{tag-0307}{I.\ref*{push:B33}}}
\expandafter\def\csname SRef@push:src:D.1.1\endcsname{\hyperlink{tag-032A}{I.\ref*{push:src:D.1.1}}}
\expandafter\def\csname SRef@push:src:D.4.1\endcsname{\hyperlink{tag-032J}{I.\ref*{push:src:D.4.1}}}
\expandafter\def\csname SRef@push:src:D.4.2\endcsname{\hyperlink{tag-032K}{I.\ref*{push:src:D.4.2}}}
\expandafter\def\csname SRef@push:src:D.4.3\endcsname{\hyperlink{tag-032L}{I.\ref*{push:src:D.4.3}}}
\expandafter\def\csname SRef@push:src:D.4.4\endcsname{\hyperlink{tag-032M}{I.\ref*{push:src:D.4.4}}}
\expandafter\def\csname SRef@push:src:D.4.5\endcsname{\hyperlink{tag-032N}{I.\ref*{push:src:D.4.5}}}
\expandafter\def\csname SRef@push:src:D.4.7\endcsname{\hyperlink{tag-032P}{I.\ref*{push:src:D.4.7}}}
\expandafter\def\csname SRef@push:src:D.4.8\endcsname{\hyperlink{tag-032Q}{I.\ref*{push:src:D.4.8}}}
\expandafter\def\csname SRef@push:src:D.2.2\endcsname{\hyperlink{tag-032E}{I.\ref*{push:src:D.2.2}}}
\expandafter\def\csname SRef@push:src:D.2.4\endcsname{\hyperlink{tag-032G}{I.\ref*{push:src:D.2.4}}}
\expandafter\def\csname SRef@inv:A3.2\endcsname{\hyperlink{tag-01RZ}{I.\ref*{inv:A3.2}}}
\expandafter\def\csname SRef@wb:lem:bcpr-eta-extension\endcsname{\hyperlink{tag-01RZ}{I.\ref*{wb:lem:bcpr-eta-extension}}}
\expandafter\def\csname SRef@inv:A4.1\endcsname{\hyperlink{tag-01S0}{I.\ref*{inv:A4.1}}}
\expandafter\def\csname SRef@inv:A4.2\endcsname{\hyperlink{tag-01S1}{I.\ref*{inv:A4.2}}}
\expandafter\def\csname SRef@inv:A1.1\endcsname{\hyperlink{tag-01RR}{I.\ref*{inv:A1.1}}}
\expandafter\def\csname SRef@inv:A1.2\endcsname{\hyperlink{tag-01RS}{I.\ref*{inv:A1.2}}}
\expandafter\def\csname SRef@wb:prop:log-relative-direct-images\endcsname{\hyperlink{tag-01RS}{I.\ref*{wb:prop:log-relative-direct-images}}}
\expandafter\def\csname SRef@inv:A1.3\endcsname{\hyperlink{tag-01RT}{I.\ref*{inv:A1.3}}}
\expandafter\def\csname SRef@wb:prop:absolute-cotangent-direct-images\endcsname{\hyperlink{tag-01RT}{I.\ref*{wb:prop:absolute-cotangent-direct-images}}}
\expandafter\def\csname SRef@inv:A2.1\endcsname{\hyperlink{tag-01RU}{I.\ref*{inv:A2.1}}}
\expandafter\def\csname SRef@wb:prop:log-ks-divisor\endcsname{\hyperlink{tag-01RU}{I.\ref*{wb:prop:log-ks-divisor}}}
\expandafter\def\csname SRef@wb:prop:period-derivative-determinant\endcsname{\hyperlink{tag-01RU}{I.\ref*{wb:prop:period-derivative-determinant}}}
\expandafter\def\csname SRef@inv:A2.2\endcsname{\hyperlink{tag-01RV}{I.\ref*{inv:A2.2}}}
\expandafter\def\csname SRef@inv:A2.3\endcsname{\hyperlink{tag-01RW}{I.\ref*{inv:A2.3}}}
\expandafter\def\csname SRef@inv:A2.4\endcsname{\hyperlink{tag-01RX}{I.\ref*{inv:A2.4}}}
\expandafter\def\csname SRef@wb:thm:h12-scalar\endcsname{\hyperlink{tag-01RX}{I.\ref*{wb:thm:h12-scalar}}}
\expandafter\def\csname SRef@inv:A3.1\endcsname{\hyperlink{tag-01RY}{I.\ref*{inv:A3.1}}}
\expandafter\def\csname SRef@wb:cor:exact-hodge-diamond\endcsname{\hyperlink{tag-01RY}{I.\ref*{wb:cor:exact-hodge-diamond}}}
\expandafter\def\csname SRef@wb:thm:parametric-hodge-diamond\endcsname{\hyperlink{tag-01RY}{I.\ref*{wb:thm:parametric-hodge-diamond}}}
\expandafter\def\csname SRef@wb:prop:workbench-holomorphic-form-vanishing\endcsname{\hyperlink{tag-01RY}{I.\ref*{wb:prop:workbench-holomorphic-form-vanishing}}}
\expandafter\def\csname SRef@inv:C1.1\endcsname{\hyperlink{tag-01SA}{I.\ref*{inv:C1.1}}}
\expandafter\def\csname SRef@wb:thm:frolicher-e2-candidate\endcsname{\hyperlink{tag-01SA}{I.\ref*{wb:thm:frolicher-e2-candidate}}}
\expandafter\def\csname SRef@inv:C3.7\endcsname{\hyperlink{tag-01SK}{I.\ref*{inv:C3.7}}}
\expandafter\def\csname SRef@inv:C1.2\endcsname{\hyperlink{tag-01SB}{I.\ref*{inv:C1.2}}}
\expandafter\def\csname SRef@inv:C2.1\endcsname{\hyperlink{tag-01SC}{I.\ref*{inv:C2.1}}}
\expandafter\def\csname SRef@wb:thm:exact-bott-chern-aeppli\endcsname{\hyperlink{tag-01SC}{I.\ref*{wb:thm:exact-bott-chern-aeppli}}}
\expandafter\def\csname SRef@inv:C2.2\endcsname{\hyperlink{tag-01SD}{I.\ref*{inv:C2.2}}}
\expandafter\def\csname SRef@inv:C3.1\endcsname{\hyperlink{tag-01SE}{I.\ref*{inv:C3.1}}}
\expandafter\def\csname SRef@wb:cor:picard-line-audit\endcsname{\hyperlink{tag-01SE}{I.\ref*{wb:cor:picard-line-audit}}}
\expandafter\def\csname SRef@inv:C3.2\endcsname{\hyperlink{tag-01SF}{I.\ref*{inv:C3.2}}}
\expandafter\def\csname SRef@wb:lem:bcpr-degree-one-input\endcsname{\hyperlink{tag-01SF}{I.\ref*{wb:lem:bcpr-degree-one-input}}}
\expandafter\def\csname SRef@inv:C3.3\endcsname{\hyperlink{tag-01SG}{I.\ref*{inv:C3.3}}}
\expandafter\def\csname SRef@wb:thm:bcpr-square-zero-algebra\endcsname{\hyperlink{tag-01SG}{I.\ref*{wb:thm:bcpr-square-zero-algebra}}}
\expandafter\def\csname SRef@wb:prop:bcpr-generators-are-bases\endcsname{\hyperlink{tag-01SG}{I.\ref*{wb:prop:bcpr-generators-are-bases}}}
\expandafter\def\csname SRef@inv:C3.4\endcsname{\hyperlink{tag-01SH}{I.\ref*{inv:C3.4}}}
\expandafter\def\csname SRef@inv:C3.5\endcsname{\hyperlink{tag-01SI}{I.\ref*{inv:C3.5}}}
\expandafter\def\csname SRef@inv:D4.1\endcsname{\hyperlink{tag-01SO}{I.\ref*{inv:D4.1}}}
\expandafter\def\csname SRef@inv:D4.2\endcsname{\hyperlink{tag-01SP}{I.\ref*{inv:D4.2}}}
\expandafter\def\csname SRef@inv:D4.3\endcsname{\hyperlink{tag-01SQ}{I.\ref*{inv:D4.3}}}
\expandafter\def\csname SRef@inv:D4.4\endcsname{\hyperlink{tag-01SR}{I.\ref*{inv:D4.4}}}
\expandafter\def\csname SRef@inv:D4.5\endcsname{\hyperlink{tag-01SS}{I.\ref*{inv:D4.5}}}
\expandafter\def\csname SRef@push:src:E.2.2\endcsname{\hyperlink{tag-0335}{I.\ref*{push:src:E.2.2}}}
\expandafter\def\csname SRef@push:src:E.2.3\endcsname{\hyperlink{tag-033B}{I.\ref*{push:src:E.2.3}}}
\expandafter\def\csname SRef@inv:B1.0\endcsname{\hyperlink{tag-01S2}{I.\ref*{inv:B1.0}}}
\expandafter\def\csname SRef@inv:B1.1\endcsname{\hyperlink{tag-01S3}{I.\ref*{inv:B1.1}}}
\expandafter\def\csname SRef@inv:B1.2\endcsname{\hyperlink{tag-01S4}{I.\ref*{inv:B1.2}}}
\expandafter\def\csname SRef@inv:B2.1\endcsname{\hyperlink{tag-01S5}{I.\ref*{inv:B2.1}}}
\expandafter\def\csname SRef@inv:B2.2\endcsname{\hyperlink{tag-01S6}{I.\ref*{inv:B2.2}}}
\expandafter\def\csname SRef@inv:B3.1\endcsname{\hyperlink{tag-01S7}{I.\ref*{inv:B3.1}}}
\expandafter\def\csname SRef@inv:B3.2\endcsname{\hyperlink{tag-01S8}{I.\ref*{inv:B3.2}}}
\expandafter\def\csname SRef@inv:B3.3\endcsname{\hyperlink{tag-01S9}{I.\ref*{inv:B3.3}}}
\expandafter\def\csname SRef@mapsI:src:A.1.1\endcsname{\hyperlink{tag-02PA}{I.\ref*{mapsI:src:A.1.1}}}
\expandafter\def\csname SRef@mapsI:quotient-eq-1\endcsname{\hyperlink{tag-02PA}{I.\ref*{mapsI:quotient-eq-1}}}
\expandafter\def\csname SRef@mapsI:src:A.1.2\endcsname{\hyperlink{tag-02PB}{I.\ref*{mapsI:src:A.1.2}}}
\expandafter\def\csname SRef@mapsI:src:A.1.3\endcsname{\hyperlink{tag-02PC}{I.\ref*{mapsI:src:A.1.3}}}
\expandafter\def\csname SRef@mapsI:quotient-eq-6\endcsname{\hyperlink{tag-02PC}{I.\ref*{mapsI:quotient-eq-6}}}
\expandafter\def\csname SRef@ari:C.2.1\endcsname{\hyperlink{tag-000L}{I.\ref*{ari:C.2.1}}}
\expandafter\def\csname SRef@ari:C.2.2\endcsname{\hyperlink{tag-000M}{I.\ref*{ari:C.2.2}}}
\expandafter\def\csname SRef@ari:C.3.1\endcsname{\hyperlink{tag-000N}{I.\ref*{ari:C.3.1}}}
\expandafter\def\csname SRef@ari:C.3.2\endcsname{\hyperlink{tag-000O}{I.\ref*{ari:C.3.2}}}
\expandafter\def\csname SRef@ari:C.6.1\endcsname{\hyperlink{tag-000V}{I.\ref*{ari:C.6.1}}}
\expandafter\def\csname SRef@ari:C.6.R5\endcsname{\hyperlink{tag-0014}{I.\ref*{ari:C.6.R5}}}
\expandafter\def\csname SRef@ari:C.6.T1\endcsname{\hyperlink{tag-0015}{I.\ref*{ari:C.6.T1}}}
\expandafter\def\csname SRef@ari:C.6.T2\endcsname{\hyperlink{tag-0016}{I.\ref*{ari:C.6.T2}}}
\expandafter\def\csname SRef@ari:C.8.6\endcsname{\hyperlink{tag-001T}{I.\ref*{ari:C.8.6}}}
\expandafter\def\csname SRef@ari:C.8.7\endcsname{\hyperlink{tag-001U}{I.\ref*{ari:C.8.7}}}
\expandafter\def\csname SRef@ari:C.8.8\endcsname{\hyperlink{tag-001V}{I.\ref*{ari:C.8.8}}}
\expandafter\def\csname SRef@ari:C.8.9\endcsname{\hyperlink{tag-001W}{I.\ref*{ari:C.8.9}}}
\expandafter\def\csname SRef@ari:C.8.10\endcsname{\hyperlink{tag-001K}{I.\ref*{ari:C.8.10}}}
\expandafter\def\csname SRef@ari:C.8.11\endcsname{\hyperlink{tag-001L}{I.\ref*{ari:C.8.11}}}
\expandafter\def\csname SRef@ari:C.6.2\endcsname{\hyperlink{tag-000W}{I.\ref*{ari:C.6.2}}}
\expandafter\def\csname SRef@ari:C.6.3\endcsname{\hyperlink{tag-000X}{I.\ref*{ari:C.6.3}}}
\expandafter\def\csname SRef@ari:C.6.4\endcsname{\hyperlink{tag-000Y}{I.\ref*{ari:C.6.4}}}
\expandafter\def\csname SRef@ari:C.6.5\endcsname{\hyperlink{tag-000Z}{I.\ref*{ari:C.6.5}}}
\expandafter\def\csname SRef@ari:C.6.R1\endcsname{\hyperlink{tag-0010}{I.\ref*{ari:C.6.R1}}}
\expandafter\def\csname SRef@ari:C.6.R2\endcsname{\hyperlink{tag-0011}{I.\ref*{ari:C.6.R2}}}
\expandafter\def\csname SRef@ari:C.6.R3\endcsname{\hyperlink{tag-0012}{I.\ref*{ari:C.6.R3}}}
\expandafter\def\csname SRef@ari:C.8.12\endcsname{\hyperlink{tag-001M}{I.\ref*{ari:C.8.12}}}
\expandafter\def\csname SRef@ari:C.8.13\endcsname{\hyperlink{tag-001N}{I.\ref*{ari:C.8.13}}}
\expandafter\def\csname SRef@ari:C.8.14\endcsname{\hyperlink{tag-001O}{I.\ref*{ari:C.8.14}}}
\expandafter\def\csname SRef@ari:C.8.15\endcsname{\hyperlink{tag-001P}{I.\ref*{ari:C.8.15}}}
\expandafter\def\csname SRef@fam:src:B.1.1\endcsname{\hyperlink{tag-01Q7}{I.\ref*{fam:src:B.1.1}}}
\expandafter\def\csname SRef@fam:st:B.1.1\endcsname{\hyperlink{tag-01Q7}{I.\ref*{fam:st:B.1.1}}}
\expandafter\def\csname SRef@fam:src:B.1.2\endcsname{\hyperlink{tag-01Q8}{I.\ref*{fam:src:B.1.2}}}
\expandafter\def\csname SRef@fam:st:B.1.2\endcsname{\hyperlink{tag-01Q8}{I.\ref*{fam:st:B.1.2}}}
\expandafter\def\csname SRef@fam:src:B.1.3\endcsname{\hyperlink{tag-01Q9}{I.\ref*{fam:src:B.1.3}}}
\expandafter\def\csname SRef@fam:st:B.1.3\endcsname{\hyperlink{tag-01Q9}{I.\ref*{fam:st:B.1.3}}}
\expandafter\def\csname SRef@fam:src:B.1.4\endcsname{\hyperlink{tag-01QA}{I.\ref*{fam:src:B.1.4}}}
\expandafter\def\csname SRef@fam:st:B.1.4\endcsname{\hyperlink{tag-01QA}{I.\ref*{fam:st:B.1.4}}}
\expandafter\def\csname SRef@fam:src:C.1.1\endcsname{\hyperlink{tag-01QH}{I.\ref*{fam:src:C.1.1}}}
\expandafter\def\csname SRef@fam:st:C.1.1\endcsname{\hyperlink{tag-01QH}{I.\ref*{fam:st:C.1.1}}}
\expandafter\def\csname SRef@fam:src:C.1.2\endcsname{\hyperlink{tag-01QI}{I.\ref*{fam:src:C.1.2}}}
\expandafter\def\csname SRef@fam:st:C.1.2\endcsname{\hyperlink{tag-01QI}{I.\ref*{fam:st:C.1.2}}}
\expandafter\def\csname SRef@fam:src:C.2.1\endcsname{\hyperlink{tag-01QJ}{I.\ref*{fam:src:C.2.1}}}
\expandafter\def\csname SRef@fam:st:C.2.1\endcsname{\hyperlink{tag-01QJ}{I.\ref*{fam:st:C.2.1}}}
\expandafter\def\csname SRef@fam:src:C.2.2\endcsname{\hyperlink{tag-01QK}{I.\ref*{fam:src:C.2.2}}}
\expandafter\def\csname SRef@fam:st:C.2.2\endcsname{\hyperlink{tag-01QK}{I.\ref*{fam:st:C.2.2}}}
\expandafter\def\csname SRef@ari:C.1.1\endcsname{\hyperlink{tag-000I}{I.\ref*{ari:C.1.1}}}
\expandafter\def\csname SRef@ari:C.1.2\endcsname{\hyperlink{tag-000J}{I.\ref*{ari:C.1.2}}}
\expandafter\def\csname SRef@ari:C.1.3\endcsname{\hyperlink{tag-000K}{I.\ref*{ari:C.1.3}}}
\expandafter\def\csname SRef@ari:C.4.2\endcsname{\hyperlink{tag-000Q}{I.\ref*{ari:C.4.2}}}
\expandafter\def\csname SRef@ari:C.4.3\endcsname{\hyperlink{tag-000R}{I.\ref*{ari:C.4.3}}}
\expandafter\def\csname SRef@ari:C.4.4\endcsname{\hyperlink{tag-000S}{I.\ref*{ari:C.4.4}}}
\expandafter\def\csname SRef@ari:C.4.5\endcsname{\hyperlink{tag-000T}{I.\ref*{ari:C.4.5}}}
\expandafter\def\csname SRef@ari:C.8.1\endcsname{\hyperlink{tag-001J}{I.\ref*{ari:C.8.1}}}
\expandafter\def\csname SRef@ari:C.8.2\endcsname{\hyperlink{tag-001Q}{I.\ref*{ari:C.8.2}}}
\expandafter\def\csname SRef@ari:C.8.3\endcsname{\hyperlink{tag-001R}{I.\ref*{ari:C.8.3}}}
\expandafter\def\csname SRef@ari:C.8.4\endcsname{\hyperlink{tag-001S}{I.\ref*{ari:C.8.4}}}
\expandafter\def\csname SRef@fam:src:A.1.1\endcsname{\hyperlink{tag-01PZ}{I.\ref*{fam:src:A.1.1}}}
\expandafter\def\csname SRef@fam:st:A.1.1\endcsname{\hyperlink{tag-01PZ}{I.\ref*{fam:st:A.1.1}}}
\expandafter\def\csname SRef@fam:src:A.1.2\endcsname{\hyperlink{tag-01Q0}{I.\ref*{fam:src:A.1.2}}}
\expandafter\def\csname SRef@fam:st:A.1.2\endcsname{\hyperlink{tag-01Q0}{I.\ref*{fam:st:A.1.2}}}
\expandafter\def\csname SRef@fam:src:A.2.1\endcsname{\hyperlink{tag-01Q1}{I.\ref*{fam:src:A.2.1}}}
\expandafter\def\csname SRef@fam:st:A.2.1\endcsname{\hyperlink{tag-01Q1}{I.\ref*{fam:st:A.2.1}}}
\expandafter\def\csname SRef@fam:src:A.3.1\endcsname{\hyperlink{tag-01Q2}{I.\ref*{fam:src:A.3.1}}}
\expandafter\def\csname SRef@fam:st:A.3.1\endcsname{\hyperlink{tag-01Q2}{I.\ref*{fam:st:A.3.1}}}
\expandafter\def\csname SRef@fam:src:A.4.1\endcsname{\hyperlink{tag-01Q3}{I.\ref*{fam:src:A.4.1}}}
\expandafter\def\csname SRef@fam:st:A.4.1\endcsname{\hyperlink{tag-01Q3}{I.\ref*{fam:st:A.4.1}}}
\expandafter\def\csname SRef@fam:src:A.4.2\endcsname{\hyperlink{tag-01Q4}{I.\ref*{fam:src:A.4.2}}}
\expandafter\def\csname SRef@fam:st:A.4.2\endcsname{\hyperlink{tag-01Q4}{I.\ref*{fam:st:A.4.2}}}
\expandafter\def\csname SRef@fam:src:A.4.3\endcsname{\hyperlink{tag-01Q5}{I.\ref*{fam:src:A.4.3}}}
\expandafter\def\csname SRef@fam:st:A.4.3\endcsname{\hyperlink{tag-01Q5}{I.\ref*{fam:st:A.4.3}}}
\expandafter\def\csname SRef@fam:src:B.2.1\endcsname{\hyperlink{tag-01QB}{I.\ref*{fam:src:B.2.1}}}
\expandafter\def\csname SRef@fam:st:B.2.1\endcsname{\hyperlink{tag-01QB}{I.\ref*{fam:st:B.2.1}}}
\expandafter\def\csname SRef@fam:src:B.2.2\endcsname{\hyperlink{tag-01QC}{I.\ref*{fam:src:B.2.2}}}
\expandafter\def\csname SRef@fam:st:B.2.2\endcsname{\hyperlink{tag-01QC}{I.\ref*{fam:st:B.2.2}}}
\expandafter\def\csname SRef@fam:src:B.2.3\endcsname{\hyperlink{tag-01QD}{I.\ref*{fam:src:B.2.3}}}
\expandafter\def\csname SRef@fam:st:B.2.3\endcsname{\hyperlink{tag-01QD}{I.\ref*{fam:st:B.2.3}}}
\expandafter\def\csname SRef@fam:src:B.3.1\endcsname{\hyperlink{tag-01QE}{I.\ref*{fam:src:B.3.1}}}
\expandafter\def\csname SRef@fam:st:B.3.1\endcsname{\hyperlink{tag-01QE}{I.\ref*{fam:st:B.3.1}}}
\expandafter\def\csname SRef@fam:src:B.4.1\endcsname{\hyperlink{tag-01QF}{I.\ref*{fam:src:B.4.1}}}
\expandafter\def\csname SRef@fam:st:B.4.1\endcsname{\hyperlink{tag-01QF}{I.\ref*{fam:st:B.4.1}}}
\expandafter\def\csname SRef@fam:src:B.4.2\endcsname{\hyperlink{tag-01QG}{I.\ref*{fam:src:B.4.2}}}
\expandafter\def\csname SRef@fam:st:B.4.2\endcsname{\hyperlink{tag-01QG}{I.\ref*{fam:st:B.4.2}}}
\expandafter\def\csname SRef@push:src:E.1.1\endcsname{\hyperlink{tag-032R}{I.\ref*{push:src:E.1.1}}}
\expandafter\def\csname SRef@inv:D1.1\endcsname{\hyperlink{tag-01SM}{I.\ref*{inv:D1.1}}}
\expandafter\def\csname SRef@inv:D1.2\endcsname{\hyperlink{tag-01SN}{I.\ref*{inv:D1.2}}}
\expandafter\def\csname SRef@ari:C.7.1\endcsname{\hyperlink{tag-0017}{I.\ref*{ari:C.7.1}}}
\expandafter\def\csname SRef@ari:C.7.11\endcsname{\hyperlink{tag-0019}{I.\ref*{ari:C.7.11}}}
\expandafter\def\csname SRef@ari:C.7.12\endcsname{\hyperlink{tag-001A}{I.\ref*{ari:C.7.12}}}
\expandafter\def\csname SRef@ari:C.7.2\endcsname{\hyperlink{tag-001B}{I.\ref*{ari:C.7.2}}}
\expandafter\def\csname SRef@ari:C.7.3\endcsname{\hyperlink{tag-001C}{I.\ref*{ari:C.7.3}}}
\expandafter\def\csname SRef@ari:C.7.4\endcsname{\hyperlink{tag-001D}{I.\ref*{ari:C.7.4}}}
\expandafter\def\csname SRef@ari:C.7.6\endcsname{\hyperlink{tag-001F}{I.\ref*{ari:C.7.6}}}
\expandafter\def\csname SRef@ari:C.7.9\endcsname{\hyperlink{tag-001I}{I.\ref*{ari:C.7.9}}}
\expandafter\def\csname SRef@ari:C.7.10\endcsname{\hyperlink{tag-0018}{I.\ref*{ari:C.7.10}}}
\expandafter\def\csname SRef@inv:C3.8\endcsname{\hyperlink{tag-01SL}{I.\ref*{inv:C3.8}}}
\expandafter\def\csname SRef@fam:src:C.3.3\endcsname{\hyperlink{tag-01QO}{I.\ref*{fam:src:C.3.3}}}
\expandafter\def\csname SRef@fam:st:C.3.3\endcsname{\hyperlink{tag-01QO}{I.\ref*{fam:st:C.3.3}}}
\expandafter\def\csname SRef@fam:src:E.1.5\endcsname{\hyperlink{tag-01R5}{I.\ref*{fam:src:E.1.5}}}
\expandafter\def\csname SRef@fam:st:E.1.5\endcsname{\hyperlink{tag-01R5}{I.\ref*{fam:st:E.1.5}}}
\expandafter\def\csname SRef@wb:rem:bcpr-literal-gauge\endcsname{\hyperlink{tag-03CT}{I.\ref*{wb:rem:bcpr-literal-gauge}}}
\expandafter\def\csname SRef@wb:rem:bcpr-global-gauge-boundary\endcsname{\hyperlink{tag-03CS}{I.\ref*{wb:rem:bcpr-global-gauge-boundary}}}
\expandafter\def\csname SRef@push:src:A.2.3\endcsname{\hyperlink{tag-02ZL}{I.\ref*{push:src:A.2.3}}}
\expandafter\def\csname SRef@push:A23\endcsname{\hyperlink{tag-02ZL}{I.\ref*{push:A23}}}
\expandafter\def\csname SRef@fam:section:F\endcsname{\hyperlink{tag-01PV}{I.\ref*{fam:section:F}}}
\expandafter\def\csname SRef@fam:section:D\endcsname{\hyperlink{tag-01PP}{I.\ref*{fam:section:D}}}
\expandafter\def\csname SRef@mapsI:real-forms\endcsname{\hyperlink{tag-02P9}{I.\ref*{mapsI:real-forms}}}
\expandafter\def\csname SRef@fam:section:E\endcsname{\hyperlink{tag-01PS}{I.\ref*{fam:section:E}}}
\expandafter\def\csname SRef@push:section:A\endcsname{\hyperlink{tag-02YZ}{I.\ref*{push:section:A}}}
\expandafter\def\csname SRef@push:section:B\endcsname{\hyperlink{tag-02Z4}{I.\ref*{push:section:B}}}
\expandafter\def\csname SRef@push:section:D\endcsname{\hyperlink{tag-02Z8}{I.\ref*{push:section:D}}}
\expandafter\def\csname SRef@inv:sec:partA\endcsname{\hyperlink{tag-01TH}{I.\ref*{inv:sec:partA}}}
\expandafter\def\csname SRef@inv:sec:partC\endcsname{\hyperlink{tag-01TI}{I.\ref*{inv:sec:partC}}}
\expandafter\def\csname SRef@def:common-conventions\endcsname{\hyperlink{tag-01O1}{I.\ref*{def:common-conventions}}}
\expandafter\def\csname SRef@mapsI:circle-quotient\endcsname{\hyperlink{tag-02OO}{I.\ref*{mapsI:circle-quotient}}}
\expandafter\def\csname SRef@inv:notation-dictionary\endcsname{\hyperlink{tag-01TG}{I.\ref*{inv:notation-dictionary}}}
\expandafter\def\csname SRef@fam:section:C\endcsname{\hyperlink{tag-01PK}{I.\ref*{fam:section:C}}}
\expandafter\def\csname SRef@fam:section:A\endcsname{\hyperlink{tag-01P9}{I.\ref*{fam:section:A}}}
\expandafter\def\csname SRef@fam:section:B\endcsname{\hyperlink{tag-01PF}{I.\ref*{fam:section:B}}}
\expandafter\def\csname SRef@push:section:E\endcsname{\hyperlink{tag-02ZD}{I.\ref*{push:section:E}}}
\expandafter\def\csname SRef@fam:section:F.1\endcsname{\hyperlink{tag-01PW}{I.\ref*{fam:section:F.1}}}
\expandafter\def\csname SRef@fam:section:F.2\endcsname{\hyperlink{tag-01PX}{I.\ref*{fam:section:F.2}}}
\expandafter\def\csname SRef@fam:section:F.3\endcsname{\hyperlink{tag-01PY}{I.\ref*{fam:section:F.3}}}
\expandafter\def\csname SRef@fam:section:D.1\endcsname{\hyperlink{tag-01PQ}{I.\ref*{fam:section:D.1}}}
\expandafter\def\csname SRef@fam:section:D.2\endcsname{\hyperlink{tag-01PR}{I.\ref*{fam:section:D.2}}}
\expandafter\def\csname SRef@fam:section:C.3\endcsname{\hyperlink{tag-01PN}{I.\ref*{fam:section:C.3}}}
\expandafter\def\csname SRef@fam:section:C.4\endcsname{\hyperlink{tag-01PO}{I.\ref*{fam:section:C.4}}}
\expandafter\def\csname SRef@fam:section:E.1\endcsname{\hyperlink{tag-01PT}{I.\ref*{fam:section:E.1}}}
\expandafter\def\csname SRef@fam:section:E.2\endcsname{\hyperlink{tag-01PU}{I.\ref*{fam:section:E.2}}}
\expandafter\def\csname SRef@push:section:A.1\endcsname{\hyperlink{tag-02Z0}{I.\ref*{push:section:A.1}}}
\expandafter\def\csname SRef@push:section:A.2\endcsname{\hyperlink{tag-02Z1}{I.\ref*{push:section:A.2}}}
\expandafter\def\csname SRef@push:section:A.3\endcsname{\hyperlink{tag-02Z2}{I.\ref*{push:section:A.3}}}
\expandafter\def\csname SRef@push:section:A.4\endcsname{\hyperlink{tag-02Z3}{I.\ref*{push:section:A.4}}}
\expandafter\def\csname SRef@push:section:B.1\endcsname{\hyperlink{tag-02Z5}{I.\ref*{push:section:B.1}}}
\expandafter\def\csname SRef@push:section:B.2\endcsname{\hyperlink{tag-02Z6}{I.\ref*{push:section:B.2}}}
\expandafter\def\csname SRef@push:section:B.3\endcsname{\hyperlink{tag-02Z7}{I.\ref*{push:section:B.3}}}
\expandafter\def\csname SRef@push:section:D.1\endcsname{\hyperlink{tag-02Z9}{I.\ref*{push:section:D.1}}}
\expandafter\def\csname SRef@push:section:D.2\endcsname{\hyperlink{tag-02ZA}{I.\ref*{push:section:D.2}}}
\expandafter\def\csname SRef@push:section:D.3\endcsname{\hyperlink{tag-02ZB}{I.\ref*{push:section:D.3}}}
\expandafter\def\csname SRef@push:section:D.4\endcsname{\hyperlink{tag-02ZC}{I.\ref*{push:section:D.4}}}
\expandafter\def\csname SRef@push:section:E.2\endcsname{\hyperlink{tag-02ZF}{I.\ref*{push:section:E.2}}}
\expandafter\def\csname SRef@push:section:E.3\endcsname{\hyperlink{tag-02ZG}{I.\ref*{push:section:E.3}}}
\expandafter\def\csname SRef@sec:e4\endcsname{\hyperlink{tag-034N}{I.\ref*{sec:e4}}}
\expandafter\def\csname SRef@push:section:E.4\endcsname{\hyperlink{tag-034N}{I.\ref*{push:section:E.4}}}
\expandafter\def\csname SRef@fam:section:B.1\endcsname{\hyperlink{tag-01PG}{I.\ref*{fam:section:B.1}}}
\expandafter\def\csname SRef@fam:section:C.1\endcsname{\hyperlink{tag-01PL}{I.\ref*{fam:section:C.1}}}
\expandafter\def\csname SRef@fam:section:C.2\endcsname{\hyperlink{tag-01PM}{I.\ref*{fam:section:C.2}}}
\expandafter\def\csname SRef@fam:section:A.1\endcsname{\hyperlink{tag-01PA}{I.\ref*{fam:section:A.1}}}
\expandafter\def\csname SRef@fam:section:A.2\endcsname{\hyperlink{tag-01PB}{I.\ref*{fam:section:A.2}}}
\expandafter\def\csname SRef@fam:section:A.3\endcsname{\hyperlink{tag-01PC}{I.\ref*{fam:section:A.3}}}
\expandafter\def\csname SRef@fam:section:A.4\endcsname{\hyperlink{tag-01PD}{I.\ref*{fam:section:A.4}}}
\expandafter\def\csname SRef@fam:section:A.5\endcsname{\hyperlink{tag-01PE}{I.\ref*{fam:section:A.5}}}
\expandafter\def\csname SRef@fam:section:B.2\endcsname{\hyperlink{tag-01PH}{I.\ref*{fam:section:B.2}}}
\expandafter\def\csname SRef@fam:section:B.3\endcsname{\hyperlink{tag-01PI}{I.\ref*{fam:section:B.3}}}
\expandafter\def\csname SRef@fam:section:B.4\endcsname{\hyperlink{tag-01PJ}{I.\ref*{fam:section:B.4}}}
\expandafter\def\csname SRef@push:section:E.1\endcsname{\hyperlink{tag-02ZE}{I.\ref*{push:section:E.1}}}
\expandafter\def\csname SRef@wb:sec:canonical-ring-recovery\endcsname{\hyperlink{tag-03CU}{I.\ref*{wb:sec:canonical-ring-recovery}}}
\expandafter\def\csname SRef@wb:sec:retained-period-deformation\endcsname{\hyperlink{tag-03CV}{I.\ref*{wb:sec:retained-period-deformation}}}
\expandafter\def\csname SRef@wb:pshift:sec:classification\endcsname{\hyperlink{tag-03CR}{I.\ref*{wb:pshift:sec:classification}}}
\expandafter\def\csname SRef@app:rederive:fam-src-A.5.1\endcsname{\hyperlink{tag-000C}{I.\ref*{app:rederive:fam-src-A.5.1}}}
\expandafter\def\csname SRef@app:rederive:fam-src-A.5.2\endcsname{\hyperlink{tag-000D}{I.\ref*{app:rederive:fam-src-A.5.2}}}
\expandafter\def\csname SRef@app:rederive:fam-src-B.4.1\endcsname{\hyperlink{tag-000E}{I.\ref*{app:rederive:fam-src-B.4.1}}}
\expandafter\def\csname SRef@fam:src:F.1.1\endcsname{\hyperlink{tag-01R9}{I.\ref*{fam:src:F.1.1}}}
\expandafter\def\csname SRef@fam:st:F.1.1\endcsname{\hyperlink{tag-01R9}{I.\ref*{fam:st:F.1.1}}}
\expandafter\def\csname SRef@cdp:st:D.1.14\endcsname{\hyperlink{tag-01R9}{I.\ref*{cdp:st:D.1.14}}}
\expandafter\def\csname SRef@fam:src:F.2.2\endcsname{\hyperlink{tag-01RC}{I.\ref*{fam:src:F.2.2}}}
\expandafter\def\csname SRef@fam:st:F.2.2\endcsname{\hyperlink{tag-01RC}{I.\ref*{fam:st:F.2.2}}}
\expandafter\def\csname SRef@fam:src:F.3.1\endcsname{\hyperlink{tag-01RE}{I.\ref*{fam:src:F.3.1}}}
\expandafter\def\csname SRef@fam:st:F.3.1\endcsname{\hyperlink{tag-01RE}{I.\ref*{fam:st:F.3.1}}}
\expandafter\def\csname SRef@fam:src:D.2.3\endcsname{\hyperlink{tag-01R0}{I.\ref*{fam:src:D.2.3}}}
\expandafter\def\csname SRef@fam:st:D.2.3\endcsname{\hyperlink{tag-01R0}{I.\ref*{fam:st:D.2.3}}}
\expandafter\def\csname SRef@fam:src:C.4.1\endcsname{\hyperlink{tag-01QQ}{I.\ref*{fam:src:C.4.1}}}
\expandafter\def\csname SRef@fam:st:C.4.1\endcsname{\hyperlink{tag-01QQ}{I.\ref*{fam:st:C.4.1}}}
\expandafter\def\csname SRef@fam:src:D.1.3\endcsname{\hyperlink{tag-01QT}{I.\ref*{fam:src:D.1.3}}}
\expandafter\def\csname SRef@fam:st:D.1.3\endcsname{\hyperlink{tag-01QT}{I.\ref*{fam:st:D.1.3}}}
\expandafter\def\csname SRef@fam:src:D.1.6\endcsname{\hyperlink{tag-01QW}{I.\ref*{fam:src:D.1.6}}}
\expandafter\def\csname SRef@fam:st:D.1.6\endcsname{\hyperlink{tag-01QW}{I.\ref*{fam:st:D.1.6}}}
\expandafter\def\csname SRef@mapsI:src:A.3.2\endcsname{\hyperlink{tag-02PF}{I.\ref*{mapsI:src:A.3.2}}}
\expandafter\def\csname SRef@mapsI:real-eq-3\endcsname{\hyperlink{tag-02PF}{I.\ref*{mapsI:real-eq-3}}}
\expandafter\def\csname SRef@mapsI:src:A.3.3\endcsname{\hyperlink{tag-02PG}{I.\ref*{mapsI:src:A.3.3}}}
\expandafter\def\csname SRef@push:src:A.4.7\endcsname{\hyperlink{tag-02ZX}{I.\ref*{push:src:A.4.7}}}
\expandafter\def\csname SRef@push:A47\endcsname{\hyperlink{tag-02ZX}{I.\ref*{push:A47}}}
\expandafter\def\csname SRef@push:src:A.4.8\endcsname{\hyperlink{tag-02ZY}{I.\ref*{push:src:A.4.8}}}
\expandafter\def\csname SRef@push:A48\endcsname{\hyperlink{tag-02ZY}{I.\ref*{push:A48}}}
\expandafter\def\csname SRef@push:src:A.1.2\endcsname{\hyperlink{tag-02ZI}{I.\ref*{push:src:A.1.2}}}
\expandafter\def\csname SRef@push:A12\endcsname{\hyperlink{tag-02ZI}{I.\ref*{push:A12}}}
\expandafter\def\csname SRef@push:src:A.2.2\endcsname{\hyperlink{tag-02ZK}{I.\ref*{push:src:A.2.2}}}
\expandafter\def\csname SRef@push:A22\endcsname{\hyperlink{tag-02ZK}{I.\ref*{push:A22}}}
\expandafter\def\csname SRef@push:src:A.2.4\endcsname{\hyperlink{tag-02ZM}{I.\ref*{push:src:A.2.4}}}
\expandafter\def\csname SRef@push:A24\endcsname{\hyperlink{tag-02ZM}{I.\ref*{push:A24}}}
\expandafter\def\csname SRef@push:src:A.3.2\endcsname{\hyperlink{tag-02ZP}{I.\ref*{push:src:A.3.2}}}
\expandafter\def\csname SRef@push:A32\endcsname{\hyperlink{tag-02ZP}{I.\ref*{push:A32}}}
\expandafter\def\csname SRef@push:src:B.4.2\endcsname{\hyperlink{tag-030O}{I.\ref*{push:src:B.4.2}}}
\expandafter\def\csname SRef@push:B42\endcsname{\hyperlink{tag-030O}{I.\ref*{push:B42}}}
\expandafter\def\csname SRef@push:src:B.4.5\endcsname{\hyperlink{tag-0315}{I.\ref*{push:src:B.4.5}}}
\expandafter\def\csname SRef@push:src:B.1.2\endcsname{\hyperlink{tag-0301}{I.\ref*{push:src:B.1.2}}}
\expandafter\def\csname SRef@push:B12\endcsname{\hyperlink{tag-0301}{I.\ref*{push:B12}}}
\expandafter\def\csname SRef@push:src:B.4.33\endcsname{\hyperlink{tag-0313}{I.\ref*{push:src:B.4.33}}}
\expandafter\def\csname SRef@push:B433\endcsname{\hyperlink{tag-0313}{I.\ref*{push:B433}}}
\expandafter\def\csname SRef@push:src:B.3.1\endcsname{\hyperlink{tag-0305}{I.\ref*{push:src:B.3.1}}}
\expandafter\def\csname SRef@push:B31\endcsname{\hyperlink{tag-0305}{I.\ref*{push:B31}}}
\expandafter\def\csname SRef@push:src:B.3.4\endcsname{\hyperlink{tag-0308}{I.\ref*{push:src:B.3.4}}}
\expandafter\def\csname SRef@push:B34\endcsname{\hyperlink{tag-0308}{I.\ref*{push:B34}}}
\expandafter\def\csname SRef@push:src:D.1.3\endcsname{\hyperlink{tag-032C}{I.\ref*{push:src:D.1.3}}}
\expandafter\def\csname SRef@push:src:D.2.3\endcsname{\hyperlink{tag-032F}{I.\ref*{push:src:D.2.3}}}
\expandafter\def\csname SRef@push:src:D.3.1\endcsname{\hyperlink{tag-032H}{I.\ref*{push:src:D.3.1}}}
\expandafter\def\csname SRef@night:P2.1\endcsname{\hyperlink{tag-02R9}{I.\ref*{night:P2.1}}}
\expandafter\def\csname SRef@stmt:P2.1\endcsname{\hyperlink{tag-02R9}{I.\ref*{stmt:P2.1}}}
\expandafter\def\csname SRef@night:P2.10\endcsname{\hyperlink{tag-02RA}{I.\ref*{night:P2.10}}}
\expandafter\def\csname SRef@stmt:P2.10\endcsname{\hyperlink{tag-02RA}{I.\ref*{stmt:P2.10}}}
\expandafter\def\csname SRef@night:P2.11\endcsname{\hyperlink{tag-02RB}{I.\ref*{night:P2.11}}}
\expandafter\def\csname SRef@stmt:P2.11\endcsname{\hyperlink{tag-02RB}{I.\ref*{stmt:P2.11}}}
\expandafter\def\csname SRef@night:P2.2\endcsname{\hyperlink{tag-02RC}{I.\ref*{night:P2.2}}}
\expandafter\def\csname SRef@stmt:P2.2\endcsname{\hyperlink{tag-02RC}{I.\ref*{stmt:P2.2}}}
\expandafter\def\csname SRef@night:P2.3\endcsname{\hyperlink{tag-02RD}{I.\ref*{night:P2.3}}}
\expandafter\def\csname SRef@stmt:P2.3\endcsname{\hyperlink{tag-02RD}{I.\ref*{stmt:P2.3}}}
\expandafter\def\csname SRef@night:P2.4\endcsname{\hyperlink{tag-02RE}{I.\ref*{night:P2.4}}}
\expandafter\def\csname SRef@stmt:P2.4\endcsname{\hyperlink{tag-02RE}{I.\ref*{stmt:P2.4}}}
\expandafter\def\csname SRef@night:P2.5\endcsname{\hyperlink{tag-02RF}{I.\ref*{night:P2.5}}}
\expandafter\def\csname SRef@stmt:P2.5\endcsname{\hyperlink{tag-02RF}{I.\ref*{stmt:P2.5}}}
\expandafter\def\csname SRef@night:P2.6\endcsname{\hyperlink{tag-02RG}{I.\ref*{night:P2.6}}}
\expandafter\def\csname SRef@stmt:P2.6\endcsname{\hyperlink{tag-02RG}{I.\ref*{stmt:P2.6}}}
\expandafter\def\csname SRef@night:P2.7\endcsname{\hyperlink{tag-02RH}{I.\ref*{night:P2.7}}}
\expandafter\def\csname SRef@stmt:P2.7\endcsname{\hyperlink{tag-02RH}{I.\ref*{stmt:P2.7}}}
\expandafter\def\csname SRef@night:P2.8\endcsname{\hyperlink{tag-02RI}{I.\ref*{night:P2.8}}}
\expandafter\def\csname SRef@stmt:P2.8\endcsname{\hyperlink{tag-02RI}{I.\ref*{stmt:P2.8}}}
\expandafter\def\csname SRef@night:P2.9\endcsname{\hyperlink{tag-02RJ}{I.\ref*{night:P2.9}}}
\expandafter\def\csname SRef@stmt:P2.9\endcsname{\hyperlink{tag-02RJ}{I.\ref*{stmt:P2.9}}}
\expandafter\def\csname SRef@push:src:E.2.1\endcsname{\hyperlink{tag-032U}{I.\ref*{push:src:E.2.1}}}
\expandafter\def\csname SRef@mapsI:src:A.2.1\endcsname{\hyperlink{tag-02PD}{I.\ref*{mapsI:src:A.2.1}}}
\expandafter\def\csname SRef@push:src:E.4.1\endcsname{\hyperlink{tag-033L}{I.\ref*{push:src:E.4.1}}}
\expandafter\def\csname SRef@push:src:E.4.2\endcsname{\hyperlink{tag-033M}{I.\ref*{push:src:E.4.2}}}
\expandafter\def\csname SRef@fam:src:C.2.3\endcsname{\hyperlink{tag-01QL}{I.\ref*{fam:src:C.2.3}}}
\expandafter\def\csname SRef@ari:C.8.5\endcsname{\hyperlink{tag-01QL}{I.\ref*{ari:C.8.5}}}
\expandafter\def\csname SRef@fam:st:C.2.3\endcsname{\hyperlink{tag-01QL}{I.\ref*{fam:st:C.2.3}}}
\expandafter\def\csname SRef@push:src:E.4.3\endcsname{\hyperlink{tag-01QL}{I.\ref*{push:src:E.4.3}}}
\expandafter\def\csname SRef@fam:src:A.5.1\endcsname{\hyperlink{tag-01Q6}{I.\ref*{fam:src:A.5.1}}}
\expandafter\def\csname SRef@wb:pshift:thm:classification\endcsname{\hyperlink{tag-01Q6}{I.\ref*{wb:pshift:thm:classification}}}
\expandafter\def\csname SRef@fam:st:A.5.1\endcsname{\hyperlink{tag-01Q6}{I.\ref*{fam:st:A.5.1}}}
\expandafter\def\csname SRef@fam:st:A.5.2\endcsname{\hyperlink{tag-01RH}{I.\ref*{fam:st:A.5.2}}}
\expandafter\def\csname SRef@fam:src:A.5.2\endcsname{\hyperlink{tag-01RH}{I.\ref*{fam:src:A.5.2}}}
\expandafter\def\csname SRef@inv:D1.3\endcsname{\hyperlink{tag-01RH}{I.\ref*{inv:D1.3}}}
\expandafter\def\csname SRef@inv:D2.1\endcsname{\hyperlink{tag-01RH}{I.\ref*{inv:D2.1}}}
\expandafter\def\csname SRef@inv:D3.1\endcsname{\hyperlink{tag-01RH}{I.\ref*{inv:D3.1}}}
\expandafter\def\csname SRef@cdp:st:E9\endcsname{\hyperlink{tag-01RH}{I.\ref*{cdp:st:E9}}}
\expandafter\def\csname SRef@push:src:E.1.2\endcsname{\hyperlink{tag-032S}{I.\ref*{push:src:E.1.2}}}
\expandafter\def\csname SRef@wb:thm:canonical-ring-recovery\endcsname{\hyperlink{tag-03DM}{I.\ref*{wb:thm:canonical-ring-recovery}}}
\expandafter\def\csname SRef@wb:thm:retained-period-deformation\endcsname{\hyperlink{tag-03DP}{I.\ref*{wb:thm:retained-period-deformation}}}
\expandafter\def\csname SRef@night:P5.1\endcsname{\hyperlink{tag-02YM}{I.\ref*{night:P5.1}}}
\expandafter\def\csname SRef@stmt:P5.1\endcsname{\hyperlink{tag-02YM}{I.\ref*{stmt:P5.1}}}
\expandafter\def\csname SRef@night:P5.2\endcsname{\hyperlink{tag-02YN}{I.\ref*{night:P5.2}}}
\expandafter\def\csname SRef@stmt:P5.2\endcsname{\hyperlink{tag-02YN}{I.\ref*{stmt:P5.2}}}
\expandafter\def\csname SRef@night:P5.3\endcsname{\hyperlink{tag-02YO}{I.\ref*{night:P5.3}}}
\expandafter\def\csname SRef@stmt:P5.3\endcsname{\hyperlink{tag-02YO}{I.\ref*{stmt:P5.3}}}
\expandafter\def\csname SRef@night:P5.4\endcsname{\hyperlink{tag-02YP}{I.\ref*{night:P5.4}}}
\expandafter\def\csname SRef@stmt:P5.4\endcsname{\hyperlink{tag-02YP}{I.\ref*{stmt:P5.4}}}
\expandafter\def\csname SRef@night:P5.5\endcsname{\hyperlink{tag-02YQ}{I.\ref*{night:P5.5}}}
\expandafter\def\csname SRef@stmt:P5.5\endcsname{\hyperlink{tag-02YQ}{I.\ref*{stmt:P5.5}}}
\expandafter\def\csname SRef@night:P5.6\endcsname{\hyperlink{tag-02YR}{I.\ref*{night:P5.6}}}
\expandafter\def\csname SRef@stmt:P5.6\endcsname{\hyperlink{tag-02YR}{I.\ref*{stmt:P5.6}}}
\expandafter\def\csname SRef@cdp:eq:A.RC.2\endcsname{\href{S6_SERIES_GRIND_II.pdf\#nameddest=tag-006V}{II.{A.RC.2} (p.~78)} }
\expandafter\def\csname SRef@cdp:eq:A.RC.3\endcsname{\href{S6_SERIES_GRIND_II.pdf\#nameddest=tag-006W}{II.{A.RC.3} (p.~78)} }
\expandafter\def\csname SRef@cdp:eq:A.RC.4\endcsname{\href{S6_SERIES_GRIND_II.pdf\#nameddest=tag-006X}{II.{A.RC.4} (p.~79)} }
\expandafter\def\csname SRef@cdp:eq:A.RC.5\endcsname{\href{S6_SERIES_GRIND_II.pdf\#nameddest=tag-006Y}{II.{A.RC.5} (p.~79)} }
\expandafter\def\csname SRef@cdp:eq:A.RC.6\endcsname{\href{S6_SERIES_GRIND_II.pdf\#nameddest=tag-006Z}{II.{A.RC.6} (p.~79)} }
\expandafter\def\csname SRef@cdp:eq:A.RC.7\endcsname{\href{S6_SERIES_GRIND_II.pdf\#nameddest=tag-0070}{II.{A.RC.7} (p.~80)} }
\expandafter\def\csname SRef@cdp:eq:A.Pi.1\endcsname{\href{S6_SERIES_GRIND_II.pdf\#nameddest=tag-006O}{II.{A.Pi.1} (p.~81)} }
\expandafter\def\csname SRef@cdp:eq:A.Pi.2\endcsname{\href{S6_SERIES_GRIND_II.pdf\#nameddest=tag-006P}{II.{A.Pi.2} (p.~81)} }
\expandafter\def\csname SRef@cdp:eq:A.Pi.3\endcsname{\href{S6_SERIES_GRIND_II.pdf\#nameddest=tag-006Q}{II.{A.Pi.3} (p.~81)} }
\expandafter\def\csname SRef@cdp:eq:A.Pi.4\endcsname{\href{S6_SERIES_GRIND_II.pdf\#nameddest=tag-006R}{II.{A.Pi.4} (p.~82)} }
\expandafter\def\csname SRef@cdp:eq:A.CD.1\endcsname{\href{S6_SERIES_GRIND_II.pdf\#nameddest=tag-0058}{II.{A.CD.1} (p.~82)} }
\expandafter\def\csname SRef@cdp:eq:A.CD.2\endcsname{\href{S6_SERIES_GRIND_II.pdf\#nameddest=tag-0059}{II.{A.CD.2} (p.~82)} }
\expandafter\def\csname SRef@cdp:eq:A.CD.3\endcsname{\href{S6_SERIES_GRIND_II.pdf\#nameddest=tag-005A}{II.{A.CD.3} (p.~83)} }
\expandafter\def\csname SRef@cdp:eq:A.CD.4\endcsname{\href{S6_SERIES_GRIND_II.pdf\#nameddest=tag-005B}{II.{A.CD.4} (p.~83)} }
\expandafter\def\csname SRef@cdp:eq:A.Q1.1\endcsname{\href{S6_SERIES_GRIND_II.pdf\#nameddest=tag-006S}{II.{A.Q1.1} (p.~83)} }
\expandafter\def\csname SRef@cdp:eq:A.Q1.2\endcsname{\href{S6_SERIES_GRIND_II.pdf\#nameddest=tag-006T}{II.{A.Q1.2} (p.~84)} }
\expandafter\def\csname SRef@cdp:eq:A.J.1\endcsname{\href{S6_SERIES_GRIND_II.pdf\#nameddest=tag-005C}{II.{A.J.1} (p.~84)} }
\expandafter\def\csname SRef@cdp:eq:A.J.2\endcsname{\href{S6_SERIES_GRIND_II.pdf\#nameddest=tag-005N}{II.{A.J.2} (p.~85)} }
\expandafter\def\csname SRef@cdp:eq:A.J.3\endcsname{\href{S6_SERIES_GRIND_II.pdf\#nameddest=tag-005Y}{II.{A.J.3} (p.~85)} }
\expandafter\def\csname SRef@cdp:eq:A.J.4\endcsname{\href{S6_SERIES_GRIND_II.pdf\#nameddest=tag-006C}{II.{A.J.4} (p.~85)} }
\expandafter\def\csname SRef@cdp:eq:A.J.5\endcsname{\href{S6_SERIES_GRIND_II.pdf\#nameddest=tag-006E}{II.{A.J.5} (p.~86)} }
\expandafter\def\csname SRef@cdp:eq:A.J.6\endcsname{\href{S6_SERIES_GRIND_II.pdf\#nameddest=tag-006F}{II.{A.J.6} (p.~86)} }
\expandafter\def\csname SRef@cdp:eq:A.J.7\endcsname{\href{S6_SERIES_GRIND_II.pdf\#nameddest=tag-006G}{II.{A.J.7} (p.~86)} }
\expandafter\def\csname SRef@cdp:eq:A.J.8\endcsname{\href{S6_SERIES_GRIND_II.pdf\#nameddest=tag-006H}{II.{A.J.8} (p.~86)} }
\expandafter\def\csname SRef@cdp:eq:A.J.9\endcsname{\href{S6_SERIES_GRIND_II.pdf\#nameddest=tag-006I}{II.{A.J.9} (p.~86)} }
\expandafter\def\csname SRef@cdp:eq:A.J.10\endcsname{\href{S6_SERIES_GRIND_II.pdf\#nameddest=tag-005D}{II.{A.J.10} (p.~87)} }
\expandafter\def\csname SRef@cdp:eq:A.J.11\endcsname{\href{S6_SERIES_GRIND_II.pdf\#nameddest=tag-005E}{II.{A.J.11} (p.~88)} }
\expandafter\def\csname SRef@cdp:eq:A.J.12\endcsname{\href{S6_SERIES_GRIND_II.pdf\#nameddest=tag-005F}{II.{A.J.12} (p.~88)} }
\expandafter\def\csname SRef@cdp:eq:A.J.13\endcsname{\href{S6_SERIES_GRIND_II.pdf\#nameddest=tag-005G}{II.{A.J.13} (p.~88)} }
\expandafter\def\csname SRef@cdp:eq:A.J.14\endcsname{\href{S6_SERIES_GRIND_II.pdf\#nameddest=tag-005H}{II.{A.J.14} (p.~89)} }
\expandafter\def\csname SRef@cdp:eq:A.J.15\endcsname{\href{S6_SERIES_GRIND_II.pdf\#nameddest=tag-005I}{II.{A.J.15} (p.~89)} }
\expandafter\def\csname SRef@cdp:eq:A.J.16\endcsname{\href{S6_SERIES_GRIND_II.pdf\#nameddest=tag-005J}{II.{A.J.16} (p.~89)} }
\expandafter\def\csname SRef@cdp:eq:A.J.17\endcsname{\href{S6_SERIES_GRIND_II.pdf\#nameddest=tag-005K}{II.{A.J.17} (p.~89)} }
\expandafter\def\csname SRef@cdp:eq:A.J.18\endcsname{\href{S6_SERIES_GRIND_II.pdf\#nameddest=tag-005L}{II.{A.J.18} (p.~89)} }
\expandafter\def\csname SRef@cdp:eq:A.J.19\endcsname{\href{S6_SERIES_GRIND_II.pdf\#nameddest=tag-005M}{II.{A.J.19} (p.~90)} }
\expandafter\def\csname SRef@cdp:eq:A.J.20\endcsname{\href{S6_SERIES_GRIND_II.pdf\#nameddest=tag-005O}{II.{A.J.20} (p.~90)} }
\expandafter\def\csname SRef@cdp:eq:A.J.30\endcsname{\href{S6_SERIES_GRIND_II.pdf\#nameddest=tag-005Z}{II.{A.J.30} (p.~90)} }
\expandafter\def\csname SRef@cdp:eq:A.J.31\endcsname{\href{S6_SERIES_GRIND_II.pdf\#nameddest=tag-0060}{II.{A.J.31} (p.~91)} }
\expandafter\def\csname SRef@cdp:eq:A.J.32\endcsname{\href{S6_SERIES_GRIND_II.pdf\#nameddest=tag-0061}{II.{A.J.32} (p.~91)} }
\expandafter\def\csname SRef@cdp:eq:A.J.21\endcsname{\href{S6_SERIES_GRIND_II.pdf\#nameddest=tag-005P}{II.{A.J.21} (p.~92)} }
\expandafter\def\csname SRef@cdp:eq:A.J.22\endcsname{\href{S6_SERIES_GRIND_II.pdf\#nameddest=tag-005Q}{II.{A.J.22} (p.~92)} }
\expandafter\def\csname SRef@cdp:eq:A.J.23\endcsname{\href{S6_SERIES_GRIND_II.pdf\#nameddest=tag-005R}{II.{A.J.23} (p.~92)} }
\expandafter\def\csname SRef@cdp:eq:A.J.24\endcsname{\href{S6_SERIES_GRIND_II.pdf\#nameddest=tag-005S}{II.{A.J.24} (p.~93)} }
\expandafter\def\csname SRef@cdp:eq:A.J.25\endcsname{\href{S6_SERIES_GRIND_II.pdf\#nameddest=tag-005T}{II.{A.J.25} (p.~93)} }
\expandafter\def\csname SRef@cdp:eq:A.J.26\endcsname{\href{S6_SERIES_GRIND_II.pdf\#nameddest=tag-005U}{II.{A.J.26} (p.~93)} }
\expandafter\def\csname SRef@cdp:eq:A.J.27\endcsname{\href{S6_SERIES_GRIND_II.pdf\#nameddest=tag-005V}{II.{A.J.27} (p.~93)} }
\expandafter\def\csname SRef@cdp:eq:A.J.28\endcsname{\href{S6_SERIES_GRIND_II.pdf\#nameddest=tag-005W}{II.{A.J.28} (p.~94)} }
\expandafter\def\csname SRef@cdp:eq:A.J.29\endcsname{\href{S6_SERIES_GRIND_II.pdf\#nameddest=tag-005X}{II.{A.J.29} (p.~94)} }
\expandafter\def\csname SRef@cdp:eq:A.J.33\endcsname{\href{S6_SERIES_GRIND_II.pdf\#nameddest=tag-0062}{II.{A.J.33} (p.~95)} }
\expandafter\def\csname SRef@cdp:eq:A.J.34\endcsname{\href{S6_SERIES_GRIND_II.pdf\#nameddest=tag-0065}{II.{A.J.34} (p.~96)} }
\expandafter\def\csname SRef@cdp:eq:A.J.33a\endcsname{\href{S6_SERIES_GRIND_II.pdf\#nameddest=tag-0063}{II.{A.J.33a} (p.~96)} }
\expandafter\def\csname SRef@cdp:eq:A.J.33b\endcsname{\href{S6_SERIES_GRIND_II.pdf\#nameddest=tag-0064}{II.{A.J.33b} (p.~97)} }
\expandafter\def\csname SRef@cdp:eq:A.J.35\endcsname{\href{S6_SERIES_GRIND_II.pdf\#nameddest=tag-0066}{II.{A.J.35} (p.~97)} }
\expandafter\def\csname SRef@cdp:eq:A.J.36\endcsname{\href{S6_SERIES_GRIND_II.pdf\#nameddest=tag-0067}{II.{A.J.36} (p.~97)} }
\expandafter\def\csname SRef@cdp:eq:A.J.37\endcsname{\href{S6_SERIES_GRIND_II.pdf\#nameddest=tag-0068}{II.{A.J.37} (p.~97)} }
\expandafter\def\csname SRef@cdp:eq:A.J.38\endcsname{\href{S6_SERIES_GRIND_II.pdf\#nameddest=tag-0069}{II.{A.J.38} (p.~98)} }
\expandafter\def\csname SRef@cdp:eq:A.J.38a\endcsname{\href{S6_SERIES_GRIND_II.pdf\#nameddest=tag-006A}{II.{A.J.38a} (p.~98)} }
\expandafter\def\csname SRef@cdp:eq:A.J.39\endcsname{\href{S6_SERIES_GRIND_II.pdf\#nameddest=tag-006B}{II.{A.J.39} (p.~99)} }
\expandafter\def\csname SRef@cdp:eq:A.J.40\endcsname{\href{S6_SERIES_GRIND_II.pdf\#nameddest=tag-006D}{II.{A.J.40} (p.~100)} }
\expandafter\def\csname SRef@cdp:eq:A.T.1\endcsname{\href{S6_SERIES_GRIND_II.pdf\#nameddest=tag-007M}{II.{A.T.1} (p.~101)} }
\expandafter\def\csname SRef@cdp:eq:A.T.2\endcsname{\href{S6_SERIES_GRIND_II.pdf\#nameddest=tag-007X}{II.{A.T.2} (p.~102)} }
\expandafter\def\csname SRef@cdp:eq:A.T.3\endcsname{\href{S6_SERIES_GRIND_II.pdf\#nameddest=tag-0082}{II.{A.T.3} (p.~102)} }
\expandafter\def\csname SRef@cdp:eq:A.T.4\endcsname{\href{S6_SERIES_GRIND_II.pdf\#nameddest=tag-0083}{II.{A.T.4} (p.~102)} }
\expandafter\def\csname SRef@cdp:eq:A.T.5\endcsname{\href{S6_SERIES_GRIND_II.pdf\#nameddest=tag-0084}{II.{A.T.5} (p.~103)} }
\expandafter\def\csname SRef@cdp:eq:A.T.6\endcsname{\href{S6_SERIES_GRIND_II.pdf\#nameddest=tag-0085}{II.{A.T.6} (p.~103)} }
\expandafter\def\csname SRef@cdp:eq:A.T.7\endcsname{\href{S6_SERIES_GRIND_II.pdf\#nameddest=tag-0086}{II.{A.T.7} (p.~103)} }
\expandafter\def\csname SRef@cdp:eq:A.T.8\endcsname{\href{S6_SERIES_GRIND_II.pdf\#nameddest=tag-0087}{II.{A.T.8} (p.~103)} }
\expandafter\def\csname SRef@cdp:eq:A.T.9\endcsname{\href{S6_SERIES_GRIND_II.pdf\#nameddest=tag-0088}{II.{A.T.9} (p.~103)} }
\expandafter\def\csname SRef@cdp:eq:A.T.10\endcsname{\href{S6_SERIES_GRIND_II.pdf\#nameddest=tag-007N}{II.{A.T.10} (p.~103)} }
\expandafter\def\csname SRef@cdp:eq:A.T.11\endcsname{\href{S6_SERIES_GRIND_II.pdf\#nameddest=tag-007O}{II.{A.T.11} (p.~104)} }
\expandafter\def\csname SRef@cdp:eq:A.T.12\endcsname{\href{S6_SERIES_GRIND_II.pdf\#nameddest=tag-007P}{II.{A.T.12} (p.~104)} }
\expandafter\def\csname SRef@cdp:eq:A.T.13\endcsname{\href{S6_SERIES_GRIND_II.pdf\#nameddest=tag-007Q}{II.{A.T.13} (p.~104)} }
\expandafter\def\csname SRef@cdp:eq:A.T.14\endcsname{\href{S6_SERIES_GRIND_II.pdf\#nameddest=tag-007R}{II.{A.T.14} (p.~105)} }
\expandafter\def\csname SRef@cdp:eq:A.T.15\endcsname{\href{S6_SERIES_GRIND_II.pdf\#nameddest=tag-007S}{II.{A.T.15} (p.~105)} }
\expandafter\def\csname SRef@cdp:eq:A.T.16\endcsname{\href{S6_SERIES_GRIND_II.pdf\#nameddest=tag-007T}{II.{A.T.16} (p.~106)} }
\expandafter\def\csname SRef@cdp:eq:A.T.17\endcsname{\href{S6_SERIES_GRIND_II.pdf\#nameddest=tag-007U}{II.{A.T.17} (p.~106)} }
\expandafter\def\csname SRef@cdp:eq:A.T.18\endcsname{\href{S6_SERIES_GRIND_II.pdf\#nameddest=tag-007V}{II.{A.T.18} (p.~106)} }
\expandafter\def\csname SRef@cdp:eq:A.T.19\endcsname{\href{S6_SERIES_GRIND_II.pdf\#nameddest=tag-007W}{II.{A.T.19} (p.~106)} }
\expandafter\def\csname SRef@cdp:eq:A.T.20\endcsname{\href{S6_SERIES_GRIND_II.pdf\#nameddest=tag-007Y}{II.{A.T.20} (p.~106)} }
\expandafter\def\csname SRef@cdp:eq:A.T.21\endcsname{\href{S6_SERIES_GRIND_II.pdf\#nameddest=tag-007Z}{II.{A.T.21} (p.~107)} }
\expandafter\def\csname SRef@cdp:eq:A.T.22\endcsname{\href{S6_SERIES_GRIND_II.pdf\#nameddest=tag-0080}{II.{A.T.22} (p.~108)} }
\expandafter\def\csname SRef@cdp:eq:A.T.23\endcsname{\href{S6_SERIES_GRIND_II.pdf\#nameddest=tag-0081}{II.{A.T.23} (p.~108)} }
\expandafter\def\csname SRef@cdp:eq:A.U.1\endcsname{\href{S6_SERIES_GRIND_II.pdf\#nameddest=tag-0089}{II.{A.U.1} (p.~109)} }
\expandafter\def\csname SRef@cdp:eq:A.U.2\endcsname{\href{S6_SERIES_GRIND_II.pdf\#nameddest=tag-008A}{II.{A.U.2} (p.~109)} }
\expandafter\def\csname SRef@cdp:eq:A.U.3\endcsname{\href{S6_SERIES_GRIND_II.pdf\#nameddest=tag-008B}{II.{A.U.3} (p.~110)} }
\expandafter\def\csname SRef@cdp:eq:A.U.4\endcsname{\href{S6_SERIES_GRIND_II.pdf\#nameddest=tag-008C}{II.{A.U.4} (p.~110)} }
\expandafter\def\csname SRef@cdp:eq:A.U.5\endcsname{\href{S6_SERIES_GRIND_II.pdf\#nameddest=tag-008D}{II.{A.U.5} (p.~110)} }
\expandafter\def\csname SRef@cdp:eq:A.U.6\endcsname{\href{S6_SERIES_GRIND_II.pdf\#nameddest=tag-008E}{II.{A.U.6} (p.~111)} }
\expandafter\def\csname SRef@cdp:eq:A.U.7\endcsname{\href{S6_SERIES_GRIND_II.pdf\#nameddest=tag-008F}{II.{A.U.7} (p.~111)} }
\expandafter\def\csname SRef@cdp:eq:A.RT.1\endcsname{\href{S6_SERIES_GRIND_II.pdf\#nameddest=tag-0071}{II.{A.RT.1} (p.~112)} }
\expandafter\def\csname SRef@cdp:eq:A.RT.2\endcsname{\href{S6_SERIES_GRIND_II.pdf\#nameddest=tag-007C}{II.{A.RT.2} (p.~112)} }
\expandafter\def\csname SRef@cdp:eq:A.RT.3\endcsname{\href{S6_SERIES_GRIND_II.pdf\#nameddest=tag-007F}{II.{A.RT.3} (p.~112)} }
\expandafter\def\csname SRef@cdp:eq:A.RT.4\endcsname{\href{S6_SERIES_GRIND_II.pdf\#nameddest=tag-007G}{II.{A.RT.4} (p.~112)} }
\expandafter\def\csname SRef@cdp:eq:A.RT.5\endcsname{\href{S6_SERIES_GRIND_II.pdf\#nameddest=tag-007H}{II.{A.RT.5} (p.~112)} }
\expandafter\def\csname SRef@cdp:eq:A.RT.6\endcsname{\href{S6_SERIES_GRIND_II.pdf\#nameddest=tag-007I}{II.{A.RT.6} (p.~112)} }
\expandafter\def\csname SRef@cdp:eq:A.RT.7\endcsname{\href{S6_SERIES_GRIND_II.pdf\#nameddest=tag-007J}{II.{A.RT.7} (p.~113)} }
\expandafter\def\csname SRef@cdp:eq:A.RT.8\endcsname{\href{S6_SERIES_GRIND_II.pdf\#nameddest=tag-007K}{II.{A.RT.8} (p.~113)} }
\expandafter\def\csname SRef@cdp:eq:A.RT.9\endcsname{\href{S6_SERIES_GRIND_II.pdf\#nameddest=tag-007L}{II.{A.RT.9} (p.~113)} }
\expandafter\def\csname SRef@cdp:eq:A.RT.10\endcsname{\href{S6_SERIES_GRIND_II.pdf\#nameddest=tag-0072}{II.{A.RT.10} (p.~113)} }
\expandafter\def\csname SRef@cdp:eq:A.RT.11\endcsname{\href{S6_SERIES_GRIND_II.pdf\#nameddest=tag-0073}{II.{A.RT.11} (p.~115)} }
\expandafter\def\csname SRef@cdp:eq:A.RT.12\endcsname{\href{S6_SERIES_GRIND_II.pdf\#nameddest=tag-0074}{II.{A.RT.12} (p.~115)} }
\expandafter\def\csname SRef@cdp:eq:A.RT.13\endcsname{\href{S6_SERIES_GRIND_II.pdf\#nameddest=tag-0075}{II.{A.RT.13} (p.~115)} }
\expandafter\def\csname SRef@cdp:eq:A.RT.14\endcsname{\href{S6_SERIES_GRIND_II.pdf\#nameddest=tag-0076}{II.{A.RT.14} (p.~115)} }
\expandafter\def\csname SRef@cdp:eq:A.RT.15\endcsname{\href{S6_SERIES_GRIND_II.pdf\#nameddest=tag-0077}{II.{A.RT.15} (p.~116)} }
\expandafter\def\csname SRef@cdp:eq:A.RT.16\endcsname{\href{S6_SERIES_GRIND_II.pdf\#nameddest=tag-0078}{II.{A.RT.16} (p.~116)} }
\expandafter\def\csname SRef@cdp:eq:A.RT.17\endcsname{\href{S6_SERIES_GRIND_II.pdf\#nameddest=tag-0079}{II.{A.RT.17} (p.~116)} }
\expandafter\def\csname SRef@cdp:eq:A.RT.18\endcsname{\href{S6_SERIES_GRIND_II.pdf\#nameddest=tag-007A}{II.{A.RT.18} (p.~116)} }
\expandafter\def\csname SRef@cdp:eq:A.RT.19\endcsname{\href{S6_SERIES_GRIND_II.pdf\#nameddest=tag-007B}{II.{A.RT.19} (p.~116)} }
\expandafter\def\csname SRef@cdp:eq:A.RT.20\endcsname{\href{S6_SERIES_GRIND_II.pdf\#nameddest=tag-007D}{II.{A.RT.20} (p.~117)} }
\expandafter\def\csname SRef@cdp:eq:A.RT.21\endcsname{\href{S6_SERIES_GRIND_II.pdf\#nameddest=tag-007E}{II.{A.RT.21} (p.~117)} }
\expandafter\def\csname SRef@cdp:eq:A.BP.1\endcsname{\href{S6_SERIES_GRIND_II.pdf\#nameddest=tag-004W}{II.{A.BP.1} (p.~118)} }
\expandafter\def\csname SRef@cdp:eq:A.BP.4\endcsname{\href{S6_SERIES_GRIND_II.pdf\#nameddest=tag-0052}{II.{A.BP.4} (p.~118)} }
\expandafter\def\csname SRef@cdp:eq:A.BP.2\endcsname{\href{S6_SERIES_GRIND_II.pdf\#nameddest=tag-0050}{II.{A.BP.2} (p.~118)} }
\expandafter\def\csname SRef@cdp:eq:A.BP.3\endcsname{\href{S6_SERIES_GRIND_II.pdf\#nameddest=tag-0051}{II.{A.BP.3} (p.~118)} }
\expandafter\def\csname SRef@cdp:eq:A.BP.5\endcsname{\href{S6_SERIES_GRIND_II.pdf\#nameddest=tag-0053}{II.{A.BP.5} (p.~119)} }
\expandafter\def\csname SRef@cdp:eq:A.BP.6\endcsname{\href{S6_SERIES_GRIND_II.pdf\#nameddest=tag-0054}{II.{A.BP.6} (p.~119)} }
\expandafter\def\csname SRef@cdp:eq:A.BP.7\endcsname{\href{S6_SERIES_GRIND_II.pdf\#nameddest=tag-0055}{II.{A.BP.7} (p.~119)} }
\expandafter\def\csname SRef@cdp:eq:A.BP.8\endcsname{\href{S6_SERIES_GRIND_II.pdf\#nameddest=tag-0056}{II.{A.BP.8} (p.~119)} }
\expandafter\def\csname SRef@cdp:eq:A.BP.9\endcsname{\href{S6_SERIES_GRIND_II.pdf\#nameddest=tag-0057}{II.{A.BP.9} (p.~120)} }
\expandafter\def\csname SRef@cdp:eq:A.BP.10\endcsname{\href{S6_SERIES_GRIND_II.pdf\#nameddest=tag-004X}{II.{A.BP.10} (p.~120)} }
\expandafter\def\csname SRef@cdp:eq:A.BP.11\endcsname{\href{S6_SERIES_GRIND_II.pdf\#nameddest=tag-004Y}{II.{A.BP.11} (p.~120)} }
\expandafter\def\csname SRef@cdp:eq:A.BP.12\endcsname{\href{S6_SERIES_GRIND_II.pdf\#nameddest=tag-004Z}{II.{A.BP.12} (p.~120)} }
\expandafter\def\csname SRef@cdp:eq:C.NC.4\endcsname{\href{S6_SERIES_GRIND_II.pdf\#nameddest=tag-00S3}{II.{C.NC.4} (p.~121)} }
\expandafter\def\csname SRef@cdp:eq:C.NC.5\endcsname{\href{S6_SERIES_GRIND_II.pdf\#nameddest=tag-00S4}{II.{C.NC.5} (p.~121)} }
\expandafter\def\csname SRef@cdp:eq:C.NC.6\endcsname{\href{S6_SERIES_GRIND_II.pdf\#nameddest=tag-00S5}{II.{C.NC.6} (p.~121)} }
\expandafter\def\csname SRef@cdp:eq:C.NC.7\endcsname{\href{S6_SERIES_GRIND_II.pdf\#nameddest=tag-00S6}{II.{C.NC.7} (p.~121)} }
\expandafter\def\csname SRef@cdp:eq:C.NC.8\endcsname{\href{S6_SERIES_GRIND_II.pdf\#nameddest=tag-00S7}{II.{C.NC.8} (p.~122)} }
\expandafter\def\csname SRef@cdp:eq:C.NC.9\endcsname{\href{S6_SERIES_GRIND_II.pdf\#nameddest=tag-00S8}{II.{C.NC.9} (p.~123)} }
\expandafter\def\csname SRef@cdp:eq:C.NC.10\endcsname{\href{S6_SERIES_GRIND_II.pdf\#nameddest=tag-00S1}{II.{C.NC.10} (p.~124)} }
\expandafter\def\csname SRef@cdp:eq:C.NC.11\endcsname{\href{S6_SERIES_GRIND_II.pdf\#nameddest=tag-00S2}{II.{C.NC.11} (p.~124)} }
\expandafter\def\csname SRef@cdp:eq:C.J.1\endcsname{\href{S6_SERIES_GRIND_II.pdf\#nameddest=tag-00OV}{II.{C.J.1} (p.~125)} }
\expandafter\def\csname SRef@cdp:eq:C.J.2\endcsname{\href{S6_SERIES_GRIND_II.pdf\#nameddest=tag-00OW}{II.{C.J.2} (p.~125)} }
\expandafter\def\csname SRef@cdp:eq:C.J.3\endcsname{\href{S6_SERIES_GRIND_II.pdf\#nameddest=tag-00OX}{II.{C.J.3} (p.~126)} }
\expandafter\def\csname SRef@cdp:eq:D1\endcsname{\href{S6_SERIES_GRIND_II.pdf\#nameddest=tag-019T}{II.{D1} (p.~21)} }
\expandafter\def\csname SRef@cdp:eq:D5\endcsname{\href{S6_SERIES_GRIND_II.pdf\#nameddest=tag-019X}{II.{D5} (p.~24)} }
\expandafter\def\csname SRef@cdp:eq:D4\endcsname{\href{S6_SERIES_GRIND_II.pdf\#nameddest=tag-019W}{II.{D4} (p.~25)} }
\expandafter\def\csname SRef@cdp:eq:D2\endcsname{\href{S6_SERIES_GRIND_II.pdf\#nameddest=tag-019U}{II.{D2} (p.~28)} }
\expandafter\def\csname SRef@cdp:eq:D3\endcsname{\href{S6_SERIES_GRIND_II.pdf\#nameddest=tag-019V}{II.{D3} (p.~28)} }
\expandafter\def\csname SRef@cdp:eq:C.SK.1\endcsname{\href{S6_SERIES_GRIND_II.pdf\#nameddest=tag-011H}{II.{C.SK.1} (p.~36)} }
\expandafter\def\csname SRef@cdp:eq:C.SK.2\endcsname{\href{S6_SERIES_GRIND_II.pdf\#nameddest=tag-011I}{II.{C.SK.2} (p.~36)} }
\expandafter\def\csname SRef@a.-one-dimensional-first-jets\endcsname{\href{S6_SERIES_GRIND_II.pdf\#nameddest=tag-0005}{II.{{$\langle 1\rangle $1.}} (p.~243)} }
\expandafter\def\csname SRef@three-smooth-branches-the-complete-excess-two-classification\endcsname{\href{S6_SERIES_GRIND_II.pdf\#nameddest=tag-0005}{II.{{$\langle 1\rangle $1.}} (p.~243)} }
\expandafter\def\csname SRef@cdp:eq:C.SK.3\endcsname{\href{S6_SERIES_GRIND_II.pdf\#nameddest=tag-011J}{II.{C.SK.3} (p.~37)} }
\expandafter\def\csname SRef@b.-two-dimensional-first-jets\endcsname{\href{S6_SERIES_GRIND_II.pdf\#nameddest=tag-003V}{II.{{$\langle 1\rangle $1.}} (p.~243)} }
\expandafter\def\csname SRef@multiplicity-four-one-cusp-and-two-smooth-branches\endcsname{\href{S6_SERIES_GRIND_II.pdf\#nameddest=tag-02QH}{II.{{$\langle 1\rangle $1.}} (p.~244)} }
\expandafter\def\csname SRef@a.-a-nonzero-coefficient-of-the-cusp-cubic\endcsname{\href{S6_SERIES_GRIND_II.pdf\#nameddest=tag-0004}{II.{{$\langle 1\rangle $1.}} (p.~245)} }
\expandafter\def\csname SRef@b.-a-zero-coefficient-of-the-cusp-cubic\endcsname{\href{S6_SERIES_GRIND_II.pdf\#nameddest=tag-003U}{II.{{$\langle 1\rangle $1.}} (p.~245)} }
\expandafter\def\csname SRef@cdp:eq:C.SK.4\endcsname{\href{S6_SERIES_GRIND_II.pdf\#nameddest=tag-011K}{II.{C.SK.4} (p.~37)} }
\expandafter\def\csname SRef@every-nonordinary-triple-exceeds-the-budget\endcsname{\href{S6_SERIES_GRIND_II.pdf\#nameddest=tag-01OO}{II.{{$\langle 1\rangle $1.}} (p.~246)} }
\expandafter\def\csname SRef@cdp:eq:C.SK.5\endcsname{\href{S6_SERIES_GRIND_II.pdf\#nameddest=tag-011L}{II.{C.SK.5} (p.~37)} }
\expandafter\def\csname SRef@cdp:eq:C.15\endcsname{\href{S6_SERIES_GRIND_II.pdf\#nameddest=tag-008T}{II.{C.15} (p.~255)} }
\expandafter\def\csname SRef@cdp:eq:C.16\endcsname{\href{S6_SERIES_GRIND_II.pdf\#nameddest=tag-008U}{II.{C.16} (p.~257)} }
\expandafter\def\csname SRef@cdp:eq:C.17\endcsname{\href{S6_SERIES_GRIND_II.pdf\#nameddest=tag-008V}{II.{C.17} (p.~257)} }
\expandafter\def\csname SRef@cdp:eq:C.18\endcsname{\href{S6_SERIES_GRIND_II.pdf\#nameddest=tag-008W}{II.{C.18} (p.~258)} }
\expandafter\def\csname SRef@cdp:eq:C.SK.6\endcsname{\href{S6_SERIES_GRIND_II.pdf\#nameddest=tag-011M}{II.{C.SK.6} (p.~37)} }
\expandafter\def\csname SRef@cdp:eq:C.19\endcsname{\href{S6_SERIES_GRIND_II.pdf\#nameddest=tag-008X}{II.{C.19} (p.~258)} }
\expandafter\def\csname SRef@cdp:eq:C.20\endcsname{\href{S6_SERIES_GRIND_II.pdf\#nameddest=tag-008Z}{II.{C.20} (p.~259)} }
\expandafter\def\csname SRef@cdp:eq:C.21\endcsname{\href{S6_SERIES_GRIND_II.pdf\#nameddest=tag-0090}{II.{C.21} (p.~260)} }
\expandafter\def\csname SRef@cdp:eq:C.22\endcsname{\href{S6_SERIES_GRIND_II.pdf\#nameddest=tag-0091}{II.{C.22} (p.~261)} }
\expandafter\def\csname SRef@cdp:eq:C.22b\endcsname{\href{S6_SERIES_GRIND_II.pdf\#nameddest=tag-0093}{II.{C.22b} (p.~261)} }
\expandafter\def\csname SRef@cdp:eq:C.22a\endcsname{\href{S6_SERIES_GRIND_II.pdf\#nameddest=tag-0092}{II.{C.22a} (p.~263)} }
\expandafter\def\csname SRef@cdp:eq:C.VC.1\endcsname{\href{S6_SERIES_GRIND_II.pdf\#nameddest=tag-015G}{II.{C.VC.1} (p.~263)} }
\expandafter\def\csname SRef@cdp:eq:C.VC.2\endcsname{\href{S6_SERIES_GRIND_II.pdf\#nameddest=tag-015H}{II.{C.VC.2} (p.~263)} }
\expandafter\def\csname SRef@cdp:eq:C.VC.3\endcsname{\href{S6_SERIES_GRIND_II.pdf\#nameddest=tag-015I}{II.{C.VC.3} (p.~264)} }
\expandafter\def\csname SRef@cdp:eq:C.VC.4\endcsname{\href{S6_SERIES_GRIND_II.pdf\#nameddest=tag-015J}{II.{C.VC.4} (p.~264)} }
\expandafter\def\csname SRef@cdp:eq:C.PK.1\endcsname{\href{S6_SERIES_GRIND_II.pdf\#nameddest=tag-00V8}{II.{C.PK.1} (p.~38)} }
\expandafter\def\csname SRef@cdp:eq:C.TEST.1\endcsname{\href{S6_SERIES_GRIND_II.pdf\#nameddest=tag-012O}{II.{C.TEST.1} (p.~265)} }
\expandafter\def\csname SRef@cdp:eq:C.TEST.2\endcsname{\href{S6_SERIES_GRIND_II.pdf\#nameddest=tag-012P}{II.{C.TEST.2} (p.~265)} }
\expandafter\def\csname SRef@cdp:eq:C.3\endcsname{\href{S6_SERIES_GRIND_II.pdf\#nameddest=tag-0094}{II.{C.3} (p.~266)} }
\expandafter\def\csname SRef@cdp:eq:C.4\endcsname{\href{S6_SERIES_GRIND_II.pdf\#nameddest=tag-0095}{II.{C.4} (p.~266)} }
\expandafter\def\csname SRef@cdp:eq:C.5\endcsname{\href{S6_SERIES_GRIND_II.pdf\#nameddest=tag-0096}{II.{C.5} (p.~267)} }
\expandafter\def\csname SRef@cdp:eq:C.6\endcsname{\href{S6_SERIES_GRIND_II.pdf\#nameddest=tag-0097}{II.{C.6} (p.~267)} }
\expandafter\def\csname SRef@cdp:eq:C.7\endcsname{\href{S6_SERIES_GRIND_II.pdf\#nameddest=tag-0098}{II.{C.7} (p.~267)} }
\expandafter\def\csname SRef@cdp:eq:C.8\endcsname{\href{S6_SERIES_GRIND_II.pdf\#nameddest=tag-0099}{II.{C.8} (p.~268)} }
\expandafter\def\csname SRef@cdp:eq:C.PK.2\endcsname{\href{S6_SERIES_GRIND_II.pdf\#nameddest=tag-00V9}{II.{C.PK.2} (p.~38)} }
\expandafter\def\csname SRef@cdp:eq:C.9\endcsname{\href{S6_SERIES_GRIND_II.pdf\#nameddest=tag-009A}{II.{C.9} (p.~269)} }
\expandafter\def\csname SRef@cdp:eq:C.10\endcsname{\href{S6_SERIES_GRIND_II.pdf\#nameddest=tag-008O}{II.{C.10} (p.~269)} }
\expandafter\def\csname SRef@cdp:eq:C.11\endcsname{\href{S6_SERIES_GRIND_II.pdf\#nameddest=tag-008P}{II.{C.11} (p.~269)} }
\expandafter\def\csname SRef@cdp:eq:C.12\endcsname{\href{S6_SERIES_GRIND_II.pdf\#nameddest=tag-008Q}{II.{C.12} (p.~269)} }
\expandafter\def\csname SRef@cdp:eq:C.13\endcsname{\href{S6_SERIES_GRIND_II.pdf\#nameddest=tag-008R}{II.{C.13} (p.~270)} }
\expandafter\def\csname SRef@cdp:eq:C.14\endcsname{\href{S6_SERIES_GRIND_II.pdf\#nameddest=tag-008S}{II.{C.14} (p.~271)} }
\expandafter\def\csname SRef@cdp:eq:D.18\endcsname{\href{S6_SERIES_GRIND_II.pdf\#nameddest=tag-0195}{II.{D.18} (p.~274)} }
\expandafter\def\csname SRef@cdp:eq:D.19\endcsname{\href{S6_SERIES_GRIND_II.pdf\#nameddest=tag-0196}{II.{D.19} (p.~275)} }
\expandafter\def\csname SRef@cdp:eq:D.S.1\endcsname{\href{S6_SERIES_GRIND_II.pdf\#nameddest=tag-019O}{II.{D.S.1} (p.~276)} }
\expandafter\def\csname SRef@cdp:eq:C.PK.3\endcsname{\href{S6_SERIES_GRIND_II.pdf\#nameddest=tag-00VA}{II.{C.PK.3} (p.~38)} }
\expandafter\def\csname SRef@cdp:eq:D.S.2\endcsname{\href{S6_SERIES_GRIND_II.pdf\#nameddest=tag-019P}{II.{D.S.2} (p.~276)} }
\expandafter\def\csname SRef@cdp:eq:D.S.3\endcsname{\href{S6_SERIES_GRIND_II.pdf\#nameddest=tag-019Q}{II.{D.S.3} (p.~277)} }
\expandafter\def\csname SRef@cdp:eq:D.S.4\endcsname{\href{S6_SERIES_GRIND_II.pdf\#nameddest=tag-019R}{II.{D.S.4} (p.~277)} }
\expandafter\def\csname SRef@cdp:eq:D.S.5\endcsname{\href{S6_SERIES_GRIND_II.pdf\#nameddest=tag-019S}{II.{D.S.5} (p.~277)} }
\expandafter\def\csname SRef@cdp:eq:D.20\endcsname{\href{S6_SERIES_GRIND_II.pdf\#nameddest=tag-0197}{II.{D.20} (p.~278)} }
\expandafter\def\csname SRef@cdp:eq:D.21\endcsname{\href{S6_SERIES_GRIND_II.pdf\#nameddest=tag-0198}{II.{D.21} (p.~279)} }
\expandafter\def\csname SRef@cdp:eq:D.22\endcsname{\href{S6_SERIES_GRIND_II.pdf\#nameddest=tag-0199}{II.{D.22} (p.~280)} }
\expandafter\def\csname SRef@cdp:eq:D.23\endcsname{\href{S6_SERIES_GRIND_II.pdf\#nameddest=tag-019A}{II.{D.23} (p.~280)} }
\expandafter\def\csname SRef@cdp:eq:D.24\endcsname{\href{S6_SERIES_GRIND_II.pdf\#nameddest=tag-019B}{II.{D.24} (p.~280)} }
\expandafter\def\csname SRef@cdp:eq:C.PK.4\endcsname{\href{S6_SERIES_GRIND_II.pdf\#nameddest=tag-00VB}{II.{C.PK.4} (p.~39)} }
\expandafter\def\csname SRef@cdp:eq:D.25\endcsname{\href{S6_SERIES_GRIND_II.pdf\#nameddest=tag-019C}{II.{D.25} (p.~280)} }
\expandafter\def\csname SRef@cdp:eq:D.26\endcsname{\href{S6_SERIES_GRIND_II.pdf\#nameddest=tag-019D}{II.{D.26} (p.~281)} }
\expandafter\def\csname SRef@cdp:eq:D.27\endcsname{\href{S6_SERIES_GRIND_II.pdf\#nameddest=tag-019E}{II.{D.27} (p.~281)} }
\expandafter\def\csname SRef@cdp:eq:D.28\endcsname{\href{S6_SERIES_GRIND_II.pdf\#nameddest=tag-019F}{II.{D.28} (p.~281)} }
\expandafter\def\csname SRef@cdp:eq:D.29\endcsname{\href{S6_SERIES_GRIND_II.pdf\#nameddest=tag-019G}{II.{D.29} (p.~282)} }
\expandafter\def\csname SRef@cdp:eq:D.30\endcsname{\href{S6_SERIES_GRIND_II.pdf\#nameddest=tag-019H}{II.{D.30} (p.~282)} }
\expandafter\def\csname SRef@cdp:eq:D.31\endcsname{\href{S6_SERIES_GRIND_II.pdf\#nameddest=tag-019I}{II.{D.31} (p.~282)} }
\expandafter\def\csname SRef@cdp:eq:D.32\endcsname{\href{S6_SERIES_GRIND_II.pdf\#nameddest=tag-019J}{II.{D.32} (p.~282)} }
\expandafter\def\csname SRef@cdp:eq:D.33\endcsname{\href{S6_SERIES_GRIND_II.pdf\#nameddest=tag-019K}{II.{D.33} (p.~283)} }
\expandafter\def\csname SRef@cdp:eq:C.PK.5\endcsname{\href{S6_SERIES_GRIND_II.pdf\#nameddest=tag-00VC}{II.{C.PK.5} (p.~40)} }
\expandafter\def\csname SRef@cdp:eq:D.34\endcsname{\href{S6_SERIES_GRIND_II.pdf\#nameddest=tag-019L}{II.{D.34} (p.~283)} }
\expandafter\def\csname SRef@cdp:eq:D.35\endcsname{\href{S6_SERIES_GRIND_II.pdf\#nameddest=tag-019M}{II.{D.35} (p.~283)} }
\expandafter\def\csname SRef@cdp:eq:D.36\endcsname{\href{S6_SERIES_GRIND_II.pdf\#nameddest=tag-019N}{II.{D.36} (p.~283)} }
\expandafter\def\csname SRef@cdp:eq:E.16.1\endcsname{\href{S6_SERIES_GRIND_II.pdf\#nameddest=tag-01A2}{II.{E.16.1} (p.~304)} }
\expandafter\def\csname SRef@cdp:eq:E.16.2\endcsname{\href{S6_SERIES_GRIND_II.pdf\#nameddest=tag-01A3}{II.{E.16.2} (p.~304)} }
\expandafter\def\csname SRef@cdp:eq:E.16.3\endcsname{\href{S6_SERIES_GRIND_II.pdf\#nameddest=tag-01A4}{II.{E.16.3} (p.~305)} }
\expandafter\def\csname SRef@cdp:eq:C.PK.6\endcsname{\href{S6_SERIES_GRIND_II.pdf\#nameddest=tag-00VD}{II.{C.PK.6} (p.~40)} }
\expandafter\def\csname SRef@cdp:eq:C.PK.7\endcsname{\href{S6_SERIES_GRIND_II.pdf\#nameddest=tag-00VE}{II.{C.PK.7} (p.~40)} }
\expandafter\def\csname SRef@cdp:eq:C.PK.8\endcsname{\href{S6_SERIES_GRIND_II.pdf\#nameddest=tag-00VF}{II.{C.PK.8} (p.~40)} }
\expandafter\def\csname SRef@cdp:eq:C.IN.1\endcsname{\href{S6_SERIES_GRIND_II.pdf\#nameddest=tag-00NK}{II.{C.IN.1} (p.~41)} }
\expandafter\def\csname SRef@invII:eq:E.2\endcsname{\href{S6_SERIES_GRIND_II.pdf\#nameddest=tag-01TR}{II.{E.2} (p.~354)} }
\expandafter\def\csname SRef@cdp:eq:C.IN.2\endcsname{\href{S6_SERIES_GRIND_II.pdf\#nameddest=tag-00NL}{II.{C.IN.2} (p.~41)} }
\expandafter\def\csname SRef@invII:eq:E.3\endcsname{\href{S6_SERIES_GRIND_II.pdf\#nameddest=tag-01TS}{II.{E.3} (p.~354)} }
\expandafter\def\csname SRef@invII:eq:E.4a\endcsname{\href{S6_SERIES_GRIND_II.pdf\#nameddest=tag-01TT}{II.{E.4a} (p.~354)} }
\expandafter\def\csname SRef@invII:eq:E.5\endcsname{\href{S6_SERIES_GRIND_II.pdf\#nameddest=tag-01TU}{II.{E.5} (p.~354)} }
\expandafter\def\csname SRef@invII:eq:E.6\endcsname{\href{S6_SERIES_GRIND_II.pdf\#nameddest=tag-01TV}{II.{E.6} (p.~356)} }
\expandafter\def\csname SRef@inv:wb:eq:9.19\endcsname{\href{S6_SERIES_GRIND_II.pdf\#nameddest=tag-01TJ}{II.{WBAPII.9.19} (p.~362)} }
\expandafter\def\csname SRef@inv:wb:eq:9.20\endcsname{\href{S6_SERIES_GRIND_II.pdf\#nameddest=tag-01TK}{II.{WBAPII.9.20} (p.~362)} }
\expandafter\def\csname SRef@inv:wb:eq:9.21\endcsname{\href{S6_SERIES_GRIND_II.pdf\#nameddest=tag-01TL}{II.{WBAPII.9.21} (p.~362)} }
\expandafter\def\csname SRef@inv:wb:eq:9.22\endcsname{\href{S6_SERIES_GRIND_II.pdf\#nameddest=tag-01TM}{II.{WBAPII.9.22} (p.~362)} }
\expandafter\def\csname SRef@inv:wb:eq:9.23\endcsname{\href{S6_SERIES_GRIND_II.pdf\#nameddest=tag-01TN}{II.{WBAPII.9.23} (p.~362)} }
\expandafter\def\csname SRef@inv:wb:eq:9.24\endcsname{\href{S6_SERIES_GRIND_II.pdf\#nameddest=tag-01TO}{II.{WBAPII.9.24} (p.~362)} }
\expandafter\def\csname SRef@cdp:eq:C.IN.3\endcsname{\href{S6_SERIES_GRIND_II.pdf\#nameddest=tag-00NM}{II.{C.IN.3} (p.~41)} }
\expandafter\def\csname SRef@inv:wb:eq:9.25\endcsname{\href{S6_SERIES_GRIND_II.pdf\#nameddest=tag-01TP}{II.{WBAPII.9.25} (p.~362)} }
\expandafter\def\csname SRef@wbII:eq:CN:2.2\endcsname{\href{S6_SERIES_GRIND_II.pdf\#nameddest=tag-03DW}{II.{W.145} (p.~364)} }
\expandafter\def\csname SRef@wbII:eq:W:5.3\endcsname{\href{S6_SERIES_GRIND_II.pdf\#nameddest=tag-03HS}{II.{W.3} (p.~365)} }
\expandafter\def\csname SRef@wbII:eq:W:5.4\endcsname{\href{S6_SERIES_GRIND_II.pdf\#nameddest=tag-03HX}{II.{W.4} (p.~365)} }
\expandafter\def\csname SRef@wbII:eq:W:5.5\endcsname{\href{S6_SERIES_GRIND_II.pdf\#nameddest=tag-03HY}{II.{W.5} (p.~366)} }
\expandafter\def\csname SRef@wbII:eq:W:5.6\endcsname{\href{S6_SERIES_GRIND_II.pdf\#nameddest=tag-03HZ}{II.{W.6} (p.~366)} }
\expandafter\def\csname SRef@cdp:eq:C.IN.4\endcsname{\href{S6_SERIES_GRIND_II.pdf\#nameddest=tag-00NN}{II.{C.IN.4} (p.~42)} }
\expandafter\def\csname SRef@wbII:eq:W:5.7\endcsname{\href{S6_SERIES_GRIND_II.pdf\#nameddest=tag-03I0}{II.{W.7} (p.~366)} }
\expandafter\def\csname SRef@wbII:eq:W:5.8\endcsname{\href{S6_SERIES_GRIND_II.pdf\#nameddest=tag-03I1}{II.{W.8} (p.~368)} }
\expandafter\def\csname SRef@wbII:eq:W:5.9\endcsname{\href{S6_SERIES_GRIND_II.pdf\#nameddest=tag-03I2}{II.{W.9} (p.~368)} }
\expandafter\def\csname SRef@wbII:eq:W:5.10\endcsname{\href{S6_SERIES_GRIND_II.pdf\#nameddest=tag-03EA}{II.{W.10} (p.~368)} }
\expandafter\def\csname SRef@wbII:eq:CN:4.3\endcsname{\href{S6_SERIES_GRIND_II.pdf\#nameddest=tag-03DX}{II.{W.160} (p.~369)} }
\expandafter\def\csname SRef@wbII:eq:CN:4.4\endcsname{\href{S6_SERIES_GRIND_II.pdf\#nameddest=tag-03DY}{II.{W.161} (p.~369)} }
\expandafter\def\csname SRef@wbII:eq:CN:4.5a\endcsname{\href{S6_SERIES_GRIND_II.pdf\#nameddest=tag-03E0}{II.{W.162} (p.~370)} }
\expandafter\def\csname SRef@cdp:eq:C.IN.5\endcsname{\href{S6_SERIES_GRIND_II.pdf\#nameddest=tag-00NO}{II.{C.IN.5} (p.~42)} }
\expandafter\def\csname SRef@wbII:eq:CN:4.5\endcsname{\href{S6_SERIES_GRIND_II.pdf\#nameddest=tag-03DZ}{II.{W.163} (p.~370)} }
\expandafter\def\csname SRef@wbII:eq:CN:4.6\endcsname{\href{S6_SERIES_GRIND_II.pdf\#nameddest=tag-03E1}{II.{W.164} (p.~370)} }
\expandafter\def\csname SRef@wbII:eq:CN:4.7\endcsname{\href{S6_SERIES_GRIND_II.pdf\#nameddest=tag-03E2}{II.{W.165} (p.~370)} }
\expandafter\def\csname SRef@wbII:eq:CN:4.8\endcsname{\href{S6_SERIES_GRIND_II.pdf\#nameddest=tag-03E3}{II.{W.166} (p.~371)} }
\expandafter\def\csname SRef@wbII:eq:CN:4.9\endcsname{\href{S6_SERIES_GRIND_II.pdf\#nameddest=tag-03E4}{II.{W.167} (p.~371)} }
\expandafter\def\csname SRef@wbII:eq:W:5.11\endcsname{\href{S6_SERIES_GRIND_II.pdf\#nameddest=tag-03EB}{II.{W.11} (p.~371)} }
\expandafter\def\csname SRef@wbII:eq:W:5.15a\endcsname{\href{S6_SERIES_GRIND_II.pdf\#nameddest=tag-03EC}{II.{W.12} (p.~372)} }
\expandafter\def\csname SRef@wbII:eq:W:5.16\endcsname{\href{S6_SERIES_GRIND_II.pdf\#nameddest=tag-03ED}{II.{W.13} (p.~374)} }
\expandafter\def\csname SRef@wbII:eq:W:5.17\endcsname{\href{S6_SERIES_GRIND_II.pdf\#nameddest=tag-03EE}{II.{W.14} (p.~374)} }
\expandafter\def\csname SRef@cdp:eq:C.IN.6\endcsname{\href{S6_SERIES_GRIND_II.pdf\#nameddest=tag-00NP}{II.{C.IN.6} (p.~43)} }
\expandafter\def\csname SRef@wbII:eq:W:5.17-V\endcsname{\href{S6_SERIES_GRIND_II.pdf\#nameddest=tag-03FU}{II.{W.15} (p.~374)} }
\expandafter\def\csname SRef@wbII:eq:W:5.17-T\endcsname{\href{S6_SERIES_GRIND_II.pdf\#nameddest=tag-03FS}{II.{W.16} (p.~374)} }
\expandafter\def\csname SRef@wbII:eq:W:5.17-U\endcsname{\href{S6_SERIES_GRIND_II.pdf\#nameddest=tag-03FT}{II.{W.17} (p.~374)} }
\expandafter\def\csname SRef@wbII:eq:W:5.17a\endcsname{\href{S6_SERIES_GRIND_II.pdf\#nameddest=tag-03FV}{II.{W.18} (p.~376)} }
\expandafter\def\csname SRef@wbII:eq:W:5.17b\endcsname{\href{S6_SERIES_GRIND_II.pdf\#nameddest=tag-03GL}{II.{W.19} (p.~376)} }
\expandafter\def\csname SRef@wbII:eq:W:5.17c\endcsname{\href{S6_SERIES_GRIND_II.pdf\#nameddest=tag-03GM}{II.{W.20} (p.~376)} }
\expandafter\def\csname SRef@wbII:eq:W:5.17d\endcsname{\href{S6_SERIES_GRIND_II.pdf\#nameddest=tag-03GN}{II.{W.21} (p.~377)} }
\expandafter\def\csname SRef@wbII:eq:W:5.17e\endcsname{\href{S6_SERIES_GRIND_II.pdf\#nameddest=tag-03GO}{II.{W.22} (p.~377)} }
\expandafter\def\csname SRef@wbII:eq:W:5.17g\endcsname{\href{S6_SERIES_GRIND_II.pdf\#nameddest=tag-03GQ}{II.{W.23} (p.~377)} }
\expandafter\def\csname SRef@wbII:eq:W:5.17h\endcsname{\href{S6_SERIES_GRIND_II.pdf\#nameddest=tag-03GR}{II.{W.24} (p.~377)} }
\expandafter\def\csname SRef@cdp:eq:C.PK.9\endcsname{\href{S6_SERIES_GRIND_II.pdf\#nameddest=tag-00VG}{II.{C.PK.9} (p.~44)} }
\expandafter\def\csname SRef@wbII:eq:W:5.17i\endcsname{\href{S6_SERIES_GRIND_II.pdf\#nameddest=tag-03GS}{II.{W.25} (p.~377)} }
\expandafter\def\csname SRef@wbII:eq:W:5.17j'\endcsname{\href{S6_SERIES_GRIND_II.pdf\#nameddest=tag-03GU}{II.{W.26} (p.~378)} }
\expandafter\def\csname SRef@wbII:eq:W:5.17j\endcsname{\href{S6_SERIES_GRIND_II.pdf\#nameddest=tag-03GT}{II.{W.27} (p.~378)} }
\expandafter\def\csname SRef@wbII:eq:W:5.17k\endcsname{\href{S6_SERIES_GRIND_II.pdf\#nameddest=tag-03GV}{II.{W.28} (p.~378)} }
\expandafter\def\csname SRef@wbII:eq:W:5.17l\endcsname{\href{S6_SERIES_GRIND_II.pdf\#nameddest=tag-03GW}{II.{W.29} (p.~379)} }
\expandafter\def\csname SRef@wbII:eq:W:5.17l'\endcsname{\href{S6_SERIES_GRIND_II.pdf\#nameddest=tag-03GX}{II.{W.30} (p.~379)} }
\expandafter\def\csname SRef@wbII:eq:W:5.17l''\endcsname{\href{S6_SERIES_GRIND_II.pdf\#nameddest=tag-03GY}{II.{W.31} (p.~379)} }
\expandafter\def\csname SRef@wbII:eq:W:5.17l'''\endcsname{\href{S6_SERIES_GRIND_II.pdf\#nameddest=tag-03GZ}{II.{W.32} (p.~379)} }
\expandafter\def\csname SRef@wbII:eq:W:5.17m\endcsname{\href{S6_SERIES_GRIND_II.pdf\#nameddest=tag-03H0}{II.{W.33} (p.~379)} }
\expandafter\def\csname SRef@wbII:eq:W:5.17m'\endcsname{\href{S6_SERIES_GRIND_II.pdf\#nameddest=tag-03H1}{II.{W.34} (p.~379)} }
\expandafter\def\csname SRef@cdp:eq:C.H.1\endcsname{\href{S6_SERIES_GRIND_II.pdf\#nameddest=tag-00L8}{II.{C.H.1} (p.~46)} }
\expandafter\def\csname SRef@wbII:eq:W:5.17m''\endcsname{\href{S6_SERIES_GRIND_II.pdf\#nameddest=tag-03H2}{II.{W.35} (p.~380)} }
\expandafter\def\csname SRef@wbII:eq:W:5.17n\endcsname{\href{S6_SERIES_GRIND_II.pdf\#nameddest=tag-03H3}{II.{W.36} (p.~380)} }
\expandafter\def\csname SRef@wbII:eq:W:5.17o\endcsname{\href{S6_SERIES_GRIND_II.pdf\#nameddest=tag-03H4}{II.{W.37} (p.~380)} }
\expandafter\def\csname SRef@wbII:eq:W:5.17p\endcsname{\href{S6_SERIES_GRIND_II.pdf\#nameddest=tag-03H5}{II.{W.38} (p.~381)} }
\expandafter\def\csname SRef@wbII:eq:W:5.17p0\endcsname{\href{S6_SERIES_GRIND_II.pdf\#nameddest=tag-03H7}{II.{W.39} (p.~381)} }
\expandafter\def\csname SRef@wbII:eq:W:5.17p'\endcsname{\href{S6_SERIES_GRIND_II.pdf\#nameddest=tag-03H6}{II.{W.40} (p.~381)} }
\expandafter\def\csname SRef@wbII:eq:W:5.17q\endcsname{\href{S6_SERIES_GRIND_II.pdf\#nameddest=tag-03H8}{II.{W.41} (p.~381)} }
\expandafter\def\csname SRef@wbII:eq:W:5.17r\endcsname{\href{S6_SERIES_GRIND_II.pdf\#nameddest=tag-03H9}{II.{W.42} (p.~382)} }
\expandafter\def\csname SRef@wbII:eq:W:5.17s\endcsname{\href{S6_SERIES_GRIND_II.pdf\#nameddest=tag-03HA}{II.{W.43} (p.~382)} }
\expandafter\def\csname SRef@wbII:eq:W:5.17t\endcsname{\href{S6_SERIES_GRIND_II.pdf\#nameddest=tag-03HB}{II.{W.61} (p.~382)} }
\expandafter\def\csname SRef@cdp:eq:A.1\endcsname{\href{S6_SERIES_GRIND_II.pdf\#nameddest=tag-003Y}{II.{A.1} (p.~48)} }
\expandafter\def\csname SRef@wbII:eq:W:5.17u\endcsname{\href{S6_SERIES_GRIND_II.pdf\#nameddest=tag-03HC}{II.{W.62} (p.~382)} }
\expandafter\def\csname SRef@wbII:eq:W:5.17v\endcsname{\href{S6_SERIES_GRIND_II.pdf\#nameddest=tag-03HD}{II.{W.63} (p.~382)} }
\expandafter\def\csname SRef@wbII:eq:W:5.17w\endcsname{\href{S6_SERIES_GRIND_II.pdf\#nameddest=tag-03HE}{II.{W.64} (p.~383)} }
\expandafter\def\csname SRef@wbII:eq:W:5.17x\endcsname{\href{S6_SERIES_GRIND_II.pdf\#nameddest=tag-03HF}{II.{W.65} (p.~383)} }
\expandafter\def\csname SRef@wbII:eq:W:5.17y\endcsname{\href{S6_SERIES_GRIND_II.pdf\#nameddest=tag-03HG}{II.{W.66} (p.~383)} }
\expandafter\def\csname SRef@wbII:eq:W:5.17z\endcsname{\href{S6_SERIES_GRIND_II.pdf\#nameddest=tag-03HH}{II.{W.67} (p.~383)} }
\expandafter\def\csname SRef@wbII:eq:W:5.17aa\endcsname{\href{S6_SERIES_GRIND_II.pdf\#nameddest=tag-03FW}{II.{W.68} (p.~383)} }
\expandafter\def\csname SRef@wbII:eq:W:5.17ab\endcsname{\href{S6_SERIES_GRIND_II.pdf\#nameddest=tag-03FX}{II.{W.69} (p.~383)} }
\expandafter\def\csname SRef@wbII:eq:W:5.17ac\endcsname{\href{S6_SERIES_GRIND_II.pdf\#nameddest=tag-03FY}{II.{W.70} (p.~384)} }
\expandafter\def\csname SRef@wbII:eq:W:5.17ad\endcsname{\href{S6_SERIES_GRIND_II.pdf\#nameddest=tag-03FZ}{II.{W.71} (p.~384)} }
\expandafter\def\csname SRef@cdp:eq:A.2\endcsname{\href{S6_SERIES_GRIND_II.pdf\#nameddest=tag-004B}{II.{A.2} (p.~48)} }
\expandafter\def\csname SRef@wbII:eq:W:5.17ae\endcsname{\href{S6_SERIES_GRIND_II.pdf\#nameddest=tag-03G0}{II.{W.72} (p.~384)} }
\expandafter\def\csname SRef@wbII:eq:W:5.17af\endcsname{\href{S6_SERIES_GRIND_II.pdf\#nameddest=tag-03G1}{II.{W.73} (p.~384)} }
\expandafter\def\csname SRef@wbII:eq:W:5.17ag\endcsname{\href{S6_SERIES_GRIND_II.pdf\#nameddest=tag-03G2}{II.{W.74} (p.~384)} }
\expandafter\def\csname SRef@wbII:eq:W:5.17ah\endcsname{\href{S6_SERIES_GRIND_II.pdf\#nameddest=tag-03G3}{II.{W.75} (p.~384)} }
\expandafter\def\csname SRef@wbII:eq:W:5.17ai\endcsname{\href{S6_SERIES_GRIND_II.pdf\#nameddest=tag-03G4}{II.{W.76} (p.~384)} }
\expandafter\def\csname SRef@wbII:eq:W:5.17aj\endcsname{\href{S6_SERIES_GRIND_II.pdf\#nameddest=tag-03G5}{II.{W.77} (p.~384)} }
\expandafter\def\csname SRef@wbII:eq:W:5.17ak\endcsname{\href{S6_SERIES_GRIND_II.pdf\#nameddest=tag-03G6}{II.{W.78} (p.~384)} }
\expandafter\def\csname SRef@wbII:eq:W:5.17al\endcsname{\href{S6_SERIES_GRIND_II.pdf\#nameddest=tag-03G7}{II.{W.79} (p.~385)} }
\expandafter\def\csname SRef@cdp:eq:A.3\endcsname{\href{S6_SERIES_GRIND_II.pdf\#nameddest=tag-004M}{II.{A.3} (p.~49)} }
\expandafter\def\csname SRef@wbII:eq:W:5.17am\endcsname{\href{S6_SERIES_GRIND_II.pdf\#nameddest=tag-03G8}{II.{W.80} (p.~385)} }
\expandafter\def\csname SRef@wbII:eq:W:5.17an\endcsname{\href{S6_SERIES_GRIND_II.pdf\#nameddest=tag-03G9}{II.{W.81} (p.~385)} }
\expandafter\def\csname SRef@wbII:eq:W:5.17ao\endcsname{\href{S6_SERIES_GRIND_II.pdf\#nameddest=tag-03GA}{II.{W.82} (p.~385)} }
\expandafter\def\csname SRef@wbII:eq:W:5.17ap\endcsname{\href{S6_SERIES_GRIND_II.pdf\#nameddest=tag-03GB}{II.{W.83} (p.~385)} }
\expandafter\def\csname SRef@wbII:eq:W:5.17aq\endcsname{\href{S6_SERIES_GRIND_II.pdf\#nameddest=tag-03GC}{II.{W.84} (p.~385)} }
\expandafter\def\csname SRef@wbII:eq:W:5.17ar\endcsname{\href{S6_SERIES_GRIND_II.pdf\#nameddest=tag-03GD}{II.{W.85} (p.~385)} }
\expandafter\def\csname SRef@wbII:eq:W:5.17as\endcsname{\href{S6_SERIES_GRIND_II.pdf\#nameddest=tag-03GE}{II.{W.86} (p.~385)} }
\expandafter\def\csname SRef@wbII:eq:W:5.17at\endcsname{\href{S6_SERIES_GRIND_II.pdf\#nameddest=tag-03GF}{II.{W.87} (p.~386)} }
\expandafter\def\csname SRef@wbII:eq:W:5.17au\endcsname{\href{S6_SERIES_GRIND_II.pdf\#nameddest=tag-03GG}{II.{W.88} (p.~386)} }
\expandafter\def\csname SRef@wbII:eq:W:5.17av\endcsname{\href{S6_SERIES_GRIND_II.pdf\#nameddest=tag-03GH}{II.{W.89} (p.~386)} }
\expandafter\def\csname SRef@cdp:eq:A.4\endcsname{\href{S6_SERIES_GRIND_II.pdf\#nameddest=tag-004O}{II.{A.4} (p.~49)} }
\expandafter\def\csname SRef@wbII:eq:W:5.17aw\endcsname{\href{S6_SERIES_GRIND_II.pdf\#nameddest=tag-03GI}{II.{W.90} (p.~386)} }
\expandafter\def\csname SRef@wbII:eq:W:5.17ax\endcsname{\href{S6_SERIES_GRIND_II.pdf\#nameddest=tag-03GJ}{II.{W.91} (p.~386)} }
\expandafter\def\csname SRef@wbII:eq:W:5.17ay\endcsname{\href{S6_SERIES_GRIND_II.pdf\#nameddest=tag-03GK}{II.{W.92} (p.~386)} }
\expandafter\def\csname SRef@wbII:eq:W:5.17-A1\endcsname{\href{S6_SERIES_GRIND_II.pdf\#nameddest=tag-03EF}{II.{W.93} (p.~387)} }
\expandafter\def\csname SRef@wbII:eq:W:5.17-A2\endcsname{\href{S6_SERIES_GRIND_II.pdf\#nameddest=tag-03EY}{II.{W.94} (p.~388)} }
\expandafter\def\csname SRef@wbII:eq:W:5.17-A3\endcsname{\href{S6_SERIES_GRIND_II.pdf\#nameddest=tag-03F9}{II.{W.95} (p.~388)} }
\expandafter\def\csname SRef@wbII:eq:W:5.17-A4\endcsname{\href{S6_SERIES_GRIND_II.pdf\#nameddest=tag-03FM}{II.{W.96} (p.~388)} }
\expandafter\def\csname SRef@wbII:eq:W:5.17-A5\endcsname{\href{S6_SERIES_GRIND_II.pdf\#nameddest=tag-03FN}{II.{W.97} (p.~388)} }
\expandafter\def\csname SRef@wbII:eq:W:5.17-A6\endcsname{\href{S6_SERIES_GRIND_II.pdf\#nameddest=tag-03FO}{II.{W.98} (p.~388)} }
\expandafter\def\csname SRef@wbII:eq:W:5.17-A7\endcsname{\href{S6_SERIES_GRIND_II.pdf\#nameddest=tag-03FP}{II.{W.99} (p.~388)} }
\expandafter\def\csname SRef@wbII:eq:W:5.17-A8\endcsname{\href{S6_SERIES_GRIND_II.pdf\#nameddest=tag-03FQ}{II.{W.100} (p.~388)} }
\expandafter\def\csname SRef@wbII:eq:W:5.17-A9\endcsname{\href{S6_SERIES_GRIND_II.pdf\#nameddest=tag-03FR}{II.{W.101} (p.~388)} }
\expandafter\def\csname SRef@wbII:eq:W:5.17-A10\endcsname{\href{S6_SERIES_GRIND_II.pdf\#nameddest=tag-03EG}{II.{W.102} (p.~388)} }
\expandafter\def\csname SRef@wbII:eq:W:5.17-A11\endcsname{\href{S6_SERIES_GRIND_II.pdf\#nameddest=tag-03EH}{II.{W.103} (p.~388)} }
\expandafter\def\csname SRef@wbII:eq:W:5.17-A12\endcsname{\href{S6_SERIES_GRIND_II.pdf\#nameddest=tag-03EI}{II.{W.104} (p.~389)} }
\expandafter\def\csname SRef@wbII:eq:W:5.17-A13\endcsname{\href{S6_SERIES_GRIND_II.pdf\#nameddest=tag-03EJ}{II.{W.105} (p.~389)} }
\expandafter\def\csname SRef@wbII:eq:W:5.17-A14\endcsname{\href{S6_SERIES_GRIND_II.pdf\#nameddest=tag-03EK}{II.{W.106} (p.~389)} }
\expandafter\def\csname SRef@wbII:eq:W:5.17-A14a\endcsname{\href{S6_SERIES_GRIND_II.pdf\#nameddest=tag-03EL}{II.{W.107} (p.~390)} }
\expandafter\def\csname SRef@cdp:eq:A.5\endcsname{\href{S6_SERIES_GRIND_II.pdf\#nameddest=tag-004P}{II.{A.5} (p.~49)} }
\expandafter\def\csname SRef@wbII:eq:W:5.17-A14c\endcsname{\href{S6_SERIES_GRIND_II.pdf\#nameddest=tag-03EM}{II.{W.108} (p.~390)} }
\expandafter\def\csname SRef@wbII:eq:W:5.17-A14d\endcsname{\href{S6_SERIES_GRIND_II.pdf\#nameddest=tag-03EN}{II.{W.109} (p.~390)} }
\expandafter\def\csname SRef@wbII:eq:W:5.17-A14e\endcsname{\href{S6_SERIES_GRIND_II.pdf\#nameddest=tag-03EO}{II.{W.110} (p.~392)} }
\expandafter\def\csname SRef@wbII:eq:W:5.17-A14f\endcsname{\href{S6_SERIES_GRIND_II.pdf\#nameddest=tag-03EP}{II.{W.111} (p.~392)} }
\expandafter\def\csname SRef@wbII:eq:W:5.17-A15\endcsname{\href{S6_SERIES_GRIND_II.pdf\#nameddest=tag-03EQ}{II.{W.112} (p.~393)} }
\expandafter\def\csname SRef@wbII:eq:W:5.17-A16\endcsname{\href{S6_SERIES_GRIND_II.pdf\#nameddest=tag-03ER}{II.{W.113} (p.~393)} }
\expandafter\def\csname SRef@wbII:eq:W:5.17-A17\endcsname{\href{S6_SERIES_GRIND_II.pdf\#nameddest=tag-03ES}{II.{W.114} (p.~393)} }
\expandafter\def\csname SRef@wbII:eq:W:5.17-A18\endcsname{\href{S6_SERIES_GRIND_II.pdf\#nameddest=tag-03ET}{II.{W.115} (p.~393)} }
\expandafter\def\csname SRef@wbII:eq:W:5.17-A19\endcsname{\href{S6_SERIES_GRIND_II.pdf\#nameddest=tag-03EU}{II.{W.116} (p.~393)} }
\expandafter\def\csname SRef@wbII:eq:W:5.17-A19a\endcsname{\href{S6_SERIES_GRIND_II.pdf\#nameddest=tag-03EV}{II.{W.117} (p.~393)} }
\expandafter\def\csname SRef@cdp:eq:A.6\endcsname{\href{S6_SERIES_GRIND_II.pdf\#nameddest=tag-004Q}{II.{A.6} (p.~50)} }
\expandafter\def\csname SRef@wbII:eq:W:5.17-A19b\endcsname{\href{S6_SERIES_GRIND_II.pdf\#nameddest=tag-03EW}{II.{W.118} (p.~393)} }
\expandafter\def\csname SRef@wbII:eq:W:5.17-A19c\endcsname{\href{S6_SERIES_GRIND_II.pdf\#nameddest=tag-03EX}{II.{W.119} (p.~394)} }
\expandafter\def\csname SRef@wbII:eq:W:5.17-A20\endcsname{\href{S6_SERIES_GRIND_II.pdf\#nameddest=tag-03EZ}{II.{W.120} (p.~394)} }
\expandafter\def\csname SRef@wbII:eq:W:5.17-A21\endcsname{\href{S6_SERIES_GRIND_II.pdf\#nameddest=tag-03F0}{II.{W.121} (p.~394)} }
\expandafter\def\csname SRef@wbII:eq:W:5.17-A22\endcsname{\href{S6_SERIES_GRIND_II.pdf\#nameddest=tag-03F1}{II.{W.122} (p.~394)} }
\expandafter\def\csname SRef@wbII:eq:W:5.17-A23\endcsname{\href{S6_SERIES_GRIND_II.pdf\#nameddest=tag-03F2}{II.{W.123} (p.~395)} }
\expandafter\def\csname SRef@wbII:eq:W:5.17-A24\endcsname{\href{S6_SERIES_GRIND_II.pdf\#nameddest=tag-03F3}{II.{W.124} (p.~395)} }
\expandafter\def\csname SRef@wbII:eq:W:5.17-A25\endcsname{\href{S6_SERIES_GRIND_II.pdf\#nameddest=tag-03F4}{II.{W.125} (p.~395)} }
\expandafter\def\csname SRef@wbII:eq:W:5.17-A26\endcsname{\href{S6_SERIES_GRIND_II.pdf\#nameddest=tag-03F5}{II.{W.126} (p.~395)} }
\expandafter\def\csname SRef@wbII:eq:W:5.17-A27\endcsname{\href{S6_SERIES_GRIND_II.pdf\#nameddest=tag-03F6}{II.{W.127} (p.~395)} }
\expandafter\def\csname SRef@cdp:eq:A.7\endcsname{\href{S6_SERIES_GRIND_II.pdf\#nameddest=tag-004R}{II.{A.7} (p.~50)} }
\expandafter\def\csname SRef@wbII:eq:W:5.17-A28\endcsname{\href{S6_SERIES_GRIND_II.pdf\#nameddest=tag-03F7}{II.{W.128} (p.~395)} }
\expandafter\def\csname SRef@wbII:eq:W:5.17-A29\endcsname{\href{S6_SERIES_GRIND_II.pdf\#nameddest=tag-03F8}{II.{W.129} (p.~395)} }
\expandafter\def\csname SRef@wbII:eq:W:5.17-A30\endcsname{\href{S6_SERIES_GRIND_II.pdf\#nameddest=tag-03FA}{II.{W.130} (p.~395)} }
\expandafter\def\csname SRef@wbII:eq:W:5.17-A31\endcsname{\href{S6_SERIES_GRIND_II.pdf\#nameddest=tag-03FB}{II.{W.131} (p.~395)} }
\expandafter\def\csname SRef@wbII:eq:W:5.17-A32\endcsname{\href{S6_SERIES_GRIND_II.pdf\#nameddest=tag-03FC}{II.{W.132} (p.~396)} }
\expandafter\def\csname SRef@wbII:eq:W:5.17-A33\endcsname{\href{S6_SERIES_GRIND_II.pdf\#nameddest=tag-03FD}{II.{W.133} (p.~396)} }
\expandafter\def\csname SRef@wbII:eq:W:5.17-A34\endcsname{\href{S6_SERIES_GRIND_II.pdf\#nameddest=tag-03FE}{II.{W.134} (p.~396)} }
\expandafter\def\csname SRef@wbII:eq:W:5.17-A34a\endcsname{\href{S6_SERIES_GRIND_II.pdf\#nameddest=tag-03FF}{II.{W.135} (p.~396)} }
\expandafter\def\csname SRef@wbII:eq:W:5.17-A34b\endcsname{\href{S6_SERIES_GRIND_II.pdf\#nameddest=tag-03FG}{II.{W.136} (p.~397)} }
\expandafter\def\csname SRef@wbII:eq:W:5.17-A34c\endcsname{\href{S6_SERIES_GRIND_II.pdf\#nameddest=tag-03FH}{II.{W.137} (p.~397)} }
\expandafter\def\csname SRef@cdp:eq:A.8\endcsname{\href{S6_SERIES_GRIND_II.pdf\#nameddest=tag-004S}{II.{A.8} (p.~50)} }
\expandafter\def\csname SRef@wbII:eq:W:5.17-A35\endcsname{\href{S6_SERIES_GRIND_II.pdf\#nameddest=tag-03FI}{II.{W.138} (p.~398)} }
\expandafter\def\csname SRef@wbII:eq:W:5.17-A36\endcsname{\href{S6_SERIES_GRIND_II.pdf\#nameddest=tag-03FJ}{II.{W.139} (p.~398)} }
\expandafter\def\csname SRef@wbII:eq:W:5.17-A37\endcsname{\href{S6_SERIES_GRIND_II.pdf\#nameddest=tag-03FK}{II.{W.140} (p.~398)} }
\expandafter\def\csname SRef@wbII:eq:W:5.17-A38\endcsname{\href{S6_SERIES_GRIND_II.pdf\#nameddest=tag-03FL}{II.{W.141} (p.~398)} }
\expandafter\def\csname SRef@wbII:eq:W:5.17f\endcsname{\href{S6_SERIES_GRIND_II.pdf\#nameddest=tag-03GP}{II.{W.44} (p.~399)} }
\expandafter\def\csname SRef@wbII:eq:W:5.18\endcsname{\href{S6_SERIES_GRIND_II.pdf\#nameddest=tag-03HI}{II.{W.45} (p.~400)} }
\expandafter\def\csname SRef@wbII:eq:W:5.19\endcsname{\href{S6_SERIES_GRIND_II.pdf\#nameddest=tag-03HJ}{II.{W.46} (p.~400)} }
\expandafter\def\csname SRef@wbII:eq:W:5.20\endcsname{\href{S6_SERIES_GRIND_II.pdf\#nameddest=tag-03HK}{II.{W.47} (p.~400)} }
\expandafter\def\csname SRef@wbII:eq:CN:5.2\endcsname{\href{S6_SERIES_GRIND_II.pdf\#nameddest=tag-03E5}{II.{W.169} (p.~401)} }
\expandafter\def\csname SRef@cdp:eq:A.9\endcsname{\href{S6_SERIES_GRIND_II.pdf\#nameddest=tag-004T}{II.{A.9} (p.~51)} }
\expandafter\def\csname SRef@wbII:eq:W:5.23\endcsname{\href{S6_SERIES_GRIND_II.pdf\#nameddest=tag-03HL}{II.{W.50} (p.~402)} }
\expandafter\def\csname SRef@wbII:eq:W:5.24\endcsname{\href{S6_SERIES_GRIND_II.pdf\#nameddest=tag-03HM}{II.{W.51} (p.~402)} }
\expandafter\def\csname SRef@wbII:eq:W:5.25\endcsname{\href{S6_SERIES_GRIND_II.pdf\#nameddest=tag-03HN}{II.{W.52} (p.~403)} }
\expandafter\def\csname SRef@wbII:eq:W:5.26\endcsname{\href{S6_SERIES_GRIND_II.pdf\#nameddest=tag-03HO}{II.{W.53} (p.~403)} }
\expandafter\def\csname SRef@wbII:eq:W:5.27\endcsname{\href{S6_SERIES_GRIND_II.pdf\#nameddest=tag-03HP}{II.{W.54} (p.~403)} }
\expandafter\def\csname SRef@wbII:eq:W:5.28\endcsname{\href{S6_SERIES_GRIND_II.pdf\#nameddest=tag-03HQ}{II.{W.55} (p.~403)} }
\expandafter\def\csname SRef@wbII:eq:W:5.29\endcsname{\href{S6_SERIES_GRIND_II.pdf\#nameddest=tag-03HR}{II.{W.56} (p.~403)} }
\expandafter\def\csname SRef@wbII:eq:W:5.30\endcsname{\href{S6_SERIES_GRIND_II.pdf\#nameddest=tag-03HT}{II.{W.57} (p.~404)} }
\expandafter\def\csname SRef@wbII:eq:W:5.31\endcsname{\href{S6_SERIES_GRIND_II.pdf\#nameddest=tag-03HU}{II.{W.58} (p.~404)} }
\expandafter\def\csname SRef@wbII:eq:W:5.32\endcsname{\href{S6_SERIES_GRIND_II.pdf\#nameddest=tag-03HV}{II.{W.59} (p.~404)} }
\expandafter\def\csname SRef@cdp:eq:A.9a\endcsname{\href{S6_SERIES_GRIND_II.pdf\#nameddest=tag-004U}{II.{A.9a} (p.~51)} }
\expandafter\def\csname SRef@wbII:eq:W:5.33\endcsname{\href{S6_SERIES_GRIND_II.pdf\#nameddest=tag-03HW}{II.{W.60} (p.~405)} }
\expandafter\def\csname SRef@wbII:eq:CN:6.1\endcsname{\href{S6_SERIES_GRIND_II.pdf\#nameddest=tag-03E6}{II.{W.176} (p.~406)} }
\expandafter\def\csname SRef@wbII:eq:CN:6.2\endcsname{\href{S6_SERIES_GRIND_II.pdf\#nameddest=tag-03E7}{II.{W.177} (p.~406)} }
\expandafter\def\csname SRef@wbII:eq:CN:6.3\endcsname{\href{S6_SERIES_GRIND_II.pdf\#nameddest=tag-03E8}{II.{W.178} (p.~406)} }
\expandafter\def\csname SRef@wbII:eq:CN:6.4\endcsname{\href{S6_SERIES_GRIND_II.pdf\#nameddest=tag-03E9}{II.{W.179} (p.~406)} }
\expandafter\def\csname SRef@cdp:eq:A.9b\endcsname{\href{S6_SERIES_GRIND_II.pdf\#nameddest=tag-004V}{II.{A.9b} (p.~51)} }
\expandafter\def\csname SRef@cdp:eq:A.NL.1\endcsname{\href{S6_SERIES_GRIND_II.pdf\#nameddest=tag-006J}{II.{A.NL.1} (p.~52)} }
\expandafter\def\csname SRef@invII:eq:E.7\endcsname{\href{S6_SERIES_GRIND_II.pdf\#nameddest=tag-01TW}{II.{E.7} (p.~795)} }
\expandafter\def\csname SRef@cdp:eq:A.NL.2\endcsname{\href{S6_SERIES_GRIND_II.pdf\#nameddest=tag-006K}{II.{A.NL.2} (p.~52)} }
\expandafter\def\csname SRef@cdp:eq:A.NL.3\endcsname{\href{S6_SERIES_GRIND_II.pdf\#nameddest=tag-006L}{II.{A.NL.3} (p.~53)} }
\expandafter\def\csname SRef@cdp:eq:A.NL.4\endcsname{\href{S6_SERIES_GRIND_II.pdf\#nameddest=tag-006M}{II.{A.NL.4} (p.~53)} }
\expandafter\def\csname SRef@cdp:eq:A.NL.5\endcsname{\href{S6_SERIES_GRIND_II.pdf\#nameddest=tag-006N}{II.{A.NL.5} (p.~54)} }
\expandafter\def\csname SRef@cdp:eq:B.1\endcsname{\href{S6_SERIES_GRIND_II.pdf\#nameddest=tag-008G}{II.{B.1} (p.~54)} }
\expandafter\def\csname SRef@cdp:eq:B.2\endcsname{\href{S6_SERIES_GRIND_II.pdf\#nameddest=tag-008H}{II.{B.2} (p.~55)} }
\expandafter\def\csname SRef@cdp:eq:B.3\endcsname{\href{S6_SERIES_GRIND_II.pdf\#nameddest=tag-008I}{II.{B.3} (p.~56)} }
\expandafter\def\csname SRef@cdp:eq:B.4\endcsname{\href{S6_SERIES_GRIND_II.pdf\#nameddest=tag-008J}{II.{B.4} (p.~56)} }
\expandafter\def\csname SRef@cdp:eq:B.5\endcsname{\href{S6_SERIES_GRIND_II.pdf\#nameddest=tag-008K}{II.{B.5} (p.~57)} }
\expandafter\def\csname SRef@cdp:eq:B.6\endcsname{\href{S6_SERIES_GRIND_II.pdf\#nameddest=tag-008L}{II.{B.6} (p.~57)} }
\expandafter\def\csname SRef@cdp:eq:B.7\endcsname{\href{S6_SERIES_GRIND_II.pdf\#nameddest=tag-008M}{II.{B.7} (p.~57)} }
\expandafter\def\csname SRef@cdp:eq:A.10\endcsname{\href{S6_SERIES_GRIND_II.pdf\#nameddest=tag-003Z}{II.{A.10} (p.~58)} }
\expandafter\def\csname SRef@cdp:eq:A.11\endcsname{\href{S6_SERIES_GRIND_II.pdf\#nameddest=tag-0040}{II.{A.11} (p.~58)} }
\expandafter\def\csname SRef@cdp:eq:A.12\endcsname{\href{S6_SERIES_GRIND_II.pdf\#nameddest=tag-0041}{II.{A.12} (p.~59)} }
\expandafter\def\csname SRef@cdp:eq:A.13\endcsname{\href{S6_SERIES_GRIND_II.pdf\#nameddest=tag-0042}{II.{A.13} (p.~60)} }
\expandafter\def\csname SRef@cdp:eq:A.14\endcsname{\href{S6_SERIES_GRIND_II.pdf\#nameddest=tag-0043}{II.{A.14} (p.~60)} }
\expandafter\def\csname SRef@cdp:eq:A.15a\endcsname{\href{S6_SERIES_GRIND_II.pdf\#nameddest=tag-0045}{II.{A.15a} (p.~60)} }
\expandafter\def\csname SRef@cdp:eq:A.15b\endcsname{\href{S6_SERIES_GRIND_II.pdf\#nameddest=tag-0046}{II.{A.15b} (p.~62)} }
\expandafter\def\csname SRef@cdp:eq:A.15\endcsname{\href{S6_SERIES_GRIND_II.pdf\#nameddest=tag-0044}{II.{A.15} (p.~63)} }
\expandafter\def\csname SRef@cdp:eq:A.16\endcsname{\href{S6_SERIES_GRIND_II.pdf\#nameddest=tag-0047}{II.{A.16} (p.~65)} }
\expandafter\def\csname SRef@cdp:eq:A.17\endcsname{\href{S6_SERIES_GRIND_II.pdf\#nameddest=tag-0048}{II.{A.17} (p.~66)} }
\expandafter\def\csname SRef@cdp:eq:A.18\endcsname{\href{S6_SERIES_GRIND_II.pdf\#nameddest=tag-0049}{II.{A.18} (p.~67)} }
\expandafter\def\csname SRef@cdp:eq:A.19\endcsname{\href{S6_SERIES_GRIND_II.pdf\#nameddest=tag-004A}{II.{A.19} (p.~67)} }
\expandafter\def\csname SRef@cdp:eq:A.20\endcsname{\href{S6_SERIES_GRIND_II.pdf\#nameddest=tag-004C}{II.{A.20} (p.~68)} }
\expandafter\def\csname SRef@cdp:eq:A.21\endcsname{\href{S6_SERIES_GRIND_II.pdf\#nameddest=tag-004D}{II.{A.21} (p.~69)} }
\expandafter\def\csname SRef@cdp:eq:A.22\endcsname{\href{S6_SERIES_GRIND_II.pdf\#nameddest=tag-004E}{II.{A.22} (p.~69)} }
\expandafter\def\csname SRef@cdp:eq:A.23\endcsname{\href{S6_SERIES_GRIND_II.pdf\#nameddest=tag-004F}{II.{A.23} (p.~70)} }
\expandafter\def\csname SRef@cdp:eq:A.24\endcsname{\href{S6_SERIES_GRIND_II.pdf\#nameddest=tag-004G}{II.{A.24} (p.~71)} }
\expandafter\def\csname SRef@cdp:eq:A.25\endcsname{\href{S6_SERIES_GRIND_II.pdf\#nameddest=tag-004H}{II.{A.25} (p.~74)} }
\expandafter\def\csname SRef@cdp:eq:A.26\endcsname{\href{S6_SERIES_GRIND_II.pdf\#nameddest=tag-004I}{II.{A.26} (p.~74)} }
\expandafter\def\csname SRef@cdp:eq:A.27\endcsname{\href{S6_SERIES_GRIND_II.pdf\#nameddest=tag-004J}{II.{A.27} (p.~76)} }
\expandafter\def\csname SRef@cdp:eq:A.28\endcsname{\href{S6_SERIES_GRIND_II.pdf\#nameddest=tag-004K}{II.{A.28} (p.~76)} }
\expandafter\def\csname SRef@cdp:eq:A.29\endcsname{\href{S6_SERIES_GRIND_II.pdf\#nameddest=tag-004L}{II.{A.29} (p.~76)} }
\expandafter\def\csname SRef@cdp:eq:A.30\endcsname{\href{S6_SERIES_GRIND_II.pdf\#nameddest=tag-004N}{II.{A.30} (p.~77)} }
\expandafter\def\csname SRef@cdp:eq:A.RC.1\endcsname{\href{S6_SERIES_GRIND_II.pdf\#nameddest=tag-006U}{II.{A.RC.1} (p.~78)} }
\expandafter\def\csname SRef@push:src:E.2.4\endcsname{\href{S6_SERIES_GRIND_II.pdf\#nameddest=tag-033C}{II.13.1 (p.~309)} }
\expandafter\def\csname SRef@push:src:E.2.24\endcsname{\href{S6_SERIES_GRIND_II.pdf\#nameddest=tag-033A}{II.13.21 (p.~348)} }
\expandafter\def\csname SRef@wbII:cor:surface-top-base-change\endcsname{\href{S6_SERIES_GRIND_II.pdf\#nameddest=tag-03DV}{II.16.10 (p.~374)} }
\expandafter\def\csname SRef@wbII:cor:candidate-pic0-route\endcsname{\href{S6_SERIES_GRIND_II.pdf\#nameddest=tag-03DS}{II.16.22 (p.~399)} }
\expandafter\def\csname SRef@wbII:cor:candidate-conductor-actual\endcsname{\href{S6_SERIES_GRIND_II.pdf\#nameddest=tag-03DR}{II.16.24 (p.~401)} }
\expandafter\def\csname SRef@wbII:cor:corrected-cdp-relative\endcsname{\href{S6_SERIES_GRIND_II.pdf\#nameddest=tag-03DT}{II.16.28 (p.~403)} }
\expandafter\def\csname SRef@wbII:cor:relative\endcsname{\href{S6_SERIES_GRIND_II.pdf\#nameddest=tag-03DT}{II.16.28 (p.~403)} }
\expandafter\def\csname SRef@wbII:cor:normal-crossing-kernel\endcsname{\href{S6_SERIES_GRIND_II.pdf\#nameddest=tag-03DU}{II.16.3 (p.~367)} }
\expandafter\def\csname SRef@wbII:test:cdp-comparison\endcsname{\href{S6_SERIES_GRIND_II.pdf\#nameddest=tag-03IV}{II.16.30 (p.~405)} }
\expandafter\def\csname SRef@cdp:A:filtration1\endcsname{\href{S6_SERIES_GRIND_II.pdf\#nameddest=tag-003W}{II.1 (p.~48)} }
\expandafter\def\csname SRef@cdp:eq:E13-residue\endcsname{\href{S6_SERIES_GRIND_II.pdf\#nameddest=tag-01A6}{II.10 (p.~296)} }
\expandafter\def\csname SRef@cdp:eq:E14-elliptic-family\endcsname{\href{S6_SERIES_GRIND_II.pdf\#nameddest=tag-01A8}{II.11 (p.~297)} }
\expandafter\def\csname SRef@cdp:eq:E14-canonical\endcsname{\href{S6_SERIES_GRIND_II.pdf\#nameddest=tag-01A7}{II.12 (p.~298)} }
\expandafter\def\csname SRef@cdp:eq:E14-log-coho\endcsname{\href{S6_SERIES_GRIND_II.pdf\#nameddest=tag-01AA}{II.13 (p.~299)} }
\expandafter\def\csname SRef@cdp:eq:E14-jet-map\endcsname{\href{S6_SERIES_GRIND_II.pdf\#nameddest=tag-01A9}{II.14 (p.~299)} }
\expandafter\def\csname SRef@cdp:eq:E15-positive\endcsname{\href{S6_SERIES_GRIND_II.pdf\#nameddest=tag-01AB}{II.15 (p.~301)} }
\expandafter\def\csname SRef@cdp:eq:E15-products\endcsname{\href{S6_SERIES_GRIND_II.pdf\#nameddest=tag-01AC}{II.16 (p.~301)} }
\expandafter\def\csname SRef@wbII:eq:analytic-toric-open-directions\endcsname{\href{S6_SERIES_GRIND_II.pdf\#nameddest=tag-03I5}{II.17 (p.~396)} }
\expandafter\def\csname SRef@wbII:eq:analytic-residual-descent-factor\endcsname{\href{S6_SERIES_GRIND_II.pdf\#nameddest=tag-03I4}{II.18 (p.~397)} }
\expandafter\def\csname SRef@wbII:eq:analytic-factor-horizontal-winding\endcsname{\href{S6_SERIES_GRIND_II.pdf\#nameddest=tag-03I3}{II.19 (p.~397)} }
\expandafter\def\csname SRef@cdp:A:filtration2\endcsname{\href{S6_SERIES_GRIND_II.pdf\#nameddest=tag-003X}{II.2 (p.~48)} }
\expandafter\def\csname SRef@eq:P4chi-resolution\endcsname{\href{S6_SERIES_GRIND_II.pdf\#nameddest=tag-01OD}{II.3 (p.~141)} }
\expandafter\def\csname SRef@eq:P4exceptional-genus\endcsname{\href{S6_SERIES_GRIND_II.pdf\#nameddest=tag-01OE}{II.4 (p.~141)} }
\expandafter\def\csname SRef@cdp:eq:E-Leray-map\endcsname{\href{S6_SERIES_GRIND_II.pdf\#nameddest=tag-019Y}{II.5 (p.~285)} }
\expandafter\def\csname SRef@cdp:eq:E-surj-embedding\endcsname{\href{S6_SERIES_GRIND_II.pdf\#nameddest=tag-01A0}{II.6 (p.~286)} }
\expandafter\def\csname SRef@cdp:eq:E-log-period\endcsname{\href{S6_SERIES_GRIND_II.pdf\#nameddest=tag-019Z}{II.7 (p.~288)} }
\expandafter\def\csname SRef@cdp:eq:E-theta-general\endcsname{\href{S6_SERIES_GRIND_II.pdf\#nameddest=tag-01A1}{II.8 (p.~289)} }
\expandafter\def\csname SRef@cdp:eq:E13-connection\endcsname{\href{S6_SERIES_GRIND_II.pdf\#nameddest=tag-01A5}{II.9 (p.~294)} }
\expandafter\def\csname SRef@wbII:lem:torsion-identification\endcsname{\href{S6_SERIES_GRIND_II.pdf\#nameddest=tag-03IF}{II.16.1 (p.~363)} }
\expandafter\def\csname SRef@wbII:lem:picard-crossed-defect\endcsname{\href{S6_SERIES_GRIND_II.pdf\#nameddest=tag-03IC}{II.16.14 (p.~377)} }
\expandafter\def\csname SRef@wbII:lem:picard-finite-norm\endcsname{\href{S6_SERIES_GRIND_II.pdf\#nameddest=tag-03IE}{II.16.15 (p.~378)} }
\expandafter\def\csname SRef@wbII:lem:picard-cusp-defect\endcsname{\href{S6_SERIES_GRIND_II.pdf\#nameddest=tag-03ID}{II.16.16 (p.~380)} }
\expandafter\def\csname SRef@wbII:lem:analytic-support-specialization\endcsname{\href{S6_SERIES_GRIND_II.pdf\#nameddest=tag-03I6}{II.16.20 (p.~394)} }
\expandafter\def\csname SRef@wbII:lem:generic-endpoints\endcsname{\href{S6_SERIES_GRIND_II.pdf\#nameddest=tag-03IA}{II.16.25 (p.~401)} }
\expandafter\def\csname SRef@wbII:lem:cdp-generic-endpoints\endcsname{\href{S6_SERIES_GRIND_II.pdf\#nameddest=tag-03I7}{II.16.26 (p.~402)} }
\expandafter\def\csname SRef@wbII:lem:gluing-bundle\endcsname{\href{S6_SERIES_GRIND_II.pdf\#nameddest=tag-03IB}{II.16.4 (p.~369)} }
\expandafter\def\csname SRef@wbII:lem:extending-gluing\endcsname{\href{S6_SERIES_GRIND_II.pdf\#nameddest=tag-03I9}{II.16.5 (p.~370)} }
\expandafter\def\csname SRef@wbII:lem:curve-flatness\endcsname{\href{S6_SERIES_GRIND_II.pdf\#nameddest=tag-03I8}{II.16.8 (p.~373)} }
\expandafter\def\csname SRef@inv:E2.3\endcsname{\href{S6_SERIES_GRIND_II.pdf\#nameddest=tag-01SW}{II.26.1 (p.~795)} }
\expandafter\def\csname SRef@cdp:st:D.2.1\endcsname{\href{S6_SERIES_GRIND_II.pdf\#nameddest=tag-01G8}{II.11.1 (p.~272)} }
\expandafter\def\csname SRef@cdp:st:D.2.3\endcsname{\href{S6_SERIES_GRIND_II.pdf\#nameddest=tag-01GA}{II.11.2 (p.~272)} }
\expandafter\def\csname SRef@cdp:st:D.2.4\endcsname{\href{S6_SERIES_GRIND_II.pdf\#nameddest=tag-01GB}{II.11.3 (p.~275)} }
\expandafter\def\csname SRef@cdp:st:D.2.5\endcsname{\href{S6_SERIES_GRIND_II.pdf\#nameddest=tag-01GC}{II.11.4 (p.~276)} }
\expandafter\def\csname SRef@cdp:st:D.2.6\endcsname{\href{S6_SERIES_GRIND_II.pdf\#nameddest=tag-01GD}{II.11.5 (p.~278)} }
\expandafter\def\csname SRef@cdp:st:D.2.7\endcsname{\href{S6_SERIES_GRIND_II.pdf\#nameddest=tag-01GE}{II.11.6 (p.~279)} }
\expandafter\def\csname SRef@cdp:st:D.2.8\endcsname{\href{S6_SERIES_GRIND_II.pdf\#nameddest=tag-01GF}{II.11.7 (p.~283)} }
\expandafter\def\csname SRef@cdp:st:D.2.2\endcsname{\href{S6_SERIES_GRIND_II.pdf\#nameddest=tag-01G9}{II.11.8 (p.~284)} }
\expandafter\def\csname SRef@cdp:st:E5\endcsname{\href{S6_SERIES_GRIND_II.pdf\#nameddest=tag-01GU}{II.12.1 (p.~285)} }
\expandafter\def\csname SRef@cdp:st:E15\endcsname{\href{S6_SERIES_GRIND_II.pdf\#nameddest=tag-01GM}{II.12.10 (p.~300)} }
\expandafter\def\csname SRef@cdp:st:E16\endcsname{\href{S6_SERIES_GRIND_II.pdf\#nameddest=tag-01GN}{II.12.11 (p.~304)} }
\expandafter\def\csname SRef@cdp:st:E17\endcsname{\href{S6_SERIES_GRIND_II.pdf\#nameddest=tag-01GO}{II.12.12 (p.~305)} }
\expandafter\def\csname SRef@cdp:st:E18\endcsname{\href{S6_SERIES_GRIND_II.pdf\#nameddest=tag-01GP}{II.12.13 (p.~307)} }
\expandafter\def\csname SRef@cdp:st:E6\endcsname{\href{S6_SERIES_GRIND_II.pdf\#nameddest=tag-01GV}{II.12.2 (p.~286)} }
\expandafter\def\csname SRef@cdp:st:E7\endcsname{\href{S6_SERIES_GRIND_II.pdf\#nameddest=tag-01GW}{II.12.3 (p.~288)} }
\expandafter\def\csname SRef@cdp:st:E8\endcsname{\href{S6_SERIES_GRIND_II.pdf\#nameddest=tag-01GX}{II.12.4 (p.~289)} }
\expandafter\def\csname SRef@cdp:st:E10\endcsname{\href{S6_SERIES_GRIND_II.pdf\#nameddest=tag-01GH}{II.12.5 (p.~290)} }
\expandafter\def\csname SRef@cdp:st:E11\endcsname{\href{S6_SERIES_GRIND_II.pdf\#nameddest=tag-01GI}{II.12.6 (p.~292)} }
\expandafter\def\csname SRef@cdp:st:E12\endcsname{\href{S6_SERIES_GRIND_II.pdf\#nameddest=tag-01GJ}{II.12.7 (p.~293)} }
\expandafter\def\csname SRef@cdp:st:E13\endcsname{\href{S6_SERIES_GRIND_II.pdf\#nameddest=tag-01GK}{II.12.8 (p.~294)} }
\expandafter\def\csname SRef@cdp:st:E14\endcsname{\href{S6_SERIES_GRIND_II.pdf\#nameddest=tag-01GL}{II.12.9 (p.~297)} }
\expandafter\def\csname SRef@push:src:E.2.13\endcsname{\href{S6_SERIES_GRIND_II.pdf\#nameddest=tag-032Y}{II.13.10 (p.~325)} }
\expandafter\def\csname SRef@push:src:E.2.14\endcsname{\href{S6_SERIES_GRIND_II.pdf\#nameddest=tag-032Z}{II.13.11 (p.~326)} }
\expandafter\def\csname SRef@push:src:E.2.15\endcsname{\href{S6_SERIES_GRIND_II.pdf\#nameddest=tag-0330}{II.13.12 (p.~328)} }
\expandafter\def\csname SRef@push:src:E.2.16\endcsname{\href{S6_SERIES_GRIND_II.pdf\#nameddest=tag-0331}{II.13.13 (p.~330)} }
\expandafter\def\csname SRef@push:src:E.2.17\endcsname{\href{S6_SERIES_GRIND_II.pdf\#nameddest=tag-0332}{II.13.14 (p.~332)} }
\expandafter\def\csname SRef@push:src:E.2.18\endcsname{\href{S6_SERIES_GRIND_II.pdf\#nameddest=tag-0333}{II.13.15 (p.~334)} }
\expandafter\def\csname SRef@push:src:E.2.19\endcsname{\href{S6_SERIES_GRIND_II.pdf\#nameddest=tag-0334}{II.13.16 (p.~335)} }
\expandafter\def\csname SRef@push:src:E.2.20\endcsname{\href{S6_SERIES_GRIND_II.pdf\#nameddest=tag-0336}{II.13.17 (p.~338)} }
\expandafter\def\csname SRef@push:src:E.2.21\endcsname{\href{S6_SERIES_GRIND_II.pdf\#nameddest=tag-0337}{II.13.18 (p.~341)} }
\expandafter\def\csname SRef@push:src:E.2.22\endcsname{\href{S6_SERIES_GRIND_II.pdf\#nameddest=tag-0338}{II.13.19 (p.~343)} }
\expandafter\def\csname SRef@push:src:E.2.5\endcsname{\href{S6_SERIES_GRIND_II.pdf\#nameddest=tag-033D}{II.13.2 (p.~310)} }
\expandafter\def\csname SRef@push:src:E.5.1\endcsname{\href{S6_SERIES_GRIND_II.pdf\#nameddest=tag-033N}{II.13.22 (p.~349)} }
\expandafter\def\csname SRef@push:src:E.5.2\endcsname{\href{S6_SERIES_GRIND_II.pdf\#nameddest=tag-033O}{II.13.23 (p.~349)} }
\expandafter\def\csname SRef@push:src:E.2.6\endcsname{\href{S6_SERIES_GRIND_II.pdf\#nameddest=tag-033E}{II.13.3 (p.~311)} }
\expandafter\def\csname SRef@push:src:E.2.7\endcsname{\href{S6_SERIES_GRIND_II.pdf\#nameddest=tag-033F}{II.13.4 (p.~313)} }
\expandafter\def\csname SRef@push:src:E.2.8\endcsname{\href{S6_SERIES_GRIND_II.pdf\#nameddest=tag-033G}{II.13.5 (p.~316)} }
\expandafter\def\csname SRef@push:src:E.2.9\endcsname{\href{S6_SERIES_GRIND_II.pdf\#nameddest=tag-033H}{II.13.6 (p.~318)} }
\expandafter\def\csname SRef@push:src:E.2.10\endcsname{\href{S6_SERIES_GRIND_II.pdf\#nameddest=tag-032V}{II.13.7 (p.~320)} }
\expandafter\def\csname SRef@push:src:E.2.11\endcsname{\href{S6_SERIES_GRIND_II.pdf\#nameddest=tag-032W}{II.13.8 (p.~321)} }
\expandafter\def\csname SRef@push:src:E.2.12\endcsname{\href{S6_SERIES_GRIND_II.pdf\#nameddest=tag-032X}{II.13.9 (p.~323)} }
\expandafter\def\csname SRef@inv:E2.2\endcsname{\href{S6_SERIES_GRIND_II.pdf\#nameddest=tag-01SV}{II.14.1 (p.~354)} }
\expandafter\def\csname SRef@inv:E2.4\endcsname{\href{S6_SERIES_GRIND_II.pdf\#nameddest=tag-01SX}{II.14.3 (p.~361)} }
\expandafter\def\csname SRef@wbII:prop:pic0-pullback\endcsname{\href{S6_SERIES_GRIND_II.pdf\#nameddest=tag-03IO}{II.16.11 (p.~376)} }
\expandafter\def\csname SRef@wbII:prop:all-pic-restriction\endcsname{\href{S6_SERIES_GRIND_II.pdf\#nameddest=tag-03IG}{II.16.12 (p.~376)} }
\expandafter\def\csname SRef@wbII:prop:leray-edge-restriction\endcsname{\href{S6_SERIES_GRIND_II.pdf\#nameddest=tag-03IM}{II.16.13 (p.~376)} }
\expandafter\def\csname SRef@wbII:prop:direct-boundary-picard\endcsname{\href{S6_SERIES_GRIND_II.pdf\#nameddest=tag-03IL}{II.16.17 (p.~381)} }
\expandafter\def\csname SRef@wbII:prop:local-cusp-picard-obstruction\endcsname{\href{S6_SERIES_GRIND_II.pdf\#nameddest=tag-03IN}{II.16.18 (p.~383)} }
\expandafter\def\csname SRef@wbII:prop:analytic-ah-boundary\endcsname{\href{S6_SERIES_GRIND_II.pdf\#nameddest=tag-03IH}{II.16.19 (p.~387)} }
\expandafter\def\csname SRef@wbII:prop:relative-local\endcsname{\href{S6_SERIES_GRIND_II.pdf\#nameddest=tag-03IP}{II.16.31 (p.~406)} }
\expandafter\def\csname SRef@wbII:prop:cartier-duality-absolute\endcsname{\href{S6_SERIES_GRIND_II.pdf\#nameddest=tag-03IJ}{II.16.7 (p.~372)} }
\expandafter\def\csname SRef@wbII:prop:analytic-base-change\endcsname{\href{S6_SERIES_GRIND_II.pdf\#nameddest=tag-03II}{II.16.9 (p.~374)} }
\expandafter\def\csname SRef@cdp:st:C.1.1\endcsname{\href{S6_SERIES_GRIND_II.pdf\#nameddest=tag-01BN}{II.3.1 (p.~18)} }
\expandafter\def\csname SRef@cdp:st:C.1.3\endcsname{\href{S6_SERIES_GRIND_II.pdf\#nameddest=tag-01BP}{II.3.2 (p.~19)} }
\expandafter\def\csname SRef@cdp:sec:E\endcsname{\href{S6_SERIES_GRIND_II.pdf\#nameddest=tag-01BP}{II.3 (p.~19)} }
\expandafter\def\csname SRef@cdp:st:E1\endcsname{\href{S6_SERIES_GRIND_II.pdf\#nameddest=tag-01GG}{II.3.3 (p.~19)} }
\expandafter\def\csname SRef@cdp:st:D.1.2\endcsname{\href{S6_SERIES_GRIND_II.pdf\#nameddest=tag-01G0}{II.4.1 (p.~20)} }
\expandafter\def\csname SRef@cdp:st:D.1.4\endcsname{\href{S6_SERIES_GRIND_II.pdf\#nameddest=tag-01G2}{II.4.10 (p.~28)} }
\expandafter\def\csname SRef@cdp:st:D.1.13\endcsname{\href{S6_SERIES_GRIND_II.pdf\#nameddest=tag-01FY}{II.4.11 (p.~28)} }
\expandafter\def\csname SRef@cdp:st:D.1.15\endcsname{\href{S6_SERIES_GRIND_II.pdf\#nameddest=tag-01FZ}{II.4.12 (p.~30)} }
\expandafter\def\csname SRef@cdp:st:D.1.5\endcsname{\href{S6_SERIES_GRIND_II.pdf\#nameddest=tag-01G3}{II.4.2 (p.~21)} }
\expandafter\def\csname SRef@cdp:st:D.1.6\endcsname{\href{S6_SERIES_GRIND_II.pdf\#nameddest=tag-01G4}{II.4.3 (p.~22)} }
\expandafter\def\csname SRef@cdp:st:D.1.7\endcsname{\href{S6_SERIES_GRIND_II.pdf\#nameddest=tag-01G5}{II.4.4 (p.~23)} }
\expandafter\def\csname SRef@cdp:st:D.1.8\endcsname{\href{S6_SERIES_GRIND_II.pdf\#nameddest=tag-01G6}{II.4.5 (p.~23)} }
\expandafter\def\csname SRef@cdp:st:D.1.9\endcsname{\href{S6_SERIES_GRIND_II.pdf\#nameddest=tag-01G7}{II.4.6 (p.~24)} }
\expandafter\def\csname SRef@cdp:st:D.1.10\endcsname{\href{S6_SERIES_GRIND_II.pdf\#nameddest=tag-01FV}{II.4.7 (p.~25)} }
\expandafter\def\csname SRef@cdp:st:D.1.11\endcsname{\href{S6_SERIES_GRIND_II.pdf\#nameddest=tag-01FW}{II.4.8 (p.~25)} }
\expandafter\def\csname SRef@cdp:st:D.1.12\endcsname{\href{S6_SERIES_GRIND_II.pdf\#nameddest=tag-01FX}{II.4.9 (p.~27)} }
\expandafter\def\csname SRef@cdp:st:C.3.2\endcsname{\href{S6_SERIES_GRIND_II.pdf\#nameddest=tag-01F8}{II.5.1 (p.~31)} }
\expandafter\def\csname SRef@cdp:st:C.3.11\endcsname{\href{S6_SERIES_GRIND_II.pdf\#nameddest=tag-01F6}{II.5.10 (p.~44)} }
\expandafter\def\csname SRef@cdp:st:C.3.5\endcsname{\href{S6_SERIES_GRIND_II.pdf\#nameddest=tag-01FB}{II.5.11 (p.~45)} }
\expandafter\def\csname SRef@cdp:st:C.3.6\endcsname{\href{S6_SERIES_GRIND_II.pdf\#nameddest=tag-01FC}{II.5.12 (p.~46)} }
\expandafter\def\csname SRef@cdp:st:C.3.3\endcsname{\href{S6_SERIES_GRIND_II.pdf\#nameddest=tag-01F9}{II.5.2 (p.~32)} }
\expandafter\def\csname SRef@cdp:st:C.3.4\endcsname{\href{S6_SERIES_GRIND_II.pdf\#nameddest=tag-01FA}{II.5.3 (p.~33)} }
\expandafter\def\csname SRef@cdp:st:C.3.7\endcsname{\href{S6_SERIES_GRIND_II.pdf\#nameddest=tag-01FD}{II.5.4 (p.~33)} }
\expandafter\def\csname SRef@cdp:st:C.3.8\endcsname{\href{S6_SERIES_GRIND_II.pdf\#nameddest=tag-01FE}{II.5.5 (p.~34)} }
\expandafter\def\csname SRef@cdp:st:C.3.9\endcsname{\href{S6_SERIES_GRIND_II.pdf\#nameddest=tag-01FF}{II.5.6 (p.~35)} }
\expandafter\def\csname SRef@cdp:st:C.3.10\endcsname{\href{S6_SERIES_GRIND_II.pdf\#nameddest=tag-01F5}{II.5.7 (p.~38)} }
\expandafter\def\csname SRef@cdp:st:C.3.12\endcsname{\href{S6_SERIES_GRIND_II.pdf\#nameddest=tag-01F7}{II.5.8 (p.~41)} }
\expandafter\def\csname SRef@cdp:st:C.3.1\endcsname{\href{S6_SERIES_GRIND_II.pdf\#nameddest=tag-01F4}{II.5.9 (p.~43)} }
\expandafter\def\csname SRef@cdp:st:A.1.1\endcsname{\href{S6_SERIES_GRIND_II.pdf\#nameddest=tag-01AD}{II.6.1 (p.~48)} }
\expandafter\def\csname SRef@wbII:lem:conormal\endcsname{\href{S6_SERIES_GRIND_II.pdf\#nameddest=tag-01AD}{II.6.1 (p.~48)} }
\expandafter\def\csname SRef@wbII:prop:criterion\endcsname{\href{S6_SERIES_GRIND_II.pdf\#nameddest=tag-01AD}{II.6.1 (p.~48)} }
\expandafter\def\csname SRef@wbII:prop:filtration\endcsname{\href{S6_SERIES_GRIND_II.pdf\#nameddest=tag-01AD}{II.6.1 (p.~48)} }
\expandafter\def\csname SRef@wbII:lem:conormal-injective\endcsname{\href{S6_SERIES_GRIND_II.pdf\#nameddest=tag-01AD}{II.6.1 (p.~48)} }
\expandafter\def\csname SRef@wbII:prop:missing-cotangent-filtration\endcsname{\href{S6_SERIES_GRIND_II.pdf\#nameddest=tag-01AD}{II.6.1 (p.~48)} }
\expandafter\def\csname SRef@wbII:prop:missing-global-section-criterion\endcsname{\href{S6_SERIES_GRIND_II.pdf\#nameddest=tag-01AD}{II.6.1 (p.~48)} }
\expandafter\def\csname SRef@cdp:st:B3\endcsname{\href{S6_SERIES_GRIND_II.pdf\#nameddest=tag-01BM}{II.6.10 (p.~58)} }
\expandafter\def\csname SRef@cdp:st:A.3.1\endcsname{\href{S6_SERIES_GRIND_II.pdf\#nameddest=tag-01AL}{II.6.11 (p.~58)} }
\expandafter\def\csname SRef@cdp:st:A.3.2\endcsname{\href{S6_SERIES_GRIND_II.pdf\#nameddest=tag-01AM}{II.6.12 (p.~59)} }
\expandafter\def\csname SRef@cdp:st:A.3.3\endcsname{\href{S6_SERIES_GRIND_II.pdf\#nameddest=tag-01AN}{II.6.13 (p.~60)} }
\expandafter\def\csname SRef@cdp:st:A.3.4\endcsname{\href{S6_SERIES_GRIND_II.pdf\#nameddest=tag-01AO}{II.6.14 (p.~60)} }
\expandafter\def\csname SRef@cdp:st:A.3.5\endcsname{\href{S6_SERIES_GRIND_II.pdf\#nameddest=tag-01AP}{II.6.15 (p.~62)} }
\expandafter\def\csname SRef@cdp:st:A.4.1\endcsname{\href{S6_SERIES_GRIND_II.pdf\#nameddest=tag-01AQ}{II.6.16 (p.~63)} }
\expandafter\def\csname SRef@cdp:st:A.4.2\endcsname{\href{S6_SERIES_GRIND_II.pdf\#nameddest=tag-01AR}{II.6.17 (p.~64)} }
\expandafter\def\csname SRef@cdp:st:A.4.3\endcsname{\href{S6_SERIES_GRIND_II.pdf\#nameddest=tag-01AS}{II.6.18 (p.~64)} }
\expandafter\def\csname SRef@cdp:st:A.4.4\endcsname{\href{S6_SERIES_GRIND_II.pdf\#nameddest=tag-01AT}{II.6.19 (p.~65)} }
\expandafter\def\csname SRef@cdp:st:A.2.1\endcsname{\href{S6_SERIES_GRIND_II.pdf\#nameddest=tag-01AF}{II.6.2 (p.~48)} }
\expandafter\def\csname SRef@cdp:st:A.4.5\endcsname{\href{S6_SERIES_GRIND_II.pdf\#nameddest=tag-01AU}{II.6.20 (p.~66)} }
\expandafter\def\csname SRef@cdp:st:A.4.6\endcsname{\href{S6_SERIES_GRIND_II.pdf\#nameddest=tag-01AV}{II.6.21 (p.~67)} }
\expandafter\def\csname SRef@cdp:st:A.4.7\endcsname{\href{S6_SERIES_GRIND_II.pdf\#nameddest=tag-01AW}{II.6.22 (p.~68)} }
\expandafter\def\csname SRef@cdp:st:A.4.8\endcsname{\href{S6_SERIES_GRIND_II.pdf\#nameddest=tag-01AX}{II.6.23 (p.~69)} }
\expandafter\def\csname SRef@cdp:st:A.4.9\endcsname{\href{S6_SERIES_GRIND_II.pdf\#nameddest=tag-01AY}{II.6.24 (p.~70)} }
\expandafter\def\csname SRef@cdp:st:A.2.2\endcsname{\href{S6_SERIES_GRIND_II.pdf\#nameddest=tag-01AG}{II.6.3 (p.~49)} }
\expandafter\def\csname SRef@wbII:lem:conductor\endcsname{\href{S6_SERIES_GRIND_II.pdf\#nameddest=tag-01AG}{II.6.3 (p.~49)} }
\expandafter\def\csname SRef@wbII:lem:conductor-square\endcsname{\href{S6_SERIES_GRIND_II.pdf\#nameddest=tag-01AG}{II.6.3 (p.~49)} }
\expandafter\def\csname SRef@cdp:st:A.1.2\endcsname{\href{S6_SERIES_GRIND_II.pdf\#nameddest=tag-01AE}{II.6.4 (p.~50)} }
\expandafter\def\csname SRef@cdp:st:A.2.4\endcsname{\href{S6_SERIES_GRIND_II.pdf\#nameddest=tag-01AI}{II.6.5 (p.~50)} }
\expandafter\def\csname SRef@cdp:st:A.2.5\endcsname{\href{S6_SERIES_GRIND_II.pdf\#nameddest=tag-01AJ}{II.6.6 (p.~51)} }
\expandafter\def\csname SRef@cdp:st:A.2.6\endcsname{\href{S6_SERIES_GRIND_II.pdf\#nameddest=tag-01AK}{II.6.7 (p.~52)} }
\expandafter\def\csname SRef@cdp:st:B1\endcsname{\href{S6_SERIES_GRIND_II.pdf\#nameddest=tag-01BK}{II.6.8 (p.~54)} }
\expandafter\def\csname SRef@wbII:thm:nodal-test\endcsname{\href{S6_SERIES_GRIND_II.pdf\#nameddest=tag-01BK}{II.6.8 (p.~54)} }
\expandafter\def\csname SRef@wbII:prop:family-geometry\endcsname{\href{S6_SERIES_GRIND_II.pdf\#nameddest=tag-01BK}{II.6.8 (p.~54)} }
\expandafter\def\csname SRef@wbII:prop:nodal-counterexample\endcsname{\href{S6_SERIES_GRIND_II.pdf\#nameddest=tag-01BK}{II.6.8 (p.~54)} }
\expandafter\def\csname SRef@cdp:st:B2\endcsname{\href{S6_SERIES_GRIND_II.pdf\#nameddest=tag-01BL}{II.6.9 (p.~56)} }
\expandafter\def\csname SRef@cdp:st:A.5.1\endcsname{\href{S6_SERIES_GRIND_II.pdf\#nameddest=tag-01AZ}{II.7.1 (p.~72)} }
\expandafter\def\csname SRef@cdp:st:A.5.5\endcsname{\href{S6_SERIES_GRIND_II.pdf\#nameddest=tag-01BF}{II.7.10 (p.~84)} }
\expandafter\def\csname SRef@cdp:st:A.5.7\endcsname{\href{S6_SERIES_GRIND_II.pdf\#nameddest=tag-01BH}{II.7.11 (p.~88)} }
\expandafter\def\csname SRef@cdp:st:A.5.8\endcsname{\href{S6_SERIES_GRIND_II.pdf\#nameddest=tag-01BI}{II.7.12 (p.~89)} }
\expandafter\def\csname SRef@cdp:st:A.5.14\endcsname{\href{S6_SERIES_GRIND_II.pdf\#nameddest=tag-01B4}{II.7.13 (p.~90)} }
\expandafter\def\csname SRef@cdp:st:A.5.10\endcsname{\href{S6_SERIES_GRIND_II.pdf\#nameddest=tag-01B0}{II.7.14 (p.~92)} }
\expandafter\def\csname SRef@cdp:st:A.5.20\endcsname{\href{S6_SERIES_GRIND_II.pdf\#nameddest=tag-01BB}{II.7.15 (p.~94)} }
\expandafter\def\csname SRef@cdp:st:A.5.15\endcsname{\href{S6_SERIES_GRIND_II.pdf\#nameddest=tag-01B5}{II.7.16 (p.~95)} }
\expandafter\def\csname SRef@cdp:st:A.5.17\endcsname{\href{S6_SERIES_GRIND_II.pdf\#nameddest=tag-01B7}{II.7.17 (p.~98)} }
\expandafter\def\csname SRef@cdp:st:A.5.18\endcsname{\href{S6_SERIES_GRIND_II.pdf\#nameddest=tag-01B8}{II.7.18 (p.~101)} }
\expandafter\def\csname SRef@cdp:st:A.5.19\endcsname{\href{S6_SERIES_GRIND_II.pdf\#nameddest=tag-01B9}{II.7.19 (p.~108)} }
\expandafter\def\csname SRef@cdp:st:A.5.2\endcsname{\href{S6_SERIES_GRIND_II.pdf\#nameddest=tag-01BA}{II.7.2 (p.~73)} }
\expandafter\def\csname SRef@cdp:st:A.5.16\endcsname{\href{S6_SERIES_GRIND_II.pdf\#nameddest=tag-01B6}{II.7.20 (p.~111)} }
\expandafter\def\csname SRef@cdp:st:A.5.21\endcsname{\href{S6_SERIES_GRIND_II.pdf\#nameddest=tag-01BC}{II.7.21 (p.~117)} }
\expandafter\def\csname SRef@cdp:st:A.5.3\endcsname{\href{S6_SERIES_GRIND_II.pdf\#nameddest=tag-01BD}{II.7.3 (p.~74)} }
\expandafter\def\csname SRef@cdp:st:A.5.4\endcsname{\href{S6_SERIES_GRIND_II.pdf\#nameddest=tag-01BE}{II.7.4 (p.~76)} }
\expandafter\def\csname SRef@cdp:st:A.5.6\endcsname{\href{S6_SERIES_GRIND_II.pdf\#nameddest=tag-01BG}{II.7.5 (p.~77)} }
\expandafter\def\csname SRef@cdp:st:A.5.9\endcsname{\href{S6_SERIES_GRIND_II.pdf\#nameddest=tag-01BJ}{II.7.6 (p.~78)} }
\expandafter\def\csname SRef@cdp:st:A.5.11\endcsname{\href{S6_SERIES_GRIND_II.pdf\#nameddest=tag-01B1}{II.7.7 (p.~80)} }
\expandafter\def\csname SRef@cdp:st:A.5.12\endcsname{\href{S6_SERIES_GRIND_II.pdf\#nameddest=tag-01B2}{II.7.8 (p.~82)} }
\expandafter\def\csname SRef@cdp:st:A.5.13\endcsname{\href{S6_SERIES_GRIND_II.pdf\#nameddest=tag-01B3}{II.7.9 (p.~83)} }
\expandafter\def\csname SRef@cdp:st:C.2.NC.2\endcsname{\href{S6_SERIES_GRIND_II.pdf\#nameddest=tag-01C9}{II.8.1 (p.~121)} }
\expandafter\def\csname SRef@cdp:st:C.2.NC.3\endcsname{\href{S6_SERIES_GRIND_II.pdf\#nameddest=tag-01CA}{II.8.2 (p.~122)} }
\expandafter\def\csname SRef@cdp:st:C.2.NC.4\endcsname{\href{S6_SERIES_GRIND_II.pdf\#nameddest=tag-01CB}{II.8.3 (p.~123)} }
\expandafter\def\csname SRef@cdp:st:C.2.NC.5\endcsname{\href{S6_SERIES_GRIND_II.pdf\#nameddest=tag-01CC}{II.8.4 (p.~124)} }
\expandafter\def\csname SRef@cdp:st:C2d\endcsname{\href{S6_SERIES_GRIND_II.pdf\#nameddest=tag-01FJ}{II.8.5 (p.~125)} }
\expandafter\def\csname SRef@night:P4.1\endcsname{\href{S6_SERIES_GRIND_II.pdf\#nameddest=tag-02WG}{II.9.1 (p.~127)} }
\expandafter\def\csname SRef@stmt:P4.1\endcsname{\href{S6_SERIES_GRIND_II.pdf\#nameddest=tag-02WG}{II.9.1 (p.~127)} }
\expandafter\def\csname SRef@night:P4.10\endcsname{\href{S6_SERIES_GRIND_II.pdf\#nameddest=tag-02WH}{II.9.10 (p.~138)} }
\expandafter\def\csname SRef@stmt:P4.10\endcsname{\href{S6_SERIES_GRIND_II.pdf\#nameddest=tag-02WH}{II.9.10 (p.~138)} }
\expandafter\def\csname SRef@night:P4.11\endcsname{\href{S6_SERIES_GRIND_II.pdf\#nameddest=tag-02WI}{II.9.11 (p.~139)} }
\expandafter\def\csname SRef@stmt:P4.11\endcsname{\href{S6_SERIES_GRIND_II.pdf\#nameddest=tag-02WI}{II.9.11 (p.~139)} }
\expandafter\def\csname SRef@night:P4.12\endcsname{\href{S6_SERIES_GRIND_II.pdf\#nameddest=tag-02WJ}{II.9.12 (p.~140)} }
\expandafter\def\csname SRef@stmt:P4.12\endcsname{\href{S6_SERIES_GRIND_II.pdf\#nameddest=tag-02WJ}{II.9.12 (p.~140)} }
\expandafter\def\csname SRef@night:P4.13\endcsname{\href{S6_SERIES_GRIND_II.pdf\#nameddest=tag-02WK}{II.9.13 (p.~142)} }
\expandafter\def\csname SRef@stmt:P4.13\endcsname{\href{S6_SERIES_GRIND_II.pdf\#nameddest=tag-02WK}{II.9.13 (p.~142)} }
\expandafter\def\csname SRef@night:P4.14\endcsname{\href{S6_SERIES_GRIND_II.pdf\#nameddest=tag-02WL}{II.9.14 (p.~143)} }
\expandafter\def\csname SRef@stmt:P4.14\endcsname{\href{S6_SERIES_GRIND_II.pdf\#nameddest=tag-02WL}{II.9.14 (p.~143)} }
\expandafter\def\csname SRef@night:P4.15\endcsname{\href{S6_SERIES_GRIND_II.pdf\#nameddest=tag-02WM}{II.9.15 (p.~144)} }
\expandafter\def\csname SRef@stmt:P4.15\endcsname{\href{S6_SERIES_GRIND_II.pdf\#nameddest=tag-02WM}{II.9.15 (p.~144)} }
\expandafter\def\csname SRef@night:P4.16\endcsname{\href{S6_SERIES_GRIND_II.pdf\#nameddest=tag-02WN}{II.9.16 (p.~145)} }
\expandafter\def\csname SRef@stmt:P4.16\endcsname{\href{S6_SERIES_GRIND_II.pdf\#nameddest=tag-02WN}{II.9.16 (p.~145)} }
\expandafter\def\csname SRef@night:P4.17\endcsname{\href{S6_SERIES_GRIND_II.pdf\#nameddest=tag-02WO}{II.9.17 (p.~147)} }
\expandafter\def\csname SRef@stmt:P4.17\endcsname{\href{S6_SERIES_GRIND_II.pdf\#nameddest=tag-02WO}{II.9.17 (p.~147)} }
\expandafter\def\csname SRef@night:P4.18\endcsname{\href{S6_SERIES_GRIND_II.pdf\#nameddest=tag-02WP}{II.9.18 (p.~148)} }
\expandafter\def\csname SRef@stmt:P4.18\endcsname{\href{S6_SERIES_GRIND_II.pdf\#nameddest=tag-02WP}{II.9.18 (p.~148)} }
\expandafter\def\csname SRef@night:P4.19\endcsname{\href{S6_SERIES_GRIND_II.pdf\#nameddest=tag-02WQ}{II.9.19 (p.~149)} }
\expandafter\def\csname SRef@stmt:P4.19\endcsname{\href{S6_SERIES_GRIND_II.pdf\#nameddest=tag-02WQ}{II.9.19 (p.~149)} }
\expandafter\def\csname SRef@night:P4.2\endcsname{\href{S6_SERIES_GRIND_II.pdf\#nameddest=tag-02WR}{II.9.2 (p.~128)} }
\expandafter\def\csname SRef@stmt:P4.2\endcsname{\href{S6_SERIES_GRIND_II.pdf\#nameddest=tag-02WR}{II.9.2 (p.~128)} }
\expandafter\def\csname SRef@night:P4.20\endcsname{\href{S6_SERIES_GRIND_II.pdf\#nameddest=tag-02WS}{II.9.20 (p.~149)} }
\expandafter\def\csname SRef@stmt:P4.20\endcsname{\href{S6_SERIES_GRIND_II.pdf\#nameddest=tag-02WS}{II.9.20 (p.~149)} }
\expandafter\def\csname SRef@night:P4.21\endcsname{\href{S6_SERIES_GRIND_II.pdf\#nameddest=tag-02WT}{II.9.21 (p.~151)} }
\expandafter\def\csname SRef@stmt:P4.21\endcsname{\href{S6_SERIES_GRIND_II.pdf\#nameddest=tag-02WT}{II.9.21 (p.~151)} }
\expandafter\def\csname SRef@night:P4.22\endcsname{\href{S6_SERIES_GRIND_II.pdf\#nameddest=tag-02WU}{II.9.22 (p.~153)} }
\expandafter\def\csname SRef@stmt:P4.22\endcsname{\href{S6_SERIES_GRIND_II.pdf\#nameddest=tag-02WU}{II.9.22 (p.~153)} }
\expandafter\def\csname SRef@night:P4.23\endcsname{\href{S6_SERIES_GRIND_II.pdf\#nameddest=tag-02WV}{II.9.23 (p.~154)} }
\expandafter\def\csname SRef@stmt:P4.23\endcsname{\href{S6_SERIES_GRIND_II.pdf\#nameddest=tag-02WV}{II.9.23 (p.~154)} }
\expandafter\def\csname SRef@night:P4.24\endcsname{\href{S6_SERIES_GRIND_II.pdf\#nameddest=tag-02WW}{II.9.24 (p.~155)} }
\expandafter\def\csname SRef@stmt:P4.24\endcsname{\href{S6_SERIES_GRIND_II.pdf\#nameddest=tag-02WW}{II.9.24 (p.~155)} }
\expandafter\def\csname SRef@night:P4.25\endcsname{\href{S6_SERIES_GRIND_II.pdf\#nameddest=tag-02WX}{II.9.25 (p.~158)} }
\expandafter\def\csname SRef@stmt:P4.25\endcsname{\href{S6_SERIES_GRIND_II.pdf\#nameddest=tag-02WX}{II.9.25 (p.~158)} }
\expandafter\def\csname SRef@night:P4.26\endcsname{\href{S6_SERIES_GRIND_II.pdf\#nameddest=tag-02WY}{II.9.26 (p.~160)} }
\expandafter\def\csname SRef@stmt:P4.26\endcsname{\href{S6_SERIES_GRIND_II.pdf\#nameddest=tag-02WY}{II.9.26 (p.~160)} }
\expandafter\def\csname SRef@night:P4.27\endcsname{\href{S6_SERIES_GRIND_II.pdf\#nameddest=tag-02WZ}{II.9.27 (p.~162)} }
\expandafter\def\csname SRef@stmt:P4.27\endcsname{\href{S6_SERIES_GRIND_II.pdf\#nameddest=tag-02WZ}{II.9.27 (p.~162)} }
\expandafter\def\csname SRef@night:P4.28\endcsname{\href{S6_SERIES_GRIND_II.pdf\#nameddest=tag-02X0}{II.9.28 (p.~164)} }
\expandafter\def\csname SRef@stmt:P4.28\endcsname{\href{S6_SERIES_GRIND_II.pdf\#nameddest=tag-02X0}{II.9.28 (p.~164)} }
\expandafter\def\csname SRef@night:P4.29\endcsname{\href{S6_SERIES_GRIND_II.pdf\#nameddest=tag-02X1}{II.9.29 (p.~166)} }
\expandafter\def\csname SRef@stmt:P4.29\endcsname{\href{S6_SERIES_GRIND_II.pdf\#nameddest=tag-02X1}{II.9.29 (p.~166)} }
\expandafter\def\csname SRef@night:P4.3\endcsname{\href{S6_SERIES_GRIND_II.pdf\#nameddest=tag-02X2}{II.9.3 (p.~129)} }
\expandafter\def\csname SRef@stmt:P4.3\endcsname{\href{S6_SERIES_GRIND_II.pdf\#nameddest=tag-02X2}{II.9.3 (p.~129)} }
\expandafter\def\csname SRef@night:P4.30\endcsname{\href{S6_SERIES_GRIND_II.pdf\#nameddest=tag-02X3}{II.9.30 (p.~167)} }
\expandafter\def\csname SRef@stmt:P4.30\endcsname{\href{S6_SERIES_GRIND_II.pdf\#nameddest=tag-02X3}{II.9.30 (p.~167)} }
\expandafter\def\csname SRef@night:P4.31\endcsname{\href{S6_SERIES_GRIND_II.pdf\#nameddest=tag-02X4}{II.9.31 (p.~168)} }
\expandafter\def\csname SRef@stmt:P4.31\endcsname{\href{S6_SERIES_GRIND_II.pdf\#nameddest=tag-02X4}{II.9.31 (p.~168)} }
\expandafter\def\csname SRef@night:P4.32\endcsname{\href{S6_SERIES_GRIND_II.pdf\#nameddest=tag-02X5}{II.9.32 (p.~170)} }
\expandafter\def\csname SRef@stmt:P4.32\endcsname{\href{S6_SERIES_GRIND_II.pdf\#nameddest=tag-02X5}{II.9.32 (p.~170)} }
\expandafter\def\csname SRef@night:P4.33\endcsname{\href{S6_SERIES_GRIND_II.pdf\#nameddest=tag-02X6}{II.9.33 (p.~171)} }
\expandafter\def\csname SRef@stmt:P4.33\endcsname{\href{S6_SERIES_GRIND_II.pdf\#nameddest=tag-02X6}{II.9.33 (p.~171)} }
\expandafter\def\csname SRef@night:P4.34\endcsname{\href{S6_SERIES_GRIND_II.pdf\#nameddest=tag-02X7}{II.9.34 (p.~172)} }
\expandafter\def\csname SRef@stmt:P4.34\endcsname{\href{S6_SERIES_GRIND_II.pdf\#nameddest=tag-02X7}{II.9.34 (p.~172)} }
\expandafter\def\csname SRef@night:P4.35\endcsname{\href{S6_SERIES_GRIND_II.pdf\#nameddest=tag-02X8}{II.9.35 (p.~174)} }
\expandafter\def\csname SRef@stmt:P4.35\endcsname{\href{S6_SERIES_GRIND_II.pdf\#nameddest=tag-02X8}{II.9.35 (p.~174)} }
\expandafter\def\csname SRef@night:P4.36\endcsname{\href{S6_SERIES_GRIND_II.pdf\#nameddest=tag-02X9}{II.9.36 (p.~176)} }
\expandafter\def\csname SRef@stmt:P4.36\endcsname{\href{S6_SERIES_GRIND_II.pdf\#nameddest=tag-02X9}{II.9.36 (p.~176)} }
\expandafter\def\csname SRef@night:P4.37\endcsname{\href{S6_SERIES_GRIND_II.pdf\#nameddest=tag-02XA}{II.9.37 (p.~178)} }
\expandafter\def\csname SRef@stmt:P4.37\endcsname{\href{S6_SERIES_GRIND_II.pdf\#nameddest=tag-02XA}{II.9.37 (p.~178)} }
\expandafter\def\csname SRef@night:P4.38\endcsname{\href{S6_SERIES_GRIND_II.pdf\#nameddest=tag-02XB}{II.9.38 (p.~179)} }
\expandafter\def\csname SRef@stmt:P4.38\endcsname{\href{S6_SERIES_GRIND_II.pdf\#nameddest=tag-02XB}{II.9.38 (p.~179)} }
\expandafter\def\csname SRef@night:P4.39\endcsname{\href{S6_SERIES_GRIND_II.pdf\#nameddest=tag-02XC}{II.9.39 (p.~181)} }
\expandafter\def\csname SRef@stmt:P4.39\endcsname{\href{S6_SERIES_GRIND_II.pdf\#nameddest=tag-02XC}{II.9.39 (p.~181)} }
\expandafter\def\csname SRef@night:P4.4\endcsname{\href{S6_SERIES_GRIND_II.pdf\#nameddest=tag-02XD}{II.9.4 (p.~130)} }
\expandafter\def\csname SRef@stmt:P4.4\endcsname{\href{S6_SERIES_GRIND_II.pdf\#nameddest=tag-02XD}{II.9.4 (p.~130)} }
\expandafter\def\csname SRef@night:P4.40\endcsname{\href{S6_SERIES_GRIND_II.pdf\#nameddest=tag-02XE}{II.9.40 (p.~183)} }
\expandafter\def\csname SRef@stmt:P4.40\endcsname{\href{S6_SERIES_GRIND_II.pdf\#nameddest=tag-02XE}{II.9.40 (p.~183)} }
\expandafter\def\csname SRef@night:P4.41\endcsname{\href{S6_SERIES_GRIND_II.pdf\#nameddest=tag-02XF}{II.9.41 (p.~184)} }
\expandafter\def\csname SRef@stmt:P4.41\endcsname{\href{S6_SERIES_GRIND_II.pdf\#nameddest=tag-02XF}{II.9.41 (p.~184)} }
\expandafter\def\csname SRef@night:P4.42\endcsname{\href{S6_SERIES_GRIND_II.pdf\#nameddest=tag-02XG}{II.9.42 (p.~187)} }
\expandafter\def\csname SRef@stmt:P4.42\endcsname{\href{S6_SERIES_GRIND_II.pdf\#nameddest=tag-02XG}{II.9.42 (p.~187)} }
\expandafter\def\csname SRef@night:P4.43\endcsname{\href{S6_SERIES_GRIND_II.pdf\#nameddest=tag-02XH}{II.9.43 (p.~188)} }
\expandafter\def\csname SRef@stmt:P4.43\endcsname{\href{S6_SERIES_GRIND_II.pdf\#nameddest=tag-02XH}{II.9.43 (p.~188)} }
\expandafter\def\csname SRef@night:P4.44\endcsname{\href{S6_SERIES_GRIND_II.pdf\#nameddest=tag-02XI}{II.9.44 (p.~189)} }
\expandafter\def\csname SRef@stmt:P4.44\endcsname{\href{S6_SERIES_GRIND_II.pdf\#nameddest=tag-02XI}{II.9.44 (p.~189)} }
\expandafter\def\csname SRef@night:P4.45\endcsname{\href{S6_SERIES_GRIND_II.pdf\#nameddest=tag-02XJ}{II.9.45 (p.~190)} }
\expandafter\def\csname SRef@stmt:P4.45\endcsname{\href{S6_SERIES_GRIND_II.pdf\#nameddest=tag-02XJ}{II.9.45 (p.~190)} }
\expandafter\def\csname SRef@night:P4.46\endcsname{\href{S6_SERIES_GRIND_II.pdf\#nameddest=tag-02XK}{II.9.46 (p.~191)} }
\expandafter\def\csname SRef@stmt:P4.46\endcsname{\href{S6_SERIES_GRIND_II.pdf\#nameddest=tag-02XK}{II.9.46 (p.~191)} }
\expandafter\def\csname SRef@night:P4.47\endcsname{\href{S6_SERIES_GRIND_II.pdf\#nameddest=tag-02XL}{II.9.47 (p.~192)} }
\expandafter\def\csname SRef@stmt:P4.47\endcsname{\href{S6_SERIES_GRIND_II.pdf\#nameddest=tag-02XL}{II.9.47 (p.~192)} }
\expandafter\def\csname SRef@night:P4.48\endcsname{\href{S6_SERIES_GRIND_II.pdf\#nameddest=tag-02XM}{II.9.48 (p.~193)} }
\expandafter\def\csname SRef@stmt:P4.48\endcsname{\href{S6_SERIES_GRIND_II.pdf\#nameddest=tag-02XM}{II.9.48 (p.~193)} }
\expandafter\def\csname SRef@night:P4.49\endcsname{\href{S6_SERIES_GRIND_II.pdf\#nameddest=tag-02XN}{II.9.49 (p.~194)} }
\expandafter\def\csname SRef@stmt:P4.49\endcsname{\href{S6_SERIES_GRIND_II.pdf\#nameddest=tag-02XN}{II.9.49 (p.~194)} }
\expandafter\def\csname SRef@night:P4.5\endcsname{\href{S6_SERIES_GRIND_II.pdf\#nameddest=tag-02XO}{II.9.5 (p.~134)} }
\expandafter\def\csname SRef@stmt:P4.5\endcsname{\href{S6_SERIES_GRIND_II.pdf\#nameddest=tag-02XO}{II.9.5 (p.~134)} }
\expandafter\def\csname SRef@night:P4.50\endcsname{\href{S6_SERIES_GRIND_II.pdf\#nameddest=tag-02XP}{II.9.50 (p.~198)} }
\expandafter\def\csname SRef@stmt:P4.50\endcsname{\href{S6_SERIES_GRIND_II.pdf\#nameddest=tag-02XP}{II.9.50 (p.~198)} }
\expandafter\def\csname SRef@night:P4.51\endcsname{\href{S6_SERIES_GRIND_II.pdf\#nameddest=tag-02XQ}{II.9.51 (p.~198)} }
\expandafter\def\csname SRef@stmt:P4.51\endcsname{\href{S6_SERIES_GRIND_II.pdf\#nameddest=tag-02XQ}{II.9.51 (p.~198)} }
\expandafter\def\csname SRef@night:P4.52\endcsname{\href{S6_SERIES_GRIND_II.pdf\#nameddest=tag-02XR}{II.9.52 (p.~199)} }
\expandafter\def\csname SRef@stmt:P4.52\endcsname{\href{S6_SERIES_GRIND_II.pdf\#nameddest=tag-02XR}{II.9.52 (p.~199)} }
\expandafter\def\csname SRef@night:P4.53\endcsname{\href{S6_SERIES_GRIND_II.pdf\#nameddest=tag-02XS}{II.9.53 (p.~201)} }
\expandafter\def\csname SRef@stmt:P4.53\endcsname{\href{S6_SERIES_GRIND_II.pdf\#nameddest=tag-02XS}{II.9.53 (p.~201)} }
\expandafter\def\csname SRef@night:P4.54\endcsname{\href{S6_SERIES_GRIND_II.pdf\#nameddest=tag-02XT}{II.9.54 (p.~204)} }
\expandafter\def\csname SRef@stmt:P4.54\endcsname{\href{S6_SERIES_GRIND_II.pdf\#nameddest=tag-02XT}{II.9.54 (p.~204)} }
\expandafter\def\csname SRef@night:P4.55\endcsname{\href{S6_SERIES_GRIND_II.pdf\#nameddest=tag-02XU}{II.9.55 (p.~210)} }
\expandafter\def\csname SRef@stmt:P4.55\endcsname{\href{S6_SERIES_GRIND_II.pdf\#nameddest=tag-02XU}{II.9.55 (p.~210)} }
\expandafter\def\csname SRef@night:P4.56\endcsname{\href{S6_SERIES_GRIND_II.pdf\#nameddest=tag-02XV}{II.9.56 (p.~213)} }
\expandafter\def\csname SRef@stmt:P4.56\endcsname{\href{S6_SERIES_GRIND_II.pdf\#nameddest=tag-02XV}{II.9.56 (p.~213)} }
\expandafter\def\csname SRef@night:P4.57\endcsname{\href{S6_SERIES_GRIND_II.pdf\#nameddest=tag-02XW}{II.9.57 (p.~216)} }
\expandafter\def\csname SRef@stmt:P4.57\endcsname{\href{S6_SERIES_GRIND_II.pdf\#nameddest=tag-02XW}{II.9.57 (p.~216)} }
\expandafter\def\csname SRef@night:P4.58\endcsname{\href{S6_SERIES_GRIND_II.pdf\#nameddest=tag-02XX}{II.9.58 (p.~218)} }
\expandafter\def\csname SRef@stmt:P4.58\endcsname{\href{S6_SERIES_GRIND_II.pdf\#nameddest=tag-02XX}{II.9.58 (p.~218)} }
\expandafter\def\csname SRef@night:P4R.1\endcsname{\href{S6_SERIES_GRIND_II.pdf\#nameddest=tag-02Y2}{II.9.59 (p.~219)} }
\expandafter\def\csname SRef@stmt:P4R.1\endcsname{\href{S6_SERIES_GRIND_II.pdf\#nameddest=tag-02Y2}{II.9.59 (p.~219)} }
\expandafter\def\csname SRef@night:P4.6\endcsname{\href{S6_SERIES_GRIND_II.pdf\#nameddest=tag-02XY}{II.9.6 (p.~135)} }
\expandafter\def\csname SRef@stmt:P4.6\endcsname{\href{S6_SERIES_GRIND_II.pdf\#nameddest=tag-02XY}{II.9.6 (p.~135)} }
\expandafter\def\csname SRef@night:P4R.2\endcsname{\href{S6_SERIES_GRIND_II.pdf\#nameddest=tag-02YD}{II.9.60 (p.~220)} }
\expandafter\def\csname SRef@stmt:P4R.2\endcsname{\href{S6_SERIES_GRIND_II.pdf\#nameddest=tag-02YD}{II.9.60 (p.~220)} }
\expandafter\def\csname SRef@night:P4R.3\endcsname{\href{S6_SERIES_GRIND_II.pdf\#nameddest=tag-02YF}{II.9.61 (p.~221)} }
\expandafter\def\csname SRef@stmt:P4R.3\endcsname{\href{S6_SERIES_GRIND_II.pdf\#nameddest=tag-02YF}{II.9.61 (p.~221)} }
\expandafter\def\csname SRef@night:P4R.4\endcsname{\href{S6_SERIES_GRIND_II.pdf\#nameddest=tag-02YG}{II.9.62 (p.~223)} }
\expandafter\def\csname SRef@stmt:P4R.4\endcsname{\href{S6_SERIES_GRIND_II.pdf\#nameddest=tag-02YG}{II.9.62 (p.~223)} }
\expandafter\def\csname SRef@night:P4R.5\endcsname{\href{S6_SERIES_GRIND_II.pdf\#nameddest=tag-02YH}{II.9.63 (p.~224)} }
\expandafter\def\csname SRef@stmt:P4R.5\endcsname{\href{S6_SERIES_GRIND_II.pdf\#nameddest=tag-02YH}{II.9.63 (p.~224)} }
\expandafter\def\csname SRef@night:P4R.6\endcsname{\href{S6_SERIES_GRIND_II.pdf\#nameddest=tag-02YI}{II.9.64 (p.~225)} }
\expandafter\def\csname SRef@stmt:P4R.6\endcsname{\href{S6_SERIES_GRIND_II.pdf\#nameddest=tag-02YI}{II.9.64 (p.~225)} }
\expandafter\def\csname SRef@night:P4R.7\endcsname{\href{S6_SERIES_GRIND_II.pdf\#nameddest=tag-02YJ}{II.9.65 (p.~227)} }
\expandafter\def\csname SRef@stmt:P4R.7\endcsname{\href{S6_SERIES_GRIND_II.pdf\#nameddest=tag-02YJ}{II.9.65 (p.~227)} }
\expandafter\def\csname SRef@night:P4R.8\endcsname{\href{S6_SERIES_GRIND_II.pdf\#nameddest=tag-02YK}{II.9.66 (p.~227)} }
\expandafter\def\csname SRef@stmt:P4R.8\endcsname{\href{S6_SERIES_GRIND_II.pdf\#nameddest=tag-02YK}{II.9.66 (p.~227)} }
\expandafter\def\csname SRef@night:P4R.9\endcsname{\href{S6_SERIES_GRIND_II.pdf\#nameddest=tag-02YL}{II.9.67 (p.~228)} }
\expandafter\def\csname SRef@stmt:P4R.9\endcsname{\href{S6_SERIES_GRIND_II.pdf\#nameddest=tag-02YL}{II.9.67 (p.~228)} }
\expandafter\def\csname SRef@night:P4R.10\endcsname{\href{S6_SERIES_GRIND_II.pdf\#nameddest=tag-02Y3}{II.9.68 (p.~229)} }
\expandafter\def\csname SRef@stmt:P4R.10\endcsname{\href{S6_SERIES_GRIND_II.pdf\#nameddest=tag-02Y3}{II.9.68 (p.~229)} }
\expandafter\def\csname SRef@night:P4R.11\endcsname{\href{S6_SERIES_GRIND_II.pdf\#nameddest=tag-02Y4}{II.9.69 (p.~231)} }
\expandafter\def\csname SRef@stmt:P4R.11\endcsname{\href{S6_SERIES_GRIND_II.pdf\#nameddest=tag-02Y4}{II.9.69 (p.~231)} }
\expandafter\def\csname SRef@night:P4.7\endcsname{\href{S6_SERIES_GRIND_II.pdf\#nameddest=tag-02XZ}{II.9.7 (p.~136)} }
\expandafter\def\csname SRef@stmt:P4.7\endcsname{\href{S6_SERIES_GRIND_II.pdf\#nameddest=tag-02XZ}{II.9.7 (p.~136)} }
\expandafter\def\csname SRef@night:P4R.12\endcsname{\href{S6_SERIES_GRIND_II.pdf\#nameddest=tag-02Y5}{II.9.70 (p.~232)} }
\expandafter\def\csname SRef@stmt:P4R.12\endcsname{\href{S6_SERIES_GRIND_II.pdf\#nameddest=tag-02Y5}{II.9.70 (p.~232)} }
\expandafter\def\csname SRef@night:P4R.13\endcsname{\href{S6_SERIES_GRIND_II.pdf\#nameddest=tag-02Y6}{II.9.71 (p.~233)} }
\expandafter\def\csname SRef@stmt:P4R.13\endcsname{\href{S6_SERIES_GRIND_II.pdf\#nameddest=tag-02Y6}{II.9.71 (p.~233)} }
\expandafter\def\csname SRef@night:P4R.14\endcsname{\href{S6_SERIES_GRIND_II.pdf\#nameddest=tag-02Y7}{II.9.72 (p.~237)} }
\expandafter\def\csname SRef@stmt:P4R.14\endcsname{\href{S6_SERIES_GRIND_II.pdf\#nameddest=tag-02Y7}{II.9.72 (p.~237)} }
\expandafter\def\csname SRef@night:P4R.15\endcsname{\href{S6_SERIES_GRIND_II.pdf\#nameddest=tag-02Y8}{II.9.73 (p.~239)} }
\expandafter\def\csname SRef@stmt:P4R.15\endcsname{\href{S6_SERIES_GRIND_II.pdf\#nameddest=tag-02Y8}{II.9.73 (p.~239)} }
\expandafter\def\csname SRef@night:P4R.16\endcsname{\href{S6_SERIES_GRIND_II.pdf\#nameddest=tag-02Y9}{II.9.74 (p.~240)} }
\expandafter\def\csname SRef@stmt:P4R.16\endcsname{\href{S6_SERIES_GRIND_II.pdf\#nameddest=tag-02Y9}{II.9.74 (p.~240)} }
\expandafter\def\csname SRef@night:P4R.17\endcsname{\href{S6_SERIES_GRIND_II.pdf\#nameddest=tag-02YA}{II.9.75 (p.~242)} }
\expandafter\def\csname SRef@stmt:P4R.17\endcsname{\href{S6_SERIES_GRIND_II.pdf\#nameddest=tag-02YA}{II.9.75 (p.~242)} }
\expandafter\def\csname SRef@bounds-without-an-ordinary-square-assumption\endcsname{\href{S6_SERIES_GRIND_II.pdf\#nameddest=tag-02YA}{II.{{$\langle 1\rangle $1.}} (p.~242)} }
\expandafter\def\csname SRef@night:P4R.18\endcsname{\href{S6_SERIES_GRIND_II.pdf\#nameddest=tag-02YB}{II.9.76 (p.~246)} }
\expandafter\def\csname SRef@stmt:P4R.18\endcsname{\href{S6_SERIES_GRIND_II.pdf\#nameddest=tag-02YB}{II.9.76 (p.~246)} }
\expandafter\def\csname SRef@night:P4R.19\endcsname{\href{S6_SERIES_GRIND_II.pdf\#nameddest=tag-02YC}{II.9.77 (p.~249)} }
\expandafter\def\csname SRef@stmt:P4R.19\endcsname{\href{S6_SERIES_GRIND_II.pdf\#nameddest=tag-02YC}{II.9.77 (p.~249)} }
\expandafter\def\csname SRef@night:P4R.20\endcsname{\href{S6_SERIES_GRIND_II.pdf\#nameddest=tag-02YE}{II.9.78 (p.~252)} }
\expandafter\def\csname SRef@stmt:P4R.20\endcsname{\href{S6_SERIES_GRIND_II.pdf\#nameddest=tag-02YE}{II.9.78 (p.~252)} }
\expandafter\def\csname SRef@night:P4.8\endcsname{\href{S6_SERIES_GRIND_II.pdf\#nameddest=tag-02Y0}{II.9.8 (p.~136)} }
\expandafter\def\csname SRef@stmt:P4.8\endcsname{\href{S6_SERIES_GRIND_II.pdf\#nameddest=tag-02Y0}{II.9.8 (p.~136)} }
\expandafter\def\csname SRef@night:P4.9\endcsname{\href{S6_SERIES_GRIND_II.pdf\#nameddest=tag-02Y1}{II.9.9 (p.~138)} }
\expandafter\def\csname SRef@stmt:P4.9\endcsname{\href{S6_SERIES_GRIND_II.pdf\#nameddest=tag-02Y1}{II.9.9 (p.~138)} }
\expandafter\def\csname SRef@cdp:st:C4a\endcsname{\href{S6_SERIES_GRIND_II.pdf\#nameddest=tag-01FK}{II.10.1 (p.~255)} }
\expandafter\def\csname SRef@cdp:st:C5c\endcsname{\href{S6_SERIES_GRIND_II.pdf\#nameddest=tag-01FT}{II.10.10 (p.~269)} }
\expandafter\def\csname SRef@cdp:st:C4b\endcsname{\href{S6_SERIES_GRIND_II.pdf\#nameddest=tag-01FL}{II.10.2 (p.~257)} }
\expandafter\def\csname SRef@cdp:st:C4c\endcsname{\href{S6_SERIES_GRIND_II.pdf\#nameddest=tag-01FM}{II.10.3 (p.~258)} }
\expandafter\def\csname SRef@cdp:st:C4d\endcsname{\href{S6_SERIES_GRIND_II.pdf\#nameddest=tag-01FN}{II.10.4 (p.~259)} }
\expandafter\def\csname SRef@cdp:st:C4e\endcsname{\href{S6_SERIES_GRIND_II.pdf\#nameddest=tag-01FO}{II.10.5 (p.~260)} }
\expandafter\def\csname SRef@cdp:st:C4f\endcsname{\href{S6_SERIES_GRIND_II.pdf\#nameddest=tag-01FP}{II.10.6 (p.~263)} }
\expandafter\def\csname SRef@cdp:st:C4g\endcsname{\href{S6_SERIES_GRIND_II.pdf\#nameddest=tag-01FQ}{II.10.7 (p.~264)} }
\expandafter\def\csname SRef@cdp:st:C5a\endcsname{\href{S6_SERIES_GRIND_II.pdf\#nameddest=tag-01FR}{II.10.8 (p.~266)} }
\expandafter\def\csname SRef@cdp:st:C5b\endcsname{\href{S6_SERIES_GRIND_II.pdf\#nameddest=tag-01FS}{II.10.9 (p.~267)} }
\expandafter\def\csname SRef@night:P3K.36\endcsname{\href{S6_SERIES_GRIND_II.pdf\#nameddest=tag-02UH}{II.17.2 (p.~411)} }
\expandafter\def\csname SRef@stmt:P3K.36\endcsname{\href{S6_SERIES_GRIND_II.pdf\#nameddest=tag-02UH}{II.17.2 (p.~411)} }
\expandafter\def\csname SRef@inv:wb:uniform-nontorsion-provenance\endcsname{\href{S6_SERIES_GRIND_II.pdf\#nameddest=tag-01TQ}{II.15.2 (p.~363)} }
\expandafter\def\csname SRef@wbII:rem:scope\endcsname{\href{S6_SERIES_GRIND_II.pdf\#nameddest=tag-03IS}{II.16.3 (p.~371)} }
\expandafter\def\csname SRef@wbII:rem:local-cusp-picard-obstruction\endcsname{\href{S6_SERIES_GRIND_II.pdf\#nameddest=tag-03IQ}{II.16.8 (p.~387)} }
\expandafter\def\csname SRef@wbII:rem:relative-local-scope\endcsname{\href{S6_SERIES_GRIND_II.pdf\#nameddest=tag-03IR}{II.16.12 (p.~407)} }
\expandafter\def\csname SRef@sec:II:conventions\endcsname{\href{S6_SERIES_GRIND_II.pdf\#nameddest=tag-034E}{II.2 (p.~18)} }
\expandafter\def\csname SRef@sec:II:13\endcsname{\href{S6_SERIES_GRIND_II.pdf\#nameddest=tag-0349}{II.11 (p.~272)} }
\expandafter\def\csname SRef@sec:II:14\endcsname{\href{S6_SERIES_GRIND_II.pdf\#nameddest=tag-034A}{II.12 (p.~285)} }
\expandafter\def\csname SRef@sec:II:push-stability\endcsname{\href{S6_SERIES_GRIND_II.pdf\#nameddest=tag-034G}{II.13 (p.~309)} }
\expandafter\def\csname SRef@sec:II:combs\endcsname{\href{S6_SERIES_GRIND_II.pdf\#nameddest=tag-034D}{II.14 (p.~354)} }
\expandafter\def\csname SRef@sec:II:uniform-nontorsion\endcsname{\href{S6_SERIES_GRIND_II.pdf\#nameddest=tag-034H}{II.15 (p.~361)} }
\expandafter\def\csname SRef@app:II:programs\endcsname{\href{S6_SERIES_GRIND_II.pdf\#nameddest=tag-000B}{II.26 (p.~795)} }
\expandafter\def\csname SRef@sec:II:01\endcsname{\href{S6_SERIES_GRIND_II.pdf\#nameddest=tag-033X}{II.3 (p.~18)} }
\expandafter\def\csname SRef@sec:II:02\endcsname{\href{S6_SERIES_GRIND_II.pdf\#nameddest=tag-033Y}{II.4 (p.~20)} }
\expandafter\def\csname SRef@sec:II:03\endcsname{\href{S6_SERIES_GRIND_II.pdf\#nameddest=tag-033Z}{II.5 (p.~31)} }
\expandafter\def\csname SRef@sec:II:04\endcsname{\href{S6_SERIES_GRIND_II.pdf\#nameddest=tag-0340}{II.6 (p.~47)} }
\expandafter\def\csname SRef@sec:II:05\endcsname{\href{S6_SERIES_GRIND_II.pdf\#nameddest=tag-0341}{II.7 (p.~72)} }
\expandafter\def\csname SRef@sec:II:08\endcsname{\href{S6_SERIES_GRIND_II.pdf\#nameddest=tag-0344}{II.8 (p.~121)} }
\expandafter\def\csname SRef@sec:II:night-sign\endcsname{\href{S6_SERIES_GRIND_II.pdf\#nameddest=tag-034F}{II.9 (p.~127)} }
\expandafter\def\csname SRef@sec:II:12\endcsname{\href{S6_SERIES_GRIND_II.pdf\#nameddest=tag-0348}{II.10 (p.~255)} }
\expandafter\def\csname SRef@wbII:subsec:local-cusp-picard-obstruction\endcsname{\href{S6_SERIES_GRIND_II.pdf\#nameddest=tag-03IU}{II.16.8 (p.~382)} }
\expandafter\def\csname SRef@wbII:subsec:analytic-all-picard\endcsname{\href{S6_SERIES_GRIND_II.pdf\#nameddest=tag-03IT}{II.16.9 (p.~387)} }
\expandafter\def\csname SRef@push:src:E.2.23\endcsname{\href{S6_SERIES_GRIND_II.pdf\#nameddest=tag-0339}{II.13.20 (p.~347)} }
\expandafter\def\csname SRef@push:src:E.5.3\endcsname{\href{S6_SERIES_GRIND_II.pdf\#nameddest=tag-033P}{II.13.24 (p.~351)} }
\expandafter\def\csname SRef@wb:thm:uniform-nontorsion-audit\endcsname{\href{S6_SERIES_GRIND_II.pdf\#nameddest=tag-03DQ}{II.15.1 (p.~361)} }
\expandafter\def\csname SRef@wbII:prop:conductor\endcsname{\href{S6_SERIES_GRIND_II.pdf\#nameddest=tag-03IK}{II.16.2 (p.~366)} }
\expandafter\def\csname SRef@wbII:prop:conductor-differential\endcsname{\href{S6_SERIES_GRIND_II.pdf\#nameddest=tag-03IK}{II.16.2 (p.~366)} }
\expandafter\def\csname SRef@wbII:thm:analytic-all-picard\endcsname{\href{S6_SERIES_GRIND_II.pdf\#nameddest=tag-03IW}{II.16.21 (p.~398)} }
\expandafter\def\csname SRef@wbII:thm:candidate-conductor\endcsname{\href{S6_SERIES_GRIND_II.pdf\#nameddest=tag-03IX}{II.16.23 (p.~400)} }
\expandafter\def\csname SRef@wbII:thm:corrected-cdp-global\endcsname{\href{S6_SERIES_GRIND_II.pdf\#nameddest=tag-03J0}{II.16.27 (p.~402)} }
\expandafter\def\csname SRef@wbII:thm:corrected-global\endcsname{\href{S6_SERIES_GRIND_II.pdf\#nameddest=tag-03J0}{II.16.27 (p.~402)} }
\expandafter\def\csname SRef@wbII:thm:corrected-cdp-fibrewise\endcsname{\href{S6_SERIES_GRIND_II.pdf\#nameddest=tag-03IZ}{II.16.29 (p.~404)} }
\expandafter\def\csname SRef@wbII:thm:cartier-duality\endcsname{\href{S6_SERIES_GRIND_II.pdf\#nameddest=tag-03IY}{II.16.6 (p.~371)} }
\expandafter\def\csname SRef@ii:LR:covering\endcsname{\href{S6_SERIES_GRIND_II.pdf\#nameddest=tag-01RO}{II.17.1 (p.~407)} }
\expandafter\def\csname SRef@maps:eq:D.O9\endcsname{\href{S6_SERIES_GRIND_III.pdf\#nameddest=tag-020U}{III.{D.O9} (p.~85)} }
\expandafter\def\csname SRef@maps:eq:D.FI1\endcsname{\href{S6_SERIES_GRIND_III.pdf\#nameddest=tag-01YT}{III.{D.FI1} (p.~87)} }
\expandafter\def\csname SRef@maps:eq:D.FI2\endcsname{\href{S6_SERIES_GRIND_III.pdf\#nameddest=tag-01YU}{III.{D.FI2} (p.~88)} }
\expandafter\def\csname SRef@maps:eq:D.FI3\endcsname{\href{S6_SERIES_GRIND_III.pdf\#nameddest=tag-01YV}{III.{D.FI3} (p.~88)} }
\expandafter\def\csname SRef@maps:eq:D.FI4\endcsname{\href{S6_SERIES_GRIND_III.pdf\#nameddest=tag-01YW}{III.{D.FI4} (p.~90)} }
\expandafter\def\csname SRef@maps:eq:D.FI5\endcsname{\href{S6_SERIES_GRIND_III.pdf\#nameddest=tag-01YX}{III.{D.FI5} (p.~91)} }
\expandafter\def\csname SRef@maps:eq:D.FG1\endcsname{\href{S6_SERIES_GRIND_III.pdf\#nameddest=tag-01YQ}{III.{D.FG1} (p.~92)} }
\expandafter\def\csname SRef@maps:eq:D.FG2\endcsname{\href{S6_SERIES_GRIND_III.pdf\#nameddest=tag-01YR}{III.{D.FG2} (p.~92)} }
\expandafter\def\csname SRef@maps:eq:D.FG3\endcsname{\href{S6_SERIES_GRIND_III.pdf\#nameddest=tag-01YS}{III.{D.FG3} (p.~92)} }
\expandafter\def\csname SRef@maps:eq:D.IS1\endcsname{\href{S6_SERIES_GRIND_III.pdf\#nameddest=tag-0207}{III.{D.IS1} (p.~93)} }
\expandafter\def\csname SRef@maps:eq:D.IS2\endcsname{\href{S6_SERIES_GRIND_III.pdf\#nameddest=tag-0208}{III.{D.IS2} (p.~93)} }
\expandafter\def\csname SRef@maps:eq:D.SC1\endcsname{\href{S6_SERIES_GRIND_III.pdf\#nameddest=tag-023Z}{III.{D.SC1} (p.~94)} }
\expandafter\def\csname SRef@maps:eq:D.SC2\endcsname{\href{S6_SERIES_GRIND_III.pdf\#nameddest=tag-0240}{III.{D.SC2} (p.~95)} }
\expandafter\def\csname SRef@maps:eq:D.SC3\endcsname{\href{S6_SERIES_GRIND_III.pdf\#nameddest=tag-0241}{III.{D.SC3} (p.~95)} }
\expandafter\def\csname SRef@maps:eq:D.SC4\endcsname{\href{S6_SERIES_GRIND_III.pdf\#nameddest=tag-0242}{III.{D.SC4} (p.~95)} }
\expandafter\def\csname SRef@maps:eq:D.SC5\endcsname{\href{S6_SERIES_GRIND_III.pdf\#nameddest=tag-0243}{III.{D.SC5} (p.~95)} }
\expandafter\def\csname SRef@maps:eq:D.SC6\endcsname{\href{S6_SERIES_GRIND_III.pdf\#nameddest=tag-0244}{III.{D.SC6} (p.~95)} }
\expandafter\def\csname SRef@maps:eq:D.SC7\endcsname{\href{S6_SERIES_GRIND_III.pdf\#nameddest=tag-0245}{III.{D.SC7} (p.~97)} }
\expandafter\def\csname SRef@maps:eq:D.SC8\endcsname{\href{S6_SERIES_GRIND_III.pdf\#nameddest=tag-0246}{III.{D.SC8} (p.~97)} }
\expandafter\def\csname SRef@maps:eq:D.HY1\endcsname{\href{S6_SERIES_GRIND_III.pdf\#nameddest=tag-01ZN}{III.{D.HY1} (p.~98)} }
\expandafter\def\csname SRef@maps:eq:D.HY2\endcsname{\href{S6_SERIES_GRIND_III.pdf\#nameddest=tag-01ZO}{III.{D.HY2} (p.~98)} }
\expandafter\def\csname SRef@maps:eq:D.HY3\endcsname{\href{S6_SERIES_GRIND_III.pdf\#nameddest=tag-01ZP}{III.{D.HY3} (p.~99)} }
\expandafter\def\csname SRef@maps:eq:D.HY4\endcsname{\href{S6_SERIES_GRIND_III.pdf\#nameddest=tag-01ZQ}{III.{D.HY4} (p.~99)} }
\expandafter\def\csname SRef@maps:eq:D.HY5\endcsname{\href{S6_SERIES_GRIND_III.pdf\#nameddest=tag-01ZR}{III.{D.HY5} (p.~99)} }
\expandafter\def\csname SRef@maps:eq:D.HY6\endcsname{\href{S6_SERIES_GRIND_III.pdf\#nameddest=tag-01ZS}{III.{D.HY6} (p.~100)} }
\expandafter\def\csname SRef@maps:eq:D.HY7\endcsname{\href{S6_SERIES_GRIND_III.pdf\#nameddest=tag-01ZT}{III.{D.HY7} (p.~100)} }
\expandafter\def\csname SRef@maps:eq:D.ST1\endcsname{\href{S6_SERIES_GRIND_III.pdf\#nameddest=tag-024V}{III.{D.ST1} (p.~102)} }
\expandafter\def\csname SRef@maps:eq:D.ST2\endcsname{\href{S6_SERIES_GRIND_III.pdf\#nameddest=tag-0250}{III.{D.ST2} (p.~103)} }
\expandafter\def\csname SRef@maps:eq:D.ST3\endcsname{\href{S6_SERIES_GRIND_III.pdf\#nameddest=tag-0251}{III.{D.ST3} (p.~103)} }
\expandafter\def\csname SRef@maps:eq:D.ST4\endcsname{\href{S6_SERIES_GRIND_III.pdf\#nameddest=tag-0252}{III.{D.ST4} (p.~103)} }
\expandafter\def\csname SRef@maps:eq:D.ST5\endcsname{\href{S6_SERIES_GRIND_III.pdf\#nameddest=tag-0253}{III.{D.ST5} (p.~103)} }
\expandafter\def\csname SRef@maps:eq:D.ST6\endcsname{\href{S6_SERIES_GRIND_III.pdf\#nameddest=tag-0254}{III.{D.ST6} (p.~103)} }
\expandafter\def\csname SRef@maps:eq:D.ST7\endcsname{\href{S6_SERIES_GRIND_III.pdf\#nameddest=tag-0255}{III.{D.ST7} (p.~104)} }
\expandafter\def\csname SRef@maps:eq:D.ST8\endcsname{\href{S6_SERIES_GRIND_III.pdf\#nameddest=tag-0256}{III.{D.ST8} (p.~104)} }
\expandafter\def\csname SRef@maps:eq:D.ST9\endcsname{\href{S6_SERIES_GRIND_III.pdf\#nameddest=tag-0257}{III.{D.ST9} (p.~104)} }
\expandafter\def\csname SRef@maps:eq:D.ST10\endcsname{\href{S6_SERIES_GRIND_III.pdf\#nameddest=tag-024W}{III.{D.ST10} (p.~105)} }
\expandafter\def\csname SRef@maps:eq:D.ST11\endcsname{\href{S6_SERIES_GRIND_III.pdf\#nameddest=tag-024X}{III.{D.ST11} (p.~105)} }
\expandafter\def\csname SRef@maps:eq:D.ST12\endcsname{\href{S6_SERIES_GRIND_III.pdf\#nameddest=tag-024Y}{III.{D.ST12} (p.~105)} }
\expandafter\def\csname SRef@maps:eq:D.ST13\endcsname{\href{S6_SERIES_GRIND_III.pdf\#nameddest=tag-024Z}{III.{D.ST13} (p.~106)} }
\expandafter\def\csname SRef@maps:eq:D.UF1\endcsname{\href{S6_SERIES_GRIND_III.pdf\#nameddest=tag-025Z}{III.{D.UF1} (p.~107)} }
\expandafter\def\csname SRef@maps:eq:D.UF2\endcsname{\href{S6_SERIES_GRIND_III.pdf\#nameddest=tag-0261}{III.{D.UF2} (p.~107)} }
\expandafter\def\csname SRef@maps:eq:D.UF3\endcsname{\href{S6_SERIES_GRIND_III.pdf\#nameddest=tag-0262}{III.{D.UF3} (p.~107)} }
\expandafter\def\csname SRef@maps:eq:D.UF4\endcsname{\href{S6_SERIES_GRIND_III.pdf\#nameddest=tag-0263}{III.{D.UF4} (p.~107)} }
\expandafter\def\csname SRef@maps:eq:D.UF5\endcsname{\href{S6_SERIES_GRIND_III.pdf\#nameddest=tag-0264}{III.{D.UF5} (p.~108)} }
\expandafter\def\csname SRef@maps:eq:D.UF6\endcsname{\href{S6_SERIES_GRIND_III.pdf\#nameddest=tag-0265}{III.{D.UF6} (p.~108)} }
\expandafter\def\csname SRef@maps:eq:D.UF7\endcsname{\href{S6_SERIES_GRIND_III.pdf\#nameddest=tag-0266}{III.{D.UF7} (p.~108)} }
\expandafter\def\csname SRef@maps:eq:D.UF8\endcsname{\href{S6_SERIES_GRIND_III.pdf\#nameddest=tag-0267}{III.{D.UF8} (p.~109)} }
\expandafter\def\csname SRef@maps:eq:D.UF9\endcsname{\href{S6_SERIES_GRIND_III.pdf\#nameddest=tag-0268}{III.{D.UF9} (p.~109)} }
\expandafter\def\csname SRef@maps:eq:D.UF10\endcsname{\href{S6_SERIES_GRIND_III.pdf\#nameddest=tag-0260}{III.{D.UF10} (p.~109)} }
\expandafter\def\csname SRef@maps:eq:D.AF1\endcsname{\href{S6_SERIES_GRIND_III.pdf\#nameddest=tag-01XI}{III.{D.AF1} (p.~110)} }
\expandafter\def\csname SRef@maps:eq:D.AF2\endcsname{\href{S6_SERIES_GRIND_III.pdf\#nameddest=tag-01XJ}{III.{D.AF2} (p.~110)} }
\expandafter\def\csname SRef@maps:eq:D.AF3\endcsname{\href{S6_SERIES_GRIND_III.pdf\#nameddest=tag-01XK}{III.{D.AF3} (p.~111)} }
\expandafter\def\csname SRef@maps:eq:D.AF4\endcsname{\href{S6_SERIES_GRIND_III.pdf\#nameddest=tag-01XL}{III.{D.AF4} (p.~111)} }
\expandafter\def\csname SRef@maps:eq:D.AF5\endcsname{\href{S6_SERIES_GRIND_III.pdf\#nameddest=tag-01XM}{III.{D.AF5} (p.~111)} }
\expandafter\def\csname SRef@maps:eq:D.AF6\endcsname{\href{S6_SERIES_GRIND_III.pdf\#nameddest=tag-01XN}{III.{D.AF6} (p.~111)} }
\expandafter\def\csname SRef@maps:eq:D.AF7\endcsname{\href{S6_SERIES_GRIND_III.pdf\#nameddest=tag-01XO}{III.{D.AF7} (p.~111)} }
\expandafter\def\csname SRef@maps:eq:D.AF8\endcsname{\href{S6_SERIES_GRIND_III.pdf\#nameddest=tag-01XP}{III.{D.AF8} (p.~111)} }
\expandafter\def\csname SRef@maps:eq:D.AF9\endcsname{\href{S6_SERIES_GRIND_III.pdf\#nameddest=tag-01XQ}{III.{D.AF9} (p.~112)} }
\expandafter\def\csname SRef@maps:eq:D.OD1\endcsname{\href{S6_SERIES_GRIND_III.pdf\#nameddest=tag-0212}{III.{D.OD1} (p.~113)} }
\expandafter\def\csname SRef@maps:eq:D.OD2\endcsname{\href{S6_SERIES_GRIND_III.pdf\#nameddest=tag-0213}{III.{D.OD2} (p.~113)} }
\expandafter\def\csname SRef@maps:eq:D.OD3\endcsname{\href{S6_SERIES_GRIND_III.pdf\#nameddest=tag-0214}{III.{D.OD3} (p.~113)} }
\expandafter\def\csname SRef@maps:eq:D.OD4\endcsname{\href{S6_SERIES_GRIND_III.pdf\#nameddest=tag-0215}{III.{D.OD4} (p.~114)} }
\expandafter\def\csname SRef@maps:eq:D.OD5\endcsname{\href{S6_SERIES_GRIND_III.pdf\#nameddest=tag-0216}{III.{D.OD5} (p.~114)} }
\expandafter\def\csname SRef@maps:eq:D.OC1\endcsname{\href{S6_SERIES_GRIND_III.pdf\#nameddest=tag-020V}{III.{D.OC1} (p.~116)} }
\expandafter\def\csname SRef@maps:eq:D.OC2\endcsname{\href{S6_SERIES_GRIND_III.pdf\#nameddest=tag-020W}{III.{D.OC2} (p.~116)} }
\expandafter\def\csname SRef@maps:eq:D.OC3\endcsname{\href{S6_SERIES_GRIND_III.pdf\#nameddest=tag-020X}{III.{D.OC3} (p.~116)} }
\expandafter\def\csname SRef@maps:eq:D.OC4\endcsname{\href{S6_SERIES_GRIND_III.pdf\#nameddest=tag-020Y}{III.{D.OC4} (p.~117)} }
\expandafter\def\csname SRef@maps:eq:D.OC5\endcsname{\href{S6_SERIES_GRIND_III.pdf\#nameddest=tag-020Z}{III.{D.OC5} (p.~117)} }
\expandafter\def\csname SRef@maps:eq:D.OC6\endcsname{\href{S6_SERIES_GRIND_III.pdf\#nameddest=tag-0210}{III.{D.OC6} (p.~117)} }
\expandafter\def\csname SRef@maps:eq:D.OC7\endcsname{\href{S6_SERIES_GRIND_III.pdf\#nameddest=tag-0211}{III.{D.OC7} (p.~118)} }
\expandafter\def\csname SRef@maps:eq:D.OM1\endcsname{\href{S6_SERIES_GRIND_III.pdf\#nameddest=tag-0222}{III.{D.OM1} (p.~119)} }
\expandafter\def\csname SRef@maps:eq:D.OM2\endcsname{\href{S6_SERIES_GRIND_III.pdf\#nameddest=tag-0224}{III.{D.OM2} (p.~119)} }
\expandafter\def\csname SRef@maps:eq:D.OM3\endcsname{\href{S6_SERIES_GRIND_III.pdf\#nameddest=tag-0225}{III.{D.OM3} (p.~120)} }
\expandafter\def\csname SRef@maps:eq:D.OM4\endcsname{\href{S6_SERIES_GRIND_III.pdf\#nameddest=tag-0226}{III.{D.OM4} (p.~120)} }
\expandafter\def\csname SRef@maps:eq:D.OM5\endcsname{\href{S6_SERIES_GRIND_III.pdf\#nameddest=tag-0227}{III.{D.OM5} (p.~120)} }
\expandafter\def\csname SRef@maps:eq:D.OM6\endcsname{\href{S6_SERIES_GRIND_III.pdf\#nameddest=tag-0228}{III.{D.OM6} (p.~120)} }
\expandafter\def\csname SRef@maps:eq:D.OM7\endcsname{\href{S6_SERIES_GRIND_III.pdf\#nameddest=tag-0229}{III.{D.OM7} (p.~120)} }
\expandafter\def\csname SRef@maps:eq:D.OM8\endcsname{\href{S6_SERIES_GRIND_III.pdf\#nameddest=tag-022A}{III.{D.OM8} (p.~120)} }
\expandafter\def\csname SRef@maps:eq:D.OM9\endcsname{\href{S6_SERIES_GRIND_III.pdf\#nameddest=tag-022B}{III.{D.OM9} (p.~121)} }
\expandafter\def\csname SRef@maps:eq:B.1.1\endcsname{\href{S6_SERIES_GRIND_III.pdf\#nameddest=tag-01UV}{III.{B.1.1} (p.~14)} }
\expandafter\def\csname SRef@maps:eq:D.OM10\endcsname{\href{S6_SERIES_GRIND_III.pdf\#nameddest=tag-0223}{III.{D.OM10} (p.~121)} }
\expandafter\def\csname SRef@maps:eq:D.WE1\endcsname{\href{S6_SERIES_GRIND_III.pdf\#nameddest=tag-026O}{III.{D.WE1} (p.~122)} }
\expandafter\def\csname SRef@maps:eq:D.WE2\endcsname{\href{S6_SERIES_GRIND_III.pdf\#nameddest=tag-026U}{III.{D.WE2} (p.~122)} }
\expandafter\def\csname SRef@maps:eq:D.WE3\endcsname{\href{S6_SERIES_GRIND_III.pdf\#nameddest=tag-026V}{III.{D.WE3} (p.~122)} }
\expandafter\def\csname SRef@maps:eq:D.WE4\endcsname{\href{S6_SERIES_GRIND_III.pdf\#nameddest=tag-026W}{III.{D.WE4} (p.~122)} }
\expandafter\def\csname SRef@maps:eq:D.WE5\endcsname{\href{S6_SERIES_GRIND_III.pdf\#nameddest=tag-026X}{III.{D.WE5} (p.~122)} }
\expandafter\def\csname SRef@maps:eq:D.WE6\endcsname{\href{S6_SERIES_GRIND_III.pdf\#nameddest=tag-026Y}{III.{D.WE6} (p.~122)} }
\expandafter\def\csname SRef@maps:eq:D.WE7\endcsname{\href{S6_SERIES_GRIND_III.pdf\#nameddest=tag-026Z}{III.{D.WE7} (p.~123)} }
\expandafter\def\csname SRef@maps:eq:D.WE8\endcsname{\href{S6_SERIES_GRIND_III.pdf\#nameddest=tag-0270}{III.{D.WE8} (p.~123)} }
\expandafter\def\csname SRef@maps:eq:D.WE9\endcsname{\href{S6_SERIES_GRIND_III.pdf\#nameddest=tag-0271}{III.{D.WE9} (p.~123)} }
\expandafter\def\csname SRef@maps:eq:B.1.2\endcsname{\href{S6_SERIES_GRIND_III.pdf\#nameddest=tag-01UW}{III.{B.1.2} (p.~14)} }
\expandafter\def\csname SRef@maps:eq:D.WE10\endcsname{\href{S6_SERIES_GRIND_III.pdf\#nameddest=tag-026P}{III.{D.WE10} (p.~123)} }
\expandafter\def\csname SRef@maps:eq:D.WE11\endcsname{\href{S6_SERIES_GRIND_III.pdf\#nameddest=tag-026Q}{III.{D.WE11} (p.~124)} }
\expandafter\def\csname SRef@maps:eq:D.WE12\endcsname{\href{S6_SERIES_GRIND_III.pdf\#nameddest=tag-026R}{III.{D.WE12} (p.~125)} }
\expandafter\def\csname SRef@maps:eq:D.WE13\endcsname{\href{S6_SERIES_GRIND_III.pdf\#nameddest=tag-026S}{III.{D.WE13} (p.~125)} }
\expandafter\def\csname SRef@maps:eq:D.WE14\endcsname{\href{S6_SERIES_GRIND_III.pdf\#nameddest=tag-026T}{III.{D.WE14} (p.~126)} }
\expandafter\def\csname SRef@maps:eq:D.C21\endcsname{\href{S6_SERIES_GRIND_III.pdf\#nameddest=tag-01Y7}{III.{D.C21} (p.~126)} }
\expandafter\def\csname SRef@maps:eq:D.C22\endcsname{\href{S6_SERIES_GRIND_III.pdf\#nameddest=tag-01YF}{III.{D.C22} (p.~126)} }
\expandafter\def\csname SRef@maps:eq:D.C23\endcsname{\href{S6_SERIES_GRIND_III.pdf\#nameddest=tag-01YG}{III.{D.C23} (p.~127)} }
\expandafter\def\csname SRef@maps:eq:D.C24\endcsname{\href{S6_SERIES_GRIND_III.pdf\#nameddest=tag-01YH}{III.{D.C24} (p.~127)} }
\expandafter\def\csname SRef@maps:eq:D.C25\endcsname{\href{S6_SERIES_GRIND_III.pdf\#nameddest=tag-01YI}{III.{D.C25} (p.~127)} }
\expandafter\def\csname SRef@maps:eq:B.1.3\endcsname{\href{S6_SERIES_GRIND_III.pdf\#nameddest=tag-01UX}{III.{B.1.3} (p.~14)} }
\expandafter\def\csname SRef@maps:eq:D.C26\endcsname{\href{S6_SERIES_GRIND_III.pdf\#nameddest=tag-01YJ}{III.{D.C26} (p.~127)} }
\expandafter\def\csname SRef@maps:eq:D.C27\endcsname{\href{S6_SERIES_GRIND_III.pdf\#nameddest=tag-01YK}{III.{D.C27} (p.~127)} }
\expandafter\def\csname SRef@maps:eq:D.C28\endcsname{\href{S6_SERIES_GRIND_III.pdf\#nameddest=tag-01YL}{III.{D.C28} (p.~127)} }
\expandafter\def\csname SRef@maps:eq:D.C29\endcsname{\href{S6_SERIES_GRIND_III.pdf\#nameddest=tag-01YM}{III.{D.C29} (p.~128)} }
\expandafter\def\csname SRef@maps:eq:D.C210\endcsname{\href{S6_SERIES_GRIND_III.pdf\#nameddest=tag-01Y8}{III.{D.C210} (p.~128)} }
\expandafter\def\csname SRef@maps:eq:D.C211\endcsname{\href{S6_SERIES_GRIND_III.pdf\#nameddest=tag-01Y9}{III.{D.C211} (p.~128)} }
\expandafter\def\csname SRef@maps:eq:D.C212\endcsname{\href{S6_SERIES_GRIND_III.pdf\#nameddest=tag-01YA}{III.{D.C212} (p.~129)} }
\expandafter\def\csname SRef@maps:eq:D.C213\endcsname{\href{S6_SERIES_GRIND_III.pdf\#nameddest=tag-01YB}{III.{D.C213} (p.~129)} }
\expandafter\def\csname SRef@maps:eq:D.C214\endcsname{\href{S6_SERIES_GRIND_III.pdf\#nameddest=tag-01YC}{III.{D.C214} (p.~129)} }
\expandafter\def\csname SRef@maps:eq:D.C215\endcsname{\href{S6_SERIES_GRIND_III.pdf\#nameddest=tag-01YD}{III.{D.C215} (p.~129)} }
\expandafter\def\csname SRef@maps:eq:B.2.1\endcsname{\href{S6_SERIES_GRIND_III.pdf\#nameddest=tag-01UY}{III.{B.2.1} (p.~15)} }
\expandafter\def\csname SRef@maps:eq:D.C216\endcsname{\href{S6_SERIES_GRIND_III.pdf\#nameddest=tag-01YE}{III.{D.C216} (p.~130)} }
\expandafter\def\csname SRef@maps:eq:D.OP1\endcsname{\href{S6_SERIES_GRIND_III.pdf\#nameddest=tag-022C}{III.{D.OP1} (p.~131)} }
\expandafter\def\csname SRef@maps:eq:D.OP2\endcsname{\href{S6_SERIES_GRIND_III.pdf\#nameddest=tag-022F}{III.{D.OP2} (p.~131)} }
\expandafter\def\csname SRef@maps:eq:D.OP3\endcsname{\href{S6_SERIES_GRIND_III.pdf\#nameddest=tag-022G}{III.{D.OP3} (p.~131)} }
\expandafter\def\csname SRef@maps:eq:D.OP4\endcsname{\href{S6_SERIES_GRIND_III.pdf\#nameddest=tag-022H}{III.{D.OP4} (p.~131)} }
\expandafter\def\csname SRef@maps:eq:D.OP5\endcsname{\href{S6_SERIES_GRIND_III.pdf\#nameddest=tag-022I}{III.{D.OP5} (p.~132)} }
\expandafter\def\csname SRef@maps:eq:D.OP6\endcsname{\href{S6_SERIES_GRIND_III.pdf\#nameddest=tag-022K}{III.{D.OP6} (p.~132)} }
\expandafter\def\csname SRef@maps:eq:D.OP5a\endcsname{\href{S6_SERIES_GRIND_III.pdf\#nameddest=tag-022J}{III.{D.OP5a} (p.~132)} }
\expandafter\def\csname SRef@maps:eq:D.OP7\endcsname{\href{S6_SERIES_GRIND_III.pdf\#nameddest=tag-022L}{III.{D.OP7} (p.~133)} }
\expandafter\def\csname SRef@maps:eq:D.OP8\endcsname{\href{S6_SERIES_GRIND_III.pdf\#nameddest=tag-022M}{III.{D.OP8} (p.~133)} }
\expandafter\def\csname SRef@maps:eq:B.3.1\endcsname{\href{S6_SERIES_GRIND_III.pdf\#nameddest=tag-01UZ}{III.{B.3.1} (p.~16)} }
\expandafter\def\csname SRef@maps:eq:D.OP9\endcsname{\href{S6_SERIES_GRIND_III.pdf\#nameddest=tag-022N}{III.{D.OP9} (p.~133)} }
\expandafter\def\csname SRef@maps:eq:D.OP10\endcsname{\href{S6_SERIES_GRIND_III.pdf\#nameddest=tag-022D}{III.{D.OP10} (p.~133)} }
\expandafter\def\csname SRef@maps:eq:D.OP11\endcsname{\href{S6_SERIES_GRIND_III.pdf\#nameddest=tag-022E}{III.{D.OP11} (p.~133)} }
\expandafter\def\csname SRef@maps:eq:D.WP1\endcsname{\href{S6_SERIES_GRIND_III.pdf\#nameddest=tag-027O}{III.{D.WP1} (p.~134)} }
\expandafter\def\csname SRef@maps:eq:D.WP2\endcsname{\href{S6_SERIES_GRIND_III.pdf\#nameddest=tag-027S}{III.{D.WP2} (p.~134)} }
\expandafter\def\csname SRef@maps:eq:D.WP3\endcsname{\href{S6_SERIES_GRIND_III.pdf\#nameddest=tag-027T}{III.{D.WP3} (p.~134)} }
\expandafter\def\csname SRef@maps:eq:D.WP4\endcsname{\href{S6_SERIES_GRIND_III.pdf\#nameddest=tag-027U}{III.{D.WP4} (p.~134)} }
\expandafter\def\csname SRef@maps:eq:D.WP5\endcsname{\href{S6_SERIES_GRIND_III.pdf\#nameddest=tag-027V}{III.{D.WP5} (p.~135)} }
\expandafter\def\csname SRef@maps:eq:D.WP6\endcsname{\href{S6_SERIES_GRIND_III.pdf\#nameddest=tag-027W}{III.{D.WP6} (p.~135)} }
\expandafter\def\csname SRef@maps:eq:D.WP7\endcsname{\href{S6_SERIES_GRIND_III.pdf\#nameddest=tag-027X}{III.{D.WP7} (p.~135)} }
\expandafter\def\csname SRef@maps:eq:B.3.2\endcsname{\href{S6_SERIES_GRIND_III.pdf\#nameddest=tag-01V0}{III.{B.3.2} (p.~16)} }
\expandafter\def\csname SRef@maps:eq:D.WP8\endcsname{\href{S6_SERIES_GRIND_III.pdf\#nameddest=tag-027Y}{III.{D.WP8} (p.~135)} }
\expandafter\def\csname SRef@maps:eq:D.WP9\endcsname{\href{S6_SERIES_GRIND_III.pdf\#nameddest=tag-027Z}{III.{D.WP9} (p.~135)} }
\expandafter\def\csname SRef@maps:eq:D.WP10\endcsname{\href{S6_SERIES_GRIND_III.pdf\#nameddest=tag-027P}{III.{D.WP10} (p.~136)} }
\expandafter\def\csname SRef@maps:eq:D.WP11\endcsname{\href{S6_SERIES_GRIND_III.pdf\#nameddest=tag-027Q}{III.{D.WP11} (p.~137)} }
\expandafter\def\csname SRef@maps:eq:D.WP12\endcsname{\href{S6_SERIES_GRIND_III.pdf\#nameddest=tag-027R}{III.{D.WP12} (p.~138)} }
\expandafter\def\csname SRef@maps:eq:D.MP1\endcsname{\href{S6_SERIES_GRIND_III.pdf\#nameddest=tag-0209}{III.{D.MP1} (p.~138)} }
\expandafter\def\csname SRef@maps:eq:D.MP2\endcsname{\href{S6_SERIES_GRIND_III.pdf\#nameddest=tag-020B}{III.{D.MP2} (p.~139)} }
\expandafter\def\csname SRef@maps:eq:D.MP3\endcsname{\href{S6_SERIES_GRIND_III.pdf\#nameddest=tag-020C}{III.{D.MP3} (p.~139)} }
\expandafter\def\csname SRef@maps:eq:D.MP4\endcsname{\href{S6_SERIES_GRIND_III.pdf\#nameddest=tag-020D}{III.{D.MP4} (p.~139)} }
\expandafter\def\csname SRef@maps:eq:D.MP5\endcsname{\href{S6_SERIES_GRIND_III.pdf\#nameddest=tag-020E}{III.{D.MP5} (p.~140)} }
\expandafter\def\csname SRef@maps:eq:B.3.3\endcsname{\href{S6_SERIES_GRIND_III.pdf\#nameddest=tag-01V1}{III.{B.3.3} (p.~16)} }
\expandafter\def\csname SRef@maps:eq:D.MP6\endcsname{\href{S6_SERIES_GRIND_III.pdf\#nameddest=tag-020F}{III.{D.MP6} (p.~140)} }
\expandafter\def\csname SRef@maps:eq:D.MP7\endcsname{\href{S6_SERIES_GRIND_III.pdf\#nameddest=tag-020G}{III.{D.MP7} (p.~140)} }
\expandafter\def\csname SRef@maps:eq:D.MP8\endcsname{\href{S6_SERIES_GRIND_III.pdf\#nameddest=tag-020H}{III.{D.MP8} (p.~140)} }
\expandafter\def\csname SRef@maps:eq:D.MP9\endcsname{\href{S6_SERIES_GRIND_III.pdf\#nameddest=tag-020I}{III.{D.MP9} (p.~141)} }
\expandafter\def\csname SRef@maps:eq:D.MP10\endcsname{\href{S6_SERIES_GRIND_III.pdf\#nameddest=tag-020A}{III.{D.MP10} (p.~141)} }
\expandafter\def\csname SRef@maps:eq:D.T31\endcsname{\href{S6_SERIES_GRIND_III.pdf\#nameddest=tag-025C}{III.{D.T31} (p.~142)} }
\expandafter\def\csname SRef@maps:eq:D.T32\endcsname{\href{S6_SERIES_GRIND_III.pdf\#nameddest=tag-025D}{III.{D.T32} (p.~142)} }
\expandafter\def\csname SRef@maps:eq:D.T33\endcsname{\href{S6_SERIES_GRIND_III.pdf\#nameddest=tag-025E}{III.{D.T33} (p.~142)} }
\expandafter\def\csname SRef@maps:eq:D.T34\endcsname{\href{S6_SERIES_GRIND_III.pdf\#nameddest=tag-025F}{III.{D.T34} (p.~143)} }
\expandafter\def\csname SRef@maps:eq:D.T35\endcsname{\href{S6_SERIES_GRIND_III.pdf\#nameddest=tag-025G}{III.{D.T35} (p.~143)} }
\expandafter\def\csname SRef@maps:eq:B.3.4\endcsname{\href{S6_SERIES_GRIND_III.pdf\#nameddest=tag-01V2}{III.{B.3.4} (p.~16)} }
\expandafter\def\csname SRef@maps:eq:D.T36\endcsname{\href{S6_SERIES_GRIND_III.pdf\#nameddest=tag-025H}{III.{D.T36} (p.~143)} }
\expandafter\def\csname SRef@maps:eq:D.T37\endcsname{\href{S6_SERIES_GRIND_III.pdf\#nameddest=tag-025I}{III.{D.T37} (p.~144)} }
\expandafter\def\csname SRef@maps:eq:D.T38\endcsname{\href{S6_SERIES_GRIND_III.pdf\#nameddest=tag-025J}{III.{D.T38} (p.~144)} }
\expandafter\def\csname SRef@maps:eq:D.T39\endcsname{\href{S6_SERIES_GRIND_III.pdf\#nameddest=tag-025K}{III.{D.T39} (p.~144)} }
\expandafter\def\csname SRef@maps:eq:D.OW1\endcsname{\href{S6_SERIES_GRIND_III.pdf\#nameddest=tag-022W}{III.{D.OW1} (p.~145)} }
\expandafter\def\csname SRef@maps:eq:D.OW2\endcsname{\href{S6_SERIES_GRIND_III.pdf\#nameddest=tag-022X}{III.{D.OW2} (p.~145)} }
\expandafter\def\csname SRef@maps:eq:D.OW3\endcsname{\href{S6_SERIES_GRIND_III.pdf\#nameddest=tag-022Y}{III.{D.OW3} (p.~145)} }
\expandafter\def\csname SRef@maps:eq:D.OW4\endcsname{\href{S6_SERIES_GRIND_III.pdf\#nameddest=tag-022Z}{III.{D.OW4} (p.~145)} }
\expandafter\def\csname SRef@maps:eq:D.OW5\endcsname{\href{S6_SERIES_GRIND_III.pdf\#nameddest=tag-0230}{III.{D.OW5} (p.~145)} }
\expandafter\def\csname SRef@maps:eq:D.OW6\endcsname{\href{S6_SERIES_GRIND_III.pdf\#nameddest=tag-0231}{III.{D.OW6} (p.~146)} }
\expandafter\def\csname SRef@maps:eq:B.3.5\endcsname{\href{S6_SERIES_GRIND_III.pdf\#nameddest=tag-01V3}{III.{B.3.5} (p.~18)} }
\expandafter\def\csname SRef@maps:eq:D.OW7\endcsname{\href{S6_SERIES_GRIND_III.pdf\#nameddest=tag-0232}{III.{D.OW7} (p.~147)} }
\expandafter\def\csname SRef@maps:eq:D.BC1\endcsname{\href{S6_SERIES_GRIND_III.pdf\#nameddest=tag-01XR}{III.{D.BC1} (p.~148)} }
\expandafter\def\csname SRef@maps:eq:D.BC2\endcsname{\href{S6_SERIES_GRIND_III.pdf\#nameddest=tag-01XY}{III.{D.BC2} (p.~148)} }
\expandafter\def\csname SRef@maps:eq:D.BC3\endcsname{\href{S6_SERIES_GRIND_III.pdf\#nameddest=tag-01XZ}{III.{D.BC3} (p.~148)} }
\expandafter\def\csname SRef@maps:eq:D.BC4\endcsname{\href{S6_SERIES_GRIND_III.pdf\#nameddest=tag-01Y0}{III.{D.BC4} (p.~148)} }
\expandafter\def\csname SRef@maps:eq:D.BC5\endcsname{\href{S6_SERIES_GRIND_III.pdf\#nameddest=tag-01Y1}{III.{D.BC5} (p.~148)} }
\expandafter\def\csname SRef@maps:eq:D.BC6\endcsname{\href{S6_SERIES_GRIND_III.pdf\#nameddest=tag-01Y2}{III.{D.BC6} (p.~149)} }
\expandafter\def\csname SRef@maps:eq:D.BC7\endcsname{\href{S6_SERIES_GRIND_III.pdf\#nameddest=tag-01Y3}{III.{D.BC7} (p.~149)} }
\expandafter\def\csname SRef@maps:eq:D.BC8\endcsname{\href{S6_SERIES_GRIND_III.pdf\#nameddest=tag-01Y4}{III.{D.BC8} (p.~149)} }
\expandafter\def\csname SRef@maps:eq:D.BC9\endcsname{\href{S6_SERIES_GRIND_III.pdf\#nameddest=tag-01Y5}{III.{D.BC9} (p.~149)} }
\expandafter\def\csname SRef@maps:eq:B.4.1\endcsname{\href{S6_SERIES_GRIND_III.pdf\#nameddest=tag-01V4}{III.{B.4.1} (p.~19)} }
\expandafter\def\csname SRef@maps:eq:D.BC10\endcsname{\href{S6_SERIES_GRIND_III.pdf\#nameddest=tag-01XS}{III.{D.BC10} (p.~149)} }
\expandafter\def\csname SRef@maps:eq:D.BC11\endcsname{\href{S6_SERIES_GRIND_III.pdf\#nameddest=tag-01XT}{III.{D.BC11} (p.~149)} }
\expandafter\def\csname SRef@maps:eq:D.BC12\endcsname{\href{S6_SERIES_GRIND_III.pdf\#nameddest=tag-01XU}{III.{D.BC12} (p.~149)} }
\expandafter\def\csname SRef@maps:eq:D.BC13\endcsname{\href{S6_SERIES_GRIND_III.pdf\#nameddest=tag-01XV}{III.{D.BC13} (p.~150)} }
\expandafter\def\csname SRef@maps:eq:D.BC14\endcsname{\href{S6_SERIES_GRIND_III.pdf\#nameddest=tag-01XW}{III.{D.BC14} (p.~150)} }
\expandafter\def\csname SRef@maps:eq:D.BC15\endcsname{\href{S6_SERIES_GRIND_III.pdf\#nameddest=tag-01XX}{III.{D.BC15} (p.~150)} }
\expandafter\def\csname SRef@maps:eq:D.MT1\endcsname{\href{S6_SERIES_GRIND_III.pdf\#nameddest=tag-020J}{III.{D.MT1} (p.~151)} }
\expandafter\def\csname SRef@maps:eq:D.MT2\endcsname{\href{S6_SERIES_GRIND_III.pdf\#nameddest=tag-020K}{III.{D.MT2} (p.~152)} }
\expandafter\def\csname SRef@maps:eq:D.MT3\endcsname{\href{S6_SERIES_GRIND_III.pdf\#nameddest=tag-020L}{III.{D.MT3} (p.~153)} }
\expandafter\def\csname SRef@maps:eq:D.GC1\endcsname{\href{S6_SERIES_GRIND_III.pdf\#nameddest=tag-01YY}{III.{D.GC1} (p.~153)} }
\expandafter\def\csname SRef@maps:eq:B.4.2\endcsname{\href{S6_SERIES_GRIND_III.pdf\#nameddest=tag-01V5}{III.{B.4.2} (p.~19)} }
\expandafter\def\csname SRef@maps:eq:D.GC2\endcsname{\href{S6_SERIES_GRIND_III.pdf\#nameddest=tag-01Z4}{III.{D.GC2} (p.~153)} }
\expandafter\def\csname SRef@maps:eq:D.GC3\endcsname{\href{S6_SERIES_GRIND_III.pdf\#nameddest=tag-01Z5}{III.{D.GC3} (p.~154)} }
\expandafter\def\csname SRef@maps:eq:D.GC4\endcsname{\href{S6_SERIES_GRIND_III.pdf\#nameddest=tag-01Z6}{III.{D.GC4} (p.~154)} }
\expandafter\def\csname SRef@maps:eq:D.GC5\endcsname{\href{S6_SERIES_GRIND_III.pdf\#nameddest=tag-01Z7}{III.{D.GC5} (p.~154)} }
\expandafter\def\csname SRef@maps:eq:D.GC6\endcsname{\href{S6_SERIES_GRIND_III.pdf\#nameddest=tag-01Z8}{III.{D.GC6} (p.~154)} }
\expandafter\def\csname SRef@maps:eq:D.GC7\endcsname{\href{S6_SERIES_GRIND_III.pdf\#nameddest=tag-01Z9}{III.{D.GC7} (p.~154)} }
\expandafter\def\csname SRef@maps:eq:D.GC8\endcsname{\href{S6_SERIES_GRIND_III.pdf\#nameddest=tag-01ZA}{III.{D.GC8} (p.~155)} }
\expandafter\def\csname SRef@maps:eq:D.GC9\endcsname{\href{S6_SERIES_GRIND_III.pdf\#nameddest=tag-01ZB}{III.{D.GC9} (p.~155)} }
\expandafter\def\csname SRef@maps:eq:D.GC10\endcsname{\href{S6_SERIES_GRIND_III.pdf\#nameddest=tag-01YZ}{III.{D.GC10} (p.~155)} }
\expandafter\def\csname SRef@maps:eq:D.GC11\endcsname{\href{S6_SERIES_GRIND_III.pdf\#nameddest=tag-01Z0}{III.{D.GC11} (p.~155)} }
\expandafter\def\csname SRef@maps:eq:B.4.3\endcsname{\href{S6_SERIES_GRIND_III.pdf\#nameddest=tag-01V6}{III.{B.4.3} (p.~20)} }
\expandafter\def\csname SRef@maps:eq:D.GC12\endcsname{\href{S6_SERIES_GRIND_III.pdf\#nameddest=tag-01Z1}{III.{D.GC12} (p.~155)} }
\expandafter\def\csname SRef@maps:eq:D.GC13\endcsname{\href{S6_SERIES_GRIND_III.pdf\#nameddest=tag-01Z2}{III.{D.GC13} (p.~155)} }
\expandafter\def\csname SRef@maps:eq:D.GC14\endcsname{\href{S6_SERIES_GRIND_III.pdf\#nameddest=tag-01Z3}{III.{D.GC14} (p.~156)} }
\expandafter\def\csname SRef@maps:eq:D44.1\endcsname{\href{S6_SERIES_GRIND_III.pdf\#nameddest=tag-0282}{III.{D44.1} (p.~157)} }
\expandafter\def\csname SRef@maps:eq:D44.0a\endcsname{\href{S6_SERIES_GRIND_III.pdf\#nameddest=tag-0280}{III.{D44.0a} (p.~157)} }
\expandafter\def\csname SRef@maps:eq:D44.0b\endcsname{\href{S6_SERIES_GRIND_III.pdf\#nameddest=tag-0281}{III.{D44.0b} (p.~157)} }
\expandafter\def\csname SRef@maps:eq:D44.2\endcsname{\href{S6_SERIES_GRIND_III.pdf\#nameddest=tag-0288}{III.{D44.2} (p.~158)} }
\expandafter\def\csname SRef@maps:eq:D44.3\endcsname{\href{S6_SERIES_GRIND_III.pdf\#nameddest=tag-0289}{III.{D44.3} (p.~158)} }
\expandafter\def\csname SRef@maps:eq:D44.4\endcsname{\href{S6_SERIES_GRIND_III.pdf\#nameddest=tag-028A}{III.{D44.4} (p.~159)} }
\expandafter\def\csname SRef@maps:eq:D44.5\endcsname{\href{S6_SERIES_GRIND_III.pdf\#nameddest=tag-028B}{III.{D44.5} (p.~159)} }
\expandafter\def\csname SRef@maps:eq:B.5.1\endcsname{\href{S6_SERIES_GRIND_III.pdf\#nameddest=tag-01V7}{III.{B.5.1} (p.~20)} }
\expandafter\def\csname SRef@maps:eq:D44.6\endcsname{\href{S6_SERIES_GRIND_III.pdf\#nameddest=tag-028C}{III.{D44.6} (p.~159)} }
\expandafter\def\csname SRef@maps:eq:D44.7\endcsname{\href{S6_SERIES_GRIND_III.pdf\#nameddest=tag-028D}{III.{D44.7} (p.~159)} }
\expandafter\def\csname SRef@maps:eq:D44.8\endcsname{\href{S6_SERIES_GRIND_III.pdf\#nameddest=tag-028E}{III.{D44.8} (p.~160)} }
\expandafter\def\csname SRef@maps:eq:D44.9\endcsname{\href{S6_SERIES_GRIND_III.pdf\#nameddest=tag-028F}{III.{D44.9} (p.~160)} }
\expandafter\def\csname SRef@maps:eq:D44.10\endcsname{\href{S6_SERIES_GRIND_III.pdf\#nameddest=tag-0283}{III.{D44.10} (p.~160)} }
\expandafter\def\csname SRef@maps:eq:D44.11\endcsname{\href{S6_SERIES_GRIND_III.pdf\#nameddest=tag-0284}{III.{D44.11} (p.~160)} }
\expandafter\def\csname SRef@maps:eq:D44.12\endcsname{\href{S6_SERIES_GRIND_III.pdf\#nameddest=tag-0285}{III.{D44.12} (p.~160)} }
\expandafter\def\csname SRef@maps:eq:D44.13\endcsname{\href{S6_SERIES_GRIND_III.pdf\#nameddest=tag-0286}{III.{D44.13} (p.~161)} }
\expandafter\def\csname SRef@maps:eq:D44.14\endcsname{\href{S6_SERIES_GRIND_III.pdf\#nameddest=tag-0287}{III.{D44.14} (p.~161)} }
\expandafter\def\csname SRef@maps:eq:D45.1\endcsname{\href{S6_SERIES_GRIND_III.pdf\#nameddest=tag-028G}{III.{D45.1} (p.~162)} }
\expandafter\def\csname SRef@maps:eq:B.5.2\endcsname{\href{S6_SERIES_GRIND_III.pdf\#nameddest=tag-01VH}{III.{B.5.2} (p.~21)} }
\expandafter\def\csname SRef@maps:eq:D45.2\endcsname{\href{S6_SERIES_GRIND_III.pdf\#nameddest=tag-028R}{III.{D45.2} (p.~162)} }
\expandafter\def\csname SRef@maps:eq:D45.3\endcsname{\href{S6_SERIES_GRIND_III.pdf\#nameddest=tag-028V}{III.{D45.3} (p.~162)} }
\expandafter\def\csname SRef@maps:eq:D45.4\endcsname{\href{S6_SERIES_GRIND_III.pdf\#nameddest=tag-028W}{III.{D45.4} (p.~162)} }
\expandafter\def\csname SRef@maps:eq:D45.5\endcsname{\href{S6_SERIES_GRIND_III.pdf\#nameddest=tag-028X}{III.{D45.5} (p.~162)} }
\expandafter\def\csname SRef@maps:eq:D45.6\endcsname{\href{S6_SERIES_GRIND_III.pdf\#nameddest=tag-028Y}{III.{D45.6} (p.~163)} }
\expandafter\def\csname SRef@maps:eq:D45.7\endcsname{\href{S6_SERIES_GRIND_III.pdf\#nameddest=tag-028Z}{III.{D45.7} (p.~163)} }
\expandafter\def\csname SRef@maps:eq:D45.8\endcsname{\href{S6_SERIES_GRIND_III.pdf\#nameddest=tag-0290}{III.{D45.8} (p.~163)} }
\expandafter\def\csname SRef@maps:eq:D45.9\endcsname{\href{S6_SERIES_GRIND_III.pdf\#nameddest=tag-0291}{III.{D45.9} (p.~163)} }
\expandafter\def\csname SRef@maps:eq:D45.10\endcsname{\href{S6_SERIES_GRIND_III.pdf\#nameddest=tag-028H}{III.{D45.10} (p.~164)} }
\expandafter\def\csname SRef@maps:eq:D45.11\endcsname{\href{S6_SERIES_GRIND_III.pdf\#nameddest=tag-028I}{III.{D45.11} (p.~164)} }
\expandafter\def\csname SRef@maps:eq:B.5.3\endcsname{\href{S6_SERIES_GRIND_III.pdf\#nameddest=tag-01VI}{III.{B.5.3} (p.~21)} }
\expandafter\def\csname SRef@maps:eq:D45.12\endcsname{\href{S6_SERIES_GRIND_III.pdf\#nameddest=tag-028J}{III.{D45.12} (p.~164)} }
\expandafter\def\csname SRef@maps:eq:D45.13\endcsname{\href{S6_SERIES_GRIND_III.pdf\#nameddest=tag-028K}{III.{D45.13} (p.~164)} }
\expandafter\def\csname SRef@maps:eq:D45.14\endcsname{\href{S6_SERIES_GRIND_III.pdf\#nameddest=tag-028L}{III.{D45.14} (p.~164)} }
\expandafter\def\csname SRef@maps:eq:D45.15\endcsname{\href{S6_SERIES_GRIND_III.pdf\#nameddest=tag-028M}{III.{D45.15} (p.~164)} }
\expandafter\def\csname SRef@maps:eq:D45.16\endcsname{\href{S6_SERIES_GRIND_III.pdf\#nameddest=tag-028N}{III.{D45.16} (p.~165)} }
\expandafter\def\csname SRef@maps:eq:D45.17\endcsname{\href{S6_SERIES_GRIND_III.pdf\#nameddest=tag-028O}{III.{D45.17} (p.~165)} }
\expandafter\def\csname SRef@maps:eq:D45.18\endcsname{\href{S6_SERIES_GRIND_III.pdf\#nameddest=tag-028P}{III.{D45.18} (p.~165)} }
\expandafter\def\csname SRef@maps:eq:D45.19\endcsname{\href{S6_SERIES_GRIND_III.pdf\#nameddest=tag-028Q}{III.{D45.19} (p.~165)} }
\expandafter\def\csname SRef@maps:eq:D45.20\endcsname{\href{S6_SERIES_GRIND_III.pdf\#nameddest=tag-028S}{III.{D45.20} (p.~165)} }
\expandafter\def\csname SRef@maps:eq:D45.21\endcsname{\href{S6_SERIES_GRIND_III.pdf\#nameddest=tag-028T}{III.{D45.21} (p.~166)} }
\expandafter\def\csname SRef@maps:eq:B.5.4\endcsname{\href{S6_SERIES_GRIND_III.pdf\#nameddest=tag-01VJ}{III.{B.5.4} (p.~21)} }
\expandafter\def\csname SRef@maps:eq:D45.22\endcsname{\href{S6_SERIES_GRIND_III.pdf\#nameddest=tag-028U}{III.{D45.22} (p.~166)} }
\expandafter\def\csname SRef@maps:eq:D.GI1\endcsname{\href{S6_SERIES_GRIND_III.pdf\#nameddest=tag-01ZC}{III.{D.GI1} (p.~167)} }
\expandafter\def\csname SRef@maps:eq:D.GI2\endcsname{\href{S6_SERIES_GRIND_III.pdf\#nameddest=tag-01ZF}{III.{D.GI2} (p.~167)} }
\expandafter\def\csname SRef@maps:eq:D.GI3\endcsname{\href{S6_SERIES_GRIND_III.pdf\#nameddest=tag-01ZG}{III.{D.GI3} (p.~167)} }
\expandafter\def\csname SRef@maps:eq:D.GI4\endcsname{\href{S6_SERIES_GRIND_III.pdf\#nameddest=tag-01ZH}{III.{D.GI4} (p.~167)} }
\expandafter\def\csname SRef@maps:eq:D.GI5\endcsname{\href{S6_SERIES_GRIND_III.pdf\#nameddest=tag-01ZI}{III.{D.GI5} (p.~168)} }
\expandafter\def\csname SRef@maps:eq:D.GI6\endcsname{\href{S6_SERIES_GRIND_III.pdf\#nameddest=tag-01ZJ}{III.{D.GI6} (p.~168)} }
\expandafter\def\csname SRef@maps:eq:D.GI7\endcsname{\href{S6_SERIES_GRIND_III.pdf\#nameddest=tag-01ZK}{III.{D.GI7} (p.~168)} }
\expandafter\def\csname SRef@maps:eq:D.GI8\endcsname{\href{S6_SERIES_GRIND_III.pdf\#nameddest=tag-01ZL}{III.{D.GI8} (p.~168)} }
\expandafter\def\csname SRef@maps:eq:D.GI9\endcsname{\href{S6_SERIES_GRIND_III.pdf\#nameddest=tag-01ZM}{III.{D.GI9} (p.~170)} }
\expandafter\def\csname SRef@maps:eq:B.5.5\endcsname{\href{S6_SERIES_GRIND_III.pdf\#nameddest=tag-01VK}{III.{B.5.5} (p.~22)} }
\expandafter\def\csname SRef@maps:eq:D.GI10\endcsname{\href{S6_SERIES_GRIND_III.pdf\#nameddest=tag-01ZD}{III.{D.GI10} (p.~171)} }
\expandafter\def\csname SRef@maps:eq:D.GI11\endcsname{\href{S6_SERIES_GRIND_III.pdf\#nameddest=tag-01ZE}{III.{D.GI11} (p.~171)} }
\expandafter\def\csname SRef@maps:eq:D47.1\endcsname{\href{S6_SERIES_GRIND_III.pdf\#nameddest=tag-0292}{III.{D47.1} (p.~172)} }
\expandafter\def\csname SRef@maps:eq:D47.2\endcsname{\href{S6_SERIES_GRIND_III.pdf\#nameddest=tag-0297}{III.{D47.2} (p.~172)} }
\expandafter\def\csname SRef@maps:eq:D47.3\endcsname{\href{S6_SERIES_GRIND_III.pdf\#nameddest=tag-0298}{III.{D47.3} (p.~172)} }
\expandafter\def\csname SRef@maps:eq:D47.4\endcsname{\href{S6_SERIES_GRIND_III.pdf\#nameddest=tag-0299}{III.{D47.4} (p.~172)} }
\expandafter\def\csname SRef@maps:eq:D47.5\endcsname{\href{S6_SERIES_GRIND_III.pdf\#nameddest=tag-029A}{III.{D47.5} (p.~173)} }
\expandafter\def\csname SRef@maps:eq:D47.6\endcsname{\href{S6_SERIES_GRIND_III.pdf\#nameddest=tag-029B}{III.{D47.6} (p.~173)} }
\expandafter\def\csname SRef@maps:eq:D47.7\endcsname{\href{S6_SERIES_GRIND_III.pdf\#nameddest=tag-029C}{III.{D47.7} (p.~173)} }
\expandafter\def\csname SRef@maps:eq:D47.8\endcsname{\href{S6_SERIES_GRIND_III.pdf\#nameddest=tag-029D}{III.{D47.8} (p.~173)} }
\expandafter\def\csname SRef@maps:eq:B.5.6\endcsname{\href{S6_SERIES_GRIND_III.pdf\#nameddest=tag-01VL}{III.{B.5.6} (p.~22)} }
\expandafter\def\csname SRef@maps:eq:D47.9\endcsname{\href{S6_SERIES_GRIND_III.pdf\#nameddest=tag-029E}{III.{D47.9} (p.~173)} }
\expandafter\def\csname SRef@maps:eq:D47.10\endcsname{\href{S6_SERIES_GRIND_III.pdf\#nameddest=tag-0293}{III.{D47.10} (p.~174)} }
\expandafter\def\csname SRef@maps:eq:D47.11\endcsname{\href{S6_SERIES_GRIND_III.pdf\#nameddest=tag-0294}{III.{D47.11} (p.~174)} }
\expandafter\def\csname SRef@maps:eq:D47.12\endcsname{\href{S6_SERIES_GRIND_III.pdf\#nameddest=tag-0295}{III.{D47.12} (p.~174)} }
\expandafter\def\csname SRef@maps:eq:D47.13\endcsname{\href{S6_SERIES_GRIND_III.pdf\#nameddest=tag-0296}{III.{D47.13} (p.~174)} }
\expandafter\def\csname SRef@maps:eq:D.WI1\endcsname{\href{S6_SERIES_GRIND_III.pdf\#nameddest=tag-0272}{III.{D.WI1} (p.~175)} }
\expandafter\def\csname SRef@maps:eq:D.WI2\endcsname{\href{S6_SERIES_GRIND_III.pdf\#nameddest=tag-027D}{III.{D.WI2} (p.~175)} }
\expandafter\def\csname SRef@maps:eq:D.WI3\endcsname{\href{S6_SERIES_GRIND_III.pdf\#nameddest=tag-027H}{III.{D.WI3} (p.~175)} }
\expandafter\def\csname SRef@maps:eq:D.WI14\endcsname{\href{S6_SERIES_GRIND_III.pdf\#nameddest=tag-0277}{III.{D.WI14} (p.~175)} }
\expandafter\def\csname SRef@maps:eq:D.WI15\endcsname{\href{S6_SERIES_GRIND_III.pdf\#nameddest=tag-0278}{III.{D.WI15} (p.~175)} }
\expandafter\def\csname SRef@maps:eq:B.5.7\endcsname{\href{S6_SERIES_GRIND_III.pdf\#nameddest=tag-01VM}{III.{B.5.7} (p.~22)} }
\expandafter\def\csname SRef@maps:eq:D.WI16\endcsname{\href{S6_SERIES_GRIND_III.pdf\#nameddest=tag-0279}{III.{D.WI16} (p.~175)} }
\expandafter\def\csname SRef@maps:eq:D.WI4\endcsname{\href{S6_SERIES_GRIND_III.pdf\#nameddest=tag-027I}{III.{D.WI4} (p.~176)} }
\expandafter\def\csname SRef@maps:eq:D.WI5\endcsname{\href{S6_SERIES_GRIND_III.pdf\#nameddest=tag-027J}{III.{D.WI5} (p.~176)} }
\expandafter\def\csname SRef@maps:eq:D.WI6\endcsname{\href{S6_SERIES_GRIND_III.pdf\#nameddest=tag-027K}{III.{D.WI6} (p.~177)} }
\expandafter\def\csname SRef@maps:eq:D.WI7\endcsname{\href{S6_SERIES_GRIND_III.pdf\#nameddest=tag-027L}{III.{D.WI7} (p.~177)} }
\expandafter\def\csname SRef@maps:eq:D.WI8\endcsname{\href{S6_SERIES_GRIND_III.pdf\#nameddest=tag-027M}{III.{D.WI8} (p.~177)} }
\expandafter\def\csname SRef@maps:eq:D.WI9\endcsname{\href{S6_SERIES_GRIND_III.pdf\#nameddest=tag-027N}{III.{D.WI9} (p.~177)} }
\expandafter\def\csname SRef@maps:eq:D.WI10\endcsname{\href{S6_SERIES_GRIND_III.pdf\#nameddest=tag-0273}{III.{D.WI10} (p.~177)} }
\expandafter\def\csname SRef@maps:eq:D.WI11\endcsname{\href{S6_SERIES_GRIND_III.pdf\#nameddest=tag-0274}{III.{D.WI11} (p.~177)} }
\expandafter\def\csname SRef@maps:eq:D.WI12\endcsname{\href{S6_SERIES_GRIND_III.pdf\#nameddest=tag-0275}{III.{D.WI12} (p.~177)} }
\expandafter\def\csname SRef@maps:eq:B.5.8\endcsname{\href{S6_SERIES_GRIND_III.pdf\#nameddest=tag-01VN}{III.{B.5.8} (p.~23)} }
\expandafter\def\csname SRef@maps:eq:D.WI13\endcsname{\href{S6_SERIES_GRIND_III.pdf\#nameddest=tag-0276}{III.{D.WI13} (p.~178)} }
\expandafter\def\csname SRef@maps:eq:D.WI17\endcsname{\href{S6_SERIES_GRIND_III.pdf\#nameddest=tag-027A}{III.{D.WI17} (p.~178)} }
\expandafter\def\csname SRef@maps:eq:D.WI18\endcsname{\href{S6_SERIES_GRIND_III.pdf\#nameddest=tag-027B}{III.{D.WI18} (p.~178)} }
\expandafter\def\csname SRef@maps:eq:D.WI19\endcsname{\href{S6_SERIES_GRIND_III.pdf\#nameddest=tag-027C}{III.{D.WI19} (p.~178)} }
\expandafter\def\csname SRef@maps:eq:D.WI20\endcsname{\href{S6_SERIES_GRIND_III.pdf\#nameddest=tag-027E}{III.{D.WI20} (p.~179)} }
\expandafter\def\csname SRef@maps:eq:D.WI21\endcsname{\href{S6_SERIES_GRIND_III.pdf\#nameddest=tag-027F}{III.{D.WI21} (p.~179)} }
\expandafter\def\csname SRef@maps:eq:D.TP1\endcsname{\href{S6_SERIES_GRIND_III.pdf\#nameddest=tag-025R}{III.{D.TP1} (p.~179)} }
\expandafter\def\csname SRef@maps:eq:D.TP2\endcsname{\href{S6_SERIES_GRIND_III.pdf\#nameddest=tag-025S}{III.{D.TP2} (p.~180)} }
\expandafter\def\csname SRef@maps:eq:D.TP3\endcsname{\href{S6_SERIES_GRIND_III.pdf\#nameddest=tag-025T}{III.{D.TP3} (p.~180)} }
\expandafter\def\csname SRef@maps:eq:B.5.9\endcsname{\href{S6_SERIES_GRIND_III.pdf\#nameddest=tag-01VO}{III.{B.5.9} (p.~24)} }
\expandafter\def\csname SRef@maps:eq:D.TP4\endcsname{\href{S6_SERIES_GRIND_III.pdf\#nameddest=tag-025U}{III.{D.TP4} (p.~180)} }
\expandafter\def\csname SRef@maps:eq:D.TP5\endcsname{\href{S6_SERIES_GRIND_III.pdf\#nameddest=tag-025V}{III.{D.TP5} (p.~180)} }
\expandafter\def\csname SRef@maps:eq:D.TP6\endcsname{\href{S6_SERIES_GRIND_III.pdf\#nameddest=tag-025W}{III.{D.TP6} (p.~180)} }
\expandafter\def\csname SRef@maps:eq:D.TP7\endcsname{\href{S6_SERIES_GRIND_III.pdf\#nameddest=tag-025X}{III.{D.TP7} (p.~180)} }
\expandafter\def\csname SRef@maps:eq:D.TP8\endcsname{\href{S6_SERIES_GRIND_III.pdf\#nameddest=tag-025Y}{III.{D.TP8} (p.~180)} }
\expandafter\def\csname SRef@maps:eq:D.WI22\endcsname{\href{S6_SERIES_GRIND_III.pdf\#nameddest=tag-027G}{III.{D.WI22} (p.~181)} }
\expandafter\def\csname SRef@maps:eq:D.SL1\endcsname{\href{S6_SERIES_GRIND_III.pdf\#nameddest=tag-0247}{III.{D.SL1} (p.~182)} }
\expandafter\def\csname SRef@maps:eq:D.SL2\endcsname{\href{S6_SERIES_GRIND_III.pdf\#nameddest=tag-024I}{III.{D.SL2} (p.~182)} }
\expandafter\def\csname SRef@maps:eq:D.SL3\endcsname{\href{S6_SERIES_GRIND_III.pdf\#nameddest=tag-024O}{III.{D.SL3} (p.~182)} }
\expandafter\def\csname SRef@maps:eq:D.SL4\endcsname{\href{S6_SERIES_GRIND_III.pdf\#nameddest=tag-024P}{III.{D.SL4} (p.~182)} }
\expandafter\def\csname SRef@maps:eq:B.5.10\endcsname{\href{S6_SERIES_GRIND_III.pdf\#nameddest=tag-01V8}{III.{B.5.10} (p.~24)} }
\expandafter\def\csname SRef@maps:eq:D.SL5\endcsname{\href{S6_SERIES_GRIND_III.pdf\#nameddest=tag-024Q}{III.{D.SL5} (p.~182)} }
\expandafter\def\csname SRef@maps:eq:D.SL6\endcsname{\href{S6_SERIES_GRIND_III.pdf\#nameddest=tag-024R}{III.{D.SL6} (p.~183)} }
\expandafter\def\csname SRef@maps:eq:D.SL7\endcsname{\href{S6_SERIES_GRIND_III.pdf\#nameddest=tag-024S}{III.{D.SL7} (p.~183)} }
\expandafter\def\csname SRef@maps:eq:D.SL8\endcsname{\href{S6_SERIES_GRIND_III.pdf\#nameddest=tag-024T}{III.{D.SL8} (p.~183)} }
\expandafter\def\csname SRef@maps:eq:D.SL9\endcsname{\href{S6_SERIES_GRIND_III.pdf\#nameddest=tag-024U}{III.{D.SL9} (p.~183)} }
\expandafter\def\csname SRef@maps:eq:D.SL10\endcsname{\href{S6_SERIES_GRIND_III.pdf\#nameddest=tag-0248}{III.{D.SL10} (p.~183)} }
\expandafter\def\csname SRef@maps:eq:D.SL11\endcsname{\href{S6_SERIES_GRIND_III.pdf\#nameddest=tag-0249}{III.{D.SL11} (p.~183)} }
\expandafter\def\csname SRef@maps:eq:D.SL12\endcsname{\href{S6_SERIES_GRIND_III.pdf\#nameddest=tag-024A}{III.{D.SL12} (p.~184)} }
\expandafter\def\csname SRef@maps:eq:D.SL13\endcsname{\href{S6_SERIES_GRIND_III.pdf\#nameddest=tag-024B}{III.{D.SL13} (p.~184)} }
\expandafter\def\csname SRef@maps:eq:D.SL14\endcsname{\href{S6_SERIES_GRIND_III.pdf\#nameddest=tag-024C}{III.{D.SL14} (p.~184)} }
\expandafter\def\csname SRef@maps:eq:B.5.11\endcsname{\href{S6_SERIES_GRIND_III.pdf\#nameddest=tag-01V9}{III.{B.5.11} (p.~24)} }
\expandafter\def\csname SRef@maps:eq:D.SL15\endcsname{\href{S6_SERIES_GRIND_III.pdf\#nameddest=tag-024D}{III.{D.SL15} (p.~185)} }
\expandafter\def\csname SRef@maps:eq:D.SL16\endcsname{\href{S6_SERIES_GRIND_III.pdf\#nameddest=tag-024E}{III.{D.SL16} (p.~185)} }
\expandafter\def\csname SRef@maps:eq:D.SL17\endcsname{\href{S6_SERIES_GRIND_III.pdf\#nameddest=tag-024F}{III.{D.SL17} (p.~185)} }
\expandafter\def\csname SRef@maps:eq:D.SL18\endcsname{\href{S6_SERIES_GRIND_III.pdf\#nameddest=tag-024G}{III.{D.SL18} (p.~185)} }
\expandafter\def\csname SRef@maps:eq:D.SL19\endcsname{\href{S6_SERIES_GRIND_III.pdf\#nameddest=tag-024H}{III.{D.SL19} (p.~185)} }
\expandafter\def\csname SRef@maps:eq:D.SL20\endcsname{\href{S6_SERIES_GRIND_III.pdf\#nameddest=tag-024J}{III.{D.SL20} (p.~186)} }
\expandafter\def\csname SRef@maps:eq:D.SL21\endcsname{\href{S6_SERIES_GRIND_III.pdf\#nameddest=tag-024K}{III.{D.SL21} (p.~186)} }
\expandafter\def\csname SRef@maps:eq:D.SL22\endcsname{\href{S6_SERIES_GRIND_III.pdf\#nameddest=tag-024L}{III.{D.SL22} (p.~186)} }
\expandafter\def\csname SRef@maps:eq:D.SL24\endcsname{\href{S6_SERIES_GRIND_III.pdf\#nameddest=tag-024N}{III.{D.SL24} (p.~186)} }
\expandafter\def\csname SRef@maps:eq:D.SL23\endcsname{\href{S6_SERIES_GRIND_III.pdf\#nameddest=tag-024M}{III.{D.SL23} (p.~187)} }
\expandafter\def\csname SRef@maps:eq:B.5.12\endcsname{\href{S6_SERIES_GRIND_III.pdf\#nameddest=tag-01VA}{III.{B.5.12} (p.~24)} }
\expandafter\def\csname SRef@maps:eq:D.VC1\endcsname{\href{S6_SERIES_GRIND_III.pdf\#nameddest=tag-026E}{III.{D.VC1} (p.~187)} }
\expandafter\def\csname SRef@maps:eq:D.VC7\endcsname{\href{S6_SERIES_GRIND_III.pdf\#nameddest=tag-026L}{III.{D.VC7} (p.~188)} }
\expandafter\def\csname SRef@maps:eq:D.VC2\endcsname{\href{S6_SERIES_GRIND_III.pdf\#nameddest=tag-026G}{III.{D.VC2} (p.~188)} }
\expandafter\def\csname SRef@maps:eq:D.VC3\endcsname{\href{S6_SERIES_GRIND_III.pdf\#nameddest=tag-026H}{III.{D.VC3} (p.~188)} }
\expandafter\def\csname SRef@maps:eq:D.VC4\endcsname{\href{S6_SERIES_GRIND_III.pdf\#nameddest=tag-026I}{III.{D.VC4} (p.~188)} }
\expandafter\def\csname SRef@maps:eq:D.VC5\endcsname{\href{S6_SERIES_GRIND_III.pdf\#nameddest=tag-026J}{III.{D.VC5} (p.~188)} }
\expandafter\def\csname SRef@maps:eq:D.VC8\endcsname{\href{S6_SERIES_GRIND_III.pdf\#nameddest=tag-026M}{III.{D.VC8} (p.~189)} }
\expandafter\def\csname SRef@maps:eq:D.VC6\endcsname{\href{S6_SERIES_GRIND_III.pdf\#nameddest=tag-026K}{III.{D.VC6} (p.~189)} }
\expandafter\def\csname SRef@maps:eq:D.VC9\endcsname{\href{S6_SERIES_GRIND_III.pdf\#nameddest=tag-026N}{III.{D.VC9} (p.~190)} }
\expandafter\def\csname SRef@maps:eq:D.VC10\endcsname{\href{S6_SERIES_GRIND_III.pdf\#nameddest=tag-026F}{III.{D.VC10} (p.~190)} }
\expandafter\def\csname SRef@maps:eq:B.5.13\endcsname{\href{S6_SERIES_GRIND_III.pdf\#nameddest=tag-01VB}{III.{B.5.13} (p.~25)} }
\expandafter\def\csname SRef@maps:eq:D.IC1\endcsname{\href{S6_SERIES_GRIND_III.pdf\#nameddest=tag-01ZY}{III.{D.IC1} (p.~191)} }
\expandafter\def\csname SRef@maps:eq:D.IC2\endcsname{\href{S6_SERIES_GRIND_III.pdf\#nameddest=tag-01ZZ}{III.{D.IC2} (p.~191)} }
\expandafter\def\csname SRef@maps:eq:D.IC3\endcsname{\href{S6_SERIES_GRIND_III.pdf\#nameddest=tag-0200}{III.{D.IC3} (p.~191)} }
\expandafter\def\csname SRef@maps:eq:D.IC4\endcsname{\href{S6_SERIES_GRIND_III.pdf\#nameddest=tag-0201}{III.{D.IC4} (p.~191)} }
\expandafter\def\csname SRef@maps:eq:D.IC5\endcsname{\href{S6_SERIES_GRIND_III.pdf\#nameddest=tag-0202}{III.{D.IC5} (p.~191)} }
\expandafter\def\csname SRef@maps:eq:D.IC6\endcsname{\href{S6_SERIES_GRIND_III.pdf\#nameddest=tag-0203}{III.{D.IC6} (p.~191)} }
\expandafter\def\csname SRef@maps:eq:D.IC7\endcsname{\href{S6_SERIES_GRIND_III.pdf\#nameddest=tag-0204}{III.{D.IC7} (p.~191)} }
\expandafter\def\csname SRef@maps:eq:D.IC8\endcsname{\href{S6_SERIES_GRIND_III.pdf\#nameddest=tag-0205}{III.{D.IC8} (p.~192)} }
\expandafter\def\csname SRef@maps:eq:D.IC9\endcsname{\href{S6_SERIES_GRIND_III.pdf\#nameddest=tag-0206}{III.{D.IC9} (p.~192)} }
\expandafter\def\csname SRef@maps:eq:D.CI1\endcsname{\href{S6_SERIES_GRIND_III.pdf\#nameddest=tag-01YN}{III.{D.CI1} (p.~195)} }
\expandafter\def\csname SRef@maps:eq:B.5.14\endcsname{\href{S6_SERIES_GRIND_III.pdf\#nameddest=tag-01VC}{III.{B.5.14} (p.~25)} }
\expandafter\def\csname SRef@maps:eq:D.OS1\endcsname{\href{S6_SERIES_GRIND_III.pdf\#nameddest=tag-022O}{III.{D.OS1} (p.~196)} }
\expandafter\def\csname SRef@maps:eq:D.OS2\endcsname{\href{S6_SERIES_GRIND_III.pdf\#nameddest=tag-022P}{III.{D.OS2} (p.~196)} }
\expandafter\def\csname SRef@maps:eq:D.OS3\endcsname{\href{S6_SERIES_GRIND_III.pdf\#nameddest=tag-022Q}{III.{D.OS3} (p.~196)} }
\expandafter\def\csname SRef@maps:eq:D.OS4\endcsname{\href{S6_SERIES_GRIND_III.pdf\#nameddest=tag-022R}{III.{D.OS4} (p.~196)} }
\expandafter\def\csname SRef@maps:eq:D.PC1\endcsname{\href{S6_SERIES_GRIND_III.pdf\#nameddest=tag-0233}{III.{D.PC1} (p.~196)} }
\expandafter\def\csname SRef@maps:eq:D.PC2\endcsname{\href{S6_SERIES_GRIND_III.pdf\#nameddest=tag-0239}{III.{D.PC2} (p.~197)} }
\expandafter\def\csname SRef@maps:eq:D.PC3\endcsname{\href{S6_SERIES_GRIND_III.pdf\#nameddest=tag-023A}{III.{D.PC3} (p.~197)} }
\expandafter\def\csname SRef@maps:eq:D.PC4\endcsname{\href{S6_SERIES_GRIND_III.pdf\#nameddest=tag-023B}{III.{D.PC4} (p.~197)} }
\expandafter\def\csname SRef@maps:eq:D.PC5\endcsname{\href{S6_SERIES_GRIND_III.pdf\#nameddest=tag-023C}{III.{D.PC5} (p.~197)} }
\expandafter\def\csname SRef@maps:eq:D.PC6\endcsname{\href{S6_SERIES_GRIND_III.pdf\#nameddest=tag-023D}{III.{D.PC6} (p.~197)} }
\expandafter\def\csname SRef@maps:eq:B.5.15\endcsname{\href{S6_SERIES_GRIND_III.pdf\#nameddest=tag-01VD}{III.{B.5.15} (p.~25)} }
\expandafter\def\csname SRef@maps:eq:D.PC7\endcsname{\href{S6_SERIES_GRIND_III.pdf\#nameddest=tag-023E}{III.{D.PC7} (p.~197)} }
\expandafter\def\csname SRef@maps:eq:D.PC8\endcsname{\href{S6_SERIES_GRIND_III.pdf\#nameddest=tag-023F}{III.{D.PC8} (p.~197)} }
\expandafter\def\csname SRef@maps:eq:D.PC9\endcsname{\href{S6_SERIES_GRIND_III.pdf\#nameddest=tag-023G}{III.{D.PC9} (p.~198)} }
\expandafter\def\csname SRef@maps:eq:D.PC10\endcsname{\href{S6_SERIES_GRIND_III.pdf\#nameddest=tag-0234}{III.{D.PC10} (p.~198)} }
\expandafter\def\csname SRef@maps:eq:D.PC11\endcsname{\href{S6_SERIES_GRIND_III.pdf\#nameddest=tag-0235}{III.{D.PC11} (p.~198)} }
\expandafter\def\csname SRef@maps:eq:D.PC12\endcsname{\href{S6_SERIES_GRIND_III.pdf\#nameddest=tag-0236}{III.{D.PC12} (p.~198)} }
\expandafter\def\csname SRef@maps:eq:D.PC13\endcsname{\href{S6_SERIES_GRIND_III.pdf\#nameddest=tag-0237}{III.{D.PC13} (p.~198)} }
\expandafter\def\csname SRef@maps:eq:D.PC14\endcsname{\href{S6_SERIES_GRIND_III.pdf\#nameddest=tag-0238}{III.{D.PC14} (p.~198)} }
\expandafter\def\csname SRef@maps:eq:D.OS5\endcsname{\href{S6_SERIES_GRIND_III.pdf\#nameddest=tag-022S}{III.{D.OS5} (p.~199)} }
\expandafter\def\csname SRef@maps:eq:B.5.16\endcsname{\href{S6_SERIES_GRIND_III.pdf\#nameddest=tag-01VE}{III.{B.5.16} (p.~25)} }
\expandafter\def\csname SRef@maps:eq:D.OS6\endcsname{\href{S6_SERIES_GRIND_III.pdf\#nameddest=tag-022T}{III.{D.OS6} (p.~199)} }
\expandafter\def\csname SRef@maps:eq:D.OS7\endcsname{\href{S6_SERIES_GRIND_III.pdf\#nameddest=tag-022U}{III.{D.OS7} (p.~199)} }
\expandafter\def\csname SRef@maps:eq:D.OS8\endcsname{\href{S6_SERIES_GRIND_III.pdf\#nameddest=tag-022V}{III.{D.OS8} (p.~200)} }
\expandafter\def\csname SRef@maps:eq:D.OF1\endcsname{\href{S6_SERIES_GRIND_III.pdf\#nameddest=tag-0217}{III.{D.OF1} (p.~200)} }
\expandafter\def\csname SRef@maps:eq:D.OF2\endcsname{\href{S6_SERIES_GRIND_III.pdf\#nameddest=tag-021I}{III.{D.OF2} (p.~200)} }
\expandafter\def\csname SRef@maps:eq:D.OF3\endcsname{\href{S6_SERIES_GRIND_III.pdf\#nameddest=tag-021J}{III.{D.OF3} (p.~201)} }
\expandafter\def\csname SRef@maps:eq:D.OF4\endcsname{\href{S6_SERIES_GRIND_III.pdf\#nameddest=tag-021K}{III.{D.OF4} (p.~201)} }
\expandafter\def\csname SRef@maps:eq:D.OF5\endcsname{\href{S6_SERIES_GRIND_III.pdf\#nameddest=tag-021L}{III.{D.OF5} (p.~201)} }
\expandafter\def\csname SRef@maps:eq:B.5.17\endcsname{\href{S6_SERIES_GRIND_III.pdf\#nameddest=tag-01VF}{III.{B.5.17} (p.~25)} }
\expandafter\def\csname SRef@maps:eq:D.OF6\endcsname{\href{S6_SERIES_GRIND_III.pdf\#nameddest=tag-021M}{III.{D.OF6} (p.~201)} }
\expandafter\def\csname SRef@maps:eq:D.OF7\endcsname{\href{S6_SERIES_GRIND_III.pdf\#nameddest=tag-021N}{III.{D.OF7} (p.~201)} }
\expandafter\def\csname SRef@maps:eq:D.OF8\endcsname{\href{S6_SERIES_GRIND_III.pdf\#nameddest=tag-021O}{III.{D.OF8} (p.~201)} }
\expandafter\def\csname SRef@maps:eq:D.OF9\endcsname{\href{S6_SERIES_GRIND_III.pdf\#nameddest=tag-021P}{III.{D.OF9} (p.~202)} }
\expandafter\def\csname SRef@maps:eq:D.OF10\endcsname{\href{S6_SERIES_GRIND_III.pdf\#nameddest=tag-0218}{III.{D.OF10} (p.~202)} }
\expandafter\def\csname SRef@maps:eq:D.OF11\endcsname{\href{S6_SERIES_GRIND_III.pdf\#nameddest=tag-0219}{III.{D.OF11} (p.~202)} }
\expandafter\def\csname SRef@maps:eq:D.OF12\endcsname{\href{S6_SERIES_GRIND_III.pdf\#nameddest=tag-021A}{III.{D.OF12} (p.~202)} }
\expandafter\def\csname SRef@maps:eq:D.OF13\endcsname{\href{S6_SERIES_GRIND_III.pdf\#nameddest=tag-021B}{III.{D.OF13} (p.~202)} }
\expandafter\def\csname SRef@maps:eq:D.OF14\endcsname{\href{S6_SERIES_GRIND_III.pdf\#nameddest=tag-021C}{III.{D.OF14} (p.~202)} }
\expandafter\def\csname SRef@maps:eq:B.5.18\endcsname{\href{S6_SERIES_GRIND_III.pdf\#nameddest=tag-01VG}{III.{B.5.18} (p.~26)} }
\expandafter\def\csname SRef@maps:eq:D.OF15\endcsname{\href{S6_SERIES_GRIND_III.pdf\#nameddest=tag-021D}{III.{D.OF15} (p.~203)} }
\expandafter\def\csname SRef@maps:eq:D.OF16\endcsname{\href{S6_SERIES_GRIND_III.pdf\#nameddest=tag-021E}{III.{D.OF16} (p.~203)} }
\expandafter\def\csname SRef@maps:eq:D.OF17\endcsname{\href{S6_SERIES_GRIND_III.pdf\#nameddest=tag-021F}{III.{D.OF17} (p.~203)} }
\expandafter\def\csname SRef@maps:eq:D.OF18\endcsname{\href{S6_SERIES_GRIND_III.pdf\#nameddest=tag-021G}{III.{D.OF18} (p.~203)} }
\expandafter\def\csname SRef@maps:eq:D.OL1\endcsname{\href{S6_SERIES_GRIND_III.pdf\#nameddest=tag-021Q}{III.{D.OL1} (p.~204)} }
\expandafter\def\csname SRef@maps:eq:D.OL2\endcsname{\href{S6_SERIES_GRIND_III.pdf\#nameddest=tag-021U}{III.{D.OL2} (p.~204)} }
\expandafter\def\csname SRef@maps:eq:D.OL3\endcsname{\href{S6_SERIES_GRIND_III.pdf\#nameddest=tag-021V}{III.{D.OL3} (p.~204)} }
\expandafter\def\csname SRef@maps:eq:D.OL4\endcsname{\href{S6_SERIES_GRIND_III.pdf\#nameddest=tag-021W}{III.{D.OL4} (p.~204)} }
\expandafter\def\csname SRef@maps:eq:D.OL5\endcsname{\href{S6_SERIES_GRIND_III.pdf\#nameddest=tag-021X}{III.{D.OL5} (p.~204)} }
\expandafter\def\csname SRef@maps:eq:C.1.1\endcsname{\href{S6_SERIES_GRIND_III.pdf\#nameddest=tag-01XE}{III.{C.1.1} (p.~27)} }
\expandafter\def\csname SRef@maps:eq:D.OL6\endcsname{\href{S6_SERIES_GRIND_III.pdf\#nameddest=tag-021Y}{III.{D.OL6} (p.~204)} }
\expandafter\def\csname SRef@maps:eq:D.OL7\endcsname{\href{S6_SERIES_GRIND_III.pdf\#nameddest=tag-021Z}{III.{D.OL7} (p.~204)} }
\expandafter\def\csname SRef@maps:eq:D.OL8\endcsname{\href{S6_SERIES_GRIND_III.pdf\#nameddest=tag-0220}{III.{D.OL8} (p.~205)} }
\expandafter\def\csname SRef@maps:eq:D.OL9\endcsname{\href{S6_SERIES_GRIND_III.pdf\#nameddest=tag-0221}{III.{D.OL9} (p.~205)} }
\expandafter\def\csname SRef@maps:eq:D.OL10\endcsname{\href{S6_SERIES_GRIND_III.pdf\#nameddest=tag-021R}{III.{D.OL10} (p.~205)} }
\expandafter\def\csname SRef@maps:eq:D.OL11\endcsname{\href{S6_SERIES_GRIND_III.pdf\#nameddest=tag-021S}{III.{D.OL11} (p.~205)} }
\expandafter\def\csname SRef@maps:eq:D.OL12\endcsname{\href{S6_SERIES_GRIND_III.pdf\#nameddest=tag-021T}{III.{D.OL12} (p.~205)} }
\expandafter\def\csname SRef@maps:eq:D.OF19\endcsname{\href{S6_SERIES_GRIND_III.pdf\#nameddest=tag-021H}{III.{D.OF19} (p.~206)} }
\expandafter\def\csname SRef@maps:eq:C.1.2\endcsname{\href{S6_SERIES_GRIND_III.pdf\#nameddest=tag-01XF}{III.{C.1.2} (p.~28)} }
\expandafter\def\csname SRef@maps:eq:C.1.3\endcsname{\href{S6_SERIES_GRIND_III.pdf\#nameddest=tag-01XG}{III.{C.1.3} (p.~30)} }
\expandafter\def\csname SRef@maps:eq:E.1\endcsname{\href{S6_SERIES_GRIND_V.pdf\#nameddest=tag-029F}{V.{E.1} (p.~14)} }
\expandafter\def\csname SRef@maps:eq:E.2\endcsname{\href{S6_SERIES_GRIND_V.pdf\#nameddest=tag-029O}{V.{E.2} (p.~15)} }
\expandafter\def\csname SRef@maps:eq:C.1.4\endcsname{\href{S6_SERIES_GRIND_III.pdf\#nameddest=tag-01XH}{III.{C.1.4} (p.~30)} }
\expandafter\def\csname SRef@maps:eq:E.3\endcsname{\href{S6_SERIES_GRIND_V.pdf\#nameddest=tag-029Q}{V.{E.3} (p.~15)} }
\expandafter\def\csname SRef@maps:eq:E.4\endcsname{\href{S6_SERIES_GRIND_V.pdf\#nameddest=tag-029R}{V.{E.4} (p.~15)} }
\expandafter\def\csname SRef@maps:eq:E.5\endcsname{\href{S6_SERIES_GRIND_V.pdf\#nameddest=tag-029S}{V.{E.5} (p.~15)} }
\expandafter\def\csname SRef@maps:eq:E.6\endcsname{\href{S6_SERIES_GRIND_V.pdf\#nameddest=tag-029T}{V.{E.6} (p.~16)} }
\expandafter\def\csname SRef@maps:eq:E.7\endcsname{\href{S6_SERIES_GRIND_V.pdf\#nameddest=tag-029U}{V.{E.7} (p.~16)} }
\expandafter\def\csname SRef@maps:eq:E.8\endcsname{\href{S6_SERIES_GRIND_V.pdf\#nameddest=tag-029V}{V.{E.8} (p.~18)} }
\expandafter\def\csname SRef@maps:eq:E.9\endcsname{\href{S6_SERIES_GRIND_V.pdf\#nameddest=tag-029W}{V.{E.9} (p.~18)} }
\expandafter\def\csname SRef@maps:eq:E.10\endcsname{\href{S6_SERIES_GRIND_V.pdf\#nameddest=tag-029G}{V.{E.10} (p.~18)} }
\expandafter\def\csname SRef@maps:eq:E.11\endcsname{\href{S6_SERIES_GRIND_V.pdf\#nameddest=tag-029H}{V.{E.11} (p.~18)} }
\expandafter\def\csname SRef@maps:eq:E.I1\endcsname{\href{S6_SERIES_GRIND_V.pdf\#nameddest=tag-02B5}{V.{E.I1} (p.~19)} }
\expandafter\def\csname SRef@maps:eq:D.C1\endcsname{\href{S6_SERIES_GRIND_III.pdf\#nameddest=tag-01Y6}{III.{D.C1} (p.~59)} }
\expandafter\def\csname SRef@maps:eq:E.I2\endcsname{\href{S6_SERIES_GRIND_V.pdf\#nameddest=tag-02B6}{V.{E.I2} (p.~20)} }
\expandafter\def\csname SRef@maps:eq:E.I3\endcsname{\href{S6_SERIES_GRIND_V.pdf\#nameddest=tag-02B7}{V.{E.I3} (p.~20)} }
\expandafter\def\csname SRef@maps:eq:E.I4\endcsname{\href{S6_SERIES_GRIND_V.pdf\#nameddest=tag-02B8}{V.{E.I4} (p.~20)} }
\expandafter\def\csname SRef@maps:eq:E.I5\endcsname{\href{S6_SERIES_GRIND_V.pdf\#nameddest=tag-02B9}{V.{E.I5} (p.~21)} }
\expandafter\def\csname SRef@maps:eq:E.12\endcsname{\href{S6_SERIES_GRIND_V.pdf\#nameddest=tag-029I}{V.{E.12} (p.~21)} }
\expandafter\def\csname SRef@maps:eq:E.13\endcsname{\href{S6_SERIES_GRIND_V.pdf\#nameddest=tag-029J}{V.{E.13} (p.~21)} }
\expandafter\def\csname SRef@maps:eq:E.14\endcsname{\href{S6_SERIES_GRIND_V.pdf\#nameddest=tag-029K}{V.{E.14} (p.~22)} }
\expandafter\def\csname SRef@maps:eq:E.16\endcsname{\href{S6_SERIES_GRIND_V.pdf\#nameddest=tag-029M}{V.{E.16} (p.~23)} }
\expandafter\def\csname SRef@maps:eq:E.17\endcsname{\href{S6_SERIES_GRIND_V.pdf\#nameddest=tag-029N}{V.{E.17} (p.~23)} }
\expandafter\def\csname SRef@maps:eq:E.20\endcsname{\href{S6_SERIES_GRIND_V.pdf\#nameddest=tag-029P}{V.{E.20} (p.~25)} }
\expandafter\def\csname SRef@maps:eq:D.R1\endcsname{\href{S6_SERIES_GRIND_III.pdf\#nameddest=tag-023H}{III.{D.R1} (p.~64)} }
\expandafter\def\csname SRef@maps:eq:E.MV1\endcsname{\href{S6_SERIES_GRIND_V.pdf\#nameddest=tag-02BA}{V.{E.MV1} (p.~26)} }
\expandafter\def\csname SRef@maps:eq:E.MV2\endcsname{\href{S6_SERIES_GRIND_V.pdf\#nameddest=tag-02BB}{V.{E.MV2} (p.~27)} }
\expandafter\def\csname SRef@maps:eq:E.MV3\endcsname{\href{S6_SERIES_GRIND_V.pdf\#nameddest=tag-02BC}{V.{E.MV3} (p.~27)} }
\expandafter\def\csname SRef@maps:eq:E.MV4\endcsname{\href{S6_SERIES_GRIND_V.pdf\#nameddest=tag-02BD}{V.{E.MV4} (p.~27)} }
\expandafter\def\csname SRef@maps:eq:E.MV5\endcsname{\href{S6_SERIES_GRIND_V.pdf\#nameddest=tag-02BE}{V.{E.MV5} (p.~27)} }
\expandafter\def\csname SRef@maps:eq:E.MV6\endcsname{\href{S6_SERIES_GRIND_V.pdf\#nameddest=tag-02BF}{V.{E.MV6} (p.~28)} }
\expandafter\def\csname SRef@maps:eq:E.MV7\endcsname{\href{S6_SERIES_GRIND_V.pdf\#nameddest=tag-02BG}{V.{E.MV7} (p.~28)} }
\expandafter\def\csname SRef@maps:eq:E.MV8\endcsname{\href{S6_SERIES_GRIND_V.pdf\#nameddest=tag-02BH}{V.{E.MV8} (p.~28)} }
\expandafter\def\csname SRef@maps:eq:E.MV9\endcsname{\href{S6_SERIES_GRIND_V.pdf\#nameddest=tag-02BI}{V.{E.MV9} (p.~29)} }
\expandafter\def\csname SRef@maps:eq:E.B1\endcsname{\href{S6_SERIES_GRIND_V.pdf\#nameddest=tag-029X}{V.{E.B1} (p.~30)} }
\expandafter\def\csname SRef@maps:eq:D.R2\endcsname{\href{S6_SERIES_GRIND_III.pdf\#nameddest=tag-023I}{III.{D.R2} (p.~64)} }
\expandafter\def\csname SRef@maps:eq:E.B2\endcsname{\href{S6_SERIES_GRIND_V.pdf\#nameddest=tag-029Y}{V.{E.B2} (p.~30)} }
\expandafter\def\csname SRef@maps:eq:E.B3\endcsname{\href{S6_SERIES_GRIND_V.pdf\#nameddest=tag-029Z}{V.{E.B3} (p.~31)} }
\expandafter\def\csname SRef@maps:eq:E.B4\endcsname{\href{S6_SERIES_GRIND_V.pdf\#nameddest=tag-02A0}{V.{E.B4} (p.~31)} }
\expandafter\def\csname SRef@maps:eq:E.B5\endcsname{\href{S6_SERIES_GRIND_V.pdf\#nameddest=tag-02A1}{V.{E.B5} (p.~31)} }
\expandafter\def\csname SRef@maps:eq:E.B6\endcsname{\href{S6_SERIES_GRIND_V.pdf\#nameddest=tag-02A2}{V.{E.B6} (p.~32)} }
\expandafter\def\csname SRef@maps:eq:E.B7\endcsname{\href{S6_SERIES_GRIND_V.pdf\#nameddest=tag-02A3}{V.{E.B7} (p.~32)} }
\expandafter\def\csname SRef@maps:eq:E.S1\endcsname{\href{S6_SERIES_GRIND_V.pdf\#nameddest=tag-02CC}{V.{E.S1} (p.~33)} }
\expandafter\def\csname SRef@maps:eq:E.S2\endcsname{\href{S6_SERIES_GRIND_V.pdf\#nameddest=tag-02CN}{V.{E.S2} (p.~33)} }
\expandafter\def\csname SRef@maps:eq:E.S3\endcsname{\href{S6_SERIES_GRIND_V.pdf\#nameddest=tag-02CY}{V.{E.S3} (p.~34)} }
\expandafter\def\csname SRef@maps:eq:E.S4\endcsname{\href{S6_SERIES_GRIND_V.pdf\#nameddest=tag-02D9}{V.{E.S4} (p.~34)} }
\expandafter\def\csname SRef@maps:eq:D.R3\endcsname{\href{S6_SERIES_GRIND_III.pdf\#nameddest=tag-023J}{III.{D.R3} (p.~64)} }
\expandafter\def\csname SRef@maps:eq:E.S5\endcsname{\href{S6_SERIES_GRIND_V.pdf\#nameddest=tag-02DK}{V.{E.S5} (p.~34)} }
\expandafter\def\csname SRef@maps:eq:E.S12\endcsname{\href{S6_SERIES_GRIND_V.pdf\#nameddest=tag-02CF}{V.{E.S12} (p.~35)} }
\expandafter\def\csname SRef@maps:eq:E.S6\endcsname{\href{S6_SERIES_GRIND_V.pdf\#nameddest=tag-02DR}{V.{E.S6} (p.~35)} }
\expandafter\def\csname SRef@maps:eq:E.S7\endcsname{\href{S6_SERIES_GRIND_V.pdf\#nameddest=tag-02DS}{V.{E.S7} (p.~35)} }
\expandafter\def\csname SRef@maps:eq:E.S8\endcsname{\href{S6_SERIES_GRIND_V.pdf\#nameddest=tag-02DT}{V.{E.S8} (p.~36)} }
\expandafter\def\csname SRef@maps:eq:E.S9\endcsname{\href{S6_SERIES_GRIND_V.pdf\#nameddest=tag-02DU}{V.{E.S9} (p.~36)} }
\expandafter\def\csname SRef@maps:eq:E.S10\endcsname{\href{S6_SERIES_GRIND_V.pdf\#nameddest=tag-02CD}{V.{E.S10} (p.~36)} }
\expandafter\def\csname SRef@maps:eq:E.S11\endcsname{\href{S6_SERIES_GRIND_V.pdf\#nameddest=tag-02CE}{V.{E.S11} (p.~36)} }
\expandafter\def\csname SRef@maps:eq:E.S13\endcsname{\href{S6_SERIES_GRIND_V.pdf\#nameddest=tag-02CG}{V.{E.S13} (p.~37)} }
\expandafter\def\csname SRef@maps:eq:E.S14\endcsname{\href{S6_SERIES_GRIND_V.pdf\#nameddest=tag-02CH}{V.{E.S14} (p.~37)} }
\expandafter\def\csname SRef@maps:eq:D.R4\endcsname{\href{S6_SERIES_GRIND_III.pdf\#nameddest=tag-023K}{III.{D.R4} (p.~64)} }
\expandafter\def\csname SRef@maps:eq:E.S15\endcsname{\href{S6_SERIES_GRIND_V.pdf\#nameddest=tag-02CI}{V.{E.S15} (p.~37)} }
\expandafter\def\csname SRef@maps:eq:E.S16\endcsname{\href{S6_SERIES_GRIND_V.pdf\#nameddest=tag-02CJ}{V.{E.S16} (p.~37)} }
\expandafter\def\csname SRef@maps:eq:E.S20\endcsname{\href{S6_SERIES_GRIND_V.pdf\#nameddest=tag-02CO}{V.{E.S20} (p.~37)} }
\expandafter\def\csname SRef@maps:eq:E.S17\endcsname{\href{S6_SERIES_GRIND_V.pdf\#nameddest=tag-02CK}{V.{E.S17} (p.~37)} }
\expandafter\def\csname SRef@maps:eq:E.S18\endcsname{\href{S6_SERIES_GRIND_V.pdf\#nameddest=tag-02CL}{V.{E.S18} (p.~37)} }
\expandafter\def\csname SRef@maps:eq:E.S19\endcsname{\href{S6_SERIES_GRIND_V.pdf\#nameddest=tag-02CM}{V.{E.S19} (p.~38)} }
\expandafter\def\csname SRef@maps:eq:E.S21\endcsname{\href{S6_SERIES_GRIND_V.pdf\#nameddest=tag-02CP}{V.{E.S21} (p.~39)} }
\expandafter\def\csname SRef@maps:eq:E.S22\endcsname{\href{S6_SERIES_GRIND_V.pdf\#nameddest=tag-02CQ}{V.{E.S22} (p.~40)} }
\expandafter\def\csname SRef@maps:eq:E.S23\endcsname{\href{S6_SERIES_GRIND_V.pdf\#nameddest=tag-02CR}{V.{E.S23} (p.~40)} }
\expandafter\def\csname SRef@maps:eq:E.S24\endcsname{\href{S6_SERIES_GRIND_V.pdf\#nameddest=tag-02CS}{V.{E.S24} (p.~40)} }
\expandafter\def\csname SRef@maps:eq:D.S1\endcsname{\href{S6_SERIES_GRIND_III.pdf\#nameddest=tag-023L}{III.{D.S1} (p.~65)} }
\expandafter\def\csname SRef@maps:eq:E.S25\endcsname{\href{S6_SERIES_GRIND_V.pdf\#nameddest=tag-02CT}{V.{E.S25} (p.~40)} }
\expandafter\def\csname SRef@maps:eq:E.S26\endcsname{\href{S6_SERIES_GRIND_V.pdf\#nameddest=tag-02CU}{V.{E.S26} (p.~40)} }
\expandafter\def\csname SRef@maps:eq:E.S27\endcsname{\href{S6_SERIES_GRIND_V.pdf\#nameddest=tag-02CV}{V.{E.S27} (p.~41)} }
\expandafter\def\csname SRef@maps:eq:E.S28\endcsname{\href{S6_SERIES_GRIND_V.pdf\#nameddest=tag-02CW}{V.{E.S28} (p.~41)} }
\expandafter\def\csname SRef@maps:eq:E.S29\endcsname{\href{S6_SERIES_GRIND_V.pdf\#nameddest=tag-02CX}{V.{E.S29} (p.~42)} }
\expandafter\def\csname SRef@maps:eq:E.S30\endcsname{\href{S6_SERIES_GRIND_V.pdf\#nameddest=tag-02CZ}{V.{E.S30} (p.~42)} }
\expandafter\def\csname SRef@maps:eq:E.S31\endcsname{\href{S6_SERIES_GRIND_V.pdf\#nameddest=tag-02D0}{V.{E.S31} (p.~42)} }
\expandafter\def\csname SRef@maps:eq:E.S32\endcsname{\href{S6_SERIES_GRIND_V.pdf\#nameddest=tag-02D1}{V.{E.S32} (p.~42)} }
\expandafter\def\csname SRef@maps:eq:E.S33\endcsname{\href{S6_SERIES_GRIND_V.pdf\#nameddest=tag-02D2}{V.{E.S33} (p.~42)} }
\expandafter\def\csname SRef@maps:eq:E.S34\endcsname{\href{S6_SERIES_GRIND_V.pdf\#nameddest=tag-02D3}{V.{E.S34} (p.~43)} }
\expandafter\def\csname SRef@maps:eq:D.S2\endcsname{\href{S6_SERIES_GRIND_III.pdf\#nameddest=tag-023R}{III.{D.S2} (p.~65)} }
\expandafter\def\csname SRef@maps:eq:E.S35\endcsname{\href{S6_SERIES_GRIND_V.pdf\#nameddest=tag-02D4}{V.{E.S35} (p.~43)} }
\expandafter\def\csname SRef@maps:eq:E.S36\endcsname{\href{S6_SERIES_GRIND_V.pdf\#nameddest=tag-02D5}{V.{E.S36} (p.~43)} }
\expandafter\def\csname SRef@maps:eq:E.S37\endcsname{\href{S6_SERIES_GRIND_V.pdf\#nameddest=tag-02D6}{V.{E.S37} (p.~43)} }
\expandafter\def\csname SRef@maps:eq:E.S38\endcsname{\href{S6_SERIES_GRIND_V.pdf\#nameddest=tag-02D7}{V.{E.S38} (p.~43)} }
\expandafter\def\csname SRef@maps:eq:E.S39\endcsname{\href{S6_SERIES_GRIND_V.pdf\#nameddest=tag-02D8}{V.{E.S39} (p.~44)} }
\expandafter\def\csname SRef@maps:eq:E.S40\endcsname{\href{S6_SERIES_GRIND_V.pdf\#nameddest=tag-02DA}{V.{E.S40} (p.~44)} }
\expandafter\def\csname SRef@maps:eq:E.S41\endcsname{\href{S6_SERIES_GRIND_V.pdf\#nameddest=tag-02DB}{V.{E.S41} (p.~44)} }
\expandafter\def\csname SRef@maps:eq:E.S42\endcsname{\href{S6_SERIES_GRIND_V.pdf\#nameddest=tag-02DC}{V.{E.S42} (p.~44)} }
\expandafter\def\csname SRef@maps:eq:E.S43\endcsname{\href{S6_SERIES_GRIND_V.pdf\#nameddest=tag-02DD}{V.{E.S43} (p.~45)} }
\expandafter\def\csname SRef@maps:eq:E.S44\endcsname{\href{S6_SERIES_GRIND_V.pdf\#nameddest=tag-02DE}{V.{E.S44} (p.~45)} }
\expandafter\def\csname SRef@maps:eq:D.S3\endcsname{\href{S6_SERIES_GRIND_III.pdf\#nameddest=tag-023S}{III.{D.S3} (p.~65)} }
\expandafter\def\csname SRef@maps:eq:E.S45\endcsname{\href{S6_SERIES_GRIND_V.pdf\#nameddest=tag-02DF}{V.{E.S45} (p.~45)} }
\expandafter\def\csname SRef@maps:eq:E.S46\endcsname{\href{S6_SERIES_GRIND_V.pdf\#nameddest=tag-02DG}{V.{E.S46} (p.~45)} }
\expandafter\def\csname SRef@maps:eq:E.S47\endcsname{\href{S6_SERIES_GRIND_V.pdf\#nameddest=tag-02DH}{V.{E.S47} (p.~45)} }
\expandafter\def\csname SRef@maps:eq:E.S48\endcsname{\href{S6_SERIES_GRIND_V.pdf\#nameddest=tag-02DI}{V.{E.S48} (p.~46)} }
\expandafter\def\csname SRef@maps:eq:E.S49\endcsname{\href{S6_SERIES_GRIND_V.pdf\#nameddest=tag-02DJ}{V.{E.S49} (p.~46)} }
\expandafter\def\csname SRef@maps:eq:E.S50\endcsname{\href{S6_SERIES_GRIND_V.pdf\#nameddest=tag-02DL}{V.{E.S50} (p.~46)} }
\expandafter\def\csname SRef@maps:eq:E.S51\endcsname{\href{S6_SERIES_GRIND_V.pdf\#nameddest=tag-02DM}{V.{E.S51} (p.~46)} }
\expandafter\def\csname SRef@maps:eq:E.S52\endcsname{\href{S6_SERIES_GRIND_V.pdf\#nameddest=tag-02DN}{V.{E.S52} (p.~46)} }
\expandafter\def\csname SRef@maps:eq:E.S53\endcsname{\href{S6_SERIES_GRIND_V.pdf\#nameddest=tag-02DO}{V.{E.S53} (p.~46)} }
\expandafter\def\csname SRef@maps:eq:E.S54\endcsname{\href{S6_SERIES_GRIND_V.pdf\#nameddest=tag-02DP}{V.{E.S54} (p.~46)} }
\expandafter\def\csname SRef@maps:eq:D.S4\endcsname{\href{S6_SERIES_GRIND_III.pdf\#nameddest=tag-023T}{III.{D.S4} (p.~65)} }
\expandafter\def\csname SRef@maps:eq:E.S55\endcsname{\href{S6_SERIES_GRIND_V.pdf\#nameddest=tag-02DQ}{V.{E.S55} (p.~47)} }
\expandafter\def\csname SRef@maps:eq:E.C1\endcsname{\href{S6_SERIES_GRIND_V.pdf\#nameddest=tag-02A4}{V.{E.C1} (p.~48)} }
\expandafter\def\csname SRef@maps:eq:E.C2\endcsname{\href{S6_SERIES_GRIND_V.pdf\#nameddest=tag-02AG}{V.{E.C2} (p.~48)} }
\expandafter\def\csname SRef@maps:eq:E.C3\endcsname{\href{S6_SERIES_GRIND_V.pdf\#nameddest=tag-02AJ}{V.{E.C3} (p.~48)} }
\expandafter\def\csname SRef@maps:eq:E.C4\endcsname{\href{S6_SERIES_GRIND_V.pdf\#nameddest=tag-02AK}{V.{E.C4} (p.~48)} }
\expandafter\def\csname SRef@maps:eq:E.C19a\endcsname{\href{S6_SERIES_GRIND_V.pdf\#nameddest=tag-02AE}{V.{E.C19a} (p.~48)} }
\expandafter\def\csname SRef@maps:eq:E.C19b\endcsname{\href{S6_SERIES_GRIND_V.pdf\#nameddest=tag-02AF}{V.{E.C19b} (p.~48)} }
\expandafter\def\csname SRef@maps:eq:E.C5\endcsname{\href{S6_SERIES_GRIND_V.pdf\#nameddest=tag-02AL}{V.{E.C5} (p.~49)} }
\expandafter\def\csname SRef@maps:eq:E.C6\endcsname{\href{S6_SERIES_GRIND_V.pdf\#nameddest=tag-02AM}{V.{E.C6} (p.~49)} }
\expandafter\def\csname SRef@maps:eq:E.C7\endcsname{\href{S6_SERIES_GRIND_V.pdf\#nameddest=tag-02AN}{V.{E.C7} (p.~49)} }
\expandafter\def\csname SRef@maps:eq:D.S5\endcsname{\href{S6_SERIES_GRIND_III.pdf\#nameddest=tag-023U}{III.{D.S5} (p.~66)} }
\expandafter\def\csname SRef@maps:eq:E.C8\endcsname{\href{S6_SERIES_GRIND_V.pdf\#nameddest=tag-02AO}{V.{E.C8} (p.~49)} }
\expandafter\def\csname SRef@maps:eq:E.C9\endcsname{\href{S6_SERIES_GRIND_V.pdf\#nameddest=tag-02AP}{V.{E.C9} (p.~49)} }
\expandafter\def\csname SRef@maps:eq:E.C10\endcsname{\href{S6_SERIES_GRIND_V.pdf\#nameddest=tag-02A5}{V.{E.C10} (p.~49)} }
\expandafter\def\csname SRef@maps:eq:E.C11\endcsname{\href{S6_SERIES_GRIND_V.pdf\#nameddest=tag-02A6}{V.{E.C11} (p.~49)} }
\expandafter\def\csname SRef@maps:eq:E.C12\endcsname{\href{S6_SERIES_GRIND_V.pdf\#nameddest=tag-02A7}{V.{E.C12} (p.~50)} }
\expandafter\def\csname SRef@maps:eq:E.C20\endcsname{\href{S6_SERIES_GRIND_V.pdf\#nameddest=tag-02AH}{V.{E.C20} (p.~50)} }
\expandafter\def\csname SRef@maps:eq:E.C13\endcsname{\href{S6_SERIES_GRIND_V.pdf\#nameddest=tag-02A8}{V.{E.C13} (p.~51)} }
\expandafter\def\csname SRef@maps:eq:E.C14\endcsname{\href{S6_SERIES_GRIND_V.pdf\#nameddest=tag-02A9}{V.{E.C14} (p.~51)} }
\expandafter\def\csname SRef@maps:eq:E.C15\endcsname{\href{S6_SERIES_GRIND_V.pdf\#nameddest=tag-02AA}{V.{E.C15} (p.~51)} }
\expandafter\def\csname SRef@maps:eq:E.C16\endcsname{\href{S6_SERIES_GRIND_V.pdf\#nameddest=tag-02AB}{V.{E.C16} (p.~51)} }
\expandafter\def\csname SRef@maps:eq:D.S6\endcsname{\href{S6_SERIES_GRIND_III.pdf\#nameddest=tag-023V}{III.{D.S6} (p.~66)} }
\expandafter\def\csname SRef@maps:eq:E.C17\endcsname{\href{S6_SERIES_GRIND_V.pdf\#nameddest=tag-02AC}{V.{E.C17} (p.~51)} }
\expandafter\def\csname SRef@maps:eq:E.C18\endcsname{\href{S6_SERIES_GRIND_V.pdf\#nameddest=tag-02AD}{V.{E.C18} (p.~52)} }
\expandafter\def\csname SRef@maps:eq:E.C21\endcsname{\href{S6_SERIES_GRIND_V.pdf\#nameddest=tag-02AI}{V.{E.C21} (p.~52)} }
\expandafter\def\csname SRef@maps:eq:E.P1\endcsname{\href{S6_SERIES_GRIND_V.pdf\#nameddest=tag-02BJ}{V.{E.P1} (p.~52)} }
\expandafter\def\csname SRef@maps:eq:E.P2\endcsname{\href{S6_SERIES_GRIND_V.pdf\#nameddest=tag-02BS}{V.{E.P2} (p.~53)} }
\expandafter\def\csname SRef@maps:eq:E.P3\endcsname{\href{S6_SERIES_GRIND_V.pdf\#nameddest=tag-02BT}{V.{E.P3} (p.~53)} }
\expandafter\def\csname SRef@maps:eq:E.P4\endcsname{\href{S6_SERIES_GRIND_V.pdf\#nameddest=tag-02BU}{V.{E.P4} (p.~53)} }
\expandafter\def\csname SRef@maps:eq:E.P5\endcsname{\href{S6_SERIES_GRIND_V.pdf\#nameddest=tag-02BV}{V.{E.P5} (p.~53)} }
\expandafter\def\csname SRef@maps:eq:E.P6\endcsname{\href{S6_SERIES_GRIND_V.pdf\#nameddest=tag-02BW}{V.{E.P6} (p.~53)} }
\expandafter\def\csname SRef@maps:eq:E.P7\endcsname{\href{S6_SERIES_GRIND_V.pdf\#nameddest=tag-02BX}{V.{E.P7} (p.~53)} }
\expandafter\def\csname SRef@maps:eq:D.S7\endcsname{\href{S6_SERIES_GRIND_III.pdf\#nameddest=tag-023W}{III.{D.S7} (p.~67)} }
\expandafter\def\csname SRef@maps:eq:E.P8\endcsname{\href{S6_SERIES_GRIND_V.pdf\#nameddest=tag-02BY}{V.{E.P8} (p.~53)} }
\expandafter\def\csname SRef@maps:eq:E.P9\endcsname{\href{S6_SERIES_GRIND_V.pdf\#nameddest=tag-02BZ}{V.{E.P9} (p.~53)} }
\expandafter\def\csname SRef@maps:eq:E.P10\endcsname{\href{S6_SERIES_GRIND_V.pdf\#nameddest=tag-02BK}{V.{E.P10} (p.~54)} }
\expandafter\def\csname SRef@maps:eq:E.P11\endcsname{\href{S6_SERIES_GRIND_V.pdf\#nameddest=tag-02BL}{V.{E.P11} (p.~54)} }
\expandafter\def\csname SRef@maps:eq:E.P12\endcsname{\href{S6_SERIES_GRIND_V.pdf\#nameddest=tag-02BM}{V.{E.P12} (p.~54)} }
\expandafter\def\csname SRef@maps:eq:E.P13\endcsname{\href{S6_SERIES_GRIND_V.pdf\#nameddest=tag-02BN}{V.{E.P13} (p.~54)} }
\expandafter\def\csname SRef@maps:eq:E.P14\endcsname{\href{S6_SERIES_GRIND_V.pdf\#nameddest=tag-02BO}{V.{E.P14} (p.~55)} }
\expandafter\def\csname SRef@maps:eq:E.P15\endcsname{\href{S6_SERIES_GRIND_V.pdf\#nameddest=tag-02BP}{V.{E.P15} (p.~55)} }
\expandafter\def\csname SRef@maps:eq:E.P16\endcsname{\href{S6_SERIES_GRIND_V.pdf\#nameddest=tag-02BQ}{V.{E.P16} (p.~55)} }
\expandafter\def\csname SRef@maps:eq:E.P17\endcsname{\href{S6_SERIES_GRIND_V.pdf\#nameddest=tag-02BR}{V.{E.P17} (p.~56)} }
\expandafter\def\csname SRef@maps:eq:D.S8\endcsname{\href{S6_SERIES_GRIND_III.pdf\#nameddest=tag-023X}{III.{D.S8} (p.~67)} }
\expandafter\def\csname SRef@maps:eq:E.Q1\endcsname{\href{S6_SERIES_GRIND_V.pdf\#nameddest=tag-02C0}{V.{E.Q1} (p.~57)} }
\expandafter\def\csname SRef@maps:eq:E.Q2\endcsname{\href{S6_SERIES_GRIND_V.pdf\#nameddest=tag-02C4}{V.{E.Q2} (p.~57)} }
\expandafter\def\csname SRef@maps:eq:E.Q3\endcsname{\href{S6_SERIES_GRIND_V.pdf\#nameddest=tag-02C5}{V.{E.Q3} (p.~57)} }
\expandafter\def\csname SRef@maps:eq:E.Q4\endcsname{\href{S6_SERIES_GRIND_V.pdf\#nameddest=tag-02C6}{V.{E.Q4} (p.~57)} }
\expandafter\def\csname SRef@maps:eq:E.Q5\endcsname{\href{S6_SERIES_GRIND_V.pdf\#nameddest=tag-02C7}{V.{E.Q5} (p.~58)} }
\expandafter\def\csname SRef@maps:eq:E.Q6\endcsname{\href{S6_SERIES_GRIND_V.pdf\#nameddest=tag-02C8}{V.{E.Q6} (p.~58)} }
\expandafter\def\csname SRef@maps:eq:E.Q7\endcsname{\href{S6_SERIES_GRIND_V.pdf\#nameddest=tag-02C9}{V.{E.Q7} (p.~58)} }
\expandafter\def\csname SRef@maps:eq:E.Q8\endcsname{\href{S6_SERIES_GRIND_V.pdf\#nameddest=tag-02CA}{V.{E.Q8} (p.~58)} }
\expandafter\def\csname SRef@maps:eq:E.Q9\endcsname{\href{S6_SERIES_GRIND_V.pdf\#nameddest=tag-02CB}{V.{E.Q9} (p.~59)} }
\expandafter\def\csname SRef@maps:eq:E.Q10\endcsname{\href{S6_SERIES_GRIND_V.pdf\#nameddest=tag-02C1}{V.{E.Q10} (p.~59)} }
\expandafter\def\csname SRef@maps:eq:D.S9\endcsname{\href{S6_SERIES_GRIND_III.pdf\#nameddest=tag-023Y}{III.{D.S9} (p.~68)} }
\expandafter\def\csname SRef@maps:eq:E.Q11\endcsname{\href{S6_SERIES_GRIND_V.pdf\#nameddest=tag-02C2}{V.{E.Q11} (p.~59)} }
\expandafter\def\csname SRef@maps:eq:E.Q12\endcsname{\href{S6_SERIES_GRIND_V.pdf\#nameddest=tag-02C3}{V.{E.Q12} (p.~60)} }
\expandafter\def\csname SRef@maps:eq:E.T1\endcsname{\href{S6_SERIES_GRIND_V.pdf\#nameddest=tag-02DV}{V.{E.T1} (p.~60)} }
\expandafter\def\csname SRef@maps:eq:E.T2\endcsname{\href{S6_SERIES_GRIND_V.pdf\#nameddest=tag-02DW}{V.{E.T2} (p.~60)} }
\expandafter\def\csname SRef@maps:eq:E.T3\endcsname{\href{S6_SERIES_GRIND_V.pdf\#nameddest=tag-02DX}{V.{E.T3} (p.~61)} }
\expandafter\def\csname SRef@maps:eq:E.T4\endcsname{\href{S6_SERIES_GRIND_V.pdf\#nameddest=tag-02DY}{V.{E.T4} (p.~61)} }
\expandafter\def\csname SRef@maps:eq:E.T5\endcsname{\href{S6_SERIES_GRIND_V.pdf\#nameddest=tag-02DZ}{V.{E.T5} (p.~61)} }
\expandafter\def\csname SRef@maps:eq:E.T6\endcsname{\href{S6_SERIES_GRIND_V.pdf\#nameddest=tag-02E0}{V.{E.T6} (p.~62)} }
\expandafter\def\csname SRef@maps:eq:E.T7\endcsname{\href{S6_SERIES_GRIND_V.pdf\#nameddest=tag-02E1}{V.{E.T7} (p.~62)} }
\expandafter\def\csname SRef@maps:eq:E.T9\endcsname{\href{S6_SERIES_GRIND_V.pdf\#nameddest=tag-02E3}{V.{E.T9} (p.~62)} }
\expandafter\def\csname SRef@maps:eq:D.S10\endcsname{\href{S6_SERIES_GRIND_III.pdf\#nameddest=tag-023M}{III.{D.S10} (p.~68)} }
\expandafter\def\csname SRef@maps:eq:E.T8\endcsname{\href{S6_SERIES_GRIND_V.pdf\#nameddest=tag-02E2}{V.{E.T8} (p.~62)} }
\expandafter\def\csname SRef@maps:eq:E.U1\endcsname{\href{S6_SERIES_GRIND_V.pdf\#nameddest=tag-02E4}{V.{E.U1} (p.~63)} }
\expandafter\def\csname SRef@maps:eq:E.U2\endcsname{\href{S6_SERIES_GRIND_V.pdf\#nameddest=tag-02E7}{V.{E.U2} (p.~63)} }
\expandafter\def\csname SRef@maps:eq:E.U3\endcsname{\href{S6_SERIES_GRIND_V.pdf\#nameddest=tag-02E8}{V.{E.U3} (p.~63)} }
\expandafter\def\csname SRef@maps:eq:E.U4\endcsname{\href{S6_SERIES_GRIND_V.pdf\#nameddest=tag-02E9}{V.{E.U4} (p.~63)} }
\expandafter\def\csname SRef@maps:eq:E.U5\endcsname{\href{S6_SERIES_GRIND_V.pdf\#nameddest=tag-02EA}{V.{E.U5} (p.~63)} }
\expandafter\def\csname SRef@maps:eq:E.U6\endcsname{\href{S6_SERIES_GRIND_V.pdf\#nameddest=tag-02EB}{V.{E.U6} (p.~64)} }
\expandafter\def\csname SRef@maps:eq:E.U7\endcsname{\href{S6_SERIES_GRIND_V.pdf\#nameddest=tag-02EC}{V.{E.U7} (p.~64)} }
\expandafter\def\csname SRef@maps:eq:E.U8\endcsname{\href{S6_SERIES_GRIND_V.pdf\#nameddest=tag-02ED}{V.{E.U8} (p.~64)} }
\expandafter\def\csname SRef@maps:eq:E.U9\endcsname{\href{S6_SERIES_GRIND_V.pdf\#nameddest=tag-02EE}{V.{E.U9} (p.~64)} }
\expandafter\def\csname SRef@maps:eq:D.S11\endcsname{\href{S6_SERIES_GRIND_III.pdf\#nameddest=tag-023N}{III.{D.S11} (p.~69)} }
\expandafter\def\csname SRef@maps:eq:E.U10\endcsname{\href{S6_SERIES_GRIND_V.pdf\#nameddest=tag-02E5}{V.{E.U10} (p.~64)} }
\expandafter\def\csname SRef@maps:eq:E.U11\endcsname{\href{S6_SERIES_GRIND_V.pdf\#nameddest=tag-02E6}{V.{E.U11} (p.~64)} }
\expandafter\def\csname SRef@maps:eq:E.H0\endcsname{\href{S6_SERIES_GRIND_V.pdf\#nameddest=tag-02AQ}{V.{E.H0} (p.~65)} }
\expandafter\def\csname SRef@maps:eq:E.H1\endcsname{\href{S6_SERIES_GRIND_V.pdf\#nameddest=tag-02AR}{V.{E.H1} (p.~65)} }
\expandafter\def\csname SRef@maps:eq:E.H2\endcsname{\href{S6_SERIES_GRIND_V.pdf\#nameddest=tag-02AW}{V.{E.H2} (p.~66)} }
\expandafter\def\csname SRef@maps:eq:E.H3\endcsname{\href{S6_SERIES_GRIND_V.pdf\#nameddest=tag-02AX}{V.{E.H3} (p.~66)} }
\expandafter\def\csname SRef@maps:eq:E.H4\endcsname{\href{S6_SERIES_GRIND_V.pdf\#nameddest=tag-02AY}{V.{E.H4} (p.~66)} }
\expandafter\def\csname SRef@maps:eq:E.H5\endcsname{\href{S6_SERIES_GRIND_V.pdf\#nameddest=tag-02AZ}{V.{E.H5} (p.~66)} }
\expandafter\def\csname SRef@maps:eq:E.H6\endcsname{\href{S6_SERIES_GRIND_V.pdf\#nameddest=tag-02B0}{V.{E.H6} (p.~66)} }
\expandafter\def\csname SRef@maps:eq:E.H7\endcsname{\href{S6_SERIES_GRIND_V.pdf\#nameddest=tag-02B1}{V.{E.H7} (p.~66)} }
\expandafter\def\csname SRef@maps:eq:D.S12\endcsname{\href{S6_SERIES_GRIND_III.pdf\#nameddest=tag-023O}{III.{D.S12} (p.~69)} }
\expandafter\def\csname SRef@maps:eq:E.H8\endcsname{\href{S6_SERIES_GRIND_V.pdf\#nameddest=tag-02B2}{V.{E.H8} (p.~67)} }
\expandafter\def\csname SRef@maps:eq:E.H9\endcsname{\href{S6_SERIES_GRIND_V.pdf\#nameddest=tag-02B3}{V.{E.H9} (p.~67)} }
\expandafter\def\csname SRef@maps:eq:E.H9a\endcsname{\href{S6_SERIES_GRIND_V.pdf\#nameddest=tag-02B4}{V.{E.H9a} (p.~67)} }
\expandafter\def\csname SRef@maps:eq:E.H10\endcsname{\href{S6_SERIES_GRIND_V.pdf\#nameddest=tag-02AS}{V.{E.H10} (p.~68)} }
\expandafter\def\csname SRef@maps:eq:E.H11\endcsname{\href{S6_SERIES_GRIND_V.pdf\#nameddest=tag-02AT}{V.{E.H11} (p.~68)} }
\expandafter\def\csname SRef@maps:eq:E.H12\endcsname{\href{S6_SERIES_GRIND_V.pdf\#nameddest=tag-02AU}{V.{E.H12} (p.~68)} }
\expandafter\def\csname SRef@maps:eq:E.H13\endcsname{\href{S6_SERIES_GRIND_V.pdf\#nameddest=tag-02AV}{V.{E.H13} (p.~69)} }
\expandafter\def\csname SRef@maps:eq:D.S13\endcsname{\href{S6_SERIES_GRIND_III.pdf\#nameddest=tag-023P}{III.{D.S13} (p.~70)} }
\expandafter\def\csname SRef@maps:eq:D.S14\endcsname{\href{S6_SERIES_GRIND_III.pdf\#nameddest=tag-023Q}{III.{D.S14} (p.~70)} }
\expandafter\def\csname SRef@maps:eq:D.I1\endcsname{\href{S6_SERIES_GRIND_III.pdf\#nameddest=tag-01ZU}{III.{D.I1} (p.~71)} }
\expandafter\def\csname SRef@maps:eq:D.I2\endcsname{\href{S6_SERIES_GRIND_III.pdf\#nameddest=tag-01ZV}{III.{D.I2} (p.~71)} }
\expandafter\def\csname SRef@maps:eq:E.15\endcsname{\href{S6_SERIES_GRIND_III.pdf\#nameddest=tag-029L}{III.{E.15} (p.~317)} }
\expandafter\def\csname SRef@maps:eq:D.I3\endcsname{\href{S6_SERIES_GRIND_III.pdf\#nameddest=tag-01ZW}{III.{D.I3} (p.~71)} }
\expandafter\def\csname SRef@maps:eq:D.I4\endcsname{\href{S6_SERIES_GRIND_III.pdf\#nameddest=tag-01ZX}{III.{D.I4} (p.~72)} }
\expandafter\def\csname SRef@maps:eq:D.T1\endcsname{\href{S6_SERIES_GRIND_III.pdf\#nameddest=tag-0258}{III.{D.T1} (p.~73)} }
\expandafter\def\csname SRef@maps:eq:D.T2\endcsname{\href{S6_SERIES_GRIND_III.pdf\#nameddest=tag-025A}{III.{D.T2} (p.~74)} }
\expandafter\def\csname SRef@maps:eq:D.T3\endcsname{\href{S6_SERIES_GRIND_III.pdf\#nameddest=tag-025B}{III.{D.T3} (p.~74)} }
\expandafter\def\csname SRef@maps:eq:D.T4\endcsname{\href{S6_SERIES_GRIND_III.pdf\#nameddest=tag-025L}{III.{D.T4} (p.~74)} }
\expandafter\def\csname SRef@maps:eq:D.T5\endcsname{\href{S6_SERIES_GRIND_III.pdf\#nameddest=tag-025M}{III.{D.T5} (p.~75)} }
\expandafter\def\csname SRef@maps:eq:D.T6\endcsname{\href{S6_SERIES_GRIND_III.pdf\#nameddest=tag-025N}{III.{D.T6} (p.~75)} }
\expandafter\def\csname SRef@maps:eq:D.T7\endcsname{\href{S6_SERIES_GRIND_III.pdf\#nameddest=tag-025O}{III.{D.T7} (p.~75)} }
\expandafter\def\csname SRef@maps:eq:D.T8\endcsname{\href{S6_SERIES_GRIND_III.pdf\#nameddest=tag-025P}{III.{D.T8} (p.~75)} }
\expandafter\def\csname SRef@maps:eq:D.T9\endcsname{\href{S6_SERIES_GRIND_III.pdf\#nameddest=tag-025Q}{III.{D.T9} (p.~76)} }
\expandafter\def\csname SRef@maps:eq:D.T10\endcsname{\href{S6_SERIES_GRIND_III.pdf\#nameddest=tag-0259}{III.{D.T10} (p.~76)} }
\expandafter\def\csname SRef@maps:eq:D.V1\endcsname{\href{S6_SERIES_GRIND_III.pdf\#nameddest=tag-0269}{III.{D.V1} (p.~78)} }
\expandafter\def\csname SRef@maps:eq:D.V2\endcsname{\href{S6_SERIES_GRIND_III.pdf\#nameddest=tag-026A}{III.{D.V2} (p.~78)} }
\expandafter\def\csname SRef@maps:eq:D.V3\endcsname{\href{S6_SERIES_GRIND_III.pdf\#nameddest=tag-026B}{III.{D.V3} (p.~78)} }
\expandafter\def\csname SRef@maps:eq:D.V4\endcsname{\href{S6_SERIES_GRIND_III.pdf\#nameddest=tag-026C}{III.{D.V4} (p.~79)} }
\expandafter\def\csname SRef@maps:eq:D.V5\endcsname{\href{S6_SERIES_GRIND_III.pdf\#nameddest=tag-026D}{III.{D.V5} (p.~79)} }
\expandafter\def\csname SRef@maps:eq:D.F1\endcsname{\href{S6_SERIES_GRIND_III.pdf\#nameddest=tag-01YO}{III.{D.F1} (p.~81)} }
\expandafter\def\csname SRef@maps:eq:D.F2\endcsname{\href{S6_SERIES_GRIND_III.pdf\#nameddest=tag-01YP}{III.{D.F2} (p.~81)} }
\expandafter\def\csname SRef@maps:eq:D.O1\endcsname{\href{S6_SERIES_GRIND_III.pdf\#nameddest=tag-020M}{III.{D.O1} (p.~82)} }
\expandafter\def\csname SRef@maps:eq:D.O2\endcsname{\href{S6_SERIES_GRIND_III.pdf\#nameddest=tag-020N}{III.{D.O2} (p.~82)} }
\expandafter\def\csname SRef@maps:eq:D.O3\endcsname{\href{S6_SERIES_GRIND_III.pdf\#nameddest=tag-020O}{III.{D.O3} (p.~82)} }
\expandafter\def\csname SRef@maps:eq:D.O4\endcsname{\href{S6_SERIES_GRIND_III.pdf\#nameddest=tag-020P}{III.{D.O4} (p.~83)} }
\expandafter\def\csname SRef@maps:eq:D.O5\endcsname{\href{S6_SERIES_GRIND_III.pdf\#nameddest=tag-020Q}{III.{D.O5} (p.~83)} }
\expandafter\def\csname SRef@maps:eq:D.O6\endcsname{\href{S6_SERIES_GRIND_III.pdf\#nameddest=tag-020R}{III.{D.O6} (p.~83)} }
\expandafter\def\csname SRef@maps:eq:D.O7\endcsname{\href{S6_SERIES_GRIND_III.pdf\#nameddest=tag-020S}{III.{D.O7} (p.~84)} }
\expandafter\def\csname SRef@maps:eq:D.O8\endcsname{\href{S6_SERIES_GRIND_III.pdf\#nameddest=tag-020T}{III.{D.O8} (p.~84)} }
\expandafter\def\csname SRef@maps:src:D.2.2\endcsname{\href{S6_SERIES_GRIND_III.pdf\#nameddest=tag-02KU}{III.5.12 (p.~42)} }
\expandafter\def\csname SRef@maps:src:D.2.17\endcsname{\href{S6_SERIES_GRIND_III.pdf\#nameddest=tag-02KR}{III.5.26 (p.~77)} }
\expandafter\def\csname SRef@maps:src:D.2.28\endcsname{\href{S6_SERIES_GRIND_III.pdf\#nameddest=tag-02L3}{III.5.37 (p.~106)} }
\expandafter\def\csname SRef@wb:cusp-cone-S12-sphere-map-null\endcsname{\href{S6_SERIES_GRIND_III.pdf\#nameddest=tag-035Z}{III.7.14 (p.~270)} }
\expandafter\def\csname SRef@wb:cusp-cone-S24-sphere-map-null\endcsname{\href{S6_SERIES_GRIND_III.pdf\#nameddest=tag-036E}{III.7.15 (p.~270)} }
\expandafter\def\csname SRef@wb:cusp-cone-literal-R4-open-cells\endcsname{\href{S6_SERIES_GRIND_III.pdf\#nameddest=tag-036V}{III.7.16 (p.~271)} }
\expandafter\def\csname SRef@wb:cusp-cone-sharp-global-R4-scope\endcsname{\href{S6_SERIES_GRIND_III.pdf\#nameddest=tag-037J}{III.7.20 (p.~278)} }
\expandafter\def\csname SRef@wb:cusp-s12-all-higher-smashes-null\endcsname{\href{S6_SERIES_GRIND_III.pdf\#nameddest=tag-038T}{III.7.22 (p.~280)} }
\expandafter\def\csname SRef@wb:cusp_s12_carrier_homotopy:L567\endcsname{\href{S6_SERIES_GRIND_III.pdf\#nameddest=tag-03A6}{III.7.24 (p.~286)} }
\expandafter\def\csname SRef@wb:sphere-metric-unit-volume\endcsname{\href{S6_SERIES_GRIND_III.pdf\#nameddest=tag-03DH}{III.7.28 (p.~294)} }
\expandafter\def\csname SRef@wb:h24:ret:q2-detected-quotient\endcsname{\href{S6_SERIES_GRIND_III.pdf\#nameddest=tag-03BE}{III.7.33 (p.~306)} }
\expandafter\def\csname SRef@wb:cusp-cone-target-slice-embeddings\endcsname{\href{S6_SERIES_GRIND_III.pdf\#nameddest=tag-037N}{III.7.9 (p.~261)} }
\expandafter\def\csname SRef@wb:cusp-cover-target-action\endcsname{\href{S6_SERIES_GRIND_III.pdf\#nameddest=tag-0388}{III.1 (p.~238)} }
\expandafter\def\csname SRef@wb:cusp-cover-normalization-diagram\endcsname{\href{S6_SERIES_GRIND_III.pdf\#nameddest=tag-0383}{III.10 (p.~239)} }
\expandafter\def\csname SRef@wb:cusp-cone-products-not-spheres\endcsname{\href{S6_SERIES_GRIND_III.pdf\#nameddest=tag-0378}{III.100 (p.~263)} }
\expandafter\def\csname SRef@wb:cusp-cone-join-dimensions\endcsname{\href{S6_SERIES_GRIND_III.pdf\#nameddest=tag-036T}{III.101 (p.~263)} }
\expandafter\def\csname SRef@wb:cusp-cone-kernel-contained\endcsname{\href{S6_SERIES_GRIND_III.pdf\#nameddest=tag-036U}{III.102 (p.~263)} }
\expandafter\def\csname SRef@wb:cusp-cone-finite-index-orbit\endcsname{\href{S6_SERIES_GRIND_III.pdf\#nameddest=tag-036Q}{III.103 (p.~263)} }
\expandafter\def\csname SRef@wb:cusp-cone-fibre-H4-nonzero\endcsname{\href{S6_SERIES_GRIND_III.pdf\#nameddest=tag-036P}{III.104 (p.~264)} }
\expandafter\def\csname SRef@wb:cusp-cone-sphere-to-sphere-maps\endcsname{\href{S6_SERIES_GRIND_III.pdf\#nameddest=tag-037L}{III.105 (p.~266)} }
\expandafter\def\csname SRef@wb:cusp-cone-S12-S4-radial-map\endcsname{\href{S6_SERIES_GRIND_III.pdf\#nameddest=tag-035R}{III.106 (p.~267)} }
\expandafter\def\csname SRef@wb:cusp-cone-S24-S4-radial-map\endcsname{\href{S6_SERIES_GRIND_III.pdf\#nameddest=tag-0361}{III.107 (p.~267)} }
\expandafter\def\csname SRef@wb:cusp-cone-S12-exact-collapsed-locus\endcsname{\href{S6_SERIES_GRIND_III.pdf\#nameddest=tag-035T}{III.108 (p.~267)} }
\expandafter\def\csname SRef@wb:cusp-cone-S24-exact-collapsed-locus\endcsname{\href{S6_SERIES_GRIND_III.pdf\#nameddest=tag-0362}{III.109 (p.~267)} }
\expandafter\def\csname SRef@wb:cusp-cover-monodromy-conjugacy\endcsname{\href{S6_SERIES_GRIND_III.pdf\#nameddest=tag-0381}{III.11 (p.~239)} }
\expandafter\def\csname SRef@wb:cusp-cone-moving-S4-R4-limits\endcsname{\href{S6_SERIES_GRIND_III.pdf\#nameddest=tag-036Y}{III.110 (p.~267)} }
\expandafter\def\csname SRef@wb:cusp-cone-transported-skeleton-collapse\endcsname{\href{S6_SERIES_GRIND_III.pdf\#nameddest=tag-037O}{III.111 (p.~268)} }
\expandafter\def\csname SRef@wb:cusp-cone-distance-to-skeleton\endcsname{\href{S6_SERIES_GRIND_III.pdf\#nameddest=tag-036M}{III.112 (p.~268)} }
\expandafter\def\csname SRef@wb:cusp-cone-S12-onto-exact-ball\endcsname{\href{S6_SERIES_GRIND_III.pdf\#nameddest=tag-035W}{III.113 (p.~269)} }
\expandafter\def\csname SRef@wb:cusp-cone-S24-onto-exact-ball\endcsname{\href{S6_SERIES_GRIND_III.pdf\#nameddest=tag-0369}{III.114 (p.~270)} }
\expandafter\def\csname SRef@wb:cusp-cone-S24-null-maps\endcsname{\href{S6_SERIES_GRIND_III.pdf\#nameddest=tag-0367}{III.115 (p.~270)} }
\expandafter\def\csname SRef@wb:cusp-cone-S24-null-special-cases\endcsname{\href{S6_SERIES_GRIND_III.pdf\#nameddest=tag-0368}{III.116 (p.~270)} }
\expandafter\def\csname SRef@wb:cusp-cone-S24-smash-square-factorization\endcsname{\href{S6_SERIES_GRIND_III.pdf\#nameddest=tag-036C}{III.117 (p.~271)} }
\expandafter\def\csname SRef@wb:cusp-cone-S24-sixfold-suspension\endcsname{\href{S6_SERIES_GRIND_III.pdf\#nameddest=tag-036B}{III.118 (p.~271)} }
\expandafter\def\csname SRef@wb:cusp-cone-fourth-stable-stem\endcsname{\href{S6_SERIES_GRIND_III.pdf\#nameddest=tag-036R}{III.119 (p.~271)} }
\expandafter\def\csname SRef@wb:cusp-cover-maximal-chart-map\endcsname{\href{S6_SERIES_GRIND_III.pdf\#nameddest=tag-0380}{III.12 (p.~241)} }
\expandafter\def\csname SRef@wb:cusp-cone-U12-homeomorphism\endcsname{\href{S6_SERIES_GRIND_III.pdf\#nameddest=tag-036G}{III.121 (p.~271)} }
\expandafter\def\csname SRef@wb:cusp-cone-U24-homeomorphism\endcsname{\href{S6_SERIES_GRIND_III.pdf\#nameddest=tag-036H}{III.122 (p.~271)} }
\expandafter\def\csname SRef@wb:cusp-cone-D12-embedding\endcsname{\href{S6_SERIES_GRIND_III.pdf\#nameddest=tag-035G}{III.123 (p.~273)} }
\expandafter\def\csname SRef@wb:cusp-cone-D24-embedding\endcsname{\href{S6_SERIES_GRIND_III.pdf\#nameddest=tag-035I}{III.124 (p.~273)} }
\expandafter\def\csname SRef@wb:cusp-cone-D12-slice\endcsname{\href{S6_SERIES_GRIND_III.pdf\#nameddest=tag-035H}{III.125 (p.~273)} }
\expandafter\def\csname SRef@wb:cusp-cone-D24-slice\endcsname{\href{S6_SERIES_GRIND_III.pdf\#nameddest=tag-035J}{III.126 (p.~273)} }
\expandafter\def\csname SRef@wb:cusp-cone-carrier-local-homology\endcsname{\href{S6_SERIES_GRIND_III.pdf\#nameddest=tag-036K}{III.127 (p.~273)} }
\expandafter\def\csname SRef@wb:cusp-cone-A-smash-wedge\endcsname{\href{S6_SERIES_GRIND_III.pdf\#nameddest=tag-0357}{III.128 (p.~274)} }
\expandafter\def\csname SRef@wb:cusp-cone-C-smash-wedge\endcsname{\href{S6_SERIES_GRIND_III.pdf\#nameddest=tag-035F}{III.129 (p.~274)} }
\expandafter\def\csname SRef@wb:cusp-cone-E12-fibres\endcsname{\href{S6_SERIES_GRIND_III.pdf\#nameddest=tag-035K}{III.130 (p.~276)} }
\expandafter\def\csname SRef@wb:cusp-cone-E24-fibres\endcsname{\href{S6_SERIES_GRIND_III.pdf\#nameddest=tag-035L}{III.131 (p.~276)} }
\expandafter\def\csname SRef@wb:cusp-s12-smash-factorization\endcsname{\href{S6_SERIES_GRIND_III.pdf\#nameddest=tag-0390}{III.132 (p.~280)} }
\expandafter\def\csname SRef@wb:cusp-s12-toda-smash-formula\endcsname{\href{S6_SERIES_GRIND_III.pdf\#nameddest=tag-0392}{III.133 (p.~280)} }
\expandafter\def\csname SRef@wb:cusp-s12-three-stem-groups\endcsname{\href{S6_SERIES_GRIND_III.pdf\#nameddest=tag-0391}{III.134 (p.~280)} }
\expandafter\def\csname SRef@wb:cusp-s12-sixfold-suspension-class\endcsname{\href{S6_SERIES_GRIND_III.pdf\#nameddest=tag-038Z}{III.135 (p.~281)} }
\expandafter\def\csname SRef@wb:cusp-s12-regular-inverse-image\endcsname{\href{S6_SERIES_GRIND_III.pdf\#nameddest=tag-038Y}{III.136 (p.~281)} }
\expandafter\def\csname SRef@wb:cusp-null:eq:periods\endcsname{\href{S6_SERIES_GRIND_III.pdf\#nameddest=tag-038N}{III.137 (p.~282)} }
\expandafter\def\csname SRef@wb:cusp-null:eq:action\endcsname{\href{S6_SERIES_GRIND_III.pdf\#nameddest=tag-038E}{III.138 (p.~282)} }
\expandafter\def\csname SRef@wb:cusp-null:eq:g\endcsname{\href{S6_SERIES_GRIND_III.pdf\#nameddest=tag-038J}{III.139 (p.~282)} }
\expandafter\def\csname SRef@wb:cusp-null:eq:chart\endcsname{\href{S6_SERIES_GRIND_III.pdf\#nameddest=tag-038H}{III.140 (p.~282)} }
\expandafter\def\csname SRef@wb:cusp-null:eq:moduli\endcsname{\href{S6_SERIES_GRIND_III.pdf\#nameddest=tag-038L}{III.141 (p.~282)} }
\expandafter\def\csname SRef@wb:cusp-null:eq:invertibility\endcsname{\href{S6_SERIES_GRIND_III.pdf\#nameddest=tag-038K}{III.142 (p.~283)} }
\expandafter\def\csname SRef@wb:cusp-null:eq:straightening\endcsname{\href{S6_SERIES_GRIND_III.pdf\#nameddest=tag-038Q}{III.144 (p.~283)} }
\expandafter\def\csname SRef@wb:cusp-null:eq:torus\endcsname{\href{S6_SERIES_GRIND_III.pdf\#nameddest=tag-038R}{III.145 (p.~283)} }
\expandafter\def\csname SRef@wb:cusp-null:eq:spatial-bound\endcsname{\href{S6_SERIES_GRIND_III.pdf\#nameddest=tag-038P}{III.146 (p.~283)} }
\expandafter\def\csname SRef@wb:cusp-null:eq:nearby-moduli\endcsname{\href{S6_SERIES_GRIND_III.pdf\#nameddest=tag-038M}{III.147 (p.~284)} }
\expandafter\def\csname SRef@wb:cusp-null:eq:J\endcsname{\href{S6_SERIES_GRIND_III.pdf\#nameddest=tag-038B}{III.148 (p.~284)} }
\expandafter\def\csname SRef@wb:cusp-null:eq:actual-frame\endcsname{\href{S6_SERIES_GRIND_III.pdf\#nameddest=tag-038F}{III.149 (p.~284)} }
\expandafter\def\csname SRef@wb:cusp-cover-equivariance-computation\endcsname{\href{S6_SERIES_GRIND_III.pdf\#nameddest=tag-037W}{III.15 (p.~241)} }
\expandafter\def\csname SRef@wb:cusp-null:eq:J0\endcsname{\href{S6_SERIES_GRIND_III.pdf\#nameddest=tag-038C}{III.150 (p.~284)} }
\expandafter\def\csname SRef@wb:cusp-null:eq:P\endcsname{\href{S6_SERIES_GRIND_III.pdf\#nameddest=tag-038D}{III.151 (p.~284)} }
\expandafter\def\csname SRef@wb:cusp-null:eq:q\endcsname{\href{S6_SERIES_GRIND_III.pdf\#nameddest=tag-038O}{III.152 (p.~285)} }
\expandafter\def\csname SRef@wb:cusp-null:eq:bordism\endcsname{\href{S6_SERIES_GRIND_III.pdf\#nameddest=tag-038G}{III.153 (p.~285)} }
\expandafter\def\csname SRef@wb:cusp-null:eq:extended-frame\endcsname{\href{S6_SERIES_GRIND_III.pdf\#nameddest=tag-038I}{III.154 (p.~285)} }
\expandafter\def\csname SRef@wb:sphere-metric-normalized-Q-limit\endcsname{\href{S6_SERIES_GRIND_III.pdf\#nameddest=tag-03DA}{III.155 (p.~287)} }
\expandafter\def\csname SRef@wb:sphere-metric-Hausdorff-collapse\endcsname{\href{S6_SERIES_GRIND_III.pdf\#nameddest=tag-03D2}{III.158 (p.~288)} }
\expandafter\def\csname SRef@wb:sphere-metric-A-fixed-metric\endcsname{\href{S6_SERIES_GRIND_III.pdf\#nameddest=tag-03CW}{III.159 (p.~288)} }
\expandafter\def\csname SRef@wb:cusp-cover-basechange-local-ring\endcsname{\href{S6_SERIES_GRIND_III.pdf\#nameddest=tag-037S}{III.16 (p.~242)} }
\expandafter\def\csname SRef@wb:sphere-metric-BC-fixed-metric\endcsname{\href{S6_SERIES_GRIND_III.pdf\#nameddest=tag-03CY}{III.160 (p.~288)} }
\expandafter\def\csname SRef@wb:sphere-metric-A-slice-exact\endcsname{\href{S6_SERIES_GRIND_III.pdf\#nameddest=tag-03CX}{III.161 (p.~289)} }
\expandafter\def\csname SRef@wb:sphere-metric-BC-slice-exact\endcsname{\href{S6_SERIES_GRIND_III.pdf\#nameddest=tag-03CZ}{III.162 (p.~289)} }
\expandafter\def\csname SRef@wb:sphere-metric-slice-fixed-blowups\endcsname{\href{S6_SERIES_GRIND_III.pdf\#nameddest=tag-03DE}{III.164 (p.~289)} }
\expandafter\def\csname SRef@wb:sphere-metric-ambient-fixed-blowups\endcsname{\href{S6_SERIES_GRIND_III.pdf\#nameddest=tag-03D3}{III.166 (p.~289)} }
\expandafter\def\csname SRef@wb:sphere-metric-scaling-trichotomy-result\endcsname{\href{S6_SERIES_GRIND_III.pdf\#nameddest=tag-03DD}{III.167 (p.~291)} }
\expandafter\def\csname SRef@wb:sphere-metric-vertex-distance\endcsname{\href{S6_SERIES_GRIND_III.pdf\#nameddest=tag-03DI}{III.168 (p.~292)} }
\expandafter\def\csname SRef@wb:sphere-metric-intrinsic-ambient-agreement\endcsname{\href{S6_SERIES_GRIND_III.pdf\#nameddest=tag-03D7}{III.169 (p.~292)} }
\expandafter\def\csname SRef@wb:cusp-cover-invariant-ring\endcsname{\href{S6_SERIES_GRIND_III.pdf\#nameddest=tag-037Z}{III.17 (p.~242)} }
\expandafter\def\csname SRef@wb:sphere-metric-intrinsic-ambient-ball-comparison\endcsname{\href{S6_SERIES_GRIND_III.pdf\#nameddest=tag-03D8}{III.170 (p.~292)} }
\expandafter\def\csname SRef@wb:sphere-metric-arbitrary-unrolling-ambient-bubbles\endcsname{\href{S6_SERIES_GRIND_III.pdf\#nameddest=tag-03D4}{III.172 (p.~292)} }
\expandafter\def\csname SRef@wb:sphere-metric-outward-shortcut-exclusion\endcsname{\href{S6_SERIES_GRIND_III.pdf\#nameddest=tag-03DB}{III.173 (p.~293)} }
\expandafter\def\csname SRef@wb:sphere-metric-C-smash-local-homology\endcsname{\href{S6_SERIES_GRIND_III.pdf\#nameddest=tag-03D1}{III.174 (p.~295)} }
\expandafter\def\csname SRef@wb:sphere-metric-C-boundary-quotient-homeomorphism\endcsname{\href{S6_SERIES_GRIND_III.pdf\#nameddest=tag-03D0}{III.175 (p.~296)} }
\expandafter\def\csname SRef@wb:sphere-metric-boundary-link-local-shift\endcsname{\href{S6_SERIES_GRIND_III.pdf\#nameddest=tag-03D5}{III.176 (p.~296)} }
\expandafter\def\csname SRef@wb:h24:ret:eq:original-periods\endcsname{\href{S6_SERIES_GRIND_III.pdf\#nameddest=tag-03B4}{III.177 (p.~298)} }
\expandafter\def\csname SRef@wb:h24:ret:eq:A\endcsname{\href{S6_SERIES_GRIND_III.pdf\#nameddest=tag-03AF}{III.178 (p.~298)} }
\expandafter\def\csname SRef@wb:h24:ret:eq:deck\endcsname{\href{S6_SERIES_GRIND_III.pdf\#nameddest=tag-03AS}{III.179 (p.~298)} }
\expandafter\def\csname SRef@wb:h24:ret:eq:logx\endcsname{\href{S6_SERIES_GRIND_III.pdf\#nameddest=tag-03B0}{III.180 (p.~298)} }
\expandafter\def\csname SRef@wb:h24:ret:eq:invertibility\endcsname{\href{S6_SERIES_GRIND_III.pdf\#nameddest=tag-03AZ}{III.181 (p.~298)} }
\expandafter\def\csname SRef@wb:h24:ret:eq:native\endcsname{\href{S6_SERIES_GRIND_III.pdf\#nameddest=tag-03B1}{III.182 (p.~298)} }
\expandafter\def\csname SRef@wb:h24:ret:eq:two-skeleton\endcsname{\href{S6_SERIES_GRIND_III.pdf\#nameddest=tag-03B9}{III.183 (p.~298)} }
\expandafter\def\csname SRef@wb:h24:ret:eq:A-target\endcsname{\href{S6_SERIES_GRIND_III.pdf\#nameddest=tag-03AH}{III.184 (p.~299)} }
\expandafter\def\csname SRef@wb:h24:ret:eq:q2\endcsname{\href{S6_SERIES_GRIND_III.pdf\#nameddest=tag-03B5}{III.185 (p.~299)} }
\expandafter\def\csname SRef@wb:h24:ret:eq:H\endcsname{\href{S6_SERIES_GRIND_III.pdf\#nameddest=tag-03AN}{III.186 (p.~299)} }
\expandafter\def\csname SRef@wb:h24:ret:eq:four-cell\endcsname{\href{S6_SERIES_GRIND_III.pdf\#nameddest=tag-03AW}{III.187 (p.~299)} }
\expandafter\def\csname SRef@wb:h24:ret:eq:chosen-value\endcsname{\href{S6_SERIES_GRIND_III.pdf\#nameddest=tag-03AR}{III.188 (p.~299)} }
\expandafter\def\csname SRef@wb:h24:ret:eq:z-chart\endcsname{\href{S6_SERIES_GRIND_III.pdf\#nameddest=tag-03BB}{III.189 (p.~299)} }
\expandafter\def\csname SRef@wb:cusp-cover-triple-fibre-length\endcsname{\href{S6_SERIES_GRIND_III.pdf\#nameddest=tag-038A}{III.19 (p.~242)} }
\expandafter\def\csname SRef@wb:h24:ret:eq:units\endcsname{\href{S6_SERIES_GRIND_III.pdf\#nameddest=tag-03BA}{III.190 (p.~299)} }
\expandafter\def\csname SRef@wb:h24:ret:eq:F\endcsname{\href{S6_SERIES_GRIND_III.pdf\#nameddest=tag-03AL}{III.191 (p.~299)} }
\expandafter\def\csname SRef@wb:h24:ret:eq:Finv\endcsname{\href{S6_SERIES_GRIND_III.pdf\#nameddest=tag-03AM}{III.192 (p.~299)} }
\expandafter\def\csname SRef@wb:h24:ret:eq:M\endcsname{\href{S6_SERIES_GRIND_III.pdf\#nameddest=tag-03AP}{III.193 (p.~300)} }
\expandafter\def\csname SRef@wb:h24:ret:eq:exact-u\endcsname{\href{S6_SERIES_GRIND_III.pdf\#nameddest=tag-03AV}{III.194 (p.~300)} }
\expandafter\def\csname SRef@wb:h24:ret:eq:full-derivative\endcsname{\href{S6_SERIES_GRIND_III.pdf\#nameddest=tag-03AX}{III.195 (p.~300)} }
\expandafter\def\csname SRef@wb:h24:ret:eq:derivative-on-M\endcsname{\href{S6_SERIES_GRIND_III.pdf\#nameddest=tag-03AT}{III.196 (p.~300)} }
\expandafter\def\csname SRef@wb:h24:ret:eq:LQ\endcsname{\href{S6_SERIES_GRIND_III.pdf\#nameddest=tag-03AO}{III.197 (p.~301)} }
\expandafter\def\csname SRef@wb:h24:ret:eq:tangents\endcsname{\href{S6_SERIES_GRIND_III.pdf\#nameddest=tag-03B8}{III.198 (p.~301)} }
\expandafter\def\csname SRef@wb:h24:ret:eq:normal-first\endcsname{\href{S6_SERIES_GRIND_III.pdf\#nameddest=tag-03B2}{III.199 (p.~301)} }
\expandafter\def\csname SRef@wb:cusp-cover-base-change\endcsname{\href{S6_SERIES_GRIND_III.pdf\#nameddest=tag-037Q}{III.2 (p.~238)} }
\expandafter\def\csname SRef@wb:cusp-cover-double-fibre-length\endcsname{\href{S6_SERIES_GRIND_III.pdf\#nameddest=tag-037V}{III.20 (p.~242)} }
\expandafter\def\csname SRef@wb:h24:ret:eq:normal-u\endcsname{\href{S6_SERIES_GRIND_III.pdf\#nameddest=tag-03B3}{III.200 (p.~301)} }
\expandafter\def\csname SRef@wb:h24:ret:eq:differential-table\endcsname{\href{S6_SERIES_GRIND_III.pdf\#nameddest=tag-03AU}{III.201 (p.~301)} }
\expandafter\def\csname SRef@wb:h24:ret:eq:A-homotopy\endcsname{\href{S6_SERIES_GRIND_III.pdf\#nameddest=tag-03AG}{III.202 (p.~301)} }
\expandafter\def\csname SRef@wb:h24:ret:eq:sign\endcsname{\href{S6_SERIES_GRIND_III.pdf\#nameddest=tag-03B7}{III.203 (p.~302)} }
\expandafter\def\csname SRef@wb:h24:ret:eq:D\endcsname{\href{S6_SERIES_GRIND_III.pdf\#nameddest=tag-03AJ}{III.204 (p.~302)} }
\expandafter\def\csname SRef@wb:h24:ret:eq:full-frame\endcsname{\href{S6_SERIES_GRIND_III.pdf\#nameddest=tag-03AY}{III.205 (p.~302)} }
\expandafter\def\csname SRef@wb:h24:ret:eq:D-homotopy\endcsname{\href{S6_SERIES_GRIND_III.pdf\#nameddest=tag-03AK}{III.206 (p.~302)} }
\expandafter\def\csname SRef@wb:h24:ret:eq:quadratic\endcsname{\href{S6_SERIES_GRIND_III.pdf\#nameddest=tag-03B6}{III.207 (p.~303)} }
\expandafter\def\csname SRef@wb:h24:ret:eq:Arf\endcsname{\href{S6_SERIES_GRIND_III.pdf\#nameddest=tag-03AI}{III.208 (p.~303)} }
\expandafter\def\csname SRef@wb:h24:ret:eq:actual-Arf\endcsname{\href{S6_SERIES_GRIND_III.pdf\#nameddest=tag-03AQ}{III.209 (p.~304)} }
\expandafter\def\csname SRef@wb:cusp-cover-additive-fibre-map\endcsname{\href{S6_SERIES_GRIND_III.pdf\#nameddest=tag-037P}{III.21 (p.~243)} }
\expandafter\def\csname SRef@mapsD:repair:cocycle-relation\endcsname{\href{S6_SERIES_GRIND_III.pdf\#nameddest=tag-02OF}{III.210 (p.~310)} }
\expandafter\def\csname SRef@mapsD:repair:canonical-matrices\endcsname{\href{S6_SERIES_GRIND_III.pdf\#nameddest=tag-02OE}{III.211 (p.~312)} }
\expandafter\def\csname SRef@mapsD:repair:free-peripherals\endcsname{\href{S6_SERIES_GRIND_III.pdf\#nameddest=tag-02OH}{III.212 (p.~312)} }
\expandafter\def\csname SRef@mapsD:repair:50-weierstrass\endcsname{\href{S6_SERIES_GRIND_III.pdf\#nameddest=tag-02OD}{III.213 (p.~312)} }
\expandafter\def\csname SRef@mapsD:repair:height-formula\endcsname{\href{S6_SERIES_GRIND_III.pdf\#nameddest=tag-02OI}{III.214 (p.~315)} }
\expandafter\def\csname SRef@wb:cusp-spec-radial-notation\endcsname{\href{S6_SERIES_GRIND_III.pdf\#nameddest=tag-03A3}{III.22 (p.~244)} }
\expandafter\def\csname SRef@wb:cusp-spec-holomorphic-errors\endcsname{\href{S6_SERIES_GRIND_III.pdf\#nameddest=tag-039V}{III.23 (p.~244)} }
\expandafter\def\csname SRef@wb:cusp-spec-full-inverse-transpose\endcsname{\href{S6_SERIES_GRIND_III.pdf\#nameddest=tag-039O}{III.24 (p.~244)} }
\expandafter\def\csname SRef@wb:cusp-spec-global-circle-map\endcsname{\href{S6_SERIES_GRIND_III.pdf\#nameddest=tag-039R}{III.25 (p.~246)} }
\expandafter\def\csname SRef@wb:cusp-spec-global-circle-quotient\endcsname{\href{S6_SERIES_GRIND_III.pdf\#nameddest=tag-039S}{III.26 (p.~246)} }
\expandafter\def\csname SRef@wb:cusp-spec-global-circle-coordinate\endcsname{\href{S6_SERIES_GRIND_III.pdf\#nameddest=tag-039Q}{III.27 (p.~246)} }
\expandafter\def\csname SRef@wb:cusp-spec-global-kernel-family\endcsname{\href{S6_SERIES_GRIND_III.pdf\#nameddest=tag-039T}{III.28 (p.~246)} }
\expandafter\def\csname SRef@wb:cusp-spec-radial-circle-character\endcsname{\href{S6_SERIES_GRIND_III.pdf\#nameddest=tag-03A2}{III.29 (p.~247)} }
\expandafter\def\csname SRef@wb:cusp-cover-source-factor\endcsname{\href{S6_SERIES_GRIND_III.pdf\#nameddest=tag-0387}{III.3 (p.~238)} }
\expandafter\def\csname SRef@wb:cusp-spec-global-circle-circumference\endcsname{\href{S6_SERIES_GRIND_III.pdf\#nameddest=tag-039P}{III.30 (p.~247)} }
\expandafter\def\csname SRef@wb:cusp-spec-fourfold-fibre-product\endcsname{\href{S6_SERIES_GRIND_III.pdf\#nameddest=tag-039N}{III.31 (p.~247)} }
\expandafter\def\csname SRef@wb:cusp-spec-fourfold-circle-map\endcsname{\href{S6_SERIES_GRIND_III.pdf\#nameddest=tag-039M}{III.32 (p.~247)} }
\expandafter\def\csname SRef@wb:cusp-spec-four-circle-fibre\endcsname{\href{S6_SERIES_GRIND_III.pdf\#nameddest=tag-039J}{III.33 (p.~247)} }
\expandafter\def\csname SRef@wb:cusp-spec-four-circle-limit\endcsname{\href{S6_SERIES_GRIND_III.pdf\#nameddest=tag-039K}{III.34 (p.~247)} }
\expandafter\def\csname SRef@wb:cusp-spec-doubled-local-quotient\endcsname{\href{S6_SERIES_GRIND_III.pdf\#nameddest=tag-0398}{III.35 (p.~247)} }
\expandafter\def\csname SRef@wb:cusp-spec-finite-affine-action\endcsname{\href{S6_SERIES_GRIND_III.pdf\#nameddest=tag-039A}{III.36 (p.~248)} }
\expandafter\def\csname SRef@wb:cusp-spec-finite-character-shift\endcsname{\href{S6_SERIES_GRIND_III.pdf\#nameddest=tag-039B}{III.37 (p.~248)} }
\expandafter\def\csname SRef@wb:cusp-spec-finite-Seifert-target\endcsname{\href{S6_SERIES_GRIND_III.pdf\#nameddest=tag-0399}{III.38 (p.~248)} }
\expandafter\def\csname SRef@wb:cusp-spec-finite-fourfold-normalization\endcsname{\href{S6_SERIES_GRIND_III.pdf\#nameddest=tag-039D}{III.39 (p.~248)} }
\expandafter\def\csname SRef@wb:cusp-cover-source-action\endcsname{\href{S6_SERIES_GRIND_III.pdf\#nameddest=tag-0386}{III.4 (p.~238)} }
\expandafter\def\csname SRef@wb:cusp-spec-finite-fourfold-target\endcsname{\href{S6_SERIES_GRIND_III.pdf\#nameddest=tag-039E}{III.40 (p.~248)} }
\expandafter\def\csname SRef@wb:cusp-spec-finite-normalization-map\endcsname{\href{S6_SERIES_GRIND_III.pdf\#nameddest=tag-039F}{III.41 (p.~248)} }
\expandafter\def\csname SRef@wb:cusp-spec-product-not-sphere\endcsname{\href{S6_SERIES_GRIND_III.pdf\#nameddest=tag-03A1}{III.42 (p.~248)} }
\expandafter\def\csname SRef@wb:cusp-spec-four-product-not-sphere\endcsname{\href{S6_SERIES_GRIND_III.pdf\#nameddest=tag-039L}{III.43 (p.~248)} }
\expandafter\def\csname SRef@wb:cusp-spec-smash-spheres\endcsname{\href{S6_SERIES_GRIND_III.pdf\#nameddest=tag-03A4}{III.44 (p.~249)} }
\expandafter\def\csname SRef@wb:cusp-spec-target-smash-map\endcsname{\href{S6_SERIES_GRIND_III.pdf\#nameddest=tag-03A5}{III.45 (p.~249)} }
\expandafter\def\csname SRef@wb:cusp-spec-cusp-kernel\endcsname{\href{S6_SERIES_GRIND_III.pdf\#nameddest=tag-0394}{III.46 (p.~249)} }
\expandafter\def\csname SRef@wb:cusp-spec-invariant-flag-calculation\endcsname{\href{S6_SERIES_GRIND_III.pdf\#nameddest=tag-039X}{III.47 (p.~249)} }
\expandafter\def\csname SRef@wb:cusp-spec-invariant-character-lattice\endcsname{\href{S6_SERIES_GRIND_III.pdf\#nameddest=tag-039W}{III.48 (p.~249)} }
\expandafter\def\csname SRef@wb:cusp-spec-circle-dual-vector\endcsname{\href{S6_SERIES_GRIND_III.pdf\#nameddest=tag-0393}{III.49 (p.~250)} }
\expandafter\def\csname SRef@wb:cusp-cover-toric-power-map\endcsname{\href{S6_SERIES_GRIND_III.pdf\#nameddest=tag-0389}{III.5 (p.~239)} }
\expandafter\def\csname SRef@wb:cusp-spec-no-finite-cover-matrix\endcsname{\href{S6_SERIES_GRIND_III.pdf\#nameddest=tag-039Z}{III.50 (p.~251)} }
\expandafter\def\csname SRef@wb:cusp-spec-no-finite-cover-orbit\endcsname{\href{S6_SERIES_GRIND_III.pdf\#nameddest=tag-03A0}{III.51 (p.~251)} }
\expandafter\def\csname SRef@wb:cusp-spec-finite-target-conjugacy\endcsname{\href{S6_SERIES_GRIND_III.pdf\#nameddest=tag-039I}{III.52 (p.~251)} }
\expandafter\def\csname SRef@wb:cusp-spec-finite-target-cocycle\endcsname{\href{S6_SERIES_GRIND_III.pdf\#nameddest=tag-039H}{III.53 (p.~251)} }
\expandafter\def\csname SRef@wb:cusp-spec-finite-overlap-identity\endcsname{\href{S6_SERIES_GRIND_III.pdf\#nameddest=tag-039G}{III.54 (p.~251)} }
\expandafter\def\csname SRef@wb:cusp-spec-finite-fibre-product-branches\endcsname{\href{S6_SERIES_GRIND_III.pdf\#nameddest=tag-039C}{III.55 (p.~252)} }
\expandafter\def\csname SRef@wb:cusp-spec-cusp-nonextension-position\endcsname{\href{S6_SERIES_GRIND_III.pdf\#nameddest=tag-0397}{III.56 (p.~252)} }
\expandafter\def\csname SRef@wb:cusp-spec-cusp-nonextension-character\endcsname{\href{S6_SERIES_GRIND_III.pdf\#nameddest=tag-0395}{III.57 (p.~252)} }
\expandafter\def\csname SRef@wb:cusp-spec-cusp-nonextension-chart\endcsname{\href{S6_SERIES_GRIND_III.pdf\#nameddest=tag-0396}{III.58 (p.~253)} }
\expandafter\def\csname SRef@wb:cusp-cone-disc\endcsname{\href{S6_SERIES_GRIND_III.pdf\#nameddest=tag-036L}{III.59 (p.~255)} }
\expandafter\def\csname SRef@wb:cusp-cover-completed-map\endcsname{\href{S6_SERIES_GRIND_III.pdf\#nameddest=tag-037T}{III.6 (p.~239)} }
\expandafter\def\csname SRef@wb:cusp-cone-quotient-lattice\endcsname{\href{S6_SERIES_GRIND_III.pdf\#nameddest=tag-037F}{III.60 (p.~255)} }
\expandafter\def\csname SRef@wb:cusp-cone-native-family-map\endcsname{\href{S6_SERIES_GRIND_III.pdf\#nameddest=tag-0371}{III.61 (p.~255)} }
\expandafter\def\csname SRef@wb:cusp-cone-radial-isometry\endcsname{\href{S6_SERIES_GRIND_III.pdf\#nameddest=tag-037G}{III.62 (p.~255)} }
\expandafter\def\csname SRef@wb:cusp-cone-native-restriction\endcsname{\href{S6_SERIES_GRIND_III.pdf\#nameddest=tag-0372}{III.63 (p.~255)} }
\expandafter\def\csname SRef@wb:cusp-cone-native-systole\endcsname{\href{S6_SERIES_GRIND_III.pdf\#nameddest=tag-0373}{III.64 (p.~255)} }
\expandafter\def\csname SRef@wb:cusp-cone-monodromy-image\endcsname{\href{S6_SERIES_GRIND_III.pdf\#nameddest=tag-036X}{III.65 (p.~255)} }
\expandafter\def\csname SRef@wb:cusp-cone-A-definition\endcsname{\href{S6_SERIES_GRIND_III.pdf\#nameddest=tag-0354}{III.66 (p.~256)} }
\expandafter\def\csname SRef@wb:cusp-cone-A-suspension\endcsname{\href{S6_SERIES_GRIND_III.pdf\#nameddest=tag-0358}{III.67 (p.~256)} }
\expandafter\def\csname SRef@wb:cusp-cone-q2\endcsname{\href{S6_SERIES_GRIND_III.pdf\#nameddest=tag-037D}{III.68 (p.~256)} }
\expandafter\def\csname SRef@wb:cusp-cone-q2-formula\endcsname{\href{S6_SERIES_GRIND_III.pdf\#nameddest=tag-037E}{III.69 (p.~256)} }
\expandafter\def\csname SRef@wb:cusp-cover-direct-ramification\endcsname{\href{S6_SERIES_GRIND_III.pdf\#nameddest=tag-037U}{III.7 (p.~239)} }
\expandafter\def\csname SRef@wb:cusp-cone-S12-map\endcsname{\href{S6_SERIES_GRIND_III.pdf\#nameddest=tag-035V}{III.70 (p.~256)} }
\expandafter\def\csname SRef@wb:cusp-cone-S12-source\endcsname{\href{S6_SERIES_GRIND_III.pdf\#nameddest=tag-035Y}{III.71 (p.~256)} }
\expandafter\def\csname SRef@wb:cusp-cone-S12-fibre-map\endcsname{\href{S6_SERIES_GRIND_III.pdf\#nameddest=tag-035U}{III.72 (p.~256)} }
\expandafter\def\csname SRef@wb:cusp-cone-S12-R4-limit\endcsname{\href{S6_SERIES_GRIND_III.pdf\#nameddest=tag-035Q}{III.73 (p.~256)} }
\expandafter\def\csname SRef@wb:cusp-cone-A-homology\endcsname{\href{S6_SERIES_GRIND_III.pdf\#nameddest=tag-0355}{III.74 (p.~256)} }
\expandafter\def\csname SRef@wb:cusp-cone-A-smash-homology\endcsname{\href{S6_SERIES_GRIND_III.pdf\#nameddest=tag-0356}{III.75 (p.~257)} }
\expandafter\def\csname SRef@wb:cusp-cone-B-definition\endcsname{\href{S6_SERIES_GRIND_III.pdf\#nameddest=tag-0359}{III.76 (p.~258)} }
\expandafter\def\csname SRef@wb:cusp-cone-q1-local\endcsname{\href{S6_SERIES_GRIND_III.pdf\#nameddest=tag-037B}{III.77 (p.~258)} }
\expandafter\def\csname SRef@wb:cusp-cone-q1-local-formula\endcsname{\href{S6_SERIES_GRIND_III.pdf\#nameddest=tag-037C}{III.78 (p.~258)} }
\expandafter\def\csname SRef@wb:cusp-cone-S24-local-fibre-map\endcsname{\href{S6_SERIES_GRIND_III.pdf\#nameddest=tag-0364}{III.79 (p.~258)} }
\expandafter\def\csname SRef@wb:cusp-cover-base-changed-target\endcsname{\href{S6_SERIES_GRIND_III.pdf\#nameddest=tag-037R}{III.8 (p.~239)} }
\expandafter\def\csname SRef@wb:cusp-cone-B-homology\endcsname{\href{S6_SERIES_GRIND_III.pdf\#nameddest=tag-035A}{III.80 (p.~258)} }
\expandafter\def\csname SRef@wb:cusp-cone-B-smash-homology\endcsname{\href{S6_SERIES_GRIND_III.pdf\#nameddest=tag-035B}{III.81 (p.~258)} }
\expandafter\def\csname SRef@wb:cusp-cone-pair-of-pants\endcsname{\href{S6_SERIES_GRIND_III.pdf\#nameddest=tag-0377}{III.82 (p.~259)} }
\expandafter\def\csname SRef@wb:cusp-cone-M-definition\endcsname{\href{S6_SERIES_GRIND_III.pdf\#nameddest=tag-035P}{III.83 (p.~259)} }
\expandafter\def\csname SRef@wb:cusp-cone-M-boundary\endcsname{\href{S6_SERIES_GRIND_III.pdf\#nameddest=tag-035N}{III.84 (p.~259)} }
\expandafter\def\csname SRef@wb:cusp-cone-C-definition\endcsname{\href{S6_SERIES_GRIND_III.pdf\#nameddest=tag-035C}{III.85 (p.~259)} }
\expandafter\def\csname SRef@wb:cusp-cone-q1\endcsname{\href{S6_SERIES_GRIND_III.pdf\#nameddest=tag-0379}{III.86 (p.~259)} }
\expandafter\def\csname SRef@wb:cusp-cone-q1-cusp-formula\endcsname{\href{S6_SERIES_GRIND_III.pdf\#nameddest=tag-037A}{III.87 (p.~259)} }
\expandafter\def\csname SRef@wb:cusp-cone-S24-map\endcsname{\href{S6_SERIES_GRIND_III.pdf\#nameddest=tag-0366}{III.88 (p.~260)} }
\expandafter\def\csname SRef@wb:cusp-cone-S24-source\endcsname{\href{S6_SERIES_GRIND_III.pdf\#nameddest=tag-036D}{III.89 (p.~260)} }
\expandafter\def\csname SRef@wb:cusp-cover-relative-factorization\endcsname{\href{S6_SERIES_GRIND_III.pdf\#nameddest=tag-0384}{III.9 (p.~239)} }
\expandafter\def\csname SRef@wb:cusp-cone-S24-fibre-map\endcsname{\href{S6_SERIES_GRIND_III.pdf\#nameddest=tag-0363}{III.90 (p.~260)} }
\expandafter\def\csname SRef@wb:cusp-cone-S24-R4-limit\endcsname{\href{S6_SERIES_GRIND_III.pdf\#nameddest=tag-0360}{III.91 (p.~260)} }
\expandafter\def\csname SRef@wb:cusp-cone-C-homology\endcsname{\href{S6_SERIES_GRIND_III.pdf\#nameddest=tag-035D}{III.92 (p.~260)} }
\expandafter\def\csname SRef@wb:cusp-cone-C-smash-homology\endcsname{\href{S6_SERIES_GRIND_III.pdf\#nameddest=tag-035E}{III.93 (p.~260)} }
\expandafter\def\csname SRef@wb:cusp-cone-Seifert-meridians\endcsname{\href{S6_SERIES_GRIND_III.pdf\#nameddest=tag-036F}{III.94 (p.~261)} }
\expandafter\def\csname SRef@wb:cusp-cone-M-H1\endcsname{\href{S6_SERIES_GRIND_III.pdf\#nameddest=tag-035M}{III.95 (p.~261)} }
\expandafter\def\csname SRef@wb:cusp-cone-M-boundary-map\endcsname{\href{S6_SERIES_GRIND_III.pdf\#nameddest=tag-035O}{III.96 (p.~261)} }
\expandafter\def\csname SRef@wb:cusp-cone-ambient-collapse\endcsname{\href{S6_SERIES_GRIND_III.pdf\#nameddest=tag-036I}{III.97 (p.~262)} }
\expandafter\def\csname SRef@wb:cusp-cone-intrinsic-expansion\endcsname{\href{S6_SERIES_GRIND_III.pdf\#nameddest=tag-036S}{III.98 (p.~262)} }
\expandafter\def\csname SRef@wb:cusp-cone-rank-bound\endcsname{\href{S6_SERIES_GRIND_III.pdf\#nameddest=tag-037H}{III.99 (p.~262)} }
\expandafter\def\csname SRef@maps:src:B.1.1\endcsname{\href{S6_SERIES_GRIND_III.pdf\#nameddest=tag-02J4}{III.3.2 (p.~13)} }
\expandafter\def\csname SRef@maps:src:C.1.2\endcsname{\href{S6_SERIES_GRIND_III.pdf\#nameddest=tag-02K3}{III.4.2 (p.~28)} }
\expandafter\def\csname SRef@maps:src:D.1.10\endcsname{\href{S6_SERIES_GRIND_III.pdf\#nameddest=tag-02K8}{III.5.10 (p.~40)} }
\expandafter\def\csname SRef@wb:h24:ret:chart-lemma\endcsname{\href{S6_SERIES_GRIND_III.pdf\#nameddest=tag-03AE}{III.7.30 (p.~300)} }
\expandafter\def\csname SRef@wb:h24:ret:orientation-lemma\endcsname{\href{S6_SERIES_GRIND_III.pdf\#nameddest=tag-03BD}{III.7.31 (p.~301)} }
\expandafter\def\csname SRef@mapsIII:carrier-three-complex-wedge\endcsname{\href{S6_SERIES_GRIND_III.pdf\#nameddest=tag-02QC}{III.7.34 (p.~306)} }
\expandafter\def\csname SRef@mapsIII:carrier-extension-lemma\endcsname{\href{S6_SERIES_GRIND_III.pdf\#nameddest=tag-02Q8}{III.7.36 (p.~308)} }
\expandafter\def\csname SRef@wb:cusp-cone-native-descent\endcsname{\href{S6_SERIES_GRIND_III.pdf\#nameddest=tag-0370}{III.7.5 (p.~255)} }
\expandafter\def\csname SRef@maps:src:D.2.4\endcsname{\href{S6_SERIES_GRIND_III.pdf\#nameddest=tag-02LG}{III.10.1 (p.~319)} }
\expandafter\def\csname SRef@maps:src:A.2.2\endcsname{\href{S6_SERIES_GRIND_III.pdf\#nameddest=tag-02J3}{III.3.1 (p.~12)} }
\expandafter\def\csname SRef@maps:src:C.1.4\endcsname{\href{S6_SERIES_GRIND_III.pdf\#nameddest=tag-02K5}{III.4.4 (p.~30)} }
\expandafter\def\csname SRef@maps:src:D.2.3\endcsname{\href{S6_SERIES_GRIND_III.pdf\#nameddest=tag-02L5}{III.5.13 (p.~43)} }
\expandafter\def\csname SRef@maps:src:D.2.19\endcsname{\href{S6_SERIES_GRIND_III.pdf\#nameddest=tag-02KT}{III.5.28 (p.~80)} }
\expandafter\def\csname SRef@maps:src:D.1.5\endcsname{\href{S6_SERIES_GRIND_III.pdf\#nameddest=tag-02KE}{III.5.5 (p.~36)} }
\expandafter\def\csname SRef@maps:src:E.1.1\endcsname{\href{S6_SERIES_GRIND_V.pdf\#nameddest=tag-02M4}{V.2.1 (p.~14)} }
\expandafter\def\csname SRef@maps:src:E.2.7\endcsname{\href{S6_SERIES_GRIND_V.pdf\#nameddest=tag-02MR}{V.2.11 (p.~29)} }
\expandafter\def\csname SRef@maps:src:E.2.8\endcsname{\href{S6_SERIES_GRIND_V.pdf\#nameddest=tag-02MS}{V.2.12 (p.~30)} }
\expandafter\def\csname SRef@maps:src:E.2.10\endcsname{\href{S6_SERIES_GRIND_V.pdf\#nameddest=tag-02MA}{V.2.14 (p.~33)} }
\expandafter\def\csname SRef@maps:src:E.2.18\endcsname{\href{S6_SERIES_GRIND_V.pdf\#nameddest=tag-02MI}{V.2.22 (p.~56)} }
\expandafter\def\csname SRef@maps:src:E.2.19\endcsname{\href{S6_SERIES_GRIND_V.pdf\#nameddest=tag-02MJ}{V.2.23 (p.~60)} }
\expandafter\def\csname SRef@maps:src:E.2.20\endcsname{\href{S6_SERIES_GRIND_V.pdf\#nameddest=tag-02ML}{V.2.24 (p.~63)} }
\expandafter\def\csname SRef@maps:src:E.2.4\endcsname{\href{S6_SERIES_GRIND_V.pdf\#nameddest=tag-02MO}{V.2.8 (p.~24)} }
\expandafter\def\csname SRef@maps:src:E.2.5\endcsname{\href{S6_SERIES_GRIND_V.pdf\#nameddest=tag-02MP}{V.2.9 (p.~25)} }
\expandafter\def\csname SRef@wb:cusp-cone-metric-scope\endcsname{\href{S6_SERIES_GRIND_III.pdf\#nameddest=tag-036W}{III.7.10 (p.~262)} }
\expandafter\def\csname SRef@wb:cusp-cone-explicit-carriers\endcsname{\href{S6_SERIES_GRIND_III.pdf\#nameddest=tag-036O}{III.7.17 (p.~273)} }
\expandafter\def\csname SRef@mapsIII:carrier-real-periods\endcsname{\href{S6_SERIES_GRIND_III.pdf\#nameddest=tag-02QB}{III.7.2 (p.~244)} }
\expandafter\def\csname SRef@wb:sphere-metric-tangent-scope\endcsname{\href{S6_SERIES_GRIND_III.pdf\#nameddest=tag-03DF}{III.7.29 (p.~294)} }
\expandafter\def\csname SRef@wb:cusp-spec-native-quotient\endcsname{\href{S6_SERIES_GRIND_III.pdf\#nameddest=tag-039Y}{III.7.3 (p.~245)} }
\expandafter\def\csname SRef@wb:h24:ret:framing-proposition\endcsname{\href{S6_SERIES_GRIND_III.pdf\#nameddest=tag-03BC}{III.7.32 (p.~302)} }
\expandafter\def\csname SRef@mapsaudit:finite-degree-three\endcsname{\href{S6_SERIES_GRIND_III.pdf\#nameddest=tag-02QE}{III.7.35 (p.~307)} }
\expandafter\def\csname SRef@mapsIII:angle-forgotten-one-copy-null\endcsname{\href{S6_SERIES_GRIND_III.pdf\#nameddest=tag-02Q7}{III.7.37 (p.~308)} }
\expandafter\def\csname SRef@wb:cusp-cone-S24-local-route\endcsname{\href{S6_SERIES_GRIND_III.pdf\#nameddest=tag-0365}{III.7.7 (p.~258)} }
\expandafter\def\csname SRef@maps:src:D.3.4\endcsname{\href{S6_SERIES_GRIND_III.pdf\#nameddest=tag-02M3}{III.9.2 (p.~317)} }
\expandafter\def\csname SRef@qa:III:layout:12\endcsname{\href{S6_SERIES_GRIND_III.pdf\#nameddest=tag-033S}{III.3 (p.~13)} }
\expandafter\def\csname SRef@wb:cusp-cover-naive-action-failure\endcsname{\href{S6_SERIES_GRIND_III.pdf\#nameddest=tag-0382}{III.7.1 (p.~243)} }
\expandafter\def\csname SRef@wb:cusp-cover-global-scope\endcsname{\href{S6_SERIES_GRIND_III.pdf\#nameddest=tag-037Y}{III.7.1 (p.~243)} }
\expandafter\def\csname SRef@qa:III:layout:11\endcsname{\href{S6_SERIES_GRIND_III.pdf\#nameddest=tag-033R}{III.{{$\langle 1\rangle $1.}} (p.~271)} }
\expandafter\def\csname SRef@wb:cusp-cone-moving-S4-null-homotopy\endcsname{\href{S6_SERIES_GRIND_III.pdf\#nameddest=tag-036Z}{III.7.3.1 (p.~272)} }
\expandafter\def\csname SRef@wb:cusp-cone-sphere-map-homotopy-status\endcsname{\href{S6_SERIES_GRIND_III.pdf\#nameddest=tag-037K}{III.7.3.1 (p.~273)} }
\expandafter\def\csname SRef@wb:cusp-s12-one-copy-invariant\endcsname{\href{S6_SERIES_GRIND_III.pdf\#nameddest=tag-038X}{III.7.4 (p.~281)} }
\expandafter\def\csname SRef@wb:cusp-s12-nullity-scope\endcsname{\href{S6_SERIES_GRIND_III.pdf\#nameddest=tag-038W}{III.7.4 (p.~281)} }
\expandafter\def\csname SRef@qa:III:layout:1\endcsname{\href{S6_SERIES_GRIND_III.pdf\#nameddest=tag-033Q}{III.{{$\langle 1\rangle $5.}} (p.~283)} }
\expandafter\def\csname SRef@qa:III:layout:2\endcsname{\href{S6_SERIES_GRIND_III.pdf\#nameddest=tag-033T}{III.{{$\langle 1\rangle $6.}} (p.~316)} }
\expandafter\def\csname SRef@qa:III:layout:3\endcsname{\href{S6_SERIES_GRIND_III.pdf\#nameddest=tag-033U}{III.{{$\langle 1\rangle $6.}} (p.~316)} }
\expandafter\def\csname SRef@qa:III:layout:4\endcsname{\href{S6_SERIES_GRIND_III.pdf\#nameddest=tag-033V}{III.8 (p.~316)} }
\expandafter\def\csname SRef@sec:III:conventions\endcsname{\href{S6_SERIES_GRIND_III.pdf\#nameddest=tag-034J}{III.2 (p.~12)} }
\expandafter\def\csname SRef@sec:III:maps\endcsname{\href{S6_SERIES_GRIND_III.pdf\#nameddest=tag-034K}{III.3 (p.~12)} }
\expandafter\def\csname SRef@sec:III:os49-night\endcsname{\href{S6_SERIES_GRIND_III.pdf\#nameddest=tag-034L}{III.6 (p.~215)} }
\expandafter\def\csname SRef@sec:III:carriers\endcsname{\href{S6_SERIES_GRIND_III.pdf\#nameddest=tag-034I}{III.7 (p.~238)} }
\expandafter\def\csname SRef@mapsD:repair:modular-collisions\endcsname{\href{S6_SERIES_GRIND_III.pdf\#nameddest=tag-02OJ}{III.8 (p.~310)} }
\expandafter\def\csname SRef@wb:cusp-cover-section\endcsname{\href{S6_SERIES_GRIND_III.pdf\#nameddest=tag-0385}{III.7.1 (p.~238)} }
\expandafter\def\csname SRef@mapsIII:carrier-radial-prerequisites\endcsname{\href{S6_SERIES_GRIND_III.pdf\#nameddest=tag-02QA}{III.7.2 (p.~244)} }
\expandafter\def\csname SRef@wb:cusp-cone-section\endcsname{\href{S6_SERIES_GRIND_III.pdf\#nameddest=tag-037I}{III.7.3 (p.~255)} }
\expandafter\def\csname SRef@wb:cusp-s12-carrier-homotopy\endcsname{\href{S6_SERIES_GRIND_III.pdf\#nameddest=tag-038U}{III.7.4 (p.~279)} }
\expandafter\def\csname SRef@wb:sphere-metric-limit-note\endcsname{\href{S6_SERIES_GRIND_III.pdf\#nameddest=tag-03D9}{III.7.6 (p.~287)} }
\expandafter\def\csname SRef@wb:h24:ret:section\endcsname{\href{S6_SERIES_GRIND_III.pdf\#nameddest=tag-03BF}{III.7.7 (p.~298)} }
\expandafter\def\csname SRef@mapsIII:carrier-extension-section\endcsname{\href{S6_SERIES_GRIND_III.pdf\#nameddest=tag-02Q9}{III.7.8 (p.~308)} }
\expandafter\def\csname SRef@the-source-section-and-component-divisor\endcsname{\href{S6_SERIES_GRIND_III.pdf\#nameddest=tag-0351}{III.6.0.1 (p.~215)} }
\expandafter\def\csname SRef@the-page-including-the-finite-cap-line\endcsname{\href{S6_SERIES_GRIND_III.pdf\#nameddest=tag-034Z}{III.6.0.10 (p.~231)} }
\expandafter\def\csname SRef@the-actual-graph-line-and-return\endcsname{\href{S6_SERIES_GRIND_III.pdf\#nameddest=tag-034Q}{III.6.0.11 (p.~232)} }
\expandafter\def\csname SRef@all-stabilizers-on-the-actual-root-chart\endcsname{\href{S6_SERIES_GRIND_III.pdf\#nameddest=tag-0008}{III.6.0.12 (p.~233)} }
\expandafter\def\csname SRef@the-full-fundamental-group-presentation\endcsname{\href{S6_SERIES_GRIND_III.pdf\#nameddest=tag-034U}{III.6.0.13 (p.~234)} }
\expandafter\def\csname SRef@the-fixed-exterior-and-the-actual-boundary-circle\endcsname{\href{S6_SERIES_GRIND_III.pdf\#nameddest=tag-034T}{III.6.0.14 (p.~235)} }
\expandafter\def\csname SRef@an-integral-basis-of-the-boundary-two-cohomology\endcsname{\href{S6_SERIES_GRIND_III.pdf\#nameddest=tag-000A}{III.6.0.15 (p.~236)} }
\expandafter\def\csname SRef@determine-the-exterior-restriction-without-guessing-a-splitting\endcsname{\href{S6_SERIES_GRIND_III.pdf\#nameddest=tag-01O2}{III.6.0.16 (p.~236)} }
\expandafter\def\csname SRef@the-two-smith-invariants-and-every-homology-group\endcsname{\href{S6_SERIES_GRIND_III.pdf\#nameddest=tag-0353}{III.6.0.17 (p.~237)} }
\expandafter\def\csname SRef@actual-local-fillings-and-their-gluing\endcsname{\href{S6_SERIES_GRIND_III.pdf\#nameddest=tag-0007}{III.6.0.2 (p.~216)} }
\expandafter\def\csname SRef@marked-monodromy-and-its-completeness\endcsname{\href{S6_SERIES_GRIND_III.pdf\#nameddest=tag-02QG}{III.6.0.3 (p.~217)} }
\expandafter\def\csname SRef@integral-topology\endcsname{\href{S6_SERIES_GRIND_III.pdf\#nameddest=tag-01RP}{III.6.0.4 (p.~218)} }
\expandafter\def\csname SRef@the-exact-line-subfan\endcsname{\href{S6_SERIES_GRIND_III.pdf\#nameddest=tag-034S}{III.6.0.5 (p.~225)} }
\expandafter\def\csname SRef@a-page-carrying-the-actual-source-line\endcsname{\href{S6_SERIES_GRIND_III.pdf\#nameddest=tag-0003}{III.6.0.6 (p.~227)} }
\expandafter\def\csname SRef@the-full-graph-to-page-comparison-by-regular-bands\endcsname{\href{S6_SERIES_GRIND_III.pdf\#nameddest=tag-034V}{III.6.0.7 (p.~228)} }
\expandafter\def\csname SRef@characteristic-calculation-and-the-scope-of-the-result\endcsname{\href{S6_SERIES_GRIND_III.pdf\#nameddest=tag-01GY}{III.6.0.8 (p.~229)} }
\expandafter\def\csname SRef@the-actual-closing-map-and-its-origins\endcsname{\href{S6_SERIES_GRIND_III.pdf\#nameddest=tag-034P}{III.6.0.9 (p.~231)} }
\expandafter\def\csname SRef@wb:cusp-cone-optimal-strengthening\endcsname{\href{S6_SERIES_GRIND_III.pdf\#nameddest=tag-0376}{III.7.3.1 (p.~264)} }
\expandafter\def\csname SRef@maps:src:B.5.2\endcsname{\href{S6_SERIES_GRIND_III.pdf\#nameddest=tag-02JC}{III.3.10 (p.~23)} }
\expandafter\def\csname SRef@maps:src:B.1.2\endcsname{\href{S6_SERIES_GRIND_III.pdf\#nameddest=tag-02J5}{III.3.3 (p.~14)} }
\expandafter\def\csname SRef@maps:src:B.2.1\endcsname{\href{S6_SERIES_GRIND_III.pdf\#nameddest=tag-02J6}{III.3.4 (p.~15)} }
\expandafter\def\csname SRef@maps:src:B.3.1\endcsname{\href{S6_SERIES_GRIND_III.pdf\#nameddest=tag-02J7}{III.3.5 (p.~16)} }
\expandafter\def\csname SRef@maps:src:B.3.2\endcsname{\href{S6_SERIES_GRIND_III.pdf\#nameddest=tag-02J8}{III.3.6 (p.~17)} }
\expandafter\def\csname SRef@maps:src:B.3.3\endcsname{\href{S6_SERIES_GRIND_III.pdf\#nameddest=tag-02J9}{III.3.7 (p.~18)} }
\expandafter\def\csname SRef@maps:src:B.4.1\endcsname{\href{S6_SERIES_GRIND_III.pdf\#nameddest=tag-02JA}{III.3.8 (p.~18)} }
\expandafter\def\csname SRef@maps:src:B.5.1\endcsname{\href{S6_SERIES_GRIND_III.pdf\#nameddest=tag-02JB}{III.3.9 (p.~20)} }
\expandafter\def\csname SRef@maps:src:C.1.1\endcsname{\href{S6_SERIES_GRIND_III.pdf\#nameddest=tag-02K2}{III.4.1 (p.~27)} }
\expandafter\def\csname SRef@maps:src:C.1.3\endcsname{\href{S6_SERIES_GRIND_III.pdf\#nameddest=tag-02K4}{III.4.3 (p.~29)} }
\expandafter\def\csname SRef@maps:src:C.1.5\endcsname{\href{S6_SERIES_GRIND_III.pdf\#nameddest=tag-02K6}{III.4.5 (p.~30)} }
\expandafter\def\csname SRef@maps:src:D.1.1\endcsname{\href{S6_SERIES_GRIND_III.pdf\#nameddest=tag-02K7}{III.5.1 (p.~32)} }
\expandafter\def\csname SRef@maps:src:D.1.12\endcsname{\href{S6_SERIES_GRIND_III.pdf\#nameddest=tag-02KA}{III.5.11 (p.~41)} }
\expandafter\def\csname SRef@maps:src:D.2.5\endcsname{\href{S6_SERIES_GRIND_III.pdf\#nameddest=tag-02LR}{III.5.14 (p.~44)} }
\expandafter\def\csname SRef@maps:src:D.2.6\endcsname{\href{S6_SERIES_GRIND_III.pdf\#nameddest=tag-02LW}{III.5.15 (p.~46)} }
\expandafter\def\csname SRef@maps:src:D.2.7\endcsname{\href{S6_SERIES_GRIND_III.pdf\#nameddest=tag-02LX}{III.5.16 (p.~48)} }
\expandafter\def\csname SRef@maps:src:D.2.8\endcsname{\href{S6_SERIES_GRIND_III.pdf\#nameddest=tag-02LY}{III.5.17 (p.~50)} }
\expandafter\def\csname SRef@maps:src:D.2.9\endcsname{\href{S6_SERIES_GRIND_III.pdf\#nameddest=tag-02LZ}{III.5.18 (p.~53)} }
\expandafter\def\csname SRef@maps:src:D.2.10\endcsname{\href{S6_SERIES_GRIND_III.pdf\#nameddest=tag-02KK}{III.5.19 (p.~55)} }
\expandafter\def\csname SRef@maps:src:D.1.2\endcsname{\href{S6_SERIES_GRIND_III.pdf\#nameddest=tag-02KB}{III.5.2 (p.~32)} }
\expandafter\def\csname SRef@maps:src:D.2.11\endcsname{\href{S6_SERIES_GRIND_III.pdf\#nameddest=tag-02KL}{III.5.20 (p.~59)} }
\expandafter\def\csname SRef@maps:src:D.2.12\endcsname{\href{S6_SERIES_GRIND_III.pdf\#nameddest=tag-02KM}{III.5.21 (p.~62)} }
\expandafter\def\csname SRef@maps:src:D.2.13\endcsname{\href{S6_SERIES_GRIND_III.pdf\#nameddest=tag-02KN}{III.5.22 (p.~65)} }
\expandafter\def\csname SRef@maps:src:D.2.14\endcsname{\href{S6_SERIES_GRIND_III.pdf\#nameddest=tag-02KO}{III.5.23 (p.~68)} }
\expandafter\def\csname SRef@maps:src:D.2.15\endcsname{\href{S6_SERIES_GRIND_III.pdf\#nameddest=tag-02KP}{III.5.24 (p.~70)} }
\expandafter\def\csname SRef@maps:src:D.2.16\endcsname{\href{S6_SERIES_GRIND_III.pdf\#nameddest=tag-02KQ}{III.5.25 (p.~73)} }
\expandafter\def\csname SRef@maps:src:D.2.18\endcsname{\href{S6_SERIES_GRIND_III.pdf\#nameddest=tag-02KS}{III.5.27 (p.~78)} }
\expandafter\def\csname SRef@maps:src:D.2.20\endcsname{\href{S6_SERIES_GRIND_III.pdf\#nameddest=tag-02KV}{III.5.29 (p.~82)} }
\expandafter\def\csname SRef@maps:src:D.1.3\endcsname{\href{S6_SERIES_GRIND_III.pdf\#nameddest=tag-02KC}{III.5.3 (p.~33)} }
\expandafter\def\csname SRef@maps:src:D.2.21\endcsname{\href{S6_SERIES_GRIND_III.pdf\#nameddest=tag-02KW}{III.5.30 (p.~86)} }
\expandafter\def\csname SRef@maps:src:D.2.22\endcsname{\href{S6_SERIES_GRIND_III.pdf\#nameddest=tag-02KX}{III.5.31 (p.~87)} }
\expandafter\def\csname SRef@maps:src:D.2.23\endcsname{\href{S6_SERIES_GRIND_III.pdf\#nameddest=tag-02KY}{III.5.32 (p.~91)} }
\expandafter\def\csname SRef@maps:src:D.2.24\endcsname{\href{S6_SERIES_GRIND_III.pdf\#nameddest=tag-02KZ}{III.5.33 (p.~93)} }
\expandafter\def\csname SRef@maps:src:D.2.25\endcsname{\href{S6_SERIES_GRIND_III.pdf\#nameddest=tag-02L0}{III.5.34 (p.~94)} }
\expandafter\def\csname SRef@maps:src:D.2.26\endcsname{\href{S6_SERIES_GRIND_III.pdf\#nameddest=tag-02L1}{III.5.35 (p.~98)} }
\expandafter\def\csname SRef@maps:src:D.2.27\endcsname{\href{S6_SERIES_GRIND_III.pdf\#nameddest=tag-02L2}{III.5.36 (p.~102)} }
\expandafter\def\csname SRef@maps:src:D.2.29\endcsname{\href{S6_SERIES_GRIND_III.pdf\#nameddest=tag-02L4}{III.5.38 (p.~107)} }
\expandafter\def\csname SRef@maps:src:D.2.30\endcsname{\href{S6_SERIES_GRIND_III.pdf\#nameddest=tag-02L6}{III.5.39 (p.~110)} }
\expandafter\def\csname SRef@maps:src:D.1.4\endcsname{\href{S6_SERIES_GRIND_III.pdf\#nameddest=tag-02KD}{III.5.4 (p.~34)} }
\expandafter\def\csname SRef@maps:src:D.2.31\endcsname{\href{S6_SERIES_GRIND_III.pdf\#nameddest=tag-02L7}{III.5.40 (p.~113)} }
\expandafter\def\csname SRef@maps:src:D.2.32\endcsname{\href{S6_SERIES_GRIND_III.pdf\#nameddest=tag-02L8}{III.5.41 (p.~116)} }
\expandafter\def\csname SRef@maps:src:D.2.33\endcsname{\href{S6_SERIES_GRIND_III.pdf\#nameddest=tag-02L9}{III.5.42 (p.~119)} }
\expandafter\def\csname SRef@maps:src:D.2.34\endcsname{\href{S6_SERIES_GRIND_III.pdf\#nameddest=tag-02LA}{III.5.43 (p.~121)} }
\expandafter\def\csname SRef@maps:src:D.2.35\endcsname{\href{S6_SERIES_GRIND_III.pdf\#nameddest=tag-02LB}{III.5.44 (p.~126)} }
\expandafter\def\csname SRef@maps:src:D.2.36\endcsname{\href{S6_SERIES_GRIND_III.pdf\#nameddest=tag-02LC}{III.5.45 (p.~130)} }
\expandafter\def\csname SRef@maps:src:D.2.37\endcsname{\href{S6_SERIES_GRIND_III.pdf\#nameddest=tag-02LD}{III.5.46 (p.~134)} }
\expandafter\def\csname SRef@maps:src:D.2.38\endcsname{\href{S6_SERIES_GRIND_III.pdf\#nameddest=tag-02LE}{III.5.47 (p.~138)} }
\expandafter\def\csname SRef@maps:src:D.2.39\endcsname{\href{S6_SERIES_GRIND_III.pdf\#nameddest=tag-02LF}{III.5.48 (p.~141)} }
\expandafter\def\csname SRef@maps:src:D.2.40\endcsname{\href{S6_SERIES_GRIND_III.pdf\#nameddest=tag-02LH}{III.5.49 (p.~145)} }
\expandafter\def\csname SRef@maps:src:D.2.41\endcsname{\href{S6_SERIES_GRIND_III.pdf\#nameddest=tag-02LI}{III.5.50 (p.~148)} }
\expandafter\def\csname SRef@maps:src:D.2.42\endcsname{\href{S6_SERIES_GRIND_III.pdf\#nameddest=tag-02LJ}{III.5.51 (p.~151)} }
\expandafter\def\csname SRef@maps:src:D.2.43\endcsname{\href{S6_SERIES_GRIND_III.pdf\#nameddest=tag-02LK}{III.5.52 (p.~153)} }
\expandafter\def\csname SRef@maps:src:D.2.44\endcsname{\href{S6_SERIES_GRIND_III.pdf\#nameddest=tag-02LL}{III.5.53 (p.~156)} }
\expandafter\def\csname SRef@maps:src:D.2.45\endcsname{\href{S6_SERIES_GRIND_III.pdf\#nameddest=tag-02LM}{III.5.54 (p.~161)} }
\expandafter\def\csname SRef@maps:src:D.2.46\endcsname{\href{S6_SERIES_GRIND_III.pdf\#nameddest=tag-02LN}{III.5.55 (p.~167)} }
\expandafter\def\csname SRef@maps:src:D.2.47\endcsname{\href{S6_SERIES_GRIND_III.pdf\#nameddest=tag-02LO}{III.5.56 (p.~172)} }
\expandafter\def\csname SRef@maps:src:D.2.48\endcsname{\href{S6_SERIES_GRIND_III.pdf\#nameddest=tag-02LP}{III.5.57 (p.~175)} }
\expandafter\def\csname SRef@maps:src:D.2.49\endcsname{\href{S6_SERIES_GRIND_III.pdf\#nameddest=tag-02LQ}{III.5.58 (p.~182)} }
\expandafter\def\csname SRef@maps:src:D.2.50\endcsname{\href{S6_SERIES_GRIND_III.pdf\#nameddest=tag-02LS}{III.5.59 (p.~187)} }
\expandafter\def\csname SRef@maps:src:D.1.6\endcsname{\href{S6_SERIES_GRIND_III.pdf\#nameddest=tag-02KF}{III.5.6 (p.~36)} }
\expandafter\def\csname SRef@maps:src:D.2.51\endcsname{\href{S6_SERIES_GRIND_III.pdf\#nameddest=tag-02LT}{III.5.60 (p.~191)} }
\expandafter\def\csname SRef@maps:src:D.2.52\endcsname{\href{S6_SERIES_GRIND_III.pdf\#nameddest=tag-02LU}{III.5.61 (p.~194)} }
\expandafter\def\csname SRef@maps:src:D.2.53\endcsname{\href{S6_SERIES_GRIND_III.pdf\#nameddest=tag-02LV}{III.5.62 (p.~196)} }
\expandafter\def\csname SRef@maps:src:D.3.1\endcsname{\href{S6_SERIES_GRIND_III.pdf\#nameddest=tag-02M0}{III.5.63 (p.~211)} }
\expandafter\def\csname SRef@maps:src:D.3.3\endcsname{\href{S6_SERIES_GRIND_III.pdf\#nameddest=tag-02M2}{III.5.64 (p.~213)} }
\expandafter\def\csname SRef@maps:src:D.1.11\endcsname{\href{S6_SERIES_GRIND_III.pdf\#nameddest=tag-02K9}{III.5.65 (p.~214)} }
\expandafter\def\csname SRef@maps:src:D.1.7\endcsname{\href{S6_SERIES_GRIND_III.pdf\#nameddest=tag-02KG}{III.5.7 (p.~37)} }
\expandafter\def\csname SRef@maps:src:D.1.8\endcsname{\href{S6_SERIES_GRIND_III.pdf\#nameddest=tag-02KH}{III.5.8 (p.~38)} }
\expandafter\def\csname SRef@maps:src:D.1.9\endcsname{\href{S6_SERIES_GRIND_III.pdf\#nameddest=tag-02KI}{III.5.9 (p.~39)} }
\expandafter\def\csname SRef@night:P3F.26\endcsname{\href{S6_SERIES_GRIND_III.pdf\#nameddest=tag-02S2}{III.6.1 (p.~215)} }
\expandafter\def\csname SRef@stmt:P3F.26\endcsname{\href{S6_SERIES_GRIND_III.pdf\#nameddest=tag-02S2}{III.6.1 (p.~215)} }
\expandafter\def\csname SRef@night:P3F.27\endcsname{\href{S6_SERIES_GRIND_III.pdf\#nameddest=tag-02S3}{III.6.2 (p.~220)} }
\expandafter\def\csname SRef@stmt:P3F.27\endcsname{\href{S6_SERIES_GRIND_III.pdf\#nameddest=tag-02S3}{III.6.2 (p.~220)} }
\expandafter\def\csname SRef@night:P3F.28\endcsname{\href{S6_SERIES_GRIND_III.pdf\#nameddest=tag-02S4}{III.6.3 (p.~221)} }
\expandafter\def\csname SRef@stmt:P3F.28\endcsname{\href{S6_SERIES_GRIND_III.pdf\#nameddest=tag-02S4}{III.6.3 (p.~221)} }
\expandafter\def\csname SRef@night:P3F.29\endcsname{\href{S6_SERIES_GRIND_III.pdf\#nameddest=tag-02S5}{III.6.4 (p.~223)} }
\expandafter\def\csname SRef@stmt:P3F.29\endcsname{\href{S6_SERIES_GRIND_III.pdf\#nameddest=tag-02S5}{III.6.4 (p.~223)} }
\expandafter\def\csname SRef@night:P3F.31\endcsname{\href{S6_SERIES_GRIND_III.pdf\#nameddest=tag-02S8}{III.6.5 (p.~225)} }
\expandafter\def\csname SRef@stmt:P3F.31\endcsname{\href{S6_SERIES_GRIND_III.pdf\#nameddest=tag-02S8}{III.6.5 (p.~225)} }
\expandafter\def\csname SRef@night:P3F.32\endcsname{\href{S6_SERIES_GRIND_III.pdf\#nameddest=tag-02S9}{III.6.6 (p.~230)} }
\expandafter\def\csname SRef@stmt:P3F.32\endcsname{\href{S6_SERIES_GRIND_III.pdf\#nameddest=tag-02S9}{III.6.6 (p.~230)} }
\expandafter\def\csname SRef@night:P3F.33\endcsname{\href{S6_SERIES_GRIND_III.pdf\#nameddest=tag-02SA}{III.6.7 (p.~233)} }
\expandafter\def\csname SRef@stmt:P3F.33\endcsname{\href{S6_SERIES_GRIND_III.pdf\#nameddest=tag-02SA}{III.6.7 (p.~233)} }
\expandafter\def\csname SRef@night:P3F.34\endcsname{\href{S6_SERIES_GRIND_III.pdf\#nameddest=tag-02SB}{III.6.8 (p.~234)} }
\expandafter\def\csname SRef@stmt:P3F.34\endcsname{\href{S6_SERIES_GRIND_III.pdf\#nameddest=tag-02SB}{III.6.8 (p.~234)} }
\expandafter\def\csname SRef@maps:src:E.2.6\endcsname{\href{S6_SERIES_GRIND_V.pdf\#nameddest=tag-02MQ}{V.2.10 (p.~26)} }
\expandafter\def\csname SRef@maps:src:E.2.9\endcsname{\href{S6_SERIES_GRIND_V.pdf\#nameddest=tag-02MT}{V.2.13 (p.~31)} }
\expandafter\def\csname SRef@maps:src:E.2.11\endcsname{\href{S6_SERIES_GRIND_V.pdf\#nameddest=tag-02MB}{V.2.15 (p.~35)} }
\expandafter\def\csname SRef@maps:src:E.2.12\endcsname{\href{S6_SERIES_GRIND_V.pdf\#nameddest=tag-02MC}{V.2.16 (p.~37)} }
\expandafter\def\csname SRef@maps:src:E.2.13\endcsname{\href{S6_SERIES_GRIND_V.pdf\#nameddest=tag-02MD}{V.2.17 (p.~40)} }
\expandafter\def\csname SRef@maps:src:E.2.14\endcsname{\href{S6_SERIES_GRIND_V.pdf\#nameddest=tag-02ME}{V.2.18 (p.~42)} }
\expandafter\def\csname SRef@maps:src:E.2.15\endcsname{\href{S6_SERIES_GRIND_V.pdf\#nameddest=tag-02MF}{V.2.19 (p.~45)} }
\expandafter\def\csname SRef@maps:src:E.1.2\endcsname{\href{S6_SERIES_GRIND_V.pdf\#nameddest=tag-02M5}{V.2.2 (p.~14)} }
\expandafter\def\csname SRef@maps:src:E.2.16\endcsname{\href{S6_SERIES_GRIND_V.pdf\#nameddest=tag-02MG}{V.2.20 (p.~48)} }
\expandafter\def\csname SRef@maps:src:E.2.17\endcsname{\href{S6_SERIES_GRIND_V.pdf\#nameddest=tag-02MH}{V.2.21 (p.~52)} }
\expandafter\def\csname SRef@maps:src:E.2.21\endcsname{\href{S6_SERIES_GRIND_V.pdf\#nameddest=tag-02MM}{V.2.25 (p.~65)} }
\expandafter\def\csname SRef@maps:src:E.1.3\endcsname{\href{S6_SERIES_GRIND_V.pdf\#nameddest=tag-02M6}{V.2.3 (p.~16)} }
\expandafter\def\csname SRef@maps:src:E.1.4\endcsname{\href{S6_SERIES_GRIND_V.pdf\#nameddest=tag-02M7}{V.2.4 (p.~18)} }
\expandafter\def\csname SRef@maps:src:E.1.5\endcsname{\href{S6_SERIES_GRIND_V.pdf\#nameddest=tag-02M8}{V.2.5 (p.~19)} }
\expandafter\def\csname SRef@maps:src:E.2.1\endcsname{\href{S6_SERIES_GRIND_V.pdf\#nameddest=tag-02M9}{V.2.6 (p.~21)} }
\expandafter\def\csname SRef@maps:src:E.2.3\endcsname{\href{S6_SERIES_GRIND_V.pdf\#nameddest=tag-02MN}{V.2.7 (p.~23)} }
\expandafter\def\csname SRef@wb:cusp-cover-extension-theorem\endcsname{\href{S6_SERIES_GRIND_III.pdf\#nameddest=tag-037X}{III.7.1 (p.~238)} }
\expandafter\def\csname SRef@wb:cusp-cone-obstructions\endcsname{\href{S6_SERIES_GRIND_III.pdf\#nameddest=tag-0374}{III.7.11 (p.~262)} }
\expandafter\def\csname SRef@wb:cusp-cone-optimal-exponents\endcsname{\href{S6_SERIES_GRIND_III.pdf\#nameddest=tag-0375}{III.7.12 (p.~265)} }
\expandafter\def\csname SRef@wb:cusp-cone-angle-forgotten-sphere-carriers\endcsname{\href{S6_SERIES_GRIND_III.pdf\#nameddest=tag-036J}{III.7.13 (p.~266)} }
\expandafter\def\csname SRef@wb:cusp-cone-target-homotopy-types\endcsname{\href{S6_SERIES_GRIND_III.pdf\#nameddest=tag-037M}{III.7.18 (p.~274)} }
\expandafter\def\csname SRef@wb:cusp-cone-exact-decompactification-space\endcsname{\href{S6_SERIES_GRIND_III.pdf\#nameddest=tag-036N}{III.7.19 (p.~276)} }
\expandafter\def\csname SRef@wb:cusp-s12-full-carrier-null\endcsname{\href{S6_SERIES_GRIND_III.pdf\#nameddest=tag-038V}{III.7.21 (p.~280)} }
\expandafter\def\csname SRef@wb:cusp-null:thm\endcsname{\href{S6_SERIES_GRIND_III.pdf\#nameddest=tag-038S}{III.7.23 (p.~282)} }
\expandafter\def\csname SRef@wb:sphere-metric-uniform-collapse\endcsname{\href{S6_SERIES_GRIND_III.pdf\#nameddest=tag-03DG}{III.7.25 (p.~288)} }
\expandafter\def\csname SRef@wb:sphere-metric-fixed-realization\endcsname{\href{S6_SERIES_GRIND_III.pdf\#nameddest=tag-03D6}{III.7.26 (p.~288)} }
\expandafter\def\csname SRef@wb:sphere-metric-scaling-trichotomy\endcsname{\href{S6_SERIES_GRIND_III.pdf\#nameddest=tag-03DC}{III.7.27 (p.~291)} }
\expandafter\def\csname SRef@wb:cusp-spec-globalization-obstruction\endcsname{\href{S6_SERIES_GRIND_III.pdf\#nameddest=tag-039U}{III.7.4 (p.~246)} }
\expandafter\def\csname SRef@wb:cusp-cone-S12-route\endcsname{\href{S6_SERIES_GRIND_III.pdf\#nameddest=tag-035X}{III.7.6 (p.~256)} }
\expandafter\def\csname SRef@wb:cusp-cone-S24-route\endcsname{\href{S6_SERIES_GRIND_III.pdf\#nameddest=tag-036A}{III.7.8 (p.~259)} }
\expandafter\def\csname SRef@mapsD:repair:collision-realization\endcsname{\href{S6_SERIES_GRIND_III.pdf\#nameddest=tag-02OG}{III.8.1 (p.~310)} }
\expandafter\def\csname SRef@maps:src:D.3.2\endcsname{\href{S6_SERIES_GRIND_III.pdf\#nameddest=tag-02M1}{III.9.1 (p.~316)} }
\expandafter\def\csname SRef@maps:src:E.2.2\endcsname{\href{S6_SERIES_GRIND_III.pdf\#nameddest=tag-02MK}{III.9.3 (p.~317)} }
\expandafter\def\csname SRef@maps:eq:B.6.2\endcsname{\href{S6_SERIES_GRIND_IV.pdf\#nameddest=tag-01VZ}{IV.{B.6.2} (p.~68)} }
\expandafter\def\csname SRef@maps:eq:B.6.3\endcsname{\href{S6_SERIES_GRIND_IV.pdf\#nameddest=tag-01WA}{IV.{B.6.3} (p.~68)} }
\expandafter\def\csname SRef@maps:eq:B.6.4\endcsname{\href{S6_SERIES_GRIND_IV.pdf\#nameddest=tag-01WL}{IV.{B.6.4} (p.~69)} }
\expandafter\def\csname SRef@maps:eq:B.6.5\endcsname{\href{S6_SERIES_GRIND_IV.pdf\#nameddest=tag-01WR}{IV.{B.6.5} (p.~69)} }
\expandafter\def\csname SRef@maps:eq:B.6.6\endcsname{\href{S6_SERIES_GRIND_IV.pdf\#nameddest=tag-01WS}{IV.{B.6.6} (p.~69)} }
\expandafter\def\csname SRef@maps:eq:B.6.7\endcsname{\href{S6_SERIES_GRIND_IV.pdf\#nameddest=tag-01WT}{IV.{B.6.7} (p.~70)} }
\expandafter\def\csname SRef@maps:eq:B.6.8\endcsname{\href{S6_SERIES_GRIND_IV.pdf\#nameddest=tag-01WU}{IV.{B.6.8} (p.~70)} }
\expandafter\def\csname SRef@maps:eq:B.6.9\endcsname{\href{S6_SERIES_GRIND_IV.pdf\#nameddest=tag-01WV}{IV.{B.6.9} (p.~70)} }
\expandafter\def\csname SRef@maps:eq:B.6.10\endcsname{\href{S6_SERIES_GRIND_IV.pdf\#nameddest=tag-01VQ}{IV.{B.6.10} (p.~70)} }
\expandafter\def\csname SRef@maps:eq:B.6.11\endcsname{\href{S6_SERIES_GRIND_IV.pdf\#nameddest=tag-01VR}{IV.{B.6.11} (p.~72)} }
\expandafter\def\csname SRef@maps:eq:B.6.12\endcsname{\href{S6_SERIES_GRIND_IV.pdf\#nameddest=tag-01VS}{IV.{B.6.12} (p.~72)} }
\expandafter\def\csname SRef@maps:eq:B.6.13\endcsname{\href{S6_SERIES_GRIND_IV.pdf\#nameddest=tag-01VT}{IV.{B.6.13} (p.~73)} }
\expandafter\def\csname SRef@maps:eq:B.6.14\endcsname{\href{S6_SERIES_GRIND_IV.pdf\#nameddest=tag-01VU}{IV.{B.6.14} (p.~73)} }
\expandafter\def\csname SRef@maps:eq:B.6.15\endcsname{\href{S6_SERIES_GRIND_IV.pdf\#nameddest=tag-01VV}{IV.{B.6.15} (p.~75)} }
\expandafter\def\csname SRef@maps:eq:B.6.16\endcsname{\href{S6_SERIES_GRIND_IV.pdf\#nameddest=tag-01VW}{IV.{B.6.16} (p.~75)} }
\expandafter\def\csname SRef@maps:eq:B.6.17\endcsname{\href{S6_SERIES_GRIND_IV.pdf\#nameddest=tag-01VX}{IV.{B.6.17} (p.~76)} }
\expandafter\def\csname SRef@maps:eq:B.C1\endcsname{\href{S6_SERIES_GRIND_IV.pdf\#nameddest=tag-01WW}{IV.{B.C1} (p.~77)} }
\expandafter\def\csname SRef@maps:eq:B.C2\endcsname{\href{S6_SERIES_GRIND_IV.pdf\#nameddest=tag-01WX}{IV.{B.C2} (p.~77)} }
\expandafter\def\csname SRef@maps:eq:B.C3\endcsname{\href{S6_SERIES_GRIND_IV.pdf\#nameddest=tag-01WY}{IV.{B.C3} (p.~77)} }
\expandafter\def\csname SRef@maps:eq:B.C4\endcsname{\href{S6_SERIES_GRIND_IV.pdf\#nameddest=tag-01WZ}{IV.{B.C4} (p.~78)} }
\expandafter\def\csname SRef@maps:eq:B.P1\endcsname{\href{S6_SERIES_GRIND_IV.pdf\#nameddest=tag-01X0}{IV.{B.P1} (p.~78)} }
\expandafter\def\csname SRef@maps:eq:B.P2\endcsname{\href{S6_SERIES_GRIND_IV.pdf\#nameddest=tag-01X1}{IV.{B.P2} (p.~79)} }
\expandafter\def\csname SRef@maps:eq:B.P3\endcsname{\href{S6_SERIES_GRIND_IV.pdf\#nameddest=tag-01X2}{IV.{B.P3} (p.~79)} }
\expandafter\def\csname SRef@maps:eq:B.P4\endcsname{\href{S6_SERIES_GRIND_IV.pdf\#nameddest=tag-01X3}{IV.{B.P4} (p.~79)} }
\expandafter\def\csname SRef@maps:eq:B.6.19\endcsname{\href{S6_SERIES_GRIND_IV.pdf\#nameddest=tag-01VY}{IV.{B.6.19} (p.~80)} }
\expandafter\def\csname SRef@maps:eq:B.6.20\endcsname{\href{S6_SERIES_GRIND_IV.pdf\#nameddest=tag-01W0}{IV.{B.6.20} (p.~81)} }
\expandafter\def\csname SRef@maps:eq:B.6.21\endcsname{\href{S6_SERIES_GRIND_IV.pdf\#nameddest=tag-01W1}{IV.{B.6.21} (p.~83)} }
\expandafter\def\csname SRef@maps:eq:B.6.22\endcsname{\href{S6_SERIES_GRIND_IV.pdf\#nameddest=tag-01W2}{IV.{B.6.22} (p.~83)} }
\expandafter\def\csname SRef@maps:eq:B.6.23\endcsname{\href{S6_SERIES_GRIND_IV.pdf\#nameddest=tag-01W3}{IV.{B.6.23} (p.~83)} }
\expandafter\def\csname SRef@maps:eq:B.6.24\endcsname{\href{S6_SERIES_GRIND_IV.pdf\#nameddest=tag-01W4}{IV.{B.6.24} (p.~84)} }
\expandafter\def\csname SRef@maps:eq:B.6.25\endcsname{\href{S6_SERIES_GRIND_IV.pdf\#nameddest=tag-01W5}{IV.{B.6.25} (p.~84)} }
\expandafter\def\csname SRef@maps:eq:B.T1\endcsname{\href{S6_SERIES_GRIND_IV.pdf\#nameddest=tag-01XD}{IV.{B.T1} (p.~87)} }
\expandafter\def\csname SRef@maps:eq:B.6.26\endcsname{\href{S6_SERIES_GRIND_IV.pdf\#nameddest=tag-01W6}{IV.{B.6.26} (p.~87)} }
\expandafter\def\csname SRef@maps:eq:B.6.27\endcsname{\href{S6_SERIES_GRIND_IV.pdf\#nameddest=tag-01W7}{IV.{B.6.27} (p.~88)} }
\expandafter\def\csname SRef@maps:eq:B.6.28\endcsname{\href{S6_SERIES_GRIND_IV.pdf\#nameddest=tag-01W8}{IV.{B.6.28} (p.~89)} }
\expandafter\def\csname SRef@maps:eq:B.6.29\endcsname{\href{S6_SERIES_GRIND_IV.pdf\#nameddest=tag-01W9}{IV.{B.6.29} (p.~89)} }
\expandafter\def\csname SRef@maps:eq:B.6.30\endcsname{\href{S6_SERIES_GRIND_IV.pdf\#nameddest=tag-01WB}{IV.{B.6.30} (p.~89)} }
\expandafter\def\csname SRef@maps:eq:B.6.31\endcsname{\href{S6_SERIES_GRIND_IV.pdf\#nameddest=tag-01WC}{IV.{B.6.31} (p.~90)} }
\expandafter\def\csname SRef@maps:eq:B.Q1\endcsname{\href{S6_SERIES_GRIND_IV.pdf\#nameddest=tag-01X4}{IV.{B.Q1} (p.~91)} }
\expandafter\def\csname SRef@maps:eq:B.Q2\endcsname{\href{S6_SERIES_GRIND_IV.pdf\#nameddest=tag-01X5}{IV.{B.Q2} (p.~91)} }
\expandafter\def\csname SRef@maps:eq:B.Q3\endcsname{\href{S6_SERIES_GRIND_IV.pdf\#nameddest=tag-01X6}{IV.{B.Q3} (p.~91)} }
\expandafter\def\csname SRef@maps:eq:B.Q4\endcsname{\href{S6_SERIES_GRIND_IV.pdf\#nameddest=tag-01X7}{IV.{B.Q4} (p.~92)} }
\expandafter\def\csname SRef@maps:eq:B.Q5\endcsname{\href{S6_SERIES_GRIND_IV.pdf\#nameddest=tag-01X8}{IV.{B.Q5} (p.~92)} }
\expandafter\def\csname SRef@maps:eq:B.Q6\endcsname{\href{S6_SERIES_GRIND_IV.pdf\#nameddest=tag-01X9}{IV.{B.Q6} (p.~92)} }
\expandafter\def\csname SRef@maps:eq:B.Q7\endcsname{\href{S6_SERIES_GRIND_IV.pdf\#nameddest=tag-01XA}{IV.{B.Q7} (p.~92)} }
\expandafter\def\csname SRef@maps:eq:B.Q8\endcsname{\href{S6_SERIES_GRIND_IV.pdf\#nameddest=tag-01XB}{IV.{B.Q8} (p.~92)} }
\expandafter\def\csname SRef@maps:eq:B.Q9\endcsname{\href{S6_SERIES_GRIND_IV.pdf\#nameddest=tag-01XC}{IV.{B.Q9} (p.~93)} }
\expandafter\def\csname SRef@maps:eq:B.6.32\endcsname{\href{S6_SERIES_GRIND_IV.pdf\#nameddest=tag-01WD}{IV.{B.6.32} (p.~93)} }
\expandafter\def\csname SRef@maps:eq:B.6.33\endcsname{\href{S6_SERIES_GRIND_IV.pdf\#nameddest=tag-01WE}{IV.{B.6.33} (p.~94)} }
\expandafter\def\csname SRef@maps:eq:B.6.34\endcsname{\href{S6_SERIES_GRIND_IV.pdf\#nameddest=tag-01WF}{IV.{B.6.34} (p.~94)} }
\expandafter\def\csname SRef@famIII:eq:E9\endcsname{\href{S6_SERIES_GRIND_IV.pdf\#nameddest=tag-01RJ}{IV.{E9} (p.~20)} }
\expandafter\def\csname SRef@maps:eq:B.6.35\endcsname{\href{S6_SERIES_GRIND_IV.pdf\#nameddest=tag-01WG}{IV.{B.6.35} (p.~95)} }
\expandafter\def\csname SRef@maps:eq:B.6.36\endcsname{\href{S6_SERIES_GRIND_IV.pdf\#nameddest=tag-01WH}{IV.{B.6.36} (p.~95)} }
\expandafter\def\csname SRef@maps:eq:B.6.37\endcsname{\href{S6_SERIES_GRIND_IV.pdf\#nameddest=tag-01WI}{IV.{B.6.37} (p.~95)} }
\expandafter\def\csname SRef@maps:eq:B.6.38\endcsname{\href{S6_SERIES_GRIND_IV.pdf\#nameddest=tag-01WJ}{IV.{B.6.38} (p.~99)} }
\expandafter\def\csname SRef@maps:eq:B.6.39\endcsname{\href{S6_SERIES_GRIND_IV.pdf\#nameddest=tag-01WK}{IV.{B.6.39} (p.~100)} }
\expandafter\def\csname SRef@famIII:eq:E10\endcsname{\href{S6_SERIES_GRIND_IV.pdf\#nameddest=tag-01RI}{IV.{E10} (p.~20)} }
\expandafter\def\csname SRef@maps:eq:B.6.40\endcsname{\href{S6_SERIES_GRIND_IV.pdf\#nameddest=tag-01WM}{IV.{B.6.40} (p.~100)} }
\expandafter\def\csname SRef@maps:eq:B.6.41\endcsname{\href{S6_SERIES_GRIND_IV.pdf\#nameddest=tag-01WN}{IV.{B.6.41} (p.~100)} }
\expandafter\def\csname SRef@maps:eq:B.6.42\endcsname{\href{S6_SERIES_GRIND_IV.pdf\#nameddest=tag-01WO}{IV.{B.6.42} (p.~100)} }
\expandafter\def\csname SRef@maps:eq:B.6.43\endcsname{\href{S6_SERIES_GRIND_IV.pdf\#nameddest=tag-01WP}{IV.{B.6.43} (p.~103)} }
\expandafter\def\csname SRef@maps:eq:B.6.44\endcsname{\href{S6_SERIES_GRIND_IV.pdf\#nameddest=tag-01WQ}{IV.{B.6.44} (p.~103)} }
\expandafter\def\csname SRef@maps:eq:B.6.1\endcsname{\href{S6_SERIES_GRIND_IV.pdf\#nameddest=tag-01VP}{IV.{B.6.1} (p.~68)} }
\expandafter\def\csname SRef@push:src:C.3.15\endcsname{\href{S6_SERIES_GRIND_IV.pdf\#nameddest=tag-031M}{IV.2.16 (p.~40)} }
\expandafter\def\csname SRef@push:src:C.3.27\endcsname{\href{S6_SERIES_GRIND_IV.pdf\#nameddest=tag-031Z}{IV.2.28 (p.~63)} }
\expandafter\def\csname SRef@push:src:C.5.1\endcsname{\href{S6_SERIES_GRIND_IV.pdf\#nameddest=tag-0329}{IV.2.31 (p.~65)} }
\expandafter\def\csname SRef@maps:src:B.6.8\endcsname{\href{S6_SERIES_GRIND_IV.pdf\#nameddest=tag-02K0}{IV.2.39 (p.~76)} }
\expandafter\def\csname SRef@maps:src:B.6.25\endcsname{\href{S6_SERIES_GRIND_IV.pdf\#nameddest=tag-02JU}{IV.2.55 (p.~104)} }
\expandafter\def\csname SRef@push:eq:C12-u-flow\endcsname{\href{S6_SERIES_GRIND_IV.pdf\#nameddest=tag-02YV}{IV.1 (p.~33)} }
\expandafter\def\csname SRef@push:eq:C12-limiting-identification\endcsname{\href{S6_SERIES_GRIND_IV.pdf\#nameddest=tag-02YU}{IV.2 (p.~34)} }
\expandafter\def\csname SRef@push:eq:C13-real-overlap\endcsname{\href{S6_SERIES_GRIND_IV.pdf\#nameddest=tag-02YX}{IV.3 (p.~35)} }
\expandafter\def\csname SRef@push:eq:C13-weighted-map\endcsname{\href{S6_SERIES_GRIND_IV.pdf\#nameddest=tag-02YY}{IV.4 (p.~35)} }
\expandafter\def\csname SRef@push:eq:C13-cusp-gauge\endcsname{\href{S6_SERIES_GRIND_IV.pdf\#nameddest=tag-02YW}{IV.5 (p.~37)} }
\expandafter\def\csname SRef@push:src:C.3.9\endcsname{\href{S6_SERIES_GRIND_IV.pdf\#nameddest=tag-0326}{IV.2.10 (p.~29)} }
\expandafter\def\csname SRef@push:src:C.3.22\endcsname{\href{S6_SERIES_GRIND_IV.pdf\#nameddest=tag-031U}{IV.2.23 (p.~54)} }
\expandafter\def\csname SRef@push:src:C.3.8\endcsname{\href{S6_SERIES_GRIND_IV.pdf\#nameddest=tag-0325}{IV.2.9 (p.~28)} }
\expandafter\def\csname SRef@iv:rank_two_specialization\endcsname{\href{S6_SERIES_GRIND_IV.pdf\#nameddest=tag-01UU}{IV.3.20 (p.~131)} }
\expandafter\def\csname SRef@fam:src:E.2.2\endcsname{\href{S6_SERIES_GRIND_IV.pdf\#nameddest=tag-01R7}{IV.2.1 (p.~20)} }
\expandafter\def\csname SRef@push:src:C.3.10\endcsname{\href{S6_SERIES_GRIND_IV.pdf\#nameddest=tag-031H}{IV.2.11 (p.~30)} }
\expandafter\def\csname SRef@push:src:C.3.11\endcsname{\href{S6_SERIES_GRIND_IV.pdf\#nameddest=tag-031I}{IV.2.12 (p.~32)} }
\expandafter\def\csname SRef@push:src:C.3.12\endcsname{\href{S6_SERIES_GRIND_IV.pdf\#nameddest=tag-031J}{IV.2.13 (p.~32)} }
\expandafter\def\csname SRef@push:src:C.3.13\endcsname{\href{S6_SERIES_GRIND_IV.pdf\#nameddest=tag-031K}{IV.2.14 (p.~34)} }
\expandafter\def\csname SRef@push:src:C.3.14\endcsname{\href{S6_SERIES_GRIND_IV.pdf\#nameddest=tag-031L}{IV.2.15 (p.~37)} }
\expandafter\def\csname SRef@push:src:C.3.16\endcsname{\href{S6_SERIES_GRIND_IV.pdf\#nameddest=tag-031N}{IV.2.17 (p.~40)} }
\expandafter\def\csname SRef@push:src:C.3.17\endcsname{\href{S6_SERIES_GRIND_IV.pdf\#nameddest=tag-031O}{IV.2.18 (p.~41)} }
\expandafter\def\csname SRef@push:src:C.3.18\endcsname{\href{S6_SERIES_GRIND_IV.pdf\#nameddest=tag-031P}{IV.2.19 (p.~44)} }
\expandafter\def\csname SRef@push:src:C.3.19\endcsname{\href{S6_SERIES_GRIND_IV.pdf\#nameddest=tag-031Q}{IV.2.20 (p.~47)} }
\expandafter\def\csname SRef@push:src:C.3.20\endcsname{\href{S6_SERIES_GRIND_IV.pdf\#nameddest=tag-031S}{IV.2.21 (p.~49)} }
\expandafter\def\csname SRef@push:src:C.3.21\endcsname{\href{S6_SERIES_GRIND_IV.pdf\#nameddest=tag-031T}{IV.2.22 (p.~51)} }
\expandafter\def\csname SRef@push:src:C.3.23\endcsname{\href{S6_SERIES_GRIND_IV.pdf\#nameddest=tag-031V}{IV.2.24 (p.~55)} }
\expandafter\def\csname SRef@push:src:C.2.2\endcsname{\href{S6_SERIES_GRIND_IV.pdf\#nameddest=tag-031E}{IV.2.3 (p.~22)} }
\expandafter\def\csname SRef@push:src:C.3.25\endcsname{\href{S6_SERIES_GRIND_IV.pdf\#nameddest=tag-031X}{IV.2.26 (p.~58)} }
\expandafter\def\csname SRef@push:src:C.3.26\endcsname{\href{S6_SERIES_GRIND_IV.pdf\#nameddest=tag-031Y}{IV.2.27 (p.~61)} }
\expandafter\def\csname SRef@push:src:C.4.1\endcsname{\href{S6_SERIES_GRIND_IV.pdf\#nameddest=tag-0327}{IV.2.29 (p.~64)} }
\expandafter\def\csname SRef@push:src:C.4.2\endcsname{\href{S6_SERIES_GRIND_IV.pdf\#nameddest=tag-0328}{IV.2.30 (p.~65)} }
\expandafter\def\csname SRef@push:src:C.2.3\endcsname{\href{S6_SERIES_GRIND_IV.pdf\#nameddest=tag-031F}{IV.2.4 (p.~23)} }
\expandafter\def\csname SRef@maps:src:B.6.1\endcsname{\href{S6_SERIES_GRIND_IV.pdf\#nameddest=tag-02JD}{IV.2.32 (p.~66)} }
\expandafter\def\csname SRef@maps:src:B.6.3\endcsname{\href{S6_SERIES_GRIND_IV.pdf\#nameddest=tag-02JV}{IV.2.34 (p.~67)} }
\expandafter\def\csname SRef@maps:src:B.6.4\endcsname{\href{S6_SERIES_GRIND_IV.pdf\#nameddest=tag-02JW}{IV.2.35 (p.~69)} }
\expandafter\def\csname SRef@maps:src:B.6.5\endcsname{\href{S6_SERIES_GRIND_IV.pdf\#nameddest=tag-02JX}{IV.2.36 (p.~71)} }
\expandafter\def\csname SRef@maps:src:B.6.6\endcsname{\href{S6_SERIES_GRIND_IV.pdf\#nameddest=tag-02JY}{IV.2.37 (p.~73)} }
\expandafter\def\csname SRef@maps:src:B.6.7\endcsname{\href{S6_SERIES_GRIND_IV.pdf\#nameddest=tag-02JZ}{IV.2.38 (p.~75)} }
\expandafter\def\csname SRef@maps:src:B.6.10\endcsname{\href{S6_SERIES_GRIND_IV.pdf\#nameddest=tag-02JE}{IV.2.40 (p.~77)} }
\expandafter\def\csname SRef@maps:src:B.6.11\endcsname{\href{S6_SERIES_GRIND_IV.pdf\#nameddest=tag-02JF}{IV.2.41 (p.~78)} }
\expandafter\def\csname SRef@push:src:C.3.2\endcsname{\href{S6_SERIES_GRIND_IV.pdf\#nameddest=tag-031R}{IV.2.5 (p.~25)} }
\expandafter\def\csname SRef@maps:src:B.6.12\endcsname{\href{S6_SERIES_GRIND_IV.pdf\#nameddest=tag-02JG}{IV.2.42 (p.~80)} }
\expandafter\def\csname SRef@maps:src:B.6.13\endcsname{\href{S6_SERIES_GRIND_IV.pdf\#nameddest=tag-02JH}{IV.2.43 (p.~83)} }
\expandafter\def\csname SRef@maps:src:B.6.14\endcsname{\href{S6_SERIES_GRIND_IV.pdf\#nameddest=tag-02JI}{IV.2.44 (p.~86)} }
\expandafter\def\csname SRef@maps:src:B.6.15\endcsname{\href{S6_SERIES_GRIND_IV.pdf\#nameddest=tag-02JJ}{IV.2.45 (p.~87)} }
\expandafter\def\csname SRef@maps:src:B.6.16\endcsname{\href{S6_SERIES_GRIND_IV.pdf\#nameddest=tag-02JK}{IV.2.46 (p.~88)} }
\expandafter\def\csname SRef@maps:src:B.6.17\endcsname{\href{S6_SERIES_GRIND_IV.pdf\#nameddest=tag-02JL}{IV.2.47 (p.~91)} }
\expandafter\def\csname SRef@maps:src:B.6.18\endcsname{\href{S6_SERIES_GRIND_IV.pdf\#nameddest=tag-02JM}{IV.2.48 (p.~93)} }
\expandafter\def\csname SRef@maps:src:B.6.19\endcsname{\href{S6_SERIES_GRIND_IV.pdf\#nameddest=tag-02JN}{IV.2.49 (p.~94)} }
\expandafter\def\csname SRef@push:src:C.3.4\endcsname{\href{S6_SERIES_GRIND_IV.pdf\#nameddest=tag-0321}{IV.2.6 (p.~25)} }
\expandafter\def\csname SRef@maps:src:B.6.22\endcsname{\href{S6_SERIES_GRIND_IV.pdf\#nameddest=tag-02JR}{IV.2.52 (p.~99)} }
\expandafter\def\csname SRef@push:src:C.3.6\endcsname{\href{S6_SERIES_GRIND_IV.pdf\#nameddest=tag-0323}{IV.2.7 (p.~25)} }
\expandafter\def\csname SRef@push:src:C.3.7\endcsname{\href{S6_SERIES_GRIND_IV.pdf\#nameddest=tag-0324}{IV.2.8 (p.~26)} }
\expandafter\def\csname SRef@iv:correction:P3F.25\endcsname{\href{S6_SERIES_GRIND_IV.pdf\#nameddest=tag-01TX}{IV.8.1 (p.~933)} }
\expandafter\def\csname SRef@iv:correction:P3K.29\endcsname{\href{S6_SERIES_GRIND_IV.pdf\#nameddest=tag-01U6}{IV.8.10 (p.~937)} }
\expandafter\def\csname SRef@iv:correction:P3K.30\endcsname{\href{S6_SERIES_GRIND_IV.pdf\#nameddest=tag-01U7}{IV.8.11 (p.~938)} }
\expandafter\def\csname SRef@iv:correction:P3K.31\endcsname{\href{S6_SERIES_GRIND_IV.pdf\#nameddest=tag-01U8}{IV.8.12 (p.~938)} }
\expandafter\def\csname SRef@iv:correction:P3K.32\endcsname{\href{S6_SERIES_GRIND_IV.pdf\#nameddest=tag-01U9}{IV.8.13 (p.~939)} }
\expandafter\def\csname SRef@iv:correction:P3K.33\endcsname{\href{S6_SERIES_GRIND_IV.pdf\#nameddest=tag-01UA}{IV.8.14 (p.~939)} }
\expandafter\def\csname SRef@iv:correction:P3K.34\endcsname{\href{S6_SERIES_GRIND_IV.pdf\#nameddest=tag-01UB}{IV.8.15 (p.~940)} }
\expandafter\def\csname SRef@iv:correction:P3F.53\endcsname{\href{S6_SERIES_GRIND_IV.pdf\#nameddest=tag-01TY}{IV.8.16 (p.~940)} }
\expandafter\def\csname SRef@iv:correction:P3K.35\endcsname{\href{S6_SERIES_GRIND_IV.pdf\#nameddest=tag-01UC}{IV.8.17 (p.~940)} }
\expandafter\def\csname SRef@iv:correction:P3K.36\endcsname{\href{S6_SERIES_GRIND_IV.pdf\#nameddest=tag-01UD}{IV.8.18 (p.~941)} }
\expandafter\def\csname SRef@iv:correction:P3K.37\endcsname{\href{S6_SERIES_GRIND_IV.pdf\#nameddest=tag-01UE}{IV.8.19 (p.~941)} }
\expandafter\def\csname SRef@iv:correction:P3K.6\endcsname{\href{S6_SERIES_GRIND_IV.pdf\#nameddest=tag-01US}{IV.8.2 (p.~934)} }
\expandafter\def\csname SRef@iv:correction:P3K.40\endcsname{\href{S6_SERIES_GRIND_IV.pdf\#nameddest=tag-01UF}{IV.8.20 (p.~942)} }
\expandafter\def\csname SRef@iv:correction:P3K.41\endcsname{\href{S6_SERIES_GRIND_IV.pdf\#nameddest=tag-01UG}{IV.8.21 (p.~942)} }
\expandafter\def\csname SRef@iv:correction:P3K.43\endcsname{\href{S6_SERIES_GRIND_IV.pdf\#nameddest=tag-01UH}{IV.8.22 (p.~943)} }
\expandafter\def\csname SRef@iv:correction:P3K.44\endcsname{\href{S6_SERIES_GRIND_IV.pdf\#nameddest=tag-01UI}{IV.8.23 (p.~943)} }
\expandafter\def\csname SRef@iv:correction:P3K.45\endcsname{\href{S6_SERIES_GRIND_IV.pdf\#nameddest=tag-01UJ}{IV.8.24 (p.~944)} }
\expandafter\def\csname SRef@iv:correction:P3K.46\endcsname{\href{S6_SERIES_GRIND_IV.pdf\#nameddest=tag-01UK}{IV.8.25 (p.~944)} }
\expandafter\def\csname SRef@iv:correction:P3K.50\endcsname{\href{S6_SERIES_GRIND_IV.pdf\#nameddest=tag-01UL}{IV.8.26 (p.~944)} }
\expandafter\def\csname SRef@iv:correction:P3K.52\endcsname{\href{S6_SERIES_GRIND_IV.pdf\#nameddest=tag-01UM}{IV.8.27 (p.~944)} }
\expandafter\def\csname SRef@iv:correction:P3K.54\endcsname{\href{S6_SERIES_GRIND_IV.pdf\#nameddest=tag-01UN}{IV.8.28 (p.~945)} }
\expandafter\def\csname SRef@iv:correction:P3K.56\endcsname{\href{S6_SERIES_GRIND_IV.pdf\#nameddest=tag-01UO}{IV.8.29 (p.~945)} }
\expandafter\def\csname SRef@iv:correction:P3K.10\endcsname{\href{S6_SERIES_GRIND_IV.pdf\#nameddest=tag-01TZ}{IV.8.3 (p.~934)} }
\expandafter\def\csname SRef@iv:correction:P3K.57\endcsname{\href{S6_SERIES_GRIND_IV.pdf\#nameddest=tag-01UP}{IV.8.30 (p.~945)} }
\expandafter\def\csname SRef@iv:correction:P3K.58\endcsname{\href{S6_SERIES_GRIND_IV.pdf\#nameddest=tag-01UQ}{IV.8.31 (p.~946)} }
\expandafter\def\csname SRef@iv:correction:P3K.59\endcsname{\href{S6_SERIES_GRIND_IV.pdf\#nameddest=tag-01UR}{IV.8.32 (p.~946)} }
\expandafter\def\csname SRef@iv:correction:P3K.60\endcsname{\href{S6_SERIES_GRIND_IV.pdf\#nameddest=tag-01UT}{IV.8.33 (p.~947)} }
\expandafter\def\csname SRef@iv:correction:P3K.17\endcsname{\href{S6_SERIES_GRIND_IV.pdf\#nameddest=tag-01U0}{IV.8.4 (p.~935)} }
\expandafter\def\csname SRef@iv:correction:P3K.22\endcsname{\href{S6_SERIES_GRIND_IV.pdf\#nameddest=tag-01U1}{IV.8.5 (p.~935)} }
\expandafter\def\csname SRef@iv:correction:P3K.23\endcsname{\href{S6_SERIES_GRIND_IV.pdf\#nameddest=tag-01U2}{IV.8.6 (p.~936)} }
\expandafter\def\csname SRef@iv:correction:P3K.24\endcsname{\href{S6_SERIES_GRIND_IV.pdf\#nameddest=tag-01U3}{IV.8.7 (p.~936)} }
\expandafter\def\csname SRef@iv:correction:P3K.25\endcsname{\href{S6_SERIES_GRIND_IV.pdf\#nameddest=tag-01U4}{IV.8.8 (p.~937)} }
\expandafter\def\csname SRef@iv:correction:P3K.27\endcsname{\href{S6_SERIES_GRIND_IV.pdf\#nameddest=tag-01U5}{IV.8.9 (p.~937)} }
\expandafter\def\csname SRef@fam:src:E.2.3\endcsname{\href{S6_SERIES_GRIND_IV.pdf\#nameddest=tag-01R8}{IV.2.2 (p.~21)} }
\expandafter\def\csname SRef@construction-and-proof\endcsname{\href{S6_SERIES_GRIND_IV.pdf\#nameddest=tag-01GZ}{IV.3.0.1 (p.~127)} }
\expandafter\def\csname SRef@the-marked-double-and-its-bordism\endcsname{\href{S6_SERIES_GRIND_IV.pdf\#nameddest=tag-034Y}{IV.{{$\langle 1\rangle $1.}} (p.~144)} }
\expandafter\def\csname SRef@removal-of-the-cubic-coordinate-by-an-extending-map\endcsname{\href{S6_SERIES_GRIND_IV.pdf\#nameddest=tag-033W}{IV.{{$\langle 1\rangle $2.}} (p.~144)} }
\expandafter\def\csname SRef@the-surgery-and-its-bounding-eight-manifold\endcsname{\href{S6_SERIES_GRIND_IV.pdf\#nameddest=tag-0352}{IV.{{$\langle 1\rangle $3.}} (p.~144)} }
\expandafter\def\csname SRef@the-invariant-and-the-resulting-homomorphism\endcsname{\href{S6_SERIES_GRIND_IV.pdf\#nameddest=tag-034W}{IV.{{$\langle 1\rangle $4.}} (p.~145)} }
\expandafter\def\csname SRef@proof\endcsname{\href{S6_SERIES_GRIND_IV.pdf\#nameddest=tag-02YT}{IV.3.0.2 (p.~129)} }
\expandafter\def\csname SRef@an-explicit-rational-elliptic-surface-and-primitive-section\endcsname{\href{S6_SERIES_GRIND_IV.pdf\#nameddest=tag-0009}{IV.3.0.3 (p.~135)} }
\expandafter\def\csname SRef@the-linearization-and-every-local-filling\endcsname{\href{S6_SERIES_GRIND_IV.pdf\#nameddest=tag-034X}{IV.3.0.4 (p.~136)} }
\expandafter\def\csname SRef@actual-integral-monodromy-not-an-abstract-substitute-tuple\endcsname{\href{S6_SERIES_GRIND_IV.pdf\#nameddest=tag-0006}{IV.3.0.5 (p.~138)} }
\expandafter\def\csname SRef@fundamental-group-and-every-integral-cohomology-differential\endcsname{\href{S6_SERIES_GRIND_IV.pdf\#nameddest=tag-01RN}{IV.3.0.6 (p.~139)} }
\expandafter\def\csname SRef@the-punctured-holomorphic-family-and-its-closing-map\endcsname{\href{S6_SERIES_GRIND_IV.pdf\#nameddest=tag-0350}{IV.3.0.7 (p.~142)} }
\expandafter\def\csname SRef@the-actual-page-and-its-section-line\endcsname{\href{S6_SERIES_GRIND_IV.pdf\#nameddest=tag-034R}{IV.3.0.8 (p.~142)} }
\expandafter\def\csname SRef@executed-verification\endcsname{\href{S6_SERIES_GRIND_IV.pdf\#nameddest=tag-01OP}{IV.3.0.9 (p.~143)} }
\expandafter\def\csname SRef@night:P3K.75\endcsname{\href{S6_SERIES_GRIND_IV.pdf\#nameddest=tag-02VN}{IV.11.1 (p.~1392)} }
\expandafter\def\csname SRef@night:P3K.117\endcsname{\href{S6_SERIES_GRIND_IV.pdf\#nameddest=tag-02TO}{IV.10.7 (p.~1213)} }
\expandafter\def\csname SRef@night:P3K.118\endcsname{\href{S6_SERIES_GRIND_IV.pdf\#nameddest=tag-02TP}{IV.10.8 (p.~1250)} }
\expandafter\def\csname SRef@night:P3K.119\endcsname{\href{S6_SERIES_GRIND_IV.pdf\#nameddest=tag-02TQ}{IV.10.9 (p.~1278)} }
\expandafter\def\csname SRef@night:P3K.109\endcsname{\href{S6_SERIES_GRIND_IV.pdf\#nameddest=tag-02TF}{IV.10.1 (p.~1116)} }
\expandafter\def\csname SRef@night:P3K.110\endcsname{\href{S6_SERIES_GRIND_IV.pdf\#nameddest=tag-02TH}{IV.10.2 (p.~1127)} }
\expandafter\def\csname SRef@night:P3K.111\endcsname{\href{S6_SERIES_GRIND_IV.pdf\#nameddest=tag-02TI}{IV.10.10 (p.~1305)} }
\expandafter\def\csname SRef@night:P3K.112\endcsname{\href{S6_SERIES_GRIND_IV.pdf\#nameddest=tag-02TJ}{IV.10.11 (p.~1369)} }
\expandafter\def\csname SRef@night:P3K.113\endcsname{\href{S6_SERIES_GRIND_IV.pdf\#nameddest=tag-02TK}{IV.10.3 (p.~1135)} }
\expandafter\def\csname SRef@night:P3K.114\endcsname{\href{S6_SERIES_GRIND_IV.pdf\#nameddest=tag-02TL}{IV.10.4 (p.~1147)} }
\expandafter\def\csname SRef@night:P3K.115\endcsname{\href{S6_SERIES_GRIND_IV.pdf\#nameddest=tag-02TM}{IV.10.5 (p.~1177)} }
\expandafter\def\csname SRef@night:P3K.116\endcsname{\href{S6_SERIES_GRIND_IV.pdf\#nameddest=tag-02TN}{IV.10.6 (p.~1193)} }
\expandafter\def\csname SRef@push:src:C.3.24\endcsname{\href{S6_SERIES_GRIND_IV.pdf\#nameddest=tag-031W}{IV.2.25 (p.~57)} }
\expandafter\def\csname SRef@maps:src:B.6.2\endcsname{\href{S6_SERIES_GRIND_IV.pdf\#nameddest=tag-02JO}{IV.2.33 (p.~67)} }
\expandafter\def\csname SRef@maps:src:B.6.20\endcsname{\href{S6_SERIES_GRIND_IV.pdf\#nameddest=tag-02JP}{IV.2.50 (p.~96)} }
\expandafter\def\csname SRef@maps:src:B.6.21\endcsname{\href{S6_SERIES_GRIND_IV.pdf\#nameddest=tag-02JQ}{IV.2.51 (p.~99)} }
\expandafter\def\csname SRef@maps:src:B.6.23\endcsname{\href{S6_SERIES_GRIND_IV.pdf\#nameddest=tag-02JS}{IV.2.53 (p.~101)} }
\expandafter\def\csname SRef@maps:src:B.6.24\endcsname{\href{S6_SERIES_GRIND_IV.pdf\#nameddest=tag-02JT}{IV.2.54 (p.~102)} }
\expandafter\def\csname SRef@night:P3F.1\endcsname{\href{S6_SERIES_GRIND_IV.pdf\#nameddest=tag-02RK}{IV.3.1 (p.~104)} }
\expandafter\def\csname SRef@night:P3F.10\endcsname{\href{S6_SERIES_GRIND_IV.pdf\#nameddest=tag-02RL}{IV.3.10 (p.~117)} }
\expandafter\def\csname SRef@night:P3F.11\endcsname{\href{S6_SERIES_GRIND_IV.pdf\#nameddest=tag-02RM}{IV.3.11 (p.~119)} }
\expandafter\def\csname SRef@night:P3F.12\endcsname{\href{S6_SERIES_GRIND_IV.pdf\#nameddest=tag-02RN}{IV.3.12 (p.~121)} }
\expandafter\def\csname SRef@night:P3F.13\endcsname{\href{S6_SERIES_GRIND_IV.pdf\#nameddest=tag-02RO}{IV.3.13 (p.~123)} }
\expandafter\def\csname SRef@night:P3F.14\endcsname{\href{S6_SERIES_GRIND_IV.pdf\#nameddest=tag-02RP}{IV.3.14 (p.~124)} }
\expandafter\def\csname SRef@night:P3F.15\endcsname{\href{S6_SERIES_GRIND_IV.pdf\#nameddest=tag-02RQ}{IV.3.15 (p.~125)} }
\expandafter\def\csname SRef@night:P3F.16\endcsname{\href{S6_SERIES_GRIND_IV.pdf\#nameddest=tag-02RR}{IV.3.16 (p.~126)} }
\expandafter\def\csname SRef@night:P3F.17\endcsname{\href{S6_SERIES_GRIND_IV.pdf\#nameddest=tag-02RS}{IV.3.17 (p.~128)} }
\expandafter\def\csname SRef@night:P3F.18\endcsname{\href{S6_SERIES_GRIND_IV.pdf\#nameddest=tag-02RT}{IV.3.18 (p.~129)} }
\expandafter\def\csname SRef@night:P3F.19\endcsname{\href{S6_SERIES_GRIND_IV.pdf\#nameddest=tag-02RU}{IV.3.19 (p.~130)} }
\expandafter\def\csname SRef@night:P3F.2\endcsname{\href{S6_SERIES_GRIND_IV.pdf\#nameddest=tag-02RV}{IV.3.2 (p.~107)} }
\expandafter\def\csname SRef@night:P3F.20\endcsname{\href{S6_SERIES_GRIND_IV.pdf\#nameddest=tag-02RW}{IV.3.21 (p.~134)} }
\expandafter\def\csname SRef@night:P3F.21\endcsname{\href{S6_SERIES_GRIND_IV.pdf\#nameddest=tag-02RX}{IV.3.22 (p.~134)} }
\expandafter\def\csname SRef@night:P3F.22\endcsname{\href{S6_SERIES_GRIND_IV.pdf\#nameddest=tag-02RY}{IV.3.23 (p.~144)} }
\expandafter\def\csname SRef@night:P3F.23\endcsname{\href{S6_SERIES_GRIND_IV.pdf\#nameddest=tag-02RZ}{IV.3.24 (p.~145)} }
\expandafter\def\csname SRef@night:P3F.24\endcsname{\href{S6_SERIES_GRIND_IV.pdf\#nameddest=tag-02S0}{IV.3.25 (p.~147)} }
\expandafter\def\csname SRef@night:P3F.30\endcsname{\href{S6_SERIES_GRIND_IV.pdf\#nameddest=tag-02S7}{IV.3.26 (p.~148)} }
\expandafter\def\csname SRef@night:P3F.35\endcsname{\href{S6_SERIES_GRIND_IV.pdf\#nameddest=tag-02SC}{IV.3.28 (p.~152)} }
\expandafter\def\csname SRef@night:P3F.36\endcsname{\href{S6_SERIES_GRIND_IV.pdf\#nameddest=tag-02SD}{IV.3.29 (p.~154)} }
\expandafter\def\csname SRef@night:P3F.37\endcsname{\href{S6_SERIES_GRIND_IV.pdf\#nameddest=tag-02SE}{IV.3.30 (p.~159)} }
\expandafter\def\csname SRef@night:P3F.3\endcsname{\href{S6_SERIES_GRIND_IV.pdf\#nameddest=tag-02S6}{IV.3.3 (p.~108)} }
\expandafter\def\csname SRef@night:P3F.38\endcsname{\href{S6_SERIES_GRIND_IV.pdf\#nameddest=tag-02SF}{IV.3.31 (p.~164)} }
\expandafter\def\csname SRef@night:P3F.39\endcsname{\href{S6_SERIES_GRIND_IV.pdf\#nameddest=tag-02SG}{IV.3.32 (p.~169)} }
\expandafter\def\csname SRef@night:P3F.40\endcsname{\href{S6_SERIES_GRIND_IV.pdf\#nameddest=tag-02SI}{IV.3.33 (p.~172)} }
\expandafter\def\csname SRef@night:P3F.41\endcsname{\href{S6_SERIES_GRIND_IV.pdf\#nameddest=tag-02SJ}{IV.3.34 (p.~175)} }
\expandafter\def\csname SRef@night:P3F.42\endcsname{\href{S6_SERIES_GRIND_IV.pdf\#nameddest=tag-02SK}{IV.3.35 (p.~180)} }
\expandafter\def\csname SRef@night:P3F.43\endcsname{\href{S6_SERIES_GRIND_IV.pdf\#nameddest=tag-02SL}{IV.3.36 (p.~186)} }
\expandafter\def\csname SRef@night:P3F.44\endcsname{\href{S6_SERIES_GRIND_IV.pdf\#nameddest=tag-02SM}{IV.3.37 (p.~188)} }
\expandafter\def\csname SRef@night:P3F.45\endcsname{\href{S6_SERIES_GRIND_IV.pdf\#nameddest=tag-02SN}{IV.3.38 (p.~191)} }
\expandafter\def\csname SRef@night:P3F.46\endcsname{\href{S6_SERIES_GRIND_IV.pdf\#nameddest=tag-02SO}{IV.3.39 (p.~195)} }
\expandafter\def\csname SRef@night:P3F.47\endcsname{\href{S6_SERIES_GRIND_IV.pdf\#nameddest=tag-02SP}{IV.3.40 (p.~202)} }
\expandafter\def\csname SRef@night:P3F.4\endcsname{\href{S6_SERIES_GRIND_IV.pdf\#nameddest=tag-02SH}{IV.3.4 (p.~109)} }
\expandafter\def\csname SRef@night:P3F.48\endcsname{\href{S6_SERIES_GRIND_IV.pdf\#nameddest=tag-02SQ}{IV.3.41 (p.~207)} }
\expandafter\def\csname SRef@night:P3F.49\endcsname{\href{S6_SERIES_GRIND_IV.pdf\#nameddest=tag-02SR}{IV.3.42 (p.~210)} }
\expandafter\def\csname SRef@night:P3F.50\endcsname{\href{S6_SERIES_GRIND_IV.pdf\#nameddest=tag-02ST}{IV.3.43 (p.~212)} }
\expandafter\def\csname SRef@night:P3F.51\endcsname{\href{S6_SERIES_GRIND_IV.pdf\#nameddest=tag-02SU}{IV.3.44 (p.~216)} }
\expandafter\def\csname SRef@night:P3F.52\endcsname{\href{S6_SERIES_GRIND_IV.pdf\#nameddest=tag-02SV}{IV.3.45 (p.~222)} }
\expandafter\def\csname SRef@night:P3F.54\endcsname{\href{S6_SERIES_GRIND_IV.pdf\#nameddest=tag-02SX}{IV.3.46 (p.~225)} }
\expandafter\def\csname SRef@night:P3F.55\endcsname{\href{S6_SERIES_GRIND_IV.pdf\#nameddest=tag-02SY}{IV.3.47 (p.~229)} }
\expandafter\def\csname SRef@night:P3F.56\endcsname{\href{S6_SERIES_GRIND_IV.pdf\#nameddest=tag-02SZ}{IV.3.48 (p.~234)} }
\expandafter\def\csname SRef@night:P3F.5\endcsname{\href{S6_SERIES_GRIND_IV.pdf\#nameddest=tag-02SS}{IV.3.5 (p.~110)} }
\expandafter\def\csname SRef@night:P3F.6\endcsname{\href{S6_SERIES_GRIND_IV.pdf\#nameddest=tag-02T0}{IV.3.6 (p.~111)} }
\expandafter\def\csname SRef@night:P3F.7\endcsname{\href{S6_SERIES_GRIND_IV.pdf\#nameddest=tag-02T1}{IV.3.7 (p.~112)} }
\expandafter\def\csname SRef@night:P3F.8\endcsname{\href{S6_SERIES_GRIND_IV.pdf\#nameddest=tag-02T2}{IV.3.8 (p.~114)} }
\expandafter\def\csname SRef@night:P3F.9\endcsname{\href{S6_SERIES_GRIND_IV.pdf\#nameddest=tag-02T3}{IV.3.9 (p.~115)} }
\expandafter\def\csname SRef@night:P3R.1\endcsname{\href{S6_SERIES_GRIND_IV.pdf\#nameddest=tag-02WE}{IV.4.1 (p.~243)} }
\expandafter\def\csname SRef@night:P3K.9\endcsname{\href{S6_SERIES_GRIND_IV.pdf\#nameddest=tag-02W3}{IV.4.10 (p.~254)} }
\expandafter\def\csname SRef@night:P3K.11\endcsname{\href{S6_SERIES_GRIND_IV.pdf\#nameddest=tag-02TG}{IV.4.11 (p.~257)} }
\expandafter\def\csname SRef@night:P3K.12\endcsname{\href{S6_SERIES_GRIND_IV.pdf\#nameddest=tag-02TR}{IV.4.12 (p.~260)} }
\expandafter\def\csname SRef@night:P3K.13\endcsname{\href{S6_SERIES_GRIND_IV.pdf\#nameddest=tag-02TS}{IV.4.13 (p.~261)} }
\expandafter\def\csname SRef@night:P3K.14\endcsname{\href{S6_SERIES_GRIND_IV.pdf\#nameddest=tag-02TT}{IV.4.14 (p.~263)} }
\expandafter\def\csname SRef@night:P3K.15\endcsname{\href{S6_SERIES_GRIND_IV.pdf\#nameddest=tag-02TU}{IV.4.15 (p.~264)} }
\expandafter\def\csname SRef@night:P3K.16\endcsname{\href{S6_SERIES_GRIND_IV.pdf\#nameddest=tag-02TV}{IV.4.16 (p.~268)} }
\expandafter\def\csname SRef@night:P3K.18\endcsname{\href{S6_SERIES_GRIND_IV.pdf\#nameddest=tag-02TX}{IV.4.17 (p.~270)} }
\expandafter\def\csname SRef@night:P3K.19\endcsname{\href{S6_SERIES_GRIND_IV.pdf\#nameddest=tag-02TY}{IV.4.18 (p.~271)} }
\expandafter\def\csname SRef@night:P3K.20\endcsname{\href{S6_SERIES_GRIND_IV.pdf\#nameddest=tag-02U0}{IV.4.19 (p.~273)} }
\expandafter\def\csname SRef@night:P3R.2\endcsname{\href{S6_SERIES_GRIND_IV.pdf\#nameddest=tag-02WF}{IV.4.2 (p.~243)} }
\expandafter\def\csname SRef@night:P3K.21\endcsname{\href{S6_SERIES_GRIND_IV.pdf\#nameddest=tag-02U1}{IV.4.20 (p.~276)} }
\expandafter\def\csname SRef@night:P3K.26\endcsname{\href{S6_SERIES_GRIND_IV.pdf\#nameddest=tag-02U6}{IV.4.21 (p.~278)} }
\expandafter\def\csname SRef@night:P3K.28\endcsname{\href{S6_SERIES_GRIND_IV.pdf\#nameddest=tag-02U8}{IV.4.22 (p.~283)} }
\expandafter\def\csname SRef@night:P3K.1\endcsname{\href{S6_SERIES_GRIND_IV.pdf\#nameddest=tag-02T4}{IV.4.3 (p.~244)} }
\expandafter\def\csname SRef@night:P3K.2\endcsname{\href{S6_SERIES_GRIND_IV.pdf\#nameddest=tag-02TZ}{IV.4.4 (p.~245)} }
\expandafter\def\csname SRef@night:P3K.3\endcsname{\href{S6_SERIES_GRIND_IV.pdf\#nameddest=tag-02UA}{IV.4.5 (p.~246)} }
\expandafter\def\csname SRef@night:P3K.4\endcsname{\href{S6_SERIES_GRIND_IV.pdf\#nameddest=tag-02UK}{IV.4.6 (p.~247)} }
\expandafter\def\csname SRef@night:P3K.5\endcsname{\href{S6_SERIES_GRIND_IV.pdf\#nameddest=tag-02UV}{IV.4.7 (p.~248)} }
\expandafter\def\csname SRef@night:P3K.7\endcsname{\href{S6_SERIES_GRIND_IV.pdf\#nameddest=tag-02VH}{IV.4.8 (p.~249)} }
\expandafter\def\csname SRef@night:P3K.8\endcsname{\href{S6_SERIES_GRIND_IV.pdf\#nameddest=tag-02VS}{IV.4.9 (p.~251)} }
\expandafter\def\csname SRef@night:P3K.38\endcsname{\href{S6_SERIES_GRIND_IV.pdf\#nameddest=tag-02UJ}{IV.5.1 (p.~291)} }
\expandafter\def\csname SRef@night:P3K.62\endcsname{\href{S6_SERIES_GRIND_IV.pdf\#nameddest=tag-02V9}{IV.5.10 (p.~333)} }
\expandafter\def\csname SRef@night:P3K.63\endcsname{\href{S6_SERIES_GRIND_IV.pdf\#nameddest=tag-02VA}{IV.5.11 (p.~340)} }
\expandafter\def\csname SRef@night:P3K.64\endcsname{\href{S6_SERIES_GRIND_IV.pdf\#nameddest=tag-02VB}{IV.5.12 (p.~348)} }
\expandafter\def\csname SRef@night:P3K.65\endcsname{\href{S6_SERIES_GRIND_IV.pdf\#nameddest=tag-02VC}{IV.5.13 (p.~352)} }
\expandafter\def\csname SRef@night:P3K.66\endcsname{\href{S6_SERIES_GRIND_IV.pdf\#nameddest=tag-02VD}{IV.5.14 (p.~362)} }
\expandafter\def\csname SRef@night:P3K.67\endcsname{\href{S6_SERIES_GRIND_IV.pdf\#nameddest=tag-02VE}{IV.5.15 (p.~369)} }
\expandafter\def\csname SRef@night:P3K.42\endcsname{\href{S6_SERIES_GRIND_IV.pdf\#nameddest=tag-02UN}{IV.5.2 (p.~298)} }
\expandafter\def\csname SRef@night:P3K.47\endcsname{\href{S6_SERIES_GRIND_IV.pdf\#nameddest=tag-02US}{IV.5.3 (p.~304)} }
\expandafter\def\csname SRef@night:P3K.48\endcsname{\href{S6_SERIES_GRIND_IV.pdf\#nameddest=tag-02UT}{IV.5.4 (p.~307)} }
\expandafter\def\csname SRef@night:P3K.49\endcsname{\href{S6_SERIES_GRIND_IV.pdf\#nameddest=tag-02UU}{IV.5.5 (p.~311)} }
\expandafter\def\csname SRef@night:P3K.51\endcsname{\href{S6_SERIES_GRIND_IV.pdf\#nameddest=tag-02UX}{IV.5.6 (p.~318)} }
\expandafter\def\csname SRef@night:P3K.53\endcsname{\href{S6_SERIES_GRIND_IV.pdf\#nameddest=tag-02UZ}{IV.5.7 (p.~321)} }
\expandafter\def\csname SRef@night:P3K.55\endcsname{\href{S6_SERIES_GRIND_IV.pdf\#nameddest=tag-02V1}{IV.5.8 (p.~325)} }
\expandafter\def\csname SRef@night:P3K.61\endcsname{\href{S6_SERIES_GRIND_IV.pdf\#nameddest=tag-02V8}{IV.5.9 (p.~330)} }
\expandafter\def\csname SRef@night:P3K.68\endcsname{\href{S6_SERIES_GRIND_IV.pdf\#nameddest=tag-02VF}{IV.6.1 (p.~374)} }
\expandafter\def\csname SRef@night:P3K.69\endcsname{\href{S6_SERIES_GRIND_IV.pdf\#nameddest=tag-02VG}{IV.6.2 (p.~380)} }
\expandafter\def\csname SRef@night:P3K.70\endcsname{\href{S6_SERIES_GRIND_IV.pdf\#nameddest=tag-02VI}{IV.6.3 (p.~384)} }
\expandafter\def\csname SRef@night:P3K.71\endcsname{\href{S6_SERIES_GRIND_IV.pdf\#nameddest=tag-02VJ}{IV.6.4 (p.~389)} }
\expandafter\def\csname SRef@night:P3K.72\endcsname{\href{S6_SERIES_GRIND_IV.pdf\#nameddest=tag-02VK}{IV.6.5 (p.~393)} }
\expandafter\def\csname SRef@night:P3K.73\endcsname{\href{S6_SERIES_GRIND_IV.pdf\#nameddest=tag-02VL}{IV.6.6 (p.~400)} }
\expandafter\def\csname SRef@night:P3K.74\endcsname{\href{S6_SERIES_GRIND_IV.pdf\#nameddest=tag-02VM}{IV.6.7 (p.~405)} }
\expandafter\def\csname SRef@night:P3K.76\endcsname{\href{S6_SERIES_GRIND_IV.pdf\#nameddest=tag-02VO}{IV.7.1 (p.~407)} }
\expandafter\def\csname SRef@night:P3K.85\endcsname{\href{S6_SERIES_GRIND_IV.pdf\#nameddest=tag-02VY}{IV.7.10 (p.~494)} }
\expandafter\def\csname SRef@night:P3K.86\endcsname{\href{S6_SERIES_GRIND_IV.pdf\#nameddest=tag-02VZ}{IV.7.11 (p.~499)} }
\expandafter\def\csname SRef@night:P3K.87\endcsname{\href{S6_SERIES_GRIND_IV.pdf\#nameddest=tag-02W0}{IV.7.12 (p.~505)} }
\expandafter\def\csname SRef@night:P3K.88\endcsname{\href{S6_SERIES_GRIND_IV.pdf\#nameddest=tag-02W1}{IV.7.13 (p.~509)} }
\expandafter\def\csname SRef@night:P3K.89\endcsname{\href{S6_SERIES_GRIND_IV.pdf\#nameddest=tag-02W2}{IV.7.14 (p.~513)} }
\expandafter\def\csname SRef@night:P3K.90\endcsname{\href{S6_SERIES_GRIND_IV.pdf\#nameddest=tag-02W4}{IV.7.15 (p.~521)} }
\expandafter\def\csname SRef@night:P3K.91\endcsname{\href{S6_SERIES_GRIND_IV.pdf\#nameddest=tag-02W5}{IV.7.16 (p.~525)} }
\expandafter\def\csname SRef@night:P3K.92\endcsname{\href{S6_SERIES_GRIND_IV.pdf\#nameddest=tag-02W6}{IV.7.17 (p.~531)} }
\expandafter\def\csname SRef@night:P3K.93\endcsname{\href{S6_SERIES_GRIND_IV.pdf\#nameddest=tag-02W7}{IV.7.18 (p.~538)} }
\expandafter\def\csname SRef@night:P3K.94\endcsname{\href{S6_SERIES_GRIND_IV.pdf\#nameddest=tag-02W8}{IV.7.19 (p.~543)} }
\expandafter\def\csname SRef@night:P3K.77\endcsname{\href{S6_SERIES_GRIND_IV.pdf\#nameddest=tag-02VP}{IV.7.2 (p.~414)} }
\expandafter\def\csname SRef@night:P3K.95\endcsname{\href{S6_SERIES_GRIND_IV.pdf\#nameddest=tag-02W9}{IV.7.20 (p.~551)} }
\expandafter\def\csname SRef@night:P3K.96\endcsname{\href{S6_SERIES_GRIND_IV.pdf\#nameddest=tag-02WA}{IV.7.21 (p.~562)} }
\expandafter\def\csname SRef@night:P3K.97\endcsname{\href{S6_SERIES_GRIND_IV.pdf\#nameddest=tag-02WB}{IV.7.22 (p.~569)} }
\expandafter\def\csname SRef@night:P3K.98\endcsname{\href{S6_SERIES_GRIND_IV.pdf\#nameddest=tag-02WC}{IV.7.23 (p.~578)} }
\expandafter\def\csname SRef@night:P3K.99\endcsname{\href{S6_SERIES_GRIND_IV.pdf\#nameddest=tag-02WD}{IV.7.24 (p.~587)} }
\expandafter\def\csname SRef@night:P3K.100\endcsname{\href{S6_SERIES_GRIND_IV.pdf\#nameddest=tag-02T6}{IV.7.25 (p.~662)} }
\expandafter\def\csname SRef@night:P3K.101\endcsname{\href{S6_SERIES_GRIND_IV.pdf\#nameddest=tag-02T7}{IV.7.26 (p.~669)} }
\expandafter\def\csname SRef@night:P3K.102\endcsname{\href{S6_SERIES_GRIND_IV.pdf\#nameddest=tag-02T8}{IV.7.27 (p.~758)} }
\expandafter\def\csname SRef@night:P3K.103\endcsname{\href{S6_SERIES_GRIND_IV.pdf\#nameddest=tag-02T9}{IV.7.28 (p.~768)} }
\expandafter\def\csname SRef@night:P3K.104\endcsname{\href{S6_SERIES_GRIND_IV.pdf\#nameddest=tag-02TA}{IV.7.29 (p.~774)} }
\expandafter\def\csname SRef@night:P3K.78\endcsname{\href{S6_SERIES_GRIND_IV.pdf\#nameddest=tag-02VQ}{IV.7.3 (p.~420)} }
\expandafter\def\csname SRef@night:P3K.105\endcsname{\href{S6_SERIES_GRIND_IV.pdf\#nameddest=tag-02TB}{IV.7.30 (p.~785)} }
\expandafter\def\csname SRef@night:P3K.106\endcsname{\href{S6_SERIES_GRIND_IV.pdf\#nameddest=tag-02TC}{IV.7.31 (p.~798)} }
\expandafter\def\csname SRef@night:P3K.108\endcsname{\href{S6_SERIES_GRIND_IV.pdf\#nameddest=tag-02TE}{IV.7.32 (p.~861)} }
\expandafter\def\csname SRef@iv:i3_all_reduced_caps\endcsname{\href{S6_SERIES_GRIND_IV.pdf\#nameddest=tag-02TE}{IV.7.32 (p.~861)} }
\expandafter\def\csname SRef@night:P3K.107\endcsname{\href{S6_SERIES_GRIND_IV.pdf\#nameddest=tag-02TD}{IV.7.33 (p.~867)} }
\expandafter\def\csname SRef@night:P3K.79\endcsname{\href{S6_SERIES_GRIND_IV.pdf\#nameddest=tag-02VR}{IV.7.4 (p.~426)} }
\expandafter\def\csname SRef@night:P3K.80\endcsname{\href{S6_SERIES_GRIND_IV.pdf\#nameddest=tag-02VT}{IV.7.5 (p.~431)} }
\expandafter\def\csname SRef@night:P3K.81\endcsname{\href{S6_SERIES_GRIND_IV.pdf\#nameddest=tag-02VU}{IV.7.6 (p.~467)} }
\expandafter\def\csname SRef@night:P3K.82\endcsname{\href{S6_SERIES_GRIND_IV.pdf\#nameddest=tag-02VV}{IV.7.7 (p.~475)} }
\expandafter\def\csname SRef@night:P3K.83\endcsname{\href{S6_SERIES_GRIND_IV.pdf\#nameddest=tag-02VW}{IV.7.8 (p.~478)} }
\expandafter\def\csname SRef@night:P3K.84\endcsname{\href{S6_SERIES_GRIND_IV.pdf\#nameddest=tag-02VX}{IV.7.9 (p.~490)} }
\expandafter\def\csname SRef@night:P3F.25\endcsname{\href{S6_SERIES_GRIND_IV.pdf\#nameddest=tag-02S1}{IV.9.1 (p.~947)} }
\expandafter\def\csname SRef@night:P3K.29\endcsname{\href{S6_SERIES_GRIND_IV.pdf\#nameddest=tag-02U9}{IV.9.10 (p.~988)} }
\expandafter\def\csname SRef@night:P3K.30\endcsname{\href{S6_SERIES_GRIND_IV.pdf\#nameddest=tag-02UB}{IV.9.11 (p.~992)} }
\expandafter\def\csname SRef@night:P3K.31\endcsname{\href{S6_SERIES_GRIND_IV.pdf\#nameddest=tag-02UC}{IV.9.12 (p.~997)} }
\expandafter\def\csname SRef@night:P3K.32\endcsname{\href{S6_SERIES_GRIND_IV.pdf\#nameddest=tag-02UD}{IV.9.13 (p.~1000)} }
\expandafter\def\csname SRef@night:P3K.33\endcsname{\href{S6_SERIES_GRIND_IV.pdf\#nameddest=tag-02UE}{IV.9.14 (p.~1010)} }
\expandafter\def\csname SRef@night:P3K.34\endcsname{\href{S6_SERIES_GRIND_IV.pdf\#nameddest=tag-02UF}{IV.9.15 (p.~1015)} }
\expandafter\def\csname SRef@night:P3F.53\endcsname{\href{S6_SERIES_GRIND_IV.pdf\#nameddest=tag-02SW}{IV.9.16 (p.~1020)} }
\expandafter\def\csname SRef@night:P3K.35\endcsname{\href{S6_SERIES_GRIND_IV.pdf\#nameddest=tag-02UG}{IV.9.17 (p.~1025)} }
\expandafter\def\csname SRef@night:P3K.37\endcsname{\href{S6_SERIES_GRIND_IV.pdf\#nameddest=tag-02UI}{IV.9.18 (p.~1032)} }
\expandafter\def\csname SRef@night:P3K.40\endcsname{\href{S6_SERIES_GRIND_IV.pdf\#nameddest=tag-02UL}{IV.9.19 (p.~1036)} }
\expandafter\def\csname SRef@night:P3K.6\endcsname{\href{S6_SERIES_GRIND_IV.pdf\#nameddest=tag-02V6}{IV.9.2 (p.~950)} }
\expandafter\def\csname SRef@night:P3K.41\endcsname{\href{S6_SERIES_GRIND_IV.pdf\#nameddest=tag-02UM}{IV.9.20 (p.~1040)} }
\expandafter\def\csname SRef@night:P3K.43\endcsname{\href{S6_SERIES_GRIND_IV.pdf\#nameddest=tag-02UO}{IV.9.21 (p.~1047)} }
\expandafter\def\csname SRef@night:P3K.44\endcsname{\href{S6_SERIES_GRIND_IV.pdf\#nameddest=tag-02UP}{IV.9.22 (p.~1050)} }
\expandafter\def\csname SRef@night:P3K.45\endcsname{\href{S6_SERIES_GRIND_IV.pdf\#nameddest=tag-02UQ}{IV.9.23 (p.~1055)} }
\expandafter\def\csname SRef@night:P3K.46\endcsname{\href{S6_SERIES_GRIND_IV.pdf\#nameddest=tag-02UR}{IV.9.24 (p.~1058)} }
\expandafter\def\csname SRef@night:P3K.50\endcsname{\href{S6_SERIES_GRIND_IV.pdf\#nameddest=tag-02UW}{IV.9.25 (p.~1061)} }
\expandafter\def\csname SRef@night:P3K.52\endcsname{\href{S6_SERIES_GRIND_IV.pdf\#nameddest=tag-02UY}{IV.9.26 (p.~1064)} }
\expandafter\def\csname SRef@night:P3K.54\endcsname{\href{S6_SERIES_GRIND_IV.pdf\#nameddest=tag-02V0}{IV.9.27 (p.~1067)} }
\expandafter\def\csname SRef@night:P3K.56\endcsname{\href{S6_SERIES_GRIND_IV.pdf\#nameddest=tag-02V2}{IV.9.28 (p.~1069)} }
\expandafter\def\csname SRef@night:P3K.57\endcsname{\href{S6_SERIES_GRIND_IV.pdf\#nameddest=tag-02V3}{IV.9.29 (p.~1072)} }
\expandafter\def\csname SRef@night:P3K.10\endcsname{\href{S6_SERIES_GRIND_IV.pdf\#nameddest=tag-02T5}{IV.9.3 (p.~952)} }
\expandafter\def\csname SRef@night:P3K.58\endcsname{\href{S6_SERIES_GRIND_IV.pdf\#nameddest=tag-02V4}{IV.9.30 (p.~1095)} }
\expandafter\def\csname SRef@night:P3K.59\endcsname{\href{S6_SERIES_GRIND_IV.pdf\#nameddest=tag-02V5}{IV.9.31 (p.~1104)} }
\expandafter\def\csname SRef@night:P3K.60\endcsname{\href{S6_SERIES_GRIND_IV.pdf\#nameddest=tag-02V7}{IV.9.32 (p.~1109)} }
\expandafter\def\csname SRef@night:P3K.17\endcsname{\href{S6_SERIES_GRIND_IV.pdf\#nameddest=tag-02TW}{IV.9.4 (p.~961)} }
\expandafter\def\csname SRef@night:P3K.22\endcsname{\href{S6_SERIES_GRIND_IV.pdf\#nameddest=tag-02U2}{IV.9.5 (p.~965)} }
\expandafter\def\csname SRef@night:P3K.23\endcsname{\href{S6_SERIES_GRIND_IV.pdf\#nameddest=tag-02U3}{IV.9.6 (p.~971)} }
\expandafter\def\csname SRef@night:P3K.24\endcsname{\href{S6_SERIES_GRIND_IV.pdf\#nameddest=tag-02U4}{IV.9.7 (p.~975)} }
\expandafter\def\csname SRef@night:P3K.25\endcsname{\href{S6_SERIES_GRIND_IV.pdf\#nameddest=tag-02U5}{IV.9.8 (p.~980)} }
\expandafter\def\csname SRef@night:P3K.27\endcsname{\href{S6_SERIES_GRIND_IV.pdf\#nameddest=tag-02U7}{IV.9.9 (p.~983)} }
\expandafter\def\csname SRef@maps:eq:F.6\endcsname{\href{S6_SERIES_GRIND_V.pdf\#nameddest=tag-02FA}{V.{F.6} (p.~85)} }
\expandafter\def\csname SRef@maps:eq:F.L2\endcsname{\href{S6_SERIES_GRIND_V.pdf\#nameddest=tag-02GV}{V.{F.L2} (p.~174)} }
\expandafter\def\csname SRef@maps:eq:F.L3\endcsname{\href{S6_SERIES_GRIND_V.pdf\#nameddest=tag-02GW}{V.{F.L3} (p.~174)} }
\expandafter\def\csname SRef@maps:eq:F.L4\endcsname{\href{S6_SERIES_GRIND_V.pdf\#nameddest=tag-02GX}{V.{F.L4} (p.~175)} }
\expandafter\def\csname SRef@maps:eq:F.L5\endcsname{\href{S6_SERIES_GRIND_V.pdf\#nameddest=tag-02GY}{V.{F.L5} (p.~175)} }
\expandafter\def\csname SRef@maps:eq:F.L6\endcsname{\href{S6_SERIES_GRIND_V.pdf\#nameddest=tag-02GZ}{V.{F.L6} (p.~175)} }
\expandafter\def\csname SRef@maps:eq:F.L7\endcsname{\href{S6_SERIES_GRIND_V.pdf\#nameddest=tag-02H0}{V.{F.L7} (p.~175)} }
\expandafter\def\csname SRef@maps:eq:F.L8\endcsname{\href{S6_SERIES_GRIND_V.pdf\#nameddest=tag-02H1}{V.{F.L8} (p.~175)} }
\expandafter\def\csname SRef@maps:eq:F.L9\endcsname{\href{S6_SERIES_GRIND_V.pdf\#nameddest=tag-02H2}{V.{F.L9} (p.~176)} }
\expandafter\def\csname SRef@maps:eq:F.L10\endcsname{\href{S6_SERIES_GRIND_V.pdf\#nameddest=tag-02GU}{V.{F.L10} (p.~176)} }
\expandafter\def\csname SRef@maps:eq:F.T1\endcsname{\href{S6_SERIES_GRIND_V.pdf\#nameddest=tag-02I7}{V.{F.T1} (p.~177)} }
\expandafter\def\csname SRef@maps:eq:F.T2\endcsname{\href{S6_SERIES_GRIND_V.pdf\#nameddest=tag-02IC}{V.{F.T2} (p.~177)} }
\expandafter\def\csname SRef@maps:eq:F.T3\endcsname{\href{S6_SERIES_GRIND_V.pdf\#nameddest=tag-02ID}{V.{F.T3} (p.~177)} }
\expandafter\def\csname SRef@maps:eq:F.T4\endcsname{\href{S6_SERIES_GRIND_V.pdf\#nameddest=tag-02IE}{V.{F.T4} (p.~178)} }
\expandafter\def\csname SRef@maps:eq:F.T5\endcsname{\href{S6_SERIES_GRIND_V.pdf\#nameddest=tag-02IF}{V.{F.T5} (p.~178)} }
\expandafter\def\csname SRef@maps:eq:F.T6\endcsname{\href{S6_SERIES_GRIND_V.pdf\#nameddest=tag-02IG}{V.{F.T6} (p.~178)} }
\expandafter\def\csname SRef@maps:eq:F.T7\endcsname{\href{S6_SERIES_GRIND_V.pdf\#nameddest=tag-02IH}{V.{F.T7} (p.~178)} }
\expandafter\def\csname SRef@maps:eq:F.T8\endcsname{\href{S6_SERIES_GRIND_V.pdf\#nameddest=tag-02II}{V.{F.T8} (p.~178)} }
\expandafter\def\csname SRef@maps:eq:F.T9\endcsname{\href{S6_SERIES_GRIND_V.pdf\#nameddest=tag-02IJ}{V.{F.T9} (p.~179)} }
\expandafter\def\csname SRef@maps:eq:F.T10\endcsname{\href{S6_SERIES_GRIND_V.pdf\#nameddest=tag-02I8}{V.{F.T10} (p.~179)} }
\expandafter\def\csname SRef@maps:eq:F.S1\endcsname{\href{S6_SERIES_GRIND_V.pdf\#nameddest=tag-02HX}{V.{F.S1} (p.~180)} }
\expandafter\def\csname SRef@maps:eq:F.S2\endcsname{\href{S6_SERIES_GRIND_V.pdf\#nameddest=tag-02HY}{V.{F.S2} (p.~180)} }
\expandafter\def\csname SRef@maps:eq:F.S3\endcsname{\href{S6_SERIES_GRIND_V.pdf\#nameddest=tag-02HZ}{V.{F.S3} (p.~180)} }
\expandafter\def\csname SRef@maps:eq:F.S4\endcsname{\href{S6_SERIES_GRIND_V.pdf\#nameddest=tag-02I0}{V.{F.S4} (p.~180)} }
\expandafter\def\csname SRef@maps:eq:F.S5\endcsname{\href{S6_SERIES_GRIND_V.pdf\#nameddest=tag-02I1}{V.{F.S5} (p.~180)} }
\expandafter\def\csname SRef@maps:eq:F.S6\endcsname{\href{S6_SERIES_GRIND_V.pdf\#nameddest=tag-02I2}{V.{F.S6} (p.~181)} }
\expandafter\def\csname SRef@maps:eq:F.S7\endcsname{\href{S6_SERIES_GRIND_V.pdf\#nameddest=tag-02I3}{V.{F.S7} (p.~181)} }
\expandafter\def\csname SRef@maps:eq:F.S8\endcsname{\href{S6_SERIES_GRIND_V.pdf\#nameddest=tag-02I4}{V.{F.S8} (p.~181)} }
\expandafter\def\csname SRef@maps:eq:F.S9\endcsname{\href{S6_SERIES_GRIND_V.pdf\#nameddest=tag-02I5}{V.{F.S9} (p.~181)} }
\expandafter\def\csname SRef@maps:eq:F.R10\endcsname{\href{S6_SERIES_GRIND_V.pdf\#nameddest=tag-02HK}{V.{F.R10} (p.~182)} }
\expandafter\def\csname SRef@maps:eq:F.R11\endcsname{\href{S6_SERIES_GRIND_V.pdf\#nameddest=tag-02HL}{V.{F.R11} (p.~182)} }
\expandafter\def\csname SRef@maps:eq:F.R12\endcsname{\href{S6_SERIES_GRIND_V.pdf\#nameddest=tag-02HM}{V.{F.R12} (p.~182)} }
\expandafter\def\csname SRef@maps:eq:F.R13\endcsname{\href{S6_SERIES_GRIND_V.pdf\#nameddest=tag-02HN}{V.{F.R13} (p.~182)} }
\expandafter\def\csname SRef@maps:eq:F.R14\endcsname{\href{S6_SERIES_GRIND_V.pdf\#nameddest=tag-02HO}{V.{F.R14} (p.~182)} }
\expandafter\def\csname SRef@maps:eq:F.X1\endcsname{\href{S6_SERIES_GRIND_V.pdf\#nameddest=tag-02IV}{V.{F.X1} (p.~183)} }
\expandafter\def\csname SRef@maps:eq:F.X2\endcsname{\href{S6_SERIES_GRIND_V.pdf\#nameddest=tag-02IW}{V.{F.X2} (p.~183)} }
\expandafter\def\csname SRef@maps:eq:F.X3\endcsname{\href{S6_SERIES_GRIND_V.pdf\#nameddest=tag-02IX}{V.{F.X3} (p.~184)} }
\expandafter\def\csname SRef@maps:eq:F.X4\endcsname{\href{S6_SERIES_GRIND_V.pdf\#nameddest=tag-02IY}{V.{F.X4} (p.~184)} }
\expandafter\def\csname SRef@maps:eq:F.X5\endcsname{\href{S6_SERIES_GRIND_V.pdf\#nameddest=tag-02IZ}{V.{F.X5} (p.~184)} }
\expandafter\def\csname SRef@maps:eq:F.T11\endcsname{\href{S6_SERIES_GRIND_V.pdf\#nameddest=tag-02I9}{V.{F.T11} (p.~185)} }
\expandafter\def\csname SRef@maps:eq:F.T12\endcsname{\href{S6_SERIES_GRIND_V.pdf\#nameddest=tag-02IA}{V.{F.T12} (p.~185)} }
\expandafter\def\csname SRef@maps:eq:F.T13\endcsname{\href{S6_SERIES_GRIND_V.pdf\#nameddest=tag-02IB}{V.{F.T13} (p.~185)} }
\expandafter\def\csname SRef@maps:eq:F.HC1\endcsname{\href{S6_SERIES_GRIND_V.pdf\#nameddest=tag-02G9}{V.{F.HC1} (p.~186)} }
\expandafter\def\csname SRef@maps:eq:F.HC2\endcsname{\href{S6_SERIES_GRIND_V.pdf\#nameddest=tag-02GE}{V.{F.HC2} (p.~186)} }
\expandafter\def\csname SRef@maps:eq:F.HC3\endcsname{\href{S6_SERIES_GRIND_V.pdf\#nameddest=tag-02GF}{V.{F.HC3} (p.~186)} }
\expandafter\def\csname SRef@maps:eq:F.HC4\endcsname{\href{S6_SERIES_GRIND_V.pdf\#nameddest=tag-02GG}{V.{F.HC4} (p.~187)} }
\expandafter\def\csname SRef@maps:eq:F.HC5\endcsname{\href{S6_SERIES_GRIND_V.pdf\#nameddest=tag-02GH}{V.{F.HC5} (p.~187)} }
\expandafter\def\csname SRef@maps:eq:F.HC6\endcsname{\href{S6_SERIES_GRIND_V.pdf\#nameddest=tag-02GI}{V.{F.HC6} (p.~187)} }
\expandafter\def\csname SRef@maps:eq:F.HC7\endcsname{\href{S6_SERIES_GRIND_V.pdf\#nameddest=tag-02GJ}{V.{F.HC7} (p.~188)} }
\expandafter\def\csname SRef@maps:eq:F.HC8\endcsname{\href{S6_SERIES_GRIND_V.pdf\#nameddest=tag-02GK}{V.{F.HC8} (p.~188)} }
\expandafter\def\csname SRef@maps:eq:F.HC9\endcsname{\href{S6_SERIES_GRIND_V.pdf\#nameddest=tag-02GL}{V.{F.HC9} (p.~189)} }
\expandafter\def\csname SRef@maps:eq:F.HC10\endcsname{\href{S6_SERIES_GRIND_V.pdf\#nameddest=tag-02GA}{V.{F.HC10} (p.~189)} }
\expandafter\def\csname SRef@maps:eq:F.HC11\endcsname{\href{S6_SERIES_GRIND_V.pdf\#nameddest=tag-02GB}{V.{F.HC11} (p.~189)} }
\expandafter\def\csname SRef@maps:eq:F.HC12\endcsname{\href{S6_SERIES_GRIND_V.pdf\#nameddest=tag-02GC}{V.{F.HC12} (p.~189)} }
\expandafter\def\csname SRef@maps:eq:F.HC13\endcsname{\href{S6_SERIES_GRIND_V.pdf\#nameddest=tag-02GD}{V.{F.HC13} (p.~190)} }
\expandafter\def\csname SRef@maps:eq:F.G1\endcsname{\href{S6_SERIES_GRIND_V.pdf\#nameddest=tag-02FO}{V.{F.G1} (p.~191)} }
\expandafter\def\csname SRef@maps:eq:F.G2\endcsname{\href{S6_SERIES_GRIND_V.pdf\#nameddest=tag-02FR}{V.{F.G2} (p.~191)} }
\expandafter\def\csname SRef@maps:eq:F.G3\endcsname{\href{S6_SERIES_GRIND_V.pdf\#nameddest=tag-02FS}{V.{F.G3} (p.~191)} }
\expandafter\def\csname SRef@maps:eq:F.G4\endcsname{\href{S6_SERIES_GRIND_V.pdf\#nameddest=tag-02FT}{V.{F.G4} (p.~191)} }
\expandafter\def\csname SRef@maps:eq:F.G5\endcsname{\href{S6_SERIES_GRIND_V.pdf\#nameddest=tag-02FU}{V.{F.G5} (p.~191)} }
\expandafter\def\csname SRef@maps:eq:F.G6\endcsname{\href{S6_SERIES_GRIND_V.pdf\#nameddest=tag-02FV}{V.{F.G6} (p.~191)} }
\expandafter\def\csname SRef@maps:eq:F.G7\endcsname{\href{S6_SERIES_GRIND_V.pdf\#nameddest=tag-02FW}{V.{F.G7} (p.~192)} }
\expandafter\def\csname SRef@maps:eq:F.G8\endcsname{\href{S6_SERIES_GRIND_V.pdf\#nameddest=tag-02FX}{V.{F.G8} (p.~192)} }
\expandafter\def\csname SRef@maps:eq:F.G9\endcsname{\href{S6_SERIES_GRIND_V.pdf\#nameddest=tag-02FY}{V.{F.G9} (p.~192)} }
\expandafter\def\csname SRef@maps:eq:F.G10\endcsname{\href{S6_SERIES_GRIND_V.pdf\#nameddest=tag-02FP}{V.{F.G10} (p.~192)} }
\expandafter\def\csname SRef@maps:eq:F.G11\endcsname{\href{S6_SERIES_GRIND_V.pdf\#nameddest=tag-02FQ}{V.{F.G11} (p.~194)} }
\expandafter\def\csname SRef@maps:eq:F.LT1\endcsname{\href{S6_SERIES_GRIND_V.pdf\#nameddest=tag-02H3}{V.{F.LT1} (p.~195)} }
\expandafter\def\csname SRef@maps:eq:F.LT2\endcsname{\href{S6_SERIES_GRIND_V.pdf\#nameddest=tag-02H4}{V.{F.LT2} (p.~195)} }
\expandafter\def\csname SRef@maps:eq:F.LT3\endcsname{\href{S6_SERIES_GRIND_V.pdf\#nameddest=tag-02H5}{V.{F.LT3} (p.~195)} }
\expandafter\def\csname SRef@maps:eq:F.LT4\endcsname{\href{S6_SERIES_GRIND_V.pdf\#nameddest=tag-02H6}{V.{F.LT4} (p.~195)} }
\expandafter\def\csname SRef@maps:eq:F.LT5\endcsname{\href{S6_SERIES_GRIND_V.pdf\#nameddest=tag-02H7}{V.{F.LT5} (p.~195)} }
\expandafter\def\csname SRef@maps:eq:F.LT6\endcsname{\href{S6_SERIES_GRIND_V.pdf\#nameddest=tag-02H8}{V.{F.LT6} (p.~196)} }
\expandafter\def\csname SRef@maps:eq:F.LT7\endcsname{\href{S6_SERIES_GRIND_V.pdf\#nameddest=tag-02H9}{V.{F.LT7} (p.~196)} }
\expandafter\def\csname SRef@maps:eq:F.LT8\endcsname{\href{S6_SERIES_GRIND_V.pdf\#nameddest=tag-02HA}{V.{F.LT8} (p.~196)} }
\expandafter\def\csname SRef@maps:eq:F.LT9\endcsname{\href{S6_SERIES_GRIND_V.pdf\#nameddest=tag-02HB}{V.{F.LT9} (p.~197)} }
\expandafter\def\csname SRef@maps:eq:F.CA1\endcsname{\href{S6_SERIES_GRIND_V.pdf\#nameddest=tag-02FE}{V.{F.CA1} (p.~197)} }
\expandafter\def\csname SRef@maps:eq:F.CA2\endcsname{\href{S6_SERIES_GRIND_V.pdf\#nameddest=tag-02FF}{V.{F.CA2} (p.~198)} }
\expandafter\def\csname SRef@maps:eq:F.CA3\endcsname{\href{S6_SERIES_GRIND_V.pdf\#nameddest=tag-02FG}{V.{F.CA3} (p.~198)} }
\expandafter\def\csname SRef@maps:eq:F.CA4\endcsname{\href{S6_SERIES_GRIND_V.pdf\#nameddest=tag-02FH}{V.{F.CA4} (p.~198)} }
\expandafter\def\csname SRef@maps:eq:F.CA5\endcsname{\href{S6_SERIES_GRIND_V.pdf\#nameddest=tag-02FI}{V.{F.CA5} (p.~198)} }
\expandafter\def\csname SRef@maps:eq:F.CK1\endcsname{\href{S6_SERIES_GRIND_V.pdf\#nameddest=tag-02FJ}{V.{F.CK1} (p.~200)} }
\expandafter\def\csname SRef@maps:eq:F.CK2\endcsname{\href{S6_SERIES_GRIND_V.pdf\#nameddest=tag-02FK}{V.{F.CK2} (p.~201)} }
\expandafter\def\csname SRef@maps:eq:F.GC1\endcsname{\href{S6_SERIES_GRIND_V.pdf\#nameddest=tag-02FZ}{V.{F.GC1} (p.~202)} }
\expandafter\def\csname SRef@maps:eq:F.GC2\endcsname{\href{S6_SERIES_GRIND_V.pdf\#nameddest=tag-02G0}{V.{F.GC2} (p.~203)} }
\expandafter\def\csname SRef@maps:eq:F.GC3\endcsname{\href{S6_SERIES_GRIND_V.pdf\#nameddest=tag-02G1}{V.{F.GC3} (p.~203)} }
\expandafter\def\csname SRef@maps:eq:F.GC4\endcsname{\href{S6_SERIES_GRIND_V.pdf\#nameddest=tag-02G2}{V.{F.GC4} (p.~203)} }
\expandafter\def\csname SRef@maps:eq:F.GC5\endcsname{\href{S6_SERIES_GRIND_V.pdf\#nameddest=tag-02G3}{V.{F.GC5} (p.~204)} }
\expandafter\def\csname SRef@maps:eq:F.GC5a\endcsname{\href{S6_SERIES_GRIND_V.pdf\#nameddest=tag-02G4}{V.{F.GC5a} (p.~204)} }
\expandafter\def\csname SRef@maps:eq:F.GC6\endcsname{\href{S6_SERIES_GRIND_V.pdf\#nameddest=tag-02G5}{V.{F.GC6} (p.~204)} }
\expandafter\def\csname SRef@maps:eq:F.GC7\endcsname{\href{S6_SERIES_GRIND_V.pdf\#nameddest=tag-02G6}{V.{F.GC7} (p.~205)} }
\expandafter\def\csname SRef@maps:eq:F.GC8\endcsname{\href{S6_SERIES_GRIND_V.pdf\#nameddest=tag-02G7}{V.{F.GC8} (p.~205)} }
\expandafter\def\csname SRef@maps:eq:F.GC9\endcsname{\href{S6_SERIES_GRIND_V.pdf\#nameddest=tag-02G8}{V.{F.GC9} (p.~205)} }
\expandafter\def\csname SRef@maps:eq:F.IC1\endcsname{\href{S6_SERIES_GRIND_V.pdf\#nameddest=tag-02GM}{V.{F.IC1} (p.~207)} }
\expandafter\def\csname SRef@maps:eq:F.IC2\endcsname{\href{S6_SERIES_GRIND_V.pdf\#nameddest=tag-02GN}{V.{F.IC2} (p.~207)} }
\expandafter\def\csname SRef@maps:eq:F.IC3\endcsname{\href{S6_SERIES_GRIND_V.pdf\#nameddest=tag-02GO}{V.{F.IC3} (p.~207)} }
\expandafter\def\csname SRef@maps:eq:F.CN1\endcsname{\href{S6_SERIES_GRIND_V.pdf\#nameddest=tag-02FL}{V.{F.CN1} (p.~209)} }
\expandafter\def\csname SRef@maps:eq:F.CN2\endcsname{\href{S6_SERIES_GRIND_V.pdf\#nameddest=tag-02FM}{V.{F.CN2} (p.~209)} }
\expandafter\def\csname SRef@maps:eq:F.CN3\endcsname{\href{S6_SERIES_GRIND_V.pdf\#nameddest=tag-02FN}{V.{F.CN3} (p.~210)} }
\expandafter\def\csname SRef@maps:eq:F.IV1\endcsname{\href{S6_SERIES_GRIND_V.pdf\#nameddest=tag-02GP}{V.{F.IV1} (p.~210)} }
\expandafter\def\csname SRef@maps:eq:F.IV2\endcsname{\href{S6_SERIES_GRIND_V.pdf\#nameddest=tag-02GQ}{V.{F.IV2} (p.~210)} }
\expandafter\def\csname SRef@maps:eq:F.IV3\endcsname{\href{S6_SERIES_GRIND_V.pdf\#nameddest=tag-02GR}{V.{F.IV3} (p.~211)} }
\expandafter\def\csname SRef@maps:eq:F.IV4\endcsname{\href{S6_SERIES_GRIND_V.pdf\#nameddest=tag-02GS}{V.{F.IV4} (p.~211)} }
\expandafter\def\csname SRef@maps:eq:F.SY1\endcsname{\href{S6_SERIES_GRIND_V.pdf\#nameddest=tag-02I6}{V.{F.SY1} (p.~212)} }
\expandafter\def\csname SRef@maps:eq:F.O1\endcsname{\href{S6_SERIES_GRIND_V.pdf\#nameddest=tag-02HC}{V.{F.O1} (p.~216)} }
\expandafter\def\csname SRef@maps:eq:F.O2\endcsname{\href{S6_SERIES_GRIND_V.pdf\#nameddest=tag-02HD}{V.{F.O2} (p.~216)} }
\expandafter\def\csname SRef@maps:eq:F.O3\endcsname{\href{S6_SERIES_GRIND_V.pdf\#nameddest=tag-02HE}{V.{F.O3} (p.~216)} }
\expandafter\def\csname SRef@maps:eq:F.O4\endcsname{\href{S6_SERIES_GRIND_V.pdf\#nameddest=tag-02HF}{V.{F.O4} (p.~217)} }
\expandafter\def\csname SRef@maps:eq:F.O5\endcsname{\href{S6_SERIES_GRIND_V.pdf\#nameddest=tag-02HG}{V.{F.O5} (p.~218)} }
\expandafter\def\csname SRef@maps:eq:F.O6\endcsname{\href{S6_SERIES_GRIND_V.pdf\#nameddest=tag-02HH}{V.{F.O6} (p.~219)} }
\expandafter\def\csname SRef@maps:eq:F.O7\endcsname{\href{S6_SERIES_GRIND_V.pdf\#nameddest=tag-02HI}{V.{F.O7} (p.~220)} }
\expandafter\def\csname SRef@maps:eq:F.7\endcsname{\href{S6_SERIES_GRIND_V.pdf\#nameddest=tag-02FB}{V.{F.7} (p.~104)} }
\expandafter\def\csname SRef@maps:eq:F.8\endcsname{\href{S6_SERIES_GRIND_V.pdf\#nameddest=tag-02FC}{V.{F.8} (p.~104)} }
\expandafter\def\csname SRef@maps:eq:F.9\endcsname{\href{S6_SERIES_GRIND_V.pdf\#nameddest=tag-02FD}{V.{F.9} (p.~105)} }
\expandafter\def\csname SRef@maps:eq:F.10\endcsname{\href{S6_SERIES_GRIND_V.pdf\#nameddest=tag-02EG}{V.{F.10} (p.~105)} }
\expandafter\def\csname SRef@maps:eq:F.11\endcsname{\href{S6_SERIES_GRIND_V.pdf\#nameddest=tag-02EH}{V.{F.11} (p.~106)} }
\expandafter\def\csname SRef@maps:eq:F.12\endcsname{\href{S6_SERIES_GRIND_V.pdf\#nameddest=tag-02EI}{V.{F.12} (p.~106)} }
\expandafter\def\csname SRef@maps:eq:F.13\endcsname{\href{S6_SERIES_GRIND_V.pdf\#nameddest=tag-02EJ}{V.{F.13} (p.~110)} }
\expandafter\def\csname SRef@maps:eq:F.14\endcsname{\href{S6_SERIES_GRIND_V.pdf\#nameddest=tag-02EK}{V.{F.14} (p.~110)} }
\expandafter\def\csname SRef@maps:eq:F.15\endcsname{\href{S6_SERIES_GRIND_V.pdf\#nameddest=tag-02EL}{V.{F.15} (p.~110)} }
\expandafter\def\csname SRef@maps:eq:F.16\endcsname{\href{S6_SERIES_GRIND_V.pdf\#nameddest=tag-02EM}{V.{F.16} (p.~110)} }
\expandafter\def\csname SRef@maps:eq:F.17\endcsname{\href{S6_SERIES_GRIND_V.pdf\#nameddest=tag-02EN}{V.{F.17} (p.~110)} }
\expandafter\def\csname SRef@maps:eq:F.18\endcsname{\href{S6_SERIES_GRIND_V.pdf\#nameddest=tag-02EO}{V.{F.18} (p.~110)} }
\expandafter\def\csname SRef@maps:eq:F.19\endcsname{\href{S6_SERIES_GRIND_V.pdf\#nameddest=tag-02EP}{V.{F.19} (p.~111)} }
\expandafter\def\csname SRef@maps:eq:F.20\endcsname{\href{S6_SERIES_GRIND_V.pdf\#nameddest=tag-02ER}{V.{F.20} (p.~111)} }
\expandafter\def\csname SRef@maps:eq:F.21\endcsname{\href{S6_SERIES_GRIND_V.pdf\#nameddest=tag-02ES}{V.{F.21} (p.~111)} }
\expandafter\def\csname SRef@maps:eq:F.22\endcsname{\href{S6_SERIES_GRIND_V.pdf\#nameddest=tag-02ET}{V.{F.22} (p.~111)} }
\expandafter\def\csname SRef@maps:eq:F.23\endcsname{\href{S6_SERIES_GRIND_V.pdf\#nameddest=tag-02EU}{V.{F.23} (p.~111)} }
\expandafter\def\csname SRef@maps:eq:F.24\endcsname{\href{S6_SERIES_GRIND_V.pdf\#nameddest=tag-02EV}{V.{F.24} (p.~115)} }
\expandafter\def\csname SRef@maps:eq:F.25\endcsname{\href{S6_SERIES_GRIND_V.pdf\#nameddest=tag-02EW}{V.{F.25} (p.~117)} }
\expandafter\def\csname SRef@maps:eq:F.3\endcsname{\href{S6_SERIES_GRIND_V.pdf\#nameddest=tag-02EZ}{V.{F.3} (p.~73)} }
\expandafter\def\csname SRef@maps:eq:F.27\endcsname{\href{S6_SERIES_GRIND_V.pdf\#nameddest=tag-02EX}{V.{F.27} (p.~147)} }
\expandafter\def\csname SRef@maps:eq:F.28\endcsname{\href{S6_SERIES_GRIND_V.pdf\#nameddest=tag-02EY}{V.{F.28} (p.~147)} }
\expandafter\def\csname SRef@maps:eq:F.4\endcsname{\href{S6_SERIES_GRIND_V.pdf\#nameddest=tag-02F7}{V.{F.4} (p.~75)} }
\expandafter\def\csname SRef@maps:eq:F.30\endcsname{\href{S6_SERIES_GRIND_V.pdf\#nameddest=tag-02F0}{V.{F.30} (p.~150)} }
\expandafter\def\csname SRef@maps:eq:F.31\endcsname{\href{S6_SERIES_GRIND_V.pdf\#nameddest=tag-02F1}{V.{F.31} (p.~150)} }
\expandafter\def\csname SRef@maps:eq:F.32\endcsname{\href{S6_SERIES_GRIND_V.pdf\#nameddest=tag-02F2}{V.{F.32} (p.~151)} }
\expandafter\def\csname SRef@maps:eq:F.33\endcsname{\href{S6_SERIES_GRIND_V.pdf\#nameddest=tag-02F3}{V.{F.33} (p.~152)} }
\expandafter\def\csname SRef@maps:eq:F.34\endcsname{\href{S6_SERIES_GRIND_V.pdf\#nameddest=tag-02F4}{V.{F.34} (p.~152)} }
\expandafter\def\csname SRef@maps:eq:F.36\endcsname{\href{S6_SERIES_GRIND_V.pdf\#nameddest=tag-02F5}{V.{F.36} (p.~155)} }
\expandafter\def\csname SRef@maps:eq:F.39\endcsname{\href{S6_SERIES_GRIND_V.pdf\#nameddest=tag-02F6}{V.{F.39} (p.~161)} }
\expandafter\def\csname SRef@maps:eq:F.40\endcsname{\href{S6_SERIES_GRIND_V.pdf\#nameddest=tag-02F8}{V.{F.40} (p.~162)} }
\expandafter\def\csname SRef@maps:eq:F.U1\endcsname{\href{S6_SERIES_GRIND_V.pdf\#nameddest=tag-02IK}{V.{F.U1} (p.~163)} }
\expandafter\def\csname SRef@maps:eq:F.5\endcsname{\href{S6_SERIES_GRIND_V.pdf\#nameddest=tag-02F9}{V.{F.5} (p.~76)} }
\expandafter\def\csname SRef@maps:eq:F.U2\endcsname{\href{S6_SERIES_GRIND_V.pdf\#nameddest=tag-02IO}{V.{F.U2} (p.~163)} }
\expandafter\def\csname SRef@maps:eq:F.U3\endcsname{\href{S6_SERIES_GRIND_V.pdf\#nameddest=tag-02IP}{V.{F.U3} (p.~163)} }
\expandafter\def\csname SRef@maps:eq:F.U4\endcsname{\href{S6_SERIES_GRIND_V.pdf\#nameddest=tag-02IQ}{V.{F.U4} (p.~164)} }
\expandafter\def\csname SRef@maps:eq:F.U5\endcsname{\href{S6_SERIES_GRIND_V.pdf\#nameddest=tag-02IR}{V.{F.U5} (p.~164)} }
\expandafter\def\csname SRef@maps:eq:F.U6\endcsname{\href{S6_SERIES_GRIND_V.pdf\#nameddest=tag-02IS}{V.{F.U6} (p.~164)} }
\expandafter\def\csname SRef@maps:eq:F.U7\endcsname{\href{S6_SERIES_GRIND_V.pdf\#nameddest=tag-02IT}{V.{F.U7} (p.~166)} }
\expandafter\def\csname SRef@maps:eq:F.U9\endcsname{\href{S6_SERIES_GRIND_V.pdf\#nameddest=tag-02IU}{V.{F.U9} (p.~166)} }
\expandafter\def\csname SRef@maps:eq:F.U10\endcsname{\href{S6_SERIES_GRIND_V.pdf\#nameddest=tag-02IL}{V.{F.U10} (p.~166)} }
\expandafter\def\csname SRef@maps:eq:F.U11\endcsname{\href{S6_SERIES_GRIND_V.pdf\#nameddest=tag-02IM}{V.{F.U11} (p.~167)} }
\expandafter\def\csname SRef@maps:eq:F.U12\endcsname{\href{S6_SERIES_GRIND_V.pdf\#nameddest=tag-02IN}{V.{F.U12} (p.~169)} }
\expandafter\def\csname SRef@maps:eq:FT.1\endcsname{\href{S6_SERIES_GRIND_V.pdf\#nameddest=tag-02J0}{V.{FT.1} (p.~81)} }
\expandafter\def\csname SRef@maps:eq:F.R1\endcsname{\href{S6_SERIES_GRIND_V.pdf\#nameddest=tag-02HJ}{V.{F.R1} (p.~170)} }
\expandafter\def\csname SRef@maps:eq:F.R2\endcsname{\href{S6_SERIES_GRIND_V.pdf\#nameddest=tag-02HP}{V.{F.R2} (p.~170)} }
\expandafter\def\csname SRef@maps:eq:F.R3\endcsname{\href{S6_SERIES_GRIND_V.pdf\#nameddest=tag-02HQ}{V.{F.R3} (p.~171)} }
\expandafter\def\csname SRef@maps:eq:F.R4\endcsname{\href{S6_SERIES_GRIND_V.pdf\#nameddest=tag-02HR}{V.{F.R4} (p.~171)} }
\expandafter\def\csname SRef@maps:eq:F.R5\endcsname{\href{S6_SERIES_GRIND_V.pdf\#nameddest=tag-02HS}{V.{F.R5} (p.~171)} }
\expandafter\def\csname SRef@maps:eq:F.R6\endcsname{\href{S6_SERIES_GRIND_V.pdf\#nameddest=tag-02HT}{V.{F.R6} (p.~171)} }
\expandafter\def\csname SRef@maps:eq:F.R7\endcsname{\href{S6_SERIES_GRIND_V.pdf\#nameddest=tag-02HU}{V.{F.R7} (p.~172)} }
\expandafter\def\csname SRef@maps:eq:F.R8\endcsname{\href{S6_SERIES_GRIND_V.pdf\#nameddest=tag-02HV}{V.{F.R8} (p.~172)} }
\expandafter\def\csname SRef@maps:eq:F.R9\endcsname{\href{S6_SERIES_GRIND_V.pdf\#nameddest=tag-02HW}{V.{F.R9} (p.~172)} }
\expandafter\def\csname SRef@maps:eq:F.L1\endcsname{\href{S6_SERIES_GRIND_V.pdf\#nameddest=tag-02GT}{V.{F.L1} (p.~174)} }
\expandafter\def\csname SRef@maps:mapsF:repair:cyclic-inclusion\endcsname{\href{S6_SERIES_GRIND_V.pdf\#nameddest=tag-02J1}{V.{{$\langle 1\rangle $6.}} (p.~199)} }
\expandafter\def\csname SRef@maps:mapsF:repair:log-cylinder-scope\endcsname{\href{S6_SERIES_GRIND_V.pdf\#nameddest=tag-02J2}{V.{{$\langle 1\rangle $5.}} (p.~201)} }
\expandafter\def\csname SRef@sec:fourfold\endcsname{\href{S6_SERIES_GRIND_V.pdf\#nameddest=tag-034O}{V.5 (p.~80)} }
\expandafter\def\csname SRef@maps:src:F.2.21\endcsname{\href{S6_SERIES_GRIND_V.pdf\#nameddest=tag-02NB}{V.11.15 (p.~173)} }
\expandafter\def\csname SRef@maps:src:F.2.30\endcsname{\href{S6_SERIES_GRIND_V.pdf\#nameddest=tag-02NL}{V.11.24 (p.~190)} }
\expandafter\def\csname SRef@maps:eq:F.1\endcsname{\href{S6_SERIES_GRIND_V.pdf\#nameddest=tag-02EF}{V.1 (p.~72)} }
\expandafter\def\csname SRef@maps:eq:F.2\endcsname{\href{S6_SERIES_GRIND_V.pdf\#nameddest=tag-02EQ}{V.2 (p.~72)} }
\expandafter\def\csname SRef@eq:P1Adegree\endcsname{\href{S6_SERIES_GRIND_V.pdf\#nameddest=tag-01O3}{V.3 (p.~134)} }
\expandafter\def\csname SRef@eq:P1Alogcotangent\endcsname{\href{S6_SERIES_GRIND_V.pdf\#nameddest=tag-01O4}{V.4 (p.~135)} }
\expandafter\def\csname SRef@maps:src:F.3.10\endcsname{\href{S6_SERIES_GRIND_V.pdf\#nameddest=tag-02O2}{V.12.5 (p.~218)} }
\expandafter\def\csname SRef@mapsaudit:order-four-torsion-transgression\endcsname{\href{S6_SERIES_GRIND_V.pdf\#nameddest=tag-02QF}{V.5.4 (p.~81)} }
\expandafter\def\csname SRef@maps:src:F.2.17\endcsname{\href{S6_SERIES_GRIND_V.pdf\#nameddest=tag-02N6}{V.11.11 (p.~166)} }
\expandafter\def\csname SRef@maps:src:F.2.8\endcsname{\href{S6_SERIES_GRIND_V.pdf\#nameddest=tag-02NZ}{V.11.2 (p.~147)} }
\expandafter\def\csname SRef@maps:src:F.2.9\endcsname{\href{S6_SERIES_GRIND_V.pdf\#nameddest=tag-02O0}{V.11.3 (p.~148)} }
\expandafter\def\csname SRef@maps:src:F.2.11\endcsname{\href{S6_SERIES_GRIND_V.pdf\#nameddest=tag-02N0}{V.11.5 (p.~151)} }
\expandafter\def\csname SRef@maps:src:F.2.15\endcsname{\href{S6_SERIES_GRIND_V.pdf\#nameddest=tag-02N4}{V.11.9 (p.~162)} }
\expandafter\def\csname SRef@maps:src:F.3.3\endcsname{\href{S6_SERIES_GRIND_V.pdf\#nameddest=tag-02O5}{V.12.1 (p.~215)} }
\expandafter\def\csname SRef@maps:src:F.1.4\endcsname{\href{S6_SERIES_GRIND_V.pdf\#nameddest=tag-02MX}{V.4.2 (p.~72)} }
\expandafter\def\csname SRef@maps:src:F.2.1\endcsname{\href{S6_SERIES_GRIND_V.pdf\#nameddest=tag-02MY}{V.4.3 (p.~73)} }
\expandafter\def\csname SRef@maps:src:F.2.2\endcsname{\href{S6_SERIES_GRIND_V.pdf\#nameddest=tag-02N9}{V.4.4 (p.~73)} }
\expandafter\def\csname SRef@maps:src:F.2.3\endcsname{\href{S6_SERIES_GRIND_V.pdf\#nameddest=tag-02NK}{V.4.5 (p.~74)} }
\expandafter\def\csname SRef@maps:src:F.3.2\endcsname{\href{S6_SERIES_GRIND_V.pdf\#nameddest=tag-02O4}{V.4.7 (p.~75)} }
\expandafter\def\csname SRef@maps:src:F.2.4\endcsname{\href{S6_SERIES_GRIND_V.pdf\#nameddest=tag-02NV}{V.4.8 (p.~75)} }
\expandafter\def\csname SRef@sec:V-known-inputs\endcsname{\href{S6_SERIES_GRIND_V.pdf\#nameddest=tag-034M}{V.3 (p.~72)} }
\expandafter\def\csname SRef@maps:src:F.2.7\endcsname{\href{S6_SERIES_GRIND_V.pdf\#nameddest=tag-02NY}{V.11.1 (p.~146)} }
\expandafter\def\csname SRef@maps:src:F.2.16\endcsname{\href{S6_SERIES_GRIND_V.pdf\#nameddest=tag-02N5}{V.11.10 (p.~163)} }
\expandafter\def\csname SRef@maps:src:F.2.18\endcsname{\href{S6_SERIES_GRIND_V.pdf\#nameddest=tag-02N7}{V.11.12 (p.~167)} }
\expandafter\def\csname SRef@maps:src:F.2.19\endcsname{\href{S6_SERIES_GRIND_V.pdf\#nameddest=tag-02N8}{V.11.13 (p.~169)} }
\expandafter\def\csname SRef@maps:src:F.2.20\endcsname{\href{S6_SERIES_GRIND_V.pdf\#nameddest=tag-02NA}{V.11.14 (p.~170)} }
\expandafter\def\csname SRef@maps:src:F.2.22\endcsname{\href{S6_SERIES_GRIND_V.pdf\#nameddest=tag-02NC}{V.11.16 (p.~174)} }
\expandafter\def\csname SRef@maps:src:F.2.23\endcsname{\href{S6_SERIES_GRIND_V.pdf\#nameddest=tag-02ND}{V.11.17 (p.~176)} }
\expandafter\def\csname SRef@maps:src:F.2.24\endcsname{\href{S6_SERIES_GRIND_V.pdf\#nameddest=tag-02NE}{V.11.18 (p.~177)} }
\expandafter\def\csname SRef@maps:src:F.2.25\endcsname{\href{S6_SERIES_GRIND_V.pdf\#nameddest=tag-02NF}{V.11.19 (p.~179)} }
\expandafter\def\csname SRef@maps:src:F.2.26\endcsname{\href{S6_SERIES_GRIND_V.pdf\#nameddest=tag-02NG}{V.11.20 (p.~181)} }
\expandafter\def\csname SRef@maps:src:F.2.27\endcsname{\href{S6_SERIES_GRIND_V.pdf\#nameddest=tag-02NH}{V.11.21 (p.~183)} }
\expandafter\def\csname SRef@maps:src:F.2.28\endcsname{\href{S6_SERIES_GRIND_V.pdf\#nameddest=tag-02NI}{V.11.22 (p.~185)} }
\expandafter\def\csname SRef@maps:src:F.2.29\endcsname{\href{S6_SERIES_GRIND_V.pdf\#nameddest=tag-02NJ}{V.11.23 (p.~186)} }
\expandafter\def\csname SRef@maps:src:F.2.31\endcsname{\href{S6_SERIES_GRIND_V.pdf\#nameddest=tag-02NM}{V.11.25 (p.~191)} }
\expandafter\def\csname SRef@maps:src:F.2.32\endcsname{\href{S6_SERIES_GRIND_V.pdf\#nameddest=tag-02NN}{V.11.26 (p.~195)} }
\expandafter\def\csname SRef@maps:src:F.2.33\endcsname{\href{S6_SERIES_GRIND_V.pdf\#nameddest=tag-02NO}{V.11.27 (p.~197)} }
\expandafter\def\csname SRef@maps:src:F.2.34\endcsname{\href{S6_SERIES_GRIND_V.pdf\#nameddest=tag-02NP}{V.11.28 (p.~200)} }
\expandafter\def\csname SRef@maps:src:F.2.35\endcsname{\href{S6_SERIES_GRIND_V.pdf\#nameddest=tag-02NQ}{V.11.29 (p.~202)} }
\expandafter\def\csname SRef@maps:src:F.2.36\endcsname{\href{S6_SERIES_GRIND_V.pdf\#nameddest=tag-02NR}{V.11.30 (p.~206)} }
\expandafter\def\csname SRef@maps:src:F.2.37\endcsname{\href{S6_SERIES_GRIND_V.pdf\#nameddest=tag-02NS}{V.11.31 (p.~208)} }
\expandafter\def\csname SRef@maps:src:F.2.38\endcsname{\href{S6_SERIES_GRIND_V.pdf\#nameddest=tag-02NT}{V.11.32 (p.~210)} }
\expandafter\def\csname SRef@maps:src:F.2.39\endcsname{\href{S6_SERIES_GRIND_V.pdf\#nameddest=tag-02NU}{V.11.33 (p.~212)} }
\expandafter\def\csname SRef@night:P1B.1\endcsname{\href{S6_SERIES_GRIND_V.pdf\#nameddest=tag-02R7}{V.11.34 (p.~213)} }
\expandafter\def\csname SRef@night:P1B.2\endcsname{\href{S6_SERIES_GRIND_V.pdf\#nameddest=tag-02R8}{V.11.35 (p.~214)} }
\expandafter\def\csname SRef@maps:src:F.2.10\endcsname{\href{S6_SERIES_GRIND_V.pdf\#nameddest=tag-02MZ}{V.11.4 (p.~150)} }
\expandafter\def\csname SRef@maps:src:F.2.12\endcsname{\href{S6_SERIES_GRIND_V.pdf\#nameddest=tag-02N1}{V.11.6 (p.~152)} }
\expandafter\def\csname SRef@maps:src:F.2.13\endcsname{\href{S6_SERIES_GRIND_V.pdf\#nameddest=tag-02N2}{V.11.7 (p.~155)} }
\expandafter\def\csname SRef@maps:src:F.2.14\endcsname{\href{S6_SERIES_GRIND_V.pdf\#nameddest=tag-02N3}{V.11.8 (p.~161)} }
\expandafter\def\csname SRef@maps:src:F.3.4\endcsname{\href{S6_SERIES_GRIND_V.pdf\#nameddest=tag-02O6}{V.12.2 (p.~215)} }
\expandafter\def\csname SRef@maps:src:F.3.8\endcsname{\href{S6_SERIES_GRIND_V.pdf\#nameddest=tag-02OA}{V.12.3 (p.~216)} }
\expandafter\def\csname SRef@maps:src:F.3.9\endcsname{\href{S6_SERIES_GRIND_V.pdf\#nameddest=tag-02OB}{V.12.4 (p.~217)} }
\expandafter\def\csname SRef@maps:src:F.3.11\endcsname{\href{S6_SERIES_GRIND_V.pdf\#nameddest=tag-02O3}{V.12.6 (p.~219)} }
\expandafter\def\csname SRef@maps:src:F.1.1\endcsname{\href{S6_SERIES_GRIND_V.pdf\#nameddest=tag-02MU}{V.4.1 (p.~72)} }
\expandafter\def\csname SRef@maps:src:F.3.1\endcsname{\href{S6_SERIES_GRIND_V.pdf\#nameddest=tag-02O1}{V.4.6 (p.~74)} }
\expandafter\def\csname SRef@maps:src:F.2.5\endcsname{\href{S6_SERIES_GRIND_V.pdf\#nameddest=tag-02NW}{V.5.1 (p.~76)} }
\expandafter\def\csname SRef@night:P1.1\endcsname{\href{S6_SERIES_GRIND_V.pdf\#nameddest=tag-02QI}{V.5.2 (p.~80)} }
\expandafter\def\csname SRef@night:P1.3\endcsname{\href{S6_SERIES_GRIND_V.pdf\#nameddest=tag-02QU}{V.5.3 (p.~80)} }
\expandafter\def\csname SRef@maps:src:F.2.6\endcsname{\href{S6_SERIES_GRIND_V.pdf\#nameddest=tag-02NX}{V.5.5 (p.~82)} }
\expandafter\def\csname SRef@night:P1.10\endcsname{\href{S6_SERIES_GRIND_V.pdf\#nameddest=tag-02QJ}{V.5.6 (p.~86)} }
\expandafter\def\csname SRef@night:P1.11\endcsname{\href{S6_SERIES_GRIND_V.pdf\#nameddest=tag-02QK}{V.5.7 (p.~87)} }
\expandafter\def\csname SRef@night:P1.12\endcsname{\href{S6_SERIES_GRIND_V.pdf\#nameddest=tag-02QL}{V.5.8 (p.~91)} }
\expandafter\def\csname SRef@night:P1.19\endcsname{\href{S6_SERIES_GRIND_V.pdf\#nameddest=tag-02QS}{V.5.9 (p.~94)} }
\expandafter\def\csname SRef@night:P1.4\endcsname{\href{S6_SERIES_GRIND_V.pdf\#nameddest=tag-02QV}{V.6.1 (p.~103)} }
\expandafter\def\csname SRef@night:P1.5\endcsname{\href{S6_SERIES_GRIND_V.pdf\#nameddest=tag-02QW}{V.6.2 (p.~104)} }
\expandafter\def\csname SRef@maps:src:F.3.5\endcsname{\href{S6_SERIES_GRIND_V.pdf\#nameddest=tag-02O7}{V.6.3 (p.~104)} }
\expandafter\def\csname SRef@night:P1.2\endcsname{\href{S6_SERIES_GRIND_V.pdf\#nameddest=tag-02QT}{V.6.4 (p.~106)} }
\expandafter\def\csname SRef@night:P1.7\endcsname{\href{S6_SERIES_GRIND_V.pdf\#nameddest=tag-02QY}{V.6.5 (p.~108)} }
\expandafter\def\csname SRef@maps:src:F.3.6\endcsname{\href{S6_SERIES_GRIND_V.pdf\#nameddest=tag-02O8}{V.6.6 (p.~109)} }
\expandafter\def\csname SRef@night:P1.8\endcsname{\href{S6_SERIES_GRIND_V.pdf\#nameddest=tag-02QZ}{V.6.7 (p.~112)} }
\expandafter\def\csname SRef@night:P1.9\endcsname{\href{S6_SERIES_GRIND_V.pdf\#nameddest=tag-02R0}{V.6.8 (p.~113)} }
\expandafter\def\csname SRef@maps:src:F.3.7\endcsname{\href{S6_SERIES_GRIND_V.pdf\#nameddest=tag-02O9}{V.6.9 (p.~114)} }
\expandafter\def\csname SRef@night:P1.13\endcsname{\href{S6_SERIES_GRIND_V.pdf\#nameddest=tag-02QM}{V.7.1 (p.~116)} }
\expandafter\def\csname SRef@maps:src:F.4.1\endcsname{\href{S6_SERIES_GRIND_V.pdf\#nameddest=tag-02OC}{V.7.2 (p.~117)} }
\expandafter\def\csname SRef@night:P1.15\endcsname{\href{S6_SERIES_GRIND_V.pdf\#nameddest=tag-02QO}{V.8.1 (p.~119)} }
\expandafter\def\csname SRef@night:P1.16\endcsname{\href{S6_SERIES_GRIND_V.pdf\#nameddest=tag-02QP}{V.8.2 (p.~121)} }
\expandafter\def\csname SRef@night:P1.18\endcsname{\href{S6_SERIES_GRIND_V.pdf\#nameddest=tag-02QR}{V.8.3 (p.~124)} }
\expandafter\def\csname SRef@night:P1A.1\endcsname{\href{S6_SERIES_GRIND_V.pdf\#nameddest=tag-02R1}{V.9.1 (p.~134)} }
\expandafter\def\csname SRef@night:P1A.2\endcsname{\href{S6_SERIES_GRIND_V.pdf\#nameddest=tag-02R2}{V.9.2 (p.~135)} }
\expandafter\def\csname SRef@night:P1A.3\endcsname{\href{S6_SERIES_GRIND_V.pdf\#nameddest=tag-02R3}{V.9.3 (p.~137)} }
\expandafter\def\csname SRef@night:P1A.4\endcsname{\href{S6_SERIES_GRIND_V.pdf\#nameddest=tag-02R4}{V.9.4 (p.~139)} }
\expandafter\def\csname SRef@night:P1.17\endcsname{\href{S6_SERIES_GRIND_V.pdf\#nameddest=tag-02QQ}{V.9.5 (p.~141)} }
\expandafter\def\csname SRef@night:P1A.5\endcsname{\href{S6_SERIES_GRIND_V.pdf\#nameddest=tag-02R5}{V.10.1 (p.~143)} }
\expandafter\def\csname SRef@night:P1A.6\endcsname{\href{S6_SERIES_GRIND_V.pdf\#nameddest=tag-02R6}{V.10.2 (p.~146)} }
\expandafter\def\csname SRef@cdp:eq:C.UN.1\endcsname{\href{S6_SERIES_GRIND_II.pdf\#nameddest=tag-0152}{II.{C.UN.1} (p.~413)} }
\expandafter\def\csname SRef@cdp:eq:C.FL.1\endcsname{\href{S6_SERIES_GRIND_II.pdf\#nameddest=tag-00KR}{II.{C.FL.1} (p.~443)} }
\expandafter\def\csname SRef@cdp:eq:C.QT.11\endcsname{\href{S6_SERIES_GRIND_II.pdf\#nameddest=tag-00WF}{II.{C.QT.11} (p.~708)} }
\expandafter\def\csname SRef@cdp:eq:C.QT.12\endcsname{\href{S6_SERIES_GRIND_II.pdf\#nameddest=tag-00WG}{II.{C.QT.12} (p.~708)} }
\expandafter\def\csname SRef@cdp:eq:C.QT.13\endcsname{\href{S6_SERIES_GRIND_II.pdf\#nameddest=tag-00WH}{II.{C.QT.13} (p.~708)} }
\expandafter\def\csname SRef@cdp:eq:C.WT.1\endcsname{\href{S6_SERIES_GRIND_II.pdf\#nameddest=tag-017O}{II.{C.WT.1} (p.~709)} }
\expandafter\def\csname SRef@cdp:eq:C.WT.2\endcsname{\href{S6_SERIES_GRIND_II.pdf\#nameddest=tag-017Z}{II.{C.WT.2} (p.~709)} }
\expandafter\def\csname SRef@cdp:eq:C.WT.3\endcsname{\href{S6_SERIES_GRIND_II.pdf\#nameddest=tag-018B}{II.{C.WT.3} (p.~709)} }
\expandafter\def\csname SRef@cdp:eq:C.WT.4\endcsname{\href{S6_SERIES_GRIND_II.pdf\#nameddest=tag-018M}{II.{C.WT.4} (p.~709)} }
\expandafter\def\csname SRef@cdp:eq:C.WT.5\endcsname{\href{S6_SERIES_GRIND_II.pdf\#nameddest=tag-018R}{II.{C.WT.5} (p.~709)} }
\expandafter\def\csname SRef@cdp:eq:C.WT.6\endcsname{\href{S6_SERIES_GRIND_II.pdf\#nameddest=tag-018S}{II.{C.WT.6} (p.~710)} }
\expandafter\def\csname SRef@cdp:eq:C.WT.7\endcsname{\href{S6_SERIES_GRIND_II.pdf\#nameddest=tag-018T}{II.{C.WT.7} (p.~710)} }
\expandafter\def\csname SRef@cdp:eq:C.FL.2\endcsname{\href{S6_SERIES_GRIND_II.pdf\#nameddest=tag-00KS}{II.{C.FL.2} (p.~443)} }
\expandafter\def\csname SRef@cdp:eq:C.WT.8\endcsname{\href{S6_SERIES_GRIND_II.pdf\#nameddest=tag-018U}{II.{C.WT.8} (p.~710)} }
\expandafter\def\csname SRef@cdp:eq:C.WT.9\endcsname{\href{S6_SERIES_GRIND_II.pdf\#nameddest=tag-018V}{II.{C.WT.9} (p.~710)} }
\expandafter\def\csname SRef@cdp:eq:C.WT.10\endcsname{\href{S6_SERIES_GRIND_II.pdf\#nameddest=tag-017P}{II.{C.WT.10} (p.~711)} }
\expandafter\def\csname SRef@cdp:eq:C.WT.11\endcsname{\href{S6_SERIES_GRIND_II.pdf\#nameddest=tag-017Q}{II.{C.WT.11} (p.~711)} }
\expandafter\def\csname SRef@cdp:eq:C.WT.12\endcsname{\href{S6_SERIES_GRIND_II.pdf\#nameddest=tag-017R}{II.{C.WT.12} (p.~711)} }
\expandafter\def\csname SRef@cdp:eq:C.WT.13\endcsname{\href{S6_SERIES_GRIND_II.pdf\#nameddest=tag-017S}{II.{C.WT.13} (p.~711)} }
\expandafter\def\csname SRef@cdp:eq:C.WT.14\endcsname{\href{S6_SERIES_GRIND_II.pdf\#nameddest=tag-017T}{II.{C.WT.14} (p.~712)} }
\expandafter\def\csname SRef@cdp:eq:C.WT.15\endcsname{\href{S6_SERIES_GRIND_II.pdf\#nameddest=tag-017U}{II.{C.WT.15} (p.~712)} }
\expandafter\def\csname SRef@cdp:eq:C.WT.16\endcsname{\href{S6_SERIES_GRIND_II.pdf\#nameddest=tag-017V}{II.{C.WT.16} (p.~712)} }
\expandafter\def\csname SRef@cdp:eq:C.WT.17\endcsname{\href{S6_SERIES_GRIND_II.pdf\#nameddest=tag-017W}{II.{C.WT.17} (p.~713)} }
\expandafter\def\csname SRef@cdp:eq:C.FL.3\endcsname{\href{S6_SERIES_GRIND_II.pdf\#nameddest=tag-00KT}{II.{C.FL.3} (p.~443)} }
\expandafter\def\csname SRef@cdp:eq:C.WT.18\endcsname{\href{S6_SERIES_GRIND_II.pdf\#nameddest=tag-017X}{II.{C.WT.18} (p.~713)} }
\expandafter\def\csname SRef@cdp:eq:C.WT.19\endcsname{\href{S6_SERIES_GRIND_II.pdf\#nameddest=tag-017Y}{II.{C.WT.19} (p.~713)} }
\expandafter\def\csname SRef@cdp:eq:C.WT.20\endcsname{\href{S6_SERIES_GRIND_II.pdf\#nameddest=tag-0180}{II.{C.WT.20} (p.~714)} }
\expandafter\def\csname SRef@cdp:eq:C.WT.21\endcsname{\href{S6_SERIES_GRIND_II.pdf\#nameddest=tag-0181}{II.{C.WT.21} (p.~714)} }
\expandafter\def\csname SRef@cdp:eq:C.WT.22\endcsname{\href{S6_SERIES_GRIND_II.pdf\#nameddest=tag-0182}{II.{C.WT.22} (p.~714)} }
\expandafter\def\csname SRef@cdp:eq:C.WT.23\endcsname{\href{S6_SERIES_GRIND_II.pdf\#nameddest=tag-0183}{II.{C.WT.23} (p.~714)} }
\expandafter\def\csname SRef@cdp:eq:C.WT.23a\endcsname{\href{S6_SERIES_GRIND_II.pdf\#nameddest=tag-0184}{II.{C.WT.23a} (p.~714)} }
\expandafter\def\csname SRef@cdp:eq:C.WT.24\endcsname{\href{S6_SERIES_GRIND_II.pdf\#nameddest=tag-0185}{II.{C.WT.24} (p.~714)} }
\expandafter\def\csname SRef@cdp:eq:C.WT.25\endcsname{\href{S6_SERIES_GRIND_II.pdf\#nameddest=tag-0186}{II.{C.WT.25} (p.~715)} }
\expandafter\def\csname SRef@cdp:eq:C.WT.26\endcsname{\href{S6_SERIES_GRIND_II.pdf\#nameddest=tag-0187}{II.{C.WT.26} (p.~715)} }
\expandafter\def\csname SRef@cdp:eq:C.FL.4\endcsname{\href{S6_SERIES_GRIND_II.pdf\#nameddest=tag-00KU}{II.{C.FL.4} (p.~443)} }
\expandafter\def\csname SRef@cdp:eq:C.WT.27\endcsname{\href{S6_SERIES_GRIND_II.pdf\#nameddest=tag-0188}{II.{C.WT.27} (p.~715)} }
\expandafter\def\csname SRef@cdp:eq:C.WT.28\endcsname{\href{S6_SERIES_GRIND_II.pdf\#nameddest=tag-0189}{II.{C.WT.28} (p.~715)} }
\expandafter\def\csname SRef@cdp:eq:C.WT.29\endcsname{\href{S6_SERIES_GRIND_II.pdf\#nameddest=tag-018A}{II.{C.WT.29} (p.~715)} }
\expandafter\def\csname SRef@cdp:eq:C.WT.30\endcsname{\href{S6_SERIES_GRIND_II.pdf\#nameddest=tag-018C}{II.{C.WT.30} (p.~715)} }
\expandafter\def\csname SRef@cdp:eq:C.WT.31\endcsname{\href{S6_SERIES_GRIND_II.pdf\#nameddest=tag-018D}{II.{C.WT.31} (p.~715)} }
\expandafter\def\csname SRef@cdp:eq:C.WT.32\endcsname{\href{S6_SERIES_GRIND_II.pdf\#nameddest=tag-018E}{II.{C.WT.32} (p.~716)} }
\expandafter\def\csname SRef@cdp:eq:C.WT.33\endcsname{\href{S6_SERIES_GRIND_II.pdf\#nameddest=tag-018F}{II.{C.WT.33} (p.~716)} }
\expandafter\def\csname SRef@cdp:eq:C.WT.34\endcsname{\href{S6_SERIES_GRIND_II.pdf\#nameddest=tag-018G}{II.{C.WT.34} (p.~716)} }
\expandafter\def\csname SRef@cdp:eq:C.WT.35\endcsname{\href{S6_SERIES_GRIND_II.pdf\#nameddest=tag-018H}{II.{C.WT.35} (p.~716)} }
\expandafter\def\csname SRef@cdp:eq:C.WT.36\endcsname{\href{S6_SERIES_GRIND_II.pdf\#nameddest=tag-018I}{II.{C.WT.36} (p.~716)} }
\expandafter\def\csname SRef@cdp:eq:C.FL.5\endcsname{\href{S6_SERIES_GRIND_II.pdf\#nameddest=tag-00KV}{II.{C.FL.5} (p.~444)} }
\expandafter\def\csname SRef@cdp:eq:C.WT.37\endcsname{\href{S6_SERIES_GRIND_II.pdf\#nameddest=tag-018J}{II.{C.WT.37} (p.~717)} }
\expandafter\def\csname SRef@cdp:eq:C.WT.38\endcsname{\href{S6_SERIES_GRIND_II.pdf\#nameddest=tag-018K}{II.{C.WT.38} (p.~717)} }
\expandafter\def\csname SRef@cdp:eq:C.WT.39\endcsname{\href{S6_SERIES_GRIND_II.pdf\#nameddest=tag-018L}{II.{C.WT.39} (p.~717)} }
\expandafter\def\csname SRef@cdp:eq:C.WT.40\endcsname{\href{S6_SERIES_GRIND_II.pdf\#nameddest=tag-018N}{II.{C.WT.40} (p.~717)} }
\expandafter\def\csname SRef@cdp:eq:C.WT.41\endcsname{\href{S6_SERIES_GRIND_II.pdf\#nameddest=tag-018O}{II.{C.WT.41} (p.~717)} }
\expandafter\def\csname SRef@cdp:eq:C.WT.42\endcsname{\href{S6_SERIES_GRIND_II.pdf\#nameddest=tag-018P}{II.{C.WT.42} (p.~718)} }
\expandafter\def\csname SRef@cdp:eq:C.WT.43\endcsname{\href{S6_SERIES_GRIND_II.pdf\#nameddest=tag-018Q}{II.{C.WT.43} (p.~718)} }
\expandafter\def\csname SRef@cdp:eq:C.NM.1\endcsname{\href{S6_SERIES_GRIND_II.pdf\#nameddest=tag-00SR}{II.{C.NM.1} (p.~719)} }
\expandafter\def\csname SRef@cdp:eq:C.NM.2\endcsname{\href{S6_SERIES_GRIND_II.pdf\#nameddest=tag-00ST}{II.{C.NM.2} (p.~719)} }
\expandafter\def\csname SRef@cdp:eq:C.NM.3\endcsname{\href{S6_SERIES_GRIND_II.pdf\#nameddest=tag-00SU}{II.{C.NM.3} (p.~719)} }
\expandafter\def\csname SRef@cdp:eq:C.PG.1\endcsname{\href{S6_SERIES_GRIND_II.pdf\#nameddest=tag-00UZ}{II.{C.PG.1} (p.~444)} }
\expandafter\def\csname SRef@cdp:eq:C.NM.4\endcsname{\href{S6_SERIES_GRIND_II.pdf\#nameddest=tag-00SV}{II.{C.NM.4} (p.~719)} }
\expandafter\def\csname SRef@cdp:eq:C.NM.5\endcsname{\href{S6_SERIES_GRIND_II.pdf\#nameddest=tag-00SW}{II.{C.NM.5} (p.~720)} }
\expandafter\def\csname SRef@cdp:eq:C.NM.6\endcsname{\href{S6_SERIES_GRIND_II.pdf\#nameddest=tag-00SX}{II.{C.NM.6} (p.~720)} }
\expandafter\def\csname SRef@cdp:eq:C.NM.7\endcsname{\href{S6_SERIES_GRIND_II.pdf\#nameddest=tag-00SY}{II.{C.NM.7} (p.~720)} }
\expandafter\def\csname SRef@cdp:eq:C.NM.8\endcsname{\href{S6_SERIES_GRIND_II.pdf\#nameddest=tag-00SZ}{II.{C.NM.8} (p.~720)} }
\expandafter\def\csname SRef@cdp:eq:C.NM.9\endcsname{\href{S6_SERIES_GRIND_II.pdf\#nameddest=tag-00T0}{II.{C.NM.9} (p.~721)} }
\expandafter\def\csname SRef@cdp:eq:C.NM.10\endcsname{\href{S6_SERIES_GRIND_II.pdf\#nameddest=tag-00SS}{II.{C.NM.10} (p.~721)} }
\expandafter\def\csname SRef@cdp:eq:C.NP.1\endcsname{\href{S6_SERIES_GRIND_II.pdf\#nameddest=tag-00T1}{II.{C.NP.1} (p.~722)} }
\expandafter\def\csname SRef@cdp:eq:C.NP.2\endcsname{\href{S6_SERIES_GRIND_II.pdf\#nameddest=tag-00T2}{II.{C.NP.2} (p.~722)} }
\expandafter\def\csname SRef@cdp:eq:C.NP.3\endcsname{\href{S6_SERIES_GRIND_II.pdf\#nameddest=tag-00T3}{II.{C.NP.3} (p.~723)} }
\expandafter\def\csname SRef@cdp:eq:C.PG.2\endcsname{\href{S6_SERIES_GRIND_II.pdf\#nameddest=tag-00V0}{II.{C.PG.2} (p.~444)} }
\expandafter\def\csname SRef@cdp:eq:C.NP.4\endcsname{\href{S6_SERIES_GRIND_II.pdf\#nameddest=tag-00T4}{II.{C.NP.4} (p.~723)} }
\expandafter\def\csname SRef@cdp:eq:C.NP.5\endcsname{\href{S6_SERIES_GRIND_II.pdf\#nameddest=tag-00T5}{II.{C.NP.5} (p.~723)} }
\expandafter\def\csname SRef@cdp:eq:C.NP.6\endcsname{\href{S6_SERIES_GRIND_II.pdf\#nameddest=tag-00T6}{II.{C.NP.6} (p.~724)} }
\expandafter\def\csname SRef@cdp:eq:C.NP.7\endcsname{\href{S6_SERIES_GRIND_II.pdf\#nameddest=tag-00T7}{II.{C.NP.7} (p.~724)} }
\expandafter\def\csname SRef@cdp:eq:C.NP.8\endcsname{\href{S6_SERIES_GRIND_II.pdf\#nameddest=tag-00T8}{II.{C.NP.8} (p.~725)} }
\expandafter\def\csname SRef@cdp:eq:C.NP.9\endcsname{\href{S6_SERIES_GRIND_II.pdf\#nameddest=tag-00T9}{II.{C.NP.9} (p.~725)} }
\expandafter\def\csname SRef@cdp:eq:C.TG.1\endcsname{\href{S6_SERIES_GRIND_II.pdf\#nameddest=tag-012Q}{II.{C.TG.1} (p.~726)} }
\expandafter\def\csname SRef@cdp:eq:C.TG.2\endcsname{\href{S6_SERIES_GRIND_II.pdf\#nameddest=tag-012X}{II.{C.TG.2} (p.~726)} }
\expandafter\def\csname SRef@cdp:eq:C.TG.3\endcsname{\href{S6_SERIES_GRIND_II.pdf\#nameddest=tag-012Y}{II.{C.TG.3} (p.~726)} }
\expandafter\def\csname SRef@cdp:eq:C.TG.4\endcsname{\href{S6_SERIES_GRIND_II.pdf\#nameddest=tag-012Z}{II.{C.TG.4} (p.~726)} }
\expandafter\def\csname SRef@cdp:eq:C.PG.3\endcsname{\href{S6_SERIES_GRIND_II.pdf\#nameddest=tag-00V1}{II.{C.PG.3} (p.~445)} }
\expandafter\def\csname SRef@cdp:eq:C.TG.5\endcsname{\href{S6_SERIES_GRIND_II.pdf\#nameddest=tag-0130}{II.{C.TG.5} (p.~726)} }
\expandafter\def\csname SRef@cdp:eq:C.TG.6\endcsname{\href{S6_SERIES_GRIND_II.pdf\#nameddest=tag-0131}{II.{C.TG.6} (p.~726)} }
\expandafter\def\csname SRef@cdp:eq:C.TG.7\endcsname{\href{S6_SERIES_GRIND_II.pdf\#nameddest=tag-0132}{II.{C.TG.7} (p.~727)} }
\expandafter\def\csname SRef@cdp:eq:C.TG.8\endcsname{\href{S6_SERIES_GRIND_II.pdf\#nameddest=tag-0133}{II.{C.TG.8} (p.~727)} }
\expandafter\def\csname SRef@cdp:eq:C.TG.9\endcsname{\href{S6_SERIES_GRIND_II.pdf\#nameddest=tag-0134}{II.{C.TG.9} (p.~728)} }
\expandafter\def\csname SRef@cdp:eq:C.TG.10\endcsname{\href{S6_SERIES_GRIND_II.pdf\#nameddest=tag-012R}{II.{C.TG.10} (p.~728)} }
\expandafter\def\csname SRef@cdp:eq:C.TG.11\endcsname{\href{S6_SERIES_GRIND_II.pdf\#nameddest=tag-012S}{II.{C.TG.11} (p.~728)} }
\expandafter\def\csname SRef@cdp:eq:C.TG.12\endcsname{\href{S6_SERIES_GRIND_II.pdf\#nameddest=tag-012T}{II.{C.TG.12} (p.~729)} }
\expandafter\def\csname SRef@cdp:eq:C.TG.13\endcsname{\href{S6_SERIES_GRIND_II.pdf\#nameddest=tag-012U}{II.{C.TG.13} (p.~729)} }
\expandafter\def\csname SRef@cdp:eq:C.TG.14\endcsname{\href{S6_SERIES_GRIND_II.pdf\#nameddest=tag-012V}{II.{C.TG.14} (p.~729)} }
\expandafter\def\csname SRef@cdp:eq:C.PG.4\endcsname{\href{S6_SERIES_GRIND_II.pdf\#nameddest=tag-00V2}{II.{C.PG.4} (p.~445)} }
\expandafter\def\csname SRef@cdp:eq:C.TG.15\endcsname{\href{S6_SERIES_GRIND_II.pdf\#nameddest=tag-012W}{II.{C.TG.15} (p.~729)} }
\expandafter\def\csname SRef@cdp:eq:C.TX.1\endcsname{\href{S6_SERIES_GRIND_II.pdf\#nameddest=tag-013Y}{II.{C.TX.1} (p.~730)} }
\expandafter\def\csname SRef@cdp:eq:C.TX.2\endcsname{\href{S6_SERIES_GRIND_II.pdf\#nameddest=tag-0144}{II.{C.TX.2} (p.~730)} }
\expandafter\def\csname SRef@cdp:eq:C.TX.3\endcsname{\href{S6_SERIES_GRIND_II.pdf\#nameddest=tag-0145}{II.{C.TX.3} (p.~730)} }
\expandafter\def\csname SRef@cdp:eq:C.TX.4\endcsname{\href{S6_SERIES_GRIND_II.pdf\#nameddest=tag-0146}{II.{C.TX.4} (p.~730)} }
\expandafter\def\csname SRef@cdp:eq:C.TX.5\endcsname{\href{S6_SERIES_GRIND_II.pdf\#nameddest=tag-0147}{II.{C.TX.5} (p.~731)} }
\expandafter\def\csname SRef@cdp:eq:C.TX.6\endcsname{\href{S6_SERIES_GRIND_II.pdf\#nameddest=tag-0148}{II.{C.TX.6} (p.~731)} }
\expandafter\def\csname SRef@cdp:eq:C.TX.7\endcsname{\href{S6_SERIES_GRIND_II.pdf\#nameddest=tag-0149}{II.{C.TX.7} (p.~731)} }
\expandafter\def\csname SRef@cdp:eq:C.TX.8\endcsname{\href{S6_SERIES_GRIND_II.pdf\#nameddest=tag-014A}{II.{C.TX.8} (p.~731)} }
\expandafter\def\csname SRef@cdp:eq:C.TX.9\endcsname{\href{S6_SERIES_GRIND_II.pdf\#nameddest=tag-014B}{II.{C.TX.9} (p.~731)} }
\expandafter\def\csname SRef@cdp:eq:C.PG.5\endcsname{\href{S6_SERIES_GRIND_II.pdf\#nameddest=tag-00V3}{II.{C.PG.5} (p.~445)} }
\expandafter\def\csname SRef@cdp:eq:C.TX.10\endcsname{\href{S6_SERIES_GRIND_II.pdf\#nameddest=tag-013Z}{II.{C.TX.10} (p.~731)} }
\expandafter\def\csname SRef@cdp:eq:C.TX.11\endcsname{\href{S6_SERIES_GRIND_II.pdf\#nameddest=tag-0140}{II.{C.TX.11} (p.~732)} }
\expandafter\def\csname SRef@cdp:eq:C.TX.12\endcsname{\href{S6_SERIES_GRIND_II.pdf\#nameddest=tag-0141}{II.{C.TX.12} (p.~732)} }
\expandafter\def\csname SRef@cdp:eq:C.TX.13\endcsname{\href{S6_SERIES_GRIND_II.pdf\#nameddest=tag-0142}{II.{C.TX.13} (p.~732)} }
\expandafter\def\csname SRef@cdp:eq:C.TX.14\endcsname{\href{S6_SERIES_GRIND_II.pdf\#nameddest=tag-0143}{II.{C.TX.14} (p.~732)} }
\expandafter\def\csname SRef@cdp:eq:C.WF.1\endcsname{\href{S6_SERIES_GRIND_II.pdf\#nameddest=tag-0165}{II.{C.WF.1} (p.~733)} }
\expandafter\def\csname SRef@cdp:eq:C.WF.2\endcsname{\href{S6_SERIES_GRIND_II.pdf\#nameddest=tag-016G}{II.{C.WF.2} (p.~733)} }
\expandafter\def\csname SRef@cdp:eq:C.WF.3\endcsname{\href{S6_SERIES_GRIND_II.pdf\#nameddest=tag-016Q}{II.{C.WF.3} (p.~733)} }
\expandafter\def\csname SRef@cdp:eq:C.WF.4\endcsname{\href{S6_SERIES_GRIND_II.pdf\#nameddest=tag-016R}{II.{C.WF.4} (p.~733)} }
\expandafter\def\csname SRef@cdp:eq:C.WF.5\endcsname{\href{S6_SERIES_GRIND_II.pdf\#nameddest=tag-016S}{II.{C.WF.5} (p.~733)} }
\expandafter\def\csname SRef@cdp:eq:C.UN.2\endcsname{\href{S6_SERIES_GRIND_II.pdf\#nameddest=tag-0153}{II.{C.UN.2} (p.~413)} }
\expandafter\def\csname SRef@cdp:eq:C.PG.6\endcsname{\href{S6_SERIES_GRIND_II.pdf\#nameddest=tag-00V4}{II.{C.PG.6} (p.~445)} }
\expandafter\def\csname SRef@cdp:eq:C.WF.6\endcsname{\href{S6_SERIES_GRIND_II.pdf\#nameddest=tag-016T}{II.{C.WF.6} (p.~734)} }
\expandafter\def\csname SRef@cdp:eq:C.WF.7\endcsname{\href{S6_SERIES_GRIND_II.pdf\#nameddest=tag-016U}{II.{C.WF.7} (p.~734)} }
\expandafter\def\csname SRef@cdp:eq:C.WF.8\endcsname{\href{S6_SERIES_GRIND_II.pdf\#nameddest=tag-016V}{II.{C.WF.8} (p.~734)} }
\expandafter\def\csname SRef@cdp:eq:C.WF.9\endcsname{\href{S6_SERIES_GRIND_II.pdf\#nameddest=tag-016W}{II.{C.WF.9} (p.~735)} }
\expandafter\def\csname SRef@cdp:eq:C.WF.10\endcsname{\href{S6_SERIES_GRIND_II.pdf\#nameddest=tag-0166}{II.{C.WF.10} (p.~735)} }
\expandafter\def\csname SRef@cdp:eq:C.WF.11\endcsname{\href{S6_SERIES_GRIND_II.pdf\#nameddest=tag-0167}{II.{C.WF.11} (p.~735)} }
\expandafter\def\csname SRef@cdp:eq:C.WF.12\endcsname{\href{S6_SERIES_GRIND_II.pdf\#nameddest=tag-0168}{II.{C.WF.12} (p.~735)} }
\expandafter\def\csname SRef@cdp:eq:C.WF.13\endcsname{\href{S6_SERIES_GRIND_II.pdf\#nameddest=tag-0169}{II.{C.WF.13} (p.~736)} }
\expandafter\def\csname SRef@cdp:eq:C.WF.14\endcsname{\href{S6_SERIES_GRIND_II.pdf\#nameddest=tag-016A}{II.{C.WF.14} (p.~736)} }
\expandafter\def\csname SRef@cdp:eq:C.WF.15\endcsname{\href{S6_SERIES_GRIND_II.pdf\#nameddest=tag-016B}{II.{C.WF.15} (p.~736)} }
\expandafter\def\csname SRef@cdp:eq:C.PG.7\endcsname{\href{S6_SERIES_GRIND_II.pdf\#nameddest=tag-00V5}{II.{C.PG.7} (p.~445)} }
\expandafter\def\csname SRef@cdp:eq:C.WF.16\endcsname{\href{S6_SERIES_GRIND_II.pdf\#nameddest=tag-016C}{II.{C.WF.16} (p.~736)} }
\expandafter\def\csname SRef@cdp:eq:C.WF.17\endcsname{\href{S6_SERIES_GRIND_II.pdf\#nameddest=tag-016D}{II.{C.WF.17} (p.~737)} }
\expandafter\def\csname SRef@cdp:eq:C.WF.18\endcsname{\href{S6_SERIES_GRIND_II.pdf\#nameddest=tag-016E}{II.{C.WF.18} (p.~737)} }
\expandafter\def\csname SRef@cdp:eq:C.WF.19\endcsname{\href{S6_SERIES_GRIND_II.pdf\#nameddest=tag-016F}{II.{C.WF.19} (p.~737)} }
\expandafter\def\csname SRef@cdp:eq:C.WF.20\endcsname{\href{S6_SERIES_GRIND_II.pdf\#nameddest=tag-016H}{II.{C.WF.20} (p.~737)} }
\expandafter\def\csname SRef@cdp:eq:C.WF.21\endcsname{\href{S6_SERIES_GRIND_II.pdf\#nameddest=tag-016I}{II.{C.WF.21} (p.~737)} }
\expandafter\def\csname SRef@cdp:eq:C.WF.22\endcsname{\href{S6_SERIES_GRIND_II.pdf\#nameddest=tag-016J}{II.{C.WF.22} (p.~737)} }
\expandafter\def\csname SRef@cdp:eq:C.WF.23\endcsname{\href{S6_SERIES_GRIND_II.pdf\#nameddest=tag-016K}{II.{C.WF.23} (p.~737)} }
\expandafter\def\csname SRef@cdp:eq:C.WF.24\endcsname{\href{S6_SERIES_GRIND_II.pdf\#nameddest=tag-016L}{II.{C.WF.24} (p.~737)} }
\expandafter\def\csname SRef@cdp:eq:C.WF.25\endcsname{\href{S6_SERIES_GRIND_II.pdf\#nameddest=tag-016M}{II.{C.WF.25} (p.~738)} }
\expandafter\def\csname SRef@cdp:eq:C.PG.8\endcsname{\href{S6_SERIES_GRIND_II.pdf\#nameddest=tag-00V6}{II.{C.PG.8} (p.~445)} }
\expandafter\def\csname SRef@cdp:eq:C.WF.26\endcsname{\href{S6_SERIES_GRIND_II.pdf\#nameddest=tag-016N}{II.{C.WF.26} (p.~738)} }
\expandafter\def\csname SRef@cdp:eq:C.WF.27\endcsname{\href{S6_SERIES_GRIND_II.pdf\#nameddest=tag-016O}{II.{C.WF.27} (p.~738)} }
\expandafter\def\csname SRef@cdp:eq:C.WF.28\endcsname{\href{S6_SERIES_GRIND_II.pdf\#nameddest=tag-016P}{II.{C.WF.28} (p.~738)} }
\expandafter\def\csname SRef@cdp:eq:C.RC.1\endcsname{\href{S6_SERIES_GRIND_II.pdf\#nameddest=tag-00XN}{II.{C.RC.1} (p.~739)} }
\expandafter\def\csname SRef@cdp:eq:C.RC.2\endcsname{\href{S6_SERIES_GRIND_II.pdf\#nameddest=tag-00XR}{II.{C.RC.2} (p.~739)} }
\expandafter\def\csname SRef@cdp:eq:C.RC.3\endcsname{\href{S6_SERIES_GRIND_II.pdf\#nameddest=tag-00XS}{II.{C.RC.3} (p.~739)} }
\expandafter\def\csname SRef@cdp:eq:C.RC.4\endcsname{\href{S6_SERIES_GRIND_II.pdf\#nameddest=tag-00XT}{II.{C.RC.4} (p.~740)} }
\expandafter\def\csname SRef@cdp:eq:C.RC.5\endcsname{\href{S6_SERIES_GRIND_II.pdf\#nameddest=tag-00XU}{II.{C.RC.5} (p.~740)} }
\expandafter\def\csname SRef@cdp:eq:C.RC.6\endcsname{\href{S6_SERIES_GRIND_II.pdf\#nameddest=tag-00XV}{II.{C.RC.6} (p.~740)} }
\expandafter\def\csname SRef@cdp:eq:C.RC.7\endcsname{\href{S6_SERIES_GRIND_II.pdf\#nameddest=tag-00XW}{II.{C.RC.7} (p.~740)} }
\expandafter\def\csname SRef@cdp:eq:C.PG.9\endcsname{\href{S6_SERIES_GRIND_II.pdf\#nameddest=tag-00V7}{II.{C.PG.9} (p.~446)} }
\expandafter\def\csname SRef@cdp:eq:C.RC.8\endcsname{\href{S6_SERIES_GRIND_II.pdf\#nameddest=tag-00XX}{II.{C.RC.8} (p.~740)} }
\expandafter\def\csname SRef@cdp:eq:C.RC.9\endcsname{\href{S6_SERIES_GRIND_II.pdf\#nameddest=tag-00XY}{II.{C.RC.9} (p.~741)} }
\expandafter\def\csname SRef@cdp:eq:C.RC.10\endcsname{\href{S6_SERIES_GRIND_II.pdf\#nameddest=tag-00XO}{II.{C.RC.10} (p.~741)} }
\expandafter\def\csname SRef@cdp:eq:C.RC.11\endcsname{\href{S6_SERIES_GRIND_II.pdf\#nameddest=tag-00XP}{II.{C.RC.11} (p.~742)} }
\expandafter\def\csname SRef@cdp:eq:C.RC.12\endcsname{\href{S6_SERIES_GRIND_II.pdf\#nameddest=tag-00XQ}{II.{C.RC.12} (p.~742)} }
\expandafter\def\csname SRef@cdp:eq:C.KG.1\endcsname{\href{S6_SERIES_GRIND_II.pdf\#nameddest=tag-00P8}{II.{C.KG.1} (p.~743)} }
\expandafter\def\csname SRef@cdp:eq:C.KG.2\endcsname{\href{S6_SERIES_GRIND_II.pdf\#nameddest=tag-00PA}{II.{C.KG.2} (p.~743)} }
\expandafter\def\csname SRef@cdp:eq:C.KG.3\endcsname{\href{S6_SERIES_GRIND_II.pdf\#nameddest=tag-00PB}{II.{C.KG.3} (p.~743)} }
\expandafter\def\csname SRef@cdp:eq:C.KG.4\endcsname{\href{S6_SERIES_GRIND_II.pdf\#nameddest=tag-00PC}{II.{C.KG.4} (p.~744)} }
\expandafter\def\csname SRef@cdp:eq:C.KG.5\endcsname{\href{S6_SERIES_GRIND_II.pdf\#nameddest=tag-00PD}{II.{C.KG.5} (p.~744)} }
\expandafter\def\csname SRef@cdp:eq:C.RG.1\endcsname{\href{S6_SERIES_GRIND_II.pdf\#nameddest=tag-00Z3}{II.{C.RG.1} (p.~446)} }
\expandafter\def\csname SRef@cdp:eq:C.KG.6\endcsname{\href{S6_SERIES_GRIND_II.pdf\#nameddest=tag-00PE}{II.{C.KG.6} (p.~744)} }
\expandafter\def\csname SRef@cdp:eq:C.KG.7\endcsname{\href{S6_SERIES_GRIND_II.pdf\#nameddest=tag-00PF}{II.{C.KG.7} (p.~745)} }
\expandafter\def\csname SRef@cdp:eq:C.KG.8\endcsname{\href{S6_SERIES_GRIND_II.pdf\#nameddest=tag-00PG}{II.{C.KG.8} (p.~745)} }
\expandafter\def\csname SRef@cdp:eq:C.KG.9\endcsname{\href{S6_SERIES_GRIND_II.pdf\#nameddest=tag-00PH}{II.{C.KG.9} (p.~747)} }
\expandafter\def\csname SRef@cdp:eq:C.KG.10\endcsname{\href{S6_SERIES_GRIND_II.pdf\#nameddest=tag-00P9}{II.{C.KG.10} (p.~747)} }
\expandafter\def\csname SRef@cdp:eq:C.HM.1\endcsname{\href{S6_SERIES_GRIND_II.pdf\#nameddest=tag-00MG}{II.{C.HM.1} (p.~748)} }
\expandafter\def\csname SRef@cdp:eq:C.HM.2\endcsname{\href{S6_SERIES_GRIND_II.pdf\#nameddest=tag-00MI}{II.{C.HM.2} (p.~748)} }
\expandafter\def\csname SRef@cdp:eq:C.HM.3\endcsname{\href{S6_SERIES_GRIND_II.pdf\#nameddest=tag-00MJ}{II.{C.HM.3} (p.~748)} }
\expandafter\def\csname SRef@cdp:eq:C.HM.4\endcsname{\href{S6_SERIES_GRIND_II.pdf\#nameddest=tag-00MK}{II.{C.HM.4} (p.~749)} }
\expandafter\def\csname SRef@cdp:eq:C.HM.5\endcsname{\href{S6_SERIES_GRIND_II.pdf\#nameddest=tag-00ML}{II.{C.HM.5} (p.~749)} }
\expandafter\def\csname SRef@cdp:eq:C.RG.2\endcsname{\href{S6_SERIES_GRIND_II.pdf\#nameddest=tag-00Z8}{II.{C.RG.2} (p.~446)} }
\expandafter\def\csname SRef@cdp:eq:C.HM.6\endcsname{\href{S6_SERIES_GRIND_II.pdf\#nameddest=tag-00MM}{II.{C.HM.6} (p.~749)} }
\expandafter\def\csname SRef@cdp:eq:C.HM.7\endcsname{\href{S6_SERIES_GRIND_II.pdf\#nameddest=tag-00MN}{II.{C.HM.7} (p.~749)} }
\expandafter\def\csname SRef@cdp:eq:C.HM.8\endcsname{\href{S6_SERIES_GRIND_II.pdf\#nameddest=tag-00MO}{II.{C.HM.8} (p.~750)} }
\expandafter\def\csname SRef@cdp:eq:C.HM.9\endcsname{\href{S6_SERIES_GRIND_II.pdf\#nameddest=tag-00MP}{II.{C.HM.9} (p.~750)} }
\expandafter\def\csname SRef@cdp:eq:C.HM.10\endcsname{\href{S6_SERIES_GRIND_II.pdf\#nameddest=tag-00MH}{II.{C.HM.10} (p.~751)} }
\expandafter\def\csname SRef@cdp:eq:C.BC.1\endcsname{\href{S6_SERIES_GRIND_II.pdf\#nameddest=tag-00AZ}{II.{C.BC.1} (p.~752)} }
\expandafter\def\csname SRef@cdp:eq:C.BC.2\endcsname{\href{S6_SERIES_GRIND_II.pdf\#nameddest=tag-00B4}{II.{C.BC.2} (p.~752)} }
\expandafter\def\csname SRef@cdp:eq:C.BC.3\endcsname{\href{S6_SERIES_GRIND_II.pdf\#nameddest=tag-00B5}{II.{C.BC.3} (p.~752)} }
\expandafter\def\csname SRef@cdp:eq:C.BC.4\endcsname{\href{S6_SERIES_GRIND_II.pdf\#nameddest=tag-00B6}{II.{C.BC.4} (p.~752)} }
\expandafter\def\csname SRef@cdp:eq:C.BC.5\endcsname{\href{S6_SERIES_GRIND_II.pdf\#nameddest=tag-00B7}{II.{C.BC.5} (p.~753)} }
\expandafter\def\csname SRef@cdp:eq:C.RG.3\endcsname{\href{S6_SERIES_GRIND_II.pdf\#nameddest=tag-00Z9}{II.{C.RG.3} (p.~447)} }
\expandafter\def\csname SRef@cdp:eq:C.BC.6\endcsname{\href{S6_SERIES_GRIND_II.pdf\#nameddest=tag-00B8}{II.{C.BC.6} (p.~753)} }
\expandafter\def\csname SRef@cdp:eq:C.BC.7\endcsname{\href{S6_SERIES_GRIND_II.pdf\#nameddest=tag-00B9}{II.{C.BC.7} (p.~753)} }
\expandafter\def\csname SRef@cdp:eq:C.BC.8\endcsname{\href{S6_SERIES_GRIND_II.pdf\#nameddest=tag-00BA}{II.{C.BC.8} (p.~753)} }
\expandafter\def\csname SRef@cdp:eq:C.BC.9\endcsname{\href{S6_SERIES_GRIND_II.pdf\#nameddest=tag-00BB}{II.{C.BC.9} (p.~753)} }
\expandafter\def\csname SRef@cdp:eq:C.BC.10\endcsname{\href{S6_SERIES_GRIND_II.pdf\#nameddest=tag-00B0}{II.{C.BC.10} (p.~754)} }
\expandafter\def\csname SRef@cdp:eq:C.BC.11\endcsname{\href{S6_SERIES_GRIND_II.pdf\#nameddest=tag-00B1}{II.{C.BC.11} (p.~754)} }
\expandafter\def\csname SRef@cdp:eq:C.BC.12\endcsname{\href{S6_SERIES_GRIND_II.pdf\#nameddest=tag-00B2}{II.{C.BC.12} (p.~754)} }
\expandafter\def\csname SRef@cdp:eq:C.BC.13\endcsname{\href{S6_SERIES_GRIND_II.pdf\#nameddest=tag-00B3}{II.{C.BC.13} (p.~754)} }
\expandafter\def\csname SRef@cdp:eq:C.EP.1\endcsname{\href{S6_SERIES_GRIND_II.pdf\#nameddest=tag-00I1}{II.{C.EP.1} (p.~755)} }
\expandafter\def\csname SRef@cdp:eq:C.EP.2\endcsname{\href{S6_SERIES_GRIND_II.pdf\#nameddest=tag-00I4}{II.{C.EP.2} (p.~755)} }
\expandafter\def\csname SRef@cdp:eq:C.RG.4\endcsname{\href{S6_SERIES_GRIND_II.pdf\#nameddest=tag-00ZA}{II.{C.RG.4} (p.~447)} }
\expandafter\def\csname SRef@cdp:eq:C.EP.3\endcsname{\href{S6_SERIES_GRIND_II.pdf\#nameddest=tag-00I5}{II.{C.EP.3} (p.~755)} }
\expandafter\def\csname SRef@cdp:eq:C.EP.4\endcsname{\href{S6_SERIES_GRIND_II.pdf\#nameddest=tag-00I6}{II.{C.EP.4} (p.~756)} }
\expandafter\def\csname SRef@cdp:eq:C.EP.5\endcsname{\href{S6_SERIES_GRIND_II.pdf\#nameddest=tag-00I7}{II.{C.EP.5} (p.~756)} }
\expandafter\def\csname SRef@cdp:eq:C.EP.6\endcsname{\href{S6_SERIES_GRIND_II.pdf\#nameddest=tag-00I8}{II.{C.EP.6} (p.~756)} }
\expandafter\def\csname SRef@cdp:eq:C.EP.7\endcsname{\href{S6_SERIES_GRIND_II.pdf\#nameddest=tag-00I9}{II.{C.EP.7} (p.~756)} }
\expandafter\def\csname SRef@cdp:eq:C.EP.8\endcsname{\href{S6_SERIES_GRIND_II.pdf\#nameddest=tag-00IA}{II.{C.EP.8} (p.~756)} }
\expandafter\def\csname SRef@cdp:eq:C.EP.9\endcsname{\href{S6_SERIES_GRIND_II.pdf\#nameddest=tag-00IB}{II.{C.EP.9} (p.~757)} }
\expandafter\def\csname SRef@cdp:eq:C.EP.10\endcsname{\href{S6_SERIES_GRIND_II.pdf\#nameddest=tag-00I2}{II.{C.EP.10} (p.~757)} }
\expandafter\def\csname SRef@cdp:eq:C.EP.11\endcsname{\href{S6_SERIES_GRIND_II.pdf\#nameddest=tag-00I3}{II.{C.EP.11} (p.~757)} }
\expandafter\def\csname SRef@cdp:eq:C.AP.1\endcsname{\href{S6_SERIES_GRIND_II.pdf\#nameddest=tag-009M}{II.{C.AP.1} (p.~757)} }
\expandafter\def\csname SRef@cdp:eq:C.RG.5\endcsname{\href{S6_SERIES_GRIND_II.pdf\#nameddest=tag-00ZB}{II.{C.RG.5} (p.~447)} }
\expandafter\def\csname SRef@cdp:eq:C.AP.2\endcsname{\href{S6_SERIES_GRIND_II.pdf\#nameddest=tag-009S}{II.{C.AP.2} (p.~758)} }
\expandafter\def\csname SRef@cdp:eq:C.AP.3\endcsname{\href{S6_SERIES_GRIND_II.pdf\#nameddest=tag-009T}{II.{C.AP.3} (p.~758)} }
\expandafter\def\csname SRef@cdp:eq:C.AP.4\endcsname{\href{S6_SERIES_GRIND_II.pdf\#nameddest=tag-009U}{II.{C.AP.4} (p.~758)} }
\expandafter\def\csname SRef@cdp:eq:C.AP.5\endcsname{\href{S6_SERIES_GRIND_II.pdf\#nameddest=tag-009V}{II.{C.AP.5} (p.~759)} }
\expandafter\def\csname SRef@cdp:eq:C.AP.6\endcsname{\href{S6_SERIES_GRIND_II.pdf\#nameddest=tag-009W}{II.{C.AP.6} (p.~759)} }
\expandafter\def\csname SRef@cdp:eq:C.AP.7\endcsname{\href{S6_SERIES_GRIND_II.pdf\#nameddest=tag-009X}{II.{C.AP.7} (p.~759)} }
\expandafter\def\csname SRef@cdp:eq:C.AP.8\endcsname{\href{S6_SERIES_GRIND_II.pdf\#nameddest=tag-009Y}{II.{C.AP.8} (p.~760)} }
\expandafter\def\csname SRef@cdp:eq:C.AP.9\endcsname{\href{S6_SERIES_GRIND_II.pdf\#nameddest=tag-009Z}{II.{C.AP.9} (p.~760)} }
\expandafter\def\csname SRef@cdp:eq:C.AP.10\endcsname{\href{S6_SERIES_GRIND_II.pdf\#nameddest=tag-009N}{II.{C.AP.10} (p.~760)} }
\expandafter\def\csname SRef@cdp:eq:C.AP.11\endcsname{\href{S6_SERIES_GRIND_II.pdf\#nameddest=tag-009O}{II.{C.AP.11} (p.~760)} }
\expandafter\def\csname SRef@cdp:eq:C.RG.6\endcsname{\href{S6_SERIES_GRIND_II.pdf\#nameddest=tag-00ZC}{II.{C.RG.6} (p.~448)} }
\expandafter\def\csname SRef@cdp:eq:C.AP.12\endcsname{\href{S6_SERIES_GRIND_II.pdf\#nameddest=tag-009P}{II.{C.AP.12} (p.~761)} }
\expandafter\def\csname SRef@cdp:eq:C.AP.13\endcsname{\href{S6_SERIES_GRIND_II.pdf\#nameddest=tag-009Q}{II.{C.AP.13} (p.~761)} }
\expandafter\def\csname SRef@cdp:eq:C.AP.14\endcsname{\href{S6_SERIES_GRIND_II.pdf\#nameddest=tag-009R}{II.{C.AP.14} (p.~761)} }
\expandafter\def\csname SRef@cdp:eq:C.ML.1\endcsname{\href{S6_SERIES_GRIND_II.pdf\#nameddest=tag-00QR}{II.{C.ML.1} (p.~762)} }
\expandafter\def\csname SRef@cdp:eq:C.ML.2\endcsname{\href{S6_SERIES_GRIND_II.pdf\#nameddest=tag-00QT}{II.{C.ML.2} (p.~762)} }
\expandafter\def\csname SRef@cdp:eq:C.ML.3\endcsname{\href{S6_SERIES_GRIND_II.pdf\#nameddest=tag-00QU}{II.{C.ML.3} (p.~763)} }
\expandafter\def\csname SRef@cdp:eq:C.ML.4\endcsname{\href{S6_SERIES_GRIND_II.pdf\#nameddest=tag-00QV}{II.{C.ML.4} (p.~763)} }
\expandafter\def\csname SRef@cdp:eq:C.ML.5\endcsname{\href{S6_SERIES_GRIND_II.pdf\#nameddest=tag-00QW}{II.{C.ML.5} (p.~763)} }
\expandafter\def\csname SRef@cdp:eq:C.ML.6\endcsname{\href{S6_SERIES_GRIND_II.pdf\#nameddest=tag-00QX}{II.{C.ML.6} (p.~763)} }
\expandafter\def\csname SRef@cdp:eq:C.ML.7\endcsname{\href{S6_SERIES_GRIND_II.pdf\#nameddest=tag-00QY}{II.{C.ML.7} (p.~764)} }
\expandafter\def\csname SRef@cdp:eq:C.UN.3\endcsname{\href{S6_SERIES_GRIND_II.pdf\#nameddest=tag-0154}{II.{C.UN.3} (p.~413)} }
\expandafter\def\csname SRef@cdp:eq:C.RG.7\endcsname{\href{S6_SERIES_GRIND_II.pdf\#nameddest=tag-00ZD}{II.{C.RG.7} (p.~448)} }
\expandafter\def\csname SRef@cdp:eq:C.ML.8\endcsname{\href{S6_SERIES_GRIND_II.pdf\#nameddest=tag-00QZ}{II.{C.ML.8} (p.~764)} }
\expandafter\def\csname SRef@cdp:eq:C.ML.9\endcsname{\href{S6_SERIES_GRIND_II.pdf\#nameddest=tag-00R0}{II.{C.ML.9} (p.~764)} }
\expandafter\def\csname SRef@cdp:eq:C.ML.10\endcsname{\href{S6_SERIES_GRIND_II.pdf\#nameddest=tag-00QS}{II.{C.ML.10} (p.~765)} }
\expandafter\def\csname SRef@cdp:eq:C.TI.1\endcsname{\href{S6_SERIES_GRIND_II.pdf\#nameddest=tag-0135}{II.{C.TI.1} (p.~766)} }
\expandafter\def\csname SRef@cdp:eq:C.TI.2\endcsname{\href{S6_SERIES_GRIND_II.pdf\#nameddest=tag-013B}{II.{C.TI.2} (p.~766)} }
\expandafter\def\csname SRef@cdp:eq:C.TI.3\endcsname{\href{S6_SERIES_GRIND_II.pdf\#nameddest=tag-013C}{II.{C.TI.3} (p.~766)} }
\expandafter\def\csname SRef@cdp:eq:C.TI.4\endcsname{\href{S6_SERIES_GRIND_II.pdf\#nameddest=tag-013D}{II.{C.TI.4} (p.~766)} }
\expandafter\def\csname SRef@cdp:eq:C.TI.5\endcsname{\href{S6_SERIES_GRIND_II.pdf\#nameddest=tag-013E}{II.{C.TI.5} (p.~766)} }
\expandafter\def\csname SRef@cdp:eq:C.TI.6\endcsname{\href{S6_SERIES_GRIND_II.pdf\#nameddest=tag-013F}{II.{C.TI.6} (p.~767)} }
\expandafter\def\csname SRef@cdp:eq:C.TI.7\endcsname{\href{S6_SERIES_GRIND_II.pdf\#nameddest=tag-013G}{II.{C.TI.7} (p.~767)} }
\expandafter\def\csname SRef@cdp:eq:C.RG.8\endcsname{\href{S6_SERIES_GRIND_II.pdf\#nameddest=tag-00ZE}{II.{C.RG.8} (p.~448)} }
\expandafter\def\csname SRef@cdp:eq:C.TI.8\endcsname{\href{S6_SERIES_GRIND_II.pdf\#nameddest=tag-013H}{II.{C.TI.8} (p.~767)} }
\expandafter\def\csname SRef@cdp:eq:C.TI.9\endcsname{\href{S6_SERIES_GRIND_II.pdf\#nameddest=tag-013I}{II.{C.TI.9} (p.~768)} }
\expandafter\def\csname SRef@cdp:eq:C.TI.10\endcsname{\href{S6_SERIES_GRIND_II.pdf\#nameddest=tag-0136}{II.{C.TI.10} (p.~768)} }
\expandafter\def\csname SRef@cdp:eq:C.TI.11\endcsname{\href{S6_SERIES_GRIND_II.pdf\#nameddest=tag-0137}{II.{C.TI.11} (p.~768)} }
\expandafter\def\csname SRef@cdp:eq:C.TI.12\endcsname{\href{S6_SERIES_GRIND_II.pdf\#nameddest=tag-0138}{II.{C.TI.12} (p.~768)} }
\expandafter\def\csname SRef@cdp:eq:C.TI.13\endcsname{\href{S6_SERIES_GRIND_II.pdf\#nameddest=tag-0139}{II.{C.TI.13} (p.~769)} }
\expandafter\def\csname SRef@cdp:eq:C.TI.14\endcsname{\href{S6_SERIES_GRIND_II.pdf\#nameddest=tag-013A}{II.{C.TI.14} (p.~769)} }
\expandafter\def\csname SRef@cdp:eq:C.TK.1\endcsname{\href{S6_SERIES_GRIND_II.pdf\#nameddest=tag-013J}{II.{C.TK.1} (p.~770)} }
\expandafter\def\csname SRef@cdp:eq:C.TK.2\endcsname{\href{S6_SERIES_GRIND_II.pdf\#nameddest=tag-013K}{II.{C.TK.2} (p.~770)} }
\expandafter\def\csname SRef@cdp:eq:C.TK.3\endcsname{\href{S6_SERIES_GRIND_II.pdf\#nameddest=tag-013L}{II.{C.TK.3} (p.~771)} }
\expandafter\def\csname SRef@cdp:eq:C.RG.9\endcsname{\href{S6_SERIES_GRIND_II.pdf\#nameddest=tag-00ZF}{II.{C.RG.9} (p.~448)} }
\expandafter\def\csname SRef@cdp:eq:C.TK.4\endcsname{\href{S6_SERIES_GRIND_II.pdf\#nameddest=tag-013M}{II.{C.TK.4} (p.~771)} }
\expandafter\def\csname SRef@cdp:eq:C.TK.5\endcsname{\href{S6_SERIES_GRIND_II.pdf\#nameddest=tag-013N}{II.{C.TK.5} (p.~771)} }
\expandafter\def\csname SRef@cdp:eq:C.TK.6\endcsname{\href{S6_SERIES_GRIND_II.pdf\#nameddest=tag-013O}{II.{C.TK.6} (p.~772)} }
\expandafter\def\csname SRef@cdp:eq:C.TK.7\endcsname{\href{S6_SERIES_GRIND_II.pdf\#nameddest=tag-013P}{II.{C.TK.7} (p.~772)} }
\expandafter\def\csname SRef@cdp:eq:C.RK.1\endcsname{\href{S6_SERIES_GRIND_II.pdf\#nameddest=tag-00ZQ}{II.{C.RK.1} (p.~773)} }
\expandafter\def\csname SRef@cdp:eq:C.RK.2\endcsname{\href{S6_SERIES_GRIND_II.pdf\#nameddest=tag-0101}{II.{C.RK.2} (p.~773)} }
\expandafter\def\csname SRef@cdp:eq:C.RK.3\endcsname{\href{S6_SERIES_GRIND_II.pdf\#nameddest=tag-0105}{II.{C.RK.3} (p.~773)} }
\expandafter\def\csname SRef@cdp:eq:C.RK.4\endcsname{\href{S6_SERIES_GRIND_II.pdf\#nameddest=tag-0106}{II.{C.RK.4} (p.~773)} }
\expandafter\def\csname SRef@cdp:eq:C.RK.5\endcsname{\href{S6_SERIES_GRIND_II.pdf\#nameddest=tag-0107}{II.{C.RK.5} (p.~773)} }
\expandafter\def\csname SRef@cdp:eq:C.RK.6\endcsname{\href{S6_SERIES_GRIND_II.pdf\#nameddest=tag-0108}{II.{C.RK.6} (p.~774)} }
\expandafter\def\csname SRef@cdp:eq:C.RG.10\endcsname{\href{S6_SERIES_GRIND_II.pdf\#nameddest=tag-00Z4}{II.{C.RG.10} (p.~448)} }
\expandafter\def\csname SRef@cdp:eq:C.RK.7\endcsname{\href{S6_SERIES_GRIND_II.pdf\#nameddest=tag-0109}{II.{C.RK.7} (p.~774)} }
\expandafter\def\csname SRef@cdp:eq:C.RK.8\endcsname{\href{S6_SERIES_GRIND_II.pdf\#nameddest=tag-010A}{II.{C.RK.8} (p.~774)} }
\expandafter\def\csname SRef@cdp:eq:C.RK.18\endcsname{\href{S6_SERIES_GRIND_II.pdf\#nameddest=tag-00ZZ}{II.{C.RK.18} (p.~774)} }
\expandafter\def\csname SRef@cdp:eq:C.RK.19\endcsname{\href{S6_SERIES_GRIND_II.pdf\#nameddest=tag-0100}{II.{C.RK.19} (p.~774)} }
\expandafter\def\csname SRef@cdp:eq:C.RK.20\endcsname{\href{S6_SERIES_GRIND_II.pdf\#nameddest=tag-0102}{II.{C.RK.20} (p.~774)} }
\expandafter\def\csname SRef@cdp:eq:C.RK.9\endcsname{\href{S6_SERIES_GRIND_II.pdf\#nameddest=tag-010B}{II.{C.RK.9} (p.~775)} }
\expandafter\def\csname SRef@cdp:eq:C.RK.10\endcsname{\href{S6_SERIES_GRIND_II.pdf\#nameddest=tag-00ZR}{II.{C.RK.10} (p.~775)} }
\expandafter\def\csname SRef@cdp:eq:C.RK.11\endcsname{\href{S6_SERIES_GRIND_II.pdf\#nameddest=tag-00ZS}{II.{C.RK.11} (p.~775)} }
\expandafter\def\csname SRef@cdp:eq:C.RK.12\endcsname{\href{S6_SERIES_GRIND_II.pdf\#nameddest=tag-00ZT}{II.{C.RK.12} (p.~776)} }
\expandafter\def\csname SRef@cdp:eq:C.RK.13\endcsname{\href{S6_SERIES_GRIND_II.pdf\#nameddest=tag-00ZU}{II.{C.RK.13} (p.~776)} }
\expandafter\def\csname SRef@cdp:eq:C.RG.11\endcsname{\href{S6_SERIES_GRIND_II.pdf\#nameddest=tag-00Z5}{II.{C.RG.11} (p.~448)} }
\expandafter\def\csname SRef@cdp:eq:C.RK.14\endcsname{\href{S6_SERIES_GRIND_II.pdf\#nameddest=tag-00ZV}{II.{C.RK.14} (p.~776)} }
\expandafter\def\csname SRef@cdp:eq:C.RK.15\endcsname{\href{S6_SERIES_GRIND_II.pdf\#nameddest=tag-00ZW}{II.{C.RK.15} (p.~776)} }
\expandafter\def\csname SRef@cdp:eq:C.RK.16\endcsname{\href{S6_SERIES_GRIND_II.pdf\#nameddest=tag-00ZX}{II.{C.RK.16} (p.~777)} }
\expandafter\def\csname SRef@cdp:eq:C.RK.17\endcsname{\href{S6_SERIES_GRIND_II.pdf\#nameddest=tag-00ZY}{II.{C.RK.17} (p.~777)} }
\expandafter\def\csname SRef@cdp:eq:C.RK.21\endcsname{\href{S6_SERIES_GRIND_II.pdf\#nameddest=tag-0103}{II.{C.RK.21} (p.~777)} }
\expandafter\def\csname SRef@cdp:eq:C.RK.22\endcsname{\href{S6_SERIES_GRIND_II.pdf\#nameddest=tag-0104}{II.{C.RK.22} (p.~778)} }
\expandafter\def\csname SRef@cdp:eq:C.NT.1\endcsname{\href{S6_SERIES_GRIND_II.pdf\#nameddest=tag-00TU}{II.{C.NT.1} (p.~779)} }
\expandafter\def\csname SRef@cdp:eq:C.NT.2\endcsname{\href{S6_SERIES_GRIND_II.pdf\#nameddest=tag-00TV}{II.{C.NT.2} (p.~779)} }
\expandafter\def\csname SRef@cdp:eq:C.NT.3\endcsname{\href{S6_SERIES_GRIND_II.pdf\#nameddest=tag-00TW}{II.{C.NT.3} (p.~779)} }
\expandafter\def\csname SRef@cdp:eq:C.NT.4\endcsname{\href{S6_SERIES_GRIND_II.pdf\#nameddest=tag-00TX}{II.{C.NT.4} (p.~779)} }
\expandafter\def\csname SRef@cdp:eq:C.RG.12\endcsname{\href{S6_SERIES_GRIND_II.pdf\#nameddest=tag-00Z6}{II.{C.RG.12} (p.~449)} }
\expandafter\def\csname SRef@cdp:eq:C.NT.5\endcsname{\href{S6_SERIES_GRIND_II.pdf\#nameddest=tag-00TY}{II.{C.NT.5} (p.~780)} }
\expandafter\def\csname SRef@cdp:eq:C.NT.6\endcsname{\href{S6_SERIES_GRIND_II.pdf\#nameddest=tag-00TZ}{II.{C.NT.6} (p.~780)} }
\expandafter\def\csname SRef@cdp:eq:C.NT.7\endcsname{\href{S6_SERIES_GRIND_II.pdf\#nameddest=tag-00U0}{II.{C.NT.7} (p.~780)} }
\expandafter\def\csname SRef@cdp:eq:C.NT.8\endcsname{\href{S6_SERIES_GRIND_II.pdf\#nameddest=tag-00U1}{II.{C.NT.8} (p.~781)} }
\expandafter\def\csname SRef@cdp:eq:C.GP.1\endcsname{\href{S6_SERIES_GRIND_II.pdf\#nameddest=tag-00KY}{II.{C.GP.1} (p.~782)} }
\expandafter\def\csname SRef@cdp:eq:C.GP.2\endcsname{\href{S6_SERIES_GRIND_II.pdf\#nameddest=tag-00L0}{II.{C.GP.2} (p.~783)} }
\expandafter\def\csname SRef@cdp:eq:C.GP.3\endcsname{\href{S6_SERIES_GRIND_II.pdf\#nameddest=tag-00L1}{II.{C.GP.3} (p.~783)} }
\expandafter\def\csname SRef@cdp:eq:C.GP.4\endcsname{\href{S6_SERIES_GRIND_II.pdf\#nameddest=tag-00L2}{II.{C.GP.4} (p.~783)} }
\expandafter\def\csname SRef@cdp:eq:C.GP.5\endcsname{\href{S6_SERIES_GRIND_II.pdf\#nameddest=tag-00L3}{II.{C.GP.5} (p.~783)} }
\expandafter\def\csname SRef@cdp:eq:C.GP.6\endcsname{\href{S6_SERIES_GRIND_II.pdf\#nameddest=tag-00L4}{II.{C.GP.6} (p.~784)} }
\expandafter\def\csname SRef@cdp:eq:C.RG.13\endcsname{\href{S6_SERIES_GRIND_II.pdf\#nameddest=tag-00Z7}{II.{C.RG.13} (p.~449)} }
\expandafter\def\csname SRef@cdp:eq:C.GP.7\endcsname{\href{S6_SERIES_GRIND_II.pdf\#nameddest=tag-00L5}{II.{C.GP.7} (p.~785)} }
\expandafter\def\csname SRef@cdp:eq:C.GP.8\endcsname{\href{S6_SERIES_GRIND_II.pdf\#nameddest=tag-00L6}{II.{C.GP.8} (p.~785)} }
\expandafter\def\csname SRef@cdp:eq:C.GP.9\endcsname{\href{S6_SERIES_GRIND_II.pdf\#nameddest=tag-00L7}{II.{C.GP.9} (p.~785)} }
\expandafter\def\csname SRef@cdp:eq:C.GP.10\endcsname{\href{S6_SERIES_GRIND_II.pdf\#nameddest=tag-00KZ}{II.{C.GP.10} (p.~786)} }
\expandafter\def\csname SRef@cdp:eq:C.AR.1\endcsname{\href{S6_SERIES_GRIND_II.pdf\#nameddest=tag-00A0}{II.{C.AR.1} (p.~788)} }
\expandafter\def\csname SRef@cdp:eq:C.AR.2\endcsname{\href{S6_SERIES_GRIND_II.pdf\#nameddest=tag-00A8}{II.{C.AR.2} (p.~788)} }
\expandafter\def\csname SRef@cdp:eq:C.AR.3\endcsname{\href{S6_SERIES_GRIND_II.pdf\#nameddest=tag-00A9}{II.{C.AR.3} (p.~788)} }
\expandafter\def\csname SRef@cdp:eq:C.AR.4\endcsname{\href{S6_SERIES_GRIND_II.pdf\#nameddest=tag-00AA}{II.{C.AR.4} (p.~789)} }
\expandafter\def\csname SRef@cdp:eq:C.AR.5\endcsname{\href{S6_SERIES_GRIND_II.pdf\#nameddest=tag-00AB}{II.{C.AR.5} (p.~789)} }
\expandafter\def\csname SRef@cdp:eq:C.CS.1\endcsname{\href{S6_SERIES_GRIND_II.pdf\#nameddest=tag-00DH}{II.{C.CS.1} (p.~450)} }
\expandafter\def\csname SRef@cdp:eq:C.AR.6\endcsname{\href{S6_SERIES_GRIND_II.pdf\#nameddest=tag-00AC}{II.{C.AR.6} (p.~789)} }
\expandafter\def\csname SRef@cdp:eq:C.AR.7\endcsname{\href{S6_SERIES_GRIND_II.pdf\#nameddest=tag-00AD}{II.{C.AR.7} (p.~789)} }
\expandafter\def\csname SRef@cdp:eq:C.AR.8\endcsname{\href{S6_SERIES_GRIND_II.pdf\#nameddest=tag-00AE}{II.{C.AR.8} (p.~789)} }
\expandafter\def\csname SRef@cdp:eq:C.AR.9\endcsname{\href{S6_SERIES_GRIND_II.pdf\#nameddest=tag-00AF}{II.{C.AR.9} (p.~789)} }
\expandafter\def\csname SRef@cdp:eq:C.AR.10\endcsname{\href{S6_SERIES_GRIND_II.pdf\#nameddest=tag-00A1}{II.{C.AR.10} (p.~789)} }
\expandafter\def\csname SRef@cdp:eq:C.AR.11\endcsname{\href{S6_SERIES_GRIND_II.pdf\#nameddest=tag-00A2}{II.{C.AR.11} (p.~789)} }
\expandafter\def\csname SRef@cdp:eq:C.AR.12\endcsname{\href{S6_SERIES_GRIND_II.pdf\#nameddest=tag-00A3}{II.{C.AR.12} (p.~790)} }
\expandafter\def\csname SRef@cdp:eq:C.AR.13\endcsname{\href{S6_SERIES_GRIND_II.pdf\#nameddest=tag-00A4}{II.{C.AR.13} (p.~790)} }
\expandafter\def\csname SRef@cdp:eq:C.AR.14\endcsname{\href{S6_SERIES_GRIND_II.pdf\#nameddest=tag-00A5}{II.{C.AR.14} (p.~791)} }
\expandafter\def\csname SRef@cdp:eq:C.AR.15\endcsname{\href{S6_SERIES_GRIND_II.pdf\#nameddest=tag-00A6}{II.{C.AR.15} (p.~791)} }
\expandafter\def\csname SRef@cdp:eq:C.CS.2\endcsname{\href{S6_SERIES_GRIND_II.pdf\#nameddest=tag-00DI}{II.{C.CS.2} (p.~451)} }
\expandafter\def\csname SRef@cdp:eq:C.AR.16\endcsname{\href{S6_SERIES_GRIND_II.pdf\#nameddest=tag-00A7}{II.{C.AR.16} (p.~791)} }
\expandafter\def\csname SRef@cdp:eq:C.G.1\endcsname{\href{S6_SERIES_GRIND_II.pdf\#nameddest=tag-00KW}{II.{C.G.1} (p.~793)} }
\expandafter\def\csname SRef@cdp:eq:C.G.2\endcsname{\href{S6_SERIES_GRIND_II.pdf\#nameddest=tag-00KX}{II.{C.G.2} (p.~794)} }
\expandafter\def\csname SRef@cdp:eq:C.CS.3\endcsname{\href{S6_SERIES_GRIND_II.pdf\#nameddest=tag-00DJ}{II.{C.CS.3} (p.~451)} }
\expandafter\def\csname SRef@cdp:eq:C.UN.4\endcsname{\href{S6_SERIES_GRIND_II.pdf\#nameddest=tag-0155}{II.{C.UN.4} (p.~414)} }
\expandafter\def\csname SRef@cdp:eq:C.CS.4\endcsname{\href{S6_SERIES_GRIND_II.pdf\#nameddest=tag-00DK}{II.{C.CS.4} (p.~451)} }
\expandafter\def\csname SRef@cdp:eq:C.CS.5\endcsname{\href{S6_SERIES_GRIND_II.pdf\#nameddest=tag-00DL}{II.{C.CS.5} (p.~452)} }
\expandafter\def\csname SRef@cdp:eq:C.CS.6\endcsname{\href{S6_SERIES_GRIND_II.pdf\#nameddest=tag-00DM}{II.{C.CS.6} (p.~452)} }
\expandafter\def\csname SRef@cdp:eq:C.EQ.1\endcsname{\href{S6_SERIES_GRIND_II.pdf\#nameddest=tag-00IC}{II.{C.EQ.1} (p.~453)} }
\expandafter\def\csname SRef@cdp:eq:C.EQ.2\endcsname{\href{S6_SERIES_GRIND_II.pdf\#nameddest=tag-00IF}{II.{C.EQ.2} (p.~453)} }
\expandafter\def\csname SRef@cdp:eq:C.EQ.3\endcsname{\href{S6_SERIES_GRIND_II.pdf\#nameddest=tag-00IG}{II.{C.EQ.3} (p.~454)} }
\expandafter\def\csname SRef@cdp:eq:C.EQ.4\endcsname{\href{S6_SERIES_GRIND_II.pdf\#nameddest=tag-00IH}{II.{C.EQ.4} (p.~454)} }
\expandafter\def\csname SRef@cdp:eq:C.EQ.5\endcsname{\href{S6_SERIES_GRIND_II.pdf\#nameddest=tag-00II}{II.{C.EQ.5} (p.~454)} }
\expandafter\def\csname SRef@cdp:eq:C.EQ.6\endcsname{\href{S6_SERIES_GRIND_II.pdf\#nameddest=tag-00IJ}{II.{C.EQ.6} (p.~455)} }
\expandafter\def\csname SRef@cdp:eq:C.EQ.7\endcsname{\href{S6_SERIES_GRIND_II.pdf\#nameddest=tag-00IK}{II.{C.EQ.7} (p.~455)} }
\expandafter\def\csname SRef@cdp:eq:C.UN.5\endcsname{\href{S6_SERIES_GRIND_II.pdf\#nameddest=tag-0156}{II.{C.UN.5} (p.~415)} }
\expandafter\def\csname SRef@cdp:eq:C.EQ.8\endcsname{\href{S6_SERIES_GRIND_II.pdf\#nameddest=tag-00IL}{II.{C.EQ.8} (p.~455)} }
\expandafter\def\csname SRef@cdp:eq:C.EQ.9\endcsname{\href{S6_SERIES_GRIND_II.pdf\#nameddest=tag-00IM}{II.{C.EQ.9} (p.~456)} }
\expandafter\def\csname SRef@cdp:eq:C.EQ.10\endcsname{\href{S6_SERIES_GRIND_II.pdf\#nameddest=tag-00ID}{II.{C.EQ.10} (p.~456)} }
\expandafter\def\csname SRef@cdp:eq:C.EQ.11\endcsname{\href{S6_SERIES_GRIND_II.pdf\#nameddest=tag-00IE}{II.{C.EQ.11} (p.~456)} }
\expandafter\def\csname SRef@cdp:eq:C.CL.1\endcsname{\href{S6_SERIES_GRIND_II.pdf\#nameddest=tag-00BW}{II.{C.CL.1} (p.~457)} }
\expandafter\def\csname SRef@cdp:eq:C.CL.2\endcsname{\href{S6_SERIES_GRIND_II.pdf\#nameddest=tag-00C1}{II.{C.CL.2} (p.~457)} }
\expandafter\def\csname SRef@cdp:eq:C.CL.3\endcsname{\href{S6_SERIES_GRIND_II.pdf\#nameddest=tag-00C2}{II.{C.CL.3} (p.~457)} }
\expandafter\def\csname SRef@cdp:eq:C.CL.4\endcsname{\href{S6_SERIES_GRIND_II.pdf\#nameddest=tag-00C3}{II.{C.CL.4} (p.~458)} }
\expandafter\def\csname SRef@cdp:eq:C.CL.5\endcsname{\href{S6_SERIES_GRIND_II.pdf\#nameddest=tag-00C4}{II.{C.CL.5} (p.~458)} }
\expandafter\def\csname SRef@cdp:eq:C.CL.6\endcsname{\href{S6_SERIES_GRIND_II.pdf\#nameddest=tag-00C5}{II.{C.CL.6} (p.~458)} }
\expandafter\def\csname SRef@cdp:eq:C.UN.6\endcsname{\href{S6_SERIES_GRIND_II.pdf\#nameddest=tag-0157}{II.{C.UN.6} (p.~415)} }
\expandafter\def\csname SRef@cdp:eq:C.CL.7\endcsname{\href{S6_SERIES_GRIND_II.pdf\#nameddest=tag-00C6}{II.{C.CL.7} (p.~458)} }
\expandafter\def\csname SRef@cdp:eq:C.CL.8\endcsname{\href{S6_SERIES_GRIND_II.pdf\#nameddest=tag-00C7}{II.{C.CL.8} (p.~458)} }
\expandafter\def\csname SRef@cdp:eq:C.CL.9\endcsname{\href{S6_SERIES_GRIND_II.pdf\#nameddest=tag-00C8}{II.{C.CL.9} (p.~459)} }
\expandafter\def\csname SRef@cdp:eq:C.CL.10\endcsname{\href{S6_SERIES_GRIND_II.pdf\#nameddest=tag-00BX}{II.{C.CL.10} (p.~460)} }
\expandafter\def\csname SRef@cdp:eq:C.CL.11\endcsname{\href{S6_SERIES_GRIND_II.pdf\#nameddest=tag-00BY}{II.{C.CL.11} (p.~460)} }
\expandafter\def\csname SRef@cdp:eq:C.CL.12\endcsname{\href{S6_SERIES_GRIND_II.pdf\#nameddest=tag-00BZ}{II.{C.CL.12} (p.~460)} }
\expandafter\def\csname SRef@cdp:eq:C.CL.13\endcsname{\href{S6_SERIES_GRIND_II.pdf\#nameddest=tag-00C0}{II.{C.CL.13} (p.~460)} }
\expandafter\def\csname SRef@cdp:eq:C.HB.1\endcsname{\href{S6_SERIES_GRIND_II.pdf\#nameddest=tag-00LF}{II.{C.HB.1} (p.~461)} }
\expandafter\def\csname SRef@cdp:eq:C.HB.2\endcsname{\href{S6_SERIES_GRIND_II.pdf\#nameddest=tag-00LG}{II.{C.HB.2} (p.~461)} }
\expandafter\def\csname SRef@cdp:eq:C.HB.3\endcsname{\href{S6_SERIES_GRIND_II.pdf\#nameddest=tag-00LH}{II.{C.HB.3} (p.~461)} }
\expandafter\def\csname SRef@cdp:eq:C.UC.1\endcsname{\href{S6_SERIES_GRIND_II.pdf\#nameddest=tag-014C}{II.{C.UC.1} (p.~416)} }
\expandafter\def\csname SRef@cdp:eq:C.HB.4\endcsname{\href{S6_SERIES_GRIND_II.pdf\#nameddest=tag-00LI}{II.{C.HB.4} (p.~461)} }
\expandafter\def\csname SRef@cdp:eq:C.HB.5\endcsname{\href{S6_SERIES_GRIND_II.pdf\#nameddest=tag-00LJ}{II.{C.HB.5} (p.~461)} }
\expandafter\def\csname SRef@cdp:eq:C.HB.6\endcsname{\href{S6_SERIES_GRIND_II.pdf\#nameddest=tag-00LK}{II.{C.HB.6} (p.~461)} }
\expandafter\def\csname SRef@cdp:eq:C.HB.7\endcsname{\href{S6_SERIES_GRIND_II.pdf\#nameddest=tag-00LL}{II.{C.HB.7} (p.~462)} }
\expandafter\def\csname SRef@cdp:eq:C.HB.8\endcsname{\href{S6_SERIES_GRIND_II.pdf\#nameddest=tag-00LM}{II.{C.HB.8} (p.~462)} }
\expandafter\def\csname SRef@cdp:eq:C.HB.9\endcsname{\href{S6_SERIES_GRIND_II.pdf\#nameddest=tag-00LN}{II.{C.HB.9} (p.~462)} }
\expandafter\def\csname SRef@cdp:eq:C.ZC.1\endcsname{\href{S6_SERIES_GRIND_II.pdf\#nameddest=tag-018W}{II.{C.ZC.1} (p.~463)} }
\expandafter\def\csname SRef@cdp:eq:C.ZC.2\endcsname{\href{S6_SERIES_GRIND_II.pdf\#nameddest=tag-018X}{II.{C.ZC.2} (p.~463)} }
\expandafter\def\csname SRef@cdp:eq:C.ZC.3\endcsname{\href{S6_SERIES_GRIND_II.pdf\#nameddest=tag-018Y}{II.{C.ZC.3} (p.~463)} }
\expandafter\def\csname SRef@cdp:eq:C.ZC.4\endcsname{\href{S6_SERIES_GRIND_II.pdf\#nameddest=tag-018Z}{II.{C.ZC.4} (p.~464)} }
\expandafter\def\csname SRef@cdp:eq:C.UC.2\endcsname{\href{S6_SERIES_GRIND_II.pdf\#nameddest=tag-014E}{II.{C.UC.2} (p.~416)} }
\expandafter\def\csname SRef@cdp:eq:C.ZC.5\endcsname{\href{S6_SERIES_GRIND_II.pdf\#nameddest=tag-0190}{II.{C.ZC.5} (p.~464)} }
\expandafter\def\csname SRef@cdp:eq:C.ZC.6\endcsname{\href{S6_SERIES_GRIND_II.pdf\#nameddest=tag-0191}{II.{C.ZC.6} (p.~465)} }
\expandafter\def\csname SRef@cdp:eq:C.ZC.7\endcsname{\href{S6_SERIES_GRIND_II.pdf\#nameddest=tag-0192}{II.{C.ZC.7} (p.~465)} }
\expandafter\def\csname SRef@cdp:eq:C.ZC.8\endcsname{\href{S6_SERIES_GRIND_II.pdf\#nameddest=tag-0193}{II.{C.ZC.8} (p.~465)} }
\expandafter\def\csname SRef@cdp:eq:C.ZC.9\endcsname{\href{S6_SERIES_GRIND_II.pdf\#nameddest=tag-0194}{II.{C.ZC.9} (p.~465)} }
\expandafter\def\csname SRef@cdp:eq:C.FG.1\endcsname{\href{S6_SERIES_GRIND_II.pdf\#nameddest=tag-00K8}{II.{C.FG.1} (p.~466)} }
\expandafter\def\csname SRef@cdp:eq:C.FG.2\endcsname{\href{S6_SERIES_GRIND_II.pdf\#nameddest=tag-00KJ}{II.{C.FG.2} (p.~466)} }
\expandafter\def\csname SRef@cdp:eq:C.FG.3\endcsname{\href{S6_SERIES_GRIND_II.pdf\#nameddest=tag-00KK}{II.{C.FG.3} (p.~466)} }
\expandafter\def\csname SRef@cdp:eq:C.FG.4\endcsname{\href{S6_SERIES_GRIND_II.pdf\#nameddest=tag-00KL}{II.{C.FG.4} (p.~466)} }
\expandafter\def\csname SRef@cdp:eq:C.FG.5\endcsname{\href{S6_SERIES_GRIND_II.pdf\#nameddest=tag-00KM}{II.{C.FG.5} (p.~466)} }
\expandafter\def\csname SRef@cdp:eq:C.UC.3\endcsname{\href{S6_SERIES_GRIND_II.pdf\#nameddest=tag-014F}{II.{C.UC.3} (p.~416)} }
\expandafter\def\csname SRef@cdp:eq:C.FG.6\endcsname{\href{S6_SERIES_GRIND_II.pdf\#nameddest=tag-00KN}{II.{C.FG.6} (p.~467)} }
\expandafter\def\csname SRef@cdp:eq:C.FG.7\endcsname{\href{S6_SERIES_GRIND_II.pdf\#nameddest=tag-00KO}{II.{C.FG.7} (p.~467)} }
\expandafter\def\csname SRef@cdp:eq:C.FG.8\endcsname{\href{S6_SERIES_GRIND_II.pdf\#nameddest=tag-00KP}{II.{C.FG.8} (p.~467)} }
\expandafter\def\csname SRef@cdp:eq:C.FG.9\endcsname{\href{S6_SERIES_GRIND_II.pdf\#nameddest=tag-00KQ}{II.{C.FG.9} (p.~467)} }
\expandafter\def\csname SRef@cdp:eq:C.FG.10\endcsname{\href{S6_SERIES_GRIND_II.pdf\#nameddest=tag-00K9}{II.{C.FG.10} (p.~468)} }
\expandafter\def\csname SRef@cdp:eq:C.FG.11\endcsname{\href{S6_SERIES_GRIND_II.pdf\#nameddest=tag-00KA}{II.{C.FG.11} (p.~468)} }
\expandafter\def\csname SRef@cdp:eq:C.FG.12\endcsname{\href{S6_SERIES_GRIND_II.pdf\#nameddest=tag-00KB}{II.{C.FG.12} (p.~468)} }
\expandafter\def\csname SRef@cdp:eq:C.FG.13\endcsname{\href{S6_SERIES_GRIND_II.pdf\#nameddest=tag-00KC}{II.{C.FG.13} (p.~469)} }
\expandafter\def\csname SRef@cdp:eq:C.FG.14\endcsname{\href{S6_SERIES_GRIND_II.pdf\#nameddest=tag-00KD}{II.{C.FG.14} (p.~469)} }
\expandafter\def\csname SRef@cdp:eq:C.FG.15\endcsname{\href{S6_SERIES_GRIND_II.pdf\#nameddest=tag-00KE}{II.{C.FG.15} (p.~469)} }
\expandafter\def\csname SRef@cdp:eq:C.UC.4\endcsname{\href{S6_SERIES_GRIND_II.pdf\#nameddest=tag-014G}{II.{C.UC.4} (p.~416)} }
\expandafter\def\csname SRef@cdp:eq:C.FG.16\endcsname{\href{S6_SERIES_GRIND_II.pdf\#nameddest=tag-00KF}{II.{C.FG.16} (p.~469)} }
\expandafter\def\csname SRef@cdp:eq:C.FG.17\endcsname{\href{S6_SERIES_GRIND_II.pdf\#nameddest=tag-00KG}{II.{C.FG.17} (p.~470)} }
\expandafter\def\csname SRef@cdp:eq:C.FG.18\endcsname{\href{S6_SERIES_GRIND_II.pdf\#nameddest=tag-00KH}{II.{C.FG.18} (p.~470)} }
\expandafter\def\csname SRef@cdp:eq:C.FG.19\endcsname{\href{S6_SERIES_GRIND_II.pdf\#nameddest=tag-00KI}{II.{C.FG.19} (p.~470)} }
\expandafter\def\csname SRef@cdp:eq:C.PF.1\endcsname{\href{S6_SERIES_GRIND_II.pdf\#nameddest=tag-00UN}{II.{C.PF.1} (p.~471)} }
\expandafter\def\csname SRef@cdp:eq:C.PF.2\endcsname{\href{S6_SERIES_GRIND_II.pdf\#nameddest=tag-00UR}{II.{C.PF.2} (p.~471)} }
\expandafter\def\csname SRef@cdp:eq:C.PF.3\endcsname{\href{S6_SERIES_GRIND_II.pdf\#nameddest=tag-00US}{II.{C.PF.3} (p.~471)} }
\expandafter\def\csname SRef@cdp:eq:C.PF.4\endcsname{\href{S6_SERIES_GRIND_II.pdf\#nameddest=tag-00UT}{II.{C.PF.4} (p.~471)} }
\expandafter\def\csname SRef@cdp:eq:C.PF.5\endcsname{\href{S6_SERIES_GRIND_II.pdf\#nameddest=tag-00UU}{II.{C.PF.5} (p.~472)} }
\expandafter\def\csname SRef@cdp:eq:C.PF.6\endcsname{\href{S6_SERIES_GRIND_II.pdf\#nameddest=tag-00UV}{II.{C.PF.6} (p.~472)} }
\expandafter\def\csname SRef@cdp:eq:C.UC.5\endcsname{\href{S6_SERIES_GRIND_II.pdf\#nameddest=tag-014H}{II.{C.UC.5} (p.~417)} }
\expandafter\def\csname SRef@cdp:eq:C.PF.7\endcsname{\href{S6_SERIES_GRIND_II.pdf\#nameddest=tag-00UW}{II.{C.PF.7} (p.~472)} }
\expandafter\def\csname SRef@cdp:eq:C.PF.8\endcsname{\href{S6_SERIES_GRIND_II.pdf\#nameddest=tag-00UX}{II.{C.PF.8} (p.~472)} }
\expandafter\def\csname SRef@cdp:eq:C.PF.9\endcsname{\href{S6_SERIES_GRIND_II.pdf\#nameddest=tag-00UY}{II.{C.PF.9} (p.~473)} }
\expandafter\def\csname SRef@cdp:eq:C.PF.10\endcsname{\href{S6_SERIES_GRIND_II.pdf\#nameddest=tag-00UO}{II.{C.PF.10} (p.~473)} }
\expandafter\def\csname SRef@cdp:eq:C.PF.11\endcsname{\href{S6_SERIES_GRIND_II.pdf\#nameddest=tag-00UP}{II.{C.PF.11} (p.~473)} }
\expandafter\def\csname SRef@cdp:eq:C.PF.12\endcsname{\href{S6_SERIES_GRIND_II.pdf\#nameddest=tag-00UQ}{II.{C.PF.12} (p.~473)} }
\expandafter\def\csname SRef@cdp:eq:C.UNN.1\endcsname{\href{S6_SERIES_GRIND_II.pdf\#nameddest=tag-015B}{II.{C.UNN.1} (p.~474)} }
\expandafter\def\csname SRef@cdp:eq:C.UNN.2\endcsname{\href{S6_SERIES_GRIND_II.pdf\#nameddest=tag-015C}{II.{C.UNN.2} (p.~475)} }
\expandafter\def\csname SRef@cdp:eq:C.CR.1\endcsname{\href{S6_SERIES_GRIND_II.pdf\#nameddest=tag-00CV}{II.{C.CR.1} (p.~476)} }
\expandafter\def\csname SRef@cdp:eq:C.CR.2\endcsname{\href{S6_SERIES_GRIND_II.pdf\#nameddest=tag-00D6}{II.{C.CR.2} (p.~476)} }
\expandafter\def\csname SRef@cdp:eq:C.UC.6\endcsname{\href{S6_SERIES_GRIND_II.pdf\#nameddest=tag-014I}{II.{C.UC.6} (p.~417)} }
\expandafter\def\csname SRef@cdp:eq:C.CR.3\endcsname{\href{S6_SERIES_GRIND_II.pdf\#nameddest=tag-00DA}{II.{C.CR.3} (p.~476)} }
\expandafter\def\csname SRef@cdp:eq:C.CR.4\endcsname{\href{S6_SERIES_GRIND_II.pdf\#nameddest=tag-00DB}{II.{C.CR.4} (p.~476)} }
\expandafter\def\csname SRef@cdp:eq:C.CR.5\endcsname{\href{S6_SERIES_GRIND_II.pdf\#nameddest=tag-00DC}{II.{C.CR.5} (p.~477)} }
\expandafter\def\csname SRef@cdp:eq:C.CR.6\endcsname{\href{S6_SERIES_GRIND_II.pdf\#nameddest=tag-00DD}{II.{C.CR.6} (p.~477)} }
\expandafter\def\csname SRef@cdp:eq:C.CR.7\endcsname{\href{S6_SERIES_GRIND_II.pdf\#nameddest=tag-00DE}{II.{C.CR.7} (p.~477)} }
\expandafter\def\csname SRef@cdp:eq:C.CR.8\endcsname{\href{S6_SERIES_GRIND_II.pdf\#nameddest=tag-00DF}{II.{C.CR.8} (p.~478)} }
\expandafter\def\csname SRef@cdp:eq:C.CR.9\endcsname{\href{S6_SERIES_GRIND_II.pdf\#nameddest=tag-00DG}{II.{C.CR.9} (p.~478)} }
\expandafter\def\csname SRef@cdp:eq:C.CR.10\endcsname{\href{S6_SERIES_GRIND_II.pdf\#nameddest=tag-00CW}{II.{C.CR.10} (p.~479)} }
\expandafter\def\csname SRef@cdp:eq:C.CR.11\endcsname{\href{S6_SERIES_GRIND_II.pdf\#nameddest=tag-00CX}{II.{C.CR.11} (p.~479)} }
\expandafter\def\csname SRef@cdp:eq:C.CR.12\endcsname{\href{S6_SERIES_GRIND_II.pdf\#nameddest=tag-00CY}{II.{C.CR.12} (p.~479)} }
\expandafter\def\csname SRef@cdp:eq:C.UC.7\endcsname{\href{S6_SERIES_GRIND_II.pdf\#nameddest=tag-014J}{II.{C.UC.7} (p.~417)} }
\expandafter\def\csname SRef@cdp:eq:C.CR.13\endcsname{\href{S6_SERIES_GRIND_II.pdf\#nameddest=tag-00CZ}{II.{C.CR.13} (p.~479)} }
\expandafter\def\csname SRef@cdp:eq:C.CR.14\endcsname{\href{S6_SERIES_GRIND_II.pdf\#nameddest=tag-00D0}{II.{C.CR.14} (p.~479)} }
\expandafter\def\csname SRef@cdp:eq:C.CR.15\endcsname{\href{S6_SERIES_GRIND_II.pdf\#nameddest=tag-00D1}{II.{C.CR.15} (p.~481)} }
\expandafter\def\csname SRef@cdp:eq:C.CR.16\endcsname{\href{S6_SERIES_GRIND_II.pdf\#nameddest=tag-00D2}{II.{C.CR.16} (p.~481)} }
\expandafter\def\csname SRef@cdp:eq:C.CR.17\endcsname{\href{S6_SERIES_GRIND_II.pdf\#nameddest=tag-00D3}{II.{C.CR.17} (p.~481)} }
\expandafter\def\csname SRef@cdp:eq:C.CR.18\endcsname{\href{S6_SERIES_GRIND_II.pdf\#nameddest=tag-00D4}{II.{C.CR.18} (p.~482)} }
\expandafter\def\csname SRef@cdp:eq:C.CR.19\endcsname{\href{S6_SERIES_GRIND_II.pdf\#nameddest=tag-00D5}{II.{C.CR.19} (p.~483)} }
\expandafter\def\csname SRef@cdp:eq:C.CR.20\endcsname{\href{S6_SERIES_GRIND_II.pdf\#nameddest=tag-00D7}{II.{C.CR.20} (p.~483)} }
\expandafter\def\csname SRef@cdp:eq:C.CR.21\endcsname{\href{S6_SERIES_GRIND_II.pdf\#nameddest=tag-00D8}{II.{C.CR.21} (p.~483)} }
\expandafter\def\csname SRef@cdp:eq:C.CR.22\endcsname{\href{S6_SERIES_GRIND_II.pdf\#nameddest=tag-00D9}{II.{C.CR.22} (p.~484)} }
\expandafter\def\csname SRef@cdp:eq:C.UC.8\endcsname{\href{S6_SERIES_GRIND_II.pdf\#nameddest=tag-014K}{II.{C.UC.8} (p.~417)} }
\expandafter\def\csname SRef@cdp:eq:C.UNNODE.1\endcsname{\href{S6_SERIES_GRIND_II.pdf\#nameddest=tag-015D}{II.{C.UNNODE.1} (p.~484)} }
\expandafter\def\csname SRef@cdp:eq:C.UNNODE.2\endcsname{\href{S6_SERIES_GRIND_II.pdf\#nameddest=tag-015E}{II.{C.UNNODE.2} (p.~485)} }
\expandafter\def\csname SRef@cdp:eq:C.UNNODE.3\endcsname{\href{S6_SERIES_GRIND_II.pdf\#nameddest=tag-015F}{II.{C.UNNODE.3} (p.~485)} }
\expandafter\def\csname SRef@cdp:eq:C.UNCR.1\endcsname{\href{S6_SERIES_GRIND_II.pdf\#nameddest=tag-0158}{II.{C.UNCR.1} (p.~486)} }
\expandafter\def\csname SRef@cdp:eq:C.UNCR.2\endcsname{\href{S6_SERIES_GRIND_II.pdf\#nameddest=tag-0159}{II.{C.UNCR.2} (p.~486)} }
\expandafter\def\csname SRef@cdp:eq:C.UNCR.3\endcsname{\href{S6_SERIES_GRIND_II.pdf\#nameddest=tag-015A}{II.{C.UNCR.3} (p.~486)} }
\expandafter\def\csname SRef@cdp:eq:C.IR.1\endcsname{\href{S6_SERIES_GRIND_II.pdf\#nameddest=tag-00OH}{II.{C.IR.1} (p.~487)} }
\expandafter\def\csname SRef@cdp:eq:C.IR.2\endcsname{\href{S6_SERIES_GRIND_II.pdf\#nameddest=tag-00OI}{II.{C.IR.2} (p.~487)} }
\expandafter\def\csname SRef@cdp:eq:C.IR.4\endcsname{\href{S6_SERIES_GRIND_II.pdf\#nameddest=tag-00OJ}{II.{C.IR.4} (p.~487)} }
\expandafter\def\csname SRef@cdp:eq:C.IR.5\endcsname{\href{S6_SERIES_GRIND_II.pdf\#nameddest=tag-00OK}{II.{C.IR.5} (p.~488)} }
\expandafter\def\csname SRef@cdp:eq:C.UC.9\endcsname{\href{S6_SERIES_GRIND_II.pdf\#nameddest=tag-014L}{II.{C.UC.9} (p.~418)} }
\expandafter\def\csname SRef@cdp:eq:C.IR.6\endcsname{\href{S6_SERIES_GRIND_II.pdf\#nameddest=tag-00OL}{II.{C.IR.6} (p.~489)} }
\expandafter\def\csname SRef@cdp:eq:C.IR.9\endcsname{\href{S6_SERIES_GRIND_II.pdf\#nameddest=tag-00OM}{II.{C.IR.9} (p.~489)} }
\expandafter\def\csname SRef@cdp:eq:C.VII.1\endcsname{\href{S6_SERIES_GRIND_II.pdf\#nameddest=tag-015K}{II.{C.VII.1} (p.~489)} }
\expandafter\def\csname SRef@cdp:eq:C.VII.2\endcsname{\href{S6_SERIES_GRIND_II.pdf\#nameddest=tag-015O}{II.{C.VII.2} (p.~490)} }
\expandafter\def\csname SRef@cdp:eq:C.VII.3\endcsname{\href{S6_SERIES_GRIND_II.pdf\#nameddest=tag-015P}{II.{C.VII.3} (p.~491)} }
\expandafter\def\csname SRef@cdp:eq:C.VII.4\endcsname{\href{S6_SERIES_GRIND_II.pdf\#nameddest=tag-015Q}{II.{C.VII.4} (p.~491)} }
\expandafter\def\csname SRef@cdp:eq:C.VII.5\endcsname{\href{S6_SERIES_GRIND_II.pdf\#nameddest=tag-015R}{II.{C.VII.5} (p.~491)} }
\expandafter\def\csname SRef@cdp:eq:C.VII.6\endcsname{\href{S6_SERIES_GRIND_II.pdf\#nameddest=tag-015S}{II.{C.VII.6} (p.~491)} }
\expandafter\def\csname SRef@cdp:eq:C.VII.7\endcsname{\href{S6_SERIES_GRIND_II.pdf\#nameddest=tag-015T}{II.{C.VII.7} (p.~492)} }
\expandafter\def\csname SRef@cdp:eq:C.VII.8\endcsname{\href{S6_SERIES_GRIND_II.pdf\#nameddest=tag-015U}{II.{C.VII.8} (p.~493)} }
\expandafter\def\csname SRef@cdp:eq:C.UC.10\endcsname{\href{S6_SERIES_GRIND_II.pdf\#nameddest=tag-014D}{II.{C.UC.10} (p.~418)} }
\expandafter\def\csname SRef@cdp:eq:C.VII.9\endcsname{\href{S6_SERIES_GRIND_II.pdf\#nameddest=tag-015V}{II.{C.VII.9} (p.~493)} }
\expandafter\def\csname SRef@cdp:eq:C.VII.10\endcsname{\href{S6_SERIES_GRIND_II.pdf\#nameddest=tag-015L}{II.{C.VII.10} (p.~493)} }
\expandafter\def\csname SRef@cdp:eq:C.VII.11\endcsname{\href{S6_SERIES_GRIND_II.pdf\#nameddest=tag-015M}{II.{C.VII.11} (p.~494)} }
\expandafter\def\csname SRef@cdp:eq:C.VII.12\endcsname{\href{S6_SERIES_GRIND_II.pdf\#nameddest=tag-015N}{II.{C.VII.12} (p.~494)} }
\expandafter\def\csname SRef@cdp:eq:C.RA.1\endcsname{\href{S6_SERIES_GRIND_II.pdf\#nameddest=tag-00WQ}{II.{C.RA.1} (p.~495)} }
\expandafter\def\csname SRef@cdp:eq:C.RA.2\endcsname{\href{S6_SERIES_GRIND_II.pdf\#nameddest=tag-00X1}{II.{C.RA.2} (p.~495)} }
\expandafter\def\csname SRef@cdp:eq:C.RA.3\endcsname{\href{S6_SERIES_GRIND_II.pdf\#nameddest=tag-00XC}{II.{C.RA.3} (p.~495)} }
\expandafter\def\csname SRef@cdp:eq:C.RA.4\endcsname{\href{S6_SERIES_GRIND_II.pdf\#nameddest=tag-00XH}{II.{C.RA.4} (p.~495)} }
\expandafter\def\csname SRef@cdp:eq:C.RA.5\endcsname{\href{S6_SERIES_GRIND_II.pdf\#nameddest=tag-00XI}{II.{C.RA.5} (p.~495)} }
\expandafter\def\csname SRef@cdp:eq:C.RA.6\endcsname{\href{S6_SERIES_GRIND_II.pdf\#nameddest=tag-00XJ}{II.{C.RA.6} (p.~495)} }
\expandafter\def\csname SRef@cdp:eq:C.VR.1\endcsname{\href{S6_SERIES_GRIND_II.pdf\#nameddest=tag-015W}{II.{C.VR.1} (p.~418)} }
\expandafter\def\csname SRef@cdp:eq:C.RA.7\endcsname{\href{S6_SERIES_GRIND_II.pdf\#nameddest=tag-00XK}{II.{C.RA.7} (p.~496)} }
\expandafter\def\csname SRef@cdp:eq:C.RA.8\endcsname{\href{S6_SERIES_GRIND_II.pdf\#nameddest=tag-00XL}{II.{C.RA.8} (p.~496)} }
\expandafter\def\csname SRef@cdp:eq:C.RA.9\endcsname{\href{S6_SERIES_GRIND_II.pdf\#nameddest=tag-00XM}{II.{C.RA.9} (p.~496)} }
\expandafter\def\csname SRef@cdp:eq:C.RA.10\endcsname{\href{S6_SERIES_GRIND_II.pdf\#nameddest=tag-00WR}{II.{C.RA.10} (p.~496)} }
\expandafter\def\csname SRef@cdp:eq:C.RA.11\endcsname{\href{S6_SERIES_GRIND_II.pdf\#nameddest=tag-00WS}{II.{C.RA.11} (p.~497)} }
\expandafter\def\csname SRef@cdp:eq:C.RA.12\endcsname{\href{S6_SERIES_GRIND_II.pdf\#nameddest=tag-00WT}{II.{C.RA.12} (p.~497)} }
\expandafter\def\csname SRef@cdp:eq:C.RA.13\endcsname{\href{S6_SERIES_GRIND_II.pdf\#nameddest=tag-00WU}{II.{C.RA.13} (p.~497)} }
\expandafter\def\csname SRef@cdp:eq:C.RA.14\endcsname{\href{S6_SERIES_GRIND_II.pdf\#nameddest=tag-00WV}{II.{C.RA.14} (p.~498)} }
\expandafter\def\csname SRef@cdp:eq:C.RA.15\endcsname{\href{S6_SERIES_GRIND_II.pdf\#nameddest=tag-00WW}{II.{C.RA.15} (p.~498)} }
\expandafter\def\csname SRef@cdp:eq:C.RA.16\endcsname{\href{S6_SERIES_GRIND_II.pdf\#nameddest=tag-00WX}{II.{C.RA.16} (p.~498)} }
\expandafter\def\csname SRef@cdp:eq:C.VR.2\endcsname{\href{S6_SERIES_GRIND_II.pdf\#nameddest=tag-015X}{II.{C.VR.2} (p.~419)} }
\expandafter\def\csname SRef@cdp:eq:C.RA.17\endcsname{\href{S6_SERIES_GRIND_II.pdf\#nameddest=tag-00WY}{II.{C.RA.17} (p.~498)} }
\expandafter\def\csname SRef@cdp:eq:C.RA.18\endcsname{\href{S6_SERIES_GRIND_II.pdf\#nameddest=tag-00WZ}{II.{C.RA.18} (p.~498)} }
\expandafter\def\csname SRef@cdp:eq:C.RA.19\endcsname{\href{S6_SERIES_GRIND_II.pdf\#nameddest=tag-00X0}{II.{C.RA.19} (p.~499)} }
\expandafter\def\csname SRef@cdp:eq:C.RA.20\endcsname{\href{S6_SERIES_GRIND_II.pdf\#nameddest=tag-00X2}{II.{C.RA.20} (p.~499)} }
\expandafter\def\csname SRef@cdp:eq:C.RA.21\endcsname{\href{S6_SERIES_GRIND_II.pdf\#nameddest=tag-00X3}{II.{C.RA.21} (p.~499)} }
\expandafter\def\csname SRef@cdp:eq:C.RA.22\endcsname{\href{S6_SERIES_GRIND_II.pdf\#nameddest=tag-00X4}{II.{C.RA.22} (p.~500)} }
\expandafter\def\csname SRef@cdp:eq:C.RA.23\endcsname{\href{S6_SERIES_GRIND_II.pdf\#nameddest=tag-00X5}{II.{C.RA.23} (p.~500)} }
\expandafter\def\csname SRef@cdp:eq:C.RA.24\endcsname{\href{S6_SERIES_GRIND_II.pdf\#nameddest=tag-00X6}{II.{C.RA.24} (p.~501)} }
\expandafter\def\csname SRef@cdp:eq:C.RA.25\endcsname{\href{S6_SERIES_GRIND_II.pdf\#nameddest=tag-00X7}{II.{C.RA.25} (p.~501)} }
\expandafter\def\csname SRef@cdp:eq:C.RA.26\endcsname{\href{S6_SERIES_GRIND_II.pdf\#nameddest=tag-00X8}{II.{C.RA.26} (p.~501)} }
\expandafter\def\csname SRef@cdp:eq:C.VR.3\endcsname{\href{S6_SERIES_GRIND_II.pdf\#nameddest=tag-015Y}{II.{C.VR.3} (p.~419)} }
\expandafter\def\csname SRef@cdp:eq:C.RA.27\endcsname{\href{S6_SERIES_GRIND_II.pdf\#nameddest=tag-00X9}{II.{C.RA.27} (p.~502)} }
\expandafter\def\csname SRef@cdp:eq:C.RA.28\endcsname{\href{S6_SERIES_GRIND_II.pdf\#nameddest=tag-00XA}{II.{C.RA.28} (p.~502)} }
\expandafter\def\csname SRef@cdp:eq:C.RA.29\endcsname{\href{S6_SERIES_GRIND_II.pdf\#nameddest=tag-00XB}{II.{C.RA.29} (p.~502)} }
\expandafter\def\csname SRef@cdp:eq:C.RA.30\endcsname{\href{S6_SERIES_GRIND_II.pdf\#nameddest=tag-00XD}{II.{C.RA.30} (p.~502)} }
\expandafter\def\csname SRef@cdp:eq:C.RA.31\endcsname{\href{S6_SERIES_GRIND_II.pdf\#nameddest=tag-00XE}{II.{C.RA.31} (p.~502)} }
\expandafter\def\csname SRef@cdp:eq:C.RA.31a\endcsname{\href{S6_SERIES_GRIND_II.pdf\#nameddest=tag-00XF}{II.{C.RA.31a} (p.~503)} }
\expandafter\def\csname SRef@cdp:eq:C.RA.32\endcsname{\href{S6_SERIES_GRIND_II.pdf\#nameddest=tag-00XG}{II.{C.RA.32} (p.~503)} }
\expandafter\def\csname SRef@cdp:eq:C.RD.1\endcsname{\href{S6_SERIES_GRIND_II.pdf\#nameddest=tag-00XZ}{II.{C.RD.1} (p.~503)} }
\expandafter\def\csname SRef@cdp:eq:C.RD.2\endcsname{\href{S6_SERIES_GRIND_II.pdf\#nameddest=tag-00Y9}{II.{C.RD.2} (p.~503)} }
\expandafter\def\csname SRef@cdp:eq:C.RD.3\endcsname{\href{S6_SERIES_GRIND_II.pdf\#nameddest=tag-00YK}{II.{C.RD.3} (p.~504)} }
\expandafter\def\csname SRef@cdp:eq:C.VR.4\endcsname{\href{S6_SERIES_GRIND_II.pdf\#nameddest=tag-015Z}{II.{C.VR.4} (p.~419)} }
\expandafter\def\csname SRef@cdp:eq:C.RD.4\endcsname{\href{S6_SERIES_GRIND_II.pdf\#nameddest=tag-00YT}{II.{C.RD.4} (p.~504)} }
\expandafter\def\csname SRef@cdp:eq:C.RD.11\endcsname{\href{S6_SERIES_GRIND_II.pdf\#nameddest=tag-00Y0}{II.{C.RD.11} (p.~504)} }
\expandafter\def\csname SRef@cdp:eq:C.RD.12\endcsname{\href{S6_SERIES_GRIND_II.pdf\#nameddest=tag-00Y1}{II.{C.RD.12} (p.~505)} }
\expandafter\def\csname SRef@cdp:eq:C.RD.13\endcsname{\href{S6_SERIES_GRIND_II.pdf\#nameddest=tag-00Y2}{II.{C.RD.13} (p.~505)} }
\expandafter\def\csname SRef@cdp:eq:C.RD.14\endcsname{\href{S6_SERIES_GRIND_II.pdf\#nameddest=tag-00Y3}{II.{C.RD.14} (p.~505)} }
\expandafter\def\csname SRef@cdp:eq:C.RD.15\endcsname{\href{S6_SERIES_GRIND_II.pdf\#nameddest=tag-00Y4}{II.{C.RD.15} (p.~506)} }
\expandafter\def\csname SRef@cdp:eq:C.RD.16\endcsname{\href{S6_SERIES_GRIND_II.pdf\#nameddest=tag-00Y5}{II.{C.RD.16} (p.~506)} }
\expandafter\def\csname SRef@cdp:eq:C.RD.17\endcsname{\href{S6_SERIES_GRIND_II.pdf\#nameddest=tag-00Y6}{II.{C.RD.17} (p.~506)} }
\expandafter\def\csname SRef@cdp:eq:C.RD.18\endcsname{\href{S6_SERIES_GRIND_II.pdf\#nameddest=tag-00Y7}{II.{C.RD.18} (p.~506)} }
\expandafter\def\csname SRef@cdp:eq:C.VR.5\endcsname{\href{S6_SERIES_GRIND_II.pdf\#nameddest=tag-0160}{II.{C.VR.5} (p.~419)} }
\expandafter\def\csname SRef@cdp:eq:C.RD.19\endcsname{\href{S6_SERIES_GRIND_II.pdf\#nameddest=tag-00Y8}{II.{C.RD.19} (p.~507)} }
\expandafter\def\csname SRef@cdp:eq:C.RD.20\endcsname{\href{S6_SERIES_GRIND_II.pdf\#nameddest=tag-00YA}{II.{C.RD.20} (p.~507)} }
\expandafter\def\csname SRef@cdp:eq:C.RD.21\endcsname{\href{S6_SERIES_GRIND_II.pdf\#nameddest=tag-00YB}{II.{C.RD.21} (p.~507)} }
\expandafter\def\csname SRef@cdp:eq:C.RD.22\endcsname{\href{S6_SERIES_GRIND_II.pdf\#nameddest=tag-00YC}{II.{C.RD.22} (p.~508)} }
\expandafter\def\csname SRef@cdp:eq:C.RD.23\endcsname{\href{S6_SERIES_GRIND_II.pdf\#nameddest=tag-00YD}{II.{C.RD.23} (p.~508)} }
\expandafter\def\csname SRef@cdp:eq:C.RD.24\endcsname{\href{S6_SERIES_GRIND_II.pdf\#nameddest=tag-00YE}{II.{C.RD.24} (p.~508)} }
\expandafter\def\csname SRef@cdp:eq:C.RD.25\endcsname{\href{S6_SERIES_GRIND_II.pdf\#nameddest=tag-00YF}{II.{C.RD.25} (p.~508)} }
\expandafter\def\csname SRef@cdp:eq:C.RD.26\endcsname{\href{S6_SERIES_GRIND_II.pdf\#nameddest=tag-00YG}{II.{C.RD.26} (p.~508)} }
\expandafter\def\csname SRef@cdp:eq:C.RD.27\endcsname{\href{S6_SERIES_GRIND_II.pdf\#nameddest=tag-00YH}{II.{C.RD.27} (p.~508)} }
\expandafter\def\csname SRef@cdp:eq:C.VR.6\endcsname{\href{S6_SERIES_GRIND_II.pdf\#nameddest=tag-0161}{II.{C.VR.6} (p.~419)} }
\expandafter\def\csname SRef@cdp:eq:C.RD.28\endcsname{\href{S6_SERIES_GRIND_II.pdf\#nameddest=tag-00YI}{II.{C.RD.28} (p.~509)} }
\expandafter\def\csname SRef@cdp:eq:C.RD.29\endcsname{\href{S6_SERIES_GRIND_II.pdf\#nameddest=tag-00YJ}{II.{C.RD.29} (p.~509)} }
\expandafter\def\csname SRef@cdp:eq:C.RD.30\endcsname{\href{S6_SERIES_GRIND_II.pdf\#nameddest=tag-00YL}{II.{C.RD.30} (p.~509)} }
\expandafter\def\csname SRef@cdp:eq:C.RD.31\endcsname{\href{S6_SERIES_GRIND_II.pdf\#nameddest=tag-00YM}{II.{C.RD.31} (p.~510)} }
\expandafter\def\csname SRef@cdp:eq:C.RD.32\endcsname{\href{S6_SERIES_GRIND_II.pdf\#nameddest=tag-00YN}{II.{C.RD.32} (p.~510)} }
\expandafter\def\csname SRef@cdp:eq:C.RD.33\endcsname{\href{S6_SERIES_GRIND_II.pdf\#nameddest=tag-00YO}{II.{C.RD.33} (p.~510)} }
\expandafter\def\csname SRef@cdp:eq:C.RD.34\endcsname{\href{S6_SERIES_GRIND_II.pdf\#nameddest=tag-00YP}{II.{C.RD.34} (p.~510)} }
\expandafter\def\csname SRef@cdp:eq:C.RD.35\endcsname{\href{S6_SERIES_GRIND_II.pdf\#nameddest=tag-00YQ}{II.{C.RD.35} (p.~510)} }
\expandafter\def\csname SRef@cdp:eq:C.VR.7\endcsname{\href{S6_SERIES_GRIND_II.pdf\#nameddest=tag-0162}{II.{C.VR.7} (p.~419)} }
\expandafter\def\csname SRef@cdp:eq:C.RD.36\endcsname{\href{S6_SERIES_GRIND_II.pdf\#nameddest=tag-00YR}{II.{C.RD.36} (p.~511)} }
\expandafter\def\csname SRef@cdp:eq:C.RD.37\endcsname{\href{S6_SERIES_GRIND_II.pdf\#nameddest=tag-00YS}{II.{C.RD.37} (p.~511)} }
\expandafter\def\csname SRef@cdp:eq:C.D4.1\endcsname{\href{S6_SERIES_GRIND_II.pdf\#nameddest=tag-00E1}{II.{C.D4.1} (p.~512)} }
\expandafter\def\csname SRef@cdp:eq:C.D4.2\endcsname{\href{S6_SERIES_GRIND_II.pdf\#nameddest=tag-00EC}{II.{C.D4.2} (p.~512)} }
\expandafter\def\csname SRef@cdp:eq:C.D4.3\endcsname{\href{S6_SERIES_GRIND_II.pdf\#nameddest=tag-00EE}{II.{C.D4.3} (p.~513)} }
\expandafter\def\csname SRef@cdp:eq:C.D4.4\endcsname{\href{S6_SERIES_GRIND_II.pdf\#nameddest=tag-00EF}{II.{C.D4.4} (p.~513)} }
\expandafter\def\csname SRef@cdp:eq:C.D4.5\endcsname{\href{S6_SERIES_GRIND_II.pdf\#nameddest=tag-00EG}{II.{C.D4.5} (p.~513)} }
\expandafter\def\csname SRef@cdp:eq:C.VR.8\endcsname{\href{S6_SERIES_GRIND_II.pdf\#nameddest=tag-0163}{II.{C.VR.8} (p.~420)} }
\expandafter\def\csname SRef@cdp:eq:C.D4.6\endcsname{\href{S6_SERIES_GRIND_II.pdf\#nameddest=tag-00EH}{II.{C.D4.6} (p.~513)} }
\expandafter\def\csname SRef@cdp:eq:C.D4.7\endcsname{\href{S6_SERIES_GRIND_II.pdf\#nameddest=tag-00EI}{II.{C.D4.7} (p.~513)} }
\expandafter\def\csname SRef@cdp:eq:C.D4.8\endcsname{\href{S6_SERIES_GRIND_II.pdf\#nameddest=tag-00EJ}{II.{C.D4.8} (p.~515)} }
\expandafter\def\csname SRef@cdp:eq:C.D4.9\endcsname{\href{S6_SERIES_GRIND_II.pdf\#nameddest=tag-00EK}{II.{C.D4.9} (p.~515)} }
\expandafter\def\csname SRef@cdp:eq:C.D4.10\endcsname{\href{S6_SERIES_GRIND_II.pdf\#nameddest=tag-00E2}{II.{C.D4.10} (p.~515)} }
\expandafter\def\csname SRef@cdp:eq:C.D4.11\endcsname{\href{S6_SERIES_GRIND_II.pdf\#nameddest=tag-00E3}{II.{C.D4.11} (p.~515)} }
\expandafter\def\csname SRef@cdp:eq:C.D4.12\endcsname{\href{S6_SERIES_GRIND_II.pdf\#nameddest=tag-00E4}{II.{C.D4.12} (p.~515)} }
\expandafter\def\csname SRef@cdp:eq:C.DD.1\endcsname{\href{S6_SERIES_GRIND_II.pdf\#nameddest=tag-00EL}{II.{C.DD.1} (p.~516)} }
\expandafter\def\csname SRef@cdp:eq:C.DD.2\endcsname{\href{S6_SERIES_GRIND_II.pdf\#nameddest=tag-00EW}{II.{C.DD.2} (p.~516)} }
\expandafter\def\csname SRef@cdp:eq:C.VR.9\endcsname{\href{S6_SERIES_GRIND_II.pdf\#nameddest=tag-0164}{II.{C.VR.9} (p.~420)} }
\expandafter\def\csname SRef@cdp:eq:C.DD.3\endcsname{\href{S6_SERIES_GRIND_II.pdf\#nameddest=tag-00EZ}{II.{C.DD.3} (p.~516)} }
\expandafter\def\csname SRef@cdp:eq:C.DD.4\endcsname{\href{S6_SERIES_GRIND_II.pdf\#nameddest=tag-00F0}{II.{C.DD.4} (p.~516)} }
\expandafter\def\csname SRef@cdp:eq:C.DD.5\endcsname{\href{S6_SERIES_GRIND_II.pdf\#nameddest=tag-00F1}{II.{C.DD.5} (p.~517)} }
\expandafter\def\csname SRef@cdp:eq:C.DD.6\endcsname{\href{S6_SERIES_GRIND_II.pdf\#nameddest=tag-00F2}{II.{C.DD.6} (p.~517)} }
\expandafter\def\csname SRef@cdp:eq:C.DD.7\endcsname{\href{S6_SERIES_GRIND_II.pdf\#nameddest=tag-00F3}{II.{C.DD.7} (p.~517)} }
\expandafter\def\csname SRef@cdp:eq:C.DD.8\endcsname{\href{S6_SERIES_GRIND_II.pdf\#nameddest=tag-00F4}{II.{C.DD.8} (p.~517)} }
\expandafter\def\csname SRef@cdp:eq:C.DD.9\endcsname{\href{S6_SERIES_GRIND_II.pdf\#nameddest=tag-00F5}{II.{C.DD.9} (p.~517)} }
\expandafter\def\csname SRef@cdp:eq:C.DD.10\endcsname{\href{S6_SERIES_GRIND_II.pdf\#nameddest=tag-00EM}{II.{C.DD.10} (p.~518)} }
\expandafter\def\csname SRef@cdp:eq:C.EC.1\endcsname{\href{S6_SERIES_GRIND_II.pdf\#nameddest=tag-00FO}{II.{C.EC.1} (p.~421)} }
\expandafter\def\csname SRef@cdp:eq:C.DD.11\endcsname{\href{S6_SERIES_GRIND_II.pdf\#nameddest=tag-00EN}{II.{C.DD.11} (p.~518)} }
\expandafter\def\csname SRef@cdp:eq:C.DD.12\endcsname{\href{S6_SERIES_GRIND_II.pdf\#nameddest=tag-00EO}{II.{C.DD.12} (p.~518)} }
\expandafter\def\csname SRef@cdp:eq:C.DD.13\endcsname{\href{S6_SERIES_GRIND_II.pdf\#nameddest=tag-00EP}{II.{C.DD.13} (p.~518)} }
\expandafter\def\csname SRef@cdp:eq:C.DD.14\endcsname{\href{S6_SERIES_GRIND_II.pdf\#nameddest=tag-00EQ}{II.{C.DD.14} (p.~518)} }
\expandafter\def\csname SRef@cdp:eq:C.DD.15\endcsname{\href{S6_SERIES_GRIND_II.pdf\#nameddest=tag-00ER}{II.{C.DD.15} (p.~518)} }
\expandafter\def\csname SRef@cdp:eq:C.DD.16\endcsname{\href{S6_SERIES_GRIND_II.pdf\#nameddest=tag-00ES}{II.{C.DD.16} (p.~519)} }
\expandafter\def\csname SRef@cdp:eq:C.EC.2\endcsname{\href{S6_SERIES_GRIND_II.pdf\#nameddest=tag-00FP}{II.{C.EC.2} (p.~422)} }
\expandafter\def\csname SRef@cdp:eq:C.DD.17\endcsname{\href{S6_SERIES_GRIND_II.pdf\#nameddest=tag-00ET}{II.{C.DD.17} (p.~519)} }
\expandafter\def\csname SRef@cdp:eq:C.DD.18\endcsname{\href{S6_SERIES_GRIND_II.pdf\#nameddest=tag-00EU}{II.{C.DD.18} (p.~519)} }
\expandafter\def\csname SRef@cdp:eq:C.DD.19\endcsname{\href{S6_SERIES_GRIND_II.pdf\#nameddest=tag-00EV}{II.{C.DD.19} (p.~519)} }
\expandafter\def\csname SRef@cdp:eq:C.DD.20\endcsname{\href{S6_SERIES_GRIND_II.pdf\#nameddest=tag-00EX}{II.{C.DD.20} (p.~519)} }
\expandafter\def\csname SRef@cdp:eq:C.DD.21\endcsname{\href{S6_SERIES_GRIND_II.pdf\#nameddest=tag-00EY}{II.{C.DD.21} (p.~520)} }
\expandafter\def\csname SRef@cdp:eq:C.D4.13\endcsname{\href{S6_SERIES_GRIND_II.pdf\#nameddest=tag-00E5}{II.{C.D4.13} (p.~520)} }
\expandafter\def\csname SRef@cdp:eq:C.D4.14\endcsname{\href{S6_SERIES_GRIND_II.pdf\#nameddest=tag-00E6}{II.{C.D4.14} (p.~521)} }
\expandafter\def\csname SRef@cdp:eq:C.D4.15\endcsname{\href{S6_SERIES_GRIND_II.pdf\#nameddest=tag-00E7}{II.{C.D4.15} (p.~521)} }
\expandafter\def\csname SRef@cdp:eq:C.SC.1\endcsname{\href{S6_SERIES_GRIND_II.pdf\#nameddest=tag-010L}{II.{C.SC.1} (p.~423)} }
\expandafter\def\csname SRef@cdp:eq:C.D4.16\endcsname{\href{S6_SERIES_GRIND_II.pdf\#nameddest=tag-00E8}{II.{C.D4.16} (p.~521)} }
\expandafter\def\csname SRef@cdp:eq:C.D4.17\endcsname{\href{S6_SERIES_GRIND_II.pdf\#nameddest=tag-00E9}{II.{C.D4.17} (p.~521)} }
\expandafter\def\csname SRef@cdp:eq:C.D4.18\endcsname{\href{S6_SERIES_GRIND_II.pdf\#nameddest=tag-00EA}{II.{C.D4.18} (p.~522)} }
\expandafter\def\csname SRef@cdp:eq:C.D4.19\endcsname{\href{S6_SERIES_GRIND_II.pdf\#nameddest=tag-00EB}{II.{C.D4.19} (p.~522)} }
\expandafter\def\csname SRef@cdp:eq:C.D4.20\endcsname{\href{S6_SERIES_GRIND_II.pdf\#nameddest=tag-00ED}{II.{C.D4.20} (p.~522)} }
\expandafter\def\csname SRef@cdp:eq:C.UI.1\endcsname{\href{S6_SERIES_GRIND_II.pdf\#nameddest=tag-014M}{II.{C.UI.1} (p.~522)} }
\expandafter\def\csname SRef@cdp:eq:C.UI.2\endcsname{\href{S6_SERIES_GRIND_II.pdf\#nameddest=tag-014U}{II.{C.UI.2} (p.~523)} }
\expandafter\def\csname SRef@cdp:eq:C.UI.3\endcsname{\href{S6_SERIES_GRIND_II.pdf\#nameddest=tag-014V}{II.{C.UI.3} (p.~523)} }
\expandafter\def\csname SRef@cdp:eq:C.UI.4\endcsname{\href{S6_SERIES_GRIND_II.pdf\#nameddest=tag-014W}{II.{C.UI.4} (p.~524)} }
\expandafter\def\csname SRef@cdp:eq:C.SC.2\endcsname{\href{S6_SERIES_GRIND_II.pdf\#nameddest=tag-010M}{II.{C.SC.2} (p.~423)} }
\expandafter\def\csname SRef@cdp:eq:C.UI.5\endcsname{\href{S6_SERIES_GRIND_II.pdf\#nameddest=tag-014X}{II.{C.UI.5} (p.~524)} }
\expandafter\def\csname SRef@cdp:eq:C.UI.6\endcsname{\href{S6_SERIES_GRIND_II.pdf\#nameddest=tag-014Y}{II.{C.UI.6} (p.~524)} }
\expandafter\def\csname SRef@cdp:eq:C.UI.7\endcsname{\href{S6_SERIES_GRIND_II.pdf\#nameddest=tag-014Z}{II.{C.UI.7} (p.~524)} }
\expandafter\def\csname SRef@cdp:eq:C.UI.8\endcsname{\href{S6_SERIES_GRIND_II.pdf\#nameddest=tag-0150}{II.{C.UI.8} (p.~525)} }
\expandafter\def\csname SRef@cdp:eq:C.UI.9\endcsname{\href{S6_SERIES_GRIND_II.pdf\#nameddest=tag-0151}{II.{C.UI.9} (p.~525)} }
\expandafter\def\csname SRef@cdp:eq:C.UI.10\endcsname{\href{S6_SERIES_GRIND_II.pdf\#nameddest=tag-014N}{II.{C.UI.10} (p.~526)} }
\expandafter\def\csname SRef@cdp:eq:C.UI.11\endcsname{\href{S6_SERIES_GRIND_II.pdf\#nameddest=tag-014O}{II.{C.UI.11} (p.~526)} }
\expandafter\def\csname SRef@cdp:eq:C.UI.12\endcsname{\href{S6_SERIES_GRIND_II.pdf\#nameddest=tag-014P}{II.{C.UI.12} (p.~526)} }
\expandafter\def\csname SRef@cdp:eq:C.UI.13\endcsname{\href{S6_SERIES_GRIND_II.pdf\#nameddest=tag-014Q}{II.{C.UI.13} (p.~527)} }
\expandafter\def\csname SRef@cdp:eq:C.UI.14\endcsname{\href{S6_SERIES_GRIND_II.pdf\#nameddest=tag-014R}{II.{C.UI.14} (p.~528)} }
\expandafter\def\csname SRef@cdp:eq:C.SC.3\endcsname{\href{S6_SERIES_GRIND_II.pdf\#nameddest=tag-010N}{II.{C.SC.3} (p.~423)} }
\expandafter\def\csname SRef@cdp:eq:C.UI.15\endcsname{\href{S6_SERIES_GRIND_II.pdf\#nameddest=tag-014S}{II.{C.UI.15} (p.~528)} }
\expandafter\def\csname SRef@cdp:eq:C.UI.16\endcsname{\href{S6_SERIES_GRIND_II.pdf\#nameddest=tag-014T}{II.{C.UI.16} (p.~528)} }
\expandafter\def\csname SRef@cdp:eq:C.IZ.1\endcsname{\href{S6_SERIES_GRIND_II.pdf\#nameddest=tag-00ON}{II.{C.IZ.1} (p.~529)} }
\expandafter\def\csname SRef@cdp:eq:C.IZ.2\endcsname{\href{S6_SERIES_GRIND_II.pdf\#nameddest=tag-00OO}{II.{C.IZ.2} (p.~529)} }
\expandafter\def\csname SRef@cdp:eq:C.IZ.3\endcsname{\href{S6_SERIES_GRIND_II.pdf\#nameddest=tag-00OP}{II.{C.IZ.3} (p.~530)} }
\expandafter\def\csname SRef@cdp:eq:C.IZ.4\endcsname{\href{S6_SERIES_GRIND_II.pdf\#nameddest=tag-00OQ}{II.{C.IZ.4} (p.~530)} }
\expandafter\def\csname SRef@cdp:eq:C.IZ.5\endcsname{\href{S6_SERIES_GRIND_II.pdf\#nameddest=tag-00OR}{II.{C.IZ.5} (p.~530)} }
\expandafter\def\csname SRef@cdp:eq:C.IZ.6\endcsname{\href{S6_SERIES_GRIND_II.pdf\#nameddest=tag-00OS}{II.{C.IZ.6} (p.~530)} }
\expandafter\def\csname SRef@cdp:eq:C.IZ.7\endcsname{\href{S6_SERIES_GRIND_II.pdf\#nameddest=tag-00OT}{II.{C.IZ.7} (p.~530)} }
\expandafter\def\csname SRef@cdp:eq:C.IZ.8\endcsname{\href{S6_SERIES_GRIND_II.pdf\#nameddest=tag-00OU}{II.{C.IZ.8} (p.~531)} }
\expandafter\def\csname SRef@cdp:eq:C.SC.4\endcsname{\href{S6_SERIES_GRIND_II.pdf\#nameddest=tag-010O}{II.{C.SC.4} (p.~423)} }
\expandafter\def\csname SRef@cdp:eq:C.INT.1\endcsname{\href{S6_SERIES_GRIND_II.pdf\#nameddest=tag-00NQ}{II.{C.INT.1} (p.~531)} }
\expandafter\def\csname SRef@cdp:eq:C.INT.2\endcsname{\href{S6_SERIES_GRIND_II.pdf\#nameddest=tag-00O1}{II.{C.INT.2} (p.~531)} }
\expandafter\def\csname SRef@cdp:eq:C.INT.3\endcsname{\href{S6_SERIES_GRIND_II.pdf\#nameddest=tag-00OA}{II.{C.INT.3} (p.~532)} }
\expandafter\def\csname SRef@cdp:eq:C.INT.4\endcsname{\href{S6_SERIES_GRIND_II.pdf\#nameddest=tag-00OB}{II.{C.INT.4} (p.~532)} }
\expandafter\def\csname SRef@cdp:eq:C.INT.5\endcsname{\href{S6_SERIES_GRIND_II.pdf\#nameddest=tag-00OC}{II.{C.INT.5} (p.~533)} }
\expandafter\def\csname SRef@cdp:eq:C.INT.6\endcsname{\href{S6_SERIES_GRIND_II.pdf\#nameddest=tag-00OD}{II.{C.INT.6} (p.~533)} }
\expandafter\def\csname SRef@cdp:eq:C.INT.7\endcsname{\href{S6_SERIES_GRIND_II.pdf\#nameddest=tag-00OE}{II.{C.INT.7} (p.~533)} }
\expandafter\def\csname SRef@cdp:eq:C.INT.8\endcsname{\href{S6_SERIES_GRIND_II.pdf\#nameddest=tag-00OF}{II.{C.INT.8} (p.~534)} }
\expandafter\def\csname SRef@cdp:eq:C.INT.9\endcsname{\href{S6_SERIES_GRIND_II.pdf\#nameddest=tag-00OG}{II.{C.INT.9} (p.~534)} }
\expandafter\def\csname SRef@cdp:eq:C.INT.10\endcsname{\href{S6_SERIES_GRIND_II.pdf\#nameddest=tag-00NR}{II.{C.INT.10} (p.~534)} }
\expandafter\def\csname SRef@cdp:eq:C.SC.5\endcsname{\href{S6_SERIES_GRIND_II.pdf\#nameddest=tag-010P}{II.{C.SC.5} (p.~424)} }
\expandafter\def\csname SRef@cdp:eq:C.INT.11\endcsname{\href{S6_SERIES_GRIND_II.pdf\#nameddest=tag-00NS}{II.{C.INT.11} (p.~535)} }
\expandafter\def\csname SRef@cdp:eq:C.INT.12\endcsname{\href{S6_SERIES_GRIND_II.pdf\#nameddest=tag-00NT}{II.{C.INT.12} (p.~535)} }
\expandafter\def\csname SRef@cdp:eq:C.INT.13\endcsname{\href{S6_SERIES_GRIND_II.pdf\#nameddest=tag-00NU}{II.{C.INT.13} (p.~535)} }
\expandafter\def\csname SRef@cdp:eq:C.INT.14\endcsname{\href{S6_SERIES_GRIND_II.pdf\#nameddest=tag-00NV}{II.{C.INT.14} (p.~535)} }
\expandafter\def\csname SRef@cdp:eq:C.INT.15\endcsname{\href{S6_SERIES_GRIND_II.pdf\#nameddest=tag-00NW}{II.{C.INT.15} (p.~535)} }
\expandafter\def\csname SRef@cdp:eq:C.INT.16\endcsname{\href{S6_SERIES_GRIND_II.pdf\#nameddest=tag-00NX}{II.{C.INT.16} (p.~536)} }
\expandafter\def\csname SRef@cdp:eq:C.INT.17\endcsname{\href{S6_SERIES_GRIND_II.pdf\#nameddest=tag-00NY}{II.{C.INT.17} (p.~536)} }
\expandafter\def\csname SRef@cdp:eq:C.INT.18\endcsname{\href{S6_SERIES_GRIND_II.pdf\#nameddest=tag-00NZ}{II.{C.INT.18} (p.~536)} }
\expandafter\def\csname SRef@cdp:eq:C.INT.19\endcsname{\href{S6_SERIES_GRIND_II.pdf\#nameddest=tag-00O0}{II.{C.INT.19} (p.~536)} }
\expandafter\def\csname SRef@cdp:eq:C.INT.20\endcsname{\href{S6_SERIES_GRIND_II.pdf\#nameddest=tag-00O2}{II.{C.INT.20} (p.~536)} }
\expandafter\def\csname SRef@cdp:eq:C.SC.6\endcsname{\href{S6_SERIES_GRIND_II.pdf\#nameddest=tag-010Q}{II.{C.SC.6} (p.~425)} }
\expandafter\def\csname SRef@cdp:eq:C.INT.21\endcsname{\href{S6_SERIES_GRIND_II.pdf\#nameddest=tag-00O3}{II.{C.INT.21} (p.~536)} }
\expandafter\def\csname SRef@cdp:eq:C.INT.22\endcsname{\href{S6_SERIES_GRIND_II.pdf\#nameddest=tag-00O4}{II.{C.INT.22} (p.~537)} }
\expandafter\def\csname SRef@cdp:eq:C.INT.23\endcsname{\href{S6_SERIES_GRIND_II.pdf\#nameddest=tag-00O5}{II.{C.INT.23} (p.~537)} }
\expandafter\def\csname SRef@cdp:eq:C.INT.24\endcsname{\href{S6_SERIES_GRIND_II.pdf\#nameddest=tag-00O6}{II.{C.INT.24} (p.~537)} }
\expandafter\def\csname SRef@cdp:eq:C.INT.25\endcsname{\href{S6_SERIES_GRIND_II.pdf\#nameddest=tag-00O7}{II.{C.INT.25} (p.~537)} }
\expandafter\def\csname SRef@cdp:eq:C.INT.26\endcsname{\href{S6_SERIES_GRIND_II.pdf\#nameddest=tag-00O8}{II.{C.INT.26} (p.~537)} }
\expandafter\def\csname SRef@cdp:eq:C.INT.27\endcsname{\href{S6_SERIES_GRIND_II.pdf\#nameddest=tag-00O9}{II.{C.INT.27} (p.~537)} }
\expandafter\def\csname SRef@cdp:eq:C.MH.1\endcsname{\href{S6_SERIES_GRIND_II.pdf\#nameddest=tag-00QA}{II.{C.MH.1} (p.~538)} }
\expandafter\def\csname SRef@cdp:eq:C.MH.2\endcsname{\href{S6_SERIES_GRIND_II.pdf\#nameddest=tag-00QJ}{II.{C.MH.2} (p.~538)} }
\expandafter\def\csname SRef@cdp:eq:C.MH.3\endcsname{\href{S6_SERIES_GRIND_II.pdf\#nameddest=tag-00QK}{II.{C.MH.3} (p.~539)} }
\expandafter\def\csname SRef@cdp:eq:C.CF.1\endcsname{\href{S6_SERIES_GRIND_II.pdf\#nameddest=tag-00BN}{II.{C.CF.1} (p.~425)} }
\expandafter\def\csname SRef@cdp:eq:C.MH.4\endcsname{\href{S6_SERIES_GRIND_II.pdf\#nameddest=tag-00QL}{II.{C.MH.4} (p.~540)} }
\expandafter\def\csname SRef@cdp:eq:C.MH.5\endcsname{\href{S6_SERIES_GRIND_II.pdf\#nameddest=tag-00QM}{II.{C.MH.5} (p.~540)} }
\expandafter\def\csname SRef@cdp:eq:C.MH.6\endcsname{\href{S6_SERIES_GRIND_II.pdf\#nameddest=tag-00QN}{II.{C.MH.6} (p.~540)} }
\expandafter\def\csname SRef@cdp:eq:C.MH.7\endcsname{\href{S6_SERIES_GRIND_II.pdf\#nameddest=tag-00QO}{II.{C.MH.7} (p.~540)} }
\expandafter\def\csname SRef@cdp:eq:C.MH.8\endcsname{\href{S6_SERIES_GRIND_II.pdf\#nameddest=tag-00QP}{II.{C.MH.8} (p.~541)} }
\expandafter\def\csname SRef@cdp:eq:C.MH.9\endcsname{\href{S6_SERIES_GRIND_II.pdf\#nameddest=tag-00QQ}{II.{C.MH.9} (p.~541)} }
\expandafter\def\csname SRef@cdp:eq:C.MH.10\endcsname{\href{S6_SERIES_GRIND_II.pdf\#nameddest=tag-00QB}{II.{C.MH.10} (p.~541)} }
\expandafter\def\csname SRef@cdp:eq:C.MH.11\endcsname{\href{S6_SERIES_GRIND_II.pdf\#nameddest=tag-00QC}{II.{C.MH.11} (p.~541)} }
\expandafter\def\csname SRef@cdp:eq:C.MH.12\endcsname{\href{S6_SERIES_GRIND_II.pdf\#nameddest=tag-00QD}{II.{C.MH.12} (p.~541)} }
\expandafter\def\csname SRef@cdp:eq:C.MH.13\endcsname{\href{S6_SERIES_GRIND_II.pdf\#nameddest=tag-00QE}{II.{C.MH.13} (p.~541)} }
\expandafter\def\csname SRef@cdp:eq:C.CF.2\endcsname{\href{S6_SERIES_GRIND_II.pdf\#nameddest=tag-00BO}{II.{C.CF.2} (p.~426)} }
\expandafter\def\csname SRef@cdp:eq:C.MH.14\endcsname{\href{S6_SERIES_GRIND_II.pdf\#nameddest=tag-00QF}{II.{C.MH.14} (p.~541)} }
\expandafter\def\csname SRef@cdp:eq:C.MH.15\endcsname{\href{S6_SERIES_GRIND_II.pdf\#nameddest=tag-00QG}{II.{C.MH.15} (p.~542)} }
\expandafter\def\csname SRef@cdp:eq:C.MH.16\endcsname{\href{S6_SERIES_GRIND_II.pdf\#nameddest=tag-00QH}{II.{C.MH.16} (p.~542)} }
\expandafter\def\csname SRef@cdp:eq:C.MH.17\endcsname{\href{S6_SERIES_GRIND_II.pdf\#nameddest=tag-00QI}{II.{C.MH.17} (p.~542)} }
\expandafter\def\csname SRef@cdp:eq:C.1\endcsname{\href{S6_SERIES_GRIND_II.pdf\#nameddest=tag-008N}{II.{C.1} (p.~543)} }
\expandafter\def\csname SRef@cdp:eq:C.2\endcsname{\href{S6_SERIES_GRIND_II.pdf\#nameddest=tag-008Y}{II.{C.2} (p.~544)} }
\expandafter\def\csname SRef@cdp:eq:C.MS.1\endcsname{\href{S6_SERIES_GRIND_II.pdf\#nameddest=tag-00RH}{II.{C.MS.1} (p.~546)} }
\expandafter\def\csname SRef@cdp:eq:C.MS.2\endcsname{\href{S6_SERIES_GRIND_II.pdf\#nameddest=tag-00RJ}{II.{C.MS.2} (p.~546)} }
\expandafter\def\csname SRef@cdp:eq:C.MS.3\endcsname{\href{S6_SERIES_GRIND_II.pdf\#nameddest=tag-00RK}{II.{C.MS.3} (p.~547)} }
\expandafter\def\csname SRef@cdp:eq:C.MS.10\endcsname{\href{S6_SERIES_GRIND_II.pdf\#nameddest=tag-00RI}{II.{C.MS.10} (p.~547)} }
\expandafter\def\csname SRef@cdp:eq:C.CF.3\endcsname{\href{S6_SERIES_GRIND_II.pdf\#nameddest=tag-00BP}{II.{C.CF.3} (p.~426)} }
\expandafter\def\csname SRef@cdp:eq:C.MS.4\endcsname{\href{S6_SERIES_GRIND_II.pdf\#nameddest=tag-00RL}{II.{C.MS.4} (p.~547)} }
\expandafter\def\csname SRef@cdp:eq:C.MS.5\endcsname{\href{S6_SERIES_GRIND_II.pdf\#nameddest=tag-00RM}{II.{C.MS.5} (p.~547)} }
\expandafter\def\csname SRef@cdp:eq:C.MS.6\endcsname{\href{S6_SERIES_GRIND_II.pdf\#nameddest=tag-00RN}{II.{C.MS.6} (p.~548)} }
\expandafter\def\csname SRef@cdp:eq:C.MS.7\endcsname{\href{S6_SERIES_GRIND_II.pdf\#nameddest=tag-00RO}{II.{C.MS.7} (p.~548)} }
\expandafter\def\csname SRef@cdp:eq:C.MS.8\endcsname{\href{S6_SERIES_GRIND_II.pdf\#nameddest=tag-00RP}{II.{C.MS.8} (p.~549)} }
\expandafter\def\csname SRef@cdp:eq:C.MS.9\endcsname{\href{S6_SERIES_GRIND_II.pdf\#nameddest=tag-00RQ}{II.{C.MS.9} (p.~549)} }
\expandafter\def\csname SRef@cdp:eq:C.K3.1\endcsname{\href{S6_SERIES_GRIND_II.pdf\#nameddest=tag-00OY}{II.{C.K3.1} (p.~551)} }
\expandafter\def\csname SRef@cdp:eq:C.K3.2\endcsname{\href{S6_SERIES_GRIND_II.pdf\#nameddest=tag-00P0}{II.{C.K3.2} (p.~551)} }
\expandafter\def\csname SRef@cdp:eq:C.K3.3\endcsname{\href{S6_SERIES_GRIND_II.pdf\#nameddest=tag-00P1}{II.{C.K3.3} (p.~551)} }
\expandafter\def\csname SRef@cdp:eq:C.K3.4\endcsname{\href{S6_SERIES_GRIND_II.pdf\#nameddest=tag-00P2}{II.{C.K3.4} (p.~551)} }
\expandafter\def\csname SRef@cdp:eq:C.CF.4\endcsname{\href{S6_SERIES_GRIND_II.pdf\#nameddest=tag-00BQ}{II.{C.CF.4} (p.~426)} }
\expandafter\def\csname SRef@cdp:eq:C.K3.5\endcsname{\href{S6_SERIES_GRIND_II.pdf\#nameddest=tag-00P3}{II.{C.K3.5} (p.~552)} }
\expandafter\def\csname SRef@cdp:eq:C.K3.6\endcsname{\href{S6_SERIES_GRIND_II.pdf\#nameddest=tag-00P4}{II.{C.K3.6} (p.~552)} }
\expandafter\def\csname SRef@cdp:eq:C.K3.7\endcsname{\href{S6_SERIES_GRIND_II.pdf\#nameddest=tag-00P5}{II.{C.K3.7} (p.~552)} }
\expandafter\def\csname SRef@cdp:eq:C.K3.8\endcsname{\href{S6_SERIES_GRIND_II.pdf\#nameddest=tag-00P6}{II.{C.K3.8} (p.~553)} }
\expandafter\def\csname SRef@cdp:eq:C.K3.9\endcsname{\href{S6_SERIES_GRIND_II.pdf\#nameddest=tag-00P7}{II.{C.K3.9} (p.~553)} }
\expandafter\def\csname SRef@cdp:eq:C.K3.10\endcsname{\href{S6_SERIES_GRIND_II.pdf\#nameddest=tag-00OZ}{II.{C.K3.10} (p.~553)} }
\expandafter\def\csname SRef@cdp:eq:C.KX.1\endcsname{\href{S6_SERIES_GRIND_II.pdf\#nameddest=tag-00PW}{II.{C.KX.1} (p.~554)} }
\expandafter\def\csname SRef@cdp:eq:C.KX.2\endcsname{\href{S6_SERIES_GRIND_II.pdf\#nameddest=tag-00Q2}{II.{C.KX.2} (p.~554)} }
\expandafter\def\csname SRef@cdp:eq:C.KX.3\endcsname{\href{S6_SERIES_GRIND_II.pdf\#nameddest=tag-00Q3}{II.{C.KX.3} (p.~554)} }
\expandafter\def\csname SRef@cdp:eq:C.KX.4\endcsname{\href{S6_SERIES_GRIND_II.pdf\#nameddest=tag-00Q4}{II.{C.KX.4} (p.~555)} }
\expandafter\def\csname SRef@cdp:eq:C.CF.5\endcsname{\href{S6_SERIES_GRIND_II.pdf\#nameddest=tag-00BR}{II.{C.CF.5} (p.~426)} }
\expandafter\def\csname SRef@cdp:eq:C.KX.5\endcsname{\href{S6_SERIES_GRIND_II.pdf\#nameddest=tag-00Q5}{II.{C.KX.5} (p.~555)} }
\expandafter\def\csname SRef@cdp:eq:C.KX.6\endcsname{\href{S6_SERIES_GRIND_II.pdf\#nameddest=tag-00Q6}{II.{C.KX.6} (p.~555)} }
\expandafter\def\csname SRef@cdp:eq:C.KX.7\endcsname{\href{S6_SERIES_GRIND_II.pdf\#nameddest=tag-00Q7}{II.{C.KX.7} (p.~555)} }
\expandafter\def\csname SRef@cdp:eq:C.KX.8\endcsname{\href{S6_SERIES_GRIND_II.pdf\#nameddest=tag-00Q8}{II.{C.KX.8} (p.~556)} }
\expandafter\def\csname SRef@cdp:eq:C.KX.9\endcsname{\href{S6_SERIES_GRIND_II.pdf\#nameddest=tag-00Q9}{II.{C.KX.9} (p.~556)} }
\expandafter\def\csname SRef@cdp:eq:C.KX.10\endcsname{\href{S6_SERIES_GRIND_II.pdf\#nameddest=tag-00PX}{II.{C.KX.10} (p.~556)} }
\expandafter\def\csname SRef@cdp:eq:C.KX.11\endcsname{\href{S6_SERIES_GRIND_II.pdf\#nameddest=tag-00PY}{II.{C.KX.11} (p.~556)} }
\expandafter\def\csname SRef@cdp:eq:C.KX.12\endcsname{\href{S6_SERIES_GRIND_II.pdf\#nameddest=tag-00PZ}{II.{C.KX.12} (p.~556)} }
\expandafter\def\csname SRef@cdp:eq:C.KX.13\endcsname{\href{S6_SERIES_GRIND_II.pdf\#nameddest=tag-00Q0}{II.{C.KX.13} (p.~557)} }
\expandafter\def\csname SRef@cdp:eq:C.KX.14\endcsname{\href{S6_SERIES_GRIND_II.pdf\#nameddest=tag-00Q1}{II.{C.KX.14} (p.~557)} }
\expandafter\def\csname SRef@cdp:eq:C.CF.6\endcsname{\href{S6_SERIES_GRIND_II.pdf\#nameddest=tag-00BS}{II.{C.CF.6} (p.~426)} }
\expandafter\def\csname SRef@cdp:eq:C.EL.1\endcsname{\href{S6_SERIES_GRIND_II.pdf\#nameddest=tag-00FY}{II.{C.EL.1} (p.~557)} }
\expandafter\def\csname SRef@cdp:eq:C.EL.2\endcsname{\href{S6_SERIES_GRIND_II.pdf\#nameddest=tag-00G1}{II.{C.EL.2} (p.~557)} }
\expandafter\def\csname SRef@cdp:eq:C.EL.3\endcsname{\href{S6_SERIES_GRIND_II.pdf\#nameddest=tag-00G2}{II.{C.EL.3} (p.~557)} }
\expandafter\def\csname SRef@cdp:eq:C.EL.4\endcsname{\href{S6_SERIES_GRIND_II.pdf\#nameddest=tag-00G3}{II.{C.EL.4} (p.~558)} }
\expandafter\def\csname SRef@cdp:eq:C.EL.5\endcsname{\href{S6_SERIES_GRIND_II.pdf\#nameddest=tag-00G4}{II.{C.EL.5} (p.~558)} }
\expandafter\def\csname SRef@cdp:eq:C.EL.6\endcsname{\href{S6_SERIES_GRIND_II.pdf\#nameddest=tag-00G5}{II.{C.EL.6} (p.~558)} }
\expandafter\def\csname SRef@cdp:eq:C.EL.7\endcsname{\href{S6_SERIES_GRIND_II.pdf\#nameddest=tag-00G6}{II.{C.EL.7} (p.~558)} }
\expandafter\def\csname SRef@cdp:eq:C.EL.8\endcsname{\href{S6_SERIES_GRIND_II.pdf\#nameddest=tag-00G7}{II.{C.EL.8} (p.~558)} }
\expandafter\def\csname SRef@cdp:eq:C.EL.9\endcsname{\href{S6_SERIES_GRIND_II.pdf\#nameddest=tag-00G8}{II.{C.EL.9} (p.~559)} }
\expandafter\def\csname SRef@cdp:eq:C.EL.10\endcsname{\href{S6_SERIES_GRIND_II.pdf\#nameddest=tag-00FZ}{II.{C.EL.10} (p.~560)} }
\expandafter\def\csname SRef@cdp:eq:C.CF.7\endcsname{\href{S6_SERIES_GRIND_II.pdf\#nameddest=tag-00BT}{II.{C.CF.7} (p.~427)} }
\expandafter\def\csname SRef@cdp:eq:C.EL.11\endcsname{\href{S6_SERIES_GRIND_II.pdf\#nameddest=tag-00G0}{II.{C.EL.11} (p.~560)} }
\expandafter\def\csname SRef@cdp:eq:C.CZ.1\endcsname{\href{S6_SERIES_GRIND_II.pdf\#nameddest=tag-00DN}{II.{C.CZ.1} (p.~561)} }
\expandafter\def\csname SRef@cdp:eq:C.CZ.2\endcsname{\href{S6_SERIES_GRIND_II.pdf\#nameddest=tag-00DT}{II.{C.CZ.2} (p.~561)} }
\expandafter\def\csname SRef@cdp:eq:C.CZ.3\endcsname{\href{S6_SERIES_GRIND_II.pdf\#nameddest=tag-00DU}{II.{C.CZ.3} (p.~561)} }
\expandafter\def\csname SRef@cdp:eq:C.CZ.4\endcsname{\href{S6_SERIES_GRIND_II.pdf\#nameddest=tag-00DV}{II.{C.CZ.4} (p.~561)} }
\expandafter\def\csname SRef@cdp:eq:C.CZ.5\endcsname{\href{S6_SERIES_GRIND_II.pdf\#nameddest=tag-00DW}{II.{C.CZ.5} (p.~562)} }
\expandafter\def\csname SRef@cdp:eq:C.CZ.6\endcsname{\href{S6_SERIES_GRIND_II.pdf\#nameddest=tag-00DX}{II.{C.CZ.6} (p.~562)} }
\expandafter\def\csname SRef@cdp:eq:C.CZ.7\endcsname{\href{S6_SERIES_GRIND_II.pdf\#nameddest=tag-00DY}{II.{C.CZ.7} (p.~562)} }
\expandafter\def\csname SRef@cdp:eq:C.CZ.8\endcsname{\href{S6_SERIES_GRIND_II.pdf\#nameddest=tag-00DZ}{II.{C.CZ.8} (p.~562)} }
\expandafter\def\csname SRef@cdp:eq:C.CZ.9\endcsname{\href{S6_SERIES_GRIND_II.pdf\#nameddest=tag-00E0}{II.{C.CZ.9} (p.~562)} }
\expandafter\def\csname SRef@cdp:eq:C.CF.8\endcsname{\href{S6_SERIES_GRIND_II.pdf\#nameddest=tag-00BU}{II.{C.CF.8} (p.~427)} }
\expandafter\def\csname SRef@cdp:eq:C.CZ.10\endcsname{\href{S6_SERIES_GRIND_II.pdf\#nameddest=tag-00DO}{II.{C.CZ.10} (p.~563)} }
\expandafter\def\csname SRef@cdp:eq:C.CZ.11\endcsname{\href{S6_SERIES_GRIND_II.pdf\#nameddest=tag-00DP}{II.{C.CZ.11} (p.~563)} }
\expandafter\def\csname SRef@cdp:eq:C.CZ.12\endcsname{\href{S6_SERIES_GRIND_II.pdf\#nameddest=tag-00DQ}{II.{C.CZ.12} (p.~564)} }
\expandafter\def\csname SRef@cdp:eq:C.CZ.13\endcsname{\href{S6_SERIES_GRIND_II.pdf\#nameddest=tag-00DR}{II.{C.CZ.13} (p.~564)} }
\expandafter\def\csname SRef@cdp:eq:C.CZ.14\endcsname{\href{S6_SERIES_GRIND_II.pdf\#nameddest=tag-00DS}{II.{C.CZ.14} (p.~564)} }
\expandafter\def\csname SRef@cdp:eq:C.MR.1\endcsname{\href{S6_SERIES_GRIND_II.pdf\#nameddest=tag-00R1}{II.{C.MR.1} (p.~566)} }
\expandafter\def\csname SRef@cdp:eq:C.MR.2\endcsname{\href{S6_SERIES_GRIND_II.pdf\#nameddest=tag-00R9}{II.{C.MR.2} (p.~566)} }
\expandafter\def\csname SRef@cdp:eq:C.MR.3\endcsname{\href{S6_SERIES_GRIND_II.pdf\#nameddest=tag-00RA}{II.{C.MR.3} (p.~566)} }
\expandafter\def\csname SRef@cdp:eq:C.MR.4\endcsname{\href{S6_SERIES_GRIND_II.pdf\#nameddest=tag-00RB}{II.{C.MR.4} (p.~566)} }
\expandafter\def\csname SRef@cdp:eq:C.MR.5\endcsname{\href{S6_SERIES_GRIND_II.pdf\#nameddest=tag-00RC}{II.{C.MR.5} (p.~567)} }
\expandafter\def\csname SRef@cdp:eq:C.CF.9\endcsname{\href{S6_SERIES_GRIND_II.pdf\#nameddest=tag-00BV}{II.{C.CF.9} (p.~427)} }
\expandafter\def\csname SRef@cdp:eq:C.MR.6\endcsname{\href{S6_SERIES_GRIND_II.pdf\#nameddest=tag-00RD}{II.{C.MR.6} (p.~567)} }
\expandafter\def\csname SRef@cdp:eq:C.MR.7\endcsname{\href{S6_SERIES_GRIND_II.pdf\#nameddest=tag-00RE}{II.{C.MR.7} (p.~567)} }
\expandafter\def\csname SRef@cdp:eq:C.MR.8\endcsname{\href{S6_SERIES_GRIND_II.pdf\#nameddest=tag-00RF}{II.{C.MR.8} (p.~568)} }
\expandafter\def\csname SRef@cdp:eq:C.MR.9\endcsname{\href{S6_SERIES_GRIND_II.pdf\#nameddest=tag-00RG}{II.{C.MR.9} (p.~568)} }
\expandafter\def\csname SRef@cdp:eq:C.MR.10\endcsname{\href{S6_SERIES_GRIND_II.pdf\#nameddest=tag-00R2}{II.{C.MR.10} (p.~569)} }
\expandafter\def\csname SRef@cdp:eq:C.MR.11\endcsname{\href{S6_SERIES_GRIND_II.pdf\#nameddest=tag-00R3}{II.{C.MR.11} (p.~569)} }
\expandafter\def\csname SRef@cdp:eq:C.MR.12\endcsname{\href{S6_SERIES_GRIND_II.pdf\#nameddest=tag-00R4}{II.{C.MR.12} (p.~570)} }
\expandafter\def\csname SRef@cdp:eq:C.MR.13\endcsname{\href{S6_SERIES_GRIND_II.pdf\#nameddest=tag-00R5}{II.{C.MR.13} (p.~570)} }
\expandafter\def\csname SRef@cdp:eq:C.MR.14\endcsname{\href{S6_SERIES_GRIND_II.pdf\#nameddest=tag-00R6}{II.{C.MR.14} (p.~570)} }
\expandafter\def\csname SRef@cdp:eq:C.MR.15\endcsname{\href{S6_SERIES_GRIND_II.pdf\#nameddest=tag-00R7}{II.{C.MR.15} (p.~571)} }
\expandafter\def\csname SRef@cdp:eq:C.CA.1\endcsname{\href{S6_SERIES_GRIND_II.pdf\#nameddest=tag-00BC}{II.{C.CA.1} (p.~428)} }
\expandafter\def\csname SRef@cdp:eq:C.MR.16\endcsname{\href{S6_SERIES_GRIND_II.pdf\#nameddest=tag-00R8}{II.{C.MR.16} (p.~572)} }
\expandafter\def\csname SRef@cdp:eq:C.RH.1\endcsname{\href{S6_SERIES_GRIND_II.pdf\#nameddest=tag-00ZG}{II.{C.RH.1} (p.~573)} }
\expandafter\def\csname SRef@cdp:eq:C.RH.2\endcsname{\href{S6_SERIES_GRIND_II.pdf\#nameddest=tag-00ZI}{II.{C.RH.2} (p.~573)} }
\expandafter\def\csname SRef@cdp:eq:C.RH.3\endcsname{\href{S6_SERIES_GRIND_II.pdf\#nameddest=tag-00ZJ}{II.{C.RH.3} (p.~573)} }
\expandafter\def\csname SRef@cdp:eq:C.RH.4\endcsname{\href{S6_SERIES_GRIND_II.pdf\#nameddest=tag-00ZK}{II.{C.RH.4} (p.~573)} }
\expandafter\def\csname SRef@cdp:eq:C.RH.5\endcsname{\href{S6_SERIES_GRIND_II.pdf\#nameddest=tag-00ZL}{II.{C.RH.5} (p.~573)} }
\expandafter\def\csname SRef@cdp:eq:C.RH.6\endcsname{\href{S6_SERIES_GRIND_II.pdf\#nameddest=tag-00ZM}{II.{C.RH.6} (p.~573)} }
\expandafter\def\csname SRef@cdp:eq:C.RH.7\endcsname{\href{S6_SERIES_GRIND_II.pdf\#nameddest=tag-00ZN}{II.{C.RH.7} (p.~574)} }
\expandafter\def\csname SRef@cdp:eq:C.RH.8\endcsname{\href{S6_SERIES_GRIND_II.pdf\#nameddest=tag-00ZO}{II.{C.RH.8} (p.~575)} }
\expandafter\def\csname SRef@cdp:eq:C.RH.9\endcsname{\href{S6_SERIES_GRIND_II.pdf\#nameddest=tag-00ZP}{II.{C.RH.9} (p.~575)} }
\expandafter\def\csname SRef@cdp:eq:C.CA.2\endcsname{\href{S6_SERIES_GRIND_II.pdf\#nameddest=tag-00BF}{II.{C.CA.2} (p.~428)} }
\expandafter\def\csname SRef@cdp:eq:C.RH.10\endcsname{\href{S6_SERIES_GRIND_II.pdf\#nameddest=tag-00ZH}{II.{C.RH.10} (p.~575)} }
\expandafter\def\csname SRef@cdp:eq:C.TC.1\endcsname{\href{S6_SERIES_GRIND_II.pdf\#nameddest=tag-012C}{II.{C.TC.1} (p.~576)} }
\expandafter\def\csname SRef@cdp:eq:C.TC.2\endcsname{\href{S6_SERIES_GRIND_II.pdf\#nameddest=tag-012G}{II.{C.TC.2} (p.~576)} }
\expandafter\def\csname SRef@cdp:eq:C.TC.3\endcsname{\href{S6_SERIES_GRIND_II.pdf\#nameddest=tag-012H}{II.{C.TC.3} (p.~576)} }
\expandafter\def\csname SRef@cdp:eq:C.TC.4\endcsname{\href{S6_SERIES_GRIND_II.pdf\#nameddest=tag-012I}{II.{C.TC.4} (p.~577)} }
\expandafter\def\csname SRef@cdp:eq:C.TC.5\endcsname{\href{S6_SERIES_GRIND_II.pdf\#nameddest=tag-012J}{II.{C.TC.5} (p.~577)} }
\expandafter\def\csname SRef@cdp:eq:C.TC.6\endcsname{\href{S6_SERIES_GRIND_II.pdf\#nameddest=tag-012K}{II.{C.TC.6} (p.~577)} }
\expandafter\def\csname SRef@cdp:eq:C.TC.7\endcsname{\href{S6_SERIES_GRIND_II.pdf\#nameddest=tag-012L}{II.{C.TC.7} (p.~577)} }
\expandafter\def\csname SRef@cdp:eq:C.TC.8\endcsname{\href{S6_SERIES_GRIND_II.pdf\#nameddest=tag-012M}{II.{C.TC.8} (p.~577)} }
\expandafter\def\csname SRef@cdp:eq:C.TC.9\endcsname{\href{S6_SERIES_GRIND_II.pdf\#nameddest=tag-012N}{II.{C.TC.9} (p.~578)} }
\expandafter\def\csname SRef@cdp:eq:C.CA.3\endcsname{\href{S6_SERIES_GRIND_II.pdf\#nameddest=tag-00BG}{II.{C.CA.3} (p.~429)} }
\expandafter\def\csname SRef@cdp:eq:C.TC.10\endcsname{\href{S6_SERIES_GRIND_II.pdf\#nameddest=tag-012D}{II.{C.TC.10} (p.~578)} }
\expandafter\def\csname SRef@cdp:eq:C.TC.11\endcsname{\href{S6_SERIES_GRIND_II.pdf\#nameddest=tag-012E}{II.{C.TC.11} (p.~578)} }
\expandafter\def\csname SRef@cdp:eq:C.TC.12\endcsname{\href{S6_SERIES_GRIND_II.pdf\#nameddest=tag-012F}{II.{C.TC.12} (p.~578)} }
\expandafter\def\csname SRef@cdp:eq:C.ER.1\endcsname{\href{S6_SERIES_GRIND_II.pdf\#nameddest=tag-00IN}{II.{C.ER.1} (p.~579)} }
\expandafter\def\csname SRef@cdp:eq:C.ER.2\endcsname{\href{S6_SERIES_GRIND_II.pdf\#nameddest=tag-00IX}{II.{C.ER.2} (p.~579)} }
\expandafter\def\csname SRef@cdp:eq:C.ER.3\endcsname{\href{S6_SERIES_GRIND_II.pdf\#nameddest=tag-00IY}{II.{C.ER.3} (p.~579)} }
\expandafter\def\csname SRef@cdp:eq:C.ER.4\endcsname{\href{S6_SERIES_GRIND_II.pdf\#nameddest=tag-00IZ}{II.{C.ER.4} (p.~579)} }
\expandafter\def\csname SRef@cdp:eq:C.ER.5\endcsname{\href{S6_SERIES_GRIND_II.pdf\#nameddest=tag-00J0}{II.{C.ER.5} (p.~579)} }
\expandafter\def\csname SRef@cdp:eq:C.ER.6\endcsname{\href{S6_SERIES_GRIND_II.pdf\#nameddest=tag-00J1}{II.{C.ER.6} (p.~579)} }
\expandafter\def\csname SRef@cdp:eq:C.ER.7\endcsname{\href{S6_SERIES_GRIND_II.pdf\#nameddest=tag-00J2}{II.{C.ER.7} (p.~579)} }
\expandafter\def\csname SRef@cdp:eq:C.CA.4\endcsname{\href{S6_SERIES_GRIND_II.pdf\#nameddest=tag-00BH}{II.{C.CA.4} (p.~429)} }
\expandafter\def\csname SRef@cdp:eq:C.ER.8\endcsname{\href{S6_SERIES_GRIND_II.pdf\#nameddest=tag-00J3}{II.{C.ER.8} (p.~580)} }
\expandafter\def\csname SRef@cdp:eq:C.ER.9\endcsname{\href{S6_SERIES_GRIND_II.pdf\#nameddest=tag-00J4}{II.{C.ER.9} (p.~580)} }
\expandafter\def\csname SRef@cdp:eq:C.ER.10\endcsname{\href{S6_SERIES_GRIND_II.pdf\#nameddest=tag-00IO}{II.{C.ER.10} (p.~580)} }
\expandafter\def\csname SRef@cdp:eq:C.ER.11\endcsname{\href{S6_SERIES_GRIND_II.pdf\#nameddest=tag-00IP}{II.{C.ER.11} (p.~580)} }
\expandafter\def\csname SRef@cdp:eq:C.ER.12\endcsname{\href{S6_SERIES_GRIND_II.pdf\#nameddest=tag-00IQ}{II.{C.ER.12} (p.~580)} }
\expandafter\def\csname SRef@cdp:eq:C.ER.13\endcsname{\href{S6_SERIES_GRIND_II.pdf\#nameddest=tag-00IR}{II.{C.ER.13} (p.~580)} }
\expandafter\def\csname SRef@cdp:eq:C.ER.14\endcsname{\href{S6_SERIES_GRIND_II.pdf\#nameddest=tag-00IS}{II.{C.ER.14} (p.~580)} }
\expandafter\def\csname SRef@cdp:eq:C.ER.15\endcsname{\href{S6_SERIES_GRIND_II.pdf\#nameddest=tag-00IT}{II.{C.ER.15} (p.~581)} }
\expandafter\def\csname SRef@cdp:eq:C.ER.16\endcsname{\href{S6_SERIES_GRIND_II.pdf\#nameddest=tag-00IU}{II.{C.ER.16} (p.~581)} }
\expandafter\def\csname SRef@cdp:eq:C.CA.5\endcsname{\href{S6_SERIES_GRIND_II.pdf\#nameddest=tag-00BI}{II.{C.CA.5} (p.~429)} }
\expandafter\def\csname SRef@cdp:eq:C.ER.17\endcsname{\href{S6_SERIES_GRIND_II.pdf\#nameddest=tag-00IV}{II.{C.ER.17} (p.~582)} }
\expandafter\def\csname SRef@cdp:eq:C.ER.18\endcsname{\href{S6_SERIES_GRIND_II.pdf\#nameddest=tag-00IW}{II.{C.ER.18} (p.~582)} }
\expandafter\def\csname SRef@cdp:eq:C.CP.1\endcsname{\href{S6_SERIES_GRIND_II.pdf\#nameddest=tag-00C9}{II.{C.CP.1} (p.~583)} }
\expandafter\def\csname SRef@cdp:eq:C.CP.2\endcsname{\href{S6_SERIES_GRIND_II.pdf\#nameddest=tag-00CB}{II.{C.CP.2} (p.~583)} }
\expandafter\def\csname SRef@cdp:eq:C.CP.3\endcsname{\href{S6_SERIES_GRIND_II.pdf\#nameddest=tag-00CC}{II.{C.CP.3} (p.~583)} }
\expandafter\def\csname SRef@cdp:eq:C.CP.4\endcsname{\href{S6_SERIES_GRIND_II.pdf\#nameddest=tag-00CD}{II.{C.CP.4} (p.~583)} }
\expandafter\def\csname SRef@cdp:eq:C.CP.5\endcsname{\href{S6_SERIES_GRIND_II.pdf\#nameddest=tag-00CE}{II.{C.CP.5} (p.~583)} }
\expandafter\def\csname SRef@cdp:eq:C.CP.6\endcsname{\href{S6_SERIES_GRIND_II.pdf\#nameddest=tag-00CF}{II.{C.CP.6} (p.~584)} }
\expandafter\def\csname SRef@cdp:eq:C.CP.7\endcsname{\href{S6_SERIES_GRIND_II.pdf\#nameddest=tag-00CG}{II.{C.CP.7} (p.~584)} }
\expandafter\def\csname SRef@cdp:eq:C.CA.6\endcsname{\href{S6_SERIES_GRIND_II.pdf\#nameddest=tag-00BJ}{II.{C.CA.6} (p.~429)} }
\expandafter\def\csname SRef@cdp:eq:C.CP.8\endcsname{\href{S6_SERIES_GRIND_II.pdf\#nameddest=tag-00CH}{II.{C.CP.8} (p.~584)} }
\expandafter\def\csname SRef@cdp:eq:C.CP.9\endcsname{\href{S6_SERIES_GRIND_II.pdf\#nameddest=tag-00CI}{II.{C.CP.9} (p.~584)} }
\expandafter\def\csname SRef@cdp:eq:C.CP.10\endcsname{\href{S6_SERIES_GRIND_II.pdf\#nameddest=tag-00CA}{II.{C.CP.10} (p.~585)} }
\expandafter\def\csname SRef@cdp:eq:C.KM.1\endcsname{\href{S6_SERIES_GRIND_II.pdf\#nameddest=tag-00PI}{II.{C.KM.1} (p.~586)} }
\expandafter\def\csname SRef@cdp:eq:C.KM.2\endcsname{\href{S6_SERIES_GRIND_II.pdf\#nameddest=tag-00PL}{II.{C.KM.2} (p.~586)} }
\expandafter\def\csname SRef@cdp:eq:C.KM.S1\endcsname{\href{S6_SERIES_GRIND_II.pdf\#nameddest=tag-00PT}{II.{C.KM.S1} (p.~586)} }
\expandafter\def\csname SRef@cdp:eq:C.KM.3\endcsname{\href{S6_SERIES_GRIND_II.pdf\#nameddest=tag-00PM}{II.{C.KM.3} (p.~587)} }
\expandafter\def\csname SRef@cdp:eq:C.KM.4\endcsname{\href{S6_SERIES_GRIND_II.pdf\#nameddest=tag-00PN}{II.{C.KM.4} (p.~587)} }
\expandafter\def\csname SRef@cdp:eq:C.KM.5\endcsname{\href{S6_SERIES_GRIND_II.pdf\#nameddest=tag-00PO}{II.{C.KM.5} (p.~587)} }
\expandafter\def\csname SRef@cdp:eq:C.KM.6\endcsname{\href{S6_SERIES_GRIND_II.pdf\#nameddest=tag-00PP}{II.{C.KM.6} (p.~587)} }
\expandafter\def\csname SRef@cdp:eq:C.CA.7\endcsname{\href{S6_SERIES_GRIND_II.pdf\#nameddest=tag-00BK}{II.{C.CA.7} (p.~430)} }
\expandafter\def\csname SRef@cdp:eq:C.KM.7\endcsname{\href{S6_SERIES_GRIND_II.pdf\#nameddest=tag-00PQ}{II.{C.KM.7} (p.~587)} }
\expandafter\def\csname SRef@cdp:eq:C.KM.8\endcsname{\href{S6_SERIES_GRIND_II.pdf\#nameddest=tag-00PR}{II.{C.KM.8} (p.~588)} }
\expandafter\def\csname SRef@cdp:eq:C.KM.9\endcsname{\href{S6_SERIES_GRIND_II.pdf\#nameddest=tag-00PS}{II.{C.KM.9} (p.~588)} }
\expandafter\def\csname SRef@cdp:eq:C.KM.10\endcsname{\href{S6_SERIES_GRIND_II.pdf\#nameddest=tag-00PJ}{II.{C.KM.10} (p.~588)} }
\expandafter\def\csname SRef@cdp:eq:C.KM.S2\endcsname{\href{S6_SERIES_GRIND_II.pdf\#nameddest=tag-00PU}{II.{C.KM.S2} (p.~589)} }
\expandafter\def\csname SRef@cdp:eq:C.KM.S3\endcsname{\href{S6_SERIES_GRIND_II.pdf\#nameddest=tag-00PV}{II.{C.KM.S3} (p.~589)} }
\expandafter\def\csname SRef@cdp:eq:C.KM.11\endcsname{\href{S6_SERIES_GRIND_II.pdf\#nameddest=tag-00PK}{II.{C.KM.11} (p.~589)} }
\expandafter\def\csname SRef@cdp:eq:C.PE.1\endcsname{\href{S6_SERIES_GRIND_II.pdf\#nameddest=tag-00U2}{II.{C.PE.1} (p.~590)} }
\expandafter\def\csname SRef@cdp:eq:C.PE.2\endcsname{\href{S6_SERIES_GRIND_II.pdf\#nameddest=tag-00UD}{II.{C.PE.2} (p.~590)} }
\expandafter\def\csname SRef@cdp:eq:C.PE.3\endcsname{\href{S6_SERIES_GRIND_II.pdf\#nameddest=tag-00UG}{II.{C.PE.3} (p.~590)} }
\expandafter\def\csname SRef@cdp:eq:C.CA.8\endcsname{\href{S6_SERIES_GRIND_II.pdf\#nameddest=tag-00BL}{II.{C.CA.8} (p.~430)} }
\expandafter\def\csname SRef@cdp:eq:C.PE.4\endcsname{\href{S6_SERIES_GRIND_II.pdf\#nameddest=tag-00UH}{II.{C.PE.4} (p.~590)} }
\expandafter\def\csname SRef@cdp:eq:C.PE.5\endcsname{\href{S6_SERIES_GRIND_II.pdf\#nameddest=tag-00UI}{II.{C.PE.5} (p.~590)} }
\expandafter\def\csname SRef@cdp:eq:C.PE.6\endcsname{\href{S6_SERIES_GRIND_II.pdf\#nameddest=tag-00UJ}{II.{C.PE.6} (p.~591)} }
\expandafter\def\csname SRef@cdp:eq:C.PE.7\endcsname{\href{S6_SERIES_GRIND_II.pdf\#nameddest=tag-00UK}{II.{C.PE.7} (p.~591)} }
\expandafter\def\csname SRef@cdp:eq:C.PE.8\endcsname{\href{S6_SERIES_GRIND_II.pdf\#nameddest=tag-00UL}{II.{C.PE.8} (p.~591)} }
\expandafter\def\csname SRef@cdp:eq:C.PE.9\endcsname{\href{S6_SERIES_GRIND_II.pdf\#nameddest=tag-00UM}{II.{C.PE.9} (p.~591)} }
\expandafter\def\csname SRef@cdp:eq:C.PE.10\endcsname{\href{S6_SERIES_GRIND_II.pdf\#nameddest=tag-00U3}{II.{C.PE.10} (p.~591)} }
\expandafter\def\csname SRef@cdp:eq:C.PE.11\endcsname{\href{S6_SERIES_GRIND_II.pdf\#nameddest=tag-00U4}{II.{C.PE.11} (p.~591)} }
\expandafter\def\csname SRef@cdp:eq:C.PE.12\endcsname{\href{S6_SERIES_GRIND_II.pdf\#nameddest=tag-00U5}{II.{C.PE.12} (p.~591)} }
\expandafter\def\csname SRef@cdp:eq:C.PE.13\endcsname{\href{S6_SERIES_GRIND_II.pdf\#nameddest=tag-00U6}{II.{C.PE.13} (p.~592)} }
\expandafter\def\csname SRef@cdp:eq:C.CA.9\endcsname{\href{S6_SERIES_GRIND_II.pdf\#nameddest=tag-00BM}{II.{C.CA.9} (p.~430)} }
\expandafter\def\csname SRef@cdp:eq:C.PE.14\endcsname{\href{S6_SERIES_GRIND_II.pdf\#nameddest=tag-00U7}{II.{C.PE.14} (p.~592)} }
\expandafter\def\csname SRef@cdp:eq:C.PE.15\endcsname{\href{S6_SERIES_GRIND_II.pdf\#nameddest=tag-00U8}{II.{C.PE.15} (p.~592)} }
\expandafter\def\csname SRef@cdp:eq:C.PE.16\endcsname{\href{S6_SERIES_GRIND_II.pdf\#nameddest=tag-00U9}{II.{C.PE.16} (p.~592)} }
\expandafter\def\csname SRef@cdp:eq:C.PE.17\endcsname{\href{S6_SERIES_GRIND_II.pdf\#nameddest=tag-00UA}{II.{C.PE.17} (p.~592)} }
\expandafter\def\csname SRef@cdp:eq:C.PE.18\endcsname{\href{S6_SERIES_GRIND_II.pdf\#nameddest=tag-00UB}{II.{C.PE.18} (p.~593)} }
\expandafter\def\csname SRef@cdp:eq:C.PE.19\endcsname{\href{S6_SERIES_GRIND_II.pdf\#nameddest=tag-00UC}{II.{C.PE.19} (p.~593)} }
\expandafter\def\csname SRef@cdp:eq:C.PE.20\endcsname{\href{S6_SERIES_GRIND_II.pdf\#nameddest=tag-00UE}{II.{C.PE.20} (p.~593)} }
\expandafter\def\csname SRef@cdp:eq:C.PE.21\endcsname{\href{S6_SERIES_GRIND_II.pdf\#nameddest=tag-00UF}{II.{C.PE.21} (p.~593)} }
\expandafter\def\csname SRef@cdp:eq:C.SKD.1\endcsname{\href{S6_SERIES_GRIND_II.pdf\#nameddest=tag-011N}{II.{C.SKD.1} (p.~594)} }
\expandafter\def\csname SRef@cdp:eq:C.SKD.2\endcsname{\href{S6_SERIES_GRIND_II.pdf\#nameddest=tag-011U}{II.{C.SKD.2} (p.~594)} }
\expandafter\def\csname SRef@cdp:eq:C.CA.10\endcsname{\href{S6_SERIES_GRIND_II.pdf\#nameddest=tag-00BD}{II.{C.CA.10} (p.~430)} }
\expandafter\def\csname SRef@cdp:eq:C.SKD.3\endcsname{\href{S6_SERIES_GRIND_II.pdf\#nameddest=tag-011V}{II.{C.SKD.3} (p.~594)} }
\expandafter\def\csname SRef@cdp:eq:C.SKD.4\endcsname{\href{S6_SERIES_GRIND_II.pdf\#nameddest=tag-011W}{II.{C.SKD.4} (p.~595)} }
\expandafter\def\csname SRef@cdp:eq:C.SKD.5\endcsname{\href{S6_SERIES_GRIND_II.pdf\#nameddest=tag-011X}{II.{C.SKD.5} (p.~595)} }
\expandafter\def\csname SRef@cdp:eq:C.SKD.6\endcsname{\href{S6_SERIES_GRIND_II.pdf\#nameddest=tag-011Y}{II.{C.SKD.6} (p.~595)} }
\expandafter\def\csname SRef@cdp:eq:C.SKD.7\endcsname{\href{S6_SERIES_GRIND_II.pdf\#nameddest=tag-011Z}{II.{C.SKD.7} (p.~595)} }
\expandafter\def\csname SRef@cdp:eq:C.SKD.8\endcsname{\href{S6_SERIES_GRIND_II.pdf\#nameddest=tag-0120}{II.{C.SKD.8} (p.~595)} }
\expandafter\def\csname SRef@cdp:eq:C.SKD.9\endcsname{\href{S6_SERIES_GRIND_II.pdf\#nameddest=tag-0121}{II.{C.SKD.9} (p.~596)} }
\expandafter\def\csname SRef@cdp:eq:C.SKD.10\endcsname{\href{S6_SERIES_GRIND_II.pdf\#nameddest=tag-011O}{II.{C.SKD.10} (p.~596)} }
\expandafter\def\csname SRef@cdp:eq:C.SKD.11\endcsname{\href{S6_SERIES_GRIND_II.pdf\#nameddest=tag-011P}{II.{C.SKD.11} (p.~596)} }
\expandafter\def\csname SRef@cdp:eq:C.SKD.14\endcsname{\href{S6_SERIES_GRIND_II.pdf\#nameddest=tag-011S}{II.{C.SKD.14} (p.~596)} }
\expandafter\def\csname SRef@cdp:eq:C.CA.11\endcsname{\href{S6_SERIES_GRIND_II.pdf\#nameddest=tag-00BE}{II.{C.CA.11} (p.~430)} }
\expandafter\def\csname SRef@cdp:eq:C.SKD.15\endcsname{\href{S6_SERIES_GRIND_II.pdf\#nameddest=tag-011T}{II.{C.SKD.15} (p.~597)} }
\expandafter\def\csname SRef@cdp:eq:C.SKD.12\endcsname{\href{S6_SERIES_GRIND_II.pdf\#nameddest=tag-011Q}{II.{C.SKD.12} (p.~597)} }
\expandafter\def\csname SRef@cdp:eq:C.SKD.13\endcsname{\href{S6_SERIES_GRIND_II.pdf\#nameddest=tag-011R}{II.{C.SKD.13} (p.~597)} }
\expandafter\def\csname SRef@cdp:eq:C.EN.1\endcsname{\href{S6_SERIES_GRIND_II.pdf\#nameddest=tag-00G9}{II.{C.EN.1} (p.~598)} }
\expandafter\def\csname SRef@cdp:eq:C.EN.18\endcsname{\href{S6_SERIES_GRIND_II.pdf\#nameddest=tag-00GI}{II.{C.EN.18} (p.~598)} }
\expandafter\def\csname SRef@cdp:eq:C.EN.2\endcsname{\href{S6_SERIES_GRIND_II.pdf\#nameddest=tag-00GJ}{II.{C.EN.2} (p.~598)} }
\expandafter\def\csname SRef@cdp:eq:C.EN.3\endcsname{\href{S6_SERIES_GRIND_II.pdf\#nameddest=tag-00GK}{II.{C.EN.3} (p.~598)} }
\expandafter\def\csname SRef@cdp:eq:C.EN.4\endcsname{\href{S6_SERIES_GRIND_II.pdf\#nameddest=tag-00GL}{II.{C.EN.4} (p.~598)} }
\expandafter\def\csname SRef@cdp:eq:C.EN.5\endcsname{\href{S6_SERIES_GRIND_II.pdf\#nameddest=tag-00GM}{II.{C.EN.5} (p.~598)} }
\expandafter\def\csname SRef@cdp:eq:C.EN.6\endcsname{\href{S6_SERIES_GRIND_II.pdf\#nameddest=tag-00GN}{II.{C.EN.6} (p.~598)} }
\expandafter\def\csname SRef@cdp:eq:C.HQ.1\endcsname{\href{S6_SERIES_GRIND_II.pdf\#nameddest=tag-00MQ}{II.{C.HQ.1} (p.~431)} }
\expandafter\def\csname SRef@cdp:eq:C.EN.7\endcsname{\href{S6_SERIES_GRIND_II.pdf\#nameddest=tag-00GO}{II.{C.EN.7} (p.~599)} }
\expandafter\def\csname SRef@cdp:eq:C.EN.8\endcsname{\href{S6_SERIES_GRIND_II.pdf\#nameddest=tag-00GP}{II.{C.EN.8} (p.~599)} }
\expandafter\def\csname SRef@cdp:eq:C.EN.9\endcsname{\href{S6_SERIES_GRIND_II.pdf\#nameddest=tag-00GQ}{II.{C.EN.9} (p.~599)} }
\expandafter\def\csname SRef@cdp:eq:C.EN.10\endcsname{\href{S6_SERIES_GRIND_II.pdf\#nameddest=tag-00GA}{II.{C.EN.10} (p.~599)} }
\expandafter\def\csname SRef@cdp:eq:C.EN.11\endcsname{\href{S6_SERIES_GRIND_II.pdf\#nameddest=tag-00GB}{II.{C.EN.11} (p.~599)} }
\expandafter\def\csname SRef@cdp:eq:C.EN.12\endcsname{\href{S6_SERIES_GRIND_II.pdf\#nameddest=tag-00GC}{II.{C.EN.12} (p.~599)} }
\expandafter\def\csname SRef@cdp:eq:C.EN.13\endcsname{\href{S6_SERIES_GRIND_II.pdf\#nameddest=tag-00GD}{II.{C.EN.13} (p.~599)} }
\expandafter\def\csname SRef@cdp:eq:C.EN.14\endcsname{\href{S6_SERIES_GRIND_II.pdf\#nameddest=tag-00GE}{II.{C.EN.14} (p.~600)} }
\expandafter\def\csname SRef@cdp:eq:C.EN.15\endcsname{\href{S6_SERIES_GRIND_II.pdf\#nameddest=tag-00GF}{II.{C.EN.15} (p.~600)} }
\expandafter\def\csname SRef@cdp:eq:C.EN.16\endcsname{\href{S6_SERIES_GRIND_II.pdf\#nameddest=tag-00GG}{II.{C.EN.16} (p.~600)} }
\expandafter\def\csname SRef@cdp:eq:C.HQ.2\endcsname{\href{S6_SERIES_GRIND_II.pdf\#nameddest=tag-00MS}{II.{C.HQ.2} (p.~431)} }
\expandafter\def\csname SRef@cdp:eq:C.EN.17\endcsname{\href{S6_SERIES_GRIND_II.pdf\#nameddest=tag-00GH}{II.{C.EN.17} (p.~600)} }
\expandafter\def\csname SRef@cdp:eq:C.EN.S1\endcsname{\href{S6_SERIES_GRIND_II.pdf\#nameddest=tag-00H7}{II.{C.EN.S1} (p.~601)} }
\expandafter\def\csname SRef@cdp:eq:C.EN.S2\endcsname{\href{S6_SERIES_GRIND_II.pdf\#nameddest=tag-00H8}{II.{C.EN.S2} (p.~601)} }
\expandafter\def\csname SRef@cdp:eq:C.EN.S3\endcsname{\href{S6_SERIES_GRIND_II.pdf\#nameddest=tag-00H9}{II.{C.EN.S3} (p.~601)} }
\expandafter\def\csname SRef@cdp:eq:C.EN.S4\endcsname{\href{S6_SERIES_GRIND_II.pdf\#nameddest=tag-00HA}{II.{C.EN.S4} (p.~602)} }
\expandafter\def\csname SRef@cdp:eq:C.EN.S5\endcsname{\href{S6_SERIES_GRIND_II.pdf\#nameddest=tag-00HB}{II.{C.EN.S5} (p.~602)} }
\expandafter\def\csname SRef@cdp:eq:C.EN.S6\endcsname{\href{S6_SERIES_GRIND_II.pdf\#nameddest=tag-00HC}{II.{C.EN.S6} (p.~602)} }
\expandafter\def\csname SRef@cdp:eq:C.EN.S7\endcsname{\href{S6_SERIES_GRIND_II.pdf\#nameddest=tag-00HD}{II.{C.EN.S7} (p.~603)} }
\expandafter\def\csname SRef@cdp:eq:C.EN.R1\endcsname{\href{S6_SERIES_GRIND_II.pdf\#nameddest=tag-00GV}{II.{C.EN.R1} (p.~603)} }
\expandafter\def\csname SRef@cdp:eq:C.EN.R2\endcsname{\href{S6_SERIES_GRIND_II.pdf\#nameddest=tag-00GZ}{II.{C.EN.R2} (p.~603)} }
\expandafter\def\csname SRef@cdp:eq:C.HQ.3\endcsname{\href{S6_SERIES_GRIND_II.pdf\#nameddest=tag-00MT}{II.{C.HQ.3} (p.~432)} }
\expandafter\def\csname SRef@cdp:eq:C.EN.R3\endcsname{\href{S6_SERIES_GRIND_II.pdf\#nameddest=tag-00H0}{II.{C.EN.R3} (p.~603)} }
\expandafter\def\csname SRef@cdp:eq:C.EN.R4\endcsname{\href{S6_SERIES_GRIND_II.pdf\#nameddest=tag-00H1}{II.{C.EN.R4} (p.~603)} }
\expandafter\def\csname SRef@cdp:eq:C.EN.R5\endcsname{\href{S6_SERIES_GRIND_II.pdf\#nameddest=tag-00H2}{II.{C.EN.R5} (p.~603)} }
\expandafter\def\csname SRef@cdp:eq:C.EN.R6\endcsname{\href{S6_SERIES_GRIND_II.pdf\#nameddest=tag-00H3}{II.{C.EN.R6} (p.~604)} }
\expandafter\def\csname SRef@cdp:eq:C.EN.R7\endcsname{\href{S6_SERIES_GRIND_II.pdf\#nameddest=tag-00H4}{II.{C.EN.R7} (p.~604)} }
\expandafter\def\csname SRef@cdp:eq:C.EN.R8\endcsname{\href{S6_SERIES_GRIND_II.pdf\#nameddest=tag-00H5}{II.{C.EN.R8} (p.~604)} }
\expandafter\def\csname SRef@cdp:eq:C.EN.R9\endcsname{\href{S6_SERIES_GRIND_II.pdf\#nameddest=tag-00H6}{II.{C.EN.R9} (p.~604)} }
\expandafter\def\csname SRef@cdp:eq:C.EN.R10\endcsname{\href{S6_SERIES_GRIND_II.pdf\#nameddest=tag-00GW}{II.{C.EN.R10} (p.~604)} }
\expandafter\def\csname SRef@cdp:eq:C.EN.R11\endcsname{\href{S6_SERIES_GRIND_II.pdf\#nameddest=tag-00GX}{II.{C.EN.R11} (p.~605)} }
\expandafter\def\csname SRef@cdp:eq:C.EN.R12\endcsname{\href{S6_SERIES_GRIND_II.pdf\#nameddest=tag-00GY}{II.{C.EN.R12} (p.~605)} }
\expandafter\def\csname SRef@cdp:eq:C.HQ.4\endcsname{\href{S6_SERIES_GRIND_II.pdf\#nameddest=tag-00MU}{II.{C.HQ.4} (p.~432)} }
\expandafter\def\csname SRef@cdp:eq:C.EN.G1\endcsname{\href{S6_SERIES_GRIND_II.pdf\#nameddest=tag-00GR}{II.{C.EN.G1} (p.~606)} }
\expandafter\def\csname SRef@cdp:eq:C.EN.G2\endcsname{\href{S6_SERIES_GRIND_II.pdf\#nameddest=tag-00GS}{II.{C.EN.G2} (p.~606)} }
\expandafter\def\csname SRef@cdp:eq:C.EN.G3\endcsname{\href{S6_SERIES_GRIND_II.pdf\#nameddest=tag-00GT}{II.{C.EN.G3} (p.~606)} }
\expandafter\def\csname SRef@cdp:eq:C.EN.G4\endcsname{\href{S6_SERIES_GRIND_II.pdf\#nameddest=tag-00GU}{II.{C.EN.G4} (p.~606)} }
\expandafter\def\csname SRef@cdp:eq:C.ET.1\endcsname{\href{S6_SERIES_GRIND_II.pdf\#nameddest=tag-00J5}{II.{C.ET.1} (p.~607)} }
\expandafter\def\csname SRef@cdp:eq:C.ET.2\endcsname{\href{S6_SERIES_GRIND_II.pdf\#nameddest=tag-00J9}{II.{C.ET.2} (p.~607)} }
\expandafter\def\csname SRef@cdp:eq:C.ET.3\endcsname{\href{S6_SERIES_GRIND_II.pdf\#nameddest=tag-00JA}{II.{C.ET.3} (p.~607)} }
\expandafter\def\csname SRef@cdp:eq:C.ET.4\endcsname{\href{S6_SERIES_GRIND_II.pdf\#nameddest=tag-00JB}{II.{C.ET.4} (p.~607)} }
\expandafter\def\csname SRef@cdp:eq:C.ET.5\endcsname{\href{S6_SERIES_GRIND_II.pdf\#nameddest=tag-00JC}{II.{C.ET.5} (p.~607)} }
\expandafter\def\csname SRef@cdp:eq:C.ET.6\endcsname{\href{S6_SERIES_GRIND_II.pdf\#nameddest=tag-00JD}{II.{C.ET.6} (p.~608)} }
\expandafter\def\csname SRef@cdp:eq:C.HQ.5\endcsname{\href{S6_SERIES_GRIND_II.pdf\#nameddest=tag-00MV}{II.{C.HQ.5} (p.~432)} }
\expandafter\def\csname SRef@cdp:eq:C.ET.7\endcsname{\href{S6_SERIES_GRIND_II.pdf\#nameddest=tag-00JE}{II.{C.ET.7} (p.~608)} }
\expandafter\def\csname SRef@cdp:eq:C.ET.8\endcsname{\href{S6_SERIES_GRIND_II.pdf\#nameddest=tag-00JF}{II.{C.ET.8} (p.~608)} }
\expandafter\def\csname SRef@cdp:eq:C.ET.9\endcsname{\href{S6_SERIES_GRIND_II.pdf\#nameddest=tag-00JG}{II.{C.ET.9} (p.~608)} }
\expandafter\def\csname SRef@cdp:eq:C.ET.10\endcsname{\href{S6_SERIES_GRIND_II.pdf\#nameddest=tag-00J6}{II.{C.ET.10} (p.~609)} }
\expandafter\def\csname SRef@cdp:eq:C.ET.11\endcsname{\href{S6_SERIES_GRIND_II.pdf\#nameddest=tag-00J7}{II.{C.ET.11} (p.~609)} }
\expandafter\def\csname SRef@cdp:eq:C.ET.12\endcsname{\href{S6_SERIES_GRIND_II.pdf\#nameddest=tag-00J8}{II.{C.ET.12} (p.~609)} }
\expandafter\def\csname SRef@cdp:eq:C.SF.1\endcsname{\href{S6_SERIES_GRIND_II.pdf\#nameddest=tag-010R}{II.{C.SF.1} (p.~611)} }
\expandafter\def\csname SRef@cdp:eq:C.SF.2\endcsname{\href{S6_SERIES_GRIND_II.pdf\#nameddest=tag-010W}{II.{C.SF.2} (p.~611)} }
\expandafter\def\csname SRef@cdp:eq:C.SF.3\endcsname{\href{S6_SERIES_GRIND_II.pdf\#nameddest=tag-010X}{II.{C.SF.3} (p.~611)} }
\expandafter\def\csname SRef@cdp:eq:C.SF.4\endcsname{\href{S6_SERIES_GRIND_II.pdf\#nameddest=tag-010Y}{II.{C.SF.4} (p.~611)} }
\expandafter\def\csname SRef@cdp:eq:C.HQ.6\endcsname{\href{S6_SERIES_GRIND_II.pdf\#nameddest=tag-00MW}{II.{C.HQ.6} (p.~432)} }
\expandafter\def\csname SRef@cdp:eq:C.SF.5\endcsname{\href{S6_SERIES_GRIND_II.pdf\#nameddest=tag-010Z}{II.{C.SF.5} (p.~611)} }
\expandafter\def\csname SRef@cdp:eq:C.SF.6\endcsname{\href{S6_SERIES_GRIND_II.pdf\#nameddest=tag-0110}{II.{C.SF.6} (p.~612)} }
\expandafter\def\csname SRef@cdp:eq:C.SF.7\endcsname{\href{S6_SERIES_GRIND_II.pdf\#nameddest=tag-0111}{II.{C.SF.7} (p.~612)} }
\expandafter\def\csname SRef@cdp:eq:C.SF.8\endcsname{\href{S6_SERIES_GRIND_II.pdf\#nameddest=tag-0112}{II.{C.SF.8} (p.~612)} }
\expandafter\def\csname SRef@cdp:eq:C.SF.9\endcsname{\href{S6_SERIES_GRIND_II.pdf\#nameddest=tag-0113}{II.{C.SF.9} (p.~612)} }
\expandafter\def\csname SRef@cdp:eq:C.SF.10\endcsname{\href{S6_SERIES_GRIND_II.pdf\#nameddest=tag-010S}{II.{C.SF.10} (p.~612)} }
\expandafter\def\csname SRef@cdp:eq:C.SF.11\endcsname{\href{S6_SERIES_GRIND_II.pdf\#nameddest=tag-010T}{II.{C.SF.11} (p.~613)} }
\expandafter\def\csname SRef@cdp:eq:C.SF.12\endcsname{\href{S6_SERIES_GRIND_II.pdf\#nameddest=tag-010U}{II.{C.SF.12} (p.~613)} }
\expandafter\def\csname SRef@cdp:eq:C.SF.13\endcsname{\href{S6_SERIES_GRIND_II.pdf\#nameddest=tag-010V}{II.{C.SF.13} (p.~613)} }
\expandafter\def\csname SRef@cdp:eq:C.EO.1\endcsname{\href{S6_SERIES_GRIND_II.pdf\#nameddest=tag-00HE}{II.{C.EO.1} (p.~614)} }
\expandafter\def\csname SRef@cdp:eq:C.HQ.7\endcsname{\href{S6_SERIES_GRIND_II.pdf\#nameddest=tag-00MX}{II.{C.HQ.7} (p.~432)} }
\expandafter\def\csname SRef@cdp:eq:C.EO.2\endcsname{\href{S6_SERIES_GRIND_II.pdf\#nameddest=tag-00HM}{II.{C.EO.2} (p.~614)} }
\expandafter\def\csname SRef@cdp:eq:C.EO.3\endcsname{\href{S6_SERIES_GRIND_II.pdf\#nameddest=tag-00HN}{II.{C.EO.3} (p.~614)} }
\expandafter\def\csname SRef@cdp:eq:C.EO.4\endcsname{\href{S6_SERIES_GRIND_II.pdf\#nameddest=tag-00HO}{II.{C.EO.4} (p.~614)} }
\expandafter\def\csname SRef@cdp:eq:C.EO.5\endcsname{\href{S6_SERIES_GRIND_II.pdf\#nameddest=tag-00HP}{II.{C.EO.5} (p.~615)} }
\expandafter\def\csname SRef@cdp:eq:C.EO.6\endcsname{\href{S6_SERIES_GRIND_II.pdf\#nameddest=tag-00HQ}{II.{C.EO.6} (p.~615)} }
\expandafter\def\csname SRef@cdp:eq:C.EO.H1\endcsname{\href{S6_SERIES_GRIND_II.pdf\#nameddest=tag-00HU}{II.{C.EO.H1} (p.~615)} }
\expandafter\def\csname SRef@cdp:eq:C.EO.H2\endcsname{\href{S6_SERIES_GRIND_II.pdf\#nameddest=tag-00HV}{II.{C.EO.H2} (p.~616)} }
\expandafter\def\csname SRef@cdp:eq:C.EO.H3\endcsname{\href{S6_SERIES_GRIND_II.pdf\#nameddest=tag-00HW}{II.{C.EO.H3} (p.~616)} }
\expandafter\def\csname SRef@cdp:eq:C.EO.H4\endcsname{\href{S6_SERIES_GRIND_II.pdf\#nameddest=tag-00HX}{II.{C.EO.H4} (p.~616)} }
\expandafter\def\csname SRef@cdp:eq:C.EO.H5\endcsname{\href{S6_SERIES_GRIND_II.pdf\#nameddest=tag-00HY}{II.{C.EO.H5} (p.~616)} }
\expandafter\def\csname SRef@cdp:eq:C.HQ.8\endcsname{\href{S6_SERIES_GRIND_II.pdf\#nameddest=tag-00MY}{II.{C.HQ.8} (p.~432)} }
\expandafter\def\csname SRef@cdp:eq:C.EO.H6\endcsname{\href{S6_SERIES_GRIND_II.pdf\#nameddest=tag-00HZ}{II.{C.EO.H6} (p.~617)} }
\expandafter\def\csname SRef@cdp:eq:C.EO.H7\endcsname{\href{S6_SERIES_GRIND_II.pdf\#nameddest=tag-00I0}{II.{C.EO.H7} (p.~617)} }
\expandafter\def\csname SRef@cdp:eq:C.EO.7\endcsname{\href{S6_SERIES_GRIND_II.pdf\#nameddest=tag-00HR}{II.{C.EO.7} (p.~617)} }
\expandafter\def\csname SRef@cdp:eq:C.EO.8\endcsname{\href{S6_SERIES_GRIND_II.pdf\#nameddest=tag-00HS}{II.{C.EO.8} (p.~617)} }
\expandafter\def\csname SRef@cdp:eq:C.EO.9\endcsname{\href{S6_SERIES_GRIND_II.pdf\#nameddest=tag-00HT}{II.{C.EO.9} (p.~617)} }
\expandafter\def\csname SRef@cdp:eq:C.EO.10\endcsname{\href{S6_SERIES_GRIND_II.pdf\#nameddest=tag-00HF}{II.{C.EO.10} (p.~618)} }
\expandafter\def\csname SRef@cdp:eq:C.EO.11\endcsname{\href{S6_SERIES_GRIND_II.pdf\#nameddest=tag-00HG}{II.{C.EO.11} (p.~618)} }
\expandafter\def\csname SRef@cdp:eq:C.EO.12\endcsname{\href{S6_SERIES_GRIND_II.pdf\#nameddest=tag-00HH}{II.{C.EO.12} (p.~618)} }
\expandafter\def\csname SRef@cdp:eq:C.EO.13\endcsname{\href{S6_SERIES_GRIND_II.pdf\#nameddest=tag-00HI}{II.{C.EO.13} (p.~618)} }
\expandafter\def\csname SRef@cdp:eq:C.EO.14\endcsname{\href{S6_SERIES_GRIND_II.pdf\#nameddest=tag-00HJ}{II.{C.EO.14} (p.~619)} }
\expandafter\def\csname SRef@cdp:eq:C.HQ.9\endcsname{\href{S6_SERIES_GRIND_II.pdf\#nameddest=tag-00MZ}{II.{C.HQ.9} (p.~433)} }
\expandafter\def\csname SRef@cdp:eq:C.EO.15\endcsname{\href{S6_SERIES_GRIND_II.pdf\#nameddest=tag-00HK}{II.{C.EO.15} (p.~619)} }
\expandafter\def\csname SRef@cdp:eq:C.EO.16\endcsname{\href{S6_SERIES_GRIND_II.pdf\#nameddest=tag-00HL}{II.{C.EO.16} (p.~620)} }
\expandafter\def\csname SRef@cdp:eq:C.EH.1\endcsname{\href{S6_SERIES_GRIND_II.pdf\#nameddest=tag-00FQ}{II.{C.EH.1} (p.~620)} }
\expandafter\def\csname SRef@cdp:eq:C.EH.2\endcsname{\href{S6_SERIES_GRIND_II.pdf\#nameddest=tag-00FR}{II.{C.EH.2} (p.~620)} }
\expandafter\def\csname SRef@cdp:eq:C.EH.3\endcsname{\href{S6_SERIES_GRIND_II.pdf\#nameddest=tag-00FS}{II.{C.EH.3} (p.~621)} }
\expandafter\def\csname SRef@cdp:eq:C.EH.4\endcsname{\href{S6_SERIES_GRIND_II.pdf\#nameddest=tag-00FT}{II.{C.EH.4} (p.~621)} }
\expandafter\def\csname SRef@cdp:eq:C.EH.5\endcsname{\href{S6_SERIES_GRIND_II.pdf\#nameddest=tag-00FU}{II.{C.EH.5} (p.~621)} }
\expandafter\def\csname SRef@cdp:eq:C.EH.6\endcsname{\href{S6_SERIES_GRIND_II.pdf\#nameddest=tag-00FV}{II.{C.EH.6} (p.~621)} }
\expandafter\def\csname SRef@cdp:eq:C.EH.7\endcsname{\href{S6_SERIES_GRIND_II.pdf\#nameddest=tag-00FW}{II.{C.EH.7} (p.~621)} }
\expandafter\def\csname SRef@cdp:eq:C.EH.8\endcsname{\href{S6_SERIES_GRIND_II.pdf\#nameddest=tag-00FX}{II.{C.EH.8} (p.~621)} }
\expandafter\def\csname SRef@cdp:eq:C.HQ.10\endcsname{\href{S6_SERIES_GRIND_II.pdf\#nameddest=tag-00MR}{II.{C.HQ.10} (p.~433)} }
\expandafter\def\csname SRef@cdp:eq:C.DH.1\endcsname{\href{S6_SERIES_GRIND_II.pdf\#nameddest=tag-00F6}{II.{C.DH.1} (p.~622)} }
\expandafter\def\csname SRef@cdp:eq:C.DH.2\endcsname{\href{S6_SERIES_GRIND_II.pdf\#nameddest=tag-00FG}{II.{C.DH.2} (p.~622)} }
\expandafter\def\csname SRef@cdp:eq:C.DH.3\endcsname{\href{S6_SERIES_GRIND_II.pdf\#nameddest=tag-00FH}{II.{C.DH.3} (p.~622)} }
\expandafter\def\csname SRef@cdp:eq:C.DH.4\endcsname{\href{S6_SERIES_GRIND_II.pdf\#nameddest=tag-00FI}{II.{C.DH.4} (p.~622)} }
\expandafter\def\csname SRef@cdp:eq:C.DH.5\endcsname{\href{S6_SERIES_GRIND_II.pdf\#nameddest=tag-00FJ}{II.{C.DH.5} (p.~623)} }
\expandafter\def\csname SRef@cdp:eq:C.DH.6\endcsname{\href{S6_SERIES_GRIND_II.pdf\#nameddest=tag-00FK}{II.{C.DH.6} (p.~623)} }
\expandafter\def\csname SRef@cdp:eq:C.DH.7\endcsname{\href{S6_SERIES_GRIND_II.pdf\#nameddest=tag-00FL}{II.{C.DH.7} (p.~623)} }
\expandafter\def\csname SRef@cdp:eq:C.DH.8\endcsname{\href{S6_SERIES_GRIND_II.pdf\#nameddest=tag-00FM}{II.{C.DH.8} (p.~623)} }
\expandafter\def\csname SRef@cdp:eq:C.DH.9\endcsname{\href{S6_SERIES_GRIND_II.pdf\#nameddest=tag-00FN}{II.{C.DH.9} (p.~623)} }
\expandafter\def\csname SRef@cdp:eq:C.DH.10\endcsname{\href{S6_SERIES_GRIND_II.pdf\#nameddest=tag-00F7}{II.{C.DH.10} (p.~624)} }
\expandafter\def\csname SRef@cdp:eq:C.FD.1\endcsname{\href{S6_SERIES_GRIND_II.pdf\#nameddest=tag-00K2}{II.{C.FD.1} (p.~433)} }
\expandafter\def\csname SRef@cdp:eq:C.DH.11\endcsname{\href{S6_SERIES_GRIND_II.pdf\#nameddest=tag-00F8}{II.{C.DH.11} (p.~624)} }
\expandafter\def\csname SRef@cdp:eq:C.DH.12\endcsname{\href{S6_SERIES_GRIND_II.pdf\#nameddest=tag-00F9}{II.{C.DH.12} (p.~624)} }
\expandafter\def\csname SRef@cdp:eq:C.DH.13\endcsname{\href{S6_SERIES_GRIND_II.pdf\#nameddest=tag-00FA}{II.{C.DH.13} (p.~625)} }
\expandafter\def\csname SRef@cdp:eq:C.DH.14\endcsname{\href{S6_SERIES_GRIND_II.pdf\#nameddest=tag-00FB}{II.{C.DH.14} (p.~625)} }
\expandafter\def\csname SRef@cdp:eq:C.DH.14a\endcsname{\href{S6_SERIES_GRIND_II.pdf\#nameddest=tag-00FC}{II.{C.DH.14a} (p.~625)} }
\expandafter\def\csname SRef@cdp:eq:C.DH.14b\endcsname{\href{S6_SERIES_GRIND_II.pdf\#nameddest=tag-00FD}{II.{C.DH.14b} (p.~625)} }
\expandafter\def\csname SRef@cdp:eq:C.DH.15\endcsname{\href{S6_SERIES_GRIND_II.pdf\#nameddest=tag-00FE}{II.{C.DH.15} (p.~625)} }
\expandafter\def\csname SRef@cdp:eq:C.DH.16\endcsname{\href{S6_SERIES_GRIND_II.pdf\#nameddest=tag-00FF}{II.{C.DH.16} (p.~626)} }
\expandafter\def\csname SRef@cdp:eq:C.NH.1\endcsname{\href{S6_SERIES_GRIND_II.pdf\#nameddest=tag-00S9}{II.{C.NH.1} (p.~626)} }
\expandafter\def\csname SRef@cdp:eq:C.NH.2\endcsname{\href{S6_SERIES_GRIND_II.pdf\#nameddest=tag-00SJ}{II.{C.NH.2} (p.~627)} }
\expandafter\def\csname SRef@cdp:eq:C.FD.2\endcsname{\href{S6_SERIES_GRIND_II.pdf\#nameddest=tag-00K3}{II.{C.FD.2} (p.~434)} }
\expandafter\def\csname SRef@cdp:eq:C.NH.3\endcsname{\href{S6_SERIES_GRIND_II.pdf\#nameddest=tag-00SK}{II.{C.NH.3} (p.~627)} }
\expandafter\def\csname SRef@cdp:eq:C.NH.4\endcsname{\href{S6_SERIES_GRIND_II.pdf\#nameddest=tag-00SL}{II.{C.NH.4} (p.~627)} }
\expandafter\def\csname SRef@cdp:eq:C.NH.5\endcsname{\href{S6_SERIES_GRIND_II.pdf\#nameddest=tag-00SM}{II.{C.NH.5} (p.~627)} }
\expandafter\def\csname SRef@cdp:eq:C.NH.6\endcsname{\href{S6_SERIES_GRIND_II.pdf\#nameddest=tag-00SN}{II.{C.NH.6} (p.~627)} }
\expandafter\def\csname SRef@cdp:eq:C.NH.7\endcsname{\href{S6_SERIES_GRIND_II.pdf\#nameddest=tag-00SO}{II.{C.NH.7} (p.~627)} }
\expandafter\def\csname SRef@cdp:eq:C.NH.8\endcsname{\href{S6_SERIES_GRIND_II.pdf\#nameddest=tag-00SP}{II.{C.NH.8} (p.~628)} }
\expandafter\def\csname SRef@cdp:eq:C.NH.9\endcsname{\href{S6_SERIES_GRIND_II.pdf\#nameddest=tag-00SQ}{II.{C.NH.9} (p.~628)} }
\expandafter\def\csname SRef@cdp:eq:C.NH.10\endcsname{\href{S6_SERIES_GRIND_II.pdf\#nameddest=tag-00SA}{II.{C.NH.10} (p.~628)} }
\expandafter\def\csname SRef@cdp:eq:C.NH.11\endcsname{\href{S6_SERIES_GRIND_II.pdf\#nameddest=tag-00SB}{II.{C.NH.11} (p.~628)} }
\expandafter\def\csname SRef@cdp:eq:C.NH.12\endcsname{\href{S6_SERIES_GRIND_II.pdf\#nameddest=tag-00SC}{II.{C.NH.12} (p.~629)} }
\expandafter\def\csname SRef@cdp:eq:C.FD.3\endcsname{\href{S6_SERIES_GRIND_II.pdf\#nameddest=tag-00K4}{II.{C.FD.3} (p.~434)} }
\expandafter\def\csname SRef@cdp:eq:C.NH.13\endcsname{\href{S6_SERIES_GRIND_II.pdf\#nameddest=tag-00SD}{II.{C.NH.13} (p.~629)} }
\expandafter\def\csname SRef@cdp:eq:C.NH.14\endcsname{\href{S6_SERIES_GRIND_II.pdf\#nameddest=tag-00SE}{II.{C.NH.14} (p.~629)} }
\expandafter\def\csname SRef@cdp:eq:C.NH.15\endcsname{\href{S6_SERIES_GRIND_II.pdf\#nameddest=tag-00SF}{II.{C.NH.15} (p.~629)} }
\expandafter\def\csname SRef@cdp:eq:C.NH.16\endcsname{\href{S6_SERIES_GRIND_II.pdf\#nameddest=tag-00SG}{II.{C.NH.16} (p.~629)} }
\expandafter\def\csname SRef@cdp:eq:C.NH.17\endcsname{\href{S6_SERIES_GRIND_II.pdf\#nameddest=tag-00SH}{II.{C.NH.17} (p.~629)} }
\expandafter\def\csname SRef@cdp:eq:C.NH.18\endcsname{\href{S6_SERIES_GRIND_II.pdf\#nameddest=tag-00SI}{II.{C.NH.18} (p.~630)} }
\expandafter\def\csname SRef@cdp:eq:C.SH.1\endcsname{\href{S6_SERIES_GRIND_II.pdf\#nameddest=tag-0114}{II.{C.SH.1} (p.~630)} }
\expandafter\def\csname SRef@cdp:eq:C.SH.2\endcsname{\href{S6_SERIES_GRIND_II.pdf\#nameddest=tag-0119}{II.{C.SH.2} (p.~631)} }
\expandafter\def\csname SRef@cdp:eq:C.SH.3\endcsname{\href{S6_SERIES_GRIND_II.pdf\#nameddest=tag-011A}{II.{C.SH.3} (p.~631)} }
\expandafter\def\csname SRef@cdp:eq:C.SH.4\endcsname{\href{S6_SERIES_GRIND_II.pdf\#nameddest=tag-011B}{II.{C.SH.4} (p.~631)} }
\expandafter\def\csname SRef@cdp:eq:C.FD.4\endcsname{\href{S6_SERIES_GRIND_II.pdf\#nameddest=tag-00K5}{II.{C.FD.4} (p.~434)} }
\expandafter\def\csname SRef@cdp:eq:C.SH.5\endcsname{\href{S6_SERIES_GRIND_II.pdf\#nameddest=tag-011C}{II.{C.SH.5} (p.~632)} }
\expandafter\def\csname SRef@cdp:eq:C.SH.6\endcsname{\href{S6_SERIES_GRIND_II.pdf\#nameddest=tag-011D}{II.{C.SH.6} (p.~632)} }
\expandafter\def\csname SRef@cdp:eq:C.SH.7\endcsname{\href{S6_SERIES_GRIND_II.pdf\#nameddest=tag-011E}{II.{C.SH.7} (p.~632)} }
\expandafter\def\csname SRef@cdp:eq:C.SH.8\endcsname{\href{S6_SERIES_GRIND_II.pdf\#nameddest=tag-011F}{II.{C.SH.8} (p.~632)} }
\expandafter\def\csname SRef@cdp:eq:C.SH.9\endcsname{\href{S6_SERIES_GRIND_II.pdf\#nameddest=tag-011G}{II.{C.SH.9} (p.~632)} }
\expandafter\def\csname SRef@cdp:eq:C.SH.10\endcsname{\href{S6_SERIES_GRIND_II.pdf\#nameddest=tag-0115}{II.{C.SH.10} (p.~632)} }
\expandafter\def\csname SRef@cdp:eq:C.SH.11\endcsname{\href{S6_SERIES_GRIND_II.pdf\#nameddest=tag-0116}{II.{C.SH.11} (p.~633)} }
\expandafter\def\csname SRef@cdp:eq:C.SH.12\endcsname{\href{S6_SERIES_GRIND_II.pdf\#nameddest=tag-0117}{II.{C.SH.12} (p.~633)} }
\expandafter\def\csname SRef@cdp:eq:C.SH.13\endcsname{\href{S6_SERIES_GRIND_II.pdf\#nameddest=tag-0118}{II.{C.SH.13} (p.~633)} }
\expandafter\def\csname SRef@cdp:eq:C.RU.1\endcsname{\href{S6_SERIES_GRIND_II.pdf\#nameddest=tag-010C}{II.{C.RU.1} (p.~634)} }
\expandafter\def\csname SRef@cdp:eq:C.FD.5\endcsname{\href{S6_SERIES_GRIND_II.pdf\#nameddest=tag-00K6}{II.{C.FD.5} (p.~434)} }
\expandafter\def\csname SRef@cdp:eq:C.RU.2\endcsname{\href{S6_SERIES_GRIND_II.pdf\#nameddest=tag-010D}{II.{C.RU.2} (p.~634)} }
\expandafter\def\csname SRef@cdp:eq:C.RU.3\endcsname{\href{S6_SERIES_GRIND_II.pdf\#nameddest=tag-010E}{II.{C.RU.3} (p.~634)} }
\expandafter\def\csname SRef@cdp:eq:C.RU.4\endcsname{\href{S6_SERIES_GRIND_II.pdf\#nameddest=tag-010F}{II.{C.RU.4} (p.~634)} }
\expandafter\def\csname SRef@cdp:eq:C.RU.5\endcsname{\href{S6_SERIES_GRIND_II.pdf\#nameddest=tag-010G}{II.{C.RU.5} (p.~634)} }
\expandafter\def\csname SRef@cdp:eq:C.RU.6\endcsname{\href{S6_SERIES_GRIND_II.pdf\#nameddest=tag-010H}{II.{C.RU.6} (p.~635)} }
\expandafter\def\csname SRef@cdp:eq:C.RU.7\endcsname{\href{S6_SERIES_GRIND_II.pdf\#nameddest=tag-010I}{II.{C.RU.7} (p.~635)} }
\expandafter\def\csname SRef@cdp:eq:C.RU.8\endcsname{\href{S6_SERIES_GRIND_II.pdf\#nameddest=tag-010J}{II.{C.RU.8} (p.~635)} }
\expandafter\def\csname SRef@cdp:eq:C.RU.9\endcsname{\href{S6_SERIES_GRIND_II.pdf\#nameddest=tag-010K}{II.{C.RU.9} (p.~636)} }
\expandafter\def\csname SRef@cdp:eq:C.HA.1\endcsname{\href{S6_SERIES_GRIND_II.pdf\#nameddest=tag-00L9}{II.{C.HA.1} (p.~636)} }
\expandafter\def\csname SRef@cdp:eq:C.HA.2\endcsname{\href{S6_SERIES_GRIND_II.pdf\#nameddest=tag-00LA}{II.{C.HA.2} (p.~636)} }
\expandafter\def\csname SRef@cdp:eq:C.FD.6\endcsname{\href{S6_SERIES_GRIND_II.pdf\#nameddest=tag-00K7}{II.{C.FD.6} (p.~434)} }
\expandafter\def\csname SRef@cdp:eq:C.HA.3\endcsname{\href{S6_SERIES_GRIND_II.pdf\#nameddest=tag-00LB}{II.{C.HA.3} (p.~637)} }
\expandafter\def\csname SRef@cdp:eq:C.HA.4\endcsname{\href{S6_SERIES_GRIND_II.pdf\#nameddest=tag-00LC}{II.{C.HA.4} (p.~637)} }
\expandafter\def\csname SRef@cdp:eq:C.HA.5\endcsname{\href{S6_SERIES_GRIND_II.pdf\#nameddest=tag-00LD}{II.{C.HA.5} (p.~638)} }
\expandafter\def\csname SRef@cdp:eq:C.HA.6\endcsname{\href{S6_SERIES_GRIND_II.pdf\#nameddest=tag-00LE}{II.{C.HA.6} (p.~638)} }
\expandafter\def\csname SRef@cdp:eq:C.HC.1\endcsname{\href{S6_SERIES_GRIND_II.pdf\#nameddest=tag-00LO}{II.{C.HC.1} (p.~638)} }
\expandafter\def\csname SRef@cdp:eq:C.HC.2\endcsname{\href{S6_SERIES_GRIND_II.pdf\#nameddest=tag-00LR}{II.{C.HC.2} (p.~638)} }
\expandafter\def\csname SRef@cdp:eq:C.HC.3\endcsname{\href{S6_SERIES_GRIND_II.pdf\#nameddest=tag-00LS}{II.{C.HC.3} (p.~639)} }
\expandafter\def\csname SRef@cdp:eq:C.HC.4\endcsname{\href{S6_SERIES_GRIND_II.pdf\#nameddest=tag-00LT}{II.{C.HC.4} (p.~639)} }
\expandafter\def\csname SRef@cdp:eq:C.HC.5\endcsname{\href{S6_SERIES_GRIND_II.pdf\#nameddest=tag-00LU}{II.{C.HC.5} (p.~640)} }
\expandafter\def\csname SRef@cdp:eq:C.HC.6\endcsname{\href{S6_SERIES_GRIND_II.pdf\#nameddest=tag-00LV}{II.{C.HC.6} (p.~640)} }
\expandafter\def\csname SRef@cdp:eq:C.RF.5\endcsname{\href{S6_SERIES_GRIND_II.pdf\#nameddest=tag-00YY}{II.{C.RF.5} (p.~435)} }
\expandafter\def\csname SRef@cdp:eq:C.HC.7\endcsname{\href{S6_SERIES_GRIND_II.pdf\#nameddest=tag-00LW}{II.{C.HC.7} (p.~640)} }
\expandafter\def\csname SRef@cdp:eq:C.HC.8\endcsname{\href{S6_SERIES_GRIND_II.pdf\#nameddest=tag-00LX}{II.{C.HC.8} (p.~641)} }
\expandafter\def\csname SRef@cdp:eq:C.HC.9\endcsname{\href{S6_SERIES_GRIND_II.pdf\#nameddest=tag-00LY}{II.{C.HC.9} (p.~641)} }
\expandafter\def\csname SRef@cdp:eq:C.HC.10\endcsname{\href{S6_SERIES_GRIND_II.pdf\#nameddest=tag-00LP}{II.{C.HC.10} (p.~641)} }
\expandafter\def\csname SRef@cdp:eq:C.HC.11\endcsname{\href{S6_SERIES_GRIND_II.pdf\#nameddest=tag-00LQ}{II.{C.HC.11} (p.~642)} }
\expandafter\def\csname SRef@cdp:eq:C.HG.1\endcsname{\href{S6_SERIES_GRIND_II.pdf\#nameddest=tag-00LZ}{II.{C.HG.1} (p.~642)} }
\expandafter\def\csname SRef@cdp:eq:C.HG.2\endcsname{\href{S6_SERIES_GRIND_II.pdf\#nameddest=tag-00M8}{II.{C.HG.2} (p.~643)} }
\expandafter\def\csname SRef@cdp:eq:C.HG.3\endcsname{\href{S6_SERIES_GRIND_II.pdf\#nameddest=tag-00M9}{II.{C.HG.3} (p.~643)} }
\expandafter\def\csname SRef@cdp:eq:C.HG.4\endcsname{\href{S6_SERIES_GRIND_II.pdf\#nameddest=tag-00MA}{II.{C.HG.4} (p.~643)} }
\expandafter\def\csname SRef@cdp:eq:C.HG.5\endcsname{\href{S6_SERIES_GRIND_II.pdf\#nameddest=tag-00MB}{II.{C.HG.5} (p.~643)} }
\expandafter\def\csname SRef@cdp:eq:C.RF.1\endcsname{\href{S6_SERIES_GRIND_II.pdf\#nameddest=tag-00YU}{II.{C.RF.1} (p.~435)} }
\expandafter\def\csname SRef@cdp:eq:C.HG.6\endcsname{\href{S6_SERIES_GRIND_II.pdf\#nameddest=tag-00MC}{II.{C.HG.6} (p.~643)} }
\expandafter\def\csname SRef@cdp:eq:C.HG.7\endcsname{\href{S6_SERIES_GRIND_II.pdf\#nameddest=tag-00MD}{II.{C.HG.7} (p.~643)} }
\expandafter\def\csname SRef@cdp:eq:C.HG.8\endcsname{\href{S6_SERIES_GRIND_II.pdf\#nameddest=tag-00ME}{II.{C.HG.8} (p.~644)} }
\expandafter\def\csname SRef@cdp:eq:C.HG.9\endcsname{\href{S6_SERIES_GRIND_II.pdf\#nameddest=tag-00MF}{II.{C.HG.9} (p.~644)} }
\expandafter\def\csname SRef@cdp:eq:C.HG.10\endcsname{\href{S6_SERIES_GRIND_II.pdf\#nameddest=tag-00M0}{II.{C.HG.10} (p.~644)} }
\expandafter\def\csname SRef@cdp:eq:C.HG.11\endcsname{\href{S6_SERIES_GRIND_II.pdf\#nameddest=tag-00M1}{II.{C.HG.11} (p.~644)} }
\expandafter\def\csname SRef@cdp:eq:C.HG.12\endcsname{\href{S6_SERIES_GRIND_II.pdf\#nameddest=tag-00M2}{II.{C.HG.12} (p.~644)} }
\expandafter\def\csname SRef@cdp:eq:C.HG.13\endcsname{\href{S6_SERIES_GRIND_II.pdf\#nameddest=tag-00M3}{II.{C.HG.13} (p.~645)} }
\expandafter\def\csname SRef@cdp:eq:C.HG.14\endcsname{\href{S6_SERIES_GRIND_II.pdf\#nameddest=tag-00M4}{II.{C.HG.14} (p.~645)} }
\expandafter\def\csname SRef@cdp:eq:C.HG.15\endcsname{\href{S6_SERIES_GRIND_II.pdf\#nameddest=tag-00M5}{II.{C.HG.15} (p.~646)} }
\expandafter\def\csname SRef@cdp:eq:C.RF.2\endcsname{\href{S6_SERIES_GRIND_II.pdf\#nameddest=tag-00YV}{II.{C.RF.2} (p.~436)} }
\expandafter\def\csname SRef@cdp:eq:C.HG.16\endcsname{\href{S6_SERIES_GRIND_II.pdf\#nameddest=tag-00M6}{II.{C.HG.16} (p.~646)} }
\expandafter\def\csname SRef@cdp:eq:C.HG.17\endcsname{\href{S6_SERIES_GRIND_II.pdf\#nameddest=tag-00M7}{II.{C.HG.17} (p.~646)} }
\expandafter\def\csname SRef@cdp:eq:C.NS.1\endcsname{\href{S6_SERIES_GRIND_II.pdf\#nameddest=tag-00TA}{II.{C.NS.1} (p.~647)} }
\expandafter\def\csname SRef@cdp:eq:C.NS.2\endcsname{\href{S6_SERIES_GRIND_II.pdf\#nameddest=tag-00TL}{II.{C.NS.2} (p.~647)} }
\expandafter\def\csname SRef@cdp:eq:C.NS.3\endcsname{\href{S6_SERIES_GRIND_II.pdf\#nameddest=tag-00TN}{II.{C.NS.3} (p.~647)} }
\expandafter\def\csname SRef@cdp:eq:C.NS.4\endcsname{\href{S6_SERIES_GRIND_II.pdf\#nameddest=tag-00TO}{II.{C.NS.4} (p.~648)} }
\expandafter\def\csname SRef@cdp:eq:C.NS.5\endcsname{\href{S6_SERIES_GRIND_II.pdf\#nameddest=tag-00TP}{II.{C.NS.5} (p.~648)} }
\expandafter\def\csname SRef@cdp:eq:C.NS.6\endcsname{\href{S6_SERIES_GRIND_II.pdf\#nameddest=tag-00TQ}{II.{C.NS.6} (p.~648)} }
\expandafter\def\csname SRef@cdp:eq:C.NS.7\endcsname{\href{S6_SERIES_GRIND_II.pdf\#nameddest=tag-00TR}{II.{C.NS.7} (p.~648)} }
\expandafter\def\csname SRef@cdp:eq:C.NS.8\endcsname{\href{S6_SERIES_GRIND_II.pdf\#nameddest=tag-00TS}{II.{C.NS.8} (p.~649)} }
\expandafter\def\csname SRef@cdp:eq:C.RF.3\endcsname{\href{S6_SERIES_GRIND_II.pdf\#nameddest=tag-00YW}{II.{C.RF.3} (p.~436)} }
\expandafter\def\csname SRef@cdp:eq:C.NS.9\endcsname{\href{S6_SERIES_GRIND_II.pdf\#nameddest=tag-00TT}{II.{C.NS.9} (p.~649)} }
\expandafter\def\csname SRef@cdp:eq:C.NS.10\endcsname{\href{S6_SERIES_GRIND_II.pdf\#nameddest=tag-00TB}{II.{C.NS.10} (p.~649)} }
\expandafter\def\csname SRef@cdp:eq:C.NS.11\endcsname{\href{S6_SERIES_GRIND_II.pdf\#nameddest=tag-00TC}{II.{C.NS.11} (p.~649)} }
\expandafter\def\csname SRef@cdp:eq:C.NS.12\endcsname{\href{S6_SERIES_GRIND_II.pdf\#nameddest=tag-00TD}{II.{C.NS.12} (p.~649)} }
\expandafter\def\csname SRef@cdp:eq:C.NS.13\endcsname{\href{S6_SERIES_GRIND_II.pdf\#nameddest=tag-00TE}{II.{C.NS.13} (p.~650)} }
\expandafter\def\csname SRef@cdp:eq:C.NS.14\endcsname{\href{S6_SERIES_GRIND_II.pdf\#nameddest=tag-00TF}{II.{C.NS.14} (p.~650)} }
\expandafter\def\csname SRef@cdp:eq:C.NS.15\endcsname{\href{S6_SERIES_GRIND_II.pdf\#nameddest=tag-00TG}{II.{C.NS.15} (p.~650)} }
\expandafter\def\csname SRef@cdp:eq:C.NS.16\endcsname{\href{S6_SERIES_GRIND_II.pdf\#nameddest=tag-00TH}{II.{C.NS.16} (p.~650)} }
\expandafter\def\csname SRef@cdp:eq:C.NS.17\endcsname{\href{S6_SERIES_GRIND_II.pdf\#nameddest=tag-00TI}{II.{C.NS.17} (p.~650)} }
\expandafter\def\csname SRef@cdp:eq:C.NS.18\endcsname{\href{S6_SERIES_GRIND_II.pdf\#nameddest=tag-00TJ}{II.{C.NS.18} (p.~650)} }
\expandafter\def\csname SRef@cdp:eq:C.RF.4\endcsname{\href{S6_SERIES_GRIND_II.pdf\#nameddest=tag-00YX}{II.{C.RF.4} (p.~437)} }
\expandafter\def\csname SRef@cdp:eq:C.NS.19\endcsname{\href{S6_SERIES_GRIND_II.pdf\#nameddest=tag-00TK}{II.{C.NS.19} (p.~651)} }
\expandafter\def\csname SRef@cdp:eq:C.NS.20\endcsname{\href{S6_SERIES_GRIND_II.pdf\#nameddest=tag-00TM}{II.{C.NS.20} (p.~651)} }
\expandafter\def\csname SRef@cdp:eq:C.SN.1\endcsname{\href{S6_SERIES_GRIND_II.pdf\#nameddest=tag-0122}{II.{C.SN.1} (p.~652)} }
\expandafter\def\csname SRef@cdp:eq:C.SN.2\endcsname{\href{S6_SERIES_GRIND_II.pdf\#nameddest=tag-0124}{II.{C.SN.2} (p.~652)} }
\expandafter\def\csname SRef@cdp:eq:C.SN.3\endcsname{\href{S6_SERIES_GRIND_II.pdf\#nameddest=tag-0125}{II.{C.SN.3} (p.~652)} }
\expandafter\def\csname SRef@cdp:eq:C.SN.4\endcsname{\href{S6_SERIES_GRIND_II.pdf\#nameddest=tag-0126}{II.{C.SN.4} (p.~652)} }
\expandafter\def\csname SRef@cdp:eq:C.SN.5\endcsname{\href{S6_SERIES_GRIND_II.pdf\#nameddest=tag-0127}{II.{C.SN.5} (p.~653)} }
\expandafter\def\csname SRef@cdp:eq:C.SN.6\endcsname{\href{S6_SERIES_GRIND_II.pdf\#nameddest=tag-0128}{II.{C.SN.6} (p.~653)} }
\expandafter\def\csname SRef@cdp:eq:C.SN.7\endcsname{\href{S6_SERIES_GRIND_II.pdf\#nameddest=tag-0129}{II.{C.SN.7} (p.~653)} }
\expandafter\def\csname SRef@cdp:eq:C.SN.8\endcsname{\href{S6_SERIES_GRIND_II.pdf\#nameddest=tag-012A}{II.{C.SN.8} (p.~653)} }
\expandafter\def\csname SRef@cdp:eq:C.RF.9\endcsname{\href{S6_SERIES_GRIND_II.pdf\#nameddest=tag-00Z2}{II.{C.RF.9} (p.~437)} }
\expandafter\def\csname SRef@cdp:eq:C.SN.9\endcsname{\href{S6_SERIES_GRIND_II.pdf\#nameddest=tag-012B}{II.{C.SN.9} (p.~653)} }
\expandafter\def\csname SRef@cdp:eq:C.SN.10\endcsname{\href{S6_SERIES_GRIND_II.pdf\#nameddest=tag-0123}{II.{C.SN.10} (p.~654)} }
\expandafter\def\csname SRef@cdp:eq:C.MX.1\endcsname{\href{S6_SERIES_GRIND_II.pdf\#nameddest=tag-00RR}{II.{C.MX.1} (p.~654)} }
\expandafter\def\csname SRef@cdp:eq:C.MX.2\endcsname{\href{S6_SERIES_GRIND_II.pdf\#nameddest=tag-00RT}{II.{C.MX.2} (p.~654)} }
\expandafter\def\csname SRef@cdp:eq:C.MX.3\endcsname{\href{S6_SERIES_GRIND_II.pdf\#nameddest=tag-00RU}{II.{C.MX.3} (p.~655)} }
\expandafter\def\csname SRef@cdp:eq:C.MX.4\endcsname{\href{S6_SERIES_GRIND_II.pdf\#nameddest=tag-00RV}{II.{C.MX.4} (p.~655)} }
\expandafter\def\csname SRef@cdp:eq:C.MX.5\endcsname{\href{S6_SERIES_GRIND_II.pdf\#nameddest=tag-00RW}{II.{C.MX.5} (p.~656)} }
\expandafter\def\csname SRef@cdp:eq:C.MX.6\endcsname{\href{S6_SERIES_GRIND_II.pdf\#nameddest=tag-00RX}{II.{C.MX.6} (p.~656)} }
\expandafter\def\csname SRef@cdp:eq:C.MX.7\endcsname{\href{S6_SERIES_GRIND_II.pdf\#nameddest=tag-00RY}{II.{C.MX.7} (p.~656)} }
\expandafter\def\csname SRef@cdp:eq:C.MX.8\endcsname{\href{S6_SERIES_GRIND_II.pdf\#nameddest=tag-00RZ}{II.{C.MX.8} (p.~656)} }
\expandafter\def\csname SRef@cdp:eq:C.RF.6\endcsname{\href{S6_SERIES_GRIND_II.pdf\#nameddest=tag-00YZ}{II.{C.RF.6} (p.~438)} }
\expandafter\def\csname SRef@cdp:eq:C.MX.9\endcsname{\href{S6_SERIES_GRIND_II.pdf\#nameddest=tag-00S0}{II.{C.MX.9} (p.~656)} }
\expandafter\def\csname SRef@cdp:eq:C.MX.10\endcsname{\href{S6_SERIES_GRIND_II.pdf\#nameddest=tag-00RS}{II.{C.MX.10} (p.~656)} }
\expandafter\def\csname SRef@cdp:eq:C.WH.0\endcsname{\href{S6_SERIES_GRIND_II.pdf\#nameddest=tag-016X}{II.{C.WH.0} (p.~658)} }
\expandafter\def\csname SRef@cdp:eq:C.WH.1\endcsname{\href{S6_SERIES_GRIND_II.pdf\#nameddest=tag-016Y}{II.{C.WH.1} (p.~658)} }
\expandafter\def\csname SRef@cdp:eq:C.WH.2\endcsname{\href{S6_SERIES_GRIND_II.pdf\#nameddest=tag-0177}{II.{C.WH.2} (p.~658)} }
\expandafter\def\csname SRef@cdp:eq:C.WH.3\endcsname{\href{S6_SERIES_GRIND_II.pdf\#nameddest=tag-0178}{II.{C.WH.3} (p.~659)} }
\expandafter\def\csname SRef@cdp:eq:C.WH.4\endcsname{\href{S6_SERIES_GRIND_II.pdf\#nameddest=tag-0179}{II.{C.WH.4} (p.~659)} }
\expandafter\def\csname SRef@cdp:eq:C.WH.5\endcsname{\href{S6_SERIES_GRIND_II.pdf\#nameddest=tag-017A}{II.{C.WH.5} (p.~660)} }
\expandafter\def\csname SRef@cdp:eq:C.WH.6\endcsname{\href{S6_SERIES_GRIND_II.pdf\#nameddest=tag-017B}{II.{C.WH.6} (p.~660)} }
\expandafter\def\csname SRef@cdp:eq:C.WH.7\endcsname{\href{S6_SERIES_GRIND_II.pdf\#nameddest=tag-017C}{II.{C.WH.7} (p.~661)} }
\expandafter\def\csname SRef@cdp:eq:C.RF.7\endcsname{\href{S6_SERIES_GRIND_II.pdf\#nameddest=tag-00Z0}{II.{C.RF.7} (p.~439)} }
\expandafter\def\csname SRef@cdp:eq:C.WH.8\endcsname{\href{S6_SERIES_GRIND_II.pdf\#nameddest=tag-017D}{II.{C.WH.8} (p.~661)} }
\expandafter\def\csname SRef@cdp:eq:C.WH.9\endcsname{\href{S6_SERIES_GRIND_II.pdf\#nameddest=tag-017E}{II.{C.WH.9} (p.~661)} }
\expandafter\def\csname SRef@cdp:eq:C.WH.10\endcsname{\href{S6_SERIES_GRIND_II.pdf\#nameddest=tag-016Z}{II.{C.WH.10} (p.~661)} }
\expandafter\def\csname SRef@cdp:eq:C.WH.11\endcsname{\href{S6_SERIES_GRIND_II.pdf\#nameddest=tag-0170}{II.{C.WH.11} (p.~662)} }
\expandafter\def\csname SRef@cdp:eq:C.WH.12\endcsname{\href{S6_SERIES_GRIND_II.pdf\#nameddest=tag-0171}{II.{C.WH.12} (p.~662)} }
\expandafter\def\csname SRef@cdp:eq:C.WH.13\endcsname{\href{S6_SERIES_GRIND_II.pdf\#nameddest=tag-0172}{II.{C.WH.13} (p.~663)} }
\expandafter\def\csname SRef@cdp:eq:C.WH.14\endcsname{\href{S6_SERIES_GRIND_II.pdf\#nameddest=tag-0173}{II.{C.WH.14} (p.~663)} }
\expandafter\def\csname SRef@cdp:eq:C.WH.15\endcsname{\href{S6_SERIES_GRIND_II.pdf\#nameddest=tag-0174}{II.{C.WH.15} (p.~663)} }
\expandafter\def\csname SRef@cdp:eq:C.WH.16\endcsname{\href{S6_SERIES_GRIND_II.pdf\#nameddest=tag-0175}{II.{C.WH.16} (p.~663)} }
\expandafter\def\csname SRef@cdp:eq:C.WH.17\endcsname{\href{S6_SERIES_GRIND_II.pdf\#nameddest=tag-0176}{II.{C.WH.17} (p.~663)} }
\expandafter\def\csname SRef@cdp:eq:C.RF.8\endcsname{\href{S6_SERIES_GRIND_II.pdf\#nameddest=tag-00Z1}{II.{C.RF.8} (p.~439)} }
\expandafter\def\csname SRef@cdp:eq:C.AS.1\endcsname{\href{S6_SERIES_GRIND_II.pdf\#nameddest=tag-00AG}{II.{C.AS.1} (p.~664)} }
\expandafter\def\csname SRef@cdp:eq:C.AS.2\endcsname{\href{S6_SERIES_GRIND_II.pdf\#nameddest=tag-00AR}{II.{C.AS.2} (p.~664)} }
\expandafter\def\csname SRef@cdp:eq:C.AS.3\endcsname{\href{S6_SERIES_GRIND_II.pdf\#nameddest=tag-00AS}{II.{C.AS.3} (p.~664)} }
\expandafter\def\csname SRef@cdp:eq:C.AS.4\endcsname{\href{S6_SERIES_GRIND_II.pdf\#nameddest=tag-00AT}{II.{C.AS.4} (p.~664)} }
\expandafter\def\csname SRef@cdp:eq:C.AS.5\endcsname{\href{S6_SERIES_GRIND_II.pdf\#nameddest=tag-00AU}{II.{C.AS.5} (p.~664)} }
\expandafter\def\csname SRef@cdp:eq:C.AS.6\endcsname{\href{S6_SERIES_GRIND_II.pdf\#nameddest=tag-00AV}{II.{C.AS.6} (p.~665)} }
\expandafter\def\csname SRef@cdp:eq:C.AS.7\endcsname{\href{S6_SERIES_GRIND_II.pdf\#nameddest=tag-00AW}{II.{C.AS.7} (p.~665)} }
\expandafter\def\csname SRef@cdp:eq:C.AS.8\endcsname{\href{S6_SERIES_GRIND_II.pdf\#nameddest=tag-00AX}{II.{C.AS.8} (p.~665)} }
\expandafter\def\csname SRef@cdp:eq:C.AS.9\endcsname{\href{S6_SERIES_GRIND_II.pdf\#nameddest=tag-00AY}{II.{C.AS.9} (p.~665)} }
\expandafter\def\csname SRef@cdp:eq:C.AS.10\endcsname{\href{S6_SERIES_GRIND_II.pdf\#nameddest=tag-00AH}{II.{C.AS.10} (p.~666)} }
\expandafter\def\csname SRef@cdp:eq:C.CQ.1\endcsname{\href{S6_SERIES_GRIND_II.pdf\#nameddest=tag-00CJ}{II.{C.CQ.1} (p.~440)} }
\expandafter\def\csname SRef@cdp:eq:C.AS.11\endcsname{\href{S6_SERIES_GRIND_II.pdf\#nameddest=tag-00AI}{II.{C.AS.11} (p.~666)} }
\expandafter\def\csname SRef@cdp:eq:C.AS.12\endcsname{\href{S6_SERIES_GRIND_II.pdf\#nameddest=tag-00AJ}{II.{C.AS.12} (p.~666)} }
\expandafter\def\csname SRef@cdp:eq:C.AS.13\endcsname{\href{S6_SERIES_GRIND_II.pdf\#nameddest=tag-00AK}{II.{C.AS.13} (p.~666)} }
\expandafter\def\csname SRef@cdp:eq:C.AS.14\endcsname{\href{S6_SERIES_GRIND_II.pdf\#nameddest=tag-00AL}{II.{C.AS.14} (p.~667)} }
\expandafter\def\csname SRef@cdp:eq:C.AS.15\endcsname{\href{S6_SERIES_GRIND_II.pdf\#nameddest=tag-00AM}{II.{C.AS.15} (p.~667)} }
\expandafter\def\csname SRef@cdp:eq:C.AS.16\endcsname{\href{S6_SERIES_GRIND_II.pdf\#nameddest=tag-00AN}{II.{C.AS.16} (p.~667)} }
\expandafter\def\csname SRef@cdp:eq:C.AS.17\endcsname{\href{S6_SERIES_GRIND_II.pdf\#nameddest=tag-00AO}{II.{C.AS.17} (p.~668)} }
\expandafter\def\csname SRef@cdp:eq:C.AS.18\endcsname{\href{S6_SERIES_GRIND_II.pdf\#nameddest=tag-00AP}{II.{C.AS.18} (p.~668)} }
\expandafter\def\csname SRef@cdp:eq:C.AS.19\endcsname{\href{S6_SERIES_GRIND_II.pdf\#nameddest=tag-00AQ}{II.{C.AS.19} (p.~668)} }
\expandafter\def\csname SRef@cdp:eq:C.HR.1\endcsname{\href{S6_SERIES_GRIND_II.pdf\#nameddest=tag-00N0}{II.{C.HR.1} (p.~669)} }
\expandafter\def\csname SRef@cdp:eq:C.CQ.2\endcsname{\href{S6_SERIES_GRIND_II.pdf\#nameddest=tag-00CN}{II.{C.CQ.2} (p.~440)} }
\expandafter\def\csname SRef@cdp:eq:C.HR.2\endcsname{\href{S6_SERIES_GRIND_II.pdf\#nameddest=tag-00N1}{II.{C.HR.2} (p.~669)} }
\expandafter\def\csname SRef@cdp:eq:C.HR.3\endcsname{\href{S6_SERIES_GRIND_II.pdf\#nameddest=tag-00N2}{II.{C.HR.3} (p.~669)} }
\expandafter\def\csname SRef@cdp:eq:C.HR.4\endcsname{\href{S6_SERIES_GRIND_II.pdf\#nameddest=tag-00N3}{II.{C.HR.4} (p.~670)} }
\expandafter\def\csname SRef@cdp:eq:C.HR.5\endcsname{\href{S6_SERIES_GRIND_II.pdf\#nameddest=tag-00N4}{II.{C.HR.5} (p.~670)} }
\expandafter\def\csname SRef@cdp:eq:C.HR.6\endcsname{\href{S6_SERIES_GRIND_II.pdf\#nameddest=tag-00N5}{II.{C.HR.6} (p.~670)} }
\expandafter\def\csname SRef@cdp:eq:C.HR.7\endcsname{\href{S6_SERIES_GRIND_II.pdf\#nameddest=tag-00N6}{II.{C.HR.7} (p.~671)} }
\expandafter\def\csname SRef@cdp:eq:C.HR.8\endcsname{\href{S6_SERIES_GRIND_II.pdf\#nameddest=tag-00N7}{II.{C.HR.8} (p.~671)} }
\expandafter\def\csname SRef@cdp:eq:C.TL.1\endcsname{\href{S6_SERIES_GRIND_II.pdf\#nameddest=tag-013Q}{II.{C.TL.1} (p.~673)} }
\expandafter\def\csname SRef@cdp:eq:C.TL.2\endcsname{\href{S6_SERIES_GRIND_II.pdf\#nameddest=tag-013R}{II.{C.TL.2} (p.~673)} }
\expandafter\def\csname SRef@cdp:eq:C.TL.3\endcsname{\href{S6_SERIES_GRIND_II.pdf\#nameddest=tag-013S}{II.{C.TL.3} (p.~673)} }
\expandafter\def\csname SRef@cdp:eq:C.CQ.3\endcsname{\href{S6_SERIES_GRIND_II.pdf\#nameddest=tag-00CO}{II.{C.CQ.3} (p.~440)} }
\expandafter\def\csname SRef@cdp:eq:C.TL.4\endcsname{\href{S6_SERIES_GRIND_II.pdf\#nameddest=tag-013T}{II.{C.TL.4} (p.~673)} }
\expandafter\def\csname SRef@cdp:eq:C.TL.5\endcsname{\href{S6_SERIES_GRIND_II.pdf\#nameddest=tag-013U}{II.{C.TL.5} (p.~673)} }
\expandafter\def\csname SRef@cdp:eq:C.TL.6\endcsname{\href{S6_SERIES_GRIND_II.pdf\#nameddest=tag-013V}{II.{C.TL.6} (p.~673)} }
\expandafter\def\csname SRef@cdp:eq:C.TL.7\endcsname{\href{S6_SERIES_GRIND_II.pdf\#nameddest=tag-013W}{II.{C.TL.7} (p.~674)} }
\expandafter\def\csname SRef@cdp:eq:C.TL.8\endcsname{\href{S6_SERIES_GRIND_II.pdf\#nameddest=tag-013X}{II.{C.TL.8} (p.~674)} }
\expandafter\def\csname SRef@cdp:eq:C.WM.1\endcsname{\href{S6_SERIES_GRIND_II.pdf\#nameddest=tag-017F}{II.{C.WM.1} (p.~675)} }
\expandafter\def\csname SRef@cdp:eq:C.WM.2\endcsname{\href{S6_SERIES_GRIND_II.pdf\#nameddest=tag-017G}{II.{C.WM.2} (p.~676)} }
\expandafter\def\csname SRef@cdp:eq:C.WM.3\endcsname{\href{S6_SERIES_GRIND_II.pdf\#nameddest=tag-017H}{II.{C.WM.3} (p.~676)} }
\expandafter\def\csname SRef@cdp:eq:C.WM.4\endcsname{\href{S6_SERIES_GRIND_II.pdf\#nameddest=tag-017I}{II.{C.WM.4} (p.~676)} }
\expandafter\def\csname SRef@cdp:eq:C.WM.5\endcsname{\href{S6_SERIES_GRIND_II.pdf\#nameddest=tag-017J}{II.{C.WM.5} (p.~676)} }
\expandafter\def\csname SRef@cdp:eq:C.CQ.4\endcsname{\href{S6_SERIES_GRIND_II.pdf\#nameddest=tag-00CP}{II.{C.CQ.4} (p.~440)} }
\expandafter\def\csname SRef@cdp:eq:C.WM.6\endcsname{\href{S6_SERIES_GRIND_II.pdf\#nameddest=tag-017K}{II.{C.WM.6} (p.~677)} }
\expandafter\def\csname SRef@cdp:eq:C.WM.7\endcsname{\href{S6_SERIES_GRIND_II.pdf\#nameddest=tag-017L}{II.{C.WM.7} (p.~677)} }
\expandafter\def\csname SRef@cdp:eq:C.WM.8\endcsname{\href{S6_SERIES_GRIND_II.pdf\#nameddest=tag-017M}{II.{C.WM.8} (p.~677)} }
\expandafter\def\csname SRef@cdp:eq:C.WM.9\endcsname{\href{S6_SERIES_GRIND_II.pdf\#nameddest=tag-017N}{II.{C.WM.9} (p.~677)} }
\expandafter\def\csname SRef@cdp:eq:C.AD.1\endcsname{\href{S6_SERIES_GRIND_II.pdf\#nameddest=tag-009B}{II.{C.AD.1} (p.~679)} }
\expandafter\def\csname SRef@cdp:eq:C.AD.2\endcsname{\href{S6_SERIES_GRIND_II.pdf\#nameddest=tag-009E}{II.{C.AD.2} (p.~680)} }
\expandafter\def\csname SRef@cdp:eq:C.AD.3\endcsname{\href{S6_SERIES_GRIND_II.pdf\#nameddest=tag-009F}{II.{C.AD.3} (p.~680)} }
\expandafter\def\csname SRef@cdp:eq:C.AD.4\endcsname{\href{S6_SERIES_GRIND_II.pdf\#nameddest=tag-009G}{II.{C.AD.4} (p.~680)} }
\expandafter\def\csname SRef@cdp:eq:C.AD.5\endcsname{\href{S6_SERIES_GRIND_II.pdf\#nameddest=tag-009H}{II.{C.AD.5} (p.~680)} }
\expandafter\def\csname SRef@cdp:eq:C.AD.6\endcsname{\href{S6_SERIES_GRIND_II.pdf\#nameddest=tag-009I}{II.{C.AD.6} (p.~680)} }
\expandafter\def\csname SRef@cdp:eq:C.CQ.5\endcsname{\href{S6_SERIES_GRIND_II.pdf\#nameddest=tag-00CQ}{II.{C.CQ.5} (p.~441)} }
\expandafter\def\csname SRef@cdp:eq:C.AD.7\endcsname{\href{S6_SERIES_GRIND_II.pdf\#nameddest=tag-009J}{II.{C.AD.7} (p.~681)} }
\expandafter\def\csname SRef@cdp:eq:C.AD.8\endcsname{\href{S6_SERIES_GRIND_II.pdf\#nameddest=tag-009K}{II.{C.AD.8} (p.~681)} }
\expandafter\def\csname SRef@cdp:eq:C.AD.9\endcsname{\href{S6_SERIES_GRIND_II.pdf\#nameddest=tag-009L}{II.{C.AD.9} (p.~681)} }
\expandafter\def\csname SRef@cdp:eq:C.AD.10\endcsname{\href{S6_SERIES_GRIND_II.pdf\#nameddest=tag-009C}{II.{C.AD.10} (p.~681)} }
\expandafter\def\csname SRef@cdp:eq:C.AD.11\endcsname{\href{S6_SERIES_GRIND_II.pdf\#nameddest=tag-009D}{II.{C.AD.11} (p.~682)} }
\expandafter\def\csname SRef@cdp:eq:C.HT.1\endcsname{\href{S6_SERIES_GRIND_II.pdf\#nameddest=tag-00N8}{II.{C.HT.1} (p.~682)} }
\expandafter\def\csname SRef@cdp:eq:C.HT.2\endcsname{\href{S6_SERIES_GRIND_II.pdf\#nameddest=tag-00NC}{II.{C.HT.2} (p.~683)} }
\expandafter\def\csname SRef@cdp:eq:C.HT.3\endcsname{\href{S6_SERIES_GRIND_II.pdf\#nameddest=tag-00ND}{II.{C.HT.3} (p.~683)} }
\expandafter\def\csname SRef@cdp:eq:C.HT.4\endcsname{\href{S6_SERIES_GRIND_II.pdf\#nameddest=tag-00NE}{II.{C.HT.4} (p.~684)} }
\expandafter\def\csname SRef@cdp:eq:C.HT.5\endcsname{\href{S6_SERIES_GRIND_II.pdf\#nameddest=tag-00NF}{II.{C.HT.5} (p.~684)} }
\expandafter\def\csname SRef@cdp:eq:C.CQ.6\endcsname{\href{S6_SERIES_GRIND_II.pdf\#nameddest=tag-00CR}{II.{C.CQ.6} (p.~441)} }
\expandafter\def\csname SRef@cdp:eq:C.HT.6\endcsname{\href{S6_SERIES_GRIND_II.pdf\#nameddest=tag-00NG}{II.{C.HT.6} (p.~684)} }
\expandafter\def\csname SRef@cdp:eq:C.HT.7\endcsname{\href{S6_SERIES_GRIND_II.pdf\#nameddest=tag-00NH}{II.{C.HT.7} (p.~684)} }
\expandafter\def\csname SRef@cdp:eq:C.HT.8\endcsname{\href{S6_SERIES_GRIND_II.pdf\#nameddest=tag-00NI}{II.{C.HT.8} (p.~684)} }
\expandafter\def\csname SRef@cdp:eq:C.HT.9\endcsname{\href{S6_SERIES_GRIND_II.pdf\#nameddest=tag-00NJ}{II.{C.HT.9} (p.~685)} }
\expandafter\def\csname SRef@cdp:eq:C.HT.10\endcsname{\href{S6_SERIES_GRIND_II.pdf\#nameddest=tag-00N9}{II.{C.HT.10} (p.~685)} }
\expandafter\def\csname SRef@cdp:eq:C.HT.11\endcsname{\href{S6_SERIES_GRIND_II.pdf\#nameddest=tag-00NA}{II.{C.HT.11} (p.~686)} }
\expandafter\def\csname SRef@cdp:eq:C.HT.12\endcsname{\href{S6_SERIES_GRIND_II.pdf\#nameddest=tag-00NB}{II.{C.HT.12} (p.~686)} }
\expandafter\def\csname SRef@cdp:eq:C.EW.1\endcsname{\href{S6_SERIES_GRIND_II.pdf\#nameddest=tag-00JH}{II.{C.EW.1} (p.~687)} }
\expandafter\def\csname SRef@cdp:eq:C.EW.2\endcsname{\href{S6_SERIES_GRIND_II.pdf\#nameddest=tag-00JS}{II.{C.EW.2} (p.~688)} }
\expandafter\def\csname SRef@cdp:eq:C.EW.3\endcsname{\href{S6_SERIES_GRIND_II.pdf\#nameddest=tag-00JV}{II.{C.EW.3} (p.~688)} }
\expandafter\def\csname SRef@cdp:eq:C.CQ.7\endcsname{\href{S6_SERIES_GRIND_II.pdf\#nameddest=tag-00CS}{II.{C.CQ.7} (p.~441)} }
\expandafter\def\csname SRef@cdp:eq:C.EW.4\endcsname{\href{S6_SERIES_GRIND_II.pdf\#nameddest=tag-00JW}{II.{C.EW.4} (p.~688)} }
\expandafter\def\csname SRef@cdp:eq:C.EW.5\endcsname{\href{S6_SERIES_GRIND_II.pdf\#nameddest=tag-00JX}{II.{C.EW.5} (p.~688)} }
\expandafter\def\csname SRef@cdp:eq:C.EW.6\endcsname{\href{S6_SERIES_GRIND_II.pdf\#nameddest=tag-00JY}{II.{C.EW.6} (p.~688)} }
\expandafter\def\csname SRef@cdp:eq:C.EW.7\endcsname{\href{S6_SERIES_GRIND_II.pdf\#nameddest=tag-00JZ}{II.{C.EW.7} (p.~689)} }
\expandafter\def\csname SRef@cdp:eq:C.EW.8\endcsname{\href{S6_SERIES_GRIND_II.pdf\#nameddest=tag-00K0}{II.{C.EW.8} (p.~689)} }
\expandafter\def\csname SRef@cdp:eq:C.EW.9\endcsname{\href{S6_SERIES_GRIND_II.pdf\#nameddest=tag-00K1}{II.{C.EW.9} (p.~689)} }
\expandafter\def\csname SRef@cdp:eq:C.EW.10\endcsname{\href{S6_SERIES_GRIND_II.pdf\#nameddest=tag-00JI}{II.{C.EW.10} (p.~689)} }
\expandafter\def\csname SRef@cdp:eq:C.EW.11\endcsname{\href{S6_SERIES_GRIND_II.pdf\#nameddest=tag-00JJ}{II.{C.EW.11} (p.~689)} }
\expandafter\def\csname SRef@cdp:eq:C.EW.12\endcsname{\href{S6_SERIES_GRIND_II.pdf\#nameddest=tag-00JK}{II.{C.EW.12} (p.~689)} }
\expandafter\def\csname SRef@cdp:eq:C.EW.13\endcsname{\href{S6_SERIES_GRIND_II.pdf\#nameddest=tag-00JL}{II.{C.EW.13} (p.~689)} }
\expandafter\def\csname SRef@cdp:eq:C.CQ.8\endcsname{\href{S6_SERIES_GRIND_II.pdf\#nameddest=tag-00CT}{II.{C.CQ.8} (p.~441)} }
\expandafter\def\csname SRef@cdp:eq:C.EW.14\endcsname{\href{S6_SERIES_GRIND_II.pdf\#nameddest=tag-00JM}{II.{C.EW.14} (p.~689)} }
\expandafter\def\csname SRef@cdp:eq:C.EW.15\endcsname{\href{S6_SERIES_GRIND_II.pdf\#nameddest=tag-00JN}{II.{C.EW.15} (p.~690)} }
\expandafter\def\csname SRef@cdp:eq:C.EW.16\endcsname{\href{S6_SERIES_GRIND_II.pdf\#nameddest=tag-00JO}{II.{C.EW.16} (p.~690)} }
\expandafter\def\csname SRef@cdp:eq:C.EW.17\endcsname{\href{S6_SERIES_GRIND_II.pdf\#nameddest=tag-00JP}{II.{C.EW.17} (p.~690)} }
\expandafter\def\csname SRef@cdp:eq:C.EW.18\endcsname{\href{S6_SERIES_GRIND_II.pdf\#nameddest=tag-00JQ}{II.{C.EW.18} (p.~691)} }
\expandafter\def\csname SRef@cdp:eq:C.EW.19\endcsname{\href{S6_SERIES_GRIND_II.pdf\#nameddest=tag-00JR}{II.{C.EW.19} (p.~691)} }
\expandafter\def\csname SRef@cdp:eq:C.EW.20\endcsname{\href{S6_SERIES_GRIND_II.pdf\#nameddest=tag-00JT}{II.{C.EW.20} (p.~691)} }
\expandafter\def\csname SRef@cdp:eq:C.EW.21\endcsname{\href{S6_SERIES_GRIND_II.pdf\#nameddest=tag-00JU}{II.{C.EW.21} (p.~691)} }
\expandafter\def\csname SRef@cdp:eq:C.PR.1\endcsname{\href{S6_SERIES_GRIND_II.pdf\#nameddest=tag-00VH}{II.{C.PR.1} (p.~693)} }
\expandafter\def\csname SRef@cdp:eq:C.PR.2\endcsname{\href{S6_SERIES_GRIND_II.pdf\#nameddest=tag-00VR}{II.{C.PR.2} (p.~693)} }
\expandafter\def\csname SRef@cdp:eq:C.CQ.9\endcsname{\href{S6_SERIES_GRIND_II.pdf\#nameddest=tag-00CU}{II.{C.CQ.9} (p.~441)} }
\expandafter\def\csname SRef@cdp:eq:C.PR.3\endcsname{\href{S6_SERIES_GRIND_II.pdf\#nameddest=tag-00VS}{II.{C.PR.3} (p.~693)} }
\expandafter\def\csname SRef@cdp:eq:C.PR.4\endcsname{\href{S6_SERIES_GRIND_II.pdf\#nameddest=tag-00VT}{II.{C.PR.4} (p.~693)} }
\expandafter\def\csname SRef@cdp:eq:C.PR.5\endcsname{\href{S6_SERIES_GRIND_II.pdf\#nameddest=tag-00VU}{II.{C.PR.5} (p.~693)} }
\expandafter\def\csname SRef@cdp:eq:C.PR.6\endcsname{\href{S6_SERIES_GRIND_II.pdf\#nameddest=tag-00VV}{II.{C.PR.6} (p.~694)} }
\expandafter\def\csname SRef@cdp:eq:C.PR.7\endcsname{\href{S6_SERIES_GRIND_II.pdf\#nameddest=tag-00VW}{II.{C.PR.7} (p.~694)} }
\expandafter\def\csname SRef@cdp:eq:C.PR.8\endcsname{\href{S6_SERIES_GRIND_II.pdf\#nameddest=tag-00VX}{II.{C.PR.8} (p.~695)} }
\expandafter\def\csname SRef@cdp:eq:C.PR.9\endcsname{\href{S6_SERIES_GRIND_II.pdf\#nameddest=tag-00VY}{II.{C.PR.9} (p.~695)} }
\expandafter\def\csname SRef@cdp:eq:C.PR.10\endcsname{\href{S6_SERIES_GRIND_II.pdf\#nameddest=tag-00VI}{II.{C.PR.10} (p.~695)} }
\expandafter\def\csname SRef@cdp:eq:C.PR.11\endcsname{\href{S6_SERIES_GRIND_II.pdf\#nameddest=tag-00VJ}{II.{C.PR.11} (p.~695)} }
\expandafter\def\csname SRef@cdp:eq:C.PR.12\endcsname{\href{S6_SERIES_GRIND_II.pdf\#nameddest=tag-00VK}{II.{C.PR.12} (p.~696)} }
\expandafter\def\csname SRef@cdp:eq:C.CQ.10\endcsname{\href{S6_SERIES_GRIND_II.pdf\#nameddest=tag-00CK}{II.{C.CQ.10} (p.~441)} }
\expandafter\def\csname SRef@cdp:eq:C.PR.13\endcsname{\href{S6_SERIES_GRIND_II.pdf\#nameddest=tag-00VL}{II.{C.PR.13} (p.~696)} }
\expandafter\def\csname SRef@cdp:eq:C.PR.14\endcsname{\href{S6_SERIES_GRIND_II.pdf\#nameddest=tag-00VM}{II.{C.PR.14} (p.~696)} }
\expandafter\def\csname SRef@cdp:eq:C.PR.15\endcsname{\href{S6_SERIES_GRIND_II.pdf\#nameddest=tag-00VN}{II.{C.PR.15} (p.~696)} }
\expandafter\def\csname SRef@cdp:eq:C.PR.16\endcsname{\href{S6_SERIES_GRIND_II.pdf\#nameddest=tag-00VO}{II.{C.PR.16} (p.~697)} }
\expandafter\def\csname SRef@cdp:eq:C.PR.17\endcsname{\href{S6_SERIES_GRIND_II.pdf\#nameddest=tag-00VP}{II.{C.PR.17} (p.~697)} }
\expandafter\def\csname SRef@cdp:eq:C.PR.18\endcsname{\href{S6_SERIES_GRIND_II.pdf\#nameddest=tag-00VQ}{II.{C.PR.18} (p.~697)} }
\expandafter\def\csname SRef@cdp:eq:C.PT.1\endcsname{\href{S6_SERIES_GRIND_II.pdf\#nameddest=tag-00VZ}{II.{C.PT.1} (p.~699)} }
\expandafter\def\csname SRef@cdp:eq:C.PT.2\endcsname{\href{S6_SERIES_GRIND_II.pdf\#nameddest=tag-00W5}{II.{C.PT.2} (p.~699)} }
\expandafter\def\csname SRef@cdp:eq:C.PT.3\endcsname{\href{S6_SERIES_GRIND_II.pdf\#nameddest=tag-00W6}{II.{C.PT.3} (p.~699)} }
\expandafter\def\csname SRef@cdp:eq:C.PT.4\endcsname{\href{S6_SERIES_GRIND_II.pdf\#nameddest=tag-00W7}{II.{C.PT.4} (p.~699)} }
\expandafter\def\csname SRef@cdp:eq:C.CQ.11\endcsname{\href{S6_SERIES_GRIND_II.pdf\#nameddest=tag-00CL}{II.{C.CQ.11} (p.~442)} }
\expandafter\def\csname SRef@cdp:eq:C.PT.5\endcsname{\href{S6_SERIES_GRIND_II.pdf\#nameddest=tag-00W8}{II.{C.PT.5} (p.~700)} }
\expandafter\def\csname SRef@cdp:eq:C.PT.6\endcsname{\href{S6_SERIES_GRIND_II.pdf\#nameddest=tag-00W9}{II.{C.PT.6} (p.~700)} }
\expandafter\def\csname SRef@cdp:eq:C.PT.7\endcsname{\href{S6_SERIES_GRIND_II.pdf\#nameddest=tag-00WA}{II.{C.PT.7} (p.~700)} }
\expandafter\def\csname SRef@cdp:eq:C.PT.8\endcsname{\href{S6_SERIES_GRIND_II.pdf\#nameddest=tag-00WB}{II.{C.PT.8} (p.~700)} }
\expandafter\def\csname SRef@cdp:eq:C.PT.9\endcsname{\href{S6_SERIES_GRIND_II.pdf\#nameddest=tag-00WC}{II.{C.PT.9} (p.~701)} }
\expandafter\def\csname SRef@cdp:eq:C.PT.10\endcsname{\href{S6_SERIES_GRIND_II.pdf\#nameddest=tag-00W0}{II.{C.PT.10} (p.~702)} }
\expandafter\def\csname SRef@cdp:eq:C.PT.11\endcsname{\href{S6_SERIES_GRIND_II.pdf\#nameddest=tag-00W1}{II.{C.PT.11} (p.~702)} }
\expandafter\def\csname SRef@cdp:eq:C.PT.12\endcsname{\href{S6_SERIES_GRIND_II.pdf\#nameddest=tag-00W2}{II.{C.PT.12} (p.~702)} }
\expandafter\def\csname SRef@cdp:eq:C.PT.13\endcsname{\href{S6_SERIES_GRIND_II.pdf\#nameddest=tag-00W3}{II.{C.PT.13} (p.~702)} }
\expandafter\def\csname SRef@cdp:eq:C.PT.14\endcsname{\href{S6_SERIES_GRIND_II.pdf\#nameddest=tag-00W4}{II.{C.PT.14} (p.~703)} }
\expandafter\def\csname SRef@cdp:eq:C.CQ.12\endcsname{\href{S6_SERIES_GRIND_II.pdf\#nameddest=tag-00CM}{II.{C.CQ.12} (p.~442)} }
\expandafter\def\csname SRef@cdp:eq:C.QT.1\endcsname{\href{S6_SERIES_GRIND_II.pdf\#nameddest=tag-00WD}{II.{C.QT.1} (p.~704)} }
\expandafter\def\csname SRef@cdp:eq:C.QT.2\endcsname{\href{S6_SERIES_GRIND_II.pdf\#nameddest=tag-00WI}{II.{C.QT.2} (p.~704)} }
\expandafter\def\csname SRef@cdp:eq:C.QT.3\endcsname{\href{S6_SERIES_GRIND_II.pdf\#nameddest=tag-00WJ}{II.{C.QT.3} (p.~705)} }
\expandafter\def\csname SRef@cdp:eq:C.QT.4\endcsname{\href{S6_SERIES_GRIND_II.pdf\#nameddest=tag-00WK}{II.{C.QT.4} (p.~705)} }
\expandafter\def\csname SRef@cdp:eq:C.QT.5\endcsname{\href{S6_SERIES_GRIND_II.pdf\#nameddest=tag-00WL}{II.{C.QT.5} (p.~705)} }
\expandafter\def\csname SRef@cdp:eq:C.QT.6\endcsname{\href{S6_SERIES_GRIND_II.pdf\#nameddest=tag-00WM}{II.{C.QT.6} (p.~705)} }
\expandafter\def\csname SRef@cdp:eq:C.QT.7\endcsname{\href{S6_SERIES_GRIND_II.pdf\#nameddest=tag-00WN}{II.{C.QT.7} (p.~706)} }
\expandafter\def\csname SRef@cdp:eq:C.QT.8\endcsname{\href{S6_SERIES_GRIND_II.pdf\#nameddest=tag-00WO}{II.{C.QT.8} (p.~706)} }
\expandafter\def\csname SRef@cdp:eq:C.QT.9\endcsname{\href{S6_SERIES_GRIND_II.pdf\#nameddest=tag-00WP}{II.{C.QT.9} (p.~706)} }
\expandafter\def\csname SRef@cdp:eq:C.QT.10\endcsname{\href{S6_SERIES_GRIND_II.pdf\#nameddest=tag-00WE}{II.{C.QT.10} (p.~706)} }
\expandafter\def\csname SRef@cdp:st:C.2.UN.1\endcsname{\href{S6_SERIES_GRIND_II.pdf\#nameddest=tag-01CM}{II.19.1 (p.~413)} }
\expandafter\def\csname SRef@cdp:st:C.2.UN.20\endcsname{\href{S6_SERIES_GRIND_II.pdf\#nameddest=tag-01CY}{II.19.10 (p.~433)} }
\expandafter\def\csname SRef@cdp:st:C.2.UN.22\endcsname{\href{S6_SERIES_GRIND_II.pdf\#nameddest=tag-01D0}{II.19.11 (p.~435)} }
\expandafter\def\csname SRef@cdp:st:C.2.UN.23\endcsname{\href{S6_SERIES_GRIND_II.pdf\#nameddest=tag-01D1}{II.19.12 (p.~439)} }
\expandafter\def\csname SRef@cdp:st:C.2.UN.26\endcsname{\href{S6_SERIES_GRIND_II.pdf\#nameddest=tag-01D4}{II.19.13 (p.~443)} }
\expandafter\def\csname SRef@cdp:st:C.2.UN.28\endcsname{\href{S6_SERIES_GRIND_II.pdf\#nameddest=tag-01D6}{II.19.14 (p.~444)} }
\expandafter\def\csname SRef@cdp:st:C.2.UN.29\endcsname{\href{S6_SERIES_GRIND_II.pdf\#nameddest=tag-01D7}{II.19.15 (p.~446)} }
\expandafter\def\csname SRef@cdp:st:C.2.UN.31\endcsname{\href{S6_SERIES_GRIND_II.pdf\#nameddest=tag-01DA}{II.19.16 (p.~450)} }
\expandafter\def\csname SRef@cdp:st:C.2.UN.33\endcsname{\href{S6_SERIES_GRIND_II.pdf\#nameddest=tag-01DC}{II.19.17 (p.~453)} }
\expandafter\def\csname SRef@cdp:st:C.2.UN.36\endcsname{\href{S6_SERIES_GRIND_II.pdf\#nameddest=tag-01DF}{II.19.18 (p.~457)} }
\expandafter\def\csname SRef@cdp:st:C.2.UN.38\endcsname{\href{S6_SERIES_GRIND_II.pdf\#nameddest=tag-01DH}{II.19.19 (p.~460)} }
\expandafter\def\csname SRef@cdp:st:C.2.UN.2\endcsname{\href{S6_SERIES_GRIND_II.pdf\#nameddest=tag-01CX}{II.19.2 (p.~414)} }
\expandafter\def\csname SRef@cdp:st:C.2.UN.39\endcsname{\href{S6_SERIES_GRIND_II.pdf\#nameddest=tag-01DI}{II.19.20 (p.~463)} }
\expandafter\def\csname SRef@cdp:st:C.2.UN.44\endcsname{\href{S6_SERIES_GRIND_II.pdf\#nameddest=tag-01DO}{II.19.21 (p.~466)} }
\expandafter\def\csname SRef@cdp:st:C.2.UN.47\endcsname{\href{S6_SERIES_GRIND_II.pdf\#nameddest=tag-01DR}{II.19.22 (p.~471)} }
\expandafter\def\csname SRef@cdp:st:C.2.UN.3\endcsname{\href{S6_SERIES_GRIND_II.pdf\#nameddest=tag-01D8}{II.19.3 (p.~415)} }
\expandafter\def\csname SRef@cdp:st:C.2.UN.4\endcsname{\href{S6_SERIES_GRIND_II.pdf\#nameddest=tag-01DJ}{II.19.4 (p.~418)} }
\expandafter\def\csname SRef@cdp:st:C.2.UN.14\endcsname{\href{S6_SERIES_GRIND_II.pdf\#nameddest=tag-01CR}{II.19.5 (p.~421)} }
\expandafter\def\csname SRef@cdp:st:C.2.UN.15\endcsname{\href{S6_SERIES_GRIND_II.pdf\#nameddest=tag-01CS}{II.19.6 (p.~423)} }
\expandafter\def\csname SRef@cdp:st:C.2.UN.16\endcsname{\href{S6_SERIES_GRIND_II.pdf\#nameddest=tag-01CT}{II.19.7 (p.~425)} }
\expandafter\def\csname SRef@cdp:st:C.2.UN.17\endcsname{\href{S6_SERIES_GRIND_II.pdf\#nameddest=tag-01CU}{II.19.8 (p.~428)} }
\expandafter\def\csname SRef@cdp:st:C.2.UN.19\endcsname{\href{S6_SERIES_GRIND_II.pdf\#nameddest=tag-01CW}{II.19.9 (p.~431)} }
\expandafter\def\csname SRef@cdp:st:C.2.UN.5\endcsname{\href{S6_SERIES_GRIND_II.pdf\#nameddest=tag-01DU}{II.20.1 (p.~474)} }
\expandafter\def\csname SRef@cdp:st:C.2.IRR.1\endcsname{\href{S6_SERIES_GRIND_II.pdf\#nameddest=tag-01C3}{II.20.10 (p.~487)} }
\expandafter\def\csname SRef@cdp:st:C.2.VII.1\endcsname{\href{S6_SERIES_GRIND_II.pdf\#nameddest=tag-01F0}{II.20.11 (p.~489)} }
\expandafter\def\csname SRef@cdp:st:C.2.VII.2\endcsname{\href{S6_SERIES_GRIND_II.pdf\#nameddest=tag-01F1}{II.20.12 (p.~491)} }
\expandafter\def\csname SRef@cdp:st:C.2.VII.3\endcsname{\href{S6_SERIES_GRIND_II.pdf\#nameddest=tag-01F2}{II.20.13 (p.~492)} }
\expandafter\def\csname SRef@cdp:st:C.2.VII.4\endcsname{\href{S6_SERIES_GRIND_II.pdf\#nameddest=tag-01F3}{II.20.14 (p.~493)} }
\expandafter\def\csname SRef@cdp:st:C.2.RAT.1\endcsname{\href{S6_SERIES_GRIND_II.pdf\#nameddest=tag-01CD}{II.20.15 (p.~494)} }
\expandafter\def\csname SRef@cdp:st:C.2.RAT.2\endcsname{\href{S6_SERIES_GRIND_II.pdf\#nameddest=tag-01CE}{II.20.16 (p.~498)} }
\expandafter\def\csname SRef@cdp:st:C.2.RAT.3\endcsname{\href{S6_SERIES_GRIND_II.pdf\#nameddest=tag-01CF}{II.20.17 (p.~499)} }
\expandafter\def\csname SRef@cdp:st:C.2.RAT.4\endcsname{\href{S6_SERIES_GRIND_II.pdf\#nameddest=tag-01CG}{II.20.18 (p.~501)} }
\expandafter\def\csname SRef@cdp:st:C.2.RAT.5\endcsname{\href{S6_SERIES_GRIND_II.pdf\#nameddest=tag-01CH}{II.20.19 (p.~503)} }
\expandafter\def\csname SRef@cdp:st:C.2.CR.1\endcsname{\href{S6_SERIES_GRIND_II.pdf\#nameddest=tag-01BQ}{II.20.2 (p.~476)} }
\expandafter\def\csname SRef@cdp:st:C.2.RAT.7\endcsname{\href{S6_SERIES_GRIND_II.pdf\#nameddest=tag-01CJ}{II.20.20 (p.~504)} }
\expandafter\def\csname SRef@cdp:st:C.2.RAT.8\endcsname{\href{S6_SERIES_GRIND_II.pdf\#nameddest=tag-01CK}{II.20.21 (p.~506)} }
\expandafter\def\csname SRef@cdp:st:C.2.RAT.9\endcsname{\href{S6_SERIES_GRIND_II.pdf\#nameddest=tag-01CL}{II.20.22 (p.~509)} }
\expandafter\def\csname SRef@cdp:st:C.2.D4.1\endcsname{\href{S6_SERIES_GRIND_II.pdf\#nameddest=tag-01BW}{II.20.23 (p.~511)} }
\expandafter\def\csname SRef@cdp:st:C.2.D4.2\endcsname{\href{S6_SERIES_GRIND_II.pdf\#nameddest=tag-01BX}{II.20.24 (p.~516)} }
\expandafter\def\csname SRef@cdp:st:C.2.D4.3\endcsname{\href{S6_SERIES_GRIND_II.pdf\#nameddest=tag-01BY}{II.20.25 (p.~520)} }
\expandafter\def\csname SRef@cdp:st:C.2.UN.7\endcsname{\href{S6_SERIES_GRIND_II.pdf\#nameddest=tag-01EG}{II.20.26 (p.~522)} }
\expandafter\def\csname SRef@cdp:st:C.2.UN.12\endcsname{\href{S6_SERIES_GRIND_II.pdf\#nameddest=tag-01CP}{II.20.27 (p.~528)} }
\expandafter\def\csname SRef@cdp:st:C.2.INT.1\endcsname{\href{S6_SERIES_GRIND_II.pdf\#nameddest=tag-01C0}{II.20.28 (p.~531)} }
\expandafter\def\csname SRef@cdp:st:C.2.INT.2\endcsname{\href{S6_SERIES_GRIND_II.pdf\#nameddest=tag-01C1}{II.20.29 (p.~533)} }
\expandafter\def\csname SRef@cdp:st:C.2.CR.2\endcsname{\href{S6_SERIES_GRIND_II.pdf\#nameddest=tag-01BR}{II.20.3 (p.~478)} }
\expandafter\def\csname SRef@cdp:st:C.2.INT.3\endcsname{\href{S6_SERIES_GRIND_II.pdf\#nameddest=tag-01C2}{II.20.30 (p.~535)} }
\expandafter\def\csname SRef@cdp:st:C.2.MH.1\endcsname{\href{S6_SERIES_GRIND_II.pdf\#nameddest=tag-01C4}{II.20.31 (p.~538)} }
\expandafter\def\csname SRef@cdp:st:C.2.MH.2\endcsname{\href{S6_SERIES_GRIND_II.pdf\#nameddest=tag-01C5}{II.20.32 (p.~540)} }
\expandafter\def\csname SRef@cdp:st:C.2.MH.3\endcsname{\href{S6_SERIES_GRIND_II.pdf\#nameddest=tag-01C6}{II.20.33 (p.~542)} }
\expandafter\def\csname SRef@cdp:st:C.2.MH.4\endcsname{\href{S6_SERIES_GRIND_II.pdf\#nameddest=tag-01C7}{II.20.34 (p.~542)} }
\expandafter\def\csname SRef@cdp:st:C.2.CR.3\endcsname{\href{S6_SERIES_GRIND_II.pdf\#nameddest=tag-01BS}{II.20.4 (p.~480)} }
\expandafter\def\csname SRef@cdp:st:C.2.CR.4\endcsname{\href{S6_SERIES_GRIND_II.pdf\#nameddest=tag-01BT}{II.20.5 (p.~481)} }
\expandafter\def\csname SRef@cdp:st:C.2.CR.5\endcsname{\href{S6_SERIES_GRIND_II.pdf\#nameddest=tag-01BU}{II.20.6 (p.~482)} }
\expandafter\def\csname SRef@cdp:st:C.2.CR.6\endcsname{\href{S6_SERIES_GRIND_II.pdf\#nameddest=tag-01BV}{II.20.7 (p.~483)} }
\expandafter\def\csname SRef@cdp:st:C.2.UN.9\endcsname{\href{S6_SERIES_GRIND_II.pdf\#nameddest=tag-01EZ}{II.20.8 (p.~484)} }
\expandafter\def\csname SRef@cdp:st:C.2.UN.10\endcsname{\href{S6_SERIES_GRIND_II.pdf\#nameddest=tag-01CN}{II.20.9 (p.~486)} }
\expandafter\def\csname SRef@cdp:st:C2a\endcsname{\href{S6_SERIES_GRIND_II.pdf\#nameddest=tag-01FG}{II.21.1 (p.~543)} }
\expandafter\def\csname SRef@cdp:st:C.2.UN.67\endcsname{\href{S6_SERIES_GRIND_II.pdf\#nameddest=tag-01ED}{II.21.10 (p.~573)} }
\expandafter\def\csname SRef@cdp:st:C.2.UN.68\endcsname{\href{S6_SERIES_GRIND_II.pdf\#nameddest=tag-01EE}{II.21.11 (p.~576)} }
\expandafter\def\csname SRef@cdp:st:C2b\endcsname{\href{S6_SERIES_GRIND_II.pdf\#nameddest=tag-01FH}{II.21.2 (p.~543)} }
\expandafter\def\csname SRef@cdp:st:C2c\endcsname{\href{S6_SERIES_GRIND_II.pdf\#nameddest=tag-01FI}{II.21.3 (p.~544)} }
\expandafter\def\csname SRef@cdp:st:C.2.UN.11\endcsname{\href{S6_SERIES_GRIND_II.pdf\#nameddest=tag-01CO}{II.21.4 (p.~546)} }
\expandafter\def\csname SRef@cdp:st:C.2.UN.24\endcsname{\href{S6_SERIES_GRIND_II.pdf\#nameddest=tag-01D2}{II.21.5 (p.~550)} }
\expandafter\def\csname SRef@cdp:st:C.2.UN.32\endcsname{\href{S6_SERIES_GRIND_II.pdf\#nameddest=tag-01DB}{II.21.6 (p.~554)} }
\expandafter\def\csname SRef@cdp:st:C.2.UN.64\endcsname{\href{S6_SERIES_GRIND_II.pdf\#nameddest=tag-01EA}{II.21.7 (p.~557)} }
\expandafter\def\csname SRef@cdp:st:C.2.UN.65\endcsname{\href{S6_SERIES_GRIND_II.pdf\#nameddest=tag-01EB}{II.21.8 (p.~561)} }
\expandafter\def\csname SRef@cdp:st:C.2.UN.66\endcsname{\href{S6_SERIES_GRIND_II.pdf\#nameddest=tag-01EC}{II.21.9 (p.~565)} }
\expandafter\def\csname SRef@cdp:st:C.2.UN.6\endcsname{\href{S6_SERIES_GRIND_II.pdf\#nameddest=tag-01E5}{II.22.1 (p.~579)} }
\expandafter\def\csname SRef@cdp:st:C.2.UN.35\endcsname{\href{S6_SERIES_GRIND_II.pdf\#nameddest=tag-01DE}{II.22.10 (p.~613)} }
\expandafter\def\csname SRef@cdp:st:C.2.UN.37\endcsname{\href{S6_SERIES_GRIND_II.pdf\#nameddest=tag-01DG}{II.22.11 (p.~620)} }
\expandafter\def\csname SRef@cdp:st:C.2.UN.41\endcsname{\href{S6_SERIES_GRIND_II.pdf\#nameddest=tag-01DL}{II.22.12 (p.~622)} }
\expandafter\def\csname SRef@cdp:st:C.2.UN.42\endcsname{\href{S6_SERIES_GRIND_II.pdf\#nameddest=tag-01DM}{II.22.13 (p.~626)} }
\expandafter\def\csname SRef@cdp:st:C.2.UN.43\endcsname{\href{S6_SERIES_GRIND_II.pdf\#nameddest=tag-01DN}{II.22.14 (p.~630)} }
\expandafter\def\csname SRef@cdp:st:C.2.UN.45\endcsname{\href{S6_SERIES_GRIND_II.pdf\#nameddest=tag-01DP}{II.22.15 (p.~634)} }
\expandafter\def\csname SRef@cdp:st:C.2.UN.8\endcsname{\href{S6_SERIES_GRIND_II.pdf\#nameddest=tag-01ER}{II.22.2 (p.~583)} }
\expandafter\def\csname SRef@cdp:st:C.2.UN.13\endcsname{\href{S6_SERIES_GRIND_II.pdf\#nameddest=tag-01CQ}{II.22.3 (p.~586)} }
\expandafter\def\csname SRef@cdp:st:C.2.UN.18\endcsname{\href{S6_SERIES_GRIND_II.pdf\#nameddest=tag-01CV}{II.22.4 (p.~590)} }
\expandafter\def\csname SRef@cdp:st:C.2.UN.21\endcsname{\href{S6_SERIES_GRIND_II.pdf\#nameddest=tag-01CZ}{II.22.5 (p.~594)} }
\expandafter\def\csname SRef@cdp:st:C.2.UN.25\endcsname{\href{S6_SERIES_GRIND_II.pdf\#nameddest=tag-01D3}{II.22.6 (p.~598)} }
\expandafter\def\csname SRef@cdp:st:C.2.UN.27\endcsname{\href{S6_SERIES_GRIND_II.pdf\#nameddest=tag-01D5}{II.22.7 (p.~603)} }
\expandafter\def\csname SRef@cdp:st:C.2.UN.30\endcsname{\href{S6_SERIES_GRIND_II.pdf\#nameddest=tag-01D9}{II.22.8 (p.~607)} }
\expandafter\def\csname SRef@cdp:st:C.2.UN.34\endcsname{\href{S6_SERIES_GRIND_II.pdf\#nameddest=tag-01DD}{II.22.9 (p.~610)} }
\expandafter\def\csname SRef@cdp:st:C.2.UN.40\endcsname{\href{S6_SERIES_GRIND_II.pdf\#nameddest=tag-01DK}{II.23.1 (p.~636)} }
\expandafter\def\csname SRef@cdp:st:C.2.UN.55\endcsname{\href{S6_SERIES_GRIND_II.pdf\#nameddest=tag-01E0}{II.23.10 (p.~672)} }
\expandafter\def\csname SRef@cdp:st:C.2.UN.56\endcsname{\href{S6_SERIES_GRIND_II.pdf\#nameddest=tag-01E1}{II.23.11 (p.~675)} }
\expandafter\def\csname SRef@cdp:st:C.2.UN.57\endcsname{\href{S6_SERIES_GRIND_II.pdf\#nameddest=tag-01E2}{II.23.12 (p.~679)} }
\expandafter\def\csname SRef@cdp:st:C.2.UN.58\endcsname{\href{S6_SERIES_GRIND_II.pdf\#nameddest=tag-01E3}{II.23.13 (p.~682)} }
\expandafter\def\csname SRef@cdp:st:C.2.UN.59\endcsname{\href{S6_SERIES_GRIND_II.pdf\#nameddest=tag-01E4}{II.23.14 (p.~686)} }
\expandafter\def\csname SRef@cdp:st:C.2.UN.60\endcsname{\href{S6_SERIES_GRIND_II.pdf\#nameddest=tag-01E6}{II.23.15 (p.~692)} }
\expandafter\def\csname SRef@cdp:st:C.2.UN.61\endcsname{\href{S6_SERIES_GRIND_II.pdf\#nameddest=tag-01E7}{II.23.16 (p.~698)} }
\expandafter\def\csname SRef@cdp:st:C.2.UN.62\endcsname{\href{S6_SERIES_GRIND_II.pdf\#nameddest=tag-01E8}{II.23.17 (p.~703)} }
\expandafter\def\csname SRef@cdp:st:C.2.UN.63\endcsname{\href{S6_SERIES_GRIND_II.pdf\#nameddest=tag-01E9}{II.23.18 (p.~708)} }
\expandafter\def\csname SRef@cdp:st:C.2.UN.46\endcsname{\href{S6_SERIES_GRIND_II.pdf\#nameddest=tag-01DQ}{II.23.2 (p.~638)} }
\expandafter\def\csname SRef@cdp:st:C.2.UN.48\endcsname{\href{S6_SERIES_GRIND_II.pdf\#nameddest=tag-01DS}{II.23.3 (p.~642)} }
\expandafter\def\csname SRef@cdp:st:C.2.UN.49\endcsname{\href{S6_SERIES_GRIND_II.pdf\#nameddest=tag-01DT}{II.23.4 (p.~646)} }
\expandafter\def\csname SRef@cdp:st:C.2.UN.50\endcsname{\href{S6_SERIES_GRIND_II.pdf\#nameddest=tag-01DV}{II.23.5 (p.~652)} }
\expandafter\def\csname SRef@cdp:st:C.2.UN.51\endcsname{\href{S6_SERIES_GRIND_II.pdf\#nameddest=tag-01DW}{II.23.6 (p.~654)} }
\expandafter\def\csname SRef@cdp:st:C.2.UN.52\endcsname{\href{S6_SERIES_GRIND_II.pdf\#nameddest=tag-01DX}{II.23.7 (p.~657)} }
\expandafter\def\csname SRef@cdp:st:C.2.UN.53\endcsname{\href{S6_SERIES_GRIND_II.pdf\#nameddest=tag-01DY}{II.23.8 (p.~664)} }
\expandafter\def\csname SRef@cdp:st:C.2.UN.54\endcsname{\href{S6_SERIES_GRIND_II.pdf\#nameddest=tag-01DZ}{II.23.9 (p.~669)} }
\expandafter\def\csname SRef@cdp:st:C.2.UN.69\endcsname{\href{S6_SERIES_GRIND_II.pdf\#nameddest=tag-01EF}{II.24.1 (p.~719)} }
\expandafter\def\csname SRef@cdp:st:C.2.UN.80\endcsname{\href{S6_SERIES_GRIND_II.pdf\#nameddest=tag-01ES}{II.24.10 (p.~755)} }
\expandafter\def\csname SRef@cdp:st:C.2.UN.81\endcsname{\href{S6_SERIES_GRIND_II.pdf\#nameddest=tag-01ET}{II.24.11 (p.~757)} }
\expandafter\def\csname SRef@cdp:st:C.2.UN.82\endcsname{\href{S6_SERIES_GRIND_II.pdf\#nameddest=tag-01EU}{II.24.12 (p.~762)} }
\expandafter\def\csname SRef@cdp:st:C.2.UN.83\endcsname{\href{S6_SERIES_GRIND_II.pdf\#nameddest=tag-01EV}{II.24.13 (p.~765)} }
\expandafter\def\csname SRef@cdp:st:C.2.UN.84\endcsname{\href{S6_SERIES_GRIND_II.pdf\#nameddest=tag-01EW}{II.24.14 (p.~770)} }
\expandafter\def\csname SRef@cdp:st:C.2.UN.78\endcsname{\href{S6_SERIES_GRIND_II.pdf\#nameddest=tag-01EP}{II.24.15 (p.~773)} }
\expandafter\def\csname SRef@cdp:st:C.2.UN.72\endcsname{\href{S6_SERIES_GRIND_II.pdf\#nameddest=tag-01EJ}{II.24.16 (p.~778)} }
\expandafter\def\csname SRef@cdp:st:C.2.UN.85\endcsname{\href{S6_SERIES_GRIND_II.pdf\#nameddest=tag-01EX}{II.24.17 (p.~782)} }
\expandafter\def\csname SRef@cdp:st:C.2.UN.86\endcsname{\href{S6_SERIES_GRIND_II.pdf\#nameddest=tag-01EY}{II.24.18 (p.~788)} }
\expandafter\def\csname SRef@cdp:st:C.2.UN.70\endcsname{\href{S6_SERIES_GRIND_II.pdf\#nameddest=tag-01EH}{II.24.2 (p.~722)} }
\expandafter\def\csname SRef@cdp:st:C.2.UN.71\endcsname{\href{S6_SERIES_GRIND_II.pdf\#nameddest=tag-01EI}{II.24.3 (p.~726)} }
\expandafter\def\csname SRef@cdp:st:C.2.UN.73\endcsname{\href{S6_SERIES_GRIND_II.pdf\#nameddest=tag-01EK}{II.24.4 (p.~730)} }
\expandafter\def\csname SRef@cdp:st:C.2.UN.76\endcsname{\href{S6_SERIES_GRIND_II.pdf\#nameddest=tag-01EN}{II.24.5 (p.~733)} }
\expandafter\def\csname SRef@cdp:st:C.2.UN.74\endcsname{\href{S6_SERIES_GRIND_II.pdf\#nameddest=tag-01EL}{II.24.6 (p.~739)} }
\expandafter\def\csname SRef@cdp:st:C.2.UN.75\endcsname{\href{S6_SERIES_GRIND_II.pdf\#nameddest=tag-01EM}{II.24.7 (p.~743)} }
\expandafter\def\csname SRef@cdp:st:C.2.UN.77\endcsname{\href{S6_SERIES_GRIND_II.pdf\#nameddest=tag-01EO}{II.24.8 (p.~747)} }
\expandafter\def\csname SRef@cdp:st:C.2.UN.79\endcsname{\href{S6_SERIES_GRIND_II.pdf\#nameddest=tag-01EQ}{II.24.9 (p.~752)} }
\expandafter\def\csname SRef@sec:II:06\endcsname{\href{S6_SERIES_GRIND_II.pdf\#nameddest=tag-0342}{II.19 (p.~413)} }
\expandafter\def\csname SRef@sec:II:07\endcsname{\href{S6_SERIES_GRIND_II.pdf\#nameddest=tag-0343}{II.20 (p.~474)} }
\expandafter\def\csname SRef@sec:II:all:08\endcsname{\href{S6_SERIES_GRIND_II.pdf\#nameddest=tag-034B}{II.21 (p.~543)} }
\expandafter\def\csname SRef@sec:II:09\endcsname{\href{S6_SERIES_GRIND_II.pdf\#nameddest=tag-0345}{II.22 (p.~579)} }
\expandafter\def\csname SRef@sec:II:10\endcsname{\href{S6_SERIES_GRIND_II.pdf\#nameddest=tag-0346}{II.23 (p.~636)} }
\expandafter\def\csname SRef@sec:II:11\endcsname{\href{S6_SERIES_GRIND_II.pdf\#nameddest=tag-0347}{II.24 (p.~719)} }
\expandafter\def\csname SRef@sec:II:all:12\endcsname{\href{S6_SERIES_GRIND_II.pdf\#nameddest=tag-034C}{II.25 (p.~792)} }
\expandafter\def\csname SRef@cdp:st:C.2.G.1\endcsname{\href{S6_SERIES_GRIND_II.pdf\#nameddest=tag-01BZ}{II.25.1 (p.~792)} }
\expandafter\def\csname SRef@inv:E1.1\endcsname{\href{S6_SERIES_GRIND_II.pdf\#nameddest=tag-01ST}{II.? (p.~?)} }
\expandafter\def\csname SRef@inv:E2.1\endcsname{\href{S6_SERIES_GRIND_II.pdf\#nameddest=tag-01SU}{II.? (p.~?)} }
\expandafter\def\csname SRef@push:src:C.1.1\endcsname{\href{S6_SERIES_GRIND_III.pdf\#nameddest=tag-031B}{III.? (p.~?)} }
\expandafter\def\csname SRef@push:src:C.1.2\endcsname{\href{S6_SERIES_GRIND_III.pdf\#nameddest=tag-031C}{III.? (p.~?)} }
\expandafter\def\csname SRef@push:src:C.2.1\endcsname{\href{S6_SERIES_GRIND_III.pdf\#nameddest=tag-031D}{III.? (p.~?)} }
\expandafter\def\csname SRef@push:src:C.3.1\endcsname{\href{S6_SERIES_GRIND_III.pdf\#nameddest=tag-031G}{III.? (p.~?)} }
\expandafter\def\csname SRef@push:src:C.3.3\endcsname{\href{S6_SERIES_GRIND_III.pdf\#nameddest=tag-0320}{III.? (p.~?)} }
\expandafter\def\csname SRef@push:src:E.3.2\endcsname{\hyperlink{tag-033J}{I.\ref*{push:src:E.3.2}}}
\expandafter\def\csname SRef@push:src:E.3.3\endcsname{\hyperlink{tag-033K}{I.\ref*{push:src:E.3.3}}}
\expandafter\def\csname SRef@cdp:st:A.2.3\endcsname{\href{S6_SERIES_GRIND_II.pdf\#nameddest=tag-01AH}{II.? (p.~?)} }
\expandafter\def\csname SRef@cdp:st:C.1.2\endcsname{\href{S6_SERIES_GRIND_II.pdf\#nameddest=tag-01BO}{II.? (p.~?)} }
\expandafter\def\csname SRef@cdp:st:C.2.NC.1\endcsname{\href{S6_SERIES_GRIND_II.pdf\#nameddest=tag-01C8}{II.? (p.~?)} }
\expandafter\def\csname SRef@cdp:st:C.2.RAT.6\endcsname{\href{S6_SERIES_GRIND_II.pdf\#nameddest=tag-01CI}{II.? (p.~?)} }
\expandafter\def\csname SRef@cdp:st:D.1.1\endcsname{\href{S6_SERIES_GRIND_II.pdf\#nameddest=tag-01FU}{II.? (p.~?)} }
\expandafter\def\csname SRef@cdp:st:D.1.3\endcsname{\href{S6_SERIES_GRIND_II.pdf\#nameddest=tag-01G1}{II.? (p.~?)} }
\expandafter\def\csname SRef@cdp:st:E2\endcsname{\href{S6_SERIES_GRIND_II.pdf\#nameddest=tag-01GR}{II.? (p.~?)} }
\expandafter\def\csname SRef@cdp:st:E3\endcsname{\href{S6_SERIES_GRIND_II.pdf\#nameddest=tag-01GS}{II.? (p.~?)} }
\expandafter\def\csname SRef@cdp:st:E4\endcsname{\href{S6_SERIES_GRIND_II.pdf\#nameddest=tag-01GT}{II.? (p.~?)} }
\expandafter\def\csname SRef@cdp:st:E19\endcsname{\href{S6_SERIES_GRIND_II.pdf\#nameddest=tag-01GQ}{II.? (p.~?)} }
\expandafter\def\csname SRef@maps:src:B.6.9\endcsname{\href{S6_SERIES_GRIND_IV.pdf\#nameddest=tag-02K1}{IV.? (p.~?)} }
\expandafter\def\csname SRef@maps:src:D.2.1\endcsname{\href{S6_SERIES_GRIND_III.pdf\#nameddest=tag-02KJ}{III.? (p.~?)} }
\expandafter\def\csname SRef@maps:src:F.1.2\endcsname{\href{S6_SERIES_GRIND_V.pdf\#nameddest=tag-02MV}{V.? (p.~?)} }
\expandafter\def\csname SRef@maps:src:F.1.3\endcsname{\href{S6_SERIES_GRIND_V.pdf\#nameddest=tag-02MW}{V.? (p.~?)} }
\expandafter\def\csname SRef@ari:C.7.5\endcsname{\hyperlink{tag-001E}{I.\ref*{ari:C.7.5}}}
\expandafter\def\csname SRef@wb:lem:triangle-generator-orientations\endcsname{\hyperlink{tag-03BM}{I.\ref*{wb:lem:triangle-generator-orientations}}}
\expandafter\def\csname SRef@wb:prop:period-input-audit\endcsname{\hyperlink{tag-03BS}{I.\ref*{wb:prop:period-input-audit}}}
\expandafter\def\csname SRef@wb:prop:period-exact-admissible-constants\endcsname{\hyperlink{tag-03BR}{I.\ref*{wb:prop:period-exact-admissible-constants}}}
\expandafter\def\csname SRef@wb:lem:toric-chart-control\endcsname{\hyperlink{tag-03BK}{I.\ref*{wb:lem:toric-chart-control}}}
\expandafter\def\csname SRef@wb:thm:candidate-cusp-geometry\endcsname{\hyperlink{tag-03DK}{I.\ref*{wb:thm:candidate-cusp-geometry}}}
\expandafter\def\csname SRef@wb:prop:candidate-cusp-overlap\endcsname{\hyperlink{tag-03BN}{I.\ref*{wb:prop:candidate-cusp-overlap}}}
\expandafter\def\csname SRef@wb:lem:middle-period-proper\endcsname{\hyperlink{tag-03BI}{I.\ref*{wb:lem:middle-period-proper}}}
\expandafter\def\csname SRef@wb:prop:candidate-finite-fillings\endcsname{\hyperlink{tag-03BO}{I.\ref*{wb:prop:candidate-finite-fillings}}}
\expandafter\def\csname SRef@wb:lem:finite-filling-stalk-map\endcsname{\hyperlink{tag-03BG}{I.\ref*{wb:lem:finite-filling-stalk-map}}}
\expandafter\def\csname SRef@wb:thm:candidate-global-geometry\endcsname{\hyperlink{tag-03DL}{I.\ref*{wb:thm:candidate-global-geometry}}}
\expandafter\def\csname SRef@wb:prop:monodromy-certificate\endcsname{\hyperlink{tag-03BQ}{I.\ref*{wb:prop:monodromy-certificate}}}
\expandafter\def\csname SRef@wb:thm:cusp-lattices\endcsname{\hyperlink{tag-03DO}{I.\ref*{wb:thm:cusp-lattices}}}
\expandafter\def\csname SRef@wb:thm:coinvariants\endcsname{\hyperlink{tag-03DN}{I.\ref*{wb:thm:coinvariants}}}
\expandafter\def\csname SRef@wb:prop:finite-indices\endcsname{\hyperlink{tag-03BP}{I.\ref*{wb:prop:finite-indices}}}
\expandafter\def\csname SRef@wb:prop:terminal-group\endcsname{\hyperlink{tag-03BV}{I.\ref*{wb:prop:terminal-group}}}
\expandafter\def\csname SRef@wb:prop:terminal-degree-two\endcsname{\hyperlink{tag-03BU}{I.\ref*{wb:prop:terminal-degree-two}}}
\expandafter\def\csname SRef@wb:prop:terminal-all-degree\endcsname{\hyperlink{tag-03BT}{I.\ref*{wb:prop:terminal-all-degree}}}
\expandafter\def\csname SRef@wb:ffcert:prop:integral-summands\endcsname{\hyperlink{tag-03A8}{I.\ref*{wb:ffcert:prop:integral-summands}}}
\expandafter\def\csname SRef@wb:ffcert:thm:product-covers\endcsname{\hyperlink{tag-03AD}{I.\ref*{wb:ffcert:thm:product-covers}}}
\expandafter\def\csname SRef@wb:ffcert:prop:quotient-charts\endcsname{\hyperlink{tag-03A9}{I.\ref*{wb:ffcert:prop:quotient-charts}}}
\expandafter\def\csname SRef@wb:ffcert:thm:normalization\endcsname{\hyperlink{tag-03AB}{I.\ref*{wb:ffcert:thm:normalization}}}
\expandafter\def\csname SRef@wb:ffcert:thm:peripheral-kernel\endcsname{\hyperlink{tag-03AC}{I.\ref*{wb:ffcert:thm:peripheral-kernel}}}
\expandafter\def\csname SRef@wb:ffcert:cor:covering-h1\endcsname{\hyperlink{tag-03A7}{I.\ref*{wb:ffcert:cor:covering-h1}}}
\expandafter\def\csname SRef@wb:ffcert:thm:normal-trivialization\endcsname{\hyperlink{tag-03AA}{I.\ref*{wb:ffcert:thm:normal-trivialization}}}
\expandafter\def\csname SRef@wb:lem:torus-algebraic-reduction-audit\endcsname{\hyperlink{tag-03BL}{I.\ref*{wb:lem:torus-algebraic-reduction-audit}}}
\expandafter\def\csname SRef@wb:prop:very-general-torus-audit\endcsname{\hyperlink{tag-03BW}{I.\ref*{wb:prop:very-general-torus-audit}}}
\expandafter\def\csname SRef@wb:lem:meromorphic-descent-audit\endcsname{\hyperlink{tag-03BH}{I.\ref*{wb:lem:meromorphic-descent-audit}}}
\expandafter\def\csname SRef@wb:thm:algebraic-dimension-audit\endcsname{\hyperlink{tag-03DJ}{I.\ref*{wb:thm:algebraic-dimension-audit}}}
\expandafter\def\csname SRef@wb:lem:semidirect-abelianization\endcsname{\hyperlink{tag-03BJ}{I.\ref*{wb:lem:semidirect-abelianization}}}
\expandafter\def\csname SRef@mapsI:top:thm:p0-integral-specialization\endcsname{\hyperlink{tag-02Q5}{I.\ref*{mapsI:top:thm:p0-integral-specialization}}}
\expandafter\def\csname SRef@mapsI:top:prop:leray-e2\endcsname{\hyperlink{tag-02PU}{I.\ref*{mapsI:top:prop:leray-e2}}}
\expandafter\def\csname SRef@mapsI:top:lem:first-transgression\endcsname{\hyperlink{tag-02PN}{I.\ref*{mapsI:top:lem:first-transgression}}}
\expandafter\def\csname SRef@mapsI:top:lem:higher-transgressions\endcsname{\hyperlink{tag-02PO}{I.\ref*{mapsI:top:lem:higher-transgressions}}}
\expandafter\def\csname SRef@mapsI:top:thm:h2-vanishing\endcsname{\hyperlink{tag-02Q3}{I.\ref*{mapsI:top:thm:h2-vanishing}}}
\expandafter\def\csname SRef@mapsI:top:prop:geometric-pi1\endcsname{\hyperlink{tag-02PT}{I.\ref*{mapsI:top:prop:geometric-pi1}}}
\expandafter\def\csname SRef@mapsI:top:cor:candidate-sphere\endcsname{\hyperlink{tag-02PK}{I.\ref*{mapsI:top:cor:candidate-sphere}}}
\expandafter\def\csname SRef@mapsI:top:prop:topological-spine\endcsname{\hyperlink{tag-02PV}{I.\ref*{mapsI:top:prop:topological-spine}}}
\expandafter\def\csname SRef@mapsI:top:lem:collapse-straightening\endcsname{\hyperlink{tag-02PM}{I.\ref*{mapsI:top:lem:collapse-straightening}}}
\expandafter\def\csname SRef@mapsI:top:lem:positive-collapse\endcsname{\hyperlink{tag-02PQ}{I.\ref*{mapsI:top:lem:positive-collapse}}}
\expandafter\def\csname SRef@mapsI:top:lem:phase-spread-collapse\endcsname{\hyperlink{tag-02PP}{I.\ref*{mapsI:top:lem:phase-spread-collapse}}}
\expandafter\def\csname SRef@mapsI:top:prop:verified-cellular-collapse\endcsname{\hyperlink{tag-02PW}{I.\ref*{mapsI:top:prop:verified-cellular-collapse}}}
\expandafter\def\csname SRef@mapsI:top:prop:verified-collapse-homology\endcsname{\hyperlink{tag-02PX}{I.\ref*{mapsI:top:prop:verified-collapse-homology}}}
\expandafter\def\csname SRef@mapsI:top:thm:collapse-equals-specialization\endcsname{\hyperlink{tag-02Q2}{I.\ref*{mapsI:top:thm:collapse-equals-specialization}}}
\expandafter\def\csname SRef@mapsI:top:remark:L600\endcsname{\hyperlink{tag-02Q0}{I.\ref*{mapsI:top:remark:L600}}}
\expandafter\def\csname SRef@mapsI:top:record:vankampen\endcsname{\hyperlink{tag-02PY}{I.\ref*{mapsI:top:record:vankampen}}}
\expandafter\def\csname SRef@mapsI:top:prop:actual-finite-band-maps\endcsname{\hyperlink{tag-02PR}{I.\ref*{mapsI:top:prop:actual-finite-band-maps}}}
\expandafter\def\csname SRef@mapsI:top:prop:cusp-mv-blocks\endcsname{\hyperlink{tag-02PS}{I.\ref*{mapsI:top:prop:cusp-mv-blocks}}}
\expandafter\def\csname SRef@mapsI:top:remark:L920\endcsname{\hyperlink{tag-02Q1}{I.\ref*{mapsI:top:remark:L920}}}
\expandafter\def\csname SRef@mapsI:top:thm:mv-leray-comparison\endcsname{\hyperlink{tag-02Q4}{I.\ref*{mapsI:top:thm:mv-leray-comparison}}}
\expandafter\def\csname SRef@mapsI:top:remark:L1113\endcsname{\hyperlink{tag-02PZ}{I.\ref*{mapsI:top:remark:L1113}}}
\expandafter\def\csname SRef@mapsI:top:lem:cohomological-closure\endcsname{\hyperlink{tag-02PL}{I.\ref*{mapsI:top:lem:cohomological-closure}}}
\expandafter\def\csname SRef@mapsI:top:thm:sphere-recognition\endcsname{\hyperlink{tag-02Q6}{I.\ref*{mapsI:top:thm:sphere-recognition}}}
\expandafter\def\csname SRef@mapsI:top:cor:abstract-s6\endcsname{\hyperlink{tag-02PI}{I.\ref*{mapsI:top:cor:abstract-s6}}}
\expandafter\def\csname SRef@mapsI:top:cor:candidate-mv-sphere\endcsname{\hyperlink{tag-02PJ}{I.\ref*{mapsI:top:cor:candidate-mv-sphere}}}
\expandafter\def\csname SRef@wb:cusp-cone-S12-almost-complex-obstruction\endcsname{\href{S6_SERIES_GRIND_III.pdf\#nameddest=tag-035S}{III.? (p.~?)} }
\expandafter\def\csname SRef@courseIII:almost-complex-structures-and-integrability:section-1\endcsname{\href{S6_SERIES_GRIND_III.pdf\#nameddest=tag-01I4}{III.? (p.~?)} }
\expandafter\def\csname SRef@courseIII:almost-complex-structures-and-integrability:section-2\endcsname{\href{S6_SERIES_GRIND_III.pdf\#nameddest=tag-01I7}{III.? (p.~?)} }
\expandafter\def\csname SRef@courseIII:almost-complex-structures-and-integrability:section-3\endcsname{\href{S6_SERIES_GRIND_III.pdf\#nameddest=tag-01I8}{III.? (p.~?)} }
\expandafter\def\csname SRef@courseIII:almost-complex-structures-and-integrability:section-4\endcsname{\href{S6_SERIES_GRIND_III.pdf\#nameddest=tag-01I9}{III.? (p.~?)} }
\expandafter\def\csname SRef@courseIII:almost-complex-structures-and-integrability:section-5\endcsname{\href{S6_SERIES_GRIND_III.pdf\#nameddest=tag-01IA}{III.? (p.~?)} }
\expandafter\def\csname SRef@courseIII:almost-complex-structures-and-integrability:section-6\endcsname{\href{S6_SERIES_GRIND_III.pdf\#nameddest=tag-01IB}{III.? (p.~?)} }
\expandafter\def\csname SRef@courseIII:almost-complex-structures-and-integrability:section-7\endcsname{\href{S6_SERIES_GRIND_III.pdf\#nameddest=tag-01IC}{III.? (p.~?)} }
\expandafter\def\csname SRef@courseIII:almost-complex-structures-and-integrability:section-8\endcsname{\href{S6_SERIES_GRIND_III.pdf\#nameddest=tag-01ID}{III.? (p.~?)} }
\expandafter\def\csname SRef@courseIII:almost-complex-structures-and-integrability:section-9\endcsname{\href{S6_SERIES_GRIND_III.pdf\#nameddest=tag-01IE}{III.? (p.~?)} }
\expandafter\def\csname SRef@courseIII:almost-complex-structures-and-integrability:section-10\endcsname{\href{S6_SERIES_GRIND_III.pdf\#nameddest=tag-01I5}{III.? (p.~?)} }
\expandafter\def\csname SRef@courseIII:almost-complex-structures-and-integrability:section-11\endcsname{\href{S6_SERIES_GRIND_III.pdf\#nameddest=tag-01I6}{III.? (p.~?)} }
\expandafter\def\csname SRef@mapsI:course:CG-S6-07-HOMOTOPY\endcsname{\hyperlink{tag-02OR}{I.\ref*{mapsI:course:CG-S6-07-HOMOTOPY}}}
\expandafter\def\csname SRef@courseIII:integral-homology-and-sphere-recognition:section\endcsname{\href{S6_SERIES_GRIND_III.pdf\#nameddest=tag-01K9}{III.? (p.~?)} }
\expandafter\def\csname SRef@courseIII:sphere-twists-and-the-original-smooth-six-sphere-comparison:twist-actual-tube\endcsname{\href{S6_SERIES_GRIND_III.pdf\#nameddest=tag-01MN}{III.? (p.~?)} }
\expandafter\def\csname SRef@courseIII:sphere-twists-and-the-original-smooth-six-sphere-comparison:twist-signed-action\endcsname{\href{S6_SERIES_GRIND_III.pdf\#nameddest=tag-01MT}{III.? (p.~?)} }
\expandafter\def\csname SRef@courseIII:sphere-twists-and-the-original-smooth-six-sphere-comparison:twist-integral-lagrangians\endcsname{\href{S6_SERIES_GRIND_III.pdf\#nameddest=tag-01MP}{III.? (p.~?)} }
\expandafter\def\csname SRef@courseIII:sphere-twists-and-the-original-smooth-six-sphere-comparison:twist-original-middle-level\endcsname{\href{S6_SERIES_GRIND_III.pdf\#nameddest=tag-01MR}{III.? (p.~?)} }
\expandafter\def\csname SRef@courseIII:sphere-twists-and-the-original-smooth-six-sphere-comparison:twist-original-gluing\endcsname{\href{S6_SERIES_GRIND_III.pdf\#nameddest=tag-01MQ}{III.? (p.~?)} }
\expandafter\def\csname SRef@courseIII:sphere-twists-and-the-original-smooth-six-sphere-comparison:twist-actual-cancellation\endcsname{\href{S6_SERIES_GRIND_III.pdf\#nameddest=tag-01MM}{III.? (p.~?)} }
\expandafter\def\csname SRef@courseIII:sphere-twists-and-the-original-smooth-six-sphere-comparison:twist-original-smooth-map\endcsname{\href{S6_SERIES_GRIND_III.pdf\#nameddest=tag-01MS}{III.? (p.~?)} }
\expandafter\def\csname SRef@courseIII:sphere-twists-and-the-original-smooth-six-sphere-comparison:twist-general-stable-framing\endcsname{\href{S6_SERIES_GRIND_III.pdf\#nameddest=tag-01MO}{III.? (p.~?)} }
\expandafter\def\csname SRef@courseIII:sphere-twists-and-the-original-smooth-six-sphere-comparison:twist-worked-exercises\endcsname{\href{S6_SERIES_GRIND_III.pdf\#nameddest=tag-01MU}{III.? (p.~?)} }
\expandafter\def\csname SRef@mapsaudit:almost-complex-homotopy\endcsname{\href{S6_SERIES_GRIND_III.pdf\#nameddest=tag-02QD}{III.? (p.~?)} }
\expandafter\def\csname SRef@mapsI:hodge:line\endcsname{\hyperlink{tag-02OV}{I.\ref*{mapsI:hodge:line}}}
\expandafter\def\csname SRef@mapsI:chain:lemma:1.1\endcsname{\hyperlink{tag-02OL}{I.\ref*{mapsI:chain:lemma:1.1}}}
\expandafter\def\csname SRef@mapsI:chain:lemma:2.1\endcsname{\hyperlink{tag-02OM}{I.\ref*{mapsI:chain:lemma:2.1}}}
\expandafter\def\csname SRef@mapsI:chain:corollary:2.2\endcsname{\hyperlink{tag-02OK}{I.\ref*{mapsI:chain:corollary:2.2}}}
\expandafter\def\csname SRef@mapsI:chain:lemma:4.1\endcsname{\hyperlink{tag-02ON}{I.\ref*{mapsI:chain:lemma:4.1}}}
\expandafter\def\csname SRef@mapsI:cw:hurewicz-theorem-result\endcsname{\hyperlink{tag-02OU}{I.\ref*{mapsI:cw:hurewicz-theorem-result}}}
\expandafter\def\csname SRef@mapsI:cw:exercise:5.1\endcsname{\hyperlink{tag-02OS}{I.\ref*{mapsI:cw:exercise:5.1}}}
\expandafter\def\csname SRef@mapsI:cw:exercise:5.2\endcsname{\hyperlink{tag-02OT}{I.\ref*{mapsI:cw:exercise:5.2}}}
\expandafter\def\csname SRef@mapsI:pf:theorem:5.1\endcsname{\hyperlink{tag-02OY}{I.\ref*{mapsI:pf:theorem:5.1}}}
\expandafter\def\csname SRef@mapsI:pf:exercise:6.1\endcsname{\hyperlink{tag-02OW}{I.\ref*{mapsI:pf:exercise:6.1}}}
\expandafter\def\csname SRef@mapsI:pf:exercise:6.2\endcsname{\hyperlink{tag-02OX}{I.\ref*{mapsI:pf:exercise:6.2}}}
\expandafter\def\csname SRef@courseIII:thom-euler-for-lesson11:CG-S6-CHAIN-LEMMA-5.1\endcsname{\href{S6_SERIES_GRIND_III.pdf\#nameddest=tag-01N3}{III.? (p.~?)} }
\expandafter\def\csname SRef@courseIII:thom-euler-for-lesson11:CG-S6-CHAIN-THEOREM-5.2\endcsname{\href{S6_SERIES_GRIND_III.pdf\#nameddest=tag-01N7}{III.? (p.~?)} }
\expandafter\def\csname SRef@courseIII:thom-euler-for-lesson11:CG-S6-CHAIN-PROPOSITION-6.1\endcsname{\href{S6_SERIES_GRIND_III.pdf\#nameddest=tag-01N4}{III.? (p.~?)} }
\expandafter\def\csname SRef@courseIII:thom-euler-for-lesson11:CG-S6-CHAIN-THEOREM-7.1\endcsname{\href{S6_SERIES_GRIND_III.pdf\#nameddest=tag-01N8}{III.? (p.~?)} }
\expandafter\def\csname SRef@courseIII:thom-euler-for-lesson11:CG-S6-CHAIN-PROPOSITION-8.1\endcsname{\href{S6_SERIES_GRIND_III.pdf\#nameddest=tag-01N5}{III.? (p.~?)} }
\expandafter\def\csname SRef@courseIII:thom-euler-for-lesson11:CG-S6-CHAIN-PROPOSITION-8.2\endcsname{\href{S6_SERIES_GRIND_III.pdf\#nameddest=tag-01N6}{III.? (p.~?)} }
\expandafter\def\csname SRef@courseIII:projective-splitting-for-lesson11:CG-S6-PROJECTIVE-1\endcsname{\href{S6_SERIES_GRIND_III.pdf\#nameddest=tag-01M7}{III.? (p.~?)} }
\expandafter\def\csname SRef@courseIII:projective-splitting-for-lesson11:CG-S6-PROJECTIVE-2\endcsname{\href{S6_SERIES_GRIND_III.pdf\#nameddest=tag-01M8}{III.? (p.~?)} }
\expandafter\def\csname SRef@courseIII:projective-splitting-for-lesson11:CG-S6-PROJECTIVE-3\endcsname{\href{S6_SERIES_GRIND_III.pdf\#nameddest=tag-01M9}{III.? (p.~?)} }
\expandafter\def\csname SRef@courseIII:projective-splitting-for-lesson11:CG-S6-PROJECTIVE-4\endcsname{\href{S6_SERIES_GRIND_III.pdf\#nameddest=tag-01MA}{III.? (p.~?)} }
\expandafter\def\csname SRef@courseIII:projective-splitting-for-lesson11:CG-S6-PROJECTIVE-5\endcsname{\href{S6_SERIES_GRIND_III.pdf\#nameddest=tag-01MB}{III.? (p.~?)} }
\expandafter\def\csname SRef@courseIII:integral-chern-classes-for-lesson11:CG-S6-CHERN-1\endcsname{\href{S6_SERIES_GRIND_III.pdf\#nameddest=tag-01K6}{III.? (p.~?)} }
\expandafter\def\csname SRef@courseIII:integral-chern-classes-for-lesson11:CG-S6-CHERN-2\endcsname{\href{S6_SERIES_GRIND_III.pdf\#nameddest=tag-01K7}{III.? (p.~?)} }
\expandafter\def\csname SRef@courseIII:integral-chern-classes-for-lesson11:CG-S6-CHERN-3\endcsname{\href{S6_SERIES_GRIND_III.pdf\#nameddest=tag-01K8}{III.? (p.~?)} }
\expandafter\def\csname SRef@courseIII:framed6:1\endcsname{\href{S6_SERIES_GRIND_III.pdf\#nameddest=tag-01J8}{III.? (p.~?)} }
\expandafter\def\csname SRef@courseIII:framed6:2\endcsname{\href{S6_SERIES_GRIND_III.pdf\#nameddest=tag-01JB}{III.? (p.~?)} }
\expandafter\def\csname SRef@courseIII:framed6:3\endcsname{\href{S6_SERIES_GRIND_III.pdf\#nameddest=tag-01JC}{III.? (p.~?)} }
\expandafter\def\csname SRef@courseIII:framed6:4\endcsname{\href{S6_SERIES_GRIND_III.pdf\#nameddest=tag-01JD}{III.? (p.~?)} }
\expandafter\def\csname SRef@courseIII:framed6:5\endcsname{\href{S6_SERIES_GRIND_III.pdf\#nameddest=tag-01JE}{III.? (p.~?)} }
\expandafter\def\csname SRef@courseIII:framed6:6\endcsname{\href{S6_SERIES_GRIND_III.pdf\#nameddest=tag-01JF}{III.? (p.~?)} }
\expandafter\def\csname SRef@courseIII:framed6:7\endcsname{\href{S6_SERIES_GRIND_III.pdf\#nameddest=tag-01JG}{III.? (p.~?)} }
\expandafter\def\csname SRef@courseIII:framed6:8\endcsname{\href{S6_SERIES_GRIND_III.pdf\#nameddest=tag-01JH}{III.? (p.~?)} }
\expandafter\def\csname SRef@courseIII:framed6:9\endcsname{\href{S6_SERIES_GRIND_III.pdf\#nameddest=tag-01JI}{III.? (p.~?)} }
\expandafter\def\csname SRef@courseIII:framed6:10.1\endcsname{\href{S6_SERIES_GRIND_III.pdf\#nameddest=tag-01J9}{III.? (p.~?)} }
\expandafter\def\csname SRef@courseIII:framed6:10.2\endcsname{\href{S6_SERIES_GRIND_III.pdf\#nameddest=tag-01JA}{III.? (p.~?)} }
\expandafter\def\csname SRef@courseIII:whitney-receiving:1\endcsname{\href{S6_SERIES_GRIND_III.pdf\#nameddest=tag-01NG}{III.? (p.~?)} }
\expandafter\def\csname SRef@courseIII:whitney-receiving:2\endcsname{\href{S6_SERIES_GRIND_III.pdf\#nameddest=tag-01NH}{III.? (p.~?)} }
\expandafter\def\csname SRef@courseIII:whitney-receiving:3\endcsname{\href{S6_SERIES_GRIND_III.pdf\#nameddest=tag-01NI}{III.? (p.~?)} }
\expandafter\def\csname SRef@courseIII:whitney-receiving:4\endcsname{\href{S6_SERIES_GRIND_III.pdf\#nameddest=tag-01NJ}{III.? (p.~?)} }
\expandafter\def\csname SRef@courseIII:whitney-receiving:5.1\endcsname{\href{S6_SERIES_GRIND_III.pdf\#nameddest=tag-01NK}{III.? (p.~?)} }
\expandafter\def\csname SRef@courseIII:whitney-receiving:5.2\endcsname{\href{S6_SERIES_GRIND_III.pdf\#nameddest=tag-01NL}{III.? (p.~?)} }
\expandafter\def\csname SRef@courseIII:higher-slides:1\endcsname{\href{S6_SERIES_GRIND_III.pdf\#nameddest=tag-01JY}{III.? (p.~?)} }
\expandafter\def\csname SRef@courseIII:higher-slides:2\endcsname{\href{S6_SERIES_GRIND_III.pdf\#nameddest=tag-01JZ}{III.? (p.~?)} }
\expandafter\def\csname SRef@courseIII:higher-slides:3\endcsname{\href{S6_SERIES_GRIND_III.pdf\#nameddest=tag-01K0}{III.? (p.~?)} }
\expandafter\def\csname SRef@courseIII:higher-slides:4\endcsname{\href{S6_SERIES_GRIND_III.pdf\#nameddest=tag-01K1}{III.? (p.~?)} }
\expandafter\def\csname SRef@courseIII:higher-slides:5\endcsname{\href{S6_SERIES_GRIND_III.pdf\#nameddest=tag-01K2}{III.? (p.~?)} }
\expandafter\def\csname SRef@courseIII:higher-slides:6\endcsname{\href{S6_SERIES_GRIND_III.pdf\#nameddest=tag-01K3}{III.? (p.~?)} }
\expandafter\def\csname SRef@courseIII:higher-slides:8.1\endcsname{\href{S6_SERIES_GRIND_III.pdf\#nameddest=tag-01K4}{III.? (p.~?)} }
\expandafter\def\csname SRef@courseIII:higher-slides:8.2\endcsname{\href{S6_SERIES_GRIND_III.pdf\#nameddest=tag-01K5}{III.? (p.~?)} }
\expandafter\def\csname SRef@courseIII:middle-cancel:1\endcsname{\href{S6_SERIES_GRIND_III.pdf\#nameddest=tag-01KT}{III.? (p.~?)} }
\expandafter\def\csname SRef@courseIII:middle-cancel:2\endcsname{\href{S6_SERIES_GRIND_III.pdf\#nameddest=tag-01KU}{III.? (p.~?)} }
\expandafter\def\csname SRef@courseIII:middle-cancel:3\endcsname{\href{S6_SERIES_GRIND_III.pdf\#nameddest=tag-01KV}{III.? (p.~?)} }
\expandafter\def\csname SRef@courseIII:middle-cancel:4\endcsname{\href{S6_SERIES_GRIND_III.pdf\#nameddest=tag-01KW}{III.? (p.~?)} }
\expandafter\def\csname SRef@courseIII:middle-cancel:5\endcsname{\href{S6_SERIES_GRIND_III.pdf\#nameddest=tag-01KX}{III.? (p.~?)} }
\expandafter\def\csname SRef@courseIII:middle-cancel:6\endcsname{\href{S6_SERIES_GRIND_III.pdf\#nameddest=tag-01KY}{III.? (p.~?)} }
\expandafter\def\csname SRef@courseIII:middle-cancel:7.1\endcsname{\href{S6_SERIES_GRIND_III.pdf\#nameddest=tag-01KZ}{III.? (p.~?)} }
\expandafter\def\csname SRef@courseIII:middle-cancel:7.2\endcsname{\href{S6_SERIES_GRIND_III.pdf\#nameddest=tag-01L0}{III.? (p.~?)} }
\expandafter\def\csname SRef@courseIII:morse-cancel:1\endcsname{\href{S6_SERIES_GRIND_III.pdf\#nameddest=tag-01L5}{III.? (p.~?)} }
\expandafter\def\csname SRef@courseIII:morse-cancel:2.1\endcsname{\href{S6_SERIES_GRIND_III.pdf\#nameddest=tag-01L6}{III.? (p.~?)} }
\expandafter\def\csname SRef@courseIII:morse-cancel:6\endcsname{\href{S6_SERIES_GRIND_III.pdf\#nameddest=tag-01L7}{III.? (p.~?)} }
\expandafter\def\csname SRef@courseIII:morse-cancel:7.1\endcsname{\href{S6_SERIES_GRIND_III.pdf\#nameddest=tag-01L8}{III.? (p.~?)} }
\expandafter\def\csname SRef@courseIII:morse-cancel:7.2\endcsname{\href{S6_SERIES_GRIND_III.pdf\#nameddest=tag-01L9}{III.? (p.~?)} }
\expandafter\def\csname SRef@courseIII:spin6:1\endcsname{\href{S6_SERIES_GRIND_III.pdf\#nameddest=tag-01MV}{III.? (p.~?)} }
\expandafter\def\csname SRef@courseIII:spin6:2\endcsname{\href{S6_SERIES_GRIND_III.pdf\#nameddest=tag-01MW}{III.? (p.~?)} }
\expandafter\def\csname SRef@courseIII:spin6:3\endcsname{\href{S6_SERIES_GRIND_III.pdf\#nameddest=tag-01MX}{III.? (p.~?)} }
\expandafter\def\csname SRef@courseIII:spin6:4\endcsname{\href{S6_SERIES_GRIND_III.pdf\#nameddest=tag-01MY}{III.? (p.~?)} }
\expandafter\def\csname SRef@courseIII:spin6:5\endcsname{\href{S6_SERIES_GRIND_III.pdf\#nameddest=tag-01MZ}{III.? (p.~?)} }
\expandafter\def\csname SRef@courseIII:spin6:6\endcsname{\href{S6_SERIES_GRIND_III.pdf\#nameddest=tag-01N0}{III.? (p.~?)} }
\expandafter\def\csname SRef@courseIII:spin6:7.1\endcsname{\href{S6_SERIES_GRIND_III.pdf\#nameddest=tag-01N1}{III.? (p.~?)} }
\expandafter\def\csname SRef@courseIII:spin6:7.2\endcsname{\href{S6_SERIES_GRIND_III.pdf\#nameddest=tag-01N2}{III.? (p.~?)} }
\expandafter\def\csname SRef@courseIII:integral-stable:1\endcsname{\href{S6_SERIES_GRIND_III.pdf\#nameddest=tag-01KA}{III.? (p.~?)} }
\expandafter\def\csname SRef@courseIII:integral-stable:2\endcsname{\href{S6_SERIES_GRIND_III.pdf\#nameddest=tag-01KB}{III.? (p.~?)} }
\expandafter\def\csname SRef@courseIII:integral-stable:3\endcsname{\href{S6_SERIES_GRIND_III.pdf\#nameddest=tag-01KC}{III.? (p.~?)} }
\expandafter\def\csname SRef@courseIII:integral-stable:4\endcsname{\href{S6_SERIES_GRIND_III.pdf\#nameddest=tag-01KD}{III.? (p.~?)} }
\expandafter\def\csname SRef@courseIII:integral-stable:5\endcsname{\href{S6_SERIES_GRIND_III.pdf\#nameddest=tag-01KE}{III.? (p.~?)} }
\expandafter\def\csname SRef@courseIII:integral-stable:6\endcsname{\href{S6_SERIES_GRIND_III.pdf\#nameddest=tag-01KF}{III.? (p.~?)} }
\expandafter\def\csname SRef@courseIII:integral-stable:7\endcsname{\href{S6_SERIES_GRIND_III.pdf\#nameddest=tag-01KG}{III.? (p.~?)} }
\expandafter\def\csname SRef@courseIII:integral-stable:8.1\endcsname{\href{S6_SERIES_GRIND_III.pdf\#nameddest=tag-01KH}{III.? (p.~?)} }
\expandafter\def\csname SRef@courseIII:integral-stable:8.2\endcsname{\href{S6_SERIES_GRIND_III.pdf\#nameddest=tag-01KI}{III.? (p.~?)} }
\expandafter\def\csname SRef@courseIII:pontryagin-thom:1\endcsname{\href{S6_SERIES_GRIND_III.pdf\#nameddest=tag-01LQ}{III.? (p.~?)} }
\expandafter\def\csname SRef@courseIII:pontryagin-thom:2\endcsname{\href{S6_SERIES_GRIND_III.pdf\#nameddest=tag-01LR}{III.? (p.~?)} }
\expandafter\def\csname SRef@courseIII:pontryagin-thom:3\endcsname{\href{S6_SERIES_GRIND_III.pdf\#nameddest=tag-01LS}{III.? (p.~?)} }
\expandafter\def\csname SRef@courseIII:pontryagin-thom:4\endcsname{\href{S6_SERIES_GRIND_III.pdf\#nameddest=tag-01LT}{III.? (p.~?)} }
\expandafter\def\csname SRef@courseIII:pontryagin-thom:5\endcsname{\href{S6_SERIES_GRIND_III.pdf\#nameddest=tag-01LU}{III.? (p.~?)} }
\expandafter\def\csname SRef@courseIII:pontryagin-thom:6\endcsname{\href{S6_SERIES_GRIND_III.pdf\#nameddest=tag-01LV}{III.? (p.~?)} }
\expandafter\def\csname SRef@courseIII:pontryagin-thom:7\endcsname{\href{S6_SERIES_GRIND_III.pdf\#nameddest=tag-01LW}{III.? (p.~?)} }
\expandafter\def\csname SRef@courseIII:pontryagin-thom:8\endcsname{\href{S6_SERIES_GRIND_III.pdf\#nameddest=tag-01LX}{III.? (p.~?)} }
\expandafter\def\csname SRef@courseIII:pontryagin-thom:9.1\endcsname{\href{S6_SERIES_GRIND_III.pdf\#nameddest=tag-01LY}{III.? (p.~?)} }
\expandafter\def\csname SRef@courseIII:pontryagin-thom:9.2\endcsname{\href{S6_SERIES_GRIND_III.pdf\#nameddest=tag-01LZ}{III.? (p.~?)} }
\expandafter\def\csname SRef@courseIII:relative-morse:1\endcsname{\href{S6_SERIES_GRIND_III.pdf\#nameddest=tag-01MG}{III.? (p.~?)} }
\expandafter\def\csname SRef@courseIII:relative-morse:2\endcsname{\href{S6_SERIES_GRIND_III.pdf\#nameddest=tag-01MH}{III.? (p.~?)} }
\expandafter\def\csname SRef@courseIII:relative-morse:3\endcsname{\href{S6_SERIES_GRIND_III.pdf\#nameddest=tag-01MI}{III.? (p.~?)} }
\expandafter\def\csname SRef@courseIII:relative-morse:4\endcsname{\href{S6_SERIES_GRIND_III.pdf\#nameddest=tag-01MJ}{III.? (p.~?)} }
\expandafter\def\csname SRef@courseIII:relative-morse:5.1\endcsname{\href{S6_SERIES_GRIND_III.pdf\#nameddest=tag-01MK}{III.? (p.~?)} }
\expandafter\def\csname SRef@courseIII:relative-morse:5.2\endcsname{\href{S6_SERIES_GRIND_III.pdf\#nameddest=tag-01ML}{III.? (p.~?)} }
\expandafter\def\csname SRef@courseIII:handle-arrangement:1\endcsname{\href{S6_SERIES_GRIND_III.pdf\#nameddest=tag-01JJ}{III.? (p.~?)} }
\expandafter\def\csname SRef@courseIII:handle-arrangement:2\endcsname{\href{S6_SERIES_GRIND_III.pdf\#nameddest=tag-01JK}{III.? (p.~?)} }
\expandafter\def\csname SRef@courseIII:handle-arrangement:3\endcsname{\href{S6_SERIES_GRIND_III.pdf\#nameddest=tag-01JL}{III.? (p.~?)} }
\expandafter\def\csname SRef@courseIII:handle-arrangement:4\endcsname{\href{S6_SERIES_GRIND_III.pdf\#nameddest=tag-01JM}{III.? (p.~?)} }
\expandafter\def\csname SRef@courseIII:handle-arrangement:5.1\endcsname{\href{S6_SERIES_GRIND_III.pdf\#nameddest=tag-01JN}{III.? (p.~?)} }
\expandafter\def\csname SRef@courseIII:handle-arrangement:5.2\endcsname{\href{S6_SERIES_GRIND_III.pdf\#nameddest=tag-01JO}{III.? (p.~?)} }
\expandafter\def\csname SRef@courseIII:zero-top:1\endcsname{\href{S6_SERIES_GRIND_III.pdf\#nameddest=tag-01NU}{III.? (p.~?)} }
\expandafter\def\csname SRef@courseIII:zero-top:2\endcsname{\href{S6_SERIES_GRIND_III.pdf\#nameddest=tag-01NV}{III.? (p.~?)} }
\expandafter\def\csname SRef@courseIII:zero-top:3\endcsname{\href{S6_SERIES_GRIND_III.pdf\#nameddest=tag-01NW}{III.? (p.~?)} }
\expandafter\def\csname SRef@courseIII:zero-top:4\endcsname{\href{S6_SERIES_GRIND_III.pdf\#nameddest=tag-01NX}{III.? (p.~?)} }
\expandafter\def\csname SRef@courseIII:zero-top:5\endcsname{\href{S6_SERIES_GRIND_III.pdf\#nameddest=tag-01NY}{III.? (p.~?)} }
\expandafter\def\csname SRef@courseIII:zero-top:6.1\endcsname{\href{S6_SERIES_GRIND_III.pdf\#nameddest=tag-01NZ}{III.? (p.~?)} }
\expandafter\def\csname SRef@courseIII:zero-top:6.2\endcsname{\href{S6_SERIES_GRIND_III.pdf\#nameddest=tag-01O0}{III.? (p.~?)} }
\expandafter\def\csname SRef@courseIII:trajectory-lowering:1\endcsname{\href{S6_SERIES_GRIND_III.pdf\#nameddest=tag-01N9}{III.? (p.~?)} }
\expandafter\def\csname SRef@courseIII:trajectory-lowering:2\endcsname{\href{S6_SERIES_GRIND_III.pdf\#nameddest=tag-01NA}{III.? (p.~?)} }
\expandafter\def\csname SRef@courseIII:trajectory-lowering:3\endcsname{\href{S6_SERIES_GRIND_III.pdf\#nameddest=tag-01NB}{III.? (p.~?)} }
\expandafter\def\csname SRef@courseIII:trajectory-lowering:4\endcsname{\href{S6_SERIES_GRIND_III.pdf\#nameddest=tag-01NC}{III.? (p.~?)} }
\expandafter\def\csname SRef@courseIII:trajectory-lowering:5\endcsname{\href{S6_SERIES_GRIND_III.pdf\#nameddest=tag-01ND}{III.? (p.~?)} }
\expandafter\def\csname SRef@courseIII:trajectory-lowering:6.1\endcsname{\href{S6_SERIES_GRIND_III.pdf\#nameddest=tag-01NE}{III.? (p.~?)} }
\expandafter\def\csname SRef@courseIII:trajectory-lowering:6.2\endcsname{\href{S6_SERIES_GRIND_III.pdf\#nameddest=tag-01NF}{III.? (p.~?)} }
\expandafter\def\csname SRef@courseIII:belt-complement:1\endcsname{\href{S6_SERIES_GRIND_III.pdf\#nameddest=tag-01IF}{III.? (p.~?)} }
\expandafter\def\csname SRef@courseIII:belt-complement:2\endcsname{\href{S6_SERIES_GRIND_III.pdf\#nameddest=tag-01IG}{III.? (p.~?)} }
\expandafter\def\csname SRef@courseIII:belt-complement:3\endcsname{\href{S6_SERIES_GRIND_III.pdf\#nameddest=tag-01IH}{III.? (p.~?)} }
\expandafter\def\csname SRef@courseIII:belt-complement:4\endcsname{\href{S6_SERIES_GRIND_III.pdf\#nameddest=tag-01II}{III.? (p.~?)} }
\expandafter\def\csname SRef@courseIII:belt-complement:5\endcsname{\href{S6_SERIES_GRIND_III.pdf\#nameddest=tag-01IJ}{III.? (p.~?)} }
\expandafter\def\csname SRef@courseIII:belt-complement:6\endcsname{\href{S6_SERIES_GRIND_III.pdf\#nameddest=tag-01IK}{III.? (p.~?)} }
\expandafter\def\csname SRef@courseIII:belt-complement:7\endcsname{\href{S6_SERIES_GRIND_III.pdf\#nameddest=tag-01IL}{III.? (p.~?)} }
\expandafter\def\csname SRef@courseIII:whitney-support:1\endcsname{\href{S6_SERIES_GRIND_III.pdf\#nameddest=tag-01NM}{III.? (p.~?)} }
\expandafter\def\csname SRef@courseIII:whitney-support:2\endcsname{\href{S6_SERIES_GRIND_III.pdf\#nameddest=tag-01NN}{III.? (p.~?)} }
\expandafter\def\csname SRef@courseIII:whitney-support:3\endcsname{\href{S6_SERIES_GRIND_III.pdf\#nameddest=tag-01NO}{III.? (p.~?)} }
\expandafter\def\csname SRef@courseIII:whitney-support:4\endcsname{\href{S6_SERIES_GRIND_III.pdf\#nameddest=tag-01NP}{III.? (p.~?)} }
\expandafter\def\csname SRef@courseIII:whitney-support:5\endcsname{\href{S6_SERIES_GRIND_III.pdf\#nameddest=tag-01NQ}{III.? (p.~?)} }
\expandafter\def\csname SRef@courseIII:whitney-support:6\endcsname{\href{S6_SERIES_GRIND_III.pdf\#nameddest=tag-01NR}{III.? (p.~?)} }
\expandafter\def\csname SRef@courseIII:whitney-support:7.1\endcsname{\href{S6_SERIES_GRIND_III.pdf\#nameddest=tag-01NS}{III.? (p.~?)} }
\expandafter\def\csname SRef@courseIII:whitney-support:7.2\endcsname{\href{S6_SERIES_GRIND_III.pdf\#nameddest=tag-01NT}{III.? (p.~?)} }
\expandafter\def\csname SRef@courseIII:handle-chains:1\endcsname{\href{S6_SERIES_GRIND_III.pdf\#nameddest=tag-01JP}{III.? (p.~?)} }
\expandafter\def\csname SRef@courseIII:handle-chains:2\endcsname{\href{S6_SERIES_GRIND_III.pdf\#nameddest=tag-01JQ}{III.? (p.~?)} }
\expandafter\def\csname SRef@courseIII:handle-chains:3\endcsname{\href{S6_SERIES_GRIND_III.pdf\#nameddest=tag-01JR}{III.? (p.~?)} }
\expandafter\def\csname SRef@courseIII:handle-chains:4\endcsname{\href{S6_SERIES_GRIND_III.pdf\#nameddest=tag-01JS}{III.? (p.~?)} }
\expandafter\def\csname SRef@courseIII:handle-chains:5\endcsname{\href{S6_SERIES_GRIND_III.pdf\#nameddest=tag-01JT}{III.? (p.~?)} }
\expandafter\def\csname SRef@courseIII:handle-chains:6.1\endcsname{\href{S6_SERIES_GRIND_III.pdf\#nameddest=tag-01JU}{III.? (p.~?)} }
\expandafter\def\csname SRef@courseIII:handle-chains:6.2\endcsname{\href{S6_SERIES_GRIND_III.pdf\#nameddest=tag-01JV}{III.? (p.~?)} }
\expandafter\def\csname SRef@courseIII:handle-chains:8.1\endcsname{\href{S6_SERIES_GRIND_III.pdf\#nameddest=tag-01JW}{III.? (p.~?)} }
\expandafter\def\csname SRef@courseIII:handle-chains:8.2\endcsname{\href{S6_SERIES_GRIND_III.pdf\#nameddest=tag-01JX}{III.? (p.~?)} }
\expandafter\def\csname SRef@courseIII:pair-creation:1\endcsname{\href{S6_SERIES_GRIND_III.pdf\#nameddest=tag-01LI}{III.? (p.~?)} }
\expandafter\def\csname SRef@courseIII:pair-creation:2\endcsname{\href{S6_SERIES_GRIND_III.pdf\#nameddest=tag-01LJ}{III.? (p.~?)} }
\expandafter\def\csname SRef@courseIII:pair-creation:3\endcsname{\href{S6_SERIES_GRIND_III.pdf\#nameddest=tag-01LK}{III.? (p.~?)} }
\expandafter\def\csname SRef@courseIII:pair-creation:4\endcsname{\href{S6_SERIES_GRIND_III.pdf\#nameddest=tag-01LL}{III.? (p.~?)} }
\expandafter\def\csname SRef@courseIII:pair-creation:5\endcsname{\href{S6_SERIES_GRIND_III.pdf\#nameddest=tag-01LM}{III.? (p.~?)} }
\expandafter\def\csname SRef@courseIII:pair-creation:6\endcsname{\href{S6_SERIES_GRIND_III.pdf\#nameddest=tag-01LN}{III.? (p.~?)} }
\expandafter\def\csname SRef@courseIII:pair-creation:7.1\endcsname{\href{S6_SERIES_GRIND_III.pdf\#nameddest=tag-01LO}{III.? (p.~?)} }
\expandafter\def\csname SRef@courseIII:pair-creation:7.2\endcsname{\href{S6_SERIES_GRIND_III.pdf\#nameddest=tag-01LP}{III.? (p.~?)} }
\expandafter\def\csname SRef@courseIII:low-index-slides:1\endcsname{\href{S6_SERIES_GRIND_III.pdf\#nameddest=tag-01KJ}{III.? (p.~?)} }
\expandafter\def\csname SRef@courseIII:low-index-slides:2\endcsname{\href{S6_SERIES_GRIND_III.pdf\#nameddest=tag-01KK}{III.? (p.~?)} }
\expandafter\def\csname SRef@courseIII:low-index-slides:3\endcsname{\href{S6_SERIES_GRIND_III.pdf\#nameddest=tag-01KL}{III.? (p.~?)} }
\expandafter\def\csname SRef@courseIII:low-index-slides:4\endcsname{\href{S6_SERIES_GRIND_III.pdf\#nameddest=tag-01KM}{III.? (p.~?)} }
\expandafter\def\csname SRef@courseIII:low-index-slides:5\endcsname{\href{S6_SERIES_GRIND_III.pdf\#nameddest=tag-01KN}{III.? (p.~?)} }
\expandafter\def\csname SRef@courseIII:low-index-slides:6\endcsname{\href{S6_SERIES_GRIND_III.pdf\#nameddest=tag-01KO}{III.? (p.~?)} }
\expandafter\def\csname SRef@courseIII:low-index-slides:7\endcsname{\href{S6_SERIES_GRIND_III.pdf\#nameddest=tag-01KP}{III.? (p.~?)} }
\expandafter\def\csname SRef@courseIII:low-index-slides:8\endcsname{\href{S6_SERIES_GRIND_III.pdf\#nameddest=tag-01KQ}{III.? (p.~?)} }
\expandafter\def\csname SRef@courseIII:low-index-slides:9.1\endcsname{\href{S6_SERIES_GRIND_III.pdf\#nameddest=tag-01KR}{III.? (p.~?)} }
\expandafter\def\csname SRef@courseIII:low-index-slides:9.2\endcsname{\href{S6_SERIES_GRIND_III.pdf\#nameddest=tag-01KS}{III.? (p.~?)} }
\expandafter\def\csname SRef@courseIII:adams-arithmetic:1\endcsname{\href{S6_SERIES_GRIND_III.pdf\#nameddest=tag-01HJ}{III.? (p.~?)} }
\expandafter\def\csname SRef@courseIII:adams-arithmetic:2\endcsname{\href{S6_SERIES_GRIND_III.pdf\#nameddest=tag-01HK}{III.? (p.~?)} }
\expandafter\def\csname SRef@courseIII:adams-arithmetic:3\endcsname{\href{S6_SERIES_GRIND_III.pdf\#nameddest=tag-01HL}{III.? (p.~?)} }
\expandafter\def\csname SRef@courseIII:adams-arithmetic:4\endcsname{\href{S6_SERIES_GRIND_III.pdf\#nameddest=tag-01HM}{III.? (p.~?)} }
\expandafter\def\csname SRef@courseIII:adams-arithmetic:5\endcsname{\href{S6_SERIES_GRIND_III.pdf\#nameddest=tag-01HN}{III.? (p.~?)} }
\expandafter\def\csname SRef@courseIII:adams-arithmetic:6\endcsname{\href{S6_SERIES_GRIND_III.pdf\#nameddest=tag-01HO}{III.? (p.~?)} }
\expandafter\def\csname SRef@courseIII:adams-arithmetic:7\endcsname{\href{S6_SERIES_GRIND_III.pdf\#nameddest=tag-01HP}{III.? (p.~?)} }
\expandafter\def\csname SRef@courseIII:adams-arithmetic:8\endcsname{\href{S6_SERIES_GRIND_III.pdf\#nameddest=tag-01HQ}{III.? (p.~?)} }
\expandafter\def\csname SRef@courseIII:adams-arithmetic:9\endcsname{\href{S6_SERIES_GRIND_III.pdf\#nameddest=tag-01HR}{III.? (p.~?)} }
\expandafter\def\csname SRef@courseIII:adams-tower:1\endcsname{\href{S6_SERIES_GRIND_III.pdf\#nameddest=tag-01HS}{III.? (p.~?)} }
\expandafter\def\csname SRef@courseIII:adams-tower:2\endcsname{\href{S6_SERIES_GRIND_III.pdf\#nameddest=tag-01HV}{III.? (p.~?)} }
\expandafter\def\csname SRef@courseIII:adams-tower:3\endcsname{\href{S6_SERIES_GRIND_III.pdf\#nameddest=tag-01HW}{III.? (p.~?)} }
\expandafter\def\csname SRef@courseIII:adams-tower:4\endcsname{\href{S6_SERIES_GRIND_III.pdf\#nameddest=tag-01HX}{III.? (p.~?)} }
\expandafter\def\csname SRef@courseIII:adams-tower:5\endcsname{\href{S6_SERIES_GRIND_III.pdf\#nameddest=tag-01HY}{III.? (p.~?)} }
\expandafter\def\csname SRef@courseIII:adams-tower:6\endcsname{\href{S6_SERIES_GRIND_III.pdf\#nameddest=tag-01HZ}{III.? (p.~?)} }
\expandafter\def\csname SRef@courseIII:adams-tower:7.1\endcsname{\href{S6_SERIES_GRIND_III.pdf\#nameddest=tag-01I0}{III.? (p.~?)} }
\expandafter\def\csname SRef@courseIII:adams-tower:7.2\endcsname{\href{S6_SERIES_GRIND_III.pdf\#nameddest=tag-01I1}{III.? (p.~?)} }
\expandafter\def\csname SRef@courseIII:adams-tower:8\endcsname{\href{S6_SERIES_GRIND_III.pdf\#nameddest=tag-01I2}{III.? (p.~?)} }
\expandafter\def\csname SRef@courseIII:adams-tower:9\endcsname{\href{S6_SERIES_GRIND_III.pdf\#nameddest=tag-01I3}{III.? (p.~?)} }
\expandafter\def\csname SRef@courseIII:adams-tower:10\endcsname{\href{S6_SERIES_GRIND_III.pdf\#nameddest=tag-01HT}{III.? (p.~?)} }
\expandafter\def\csname SRef@courseIII:adams-tower:11\endcsname{\href{S6_SERIES_GRIND_III.pdf\#nameddest=tag-01HU}{III.? (p.~?)} }
\expandafter\def\csname SRef@courseIII:formal-char2:1\endcsname{\href{S6_SERIES_GRIND_III.pdf\#nameddest=tag-01J2}{III.? (p.~?)} }
\expandafter\def\csname SRef@courseIII:formal-char2:2\endcsname{\href{S6_SERIES_GRIND_III.pdf\#nameddest=tag-01J3}{III.? (p.~?)} }
\expandafter\def\csname SRef@courseIII:formal-char2:3\endcsname{\href{S6_SERIES_GRIND_III.pdf\#nameddest=tag-01J4}{III.? (p.~?)} }
\expandafter\def\csname SRef@courseIII:formal-char2:4\endcsname{\href{S6_SERIES_GRIND_III.pdf\#nameddest=tag-01J5}{III.? (p.~?)} }
\expandafter\def\csname SRef@courseIII:formal-char2:4.1\endcsname{\href{S6_SERIES_GRIND_III.pdf\#nameddest=tag-01J6}{III.? (p.~?)} }
\expandafter\def\csname SRef@courseIII:formal-char2:4.2\endcsname{\href{S6_SERIES_GRIND_III.pdf\#nameddest=tag-01J7}{III.? (p.~?)} }
\expandafter\def\csname SRef@courseIII:projective-coop:1\endcsname{\href{S6_SERIES_GRIND_III.pdf\#nameddest=tag-01M0}{III.? (p.~?)} }
\expandafter\def\csname SRef@courseIII:projective-coop:2\endcsname{\href{S6_SERIES_GRIND_III.pdf\#nameddest=tag-01M1}{III.? (p.~?)} }
\expandafter\def\csname SRef@courseIII:projective-coop:3\endcsname{\href{S6_SERIES_GRIND_III.pdf\#nameddest=tag-01M2}{III.? (p.~?)} }
\expandafter\def\csname SRef@courseIII:projective-coop:4\endcsname{\href{S6_SERIES_GRIND_III.pdf\#nameddest=tag-01M3}{III.? (p.~?)} }
\expandafter\def\csname SRef@courseIII:projective-coop:5\endcsname{\href{S6_SERIES_GRIND_III.pdf\#nameddest=tag-01M4}{III.? (p.~?)} }
\expandafter\def\csname SRef@courseIII:projective-coop:6.1\endcsname{\href{S6_SERIES_GRIND_III.pdf\#nameddest=tag-01M5}{III.? (p.~?)} }
\expandafter\def\csname SRef@courseIII:projective-coop:6.2\endcsname{\href{S6_SERIES_GRIND_III.pdf\#nameddest=tag-01M6}{III.? (p.~?)} }
\expandafter\def\csname SRef@courseIII:actual-degree8:1\endcsname{\href{S6_SERIES_GRIND_III.pdf\#nameddest=tag-01H0}{III.? (p.~?)} }
\expandafter\def\csname SRef@courseIII:actual-degree8:2\endcsname{\href{S6_SERIES_GRIND_III.pdf\#nameddest=tag-01H1}{III.? (p.~?)} }
\expandafter\def\csname SRef@courseIII:actual-degree8:3\endcsname{\href{S6_SERIES_GRIND_III.pdf\#nameddest=tag-01H2}{III.? (p.~?)} }
\expandafter\def\csname SRef@courseIII:actual-degree8:4\endcsname{\href{S6_SERIES_GRIND_III.pdf\#nameddest=tag-01H3}{III.? (p.~?)} }
\expandafter\def\csname SRef@courseIII:actual-degree8:5\endcsname{\href{S6_SERIES_GRIND_III.pdf\#nameddest=tag-01H4}{III.? (p.~?)} }
\expandafter\def\csname SRef@courseIII:actual-degree8:6\endcsname{\href{S6_SERIES_GRIND_III.pdf\#nameddest=tag-01H5}{III.? (p.~?)} }
\expandafter\def\csname SRef@courseIII:actual-degree8:7\endcsname{\href{S6_SERIES_GRIND_III.pdf\#nameddest=tag-01H6}{III.? (p.~?)} }
\expandafter\def\csname SRef@courseIII:actual-degree8:8\endcsname{\href{S6_SERIES_GRIND_III.pdf\#nameddest=tag-01H7}{III.? (p.~?)} }
\expandafter\def\csname SRef@courseIII:actual-degree8:9.1\endcsname{\href{S6_SERIES_GRIND_III.pdf\#nameddest=tag-01H8}{III.? (p.~?)} }
\expandafter\def\csname SRef@courseIII:actual-degree8:9.2\endcsname{\href{S6_SERIES_GRIND_III.pdf\#nameddest=tag-01H9}{III.? (p.~?)} }
\expandafter\def\csname SRef@courseIII:pt-relative-embedding\endcsname{\href{S6_SERIES_GRIND_III.pdf\#nameddest=tag-01ME}{III.? (p.~?)} }
\expandafter\def\csname SRef@courseIII:pt-relative-normal-frame\endcsname{\href{S6_SERIES_GRIND_III.pdf\#nameddest=tag-01MF}{III.? (p.~?)} }
\expandafter\def\csname SRef@courseIII:pt-normal-inverse-comparison\endcsname{\href{S6_SERIES_GRIND_III.pdf\#nameddest=tag-01MD}{III.? (p.~?)} }
\expandafter\def\csname SRef@courseIII:pt-abstract-framed-comparison\endcsname{\href{S6_SERIES_GRIND_III.pdf\#nameddest=tag-01MC}{III.? (p.~?)} }
\expandafter\def\csname SRef@courseIII:actual-resolution:1\endcsname{\href{S6_SERIES_GRIND_III.pdf\#nameddest=tag-01HA}{III.? (p.~?)} }
\expandafter\def\csname SRef@courseIII:actual-resolution:2\endcsname{\href{S6_SERIES_GRIND_III.pdf\#nameddest=tag-01HB}{III.? (p.~?)} }
\expandafter\def\csname SRef@courseIII:actual-resolution:3\endcsname{\href{S6_SERIES_GRIND_III.pdf\#nameddest=tag-01HC}{III.? (p.~?)} }
\expandafter\def\csname SRef@courseIII:actual-resolution:4\endcsname{\href{S6_SERIES_GRIND_III.pdf\#nameddest=tag-01HD}{III.? (p.~?)} }
\expandafter\def\csname SRef@courseIII:actual-resolution:5\endcsname{\href{S6_SERIES_GRIND_III.pdf\#nameddest=tag-01HE}{III.? (p.~?)} }
\expandafter\def\csname SRef@courseIII:actual-resolution:6\endcsname{\href{S6_SERIES_GRIND_III.pdf\#nameddest=tag-01HF}{III.? (p.~?)} }
\expandafter\def\csname SRef@courseIII:actual-resolution:7\endcsname{\href{S6_SERIES_GRIND_III.pdf\#nameddest=tag-01HG}{III.? (p.~?)} }
\expandafter\def\csname SRef@courseIII:actual-resolution:8.1\endcsname{\href{S6_SERIES_GRIND_III.pdf\#nameddest=tag-01HH}{III.? (p.~?)} }
\expandafter\def\csname SRef@courseIII:actual-resolution:8.2\endcsname{\href{S6_SERIES_GRIND_III.pdf\#nameddest=tag-01HI}{III.? (p.~?)} }
\expandafter\def\csname SRef@courseIII:milnor-cobar:1.1\endcsname{\href{S6_SERIES_GRIND_III.pdf\#nameddest=tag-01L1}{III.? (p.~?)} }
\expandafter\def\csname SRef@courseIII:milnor-cobar:6.1\endcsname{\href{S6_SERIES_GRIND_III.pdf\#nameddest=tag-01L2}{III.? (p.~?)} }
\expandafter\def\csname SRef@courseIII:milnor-cobar:7.1\endcsname{\href{S6_SERIES_GRIND_III.pdf\#nameddest=tag-01L3}{III.? (p.~?)} }
\expandafter\def\csname SRef@courseIII:milnor-cobar:7.2\endcsname{\href{S6_SERIES_GRIND_III.pdf\#nameddest=tag-01L4}{III.? (p.~?)} }
\expandafter\def\csname SRef@courseIII:odd-coop:1\endcsname{\href{S6_SERIES_GRIND_III.pdf\#nameddest=tag-01LA}{III.? (p.~?)} }
\expandafter\def\csname SRef@courseIII:odd-coop:2\endcsname{\href{S6_SERIES_GRIND_III.pdf\#nameddest=tag-01LB}{III.? (p.~?)} }
\expandafter\def\csname SRef@courseIII:odd-coop:3\endcsname{\href{S6_SERIES_GRIND_III.pdf\#nameddest=tag-01LC}{III.? (p.~?)} }
\expandafter\def\csname SRef@courseIII:odd-coop:4\endcsname{\href{S6_SERIES_GRIND_III.pdf\#nameddest=tag-01LD}{III.? (p.~?)} }
\expandafter\def\csname SRef@courseIII:odd-coop:5\endcsname{\href{S6_SERIES_GRIND_III.pdf\#nameddest=tag-01LE}{III.? (p.~?)} }
\expandafter\def\csname SRef@courseIII:odd-coop:6\endcsname{\href{S6_SERIES_GRIND_III.pdf\#nameddest=tag-01LF}{III.? (p.~?)} }
\expandafter\def\csname SRef@courseIII:odd-coop:7.1\endcsname{\href{S6_SERIES_GRIND_III.pdf\#nameddest=tag-01LG}{III.? (p.~?)} }
\expandafter\def\csname SRef@courseIII:odd-coop:7.2\endcsname{\href{S6_SERIES_GRIND_III.pdf\#nameddest=tag-01LH}{III.? (p.~?)} }
\expandafter\def\csname SRef@courseIII:coefficient-rings:1\endcsname{\href{S6_SERIES_GRIND_III.pdf\#nameddest=tag-01IM}{III.? (p.~?)} }
\expandafter\def\csname SRef@courseIII:coefficient-rings:2\endcsname{\href{S6_SERIES_GRIND_III.pdf\#nameddest=tag-01IN}{III.? (p.~?)} }
\expandafter\def\csname SRef@courseIII:coefficient-rings:3\endcsname{\href{S6_SERIES_GRIND_III.pdf\#nameddest=tag-01IO}{III.? (p.~?)} }
\expandafter\def\csname SRef@courseIII:coefficient-rings:4\endcsname{\href{S6_SERIES_GRIND_III.pdf\#nameddest=tag-01IP}{III.? (p.~?)} }
\expandafter\def\csname SRef@courseIII:coefficient-rings:5\endcsname{\href{S6_SERIES_GRIND_III.pdf\#nameddest=tag-01IQ}{III.? (p.~?)} }
\expandafter\def\csname SRef@courseIII:coefficient-rings:6\endcsname{\href{S6_SERIES_GRIND_III.pdf\#nameddest=tag-01IR}{III.? (p.~?)} }
\expandafter\def\csname SRef@courseIII:coefficient-rings:7.1\endcsname{\href{S6_SERIES_GRIND_III.pdf\#nameddest=tag-01IS}{III.? (p.~?)} }
\expandafter\def\csname SRef@courseIII:coefficient-rings:7.2\endcsname{\href{S6_SERIES_GRIND_III.pdf\#nameddest=tag-01IT}{III.? (p.~?)} }
\expandafter\def\csname SRef@courseIII:duality:2.1\endcsname{\href{S6_SERIES_GRIND_III.pdf\#nameddest=tag-01IU}{III.? (p.~?)} }
\expandafter\def\csname SRef@courseIII:duality:2.2\endcsname{\href{S6_SERIES_GRIND_III.pdf\#nameddest=tag-01IV}{III.? (p.~?)} }
\expandafter\def\csname SRef@courseIII:duality:2.3\endcsname{\href{S6_SERIES_GRIND_III.pdf\#nameddest=tag-01IW}{III.? (p.~?)} }
\expandafter\def\csname SRef@courseIII:duality:3.1\endcsname{\href{S6_SERIES_GRIND_III.pdf\#nameddest=tag-01IX}{III.? (p.~?)} }
\expandafter\def\csname SRef@courseIII:duality:3.2\endcsname{\href{S6_SERIES_GRIND_III.pdf\#nameddest=tag-01IY}{III.? (p.~?)} }
\expandafter\def\csname SRef@courseIII:duality:3.3\endcsname{\href{S6_SERIES_GRIND_III.pdf\#nameddest=tag-01IZ}{III.? (p.~?)} }
\expandafter\def\csname SRef@courseIII:duality:4.1\endcsname{\href{S6_SERIES_GRIND_III.pdf\#nameddest=tag-01J0}{III.? (p.~?)} }
\expandafter\def\csname SRef@courseIII:duality:4.2\endcsname{\href{S6_SERIES_GRIND_III.pdf\#nameddest=tag-01J1}{III.? (p.~?)} }
\expandafter\def\csname SRef@night:P1.6\endcsname{\href{S6_SERIES_GRIND_V.pdf\#nameddest=tag-02QX}{V.? (p.~?)} }
\expandafter\def\csname SRef@night:P1.14\endcsname{\href{S6_SERIES_GRIND_V.pdf\#nameddest=tag-02QN}{V.? (p.~?)} }
\expandafter\def\csname SRef@night:P5.7\endcsname{\hyperlink{tag-02YS}{I.\ref*{night:P5.7}}}
\expandafter\def\csname SRef@ari:A.1.1\endcsname{\hyperlink{tag-03J1}{I.\ref*{ari:A.1.1}}}
\expandafter\def\csname SRef@ari:A.1.2\endcsname{\hyperlink{tag-03J2}{I.\ref*{ari:A.1.2}}}
\expandafter\def\csname SRef@ari:A.1.3\endcsname{\hyperlink{tag-03J3}{I.\ref*{ari:A.1.3}}}
\expandafter\def\csname SRef@ari:A.1.4\endcsname{\hyperlink{tag-03J4}{I.\ref*{ari:A.1.4}}}
\expandafter\def\csname SRef@ari:A.2.1\endcsname{\hyperlink{tag-03J5}{I.\ref*{ari:A.2.1}}}
\expandafter\def\csname SRef@ari:A.2.2\endcsname{\hyperlink{tag-03J6}{I.\ref*{ari:A.2.2}}}
\expandafter\def\csname SRef@ari:A.2.3\endcsname{\hyperlink{tag-03J7}{I.\ref*{ari:A.2.3}}}
\expandafter\def\csname SRef@ari:A.2.4\endcsname{\hyperlink{tag-03J8}{I.\ref*{ari:A.2.4}}}
\expandafter\def\csname SRef@ari:A.2.5\endcsname{\hyperlink{tag-03J9}{I.\ref*{ari:A.2.5}}}
\expandafter\def\csname SRef@ari:A.3.1\endcsname{\hyperlink{tag-03JA}{I.\ref*{ari:A.3.1}}}
\expandafter\def\csname SRef@ari:A.3.2\endcsname{\hyperlink{tag-03JB}{I.\ref*{ari:A.3.2}}}
\expandafter\def\csname SRef@ari:A.4.1\endcsname{\hyperlink{tag-03JC}{I.\ref*{ari:A.4.1}}}
\expandafter\def\csname SRef@ari:A.4.2\endcsname{\hyperlink{tag-03JD}{I.\ref*{ari:A.4.2}}}
\expandafter\def\csname SRef@ari:A.4.3\endcsname{\hyperlink{tag-03JE}{I.\ref*{ari:A.4.3}}}
\expandafter\def\csname SRef@ari:A.4.4\endcsname{\hyperlink{tag-03JF}{I.\ref*{ari:A.4.4}}}
\expandafter\def\csname SRef@ari:A.4.5\endcsname{\hyperlink{tag-03JG}{I.\ref*{ari:A.4.5}}}
\expandafter\def\csname SRef@ari:A.4.6\endcsname{\hyperlink{tag-03JH}{I.\ref*{ari:A.4.6}}}
\expandafter\def\csname SRef@ari:A.4.7\endcsname{\hyperlink{tag-03JI}{I.\ref*{ari:A.4.7}}}
\expandafter\def\csname SRef@ari:A.4.8\endcsname{\hyperlink{tag-03JJ}{I.\ref*{ari:A.4.8}}}
\expandafter\def\csname SRef@ari:A.5.1\endcsname{\hyperlink{tag-03JK}{I.\ref*{ari:A.5.1}}}
\expandafter\def\csname SRef@ari:A.5.2\endcsname{\hyperlink{tag-03JL}{I.\ref*{ari:A.5.2}}}
\expandafter\def\csname SRef@ari:A.5.3\endcsname{\hyperlink{tag-03JM}{I.\ref*{ari:A.5.3}}}
\expandafter\def\csname SRef@ari:A.5.4\endcsname{\hyperlink{tag-03JN}{I.\ref*{ari:A.5.4}}}
\expandafter\def\csname SRef@ari:A.5.5\endcsname{\hyperlink{tag-03JO}{I.\ref*{ari:A.5.5}}}
\expandafter\def\csname SRef@ari:A.5.6\endcsname{\hyperlink{tag-03JP}{I.\ref*{ari:A.5.6}}}
\expandafter\def\csname SRef@ari:A.5.8\endcsname{\hyperlink{tag-03JR}{I.\ref*{ari:A.5.8}}}
\expandafter\def\csname SRef@ari:A.5.7\endcsname{\hyperlink{tag-03JQ}{I.\ref*{ari:A.5.7}}}
\expandafter\def\csname SRef@series:I:main-theorems\endcsname{\hyperlink{tag-03JS}{I.\ref*{series:I:main-theorems}}}
\expandafter\def\csname SRef@series:I:A\endcsname{\hyperlink{tag-03JT}{I.\ref*{series:I:A}}}
\expandafter\def\csname SRef@series:I:B\endcsname{\hyperlink{tag-03JU}{I.\ref*{series:I:B}}}
\expandafter\def\csname SRef@series:I:C\endcsname{\hyperlink{tag-03JV}{I.\ref*{series:I:C}}}
\expandafter\def\csname SRef@series:I:D\endcsname{\hyperlink{tag-03JW}{I.\ref*{series:I:D}}}
\expandafter\def\csname SRef@series:I:E\endcsname{\hyperlink{tag-03JX}{I.\ref*{series:I:E}}}
\expandafter\def\csname SRef@series:I:F\endcsname{\hyperlink{tag-03JY}{I.\ref*{series:I:F}}}
\expandafter\def\csname SRef@series:I:G\endcsname{\hyperlink{tag-03JZ}{I.\ref*{series:I:G}}}
\expandafter\def\csname SRef@series:I:outlines\endcsname{\hyperlink{tag-03K0}{I.\ref*{series:I:outlines}}}
\expandafter\def\csname SRef@series:II:main-theorems\endcsname{\href{S6_SERIES_GRIND_II.pdf\#nameddest=tag-03K1}{II.1.2 (p.~5)} }
\expandafter\def\csname SRef@series:II:A\endcsname{\href{S6_SERIES_GRIND_II.pdf\#nameddest=tag-03K2}{II.A (p.~5)} }
\expandafter\def\csname SRef@series:II:B\endcsname{\href{S6_SERIES_GRIND_II.pdf\#nameddest=tag-03K3}{II.B (p.~6)} }
\expandafter\def\csname SRef@series:II:C\endcsname{\href{S6_SERIES_GRIND_II.pdf\#nameddest=tag-03K4}{II.C (p.~6)} }
\expandafter\def\csname SRef@series:II:D\endcsname{\href{S6_SERIES_GRIND_II.pdf\#nameddest=tag-03K5}{II.D (p.~7)} }
\expandafter\def\csname SRef@series:II:E\endcsname{\href{S6_SERIES_GRIND_II.pdf\#nameddest=tag-03K6}{II.E (p.~7)} }
\expandafter\def\csname SRef@series:II:F\endcsname{\href{S6_SERIES_GRIND_II.pdf\#nameddest=tag-03K7}{II.F (p.~8)} }
\expandafter\def\csname SRef@series:II:G\endcsname{\href{S6_SERIES_GRIND_II.pdf\#nameddest=tag-03K8}{II.G (p.~9)} }
\expandafter\def\csname SRef@series:II:outline\endcsname{\href{S6_SERIES_GRIND_II.pdf\#nameddest=tag-03K9}{II.1.3 (p.~13)} }
\expandafter\def\csname SRef@series:III:main-theorems\endcsname{\href{S6_SERIES_GRIND_III.pdf\#nameddest=tag-03KA}{III.1.2 (p.~3)} }
\expandafter\def\csname SRef@series:III:A\endcsname{\href{S6_SERIES_GRIND_III.pdf\#nameddest=tag-03KB}{III.A (p.~3)} }
\expandafter\def\csname SRef@series:III:B\endcsname{\href{S6_SERIES_GRIND_III.pdf\#nameddest=tag-03KC}{III.B (p.~3)} }
\expandafter\def\csname SRef@series:III:C\endcsname{\href{S6_SERIES_GRIND_III.pdf\#nameddest=tag-03KD}{III.C (p.~4)} }
\expandafter\def\csname SRef@series:III:D\endcsname{\href{S6_SERIES_GRIND_III.pdf\#nameddest=tag-03KE}{III.D (p.~4)} }
\expandafter\def\csname SRef@series:III:E\endcsname{\href{S6_SERIES_GRIND_III.pdf\#nameddest=tag-03KF}{III.E (p.~4)} }
\expandafter\def\csname SRef@series:III:F\endcsname{\href{S6_SERIES_GRIND_III.pdf\#nameddest=tag-03KG}{III.F (p.~5)} }
\expandafter\def\csname SRef@series:III:G\endcsname{\href{S6_SERIES_GRIND_III.pdf\#nameddest=tag-03KH}{III.G (p.~5)} }
\expandafter\def\csname SRef@series:IV:main-theorems\endcsname{\href{S6_SERIES_GRIND_IV.pdf\#nameddest=tag-03KI}{IV.1.2 (p.~3)} }
\expandafter\def\csname SRef@series:IV:A\endcsname{\href{S6_SERIES_GRIND_IV.pdf\#nameddest=tag-03KJ}{IV.A (p.~3)} }
\expandafter\def\csname SRef@series:IV:B\endcsname{\href{S6_SERIES_GRIND_IV.pdf\#nameddest=tag-03KK}{IV.B (p.~4)} }
\expandafter\def\csname SRef@series:IV:C\endcsname{\href{S6_SERIES_GRIND_IV.pdf\#nameddest=tag-03KL}{IV.C (p.~5)} }
\expandafter\def\csname SRef@series:IV:D\endcsname{\href{S6_SERIES_GRIND_IV.pdf\#nameddest=tag-03KM}{IV.D (p.~6)} }
\expandafter\def\csname SRef@series:IV:E\endcsname{\href{S6_SERIES_GRIND_IV.pdf\#nameddest=tag-03KN}{IV.E (p.~7)} }
\expandafter\def\csname SRef@series:IV:F\endcsname{\href{S6_SERIES_GRIND_IV.pdf\#nameddest=tag-03KO}{IV.F (p.~8)} }
\expandafter\def\csname SRef@series:IV:G\endcsname{\href{S6_SERIES_GRIND_IV.pdf\#nameddest=tag-03KP}{IV.G (p.~8)} }
\expandafter\def\csname SRef@series:V:main-theorems\endcsname{\href{S6_SERIES_GRIND_V.pdf\#nameddest=tag-03KQ}{V.1.2 (p.~3)} }
\expandafter\def\csname SRef@series:V:A\endcsname{\href{S6_SERIES_GRIND_V.pdf\#nameddest=tag-03KR}{V.A (p.~3)} }
\expandafter\def\csname SRef@series:V:B\endcsname{\href{S6_SERIES_GRIND_V.pdf\#nameddest=tag-03KS}{V.B (p.~3)} }
\expandafter\def\csname SRef@series:V:C\endcsname{\href{S6_SERIES_GRIND_V.pdf\#nameddest=tag-03KT}{V.C (p.~4)} }
\expandafter\def\csname SRef@series:V:D\endcsname{\href{S6_SERIES_GRIND_V.pdf\#nameddest=tag-03KU}{V.D (p.~4)} }
\expandafter\def\csname SRef@series:V:E\endcsname{\href{S6_SERIES_GRIND_V.pdf\#nameddest=tag-03KV}{V.E (p.~5)} }
\expandafter\def\csname SRef@series:V:F\endcsname{\href{S6_SERIES_GRIND_V.pdf\#nameddest=tag-03KW}{V.F (p.~5)} }
\expandafter\def\csname SRef@layout:V:fresh-gauges\endcsname{\href{S6_SERIES_GRIND_V.pdf\#nameddest=tag-03KX}{V.{{$\langle 2\rangle $1.}} (p.~126)} }
\expandafter\def\csname SRef@layout:V:fresh-minimum\endcsname{\href{S6_SERIES_GRIND_V.pdf\#nameddest=tag-03KY}{V.{{$\langle 1\rangle $4.}} (p.~130)} }
\expandafter\def\csname SRef@layout:V:fresh-bounds\endcsname{\href{S6_SERIES_GRIND_V.pdf\#nameddest=tag-03KZ}{V.{1.18.4} (p.~130)} }
\expandafter\def\csname SRef@layout:V:fresh-marking\endcsname{\href{S6_SERIES_GRIND_V.pdf\#nameddest=tag-03L0}{V.{{$\langle 1\rangle $6.}} (p.~132)} }
\expandafter\def\csname SRef@layout:V:fresh-splitting\endcsname{\href{S6_SERIES_GRIND_V.pdf\#nameddest=tag-03L1}{V.{{$\langle 1\rangle $7.}} (p.~133)} }
\expandafter\def\csname SRef@iv:source_transports\endcsname{\href{S6_SERIES_GRIND_IV.pdf\#nameddest=tag-03L2}{IV.3.27 (p.~150)} }
\expandafter\def\csname SRef@inv:section:coefficient-atlas\endcsname{\hyperlink{tag-03L3}{I.\ref*{inv:section:coefficient-atlas}}}
\expandafter\def\csname SRef@inv:eq:atlas-theta-factor\endcsname{\hyperlink{tag-03L4}{I.\ref*{inv:eq:atlas-theta-factor}}}
\expandafter\def\csname SRef@inv:eq:atlas-semiflat\endcsname{\hyperlink{tag-03L5}{I.\ref*{inv:eq:atlas-semiflat}}}
\expandafter\def\csname SRef@inv:eq:atlas-affine-lift\endcsname{\hyperlink{tag-03L6}{I.\ref*{inv:eq:atlas-affine-lift}}}
\expandafter\def\csname SRef@inv:eq:atlas-cusp-factor\endcsname{\hyperlink{tag-03L7}{I.\ref*{inv:eq:atlas-cusp-factor}}}
\expandafter\def\csname SRef@inv:eq:atlas-cusp-unit\endcsname{\hyperlink{tag-03L8}{I.\ref*{inv:eq:atlas-cusp-unit}}}
\expandafter\def\csname SRef@inv:eq:atlas-connection-inverse\endcsname{\hyperlink{tag-03L9}{I.\ref*{inv:eq:atlas-connection-inverse}}}
\expandafter\def\csname SRef@inv:eq:atlas-cusp-connection\endcsname{\hyperlink{tag-03LA}{I.\ref*{inv:eq:atlas-cusp-connection}}}
\expandafter\def\csname SRef@inv:eq:atlas-finite-connection\endcsname{\hyperlink{tag-03LB}{I.\ref*{inv:eq:atlas-finite-connection}}}
\expandafter\def\csname SRef@inv:eq:atlas-averaged-connection\endcsname{\hyperlink{tag-03LC}{I.\ref*{inv:eq:atlas-averaged-connection}}}
\expandafter\def\csname SRef@inv:eq:atlas-cocycle\endcsname{\hyperlink{tag-03LD}{I.\ref*{inv:eq:atlas-cocycle}}}
\expandafter\def\csname SRef@inv:eq:atlas-weights\endcsname{\hyperlink{tag-03LE}{I.\ref*{inv:eq:atlas-weights}}}
\expandafter\def\csname SRef@inv:eq:atlas-weight-derivative\endcsname{\hyperlink{tag-03LF}{I.\ref*{inv:eq:atlas-weight-derivative}}}
\expandafter\def\csname SRef@inv:eq:atlas-final-forms\endcsname{\hyperlink{tag-03LG}{I.\ref*{inv:eq:atlas-final-forms}}}
\expandafter\def\csname SRef@inv:eq:atlas-final-coefficients\endcsname{\hyperlink{tag-03LH}{I.\ref*{inv:eq:atlas-final-coefficients}}}
\expandafter\def\csname SRef@inv:C3.8:operator-comparison\endcsname{\hyperlink{tag-03LI}{I.\ref*{inv:C3.8:operator-comparison}}}
\expandafter\def\csname SRef@iv:open:K111:original_auxiliary_arrays\endcsname{\href{S6_SERIES_GRIND_IV.pdf\#nameddest=tag-03LJ}{IV.12.1 (p.~1396)} }
\expandafter\def\csname SRef@maps:proof:D.OS.refined\endcsname{\href{S6_SERIES_GRIND_III.pdf\#nameddest=tag-03LK}{III.5.2 (p.~207)} }
\expandafter\def\csname SRef@maps:eq:D.OR1\endcsname{\href{S6_SERIES_GRIND_III.pdf\#nameddest=tag-03LL}{III.{D.OR1} (p.~207)} }
\expandafter\def\csname SRef@maps:eq:D.OR2\endcsname{\href{S6_SERIES_GRIND_III.pdf\#nameddest=tag-03LM}{III.{D.OR2} (p.~207)} }
\expandafter\def\csname SRef@maps:eq:D.OR3\endcsname{\href{S6_SERIES_GRIND_III.pdf\#nameddest=tag-03LN}{III.{D.OR3} (p.~207)} }
\expandafter\def\csname SRef@maps:eq:D.OR4\endcsname{\href{S6_SERIES_GRIND_III.pdf\#nameddest=tag-03LO}{III.{D.OR4} (p.~207)} }
\expandafter\def\csname SRef@maps:eq:D.OR5\endcsname{\href{S6_SERIES_GRIND_III.pdf\#nameddest=tag-03LP}{III.{D.OR5} (p.~208)} }
\expandafter\def\csname SRef@maps:eq:D.OR6\endcsname{\href{S6_SERIES_GRIND_III.pdf\#nameddest=tag-03LQ}{III.{D.OR6} (p.~208)} }
\expandafter\def\csname SRef@maps:eq:D.OR7\endcsname{\href{S6_SERIES_GRIND_III.pdf\#nameddest=tag-03LR}{III.{D.OR7} (p.~208)} }
\expandafter\def\csname SRef@maps:eq:D.OR8\endcsname{\href{S6_SERIES_GRIND_III.pdf\#nameddest=tag-03LS}{III.{D.OR8} (p.~208)} }
\expandafter\def\csname SRef@maps:eq:D.OR9\endcsname{\href{S6_SERIES_GRIND_III.pdf\#nameddest=tag-03LT}{III.{D.OR9} (p.~208)} }
\expandafter\def\csname SRef@maps:eq:D.OR10\endcsname{\href{S6_SERIES_GRIND_III.pdf\#nameddest=tag-03LU}{III.{D.OR10} (p.~209)} }
\expandafter\def\csname SRef@maps:eq:D.OR11\endcsname{\href{S6_SERIES_GRIND_III.pdf\#nameddest=tag-03LV}{III.{D.OR11} (p.~209)} }
\expandafter\def\csname SRef@maps:eq:D.OR12\endcsname{\href{S6_SERIES_GRIND_III.pdf\#nameddest=tag-03LW}{III.{D.OR12} (p.~210)} }
\expandafter\def\csname SRef@maps:eq:D.OR13\endcsname{\href{S6_SERIES_GRIND_III.pdf\#nameddest=tag-03LX}{III.{D.OR13} (p.~210)} }
\expandafter\def\csname SRef@maps:eq:D.OR14\endcsname{\href{S6_SERIES_GRIND_III.pdf\#nameddest=tag-03LY}{III.{D.OR14} (p.~211)} }
\expandafter\def\csname SRef@inv:E2.5\endcsname{\href{S6_SERIES_GRIND_II.pdf\#nameddest=tag-03LZ}{II.14.2 (p.~356)} }
\expandafter\def\csname SRef@invII:eq:E12-node\endcsname{\href{S6_SERIES_GRIND_II.pdf\#nameddest=tag-03M0}{II.{E.12.1} (p.~358)} }
\expandafter\def\csname SRef@invII:eq:E12-jets\endcsname{\href{S6_SERIES_GRIND_II.pdf\#nameddest=tag-03M1}{II.{E.12.2} (p.~359)} }
\expandafter\def\csname SRef@invII:eq:E12-open\endcsname{\href{S6_SERIES_GRIND_II.pdf\#nameddest=tag-03M2}{II.{E.12.3} (p.~360)} }
\expandafter\def\csname SBack@I:weierstrass-comparison\endcsname{\hyperlink{step-I-A-1}{proves step 1 of Theorem A}}
\expandafter\def\csname SBack@fam:src:B.1.1\endcsname{\hyperlink{step-I-A-1}{proves step 1 of Theorem A}}
\expandafter\def\csname SBack@fam:src:C.1.1\endcsname{\hyperlink{step-I-A-2}{proves step 2 of Theorem A}}
\expandafter\def\csname SBack@ari:C.4.2\endcsname{\hyperlink{step-I-A-3}{proves step 3 of Theorem A}}
\expandafter\def\csname SBack@ari:C.4.3\endcsname{\hyperlink{step-I-A-3}{proves step 3 of Theorem A}}
\expandafter\def\csname SBack@ari:C.4.4\endcsname{\hyperlink{step-I-A-3}{proves step 3 of Theorem A}}
\expandafter\def\csname SBack@ari:C.8.1\endcsname{\hyperlink{step-I-A-4}{proves step 4 of Theorem A}}
\expandafter\def\csname SBack@ari:C.8.2\endcsname{\hyperlink{step-I-A-4}{proves step 4 of Theorem A}}
\expandafter\def\csname SBack@ari:C.8.3\endcsname{\hyperlink{step-I-A-4}{proves step 4 of Theorem A}}
\expandafter\def\csname SBack@ari:C.8.4\endcsname{\hyperlink{step-I-A-4}{proves step 4 of Theorem A}}
\expandafter\def\csname SBack@fam:src:C.2.3\endcsname{\hyperlink{step-I-A-4}{proves step 4 of Theorem A}}
\expandafter\def\csname SBack@ari:C.8.15\endcsname{\hyperlink{step-I-A-5}{proves step 5 of Theorem A}}
\expandafter\def\csname SBack@push:src:D.4.8\endcsname{\hyperlink{step-I-A-5}{proves step 5 of Theorem A}}
\expandafter\def\csname SBack@ari:C.6.1\endcsname{\hyperlink{step-I-A-6}{proves step 6 of Theorem A}}
\expandafter\def\csname SBack@ari:C.6.3\endcsname{\hyperlink{step-I-A-6}{proves step 6 of Theorem A}}
\expandafter\def\csname SBack@fam:src:A.1.1\endcsname{\hyperlink{step-I-B-1}{proves step 1 of Theorem B}}
\expandafter\def\csname SBack@fam:src:A.5.1\endcsname{\hyperlink{step-I-B-1}{proves step 1 of Theorem B}}
\expandafter\def\csname SBack@fam:src:B.2.1\endcsname{\hyperlink{step-I-B-2}{proves step 2 of Theorem B}}
\expandafter\def\csname SBack@fam:src:B.2.2\endcsname{\hyperlink{step-I-B-2}{proves step 2 of Theorem B}}
\expandafter\def\csname SBack@fam:src:B.2.3\endcsname{\hyperlink{step-I-B-2}{proves step 2 of Theorem B}}
\expandafter\def\csname SBack@fam:src:B.3.1\endcsname{\hyperlink{step-I-B-3}{proves step 3 of Theorem B}}
\expandafter\def\csname SBack@push:src:E.1.2\endcsname{\hyperlink{step-I-B-3}{proves step 3 of Theorem B}; \hyperlink{step-I-E-5}{proves step 5 of Theorem E}}
\expandafter\def\csname SBack@ari:C.2.1\endcsname{\hyperlink{step-I-B-4}{proves step 4 of Theorem B}}
\expandafter\def\csname SBack@fam:src:B.4.2\endcsname{\hyperlink{step-I-B-5}{proves step 5 of Theorem B}}
\expandafter\def\csname SBack@mapsI:src:A.3.1\endcsname{\hyperlink{step-I-B-5}{proves step 5 of Theorem B}}
\expandafter\def\csname SBack@mapsI:src:A.3.2\endcsname{\hyperlink{step-I-B-6}{proves step 6 of Theorem B}}
\expandafter\def\csname SBack@mapsI:src:A.3.3\endcsname{\hyperlink{step-I-B-6}{proves step 6 of Theorem B}}
\expandafter\def\csname SBack@fam:src:F.1.1\endcsname{\hyperlink{step-I-C-1}{proves step 1 of Theorem C}}
\expandafter\def\csname SBack@I:monodromy-fundamental-set\endcsname{\hyperlink{step-I-C-2}{proves step 2 of Theorem C}}
\expandafter\def\csname SBack@fam:src:F.2.2\endcsname{\hyperlink{step-I-C-3}{proves step 3 of Theorem C}}
\expandafter\def\csname SBack@fam:src:F.2.3\endcsname{\hyperlink{step-I-C-3}{proves step 3 of Theorem C}}
\expandafter\def\csname SBack@fam:src:F.3.1\endcsname{\hyperlink{step-I-C-4}{proves step 4 of Theorem C}}
\expandafter\def\csname SBack@fam:src:D.1.4\endcsname{\hyperlink{step-I-C-5}{proves step 5 of Theorem C}}
\expandafter\def\csname SBack@fam:src:D.1.6\endcsname{\hyperlink{step-I-C-5}{proves step 5 of Theorem C}}
\expandafter\def\csname SBack@inv:A2.2\endcsname{\hyperlink{step-I-D-1}{proves step 1 of Theorem D}}
\expandafter\def\csname SBack@inv:A2.4\endcsname{\hyperlink{step-I-D-1}{proves step 1 of Theorem D}}
\expandafter\def\csname SBack@inv:A1.0\endcsname{\hyperlink{step-I-D-2}{proves step 2 of Theorem D}}
\expandafter\def\csname SBack@inv:A3.1\endcsname{\hyperlink{step-I-D-2}{proves step 2 of Theorem D}}
\expandafter\def\csname SBack@inv:C1.1\endcsname{\hyperlink{step-I-D-3}{proves step 3 of Theorem D}}
\expandafter\def\csname SBack@inv:C2.1\endcsname{\hyperlink{step-I-D-3}{proves step 3 of Theorem D}}
\expandafter\def\csname SBack@inv:C3.3\endcsname{\hyperlink{step-I-D-3}{proves step 3 of Theorem D}}
\expandafter\def\csname SBack@inv:A4.1\endcsname{\hyperlink{step-I-D-4}{proves step 4 of Theorem D}}
\expandafter\def\csname SBack@inv:A4.2\endcsname{\hyperlink{step-I-D-4}{proves step 4 of Theorem D}}
\expandafter\def\csname SBack@inv:C3.8\endcsname{\hyperlink{step-I-D-5}{proves step 5 of Theorem D}}
\expandafter\def\csname SBack@inv:C3.2\endcsname{\hyperlink{step-I-E-1}{proves step 1 of Theorem E}; \hyperlink{step-I-E-4}{proves step 4 of Theorem E}}
\expandafter\def\csname SBack@inv:D4.1\endcsname{\hyperlink{step-I-E-1}{proves step 1 of Theorem E}}
\expandafter\def\csname SBack@inv:D4.2\endcsname{\hyperlink{step-I-E-2}{proves step 2 of Theorem E}}
\expandafter\def\csname SBack@inv:D4.3\endcsname{\hyperlink{step-I-E-2}{proves step 2 of Theorem E}}
\expandafter\def\csname SBack@inv:D4.4\endcsname{\hyperlink{step-I-E-3}{proves step 3 of Theorem E}}
\expandafter\def\csname SBack@inv:D4.5\endcsname{\hyperlink{step-I-E-3}{proves step 3 of Theorem E}; \hyperlink{step-I-E-4}{proves step 4 of Theorem E}}
\expandafter\def\csname SBack@inv:C3.1\endcsname{\hyperlink{step-I-E-4}{proves step 4 of Theorem E}}
\expandafter\def\csname SBack@push:src:E.2.1\endcsname{\hyperlink{step-I-E-5}{proves step 5 of Theorem E}; \hyperlink{step-I-E-6}{proves step 6 of Theorem E}}
\expandafter\def\csname SBack@fam:st:A.5.2\endcsname{\hyperlink{step-I-E-6}{proves step 6 of Theorem E}}
\expandafter\def\csname SBack@fam:src:B.4.1\endcsname{\hyperlink{step-I-E-6}{proves step 6 of Theorem E}}
\expandafter\def\csname SBack@night:P2.2\endcsname{\hyperlink{step-I-F-1}{proves step 1 of Theorem F}}
\expandafter\def\csname SBack@night:P2.3\endcsname{\hyperlink{step-I-F-1}{proves step 1 of Theorem F}}
\expandafter\def\csname SBack@night:P2.4\endcsname{\hyperlink{step-I-F-2}{proves step 2 of Theorem F}}
\expandafter\def\csname SBack@night:P2.5\endcsname{\hyperlink{step-I-F-3}{proves step 3 of Theorem F}}
\expandafter\def\csname SBack@night:P2.6\endcsname{\hyperlink{step-I-F-3}{proves step 3 of Theorem F}}
\expandafter\def\csname SBack@night:P2.7\endcsname{\hyperlink{step-I-F-4}{proves step 4 of Theorem F}; \hyperlink{step-I-F-5}{proves step 5 of Theorem F}}
\expandafter\def\csname SBack@night:P2.9\endcsname{\hyperlink{step-I-F-6}{proves step 6 of Theorem F}}
\expandafter\def\csname SBack@night:P2.10\endcsname{\hyperlink{step-I-F-7}{proves step 7 of Theorem F}}
\expandafter\def\csname SBack@night:P2.11\endcsname{\hyperlink{step-I-F-7}{proves step 7 of Theorem F}}
\expandafter\def\csname SBack@push:src:D.1.3\endcsname{\hyperlink{step-I-G-1}{proves step 1 of Theorem G}}
\expandafter\def\csname SBack@push:src:D.2.1\endcsname{\hyperlink{step-I-G-2}{proves step 2 of Theorem G}}
\expandafter\def\csname SBack@eq:d-dual-exact\endcsname{\hyperlink{step-I-G-2}{proves step 2 of Theorem G}}
\expandafter\def\csname SBack@push:src:D.3.1\endcsname{\hyperlink{step-I-G-3}{proves step 3 of Theorem G}; \hyperlink{step-I-G-4}{proves step 4 of Theorem G}}
\expandafter\def\csname SBack@push:src:D.2.3\endcsname{\hyperlink{step-I-G-5}{proves step 5 of Theorem G}}
\expandafter\def\csname SBack@mapsI:src:A.1.1\endcsname{\href{S6_SERIES_GRIND_III.pdf\#nameddest=step-III-A-1}{proves step 1 of Theorem III.A}}
\expandafter\def\csname SBack@mapsI:src:A.2.1\endcsname{\href{S6_SERIES_GRIND_III.pdf\#nameddest=step-III-A-2}{proves step 2 of Theorem III.A}}

\begin{document}
\title{The $(3,4,\infty)$ family of complex six-spheres\\[.5em]I. Periods, moduli and the two ends}
\author{AI-produced mathematical text}
\date{19 November 2026}
\maketitle
\begin{abstract}
We study the holomorphic $(3,4,\infty)$ family through its normalized
elliptic periods, integral monodromy and cusp charts. The exact parameter
disc is $0<|w|<216$. The period norm determines its boundary, and the
marked classification gives the parameter equivalence, fibre strata,
automorphisms and real structures. We compute the differential-sheaf
and cohomological invariants with their explicit maps and compare the
two limiting regimes. At the puncture the deck and radial monodromies
give a proper-filling obstruction after every finite cover. At the
outer boundary the period calculations determine the limiting groups
and the failure of a global continuation. A credited appendix gives
our parameter-equivalence proof and the re-derived tangent and
deformation results in the source conventions. A holomorphic \v Cech
cocycle and a specified locally finite coordinate partition give global
Dolbeault representatives through all three special fibres. The
arithmetic period question is stated with its established reductions
and precise remaining exclusion.
\end{abstract}
\clearpage
\tableofcontents
\clearpage
\section{Results and their proofs in outline}\label{sec:results}
\subsection{The series and this paper}
We study the holomorphic families supplied by the marked
$(3,4,\infty)$ torus-fibred construction and the maps, degenerations
and manifolds obtained from its explicit periods and gluing data.
The source lattice, finite twists, cusp fan, line convention and
parameter normalization remain part of each construction. The five
papers form one series; the table gives their main conclusions and
their cross-paper inputs.
\begin{center}\small
\begin{tabularx}{\textwidth}{@{}lXl@{}}\toprule
Paper & Main result & Uses\\\midrule
I & Exact periods, moduli, fibre strata and the two ends & ---\\
II & Corrected cotangent argument and the sign for all degenerations & I, IV\\
III & Marked map classes, integral covers and their smooth outputs & I, II, IV, V\\
IV & Physical-circle return classes and the complete retained fan calculations & I, II, III, V\\
V & Resolved fibre products and the marked rank-six fourfold family & I, III, IV\\
\bottomrule\end{tabularx}\end{center}
\paragraph{The whole series: conclusions and mechanisms.}
In Paper~I the pointed Weierstrass map transports the section,
invariant differential and period lattice to a theta expression;
its exact norm gives the disc $0<|w|<216$, and its marked finite
and cusp quotients give the moduli, real structures and two endpoint
obstructions. Paper~II retains the cotangent-torsion connecting map
in the Campana--Demailly--Peternell calculation and uses normalization,
conductor, monodromy and parity maps to establish the sign for every
torus degeneration in its stated scope. Paper~III compares the actual
normal frames and integral attaching maps, determining the quotient
homotopy classes and the smooth invariants of the constructed covers
and elliptic-profile outputs. Paper~IV carries the physical circle
line through its Hermitian dual and cap maps; the resulting integral
characteristic formula gives the return classes in $\Theta_7$ and
the classes reached by positive covers. Paper~V combines a marked
Kuhn resolution with a rank-six exterior-square period construction;
the integral boundary calculation and oriented smooth classification
give the fourfolds, while the certified period minimum determines
the exact family half-plane. Every conclusion retains the source
section, finite twists, fan, orientation and parameter marking named
in its full statement.
% Proposed Paper I frontmatter.  This file reorganizes established body
% statements; it does not add verification events to those statements.
% Shared series overview, methodology, navigation and dependency graphic
% are supplied by the series assembler.

\subsection{The main theorems}
\hypertarget{tag-03JS}{}\label{series:I:main-theorems}

We use the Claude--Alp\"oge construction with its specified periodic
cusp fan, clockwise finite root generators and affine translations.
Its base is $B=\Pp^1_t$, with special values $0,1,\infty$, and
$B^\circ=B\setminus\{0,1,\infty\}$.  The ordered homology lattice and
period matrix are
\[
 \Lambda=\Z\langle\widehat\gamma,\widehat u,
                  \widehat w,\widehat\delta\rangle,
 \qquad
 \Pi_c=\begin{pmatrix}6\mu&\tau&1&0\\
                       \beta_{\rm part}+c&\mu&0&1\end{pmatrix}.
\]
Write $d=\widehat\delta$, and use $\gamma,u,w,\delta$ for the dual
integral covectors when they occur in lattice forms.  The parameter
$w=e^{-\pi i c}$ and its square $u=w^2$ are distinguished by their
arguments or by the words ``parameter'' and ``covector.''  The
particular solution is fixed by
$(\beta_{\rm part}+\tau)(\infty)=0$; $\mu$ is the cusp-selected
zero of $\wp_\tau$.  A statement about $t=0$ or $t=1$ uses the
specified smooth finite root model.  Unless further data are
displayed, $X_c$ denotes the standard $(3,4)$ sphere member with
$v_1=\varepsilon=(1,2,-4,0)^{\mathsf t}$,
$v_2=-\varepsilon'=(-1,-3,3,0)^{\mathsf t}$ and the untwisted
cusp collar.  We write $e=(0,1)^{\mathsf t}=\Pi_c d$ for the
invariant vertical vector.  The broader admissible translation and higher
multiplicity data are stated in Theorems B and F.

\hypertarget{main-I-A}{}
\begin{maintheorem}[The period, its exact domain, and its arithmetic norm]
\hypertarget{tag-03JT}{}\label{series:I:A}
\StatementRecord{series:I:A}{proved}{03JT}{\SeriesRef{I:weierstrass-comparison}, \SeriesRef{fam:src:B.1.1}, \SeriesRef{fam:src:C.1.1}, \SeriesRef{fam:src:C.2.3}, \SeriesRef{ari:C.4.2}, \SeriesRef{ari:C.4.3}, \SeriesRef{ari:C.4.4}, \SeriesRef{ari:C.6.1}, \SeriesRef{ari:C.6.3}, \SeriesRef{ari:C.8.15}, \SeriesRef{push:src:D.4.8}, \SeriesRef{ari:C.8.1}, \SeriesRef{ari:C.8.2}, \SeriesRef{ari:C.8.3}, \SeriesRef{ari:C.8.4}}
The missing period and the normalized norm are
\[
 \beta_{\rm part}=\frac{6i}{\pi}\Log
 \frac{\vartheta_1(\mu,\tau)}{2\sqrt3\,\eta_D(\tau)},
 \qquad
 N(P)=\left|\frac{\vartheta_1(\mu,\tau)}{\eta_D(\tau)}\right|
           e^{-\pi(\Im\mu)^2/\Im\tau},
 \quad P=[\mu].
\]
The displayed logarithm is the branch fixed by the cusp condition.
Putting $L=\log N$ and
$D_c=\Im(\beta_{\rm part}+c)-6(\Im\mu)^2/\Im\tau$, one has
\[
 D_{
 \rm part}=\frac6\pi\log\frac{N}{2\sqrt3},\qquad
 \max N=\sqrt2\quad(t=1),\qquad
 \max D_{\rm part}=-\frac3\pi\log6.
\]
All ordinary fibres and both normalized finite root fibres are
compact complex tori exactly on
\[
 \Im c<H_*:=\frac3\pi\log6,
 \qquad 0<|w|<216.
\]
There is no coarse critical point of $L$ away from $t=1$;
$t=0$ is a noncritical $C^1$ point with $L_t(0)=7/384$.
For every angle $\theta$, $L(1+re^{i\theta})$ decreases strictly
from $\log\sqrt2$ to $-\infty$ as $r$ increases from zero.
Every level below the maximum is one Jordan curve strictly
star-shaped about $1$, real analytic except for a possible
$C^1$ point with nonzero tangent at $t=0$ on the level
$L=\tfrac14\log3$.  Its bounded complement is the strict
superlevel set; these sets form nested discs.

In any of the explicitly transported pointed Weierstrass models,
for $m\ge2$ and $mP\ne O$,
\[
 N(P)^{m^2}=\frac{N(mP)|\Delta_E|^{(m^2-1)/12}}{|\psi_m(P)|}.
\]
For torsion $P$ of order $n$ and $m\equiv\pm1\pmod n$ this gives
\[
 N(P)^{12}=\frac{|\Delta_E|}{|\psi_m(P)|^{12/(m^2-1)}}.
\]
In particular the five specified values are
\[
\begin{array}{c|c|c}
t&\operatorname{ord}(P)&N(P)^{12}\\\hline
0&3&27\\1&2&64\\32/5&4&10\\
64(47-27\sqrt3)/11&6&17+10\sqrt3\\
64(47+27\sqrt3)/11&6&10\sqrt3-17
\end{array}
\]
and the body gives the complete exact division-polynomial lists
for orders four through eight.  For every number field $K$ and
$t\in K\setminus\{0,1\}$, in the arithmetic model
$y^2=x^3-3t^3(t-1)x/4+t^4(t-1)^2/4$ with
$P=(0,t^2(t-1)/2)$,
\[
 \sum_{\sigma:K\hookrightarrow\C}\log N(P_{\sigma(t)})
 =\sum_{v\nmid\infty}n_v\lambda_v(P_t)
       -[K:\Q]\widehat h_O(P_t).
\]
The normalization is $\lambda_\sigma(P)=-\log N(P)$, with complex
embeddings counted separately.  For rational $t$, the positive
algebraic factor $A_t=\exp(\sum_p\lambda_p(P_t))$ gives
$N(P_t)=A_t e^{-\widehat h_O(P_t)}$.  At the eleven nonspecial
rational CM parameters $t=j/1728$ its values are
\[
\begin{array}{r|c@{\qquad}r|c}
j&A_t&j&A_t\\\hline
8000&\sqrt6&-3375&2\\
54000&2^{2/3}3^{1/4}&-32768&\sqrt3\\
287496&2^{3/4}&-884736&1\\
16581375&2&-12288000&3^{5/12}\\
-884736000&1&-147197952000&1\\
-262537412640768000&1&&
\end{array}
\]
These are exact height identities; no transcendence conclusion
at the stated non-torsion points is included.

\end{maintheorem}

\textit{Mechanism.} Theta transformation laws fix the additive
constant, division identities control the norm, and the local
height product formula fixes its arithmetic normalization.
\textit{Relation to literature.} The elliptic surface and section
come from \cite[Remark~3.12, p.~19]{AC2026} and
\cite[Definition~2.3]{Engel2026}; the exact coordinate transports
are retained in \ref{I:weierstrass-comparison} and
\ref{fam:src:B.1.1}.  The norm and height conclusions have the
specific scopes of the body results listed here.
\hyperlink{outline-I-A}{Proof in outline}; full statements:
\ref{fam:src:C.1.1}, \ref{fam:src:C.2.3},
\ref{ari:C.4.2}, \ref{ari:C.4.3}, \ref{ari:C.4.4},
\ref{ari:C.6.1}, \ref{ari:C.6.3}, \ref{ari:C.8.15},
\ref{push:src:D.4.8}.

\hypertarget{main-I-B}{}
\begin{maintheorem}[All discrete data, the parameter, and real forms]
\hypertarget{tag-03JU}{}\label{series:I:B}
\StatementRecord{series:I:B}{proved}{03JU}{\SeriesRef{fam:src:A.5.1}, \SeriesRef{fam:src:B.2.1}, \SeriesRef{fam:src:B.2.2}, \SeriesRef{fam:src:B.2.3}, \SeriesRef{fam:src:B.3.1}, \SeriesRef{ari:C.2.1}, \SeriesRef{fam:src:B.4.2}, \SeriesRef{mapsI:src:A.3.1}, \SeriesRef{mapsI:src:A.3.2}, \SeriesRef{mapsI:src:A.3.3}, \SeriesRef{fam:src:A.1.1}, \SeriesRef{push:src:E.1.2}}
For the fixed period solution and cusp fan, take
\[
 w_1=r_1\varepsilon+k_1d,\qquad
 w_2=r_2\varepsilon'+k_2d,\qquad
 3\nmid r_1,\quad2\nmid r_2,
 \quad p=12l_0-4r_1-3r_2,
\]
where $l_0$ is the first integral winding of the full holomorphic
cusp collar in the stated positive convention.  Two such members
are biholomorphic exactly when
\[
 n=c'-c\in\Z,\qquad p'=ap\ (a\in\{1,-1\}),\qquad
 n\equiv k_2-k'_2\pmod2.
\]
Every biholomorphism is over $B$, and the map with sign $a$ has
fibre-linear part $aI_2$.  The complete invariants are
\[
 \bigl(|p|,e^{-\pi i(c+k_2)}\bigr),\qquad
 |p|\ge1,\quad\gcd(|p|,6)=1,
\]
and every displayed integer occurs.  The additional analytic
collar part and its second integral winding contribute no further
invariant.  For sphere members, $|p|=1$, a fixed datum has period
subgroup $2\Z$.  At each fixed admissible $u=w^2$ there are
exactly two sphere isomorphism classes, exchanged by a positive
circuit of $u$ about zero.

For positive $k,l$ with $\gcd(3k,4l)=1$, put $m_1=3k,m_2=4l$.
Use Viaclovsky's sphere data
$w_j=r_j\varepsilon_j+s_jd$, $\gcd(r_j,m_j)=1$,
$\varepsilon_1=\varepsilon$, $\varepsilon_2=\varepsilon'$, and
$p=m_1m_2l_0-m_2r_1-m_1r_2=\pm1$.
Two members are biholomorphic precisely when their multiplicity
pairs agree and, for some $a=\pm1$ and integer $n=c'-c$,
\[
 p'=ap,\quad r'_j\equiv ar_j\pmod{m_j},\quad
 nr'_1\equiv as_1-s'_1\pmod{k},\quad
 nr'_2\equiv as_2-s'_2\pmod{2l}.
\]
The fixed-data period subgroup is $2kl\Z$.  At fixed admissible
$u$ there are exactly $2kl$ classes for this multiplicity pair,
and a parameter circuit permutes them in one cycle.

For the standard family the reduced finite fibres have Albanese
nomes
\[
 q_2=iw/216,\qquad q_1=-w^2/48^3=(27/64)q_2^2.
\]
An antiholomorphic diffeomorphism $X_c\to X_{c'}$ exists exactly
when $c'-(1-\bar c)\in2\Z$, equivalently $w'=-\bar w$.
The members with real structures are precisely
$\Re c\in\tfrac12+\Z$.  On each,
\[
 \operatorname{Aut}^{\pm}(X_c)=\C^*\rtimes\Z/2,
 \qquad z\longmapsto1/\bar z,
\]
and there are two holomorphic-conjugacy classes of antiholomorphic
involutions.  For negative imaginary $w$ both fixed manifolds
are $S^3$, with the semifree circle action having the circle
$D_2^{\sigma_c}$ as fixed set.  For positive imaginary $w$ one
class is fixed-point free and the other has fixed manifold $S^3$
with the free Hopf circle action, up to the stated circle and
orientation conventions.

\end{maintheorem}

\textit{Mechanism.} Integral centralizers and finite norm maps
retain the cusp winding and the parity obstruction; the marked
reflection and cap-meridian calculation determine the real forms.
\textit{Relation to literature.} The construction is that of
\cite{AC2026,Engel2026}; the higher input is Viaclovsky's specified
root construction.  The standard exact $2\Z$ classification is
also Theorem~1.4 of \cite{DuEtAl2026v1}; the version and parameter
comparison is retained in \ref{fam:src:A.5.1} and the credited
appendix, including the constant translation between particular
period solutions.
\hyperlink{outline-I-B}{Proof in outline}; full statements:
\ref{fam:src:B.2.1}, \ref{fam:src:B.2.2}, \ref{fam:src:B.2.3},
\ref{fam:src:B.3.1}, \ref{ari:C.2.1}, \ref{fam:src:B.4.2},
\ref{mapsI:src:A.3.1}, \ref{mapsI:src:A.3.2},
\ref{mapsI:src:A.3.3}.

\hypertarget{main-I-C}{}
\begin{maintheorem}[The monodromy quotient and every fibre's algebraic dimension]
\hypertarget{tag-03JV}{}\label{series:I:C}
\StatementRecord{series:I:C}{proved}{03JV}{\SeriesRef{fam:src:F.1.1}, \SeriesRef{I:monodromy-fundamental-set}, \SeriesRef{fam:src:F.2.2}, \SeriesRef{fam:src:F.2.3}, \SeriesRef{fam:src:F.3.1}, \SeriesRef{fam:src:D.1.4}, \SeriesRef{fam:src:D.1.6}}
Let $\psi=6\gamma\wedge\delta+u\wedge w$ on $\Lambda$, and
let $P_{d,\Z}$ fix $d$ as a vector in its integral symplectic
group.  Write $J_2=\left(\begin{smallmatrix}0&1\\-1&0\end{smallmatrix}\right)$.
Every member has unique block coordinates
\[
 g(B,x,t)=\begin{pmatrix}1&0&0\\6x&B&0\\
                   t&-x^{\mathsf t}J_2B&1\end{pmatrix},
 \quad B\in\operatorname{SL}_2(\Z),\ x\in\Z^2,\ t\in\Z.
\]
For $S=\left(\begin{smallmatrix}0&-1\\1&0\end{smallmatrix}\right)$,
$T=\left(\begin{smallmatrix}1&1\\0&1\end{smallmatrix}\right)$,
let $\alpha:\operatorname{SL}_2(\Z)\to\Z/12$ be the homomorphism
with $\alpha(S)=9$, $\alpha(T)=1$.  The actual monodromy group is
\[
 \Gamma=\{g(B,x,t):t+6(x_1+x_2+x_1x_2)-\alpha(B)=0\pmod{12}\}.
\]
It has index 12 in $P_{d,\Z}$ and index 24 in $P_\Z$, the group preserving
$\Z d\subset\ker\gamma\subset\Lambda$.  In the real parabolic,
a fundamental set is given by the classical modular region for
$\tau$, the Iwasawa angle $0\le\theta<\pi$,
$0\le x_1,x_2<1$ and $0\le t<12$, with the actual modular
lifts $g(S,0,9)$, $g(T,0,1)$, $g(-I,0,6)$ and the Heisenberg
face identifications
\[
 (x_1,x_2,t)\mapsto(x_1+1,x_2,t-6-6x_2),\quad
 (x_1,x_2,t)\mapsto(x_1,x_2+1,t-6+6x_1).
\]

On $\mathcal D=\operatorname{Sp}_4(\R)/\operatorname U(1,1)$ put
$q_d(J)=\psi(d,Jd)$.  The largest $\Gamma$-invariant open domain
of proper discontinuity is $\mathcal O=\{q_d\ne0\}$;
the same assertion holds for $P_\Z$-invariant open domains and
its action.  The null locus has a dense set of
points with infinite monodromy stabilizer.  Each sign region has
a Hausdorff complex analytic orbifold quotient.  In the family
chart, $\mathcal O_+$ has $\Im\tau>0$, $D_c<0$ and
$q_d=-6/D_c$.  For every admissible $c$ the period map is
injective and immersive on $B^\circ$ and extends to a proper
holomorphic embedding
\[
 \C=B\setminus\{\infty\}\lhook\joinrel\longrightarrow
 \Gamma\backslash\mathcal O_+.
\]

For every smooth fibre its N\'eron--Severi group is the integral
kernel in $\Z^6$ of
\[
 E_{12}-E_{13}\tau+(6E_{23}-E_{14})\mu+E_{24}\beta_c
             +E_{34}(6\mu^2-\tau\beta_c)=0.
\]
With its complex orientation,
$E^2=-2\operatorname{Pf}(E)$, where
$\operatorname{Pf}(E)=E_{12}E_{34}-E_{13}E_{24}+E_{14}E_{23}$.
Let $\mathcal E_c$ be the image on $B^\circ$ of the union of
these zero sets for $E\ne0$ with $\operatorname{Pf}(E)=0$,
and let $\mathcal A_c$ use $\operatorname{Pf}(E)<0$.
They are exactly the elliptic-quotient locus and abelian-surface
locus, respectively, and
\[
 a(F_t)=2\ (t\in\mathcal A_c),\quad
 a(F_t)=1\ (t\in\mathcal E_c\setminus\mathcal A_c),\quad
 a(F_t)=0\ (t\notin\mathcal E_c\cup\mathcal A_c).
\]
These countable unions classify all quotient maps, including
those that mix the two displayed coordinates.  The invariant
class $\eta=6\gamma\wedge\delta+u\wedge w$ has square $-12$;
a fibre with $\operatorname{NS}=\Z\eta$ has no compact curves.

\end{maintheorem}

\textit{Mechanism.} The exact congruence kernel controls the real
parabolic quotient, while the integral $(1,1)$ equation and its
Hermitian form determine every fibre stratum with the actual
orientation.
\textit{Relation to literature.} The proper-action question is
posed in \cite[Remark~8.4, pp.~23--24]{Engel2026}.  The stated
answer retains the precise monodromy subgroup and full invariant
open set; it does not rely on replacing that subgroup by its
rational or real closure.
\hyperlink{outline-I-C}{Proof in outline}; full statements:
\ref{fam:src:F.1.1}, \ref{I:monodromy-fundamental-set},
\ref{fam:src:F.2.2}, \ref{fam:src:F.2.3}, \ref{fam:src:F.3.1},
\ref{fam:src:D.1.4}, \ref{fam:src:D.1.6}.

\hypertarget{main-I-D}{}
\begin{maintheorem}[Cohomology and the ordinary relative differential modules]
\hypertarget{tag-03JW}{}\label{series:I:D}
\StatementRecord{series:I:D}{proved}{03JW}{\SeriesRef{inv:C3.8}, \SeriesRef{inv:A2.2}, \SeriesRef{inv:A3.1}, \SeriesRef{inv:C1.1}, \SeriesRef{inv:C2.1}, \SeriesRef{inv:C3.3}, \SeriesRef{inv:A4.1}, \SeriesRef{inv:A4.2}, \SeriesRef{inv:A2.4}, \SeriesRef{inv:A1.0}}
Every admissible standard sphere member has the following
Dolbeault, Bott--Chern and Aeppli dimensions, with columns
$p=0,1,2,3$ and rows $q=3,2,1,0$:
\[
\begin{array}{c|rrrr}q\backslash p&0&1&2&3\\\hline
3&0&0&0&1\\2&0&1&2&1\\1&1&2&1&0\\0&1&0&0&0
\end{array}_{\bar\partial}\quad
\begin{array}{c|rrrr}q\backslash p&0&1&2&3\\\hline
3&0&0&1&1\\2&0&1&2&1\\1&0&2&1&0\\0&1&0&0&0
\end{array}_{BC}\quad
\begin{array}{c|rrrr}q\backslash p&0&1&2&3\\\hline
3&0&0&0&1\\2&0&1&2&0\\1&1&2&1&0\\0&1&1&0&0
\end{array}_{A}.
\]
The Fr\"olicher sequence satisfies $E_1\ne E_2=E_\infty$;
the only nonzero $d_1$ have rank one at
$(p,q)=(0,1),(1,1),(1,2),(2,2)$.  The Bott--Chern algebra is
$\C\cdot1\oplus\mathcal N$ with $\mathcal N^2=0$.

For the ordinary relative differentials the local presentations
retain both the degree and the full fibre multiplicity.  For
$R=\C\{x,y,z\}$ and $t=xy$,
\[
 \Omega^1_{R/\C\{t\}}=R^3/R(y,x,0)\simeq(x,y)\oplus R,
 \quad\Omega^2\simeq(x,y)\oplus R/(x,y),
 \quad\Omega^3\simeq R/(x,y).
\]
For $t=xyz$, $\Omega^1=R^3/R(yz,xz,xy)$ is torsion-free of
rank two and not locally free on the double locus.  With
$I=(xy,xz,yz)$,
\[
 0\longrightarrow
 \frac{Ru\oplus Rv}{(xu,\,y(u-v),\,zv)}
 \longrightarrow\Omega^2\xrightarrow{\wedge dt}I
 \longrightarrow0,
 \qquad \Omega^3=R/I.
\]
The left term is exactly the torsion and is annihilated by $I$.
For $t=s^m$ in $R=\C\{s,y_1,y_2\}$, the three modules are
\[
 R^2\oplus R/(s^{m-1})\,ds,\qquad
 R\,dy_1\wedge dy_2\oplus(R/(s^{m-1}))^2,
 \qquad R/(s^{m-1})\,ds\wedge dy_1\wedge dy_2.
\]
Consequently the degree-one torsion of $\Omega^1_{X/B}$ is
supported exactly on $S_1\cup S_2$, while degree two also has
torsion along the cusp double locus.

There are specified global forms
$G\in\mathcal A^{0,1}(X)$ and $Q\in\mathcal A^{1,1}(X)$ satisfying
\[
 \bar\partial G=\bar\partial Q=0,\qquad
 [G]=g_\gamma,\qquad i_b^*[Q]=\eta,\qquad
 \exp(12[G])=f^*\mathcal O_{\mathbb P^1}(1).
\]
On the specified common refinement of the original atlas they are
$G=\bar\partial\sum w_nf_{a,n}$ and
$Q=\tfrac1{12}\bar\partial\sum w_nc_{a,n}$, with the locally finite
weights, holomorphic scalar cocycle and holomorphic one-form cocycle
of \eqref{inv:eq:atlas-final-forms}.  The elementary derivative
recursion specifies every coefficient through both multiple fibres
and every cusp stratum.  The form $Q$ is complex; a real
$\bar\partial$-closed form cannot have this fibre class.

\end{maintheorem}

\textit{Mechanism.} Logarithmic residues and root characters
compute the direct images; the evaluated connecting scalar fixes
the remaining Hodge integer.  The ordinary modules follow from
the exact quotient by $dt$ and its exterior powers.
Twelfth-power theta lines, cusp Cartier frames and annular
holomorphic connections give the cocycles for the coefficient
construction; the specified locally finite weights evaluate them.
\textit{Relation to literature.} The holomorphic-form vanishing
and Hodge pattern with one undetermined integer are credited to
\cite[Theorem~9.1(vi), p.~65; Proposition~9.20, p.~80]{AC2026}.
The local presentations correct the degree-one torsion assertion
in its Remark~9.21, p.~81, and qualify Remark~1.3, p.~5, by
degree.  The coefficient construction in \ref{inv:C3.8} uses the
original period, finite affine and toric charts; its source comparison
and exact cocycle maps are given there.
\hyperlink{outline-I-D}{Proof in outline}; full statements:
\ref{inv:A2.2}, \ref{inv:A3.1}, \ref{inv:C1.1},
\ref{inv:C2.1}, \ref{inv:C3.3}, \ref{inv:A4.1},
\ref{inv:A4.2}, \ref{inv:C3.8}.

\hypertarget{main-I-E}{}
\begin{maintheorem}[Every Picard twist and the deformation parameter]
\hypertarget{tag-03JX}{}\label{series:I:E}
\StatementRecord{series:I:E}{proved}{03JX}{\SeriesRef{inv:D4.1}, \SeriesRef{inv:D4.2}, \SeriesRef{inv:D4.3}, \SeriesRef{inv:D4.4}, \SeriesRef{inv:D4.5}, \SeriesRef{fam:st:A.5.2}, \SeriesRef{fam:src:B.4.1}, \SeriesRef{push:src:E.2.1}, \SeriesRef{inv:C3.2}, \SeriesRef{inv:C3.1}, \SeriesRef{push:src:E.1.2}}
Use the primitive coordinate $M_z=\exp(z\mathfrak g)$ on
$\operatorname{Pic}(X)=\C$, so
$M_n=\mathcal O(n(S_1-S_2))$ for $n\in\Z$.
Let $r_1(n)\in\{0,1,2\}$ and $r_2(n)\in\{0,1,2,3\}$ be the
residues of $n$ modulo 3 and $3n$ modulo 4, respectively, and set
\label{layout:I:1}
\[
\begin{gathered}
d(n)=\frac{n-4r_1(n)-3r_2(n)}{12},\quad
 \\
U(n)=\max\{d(n)+1,0\},\quad V(n)=\max\{-d(n)-1,0\},
 \\
\quad\epsilon_n=\begin{cases}0&3\mid n,\\1&3\nmid n.\end{cases}
\end{gathered}
\]
Then
\[
\begin{aligned}
h^0(T_X\otimes M_n)&=U(n)+U(n-1),\\
h^1(T_X\otimes M_n)&=V(n)+V(n-1)+U(n)+U(n-1)+U(n+2),\\
h^2(T_X\otimes M_n)&=V(n)+V(n-1)+V(n+2)+U(n+6)+\epsilon_n,\\
h^3(T_X\otimes M_n)&=V(n+6).
\end{aligned}
\]
For $z\notin\Z$, let $\mathcal R_z$ be the set of distinct
regular base points for which integers $n,p,q,r$ satisfy
\[
 q\tau+r\mu=p,\qquad z=n-6q\mu-r\beta_c.
\]
It is finite; write $N_z=|\mathcal R_z|$.  Put
$\rho=e^{\pi i/3}$, $\beta_j=\beta_c(z_j)$ at the specified
order-three and order-four marked root centers.  Define
$\ell_1(z)=1+\epsilon_n$ at
$z=n+k(3\beta_1+4-2\rho)$, $k\ne0$, and zero elsewhere;
define $\ell_2(z)=1$ at
$z=n+k(2\beta_2+3-3i)$, $k\ne0$, $n\equiv1,2\pmod4$,
and zero elsewhere.  With
$J_z=2N_z+\ell_1(z)+\ell_2(z)$,
\[
 (h^0,h^1,h^2,h^3)(T_X\otimes M_z)=(0,J_z,1+J_z,0).
\]
The actual cohomology jump loci are locally discrete in this
Picard line, and the generic tuple is $(0,0,1,0)$.

The untwisted tangent tuple is $(1,1,1,0)$, the holomorphic
automorphism group of the standard member is $\C^*$, and its
Kuranishi germ is a smooth reduced line.  For every sphere member
with coprime multiplicities $(3k,4l)$ and the data of Theorem B,
put $z_0=kl$, $d_0=3k+4l$.  On the analytic weighted projective root stack
$\mathfrak B_{k,l}=\Pp(3k,4l)$ the actual root lift satisfies
\[
 \mathcal E=\mathcal O\oplus\mathcal O(-z_0),\quad
 R^q\bar f_*\mathcal O_X=
 \mathcal O,\ \mathcal O\oplus\mathcal O(-z_0),\ \mathcal O(-z_0),
 \quad K_X=\bar f^*\mathcal O(z_0-d_0).
\]
Its tangent tuple is again $(1,1,1,0)$ and its Kuranishi germ is
a smooth reduced line.  For the groupoid derivative
$\kappa(\Phi)=(\partial_c\Phi)\Phi^{-1}$ the classes of
$\partial_c$ and $\partial_u$ are
$e\otimes[\gamma]$ and $-(2\pi i u)^{-1}e\otimes[\gamma]$,
respectively, and are nonzero.  Every stated fixed-data period
family is locally miniversal.

\end{maintheorem}

\textit{Mechanism.} Root-line section counts and the exact local
Koszul modules calculate every Picard twist.  The nonzero
Kodaira--Spencer class supplies the actual coordinate on the
one-dimensional Kuranishi germ.
\textit{Relation to literature.} The standard tangent, automorphism
and Kuranishi conclusions are credited to
\cite[Theorems~1.3--1.4, p.~3; Theorem~6.12, pp.~81--82;
Theorem~6.16, p.~85]{DuEtAl2026v1}.  The version-specific
re-derivations and parameter transport are preserved; version~2's
omission of the deformation section changes none of that credit.
\hyperlink{outline-I-E}{Proof in outline}; full statements:
\ref{inv:D4.1}, \ref{inv:D4.2}, \ref{inv:D4.3},
\ref{inv:D4.4}, \ref{inv:D4.5}, \ref{fam:st:A.5.2},
\ref{fam:src:B.4.1}, \ref{push:src:E.2.1}.

\hypertarget{main-I-F}{}
\begin{maintheorem}[The puncture and the obstruction to every proper filling]
\hypertarget{tag-03JY}{}\label{series:I:F}
\StatementRecord{series:I:F}{proved}{03JY}{\SeriesRef{night:P2.2}, \SeriesRef{night:P2.3}, \SeriesRef{night:P2.4}, \SeriesRef{night:P2.5}, \SeriesRef{night:P2.6}, \SeriesRef{night:P2.7}, \SeriesRef{night:P2.9}, \SeriesRef{night:P2.10}, \SeriesRef{night:P2.11}}
For an admissible $(3,4)$ logarithmic datum put
$b=4l_1+3l_2$, $p=12l_0-b$, $a=-p$, with
$3\nmid l_1$ and $l_2$ odd.  After $w=r^6$ there is a proper
holomorphic partial filling over
\[
 (\Pp^1\times\Delta_r)\setminus\{(\infty,0)\}
\]
with smooth total space and reduced normal-crossing parameter
central divisor.  Over a regular finite base point the central
fibre is a cycle of twelve ruled surfaces, with component cover
\[
 y^{12}=t^{4l_1}(t-1)^{3l_2}
\]
and rotations $4l_1,3l_2,-b$ modulo 12.
On $q=Q^{12m}$ and $r=R^n$, the difference between the actual
cusp-extending and parameter-extending origins has deck
differences
\[
 am\gamma-6am^2\delta,\qquad
 an\gamma-6an(n-1)\delta,
\]
so the radial class modulo $\Z\widehat\delta$ is
$-pn\widehat\gamma$ in these fixed conventions.

The punctured family has no proper Hausdorff complex-analytic
filling at $w=0$ after any finite parameter cover.  The conclusion
permits arbitrary singularities, multiplicities and extra central
components, and holds for every admissible datum, including every
sphere member.  It also holds for every admissible datum of the
coprime $(3k,4l)$ construction.  In that case
\[
 N=12kl,\quad b=4l\,l_1+3k\,l_2,\quad a=b-Nl_0,
 \qquad \gcd(l_1,3k)=\gcd(l_2,4l)=1,
\]
the corresponding partial model has $N$ components, and the two
displayed deck differences have $N/2$ in place of 6.

For the untwisted elliptic-surface quotient $X(u)$, whose
logarithmic translations are zero, the actual origins have zero
corner winding.  That family admits a proper flat Hausdorff
analytic filling at $u=0$, obtained by the two-projective-line-bundle
plumbing, proper modification and finite quotient in the body.
The particular one-component central model proposed in
\cite[Conjecture~8.2, p.~22]{Engel2026} cannot be the central
fibre of a proper flat smoothing to these $X(u)$: its virtual
Euler characteristic is zero, whereas $e(X(u))=2$.

\end{maintheorem}

\textit{Mechanism.} The intrinsic canonical system recovers the
marked base over a proposed filling.  Finite-orbit vanishing kills
horizontal meridians, while commuting meridians propagate their
ratio across the connected compact corner to the nonzero cusp
coefficient $a$.
\textit{Relation to literature.} The comparison retains Engel's
conjectural status, its specified central model, and the actual
maps between the elliptic-surface and period descriptions.  The
higher family is Viaclovsky's construction; the obstruction is
applied to its precise primitive logarithmic data.
\hyperlink{outline-I-F}{Proof in outline}; full statements:
\ref{night:P2.2}, \ref{night:P2.3}, \ref{night:P2.4},
\ref{night:P2.5}, \ref{night:P2.6}, \ref{night:P2.7},
\ref{night:P2.9}, \ref{night:P2.10}, \ref{night:P2.11}.

\hypertarget{main-I-G}{}
\begin{maintheorem}[The outer boundary and all marked additive limits]
\hypertarget{tag-03JZ}{}\label{series:I:G}
\StatementRecord{series:I:G}{proved}{03JZ}{\SeriesRef{push:src:D.1.3}, \SeriesRef{push:src:D.2.1}, \SeriesRef{push:src:D.2.3}, \SeriesRef{push:src:D.3.1}, \SeriesRef{eq:d-dual-exact}}
For every $w_0$ with $|w_0|=216$ and every connected open
neighbourhood $U$ of $w_0$, there is no flat proper holomorphic
map from a reduced complex space to $U$ whose fibre at each
$w\in U\cap\{0<|w|<216\}$ is biholomorphic to $X_w$.
The intrinsic divisor $2S_2$ and its Albanese nome $iw/216$
retain the invariant $J(iw/216)$, where
$J(q)=q^{-1}+744+196884q+\cdots$ has the unit circle as a
natural boundary.

Fix $c_*$ with $\Im c_*=H_*$, put $a=\Re c_*+3/2$, and use the
specified order-four root coordinate $h$ at $t=1$ and its fixed
complex-linear fibre coordinates $(\eta,\xi)$.  In the transformed
integral period basis the four vectors are
\[
 (1,1/2),\quad(\tau,1/2+h),\quad(0,1),
 \quad(6h,a+B(h)+c-c_*),
 \qquad B(h)=-3ih^2+O(h^4).
\]
For the total root quotient at fixed $c=c_*$, two central points
are nonseparated exactly when
\[
 \Re(\xi'-\xi)\in G_a=\tfrac12\Z+a\Z\quad(a\in\Q);
\]
every central pair is nonseparated for $a\notin\Q$.  Distinct
noncentral points and distinct base fibres are separated.
The Hausdorff reflection keeps the noncentral complex tori and
has central fibre $\R/G_a\simeq S^1$ for rational $a$, and one
point for irrational $a$.  For reduced $a=m/n$, $n>0$,
$G_a=\operatorname{lcm}(2,n)^{-1}\Z$.

For sequences $c_n\to c_*$ with $\Im c_n<H_*$ and $h_n\to0$,
write
\[
 \delta_n=-D_{c_n}(h_n)>0,\quad
 V_n=\delta_n^{-1}(6\Re h_n,6\Im h_n,\Re(c_n-c_*)),
\]
\[
 L_a=\{(p,q,r):r\in\Z,\ p+r/2\in\Z,
                      \ q+r/2\in\Z,\ ar\in\Z\}.
\]
It has rank three for rational $a$, and equals
$\Z^2\times\{0\}$ for irrational $a$.  Each convergent
subsequence of the closed period lattices has limit
\[
 \Gamma_\infty=\{(\eta,\xi):p\Re\eta+q\Im\eta+r\Re\xi
       +\ell(p,q,r)\Im\xi\in\Z\text{ for every }(p,q,r)\in K\},
\]
where $K\subset L_a$ is primitive,
$K=\{\lambda:\lambda\cdot V_n\text{ has a finite limit}\}$,
and $\ell(\lambda)=\lim\lambda\cdot V_n$.
For labels outside $K$ the absolute values tend to infinity.
Conversely every such primitive $K$ and homomorphism
$\ell:K\to\R$ occurs.  The resulting pointed additive quotient
limits are real tori of every dimension zero through three for
rational $a$, and zero through two for irrational $a$.

\end{maintheorem}

\textit{Mechanism.} The Albanese invariant detects the outer
natural boundary.  At a boundary root chart, the exact dual
period matrix and its quadratic cancellation describe every
surviving integral covector, hence both nonseparation and all
closed-group limits.
\textit{Relation to literature.} The elliptic-surface description
is transported through the marked quotient dictionary; the body
states the precise Pontryagin--Chabauty duality input used for
closed-group convergence.  The Hausdorff reflection and additive
limits retain their stated topological categories and maps.
\hyperlink{outline-I-G}{Proof in outline}; full statements:
\ref{push:src:D.1.3}, \ref{push:src:D.2.1},
\ref{push:src:D.2.3}, \ref{push:src:D.3.1}.


\paragraph{Further results.}
\begingroup\small
\begin{longtable}{@{}ll@{\hspace{1cm}}ll@{}}
\toprule Result & Tag & Result & Tag\\\midrule\endhead
\SeriesRef{fam:src:B.1.2} & \texttt{01Q8} & \SeriesRef{push:src:B.4.25} & \texttt{030U}\\
\SeriesRef{fam:src:B.1.3} & \texttt{01Q9} & \SeriesRef{push:src:B.4.26} & \texttt{030V}\\
\SeriesRef{fam:src:B.1.4} & \texttt{01QA} & \SeriesRef{push:src:B.4.27} & \texttt{030W}\\
\SeriesRef{fam:src:C.1.2} & \texttt{01QI} & \SeriesRef{push:src:B.4.28} & \texttt{030X}\\
\SeriesRef{fam:src:C.2.1} & \texttt{01QJ} & \SeriesRef{push:src:B.4.29} & \texttt{030Y}\\
\SeriesRef{fam:src:C.2.2} & \texttt{01QK} & \SeriesRef{push:src:B.4.30} & \texttt{0310}\\
\SeriesRef{ari:C.1.1} & \texttt{000I} & \SeriesRef{push:src:B.4.31} & \texttt{0311}\\
\SeriesRef{ari:C.1.2} & \texttt{000J} & \SeriesRef{push:src:B.4.32} & \texttt{0312}\\
\SeriesRef{ari:C.1.3} & \texttt{000K} & \SeriesRef{push:src:B.4.33} & \texttt{0313}\\
\SeriesRef{ari:C.4.1} & \texttt{000P} & \SeriesRef{push:src:B.5.1} & \texttt{031A}\\
\SeriesRef{ari:C.4.5} & \texttt{000T} & \SeriesRef{push:src:D.1.1} & \texttt{032A}\\
\SeriesRef{fam:src:A.1.2} & \texttt{01Q0} & \SeriesRef{push:src:D.1.2} & \texttt{032B}\\
\SeriesRef{fam:src:A.2.1} & \texttt{01Q1} & \SeriesRef{push:src:D.2.2} & \texttt{032E}\\
\SeriesRef{fam:src:A.3.1} & \texttt{01Q2} & \SeriesRef{push:src:D.2.4} & \texttt{032G}\\
\SeriesRef{fam:src:A.4.1} & \texttt{01Q3} & \SeriesRef{push:src:D.3.2} & \texttt{032I}\\
\SeriesRef{fam:src:A.4.2} & \texttt{01Q4} & \SeriesRef{push:src:D.4.1} & \texttt{032J}\\
\SeriesRef{fam:src:A.4.3} & \texttt{01Q5} & \SeriesRef{push:src:D.4.2} & \texttt{032K}\\
\SeriesRef{push:src:E.1.1} & \texttt{032R} & \SeriesRef{push:src:D.4.3} & \texttt{032L}\\
\SeriesRef{push:src:E.1.3} & \texttt{032T} & \SeriesRef{push:src:D.4.4} & \texttt{032M}\\
\SeriesRef{fam:src:F.1.2} & \texttt{01RA} & \SeriesRef{push:src:D.4.5} & \texttt{032N}\\
\SeriesRef{fam:src:F.2.1} & \texttt{01RB} & \SeriesRef{push:src:D.4.6} & \texttt{032O}\\
\SeriesRef{fam:src:F.3.2} & \texttt{01RF} & \SeriesRef{push:src:D.4.7} & \texttt{032P}\\
\SeriesRef{fam:src:F.3.3} & \texttt{01RG} & \SeriesRef{night:P2.1} & \texttt{02R9}\\
\SeriesRef{fam:src:D.1.1} & \texttt{01QR} & \SeriesRef{night:P2.8} & \texttt{02RI}\\
\SeriesRef{fam:src:D.1.2} & \texttt{01QS} & \SeriesRef{inv:A1.1} & \texttt{01RR}\\
\SeriesRef{fam:src:D.1.3} & \texttt{01QT} & \SeriesRef{inv:A1.2} & \texttt{01RS}\\
\SeriesRef{fam:src:D.1.5} & \texttt{01QV} & \SeriesRef{inv:A1.3} & \texttt{01RT}\\
\SeriesRef{fam:src:D.1.7} & \texttt{01QX} & \SeriesRef{inv:A2.1} & \texttt{01RU}\\
\SeriesRef{fam:src:D.2.1} & \texttt{01QY} & \SeriesRef{inv:A2.3} & \texttt{01RW}\\
\SeriesRef{fam:src:D.2.2} & \texttt{01QZ} & \SeriesRef{inv:A3.2} & \texttt{01RZ}\\
\SeriesRef{fam:src:D.2.3} & \texttt{01R0} & \SeriesRef{inv:C1.2} & \texttt{01SB}\\
\SeriesRef{fam:src:C.3.1} & \texttt{01QM} & \SeriesRef{inv:C2.2} & \texttt{01SD}\\
\SeriesRef{fam:src:C.3.2} & \texttt{01QN} & \SeriesRef{inv:C3.4} & \texttt{01SH}\\
\SeriesRef{fam:src:C.3.3} & \texttt{01QO} & \SeriesRef{inv:C3.5} & \texttt{01SI}\\
\SeriesRef{fam:src:C.3.4} & \texttt{01QP} & \SeriesRef{inv:C3.6} & \texttt{01SJ}\\
\SeriesRef{fam:src:C.4.1} & \texttt{01QQ} & \SeriesRef{inv:C3.7} & \texttt{01SK}\\
\SeriesRef{mapsI:src:A.3.4} & \texttt{02PH} & \SeriesRef{push:src:E.2.2} & \texttt{0335}\\
\SeriesRef{push:src:C.3.5} & \texttt{0322} & \SeriesRef{push:src:E.2.3} & \texttt{033B}\\
\SeriesRef{fam:src:E.1.1} & \texttt{01R1} & \SeriesRef{push:src:E.3.1} & \texttt{033I}\\
\SeriesRef{fam:src:E.1.2} & \texttt{01R2} & \SeriesRef{inv:B1.0} & \texttt{01S2}\\
\SeriesRef{fam:src:E.1.3} & \texttt{01R3} & \SeriesRef{inv:B1.1} & \texttt{01S3}\\
\SeriesRef{fam:src:E.1.4} & \texttt{01R4} & \SeriesRef{inv:B1.2} & \texttt{01S4}\\
\SeriesRef{fam:src:E.1.5} & \texttt{01R5} & \SeriesRef{inv:B2.1} & \texttt{01S5}\\
\SeriesRef{fam:src:E.2.1} & \texttt{01R6} & \SeriesRef{inv:B2.2} & \texttt{01S6}\\
\SeriesRef{push:src:A.1.1} & \texttt{02ZH} & \SeriesRef{inv:B3.1} & \texttt{01S7}\\
\SeriesRef{push:src:A.1.2} & \texttt{02ZI} & \SeriesRef{inv:B3.2} & \texttt{01S8}\\
\SeriesRef{push:src:A.2.1} & \texttt{02ZJ} & \SeriesRef{inv:B3.3} & \texttt{01S9}\\
\SeriesRef{push:src:A.2.2} & \texttt{02ZK} & \SeriesRef{mapsI:src:A.1.1} & \texttt{02PA}\\
\SeriesRef{push:src:A.2.3} & \texttt{02ZL} & \SeriesRef{mapsI:src:A.1.2} & \texttt{02PB}\\
\SeriesRef{push:src:A.2.4} & \texttt{02ZM} & \SeriesRef{mapsI:src:A.1.3} & \texttt{02PC}\\
\SeriesRef{push:src:A.2.5} & \texttt{02ZN} & \SeriesRef{mapsI:src:A.2.1} & \texttt{02PD}\\
\SeriesRef{push:src:A.3.1} & \texttt{02ZO} & \SeriesRef{ari:C.2.2} & \texttt{000M}\\
\SeriesRef{push:src:A.3.2} & \texttt{02ZP} & \SeriesRef{ari:C.3.1} & \texttt{000N}\\
\SeriesRef{push:src:A.3.3} & \texttt{02ZQ} & \SeriesRef{ari:C.3.2} & \texttt{000O}\\
\SeriesRef{push:src:A.4.1} & \texttt{02ZR} & \SeriesRef{ari:C.5.1} & \texttt{000U}\\
\SeriesRef{push:src:A.4.2} & \texttt{02ZS} & \SeriesRef{ari:C.6.2} & \texttt{000W}\\
\SeriesRef{push:src:A.4.3} & \texttt{02ZT} & \SeriesRef{ari:C.6.4} & \texttt{000Y}\\
\SeriesRef{push:src:A.4.4} & \texttt{02ZU} & \SeriesRef{ari:C.6.5} & \texttt{000Z}\\
\SeriesRef{push:src:A.4.5} & \texttt{02ZV} & \SeriesRef{ari:C.6.R1} & \texttt{0010}\\
\SeriesRef{push:src:A.4.6} & \texttt{02ZW} & \SeriesRef{ari:C.6.R2} & \texttt{0011}\\
\SeriesRef{push:src:A.4.7} & \texttt{02ZX} & \SeriesRef{ari:C.6.R3} & \texttt{0012}\\
\SeriesRef{push:src:A.4.8} & \texttt{02ZY} & \SeriesRef{ari:C.6.R4} & \texttt{0013}\\
\SeriesRef{push:src:A.4.9} & \texttt{02ZZ} & \SeriesRef{ari:C.6.R5} & \texttt{0014}\\
\SeriesRef{push:src:B.1.1} & \texttt{0300} & \SeriesRef{ari:C.6.T1} & \texttt{0015}\\
\SeriesRef{push:src:B.1.2} & \texttt{0301} & \SeriesRef{ari:C.6.T2} & \texttt{0016}\\
\SeriesRef{push:src:B.2.1} & \texttt{0302} & \SeriesRef{ari:C.8.6} & \texttt{001T}\\
\SeriesRef{push:src:B.2.2} & \texttt{0303} & \SeriesRef{ari:C.8.7} & \texttt{001U}\\
\SeriesRef{push:src:B.2.3} & \texttt{0304} & \SeriesRef{ari:C.8.8} & \texttt{001V}\\
\SeriesRef{push:src:B.3.1} & \texttt{0305} & \SeriesRef{ari:C.8.9} & \texttt{001W}\\
\SeriesRef{push:src:B.3.2} & \texttt{0306} & \SeriesRef{ari:C.8.10} & \texttt{001K}\\
\SeriesRef{push:src:B.3.3} & \texttt{0307} & \SeriesRef{ari:C.8.11} & \texttt{001L}\\
\SeriesRef{push:src:B.3.4} & \texttt{0308} & \SeriesRef{ari:C.8.12} & \texttt{001M}\\
\SeriesRef{push:src:B.3.5} & \texttt{0309} & \SeriesRef{ari:C.8.13} & \texttt{001N}\\
\SeriesRef{push:src:B.3.6} & \texttt{030A} & \SeriesRef{ari:C.8.14} & \texttt{001O}\\
\SeriesRef{push:src:B.3.7} & \texttt{030B} & \SeriesRef{push:src:E.4.1} & \texttt{033L}\\
\SeriesRef{push:src:B.3.8} & \texttt{030C} & \SeriesRef{push:src:E.4.2} & \texttt{033M}\\
\SeriesRef{push:src:B.4.1} & \texttt{030D} & \SeriesRef{night:P5.7} & \texttt{02YS}\\
\SeriesRef{push:src:B.4.2} & \texttt{030O} & \SeriesRef{wb:thm:canonical-ring-recovery} & \texttt{03DM}\\
\SeriesRef{push:src:B.4.3} & \texttt{030Z} & \SeriesRef{wb:thm:retained-period-deformation} & \texttt{03DP}\\
\SeriesRef{push:src:B.4.4} & \texttt{0314} & \SeriesRef{wb:pshift:lem:centralizer} & \texttt{03CQ}\\
\SeriesRef{push:src:B.4.5} & \texttt{0315} & \SeriesRef{inv:D1.1} & \texttt{01SM}\\
\SeriesRef{push:src:B.4.6} & \texttt{0316} & \SeriesRef{inv:D1.2} & \texttt{01SN}\\
\SeriesRef{push:src:B.4.7} & \texttt{0317} & \SeriesRef{night:P5.1} & \texttt{02YM}\\
\SeriesRef{push:src:B.4.8} & \texttt{0318} & \SeriesRef{night:P5.2} & \texttt{02YN}\\
\SeriesRef{push:src:B.4.9} & \texttt{0319} & \SeriesRef{night:P5.3} & \texttt{02YO}\\
\SeriesRef{push:src:B.4.10} & \texttt{030E} & \SeriesRef{night:P5.4} & \texttt{02YP}\\
\SeriesRef{push:src:B.4.11} & \texttt{030F} & \SeriesRef{night:P5.5} & \texttt{02YQ}\\
\SeriesRef{push:src:B.4.12} & \texttt{030G} & \SeriesRef{night:P5.6} & \texttt{02YR}\\
\SeriesRef{push:src:B.4.13} & \texttt{030H} & \SeriesRef{ari:C.7.1} & \texttt{0017}\\
\SeriesRef{push:src:B.4.14} & \texttt{030I} & \SeriesRef{ari:C.7.2} & \texttt{001B}\\
\SeriesRef{push:src:B.4.15} & \texttt{030J} & \SeriesRef{ari:C.7.3} & \texttt{001C}\\
\SeriesRef{push:src:B.4.16} & \texttt{030K} & \SeriesRef{ari:C.7.4} & \texttt{001D}\\
\SeriesRef{push:src:B.4.17} & \texttt{030L} & \SeriesRef{ari:C.7.6} & \texttt{001F}\\
\SeriesRef{push:src:B.4.18} & \texttt{030M} & \SeriesRef{ari:C.7.7} & \texttt{001G}\\
\SeriesRef{push:src:B.4.19} & \texttt{030N} & \SeriesRef{ari:C.7.8} & \texttt{001H}\\
\SeriesRef{push:src:B.4.20} & \texttt{030P} & \SeriesRef{ari:C.7.9} & \texttt{001I}\\
\SeriesRef{push:src:B.4.21} & \texttt{030Q} & \SeriesRef{ari:C.7.10} & \texttt{0018}\\
\SeriesRef{push:src:B.4.22} & \texttt{030R} & \SeriesRef{ari:C.7.11} & \texttt{0019}\\
\SeriesRef{push:src:B.4.23} & \texttt{030S} & \SeriesRef{ari:C.7.12} & \texttt{001A}\\
\SeriesRef{push:src:B.4.24} & \texttt{030T} &  & \\
\bottomrule
\end{longtable}\endgroup

\subsection{The proofs in outline}
\hypertarget{tag-03K0}{}\label{series:I:outlines}

\hypertarget{outline-I-A}{}
\paragraph{Theorem A.}
\begin{enumerate}[label=\arabic*.]
\item\hypertarget{step-I-A-1}{} The pointed maps
$(x,y)\mapsto(2x,-2\sqrt2y)$ and
$(x_0,y_0)\mapsto(-x_0,y_0/(2i))$ transport the chosen section,
invariant differential and period lattice.  The quotient
exponential identifies the period model with the elliptic-surface
line-bundle quotient
(\hyperref[I:weierstrass-comparison]{\ref*{I:weierstrass-comparison}},
\hyperref[fam:src:B.1.1]{\ref*{fam:src:B.1.1}}).
\item\hypertarget{step-I-A-2}{} The two theta transformation
factors give the prescribed affine period laws; the cusp product
fixes the additive constant and hence the displayed logarithm
(\hyperref[fam:src:C.1.1]{\ref*{fam:src:C.1.1}}).
\item\hypertarget{step-I-A-3}{} The distribution identity gives
the exact division-polynomial norm formula, and $mP=\pm P$
cancels the remaining norm.  The stated polynomial roots and
values then determine every order-four-through-eight
specialization
(\hyperref[ari:C.4.2]{\ref*{ari:C.4.2}},
\hyperref[ari:C.4.3]{\ref*{ari:C.4.3}},
\hyperref[ari:C.4.4]{\ref*{ari:C.4.4}}).
\item\hypertarget{step-I-A-4}{} Exact polynomial bounds away
from the square root point and the controlled local theta
expansion give $\max N=\sqrt2$.  The determinant identity and
its cusp limit then give exactly $\Im c<H_*$
(\hyperref[ari:C.8.1]{\ref*{ari:C.8.1}},
\hyperref[ari:C.8.2]{\ref*{ari:C.8.2}},
\hyperref[ari:C.8.3]{\ref*{ari:C.8.3}},
\hyperref[ari:C.8.4]{\ref*{ari:C.8.4}},
\hyperref[fam:src:C.2.3]{\ref*{fam:src:C.2.3}}).
\item\hypertarget{step-I-A-5}{} The normalized critical factor
satisfies $|1-X|<1$ and
$(t-1)L_t=X(X-2)/24$.  Its radial sign, the coarse derivative
at zero and the cusp limit give the complete critical and level
classification
(\hyperref[ari:C.8.15]{\ref*{ari:C.8.15}},
\hyperref[push:src:D.4.8]{\ref*{push:src:D.4.8}}).
\item\hypertarget{step-I-A-6}{} The local multiplication identity
evaluates every finite-place contribution on the stated minimal
models.  The product formula identifies the archimedean term and
gives the eleven exact CM factors
(\hyperref[ari:C.6.1]{\ref*{ari:C.6.1}},
\hyperref[ari:C.6.3]{\ref*{ari:C.6.3}}).
\end{enumerate}
\hyperlink{main-I-A}{Return to Theorem A}.

\hypertarget{outline-I-B}{}
\paragraph{Theorem B.}
\begin{enumerate}[label=\arabic*.]
\item\hypertarget{step-I-B-1}{} The intrinsic canonical pencil
fixes the base.  Its period centralizer forces $c'-c\in\Z$ and
the fibre-linear part $aI_2$
(\hyperref[fam:src:A.1.1]{\ref*{fam:src:A.1.1}},
\hyperref[fam:src:A.5.1]{\ref*{fam:src:A.5.1}}).
\item\hypertarget{step-I-B-2}{} Solving the two finite norm
equations gives the parity condition; comparing the actual cusp
collar gives $p'=ap$.  The holomorphic correction solves the
remaining gluing, so these necessary equations are sufficient
(\hyperref[fam:src:B.2.1]{\ref*{fam:src:B.2.1}},
\hyperref[fam:src:B.2.2]{\ref*{fam:src:B.2.2}},
\hyperref[fam:src:B.2.3]{\ref*{fam:src:B.2.3}}).
\item\hypertarget{step-I-B-3}{} The higher root pullbacks multiply
the norm images by $k,l$, yielding exactly the congruences modulo
$k,2l$.  Their intersection is $2kl\Z$, and a positive $u$
circuit is the cyclic parameter translation
(\hyperref[fam:src:B.3.1]{\ref*{fam:src:B.3.1}},
\hyperref[push:src:E.1.2]{\ref*{push:src:E.1.2}}).
\item\hypertarget{step-I-B-4}{} The invariant-coordinate lattice,
including its affine half-period, and the two exact finite theta
values give the Albanese nomes
(\hyperref[ari:C.2.1]{\ref*{ari:C.2.1}}).
\item\hypertarget{step-I-B-5}{} The explicit reflected period
and affine lattice maps give $c'\equiv1-\bar c\pmod{2\Z}$.
Their square and holomorphic conjugacy identify the two
involution classes
(\hyperref[fam:src:B.4.2]{\ref*{fam:src:B.4.2}},
\hyperref[mapsI:src:A.3.1]{\ref*{mapsI:src:A.3.1}}).
\item\hypertarget{step-I-B-6}{} The fixed-lift congruences and
cusp units determine the nonempty pieces.  The exact cap-meridian
matrices close them to $S^3$ and determine the semifree or Hopf
circle action
(\hyperref[mapsI:src:A.3.2]{\ref*{mapsI:src:A.3.2}},
\hyperref[mapsI:src:A.3.3]{\ref*{mapsI:src:A.3.3}}).
\end{enumerate}
\hyperlink{main-I-B}{Return to Theorem B}.

\hypertarget{outline-I-C}{}
\paragraph{Theorem C.}
\begin{enumerate}[label=\arabic*.]
\item\hypertarget{step-I-C-1}{} The parabolic multiplication law
and the polarization of $x_1+x_2+x_1x_2$ prove the character law.
The commutator and its conjugates generate its integral kernel,
giving the two exact indices
(\hyperref[fam:src:F.1.1]{\ref*{fam:src:F.1.1}}).
\item\hypertarget{step-I-C-2}{} Actual modular lifts first reduce
the Iwasawa factor, integral Heisenberg translations reduce $x$,
and the center reduces $t$.  Interior uniqueness and the displayed
face maps give the fundamental set
(\hyperref[I:monodromy-fundamental-set]{\ref*{I:monodromy-fundamental-set}}).
\item\hypertarget{step-I-C-3}{} Real parabolic stabilizers are
compact on $q_d\ne0$.  Rational parabolic conjugates of the
explicit null unipotent give a dense infinite-stabilizer locus,
which proves maximality
(\hyperref[fam:src:F.2.2]{\ref*{fam:src:F.2.2}},
\hyperref[fam:src:F.2.3]{\ref*{fam:src:F.2.3}}).
\item\hypertarget{step-I-C-4}{} The modular function $j/1728$
is the period map's left inverse, while the finite root charts
extend it at zero and one.  Its control on compact inverse
images proves properness on the finite $t$-line
(\hyperref[fam:src:F.3.1]{\ref*{fam:src:F.3.1}}).
\item\hypertarget{step-I-C-5}{} Wedge the period rows with an
integral alternating form to obtain the exact $(1,1)$ equation.
Its Hermitian determinant and the complex orientation give
$E^2=-2\operatorname{Pf}(E)$ and the quotient/projectivity tests
(\hyperref[fam:src:D.1.4]{\ref*{fam:src:D.1.4}},
\hyperref[fam:src:D.1.6]{\ref*{fam:src:D.1.6}}).
\end{enumerate}
\hyperlink{main-I-C}{Return to Theorem C}.

\hypertarget{outline-I-D}{}
\paragraph{Theorem D.}
\begin{enumerate}[label=\arabic*.]
\item\hypertarget{step-I-D-1}{} Root characters and the cusp
normalization compute the logarithmic direct images.  The residue
connecting map is evaluated as an actual nonzero scalar, fixing
the remaining Hodge integer at one
(\hyperref[inv:A2.2]{\ref*{inv:A2.2}},
\hyperref[inv:A2.4]{\ref*{inv:A2.4}}).
\item\hypertarget{step-I-D-2}{} Coherent Leray on the curve,
the structure-sheaf inputs and Serre duality give the complete
Dolbeault table, uniformly for every admissible parameter
(\hyperref[inv:A1.0]{\ref*{inv:A1.0}},
\hyperref[inv:A3.1]{\ref*{inv:A3.1}}).
\item\hypertarget{step-I-D-3}{} The middle rows and the
$S^6$ de Rham groups force the four rank-one differentials.
The double-complex decomposition then gives the Bott--Chern and
Aeppli tables and the explicitly checked square-zero products
(\hyperref[inv:C1.1]{\ref*{inv:C1.1}},
\hyperref[inv:C2.1]{\ref*{inv:C2.1}},
\hyperref[inv:C3.3]{\ref*{inv:C3.3}}).
\item\hypertarget{step-I-D-4}{} Quotienting by $dt$ gives the
degree-one modules.  The maps $(a,b,c)\mapsto(xa-zc,yb-zc)$
and $\Omega^2\xrightarrow{\wedge dt}I$ compute their kernels
and torsion exactly; finite etale descent gives the global
support assertion
(\hyperref[inv:A4.1]{\ref*{inv:A4.1}},
\hyperref[inv:A4.2]{\ref*{inv:A4.2}}).
\item\hypertarget{step-I-D-5}{} The theta factors with their original
finite affine phases and the extending cusp Cartier units define the
local twelfth-power lines on the four pieces.  The exact determinant divisor makes
the displayed holomorphic connection inverse valid on every overlap
annulus.  Its one-form cocycle and the specified coordinate weights
give $Q$; the scalar cocycle gives the primitive $G$ and its factor
twelve normalization
(\hyperref[inv:C3.8]{\ref*{inv:C3.8}}).
\end{enumerate}
\hyperlink{main-I-D}{Return to Theorem D}.

\hypertarget{outline-I-E}{}
\paragraph{Theorem E.}
\begin{enumerate}[label=\arabic*.]
\item\hypertarget{step-I-E-1}{} The primitive Picard
normalization identifies the integer twists, and weighted
coarse degrees give their section counts and four dimensions
(\hyperref[inv:C3.2]{\ref*{inv:C3.2}},
\hyperref[inv:D4.1]{\ref*{inv:D4.1}}).
\item\hypertarget{step-I-E-2}{} The actual integral cocycle
modulo the two period rows gives the regular resonance equation
and both finite root lattices.  At the cusp the normalization
complex gives exactly the reduced degree-two skyscraper
(\hyperref[inv:D4.2]{\ref*{inv:D4.2}},
\hyperref[inv:D4.3]{\ref*{inv:D4.3}}).
\item\hypertarget{step-I-E-3}{} The joint order-three
coefficient $123ik/(64\pi)$ fixes the local normal form;
Smith modules and every stabilizer character give the complete
finite stalk lengths.  Leray then gives the formula for $J_z$
(\hyperref[inv:D4.4]{\ref*{inv:D4.4}},
\hyperref[inv:D4.5]{\ref*{inv:D4.5}}).
\item\hypertarget{step-I-E-4}{} The established
$\bar\partial$-closed form $G_\gamma=G_0/12$ defines the
Picard family by $\bar\partial-2\pi izG_\gamma$.
Proper coherent base change and a real nonintegral nonresonant
point make every actual jump locus locally discrete
(\hyperref[inv:C3.1]{\ref*{inv:C3.1}},
\hyperref[inv:C3.2]{\ref*{inv:C3.2}},
\hyperref[inv:D4.5]{\ref*{inv:D4.5}}).
\item\hypertarget{step-I-E-5}{} The root-forgetting map
$(x,y)\mapsto(x^l,y^k)$ with character
$\lambda\mapsto\lambda^{kl}$ pulls back the standard Hodge
line.  Its weighted direct images and logarithmic tangent
sequence give $(1,1,1,0)$ for every stated higher sphere member
(\hyperref[push:src:E.1.2]{\ref*{push:src:E.1.2}},
\hyperref[push:src:E.2.1]{\ref*{push:src:E.2.1}}).
\item\hypertarget{step-I-E-6}{} Differentiating the actual
period groupoid gives the nonzero class $e\otimes[\gamma]$.
The induced parameter map forces the one-variable Kuranishi
ring to be reduced and smooth.  The integral translation sheaf
gives the standard automorphism group
(\hyperref[fam:st:A.5.2]{\ref*{fam:st:A.5.2}},
\hyperref[fam:src:B.4.1]{\ref*{fam:src:B.4.1}},
\hyperref[push:src:E.2.1]{\ref*{push:src:E.2.1}}).
\end{enumerate}
\hyperlink{main-I-E}{Return to Theorem E}.

\hypertarget{outline-I-F}{}
\paragraph{Theorem F.}
\begin{enumerate}[label=\arabic*.]
\item\hypertarget{step-I-F-1}{} The node charts
$X_i=W/r^i$, $Y_i=r^{i+1}/W$, $X_iY_i=r$ and their inverse
overlaps give the proper rank-one cycle.  The finite affine
actions rotate its components freely and produce the stated
cover.  Subtraction of the actual two origins fixes both deck
differences
(\hyperref[night:P2.2]{\ref*{night:P2.2}},
\hyperref[night:P2.3]{\ref*{night:P2.3}}).
\item\hypertarget{step-I-F-2}{} Resolve a proposed filling's
dominating component.  Its relative anticanonical direct image
recovers the pencil, and the three labelled critical-value
sections recover the marked base by their meromorphic cross-ratio
(\hyperref[night:P2.4]{\ref*{night:P2.4}}).
\item\hypertarget{step-I-F-3}{} Modulo $\widehat\delta$ the
actual cross-cocycles vanish.  The three unipotent annihilator
equations have determinant 6; Stein factorization on a horizontal
central component then kills its normal-meridian class
(\hyperref[night:P2.5]{\ref*{night:P2.5}},
\hyperref[night:P2.6]{\ref*{night:P2.6}}).
\item\hypertarget{step-I-F-4}{} Principalization makes both
axes normal crossing.  The Stein finite factor is birational
over the normal bidisc, so the compact corner fibre is connected.
The lifts $Q_o=T_{-mb\widehat\gamma}Q$, $R_o=R$ commute;
local meridians therefore satisfy the constant-ratio equation
(\hyperref[night:P2.7]{\ref*{night:P2.7}}).
\item\hypertarget{step-I-F-5}{} The connected incidence graph
joins a horizontal ratio zero to the strict cusp ratio $a$.
Admissibility gives $\gcd(a,12)=1$, producing the contradiction
for every proper analytic filling and every finite cover
(\hyperref[night:P2.7]{\ref*{night:P2.7}}).
\item\hypertarget{step-I-F-6}{} For higher multiplicities,
$K_X^{-N}=f^*\mathcal O(4l+3k-kl)$ with nonzero degree.
The positive canonical power and its three labelled critical
points recover the degree-one base parameter.  The same ratio
argument and $\gcd(a,N)=1$ give the full higher conclusion
(\hyperref[night:P2.9]{\ref*{night:P2.9}}).
\item\hypertarget{step-I-F-7}{} With zero logarithmic
translations the two line-bundle neck equations have zero
corner winding and give the proper flat plumbing after its
specified modification and finite quotient.  Normalization and
the double-locus smoothing formula give virtual Euler zero for
Engel's particular one-component proposal, excluding its
smoothing to $e=2$
(\hyperref[night:P2.10]{\ref*{night:P2.10}},
\hyperref[night:P2.11]{\ref*{night:P2.11}}).
\end{enumerate}
\hyperlink{main-I-F}{Return to Theorem F}.

\hypertarget{outline-I-G}{}
\paragraph{Theorem G.}
\begin{enumerate}[label=\arabic*.]
\item\hypertarget{step-I-G-1}{} The canonical divisor recovers
$S_2$ intrinsically in a proper extension.  Its normalization
and Albanese invariant retain $J(iw/216)$, whose natural boundary
excludes crossing the outer circle
(\hyperref[push:src:D.1.3]{\ref*{push:src:D.1.3}}).
\item\hypertarget{step-I-G-2}{} The fixed complex-linear and
integral basis changes give the displayed four periods.
The local theta expansion gives $B(h)=-3ih^2+O(h^4)$ and
the exact inverse transpose gives every integral dual covector
(\hyperref[push:src:D.2.1]{\ref*{push:src:D.2.1}},
\hyperref[eq:d-dual-exact]{\ref*{eq:d-dual-exact}}).
\item\hypertarget{step-I-G-3}{} Bounded dual covectors have
constant integral labels after subselection; their labels form
$L_a$.  The quadratic real term cancels, leaving
$\lambda\cdot V_n+o(1)$.  The primitive surviving subgroup and
its limiting homomorphism give exactly $\Gamma_\infty$
(\hyperref[push:src:D.3.1]{\ref*{push:src:D.3.1}}).
\item\hypertarget{step-I-G-4}{} Choosing successively separated
scales realizes every primitive subgroup and limiting
homomorphism.  Pontryagin--Chabauty duality converts their dual
graphs to the stated real-torus additive quotients
(\hyperref[push:src:D.3.1]{\ref*{push:src:D.3.1}}).
\item\hypertarget{step-I-G-5}{} At fixed boundary parameter,
irrational root rays make all noncentral dual labels escape.
The surviving rational label is
$\operatorname{lcm}(2,n)\Re\xi$; its circle-valued function
separates precisely the asserted classes.  Collapsing those
classes proves the universal Hausdorff-reflection property
(\hyperref[push:src:D.2.3]{\ref*{push:src:D.2.3}}).
\end{enumerate}
\hyperlink{main-I-G}{Return to Theorem G}.

% RungZero block 1 belongs to the shared series overview.
% The coefficient-atlas portion of block 25 is included in Theorem D;
% its arithmetic portion remains in the final open-problems part,
% with the explicit unresolved relative-period gate at ari:C.7.4.
\subsection{Known and new}
\begingroup\small
\begin{longtable}{@{}p{.10\textwidth}p{.85\textwidth}@{}}\toprule Theorem & Exact source relationship\\\midrule\endhead
A & The elliptic surface and section
come from \cite[Remark~3.12, p.~19]{AC2026} and
\cite[Definition~2.3]{Engel2026}; the exact coordinate transports
are retained in \ref{I:weierstrass-comparison} and
\ref{fam:src:B.1.1}.  The norm and height conclusions have the
specific scopes of the body results listed here.\\
B & The construction is that of
\cite{AC2026,Engel2026}; the higher input is Viaclovsky's specified
root construction.  The standard exact $2\Z$ classification is
also Theorem~1.4 of \cite{DuEtAl2026v1}; the version and parameter
comparison is retained in \ref{fam:src:A.5.1} and the credited
appendix, including the constant translation between particular
period solutions.\\
C & The proper-action question is
posed in \cite[Remark~8.4, pp.~23--24]{Engel2026}.  The stated
answer retains the precise monodromy subgroup and full invariant
open set; it does not rely on replacing that subgroup by its
rational or real closure.\\
D & The holomorphic-form vanishing
and Hodge pattern with one undetermined integer are credited to
\cite[Theorem~9.1(vi), p.~65; Proposition~9.20, p.~80]{AC2026}.
The local presentations correct the degree-one torsion assertion
in its Remark~9.21, p.~81, and qualify Remark~1.3, p.~5, by
degree.  The coefficient construction in \ref{inv:C3.8} uses the
original period, finite affine and toric charts; its source comparison
and exact cocycle maps are given there.\\
E & The standard tangent, automorphism
and Kuranishi conclusions are credited to
\cite[Theorems~1.3--1.4, p.~3; Theorem~6.12, pp.~81--82;
Theorem~6.16, p.~85]{DuEtAl2026v1}.  The version-specific
re-derivations and parameter transport are preserved; version~2's
omission of the deformation section changes none of that credit.\\
F & The comparison retains Engel's
conjectural status, its specified central model, and the actual
maps between the elliptic-surface and period descriptions.  The
higher family is Viaclovsky's construction; the obstruction is
applied to its precise primitive logarithmic data.\\
G & The elliptic-surface description
is transported through the marked quotient dictionary; the body
states the precise Pontryagin--Chabauty duality input used for
closed-group convergence.  The Hausdorff reflection and additive
limits retain their stated topological categories and maps.\\
\bottomrule\end{longtable}\endgroup
\subsection{Methodology}
This mathematical text is AI-produced. GPT-5.6 Sol (OpenAI), in
Codex at reasoning effort ultra, produced the working edition during
24 August--9 September 2026. GPT-6 Astra (OpenAI), in ChatGPT Pro
mode at effort Pro, carried out the research runs of 12--16 November,
the S6\_06\_BLIND run of 16--18 November, the S6\_09\_PAPERS and
S6\_10\_PAPERS writing runs of 16--19 November, and this series run
on 19 November 2026. Claude Fable 5.2 (Anthropic), at effort high and
at effort max, and Claude Opus 5.5 (Anthropic), at effort max,
performed readings, intakes, sweeps and checks during 12--19 November
2026. Prompts specified the statements, marked inputs and requested
checks. Returned proofs and computations were read against those
inputs and the opened original sources. Corrections entered the next
prompt and the receiving statements with their explicit transports.

The recorded methods include derivation before reading the supplied
proof, independent re-proofs, exact and interval computations,
checks of special cases, and comparison of scope and conventions
with original sources. For this paper the received and current
record contains 853 targeted proof checks, 496 use or reconciliation events, 104 computational or formal certificates, 205 background-carry events, 9 independent re-proofs.
Each event keeps its actual scope; an editorial reconciliation
contributes no proof verification. The linked entry for every
statement records the agreeing routes and remaining risk.
No human has checked the mathematics.

We dedicate our text to the public domain under
\href{https://creativecommons.org/publicdomain/zero/1.0/}{CC0 1.0 Universal}.
The earlier runs and check records are preserved in the
\href{https://github.com/KokunoYumeto/s6-complex-structure-workbench/tree/4e8b547874c2acff6b49713ca96f84b281ab5793}{pinned workbench archive};
the register sweep is in the
\href{https://github.com/KokunoYumeto/grothendieck-montpellier-archive/tree/96b2bcaec355870a8faa9457406673b7cd897654/raw/results_register/sweep1_s6}{pinned sweep archive}.
The current statement records are
\href{S6_SERIES_GRIND_checks_I.md}{the accompanying check file},
with the complete typed ledger in \href{CHECK_EVENTS.tsv}{CHECK\_EVENTS.tsv}.
\subsection{How the paper is organized}
The body gives the marked constructions, their integral and holomorphic
comparison maps, and the consequences grouped in Theorems A, B, C, D, E, F and G.
The final part states each remaining question with its established
reduction. In the diagram an arrow records reuse of the particular
body argument identified below; a theorem group can contain several
arguments with different inputs. The statement headers give the full
local uses and the links to the corresponding outline steps.
\begin{figure}[htbp]\centering
\begin{tikzpicture}[>=stealth,every node/.style={font=\small}]
\node[draw,rounded corners,minimum width=2.2cm,minimum height=.65cm] (TA) at (0.0,0.0) {\hyperlink{main-I-A}{Theorem A}};
\node[draw,rounded corners,minimum width=2.2cm,minimum height=.65cm] (TB) at (3.5,0.0) {\hyperlink{main-I-B}{Theorem B}};
\node[draw,rounded corners,minimum width=2.2cm,minimum height=.65cm] (TC) at (7.0,0.0) {\hyperlink{main-I-C}{Theorem C}};
\node[draw,rounded corners,minimum width=2.2cm,minimum height=.65cm] (TD) at (0.0,-1.25) {\hyperlink{main-I-D}{Theorem D}};
\node[draw,rounded corners,minimum width=2.2cm,minimum height=.65cm] (TE) at (3.5,-1.25) {\hyperlink{main-I-E}{Theorem E}};
\node[draw,rounded corners,minimum width=2.2cm,minimum height=.65cm] (TF) at (7.0,-1.25) {\hyperlink{main-I-F}{Theorem F}};
\node[draw,rounded corners,minimum width=2.2cm,minimum height=.65cm] (TG) at (0.0,-2.5) {\hyperlink{main-I-G}{Theorem G}};
\draw[->,semithick] (TA) to (TB);
\draw[->,semithick] (TE) to (TB);
\draw[->,semithick] (TD) to (TE);
\draw[->,bend left=18,semithick] (TB) to (TG);
\draw[->,bend left=18,semithick] (TG) to (TA);
\end{tikzpicture}
\caption{Reuse of the indicated body arguments among the main theorem groups.}
\par\medskip
\begingroup\small
\begin{tabularx}{\textwidth}{@{}p{.10\textwidth}X@{}}
\toprule Arrow & Supplying and receiving arguments\\\midrule
$A\to B$ & \SeriesRef{fam:src:C.2.3} $\longrightarrow$ \SeriesRef{ari:C.2.1}. The exact parameter domain is used for the intrinsic Albanese nome.\\
$E\to B$ & \SeriesRef{fam:src:B.4.1} $\longrightarrow$ \SeriesRef{mapsI:src:A.3.1}. The automorphism group acts in the conjugacy classification of real involutions.\\
$D\to E$ & \SeriesRef{inv:C1.1} $\longrightarrow$ \SeriesRef{inv:C3.1}. The differential and Bott--Chern dimensions determine the Picard curvature statement.\\
$B\to G$ & \SeriesRef{fam:src:A.1.1} $\longrightarrow$ \SeriesRef{push:src:D.1.3}. The intrinsic multiplicity-four divisor supplies the Albanese continuation obstruction.\\
$G\to A$ & \SeriesRef{push:src:D.2.1} $\longrightarrow$ \SeriesRef{ari:C.8.2}. The half-period derivative identity supplies the theta-heat calculation.\\
\bottomrule\end{tabularx}\endgroup\end{figure}

\clearpage
\section{Common conventions and the marked construction}
\hypertarget{tag-01O1}{}\label{def:common-conventions}

\section*{Conventions and the notation dictionary}
We use $B=\Pp^1_t$, $p_1=0$, $p_2=1$, $p_0=\infty$, and
$B^\circ=B\setminus\{0,1,\infty\}$. Base meridians are counterclockwise
in a local holomorphic coordinate. The analytic generators are clockwise
root deck transformations, $s_j(g_jz)=e^{-2\pi i/m_j}s_j(z)$, with
$g_1g_2g_0=1$. A positive base loop ends at the inverse deck translate;
identifying its endpoint with the initial fibre applies $g_j$, giving
homology monodromy $A_j$. The retained multiplicities are $m_1=3,m_2=4$;
$f^*p_j=m_jS_j$ includes the nonreduced structure, whereas $S_j$ is reduced.
The cusp fibre $W=f^{-1}(\infty)$ is reduced.
\[
 \Lambda=\Z\langle\gh,\uh,\wh,\dh\rangle,\qquad
 \Pi_c=\begin{pmatrix}6\mu&\tau&1&0\\\beta^{\partt}+c&\mu&0&1\end{pmatrix},
 \qquad w=e^{-\pi ic}.
\]
We fix $(\beta^{\partt}+\tau)(p_0)=0$. Write
$D_{\partt}=\Im\beta^{\partt}-6(\Im\mu)^2/\Im\tau$ and
$M=\sup D_{\partt}$. Engel's line bundle is
$M_S=\OO_S(P-O)$; his indices $0,1,\infty$ correspond to our $1,2,0$.
In his common integral basis, the columns of the following matrix are
our ordered basis vectors:
\[
 P_{\rm mark}=\begin{pmatrix}
 0&-2&-1&-2\\2&-1&0&-1\\0&2&1&3\\-1&0&0&0
 \end{pmatrix},\quad
 T_0P_{\rm mark}=P_{\rm mark}A_1,\quad
 T_1P_{\rm mark}=P_{\rm mark}A_2,\quad
 \delta_{\rm E}=P_{\rm mark}\dh.
\]
For our finite twist vectors $v_1,v_2$, his positive-root affine vectors
are $a_0=-P_{\rm mark}v_1/3$, $a_1=-P_{\rm mark}v_2/4$.
In additive fibre coordinates $(\xi,\eta)$, put $v=e^{2\pi i\eta}$.
The first period acts by
$(\xi,v)\mapsto(\xi+6\mu,e^{2\pi i\beta_c}v)$; Engel's contracting
generator is its inverse, over $-6P$ with multiplier $e^{-2\pi i\beta_c}$.
We choose his reference linearization to be this transported inverse
period lift at $c=0$. Then $\lambda_0=1$ and $u=w^2$.
The comparison maps and the effect of changing this reference are
proved in \ref{fam:section:B.1} and B.2. Homology has integral coefficients unless stated
otherwise; tori and the threefold have their complex orientations.

Each statement below has its own proof or an explicit unresolved target
with its complete partial proof and the precise remaining step.


\section{Invariant and arithmetic notation in the family marking}
\hypertarget{tag-01TG}{}\label{inv:notation-dictionary}

\paragraph{Base, fibres, and the two peripheral orientations.}
We use $B=\mathbb P^1_t$, $p_1=0$, $p_2=1$, $p_0=\infty$,
$\Sigma=\{p_0,p_1,p_2\}$, and $D=W+S_1+S_2$ with reduced
components. The full fibres are $f^*p_1=3S_1$, $f^*p_2=4S_2$,
and $f^*p_0=W$. A positive base loop is counterclockwise in a
holomorphic local coordinate. The analytic root deck generator is
clockwise, $s_j(g_jz)=e^{-2\pi i/m_j}s_j(z)$, with
$g_1g_2g_0=1$. The lifted positive base loop ends at the inverse deck
translate; its endpoint identification applies $g_j$. Thus its
homology action is $A_j$, not $A_j^{-1}$.
The clockwise boundary orientation used in the collar integration
is stated separately in \cref{inv:C3.2}. References there to a
peripheral generator refer to this fixed analytic generator.

\paragraph{Integral marking and coordinate order.}
The homology basis is
$(\hat\gamma,\hat u,\hat w,\hat\delta)$, and
$(\gamma,u,w,\delta)$ denotes its dual cohomology basis.
The letter $w$ in that cohomology basis is a covector, whereas the
scalar family parameter below is $e^{-\pi ic}$.
In the ordered additive coordinates $(\zeta_1,\zeta_2)$,
\[
 \Pi_c=\begin{pmatrix}6\mu&\tau&1&0\\
             \beta_c&\mu&0&1\end{pmatrix},
 \qquad \beta_c=\beta^{\rm part}+c,
 \qquad (\beta^{\rm part}+\tau)(p_0)=0.
\]
Its rows are $\sigma_1=6\mu\gamma+\tau u+w$ and
$\sigma_2=\beta_c\gamma+\mu u+\delta$.
The actual monodromy matrices in this column basis are
\[
 A_1=\begin{pmatrix}1&0&0&0\\6&0&1&0\\-6&-1&-1&0\\-2&1&0&1\end{pmatrix},
 \quad
 A_2=\begin{pmatrix}1&0&0&0\\0&0&-1&0\\-6&1&0&0\\3&0&1&1\end{pmatrix},
 \quad A_0=(A_1A_2)^{-1}.
\]
No row interchange is implicit: if a local calculation orders the
coordinates as $(\zeta_2,\zeta_1)$, its period matrix is
$J\Pi_c$, $J=\left(\begin{smallmatrix}0&1\\1&0\end{smallmatrix}\right)$,
and its Hermitian coefficient matrix is $JH J^{\mathsf t}$.
Both the main invariant proofs and the retained working-edition
fibre formulas use the displayed $(\zeta_1,\zeta_2)$ order.

\paragraph{The exact Engel and reversed-meridian transitions.}
Engel's base indices $(0,1,\infty)$ correspond to $(1,2,0)$ here.
Writing his common integral basis as column coordinates, the family
basis is given by the columns of
\[
 P_{\rm mark}=\begin{pmatrix}
 0&-2&-1&-2\\2&-1&0&-1\\0&2&1&3\\-1&0&0&0
 \end{pmatrix},\qquad
 T_0P_{\rm mark}=P_{\rm mark}A_1,\quad
 T_1P_{\rm mark}=P_{\rm mark}A_2.
\]
The invariant source's reversed-meridian matrix is
\[
 Q=P_{\rm mark}S,\qquad
 S=\begin{pmatrix}-1&0&0&0\\-6&0&-1&0\\6&-1&0&0\\2&1&1&1\end{pmatrix},
 \qquad
 Q=\begin{pmatrix}2&-1&0&-2\\2&-1&0&-1\\0&2&1&3\\1&0&0&0\end{pmatrix}.
\]
It satisfies $T_0Q=QA_1^{-1}$ and $T_1Q=QA_2^{-1}$.
This is the explicit convention change of \cref{inv:B1.0}.
For invariant finite twist vectors $v_1,v_2$, Engel's positive-root
affine vectors are $a_0=-P_{\rm mark}v_1/3$ and
$a_1=-P_{\rm mark}v_2/4$; his distinguished cycle is
$P_{\rm mark}\hat\delta$.

\paragraph{Multiplicative parameters and special-fibre coordinates.}
Set $v=e^{2\pi i\zeta_2}$. The first period acts by
$(\zeta_1,v)\mapsto(\zeta_1+6\mu,e^{2\pi i\beta_c}v)$.
Engel's contracting generator is its inverse over $-6P$, on the
complement of the zero section of $\mathcal O(P-O)$.
With the reference inverse lift transported from $c=0$,
\[
 w=e^{-\pi ic},\qquad u_{\rm Engel}=w^2,
 \qquad 0<|w|<216,\quad 0<|u_{\rm Engel}|<46656.
\]
A different reference lift multiplies $u_{\rm Engel}$ by a fixed
$\lambda_0\in\mathbb C^*$. Thus $u_{\rm Engel}$ is distinct from
the integral covector $u$ and from an elliptic logarithm denoted $u$
in the period-transcendence section.
The special values use $\rho=e^{\pi i/3}$ rather than
$e^{2\pi i/3}=\rho-1$; the two conventions differ by the modular
translation $\tau\mapsto\tau+1$. The Albanese parameters and their
upper-half-plane representatives are
\[
 \begin{array}{c|c|c|c}
 &\kappa&\text{upper-half-plane period}&\text{nome}\\\hline
 S_1&c+\frac12-\frac{3i}{2\pi}\log48&-\kappa_1&
             q_1=e^{-2\pi i\kappa_1}=-w^2/48^3\\
 S_2&c+\frac32-\frac{3i}{\pi}\log6&-\kappa_2/2&
             q_2=e^{-\pi i\kappa_2}=iw/216.
 \end{array}
\]
The factor $1/2$ is the full order-four Albanese lattice projection.
In particular $q_1=(27/64)q_2^2$.

\paragraph{Determinants, orientations, and cohomology generators.}
Write $y=\Im\tau>0$, $m=\Im\mu$, $r=\Im\beta_c$, and
$D_c=r-6m^2/y=D_{\rm part}+\Im c$. The symbol $\Delta$ used
for this real quantity in the invariant fibre formulas means $D_c$,
not a Weierstrass discriminant. In admissible fibres $D_c<0$.
The real period determinant in the row order
$(\Re\zeta_1,\Im\zeta_1,\Re\zeta_2,\Im\zeta_2)$ is $yD_c$.
Consequently the marked real orientation is opposite to the complex
orientation, and for
$\eta_F=u\wedge w+6\gamma\wedge\delta$,
\[
 \int_{F,\operatorname{or}_{\mathbb C}}\eta_F^2=-12,
 \qquad
 \Gamma=\frac{i}{2D_c}
       \left(d\bar\zeta_2-\frac myd\bar\zeta_1\right)
       \quad\text{represents }\gamma^{0,1}.
\]
The subscript on $\eta_F$ separates this integral fibre class from
Dedekind's eta function and from Weierstrass quasi-periods.
The primitive global class is $\mathfrak g=g_\gamma$.
For standard data $(l_0,l_1,l_2)=(0,1,-1)$ and
$p=12l_0-4l_1-3l_2=-1$,
\[
 \exp(12\mathfrak g)=f^*\mathcal O_B(1),\quad
 \mathcal O_X(S_1-S_2)=\exp(\mathfrak g),\quad
 [G_0]=12\mathfrak g.
\]
In general the first equality is
$\exp(12g_\gamma)=f^*\mathcal O_B(-p)$.
Here $\exp$ means the sheaf map $z\mapsto e^{2\pi iz}$,
$\int_B\omega_{FS}=1$, $c_{1,BC}(L)=[iF_L/(2\pi)]$, and the
unnormalized Chern--Ricci form represents $2\pi c_{1,BC}(T_X)$.
Thus the prescribed Chern--Ricci form is $\pi f^*\omega_{FS}$.

\paragraph{Root and coarse line bundles.}
For the root curve $\mathcal B$ of orders $(3,4)$ put
$R=R_1R_2^{-1}$, so
$R_1=R^4$, $R_2=R^3$, $\rho^*\mathcal O_B(1)=R^{12}$,
and $K_{\mathcal B}=R^{-7}$.
The notation $\mathcal O_{\mathcal B}(n)$ means $R^n$;
$\mathcal O_B(n)$ is the ordinary coarse degree-$n$ line bundle.
They have different degrees and stabilizer characters.
The logarithmic sheaves are
$\mathcal L=\Omega_B^1(\log\Sigma)$,
$\mathcal A=\Omega_X^1(\log D)$, and
$\mathcal E=\mathcal A/f^*\mathcal L$.
Every coherent base-change assertion uses the full scheme fibre.
The finite cyclic covers and subsequent invariant functor retain its
nilpotent normal directions.

\paragraph{Theta norm, modular coordinate, and heights.}
The theta variable has periods $(1,\tau)$; software using a first
argument of period $\pi$ receives $\pi z$. With $q=e^{2\pi i\tau}$,
\[
 N(z,\tau)=\left|\frac{\vartheta_1(z,\tau)}{\eta(\tau)}\right|
       e^{-\pi(\Im z)^2/\Im\tau},\qquad
 D_{\rm part}=\frac6\pi\log\frac{N(\mu,\tau)}{2\sqrt3},
 \qquad t=\frac{j(\tau)}{1728}.
\]
The modular discriminant is
$\Delta_0=q\prod_{n\ge1}(1-q^n)^{24}$, so
$t=E_4^3/(1728\Delta_0)$ and $t-1=E_6^2/(1728\Delta_0)$.
For the rational family
\[
 E_t:Y^2=X^3-\tfrac34t^3(t-1)X+\tfrac14t^4(t-1)^2,
 \qquad P_t=(0,t^2(t-1)/2),
\]
the invariant differential is $dX/(2Y)$ and
$\Delta_E=27t^8(t-1)^3$.
In the lattice model $X=\wp(u)$ and $Y=\wp'(u)/2$;
the lattice discriminant is
$g_2^3-27g_3^2=(2\pi)^{12}\eta(\tau)^{24}$ when the
period lattice is $\mathbb Z+\tau\mathbb Z$.
We use the height of $(O)$,
$\widehat h_O(P)=\frac12\lim n^{-2}h(x(nP))$.
PARI's \texttt{ellheight}, the surface self-pairing, and the
rank-one BSD regulator use $2\widehat h_O$.
The full real period is $\Omega=\int_{E(\mathbb R)}|\omega|$,
including both components if disconnected.
The archimedean local height is exactly $-\log N$.
At finite places $|p|_v=p^{-1}$ and $n_v=[K_v:\mathbb Q_p]$;
complex embeddings are counted separately in the global sum.

For period arguments, $\eta_j$ is defined by
$\zeta(z+\omega_j)=\zeta(z)+\eta_j$, with
$\omega_1\eta_2-\omega_2\eta_1=-2\pi i$.
The real-linear extension of these quasi-periods is denoted $\eta(u)$.
In the critical-point proof, the Cayley coordinate is
$\chi=(\tau-i)/(\tau+i)$, $z=k\chi^2$, and
$\beta_{\rm cusp}=\log(5+2\sqrt6)/(2\pi)$;
neither symbol is the family parameter $w$ or the period $\beta_c$.

\paragraph{Shared statements.}
The maximum has the single label \cref{fam:src:C.2.3}, with aliases
\texttt{ari:C.8.5} and \texttt{push:src:E.4.3}; the independent
interval, duplication, and radial arguments are all retained.
Tangent cohomology, the parameter Kodaira--Spencer map, and the
one-dimensional Kuranishi germ are collected in
\cref{fam:src:A.5.2}, with aliases
\texttt{inv:D1.3}, \texttt{inv:D2.1}, and \texttt{inv:D3.1}.
The exact parameter classification is \cref{fam:src:A.5.1};
the canonical-pencil statement is \cref{fam:src:A.1.1}.
The working Hodge-line proof is \CompendiumRef{I}{mapsI:hodge:line},
and the finite normal-order calculation is
\CompendiumRef{I}{wb:ffcert:thm:normal-trivialization}.
Working structure-sheaf and cusp statements share \cref{inv:A1.0};
working Picard restriction shares \cref{inv:C3.1};
the global fibre-class extension shares \cref{inv:A3.2};
the Bott--Chern multiplication statement shares \cref{inv:C3.3}.



\begin{lemma}[The exact comparison of the two Weierstrass equations]
\hypertarget{tag-0002}{}\label{I:weierstrass-comparison}
\StatementRecord{I:weierstrass-comparison}{proved}{0002}{\SeriesRef{ari:C.1.1}}
Put $C=t^3(t-1)$ and $D=t^4(t-1)^2$, with $t\ne0,1$.
The pointed map
\[
 \Phi:E_t^{\rm ar}\longrightarrow E_t^{\rm geom},\qquad
 (x,y)\longmapsto(2x,-2\sqrt2y)
\]
is an isomorphism from
$y^2=x^3-3Cx/4+D/4$ to $Y^2=X^3-3CX+2D$.
It sends $(0,t^2(t-1)/2)$ to
$(0,-\sqrt2t^2(t-1))$ and preserves the normalized theta norm
under the induced pointed elliptic-curve identification.
The discriminants satisfy $\Delta_{\rm geom}=64\Delta_{\rm ar}$.
\end{lemma}
% ALL PRINTED PROOF MATERIAL MASKED





































\subsection{A typed fibre dictionary}\hypertarget{tag-01PG}{}\label{fam:section:B.1}

\begin{proposition}\hypertarget{tag-01Q7}{}\label{fam:src:B.1.1}\label{fam:st:B.1.1}
\StatementRecord{fam:src:B.1.1}{proved}{01Q7}{\SeriesRef{fam:src:A.3.1}}

Let $E_\tau=\mathbb C/(\mathbb Z+\tau\mathbb Z)$, $P=[\mu]$, and $L_\mu=\mathcal O_{E_\tau}(P-O)$. The holomorphic map
\[
 [\xi,\eta]\longmapsto[\xi,e^{2\pi i\eta}]
\]
identifies $\mathbb C^2/\Pi_c\Lambda$ with the quotient of $L_\mu^\times$ by the lift
$A(\xi,W)=(\xi+6\mu,e^{2\pi i\beta_c}W)$. Engel's generator, in Definition 2.3 of the attached paper, is its inverse, with translation $-6P$ and multiplier $e^{-2\pi i\beta_c}$. The section $P=[\mu]$ is a generator of the Mordell--Weil group of the stated pointed Weierstrass surface.
\end{proposition}


% ALL PRINTED PROOF MATERIAL MASKED



































\begin{proposition}\hypertarget{tag-01Q8}{}\label{fam:src:B.1.2}\label{fam:st:B.1.2}
\StatementRecord{fam:src:B.1.2}{proved}{01Q8}{the definitions and constructions stated in the proof}

The exact set of nonzero parameters for which every ordinary fibre and both normalized finite root fibres of Engel's period-reference quotient are compact complex two-tori is
\[
 0<|u|<e^{-2\pi M}.
\]
This is the parameter domain for the compact family with both smooth finite root models. The condition on ordinary fibres alone at equality is determined by whether a maximizing point of $D_{\rm part}$ is ordinary; this boundary issue is treated in Part C. For an independent reference, the radius is $|\lambda_0|e^{-2\pi M}$.
\end{proposition}


% ALL PRINTED PROOF MATERIAL MASKED






\begin{proposition}\hypertarget{tag-01Q9}{}\label{fam:src:B.1.3}\label{fam:st:B.1.3}
\StatementRecord{fam:src:B.1.3}{proved}{01Q9}{\SeriesRef{fam:src:B.1.1}}

For the original affine logarithm $\mu$, we have $\wp_\tau(\mu)=0$ identically. Equivalently the generator $P=[\mu]$ of \ref{fam:src:B.1.1} is one of the two sections $P_*=(0,\sqrt2t^2(t-1))$ and $-P_*$. This assertion does not use a numerical evaluation of $\mu$ or of $\beta$.
\end{proposition}


% ALL PRINTED PROOF MATERIAL MASKED














































\begin{proposition}\hypertarget{tag-01QA}{}\label{fam:src:B.1.4}\label{fam:st:B.1.4}
\StatementRecord{fam:src:B.1.4}{proved}{01QA}{\SeriesRef{fam:src:B.1.3}, \SeriesRef{fam:src:B.1.1}; \cite[Lemma 2.1]{Engel2026}; \cite[Corollary 2.2]{Engel2026}}

The line-bundle comparison of the period family with Engel's family over $B^\circ$ transports its inverse period lift to a global section of $\operatorname{Hom}(t_{6P}^*M_S,M_S)$ with divisor $F_\infty$. Take this transported section for $c=0$ as the reference linearization. Then $\lambda_0=1$ exactly and $u=e^{-2\pi ic}=w^2$. For another independently fixed reference the precise constant is its scalar ratio with this section. With the cusp branch fixed in Part C and $f_{\rm sc}/t\to i\sqrt3/\pi$, the section is $P=(0,-\sqrt2t^2(t-1))$.
\end{proposition}


% ALL PRINTED PROOF MATERIAL MASKED

















\section{The normalized periods, quantitative bounds, and the boundary}\hypertarget{tag-01PK}{}\label{fam:section:C}

\subsection{A theta expression for the missing period}\hypertarget{tag-01PL}{}\label{fam:section:C.1}
We use the odd Jacobi function
\[
 \vartheta_1(v,\tau)=2\sum_{n\geq0}(-1)^n
 e^{\pi i\tau(n+1/2)^2}\sin((2n+1)\pi v)
\]
and the Dedekind function
\(\eta_D(\tau)=e^{\pi i\tau/12}\prod_{n\geq1}(1-e^{2\pi in\tau})\).
The subscript distinguishes this function from the invariant integral class.
We choose the cusp lift reached from \(\tau=iT\), \(T>1\), on which
\(\mu=1/2-iy(T)\), \(0<y(T)<T/2\). Put
\(a=5+2\sqrt6\) and \(\rho=e^{\pi i/3}\).

\begin{proposition}\hypertarget{tag-01QH}{}\label{fam:src:C.1.1}\label{fam:st:C.1.1}
\StatementRecord{fam:src:C.1.1}{proved}{01QH}{\SeriesRef{fam:src:B.1.3}}

For the period section identified in Statement \ref{fam:src:B.1.3}, the particular solution
with \((\beta^{\rm part}+\tau)(p_0)=0\) is
\begin{equation}\hypertarget{tag-01P5}{}\label{fam:eq:cd-beta}
 \beta^{\rm part}=\frac{6i}{\pi}\operatorname{Log}
 \frac{\vartheta_1(\mu,\tau)}{2\sqrt3\,\eta_D(\tau)}.
\end{equation}
Here the logarithm is continued on the simply connected uniformizing
half-plane from its cusp expansion. Its special values are
\[
 \mu_0=\frac12-\frac{i\log a}{2\pi},\qquad
 \mu(z_1)=\frac{2-\rho}{3},\qquad \mu(z_2)=\frac{1-i}{2},
\]
\[
 \beta^{\rm part}(z_2)=i\left(\frac32-\frac3\pi\log6\right),
 \qquad \operatorname{Re}\beta^{\rm part}(z_1)=-\frac12.
\]
In particular \(e^{2\pi i\mu_0}=-a\).
\end{proposition}


% ALL PRINTED PROOF MATERIAL MASKED












































































\begin{proposition}\hypertarget{tag-01QI}{}\label{fam:src:C.1.2}\label{fam:st:C.1.2}
\StatementRecord{fam:src:C.1.2}{proved}{01QI}{\SeriesRef{fam:src:C.1.1}, \SeriesRef{fam:eq:cd-beta}}

The following decimal evaluations of the exact expressions in \ref{fam:src:C.1.1} agree
in two independent evaluations through the displayed twenty decimal places:
\[
\begin{array}{c|l}
\mu_0&0.5-0.36485183191106784792\,i\\
\mu(z_1)&0.5-0.28867513459481288225\,i\\
\mu(z_2)&0.5-0.5\,i\\
\beta^{\rm part}(z_1)&-0.5-1.27101206057537346287\,i\\
\beta^{\rm part}(z_2)&-0.21100425815613411418\,i
\end{array}
\]
These are evaluations of explicit convergent expressions, with working
precision ninety decimal digits. The formula \eqref{fam:eq:cd-beta}, together
with \(j(\tau)=1728t\) and \(\wp_\tau(\mu)=0\), determines the periods
throughout a fundamental domain.
\end{proposition}


% ALL PRINTED PROOF MATERIAL MASKED




























\subsection{Bounds for the sharp parameter disc}\hypertarget{tag-01PM}{}\label{fam:section:C.2}
\begin{proposition}\hypertarget{tag-01QJ}{}\label{fam:src:C.2.1}\label{fam:st:C.2.1}
\StatementRecord{fam:src:C.2.1}{proved}{01QJ}{\SeriesRef{fam:src:C.1.1}, \SeriesRef{fam:eq:cd-beta}, \SeriesRef{fam:eq:cd-norm}, \SeriesRef{fam:eq:cd-wp}, \SeriesRef{fam:src:C.2.3}}

For the normalization of \ref{fam:src:C.1.1}, the exact admissibility maximum is finite
and satisfies
\[
 -\frac3\pi\log6\ \leq M<2.
\]
More precisely the proof below gives \(M<1.635\). Consequently the sharp
radius \(R_w=e^{-\pi M}\) satisfies
\(e^{-2\pi}<R_w\leq216\), and for the reference linearization with
\(u=w^2\), \(e^{-4\pi}<R_u\leq46656\).
\end{proposition}


% ALL PRINTED PROOF MATERIAL MASKED








































































\begin{proposition}\hypertarget{tag-01QK}{}\label{fam:src:C.2.2}\label{fam:st:C.2.2}
\StatementRecord{fam:src:C.2.2}{proved}{01QK}{the definitions and constructions stated in the proof}

The point \(t=1\) is a strict local maximum of \(D_{\rm part}\).
A double-precision calculation on 5151 points of the half modular domain
\(0\leq\operatorname{Re}\tau\leq1/2\), \(|\tau|\geq1\),
\(\operatorname{Im}\tau\leq3\), has largest sampled value at \(\tau=i\).
This finite computation does not establish the global maximum.
\end{proposition}


% ALL PRINTED PROOF MATERIAL MASKED













































































































\begin{theorem}[Certified bound]\hypertarget{tag-01QL}{}\label{fam:src:C.2.3}\label{fam:st:C.2.3}\label{ari:C.8.5}\label{push:src:E.4.3}
\StatementRecord{fam:src:C.2.3}{proved}{01QL}{\SeriesRef{fam:src:C.2.2}, \SeriesRef{fam:eq:cd-norm}, \SeriesRef{fam:src:C.2.1}, \SeriesRef{fam:eq:cd-wp}, \SeriesRef{fam:src:C.1.1}, \SeriesRef{fam:src:C.1.2}}

For the cusp normalization $\beta^{\rm part}+\tau\to0$,
\[
 M=-\frac3\pi\log6=-1.71100425815613411418\ldots,
 \qquad \operatorname*{arg\,max}_{B\setminus\{p_0\}}D_{\rm part}=\{1\}.
\]
The exact parameter discs are
\[
 0<|w|<216,
 \qquad 0<|u|<46656
\]
for the reference linearization $u=w^2$; an arbitrary reference constant
$\lambda_0$ multiplies the latter radius by $|\lambda_0|$.
The proof combines an analytic cusp bound, the quantitative local proof
\ref{fam:src:C.2.2}, and a completed outward-rounded interval cover of the remaining
compact domain.
\end{theorem}


% ALL PRINTED PROOF MATERIAL MASKED
























































































































\section{Special values, local expansions and torsion}
\subsection{The normalized Green function and special arguments}
We write $\rho=(1+i\sqrt3)/2$, $T=\operatorname{Im}\tau$, and
\[
 U_\tau(z)=\log N(z,\tau)
 =\log\left|\frac{\vartheta _1(z,\tau)}{\eta(\tau)}\right|
       -\frac{\pi(\operatorname{Im}z)^2}{T}.
\]
Our theta variable has periods $1,\tau$; thus the standard theta series
is evaluated at $\pi z$ when its first argument has period $\pi$.
For a general lattice $\Lambda=\omega(\mathbb Z+\tau\mathbb Z)$ we set
$U_\Lambda(z)=U_\tau(z/\omega)$. The transformation laws of the theta
function show that this definition is independent of the oriented basis.

\begin{proposition}[distribution under isogenies]\hypertarget{tag-000I}{}\label{ari:C.1.1}
\StatementRecord{ari:C.1.1}{proved}{000I}{the definitions and constructions stated in the proof}

Let $\Lambda,\Lambda'\subset\mathbb C$ be lattices and let
$\alpha\ne0$ satisfy $\alpha\Lambda\subset\Lambda'$. Write
$\phi:\mathbb C/\Lambda\to\mathbb C/\Lambda'$ for $z\mapsto\alpha z$,
$K=\alpha^{-1}\Lambda'/\Lambda$, and
$d=[\Lambda':\alpha\Lambda]$. Then
\begin{equation}\hypertarget{tag-003O}{}\label{ari:eq:cs-isogeny}
 U_{\Lambda'}(\alpha z)=\sum_{v\in K}U_\Lambda(z+v).
\end{equation}
The equality holds away from the zeros and as an equality of locally
integrable functions at them. If $\Delta_\Lambda=g_2(\Lambda)^3-
27g_3(\Lambda)^2$, then
\[
 \prod_{0\ne v\in K}N_\Lambda(v)
 =|\alpha|\left|\frac{\Delta_{\Lambda'}}{\Delta_\Lambda}\right|^{1/12}.
\]
In particular, for every nonzero endomorphism of one elliptic curve,
$\prod_{v\in\ker\alpha,\ v\ne0}N(v)=|\alpha|$.
\end{proposition}


% ALL PRINTED PROOF MATERIAL MASKED





































\begin{proposition}[the phases and the two periods]\hypertarget{tag-000J}{}\label{ari:C.1.2}
\StatementRecord{ari:C.1.2}{proved}{000J}{\SeriesRef{ari:C.1.1}}

With the branch continued from the specified cusp, the exact values are
\begin{align*}
 \frac{\vartheta _1((2-\rho)/3,\rho)}{\eta(\rho)}
   &=e^{\pi i/12}\,3^{1/4}e^{\pi/(6\sqrt3)},\\
 \frac{\vartheta _1((1-i)/2,i)}{\eta(i)}
   &=\sqrt2\,e^{\pi/4},\\
 \beta^{\rm part}(z_1)
   &=-\frac12+i\left(\frac1{\sqrt3}-\frac3{2\pi}\log48\right),\\
 \beta^{\rm part}(z_2)
   &=i\left(\frac32-\frac3\pi\log6\right).
\end{align*}
Consequently $N(z_1)=3^{1/4}$, $N(z_2)=\sqrt2$,
$D_{\rm part}(z_1)=-3\log48/(2\pi)$, and
$D_{\rm part}(z_2)=-3\log6/\pi$.
\end{proposition}


% ALL PRINTED PROOF MATERIAL MASKED











































\begin{proposition}[numerical cross-check]\hypertarget{tag-000K}{}\label{ari:C.1.3}
\StatementRecord{ari:C.1.3}{proved}{000K}{\SeriesRef{ari:C.1.2}}

The independent analytic computation in
\path{code/Cspecial_torsion_numerics.py} gives, to fifty decimals,
\begin{align*}
 N(z_1)&=1.31607401295249246081921890179699905516006859020582\ldots,\\
 N(z_2)&=1.41421356237309504880168872420969807856967187537695\ldots,\\
 \operatorname{Im}\beta^{\rm part}(z_1)
  &=-1.27101206057537346286579673538412682778809507094949\ldots,\\
 \operatorname{Im}\beta^{\rm part}(z_2)
  &=-0.21100425815613411418178314271018934802718590587233\ldots.
\end{align*}
The theta/eta evaluations and exact expressions differ by less than
$10^{-100}$ at working precision $110$ decimal digits.
\end{proposition}


% ALL PRINTED PROOF MATERIAL MASKED













\subsection{Division identities and the entire torsion locus}
We use throughout this subsection the rational model
\[
 E_t:\quad y^2=x^3+A(t)x+B(t),\qquad
 A=-\tfrac34t^3(t-1),\quad B=\tfrac14t^4(t-1)^2,
 \quad P=(0,t^2(t-1)/2).
\]
Its discriminant is $\Delta(t)=27t^8(t-1)^3$.
The change of complex Weierstrass model does not change $N$.

\begin{cited}[Siegel functions and local heights]\hypertarget{tag-000P}{}\label{ari:C.4.1}
\StatementRecord{ari:C.4.1}{proved}{000P}{\citet[Remark~2.2 and Proposition~2.3]{Shin2015}}

For $a=(a_1,a_2)\in\mathbb R^2\setminus\mathbb Z^2$ and
$v=a_1\tau+a_2$, define
\[
 g_a(\tau)=-e^{\pi ia_2(a_1-1)}q^{B_2(a_1)/2}
 (1-e^{2\pi iv})\prod_{r\ge1}
 (1-q^re^{2\pi iv})(1-q^re^{-2\pi iv}),
 \quad B_2(x)=x^2-x+\tfrac16.
\]
Then
\[
 g_a(\tau)=i e^{\pi ia_1v}
       \frac{\vartheta_1(v,\tau)}{\eta(\tau)},\qquad
 N(v,\tau)=|g_a(\tau)|.
\]
When $a\in n^{-1}\mathbb Z^2\setminus\mathbb Z^2$, $n\ge2$,
$g_a^{12n}$ is a modular unit for $\Gamma(n)$ with Fourier coefficients
in $\mathbb Q(\zeta_n)$ in the variable $q^{1/n}$. On every complex
elliptic curve the model-independent Néron local function for $(O)$ is
\[
 \lambda_\infty(Q)=-\log N(Q).
\]
Thus $D_{\rm part}=-(6/\pi)(\lambda_\infty(P)+\log(2\sqrt3))$.
\end{cited}
The expanded source derivation is retained in the proof compendium, Paper~I, label \nolinkurl{ari:C.4.1}.




The Siegel-function identities are those of \citet[Remark~2.2 and Proposition~2.3]{Shin2015}, with the local-height normalization proved explicitly. A complete re-derivation is given in \CompendiumRef{I}{app:rederive:ari-C.4.1}.


\begin{proposition}[the full division identity]\hypertarget{tag-000Q}{}\label{ari:C.4.2}
\StatementRecord{ari:C.4.2}{proved}{000Q}{\SeriesRef{ari:C.1.1}}

Let $\psi_m$ be the standard division polynomial of a smooth
Weierstrass model, normalized by $\psi_1=1$, $\psi_2=2y$.
For $m\ge2$ and $Q,mQ\ne O$,
\begin{equation}\hypertarget{tag-003N}{}\label{ari:eq:cs-division}
 N(Q)^{m^2}
       =N(mQ)\frac{|\Delta_E|^{(m^2-1)/12}}{|\psi_m(Q)|}.
\end{equation}
For the displayed family this specializes to
\[
 N(P)^4=3^{3/4}|t-1|^{-1/4}N(2P),\qquad
 N(P)^9=16|t|^{-2/3}N(3P).
\]
If $P_t$ has exact order $n$ and $m\ge2$ satisfies
$m\equiv1$ or $-1\pmod n$, then
\begin{equation}\hypertarget{tag-003P}{}\label{ari:eq:cs-torsionvalue}
 N(P_t)^{12}
       =\frac{|\Delta(t)|}{|\psi_m(P_t)|^{12/(m^2-1)}}.
\end{equation}
The positive value on the right is independent of $m$.
\end{proposition}


% ALL PRINTED PROOF MATERIAL MASKED




































\begin{proposition}[algebraicity and explicit fields for every order]\hypertarget{tag-000R}{}\label{ari:C.4.3}
\StatementRecord{ari:C.4.3}{proved}{000R}{\SeriesRef{ari:eq:cs-torsionvalue}}

Let $t_0\in\mathbb C\setminus\{0,1\}$ and suppose $P_{t_0}$ has
exact order $n\ge2$. Then $t_0$ and $N(P_{t_0})$ are algebraic.
Put $K_0=\mathbb Q(t_0,\bar t_0)$,
\[
 W=\Delta(t_0)\Delta(\bar t_0)>0,\qquad
 V=\psi_{n+1}(P_{t_0})\psi_{n+1}(P_{\bar t_0})>0.
\]
With positive real radicals,
\[
 N(P_{t_0})^{12n}=W^{n/2}V^{-6/(n+2)}
 \ \in\ K_0\bigl(W^{1/2},V^{1/(n+2)}\bigr).
\]
This supplies an explicit field, without a claim that it is minimal.
More generally, if $m\equiv\pm1\pmod n$, then
\[
 \lambda_\infty(P_{t_0})
 =\frac{\log|\psi_m(P_{t_0})|}{m^2-1}
       -\frac1{12}\log|\Delta(t_0)|.
\]
These formulas determine the height and $N^{12}$ on the entire
torsion locus.
\end{proposition}


% ALL PRINTED PROOF MATERIAL MASKED







































\begin{proposition}[the five values and all orders through eight]\hypertarget{tag-000S}{}\label{ari:C.4.4}
\StatementRecord{ari:C.4.4}{proved}{000S}{\SeriesRef{ari:eq:cs-torsionvalue}, \SeriesRef{ari:C.1.2}, \SeriesRef{ari:C.4.3}}

The exact special values are
\[
 \begin{array}{c|c|c}
 t&\operatorname{ord}(P)&N(P)^{12}\\ \hline
 0&3&27\\
 1&2&64\\
 32/5&4&10\\
 64(47-27\sqrt3)/11&6&17+10\sqrt3\\
 64(47+27\sqrt3)/11&6&10\sqrt3-17 .
 \end{array}
\]
At $t=0,1$ the orders refer to the smooth models after the specified
base changes. On the smooth base, $F_2=F_3=1$ and
\[
 \begin{aligned}
 F_4&=5t-32,\\
 F_5&=89t^2+4736t-4096,\\
 F_6&=11t^2-6016t+8192,\\
 F_7&=881t^4+4745344t^3-12079104t^2\\
    &\hspace{13mm}+24641536t-16777216,\\
 F_8&=49t^4-2614336t^3-7864320t^2\\
    &\hspace{13mm}+19398656t-8388608.
 \end{aligned}
\]
All these polynomials are squarefree. At each root of exact order
$n=4,5,6,7,8$, the positive number $r=N^{12}$ is respectively
determined by
\[
 \begin{array}{c|l}
 n&\text{equation for }r\\ \hline
 4&r=64/|t|\\
 5&r^5=432^5/\bigl((1-t)(5t-32)^4\bigr)\\
 6&r^2=1728^2/F_5(t)\\
 7&r^{35}=2^{192}3^{81}/
       \bigl(|t|^8|t-1|^3|F_6(t)|^{12}\bigr)\\
 8&r^4=1728^4/|F_7(t)|.
 \end{array}
\]
The order-five roots are real and below one; the order-six roots
are real and above one. The two complex order-seven roots are
included in the formula.
\end{proposition}


% ALL PRINTED PROOF MATERIAL MASKED

















































\begin{proposition}[complete numerical verification through order eight]\hypertarget{tag-000T}{}\label{ari:C.4.5}
\StatementRecord{ari:C.4.5}{proved}{000T}{the definitions and constructions stated in the proof}

All $13$ smooth exact-order roots for $2\le n\le8$ were evaluated
at $110$ decimal digits independently of the algebraic norm formula.
The discrepancy in $N^{12}$ is below $10^{-90}$ at every root;
the period coordinates also verify $nP=O$ to that tolerance.
There are $0,0,1,2,2,4,4$ roots for $n=2,\ldots,8$.
\end{proposition}


% ALL PRINTED PROOF MATERIAL MASKED



























\section{An analytic proof of the elliptic maximum}

Throughout this section $q=e^{2\pi i\tau}$, $T=\Im\tau$,
$\Delta_0=q\prod_{n\ge1}(1-q^n)^{24}$, and
$t=E_4^3/(1728\Delta_0)$. Thus $t-1=E_6^2/(1728\Delta_0)$.
The discriminant $\Delta_0$ is the modular discriminant and must not
be confused with the discriminant $27t^8(t-1)^3$ of the rational
Weierstrass equation. We use the division identity
\eqref{ari:eq:cs-division}, already proved above, and the modular
invariance of $N$. No interval cover or numerical root computation
enters the following proof.

\begin{proposition}[three uniform bounds on the elliptic norm]\hypertarget{tag-001J}{}\label{ari:C.8.1}
\StatementRecord{ari:C.8.1}{proved}{001J}{\SeriesRef{ari:eq:ca-sup}, \SeriesRef{ari:eq:ca-exp}}

For every $z\in\C$ and every $\tau$ in the standard modular
fundamental region, we have
\[
 \begin{array}{c|ccc}
 \text{restriction}&T\le2&T\le7/5&T\le11/10\\ \hline
 N(z,\tau)&<17/10&<3/2&<143/100.
 \end{array}
\]
If $\wp_\tau(z)=0$ and $T\ge2$, then $N(z,\tau)<541/500$.
\end{proposition}


% ALL PRINTED PROOF MATERIAL MASKED





























































































\begin{proposition}[an explicit local maximum on a larger disc]\hypertarget{tag-001Q}{}\label{ari:C.8.2}
\StatementRecord{ari:C.8.2}{proved}{001Q}{\SeriesRef{ari:eq:ca-exp}, \SeriesRef{ari:eq:ca-derivatives}, \SeriesRef{ari:eq:ca-GHbounds}, \SeriesRef{ari:eq:ca-H}, \SeriesRef{ari:eq:ca-delta-h}}

If $|\tau-i|\le1/50$ and $\wp_\tau(z)=0$, then
$N(z,\tau)\le\sqrt2$, with equality exactly when $\tau=i$.
\end{proposition}


% ALL PRINTED PROOF MATERIAL MASKED







































































































































































\begin{proposition}[the comparison of the two local discs]\hypertarget{tag-001R}{}\label{ari:C.8.3}
\StatementRecord{ari:C.8.3}{proved}{001R}{\SeriesRef{ari:eq:ca-derivatives}, \SeriesRef{ari:C.8.2}}

Every complex number $t_0$ with $|t_0-1|\le1/500$ has a
representative $\tau$ with $|\tau-i|<1/50$ and $j(\tau)=1728t_0$.
Consequently every zero of $\wp$ for that modulus satisfies
$N\le\sqrt2$, with equality only for $t_0=1$.
\end{proposition}


% ALL PRINTED PROOF MATERIAL MASKED




























\begin{proposition}[the first three normalized duplication steps]\hypertarget{tag-001S}{}\label{ari:C.8.4}
\StatementRecord{ari:C.8.4}{proved}{001S}{\SeriesRef{ari:C.4.2}, \SeriesRef{ari:eq:cs-division}, \SeriesRef{ari:C.3.1}}

Put $u=t-1$. For $n=2^k$, $k=1,2,3$, and $nP\ne O$,
\begin{equation}\hypertarget{tag-002D}{}\label{ari:eq:ca-ladder}
 \left(\frac{N(P)}{\sqrt2}\right)^{n^2}
 =\frac{3^{3/4}}{2n}\,
       \frac{N(nP)|u|^{-1/4}}{|A_k(u)|},
\end{equation}
where
\begin{align*}
 A_1(u)&=1,\qquad A_2(u)=1-5u/27,\\
 A_3(u)&=1-\frac{217}{27}u-\frac{20486}{729}u^2
       +\frac{10910}{19683}u^3
       +\frac{161383}{177147}u^4
       -\frac{245}{14348907}u^5.
\end{align*}
On $|u|\le1/28$ we have $|A_3(u)|\ge1-10|u|$.
For every $k\ge1$, not just these first three steps, define
\[
 E_k(t)=\frac{|\Delta_E|^{(4^k-1)/12}|t-1|^{1/4}}
 {2^{4^k/2}|\psi_{2^k}(P)|}.
\]
Then $(N(P)/\sqrt2)^{4^k}=E_k(t)N(2^kP)|t-1|^{-1/4}$,
and the removable limiting value is
\[
 \lim_{t\to1}E_k(t)=\frac{3^{3/4}}{2^{k+1}}
                 =2^{1-k}\frac{3^{3/4}}4.
\]
\end{proposition}


% ALL PRINTED PROOF MATERIAL MASKED

































% ALL PRINTED PROOF MATERIAL MASKED































































Theorems~\ref{fam:src:C.2.3} and \cref{push:src:E.4.1}, \cref{push:src:E.4.2} establish this inequality by
computer-assisted proofs. Statements \ref{ari:C.8.1}--\ref{ari:C.8.5} give another
proof with a three-step division ladder and the explicit radius
$1/50$ in the $\tau$-plane. The programs
\path{code/C_ladder_estimates.py} and
\path{code/C_local_analytic_bounds.py} record numerical diagnostics
of the constants; the proof uses the displayed rational estimates
and their series remainder bounds alone.




\section{The retained family and its intrinsic classification}\hypertarget{tag-01P9}{}\label{fam:section:A}

Throughout this section $\mathcal P=\{c:\operatorname{Im}c<-M\}$, and $X_c$ means the actual glued space with $v_1=\varepsilon=(1,2,-4,0)^t$, $v_2=-\varepsilon'=(-1,-3,3,0)^t$ and the untwisted cusp collar. We write $d=\widehat\delta$, $e=(0,1)^t$, $E=d\gamma$, and use column vectors. Our peripheral generators are clockwise: $s_j(g_jz)=e^{-2\pi i/m_j}s_j(z)$ and $g_1g_2g_0=1$. Thus $s(g_0z)=s(z)-1$ at the cusp. All period comparisons use
\[
 A_1=\begin{pmatrix}1&0&0&0\\6&0&1&0\\-6&-1&-1&0\\-2&1&0&1\end{pmatrix},\quad
 A_2=\begin{pmatrix}1&0&0&0\\0&0&-1&0\\-6&1&0&0\\3&0&1&1\end{pmatrix},\qquad
 R_g(z)\Pi_c(z)=\Pi_c(gz)A_g.
\]
Here $R_1=\left(\begin{smallmatrix}-1/\tau&0\\(1-\mu)/\tau&1\end{smallmatrix}\right)$, $R_2=\left(\begin{smallmatrix}1/\tau&0\\-\mu/\tau&1\end{smallmatrix}\right)$, and $R_0=I$. The distinction between source and target markings in each equation below is retained by its displayed period equation.

\subsection{Canonical sections and the finite covers}\hypertarget{tag-01PA}{}\label{fam:section:A.1}

\begin{proposition}\hypertarget{tag-01PZ}{}\label{fam:src:A.1.1}\label{fam:st:A.1.1}
\StatementRecord{fam:src:A.1.1}{proved}{01PZ}{\citet[Theorem~9.1]{AC2026}}

For the constructed holomorphic map $f_c:X_c\to\mathbb P^1$, and also for every admissible choice of its affine translations and cusp collar, we have
\[
 \omega_{X_c}\simeq\mathcal O_{X_c}(-2S_2),\qquad
 f_{c*}\mathcal O_{X_c}(kS_2)=\mathcal O_B(\lfloor k/4\rfloor p_2)\quad(k\in\mathbb Z).
\]
The latter equality is an equality of subsheaves of meromorphic functions under pullback. The anticanonical section ring is $\mathbb C[U,V]$, with $\deg U=1$, $\deg V=2$. The complete degree-two map is $f_c$ followed by a projective coordinate isomorphism of its base. Every biholomorphism between two such threefolds is over the identity of $B$.
\end{proposition}



The canonical pencil is due to \citet[Theorem~9.1]{AC2026}. A complete re-derivation is given in \CompendiumRef{I}{app:rederive:fam-src-A.1.1}.


\begin{proposition}\hypertarget{tag-01Q0}{}\label{fam:src:A.1.2}\label{fam:st:A.1.2}
\StatementRecord{fam:src:A.1.2}{proved}{01Q0}{the definitions and constructions stated in the proof}

The normalized base change $N_j\times_{D_j}\Delta_j\to\Delta_j$, $t_j=s_j^{m_j}$, is the original smooth torus family $J'_j$. The maps $E_1\times F_1\to S_1$ and $E_2\times F_2\to S_2$ described below are finite etale covers of degrees $9$ and $8$.
\end{proposition}


% ALL PRINTED PROOF MATERIAL MASKED


































\subsection{The integral obstruction}\hypertarget{tag-01PB}{}\label{fam:section:A.2}

\begin{proposition}\hypertarget{tag-01Q1}{}\label{fam:src:A.2.1}\label{fam:st:A.2.1}
\StatementRecord{fam:src:A.2.1}{proved}{01Q1}{\SeriesRef{fam:src:A.1.2}}

The rational centralizer of $(A_1,A_2)$ is $\mathbb QI+\mathbb QE$. Every biholomorphism $X_c\to X_{c'}$ has $c'-c=n\in2\mathbb Z$, complex-linear fibre part $I_2$, and source-to-target lattice map $I-nE$. In the full fixed lattices,
\[
 \Lambda^{A_1}/N_1\Lambda\simeq\mathbb Z/3,\qquad
 \Lambda^{A_2}/N_2\Lambda\simeq\mathbb Z/4\oplus\mathbb Z/2,
\]
with $d$ generating the second factor of the latter group and vanishing in the former.
\end{proposition}


% ALL PRINTED PROOF MATERIAL MASKED






























\subsection{Extension and coherent descent}\hypertarget{tag-01PC}{}\label{fam:section:A.3}

\begin{proposition}\hypertarget{tag-01Q2}{}\label{fam:src:A.3.1}\label{fam:st:A.3.1}
\StatementRecord{fam:src:A.3.1}{proved}{01Q2}{\SeriesRef{fam:src:A.2.1}}

For each $n\in2\mathbb Z$ there is a biholomorphism $F_{n,c}:X_c\to X_{c+n}$, holomorphic jointly in $c$, with $F_{2h,c+2k}F_{2k,c}=F_{2(h+k),c}$. The quotient family over $0<|w|<e^{-\pi M}$, $w=e^{-\pi ic}$, is proper and holomorphic, with pairwise non-biholomorphic fibres.
\end{proposition}


% ALL PRINTED PROOF MATERIAL MASKED


















































\subsection{Albanese invariants and the derivative of the deformation}\hypertarget{tag-01PD}{}\label{fam:section:A.4}

\begin{proposition}\hypertarget{tag-01Q3}{}\label{fam:src:A.4.1}\label{fam:st:A.4.1}
\StatementRecord{fam:src:A.4.1}{proved}{01Q3}{\SeriesRef{fam:src:A.1.2}, \SeriesRef{fam:src:A.2.1}}

With $\kappa_j(c)=\kappa_j^{\rm part}+c$, the Albanese elliptic curves of the reduced finite fibres are
\[
 \operatorname{Alb}(S_1)=\mathbb C/(\mathbb Z+\kappa_1(c)\mathbb Z),\qquad
 \operatorname{Alb}(S_2)=\mathbb C/(\mathbb Z+\tfrac12\kappa_2(c)\mathbb Z).
\]
Writing $J(q)=q^{-1}+744+196884q+\cdots$, their invariants are
\[
 j_1(w)=J(e^{-2\pi i\kappa_1^{\rm part}}w^2),\qquad
 j_2(w)=J(e^{-\pi i\kappa_2^{\rm part}}w).
\]
These are convergent formulas because the displayed nome moduli are below one. Outside the explicitly defined countable exceptional set in the proof, equality of both Albanese invariants forces $c'-c\in2\mathbb Z$. Thus the generic Albanese invariants recover the parameter modulo its exact period subgroup, independently of the centralizer calculation.
\end{proposition}


% ALL PRINTED PROOF MATERIAL MASKED


















\begin{proposition}\hypertarget{tag-01Q4}{}\label{fam:src:A.4.2}\label{fam:st:A.4.2}
\StatementRecord{fam:src:A.4.2}{proved}{01Q4}{\SeriesRef{fam:src:A.1.1}}

For every $c\in\mathcal P$, the Kodaira--Spencer class of $\partial/\partial c$ in $H^1(X_c,T_{X_c})$ is nonzero. On any fixed ordinary fibre, put $T=\operatorname{Im}\tau$, $m=\operatorname{Im}\mu$, $a=m/T$, $D=\operatorname{Im}\beta_c-6m^2/T$. In the convention $\mathrm{KS}(\partial_c)=[\bar\partial h]$ for a smooth $(1,0)$ lift $h$, its restricted torus-deformation representative is
\[
 -\frac{1}{2iD}\,\partial_{\zeta_2}\otimes
       (d\bar\zeta_2-a\,d\bar\zeta_1).
\]
\end{proposition}


% ALL PRINTED PROOF MATERIAL MASKED


















\begin{proposition}\hypertarget{tag-01Q5}{}\label{fam:src:A.4.3}\label{fam:st:A.4.3}
\StatementRecord{fam:src:A.4.3}{proved}{01Q5}{\SeriesRef{fam:src:A.4.2}, \SeriesRef{fam:src:A.1.1}, \SeriesRef{fam:src:A.2.1}, \SeriesRef{fam:src:A.3.1}; \cite[Proposition 4.7]{ViaclovskyLog2026}; \cite[Lemmas 6.6--6.9 and Proposition 6.15]{Yang2026}}

The non-isotriviality argument of Viaclovsky, \cite[Proposition 4.7]{ViaclovskyLog2026}, follows for the present family from \ref{fam:src:A.4.2} and \ref{fam:src:A.1.1}. The cited countability argument proves local non-isotriviality and positive Kuranishi dimension; the computation in \ref{fam:src:A.4.2} additionally proves nonzero derivative at every admissible parameter. The arithmetic method of Yang, \cite[Lemmas 6.6--6.9 and Proposition 6.15]{Yang2026}, specializes to the same centralizer, finite-norm and collar-gluing operations after replacing his displayed matrices by $A_1,A_2$; those replacements are evaluated in \ref{fam:src:A.2.1} and \ref{fam:src:A.3.1}.
\end{proposition}


% ALL PRINTED PROOF MATERIAL MASKED








\subsection{The assembled theorem}\hypertarget{tag-01PE}{}\label{fam:section:A.5}

\begin{theorem}\hypertarget{tag-01Q6}{}\label{fam:src:A.5.1}\label{fam:st:A.5.1}\label{wb:pshift:thm:classification}
\StatementRecord{fam:src:A.5.1}{proved}{01Q6}{\SeriesRef{app:rederive:fam-src-A.5.1}; \citet[Theorem~1.4, p.~3]{DuEtAl2026v1}}

Fix the original cusp-regular period solution, the clockwise peripheral marking, the prescribed periodic triangular cusp fan, the affine vectors $(\varepsilon,-\varepsilon')$, and the untwisted cusp collar. For every $c,c'\in\mathcal P$, the constructed compact complex threefolds satisfy
\[
 X_c\simeq X_{c'}\quad\Longleftrightarrow\quad c'-c\in2\mathbb Z.
\]
Every biholomorphism is over the identity of the original $\mathbb P^1$. The proper holomorphic family of their isomorphism classes has parameter $0<|w|<e^{-\pi M}$ and nonzero Kodaira--Spencer class at every point.
\end{theorem}



The standard classification is \citet[Theorem~1.4, p.~3]{DuEtAl2026v1}. A complete re-derivation is given in Appendix~\ref{app:rederive:fam-src-A.5.1}.


\begin{theorem}\hypertarget{tag-01RH}{}\label{fam:src:A.5.2}\label{fam:st:A.5.2}\label{inv:D1.3}\label{inv:D2.1}\label{inv:D3.1}\label{cdp:st:E9}
\StatementRecord{fam:st:A.5.2}{proved}{01RH}{\SeriesRef{fam:src:A.5.1}, \SeriesRef{app:rederive:fam-src-A.5.2}, \SeriesRef{fam:section:B.2}, \SeriesRef{fam:section:B.3}, \SeriesRef{fam:eq:ab-normal-quotient}; \citet[Theorem~6.12, pp.~81--82, and Proposition~6.16, p.~85]{DuEtAl2026v1}; \cite[Sections 2--3 and Proposition 6.1]{ViaclovskyTwo2026}}

For every admissible $c$, the tangent cohomology of the actual compact
member is
\[
 \bigl(h^0(X_c,T_{X_c}),h^1(X_c,T_{X_c}),h^2(X_c,T_{X_c}),
 h^3(X_c,T_{X_c})\bigr)=(1,1,1,0).
\]
Its Kuranishi germ is a smooth reduced disc. The family of \ref{fam:src:A.5.1} is
miniversal at every parameter: its Kodaira--Spencer map is an isomorphism,
and every sufficiently small deformation is induced from this family.
More precisely, for the intrinsic fibration $f:X_c\to\mathbb P^1$,
\[
 f_*T_{X_c}=\mathcal O\oplus\mathcal O(-1),\qquad
 R^1f_*T_{X_c}=\mathcal O\oplus\mathcal O(-1)^2,\qquad
 R^2f_*T_{X_c}=\mathcal O.\tag{D.4}
\]
The Leray class $\xi\otimes[\gamma]$ generates $H^1(X_c,T_{X_c})$.
The primary Kuranishi bracket is zero; the one-dimensional obstruction
space $H^2(X_c,T_{X_c})$ therefore does not imply a nonzero obstruction
map.

\end{theorem}



The tangent dimensions and Kuranishi germ are re-derived from \citet[Theorem~6.12, pp.~81--82, and Proposition~6.16, p.~85]{DuEtAl2026v1}. A complete re-derivation is given in Appendix~\ref{app:rederive:fam-src-A.5.2}.


\paragraph{Scope and a precise changed-input test.}
The hypotheses of \ref{fam:src:A.5.1}--\ref{fam:src:A.5.2} fix the original period system and
multiplicities $(3,4)$; \ref{fam:section:B.2} treats its arbitrary affine and cusp data,
and \ref{fam:section:B.3} classifies its stipulated higher-multiplicity variants.
A concrete extension beyond these hypotheses begins with
$y^2=x^3+tx+1$ in \cite[Sections 2--3 and Proposition 6.1]{ViaclovskyTwo2026}.
Already one rank-one filling over a specified root $p_1$ of
$4t^3+27=0$ replaces the hexagonal normalization calculation by
\[
 \begin{gathered}
 E'=\mathbb C^*/\langle c_0\rangle,\qquad
 R=\mathbb P_{E'}(\mathcal O_{E'}\oplus L),\\
 R=(\mathbb P^1\times\mathbb C^*)/
       \langle(z,v)\mapsto(\gamma_0z,c_0v)\rangle.
 \end{gathered}
\]
with its boundary sections identified by $v\mapsto\beta_0v$.
Here $c_0,\gamma_0,\beta_0$ are exactly the nonzero central multipliers
in that paper's equation (3.1). The invariant nodal volume form there
trivializes the ambient canonical line, so the branch-normal map has
quotient $\nu_*K_R$ by the same local coefficient map as
\eqref{fam:eq:ab-normal-quotient}. The ruling has
$\pi_*\mathcal O_R=\mathcal O_{E'}$ and $R^1\pi_*\mathcal O_R=0$;
Leray and Serre duality therefore give $h^\bullet(K_R)=(0,1,1)$.
The first unevaluated global datum for transferring the deformation
argument is the connecting map from this single
$H^1(R,K_R)=\mathbb C$ summand into $H^2(T_X(-\log D))$, where
$D$ is the union of the three reduced singular fibres of that model.
Its direct attack is to compute the boundary-translation Cech cocycle
using the full holomorphic unit functions in that equation (3.1),
their central values $(\gamma_0,c_0,\beta_0)$, and the corresponding
vertical frames.
The global identification argument also requires the integral
centralizer and collar defects of that model's four monodromy matrices.
These are specified substitutions to perform for the stated one-fibre
test; no classification or deformation conclusion for that changed
period system is included in \ref{fam:src:A.5.1}--\ref{fam:src:A.5.2}.

\section{Discrete data and the geometric parameter}\hypertarget{tag-01PF}{}\label{fam:section:B}

\subsection{All affine vectors and cusp sections}\hypertarget{tag-01PH}{}\label{fam:section:B.2}

\begin{proposition}\hypertarget{tag-01QB}{}\label{fam:src:B.2.1}\label{fam:st:B.2.1}
\StatementRecord{fam:src:B.2.1}{proved}{01QB}{\SeriesRef{fam:src:A.1.1}, \SeriesRef{fam:src:A.3.1}; \cite[Section 6.4]{AC2026}}

Let
\[
 w_1=r_1\varepsilon+k_1d,\qquad w_2=r_2\varepsilon'+k_2d,
 \qquad 3\nmid r_1,\quad 2\nmid r_2,
\]
and use an arbitrary holomorphic cusp collar section $\sigma$ in the construction of \cite[Section 6.4]{AC2026}. Write $l_0=\gamma(c(\sigma))$ in its positive cusp-winding convention, and set $p=12l_0-4r_1-3r_2$. Let primes denote a second such datum with the same period solution and cusp fan. Then
\[
 X(c;w_1,w_2,\sigma)\simeq X(c';w'_1,w'_2,\sigma')
\]
holds exactly for the following arithmetic conditions:
\[
 n=c'-c\in\mathbb Z,\qquad
 p'=a p\text{ for some }a\in\{1,-1\},\qquad
 n\equiv k_2-k'_2\pmod2.
\]
All maps are over $B$. A map with the indicated $a$ has fibre-linear part $aI_2$. In particular the complete invariants are
\[
 \bigl(|p|,\ e^{-\pi i(c+k_2)}\bigr),\qquad
 |p|\ge1,\quad\gcd(|p|,6)=1.
\]
Every value of this positive integer occurs. The full collar section enters the construction, but its additional analytic part and its second integral winding contribute no further invariant.
\end{proposition}


% ALL PRINTED PROOF MATERIAL MASKED















































\begin{proposition}\hypertarget{tag-01QC}{}\label{fam:src:B.2.2}\label{fam:st:B.2.2}
\StatementRecord{fam:src:B.2.2}{proved}{01QC}{\SeriesRef{fam:src:A.5.1}, \SeriesRef{fam:src:B.2.1}, \SeriesRef{fam:src:A.2.1}}

For multiplicities $(3,4)$ and $|p|=1$, every admissible discrete member is biholomorphic to one member of \ref{fam:src:A.5.1}, uniquely modulo $2\mathbb Z$. At fixed nonzero admissible Engel parameter $u$ there are exactly two isomorphism classes, exchanged by a positive circuit of $u$ about zero. In the convention that Engel's root generator is the inverse of our clockwise one and his circle class is $d$, the identity of the middle quotient sends
\[
 X_{c+1}\longleftrightarrow X(u;a_0-d/3,\ a_1+d/4).
\]
Reversing both declared root generators reverses these two displayed signs.
\end{proposition}


% ALL PRINTED PROOF MATERIAL MASKED
















\begin{proposition}[counterexample, with corrected family]\hypertarget{tag-01QD}{}\label{fam:src:B.2.3}\label{fam:st:B.2.3}
\StatementRecord{fam:src:B.2.3}{proved}{01QD}{\SeriesRef{fam:src:A.3.1}, \SeriesRef{fam:src:A.2.1}, \SeriesRef{fam:src:B.2.1}, \SeriesRef{fam:src:A.4.1}; \cite[Remark 8.1]{Engel2026}}

The assertion in \cite[Remark 8.1]{Engel2026} of a global proper family over the original punctured $u$-disc for fixed sphere logarithmic data requires a double-cover qualification. The retained family with such data does not descend through $u=w^2$. It does give a proper holomorphic family on each simply connected $u$-open set, and a global proper family over the punctured $w$-disc as in \ref{fam:src:A.3.1}.
\end{proposition}


% ALL PRINTED PROOF MATERIAL MASKED








\subsection{Higher multiplicities}\hypertarget{tag-01PI}{}\label{fam:section:B.3}

\begin{proposition}\hypertarget{tag-01QE}{}\label{fam:src:B.3.1}\label{fam:st:B.3.1}
\StatementRecord{fam:src:B.3.1}{proved}{01QE}{\SeriesRef{fam:src:A.1.2}, \SeriesRef{fam:src:A.2.1}, \SeriesRef{fam:src:B.2.1}}

Let $k,l$ be positive integers with $\gcd(3k,4l)=1$, and put $m_1=3k,m_2=4l$. Consider Viaclovsky's sphere members expressed in the same clockwise marking, with $w_j=r_j\varepsilon_j+s_jd$, $\gcd(r_j,m_j)=1$, cusp winding $l_0$, and $p=m_1m_2l_0-m_2r_1-m_1r_2=\pm1$, where $\varepsilon_1=\varepsilon,\varepsilon_2=\varepsilon'$. Two such members are biholomorphic precisely when their multiplicity pair agrees and, for some $a\in\{1,-1\}$ and integer $n=c'-c$,
\[
 p'=ap,\quad r'_j\equiv ar_j\pmod{m_j},\quad
 nr'_1\equiv as_1-s'_1\pmod{k},\quad
 nr'_2\equiv as_2-s'_2\pmod{2l}.
\]
The congruences on $r'_j$ already follow from $p'=ap$. For fixed discrete vectors and cusp collar the identification subgroup is $2kl\mathbb Z$. For fixed admissible $u$ and this multiplicity pair, all the stated sphere translation data give exactly $2kl$ isomorphism classes, permuted in one cycle by a circuit of $u$ about zero.
\end{proposition}


% ALL PRINTED PROOF MATERIAL MASKED



















\subsection{Automorphisms and the remaining real-structure calculation}\hypertarget{tag-01PJ}{}\label{fam:section:B.4}

\begin{proposition}\hypertarget{tag-01QF}{}\label{fam:src:B.4.1}\label{fam:st:B.4.1}
\StatementRecord{fam:src:B.4.1}{proved}{01QF}{\SeriesRef{app:rederive:fam-src-B.4.1}; \citet[Theorem~7.14, p.~95, label \nolinkurl{p2:thm:full-automorphisms}]{DuEtAl2026v1}}

For each member of the retained $(3,4)$ family, $\operatorname{Aut}(X_c)$, including all its components, is $\mathbb C/\mathbb Z\simeq\mathbb C^*$. Its action on the ordinary fibres is $\zeta\mapsto\zeta+be$, and the exponential group isomorphism is $b\mapsto e^{2\pi ib}$.
\end{proposition}
\Uses{the hypotheses, notation and definitions stated here; \citet[Theorem~1.4, p.~3]{DuEtAl2026v1}; \cref{app:rederive:fam-src-B.4.1}}


The full automorphism group is \citet[Theorem~1.4, p.~3]{DuEtAl2026v1}. A complete re-derivation is given in Appendix~\ref{app:rederive:fam-src-B.4.1}.



\begin{proposition}\hypertarget{tag-01QG}{}\label{fam:src:B.4.2}\label{fam:st:B.4.2}
\StatementRecord{fam:src:B.4.2}{proved}{01QG}{\SeriesRef{fam:src:B.4.1}, \SeriesRef{fam:src:A.3.1}, \SeriesRef{fam:src:B.2.1}, \SeriesRef{fam:src:A.5.1}}

Use the normalization $(\beta^{\rm part}+\tau)(p_0)=0$. There is an antiholomorphic diffeomorphism $X_c\to X_{c'}$ exactly when
\[
 c'-(1-\bar c)\in2\mathbb Z,
 \qquad\text{equivalently }w'=-\bar w.
\]
Consequently $\overline{X_c}\simeq X_{1-\bar c}$. The members admitting antiholomorphic self-maps are exactly those with $\operatorname{Re}c\in\tfrac12+\mathbb Z$, equivalently $w=-\bar w$. On each such member the group of holomorphic and antiholomorphic diffeomorphisms is
\[
 \operatorname{Aut}^{\pm}(X_c)\simeq\mathbb C^*\rtimes\mathbb Z/2,
 \qquad z\longmapsto1/\bar z
\]
for the action of the second factor. Its antiholomorphic maps form the nonidentity coset. For all other members the full group is the holomorphic group of \ref{fam:src:B.4.1}.
\end{proposition}


% ALL PRINTED PROOF MATERIAL MASKED





















































\section{Classification and deformation questions}\hypertarget{tag-02ZD}{}\label{push:section:E}
\subsection{Higher logarithmic transformations}\hypertarget{tag-02ZE}{}\label{push:section:E.1}
We distinguish Viaclovsky's indices $0,1$ from the indices $1,2$ of
the question.  Put $m=3k$, $n=4l$, and assume $\gcd(m,n)=1$.
The lattice coordinates in this subsection are integral homology
coordinates, and matrices act on columns.

\begin{proposition}[the integral dictionary]\hypertarget{tag-032R}{}\label{push:src:E.1.1}
\StatementRecord{push:src:E.1.1}{proved}{032R}{\cite[Section 2, equations (2.1), (2.2), (2.7), (2.11)--(2.12)]{ViaclovskyLog2026}}

Let $(\delta,a,b,d)$ be the marking in
\cite[Section 2, equations (2.1), (2.2), (2.7), (2.11)--(2.12)]{ViaclovskyLog2026},
and put $u_0=d+2a+2b$, $u_1=d+3a+3b$.
The homomorphism from the marking of the question to this marking is
\[
 P_V=\begin{pmatrix}
 -2&1&0&1\\-6&0&-1&0\\0&-1&0&0\\-1&0&0&0
 \end{pmatrix},\qquad \det P_V=-1.
\]
It intertwines the two finite monodromies and satisfies
\[
 P_V\widehat\delta=\delta,\qquad
 P_V\varepsilon=-u_0,\qquad
 P_V\varepsilon'=\delta-u_1.
\]
With the counterclockwise deck generators of that paper, its data
$m_0=3k$, $m_1=4l$ and
$a_i=[v_i/m_i]$ with $v_i=\rho_i\delta+N_i u_i$ correspond to
\[
 w_1=N_0\varepsilon-\rho_0\widehat\delta,
 \qquad
 w_2=N_1\varepsilon'-(\rho_1+N_1)\widehat\delta.
\]
For the section-preserving cusp comparison used there, $l_0=0$ and
$p=-(nN_0+mN_1)$.  Thus its sphere condition is
$b:=nN_0+mN_1=\pm1$.
The brackets in $a_i$ mean the torsion point in the smooth
root torus, using its section-preserving origin and the
transported period vector.  Replacing $v_i$ by an integral
multiple $m_i$ times a period does not change that torsion
point, although its chosen lift fixes the collar bookkeeping.
\end{proposition}


% TARGET PROOF MASKED: push:src:E.1.1
























For a general cusp collar $\sigma$, write $c(\sigma)$ for
its residual translation defect in
$\overline\Lambda=\Lambda/(\mathbb Z\widehat w+
\mathbb Z\widehat\delta)$, using the convention
$\sigma(g_0z)-R_0\sigma(z)=\Pi_c(g_0z)\ell(\sigma)$
and $c(\sigma)=[\ell(\sigma)]$.
Set $l_0=\gamma(c(\sigma))$.  In the question's notation
$w_j=r_j\varepsilon_j+s_j\widehat\delta$, the integer is
$p=mn l_0-nr_1-mr_2$.
The section-preserving source data of \ref{push:src:E.1.1} have $c(\sigma)=0$.

\begin{theorem}[classification in Viaclovsky's coordinates]\hypertarget{tag-032S}{}\label{push:src:E.1.2}
\StatementRecord{push:src:E.1.2}{proved}{032S}{\SeriesRef{push:src:E.1.1}}

Use a local logarithm $c$ of the transported linearization,
$u=e^{-2\pi i c}$, and the same reference section and cusp comparison
on both sides.  Two sphere members with data
$(m,\rho_0,N_0;n,\rho_1,N_1;c)$ and primed data are biholomorphic
if and only if their multiplicities agree and there are
$\alpha\in\{1,-1\}$ and $h=c'-c\in\mathbb Z$ such that
\[
 b'=\alpha b,\qquad
 hN'_0\equiv\rho'_0-\alpha\rho_0\pmod{k},\qquad
 hN'_1\equiv\rho'_1-\alpha\rho_1\pmod{2l}.
\]
Every such biholomorphism covers the identity of $B$.
For fixed discrete data the period identification subgroup is
$2kl\mathbb Z$.  At fixed $u$ all sphere translation data give
exactly $2kl$ classes.  A positive $u$-circuit acts on them by a
single cycle.
\end{theorem}


% TARGET PROOF MASKED: push:src:E.1.2

















































































































































The preceding theorem answers Question 3 of
\cite[Section 1.1]{ViaclovskyLog2026}.  The original return asserted the
same arithmetic classification in \cref{fam:src:B.3.1}; the contribution here is
the explicit dictionary and the complete enlarged-root and collar
argument.

\begin{corollary}[holomorphic dependence of the comparisons]\hypertarget{tag-032T}{}\label{push:src:E.1.3}
\StatementRecord{push:src:E.1.3}{proved}{032T}{\SeriesRef{push:src:E.1.2}, \SeriesRef{fam:src:B.4.1}}

For fixed discrete data in \ref{push:src:E.1.2} and a fixed admissible shift $h$,
the constructed comparisons are holomorphic isomorphisms of families
over a sufficiently small parameter disc.  In particular, for sphere
data with multiplicities $(3,4)$ the identification with the standard
family uses the parameter
\[
 W=e^{-\pi i(c+s_2)}=(-1)^{s_2}w.
\]
The smooth monodromy conjugacy class around the positive parameter
circle is unchanged by this identification of twist data.
\end{corollary}


% ALL PRINTED PROOF MATERIAL MASKED

























\section{The monodromy group and its period domain}\hypertarget{tag-01PV}{}\label{fam:section:F}

We use the clockwise root deck generators and the lattice basis specified
on the first page.  In this section a group acts properly discontinuously when, for every
compact subset $K$ of its space, only finitely many group elements $g$ satisfy
$gK\cap K\ne\varnothing$.  This definition permits finite stabilizers.  We
write $e=\widehat\delta$, $E=e\otimes\gamma$, and
\[
 J_2=\begin{pmatrix}0&1\\-1&0\end{pmatrix},\qquad
 \Omega=\begin{pmatrix}
 0&0&0&6\\0&0&1&0\\0&-1&0&0\\-6&0&0&0
 \end{pmatrix}.
\]
Thus the alternating form is
$\psi=u\wedge w+6\gamma\wedge\delta$, an integral alternating form of
type $(1,6)$ on $\Lambda$.  We use $\Omega$ for its matrix and retain $J$ for
a real complex structure.  The domain $\mathcal D$ consists of the real
endomorphisms $J$ of $\Lambda_{\mathbb R}$ satisfying
$J^2=-I$, $J^{\mathsf t}\Omega J=\Omega$, and
$\operatorname{sign}(\Omega J)=(2,2)$.  It is the stated
$\operatorname{Sp}_4(\mathbb R)/U(1,1)$, with its complex structure as an
open subset of the Lagrangian Grassmannian.  An integral symplectic matrix
acts by $J\mapsto gJg^{-1}$.

\subsection{An exact arithmetic description}\hypertarget{tag-01PW}{}\label{fam:section:F.1}

\begin{theorem}\hypertarget{tag-01R9}{}\label{fam:src:F.1.1}\label{fam:st:F.1.1}\label{cdp:st:D.1.14}
\StatementRecord{fam:src:F.1.1}{proved}{01R9}{\SeriesRef{fam:eq:F-law}, \SeriesRef{fam:eq:F-SLpresentation}}

Let $P_{e,\mathbb Z}$ be the integral symplectic group fixing $e$ as a
vector, and let $P_{\mathbb Z}$ be the integral symplectic group preserving
$\mathbb Ze\subset\ker\gamma\subset\Lambda$.  Every element of
$P_{e,\mathbb Z}$ has a unique expression
\begin{equation}\hypertarget{tag-01OT}{}\label{fam:eq:F-g}
 g(B,x,t)=
 \begin{pmatrix}
 1&0&0\\6x&B&0\\t&-x^{\mathsf t}J_2B&1
 \end{pmatrix},\qquad
 B\in\operatorname{SL}_2(\mathbb Z),\quad x\in\mathbb Z^2,\quad t\in\mathbb Z,
\end{equation}
where the block sizes are $1,2,1$.  Put
\[
 S=\begin{pmatrix}0&-1\\1&0\end{pmatrix},\quad
 T=\begin{pmatrix}1&1\\0&1\end{pmatrix},\quad
 a=\begin{pmatrix}0&1\\-1&-1\end{pmatrix}.
\]
There is a homomorphism
$\alpha:\operatorname{SL}_2(\mathbb Z)\to\mathbb Z/12$ determined by
$\alpha(S)=9$ and $\alpha(T)=1$; in particular $\alpha(a)=4$.
With $Q(x)=x_1+x_2+x_1x_2$, the actual monodromy group is
\begin{equation}\hypertarget{tag-01OQ}{}\label{fam:eq:F-Gamma}
 \Gamma=\left\{g(B,x,t):t+6Q(x)-\alpha(B)=0\ \text{in }\mathbb Z/12\right\}.
\end{equation}
Consequently $[P_{e,\mathbb Z}:\Gamma]=12$ and
$P_{\mathbb Z}=\{\pm I\}\times P_{e,\mathbb Z}$, so
$[P_{\mathbb Z}:\Gamma]=24$.  Its image in the Levi quotient is exactly
$\operatorname{SL}_2(\mathbb Z)$, with scalar on $e$ equal to $1$.
\end{theorem}



% ALL PRINTED PROOF MATERIAL MASKED





































































































\begin{proposition}\hypertarget{tag-01RA}{}\label{fam:src:F.1.2}\label{fam:st:F.1.2}
\StatementRecord{fam:src:F.1.2}{proved}{01RA}{\SeriesRef{fam:eq:F-Gamma}, \SeriesRef{fam:eq:F-law}}

The group in \eqref{fam:eq:F-Gamma} has center
$\langle Z\rangle\cong\mathbb Z$ and is the following explicit central
extension of $\mathbb Z^2\rtimes\operatorname{SL}_2(\mathbb Z)$ by
$\mathbb Z$.  Choose the representative $a(B)\in\{0,\ldots,11\}$ of
$\alpha(B)$.  Its underlying set consists of triples $(B,x,k)$, with
product
\[
 (B,x,k)(C,y,l)=(BC,x+By,k+l+\kappa((B,x),(C,y))),
\]
where
\begin{align*}
12\kappa((B,x),(C,y))={}&a(B)+a(C)-a(BC)-6Q(x)-6Q(y)\\
 &-6x^{\mathsf t}J_2By+6Q(x+By).
\end{align*}
Its real and rational Zariski closures are respectively the full real
and rational symplectic stabilizers of the vector $e$.  They have
dimension $6$, with three-dimensional Heisenberg unipotent radical and
Levi quotient $\operatorname{SL}_2$.
\end{proposition}



% ALL PRINTED PROOF MATERIAL MASKED
































\subsection{Properness and the maximal open set}\hypertarget{tag-01PX}{}\label{fam:section:F.2}

Define the real-analytic function
\[
 q_e:\mathcal D\longrightarrow\mathbb R,\qquad q_e(J)=\psi(e,Je),
 \qquad\mathcal O=\{J:q_e(J)\ne0\}.
\]
The function $q_e$ is invariant under $P_{\mathbb Z}$: replacing $e$ by
$-e$ changes both of its arguments.  Write $\mathcal O_+$ and
$\mathcal O_-$ for its two sign regions.

\begin{proposition}[counterexample]\hypertarget{tag-01RB}{}\label{fam:src:F.2.1}\label{fam:st:F.2.1}
\StatementRecord{fam:src:F.2.1}{proved}{01RB}{\SeriesRef{fam:eq:F-Gamma}}

The assertion that $\Gamma$ acts properly discontinuously on all of
$\mathcal D$ is false.  The following point of $\mathcal D$ has an
explicit infinite stabilizer:
\[
 J_0=\begin{pmatrix}
 0&0&-1&0\\0&0&0&-6\\1&0&0&0\\0&1/6&0&0
 \end{pmatrix},\qquad
 R=\begin{pmatrix}
 1&0&0&0\\0&1&72&0\\0&0&1&0\\-12&0&0&1
 \end{pmatrix}\in\Gamma.
\]
Every power of $R$ fixes $J_0$, and all these powers are distinct.
\end{proposition}



% ALL PRINTED PROOF MATERIAL MASKED























\begin{theorem}\hypertarget{tag-01RC}{}\label{fam:src:F.2.2}\label{fam:st:F.2.2}
\StatementRecord{fam:src:F.2.2}{proved}{01RC}{\SeriesRef{fam:eq:F-g}}

The set $\mathcal O$ is the largest $\Gamma$-invariant open subset of
$\mathcal D$ on which the action is properly discontinuous.  It is also
the largest such open subset for $P_{\mathbb Z}$.  Each sign region is
a complex manifold in the analytic topology, and its quotient by either
discrete group is a Hausdorff complex analytic orbifold.
\end{theorem}



% ALL PRINTED PROOF MATERIAL MASKED













































































\begin{proposition}\hypertarget{tag-01RD}{}\label{fam:src:F.2.3}\label{fam:st:F.2.3}
\StatementRecord{fam:src:F.2.3}{proved}{01RD}{the definitions and constructions stated in the proof}

The positive sign region admits the global period chart
\[
 \mathcal O_+=\left\{
 (\tau,\mu,\beta)\in\mathbb C^3:
 \operatorname{Im}\tau>0,\quad
 \operatorname{Im}\beta-\frac{6(\operatorname{Im}\mu)^2}
 {\operatorname{Im}\tau}<0\right\},
\]
represented by the matrix $\Pi$ in the stated problem.  In this chart,
\begin{equation}\hypertarget{tag-01OW}{}\label{fam:eq:F-qD}
 q_e(J)=-\frac6D.
\end{equation}
Every admissible period $\Pi_c(z)$ belongs to $\mathcal O_+$, for every
admissible $c$.  Thus this single maximal proper-action open set contains
all the period images in question.
\end{proposition}



% ALL PRINTED PROOF MATERIAL MASKED






































\subsection{The curve in the proper quotient}\hypertarget{tag-01PY}{}\label{fam:section:F.3}

\begin{theorem}\hypertarget{tag-01RE}{}\label{fam:src:F.3.1}\label{fam:st:F.3.1}
\StatementRecord{fam:src:F.3.1}{proved}{01RE}{\SeriesRef{fam:eq:F-j}, \SeriesRef{fam:src:F.2.2}}

For every admissible $c$, the marked period map on the covering of
$B^\circ$ is a holomorphic immersion into $\mathcal O_+$.  The induced
map
\[
 p_c:B^\circ\longrightarrow\Gamma\backslash\mathcal O_+
\]
is injective and a holomorphic immersion.  It extends over $t=0,1$ to a
proper holomorphic embedding
\[
 \bar p_c:\mathbb C=B\setminus\{\infty\}
 \lhook\joinrel\longrightarrow\Gamma\backslash\mathcal O_+.
\]
The closure of $p_c(B^\circ)$ in this quotient is exactly
$\bar p_c(\mathbb C)$.
\end{theorem}



% ALL PRINTED PROOF MATERIAL MASKED
































































\begin{proposition}\hypertarget{tag-01RF}{}\label{fam:src:F.3.2}\label{fam:st:F.3.2}
\StatementRecord{fam:src:F.3.2}{proved}{01RF}{\SeriesRef{fam:src:F.3.1}}

At the two finite punctures the period map has the holomorphic
finite-cover limits
\[
 \Pi_1=
 \begin{pmatrix}2(2-\rho)&\rho&1&0\\
 \beta_c(z_1)&(2-\rho)/3&0&1\end{pmatrix},\qquad
 \Pi_2=
 \begin{pmatrix}3(1-i)&i&1&0\\
 \beta_c(z_2)&(1-i)/2&0&1\end{pmatrix},
 \quad \rho=e^{\pi i/3}.
\]
Their monodromies are the given order-three and order-four matrices.
At the cusp the curve leaves every compact subset of
$\Gamma\backslash\mathcal O_+$.  On the prescribed cusp branch its
Hodge plane in the compact dual tends to
$\mathbb C\gamma+\mathbb Cu$.  The latter plane has complex dimension
two in common with its conjugate and therefore does not represent a
point of $\mathcal D$.
\end{proposition}



% ALL PRINTED PROOF MATERIAL MASKED
































\begin{proposition}\hypertarget{tag-01RG}{}\label{fam:src:F.3.3}\label{fam:st:F.3.3}
\StatementRecord{fam:src:F.3.3}{proved}{01RG}{\SeriesRef{fam:src:F.3.1}, \SeriesRef{fam:eq:F-j}}

The determinant used in the workbench is generically nonzero, with
the stated zero at $z_1$:
\[
 \det\dot Z=6\dot\mu^2-\dot\tau\dot\beta,\qquad
 \det\dot Z(z_1)=0.
\]
The immersion established in \ref{fam:src:F.3.1} holds at every point of $B^\circ$
without requiring this determinant to be everywhere nonzero.
\end{proposition}



% ALL PRINTED PROOF MATERIAL MASKED


























\section{Exceptional fibres, Neron--Severi classes, and compact subspaces}\hypertarget{tag-01PP}{}\label{fam:section:D}
\subsection{Elliptic quotients and the torsion locus}\hypertarget{tag-01PQ}{}\label{fam:section:D.1}
\begin{proposition}\hypertarget{tag-01QR}{}\label{fam:src:D.1.1}\label{fam:st:D.1.1}
\StatementRecord{fam:src:D.1.1}{proved}{01QR}{the definitions and constructions stated in the proof}

The specified linear projection \(\operatorname{pr}_1:\mathbb C^2\to
\mathbb C\) induces a holomorphic group homomorphism
\(F_t\to E_t=\mathbb C/(\mathbb Z\tau+\mathbb Z)\) exactly on
\(6\mu=A\tau+B\), \(A,B\in\mathbb Z\). It is then surjective with
connected kernel
\(\mathbb C/(\mathbb Z+(A\mu-\beta)\mathbb Z)\).
On the smooth base the two values of \(t\) are precisely
\[
 t_\pm=\frac{64(47\pm27\sqrt3)}{11}.
\]
\end{proposition}


% ALL PRINTED PROOF MATERIAL MASKED















































\begin{proposition}\hypertarget{tag-01QS}{}\label{fam:src:D.1.2}\label{fam:st:D.1.2}
\StatementRecord{fam:src:D.1.2}{proved}{01QS}{\SeriesRef{fam:src:C.2.1}, \SeriesRef{fam:src:D.1.1}; \cite[Lemma 2.1]{Engel2026}; \cite[Theorem 4.4, p.~3]{Schnell19}}

The locus on \(B^\circ\) where \(P(t)\) is torsion is countable and dense.
Every fibre over this locus is an abelian surface. The divisor on the
smooth base cut out by \(6kP=O\) has degree \(3k^2-1\), including local
intersection multiplicity, for every positive integer \(k\).
\end{proposition}

% ALL PRINTED PROOF MATERIAL MASKED









































































\Uses{the hypotheses, notation and definitions stated here; \cref{app:rederive:fam-src-D.1.2}}


This specializes the elliptic-surface height and specialization results cited in the re-derivation. A complete re-derivation is given in Appendix~\ref{app:rederive:fam-src-D.1.2}.


\begin{theorem}[counterexample]\hypertarget{tag-01QT}{}\label{fam:src:D.1.3}\label{fam:st:D.1.3}
\StatementRecord{fam:src:D.1.3}{proved}{01QT}{\SeriesRef{fam:src:D.1.1}, \SeriesRef{fam:src:C.2.1}}

The proposed assertion that \emph{existence of a surjection onto the
specified \(E_t\) with elliptic kernel} is equivalent to \(6P(t)=O\)
is false. It remains true for the specified projection in \ref{fam:src:D.1.1}.
\end{theorem}


% ALL PRINTED PROOF MATERIAL MASKED









































\begin{proposition}\hypertarget{tag-01QU}{}\label{fam:src:D.1.4}\label{fam:st:D.1.4}
\StatementRecord{fam:src:D.1.4}{proved}{01QU}{\SeriesRef{fam:src:C.3.1}, \SeriesRef{fam:src:D.1.6}}

For every smooth fibre and every admissible \(c\), its Neron--Severi lattice
is the following explicit subgroup of \(\mathbb Z^6\):
\begin{equation}\hypertarget{tag-01P4}{}\label{fam:eq:cd-NS}
 \left\{E:\ E_{12}-E_{13}\tau+(6E_{23}-E_{14})\mu+
 E_{24}\beta+E_{34}(6\mu^2-\tau\beta)=0\right\}.
\end{equation}
The complex orientation has intersection form
\(E^2=-2\operatorname{Pf}(E)\), where
\(\operatorname{Pf}(E)=E_{12}E_{34}-E_{13}E_{24}+E_{14}E_{23}\).
In particular the invariant class \(\eta=6\gamma\wedge\delta+u\wedge w\)
has \(\int_{F_t}\eta^2=-12\). Its Hermitian form is indefinite.
The fibres with lattice exactly \(\mathbb Z\eta\) have no compact curves
and algebraic dimension zero.
\end{proposition}


% ALL PRINTED PROOF MATERIAL MASKED









































\begin{proposition}\hypertarget{tag-01QV}{}\label{fam:src:D.1.5}\label{fam:st:D.1.5}
\StatementRecord{fam:src:D.1.5}{proved}{01QV}{\SeriesRef{fam:eq:cd-NS}}

For a fixed \(c\), all but countably many smooth fibres have
\(\operatorname{NS}=\mathbb Z\eta\). Conversely, for every fixed smooth
\(t\), there are admissible parameters tending to the cusp of the
parameter disc for which \(F_t\) is an abelian surface, irrespective of
the order of \(P(t)\).
\end{proposition}


% ALL PRINTED PROOF MATERIAL MASKED





































\begin{theorem}\hypertarget{tag-01QW}{}\label{fam:src:D.1.6}\label{fam:st:D.1.6}
\StatementRecord{fam:src:D.1.6}{proved}{01QW}{\SeriesRef{fam:eq:cd-NS}, \SeriesRef{fam:src:D.1.4}; \cite[Theorem 4.1, p.~1, and Theorem 4.4, p.~3]{Schnell19}}

Fix any admissible $c$. For each integral alternating matrix $E$, let
$Z_E(c)$ be the zero set on the marking cover of the explicit holomorphic
function in \eqref{fam:eq:cd-NS}, with $\beta=\beta^{\rm part}+c$ given by
\eqref{fam:eq:cd-beta}. Let $\pi$ be the cover map to $B^\circ$, and define
\[
 \mathcal E_c=\pi\!\left(\bigcup_{E\ne0,\ \operatorname{Pf}(E)=0}Z_E(c)\right),
 \qquad
 \mathcal A_c=\pi\!\left(\bigcup_{\operatorname{Pf}(E)<0}Z_E(c)\right).
\]
These give the exact loci of all elliptic quotients and all abelian
surfaces, respectively. The algebraic dimension of every smooth fibre is
\[
 a(F_t)=
 \begin{cases}
 2,&t\in\mathcal A_c,\\
 1,&t\in\mathcal E_c\setminus\mathcal A_c,\\
 0,&t\notin\mathcal E_c\cup\mathcal A_c.
 \end{cases}
\]
Together with \eqref{fam:eq:cd-NS}, these are an exact classification for every
$t$, by countable unions of explicitly given analytic zero sets; they do
not assert that the loci have a finite algebraic description. They include
all quotient maps, without requiring preservation of the first coordinate.
\end{theorem}

% ALL PRINTED PROOF MATERIAL MASKED








































































\Uses{the hypotheses, notation and definitions stated here; \cref{fam:eq:cd-NS}; \cref{fam:eq:cd-beta}; \cref{app:rederive:fam-src-D.1.6}}


This is the Appell--Humbert and Riemann criterion in the present lattice. A complete re-derivation is given in Appendix~\ref{app:rederive:fam-src-D.1.6}.


\begin{proposition}\hypertarget{tag-01QX}{}\label{fam:src:D.1.7}\label{fam:st:D.1.7}
\StatementRecord{fam:src:D.1.7}{proved}{01QX}{\SeriesRef{fam:eq:cd-NS}, \SeriesRef{fam:src:D.1.6}, \SeriesRef{fam:src:C.2.1}, \SeriesRef{fam:src:D.1.2}}

There is an explicit nonempty locus of admissible members with ordinary
fibres of algebraic dimension one. On a sufficiently small punctured
cusp neighbourhood set
\[
 c(z)=-6\tau(z)-12\mu(z)-\beta^{\rm part}(z).
\]
Outside a countable set of $z$, the fibre of $X_{c(z)}$ over $\pi(z)$ has
Neron--Severi lattice $\mathbb Z\eta+\mathbb ZE$, where the nonzero upper
entries of $E$ are
$E_{13}=-6$, $E_{14}=-6$, $E_{23}=1$, $E_{24}=1$.
Its algebraic dimension is one. The section $P(\pi(z))$ can simultaneously
be required to be non-torsion.
\end{proposition}


% ALL PRINTED PROOF MATERIAL MASKED





























\subsection{Compact curves and surfaces}\hypertarget{tag-01PR}{}\label{fam:section:D.2}
\begin{proposition}\hypertarget{tag-01QY}{}\label{fam:src:D.2.1}\label{fam:st:D.2.1}
\StatementRecord{fam:src:D.2.1}{proved}{01QY}{the definitions and constructions stated in the proof}

For a standard-twist member \(X_c\) with underlying manifold \(S^6\),
every compact irreducible complex curve lies in one fibre of \(f\), and
every compact irreducible complex surface is one of the reduced fibres.
\end{proposition}


% ALL PRINTED PROOF MATERIAL MASKED

























\begin{proposition}\hypertarget{tag-01QZ}{}\label{fam:src:D.2.2}\label{fam:st:D.2.2}
\StatementRecord{fam:src:D.2.2}{proved}{01QZ}{\SeriesRef{fam:src:D.2.1}, \SeriesRef{fam:src:D.1.1}, \SeriesRef{fam:src:C.4.1}}

There is no rational curve in \(X_c\setminus W_c\). There are elliptic
curves outside all three special fibres, in particular in each of the
two smooth fibres \(F_{t_\pm}\). All curves in \(W_c\) are precisely the
images of irreducible algebraic curves in \(\mathrm{dP}_6\) under its
finite normalization map. Distinct curves have the same image only for
the three paired boundary lines. All curves in a bielliptic fibre
\(S_j\) are images of irreducible curves in its finite product cover.
\end{proposition}


% ALL PRINTED PROOF MATERIAL MASKED
































\begin{theorem}\hypertarget{tag-01R0}{}\label{fam:src:D.2.3}\label{fam:st:D.2.3}
\StatementRecord{fam:src:D.2.3}{proved}{01R0}{\SeriesRef{fam:src:D.1.6}, \SeriesRef{fam:src:D.2.1}, \SeriesRef{fam:src:D.2.2}; \cite[Theorem 4.1 and equation (4.3), pp.~1--2]{Schnell19}}

The following descriptions give all compact irreducible curves of $X_c$.
In $W_c$ and $S_j$ they are exactly the images described in \ref{fam:src:D.2.2}.
In a smooth fibre of algebraic dimension zero there are none. In a
smooth fibre of algebraic dimension one they are exactly the translates
of its unique elliptic kernel $K_t$. In an abelian fibre they are exactly
the translates of elliptic subtori, together with the irreducible zero
divisors of the positive Appell--Humbert theta sections described below.
These descriptions, \ref{fam:src:D.1.6}, and \ref{fam:src:D.2.1} constitute a classification including
every smooth exceptional fibre and every compact irreducible surface.
\end{theorem}

% ALL PRINTED PROOF MATERIAL MASKED






















































\Uses{the hypotheses, notation and definitions stated here; \cref{fam:src:D.2.2}; \cref{fam:src:D.1.6}; \cref{fam:src:D.2.1}; \cref{app:rederive:fam-src-D.2.3}}


This is the effective line-bundle criterion for compact complex tori in the present lattice. A complete re-derivation is given in Appendix~\ref{app:rederive:fam-src-D.2.3}.




\subsection{The rank-three period quotient}\hypertarget{tag-01PN}{}\label{fam:section:C.3}
\begin{proposition}\hypertarget{tag-01QM}{}\label{fam:src:C.3.1}\label{fam:st:C.3.1}
\StatementRecord{fam:src:C.3.1}{proved}{01QM}{the definitions and constructions stated in the proof}

At every point with \(D=0\), \(T>0\), the real-linear map
\(\Pi:\mathbb R^4\to\mathbb C^2\) has rank three. Write
\(\mu=a_0\tau+b_0\), \(a_0,b_0\in\mathbb R\), and put
\[
 k=(1,-6a_0,-6b_0,6a_0\operatorname{Re}\mu-\operatorname{Re}\beta).
\]
Its real kernel is \(\mathbb R k\). The subgroup \(\Pi\mathbb Z^4\) is
discrete exactly when \(k\in\mathbb Q^4\). In that case it is a
rank-three lattice in a real hyperplane and its quotient in \(\mathbb C^2\)
is a noncompact complex Lie group. In every other case the topological
quotient \(\mathbb C^2/\Pi\mathbb Z^4\) is non-Hausdorff.
\end{proposition}


% ALL PRINTED PROOF MATERIAL MASKED





























\begin{proposition}\hypertarget{tag-01QN}{}\label{fam:src:C.3.2}\label{fam:st:C.3.2}
\StatementRecord{fam:src:C.3.2}{proved}{01QN}{\SeriesRef{fam:src:C.1.1}, \SeriesRef{fam:src:C.2.3}}

The unique degenerate base point at the actual admissibility boundary is
\(t=1\). The boundary constants
\(c=x+i(3/\pi)\log6\) give
\[
 \ker_{\mathbb R}\Pi=\mathbb R(1,3,-3,-x-3/2).
\]
For rational $x+3/2=m/n$ in lowest terms, the quotient is a
finite holomorphic quotient of $E_i\times\mathbb C^*$, with a covering
of degree four. For irrational \(x+3/2\), the closure of the period subgroup is generated
by the real line \(\mathbb R(0,1)\) and the two periods
\((i,(1-i)/2)\), \((1,0)\). Its Hausdorff quotient is a real manifold
homeomorphic to \(\mathbb R\times(S^1)^2\).
\end{proposition}


% ALL PRINTED PROOF MATERIAL MASKED




























\paragraph{Local boundary obstruction}\hypertarget{tag-01QO}{}\label{fam:src:C.3.3}\label{fam:st:C.3.3}
\StatementRecord{fam:src:C.3.3}{proved}{01QO}{\SeriesRef{fam:src:C.3.2}, \SeriesRef{fam:src:C.3.4}, \SeriesRef{push:src:D.1.2}, \SeriesRef{push:src:D.1.3}}

The local compactification obstruction at the actual boundary is localized at the order-four fibre over $t=1$. The proposed extension would require a compact complex-space replacement of the rank-three
quotient in \ref{fam:src:C.3.2}, over a neighbourhood of one boundary parameter,
with a proof that it extends the original cyclic filling and its overlap
with the ordinary torus family. Both rational and irrational values of
$x+3/2$ are included in the exact quotient description. The naive
quotients cannot be central fibres of a proper complex-space family:
for rational $x+3/2$ they are noncompact, and for irrational values
they are non-Hausdorff. \ref{fam:src:C.3.4} supplies a further obstruction to preserving
a holomorphic Weierstrass Albanese model in this parameter.
These prove failure of the stated specific replacements. The unique base-fibre location and exact local period obstruction are thereby determined. The general proper holomorphic extension is excluded by \cref{push:src:D.1.2}, \cref{push:src:D.1.3}, which recovers the intrinsic divisor in a hypothetical extension.





\begin{proposition}\hypertarget{tag-01QP}{}\label{fam:src:C.3.4}\label{fam:st:C.3.4}
\StatementRecord{fam:src:C.3.4}{proved}{01QP}{\SeriesRef{fam:section:A.4}, \SeriesRef{fam:src:C.1.1}}

On its half-plane of definition, the $j$-invariant of
$\operatorname{Alb}(S_2)$ has no meromorphic continuation across any
point of the line
$\operatorname{Im}c=(3/\pi)\log6$. In particular the rank-three quotient
at $t=1$ cannot be replaced by a Weierstrass family of elliptic Albanese
curves whose coefficients are holomorphic in the same parameter across
that line.
\end{proposition}


% ALL PRINTED PROOF MATERIAL MASKED
























\subsection{Classification of the non-normal cusp surface}\hypertarget{tag-01PO}{}\label{fam:section:C.4}
\begin{theorem}\hypertarget{tag-01QQ}{}\label{fam:src:C.4.1}\label{fam:st:C.4.1}
\StatementRecord{fam:src:C.4.1}{proved}{01QQ}{\SeriesRef{fam:src:C.1.1}}

For the normalization of \ref{fam:src:C.1.1}, the three opposite-side identification
constants of \(W_c\), in the characters specified below, are
\[
 (g_1,g_2,g_3)=(-a^{-1},a^6,-a^{-7}e^{-2\pi ic}).
\]
Two such compact complex surfaces are biholomorphic exactly when
\(c'-c\in\mathbb Z\). The vertical \(\mathbb C^*\)-action inherited from
translation in the \(\widehat\delta\)-direction fixes exactly the double
curve represented by the pair of edges of direction \(e_2\).
\end{theorem}


% ALL PRINTED PROOF MATERIAL MASKED












































































\section{Real structures and their fixed three-spheres}\hypertarget{tag-02P9}{}\label{mapsI:real-forms}

We use the family marking
$\Lambda=\mathbb Z\langle\widehat\gamma,\widehat u,\widehat w,
\widehat\delta\rangle$ and its dual covector $\gamma$.
In this section $d=\widehat\delta$, $g=\widehat\gamma$,
$u=\widehat u$, and $w_{\Lambda}=\widehat w$; an unadorned $w$ in a
lattice expression denotes $w_{\Lambda}$, whereas the scalar parameter is
$w=\exp(-\pi i c)$. Put
\[
 e=\Pi_c d=(0,1)^t,\qquad E=d\otimes\gamma,
 \qquad \Pi_c=
 \begin{pmatrix}6\mu&\tau&1&0\\\beta^{\rm part}+c&\mu&0&1\end{pmatrix},
\]
\[
 D_c=\operatorname{Im}(\beta^{\rm part}+c)
       -6(\operatorname{Im}\mu)^2/\operatorname{Im}\tau<0.
\]
Thus $d$ is a lattice vector and $e$ is its image in the complex translation
space. The circle $\rho(s)$ is translation by $se$, $s\in\mathbb R/\mathbb Z$.
The clockwise root generators and the affine vectors are
\[
 A_1=\begin{pmatrix}1&0&0&0\\6&0&1&0\\-6&-1&-1&0\\-2&1&0&1\end{pmatrix},\quad
 A_2=\begin{pmatrix}1&0&0&0\\0&0&-1&0\\-6&1&0&0\\3&0&1&1\end{pmatrix},
\]
\[
 M_0=(A_1A_2)^{-1},\qquad
 v_1=(1,2,-4,0)^t,\qquad v_2=(-1,-3,3,0)^t.
\]
These are the same conventions as in the family construction. In cusp
coordinates $q=1/t$ and $s=(2\pi i)^{-1}\log q$, the bounded matrix $C_c(q)$
is defined by $Z_c=sB_0+C_c(q)$, with
$B_0=\left(\begin{smallmatrix}0&1\\-1&0\end{smallmatrix}\right)$.
In particular $\Pi_cg=(6\mu,-s+C_{c,21})^t$.
We write $D_1,D_2,D_{12}$ for the three cusp double curves belonging to
edges in directions $e_1,e_2,e_1-e_2$, respectively. Thus $D_2$ is the
fixed sphere of the circle action $\rho$, whereas $D_1$ is the double
curve along which the literal character has a limit. This distinction
is essential in the real calculation as well.

\begin{proposition}[Two conjugacy classes of real structures]\hypertarget{tag-02PE}{}\label{mapsI:src:A.3.1}
\StatementRecord{mapsI:src:A.3.1}{proved}{02PE}{\SeriesRef{fam:src:B.4.2}, \SeriesRef{fam:src:B.4.1}; \cite[Theorem~3.1, p.~48]{Illman1972}; \cite[Theorem 4.1(1)]{Knudson2004}}

For a real member the antiholomorphic involutions have exactly two
conjugacy classes under holomorphic automorphisms. For any fixed
involutive splitting $\sigma$ they are represented by $\sigma$ and
$\rho(1/2)\sigma$. Every nonempty fixed set is a closed connected
orientable totally real three-manifold with the mod-two homology of
$S^3$. The circle $\rho$ preserves each such fixed set and acts
semifreely on it.
\end{proposition}



% ALL PRINTED PROOF MATERIAL MASKED
































We now fix a representative of each sign class, instead of leaving the
symbol $\sigma$ unspecified. Put
\[
 L_F=(2\pi i)^{-1}\log F,\qquad L_G=(2\pi i)^{-1}\log G,
 \qquad F^3=t,\quad G^4=t-1.
 \hypertarget{tag-02P4}{}\label{mapsI:real-eq-1}
\]
The roots are positive over $t>1$. We continue from that interval
through the upper half $t$-plane. Their logarithms are taken on the
orbifold uniformizing cover. Write $\iota$ for its reflection lift.
Every point of $-i\R_{>0}$ has a representative $c=1/2+iy$, and every
point of $i\R_{>0}$ has a representative $c=-1/2+iy$, with $y$ in the
admissible range. Set
\[
\begin{array}{c|c|c}
 &\alpha&S(t)\\\hline
 c=1/2+iy&1/3&2L_F-4L_G-1/6\\
 c=-1/2+iy&1/2&-2L_G+1/4
\end{array}
\qquad b_c(z)=\alpha\Pi_c(z)g+S(t(z))e.
\label{mapsI:real-eq-2}
\]

\begin{theorem}[Explicit real involutions and all fixed fibres]\hypertarget{tag-02PF}{}\label{mapsI:src:A.3.2}
\StatementRecord{mapsI:src:A.3.2}{proved}{02PF}{\SeriesRef{fam:src:B.4.2}, \SeriesRef{mapsI:real-eq-4}, \SeriesRef{mapsI:real-eq-2}, \SeriesRef{fam:src:A.3.1}, \SeriesRef{mapsI:real-eq-1}}

On the ordinary period charts the formulas
\[
 \sigma_c[z,\zeta]=[\iota z,\bar\zeta+b_c(\iota z)],
 \qquad \sigma_c^- =\rho(1/2)\sigma_c
 \label{mapsI:real-eq-3}
\]
extend to antiholomorphic involutions of the compact threefold.
With these representatives their complete fixed-set data are as follows.
Here $2T^2$ means the disjoint union of two real two-tori; every displayed
circle in a double curve avoids its two triple-point endpoints.
\begin{center}\small
\begin{tabular}{cclccc}
\toprule
$w$&involution&image in $\RP^1$&on $S_1$&on $S_2$&on $W$\\
\midrule
$-ir$, $r>0$&$\sigma_c$&all $\RP^1$&$2T^2$&$2T^2$&$T^2\sqcup S^1_{D_2}\sqcup S^1_{D_{12}}$\\
$-ir$, $r>0$&$\sigma_c^-$&$[1,\infty]$&$\varnothing$&$T^2$&$S^1_{D_1}\sqcup S^1_{D_2}$\\
$ir$, $r>0$&$\sigma_c$&$\varnothing$&$\varnothing$&$\varnothing$&$\varnothing$\\
$ir$, $r>0$&$\sigma_c^-$&$[\infty,1]$ through $0$&$2T^2$&$T^2$&$S^1_{D_1}\sqcup S^1_{D_{12}}$\\
\bottomrule
\end{tabular}
\end{center}
The first interval is the closure of $t>1$; the second is the closure
of $t<1$. On every open real arc contained in the indicated image the
real regular fibre is $2T^2$. No triple point is real for these structures.
\end{theorem}



% ALL PRINTED PROOF MATERIAL MASKED






















































































































































































\begin{theorem}[The real three-manifolds]\hypertarget{tag-02PG}{}\label{mapsI:src:A.3.3}
\StatementRecord{mapsI:src:A.3.3}{proved}{02PG}{\SeriesRef{mapsI:src:A.3.2}}

Every nonempty fixed set in Theorem~\ref{mapsI:src:A.3.2} is diffeomorphic to $S^3$.
Its fundamental group is trivial and its integral homology is
$(\Z,0,0,\Z)$. For negative imaginary $w$, the restricted circle action
is the standard semifree action on $S^3$ with fixed circle
$D_2^{\sigma_c}$. For positive imaginary $w$, the nonempty real form
has a free circle action isomorphic, up to circle and orientation
conventions, to the Hopf action.
\end{theorem}



% ALL PRINTED PROOF MATERIAL MASKED


















































\begin{proposition}[Restrictions of the repaired circle quotient]\hypertarget{tag-02PH}{}\label{mapsI:src:A.3.4}
\StatementRecord{mapsI:src:A.3.4}{proved}{02PH}{\SeriesRef{mapsI:src:A.2.1}, \SeriesRef{mapsI:real-eq-2}}

The repaired map $e_{\rm q}$ of Theorem~\ref{mapsI:src:A.2.1} may be chosen $\rho$-invariant. Its restriction to
every nonempty real form is null-homotopic. On an ordinary real fibre
its image consists of the two quotient-circle values
\[
 \theta=\tfrac12\gamma(r_c(t)),\qquad
 \theta=\tfrac12\gamma(r_c(t))+\tfrac12\pmod\Z,
 \quad
 \gamma(r_c(t))=\alpha+\frac{\Im S(t)}{D_c(t)},
 \label{mapsI:real-eq-9}
\]
with $\alpha,S$ from \eqref{mapsI:real-eq-2}. The circle action is along each real
two-torus and is annihilated by $e_{\rm q}$.
\end{proposition}



% ALL PRINTED PROOF MATERIAL MASKED























The formulas specify the embedded fixed sets by their torus cosets and
their actual finite and toric attachments. No ambient-isotopy
classification or value of the Haefliger knot invariant of these
embedded three-spheres is asserted here.


\begin{cited}[The coherent closing translation]\hypertarget{tag-0322}{}\label{push:src:C.3.5}
\StatementRecord{push:src:C.3.5}{proved}{0322}{\SeriesRef{fam:src:A.3.1}; \citet[Proposition~7.8, pp.~91--92, and Proposition~C.2, p.~102]{DuEtAl2026v1}}

Fix the standard twists and the root and cusp markings of Proposition~\ref{fam:src:A.3.1}.
There is a coherent choice of the even-shift maps
\(\widetilde F_{h,c}:X_c\to X_{c+h}\), \(h\in2\mathbb Z\),
with the following explicit ordinary-chart translation. Choose global
roots \(F^3=t\), \(G^4=t-1\) on the uniformizing cover and put
\(\ell=(2\pi i)^{-1}\log(F/G)\pmod{\mathbb Z}\), \(e=(0,1)^t\).
Then the translation is
\begin{equation}\hypertarget{tag-02OP}{}\label{mapsI:closing:eq:C-explicit-closing}
 b_{h,c}=-\frac h{12}\Pi_c\widehat\gamma
              +\left(h\ell-\frac{h^2}{24}\right)e.
\end{equation}
For \(h=12\), it is the scalar multiplication
\begin{equation}\hypertarget{tag-02OQ}{}\label{mapsI:closing:eq:C-twelve-scalar}
 K(t)=\frac{t^4}{(t-1)^3}
\end{equation}
on the ordinary \(\mathbb C^*\)-coordinate, after removing an integral
period. On the finite root charts its scalar representatives are
\(K(t)s_1^{-12}\) and \(K(t)s_2^{12}\), respectively, both units.
In a cusp coordinate \(q=t_c\), the scalar representative is
\(q^{-1}H(q)\), with \(H\) a holomorphic nonvanishing function.
The same map has a holomorphic \((\mathbb C^*)^2\)-unit representative
on the entire cusp filling. These two representatives differ by
exactly the \(\widehat\gamma\) cusp deck transformation.

For six positive \(w\)-loops the closing map is
\(\widetilde F_{12,c_0-12}\); thus the stated cusp exponent is \(-1\)
with $w=e^{-\pi i c}$ and positive loops in the $w$-plane. This computes the closing
cocycle, not the mapping class of its composition with parallel
transport. In particular it does not imply \(6[\Phi]=0\).
\end{cited}
The expanded source derivation is retained in the proof compendium, Paper~I, label \nolinkurl{push:src:C.3.5}.




The raw even-shift formula and its coherent correction are due to \citet[Proposition~7.8, pp.~91--92, and Proposition~C.2, p.~102]{DuEtAl2026v1}. A complete re-derivation is given in \CompendiumRef{I}{app:rederive:push-src-C.3.5}.




\begin{figure}[htbp]
\centering
\begin{tikzpicture}[x=1cm,y=1cm,>=stealth,font=\small]
\draw[thick] (-5.3,0) .. controls (-1.8,-.35) and (1.8,-.35) .. (5.3,0);
\node at (0,-1.05) {$B=\mathbb P^1_t$};
\foreach \x/\lab in {-4/0,0/1,4/\infty}{
 \fill (\x,-.1) circle (2pt);
 \node[below=5pt] at (\x,-.1) {$\lab$};
 \draw[->] (\x,1.15)--(\x,.15);
}
\node[draw,rounded corners,align=center,text width=2.8cm,minimum height=1.2cm] at (-4,1.75)
 {$3S_1$\\free order-three\\root quotient};
\node[draw,rounded corners,align=center,text width=2.8cm,minimum height=1.2cm] at (0,1.75)
 {$4S_2$\\free order-four\\root quotient};
\node[draw,rounded corners,align=center,text width=2.8cm,minimum height=1.2cm] at (4,1.75)
 {$W$\\normalization $\mathrm{dP}_6$};
\draw[->] (-2,1.15)--(-2,.15);
\node[align=center] at (-2,2.85) {ordinary fibre\\$\mathbb C^2/\Pi_c\mathbb Z^4$};
\node[align=center] at (4,3.05) {three paired boundary curves\\two triple points};
\end{tikzpicture}
\caption{The fixed $(3,4,\infty)$ base data. The arrows denote the
restriction of $f$ to each indicated fibre, with its scheme
multiplicity retained. The ordinary period quotient, finite affine
quotients and the paired hexagon are the three actual local models
in the common conventions; their overlap maps are fixed there.}
\hypertarget{tag-01RL}{}\label{fig:I:orbifold}
\end{figure}

\begin{figure}[htbp]
\centering
\begin{tikzpicture}[x=1cm,y=1cm,>=stealth,font=\small]
\fill[blue!4] (0,0) circle (2.6);
\draw[thick,dashed] (0,0) circle (2.6);
\draw[->,gray] (-3.1,0)--(3.2,0) node[right] {$\Re w$};
\draw[->,gray] (0,-3.15)--(0,3.2) node[above] {$\Im w$};
\draw[very thick,blue!65!black] (0,-2.56)--(0,-.13);
\draw[very thick,blue!65!black] (0,.13)--(0,2.56);
\fill[white] (0,0) circle (.11);
\draw[thick] (0,0) circle (.11);
\draw[->] (.23,-.18)--(1.1,-.7);
\node[anchor=west,align=left] at (1.05,-.88) {puncture $w=0$\\corner winding $-pn$};
\node[align=center] at (-1.45,.8) {$0<|w|<216$};
\node[anchor=west,align=left] at (3.15,1.65)
 {outer circle $|w|=216$\\$|q_2|=1$\\$q_2=iw/216$};
\draw[->] (3.05,1.55)--(2.15,1.43);
\node[anchor=east,align=right] at (-.25,1.8)
 {real structures\\$w\in i\mathbb R$};
\node[anchor=west,align=left] at (3.15,-2.1)
 {all admissible members\\retain $3S_1,4S_2,W$};
\draw[->] (3.05,-1.95)--(1.85,-1.7);
\end{tikzpicture}
\caption{The exact parameter disc for the standard fixed data.
The coordinate map is $w=e^{-\pi i c}$ and the coarser coordinate
map is $u=w^2$. The imaginary diameter is the fixed locus of
$w\mapsto-\bar w$, which is the reflected-family map.
The two boundary obstructions use the meridian comparison at the
puncture and the intrinsic Albanese nome at the outer circle.
The special fibres are over $t=0,1,\infty$ in every member, as
indicated in Figure~\ref{fig:I:orbifold}.}
\hypertarget{tag-01RM}{}\label{fig:I:parameter-disc}
\end{figure}

\begin{figure}[htbp]
\centering
\begin{tikzpicture}[x=1cm,y=1cm,>=stealth,font=\small]
\begin{scope}[shift={(-3.3,0)}]
 \fill[blue!5] (-1,0)--(-1,2.8)--(1,2.8)--(1,0)
  arc[start angle=30,end angle=150,radius=1.1547];
 \draw[thick] (-1,0)--(-1,2.8);
 \draw[thick] (1,0)--(1,2.8);
 \draw[thick] (1,0) arc[start angle=30,end angle=150,radius=1.1547];
 \node at (0,1.3) {$\mathcal F$};
 \node[align=center] at (0,-1.05)
 {$|\Re\tau|\le\tfrac12,\ |\tau|\ge1$\\$0\le\theta<\pi$};
 \draw[->] (0,2.1)--(0,3.1) node[above] {cusp};
\end{scope}
\begin{scope}[shift={(1.0,.2)},x={(1.8cm,0cm)},y={(.75cm,.5cm)},z={(0cm,2.2cm)}]
 \fill[blue!4] (0,0,0)--(1,0,0)--(1,0,1)--(0,0,1)--cycle;
 \draw[thick] (0,0,0)--(1,0,0)--(1,1,0)--(1,1,1)--(0,1,1)--(0,0,1)--cycle;
 \draw[thick] (1,0,0)--(1,0,1)--(0,0,1);
 \draw[thick] (1,0,1)--(1,1,1);
 \draw[dashed] (0,0,0)--(0,1,0)--(1,1,0);
 \draw[dashed] (0,1,0)--(0,1,1);
 \node[below] at (.5,0,0) {$0\le x_1<1$};
 \node[right] at (1,.5,0) {$0\le x_2<1$};
 \node[left] at (0,0,.5) {$0\le t<12$};
\end{scope}
\node[align=center] at (1.95,-1.55)
 {$U:\ t\mapsto t-6-6x_2$\qquad
  $V:\ t\mapsto t-6+6x_1$\\
  $Z:\ t\mapsto t+12$};
\end{tikzpicture}
\caption{A fundamental set for the index-twelve monodromy kernel
in the real parabolic stabilizer of $d$.
The base is a modular fundamental set with its rotation coordinate;
the fibre is the indicated Heisenberg parallelepiped. The displayed
face maps include the central shears. The representative map and
the boundary identifications are proved in
\ref{I:monodromy-fundamental-set}; the index twelve occurs in the
central fibre coordinate.}
\hypertarget{tag-01RK}{}\label{fig:I:monodromy-domain}
\end{figure}

\begin{proposition}[The fundamental set of the monodromy kernel]
\hypertarget{tag-0001}{}\label{I:monodromy-fundamental-set}
\StatementRecord{I:monodromy-fundamental-set}{proved}{0001}{\SeriesRef{fam:src:F.1.1}, \SeriesRef{fam:src:F.1.2}}
In the real coordinates $g(B,x,t)$ of
\ref{fam:src:F.1.1}, a measurable fundamental set for
$\Gamma\backslash P_d(\mathbb R)$ consists of the matrices
\[
 g(B,x,t),\qquad
 B=n(\Re\tau)a(\Im\tau)k(\theta),\quad
 \tau\in\mathcal F,\quad 0\le\theta<\pi,
 \quad x_1,x_2\in[0,1),\quad t\in[0,12).
\]
Here
$n(u)=\left(\begin{smallmatrix}1&u\\0&1\end{smallmatrix}\right)$,
$a(v)=\operatorname{diag}(\sqrt v,1/\sqrt v)$,
and $k(\theta)$ is the rotation matrix. Paired modular edges
are assigned to one side, and their lower-dimensional
stabilizer identifications use the corresponding lifts in
$\Gamma$: for
$S=\left(\begin{smallmatrix}0&-1\\1&0\end{smallmatrix}\right)$
and $T=\left(\begin{smallmatrix}1&1\\0&1\end{smallmatrix}\right)$,
we use $\widetilde S=g(S,0,9)$,
$\widetilde T=g(T,0,1)$ and
$\widetilde{-I}=g(-I,0,6)$.
Their maps on the three coordinates are
$(B,x,t)\mapsto(AB,Ax,t+\alpha(A))$, with
$\alpha(S)=9$, $\alpha(T)=1$ and $\alpha(-I)=6$.
The Heisenberg face identifications are the
maps printed in Figure~\ref{fig:I:monodromy-domain}.
\end{proposition}
% ALL PRINTED PROOF MATERIAL MASKED






































\section{The limit at the puncture}\hypertarget{tag-01PS}{}\label{fam:section:E}
\subsection{The untwisted pinching and its local smoothings}\hypertarget{tag-01PT}{}\label{fam:section:E.1}

In this section $S$ is the smooth minimal rational elliptic surface with
Weierstrass equation
\[
 y^2=x^3-3t^3(t-1)x+2t^4(t-1)^2,
\]
$\pi:S\to B$ is its elliptic fibration, $P$ is the section used in the
period dictionary, $T=t_{6P}:S\to S$, and $L=M_S=\OO_S(P-O)$.
We use the convention that $\mathbb P_S(\OO\oplus L)$ contains the total
space of $L$, with fibre coordinate $a$, zero section $s_0$, and infinity
section $s_\infty$. The two normal bundles are respectively $L$ and
$L^{-1}$. We write $R=\mathbb P_S(\OO\oplus L)$.

\begin{proposition}\hypertarget{tag-01R1}{}\label{fam:src:E.1.1}\label{fam:st:E.1.1}
\StatementRecord{fam:src:E.1.1}{proved}{01R1}{\cite[Lemma 2.1]{Engel2026}}

The quotient
\[
 Q=R/\{s_0(x)\sim s_\infty(Tx):x\in S\}
\]
has a compact, reduced complex-space structure. Its normalization is
the finite holomorphic map $\nu:R\to Q$. Its singular locus is a copy
of $S$, and its local analytic equation along that locus is $ab=0$,
with two further coordinates on $S$. The induced map $Q\to B$ is
proper. Its sheaf of first-order local smoothing parameters is
\[
 \mathcal T^1_Q=\mathcal Ext^1_{\OO_Q}(\Omega_Q,\OO_Q)
       \simeq L\otimes T^*L^{-1}\simeq\pi^*\OO_B(1)
\]
on the double locus.
\end{proposition}



% ALL PRINTED PROOF MATERIAL MASKED



















































\begin{proposition}\hypertarget{tag-01R2}{}\label{fam:src:E.1.2}\label{fam:st:E.1.2}
\StatementRecord{fam:src:E.1.2}{proved}{01R2}{\SeriesRef{fam:src:E.1.1}; \cite{Engel2026}}

Over $S^\circ=S\setminus S_\infty$, the pushout in \ref{fam:src:E.1.1} is the central
fibre of a holomorphic plumbing with parameter $u$. On every relatively
compact open disc $U\Subset B\setminus\{\infty\}$ this plumbing is
proper and flat over $U\times\Delta_\epsilon$, for a sufficiently small
positive $\epsilon$. Its fibre at $u\ne0$ is, over $U$, Engel's untwisted
quotient by the contracting generator $u\phi$ covering $T^{-1}$.
This assertion includes discs containing $0$ or $1$; it is made before
the logarithmic transformations at those points.
\end{proposition}



% ALL PRINTED PROOF MATERIAL MASKED































































\begin{proposition}\hypertarget{tag-01R3}{}\label{fam:src:E.1.3}\label{fam:st:E.1.3}
\StatementRecord{fam:src:E.1.3}{proved}{01R3}{the definitions and constructions stated in the proof}

The first simultaneous local plumbing at the nodal point of $S_\infty$
has equation
\[
                 xy=u zv,\qquad q=zv,                         \tag{E2}
\]
after a holomorphic change of one normal coordinate, where $q$ is a
local coordinate on $B$ at infinity. This hypersurface is flat over
the $u$-disc. For each $u\ne0$, its three-dimensional fibre has an
ordinary double point at $x=y=z=v=0$. The projective modification
defined by the blow-up of $(x,z)$ resolves that point for $u\ne0$.
Its two strict-transform charts are
\[
       x=zr,\quad ry=uv;
       \qquad z=xa,\quad y=uav.                               \tag{E3}
\]
The modified family remains flat over $u$. At $u=0$, the normalization
of its central fibre contains the blow-up of one branch of (E2) along
the smooth curve $x=z=0$; hence this particular modification changes
the normalization of the unmodified pinching.
\end{proposition}



% ALL PRINTED PROOF MATERIAL MASKED










































\begin{proposition}\hypertarget{tag-01R4}{}\label{fam:src:E.1.4}\label{fam:st:E.1.4}
\StatementRecord{fam:src:E.1.4}{proved}{01R4}{\SeriesRef{fam:src:E.1.1}; \cite[Lemma 2.1]{Engel2026}}

The compact candidate $Q$ of \ref{fam:src:E.1.1} has
\[
 \pi_1(Q)=\Z,\qquad
 (H_i(Q;\Z))_{i=0}^6
   =(\Z,\Z,\Z^8,\Z^7,\Z^{10},0,\Z),\qquad\chi(Q)=12.
\]
These are invariants of the stated candidate; they do not identify it
as a central fibre of the prescribed compact family.
\end{proposition}



% ALL PRINTED PROOF MATERIAL MASKED



















































































\paragraph{Local obstruction and the precise gluing requirement}\hypertarget{tag-01R5}{}\label{fam:src:E.1.5}\label{fam:st:E.1.5}
\StatementRecord{fam:src:E.1.5}{proved}{01R5}{\SeriesRef{fam:src:E.1.1}, \SeriesRef{fam:src:E.1.2}, \SeriesRef{fam:src:E.1.3}, \SeriesRef{push:src:A.1.1}, \SeriesRef{push:src:A.2.1}, \SeriesRef{push:src:A.2.2}, \SeriesRef{push:src:A.3.1}, \SeriesRef{push:src:A.3.2}}

The following local construction isolates the obstruction to a flat proper filling of the prescribed compact untwisted family with central fibre exactly $Q$. The full nonexistence result is proved in \cref{push:src:A.3.1}, \cref{push:src:A.3.2}. \ref{fam:src:E.1.1} constructs the compact candidate as an actual complex
space; \ref{fam:src:E.1.2} constructs its smoothing away from infinity, including
the two other special fibres; \ref{fam:src:E.1.3} computes the two-parameter germ and a concrete flat modification of it.

The local requirement identified by this argument is the following analytic gluing
problem at one lift of the node of $S_\infty$. The local model
$xy=uzv$ must be replaced by charts whose restriction to every
$u\ne0$ is the chosen small resolution of the rank-two Mumford
filling, and these charts must be equivariant under both period
generators, with proper quotient and with central normalization $R$.
The explicit blow-up in \ref{fam:src:E.1.3} supplies the required smooth local
general fibre but changes the normalization at $u=0$. This particular blow-up does not settle whether another gluing can preserve that normalization. Thus the local computation alone is an obstruction to the unmodified plumbing argument. The global virtual-Euler argument in \cref{push:src:A.3.1}, \cref{push:src:A.3.2} supplies the exclusion. A modified compact limit
also requires an equivariant periodic fan and verification of its
quotient near the crossing $(u,q)=(0,0)$; the complete fan and quotient data are supplied in \cref{push:src:A.1.1}, \cref{push:src:A.2.1}, \cref{push:src:A.2.2}.



% ALL PRINTED PROOF MATERIAL MASKED























\Uses{the hypotheses, notation and definitions stated here; \cref{push:src:A.3.1}, \cref{push:src:A.3.2}; \cref{fam:src:E.1.1}; \cref{fam:src:E.1.2}; \cref{fam:src:E.1.3}; \cref{push:src:A.1.1}, \cref{push:src:A.2.1}, \cref{push:src:A.2.2}}


\subsection{Finite bielliptic limits and the smooth monodromy}\hypertarget{tag-01PU}{}\label{fam:section:E.2}

\begin{proposition}\hypertarget{tag-01R6}{}\label{fam:src:E.2.1}\label{fam:st:E.2.1}
\StatementRecord{fam:src:E.2.1}{proved}{01R6}{the definitions and constructions stated in the proof}

After the auxiliary base change $w=s^6$, each of the two families of
\emph{reduced} special surfaces $S_j(w)$ has a proper flat
semistable filling by explicit finite free quotients. The central
surface for $j=1$ is a cycle of four ruled surfaces over an elliptic
curve; that for $j=2$ is a cycle of three. They have
\[
\begin{array}{c|ccccc|c}
 &H_0&H_1&H_2&H_3&H_4&\chi\\\hline
 j=1&\Z&\Z\oplus\Z/3&\Z^5&\Z^9&\Z^4&0\\
 j=2&\Z&\Z\oplus\Z/2&\Z^4&\Z^7&\Z^3&0
\end{array}
\]
with integral coefficients. Their fundamental groups are
$\Z^2\rtimes_{R_3}\Z$ and $\Z^2\rtimes_{R_4}\Z$, respectively,
where
\[
 R_3=\begin{pmatrix}0&-1\\1&-1\end{pmatrix},\qquad
 R_4=\begin{pmatrix}0&-1\\1&0\end{pmatrix}.
\]
\end{proposition}



% ALL PRINTED PROOF MATERIAL MASKED


















































































































\section{The untwisted limit}\hypertarget{tag-02YZ}{}\label{push:section:A}
\subsection{The nodal plumbing}\hypertarget{tag-02Z0}{}\label{push:section:A.1}

We write $L=\mathcal O_S(P-O)$, $T=t_{6P}$,
$J=L\otimes T^*L^{-1}=\pi^*\mathcal O_B(1)$, and
$R=\mathbb P_S(\mathcal O\oplus L)$. The convention for this projective
bundle is that its zero and infinity sections have normal bundles
$L$ and $L^{-1}$. The identification defining $Q$ is
$s_0(x)\sim s_\infty(Tx)$. The equality defining $J$ is the equality
in Engel's Lemma 2.1, with this direction of $T$; its intersection
calculation uses $P\cdot O=0$, $P^2=O^2=-1$,
$6P\cdot O=6P\cdot P=2$. Thus the degree on $O$ is
$(O-P)\cdot6P+(P-O)\cdot O=1$. We use the section
$\phi\in H^0(S,J)$ whose divisor is $S_\infty$, as in his
Corollary 2.2. We use the stated version of
\cite{Engel2026}, as is shown by its explicit local equations.

\begin{proposition}\hypertarget{tag-02ZH}{}\label{push:src:A.1.1}\label{push:A11}
\StatementRecord{push:src:A.1.1}{proved}{02ZH}{the definitions and constructions stated in the proof}

For $k,l\in\mathbb Z$, let $P_{k,l}$ be the cone defined by $h_q,h_u\geq0$, $kh_q\leq y_1\leq(k+1)h_q$ and $l(h_q+h_u)\leq y_2\leq(l+1)(h_q+h_u)$. These cones and their faces form a fan. Its two
period automorphisms on the cocharacter lattice are
\[
 G_1(y_1,y_2,h_q,h_u)=(y_1+h_q,y_2,h_q,h_u),\qquad
 G_2(y_1,y_2,h_q,h_u)=(y_1,y_2+h_q+h_u,h_q,h_u).
\]
The affine chart of $P_{0,0}$ is
\[
 U_{0,0}=\operatorname{Spec}\mathbb C[x,y,z,v,u]/(xy-uzv),
 \qquad q=zv.
\]
In particular the unsubdivided chart is not a smooth simultaneous
model of the cusp.
\end{proposition}


% ALL PRINTED PROOF MATERIAL MASKED




























\begin{theorem}\hypertarget{tag-02ZI}{}\label{push:src:A.1.2}\label{push:A12}
\StatementRecord{push:src:A.1.2}{proved}{02ZI}{\SeriesRef{push:A11}; \cite{EGFS25}}

There is a reduced complex space $\mathcal X^\flat$ and a proper flat
holomorphic map
\[
 (f,\varpi):\mathcal X^\flat\longrightarrow B\times\Delta_u
\]
whose central threefold is $Q$. For $u\ne0$ its fibre is the
contraction of Engel's chosen $X(u)$ at one rational curve of the
Mumford fibre. This curve is the double curve corresponding to the
diagonal of the square in the chosen $A_2$ triangulation; in the
return's coordinates its direction is $e_1-e_2$. Its normal bundle
in $X(u)$ is $\mathcal O_{\mathbb P^1}(-1)^{\oplus2}$. Its image is
the unique ordinary double point of $\mathcal X^\flat_u$.
\end{theorem}


% ALL PRINTED PROOF MATERIAL MASKED





















































































































\subsection{Two minimal toroidal fillings and further modifications}\hypertarget{tag-02Z1}{}\label{push:section:A.2}

\begin{proposition}\hypertarget{tag-02ZJ}{}\label{push:src:A.2.1}\label{push:A21}
\StatementRecord{push:src:A.2.1}{proved}{02ZJ}{\SeriesRef{push:src:A.1.1}}

Among subdivisions with no new rays, the fan in \ref{push:src:A.1.1} has exactly
two smooth period-invariant subdivisions restricting to the chosen
$A_2$ triangulation at $h_u=0$. These are also all subdivisions
which are flat over the original $(q,u)$ bidisc, have reduced
$u$-central divisor, and retain that boundary triangulation.
Without the joint-flatness restriction there are infinitely many
smooth invariant subdivisions retaining that triangulation and
having reduced $u$-central divisor.
\end{proposition}


% ALL PRINTED PROOF MATERIAL MASKED




























































\begin{theorem}\hypertarget{tag-02ZK}{}\label{push:src:A.2.2}\label{push:A22}
\StatementRecord{push:src:A.2.2}{proved}{02ZK}{\SeriesRef{push:src:A.2.1}, \SeriesRef{push:A12}, \SeriesRef{push:src:A.1.2}}

The two minimal subdivisions of \ref{push:src:A.2.1} give proper flat holomorphic
families $\mathcal X^+,\mathcal X^-\to\Delta_u$ with smooth total
spaces. Their joint maps to $B\times\Delta_u$ are flat. Over
$\Delta_u^*$ they are isomorphic as families to the specified
untwisted family $X(u)$. Each central fibre $Q^\pm$ is reduced,
irreducible, and has normal crossings, with smooth double locus
$S$ and no triple points.

Its normalization $R^\pm$ is obtained by blowing up $R$ along
$S_\infty$ in, respectively, one of its two sections and then
taking the small resolution at the unique ordinary double point
of that blow-up prescribed by the fixed bottom diagonal.
The two proper transforms of the sections are still copies of
$S$, glued by $T$. Their normal-bundle product is trivial;
in particular $\mathcal T^1_{Q^\pm}\simeq\mathcal O_S$.
\end{theorem}


% ALL PRINTED PROOF MATERIAL MASKED















































































\paragraph{counterexample; the literal blow-up shape}\hypertarget{tag-02ZL}{}\label{push:src:A.2.3}\label{push:A23}
\StatementRecord{push:src:A.2.3}{proved}{02ZL}{\SeriesRef{push:src:A.2.1}, \SeriesRef{push:src:A.2.4}, \SeriesRef{push:src:A.2.5}, \SeriesRef{fam:src:E.1.1}}

The unrestricted request that every normalized component
be a blow-up of $R$ is false. The first maximal-cone
star already contributes a $\mathbb P^3$ component,
which admits no proper bimeromorphic map to $R$, since
such a map would inject a rank-eleven rational $H^2$
into a rank-one group. \ref{push:src:A.2.4} gives every normalization
by explicit star charts; \ref{push:src:A.2.5} supplies the terminating
blow-up and blow-down algorithm with centres, order,
exceptional bundles and gluing maps. Thus the literal
shape is refuted and its general algorithmic
replacement is proved.



% ALL PRINTED PROOF MATERIAL MASKED























\Uses{the hypotheses, notation and definitions stated here; \cref{push:src:A.2.4}; \cref{push:src:A.2.5}}


\begin{theorem}[exhaustive finite fan data]\hypertarget{tag-02ZM}{}\label{push:src:A.2.4}\label{push:A24}
\StatementRecord{push:src:A.2.4}{proved}{02ZM}{\SeriesRef{push:src:A.2.1}, \SeriesRef{push:src:B.2.1}, \SeriesRef{push:src:A.1.2}; \cite[Section 2, p.6]{DMM16}}

The following finite-data procedure parametrizes all smooth
period-invariant subdivisions of the cusp fan after
$u=r^n$, retaining the prescribed bottom triangulation and
having finitely many cones modulo the periods. Here $n=1$
gives Part A, and $n=2,12$ give the two cases of B.2.
Each output defines a proper flat family over $\Delta_r$
with smooth total space, whose punctured restriction is the
given untwisted family pulled back by $u=r^n$.

For a bound $H\geq1$, retain the four bottom rays
$A,B,C,D$ and consider all primitive vectors
\[
 \begin{gathered}
 v=(a,b,k,\ell),\qquad
 \ell\geq1,\quad k\geq0,\quad k+\ell\leq H,\\
 0\leq a\leq k,\qquad0\leq b\leq k+n\ell.
 \end{gathered}
 \tag{A.24a}
\]
Take any collection of four-element subsets of this finite
set satisfying the following exact tests:
\begin{enumerate}
\item Each four-column determinant has absolute value one.
\item After dividing each ray by $k+\ell$, the resulting
tetrahedra intersect pairwise in common faces, and their
volumes sum to $(n+2)/6$.
\item Their bottom boundary is the two triangles $ABC,BCD$.
Their boundary subdivisions on $a=0$ and $a=k$ agree under
$G_1$, and those on $b=0$ and $b=k+n\ell$ agree under $G_2$.
\end{enumerate}
All these are finite rational-arithmetic tests. Extend a
passing collection by the two periods, and discard duplicates
as $H$ increases. This is exhaustive. For a reduced central
divisor, retain only $\ell=1$ in (A.24a).

The central prime components are indexed by period orbits of
the selected rays with $\ell>0$, and the multiplicity of the
component indexed by $v$ is exactly $\ell$. The normalization
charts of that component are the star fan of $v$ in
$N/\mathbb Zv$, with the inherited bundle units and the
stabilizer action. For $k>0$ the stabilizer is trivial and this
normalization is a smooth complete toric threefold. All
multiple loci, including possible triple and quadruple loci,
are specified by the positive-$\ell$ rays belonging to each
cone. In the reduced case
$\mathcal T^1\simeq\mathcal O_{\operatorname{Sing}D}$,
so every output is d-semistable.
\end{theorem}


% ALL PRINTED PROOF MATERIAL MASKED





























































































































































































\begin{proposition}[explicit normalization algorithm and its necessary scope]\hypertarget{tag-02ZN}{}\label{push:src:A.2.5}\label{push:A25}
\StatementRecord{push:src:A.2.5}{proved}{02ZN}{\SeriesRef{push:src:A.2.4}, \SeriesRef{push:src:A.4.6}, \SeriesRef{push:src:A.4.8}, \SeriesRef{push:src:A.2.2}}

The requirement that every normalized component of every
smooth refinement be a blow-up of
$R=\mathbb P_S(\mathcal O\oplus L)$ is false, already
for a reduced refinement with $n=1$.
The following is a uniform replacement for any finite
fan certificate $\Sigma$ in \ref{push:src:A.2.4}.

The star charts and integer transition matrices of \ref{push:src:A.2.4}
give every normalized component, its gluing maps, and
all multiple loci directly. There is also a terminating
finite-data algorithm that expresses its comparison
with a reference model as an explicit sequence of
blow-ups and blow-downs along specified smooth strata.
At each blow-up the algorithm gives the normalized
strict transforms, the exceptional projective bundle,
its multiplicity, and every new intersection.
This algorithm allows blow-downs; no toric
strong-factorization assertion is used.

For a component of multiplicity $\ell$, the local
nonreduced central equation is the monomial in A.24d.
More generally, in a unimodular chart with central
equation $F=\prod_i z_i^{\ell_i}$, the cotangent-complex
first-order deformation sheaf is
\[
 \mathcal T^1_D
  =\mathcal O_D\Big/
    \left(\ell_i z_i^{\ell_i-1}
          \prod_{j\ne i}z_j^{\ell_j}:\ell_i>0\right).
 \tag{A.25a}
\]
Here $\mathcal T^1_D$ means
$\mathcal Ext^1(\mathbb L_D,\mathcal O_D)$; in the
reduced case it also equals
$\mathcal Ext^1(\Omega_D,\mathcal O_D)$.
The reduced case is the trivial line on the singular
scheme, as in \ref{push:src:A.2.4}. For a nonreduced centre the displayed
Jacobian quotient is the appropriate deformation sheaf;
the reduced normal-crossings notion of d-semistability
is not asserted for it.
\end{proposition}


% ALL PRINTED PROOF MATERIAL MASKED





















































































































































\subsection{The virtual Euler obstruction and the conjecture}\hypertarget{tag-02Z2}{}\label{push:section:A.3}

\begin{lemma}\hypertarget{tag-02ZO}{}\label{push:src:A.3.1}\label{push:A31}
\StatementRecord{push:src:A.3.1}{proved}{02ZO}{\cite[Lemmas 3.1--3.2 and Section 3.1]{Levy08}; \cite[Section~3, Theorems~1--2, pp.~463--464]{Loj64}; \cite[Theorem~1.6, p.~5; Remark~1.16, p.~11]{BM95AB}; \cite{FG10}}

Let $Z$ be a compact reduced local-complete-intersection complex
space of pure dimension $n$. In any proper flat holomorphic
family over a disc with central fibre $Z$ and smooth nearby
fibres, their Euler number is
\[
 e_{\mathrm{vir}}(Z)=\langle c_n(T_Z^{\mathrm{vir}}),[Z]\rangle.
\]
This statement needs no projectivity or K\"ahler hypothesis.
\end{lemma}


% ALL PRINTED PROOF MATERIAL MASKED







































































\begin{theorem}[Conjecture 8.2 refuted]\hypertarget{tag-02ZP}{}\label{push:src:A.3.2}\label{push:A32}
\StatementRecord{push:src:A.3.2}{proved}{02ZP}{\SeriesRef{push:src:A.1.2}, \SeriesRef{push:src:A.2.2}, \SeriesRef{push:src:A.4.3}, \SeriesRef{push:A31}; \cite{Engel2026}}

For the specified pinching $Q$ one has $e_{\mathrm{vir}}(Q)=0$.
Every smooth proper-flat smoothing of $Q$ therefore has Euler
number zero. In particular there is no proper flat holomorphic
filling of Engel's untwisted family with central fibre exactly
$Q$, even if no map to $B$ is required. Thus Conjecture 8.2
of the attached \cite{Engel2026} is false as worded. The proper nodal
filling of \ref{push:src:A.1.2} and the two natural semistable fillings of
\ref{push:src:A.2.2} are precise replacements, with different central fibres.
\end{theorem}


% ALL PRINTED PROOF MATERIAL MASKED









































































\begin{proposition}[counterexample and corrected]\hypertarget{tag-02ZQ}{}\label{push:src:A.3.3}\label{push:A33}
\StatementRecord{push:src:A.3.3}{proved}{02ZQ}{\SeriesRef{push:section:A.3}, \SeriesRef{push:src:A.2.1}}

The assertion that the two natural
fillings exhaust smooth-total-space reduced-normal-crossings
fillings is false, even when joint flatness over
$B\times\Delta$ is imposed. There are infinitely many such
fillings of the same punctured untwisted family. The
two-model classification in \ref{push:src:A.2.1} is a classification of the
specified minimal toroidal refinements, not of all fillings.
\end{proposition}


% ALL PRINTED PROOF MATERIAL MASKED














































\subsection{Integral topology and nearby cycles}\hypertarget{tag-02Z3}{}\label{push:section:A.4}

\begin{proposition}\hypertarget{tag-02ZR}{}\label{push:src:A.4.1}\label{push:A41}
\StatementRecord{push:src:A.4.1}{proved}{02ZR}{\SeriesRef{fam:src:E.1.1}, \SeriesRef{fam:src:E.1.2}, \SeriesRef{fam:src:E.1.3}, \SeriesRef{fam:src:E.1.4}, \SeriesRef{push:src:A.2.2}}

With integral coefficients, the spaces constructed above have
\[
\begin{array}{c|rrrrrrr|r}
&b_0&b_1&b_2&b_3&b_4&b_5&b_6&\chi\\\hline
R&1&0&11&0&11&0&1&24\\
R^\pm&1&0&12&0&12&0&1&26\\
Q&1&1&8&7&10&0&1&12\\
Q^\pm&1&1&9&7&11&0&1&14
\end{array}
\]
All these homology groups are torsion free. The fundamental
groups of $R,R^\pm$ are trivial and those of $Q,Q^\pm$ are
$\mathbb Z$. The dual complex of either $Q^\pm$ is a circle.
\end{proposition}


% ALL PRINTED PROOF MATERIAL MASKED





































































































\begin{proposition}\hypertarget{tag-02ZS}{}\label{push:src:A.4.2}\label{push:A42}
\StatementRecord{push:src:A.4.2}{proved}{02ZS}{\SeriesRef{push:src:A.2.2}, \SeriesRef{push:src:A.4.1}}

The untwisted smooth fibre has torsion-free integral homology
with Betti numbers
\[
 (b_0,b_1,b_2,b_3,b_4,b_5,b_6)=(1,1,8,14,8,1,1).
\]
The positive-$u$ monodromy is the identity outside degree three.
In degree three there is an integral exact sequence
\[
 0\longrightarrow \mathcal L/\langle F,v,m\rangle
 \longrightarrow H_3(X(u);\mathbb Z)
 \longrightarrow K\longrightarrow0,\qquad
 K=\{a\in\mathcal L:aF=av=am=0\},
\]
in which $T_u-I$ is zero on the left term and, on the quotient,
is the map
\[
 K\longrightarrow \mathcal L/\langle F,v,m\rangle,\qquad
 a\longmapsto a\bmod\langle F,v,m\rangle.
\]
The sign convention orients the neck circle by increasing
$\arg x$ in $xy=u$; the boundary convention for a transported
chain is chosen so that a positive $u$ loop adds this circle
times its boundary cycle. With these conventions
$(T_u-I)^2=0$, its rank is six, and its nonzero Smith factors
are $1,1,1,1,1,11$.
\end{proposition}


% ALL PRINTED PROOF MATERIAL MASKED





















































































\begin{proposition}[the pointwise nearby Euler function]\hypertarget{tag-02ZT}{}\label{push:src:A.4.3}\label{push:A43}
\StatementRecord{push:src:A.4.3}{proved}{02ZT}{\SeriesRef{push:src:A.3.2}}

For $\mathcal X^\flat\to\Delta_u$ the nearby-fibre Euler
function on the double locus of $Q$ is zero off $S_\infty$
and one at every point of $S_\infty$, including its node.
Its integral is one. In particular the assertion that the
pointwise value on the smooth part of $S_\infty$ is zero
is false; its integrated contribution is zero.
For either semistable family $\mathcal X^\pm$ the nearby
Euler function is one off the double locus and zero on
it, giving $\chi(X(u))=26-2\cdot12=2$.
\end{proposition}


% ALL PRINTED PROOF MATERIAL MASKED






























\begin{proposition}[limitation of Clemens--Schmid]\hypertarget{tag-02ZU}{}\label{push:src:A.4.4}\label{push:A44}
\StatementRecord{push:src:A.4.4}{proved}{02ZU}{\SeriesRef{push:src:A.4.1}}

The classical K\"ahler or projective Clemens--Schmid
comparison cannot be applied to these untwisted models.
Their smooth fibres have $b_1=1$ and are not K\"ahler.
Moreover the exceptional small-resolution curve in
$R^\pm$ is null-homologous, so that normalization is
not K\"ahler either. The rational monodromy weight
filtration, defined algebraically from $N=T_u-I$ and
centred at degree three, has graded ranks $6,2,6$ in
weights $2,3,4$. This is a filtration of the computed
monodromy representation; we do not identify it with
a polarized limiting mixed Hodge structure.
\end{proposition}


% ALL PRINTED PROOF MATERIAL MASKED


















\begin{corollary}[unrestricted toroidal topology]\hypertarget{tag-02ZV}{}\label{push:src:A.4.5}\label{push:A45}
\StatementRecord{push:src:A.4.5}{proved}{02ZV}{\SeriesRef{push:src:A.2.4}, \SeriesRef{push:src:A.4.4}, \SeriesRef{push:src:A.4.7}, \SeriesRef{push:src:A.4.8}, \SeriesRef{push:src:A.4.9}}

For every output of \ref{push:src:A.2.4}, \ref{push:src:A.4.7}--\ref{push:src:A.4.9} determine the
Euler number, fundamental group, all integral homology
groups with torsion, nearby Euler function and marked
monodromy. The integral normalization computation is
given by the finite algorithm and the primitive
central pushforward matrices of \ref{push:src:A.4.9}; the independent
ambient computation is \ref{push:src:A.4.8}, and the torus-stratum
Euler count is \ref{push:src:A.4.7}. The non-K\"ahler limitation on
classical Clemens--Schmid is the explicit one in
\ref{push:src:A.4.4}; no polarized limiting Hodge structure is
asserted.
\end{corollary}

% ALL PRINTED PROOF MATERIAL MASKED























\Uses{the hypotheses, notation and definitions stated here; \cref{push:src:A.2.4}; \cref{push:src:A.4.7}; \cref{push:src:A.4.9}; \cref{push:src:A.4.8}; \cref{push:src:A.4.4}}


\begin{proposition}[an infinite explicit subclass]\hypertarget{tag-02ZW}{}\label{push:src:A.4.6}\label{push:A46}
\StatementRecord{push:src:A.4.6}{proved}{02ZW}{\SeriesRef{push:src:A.2.1}, \SeriesRef{push:src:A.4.2}, \SeriesRef{push:src:A.4.1}}

Let $D_m$ be the central fibre after the $m$ maximal-cone
star subdivisions in \ref{push:src:A.2.1}, starting with $Q^+$, where
$m\geq0$. Then
\[
 \begin{gathered}
 \pi_1(D_m)=\mathbb Z,\qquad
 H_*(D_m;\mathbb Z)\text{ is free},\\
 (b_0,\ldots,b_6)=(1,1,9+m,7,11+m,0,1+m).
 \end{gathered}
\]
and $\chi(D_m)=14+3m$. For $m\geq1$ its normalization is
the disjoint union of $\operatorname{Bl}_pR^+$, $m-1$
copies of $\operatorname{Bl}_{p_i}\mathbb P^3$, and a final
$\mathbb P^3$. The original two sections are still glued
by $T$. The new components form a chain attached to the
original component. Consecutive components meet in one
$\mathbb P^2$, with no triple intersections. The dual
complex is the original circle with a chain of $m$ edges
attached. The nearby Euler function is one off the double
locus and zero on it; its integral remains two. The nearby
integral monodromy remains exactly that of \ref{push:src:A.4.2}.
\end{proposition}


% ALL PRINTED PROOF MATERIAL MASKED





























































\begin{theorem}[Euler number and fundamental group for every toroidal output]\hypertarget{tag-02ZX}{}\label{push:src:A.4.7}\label{push:A47}
\StatementRecord{push:src:A.4.7}{proved}{02ZX}{\SeriesRef{push:src:A.2.4}, \SeriesRef{push:src:B.4.2}; \cite[Theorem 1.6 and its following paragraph]{BM95AB}}

Let $\Sigma$ be any output of \ref{push:src:A.2.4} for $u=r^n$, and let
$D_\Sigma$ denote the underlying reduced central space.
Let $k(\Sigma)$ be the number of selected pure top ray
orbits, and let $t_4(\Sigma)$ be the number of maximal
four-cones in one fundamental prism. Then
\[
 \chi(D_\Sigma)=11k(\Sigma)+t_4(\Sigma),
 \qquad \pi_1(D_\Sigma)=\mathbb Z,
 \qquad H_1(D_\Sigma;\mathbb Z)=\mathbb Z.
\]
These assertions include nonreduced central divisors.
At a point with precisely one local central branch of
multiplicity $\ell$, the nearby Euler function is $\ell$;
at a point with at least two local branches it is zero.
Its integral is always two. For every reduced output the
pure top count is $k(\Sigma)=n$, so
$\chi(D_\Sigma)=11n+t_4(\Sigma)$.
The higher integral homology of every arbitrary output is
not asserted by this theorem.
\end{theorem}


% ALL PRINTED PROOF MATERIAL MASKED






























































































































\begin{theorem}[integral homology of every toroidal output]\hypertarget{tag-02ZY}{}\label{push:src:A.4.8}\label{push:A48}
\StatementRecord{push:src:A.4.8}{proved}{02ZY}{\SeriesRef{push:src:A.2.4}, \SeriesRef{push:src:A.4.2}, \SeriesRef{push:src:A.4.1}, \SeriesRef{push:src:A.4.7}; \cite[Section 5.6, pp.565--566]{AKMW02AB}; \cite[Theorems 1.6.1 and 2.7.1]{AKMW02AB}; \cite[Section 2, p.2]{Matsuki00ErrAB}; \cite[Section 5.5]{AKMW02AB}}

Let $D_\Sigma$ be the reduced central space of an output of
\ref{push:src:A.2.4} for $u=r^n$. Write $k$ for its number of pure top
ray orbits, $m$ for its number of mixed ray orbits, and
$t_4$ for its number of maximal four-cone orbits. Then
every integral homology group is free, and
\[
 \begin{split}
 (b_0,b_1,b_2,b_3,b_4,b_5,b_6)
   ={}&(1,1,8+k+m,7,\\
       &\hspace{20mm}9k+t_4-2m-1,0,k+m).
 \end{split}
\]
This includes nonreduced central divisors: the homology
is that of their underlying support. The marked nearby
fibre is $X(r^n)$, and its integral monodromy is
$I+nN$ in degree three, where $N$ is the map in \ref{push:src:A.4.2};
it is the identity in every other degree. In particular
the six nonzero Smith factors of this degree-three
monodromy minus the identity are
$n,n,n,n,n,11n$.
\end{theorem}


% ALL PRINTED PROOF MATERIAL MASKED



























































































































































































































\begin{proposition}[a second integral central-space computation]\hypertarget{tag-02ZZ}{}\label{push:src:A.4.9}\label{push:A49}
\StatementRecord{push:src:A.4.9}{proved}{02ZZ}{\SeriesRef{push:src:A.4.8}, \SeriesRef{push:src:A.2.5}, \SeriesRef{push:src:A.4.7}}

The homology formula of \ref{push:src:A.4.8} can also be computed
on the central spaces, by their normalization
resolutions and the strata of each recorded star.
More precisely, let $Y'\to Y$ be one ordinary
blow-up in the algorithm of \ref{push:src:A.2.5}, with closed
smooth centre $C$ of complex codimension $c$.
Let $D,D'$ be the reduced central supports and
$C_0=C\cap D$. In the derived category of integral
sheaves on $D$ there is an isomorphism
\[
 R(\pi_D)_*\mathbb Z_{D'}\simeq
 \mathbb Z_D\oplus
 \bigoplus_{a=1}^{c-1}(i_0)_*\mathbb Z_{C_0}[-2a],
 \tag{A.49a}
\]
where $\pi_D:D'\to D$ and $i_0:C_0\hookrightarrow D$
are the induced continuous maps. In particular
\[
 H^j(D';\mathbb Z)\simeq H^j(D;\mathbb Z)
 \oplus\bigoplus_{a=1}^{c-1}H^{j-2a}(C_0;\mathbb Z).
 \tag{A.49b}
\]
The maps in this splitting are integral, with unit
diagonal on the exceptional-fibre basis. Thus it
also determines the primitive boundary maps and
the absence of torsion in the normalization
mapping-cylinder computation along the entire
finite algorithm, not just the Euler characteristic.
\end{proposition}


% ALL PRINTED PROOF MATERIAL MASKED

































































































































\section{Local limits and the global component obstruction}\hypertarget{tag-02Z4}{}\label{push:section:B}
\subsection{The finite root neighbourhoods}\hypertarget{tag-02Z5}{}\label{push:section:B.1}

\begin{proposition}[counterexample and corrected]\hypertarget{tag-0300}{}\label{push:src:B.1.1}\label{push:B11}
\StatementRecord{push:src:B.1.1}{proved}{0300}{\SeriesRef{fam:src:B.1.1}, \SeriesRef{fam:src:B.1.4}, \SeriesRef{push:src:A.1.2}, \SeriesRef{fam:src:E.2.1}}

The assertion that the filling is the product of the smooth elliptic
family with a twelve-gon Tate model is false on the whole root disc. There is, however, a proper flat holomorphic filling of
the marked root family over
$\Delta_{s_j}\times\Delta_r$, where $u=r^{12}$, with smooth
total space and reduced normal-crossings $r$-central divisor.
Its central surface over an ordinary value of $s_j$ is a cycle
of twelve ruled surfaces over the smooth elliptic curve
$E'_{j,s_j}$. Their normalizations are projective bundles
associated with the degree-zero line bundle represented by
$P_{s_j}-O_{s_j}$, with the appropriate pullbacks by translations.
They are not generically products with $\mathbb P^1$.
\end{proposition}


% ALL PRINTED PROOF MATERIAL MASKED





































































\begin{theorem}\hypertarget{tag-0301}{}\label{push:src:B.1.2}\label{push:B12}
\StatementRecord{push:src:B.1.2}{proved}{0301}{\SeriesRef{push:src:B.1.1}, \SeriesRef{fam:src:E.2.1}, \SeriesRef{fam:src:A.5.1}}

The twisted finite deck generators extend to the models of
\ref{push:src:B.1.1}. On the twelve component labels their actions are
$i\mapsto i+4$ at $j=1$ and $i\mapsto i-3$ at $j=2$,
with labels in $\mathbb Z/12$. Every nonidentity group
element acts freely, including at every node of the central
cycle. The quotient is a smooth complex fourfold, proper
and flat over $(t,r)=(s_j^{m_j},r)$, with a reduced
normal-crossings $r$-central divisor. Its fibre divisor
over $t=p_j$ has multiplicity exactly $m_j$.

The reduction over $s_j=r=0$ is the central surface of \cref{fam:src:E.2.1}:
a cycle of four ruled surfaces at $j=1$, and a cycle of
three at $j=2$. In the root presentation the normal
character to $s_j=0$ is $e^{-2\pi i/m_j}$ and the
character on $r$ is trivial.
\end{theorem}


% ALL PRINTED PROOF MATERIAL MASKED










































































































\subsection{The base-changed cusp fan}\hypertarget{tag-02Z6}{}\label{push:section:B.2}

\begin{proposition}\hypertarget{tag-0302}{}\label{push:src:B.2.1}\label{push:B21}
\StatementRecord{push:src:B.2.1}{proved}{0302}{\SeriesRef{push:src:A.2.1}, \SeriesRef{push:src:A.1.2}, \SeriesRef{push:src:A.2.2}}

For the plain base change $u=r^n$ of the untwisted cusp
model, the smooth period-invariant refinements which
retain the prescribed bottom triangulation, are jointly
flat over $(q,r)$, and have reduced $r$-central divisor
are indexed by $k=0,\ldots,n$. In a fundamental prism
write its four bottom rays as $A,B,C,D$ as in \ref{push:src:A.2.1},
and put $\rho_i=(0,i,0,1)$ for $0\leq i\leq n$.
The maximal cones for the index $k$ are
\[
 ABC\rho_k,\quad BCD\rho_k,\quad
 AB\rho_i\rho_{i+1}\ (i<k),\quad
 CD\rho_i\rho_{i+1}\ (i\geq k).
\]
There are thirteen such models for $n=12$ and three
for $n=2$. Each has smooth total space and a reduced
normal-crossings central divisor with $n$ components
and no triple locus. Its component labels over the
punctured $q$-disc are globally $\mathbb Z/n$; in
particular their monodromy around $q=0$ is the identity.
\end{proposition}


% ALL PRINTED PROOF MATERIAL MASKED




















































































\begin{corollary}[general plain-cusp normalization]\hypertarget{tag-0303}{}\label{push:src:B.2.2}\label{push:B22}
\StatementRecord{push:src:B.2.2}{proved}{0303}{\SeriesRef{push:src:A.2.4}, \SeriesRef{push:src:A.2.5}, \SeriesRef{push:src:A.4.7}, \SeriesRef{push:src:A.4.8}, \SeriesRef{push:src:A.4.9}}

For $n=2$ and $n=12$, \ref{push:src:A.2.4}--\ref{push:src:A.2.5} give all
finite-orbit smooth invariant plain-cusp refinements,
their proper flat families, normalization charts,
finite blow-up and blow-down presentations,
multiplicities and all multiple loci. \ref{push:src:A.4.7}--\ref{push:src:A.4.9}
give their complete integral topology and nearby
Euler function. These statements concern the plain
untwisted cusp. They assert no compatibility with
the finite logarithmic-transform pieces of B.1.
\end{corollary}

% ALL PRINTED PROOF MATERIAL MASKED























\Uses{the hypotheses, notation and definitions stated here; \cref{push:src:A.2.4}; \cref{push:src:A.2.5}; \cref{push:src:A.4.7}; \cref{push:src:A.4.9}}


\begin{corollary}[unrestricted toroidal parametrization]\hypertarget{tag-0304}{}\label{push:src:B.2.3}\label{push:B23}
\StatementRecord{push:src:B.2.3}{proved}{0304}{\SeriesRef{push:src:A.2.4}, \SeriesRef{push:src:A.2.5}, \SeriesRef{push:src:A.4.7}, \SeriesRef{push:src:A.4.8}, \SeriesRef{push:src:A.4.9}}

For $u=r^{12}$ and $u=w^2$, all finite-orbit smooth invariant
subdivisions of the plain cusp fan are obtained by \ref{push:src:A.2.4}
with $n=12$ and $n=2$, respectively. The construction there
gives their proper flat smooth-total-space families,
central multiplicities, normalization charts and all
multiple loci. Reducedness is exactly the restriction
$\ell=1$. This is a classification of the plain untwisted
cusp models; it does not supply the missing transition to
the finite logarithmic-transform models in B.3.
\end{corollary}

% ALL PRINTED PROOF MATERIAL MASKED























\Uses{the hypotheses, notation and definitions stated here; \cref{push:src:A.2.4}}


\subsection{The prescribed pieces do not glue}\hypertarget{tag-02Z7}{}\label{push:section:B.3}

\begin{theorem}[counterexample to simultaneous gluing]\hypertarget{tag-0305}{}\label{push:src:B.3.1}\label{push:B31}
\StatementRecord{push:src:B.3.1}{proved}{0305}{\SeriesRef{push:section:B.1}, \SeriesRef{push:section:B.2}, \SeriesRef{push:src:B.1.2}, \SeriesRef{push:src:B.2.1}}

The finite models of \ref{push:section:B.1} and any of the plain
$u=r^{12}$ cusp models of \ref{push:section:B.2} cannot be glued to a
proper complex space over $B\times\Delta_r$ with
the prescribed central fibres on all three base
neighbourhoods. In particular the proposed simultaneous gluing is false. This conclusion does
not rule out a different cusp model of the same
punctured $S^6$ family.
\end{theorem}


% ALL PRINTED PROOF MATERIAL MASKED


































































\begin{proposition}[the marked genus-three component cover]\hypertarget{tag-0306}{}\label{push:src:B.3.2}\label{push:B32}
\StatementRecord{push:src:B.3.2}{proved}{0306}{\SeriesRef{push:src:B.3.1}}

Let $B^\circ=\mathbb P^1\setminus\{0,1,\infty\}$ and use the
clockwise deck-identification convention of \ref{push:src:B.3.1}.
The peripheral rotations $4,-3,-1$ in $\mathbb Z/12$ define a
connected degree-twelve cover of $B^\circ$. Its compact
normalization is the smooth genus-three curve obtained from
\[
 y^{12}=\frac{t^4}{(t-1)^3}.
\]
The ramification indices over $0,1,\infty$ are $3,4,12$,
respectively. The peripheral relation determines the last rotation
from the first two, so this is the marked cyclic component cover
used in the partial filling of \ref{push:src:B.4.19}.
\end{proposition}

% ALL PRINTED PROOF MATERIAL MASKED




























\begin{proposition}[obstruction to a specified repair]\hypertarget{tag-0307}{}\label{push:src:B.3.3}\label{push:B33}
\StatementRecord{push:src:B.3.3}{proved}{0307}{\SeriesRef{push:src:B.3.1}, \SeriesRef{push:src:A.1.2}}

A cyclic cusp-root construction using a
fractional translation by the two cusp
periods cannot repair \ref{push:src:B.3.1} while retaining
the reduced Mumford fibre for $r\ne0$.
More precisely, let $q=s^d$, $u=r^n$,
and suppose a cyclic deck generator has
translation $d^{-1}(aG_1+bG_2)$ in the
additive notation for the two periods.
Assume its $r$-component shift is integral.
Then its $r$-component shift is
$nb/d\bmod n$, whereas the multiplicity
of every generic quotient $q$-component
is a positive multiple of
\[
 \frac{d}{\gcd(d,a,b)}.
\]
A nonzero $r$-component shift therefore
forces a nonreduced $q$-fibre. In
particular a shift $-1\bmod12$ in this
class of constructions forces multiplicity
at least twelve.
\end{proposition}


% ALL PRINTED PROOF MATERIAL MASKED
















































\begin{theorem}[split-normalization trace obstruction]\hypertarget{tag-0308}{}\label{push:src:B.3.4}\label{push:B34}
\StatementRecord{push:src:B.3.4}{proved}{0308}{\SeriesRef{push:src:B.3.2}; \cite[Theorems 8.4 and 8.6, table (8.16)]{Shioda90}; \cite[Theorem 4.1]{Shioda90}; \cite[Theorem 1.3]{Shioda90}}

Let $C$ be the genus-three component cover
in \ref{push:src:B.3.2}, and let $\mathscr E\to C$ be
the relatively minimal elliptic surface
obtained by pulling back $S\to B$.
Write $\sigma$ for the order-twelve deck
transformation, with its natural lift to
$\mathscr E$. Then
\[
 \operatorname{Tr}_{\sigma}
   \mathscr E(\mathbb C(C))
 :=\left\{\sum_{i=0}^{11}\sigma^i Q:
                 Q\in\mathscr E(\mathbb C(C))\right\}
 =12\mathbb Z P.
\]
In particular there is no section $Q$
for which $(t_Q\circ\sigma)^{12}=t_{6P}$.
Thus the proposed repair that takes the
normalization over $C$ to be ruled over
this split elliptic surface, with adjacent
conductor map $t_Q\circ\sigma$, is
impossible.

This conclusion is restricted to the
stated split normalization. A ruled
normalization over a nontrivial genus-one
torsor need not admit a section or such
an expression for its affine deck lift;
it is not excluded by this theorem.
\end{theorem}


% ALL PRINTED PROOF MATERIAL MASKED






























































































































































































































































\begin{corollary}[exact Mordell--Weil lattice]\hypertarget{tag-0309}{}\label{push:src:B.3.5}\label{push:B35}
\StatementRecord{push:src:B.3.5}{proved}{0309}{\SeriesRef{push:src:B.3.4}}

In the notation of \ref{push:src:B.3.4},
\[
 \mathscr E(\mathbb C(C))
   =\mathbb ZP\oplus\mathbb ZR\oplus\mathbb ZR',
 \qquad \det\operatorname{NS}(\mathscr E)=-32.
\]
\end{corollary}


% ALL PRINTED PROOF MATERIAL MASKED

















































\begin{proposition}[an unramified torsor alternative]\hypertarget{tag-030A}{}\label{push:src:B.3.6}\label{push:B36}
\StatementRecord{push:src:B.3.6}{proved}{030A}{\SeriesRef{push:src:B.3.4}, \SeriesRef{push:src:B.4.1}; \cite[Chapter III, Section 4, Theorem 2
and the following remark, pp.129--130]{HS53}; \cite[Lemma 4.7.8]{CDL24}; \cite{CDL24}; \cite[Section 4, equation (4)]{Xiao91}; \cite[Sections 3--4]{Xiao91}}

Put $K=\mathbb C(t)$, $L=\mathbb C(C)$,
$G=\operatorname{Gal}(L/K)=\langle\sigma\rangle$,
and let $E/K$ be the generic elliptic
curve of $S\to B$. There is a nonzero
class
\[
 \xi\in\operatorname{Sha}(E_L/L)^G
 \quad\hbox{whose transgression in}\quad
 H^2(G,E(L))=\mathbb ZP/12\mathbb ZP
 \quad\hbox{is }6P.
\]
Here $\operatorname{Sha}$ denotes the
subgroup of genus-one torsors trivial
over the henselian local fields of all
points of $C$.
In particular the trace obstruction of
\ref{push:src:B.3.4} does not obstruct all genus-one
torsors. Such a class has a relatively
minimal projective genus-one surface
$V\to C$, locally isomorphic to
$\mathscr E\to C$, and an affine lift
$A:V\to V$ of $\sigma$ with
\[
 A^{12}=t_{6P}.
\]
This is an existence statement for a
possible normalization base, not a
construction of a threefold central
fibre or of its smoothing.
\end{proposition}


% ALL PRINTED PROOF MATERIAL MASKED








































































































































































\begin{proposition}[an auxiliary family with the wrong cusp]\hypertarget{tag-030B}{}\label{push:src:B.3.7}\label{push:B37}
\StatementRecord{push:src:B.3.7}{proved}{030B}{\SeriesRef{push:src:B.3.6}, \SeriesRef{push:src:A.1.2}, \SeriesRef{push:src:A.2.2}; \cite[Proposition 0.9.2 and the opening
of Section 4.2, p.370]{CDL24}; \cite[Proposition 8.12, equation (8.25)]{Shioda90}}

For the torsor and affine lift $(V,A)$
of \ref{push:src:B.3.6}, there is a line bundle $M$ on
$V$ representing the narrow Jacobian
section $P$, with
\[
 M^2=-2,\qquad
 M\otimes(A^*M)^{-1}
       \simeq\pi^*\mathcal O_C(\infty).
\]
These data give a proper flat nodal
plumbing $\mathcal Z^\flat\to\Delta_h$
and two uniform smooth-total-space
semistable modifications
$\mathcal Z^\pm\to\Delta_h$, together
with holomorphic maps to $B$.
For $h\ne0$, the nodal fibre has
twelve ordinary double points and
Euler number twelve; each small-resolved
fibre has Euler number twenty-four.
The cusp fibre of the latter map to
$B$ has multiplicity twelve.
Consequently this auxiliary construction
does not fill the prescribed $S^6$
family. No proper birational modification
over $B$ can turn its multiple cusp
into a reduced cusp.
\end{proposition}


% ALL PRINTED PROOF MATERIAL MASKED













































































































































































\begin{proposition}[analytic unramified torsors do not evade the Hodge obstruction]\hypertarget{tag-030C}{}\label{push:src:B.3.8}\label{push:B38}
\StatementRecord{push:src:B.3.8}{proved}{030C}{\SeriesRef{push:src:B.3.4}, \SeriesRef{push:src:B.4.5}, \SeriesRef{push:src:B.3.6}; \cite[Theorem 11, p.297]{Buchdahl1999}; \cite[Section 2.1, p.8]{Chen2026HodgeRoot}}

Let $V\to C$ be any compact smooth analytic genus-one surface
locally isomorphic, as a torsor over every base disc, to the
Jacobian surface $\mathscr E\to C$ of \ref{push:src:B.3.4}. Algebraicity
or finite order of this analytic torsor class is not assumed.
Then $b_1(V)=6$, and $V$ is K\"ahler. For every holomorphic
line bundle $M$ on $V$, the projective bundle
$\mathbb P_V(\mathcal O\oplus M)$ is K\"ahler. Every smooth
compact threefold bimeromorphic to such a bundle satisfies
the $\partial\bar\partial$-lemma.

Consequently the necessary non-$\partial\bar\partial$
normalization component in \ref{push:src:B.4.5} cannot be obtained simply
by replacing the projective torsor of \ref{push:src:B.3.6} by a nonalgebraic
locally trivial analytic torsor, followed by a ruled bundle
and bimeromorphic modifications. This excludes that specified
repair; it does not assert that every possible cusp
normalization has this form.
\end{proposition}


% ALL PRINTED PROOF MATERIAL MASKED






















































\subsection{Nonexistence of a proper analytic filling}

\begin{corollary}[counterexample; the proposed global filling]\hypertarget{tag-030D}{}\label{push:src:B.4.1}\label{push:B41}
\StatementRecord{push:src:B.4.1}{proved}{030D}{\SeriesRef{push:src:B.4.33}, \SeriesRef{push:section:B.1}, \SeriesRef{push:section:B.2}, \SeriesRef{push:src:B.4.19}}

The actual punctured $S^6$ family has no proper
Hausdorff complex-analytic filling after any finite
cover of $w=0$. In particular, a proper flat filling
with smooth total space and reduced normal-crossings
centre does not exist. There is no positive least
semistable exponent, no equivariant analytic filling
to descend, and no central multiplicity data for
such a descent.
\end{corollary}


% ALL PRINTED PROOF MATERIAL MASKED













The obstruction cannot be removed by citing
$H^1(B,\mathcal V)=0$ from \cref{fam:src:A.5.1}. That
vanishing solves an additive translation
cocycle for a fixed nonzero parameter and
its integer shifts. On a finite overlap
the translation contains
\[
 \frac{\log s_j}{2\pi i}\,\Pi_c v_j,
 \qquad c=-\frac{12}{2\pi i}\log r.
\]
Its vertical exponential contains a product
of the two logarithms. No assertion in
the cited vanishing controls its growth
at $r=0$ or makes it a unit in the central
toric charts. \ref{push:src:B.3.1} proves that at least
one of those extension requirements must
fail for the prescribed pieces. That local
failure alone would not rule out a different cap.
\ref{push:src:B.4.30}--\ref{push:src:B.4.33} supply the additional global argument
which excludes every proper analytic filling of
the actual family, including such alternatives.

\begin{theorem}[fundamental-group obstruction]\hypertarget{tag-030O}{}\label{push:src:B.4.2}\label{push:B42}
\StatementRecord{push:src:B.4.2}{proved}{030O}{\cite[Theorem 5.1, p.2226]{NO10}}

Let $f:\mathcal Y\to\Delta$ be a
proper holomorphic map from a smooth
complex manifold, with smooth connected
fibres over $\Delta^*$ and reduced
normal-crossings central fibre $D$.
There is a continuous specialization
map from a nearby fibre $Y_t$ to $D$
which induces a surjection
\[
 \pi_1(Y_t)\longrightarrow\pi_1(D).
\]
No K\"ahler hypothesis is needed.
Hence the central fibre of any such
semistable filling of $S^6$ is simply
connected. In particular it cannot
be a simple pinching of a connected
normalization along one pair of
disjoint nonempty conductor sections,
whether its elliptic normalization
base is split or is a genus-one torsor.
\end{theorem}


% ALL PRINTED PROOF MATERIAL MASKED






















































































































\subsection{4 supplement. A Hodge obstruction to the central strata}

\begin{cited}[consequences with explicit derivation]\hypertarget{tag-030Z}{}\label{push:src:B.4.3}
\StatementRecord{push:src:B.4.3}{proved}{030Z}{\SeriesRef{push:src:B.4.33}; \citet[Theorem 1.4, preprint p.~7, and Example 2.1, p.~13]{Nishiguchi1988IIRoot}; \citet[Theorem 7.1]{IstratiEtAl2022Root}}

Let $g:\mathcal Z\to\Delta_s$ be a proper holomorphic map from a
smooth complex fourfold, smooth over $\Delta_s^*$, with reduced
normal-crossings central fibre $D$. Suppose its nearby fibres are
biholomorphic to the six-sphere members of the specified family,
after any finite base change $w=s^N$. Then
\[
 H^1(D;\C)=0,\qquad h^1(D,\OO_D)\geq1.
\]
Consequently the natural map
$H^1(D;\C)\to H^1(D,\OO_D)$ is not surjective. In particular,
$D$ cannot be the analytification of a proper complex algebraic
variety with normal-crossings singularities.

If $D$ has \emph{simple} normal crossings, at least one of its
nonempty closed strata fails the $\partial\bar\partial$-lemma.
In particular, no such simple normal-crossings filling can have
all irreducible components K\"ahler or all irreducible components
Moishezon. These obstructions hold for every $N\geq1$.
\end{cited}
The expanded source derivation is retained in the proof compendium, Paper~I, label \nolinkurl{push:src:B.4.3}.




This is the stated application of the classical Euler and normal-crossing cohomology results; all hypotheses and the independent derivation are retained in the appendix. A complete re-derivation is given in \CompendiumRef{I}{app:rederive:push-src-B.4.3}.


This is an application of existing Hodge-theoretic results, not
a claim of a new degeneration theorem. Searches on 15 November
2026 for six-sphere semistable degenerations, normal-crossings
limits, Moishezon central fibres, and Conjecture 8.2 with Hodge
theory did not locate this specific application. We opened
Chen's version 2 and checked Definition 2.3
and Theorem 2.4(4), including the condition inserted before
parts (3)--(4); omitting that condition would make the
statement false. We also opened the original Graf--Kov\'acs
article and checked the proper Du Bois assertion in Section
1.D. The separate source audit records the searches.

The theorem prevents a repair made solely from compact
Moishezon components in a simple normal-crossings model.
It does not prove that an arbitrary analytic filling must have
such components. In particular, no argument here makes an
unknown three-dimensional component over the cusp Moishezon.
These Hodge arguments alone give conditional restrictions, not
unrestricted nonexistence. The latter is proved in \ref{push:src:B.4.33} for
every finite parameter cover, using the actual affine clutching.
The Hopf-component example of
\citet[Theorem 1.4, preprint p.~7, and Example 2.1, p.~13]{Nishiguchi1988IIRoot}
shows why an analytic degeneration cannot simply be assumed to
have class-$\mathcal C$ strata. The toric-Kato construction of
\citet[Theorem 7.1]{IstratiEtAl2022Root} is another relevant
non-K\"ahler analogue. Neither source supplies the required
cusp comparison for the six-sphere family.

\begin{proposition}[extension of the marked fibration]\hypertarget{tag-0314}{}\label{push:src:B.4.4}
\StatementRecord{push:src:B.4.4}{proved}{0314}{\SeriesRef{fam:src:A.1.1}, \SeriesRef{fam:src:A.1.2}; \citep[Section 6, Hauptsatz I]{Grauert1960D}}

Suppose $g:\mathcal Z\to\Delta_s$ is a proper holomorphic map
from a smooth connected complex fourfold, and its restriction
over $\Delta_s^*$ is
isomorphic to a finite base change of the specified six-sphere
family. The given marked holomorphic map
\[
 f^*:\mathcal Z|_{\Delta_s^*}\longrightarrow\PP^1_t
\]
extends to a meromorphic map $\mathcal Z\dashrightarrow\PP^1_t$.
Its graph gives a proper projective modification
$\mathcal Z'\to\mathcal Z$, an isomorphism over
$\Delta_s^*$, on which the marked map is holomorphic.
No assertion that $\mathcal Z'$ is smooth or semistable is
included in this statement.
\end{proposition}


% ALL PRINTED PROOF MATERIAL MASKED





































































































No condition on the central divisor was used: a nonzero base
germ is a nonzero divisor on the smooth dominating total space.
In particular the proposition applies after resolving the
reduced main component of any proper analytic filling, even
when the resulting central fibre has multiplicities.

The proposition removes a possible essential projective-coordinate
ambiguity. It does not extend a punctured isomorphism to any of
the auxiliary models of \ref{push:section:B.3}, and does not make a component over
$t=\infty,s=0$ Moishezon. Those would be additional statements.

\begin{theorem}[the obstruction on normalization components]\hypertarget{tag-0315}{}\label{push:src:B.4.5}
\StatementRecord{push:src:B.4.5}{proved}{0315}{\SeriesRef{push:src:B.4.2}, \SeriesRef{push:src:B.4.3}}

Let $D$ be a compact reduced complex space with normal-crossings
singularities, allowing self-intersections of global components.
Suppose that every connected component of its smooth normalization
satisfies the $\partial\bar\partial$-lemma. Then the natural map
\[
 H^1(D;\C)\longrightarrow H^1(D,\OO_D)
\]
is surjective.

Consequently every smooth-total-space reduced normal-crossings
filling of the specified six-sphere family, after any finite base
change, has a normalization component which fails the
$\partial\bar\partial$-lemma. In particular its normalization
cannot consist entirely of compact K\"ahler or Moishezon manifolds.
This conclusion includes normal crossings which are not simple.
\end{theorem}


% ALL PRINTED PROOF MATERIAL MASKED







































































































This strengthens the simple-normal-crossings obstruction in
\ref{push:src:B.4.3} by using only degree one. It also avoids an unverified
extension of an algebraic Du Bois theorem to arbitrary analytic
spaces. Searches for normal-crossings cohomology with K\"ahler
normalization, the normalization Dolbeault complex, and
degree-one surjectivity located the broader Hodge-theoretic
setting discussed by \citet{Chen2026HodgeRoot} and the proper
algebraic statement of \citet{GrafKovacs2014Root}. The proof above
states and verifies the particular analytic comparison required
here. It still does not prove that all possible analytic
normalization components satisfy its hypothesis.

\begin{proposition}[a restricted non-K\"ahler cap obstruction]\hypertarget{tag-0316}{}\label{push:src:B.4.6}\label{push:B46}
\StatementRecord{push:src:B.4.6}{proved}{0316}{\SeriesRef{push:src:B.3.8}, \SeriesRef{push:src:B.3.7}, \SeriesRef{push:src:B.4.5}, \SeriesRef{push:src:B.4.2}, \SeriesRef{push:src:B.4.3}; \cite[Section 2, p.1352]{IstratiEtAl2022Root}; \cite[Proposition 8.4 and Theorem 8.5]{IstratiEtAl2022Root}; \cite[proof of Theorem 8.8, p.1380]{IstratiEtAl2022Root}; \cite[Theorem 5.2]{IstratiEtAl2022Root}}

Let $D$ be a compact reduced normal-crossings complex
space. Suppose its irreducible components have been
partitioned into closed unions $D_0,\ldots,D_s$, with
their reduced induced structures. Assume that
\[
 H^1(D_i,\mathbb C)\longrightarrow H^1(D_i,\mathcal O)
 \quad\hbox{is surjective for every }i,
 \tag{B.46a}
\]
and that the analogous map is injective for every
pairwise intersection $D_i\cap D_j$. Then
\[
 H^1(D,\mathbb C)\longrightarrow H^1(D,\mathcal O_D)
 \quad\hbox{is surjective}.
 \tag{B.46b}
\]

In particular these hypotheses hold for the following
specified cap assembly. The horizontal block $D_0$ is
a two-section pinching of a ruled bundle over an
analytic locally trivial torsor $V\to C$ of \ref{push:src:B.3.8},
with conductor map $A$ covering $\sigma$ and
$A^{12}=t_{6P}$. The normalization may undergo the
modifications of \ref{push:src:B.3.7}, or any smooth bimeromorphic
modifications for which the two conductor restriction
maps on $H^1$ are the same pullbacks as before.
The other blocks are smooth hyperbolic toric Kato
threefolds or compact smooth spaces satisfying
(B.46a), for example class-$\mathcal C$ spaces.
Assume every pairwise overlap is either a reduced
union of invariant divisors of a hyperbolic toric
Kato threefold, or a compact normal-crossings surface
whose normalized components and normalized double
curves are rational. No such assembly can be a
semistable central fibre of the $S^6$ family.

This excludes the stated cap construction even when
an invariant Kato surface in an overlap kills a
fundamental-group generator of a Kato component.
It is not a classification of arbitrary analytic
cusp components or of their possible overlaps.
\end{proposition}


% ALL PRINTED PROOF MATERIAL MASKED






































































































































\paragraph{Why automatic analytic comparison is not used.}
We also examined whether general results on degenerations of tori
would force an arbitrary filling into the excluded toroidal class.
The distinction between a construction theorem and such a comparison
theorem is essential. \citet[Theorem 5.1.3 and Sections 5.2--5.4]
{KKN2008Root} construct proper analytic models of log complex tori
without assuming a polarization. Their input is a specified log
complex torus and an admissible pairing; the theorem does not assert
that every proper analytic filling is one of these models.
\ref{push:src:B.4.24} constructs the actual local Jacobian log input and isolates
its affine torsor obstruction. \ref{push:src:B.4.33} derives an associated torsor
from traces in a hypothetical cap; it does not identify that cap
with a proper fan model.

The classification in \citet[Theorem 3.1]{KamenovaToriRoot} assumes
a K\"ahler total space and trivial canonical bundle. Those hypotheses
have not been proved for a slice of an arbitrary filling here.
\citet[Proposition 3.1, preprint pp.~6--7]{SchroeerSiebert2002Root} explicitly
study normal-crossings degenerations with Hopf normalization, but
their smooth fibres are primary Kodaira surfaces. This is a source
of alternative analytic models to test, not a comparison with our
punctured family.

Even extending the marked map in \ref{push:src:B.4.4} does not extend a chosen
punctured isomorphism to a second model. Properness alone cannot
justify that inference: the automorphism
$(s,[z_0:z_1])\mapsto(s,[e^{1/s}z_0:z_1])$ of
$\Delta^*\times\PP^1$ has no meromorphic extension across zero,
although both product families are proper over the disc. Its
restriction to $[1:1]$ has an essential singularity. Thus an
unrestricted proof needs additional information beyond a
punctured isomorphism. \ref{push:src:B.4.30}--\ref{push:src:B.4.33} supply exactly such information:
bounded differences of extending sections and a global radial-loop
invariant. Their contradiction proves nonexistence without extending
an individual punctured isomorphism meromorphically.


\begin{proposition}[a parabolic analytic model and its comparison failure]\hypertarget{tag-0317}{}\label{push:src:B.4.7}
\StatementRecord{push:src:B.4.7}{proved}{0317}{\SeriesRef{push:src:B.4.6}, \SeriesRef{push:src:B.4.2}, \SeriesRef{push:src:A.4.7}; \citep[Definition 1.1, pp.~85--86]{Fujimoto1990AB}; \citet[Section 2.2]{AbashevaDeev2024AB}}

There is a proper semistable analytic family of compact complex
surfaces $p:\mathcal H\to\Delta$ with smooth total space such that:
\begin{enumerate}
\item its punctured restriction is a holomorphic product with the
scalar primary Hopf surface
\[
 H_-=(\C^2\setminus\{0\})/\langle -e^{-\pi}\id\rangle;
\]
\item its central fibre is $Y=H_1\cup_E H_2$, where both $H_i$
are isomorphic to
\[
 H_+=(\C^2\setminus\{0\})/\langle e^{-2\pi}\id\rangle,
 \qquad E=\C/(\Z+i\Z);
\]
\item $N_{E/H_1}\simeq N_{E/H_2}\simeq\mathcal O_E$, and
\[
 \pi_1(Y)=0,\qquad h^1(Y,\mathcal O_Y)=1,
 \qquad H_*(Y;\Z)=(\Z,0,0,\Z^3,\Z^2);
\]
\item no punctured product isomorphism has a meromorphic extension
whose graph is proper over the two total spaces near the central
fibres.
\end{enumerate}
Consequently $\mathcal H\times\PP^1\to\Delta$ is an actual
semistable family of compact threefolds with simply connected centre
and $h^1(\mathcal O)=1$, whose normalized components fail the
$\partial\bar\partial$ lemma.  Its smooth fibres have Euler number
zero and fundamental group $\Z$.  It is not a filling of the
six-sphere family.
\end{proposition}


% ALL PRINTED PROOF MATERIAL MASKED























































































































































































\paragraph{Primary literature and the scope of the auxiliary cap}
The Cartier-divisor cocycle above is established machinery, not
a new logarithmic-transformation construction.  Fujimoto's
Definition~1.1 permits an arbitrary Cartier divisor; his
Remark~1.4 attributes non-K\"ahler torus degenerations to Ueno.
His Proposition~1.14 extends principal elliptic bundles across
subsets of codimension at least three.  Example~1.15 concerns
failure of extension as a principal elliptic bundle in
codimension two and explicitly leaves general relative
compactification open there; it supplies no unrestricted
nonexistence theorem for the present cusp
\citep[pp.~85--86, 93--94]{Fujimoto1990AB}.

Hopf components in analytic semistable degenerations are already
present in the original work of
\citet[Theorem 1.4, preprint p.~7, and Example 2.1, p.~13]{Nishiguchi1988IIRoot}.
The explicit diagonal Hopf quotients, elliptic conductor gluing
and normal-line calculation in
\citet[Propositions 3.1--3.4, preprint pp.~6--9]{SS00}
give a closely related analytic model.  Those are comparisons,
not smoothing theorems for the present threefold.  The construction
above supplies its own complete proper family and does not invoke
such a theorem by analogy.

There are two concrete obstacles to directly inserting the product
cap into B.3.  The divisor $E\times\PP^1$ in
$H_i\times\PP^1$ has normal degree zero along the $\PP^1$
ruling.  The exceptional divisor obtained by blowing up an
elliptic curve in a horizontal normalization has normal degree
$-1$ on its exceptional $\PP^1$ fibres.  For a two-component
attachment with no triple divisor on this interface, these cannot
be dual normal lines.  In the general normal-crossings formula
a triple divisor must also be included: a section
$E\times\{p\}$ has degree $+1$ on the ruling and could cancel
this particular mismatch.  No obstruction to such a higher-intersection
attachment is asserted.  Moreover the actual generic two-tori
do not possess the elliptic fibration of a principal elliptic
bundle, as the following exact calculation shows.

Fix an ordinary base point for which
$\mu\notin\Q\tau+\Q$.  Such points exist: otherwise the
ordinary base would be the countable union of the analytic loci
where one rational relation $\mu=r\tau+s$ holds; Baire's
theorem would force one such relation identically, contradicting
that $P$ is a nontorsion section.  For integral vectors
$v=(a,b,c,d)$ and $v'=(A,B,C,D)$, write
\[
 \Pi v=(6a\mu+b\tau+c,\ a\beta+b\mu+d).
\]
The coefficient of $\beta$ in
$\det(\Pi v,\Pi v')$ is
\[
 (Ab-aB)\tau+(Ac-aC).
\]
If it vanishes and one of $a,A$ is nonzero, the first three
coordinates of $v,v'$ are rationally proportional.  The
remaining determinant is
$(6a\mu+b\tau+c)(D-(A/a)d)$ after choosing $a\ne0$;
the first factor is nonzero, so a zero determinant would make
the two vectors dependent.  The case $a=0,A\ne0$ is the
same statement with the vectors interchanged.  If $a=A=0$,
the determinant is
\[
 (bD-Bd)\tau+(cB-Cb)\mu+(cD-Cd),
\]
and the rational independence of $1,\tau,\mu$ again forces
all the minors to vanish.  Thus every independent integral
pair spans a complex line for at most one value of $\beta$.
There are countably many pairs.  For all other admissible
parameters the two-torus is simple and contains no elliptic
subtorus.  In particular the distinguished line $\C\delta$
meets the period lattice in $\Z\delta$, so its complexified
orbit is $\C^*$, rather than an elliptic curve.

A principal elliptic-bundle construction may therefore describe
a central cap, but cannot by itself equal the entire prescribed
punctured cusp family.  It would need an additional affine
identification that destroys its generic elliptic fibration,
together with the actual nodal $q$-interface and the finite-place
component rotations. The construction here establishes none of
those additional comparisons. \ref{push:src:B.4.33} separately rules out a proper
analytic filling of the actual global family, without denying the
auxiliary Hopf degeneration constructed above.


\Uses{the hypotheses, notation and definitions stated here; \citep[pp.~85--86, 93--94]{Fujimoto1990AB}; \citet[Theorem 1.4, preprint p.~7, and Example 2.1, p.~13]{Nishiguchi1988IIRoot}; \citet[Propositions 3.1--3.4, preprint pp.~6--9]{SS00}; \cref{push:src:B.4.33}}



\paragraph{Limits of automatic analytic comparison}
The abstract identity $\Aut(X)=\C^*$ does not supply a
meromorphic compactification of the action. On an ordinary
two-torus the distinguished orbit is a nonconstant holomorphic
map from $\C^*$. If it extended meromorphically across both
ends to $\PP^1$, it would extend holomorphically: a
meromorphic map from a smooth curve to a compact complex
manifold has a removable indeterminacy. Every holomorphic
map $\PP^1\to T^2$ is constant, since it lifts to $\C^2$.
This is a contradiction. Thus an algebraic structure on the
abstract automorphism group is insufficient here.

The compactification theorem of
\citet[Theorem 1.2]{BiliottiGhigi2021Root} assumes Fujiki
class $\mathcal C$; that hypothesis is unavailable.
The local isomorphism theorem of
\citet[Theorem 3]{Meersseman2011Root} assumes pointwise
isomorphism on the whole base, including its central point,
and constant vector-field dimension. It cannot identify the
unknown central fibre. The Hopf compactification example of
\citet[Example 1]{KatzarkovEtAl2025Root} likewise illustrates
that an analytic open part need not determine a bimeromorphic
compactification. Their definability-based comparison is not
used: no definability of our punctured comparison has been
proved. \ref{push:src:B.4.7} supplies a proper family exhibiting the relevant
failure directly. None of these automatic-comparison results
proves the conclusion of \ref{push:src:B.4.33}, whose argument instead uses
extending-section differences and the global affine loop invariant.

Non-K\"ahler modifications of elliptic degenerations are
established. The product case in
\citet[Section 3, Example 3.4(1)]{Fujimoto1988Root}, using
the preceding construction, already gives such degenerations
with elliptic abelian nearby surfaces. Its elliptic-bundle
structure prevents a direct identification with our generic
simple two-tori. For the related codimension-two extension
question, \citet[Theorem 7.1.3]{Nakayama2002Root} proves an
extension criterion for locally bimeromorphically projective
elliptic fibrations over a punctured polydisc: the fibration
must be bimeromorphically projective over that punctured base.
It rules out the sectionless elliptic example discussed there.
Our regular fibres have dimension two and no asserted elliptic
quotient, so that theorem does not establish a nonexistence
result for the present filling.

The differential-geometric smoothing theorem of
\citet[Theorem 5.5, pp.~882--883]{Doi2009Root} allows
non-K\"ahler components with dual normal bundles,
anticanonical double divisors and compatible residues.
Its total space is a smooth real six-manifold and its
map is smooth; the complex structures on the surface
fibres vary continuously in the stated sense.
This is not a holomorphic total-space filling theorem.
Example 5.3 identifies the degree-zero elliptic ruled
self-gluing with a torus and the nonzero-degree cases
with Kodaira surfaces. Neither result supplies the
prescribed holomorphic two-parameter comparison at our
missing origin.


\Uses{the hypotheses, notation and definitions stated here; \citet[Theorem 1.2]{BiliottiGhigi2021Root}; \citet[Theorem 3]{Meersseman2011Root}; \citet[Example 1]{KatzarkovEtAl2025Root}; \cref{push:src:B.4.7}; \cref{push:src:B.4.33}; \citet[Section 3, Example 3.4(1)]{Fujimoto1988Root}; \citet[Theorem 7.1.3]{Nakayama2002Root}; \citet[Theorem 5.5, pp.~882--883]{Doi2009Root}}


\begin{proposition}[a normal-line obstruction on every semistable good-reduction interface]\hypertarget{tag-0318}{}\label{push:src:B.4.8}
\StatementRecord{push:src:B.4.8}{proved}{0318}{\citet[Section 2.2]{AbashevaDeev2024AB}}

Let $a,b$ be holomorphic units on a disc $\Delta_q$, with
$0<|a(q)|<1$, and let $m\ne0$ be an integer.  Put
\[
 \mathcal E_a=(\C^*\times\Delta_q)/
       \bigl((z,q)\sim(a(q)z,q)\bigr).
\]
This is a smooth proper elliptic family over the full disc.
Over $\Delta_q^*$ define a line bundle $\mathcal L_m$ by
\[
 (z,v,q)\sim\bigl(a(q)z,\ q^m b(q)v,\ q\bigr).
\]
Its restriction to each elliptic fibre has degree zero.
After any finite base change $q=s^N\varepsilon(s)$, where
$N>0$ and $\varepsilon$ is a holomorphic unit, the following
extension is impossible: a proper semistable family
$p:\mathcal C\to\Delta_s$ of connected curves, with smooth
total surface, an isomorphism
$\mathcal C|_{\Delta_s^*}\simeq\mathcal E_a|_{\Delta_s^*}$,
and a line bundle on $\mathcal C$ restricting to
$\mathcal L_m$ under that isomorphism.

The assertion allows any finite collection of rational tails
in the central curve and arbitrary translations in the
punctured isomorphism.  It remains true after tensoring
$\mathcal L_m$ by the restriction of a line bundle extending
on the proposed interface.  It does not assert an obstruction
for arbitrary singular, non-flat, or higher-dimensional
interface networks.

In particular the axis $z_1=0$ in the diagonal Hopf surfaces
\[
 H_{q,a(q)}=(\C^2\setminus\{0\})/
       \langle(z_1,z_2)\mapsto(qz_1,a(q)z_2)\rangle
\]
cannot acquire such a semistable completion together with
its normal line, after any finite base change.
\end{proposition}


% ALL PRINTED PROOF MATERIAL MASKED









































































































































\paragraph{What a corrected cap must change}
The obstruction concerns an individual normal line,
not just the tensor product of two normal lines in
the $T^1$ condition.  Opposite windings can cancel in
that tensor product while neither line extends on its
own proposed smooth branch.  Thus cancelling the
scalar $q$ by $q^{-1}$ is insufficient for this
good-reduction interface.  A successful correction
would have to change the interface or its punctured
normal bundle outside the class specified above.
The constant-elliptic example \ref{push:src:B.4.7} has no such
winding and remains consistent with \ref{push:src:B.4.8}.


\Uses{the hypotheses, notation and definitions stated here; \cref{push:src:B.4.7}; \cref{push:src:B.4.8}}


\begin{proposition}[ordinary logarithmic torus cocycles and the remaining descent classes]\hypertarget{tag-0319}{}\label{push:src:B.4.9}
\StatementRecord{push:src:B.4.9}{proved}{0319}{\SeriesRef{push:src:B.3.1}, \SeriesRef{push:src:C.3.5}; \citep[Definition 1.1]{Fujimoto1990AB}; \citep[Section 5]{Cartan1952B_AB}}

Let $\mathcal T\to U$, $U=\Delta_q^*\times\Delta_r^*$,
be the Jacobian two-torus family of the marked cusp,
after $w=r^6$.  A holomorphic homomorphism of group
families
\[
 \C^*\times U\longrightarrow\mathcal T
\]
is determined by a flat integral cycle.  For this family
all such cycles are
\[
 c=a\widehat w+b\widehat\delta,\qquad a,b\in\Z.
\]
In the exponential cusp coordinates its homomorphism
is $z\mapsto(z^a,z^b)$.  Thus the ordinary Cartier
logarithm construction along a globally specified
cocharacter has transition maps
\[
 (Z,W)\longmapsto(g_{ij}^{\,a}Z,g_{ij}^{\,b}W).
\]
If these units extend over the prescribed fan charts,
the maps extend there and preserve each toric
boundary component.  Changing only these transition
units cannot repair the component permutations
$+4,-3,0$ refuted in \ref{push:src:B.3.1}.

This is not a classification of all torus-torsor
descent.  In fact, for the sheaf $\mathscr T$ of
holomorphic sections of $\mathcal T$,
\[
 H^1(U,\mathscr T)\simeq H^2(U,\Lambda)\simeq\Z^2.
\]
The last group is represented by the coinvariant
classes of $\widehat\gamma,\widehat u$.  The ordinary
fixed-cocharacter Cartier constructions have zero
image in this group.  The actual cusp family with
the coherent closing translations of \ref{push:src:C.3.5} also
has zero torsor class, by the explicit holomorphic
section below.  Thus nonzero classes in this
$\Z^2$ do not have its prescribed marked punctured
cusp germ.  No compactification of an arbitrary
torsor is asserted.
\end{proposition}


% ALL PRINTED PROOF MATERIAL MASKED




































































































































































\begin{proposition}[an explicit parabolic family and its nonextending exceptional correction]\hypertarget{tag-030E}{}\label{push:src:B.4.10}
\StatementRecord{push:src:B.4.10}{proved}{030E}{\SeriesRef{push:src:B.4.8}; \citet[Section 2, pp.1351--1353]{IstratiEtAl2022Root}; \citet[Introduction, pp.1750--1751]{Teleman2010AB}; \citet[Example 9.10, p.1386]{IstratiEtAl2022Root}; \citet[Proposition 5.6, pp.1369--1370]{IstratiEtAl2022Root}; \citet[Section 1.2]{DlousskyTeleman2012AB}}

Fix an integer $n\ge1$ and holomorphic units $a(s),\kappa(s)$ on a
disc, with $0<|a(s)|<1$. After shrinking the disc there is a proper
holomorphic submersion $\mathcal K\to\Delta_s$, with an embedded smooth
proper elliptic family
\[
 \mathcal E=(\C^*\times\Delta_s)/
       ((v,s)\sim(a(s)v,s)).
\]
Its central surface is a parabolic Inoue surface $S_n$, with
$\pi_1(S_n)=\Z$, $b_2(S_n)=n$, and elliptic axis of square $-n$.
For $s\ne0$, the surface $\mathcal K_s$ is obtained from the diagonal
Hopf surface $H_{\kappa(s)s^n,a(s)}$ by $n$ successive blow-ups at one
axis point and its indicated infinitely near strict-axis points.
In a fixed annulus marking of $\mathcal E$, its normal line has factor
\[
 g_s(v)=\kappa(s)(a(s)v+s)^n.
\]
This line extends across $s=0$ and is isomorphic to the holomorphic
family with factor $\kappa(s)a(s)^n v^n$.
For constant $a,\kappa$ it is a constant family of degree $-n$ lines.

The relation with the normal line of the unmodified Hopf axis is
\[
 N_{\mathcal E_s/\mathcal K_s}
   =L_{\kappa(s)s^n}\otimes\OO_{\mathcal E_s}(-n[-s]),
       \qquad s\ne0.
\]
Here $n[-s]$ is the weighted total exceptional correction, not the
intersection with the reduced final exceptional locus. It has no
Cartier extension on $\mathcal E$, even after a finite base change or
the addition of a vertical component. Thus this construction does
not contradict \ref{push:src:B.4.8}: the marked divisor used to change the normal
line fails to extend.

For $n=8$, the degree $-8$ axis normal can be chosen to be the inverse
of $\OO_E(2)$ on any chosen smooth elliptic complete intersection
$E=(2,2)\subset\PP^3$. This supplies the normal-line matching with
$\operatorname{Bl}_E\PP^3$ described below, including one positive
triple section. It is an auxiliary construction and first-order
matching result. No attachment to the prescribed punctured
$S^6$ family, compact third component, or global smoothing is asserted.
\end{proposition}


% ALL PRINTED PROOF MATERIAL MASKED



























































































































































































































\paragraph{Auxiliary normal matching and the actual filling obstruction}
Choose a smooth quadric $Q\subset\PP^3$ containing $E$.
Its strict transform in $A$ is still $Q\simeq\PP^1\times\PP^1$:
blowing up its Cartier divisor $E$ does not change $Q$.
Its normal is
$\OO_Q(2,2)(-E)\simeq\OO_Q$, and it meets $D$ in a section.
The matching surface $S_8\times\{p\}$ in $B$ also has trivial
normal. A third threefold would therefore need divisors $Q$ and
$S_8$ meeting along $E$, with normals
$\OO_Q(-E)$ and $\OO_{S_8}(-E)$ respectively.
These are exactly the normal conditions for the central surface
$Q\cup_E S_8$ in a semistable surface smoothing.
The surface smoothing is constructed in \ref{push:src:B.4.11}--\ref{push:src:B.4.12}.
Its topological Euler number is $4+8-2\chi(E)=12$ and its
holomorphic Euler number is one. Its total space over an open
disc is not a compact third threefold, so that construction
does not supply a compact threefold assembly.

A full assembly must also preserve the required coherent cohomology
from \ref{push:src:B.4.5}; a lone parabolic component attached to rational components
can lose its $H^1(\OO)$ through the elliptic overlap.
Additional conductor surfaces meeting $E$ would contribute further
triple divisors and must be included in the normal equation.
Nor is a local change of the cusp parameter automatically a permitted
global base change of the prescribed family.
These auxiliary matching statements do not produce a filling
of the actual family. \ref{push:src:B.4.33} proves that no such proper analytic
filling exists after any finite parameter cover.


\Uses{the hypotheses, notation and definitions stated here; \cref{push:src:B.4.11}; \cref{push:src:B.4.12}; \cref{push:src:B.4.5}; \cref{push:src:B.4.33}}


\begin{proposition}[globalization at the contracted parabolic elliptic axis]\hypertarget{tag-030F}{}\label{push:src:B.4.11}
\StatementRecord{push:src:B.4.11}{proved}{030F}{\SeriesRef{push:src:B.4.10}; \citep[Section~6, Hauptsatz~IIa, p.~59]{Grauert1960D}; \citet[Lemma~1, p.~93]{Manetti1991D}; \citet[Section~3.2, Proposition~3.3 and the discussion preceding
Proposition~3.4]{DlousskyTeleman2020D}}

Let $S=S_8(a,\kappa)$ be the central parabolic surface in \ref{push:src:B.4.10},
let $E\subset S$ be its elliptic axis, and put
$N=N_{E/S}$, so $\deg N=-8$.
The contraction $\pi:S\to Y$ of $E$ is a compact normal surface
with a single isolated singular point $p$.
Its germ at $p$ is the vertex of the affine cone over
$(E,N^{-1})$.  In particular,
\[
 H^2(Y,\Theta_Y)=0.
\]
The restriction morphism from a semiuniversal deformation of $Y$
to a semiuniversal deformation of $(Y,p)$ is regular.
Consequently every prescribed analytic one-parameter deformation
of this isolated singularity is induced, after shrinking its
parameter disc, by a proper deformation of the compact surface $Y$.
If the prescribed local deformation smooths the singularity, the
global deformation has smooth fibres away from the central value.

This concerns globalization of a specified local deformation of
an auxiliary compact surface.  It does not assert an attachment to
the punctured $S^6$ family or a compactification of the resulting
threefold over the open parameter disc.
\end{proposition}


% ALL PRINTED PROOF MATERIAL MASKED




















































































































































































































\begin{proposition}[a proper quadric--parabolic surface smoothing]\hypertarget{tag-030G}{}\label{push:src:B.4.12}
\StatementRecord{push:src:B.4.12}{proved}{030G}{\SeriesRef{push:src:B.4.11}, \SeriesRef{push:src:B.4.10}; \citep[Remark~7.6(iii), p.~46; Section~9, pp.~63--65]{Pinkham1974Root}}

Let $S=S_8(a,\kappa)$ be as in \ref{push:src:B.4.11}, with elliptic axis $E$
and normal line $N$ of degree $-8$.
There is an embedding $E\hookrightarrow Q=\PP^1\times\PP^1$
as a smooth divisor of bidegree $(2,2)$ with
$N_{E/Q}\simeq N^{-1}$.
There is a proper flat map from a smooth complex threefold
\[
 \varpi:\mathcal V\longrightarrow\Delta
\]
with smooth noncentral fibres and reduced simple normal-crossings
central fibre exactly $S\cup_E Q$.  Its two normal lines are
\[
 N_{S/\mathcal V}=\OO_S(-E),\qquad
 N_{Q/\mathcal V}=K_Q=\OO_Q(-2,-2).
\]
The family is obtained by globalizing a specified cone smoothing
and blowing up its vertex; no finite base change is needed.

For the elliptic-quartic configuration of \ref{push:src:B.4.10}, the three
pairwise normal-line products, including the triple-curve factors,
can consequently all be made trivial using
\[
 A=\operatorname{Bl}_E\PP^3,\qquad
 B=S\times\PP^1,\qquad C=\mathcal V.
\]
Here $C$ is a proper family over an open disc and is not a compact
threefold.  This assertion supplies the indicated surface smoothing
and normal-line products only.  It does not construct a compact
three-component central fibre, a compatible global smoothing of
such a fibre, or a comparison with the prescribed punctured
$S^6$ family.
\end{proposition}


% ALL PRINTED PROOF MATERIAL MASKED





















































































































































































\begin{proposition}[the second period does not lift to a finite marked Hopf modification]\hypertarget{tag-030H}{}\label{push:src:B.4.13}
\StatementRecord{push:src:B.4.13}{proved}{030H}{\SeriesRef{push:src:B.4.10}, \SeriesRef{push:src:A.1.2}, \SeriesRef{push:src:B.4.8}}

Fix $0<|\alpha|,|a|<1$, and put
\[
 H=(\C^2\setminus\{0\})/
       \langle(x,v)\mapsto(\alpha x,av)\rangle,
       \qquad E=\{x=0\}\simeq\C^*/a^\Z.
\]
Let $h$ be the automorphism induced by
$(x,v)\mapsto(c x,bv)$, where $b,c\in\C^*$.
Suppose $p:Y\to H$ is a finite composition of point blow-ups,
allowing infinitely near centres, and at least one point of $E$
is an image of its exceptional locus. If $[b]\in\C^*/a^\Z$ has
infinite order, there is no holomorphic automorphism of $Y$
which agrees, on the dense unmodified core, with the marked
automorphism $h$.

The same obstruction holds for a finite collection of such marked
modifications if the proposed period permutes the collection and
acts as $h$ on each marked core.
More generally, if a compact smooth surface contains a preserved
marked elliptic curve $E$ of nonzero normal degree $d$, every
holomorphic automorphism restricting to translation by $b$ on $E$
must satisfy $d\,b=0$ in $E$.
The latter assertion also excludes a finite cycle of such preserved
marked cap axes when the cycle composite is an infinite-order
translation. These assertions apply to the direct attempt
to lift the second actual cusp period to the finite Hopf
modifications in \ref{push:src:B.4.10}, for generic nonzero parameters, after every
finite base change. It concerns a holomorphic period action on
these full modifications. It does not exclude a meromorphic
correspondence, a different connected unfolding made from
noncompact pieces, or a cap present only over the central parameter.
\end{proposition}


% ALL PRINTED PROOF MATERIAL MASKED



































































































\paragraph{What an infinite unfolding must additionally prove}
There is no corresponding obstruction to merely indexing translated
modifications by all integers and letting the deck map shift the
index. That construction alone has the wrong core. Indeed suppose
$p_j:Y_j\to H$ and
$\widetilde h_j:Y_j\to Y_{j+1}$ satisfy
$p_{j+1}\widetilde h_j=hp_j$.
On the disjoint union of their unmodified cores, the quotient by
the shift is identified with the single core $H^\circ$ by
\[
 [(j,y)]\longmapsto h^{-j}p_j(y).
\]
It is not the required quotient $H^\circ/\langle h\rangle$.
Thus a useful unfolding must have additional connected overlaps
and must recover the original second-period identifications there.

Nor can one simply glue several full compact Hopf modifications
along the same dense core and assume a Hausdorff space results.
If two embedded compact spaces in a Hausdorff space share an open
subset dense in both, their images, being closed, both equal its
closure. They coincide. A construction by genuinely distinct
compact copies with only that common core consequently cannot be
such a Hausdorff gluing.
These observations alone do not classify connected unfoldings
of noncompact pieces with a new boundary geometry. Analytic
transition maps, proper discontinuity and the retained punctured
quotient are not supplied by the auxiliary compact family in
\ref{push:src:B.4.10}. Independently of that restricted classification question,
\ref{push:src:B.4.33} rules out any such construction giving a proper filling
of the actual global family.


\Uses{the hypotheses, notation and definitions stated here; \cref{push:src:B.4.10}; \cref{push:src:B.4.33}}


\begin{proposition}[tangent weights and the compact-bundle obstruction reduction]\hypertarget{tag-030I}{}\label{push:src:B.4.14}
\StatementRecord{push:src:B.4.14}{proved}{030I}{\SeriesRef{push:src:B.4.11}; \citet[Section~2, equations~(1)--(8)]{Fujiki2009D}; \citet[Theorem~1.1]{Fujiki2009D}; \citet[Theorem~(4.6), p.~566]{AtiyahSinger1968IIIE}}

For the explicit parabolic surface $S=S_8(a,\kappa)$ and its
elliptic contraction $\pi:S\to Y$ of \ref{push:src:B.4.11}, the genuine action
$\lambda\cdot(x,v)=(\lambda x,v)$ induces trivial representations
on the following cohomology groups:
\[
 H^0(S,\Theta_S(-\log E))=\C,\qquad
 H^1(S,\Theta_S(-\log E))=\C^9,
\]
\[
 H^0(Y,\Theta_Y)=\C,\qquad H^1(Y,\Theta_Y)=\C^8.
\]
The two second cohomology groups vanish.

Let $p:Z_L\to\PP^1$ be the compact holomorphic $Y$-bundle
associated to any holomorphic line bundle $L$ and this
$\C^*$ action.  Put
$\mathcal T=\Theta_{Z_L/\PP^1}$.
For every $b\in\PP^1$ one has
\[
 R^0p_*\mathcal T=\mathcal O_{\PP^1},\qquad
 R^1p_*\mathcal T=\mathcal O_{\PP^1}^{\oplus8},\qquad
 R^2p_*\mathcal T=0,
\]
and
\[
 H^2(Z_L,\mathcal T)=0,
 \qquad
 H^2\bigl(Z_L,\mathcal T\otimes
                   p^*\mathcal O_{\PP^1}(-b)\bigr)=0.
\]
These are cohomological reductions.  They do not, by themselves,
globalize a prescribed deformation along the nonisolated
singular section of $Z_L$ or construct a compact third component.
\end{proposition}


% ALL PRINTED PROOF MATERIAL MASKED








































































































































\subsection{A local insertion into the prescribed simple-torus family}

\begin{proposition}\hypertarget{tag-030J}{}\label{push:src:B.4.15}\label{push:B415}
\StatementRecord{push:src:B.4.15}{proved}{030J}{\SeriesRef{push:src:B.3.1}; \citep{Cartan1952B_AB}; \citet[Definition~1.1]{Fujimoto1990AB}}

Let $q_0$ be a regular point of the elliptic base, and let
$\Delta_z$, $z=q-q_0$, be a simply connected disc avoiding
$0,1,\infty$. After $w=r^6$, take the standard twelve-cycle
model
\[
 \mathcal X\longrightarrow\Delta_z\times\Delta_r
\]
of the prescribed local two-torus family. There is a proper
replacement $\mathcal X'\to\Delta_z\times\Delta_r$, with
a specified isomorphism to $\mathcal X$ off $(z,r)=(0,0)$,
whose total fourfold is smooth and whose divisor $r=0$ is
reduced normal crossings. The composite map to $\Delta_r$
is flat.

The replacement has twelve new compact vertical normalization
components. Six are $\PP^1$-bundles over a scalar Hopf
surface; the other six are the blow-ups of those bundles
along two disjoint elliptic curves. Every new component
has fundamental group $\Z$, and therefore fails the
$\partial\bar\partial$ lemma. The punctured family is the
actual prescribed family, so its generic simplicity is
preserved.

This is a statement about a proper model over a bidisc,
flat over the degeneration parameter. It does not assert
joint flatness over the bidisc, a one-parameter degeneration
of compact surfaces with Hopf components, or a global
$S^6$ filling. All twelve component labels, and hence the
obstruction of \ref{push:src:B.3.1}, are unchanged.
\end{proposition}



% ALL PRINTED PROOF MATERIAL MASKED






































































































































































































































\paragraph{Scope}
For a generic arc $z=kr$ through the blown-up base
origin, the strict transform meets $F$ at an ordinary
point. Its central surface remains a cycle of ruled
surfaces over the elliptic curve. The non-K\"ahler
geometry in (\ref{push:src:B.4.15}.4) is carried by the variation
around the whole exceptional $\PP^1$. Therefore a
theorem only about one-parameter limits of simple
two-tori would not, by itself, exclude these vertical
threefold components.


\Uses{the hypotheses, notation and definitions stated here; \cref{push:src:B.4.15}}


\begin{proposition}[scope of the insertion operations]\hypertarget{tag-030K}{}\label{push:src:B.4.16}
\label{push:B416}
\StatementRecord{push:src:B.4.16}{proved}{030K}{\SeriesRef{push:src:B.4.15}, \SeriesRef{push:src:B.3.1}, \SeriesRef{push:src:B.4.2}}

The local operation \ref{push:src:B.4.15} attaches a triangulated
cylinder to the original twelve-cycle in the dual
complex and admits an elementary collapse back to that
cycle. In a compact normal-crossings model, performing
it at a regular place consequently preserves the
homotopy type of the global dual complex.

The same conclusion holds for finite sequences of these
insertions, twists by a connected holomorphic group
action preserving the prescribed normal-crossings
model, and toroidal subdivisions of its central cone
complex. In particular, starting with a model whose
dual complex has a nonzero first integral cohomology
class, these operations cannot produce a semistable
$S^6$ centre. They also cannot change the component
rotations in \ref{push:src:B.3.1}.

This statement concerns the specified generated class
of operations. It makes no assertion that arbitrary
analytic caps or arbitrary changes of the punctured
comparison belong to that class.
\end{proposition}



% ALL PRINTED PROOF MATERIAL MASKED






































































































\begin{proposition}[trivial monodromy does not force equality of first Betti numbers]\hypertarget{tag-030L}{}\label{push:src:B.4.17}
\StatementRecord{push:src:B.4.17}{proved}{030L}{\SeriesRef{push:src:B.4.7}}

There is a compact connected complex threefold $Z$ and a proper
holomorphic map $f:Z\to\PP^1$ with connected fibres, with precisely
one singular fibre $Y$, such that
\[
 f^{-1}(\PP^1\setminus\{0\})
 \simeq H_-\times(\PP^1\setminus\{0\}),\qquad
 b_1(Y)=0<1=b_1(H_-).
\]
Here $H_-$ is the scalar Hopf surface with multiplier $-e^{-\pi}$,
and $Y=H_+\cup_E H_+$ is a reduced simple-normal-crossings union
with $E=\C/(\Z+i\Z)$ and $H_+$ of multiplier $e^{-2\pi}$.
In particular, trivial local monodromy on $H^1$ does not imply
$b_1(Y)=b_1(H_-)$, even with compact smooth total space.
\end{proposition}


% ALL PRINTED PROOF MATERIAL MASKED







































































\paragraph{Source audit and the valid remaining implication}
The converse just disproved is asserted in
\citet[Theorems 1.2 and 3.2, arXiv version~1]{HondaViaclovsky2025D}.
The published primary PDF excerpt, printed p.~553,
retains it as Theorem~1.1.  The two versions and the precise
access record are distinguished in
\texttt{B\_monodromy\_source\_audit.md}.
In the version read in full, equation~(3.9) establishes only
\[
 b_1(Y)\leq\dim H^1(F;\R)^{T}\leq b_1(F).
\]
Identity of $T$ gives equality in the second inequality;
it does not give equality in the first.  The inequality and
the implication from equality of the endpoint Betti numbers
to identity monodromy survive this counterexample.
The irreducible-fibre argument in its Corollary~3.7 also
survives: the top-dimensional homology of an irreducible
central fibre has rank one, forcing equality with the
invariant subspace in the displayed relative-homology
sequence.  No irreducibility is present in the counterexample.
This is a source correction relevant to possible
non-K\"ahler completions, not a nonextension theorem for
the prescribed $S^6$ family.


\Uses{the hypotheses, notation and definitions stated here; \citet[Theorems 1.2 and 3.2, arXiv version~1]{HondaViaclovsky2025D}}


\subsection{The corrected cusp on the punctured bidisc}

\begin{proposition}[a partial cusp filling]\hypertarget{tag-030M}{}\label{push:src:B.4.18}
\label{push:B418}
\StatementRecord{push:src:B.4.18}{proved}{030M}{\SeriesRef{push:src:B.3.1}, \SeriesRef{push:src:B.3.2}, \SeriesRef{push:src:B.4.9}, \SeriesRef{push:src:A.2.1}, \SeriesRef{push:src:A.1.2}, \SeriesRef{push:src:B.3.3}}

Put
\[
 B^\times=(\Delta_q\times\Delta_r)\setminus\{(0,0)\}.
\]
There is a proper Hausdorff analytic family over
$B^\times$ with smooth total fourfold, flat over $r$,
whose restriction to $r\ne0$ is the actual prescribed
cusp family after $w=r^6$, and whose reduced
normal-crossings $r$-central fibre over $q\ne0$ has
component deck permutation the shift $-1$ on $\Z/12$
for the clockwise cusp generator $g_0$ of \ref{push:src:B.3.1}.
Equivalently, the positive $q$ loop has shift $+1$.
In particular its $q$-fibre for every $r\ne0$ is the
actual reduced Mumford fibre, while its other axis has
the component transport required by \ref{push:src:B.3.2}.

The overlap is an explicit holomorphic translation,
including the retained affine marking. This proposition
does not assert a proper completion over the missing
origin, nor the global finite-place overlap maps.
\end{proposition}



% ALL PRINTED PROOF MATERIAL MASKED










































































































































































































\paragraph{Scope}
This construction fills the punctured bidisc, not its
origin. \ref{push:src:B.4.30} proves that its prescribed partial family
has no proper analytic completion at that point.
The comparison (\ref{push:src:B.4.18}.3) contains quadratic
logarithms. Neither a meromorphic graph closure
nor a common toroidal modification is thereby
available. For fixed $r\ne0$, the multiple
$q$-boundary of $\mathcal R$ cannot be substituted
for the actual reduced boundary: the valuation
argument of \ref{push:src:B.3.7} excludes a proper birational
comparison over $q$ between them. \ref{push:src:B.4.33} further
excludes every proper filling of the actual global family,
even if its other central axial model is changed. Thus
the missing origin is not left as a construction claim.


\Uses{the hypotheses, notation and definitions stated here; \cref{push:src:B.4.30}; \cref{push:src:B.4.18}; \cref{push:src:B.3.7}; \cref{push:src:B.4.33}}


\subsection{The actual global family with one base point removed}

\begin{proposition}[the global partial filling]\hypertarget{tag-030N}{}\label{push:src:B.4.19}
\label{push:B419}
\StatementRecord{push:src:B.4.19}{proved}{030N}{\SeriesRef{push:src:B.1.2}, \SeriesRef{push:src:B.4.18}, \SeriesRef{push:src:B.4.30}, \SeriesRef{push:src:B.4.33}, \SeriesRef{push:src:C.3.5}, \SeriesRef{push:src:B.1.1}, \SeriesRef{push:src:B.3.2}, \SeriesRef{push:src:B.4.9}, \SeriesRef{push:src:A.1.2}}

Let
\[
 \Omega=(\mathbb P^1\times\Delta_r)
                  \setminus\{(\infty,0)\}.
\]
There is a proper Hausdorff holomorphic map
$\mathcal Y^\circ\to\Omega$ with smooth total fourfold,
flat over $r$, whose restriction to $r\ne0$ is the
actual family pulled back by $w=r^6$. Its reduced
normal-crossings central divisor has the finite-place
models of \ref{push:src:B.1.2} and the corrected punctured cusp of
\ref{push:src:B.4.18}, with their retained markings. Over
$B^\circ=\mathbb P^1\setminus\{0,1,\infty\}$ its
relative component cover is
\[
 y^{12}=K(x),\qquad K(x)=\frac{x^4}{(x-1)^3}.
\]
Thus the global finite-place overlap maps and their
cocycle are constructed. The sole missing base point
is $(\infty,0)$. This proposition constructs only the
partial family: \ref{push:src:B.4.30} proves its noncompletion, and
\ref{push:src:B.4.33} excludes every proper filling of the actual
family over the whole $r$ disc.
\end{proposition}



% ALL PRINTED PROOF MATERIAL MASKED





























































































































































































































\subsection{An embedded Tate interface at the missing corner}

\begin{proposition}[the other Kato axis]\hypertarget{tag-030P}{}\label{push:src:B.4.20}
\label{push:B420}
\StatementRecord{push:src:B.4.20}{proved}{030P}{\SeriesRef{push:src:B.4.10}, \SeriesRef{push:src:B.4.33}; \cite[Example 9.10]{IstratiEtAl2022Root}}

In the proper smooth Kato family $\mathcal K\to\Delta_s$
of \ref{push:src:B.4.10} there is a second relative effective Cartier
divisor $\mathcal D\subset\mathcal K$, disjoint from
the persistent elliptic divisor $\mathcal E$.
The surface $\mathcal D$ is smooth and proper over
$\Delta_s$. Its general fibre and central fibre are
\[
 \mathcal D_s=\mathbb C^*/(\kappa(s)s^n)^\mathbb Z
 \quad(s\ne0),\qquad
 \mathcal D_0=\hbox{the reduced rational }n\hbox{-cycle}.
\]
At its nodes the relative local equation is $uv=s$.
Its normal line in $\mathcal K$ is the multidegree-zero
Tate line whose factor on the closing collar is $a(s)$.
In particular it has degree zero on every component
of the central cycle. Scaling the first blowdown
coordinate gives a holomorphic $\mathbb C^*$ action
on the pair $(\mathcal K,\mathcal D)$ over the base.

For $n=12$ this constructs an embedded proper copy of
the actual cusp's elliptic root interface, with its
degree-zero normal line, after choosing the units
$\kappa,a$ to be the actual Tate and normal multipliers.
In the additive period convention $\mathbf u=(\tau,\mu)$,
$W=e^{-2\pi i\zeta_2}$, these multipliers are
$\kappa(s)s^{12}=e^{2\pi i\tau}$ and
$a(s)=e^{-2\pi i\mu}$. The elliptic translation by
$6P$ is realized by first-coordinate scaling
$a(s)^{-6}$.

This is an embedded Cartier and normal-line comparison,
not an identification of relative tubular neighbourhoods.
No lift of the whole second period, attachment to the
global partial filling, or origin completion is asserted.
\end{proposition}



% ALL PRINTED PROOF MATERIAL MASKED

















































































































































































\begin{proposition}[a degree-one Kato model with a linear twelve-gon collar]\hypertarget{tag-030Q}{}\label{push:src:B.4.21}
\StatementRecord{push:src:B.4.21}{proved}{030Q}{\SeriesRef{push:src:B.4.10}, \SeriesRef{push:src:B.4.19}, \SeriesRef{push:src:B.4.33}, \SeriesRef{push:src:B.4.20}}

Start with the $n=1$ Kato family of \ref{push:src:B.4.10} over a disc
$\Delta_q$, with holomorphic units $a(q),\kappa(q)$ and
$0<|a(q)|<1$.  Pull it back by $q=s^{12}$.
There is a proper modification of this pulled-back threefold,
obtained by six point blow-ups and seventy-two smooth curve
blow-ups supported on its central rational cycle, with the
following properties.  Its total space $\mathcal K^*$ is smooth,
it is proper and flat over $\Delta_s$, and its restriction to
$s\ne0$ is the original pulled-back Kato family.

The strict transform $\mathcal D^*$ of the other Cartier axis
is the smooth semistable Tate surface with multiplier
$\kappa(s^{12})s^{12}$ and central fibre $I_{12}$.  A neighbourhood
of $\mathcal D^*$ in $\mathcal K^*$ is biholomorphic over
$\Delta_s$ to a neighbourhood of the zero section in the
multidegree-zero Tate line with closing factor $a(s^{12})$.
The comparison retains the explicit transverse coordinate
\[
 H=\frac{h}{q},\qquad h=\pi_1^*w-q,
 \tag{B.4.21.1}
\]
on the original punctured family.  Thus its normal comparison,
relative to the original punctured normal frame, has factor
$q^{-1}=s^{-12}$; its inverse has factor $q$.

The persistent elliptic axis has an unchanged neighbourhood.
In particular its normal in the strict transform of the original
central Kato surface still has degree $-1$, rather than $-12$.
The construction supplies an actual linear collar with this
marking.  It does not extend the whole second period, construct
an origin cap for \ref{push:src:B.4.19}, or identify the modified Kato family
with the prescribed torus family. Such a proper filling of
the actual family is excluded by \ref{push:src:B.4.33}; the present
statement retains its auxiliary embedded-collar scope.
\end{proposition}


% ALL PRINTED PROOF MATERIAL MASKED






































































































































































































































































\subsection{A local resolution retaining the cusp factor}

\begin{proposition}[the repaired local plumbing model]\hypertarget{tag-030R}{}\label{push:src:B.4.22}
\label{push:B422}
\StatementRecord{push:src:B.4.22}{proved}{030R}{\SeriesRef{push:src:B.4.21}, \SeriesRef{push:src:B.4.19}, \SeriesRef{push:src:B.4.20}, \SeriesRef{push:src:B.4.30}, \SeriesRef{push:src:B.4.33}}

Let
\[
 X_n=\{fuv=g s^n\}\subset\mathbb C^5,
 \qquad r=g-f,\qquad q=s^n.
\]
There is a proper bimeromorphic local modification
$\widetilde X_n\to X_n$ with smooth total fourfold
and reduced normal-crossings divisor $r=0$.
For fixed $r\ne0$ it resolves the cyclic base change
of a three-branch $q$-normal-crossings chart.
This is a local statement and does not identify its
global smooth fibres with the prescribed tori.
The divisor $s=0$ after the displayed cyclic base change
is reduced; the divisor of the original function $q=s^n$
has multiplicity $n$. No descent removing that multiplicity
is asserted.
\end{proposition}



% ALL PRINTED PROOF MATERIAL MASKED

















































































































































\begin{proposition}[exceptional-curve nonclosure and persistence in a semistable cap]\hypertarget{tag-030S}{}\label{push:src:B.4.23}
\StatementRecord{push:src:B.4.23}{proved}{030S}{\SeriesRef{push:src:B.4.21}, \SeriesRef{push:src:B.4.10}; \citet[Section 1.2]{DlousskyTeleman2012AB}; \citet[Proposition 2.3 and its proof, pp.~6--7 of the
original preprint]{DlousskyTeleman2012AB}; \citet{Merkulov1995D}}

Let $\mathcal K^*\to\Delta_s$ be the repaired degree-one
Kato family of \ref{push:src:B.4.21}, with $q=s^{12}$, and let
$\mathcal E\to\Delta_s$ be its unchanged persistent elliptic
axis. For $s\ne0$ let $C_s\simeq\PP^1$ be the exceptional
curve of the single Hopf blow-up.

The union of the $C_s$ has no proper analytic divisor
closure in $\mathcal K^*$. Further finite proper modifications
cannot give these curves a proper analytic divisor closure:
any such closure would push forward to a forbidden one.
In particular, the generic Hopf blowdown does not extend
as a relative proper holomorphic modification, or as a
relative meromorphic contraction represented by a proper
analytic graph, across $s=0$.

There is also a deformation obstruction which does not
assume a particular plumbing. Suppose a smooth complex
fourfold $\mathcal Y$ has holomorphic functions $(Q,r)$,
its divisor $r=0$ is reduced normal crossings, and one
component $K$ contains a neighbourhood of some $C_{s_0}$,
$s_0\ne0$, agreeing with the complementary Kato region
and its ordinary base coordinate $Q$. Suppose this curve
is disjoint from every conductor branch of $r=0$.
Then every sufficiently nearby fibre $(Q,r)=(Q_0,r_0)$
contains an embedded rational curve. Consequently such a
cap cannot have the prescribed ordinary two-tori as its
$r\ne0$ surface fibres. The conclusion applies to the
whole complementary Kato component with these hypotheses;
it makes no assertion about a different central component,
a curve meeting the conductor, or a replacement changing
the holomorphic base map.
\end{proposition}


% ALL PRINTED PROOF MATERIAL MASKED















































































































































\subsection{The associated logarithmic torsor at the missing origin}

\begin{proposition}[obstruction in the fixed log-torus category]\hypertarget{tag-030T}{}\label{push:src:B.4.24}
\label{push:B424}
\StatementRecord{push:src:B.4.24}{proved}{030T}{\SeriesRef{push:src:B.4.18}, \SeriesRef{push:src:B.4.19}, \SeriesRef{push:src:B.4.9}, \SeriesRef{push:src:B.4.30}, \SeriesRef{push:src:B.4.33}; \citet{KKN2008Root}; \citep[Definition 1.2.3, pp.~77--78]{KKN2008Root}; \citep[Definition 1.2.6, p.~79]{KKN2008Root}; \citep[Sections 1.3.1--1.3.4, pp.~81--82]{KKN2008Root}; \citet[Sections 7--9]{KKN2008PicardE}}

The actual local Jacobian admits a nondegenerate logarithmic
complex-torus extension over the full $(q,r)$ bidisc and hence
has proper analytic models in the sense of
\citet{KKN2008Root}. After the normalized base change
$q=Q^{12}$, the two marked axial fan presentations of
\ref{push:src:B.4.18}--\ref{push:src:B.4.19} determine a torsor under its associated
log-torus sheaf on the punctured bidisc. This torsor does
not extend across the origin as a torsor locally trivial
for the usual analytic topology. The obstruction persists
after every further finite coordinate-root base change.

Here the torsor is formed from the \emph{associated
log-torus sheaves}, not from the fixed proper fan models
themselves, which need not be group objects. The conclusion
excludes an origin cap obtained from a locally trivial
torsor of this specified kind with the given axial
identifications. No implication from an arbitrary proper
analytic cap with vertical threefold components to this
torsor category is asserted.
\end{proposition}



% ALL PRINTED PROOF MATERIAL MASKED










































































































































































































































\begin{proposition}[a one-eigenvalue Hopf boundary has no proper analytic completion]\hypertarget{tag-030U}{}\label{push:src:B.4.25}
\label{push:B425}
\StatementRecord{push:src:B.4.25}{proved}{030U}{\SeriesRef{push:src:B.4.21}, \SeriesRef{push:src:B.4.19}; \citep[Section 6, Hauptsatz I, p.~59]{Grauert1960D}; \citet[Theorem 1.6 and the noncompact analytic
extension immediately following it]{BM95AB}}

Let $\alpha(q)=q^m\kappa(q)$, where $m\ge1$ and
$\kappa$ is a holomorphic unit on a disc, and let $\beta$
be holomorphic with $0<|\beta(0)|<1$. Shrink the disc so
$\beta$ is a unit and $0<|\alpha(q)|<|\beta(q)|<1$ for
$q\ne0$. Consider the diagonal primary Hopf surfaces
\[
 H_q=(\C^2\setminus\{0\})/
       \langle(z,w)\mapsto(\alpha(q)z,\beta(q)w)\rangle.
                                                        \tag{B.4.25.1}
\]
Their holomorphic family over the punctured disc has no
proper complex-analytic completion over the full disc,
even allowing singular or nonreduced central fibres and
a singular total space. This remains true after every
finite base change ramified at zero.

In particular one cannot replace the blown-up Hopf
fibres in \ref{push:src:B.4.21} by the unblown family with one multiplier
$\kappa(q)q\to0$ and the other tending to a nonzero
contraction. This concerns that specified auxiliary
Hopf family. It is not a nonexistence theorem for a
completion of the actual two-torus family in \ref{push:src:B.4.19}.
\end{proposition}


% ALL PRINTED PROOF MATERIAL MASKED






















































































































































\paragraph{Scope and comparison with the classical toric Hopf degeneration}
There is no conflict with the degeneration in
\citet[Section 13, pp.~116--118]{Oda1978D}.
There the boundary rays are generated by $n$ and
$-n+a n'$ with $a>0$, and the period is translation
by $\lambda^{n'}$. In the universal-cover basis
\[
 v_1=n,\qquad v_2=-n+a n',\qquad
 n'=\frac{v_1+v_2}{a},
\]
the base change $\lambda=t^a$ makes this period
act with eigenvalues $(t,t)$, up to fixed torus
units in the indicated variant. Both eigenvalues
tend to zero. \ref{push:src:B.4.25} assumes instead a nonzero
limit for the second eigenvalue. The distinction
is forced by the exact period marking of the
candidate Tate interface, and cannot be removed
by citing the existence of some Hopf degeneration.
The blown-up Kato family of \ref{push:src:B.4.10} is also consistent
with \ref{push:src:B.4.25}: its generic fibres are blow-ups of
$H_q$, so the anticanonical sections must satisfy
the extra vanishing conditions at the marked
blow-up centre and are not the spaces in
(\ref{push:src:B.4.25}.5).


\Uses{the hypotheses, notation and definitions stated here; \citet[Section 13, pp.~116--118]{Oda1978D}; \cref{push:src:B.4.25}; \cref{push:src:B.4.10}}


\subsection{A necessary vertical fibre over the missing origin}

\begin{proposition}[every origin component is a threefold]\hypertarget{tag-030V}{}\label{push:src:B.4.26}
\label{push:B426}
\StatementRecord{push:src:B.4.26}{proved}{030V}{\SeriesRef{push:src:B.4.19}, \SeriesRef{push:src:B.4.30}, \SeriesRef{push:src:B.4.33}, \SeriesRef{push:src:B.4.24}; \citet[Chapter II, Sections 7--8]{Demailly2012E}}

Suppose the actual partial family of \ref{push:src:B.4.19} has a
proper analytic completion $f:\mathcal Z\to\Delta_q\times\Delta_r$
with smooth pure four-dimensional total space, with
$r\circ f$ flat, and with its prescribed identification
over the punctured bidisc. The map to this fixed bidisc
is assumed holomorphic. Then every irreducible component
of $f^{-1}(0,0)$ has dimension three. In particular,
there is no equidimensional completion, and no completion
having even one two-dimensional component of the origin
fibre.

The conclusion still permits the vertical threefold
components allowed in B.4. It is not an assertion that
an arbitrary $r$-proper filling initially has a holomorphic
extension of $q$; if $q$ is only meromorphic, a graph
modification must first be supplied and its resulting
geometry retained. \ref{push:src:B.4.30} subsequently excludes even
the permitted vertical-component completions of the
fixed partial family, and \ref{push:src:B.4.33} excludes every proper
filling of the actual global family.
\end{proposition}



% ALL PRINTED PROOF MATERIAL MASKED


































































































































































\subsection{The obstruction survives modifications of the base}

\begin{proposition}[no log-torsor repair by point blowups]\hypertarget{tag-030W}{}\label{push:src:B.4.27}
\label{push:B427}
\StatementRecord{push:src:B.4.27}{proved}{030W}{\SeriesRef{push:src:B.4.24}, \SeriesRef{push:src:B.4.30}, \SeriesRef{push:src:B.4.33}; \citep[Lemma 1.2.9, pp.~80--81]{KKN2008Root}}

Let $B=\Delta_Q\times\Delta_R$, with
$q=Q^{12m}$ and $r=R^n$, where $m,n$ are positive
integers. Let $\beta:S\to B$ be any finite sequence
of point blowups whose centres lie over the origin.
Give $S$ the divisorial log structure of the reduced
inverse image of $QR=0$, and let $\mathcal A_S$ be
the pullback of the actual Jacobian log torus of
\ref{push:src:B.4.24}. Then
\[
 H^1(S,\mathcal A_S)=0
 \tag{B.4.27.1}
\]
for the usual analytic topology. In particular the
actual associated axial log-torus torsor on the
punctured bidisc has no locally trivial
$\mathcal A_S$-torsor extension after these
modifications.

This result concerns the specified associated
log-torus sheaf. It does not assert that flattening
or resolving an arbitrary analytic origin cap
produces a torsor under that sheaf.
\end{proposition}



% ALL PRINTED PROOF MATERIAL MASKED










































































































































































































\begin{proposition}[maximal monodromy with trivial central dualizing bundle]\hypertarget{tag-030X}{}\label{push:src:B.4.28}
\StatementRecord{push:src:B.4.28}{proved}{030X}{\citet[Definition 1.1 and Lemmas 1.2, 2.2, 2.7]{Friedman2016AB}; \citet[Corollary 2.25, Lemma 2.26 and Proposition
2.27]{Dloussky2005AB}; \citet[Sections 3.2--3.3]{FujikiPontecorvo2010AB}; \citet[Section 15, Theorem 15.1 and the following
paragraph, pp.~130--131]{Oda1978D}; \citet[Chapter III, Theorem 4.12(ii)]{BS76}; \citep[Chapter III, Theorem 4.12(iii)]{BS76}}

Let $p:\mathcal X\to\Delta$ be a proper holomorphic map
from a smooth complex threefold, smooth over $\Delta^*$,
whose nonzero fibres are complex two-tori. Suppose that
$D=p^{-1}(0)$ is a reduced normal-crossing divisor and
\emph{assume in addition} that $\omega_D\simeq\OO_D$.
If its monodromy on $H_1(\mathcal X_t,\Q)$ satisfies
\[
 \operatorname{rank}(T-1)=2,                         \tag{B.4.28.1}
\]
then every normalization component of $D$ is a smooth
rational toric surface, its conductor is its toric
boundary, and the dual complex of $D$ is a topological
two-torus. In particular all normalized components are
K\"ahler.

No K\"ahler hypothesis on $\mathcal X$ is used. The
additional triviality of $\omega_D$ is not inferred from
the neighbouring two-tori. The proposition does not
produce a modification with this property from an
arbitrary analytic completion.
\end{proposition}


% ALL PRINTED PROOF MATERIAL MASKED








































































































































































































































































\paragraph{The separate canonical-model reduction}
This conditional proposition is not the noncompletion
argument for \ref{push:src:B.4.19}. Suppose a relative holomorphic two-form
on a hypothetical transverse punctured slice has first
been proved to extend meromorphically. Rescale it so
that its vertical divisor is $\sum_i a_iV_i$ with
$\min_i a_i=0$. On each component its restriction gives
\[
 K_{V_i}+C_i\sim\sum_{j\ne i}(a_j-a_i)C_{ij}.          \tag{B.4.28.11}
\]
Here $\OO_{\mathcal X}(V_i)|_{V_i}\sim
-\sum_{j\ne i}C_{ij}$ since the full fibre is principal.
The minimal coefficient does not make this right side
zero: neighbouring higher-coefficient components remain.
Contracting them, while retaining semistability and
properness in an arbitrary non-K\"ahler threefold, is
not proved here. Nor is arbitrary meromorphic extension
of a punctured comparison assumed. The K\"ahler,
$K$-trivial torus classification of
\citet[Theorem 3.1]{KamenovaToriRoot} does not supply
that missing reduction. The unrestricted nonexistence
of the actual sphere filling in \ref{push:src:B.4.33} does not require
such a canonical model; it uses extending-section traces
and the global affine loop invariant instead.


\Uses{the hypotheses, notation and definitions stated here; \cref{push:src:B.4.19}; \citet[Theorem 3.1]{KamenovaToriRoot}; \cref{push:src:B.4.33}}



\subsection{What invariant forms do not control}

\begin{proposition}[full-rank periods do not bound an affine comparison]\hypertarget{tag-030Y}{}\label{push:src:B.4.29}
\label{push:B429}
\StatementRecord{push:src:B.4.29}{proved}{030Y}{\SeriesRef{push:src:B.4.30}, \SeriesRef{push:src:B.4.31}, \SeriesRef{push:src:B.4.33}; \citet[Sections 1.1.1--1.1.3, pp.~73--74]{KKN2008Root}}

There is a proper semistable family of two-dimensional
complex tori over a disc, with smooth total threefold
and rank-two unipotent monodromy on $H^1$, and a
fibre-preserving automorphism over the punctured disc
with the following properties. It acts identically
on all relative invariant one- and two-forms and
on the Gauss--Manin variation. Their coherent
extensions supplied by the proper family therefore
give the same lattices under this identification.
An absolute holomorphic canonical form extending
across the centre is also unchanged. Nevertheless,
the automorphism has no meromorphic extension to
the total family.

Thus these invariant-form and connection data do
not alone control an affine origin or a marked
total-space comparison, even in the presence of
a proper smooth semistable cap. This statement
does not refute a separate extension theorem
for the difference of two sections which both
extend holomorphically into the same cap.
\end{proposition}



% ALL PRINTED PROOF MATERIAL MASKED











































































































































































































\subsection{Noncompletion of the prescribed partial filling}

\begin{proposition}[the fixed partial filling cannot be completed]\hypertarget{tag-0310}{}\label{push:src:B.4.30}
\label{push:B430}
\StatementRecord{push:src:B.4.30}{proved}{0310}{\SeriesRef{push:src:B.4.19}, \SeriesRef{push:src:B.4.4}, \SeriesRef{push:src:B.4.31}, \SeriesRef{push:src:B.4.33}, \SeriesRef{push:src:B.4.27}, \SeriesRef{push:src:B.4.29}, \SeriesRef{push:src:B.4.24}, \SeriesRef{push:src:B.4.26}, \SeriesRef{push:src:B.4.32}; \citet[Theorems 6.2 and 6.3,
pp.~32 and 34]{IKN2005E}; \citet{IKN2007ErratumE}; \citet[Definition 7.1(i) and Section 9.2,
pp.~164--165 and 173]{KKN2008PicardE}; \citet[Corollary 1.2 and its proof,
manuscript pp.~3--6]{Parusinski1990E}; \citet[Theorems 1.6 and 1.10]{BM95AB}; \citet[Lemma 2.14, pp.~6--7]{EH2024E}}

The partial filling of \ref{push:src:B.4.19} has no proper
Hausdorff complex-analytic completion over its
missing base point. This remains true after finite
coordinate roots of the two base coordinates and
allows arbitrary additional three-dimensional
components of the origin fibre. No K\"ahler or
projective hypothesis is imposed on a putative
completion.

This is a statement about the prescribed partial
filling, including its two punctured axial models.
This proposition alone does not exclude every filling
over the $r$ disc: a different completion along
$r=0$, $q\ne0$ need not have the axial model of
\ref{push:src:B.4.19}. Resolving the graph in \ref{push:src:B.4.4} does not by
itself supply that identification. \ref{push:src:B.4.31}--\ref{push:src:B.4.33}
remove the need for it by controlling meridians and
the global affine loop class, and prove unrestricted
nonexistence for the actual family.
\end{proposition}



% ALL PRINTED PROOF MATERIAL MASKED



































































































































































































































































































































































































































































































































\subsection{Global elliptic monodromy constrains every central meridian}

\begin{proposition}[the global meridians have only the prescribed toric direction]\hypertarget{tag-0311}{}\label{push:src:B.4.31}
\StatementRecord{push:src:B.4.31}{proved}{0311}{\SeriesRef{push:src:B.4.4}, \SeriesRef{push:src:B.4.19}, \SeriesRef{push:src:B.4.9}; \citet[Theorems 1.6 and 1.10 and the paragraph following
Theorem 1.6]{BM95AB}}

Suppose $g:\mathcal Z\to\Delta_r$ is any proper complex-analytic
filling of a finite parameter cover of the actual six-sphere
family. Retain its reduced main component, normalize and resolve
it, extend its intrinsically marked map to $\PP^1$ by \ref{push:src:B.4.4},
resolve the graph, and resolve the support of its central divisor. Denote
the resulting smooth proper modification by $\widetilde{\mathcal Z}$,
and its holomorphic marked map by $h$. All modifications can be
chosen to be isomorphisms over $r\ne0$. The central support has
normal crossings; its multiplicities need not be one.

There is a finite subset $\Sigma\subset\PP^1$, containing
$0,1,\infty$, with the following property. For every
$q\in B^\circ:=\PP^1\setminus\Sigma$, every local angular
vanishing circle of the $r$-central normal-crossings divisor,
viewed in the homology of the nearby ordinary two-torus, lies in
\[
             \Z\widehat\delta\subset H_1(X_{q,r};\Z).
 \tag{B.4.31.1}
\]
The same conclusion holds for the image of the entire local
angular torus at a higher crossing over $B^\circ$.
Here $\widehat\delta$ is the primitive cycle in the audited
period marking. The conclusion is unaffected by any further
finite cover of the degeneration parameter.

This is a constraint on every hypothetical global filling. It
does not assert that the unknown rank-one axial completion is
the prescribed model of \ref{push:src:B.4.19}, or that an affine origin in one
model extends into the other.
\end{proposition}


% ALL PRINTED PROOF MATERIAL MASKED







































































































































































































\paragraph{Scope of the affine comparison}
The finite-orbit mechanism uses the global compact marked
base and its actual elliptic monodromy. A rank-one degeneration
over an isolated ordinary base disc has no such monodromy
constraint. Nor does (\ref{push:src:B.4.31}.1) distinguish an affine origin:
even if differences of two sections extending into the same
cap have controlled logarithmic characters, a common
nonextending translation cancels from those differences.
Controlling its affine loop relative to the particular
origin of \ref{push:src:B.4.19} is an additional step, proved in \ref{push:src:B.4.32}.
The invariant-form counterexample \ref{push:src:B.4.29} is consistent with
this limitation. \ref{push:src:B.4.33} uses the loop class, rather than
an extending comparison of individual origins, to prove
unrestricted nonexistence for Part B.


\Uses{the hypotheses, notation and definitions stated here; \cref{push:src:B.4.31}; \cref{push:src:B.4.19}; \cref{push:src:B.4.32}; \cref{push:src:B.4.29}; \cref{push:src:B.4.33}}


The finite matrix calculation, including arbitrary positive
powers of the two unipotents, is reproduced in
\path{code/ab_push_global_meridians.py}; its saved output and
JSON record all checks. The selected fixed-space minor is
$-k\ell\ne0$. An initial sign-only assertion mismatch is
retained with its exact script and output; it did not affect
the nonvanishing argument.

\subsection{The radial-loop class of every global cap}

\begin{proposition}[vanishing of the global affine loop defect]\hypertarget{tag-0312}{}\label{push:src:B.4.32}
\label{push:B432}
\StatementRecord{push:src:B.4.32}{proved}{0312}{\SeriesRef{push:src:B.4.31}, \SeriesRef{push:src:B.4.19}, \SeriesRef{push:src:B.4.4}, \SeriesRef{push:src:B.1.2}, \SeriesRef{push:src:B.4.30}; \citet[Theorem 6.3]{IKN2005E}}

Suppose a proper complex-analytic filling of a finite parameter
cover of the actual six-sphere family exists. Make the
modifications in \ref{push:src:B.4.31} and delete its finite exceptional set
from the marked base. Use the ordinary local origins of
\ref{push:src:B.4.19} as reference. There is a well-defined normalized
radial-loop defect of a local extending cap section,
\[
 \lambda\in
 V:=\bigl(\Lambda\otimes\Q\bigr)/\Q\widehat\delta .
 \tag{B.4.32.1}
\]
It is independent of the section, of the finite parameter root
needed to obtain that section, and of the modifications used.
It patches to a flat section of the actual local system $V$
over the remaining ordinary base, and this section is zero.

In particular, let $t$ be the trace of an extending finite
multisection of degree $d$, and let $o$ be the reference ordinary
origin. The radial-loop defect of $t-do$ in $V$ is zero.
The same holds for an origin obtained from multiplied such
traces by bounded logarithmic differences along $r=0$.
This conclusion controls a loop class; it does not assert
that $t-do$ is bounded or extends holomorphically at $r=0$.
\end{proposition}



% ALL PRINTED PROOF MATERIAL MASKED

































































































































































































































































\paragraph{Scope of the loop calculation}
No assertion of a bounded, meromorphic or holomorphic
comparison with the reference origin has been made.
Essential translations with zero loop defect remain
possible. The conclusion uses the global marked
$\PP^1$ base and its two peripheral matrices; an isolated
rank-one degeneration over an ordinary base disc has
no such constraint. Arbitrary singular proper fillings
are covered by resolving their reduced main component
before \ref{push:src:B.4.4} and \ref{push:src:B.4.31}; their central multiplicities
are retained until the local finite-root calculations.


\Uses{the hypotheses, notation and definitions stated here; \cref{push:src:B.4.4}; \cref{push:src:B.4.31}}


The exact quotient fixed spaces, the root-transfer identity,
the finite connector cancellation and the cusp loop label
are checked in
\path{code/d_push_global_loop_defect.py}, with its saved
output and JSON record. The proof of the global geometric
comparison is the argument above.

\subsection{No proper analytic filling of the actual punctured family}

\begin{theorem}[nonexistence after every finite parameter cover]\hypertarget{tag-0313}{}\label{push:src:B.4.33}
\label{push:B433}
\StatementRecord{push:src:B.4.33}{proved}{0313}{\SeriesRef{push:src:B.4.29}, \SeriesRef{push:src:B.4.32}, \SeriesRef{push:src:B.4.31}, \SeriesRef{push:src:B.4.4}, \SeriesRef{push:src:B.4.24}, \SeriesRef{push:src:B.4.30}, \SeriesRef{push:src:B.4.27}, \SeriesRef{push:src:B.4.19}, \SeriesRef{push:src:B.4.26}; \citet[Theorems 1.6 and 1.10]{BM95AB}}

The actual punctured six-sphere family has no proper
Hausdorff complex-analytic filling after any finite
cover of its degeneration parameter. In particular,
there is no positive integer $N$ for which the cover
$w=r^N$ has a proper filling with smooth total space
and reduced normal-crossings central fibre.

The assertion permits an arbitrary singular central
space, arbitrary additional components, and arbitrary
central multiplicities. It does not impose a
K\"ahler, projective, or logarithmic hypothesis on
the putative filling. Its proof derives the specific
logarithmic torsor needed for the contradiction from
differences of sections extending into that filling.
\end{theorem}



% ALL PRINTED PROOF MATERIAL MASKED






























































































































































































































































































































































































































The final step uses two separate facts: a
cap-derived origin has zero global radial
defect by \ref{push:src:B.4.32}, and a section extending
along the actual cusp has zero radial
ghost class by \ref{push:src:B.4.24}. Neither fact
identifies the unknown central model with
\ref{push:src:B.4.19}. A common essential translation
with zero loop class, as allowed by
\ref{push:src:B.4.29}, is harmless in this comparison.
The untwisted Engel family provides a
useful scope check: its identity closing
has zero affine origins on both axial
fan models, hence $F=0$. The nonzero
clutching class used in (\ref{push:src:B.4.33}.14)
is specific to the actual twisted
family; the proof does not exclude
the untwisted family's known filling.


\begin{corollary}[counterexample; the premise of a global central fibre]\hypertarget{tag-031A}{}\label{push:src:B.5.1}\label{push:B51}
\StatementRecord{push:src:B.5.1}{proved}{031A}{\SeriesRef{push:src:B.4.33}}

There is no central fibre or descent resolution
of the requested proper analytic filling of the
actual $S^6$ family. Its proposed integral homology,
fundamental group, central Euler number and
nearby-cycle comparison are therefore undefined.
\end{corollary}


% ALL PRINTED PROOF MATERIAL MASKED















\section{The boundary of the parameter disc}\hypertarget{tag-02Z8}{}\label{push:section:D}

We use the standard twists throughout this section. Put
\(H_*=(3/\pi)\log6\), fix a boundary parameter
\(c_*=x+iH_*\), and set \(a=x+3/2\). The distinction between a
single degenerate fibre and the total quotient over its root disc is
essential below. All topological group quotients have their quotient
topology; a Hausdorff reflection means a continuous quotient universal
for continuous maps to Hausdorff spaces.

\subsection{The intrinsic obstruction to extension}\hypertarget{tag-02Z9}{}\label{push:section:D.1}

\begin{proposition}\hypertarget{tag-032A}{}\label{push:src:D.1.1}
\StatementRecord{push:src:D.1.1}{proved}{032A}{\SeriesRef{fam:src:A.4.1}, \SeriesRef{fam:src:C.1.1}, \SeriesRef{fam:src:C.2.3}, \SeriesRef{fam:src:C.3.4}}

For \(0<|w|<216\), the reduced multiplicity-four surface has
\[
 j(\operatorname{Alb}S_2(w))=J(iw/216),\qquad
 J(q)=q^{-1}+744+196884q+\cdots.
\]
This holomorphic function has no meromorphic continuation through any
point of \(|w|=216\).
\end{proposition}


% ALL PRINTED PROOF MATERIAL MASKED
































\begin{lemma}\hypertarget{tag-032B}{}\label{push:src:D.1.2}
\StatementRecord{push:src:D.1.2}{proved}{032B}{the definitions and constructions stated in the proof}

Let \(V\) be a connected Riemann surface and
\(\pi:\mathcal X\to V\) a flat proper holomorphic map of reduced
complex spaces. Suppose that over a nonempty open subset \(V_0\), the
fibres are smooth compact threefolds whose unique anticanonical divisor
is \(2S_v\), where \(S_v\) is a smooth connected surface with
\(b_1(S_v)=2\). Then there is an open subset \(V^{\rm good}\subset V\), with at
most countable complement, on which the reduced support of the
relative anticanonical divisor is a smooth proper surface family. Its
Albanese elliptic invariant is holomorphic.
\end{lemma}


% ALL PRINTED PROOF MATERIAL MASKED

















































































\begin{theorem}[nonextension across the parameter boundary]\hypertarget{tag-032C}{}\label{push:src:D.1.3}
\StatementRecord{push:src:D.1.3}{proved}{032C}{\SeriesRef{fam:src:A.1.1}, \SeriesRef{fam:src:A.1.2}, \SeriesRef{fam:src:A.4.1}, \SeriesRef{push:src:D.1.2}, \SeriesRef{push:src:D.1.1}}

For every \(w_0\) with \(|w_0|=216\) and every connected open
neighbourhood \(U\) of \(w_0\), there is no flat proper holomorphic
map \(\mathcal X\to U\) from a reduced complex space whose fibres
over \(U\cap\{0<|w|<216\}\) are biholomorphic, fibre by fibre, to
\(X_w\). In particular an extension isomorphic to the given family
over that intersection does not exist.
\end{theorem}


% TARGET PROOF MASKED: push:src:D.1.3
























\subsection{The local quotient and its Hausdorff reflection}\hypertarget{tag-02ZA}{}\label{push:section:D.2}

\begin{lemma}\hypertarget{tag-032D}{}\label{push:src:D.2.1}
\StatementRecord{push:src:D.2.1}{proved}{032D}{\SeriesRef{fam:src:C.2.2}}

Near the order-four root centre there is a holomorphic coordinate
\[
 h=\mu-(1-\tau)/2
\]
with nonzero derivative in \(s_2\). In the fixed complex-linear fibre
coordinates \(\eta=\zeta_1\), \(\xi=\zeta_2+\zeta_1/2\), and the
integral ordered basis
\((\widehat w,\widehat u,\widehat\delta,
\widehat\gamma+3\widehat u-3\widehat w)\), the periods at \(c_*\) are
\begin{equation}\hypertarget{tag-01ON}{}\label{eq:d-period-normal}
 (1,\tfrac12),\quad(\tau,\tfrac12+h),\quad(0,1),\quad(6h,a+B(h)),
\end{equation}
where
\[
 \tau=i-\frac{3\pi i}{2}h^2+O(h^4),\qquad
 B(h)=-3ih^2+O(h^4),\qquad
 D_{c_*}=-3|h|^2+O(|h|^4).
\]
For \(c\) near \(c_*\), only the last \(\xi\)-period changes, by
\(c-c_*\).
\end{lemma}


% ALL PRINTED PROOF MATERIAL MASKED


























































\begin{proposition}[nonseparation beyond the rank-drop line]\hypertarget{tag-032E}{}\label{push:src:D.2.2}
\StatementRecord{push:src:D.2.2}{proved}{032E}{\SeriesRef{push:src:D.2.1}}

The assertion that total-space nonseparation is confined to
\(\mathbb C(2,-1)^t\), modulo central periods, is false. At the
boundary value \(a=0\), the central points with \((\eta,\xi)=(0,0)\)
and \((0,-3i)\) cannot be separated in the root-family quotient.
They are not related by that complex line and the central periods.
\end{proposition}


% ALL PRINTED PROOF MATERIAL MASKED




















\begin{theorem}\hypertarget{tag-032F}{}\label{push:src:D.2.3}
\StatementRecord{push:src:D.2.3}{proved}{032F}{\SeriesRef{push:src:D.3.1}, \SeriesRef{eq:d-dual-exact}, \SeriesRef{push:src:D.2.2}}

Let \(\mathcal T=(\Delta_h\times\mathbb C^2)/\Lambda\) at the fixed
boundary parameter \(c_*\), after shrinking \(\Delta_h\). Two
central points represented by \((\eta,\xi)\) and
\((\eta',\xi')\) are nonseparated precisely when
\[
 \operatorname{Re}(\xi'-\xi)\in G_a
       :=\tfrac12\mathbb Z+a\mathbb Z\quad(a\in\mathbb Q),
\]
and every pair is nonseparated when \(a\notin\mathbb Q\).
There is no nonseparation between distinct base fibres, or between
distinct points of a noncentral fibre. The Hausdorff reflection
\(\mathcal T^{\mathrm H}\to\Delta_h\) has the original complex tori
over \(h\ne0\), and its central fibre is
\[
 \mathbb R/G_a\simeq S^1\quad(a\in\mathbb Q),\qquad
 \{*\}\quad(a\notin\mathbb Q).
\]
If \(a=m/n\) is in lowest terms, \(n>0\), then
\(G_a=\operatorname{lcm}(2,n)^{-1}\mathbb Z\).
\end{theorem}


% TARGET PROOF MASKED: push:src:D.2.3
























































\begin{proposition}\hypertarget{tag-032G}{}\label{push:src:D.2.4}
\StatementRecord{push:src:D.2.4}{proved}{032G}{\SeriesRef{fam:src:C.1.1}, \SeriesRef{fam:src:C.2.3}, \SeriesRef{fam:src:C.3.1}, \SeriesRef{fam:src:C.3.2}}

At the boundary, the construction over \(B\setminus\{1\}\) is a
complex manifold, since \cref{fam:src:C.1.1}, \cref{fam:src:C.2.3} makes every remaining finite period map
invertible and the cusp completion is unchanged.
The order-four deck action descends continuously to
\(\mathcal T^{\mathrm H}\). On its rational central circle it is
translation by \(-a/4\); on the irrational central point it is
trivial. Its finite quotient is Hausdorff with central fibre a circle
or a point, respectively. Neither quotient is a complex space over
the base compatible with the given holomorphic family off the centre.
The linear foliation by \(\mathbb C(2,-1)^t\) on the original central
torus cover has compact elliptic leaves. For rational \(a\) the
leaves are closed; for irrational \(a\) the closure of a leaf is the
set with its fixed value of \(\operatorname{Im}\xi\).
\end{proposition}


% ALL PRINTED PROOF MATERIAL MASKED


















































\subsection{All local limiting period groups}\hypertarget{tag-02ZB}{}\label{push:section:D.3}

\begin{theorem}\hypertarget{tag-032H}{}\label{push:src:D.3.1}
\StatementRecord{push:src:D.3.1}{proved}{032H}{\SeriesRef{push:src:D.2.1}, \SeriesRef{eq:d-dual-exact}}

Let \(c_n\to c_*\) with \(\operatorname{Im}c_n<H_*\), and let
\(h_n\to0\). Set
\[
 \delta_n=-D_{c_n}(h_n)>0,\qquad
 V_n=\left(\frac{6\operatorname{Re}h_n}{\delta_n},
           \frac{6\operatorname{Im}h_n}{\delta_n},
           \frac{\operatorname{Re}(c_n-c_*)}{\delta_n}\right).
\]
Define the discrete group
\[
 L_a=\{(p,q,r)\in\mathbb R^3:
    r\in\mathbb Z,\ p+r/2\in\mathbb Z,
    \ q+r/2\in\mathbb Z,\ ar\in\mathbb Z\}.
\]
It has rank three for rational \(a\), and is
\(\mathbb Z^2\times\{0\}\) for irrational \(a\).
Every convergent subsequence of the closed lattices
\(\Gamma_n=\Pi_{c_n}(h_n)\mathbb Z^4\subset\mathbb C^2\) has the
following form. There is a primitive subgroup \(K\subset L_a\) and
a homomorphism \(\ell:K\to\mathbb R\) such that
\begin{equation}\hypertarget{tag-01OM}{}\label{eq:d-limit-group}
 \Gamma_\infty=
 \{(\eta,\xi):p\operatorname{Re}\eta+q\operatorname{Im}\eta
       +r\operatorname{Re}\xi+\ell(p,q,r)\operatorname{Im}\xi
       \in\mathbb Z\text{ for every }(p,q,r)\in K\}.
\end{equation}
Here
\[
 K=\{\lambda\in L_a:\lambda\cdot V_n
                      \text{ has a finite limit}\},\qquad
 \ell(\lambda)=\lim_n\lambda\cdot V_n,
\]
and for \(\lambda\notin K\) the absolute value tends to infinity.
Conversely every primitive \(K\subset L_a\) and every such
\(\ell\) occur for sequences with \(\operatorname{Im}c_n<H_*\).
Thus the possible pointed additive quotient limits are real tori of
any dimension from zero to three for rational \(a\), and from zero
to two for irrational \(a\), with the marked groups given explicitly
by \eqref{eq:d-limit-group}.
\end{theorem}


% ALL PRINTED PROOF MATERIAL MASKED










































































































\begin{corollary}\hypertarget{tag-032I}{}\label{push:src:D.3.2}
\StatementRecord{push:src:D.3.2}{proved}{032I}{\SeriesRef{push:src:D.2.3}, \SeriesRef{push:src:D.3.1}, \SeriesRef{push:src:D.1.3}}

There is no path-independent limiting group at the boundary even
when \(c_n\to c_*\) from the admissible side. At \(h_n=0\) and
\(c_n=c_*-i\epsilon_n\), all of \(L_a\) survives, giving a
three-dimensional real torus if \(a\) is rational and a
two-dimensional real torus if \(a\) is irrational. Other admissible
paths give a point. The fixed-boundary Hausdorff reflection of \ref{push:src:D.2.3}
therefore cannot be identified with all pointed fibre limits.
\end{corollary}


% ALL PRINTED PROOF MATERIAL MASKED













\subsection{The supercritical locus}\hypertarget{tag-02ZC}{}\label{push:section:D.4}

\begin{proposition}\hypertarget{tag-032J}{}\label{push:src:D.4.1}
\StatementRecord{push:src:D.4.1}{proved}{032J}{\SeriesRef{push:src:D.2.1}, \SeriesRef{fam:src:C.1.1}, \SeriesRef{fam:src:C.2.3}}

For \(\operatorname{Im}c=H_*+\epsilon\), with
\(0<\epsilon\ll1\), the zero locus \(D_c=0\) is one real-analytic
simple closed curve surrounding \(t=1\). In the coordinate of \ref{push:src:D.2.1},
\[
 |h|^2=\epsilon/3+O(\epsilon^2).
\]
In the fixed coarse coordinate it has the leading shape
\[
 |t-1|=\frac{\pi^4E_4(i)}4\epsilon^2
                         \bigl(1+O(\epsilon)\bigr).
\]
The error is uniform along the curve. The region containing \(t=1\)
has \(D_c>0\); the region containing the cusp has \(D_c<0\).
\end{proposition}


% TARGET PROOF MASKED: push:src:D.4.1





























\begin{proposition}\hypertarget{tag-032K}{}\label{push:src:D.4.2}
\StatementRecord{push:src:D.4.2}{proved}{032K}{\SeriesRef{fam:src:C.3.1}, \SeriesRef{fam:src:C.3.2}, \SeriesRef{fam:src:C.1.1}, \SeriesRef{fam:src:C.2.3}, \SeriesRef{fam:src:A.1.1}, \SeriesRef{fam:src:A.1.2}}

For any height \(H=\operatorname{Im}c>H_*\), the zero locus is the
well-defined modular-norm level set
\[
 Z_H=\left\{t:\mathcal N(\mu(t),\tau(t))
                    =2\sqrt3\,e^{-\pi H/6}\right\},\qquad
 \wp_{\tau(t)}(\mu(t))=0.
\]
It is compact away from the cusp, and is a real-analytic set in the
finite root charts. At every regular point it is a smooth real curve.
The construction gives a complex manifold precisely over
\(B\setminus Z_H\), retaining the original Mumford fibre at the cusp.
The open sublevel set \(\{D_c<0\}\), with the cusp adjoined, is
connected. At every point of \(Z_H\), the period quotient is the
rank-three quotient of \cref{fam:src:C.3.1}, \cref{fam:src:C.3.2}; its exact Hausdorffization is described
below.
\end{proposition}


% TARGET PROOF MASKED: push:src:D.4.2













































\begin{proposition}\hypertarget{tag-032L}{}\label{push:src:D.4.3}
\StatementRecord{push:src:D.4.3}{proved}{032L}{\SeriesRef{fam:src:C.3.1}, \SeriesRef{fam:src:C.3.2}}

At a zero of \(D_c\), let \(A:\mathbb R^4\to\mathbb C^2\) be
the real period map and \(\ker A=\mathbb R k\), with the kernel
vector stated in \cref{fam:src:C.3.1}, \cref{fam:src:C.3.2}. Define
\[
 L_k=\{m\in\mathbb Z^4:m\cdot k=0\},\qquad r=\operatorname{rank}L_k,
 \qquad V_k=\bigcap_{m\in L_k}\ker m.
\]
Then
\[
 \overline{A\mathbb Z^4}=A(\mathbb Z^4+V_k),\qquad
 \mathbb C^2/\overline{A\mathbb Z^4}\simeq
                         \mathbb R\times(S^1)^r
\]
as real topological groups. The original quotient is a complex Lie
group precisely when \(k\) is rational, in which case \(r=3\).
Otherwise it is non-Hausdorff. Every Chabauty limit of nearby full
lattices contains \(\overline{A\mathbb Z^4}\); the quotient topology
at the fixed zero and a limit of full lattices need not coincide.
\end{proposition}


% TARGET PROOF MASKED: push:src:D.4.3






























\begin{proposition}\hypertarget{tag-032M}{}\label{push:src:D.4.4}
\StatementRecord{push:src:D.4.4}{proved}{032M}{the definitions and constructions stated in the proof}

On an ordinary modulus chart where \(P=\wp'_\tau(\mu)\ne0\), put
\[
 L=\partial_\mu\log\vartheta_1(\mu,\tau),\qquad
 Z=L+2\pi i\frac{\operatorname{Im}\mu}{\operatorname{Im}\tau}.
\]
Then the exact real-analytic derivative is
\[
 \partial_\tau D_{\rm part}
       =\frac{iZ}{4\pi^2P}(3ZP-2g_2).
\]
Away also from the ramification of \(j\), the critical-point equation
on the coarse base is therefore \(Z=0\) or \(3ZP=2g_2\).
\end{proposition}


% ALL PRINTED PROOF MATERIAL MASKED





































\begin{proposition}\hypertarget{tag-032N}{}\label{push:src:D.4.5}
\StatementRecord{push:src:D.4.5}{proved}{032N}{\SeriesRef{push:section:D.4}, \SeriesRef{fam:src:C.1.1}, \SeriesRef{fam:src:C.2.3}, \SeriesRef{push:src:D.4.4}}

At the equianharmonic point \(\tau=\rho=e^{\pi i/3}\), with the
original period representative \(\mu=(2-\rho)/3\), one has
\[
 \mathcal N(\mu,\rho)=3^{1/4},\qquad
 D_{\rm part}(z_1)=-\frac3{2\pi}\log48.
\]
Consequently the zero locus passes through the coarse point \(t=0\)
exactly at the height
\[
 H_0=\frac3{2\pi}\log48
 =H_*+\frac3{2\pi}\log\frac43.
\]
It has \(D_c(0)<0\) below this height and \(D_c(0)>0\) above it.
In the fixed coarse coordinate the function is \(C^1\) at zero and
\[
 D_{\rm part}(t)=D_{\rm part}(0)
       +\frac7{32\pi}\operatorname{Re}t+O(|t|^{4/3}),\qquad
 \partial_tD_{\rm part}(t)=\frac7{64\pi}+O(|t|^{1/3}).
\]
Thus at \(H=H_0\) the level set is a regular \(C^1\) arc through
zero, tangent to the imaginary axis. This proves the crossing value
and its local regularity requested in \ref{push:section:D.4}, without using an unproved
closed form for \(\beta^{\rm part}(z_1)\).
\end{proposition}


% ALL PRINTED PROOF MATERIAL MASKED











































































































































\begin{cited}\hypertarget{tag-032O}{}\label{push:src:D.4.6}
\StatementRecord{push:src:D.4.6}{proved}{032O}{\SeriesRef{push:src:D.4.2}, \SeriesRef{push:src:D.4.5}, \SeriesRef{push:src:D.4.7}, \SeriesRef{push:src:D.4.8}}

The minimum principle in \ref{push:src:D.4.2} and the crossing
in \ref{push:src:D.4.5} enter the complete level classification through
the nonvanishing and boundary estimates of
\ref{push:src:D.4.7}--\ref{push:src:D.4.8}. Lin and Wang (2010,
Theorem~1.3) count critical points of the Green function on a fixed
elliptic curve, where \(Z=0\) is its critical equation. That theorem
does not count zeros along \(\wp_\tau(\mu(\tau))=0\) as the modulus
varies, nor does it address \(3ZP-2g_2\). Our completed argument
therefore supplies a separate nonvanishing certificate for both
factors, with analytic local and cusp estimates.
\end{cited}
The expanded source derivation is retained in the proof compendium, Paper~I, label \nolinkurl{push:src:D.4.6}.




\begin{proposition}\hypertarget{tag-032P}{}\label{push:src:D.4.7}
\StatementRecord{push:src:D.4.7}{proved}{032P}{\SeriesRef{ari:C.8.15}, \SeriesRef{push:src:D.4.4}, \SeriesRef{push:src:D.4.5}, \SeriesRef{ari:C.1.1}}

The function $D_{\rm part}$ has no critical point on the punctured
root neighbourhood $0<|\tau-i|\leq3/20000$, no critical point on the
coarse neighbourhood represented by $|\tau-\rho|\leq1/50$, and no
critical point where $\Im\tau\geq2$. At $\rho$ the two factors
$Z$ and $F=3ZP-2g_2$ of \ref{push:src:D.4.4} vanish, but the coarse gradient is
nonzero, as in \ref{push:src:D.4.5}. The earlier interval proof and its 400 local rectangles remain in
the S6\_09 proof compendium under this label. We give the consequence
of the full critical-point classification here.
\end{proposition}

% ALL PRINTED PROOF MATERIAL MASKED































\begin{proposition}\hypertarget{tag-032Q}{}\label{push:src:D.4.8}
\StatementRecord{push:src:D.4.8}{proved}{032Q}{\SeriesRef{ari:C.8.15}, \SeriesRef{push:src:D.4.1}, \SeriesRef{push:src:D.4.2}, \SeriesRef{push:src:D.4.3}, \SeriesRef{push:src:D.4.5}, \SeriesRef{ari:C.1.1}}

For every \(H>H_*\), the degeneracy locus \(Z_H\) is exactly one
embedded \(C^1\) Jordan curve in the finite \(t\)-plane. It is
real analytic unless \(H=H_0=3\log48/(2\pi)\); at that height it
is real analytic away from \(t=0\) and has the regular crossing
described in \ref{push:src:D.4.5}. There are no singular coarse levels at any
supercritical height. The bounded complementary disc contains
\(t=1\) and has \(D_c>0\); the complementary disc in
\(\mathbb P^1\), containing the cusp, has \(D_c<0\). The point
\(t=0\) is in the latter disc for \(H<H_0\), on the curve for
\(H=H_0\), and in the former disc for \(H>H_0\).

The bounded discs are strictly nested as \(H\) increases, shrink
to \(\{1\}\) as \(H\downarrow H_*\), and exhaust the finite
plane as \(H\to\infty\). Their small-height boundary shape is
the expansion in \ref{push:src:D.4.1}. The maximal open base locus of the complex
manifold construction consists of these two complementary discs.
Approaching their common boundary, the covolume of the marked
root torus lattice is
\[
 \operatorname{covol}(\Pi_c\mathbb Z^4)
       =\operatorname{Im}\tau\,|D_c|\longrightarrow0.
\]
At the boundary the period map has rank three and the exact quotient
and Hausdorff reflection are those of \ref{push:src:D.4.3}. Thus this statement
classifies the locus and the local degeneration requested in D.4.
\end{proposition}

% TARGET PROOF MASKED: push:src:D.4.8















































\section{The uniform corner obstruction}

\subsection{Period transport and the actual parameter descent}

\begin{theorem}\hypertarget{tag-02R9}{}\label{night:P2.1}\label{stmt:P2.1}
\StatementRecord{night:P2.1}{proved}{02R9}{the definitions and constructions stated in the proof}

For the prescribed $(3,4)$ construction with integral twists $(l_0,l_1,l_2)$, put
\[
 b=4l_1+3l_2,\qquad p=12l_0-4l_1-3l_2,\qquad a=b-12l_0=-p.
\]
On the logarithmic cover of the cusp bidisc the period columns are
\[
 \gamma=\binom{6\mu}{C-s-12v},\quad
 u=\binom{s+h}{\mu},\quad w=\binom10,\quad\delta=\binom01,
\]
where $h,\mu,C$ are holomorphic at $q=0$. Their positive-loop column comparisons are
\[
 K_q=\begin{pmatrix}1&0&0&0\\0&1&0&0\\0&1&1&0\\-1&0&0&1\end{pmatrix},\qquad
 K_r=\begin{pmatrix}1&0&0&0\\0&1&0&0\\0&0&1&0\\-12&0&0&1\end{pmatrix}.
\]
The parameter descent on the ordinary Jacobian chart is
\begin{equation}\hypertarget{tag-01O8}{}\label{eq:P2H}
 R:(t,v,\zeta)\longmapsto(t,v+1,\zeta+H(t)\delta),\qquad
 H(t)=-\frac{4l_1\log t+3l_2\log(t-1)}{2\pi i}.
\end{equation}
In the cusp collar it extends across $q=0$ and has coefficient $a$ on $s\delta$.
\end{theorem}
% ALL PRINTED PROOF MATERIAL MASKED



























\begin{theorem}\hypertarget{tag-02RC}{}\label{night:P2.2}\label{stmt:P2.2}
\StatementRecord{night:P2.2}{proved}{02RC}{the definitions and constructions stated in the proof}

On $q=Q^{12}$, actual extending origins have lifts
\begin{align*}
 o&=\binom{-b\mu s/2}{bs^2/24-bCs/12+l_0s/2+v(bs+U)},\\
 \sigma&=-l_0z_\gamma+\binom{6a\mu v}{vU+aCv-6av^2+6av}.
\end{align*}
The section $\sigma$ extends across the cusp for $r\ne0$, and $o$ extends across the parameter central fibre of the partial filling. Their difference is
\begin{equation}\hypertarget{tag-01O5}{}\label{eq:P2F}
 F=\sigma-o=a\binom{\mu(s+12v)/2}{C(s+12v)/12-(s+12v)^2/24+6v}.
\end{equation}
On $q=Q^{12m}$, $r=R^n$, with $m,n$ positive integers,
\begin{align}
 \Delta_QF&=am\gamma-6am^2\delta,\hypertarget{tag-01O6}{}\label{eq:P2FQ}\\
 \Delta_RF&=an\gamma-6an(n-1)\delta.\hypertarget{tag-01O7}{}\label{eq:P2FR}
\end{align}
The radial class modulo $\Z\widehat\delta$ is $-pn\widehat\gamma$. The proposed $+pn\widehat\gamma$ has the opposite sign in the fixed conventions.
\end{theorem}
% ALL PRINTED PROOF MATERIAL MASKED














\subsection{The partial parameter filling}
\begin{theorem}\hypertarget{tag-02RD}{}\label{night:P2.3}\label{stmt:P2.3}
\StatementRecord{night:P2.3}{proved}{02RD}{the definitions and constructions stated in the proof}

There is a proper holomorphic partial filling over $(\Pp^1\times\Delta_r)\setminus\{(\infty,0)\}$ with smooth total space and reduced normal-crossing parameter central divisor. Over a regular finite base point its central fibre is a cycle of twelve ruled surfaces. Its component rotations at $0,1,\infty$ are $4l_1,3l_2,-b$ modulo $12$, with component cover
\[
 y^{12}=t^{4l_1}(t-1)^{3l_2}.
\]
\end{theorem}
% TARGET PROOF MASKED: night:P2.3






















\subsection{Extending the marked base map}
\begin{theorem}\hypertarget{tag-02RE}{}\label{night:P2.4}\label{stmt:P2.4}
\StatementRecord{night:P2.4}{proved}{02RE}{the definitions and constructions stated in the proof}

For every admissible $(3,4)$ twist, the marked map $f:X_c\to\Pp^1$ extends meromorphically over every hypothetical proper Hausdorff analytic filling of a finite parameter cover. A proper modification carries a holomorphic marked map to $\Pp^1$, and hence a proper map to $\Pp^1\times\Delta$.
\end{theorem}
% ALL PRINTED PROOF MATERIAL MASKED












\begin{theorem}\hypertarget{tag-02RF}{}\label{night:P2.5}\label{stmt:P2.5}
\StatementRecord{night:P2.5}{proved}{02RF}{the definitions and constructions stated in the proof}

The global monodromy on $\Lambda_\Q/\Q\widehat\delta$ has no nonzero vector with finite orbit.
\end{theorem}
% ALL PRINTED PROOF MATERIAL MASKED












\begin{theorem}\hypertarget{tag-02RG}{}\label{night:P2.6}\label{stmt:P2.6}
\StatementRecord{night:P2.6}{proved}{02RG}{the definitions and constructions stated in the proof}

In a smooth resolved proper filling carrying the marked base map, every parameter-central component dominating $\Pp^1$ has zero normal-meridian lattice class modulo $\Z\widehat\delta$, relative to the ordinary parameter origin.
\end{theorem}
% ALL PRINTED PROOF MATERIAL MASKED
















\subsection{The commuting-meridian obstruction}
\begin{theorem}\hypertarget{tag-02RH}{}\label{night:P2.7}\label{stmt:P2.7}
\StatementRecord{night:P2.7}{proved}{02RH}{the definitions and constructions stated in the proof}

For every admissible $(3,4)$ twist, the punctured family has no proper Hausdorff complex-analytic filling at $w=0$ after any finite parameter cover. This includes every sphere member and permits arbitrary singularities, multiplicities and extra central components.
\end{theorem}
% ALL PRINTED PROOF MATERIAL MASKED







































\begin{theorem}\hypertarget{tag-02RI}{}\label{night:P2.8}\label{stmt:P2.8}
\StatementRecord{night:P2.8}{proved}{02RI}{the definitions and constructions stated in the proof}

For $(l_0,l_1,l_2)=(0,2,-1)$ we have $p=-5$, $a=b=5$, component cover $y^{12}=t^8/(t-1)^3$, and rotations $8,-3,-5$ modulo $12$. The descent and origins are
\begin{align*}
 H&=5s+U,& U&=\frac{3\log(1-q)}{2\pi i},\\
 o&=\binom{-5\mu s/2}{5s^2/24-5Cs/12+v(5s+U)},\\
 \sigma&=\binom{30\mu v}{vU+5Cv-30v^2+30v},& F&=5F_0,
\end{align*}
where $F_0$ is the vector in \eqref{eq:P2F} without $a$. Its two differences are $5m\gamma-30m^2\delta$ and $5n\gamma-30n(n-1)\delta$. No finite parameter cover admits a proper filling.
\end{theorem}
% ALL PRINTED PROOF MATERIAL MASKED







\subsection{Higher coprime multiplicities}
\begin{theorem}\hypertarget{tag-02RJ}{}\label{night:P2.9}\label{stmt:P2.9}
\StatementRecord{night:P2.9}{proved}{02RJ}{the definitions and constructions stated in the proof}

Every admissible logarithmic datum in the coprime $(3k,4l)$ family of the supplied Viaclovsky construction has no proper Hausdorff analytic filling at parameter zero after any finite cover.
\end{theorem}
% ALL PRINTED PROOF MATERIAL MASKED














































\subsection{The untwisted control}
\begin{theorem}\hypertarget{tag-02RA}{}\label{night:P2.10}\label{stmt:P2.10}
\StatementRecord{night:P2.10}{proved}{02RA}{the definitions and constructions stated in the proof}

The untwisted family $X(u)$ has a proper flat Hausdorff analytic filling at $u=0$. Its corner winding is zero. The filling is obtained from two projective-line bundles by explicit plumbing, a proper modification and a finite quotient.
\end{theorem}
% ALL PRINTED PROOF MATERIAL MASKED








































\begin{theorem}\hypertarget{tag-02RB}{}\label{night:P2.11}\label{stmt:P2.11}
\StatementRecord{night:P2.11}{proved}{02RB}{the definitions and constructions stated in the proof}

The one-component central space specified in Engel's Conjecture 8.2 cannot be the central fibre of a proper flat smoothing with smooth fibres $X(u)$. Its intrinsic virtual Euler number is zero, whereas $e(X(u))=2$.
\end{theorem}
% TARGET PROOF MASKED: night:P2.11























\section{The cotangent direct images and the decisive period scalar}
\hypertarget{tag-01TH}{}\label{inv:sec:partA}

\subsection{Local logarithmic bundles and a duality calculation}

We use the marked threefold, with its arbitrary admissible
constant $c$.  Coefficients of coherent cohomology are complex.  In this
section $\omega_f=K_X\otimes f^*K_B^{-1}$,
$\mathcal L=\Omega_B^1(\log\Sigma)=\OO_B(1)$,
$\mathcal A=\Omega_X^1(\log D)$, and
$\mathcal E=\mathcal A/f^*\mathcal L$.  We write $\gamma,u$ also for
their images in $R^1f_*\OO_X$; this convention concerns cohomology
classes and does not identify them with differential-form representatives.

\begin{cited}[The structure-sheaf input in the precise form used below]\hypertarget{tag-01RQ}{}\label{inv:A1.0}\label{wb:thm:structure-direct-images-audit}\label{wb:lem:cusp-structure-sheaf-audit}
\StatementRecord{inv:A1.0}{proved}{01RQ}{\cite[Lemma~9.12, pp.~72--74; Proposition~9.13, pp.~74--76; p.~92]{AC2026}}

This is the structure-sheaf calculation of Claude and Alp\"oge: \cite[Lemma~9.12, pp.~72--74; Proposition~9.13, pp.~74--76, and the base-change clarification on p.~92]{AC2026}. The direct images $R^qf_*\OO_X$, for $q=0,1,2$, are
$\OO$, $G=\OO\gamma\oplus\OO(-1)$, and $\OO(-1)$.
The inclusion $\OO\gamma\subset G$ is a subbundle.  At the cusp,
$(\gamma,u)$ is a holomorphic frame, and cup product identifies
$\bigwedge^2G$ with $R^2f_*\OO_X$.  In the cyclic covering coordinate
$s$, the action on $G$ has characters $(1,\zeta^{-a})$, where
$(m,a)=(3,2)$ or $(4,1)$ and $\zeta=e^{-2\pi i/m}$.
These assertions concern full scheme fibres, including $3S_1$ and $4S_2$.
At the cusp, the natural map $H^1(W;\mathbb C)\to H^1(W,\mathcal O_W)$ is an isomorphism. In degree two the comparison is surjective and is injective on $\ker[H^2(W;\mathbb C)\to H^2(\mathrm{dP}_6;\mathbb C)]$. The cup product $\bigwedge^2H^1(W,\mathcal O_W)\to H^2(W,\mathcal O_W)$ is an isomorphism. Every base-change map $(R^qf_*\mathcal O_X)\otimes k(b)\to H^q(X_b,\mathcal O_{X_b})$ is an isomorphism, including at the nonreduced fibres.
\end{cited}


The expanded source derivation is retained in the proof compendium, Paper~I, label \nolinkurl{inv:A1.0}.

\Uses{the hypotheses, notation and definitions stated here; \citet[Lemma~9.12 and Proposition~9.13]{AC2026}; \CompendiumRef{I}{app:rederive:inv-A1.0}}


The structure-sheaf and cusp statements are due to \citet[Lemma~9.12 and Proposition~9.13]{AC2026}; we retain the independent root-line derivation. A complete re-derivation is given in \CompendiumRef{I}{app:rederive:inv-A1.0}.


\begin{proposition}[Local freeness and the determinant line]\hypertarget{tag-01RR}{}\label{inv:A1.1}
\StatementRecord{inv:A1.1}{proved}{01RR}{the definitions and constructions stated in the proof}

Both $\mathcal A$ and $\mathcal E$ are locally free, of ranks three and
two.  The vertical vector $\xi=\partial/\partial\zeta_2$ is nowhere
zero as a section of $\mathcal E^\vee$.  There is an exact sequence
of holomorphic vector bundles
\begin{equation}\hypertarget{tag-01SZ}{}\label{inv:eq:A-E-line}
0\longrightarrow N\longrightarrow\mathcal E
 \xrightarrow{\langle\cdot,\xi\rangle}\OO_X\longrightarrow0,
\qquad N=\det\mathcal E=\omega_f(-2S_1-3S_2).
\end{equation}
\end{proposition}


% ALL PRINTED PROOF MATERIAL MASKED










































\begin{proposition}[The logarithmic direct images by vertical twisting and duality]\hypertarget{tag-01RS}{}\label{inv:A1.2}\label{wb:prop:log-relative-direct-images}
\StatementRecord{inv:A1.2}{proved}{01RS}{\SeriesRef{inv:eq:A-H-images}, \SeriesRef{inv:eq:A-E-line}, \SeriesRef{inv:eq:A-N-images}, \SeriesRef{inv:A1.0}}

One has, as holomorphic vector bundles on $B$,
\[
\begin{array}{c|ccc}
q&0&1&2\\\hline
R^qf_*\mathcal E&\OO\oplus\OO(-1)&
 \OO^2\oplus\OO(-1)^2&\OO\oplus\OO(-1).
\end{array}
\]
They commute with base change to every full scheme fibre.
\end{proposition}


% ALL PRINTED PROOF MATERIAL MASKED































































\begin{proposition}[All absolute cotangent direct images]\hypertarget{tag-01RT}{}\label{inv:A1.3}\label{wb:prop:absolute-cotangent-direct-images}
\StatementRecord{inv:A1.3}{proved}{01RT}{\SeriesRef{inv:A2.1}, \SeriesRef{inv:A2.2}, \SeriesRef{inv:A2.4}, \SeriesRef{inv:A1.1}}

For every admissible $c$, the six direct images are
\begin{equation}\hypertarget{tag-01T2}{}\label{inv:eq:A-absolute-table}
\begin{array}{c|cc}
q&R^qf_*\Omega_X^1&R^qf_*\Omega_X^2\\\hline
0&\OO(-2)&\OO(-2)\oplus\OO(-1)\\
1&\OO\oplus\OO(-1)^2&\OO(-2)\oplus\OO(-1)^2\\
2&\OO\oplus\OO(-1)&\OO.
\end{array}
\end{equation}
In particular every listed sheaf is torsion-free.
\end{proposition}


% ALL PRINTED PROOF MATERIAL MASKED























































\subsection{An exact modular determinant and two numerical routes}

The scalar depends on the fourth root chosen in $t-1=s^4$.  We make
that ambiguity explicit.  We choose the local branch for which, on
the positive $t>1$ arc leading to the distinguished cusp,
$\tau=iy$ with $y>1$, and
\begin{equation}\hypertarget{tag-01T5}{}\label{inv:eq:A-s-branch}
\tau=i+\frac{i}{\pi\sqrt{E_4(i)}}s^2+O(s^4),\qquad
\operatorname{Im}\mu_s(0)>0.
\end{equation}
The nonvanishing proved below makes the second condition meaningful;
it fixes the remaining sign of $s$.  Equivalently it is the lower
cusp-regular zero branch of the Weierstrass function in Statement \ref{inv:A2.3}.
Replacing $s$ by $i^js$ multiplies the reported derivative by $i^{-j}$.
The workbench instead normalizes $ds/dz(z_2)=1$; its number is obtained
by the chain rule and is not numerically the same unspecified coordinate.

\begin{proposition}[The complete raw period determinant]\hypertarget{tag-01RU}{}\label{inv:A2.1}\label{wb:prop:log-ks-divisor}\label{wb:prop:period-derivative-determinant}
\StatementRecord{inv:A2.1}{proved}{01RU}{the definitions and constructions stated in the proof}

On the marked upper half-plane,
\begin{equation}\hypertarget{tag-01T3}{}\label{inv:eq:A-modular-K}
6\mu_z^2-\tau_z\beta_z
=\tau_z^2\frac{E_4(\tau)^2}{E_6(\tau)}.
\end{equation}
The right side extends holomorphically through the apparent poles.
Its exact orders are two at $z_1$ and zero at $z_2$; it has no other
zero away from their translates.  The logarithmic connecting map
$\kappa_0:f_*\mathcal E\to\mathcal L\otimes G$ has divisor
$\operatorname{divisor}(\det\kappa_0)=p_1+p_2$.
\end{proposition}


% ALL PRINTED PROOF MATERIAL MASKED










































































\begin{proposition}[The evaluated scalar and the two local leading terms]\hypertarget{tag-01RV}{}\label{inv:A2.2}
\StatementRecord{inv:A2.2}{proved}{01RV}{\SeriesRef{inv:eq:A-s-branch}, \SeriesRef{inv:eq:A-modular-K}, \SeriesRef{inv:A2.3}}

In the coordinate \eqref{inv:eq:A-s-branch},
\begin{equation}\hypertarget{tag-01T4}{}\label{inv:eq:A-mu-value}
\mu_s(0)=i\,C_4,
\qquad
C_4=\frac{\sqrt{2/3}}{\pi E_4(i)^{1/4}}
=\frac{4\sqrt\pi}{3^{3/4}\Gamma(1/4)^2}>0.
\end{equation}
Numerically,
\[
C_4=0.2366093140807852246877051134504496062943051340858243372192243072\ldots.
\]
For $r^3=t$, with $r>0$ on the arc $0<t<1$ and $\tau$ travelling
from $\rho$ to $i$ along the unit circle, set
\[
\lambda=\frac{3e^{5\pi i/6}}{2\pi E_6(\rho)^{1/3}}.
\]
Then $\tau=\rho+\lambda r+O(r^2)$ and
\[
K_r=-\frac{3i\lambda}{2\pi}r^2+O(r^3).
\]
The leading coefficient has real and imaginary parts
\[
\begin{aligned}
\operatorname{Re}K_2&=0.08010221108753827140763650309934451243\ldots,\\
\operatorname{Im}K_2&=0.13874109940222333968402480721280640441\ldots.
\end{aligned}
\]
\end{proposition}


% ALL PRINTED PROOF MATERIAL MASKED




















































































\begin{proposition}[Identification of the elliptic logarithm and the verified zero formula]\hypertarget{tag-01RW}{}\label{inv:A2.3}
\StatementRecord{inv:A2.3}{proved}{01RW}{\SeriesRef{inv:eq:A-s-branch}, \SeriesRef{inv:eq:A-mu-value}; \citet{EichlerZagier1982}; \citet{DukeImamoglu2008}}

For the Weierstrass equation
\[
y^2=x^3-3t^3(t-1)x+2t^4(t-1)^2,
\]
the section $[\mu]$ is a generator of its rational section group and
equals one of $P_*=(0,\sqrt2t^2(t-1))$ and $-P_*$.
In particular $\wp(\mu,\tau)=0$ on every regular fibre.
Along the branch in \eqref{inv:eq:A-s-branch}, with
$x_0=1-1/t$, the exact branch is
\begin{equation}\hypertarget{tag-01SY}{}\label{inv:eq:A-DI-branch}
\mu=\frac{1-\tau}{2}
 +\frac{i\sqrt6}{3\pi}
 \frac{x_0^{1/4}\,{}_3F_2(1/3,2/3,1;3/4,5/4;x_0)}
 {{}_2F_1(1/12,5/12;1;1-x_0)}.
\end{equation}
Consequently $\mu-(2-\tau)/3$ has exact order two at the
order-three point in $r^3=t$.
\end{proposition}


% ALL PRINTED PROOF MATERIAL MASKED











































































































































































\begin{proposition}[The global rank and the decisive scalar]\hypertarget{tag-01RX}{}\label{inv:A2.4}\label{wb:thm:h12-scalar}
\StatementRecord{inv:A2.4}{proved}{01RX}{\SeriesRef{inv:A2.1}, \SeriesRef{inv:A1.3}}

The map
$\kappa_1:R^1f_*\mathcal E\to\mathcal L\otimes R^2f_*\OO=\OO$
is surjective, and $H^0(\kappa_1):\C^2\to\C$ has rank one.
Let $h=a\,d\zeta_1+d\zeta_2$ be the unique global section of
$f_*\mathcal E$ with $h(\xi)=1$.  In a genuine invariant target
frame normalized as below,
\[
\kappa_1(h\gamma)(p_2)=-\tfrac14\mu_s(0)=-\tfrac i4 C_4.
\]
The rank-one condition is equivalent to
$a\tau_z+\mu_z\not\equiv0$.  Under the same direct-image and
sheaf-surjectivity hypotheses a rank-zero global map would instead
give $h^{1,2}=2$, $h^{1,1}=3$; that alternative does not occur here.
\end{proposition}


% ALL PRINTED PROOF MATERIAL MASKED





























































\subsection{The Hodge table and its Leray generators}

\begin{proposition}[The complete Hodge table for every admissible parameter]\hypertarget{tag-01RY}{}\label{inv:A3.1}\label{wb:cor:exact-hodge-diamond}\label{wb:thm:parametric-hodge-diamond}\label{wb:prop:workbench-holomorphic-form-vanishing}
\StatementRecord{inv:A3.1}{proved}{01RY}{\SeriesRef{inv:eq:A-absolute-table}, \SeriesRef{inv:A1.0}}

For every $\operatorname{Im}c<-M$, the Hodge numbers are
\[
\begin{array}{c|rrrr}
q\backslash p&0&1&2&3\\\hline
3&0&0&0&1\\
2&0&1&2&1\\
1&1&2&1&0\\
0&1&0&0&0.
\end{array}
\]
\end{proposition}



The vanishing of holomorphic forms and the Hodge-diamond
pattern with one undetermined integer are due to
\citet[Proposition~9.20, p.~80; Theorem~9.1(vi), p.~65]{AC2026}.
The value of that integer is determined here by the nonzero direct-image
scalar. The full calculations of both the pattern and its value are retained.

% ALL PRINTED PROOF MATERIAL MASKED



















\begin{proposition}[Specified associated-graded generators and a global Dolbeault lift]\hypertarget{tag-01RZ}{}\label{inv:A3.2}\label{wb:lem:bcpr-eta-extension}
\StatementRecord{inv:A3.2}{proved}{01RZ}{\SeriesRef{inv:A2.4}, \SeriesRef{inv:eq:A-E-line}, \SeriesRef{inv:A1.3}}

The Leray associated-graded spaces for the four middle groups are
\[
\begin{array}{c|cc}
&E_2^{1,q-1}&E_2^{0,q}\\\hline
H^{1,1}&H^1(\OO(-2))=\C\alpha&\C\eta_{\rm rel}\\
H^{1,2}&0&\C\upsilon\\
H^{2,1}&\C\upsilon^\vee&0\\
H^{2,2}&\C\eta_{\rm rel}^\vee&\C\alpha^\vee.
\end{array}
\]
Here $\alpha$ is the Cech class of $dt/t$ on the two standard affine
base charts, $\eta_{\rm rel}=-q_E$ is the invariant relative class
from \ref{inv:A2.4}, and $\upsilon$ is the image of the trace-dual of
$1\in f_*\OO(H)$ under $R^2f_*N\to R^2f_*\mathcal E$ and the
isomorphism $R^2f_*\Omega_X^1=R^2f_*\mathcal E$.
The superscript $\vee$ denotes the bases dual under the relative
trace and the base Serre pairing.  There exists a Dolbeault class
$Q\in H^1(X,\Omega_X^1)$ whose restriction to every regular fibre
is $\eta=u\wedge w+6\gamma\wedge\delta$.
\end{proposition}


% ALL PRINTED PROOF MATERIAL MASKED























\subsection{Ordinary relative differentials and the correction at the cusp}

\begin{proposition}[All local modules needed to distinguish degrees one and two]\hypertarget{tag-01S0}{}\label{inv:A4.1}
\StatementRecord{inv:A4.1}{proved}{01S0}{the definitions and constructions stated in the proof}

Let $R=\C\{x,y,z\}$.  For $t=xy$,
\[
\Omega_{R/\C\{t\}}^1=R^3/R(y,x,0)
\simeq(x,y)\oplus R,
\]
and
\[
\Omega_{R/\C\{t\}}^2\simeq
(x,y)\oplus R/(x,y),\qquad
\Omega_{R/\C\{t\}}^3\simeq R/(x,y).
\]
For $t=xyz$,
$\Omega^1=R^3/R(yz,xz,xy)$ is torsion-free of rank two and is
not free at the triple point.  If $I=(xy,xz,yz)$, there is an exact
sequence
\[
0\to T\to\Omega^2\xrightarrow{\wedge dt} I\to0,
\qquad
T=\frac{R u\oplus R v}{(xu,\ y(u-v),\ zv)},
\quad IT=0,
\]
and $\Omega^3=R/I$.  Finally, for $t=s^m$ in
$R=\C\{s,y_1,y_2\}$,
\[
\begin{aligned}
\Omega^1&=R^2\oplus R/(s^{m-1})\,ds,\\
\Omega^2&=R\,dy_1\wedge dy_2
 \oplus\bigl(R/(s^{m-1})\bigr)^2,\\
\Omega^3&=R/(s^{m-1})\,ds\wedge dy_1\wedge dy_2.
\end{aligned}
\]
\end{proposition}


% ALL PRINTED PROOF MATERIAL MASKED












































\begin{proposition}[Refuted: The claimed degree-one torsion along the cusp]\hypertarget{tag-01S1}{}\label{inv:A4.2}
\StatementRecord{inv:A4.2}{proved}{01S1}{\SeriesRef{inv:A4.1}}

The assertion in AC26 Remark 9.21, PDF page 81, that
$\Omega_{X/B}^1$ has torsion along $W\cup S_1\cup S_2$ is false.
Its torsion is supported precisely on $S_1\cup S_2$; its degree-one
cusp modules are torsion-free but not locally free on the double
locus.  The plural statement in Remark 1.3, PDF page 5, must be
qualified by degree: $\Omega^2_{X/B}$ does have torsion on the cusp
double locus.  The workbench's degree-one correction is correct;
it must not be extended to degree two.
\end{proposition}


% ALL PRINTED PROOF MATERIAL MASKED






















\section{Fr\"olicher, Bott--Chern, Aeppli, and Hermitian metrics}
\hypertarget{tag-01TI}{}\label{inv:sec:partC}

\subsection{The differentials and their dependence on the Hodge number}

In this section $X$ is the standard marked $(3,4)$ threefold and
$h=h^{1,2}(X)=1$ by Part A.  We retain $h$ in the formal calculations
to specify exactly which conclusions depend on its evaluated value.
Every assertion involving $h$ is valid for every value
consistent with the Hodge table
\begin{equation}\hypertarget{tag-01TF}{}\label{inv:eq:C-hodge-input}
\begin{array}{c|rrrr}
q\backslash p&0&1&2&3\\\hline
3&0&0&0&1\\
2&0&h&h+1&1\\
1&1&h+1&h&0\\
0&1&0&0&0.
\end{array}
\end{equation}
Thus no conclusion below silently uses $h=1$.  The geometric inputs for
this part are the structure-sheaf calculation, the vanishing of
holomorphic forms, and the sphere topology, which give
\eqref{inv:eq:C-hodge-input}; Statements \ref{inv:A1.3} and \ref{inv:A3.1} determine $h$.
We use $d=\partial+\bar\partial$ on complex-valued
smooth forms and the complex orientation for integration.

\begin{proposition}[The complete Fr\"olicher differential in terms of $h$]\hypertarget{tag-01SA}{}\label{inv:C1.1}\label{wb:thm:frolicher-e2-candidate}
\StatementRecord{inv:C1.1}{proved}{01SA}{\SeriesRef{inv:eq:C-hodge-input}}

For a compact complex threefold with \eqref{inv:eq:C-hodge-input} and the
complex de Rham cohomology of $S^6$, the Fr\"olicher spectral sequence
satisfies $E_1\ne E_2=E_\infty$.  Its only nonzero $d_1$ ranks are
\[
\rank d_1^{0,1}=1,\quad \rank d_1^{1,1}=h,\quad
\rank d_1^{1,2}=h,\quad \rank d_1^{2,2}=1.
\]
In particular its two middle rows are exact.
\end{proposition}


% ALL PRINTED PROOF MATERIAL MASKED



























\begin{proposition}[A basis on which all first differentials are specified]\hypertarget{tag-01SB}{}\label{inv:C1.2}
\StatementRecord{inv:C1.2}{proved}{01SB}{\SeriesRef{inv:A1.3}, \SeriesRef{inv:A3.2}, \SeriesRef{inv:C1.1}, \SeriesRef{inv:C3.7}}

Choose a nonzero $g\in H^{0,1}_{\bar\partial}(X)$, a basis
$q_1,\ldots,q_h$ of a complement to $\C d_1g$ in
$H^{1,1}_{\bar\partial}(X)$, a basis $v_1,\ldots,v_h$ of
$H^{1,2}_{\bar\partial}(X)$, and a class $e$ complementary to the
image of $d_1:H^{1,2}\to H^{2,2}$.  With
$a=d_1g$, $c_j=d_1q_j$, $d_j=d_1v_j$, and $t=d_1e$, the first
differential is
\[
g\mapsto a\mapsto0,\quad q_j\mapsto c_j\mapsto0,
\qquad v_j\mapsto d_j\mapsto0,\quad e\mapsto t\mapsto0.
\]
These are bases in the corresponding bidegrees.  If $h=1$ and a
class $q$ restricts to the invariant fibre class $\eta$ described below,
we may take $q_1=q$, $v_1=gq$, and $e=q^2$, after scaling $g$ if
necessary.  Statements \ref{inv:A1.3} and \ref{inv:A3.2} prove these hypotheses for $X$,
so this geometric choice applies to the threefold in the stated problem.
\end{proposition}


% ALL PRINTED PROOF MATERIAL MASKED



















\subsection{A computation from the two definitions}

\begin{proposition}[Bott--Chern and Aeppli tables as functions of $h$]\hypertarget{tag-01SC}{}\label{inv:C2.1}\label{wb:thm:exact-bott-chern-aeppli}
\StatementRecord{inv:C2.1}{proved}{01SC}{\SeriesRef{inv:A3.1}, \SeriesRef{inv:C1.1}, \SeriesRef{inv:eq:C-arrows}; \citet[Theorem 3]{Stelzig2021}; \citet[Proposition 6(1)]{Stelzig2021}}

With
\[
H^{p,q}_{BC}=
\frac{\ker\partial\cap\ker\bar\partial\subset A^{p,q}}
{\im\partial\bar\partial},\qquad
H^{p,q}_{A}=
\frac{\ker\partial\bar\partial\subset A^{p,q}}
{\im\partial+\im\bar\partial},
\]
the complete tables are
\begin{equation}\hypertarget{tag-01T6}{}\label{inv:eq:C-BC-A}
\begin{array}{c|rrrr}
q\backslash p&0&1&2&3\\\hline
3&0&0&1&1\\
2&0&h&2h&1\\
1&0&2&h&0\\
0&1&0&0&0
\end{array}_{BC}
\qquad
\begin{array}{c|rrrr}
q\backslash p&0&1&2&3\\\hline
3&0&0&0&1\\
2&0&h&2&0\\
1&1&2h&h&0\\
0&1&1&0&0.
\end{array}_{A}
\end{equation}
In particular, substitution of the value $h=1$ proved in \ref{inv:A3.1} gives
the workbench's numerical tables.  The totals are
$(1,0,2,2,2,2,1)$ and $(1,2,2,2,2,0,1)$, respectively.
\end{proposition}


% ALL PRINTED PROOF MATERIAL MASKED











































\begin{proposition}[Comparison with McHugh and Angella]\hypertarget{tag-01SD}{}\label{inv:C2.2}
\StatementRecord{inv:C2.2}{proved}{01SD}{\SeriesRef{inv:eq:C-hodge-input}, \SeriesRef{inv:C2.1}, \SeriesRef{inv:C1.1}, \SeriesRef{inv:eq:C-BC-A}; \citet[Theorem 6.3, Proposition 6.4, (7), pp. 20--25]{McHugh2018}; \citet[Section 4, pp. 13--16]{Angella2018}}

The specialization of McHugh's Theorem 6.3, Proposition 6.4, equation
(7), and Section 7.1 to \eqref{inv:eq:C-hodge-input} agrees with \ref{inv:C2.1}.
The corresponding relations and Figure 5 in Angella's Section 4 also
agree.  No correction to these specific formulas is required.
\end{proposition}


% ALL PRINTED PROOF MATERIAL MASKED



























\subsection{The algebra, line bundles, and metric consequences}

\begin{proposition}[The real Bott--Chern Chern-class map and a global $G$]\hypertarget{tag-01SE}{}\label{inv:C3.1}\label{wb:cor:picard-line-audit}
\StatementRecord{inv:C3.1}{proved}{01SE}{\SeriesRef{inv:C1.1}, \SeriesRef{inv:C2.1}, \SeriesRef{inv:A1.0}; \cite[Theorem~10.5(b), proof Steps~1 and~4, p.~88]{AC2026}}

The Picard-group identification and fibrewise injectivity below are due to Claude and Alp\"oge \cite[Theorem~10.5(b), proof Steps~1 and~4, p.~88]{AC2026}. We compute the real Chern map and a specified Green-operator representative in that coordinate. Let $G$ be a $\bar\partial$-closed smooth $(0,1)$ representative of
a nonzero $g\in H^1(X,\OO_X)$.  Set
\[
\mathsf A=[\partial G]_{BC},\qquad
\mathsf B=[\bar\partial\bar G]_{BC}.
\]
These classes form a complex basis of $H^{1,1}_{BC}(X)$.
With the exponential $z\mapsto e^{2\pi iz}$ and the Chern convention
$c_{1,BC}=[iF/(2\pi)]$, the map of real Lie groups is
\begin{equation}\hypertarget{tag-01T8}{}\label{inv:eq:C-Pic-map}
\Pic(X)=\{\exp(zg):z\in\C\}\longrightarrow H^{1,1}_{BC}(X),
\qquad \exp(zg)\longmapsto z\mathsf A+\bar z\mathsf B.
\end{equation}
It is injective and its image is the real subspace fixed by conjugation.

There is an operator-explicit global representative.  Normalize the
Fubini--Study form by $\int_{\Pone}\omega_{FS}=1$, put
$\theta=f^*\omega_{FS}$, choose any smooth Riemannian metric on $X$,
and let $\mathcal G_d$ be its Green operator for the real de Rham
Laplacian.  Then
\begin{equation}\hypertarget{tag-01T7}{}\label{inv:eq:C-G-Green}
\alpha=d^*\mathcal G_d\theta,\qquad G_0=\alpha^{0,1}
\end{equation}
satisfy $d\alpha=\theta$, $\bar\partial G_0=0$,
$[G_0]\ne0$, and $\exp([G_0])=f^*\OO_{\Pone}(1)$.
Restriction $H^1(X,\mathcal O_X)\to H^1(X_b,\mathcal O_{X_b})$ is injective for every scheme-theoretic fibre; on a smooth fibre its image is $\mathbb C\gamma^{0,1}$.
\end{proposition}

% ALL PRINTED PROOF MATERIAL MASKED



































































\Uses{the hypotheses, notation and definitions stated here; \citet[Theorem~10.5]{AC2026}; \CompendiumRef{I}{app:rederive:inv-C3.1}}


The Picard and fibre-restriction statement is due to \citet[Theorem~10.5]{AC2026}. A complete re-derivation is given in \CompendiumRef{I}{app:rederive:inv-C3.1}.


\begin{proposition}[The factor $12$ with a fixed collar orientation]\hypertarget{tag-01SF}{}\label{inv:C3.2}\label{wb:lem:bcpr-degree-one-input}
\StatementRecord{inv:C3.2}{proved}{01SF}{\SeriesRef{inv:eq:C-collar-sum}, \SeriesRef{inv:C3.1}}

Use the marked cohomology class $\gamma$ dual to
$\hat\gamma$, the analytic peripheral generators and the lift relation of
[AC26, Section 7.5], and put
$p=12l_0-4l_1-3l_2=\pm1$.  Let $g_\gamma$ be the unique coherent
Leray class whose fibre image is $\gamma^{0,1}$.  Then
\begin{equation}\hypertarget{tag-01TE}{}\label{inv:eq:C-factor12}
\exp(12g_\gamma)=f^*\OO_{\Pone}(-p).
\end{equation}
For the standard data $(l_0,l_1,l_2)=(0,1,-1)$, this is
$\exp(12g_\gamma)=f^*\OO(1)$.  For the opposite value $p=1$,
the exponent representing $f^*\OO(1)$ is $-12g_\gamma$.
\end{proposition}


% ALL PRINTED PROOF MATERIAL MASKED


















































































\begin{proposition}[The complete Bott--Chern algebra, for every $h$]\hypertarget{tag-01SG}{}\label{inv:C3.3}\label{wb:thm:bcpr-square-zero-algebra}\label{wb:prop:bcpr-generators-are-bases}
\StatementRecord{inv:C3.3}{proved}{01SG}{\SeriesRef{inv:C2.1}, \SeriesRef{inv:C1.2}, \SeriesRef{inv:C3.1}, \SeriesRef{inv:C1.1}}

The positive-total-degree ideal $N\subset H^{\bullet,\bullet}_{BC}(X)$
satisfies
\[
H^{\bullet,\bullet}_{BC}(X)=\C\cdot1\oplus N,\qquad N^2=0.
\]
Its nonzero bidegrees and dimensions are those of \ref{inv:C2.1}.  Thus the
square-zero conclusion is independent of the evaluated value of $h$.
For $h=1$, the nine classes specified in \CompendiumRef{I}{wb:eq:bcpr-generators}, together with the unit, give bases of every nonzero Bott--Chern group.
\end{proposition}


% ALL PRINTED PROOF MATERIAL MASKED





































































\begin{proposition}[Non-K\"ahler degrees and the precise degreewise statement]\hypertarget{tag-01SH}{}\label{inv:C3.4}
\StatementRecord{inv:C3.4}{proved}{01SH}{\SeriesRef{inv:C2.1}}

For $k=0,\ldots,6$ the non-K\"ahler degrees are
\begin{equation}\hypertarget{tag-01TC}{}\label{inv:eq:C-delta}
(\Delta^k)=(0,2,2+2h,4h,2+2h,2,0).
\end{equation}
If $h\ge1$, the two maps
$\bigoplus_{p+q=k}H^{p,q}_{BC}\to H^k_{\mathrm{dR}}(X;\C)
\to\bigoplus_{p+q=k}H^{p,q}_A$
are both isomorphisms exactly for $k=0,6$.
For the weaker assertion that a $d$-exact pure $d$-closed form of
degree $k$ is $\partial\bar\partial$-exact, the valid degrees are
$k=0,1,6$, and also $k=3,4$ if $h=0$.
\end{proposition}


% ALL PRINTED PROOF MATERIAL MASKED





















\begin{proposition}[Metric obstructions that do not require $h$]\hypertarget{tag-01SI}{}\label{inv:C3.5}
\StatementRecord{inv:C3.5}{proved}{01SI}{the definitions and constructions stated in the proof}

The threefold $X$ admits no balanced, strongly Gauduchon, K\"ahler,
locally conformally K\"ahler, or Hermitian-symplectic metric.
\end{proposition}


% ALL PRINTED PROOF MATERIAL MASKED
























\begin{cited}[Gauduchon metrics with specified Chern--Ricci form]\hypertarget{tag-01SJ}{}\label{inv:C3.6}
\StatementRecord{inv:C3.6}{proved}{01SJ}{\citet[Theorem~2]{Gauduchon1985}; \citet[Theorem~1.3]{STW2017}}

There exists a Gauduchon metric on $X$.  More precisely, for every
Gauduchon metric $\omega_0$ there is another Gauduchon metric
$\omega$ with $[\omega^2]_A=[\omega_0^2]_A$ and
\begin{equation}\hypertarget{tag-01T9}{}\label{inv:eq:C-Ricci}
\operatorname{Ric}^{\mathrm{Ch}}(\omega)=\pi f^*\omega_{FS}.
\end{equation}
This Ricci form is semipositive and has rank one at regular points
of $f$.  No Chern--Ricci-flat Hermitian metric exists on $X$.
\end{cited}
The expanded source derivation is retained in the proof compendium, Paper~I, label \nolinkurl{inv:C3.6}.




This is the application of \citet[Theorem~2]{Gauduchon1985} and \citet[Theorem~1.3]{STW2017} with the computed Bott--Chern class. A complete re-derivation is given in \CompendiumRef{I}{app:rederive:inv-C3.6}.


\begin{proposition}[The fibre obstruction to SKT and astheno-K\"ahler metrics]\hypertarget{tag-01SK}{}\label{inv:C3.7}
\StatementRecord{inv:C3.7}{proved}{01SK}{\SeriesRef{inv:A3.2}, \SeriesRef{inv:C1.2}, \SeriesRef{inv:C2.1}, \SeriesRef{inv:eq:C-eta-negative}}

The threefold $X$ admits no pluriclosed (SKT) metric and no
astheno-K\"ahler metric.  The geometric input is the value $h=1$
and the global Dolbeault class $q$ in \ref{inv:A3.2}, whose restriction to a
regular fibre is $\eta=u\wedge w+6\gamma\wedge\delta$.
More generally, it is enough that the real image of
$H_A^{1,1}(X;\R)\to H_A^{1,1}(F;\R)$ be contained in
$\R\eta$ for one regular fibre $F$.
\end{proposition}


% ALL PRINTED PROOF MATERIAL MASKED











































































\section{Twisted tangent cohomology}
\subsection{All integral twists and the generic nonintegral twist}

\begin{proposition}[Integral Picard twists, including their torsion]\hypertarget{tag-01SO}{}\label{inv:D4.1}
\StatementRecord{inv:D4.1}{proved}{01SO}{the definitions and constructions stated in the proof}

For $M_n=\OO_X(n(S_1-S_2))$, define
$u(n)=\max\{d(n)+1,0\}$, $v(n)=\max\{-d(n)-1,0\}$ and
$\epsilon_n=0$ if $3\mid n$, $\epsilon_n=1$ otherwise. Then
\[
 f_*(T_X\otimes M_n)=\OO(d(n))\oplus\OO(d(n-1)),
\]
\[
 R^1f_*(T_X\otimes M_n)=
 \OO(d(n))\oplus\OO(d(n-1))\oplus\OO(d(n+2)),           \tag{D.12}
\]
\[
 R^2f_*(T_X\otimes M_n)=
 \OO(d(n+6))\oplus\C_{p_1}^{\epsilon_n}.
\]
In particular its cusp stalk is free of rank one. The four global
dimensions, for every integer $n$, are
\begin{align*}
 h^0&=u(n)+u(n-1),\\
 h^1&=v(n)+v(n-1)+u(n)+u(n-1)+u(n+2),\\
 h^2&=v(n)+v(n-1)+v(n+2)+u(n+6)+\epsilon_n,\\
 h^3&=v(n+6).                                          \tag{D.13}
\end{align*}
\end{proposition}


% ALL PRINTED PROOF MATERIAL MASKED


















\begin{proposition}[The cusp length and the generic twist]\hypertarget{tag-01SP}{}\label{inv:D4.2}
\StatementRecord{inv:D4.2}{proved}{01SP}{\SeriesRef{inv:C3.2}, \SeriesRef{inv:D4.3}, \SeriesRef{inv:D1.2}}

Use the coordinate $M_z=\exp(z\mathfrak g)$ of \ref{inv:C3.2}. For every
$z\notin\Z$ the cusp contribution is exactly
\[
 (R^0f_*(T_X\otimes M_z))_{p_0}=0,\quad
 (R^1f_*(T_X\otimes M_z))_{p_0}=0,\quad
 (R^2f_*(T_X\otimes M_z))_{p_0}=\C\{q\}/(q).           \tag{D.14}
\]
Outside the countable exceptional set explicitly defined in \ref{inv:D4.3},
these are the only nonzero direct-image stalks. Hence the generic
twist has $(h^0,h^1,h^2,h^3)=(0,0,1,0)$, and the cusp nilpotence exponent $r$ equals one there.
\end{proposition}


% ALL PRINTED PROOF MATERIAL MASKED







































\begin{proposition}[The exceptional nonintegral twists cannot be omitted]\hypertarget{tag-01SQ}{}\label{inv:D4.3}
\StatementRecord{inv:D4.3}{proved}{01SQ}{\SeriesRef{inv:D4.2}}

On a regular marked fibre, $M_z$ is trivial if and only if there are
integers $n,p,q,r$ such that
\[
 q\tau(t)+r\mu(t)=p,\qquad
 z=n-6q\mu(t)-r\beta_c(t).                             \tag{D.15}
\]
For $z\notin\Z$ these resonances occur in a countable set of parameter
values $z$; every resonant $z$ is nonreal. On the finite root covers
the exceptional lattices are
\[
 \Lambda_1=\Z+(3\beta_c(z_1)+4-2\rho)\Z,\qquad
 \Lambda_2=\Z+(2\beta_c(z_2)+3-3i)\Z.                  \tag{D.16}
\]
Take the union of $\Z$, these two lattices, and all values from (D.15)
over regular fibres; this is the countable exceptional set used in
\ref{inv:D4.2}. The set is larger than $\Z$: the regular value $t=32/5$
supplies nonintegral twists with $h^1(X,T_X\otimes M_z)\ge2$.
\end{proposition}


% ALL PRINTED PROOF MATERIAL MASKED























































\begin{proposition}[Local complexes and every noncusp torsion length]\hypertarget{tag-01SR}{}\label{inv:D4.4}
\StatementRecord{inv:D4.4}{proved}{01SR}{\SeriesRef{inv:D4.1}, \SeriesRef{inv:A2.3}}

Let $z\notin\Z$. At a regular resonance $t_*$ in (D.15), with
$A=\C\{t-t_*\}$, the two stalks are
\[
 (R^1f_*(T_X\otimes M_z))_{t_*}=(A/(t-t_*))^2,
 \qquad (R^2f_*(T_X\otimes M_z))_{t_*}=A/((t-t_*)^2).    \tag{D.18}
\]
At $p_1$, write $z=n+k(3\beta_1+4-2\rho)$, $k\ne0$, when this is
possible. Both coarse stalks are killed by $t$ and have length
$1+\epsilon_n$, with $\epsilon_n$ as in \ref{inv:D4.1}.
At $p_2$, write $z=n+k(2\beta_2+3-3i)$, $k\ne0$, when possible.
Both coarse stalks are killed by $t-1$ and have length one for
$n\equiv1,2\pmod4$ and length zero for $n\equiv0,3\pmod4$.
If the relevant lattice representation does not exist, both stalks
at that finite fibre are zero.
\end{proposition}


% ALL PRINTED PROOF MATERIAL MASKED







































































































































































































\begin{proposition}[The full twisted cohomology and its jump loci]\hypertarget{tag-01SS}{}\label{inv:D4.5}
\StatementRecord{inv:D4.5}{proved}{01SS}{\SeriesRef{inv:D4.1}, \SeriesRef{inv:D4.2}, \SeriesRef{inv:D4.4}, \SeriesRef{inv:C3.1}, \SeriesRef{inv:C3.2}}

For $z\notin\Z$, let $\mathcal R_z$ be the set of regular base
points satisfying (D.15), counted as points, not as marked integer
tuples. It is finite. Put $N_z=|\mathcal R_z|$. Define
$\ell_1(z)=0$ off $\Lambda_1$, and
$\ell_1(z)=1+\epsilon_n$ if
$z=n+k(3\beta_1+4-2\rho)$, $k\ne0$.
Define $\ell_2(z)=1$ if
$z=n+k(2\beta_2+3-3i)$, $k\ne0$, with
$n\equiv1,2\pmod4$, and $\ell_2(z)=0$ otherwise.
Then, setting $J_z=2N_z+\ell_1(z)+\ell_2(z)$,
\[
 \boxed{(h^0,h^1,h^2,h^3)(T_X\otimes M_z)=(0,J_z,1+J_z,0).} \tag{D.25}
\]
More precisely, with $\mathfrak m_t$ the maximal ideal at $t$,
\[
 R^1f_*(T_X\otimes M_z)=
   \bigoplus_{t\in\mathcal R_z}\C_t^2
      \oplus\C_{p_1}^{\ell_1(z)}\oplus\C_{p_2}^{\ell_2(z)},
\]
\[
 R^2f_*(T_X\otimes M_z)=\C_{p_0}\oplus
   \bigoplus_{t\in\mathcal R_z}\OO_{B,t}/\mathfrak m_t^2
      \oplus\C_{p_1}^{\ell_1(z)}\oplus\C_{p_2}^{\ell_2(z)}. \tag{D.26}
\]
Together with \ref{inv:D4.1} these formulas cover every element of
$\operatorname{Pic}(X)=\C$. The actual cohomology jump loci are
locally discrete subsets of this complex line.
\end{proposition}


% ALL PRINTED PROOF MATERIAL MASKED



























































\paragraph{Source and re-derivation of the tangent theorem.}
The tangent cohomology $(1,1,1,0)$ and the smooth one-dimensional
Kuranishi germ are theorems of Du, Guo, Kusakabe, and Xie,
arXiv:2610.00250v1, Section~6 \citep{DuEtAl2026v1}.
The invariant run re-derived these results after reading that paper.
The present proof prints the root-base calculation, its local
comparisons, and the additional coarse direct images and all Picard
twists. Version~2 \citep{DuEtAl2026v2} does not contain the deformation
section; this does not alter the credit for the version~1 theorems.



\subsection{Tangent cohomology for arbitrary coprime multiplicities}\hypertarget{tag-02ZF}{}\label{push:section:E.2}
\begin{theorem}[tangent cohomology and local miniversality]\hypertarget{tag-032U}{}\label{push:src:E.2.1}
\StatementRecord{push:src:E.2.1}{proved}{032U}{\SeriesRef{push:section:E.1}, \SeriesRef{push:src:E.3.1}}

For every sphere member with multiplicities $(3k,4l)$ and the
admissible logarithmic data of \ref{push:section:E.1}, put $z=kl$ and $d=3k+4l$.
On $\mathfrak B_{k,l}=\mathbb P(3k,4l)$ the root lift
$\bar f:X\to\mathfrak B_{k,l}$ satisfies
\[
 \begin{split}
 \mathcal E&=\mathcal O\oplus\mathcal O(-z),\\
 R^q\bar f_*\mathcal O_X&=
 \mathcal O,\quad\mathcal O\oplus\mathcal O(-z),\quad\mathcal O(-z),\\
 K_X&=\bar f^*\mathcal O(z-d).
 \end{split}
\]
The tangent cohomology is $(h^0,h^1,h^2,h^3)(X,T_X)=(1,1,1,0)$.
With the groupoid derivative convention
$\kappa(\Phi)=(\partial_c\Phi)\Phi^{-1}$, the class of
$\partial/\partial c$ is $e\otimes[\gamma]$.  The class of
$\partial/\partial u$ is
$-(2\pi i u)^{-1}e\otimes[\gamma]$.  Both are nonzero.
The Kuranishi germ is a smooth reduced line, and the fixed-data
period family is locally miniversal.
\end{theorem}


% TARGET PROOF MASKED: push:src:E.2.1

















































































































































































\begin{proposition}[local separation and canonical dimension]\hypertarget{tag-0335}{}\label{push:src:E.2.2}
\StatementRecord{push:src:E.2.2}{proved}{0335}{\SeriesRef{push:src:E.1.2}, \SeriesRef{push:src:E.2.1}}

The multiplicity pair $(3k,4l)$ is constant on every sufficiently
small deformation of one of the preceding members.  Distinct
pairs therefore determine disjoint open loci among these
isomorphism classes.  Moreover
\[
 \kappa(X)=\begin{cases}-\infty,&kl<3k+4l,\\1,&kl>3k+4l.
 \end{cases}
\]
Equality $kl=3k+4l$ never occurs under $\gcd(3k,4l)=1$.
\end{proposition}


% ALL PRINTED PROOF MATERIAL MASKED














\begin{proposition}[separation in irreducible analytic families]\hypertarget{tag-033B}{}\label{push:src:E.2.3}
\StatementRecord{push:src:E.2.3}{proved}{033B}{\SeriesRef{push:src:E.2.1}; \cite[Chapter III, Theorem 4.1]{BS76}}

Suppose a proper holomorphic submersion over a connected
irreducible reduced complex base has two fibres of the type
in \ref{push:src:E.2.1}.  Then their multiplicity pairs agree, even if some
other fibres have not been constructed by this method.
For these members the asymptotic number
\[
 s(X)=\lim_{r\to\infty}
 \frac{h^0(X,K_X^r)-h^0(X,K_X^{-r})}{r}
 =\frac1{12}-\frac1{3k}-\frac1{4l}
\]
recovers the ordered multiplicity pair uniquely.
\end{proposition}


% ALL PRINTED PROOF MATERIAL MASKED

































\subsection{Analytic foundations and the version 1}\hypertarget{tag-02ZG}{}\label{push:section:E.3}
\begin{cited}[base change and duality]\hypertarget{tag-033I}{}\label{push:src:E.3.1}
\StatementRecord{push:src:E.3.1}{proved}{033I}{\SeriesRef{push:src:E.2.1}, \SeriesRef{fam:src:A.5.2}; \cite{DuEtAl2026v1}; \cite[Section 6, Hauptsatz I, and Section 7, Satz 5]{Grauert1960D}; \cite{Grauert1963ErratumD}; \cite[Chapter III, Theorem 4.1]{BS76}}

For the proper flat root-curve morphisms in \ref{push:src:E.2.1}, the canonical
geometric-point comparison
\[
 (R^q\bar f_*\mathcal O_X)\otimes k(b)
 \longrightarrow H^q(X_b,\mathcal O_{X_b})
\]
is an isomorphism for $q=0,1,2$, and the direct images are locally
free.  The trace induces the canonical isomorphism
\[
 (R^2\bar f_*\mathcal O_X)^\vee
 \simeq\bar f_*\omega_{X/\mathfrak B}.
\]
Both comparisons descend on the analytic root-curve groupoid.
These conclusions, and hence the use made of them in \cref{fam:src:A.5.2}, require
no result specific to the version 1 of
\cite{DuEtAl2026v1}.
\end{cited}
The expanded source derivation is retained in the proof compendium, Paper~I, label \nolinkurl{push:src:E.3.1}.




This is analytic base change and duality, with the source hypotheses specified in the re-derivation. A complete re-derivation is given in \CompendiumRef{I}{app:rederive:push-src-E.3.1}.


Grauert's original direct-image theorem
\cite[Section 6, Hauptsatz I, and Section 7, Satz 5]{Grauert1960D}
was also opened together with its correction \cite{Grauert1963ErratumD}.
The correction explicitly warns that the exceptional locus in
Section 7, Satz 2 need not be closed.  \ref{push:src:E.3.1} uses the
constant-cohomology theorem on every chart and the later
base-change formulation, so it does not make that erroneous
closedness inference.  For a generic argument over a smooth
curve one can instead use the finite free complex of
\cite[Chapter III, Theorem 4.1]{BS76}; its rank-jump minors
give the required local analytic exceptional sets.


% Compendium-only push:src:E.3.2: apparatus: historical/current dependency audit of the five credited DGKX version1 theorems; complete mathematical applications retained at their actual receiving results and in the credited-result appendix



% Compendium-only push:src:E.3.3: apparatus: exact scope comparison for the five DGKX version1 standard(3,4) results; the source-specific receiving results are separately retained



\section{Other logarithmic transformations and finite covers}

\subsection{Comparison under logarithmic transformation}

\begin{proposition}[The marked geometric dictionary on the first page]\hypertarget{tag-01S2}{}\label{inv:B1.0}
\StatementRecord{inv:B1.0}{proved}{01S2}{the definitions and constructions stated in the proof}

The theta quotient identifies the period family with the line-bundle
quotient in Engel's construction. In the marking and reference
linearization fixed in the conventions, $u=w^2$. Replacing that
reference multiplies $u$ by the corresponding nonzero constant.
\end{proposition}


% ALL PRINTED PROOF MATERIAL MASKED



















































\begin{proposition}[The relative cohomology comparison]\hypertarget{tag-01S3}{}\label{inv:B1.1}
\StatementRecord{inv:B1.1}{proved}{01S3}{the definitions and constructions stated in the proof}

Fix the period family and cusp filling. Let $X'$ be obtained by any
admissible translation twists and multiplicities $m_1=3k$, $m_2=4l$,
where $k,l$ are positive integers. Write $D'=W+S'_1+S'_2$ and
\[
E'=\Omega^1_{X'}(\log D')/f'^*\Omega^1_B(\log\Sigma).
\]
The identity on the common smooth family induces isomorphisms on $B$
\[
R^qf_*\OO_X\simeq R^qf'_*\OO_{X'},\qquad
R^qf_*E\simeq R^qf'_*E'\quad(0\le q\le2).
\]
These isomorphisms identify the logarithmic connecting maps $\kappa_q$.
They also hold when the cusp collar is changed by an admissible translation.
\end{proposition}


% ALL PRINTED PROOF MATERIAL MASKED





















































\begin{proposition}[All admissible twists have the same Hodge numbers]\hypertarget{tag-01S4}{}\label{inv:B1.2}
\StatementRecord{inv:B1.2}{proved}{01S4}{\SeriesRef{inv:B1.1}}

For every family in \ref{inv:B1.1} whose cusp is the one-component hexagonal filling,
the direct images of the absolute cotangent sheaves are
\[
\begin{array}{c|ccc}
 &R^0f'_*&R^1f'_*&R^2f'_*\\\hline
\Omega^1_{X'}&\OO(-2)&\OO\oplus\OO(-1)^2&\OO\oplus\OO(-1)\\
\Omega^2_{X'}&\OO(-2)\oplus\OO(-1)&\OO(-2)\oplus\OO(-1)^2&\OO
\end{array}
\]
and $R^qf'_*\OO_{X'}=(\OO,\OO\oplus\OO(-1),\OO(-1))$.
In particular every admissible $(3,4)$ member with
$|12l_0-4l_1-3l_2|=1$ has the Hodge diamond of Part A.
\end{proposition}


% ALL PRINTED PROOF MATERIAL MASKED
































\subsection{Higher multiplicities and the anticanonical ring}

\begin{proposition}[Canonical bundles for all coprime multiplicities]\hypertarget{tag-01S5}{}\label{inv:B2.1}
\StatementRecord{inv:B2.1}{proved}{01S5}{\SeriesRef{inv:B1.1}, \SeriesRef{inv:B1.2}; \citet[Theorem 1.1 and Proposition 2.3]{ViaclovskyLog2026}}

Let $X=X_{3k,4l}(u)$ with admissible translation data as in
\citet[Theorem 1.1 and Proposition 2.3]{ViaclovskyLog2026},
$\gcd(3k,4l)=1$, and any nonzero $u$ in its construction domain.
Then
\[
K_X=f^*\OO_B(-1)\otimes
\OO_X\bigl((k-1)S_1+(3l-1)S_2\bigr).
\]
Moreover $(h^{0,q})_{q=0}^3=(1,1,0,0)$, $h^{p,0}=0$ for $p>0$,
and $(h^{1,1},h^{1,2})=(2,1)$ for every such pair.
Thus these Hodge numbers do not distinguish the multiplicity families.
\end{proposition}


% ALL PRINTED PROOF MATERIAL MASKED




























\begin{proposition}[The complete graded anticanonical ring]\hypertarget{tag-01S6}{}\label{inv:B2.2}
\StatementRecord{inv:B2.2}{proved}{01S6}{\SeriesRef{inv:B2.1}}

Put $m_1=3k$, $m_2=4l$, and $d=4l+3k-kl$.
Choose integers $r,s$ with $m_2r+m_1s=1$ and set
$\mathcal L=\OO_X(rS_1+sS_2)$.
There is a graded algebra isomorphism
\[
\bigoplus_{n\ge0}H^0(X,-nK_X)
\simeq\bigoplus_{n\ge0}\C[x,y]_{dn},
\qquad\deg x=m_2,\quad\deg y=m_1,
\]
where the right-hand summand is assigned degree $n$.
For negative weights the polynomial space is zero. Explicitly
\[
h^0(X,-nK_X)=
\max\left\{1+n-\left\lceil\frac{n(k-1)}{3k}\right\rceil
-\left\lceil\frac{n(3l-1)}{4l}\right\rceil,\,0\right\}.
\]
The integer $d$ is nonzero for admissible pairs. If $d<0$ the
anticanonical ring is $\C$; if $d>0$ it is the indicated weighted
Veronese ring. For $(k,l)=(1,1)$ it is
$\C[U,V]$, $\deg U=1$, $\deg V=2$.
\end{proposition}


% ALL PRINTED PROOF MATERIAL MASKED


















































For completeness the requested small pairs give the following exact data.
Here $r_1,r_2$ are the coefficients in \ref{inv:B2.1} and $P_n^-=h^0(-nK_X)$.
\[
\begin{array}{c|c|c|c|c|ccc}
(k,l)&(m_1,m_2)&(r_1,r_2)&(\deg x,\deg y)&d&P^-_1&P^-_2&P^-_3\\\hline
(1,1)&(3,4)&(0,2)&(4,3)&6&1&2&2\\
(1,2)&(3,8)&(0,5)&(8,3)&9&1&1&2\\
(3,1)&(9,4)&(2,2)&(4,9)&10&0&1&1\\
(1,4)&(3,16)&(0,11)&(16,3)&15&1&1&1\\
(5,1)&(15,4)&(4,2)&(4,15)&14&0&1&1\\
(1,5)&(3,20)&(0,14)&(20,3)&18&1&1&1
\end{array}
\]
The program \texttt{b\_families.py} enumerates all coprime pairs with
$m_1m_2\le60$ and checks the floor formula against exact weighted
monomial counting for $0\le n\le12$. These checks accompany the
all-degree proof above.

\subsection{Finite fundamental groups and universal covers}

\begin{proposition}[Connected cyclic covers and the cusp comparison]\hypertarget{tag-01S7}{}\label{inv:B3.1}
\StatementRecord{inv:B3.1}{proved}{01S7}{\SeriesRef{inv:B3.3}, \SeriesRef{inv:B1.1}}

For an admissible $(3,4)$ member with
$n=|12l_0-4l_1-3l_2|>1$, let $\pi:Y\to X$ be the universal
cover and $g=f\pi$. It is a finite etale holomorphic map of degree $n$.
Every fibre of $g$ is connected; its finite reductions are bielliptic,
its multiplicities are $3,4$, and its cusp $W_n$ normalizes to a
disjoint union of $n$ copies of $\mathrm{dP}_6$.
The pullback maps induce isomorphisms
\[
R^qf_*\OO_X\simeq R^qg_*\OO_Y,\qquad
R^qf_*E_X\simeq R^qg_*E_Y.
\]
\end{proposition}


% ALL PRINTED PROOF MATERIAL MASKED

































































































\begin{proposition}[Cotangent direct images and Hodge numbers of the covers]\hypertarget{tag-01S8}{}\label{inv:B3.2}
\StatementRecord{inv:B3.2}{proved}{01S8}{\SeriesRef{inv:B3.1}, \SeriesRef{inv:B1.2}}

Write $T_n=\C_{p_0}^{\,n-1}$. Then
\[
\begin{array}{c|ccc}
 &R^0g_*&R^1g_*&R^2g_*\\\hline
\Omega^1_Y&\OO(-2)&\OO\oplus\OO(-1)^2\oplus T_n&\OO\oplus\OO(-1)\\
\Omega^2_Y&\OO(-2)\oplus\OO(-1)&\OO(-2)\oplus\OO(-1)^2&\OO\oplus T_n
\end{array}
\]
The splittings involving $T_n$ are not canonical.
The only nonzero Hodge numbers of $Y$ are
\[
h^{0,0}=h^{0,1}=h^{3,2}=h^{3,3}=1,\quad
h^{1,2}=h^{2,1}=1,\quad h^{1,1}=h^{2,2}=n+1.
\]
The original finite-fundamental-group threefold $X$ itself has the
same Hodge diamond as the simply connected $(3,4)$ member.
\end{proposition}


% ALL PRINTED PROOF MATERIAL MASKED





































\begin{proposition}[Canonical rings and rational topology of the covers]\hypertarget{tag-01S9}{}\label{inv:B3.3}
\StatementRecord{inv:B3.3}{proved}{01S9}{\SeriesRef{inv:B3.1}, \SeriesRef{inv:B2.2}, \SeriesRef{inv:B3.2}}

For $X$ and $Y$ as above,
\[
K_Y=\pi^*K_X=\OO_Y(-2S'_2),\quad
\bigoplus_{m\ge0}H^0(Y,-mK_Y)=\C[U,V],\quad
\deg(U,V)=(1,2).
\]
The same ring holds for $X$.
Over $\C$, the Betti numbers of $Y$ are
$(1,0,n-1,0,n-1,0,1)$; its Fr\"olicher sequence degenerates at $E_2$.
For every nontrivial character of the deck group, the character
part of Dolbeault cohomology consists of one class in $(1,1)$ and
one in $(2,2)$ and no others.
\end{proposition}


% ALL PRINTED PROOF MATERIAL MASKED































































































































































































\section{The integral circle quotient and its cusp replacement}\hypertarget{tag-02OO}{}\label{mapsI:circle-quotient}

We write $g,u,w,d$ for the ordered homology basis
$(\widehat\gamma,\widehat u,\widehat w,\widehat\delta)$ in this
section, and $\gamma$ for its first dual covector.  Thus $d$ is a
vector, whereas $\gamma$ is a covector.  The cusp coordinate is
$q=1/t$.  Positive target orientations are the complex orientation
on the base followed by the positive $\gamma$-circle orientation.
The complex orientation on $X$ fixes the source orientation.

\begin{proposition}[The primitive integral quotient and its finite fillings]\hypertarget{tag-02PA}{}\label{mapsI:src:A.1.1}
\StatementRecord{mapsI:src:A.1.1}{proved}{02PA}{\cite{AC2026}}

For the markings and the clockwise root generators of the family construction,
the homological matrices and twisting vectors are
\[
 A_1=\begin{pmatrix}1&0&0&0\\6&0&1&0\\-6&-1&-1&0\\-2&1&0&1\end{pmatrix},
 \qquad
 A_2=\begin{pmatrix}1&0&0&0\\0&0&-1&0\\-6&1&0&0\\3&0&1&1\end{pmatrix},
\]
\[
 M_0=\begin{pmatrix}1&0&0&0\\0&1&0&0\\0&1&1&0\\-1&0&0&1\end{pmatrix},
 \qquad v_1=g+2u-4w,\qquad v_2=-g-3u+3w.
\]
The common invariant vectors and covectors are $\mathbb Zd$ and
$\mathbb Z\gamma$, respectively.  The smallest invariant primitive
sublattice containing the cusp vanishing lattice is
$L=\mathbb Z\langle u,w,d\rangle=\ker\gamma$.
The formula
\[
 [z,\Pi(z)x]\longmapsto([z],\gamma(x)\bmod\mathbb Z)
 \label{mapsI:quotient-eq-1}
\]
defines a real-analytic fibrewise group epimorphism
$J\longrightarrow B^\circ\times S^1$ with connected kernel $T_L$.
At the finite fillings it descends to a smooth proper $T^3$-bundle
over a Seifert solid torus, of multiplicities $3$ and $4$ and
translation residues $1$ and $-1$.
\end{proposition}



% ALL PRINTED PROOF MATERIAL MASKED















































Engel's primitive invariant covector is named $\psi$ in his
Lemma~3.8 \cite[Lemma~3.8]{Engel2026}.  It is not named $\varepsilon$ there.  In his
common marking $\mathcal B_0$, the integral comparison is
\[
 P_E=\begin{pmatrix}0&-2&-1&-2\\2&-1&0&-1\\0&2&1&3\\-1&0&0&0\end{pmatrix},
 \qquad \psi P_E=-\gamma.
 \hypertarget{tag-02OZ}{}\label{mapsI:quotient-eq-2}
\]
Thus the quotient characters agree after reversing the target
circle orientation.  With the inverse root generators the
translation vectors are $a_{j-1}=-P_Ev_j/m_j$.
This is the precise sign convention in Proposition~\ref{fam:src:B.2.2}; it explains why equality of the unnamed primitive
covectors without a sign convention would be insufficient.

\begin{proposition}[Pointwise cusp obstruction and the corrected double curve]\hypertarget{tag-02PB}{}\label{mapsI:src:A.1.2}
\StatementRecord{mapsI:src:A.1.2}{proved}{02PB}{\SeriesRef{mapsI:quotient-eq-4}}

The literal character has no continuous extension at either
triple point or at generic points of two of the three double
curves.  It has the constant limit $0$ on the smooth locus of
$W$ and on the third double curve away from the triple points.
In the chart $q=z_0z_1z_2$, $x_1=z_1$, $x_2=z_2$, the extendible
double curve is $z_0=z_1=0$, $z_2\ne0$.  The fixed curve of the
$\mathbb C^*$-action in direction $d=e_2$ is instead
$z_0=z_2=0$, $z_1\ne0$, and is one of the nonextendible curves.
Consequently the extendible double curve and the fixed double curve are distinct.
\end{proposition}



% ALL PRINTED PROOF MATERIAL MASKED




























































The nonextension above is local and pointwise.  It does not
contradict extension of the cohomology class over a cusp
neighbourhood.  The toric covering has rays $(v,1)$ for all
$v\in\mathbb Z^2$.  Meridians about its toric divisors kill
the whole cocharacter lattice: $(0,1)$ is a ray and
$(e_i,1)-(0,1)=(e_i,0)$.  The toric cover is therefore simply
connected.  Its free period deck group is $\overline\Lambda$
and hence
\[
 \pi_1(N_0)=\overline\Lambda=\Lambda/K.
 \hypertarget{tag-02P2}{}\label{mapsI:quotient-eq-5}
\]
For the punctured cusp the map to this group kills $K$ and
the base meridian.  The latter bounds the disc given by a
point of the dense orbit on a central component.  Both
long-period covectors extend in integral $H^1$, since they
kill precisely these relations.  In contrast, an extension
of $\theta_\gamma$ to all of $X\simeq S^6$ would induce a
nonzero map $H_1(F)\to\mathbb Z$ factoring through
$H_1(X)=0$, which is impossible.  This last argument is the
global cohomological obstruction; it says nothing about
pointwise cusp limits.

\begin{proposition}[The filled three-dimensional target]\hypertarget{tag-02PC}{}\label{mapsI:src:A.1.3}
\StatementRecord{mapsI:src:A.1.3}{proved}{02PC}{the definitions and constructions stated in the proof}

For the original data the quotient over the complement of
the cusp disc is the exterior of the $(3,4)$ torus knot,
with the specified Seifert structure.  For general admissible
$(l_0,l_1,l_2)$ the filled quotient is the lens space determined
by the following meridian slopes in one oriented integral
torus basis $(\rho_1,h)$:
\[
 \mu_1=(3,-l_1),\qquad
 \mu_2=(-4,4l_0-l_2),\qquad
 \det(\mu_1,\mu_2)=p=12l_0-4l_1-3l_2.
 \label{mapsI:quotient-eq-6}
\]
In the lens-space convention that the second meridian is
$-q_0\mu_1+p\lambda_1$, where
$\det(\mu_1,\lambda_1)=1$, the lens parameter satisfies
$3q_0\equiv4\pmod p$.  In particular the target depends only
on the signed integer $p$, and is a three-sphere exactly
when $|p|=1$.
\end{proposition}



% ALL PRINTED PROOF MATERIAL MASKED









































\subsection{The repaired map and its relative uniqueness}

\begin{theorem}[The repaired map and its relative homotopy class]\hypertarget{tag-02PD}{}\label{mapsI:src:A.2.1}
\StatementRecord{mapsI:src:A.2.1}{proved}{02PD}{\SeriesRef{mapsI:quotient-eq-5}}

For every admissible member, there is a smooth map
$e_{\rm q}:X\to Q$ over $B$ whose restriction outside an arbitrarily
small cusp neighbourhood is the literal circle quotient.
For fixed quotient gluing, its homotopy class is independent
of the cusp replacement, provided the replacement maps the
cusp neighbourhood into its designated solid torus and is
kept fixed on its outer collar.  In particular there is a
well-defined class for the original member $X\simeq S^6$.
\end{theorem}



% ALL PRINTED PROOF MATERIAL MASKED



































\subsection{The Albanese coordinates and real structures}
\begin{proposition}[exact nomes and continuation radius]\hypertarget{tag-000L}{}\label{ari:C.2.1}
\StatementRecord{ari:C.2.1}{proved}{000L}{\SeriesRef{fam:src:A.4.1}, \SeriesRef{ari:C.1.2}, \SeriesRef{fam:src:C.2.3}, \SeriesRef{fam:src:A.5.1}; \cite{DuGuoKusakabeXie2026v1}}

For $w=e^{-\pi ic}$, the fixed branches give
\begin{align*}
 \kappa_1&=c+\tfrac12-\frac{3i}{2\pi}\log48,&
 q_1&=-w^2/48^3,\\
 \kappa_2&=c+\tfrac32-\frac{3i}{\pi}\log6,&
 q_2&=iw/216,\qquad q_1=(27/64)q_2^2.
\end{align*}
The family is parameterized by $0<|q_2|<1$ and exists precisely in
that domain. Writing $J(q)=q^{-1}+744+\cdots$, one has
\[
 j(\operatorname{Alb}S_1)=J(-w^2/48^3).
\]
This function is holomorphic on $0<|w|<48^{3/2}$, has a pole of order
two at zero, and has a natural boundary at $|w|=48^{3/2}$.
In particular it continues holomorphically through every point of the
family boundary $|w|=216$.
\end{proposition}


% ALL PRINTED PROOF MATERIAL MASKED







































\begin{proposition}[the real locus]\hypertarget{tag-000M}{}\label{ari:C.2.2}
\StatementRecord{ari:C.2.2}{proved}{000M}{\SeriesRef{fam:src:B.4.2}, \SeriesRef{fam:src:A.5.1}; \cite{DuGuoKusakabeXie2026v1}}

An antiholomorphic equivalence $X_c\to X_{c'}$ exists exactly when
$c'-(1-\bar c)\in2\mathbb Z$. Consequently $X_c$ admits an
antiholomorphic involution exactly when $q_2\in\mathbb R\setminus\{0\}$.
\end{proposition}


% ALL PRINTED PROOF MATERIAL MASKED





































































\subsection{The next local terms}
Put $G=g_2(i)=\Gamma(1/4)^8/(16\pi^2)$ and
$D_*=-3\log6/\pi$.
Our theta convention is
$\vartheta_j(z,\tau)=\theta_j(\pi z\mid\tau)$, with lattice
$\mathbb Z+\tau\mathbb Z$.  The heat identity and modular
transformations are \cite[Equations 20.13.1, 20.7.26, 20.7.30
and 23.18.5]{DLMFTheta2026}.  The classical differential system,
in the notation $P=E_2$, $Q=E_4$, $R=E_6$, is
\[
 xP_x=(P^2-Q)/12,\qquad
 xQ_x=(PQ-R)/3,\qquad
 xR_x=(PR-Q^2)/2;
\]
see \cite[Equations (1.1)--(1.4), p.~1]{ErbanVanGorder2021}.
Here $x=e^{2\pi i\tau}$, so
$d/d\tau=2\pi i x\,d/dx$.
For $E=\pi^2E_2/3$, $g_2=4\pi^4E_4/3$ and
$g_3=8\pi^6E_6/27$, this gives the exact normalization
\[
 \begin{split}
 E'&=\frac{i}{2\pi}(E^2-g_2/12),\\
 g_2'&=\frac{i}{\pi}(2Eg_2-3g_3),\qquad
 g_3'=\frac{i}{\pi}(3Eg_3-g_2^2/6).
 \end{split}
\]

\begin{proposition}[fourth order at the maximum]\hypertarget{tag-000N}{}\label{ari:C.3.1}
\StatementRecord{ari:C.3.1}{proved}{000N}{the definitions and constructions stated in the proof}

On either root branch at $\tau=i$, let
$h=\mu-(1-\tau)/2$ and $\delta=\tau-i$. Then
\begin{align*}
 \tau&=i-\frac{3\pi i}{2}h^2-\frac{9\pi^2i}{8}h^4+O(h^6),\\
 h^2&=\frac{2i}{3\pi}\delta+\frac1{3\pi}\delta^2+O(\delta^3),\\
 \mu&=\frac{1-\tau}{2}
  \mathbin{\pm}\left(\frac{2i\delta}{3\pi}\right)^{1/2}
          \left(1-\frac{i}{4}\delta+O(\delta^2)\right).
\end{align*}
The square root is taken holomorphically on the chosen root chart.
Furthermore
\begin{align}\hypertarget{tag-003M}{}\label{ari:eq:cs-Li}
 \log N={}&\tfrac12\log2-\frac\pi2|h|^2
   +\frac{3\pi^2}{8}|h|^4
   -\frac{3\pi^2}{4}|h|^2\operatorname{Re}(h^2)
   -\frac{5G}{192}\operatorname{Re}(h^4)+O(|h|^6),\\
 D_{\rm part}={}&D_* -\frac2\pi|\delta|
  +\frac1\pi\bigl(|\delta|^2+|\delta|\operatorname{Im}\delta\bigr)
  +\frac{5G}{72\pi^3}\operatorname{Re}(\delta^2)
  +O(|\delta|^3).\nonumber
\end{align}
In particular the leading terms agree with the boundary expansion in Lemma~\ref{push:src:D.2.1}.
\end{proposition}


% ALL PRINTED PROOF MATERIAL MASKED



































































\begin{proposition}[the equianharmonic point is neither an extremum]\hypertarget{tag-000O}{}\label{ari:C.3.2}
\StatementRecord{ari:C.3.2}{proved}{000O}{the definitions and constructions stated in the proof}

Set $T_0=\sqrt3/2$, $\mu_0=(2-\rho)/3$, and
$p=\wp'_\rho(\mu_0)$. Then $p\ne0$, $p^2=-g_3(\rho)$, and, for
$\delta=\tau-\rho$ on the order-three root chart,
\begin{align*}
 \mu&=\frac{2-\tau}{3}
      +\frac{3p}{16\pi^2}\delta^2
      +\frac{3ip}{32\pi^2T_0}\delta^3+O(\delta^4),\\
 t&=\frac{ip^2}{\pi^3}\delta^3
        \left(1+\frac{3i}{2T_0}\delta+O(\delta^2)\right),\\
 \log N&=\tfrac14\log3+\frac7{192}\operatorname{Re}t
       -\frac{9\pi}{512T_0|p|^{2/3}}|t|^{4/3}
       +O(|t|^{5/3}).
\end{align*}
Thus
\[
 D_{\rm part}(t)=-\frac3{2\pi}\log48
       +\frac7{32\pi}\operatorname{Re}t
       -\frac{27}{256T_0|p|^{2/3}}|t|^{4/3}
       +O(|t|^{5/3}).
\]
The point $t=0$ is neither a local maximum nor a local minimum.
\end{proposition}


% ALL PRINTED PROOF MATERIAL MASKED





































































\subsection{CM and torsion}
\begin{cited}[no smooth CM-torsion parameter]\hypertarget{tag-000U}{}\label{ari:C.5.1}
\StatementRecord{ari:C.5.1}{proved}{000U}{\citet[Theorem~4.15]{BourdonClarkStankewicz2017}; \citet[Theorem~20.12 and Corollary~20.13]{Sutherland2023CM}}

There is no $t\in\mathbb C\setminus\{0,1\}$ for which $E_t$ has
complex multiplication and $P_t$ is torsion. On the smooth models
at $t=0,1$ after base change of orders $3,4$, respectively, the
limiting points have orders $3,2$. Thus these are the only CM
torsion points on the completed family.
\end{cited}
The expanded source derivation is retained in the proof compendium, Paper~I, label \nolinkurl{ari:C.5.1}.




This applies \citet[Theorem~4.15]{BourdonClarkStankewicz2017} and \citet[Theorem~20.12 and Corollary~20.13]{Sutherland2023CM}; the exact finite enumeration is retained. A complete re-derivation is given in \CompendiumRef{I}{app:rederive:ari-C.5.1}.



\section{Global heights and the eleven rational CM fibres}

We use the equation
\[
 E_t:Y^2=X^3-\frac34t^3(t-1)X+\frac14t^4(t-1)^2,
 \qquad P_t=(0,t^2(t-1)/2).
\]
Its discriminant is $27t^8(t-1)^3$. Our canonical height is
$\widehat h_O(P)=\frac12\lim_{n\to\infty}n^{-2}h(x(nP))$.
In particular the bilinear Mordell--Weil self-pairing used in the
supplied surface calculation is $2\widehat h_O$.
For a finite place $v$ above the rational prime $p$, we normalize
$|p|_v=p^{-1}$ and set $n_v=[K_v:\Q_p]$. At an archimedean
place $n_v$ is one or two, according as the place is real or complex.

\begin{proposition}[the global formula]\hypertarget{tag-000V}{}\label{ari:C.6.1}
\StatementRecord{ari:C.6.1}{proved}{000V}{\SeriesRef{ari:C.7.1}; \citet[Section 2.1]{demarco2017}}

For every number field $K$, every $t\in K\setminus\{0,1\}$,
and the section just specified,
\begin{equation}\hypertarget{tag-0039}{}\label{ari:eq:ch-global}
 \sum_{\sigma:K\hookrightarrow\C}\log N(P_{\sigma(t)})
 =\sum_{v\nmid\infty}n_v\lambda_v(P_t)
          -[K:\Q]\widehat h_O(P_t).
\end{equation}
Here the local heights are invariant under change of Weierstrass
model and satisfy, when $nP\ne O$,
\begin{equation}\hypertarget{tag-003B}{}\label{ari:eq:ch-local-mult}
 \lambda_v(nP)=n^2\lambda_v(P)-\log|\psi_n(P)|_v
                   +\frac{n^2-1}{12}\log|\Delta|_v.
\end{equation}
At a complex embedding $\lambda_v(P)=-\log N(P)$; complex
embeddings are counted separately in \eqref{ari:eq:ch-global}.
For rational $t$, put $A_t=\exp(\sum_p\lambda_p(P_t))$.
Then $A_t$ is a positive real algebraic number and
\[
 N(P_t)=A_t e^{-\widehat h_O(P_t)},\qquad
 D_{\rm part}(t)=\frac6\pi
       \left(\log\frac{A_t}{2\sqrt3}-\widehat h_O(P_t)\right).
\]
Thus $N(P_t)$ is algebraic if and only if
$e^{\widehat h_O(P_t)}$ is algebraic.
\end{proposition}


% ALL PRINTED PROOF MATERIAL MASKED




































\begin{proposition}[exact finite local-height calculation]\hypertarget{tag-000W}{}\label{ari:C.6.2}
\StatementRecord{ari:C.6.2}{proved}{000W}{\SeriesRef{ari:eq:ch-local-mult}}

Let $E/\Q$ be a global minimal Weierstrass model. If $mP$ is
in the identity component at $p$, and $mP\ne O$, then
\begin{equation}\hypertarget{tag-003A}{}\label{ari:eq:ch-local-exact}
 \frac{\lambda_p(P)}{\log p}
 =\frac{v_p(\Delta)}{12}
   +\frac{\frac12\max(0,-v_p(x(mP)))-v_p(\psi_m(P))}{m^2}.
\end{equation}
For $m=1$ we read $\psi_1=1$. This formula, with exact rational
coordinates and valuations, gives all the finite contributions
in the tables below.
\end{proposition}


% ALL PRINTED PROOF MATERIAL MASKED


















\begin{proposition}[the CM models and their height identities]\hypertarget{tag-000X}{}\label{ari:C.6.3}
\StatementRecord{ari:C.6.3}{proved}{000X}{\SeriesRef{ari:eq:ch-local-exact}}

The eleven rational CM fibres have the global minimal models,
section coordinates, component classes, and finite local heights
in the following tables. On each model $a_1=a_2=0$, and we
write $[a_3,a_4,a_6]$ for
$y^2+a_3y=x^3+a_4x+a_6$. The displayed point is $(0,y_P)$.
The sign of $P$, which is immaterial to all these heights,
is chosen by the recorded change of variables.
The constants $A_t$ in increasing row order are
\[
 \sqrt6,\quad 2^{2/3}3^{1/4},\quad 2^{3/4},\quad
 2,\quad 2,\quad \sqrt3,\quad 1,\quad 3^{5/12},\quad
 1,\quad1,\quad1.
\]
\end{proposition}


% ALL PRINTED PROOF MATERIAL MASKED













































\begingroup\footnotesize\setlength{\tabcolsep}{4pt}
\begin{longtable}{r r r r r r}
\toprule Row & $j$ & $a_3$ & $a_4$ & $a_6$ & $y_P$\\\midrule\endhead
1 & 8000 & 0 & -5880 & 153664 & 392\\
2 & 54000 & 0 & -7260 & 234256 & 484\\
3 & 287496 & 0 & -2156 & 38416 & 196\\
4 & 16581375 & 0 & -7732620 & 8275904784 & 90972\\
5 & -3375 & 0 & -140 & 784 & -28\\
6 & -32768 & 1 & -3234 & 72630 & -270\\
7 & -884736 & 1 & -38 & 90 & -10\\
8 & -12288000 & 1 & -1920270 & 1024288020 & -32005\\
9 & -884736000 & 1 & -379260 & 89898842 & -9482\\
10 & -147197952000 & 1 & -347045930 & 2488447882842 & -1577482\\
11 & -262537412640768000 & 1 & -75065352465446420 & 7916029533155605844266842 & -2813543945482\\
\bottomrule\end{longtable}
\begin{longtable}{r r l c c l}
\toprule Row & $p$ & Type & Component & Order & $\lambda_p/\log p$\\\midrule\endhead
1 & 2 & $III^*$ & $c_P$ & 2 & $\frac{1}{2}$\\
1 & 3 & $I_0^*$ & $0$ & 1 & $\frac{1}{2}$\\
1 & 7 & $I_0^*$ & $c_P$ & 2 & $0$\\
2 & 2 & $I_{4}^*$ & $c_P$ & 2 & $\frac{2}{3}$\\
2 & 3 & $III$ & $0$ & 1 & $\frac{1}{4}$\\
2 & 11 & $I_0^*$ & $c_P$ & 2 & $0$\\
3 & 2 & $I_{5}^*$ & $c_P$ & 2 & $\frac{3}{4}$\\
3 & 7 & $I_0^*$ & $c_P$ & 2 & $0$\\
4 & 2 & $I_{8}^*$ & $c_P$ & 2 & $1$\\
4 & 3 & $I_0^*$ & $c_P$ & 2 & $0$\\
4 & 7 & $III$ & $c_P$ & 2 & $0$\\
4 & 19 & $I_0^*$ & $c_P$ & 2 & $0$\\
5 & 2 & $I_{8}^*$ & $c_P$ & 2 & $1$\\
5 & 7 & $III$ & $c_P$ & 2 & $0$\\
6 & 3 & $I_0^*$ & $0$ & 1 & $\frac{1}{2}$\\
6 & 7 & $I_0^*$ & $c_P$ & 2 & $0$\\
6 & 11 & $III$ & $c_P$ & 2 & $0$\\
7 & 19 & $III$ & $c_P$ & 2 & $0$\\
8 & 3 & $IV$ & $0$ & 1 & $\frac{5}{12}$\\
8 & 11 & $I_0^*$ & $c_P$ & 2 & $0$\\
8 & 23 & $I_0^*$ & $c_P$ & 2 & $0$\\
9 & 3 & $I_0^*$ & $c_P$ & 2 & $0$\\
9 & 7 & $I_0^*$ & $c_P$ & 2 & $0$\\
9 & 43 & $III$ & $c_P$ & 2 & $0$\\
10 & 7 & $I_0^*$ & $c_P$ & 2 & $0$\\
10 & 31 & $I_0^*$ & $c_P$ & 2 & $0$\\
10 & 67 & $III$ & $c_P$ & 2 & $0$\\
11 & 7 & $I_0^*$ & $c_P$ & 2 & $0$\\
11 & 11 & $I_0^*$ & $c_P$ & 2 & $0$\\
11 & 19 & $I_0^*$ & $c_P$ & 2 & $0$\\
11 & 127 & $I_0^*$ & $c_P$ & 2 & $0$\\
11 & 163 & $III$ & $c_P$ & 2 & $0$\\
\bottomrule\end{longtable}
\endgroup
\noindent The next table prints 52 significant digits. The last column is both the directly computed $\log N$ and the independently computed $\log A_t-\widehat h_O(P_t)$; the agreement is more precise than the digits printed.
\begingroup\small
\begin{longtable}{r l}
\toprule Row & Height and archimedean logarithm\\\midrule\endhead
1 & $\widehat h_O=0.68444758838339190721754413638128047178526153154571$\\
 & $\log N=0.2114321462306355931886945428090706645762338145457924$\\[3pt]
2 & $\widehat h_O=0.6417998032172362080432024005106271224631326545214062$\\
 & $\log N=0.09495138932308808775043032302545534908240674117445132$\\[3pt]
3 & $\widehat h_O=0.5335302122611885050129127004864161587766141055417204$\\
 & $\log N=-0.01366982684122952294998860939278373271998900477152897$\\[3pt]
4 & $\widehat h_O=0.9940668685687274300038154462317330832788241518412543$\\
 & $\log N=-0.300919688008782120586583324773556515203324017480999$\\[3pt]
5 & $\widehat h_O=0.4673656771814712800075202012628467836356203437033082$\\
 & $\log N=0.2257815033784740294097119201953297844398797906569471$\\[3pt]
6 & $\widehat h_O=0.4286091723813209970260194687110167110834518087280487$\\
 & $\log N=0.120696971952733848671603149750246141240293470183326$\\[3pt]
7 & $\widehat h_O=0.09064583701789963707894077069188164993051948122321821$\\
 & $\log N=-0.09064583701789963707894077069188164993051948122321821$\\[3pt]
8 & $\widehat h_O=0.7366208256267509672496883390906560396314066510934334$\\
 & $\log N=-0.2788657053483719291683361570396036626949522520006211$\\[3pt]
9 & $\widehat h_O=0.6018296198896332298783588845164397135769872973598109$\\
 & $\log N=-0.6018296198896332298783588845164397135769872973598109$\\[3pt]
10 & $\widehat h_O=1.002649286144114870320961646648539589847474352238484$\\
 & $\log N=-1.002649286144114870320961646648539589847474352238484$\\[3pt]
11 & $\widehat h_O=2.165489246796061077105970244492001901750597685962121$\\
 & $\log N=-2.165489246796061077105970244492001901750597685962121$\\[3pt]
\bottomrule\end{longtable}\endgroup


\begin{proposition}[the maximum as a height inequality]\hypertarget{tag-000Y}{}\label{ari:C.6.4}
\StatementRecord{ari:C.6.4}{proved}{000Y}{\SeriesRef{ari:C.8.5}, \SeriesRef{ari:eq:ch-global}; \citet{demarco2017}; \citet[pp.~183--184]{callSilverman1993}; \citet{silverman2006heights}}

For $t\in K\setminus\{0,1\}$,
\begin{equation}\hypertarget{tag-0038}{}\label{ari:eq:ch-bound}
 \widehat h_O(P_t)\geq
 \frac1{[K:\Q]}\sum_{v\nmid\infty}n_v\lambda_v(P_t)
                    -\frac12\log2.
\end{equation}
The inequality is strict on this smooth base. Over $\Q$ it
is $\widehat h_O(P_t)\geq\log A_t-\frac12\log2$.
The generic height is $\widehat h_O(P)=1/12$, and the
variation divisor is $D_P=(\infty)/12$. Consequently
\[
 \widehat h_O(P_t)=\frac1{12}h(t)+O(1)
 \qquad(t\in\overline\Q\setminus\{0,1\}),
\]
where the constant is uniform over algebraic $t$ and is for
the fixed family and section.
\end{proposition}


% ALL PRINTED PROOF MATERIAL MASKED




































































\begin{proposition}[rigorous high-precision height enclosures]\hypertarget{tag-000Z}{}\label{ari:C.6.5}
\StatementRecord{ari:C.6.5}{proved}{000Z}{\SeriesRef{ari:C.6.3}, \SeriesRef{ari:C.5.1}}

For every one of the eleven rational CM parameters in \ref{ari:C.6.3},
the files \path{records/C_cm_height_arb.json} and
\path{code/C_cm_height_arb.out} give certified enclosures of
$\log N(P)$ and $\widehat h_O(P)$ of radius less than
$10^{-236}$. These enclosures use no floating-point height
value as a proof step.
\end{proposition}


% ALL PRINTED PROOF MATERIAL MASKED























































































\subsection{Supplement: ranks, saturation, and central derivatives}

We retain the order of the eleven rational CM values in the stated conventions, and
write $E_i$ for the global minimal models recorded in the preceding
tables. The real period in this subsection is
\[
 \Omega_i=\int_{E_i(\R)}|\omega_i|,
\]
including both connected components when $E_i(\R)$ is disconnected.
The canonical height $\widehat h$ is relative to $(O)$; PARI's
\texttt{ellheight} and the regulator convention used in the BSD formula
are relative to $2(O)$, and hence equal $2\widehat h$ in rank one.
This factor of two is essential in every comparison below.

\begin{proposition}[ARI \ref{ari:C.6.R1}; Proved by descent and full saturation]\hypertarget{tag-0010}{}\label{ari:C.6.R1}
\StatementRecord{ari:C.6.R1}{proved}{0010}{\citet[Theorem 1]{CrCremonaSiksek2006}}

The algebraic and analytic ranks agree for the eleven curves. They are
two for $j=8000$ and $j=-12288000$, and one for all nine other values.
For every rank-one curve the section $P$ is a primitive generator of
$E(\Q)/E(\Q)_{\rm tors}$. Its index is therefore one. For the two
rank-two curves the index of $\Z P$ in this free quotient is infinite.
\[
\begin{array}{r|r|r|r}
i&j&\operatorname{rank}E(\Q)&[E(\Q)/E(\Q)_{\rm tors}:\Z P]\\\hline
1&8000&2&\infty\\
2&54000&1&1\\
3&287496&1&1\\
4&16581375&1&1\\
5&-3375&1&1\\
6&-32768&1&1\\
7&-884736&1&1\\
8&-12288000&2&\infty\\
9&-884736000&1&1\\
10&-147197952000&1&1\\
11&-262537412640768000&1&1
\end{array}
\]
\end{proposition}


% ALL PRINTED PROOF MATERIAL MASKED







































\begin{proposition}[ARI \ref{ari:C.6.R2}; Proved: central derivatives with explicit tails]\hypertarget{tag-0011}{}\label{ari:C.6.R2}
\StatementRecord{ari:C.6.R2}{proved}{0011}{\SeriesRef{ari:eq:Cr-tails}; \citet[Section 3]{CrDarmon2009}; \citet[Theorem 10.1]{CrRubin1999}}

The following are the values of $L'(E_i,1)/\Omega_i$; zeros in rows
one and eight are exact. More digits, the periods, derivatives and
tail bounds are in \path{results_Crank.csv}.
\begin{center}\small
\begin{tabular}{r|l}
$i$&$L'(E_i,1)/\Omega_i$\\\hline
1&$0$\\
2&$5.134398425737889664345619204085016979705061236171249722654226\ldots$\\
3&$2.134120849044754020051650801945664635106456422166881627967598\ldots$\\
4&$31.81013979419927776012209427941545866492237285892013668350005\ldots$\\
5&$1.869462708725885120030080805051387134542481374813232701178304\ldots$\\
6&$3.428873379050567976208155749688133688667614469824389504923957\ldots$\\
7&$0.362583348071598548315763082767526599722077924892872839203426\ldots$\\
8&$0$\\
9&$9.629273918234131678053742152263035417231796757756973886621556\ldots$\\
10&$16.04238857830583792513538634637663343755958963581574670348088\ldots$\\
11&$1247.321806154531180413038860827393095408344267114181914400926\ldots$
\end{tabular}
\end{center}
The two second derivatives, divided by the same real periods, are
\begin{align*}
 L''(E_1,1)/\Omega_1
 &=28.71408722758099118128474692322079946598060512912402172882525\ldots,\\
 L''(E_8,1)/\Omega_8
 &=78.03222041637032223732443782051140882655135939512581188900823\ldots.
\end{align*}
In particular the leading Taylor coefficients for these two curves
are half the displayed second derivatives.
\end{proposition}


% ALL PRINTED PROOF MATERIAL MASKED



















































































































































































\begin{proposition}[ARI \ref{ari:C.6.R3}; Proved: the consequences of Gross--Zagier and Kolyvagin]\hypertarget{tag-0012}{}\label{ari:C.6.R3}
\StatementRecord{ari:C.6.R3}{proved}{0012}{\citep{CrGrossZagier1986}; \citet[Theorem 1.3]{CrGross1991}; \citet[Theorem 1.1]{CrJetchev2008}}

For each of the nine rank-one curves there is a positive rational
number $q_i$ such that
\begin{equation}\hypertarget{tag-0021}{}\label{ari:eq:Cr-rational-factor}
 \frac{L'(E_i,1)}{\Omega_i}=q_i\widehat h(P_i).
\end{equation}
Their Tate--Shafarevich groups over $\Q$ are finite. Consequently
the corresponding $D_{\rm part}$ is an explicit logarithm minus
$(6/(\pi q_i))L'(E_i,1)/\Omega_i$. For the two rank-two curves, $L'(E_i,1)=0$ while
$\widehat h(P_i)>0$.
\end{proposition}


% ALL PRINTED PROOF MATERIAL MASKED




















































\paragraph{What the numerical BSD comparison does and does not determine.}
For a primitive rank-one point, the BSD prediction in these conventions is
\[
 \frac{L'(E,1)}{\Omega_E}
 =\frac{2\#\operatorname{Sha}(E/\Q)\prod_pc_p}
        {\#E(\Q)_{\rm tors}^{\,2}}\,\widehat h(P).
\]
The independently evaluated quotient after dividing by
$2\widehat h(P)\prod_pc_p/\#E(\Q)_{\rm tors}^2$ is one, to
more than 69 digits, for $i=2,3,4,5,6,7,9,10$, and nine for
$i=11$. Thus the numerically predicted rational factors $q_i$ are
\[
 8,4,32,4,8,4,16,16,576,
\]
in that order. This initial numerical comparison alone proves neither
these exact factors nor the orders $1$ or $9$ of $\operatorname{Sha}$.
The exact factors are subsequently proved in \ref{ari:C.6.R4} and \ref{ari:C.6.R5}:
the latter supplies an independent Gross--Zagier denominator bound
and rigorous intervals for all remaining ratios. This determines
the analytic BSD quotients exactly, while the orders of
$\operatorname{Sha}$ are proved here only in the three cases of
\ref{ari:C.6.R4}. In particular primitivity alone does not determine
$\#\operatorname{Sha}$ for the last curve.

\begin{cited}[ARI \ref{ari:C.6.R4}; Known input: three exact leading factors]\hypertarget{tag-0013}{}\label{ari:C.6.R4}
\StatementRecord{ari:C.6.R4}{proved}{0013}{\SeriesRef{ari:C.6.R1}, \SeriesRef{ari:C.6.R2}; \citet[Definition~1.1, Theorem~1.2, and Table~9]{miller2011}}

For $i=3,5,7$, namely $j=287496,-3375,-884736$, one has
\[
 \operatorname{Sha}(E_i/\Q)=0,\qquad
 \frac{L'(E_i,1)}{\Omega_i}=4\widehat h_O(P_i).
\]
Thus the factors $q_3=q_5=q_7=4$ are determined exactly by the integer bounds in the cited application. For the smallest of these examples,
$E:y^2+y=x^3-38x+90$ and $P=(0,9)$, this gives the exact
reduction
\[
 N(P)=\exp\left(-\frac{L'(E,1)}{4\Omega_E}\right).
\]
\end{cited}

The expanded source derivation is retained in the proof compendium, Paper~I, label \nolinkurl{ari:C.6.R4}.




The three exact cases specialize \citet[Definition~1.1, Theorem~1.2 and Table~9]{miller2011}. The complete application is retained in the proof compendium, Paper~I, label \nolinkurl{app:rederive:ari-C.6.R4}.


\paragraph{The scope of Rubin's theorem.}
The applicable ordinary CM results concern a prime $p$ at which
$E$ has good ordinary reduction, equivalently $p$ splits in the CM
field. One precise opened formulation of Rubin's rank-one formula
is \citet[Theorem 1]{CrBDP2012}: for a curve over $\Q$ with CM
by the maximal order, an ordinary prime prime to the conductor,
and a simple central zero, an appropriate Katz $p$-adic $L$-value
is, up to a nonzero algebraic factor, the square of a formal
logarithm of a nontorsion rational point divided by a $p$-adic
period. Its input is analytic rank one, already verified here.
It is a $p$-adic identity; it does not identify the real exponential
$\exp\widehat h(P)$ as algebraic.

For the $p$-part of the real leading-term formula, a modern precise
version is \citet[Theorem A]{CrCastella2025}. In addition to analytic
rank one it requires full CM, abelian torsion field over the CM
field, $p\nmid6h_K$, splitting of $p$, conductor prime to $p$,
and that the different of $K$ divide the Hecke conductor exactly.
For the maximal odd-discriminant cases here, the latter follows
from the exponent two of the ramified prime in the rational
conductor. The nonmaximal cases of discriminant $-12$ and $-28$
can be passed through a $2$-isogeny at $p>3$. The discriminant
$-16$ case requires care: the direct application of this version
has a different-conductor hypothesis at $2$ that fails, and we
do not claim that this particular formulation proves its full
leading-term formula. The asserted application retains the exact hypotheses of
\citet[Theorem~A]{CrCastella2025}.

In particular these ordinary-prime statements do not prove that
the $3$-primary part of $\operatorname{Sha}$ has order $9$ for the
last curve: $3$ is inert in $\Q(\sqrt{-163})$ and also excluded
from the stated $p>3$ hypotheses. \ref{ari:C.6.R5} proves the corresponding
analytic quotient to be exactly $9$ by a separate denominator
argument. The identification with the actual Tate--Shafarevich
group remains a separate arithmetic question in that row.


\subsection{Supplement: exact leading factors from a genus-character denominator bound}

We retain the full real period $\Omega_i$ and the absolute canonical
height $\widehat h_O$ of the preceding subsection. The argument below
uses the Gross--Zagier formula to bound a rational denominator; it does
not assume the leading-term Birch--Swinnerton-Dyer formula.

\begin{proposition}[ARI \ref{ari:C.6.R5}; Proved: all remaining rational leading factors]\hypertarget{tag-0014}{}\label{ari:C.6.R5}
\StatementRecord{ari:C.6.R5}{proved}{0014}{\SeriesRef{ari:C.6.R4}, \SeriesRef{ari:C.6.R2}, \SeriesRef{ari:C.6.R3}, \SeriesRef{ari:C.6.R1}, \SeriesRef{ari:eq:Cgz-product}, \SeriesRef{ari:eq:Cr-first}, \SeriesRef{ari:C.6.5}, \SeriesRef{ari:eq:Cgz-central-integrality}, \SeriesRef{ari:eq:Cgz-denominator}; \citet[Appendix B, p.~102, unnumbered $26569\mathrm a1/26569\mathrm a2$
block]{CgzBelabas2019}; \citet[Theorem 1.1, pp.~2524--2525]{CgzCai2014}; \citet[\S7.1.1, pp.~218--219, Theorem 7.1.2]{CgzYZZ2011}; \citet[\S3, pp.~4--5]{CrDarmon2009}; \citet[Theorem 9.18]{CgzStein2007}}

For the six rank-one rows $i=2,4,6,9,10,11$, one has
\[
 \frac{L'(E_i,1)}{\Omega_i}
   =q_i\widehat h_O(P_i),\qquad
 (q_2,q_4,q_6,q_9,q_{10},q_{11})=(8,32,8,16,16,576).
\]
Together with \ref{ari:C.6.R4}, this determines every rank-one factor:
\[
\begin{array}{c|rrrrrrrrr}
 i&2&3&4&5&6&7&9&10&11\\\hline
 q_i&8&4&32&4&8&4&16&16&576.
\end{array}
\]
The certified auxiliary ratios and denominator bounds used to identify
the six additional factors are as follows. Here $D$ is the auxiliary
negative fundamental discriminant,
$r_i=L(E_i^D,1)/\Omega_{i,D}$, and the denominator of $q_i$ is at most
$B_i$.
\[
\begin{array}{c|r|r|r|r|r}
 i&D&r_i&B_i&q_i&|a_i|\\\hline
 2&-23&64&256&8&32\\
 4&-31&64&256&32&64\\
 6&-19&8&32&8&8\\
 9&-19&16&64&16&16\\
 10&-3&400&1600&16&80\\
 11&-3&64&256&576&192
\end{array}
\]
The last column is the absolute genus-trace integer in the proof. These
identities determine the analytic BSD quotients exactly: they are one
for $i=2,4,6,9,10$ and nine for $i=11$. They do not by themselves
identify those numbers with the orders of the Tate--Shafarevich groups.
\end{proposition}



% ALL PRINTED PROOF MATERIAL MASKED





































































































































































































































































































































































































\subsection{Supplement: the three specified smooth torsion fibres}

\begin{proposition}[T1; Proved, with exact local and descent computations]\hypertarget{tag-0015}{}\label{ari:C.6.T1}
\StatementRecord{ari:C.6.T1}{proved}{0015}{the definitions and constructions stated in the proof}

At $t=32/5$ a global minimal model and the image of the section are
\[
 E:y^2=x^3-2x+1,\qquad P=(0,1),\qquad \Delta=80.
\]
The conductor is $40$, $E(\Q)=\Z/4\Z$ is generated by $P$,
and $\widehat h(P)=0$. The index of $\Z P$ in the quotient
modulo torsion is one. The complete bad-prime data are
\[
\begin{array}{c|c|c|c|c}
p&\operatorname{ord}_p\Delta&\text{type}&\text{component of }P&\lambda_p(P)\\\hline
2&4&\mathrm{III}&\text{unique nonzero class, order }2&\log2/12\\
5&1&\mathrm I_1&\text{identity}&\log5/12.
\end{array}
\]
All other finite local heights vanish. Consequently
\[
 \log N=\frac{\log10}{12},\qquad
 D_{\rm part}=\frac6\pi\log\frac{10^{1/12}}{2\sqrt3}.
\]
Numerically the two sides of the global height identity equal
\[
 0.191882091082837140334832621223697017300091790719064414669444\ldots.
\]
With the full real period $\Omega=5.9376498936895458116293755517\ldots$,
the central derivative is
\[
 \frac{L'(E,1)}\Omega=
 0.071331625531738774092617964261659879861112925567606085287841\ldots.
\]
The analytic rank is zero; a vanishing first derivative is therefore
not expected even though the point $P$ is torsion.
\end{proposition}


% ALL PRINTED PROOF MATERIAL MASKED


























\begin{proposition}[T2; Proved, with exact number-field local arithmetic]\hypertarget{tag-0016}{}\label{ari:C.6.T2}
\StatementRecord{ari:C.6.T2}{proved}{0016}{\SeriesRef{ari:C.4.4}}

Put $s=\sqrt3$ and
$t_\pm=(3008\pm1728s)/11$. Over $K=\Q(s)$, a global minimal
model for the fibre at $t_-$ is
\[
 E_-:y^2=x^3-(18+10s)x+21+12s,\qquad P_-=(0,3+2s).
\]
The model and point for $t_+$ are its conjugates. The
discriminant of this model is $1032912+596352s$. Define
\[
 \mathfrak p_2=(2,1+s),\qquad\mathfrak p_3=(s),\qquad
 \mathfrak p_{11,-}=(11,s+5),\quad
 \mathfrak p_{11,+}=(11,s-5).
\]
The bad-prime data of $E_-$ are
\[
\begin{array}{c|c|c|c|c|c}
v&e_v&v(\Delta)&\text{type}&\text{component order of }P_-&\lambda_v(P_-)\\\hline
\mathfrak p_2&2&8&\mathrm{IV}^*&3&0\\
\mathfrak p_3&2&3&\mathrm{III}&2&0\\
\mathfrak p_{11,-}&1&1&\mathrm I_1&1&\log11/12.
\end{array}
\]
For $E_+$ replace $\mathfrak p_{11,-}$ by
$\mathfrak p_{11,+}$. The nonzero component at
$\mathfrak p_3$ is unique; the two nonzero classes at
$\mathfrak p_2$ are interchanged by $P\mapsto-P$ and the
class of $P$ is a generator. All omitted local heights vanish.
Both points have order six, and
\begin{align*}
 \log N(P_-)&=\frac1{12}\log(17+10\sqrt3),\\
 \log N(P_+)&=\frac1{12}\log(10\sqrt3-17),\\
 \sum_{\sigma:K\hookrightarrow\R}\log N(P_\sigma)
 &=\frac{\log11}{12}
 =\sum_{v<\infty}n_v\lambda_v(P)-2\widehat h(P).
\end{align*}
The common value in the last line is
\[
 0.199824606066530878671828631497094108318475571161451431268214\ldots.
\]
In particular $A_\mp=(17+10\sqrt3)^{1/12}$ and
$(10\sqrt3-17)^{1/12}$ respectively give
$D_{\rm part}(t_\mp)=(6/\pi)\log(A_\mp/(2\sqrt3))$.

The natural Mordell--Weil group here is $E_\pm(K)$; the
notation $E_{t_\pm}(\Q)$ requested in the stated conventions does not
define an elliptic curve over $\Q$, since its $j$-invariant
is irrational. Over $K$, both curves have rank one and torsion
group $\Z/6\Z$. Thus the index of the torsion section in the
free quotient is infinite. For example $(1,-1-s)$ has infinite
order on $E_-$.
\end{proposition}


% ALL PRINTED PROOF MATERIAL MASKED

















































\paragraph{The corresponding number-field $L$-derivative.}
For completeness, the natural replacement for the nonexistent
rational-curve quantity at $t_\pm$ can also be computed. Put
\[
 \Omega_K=\frac1{\sqrt{12}}
  \prod_{\sigma:K\hookrightarrow\R}
       \int_{E_\sigma(\R)}|\omega_\sigma|.
\]
The Hasse--Weil $L$-function over $K$ has analytic conductor
$12^2\cdot396=57024$, gamma parameters $(0,0,1,1)$, and sign
$-1$. Direct 100-digit evaluation gives
\begin{align*}
 L'(E/K,1)&=2.046708855992864213402012895411307398829979893338038716037006\ldots,\\
 L'(E/K,1)/\Omega_K
 &=0.260797680551692455136263550519918387218960539765929916318910\ldots.
\end{align*}
This is a value for the curve over $K$, common to the two
conjugates. If the period is instead defined as the product of
the two full real periods, omitting the discriminant factor,
the ratio is
$0.075285805535274833467509818428226475613981023404668273560620\ldots$.
Both conventions and all input periods are explicitly recorded
in \path{code/Crank_torsion_final.out}; neither value enters
the exact torsion height proof.



\subsection{Integral and Gaussian descriptions}

\begin{proposition}[the Mahler-measure identities]\hypertarget{tag-001T}{}\label{ari:C.8.6}
\StatementRecord{ari:C.8.6}{proved}{001T}{\SeriesRef{ari:C.8.5}}

Let $\Lambda\subset\C$ be a lattice of area $A$, let $v$ be a zero
of $\wp_\Lambda$, and let $d\mu=A^{-1}dA_z$ be Haar probability
measure on $\C/\Lambda$. Put $U(z)=\log N(z)$ and
$L=U(v)$. Then
\begin{align}
 6L&=\int_{\C/\Lambda}\log\left|\frac{\wp'(z)^2}{\wp(z)^3}\right|d\mu(z),
 \hypertarget{tag-002E}{}\label{ari:eq:ca-mahler}\\
 N(v)^2&=|\Delta(\Lambda)|^{1/6}
      \exp\left(-\int_{\C/\Lambda}\log|\wp(z)|d\mu(z)\right).
 \hypertarget{tag-002H}{}\label{ari:eq:ca-wpmean}
\end{align}
If $C(u)=1-(g_2/4)u^2-(g_3/4)u^3$, the pushforward of $d\mu$
under $u=1/\wp(z)$ is
\begin{equation}\hypertarget{tag-002F}{}\label{ari:eq:ca-nu}
 d\nu(u)=\frac1{2A}\frac{dA_u}{|u|\,|C(u)|}.
\end{equation}
In particular its unnormalized integral is exactly $2A$, and
\[
 N(v)^6=4\exp\left(\int_{\C}\log|C(u)|d\nu(u)\right).
\]
Writing $e_1,e_2,e_3$ for the three branch values and
$p(x)=\int\log|x-\wp(z)|d\mu(z)$, the maximum theorem is equivalently
\[
 \int\log|C|d\nu\leq\log2,\qquad
 p(0)\geq\frac{p(e_1)+p(e_2)+p(e_3)}3-\frac{\log2}3.
\]
Equality in either inequality occurs exactly when $g_3=0$.
\end{proposition}


% ALL PRINTED PROOF MATERIAL MASKED












































\begin{proposition}[the two complex beta integrals]\hypertarget{tag-001U}{}\label{ari:C.8.7}
\StatementRecord{ari:C.8.7}{proved}{001U}{\SeriesRef{ari:C.8.6}, \SeriesRef{ari:eq:ca-mahler}}

For the square and equianharmonic lattices respectively, the
integrals in \ref{ari:C.8.6} are
\[
 \int\log|C|\,d\nu=\log2,\qquad
 \int\log|C|\,d\nu=\frac32\log3-2\log2.
\]
They recover $N^2=2$ and $N^4=3$ with the same normalization.
\end{proposition}


% ALL PRINTED PROOF MATERIAL MASKED













































\begin{proposition}[the signed Gaussian identity]\hypertarget{tag-001V}{}\label{ari:C.8.8}
\StatementRecord{ari:C.8.8}{proved}{001V}{the definitions and constructions stated in the proof}

For $\Lambda=\Z+\tau\Z$, $T=\Im\tau$, and every $z\in\C$,
\[
 \sqrt{2T}\,|\eta(\tau)|^2N(z,\tau)^2
 =\sum_{m\in\frac12\Lambda}\epsilon(m)
                         e^{-2\pi|z-m|^2/T},
\]
where $\epsilon=-1$ on $\Lambda$ and $+1$ on the three other
cosets of $\Lambda$ in $\frac12\Lambda$.
Both sides and all their derivatives converge locally uniformly.
\end{proposition}


% ALL PRINTED PROOF MATERIAL MASKED







































\subsection{The Betti measure and torsion equidistribution}

We take $dd^c=(i/\pi)\partial\bar\partial$, so that
$dd^c\log|w|=\delta_0$. For a local complex coordinate $w$,
$dA_w=(i/2)dw\wedge d\bar w$.

\begin{proposition}[Refuted sign and normalization, with the corrected current identity proved]\hypertarget{tag-001W}{}\label{ari:C.8.9}
\StatementRecord{ari:C.8.9}{proved}{001W}{\SeriesRef{ari:C.3.2}, \SeriesRef{ari:eq:ca-cuspexp}}

On a local period chart write $\mu=a\tau+b$, with $a,b$ real,
and put $T=\Im\tau$. The positive Betti measure is
\[
 m_P=db\wedge da
 =\frac{|\mu_w-a\tau_w|^2}{T}\,dA_w.
\]
Thus the proposed form $da\wedge db$ has the wrong sign.
The globally single-valued function $L(t)=\log N(P_t)$ satisfies
on $\mathbb P^1(\C)$
\begin{equation}\hypertarget{tag-0026}{}\label{ari:eq:ca-current}
 dd^cL=\frac1{12}\delta_\infty-m_P,\qquad
               \int_{\mathbb P^1}m_P=\frac1{12}.
\end{equation}
There are no atoms of $m_P$ at $0$, $1$ or $\infty$.
Let $T_0=\sqrt3/2$, $\rho=(1+i\sqrt3)/2$, and
$p=\wp'_\rho((2-\rho)/3)$. The density in the ordinary
$t$ coordinate satisfies
\begin{align}
 \frac{dm_P}{dA_t}
 &=\frac{|t|^{-2/3}}{64T_0|p|^{2/3}}
                      \bigl(1+O(|t|^{1/3})\bigr),
                         &&t\longrightarrow0,\hypertarget{tag-0029}{}\label{ari:eq:ca-density0}\\
 \frac{dm_P}{dA_t}
 &=\frac{|t-1|^{-3/2}}{24\pi^2\sqrt{E_4(i)}}
                      \bigl(1+O(|t-1|^{1/2})\bigr),
                         &&t\longrightarrow1.\hypertarget{tag-002A}{}\label{ari:eq:ca-density1}
\end{align}
\end{proposition}


% ALL PRINTED PROOF MATERIAL MASKED










































































\begin{proposition}[the correctly normalized equidistribution]\hypertarget{tag-001K}{}\label{ari:C.8.10}
\StatementRecord{ari:C.8.10}{proved}{001K}{\SeriesRef{ari:C.6.4}, \SeriesRef{ari:C.8.9}; \citet[Theorem 1.1 and Corollary 1.2]{demarco2017}}

Let $D_n$ be the divisor, with intersection multiplicity, of
parameters $t\in\mathbb P^1\setminus\{0,1,\infty\}$ for which
$nP_t=O$. Then
\[
 \deg D_n=\frac{n^2}{12}+O(1),\qquad
 \frac1{n^2}[D_n]\ \longrightarrow\ m_P,\qquad
 \frac1{\deg D_n}[D_n]\ \longrightarrow\ 12m_P
\]
weakly as $n\to\infty$ through indices for which $D_n\ne0$.
For any nonrepeating sequence of individual algebraic torsion
parameters, their Galois-orbit probability measures likewise
converge to $12m_P$. Thus the probability limit requires the factor $12$.
\end{proposition}


% ALL PRINTED PROOF MATERIAL MASKED


























































\begin{proposition}[the two first-derivative identities]\hypertarget{tag-001L}{}\label{ari:C.8.11}
\StatementRecord{ari:C.8.11}{proved}{001L}{the definitions and constructions stated in the proof}

On an ordinary modulus chart put
\[
 a=\frac{\Im\mu}{T},\qquad
 W=\frac{\vartheta_1'(\mu,\tau)}{\vartheta_1(\mu,\tau)}
                         +2\pi ia,\qquad
 R=\frac{d\mu}{d\tau}-a.
\]
Then
\[
 L_\tau=\frac{W}{8\pi i}(W+4\pi iR),\qquad
 \det\frac{d}{d\tau}
 \begin{pmatrix}\tau&\mu\\\mu&\beta^{\rm part}/6\end{pmatrix}
       =\frac{(W+2\pi iR)^2}{4\pi^2}.
\]
Moreover
\[
 R=\frac{g_2}{6\pi i\wp'(\mu)}-\frac{W}{2\pi i},
 \qquad
 L_\tau=\frac{iW(3W\wp'(\mu)-2g_2)}
                         {24\pi\wp'(\mu)}.
\]
Multiplication by $6/\pi$ gives the critical-point factorization for
$D_{\rm part}$ in Proposition~\ref{push:src:D.4.4}.
\end{proposition}


% ALL PRINTED PROOF MATERIAL MASKED






































\section{Critical points and the global level curves}

We continue to write $L(t)=\log N(\mu(\tau),\tau)$, where
$t=j(\tau)/1728$ and $\wp_\tau(\mu)=0$. The sign and lattice
translation of $\mu$ do not affect any expression below. We use
Eichler and Zagier's unnumbered theorem on p.~399 of
\cite{EichlerZagier1982}, also stated as equation (6) and the
unnumbered theorem in Section~2 of \cite{DukeImamoglu2008}.
For the cusp branch in this problem their theorem says
\begin{equation}\hypertarget{tag-002I}{}\label{ari:eq:cc-EZ}
 \begin{aligned}
 \mu&=\frac12-i\beta_{\rm cusp}-
       \frac{36i\sqrt6}{\pi}\sum_{n\ge1}\frac{a_nq^n}{n^2},
 &\beta_{\rm cusp}&=\frac{\log(5+2\sqrt6)}{2\pi},\\
 \sum_{n\ge1}a_nq^n&=\frac{\Delta_0}{E_6^{3/2}},
 &\mu''&=144\pi i\sqrt6\frac{\Delta_0}{E_6^{3/2}}.
 \end{aligned}
\end{equation}
The branch of $E_6^{3/2}$ tends to one at the cusp. Thus the
formula applies directly on $\Im\tau>1$ and then by analytic
continuation away from the half-period branch points. Our use at $i$ is on the square-root chart specified explicitly below.
The original theorem and the hypergeometric reformulation concern
the logarithm $\mu$; neither asserts the global inequality for
$L$ proved here. Searches through these papers and their cited
work, and through the torus Green-function literature, did not
locate the radial monotonicity statement below.

\begin{proposition}[a radial derivative identity]\hypertarget{tag-001M}{}\label{ari:C.8.12}
\StatementRecord{ari:C.8.12}{proved}{001M}{\SeriesRef{ari:C.8.11}, \SeriesRef{ari:C.3.2}}

Put $p=\wp'_\tau(\mu)$ and
\[
 X=\frac{3Wp}{g_2},\qquad F=1-X,
 \qquad W=\frac{\vartheta_1'(\mu,\tau)}{\vartheta_1(\mu,\tau)}
              +\frac{2\pi i\Im\mu}{\Im\tau}.
\]
These are single-valued modular invariants wherever $g_2\ne0$.
They have continuous extensions at $t=0,1$, with
$X(0)=1/4$ and $X(1)=0$. At every $t\ne0,1$,
\begin{equation}\hypertarget{tag-002W}{}\label{ari:eq:cc-radial}
 (t-1)L_t=\frac{X(X-2)}{24},\qquad
 r\frac{d}{dr}L(1+re^{i\theta})
       =\frac{\Re(X(X-2))}{12}.
\end{equation}
Also, if $a=\Im\mu/\Im\tau$ and $R=\mu'-a$, then
\begin{equation}\hypertarget{tag-002J}{}\label{ari:eq:cc-FR}
 F=\frac{6\pi ip}{g_2}R.
\end{equation}
Consequently $|F|<1$ implies strict decrease along the ray
from $1$ through $t$.
\end{proposition}


% ALL PRINTED PROOF MATERIAL MASKED


































\begin{proposition}[the cusp exclusion]\hypertarget{tag-001N}{}\label{ari:C.8.13}
\StatementRecord{ari:C.8.13}{proved}{001N}{\SeriesRef{ari:eq:cc-EZ}, \SeriesRef{ari:eq:ca-exp}}

If $\Im\tau\ge7/5$, then $|F|<17/20$.
\end{proposition}


% ALL PRINTED PROOF MATERIAL MASKED

















































\begin{proposition}[the compact exclusion by one Cayley disc]\hypertarget{tag-001O}{}\label{ari:C.8.14}
\StatementRecord{ari:C.8.14}{proved}{001O}{\SeriesRef{ari:eq:cc-EZ}, \SeriesRef{ari:eq:cc-Fseries}, \SeriesRef{ari:eq:cc-ode}, \SeriesRef{ari:eq:cc-recurrence}, \SeriesRef{ari:eq:cc-power}, \SeriesRef{ari:eq:cc-finite-table}, \SeriesRef{ari:eq:cc-norms}, \SeriesRef{ari:eq:cc-Ffactored}, \SeriesRef{ari:eq:cc-scalar}, \SeriesRef{ari:C.8.12}}

Let $\chi=(\tau-i)/(\tau+i)$. On
$0<|\chi|\le2-\sqrt3$, the invariant $F$ satisfies $|F|<1$.
At the two equianharmonic points on the boundary its continuous
value is $3/4$.
\end{proposition}


% ALL PRINTED PROOF MATERIAL MASKED






















































































































































































































































































































































\begin{proposition}[every critical point and every level set]\hypertarget{tag-001P}{}\label{ari:C.8.15}
\StatementRecord{ari:C.8.15}{proved}{001P}{\SeriesRef{ari:C.8.13}, \SeriesRef{ari:C.8.14}, \SeriesRef{ari:eq:cc-radial}, \SeriesRef{ari:eq:cc-EZ}, \SeriesRef{fam:src:C.2.3}}

For every zero $\mu$ of $\wp_\tau$, the zero sets of the two
critical factors are exactly
\[
 \begin{aligned}
 W=0&\quad\Longleftrightarrow\quad
       \tau\in\operatorname{SL}_2(\Z)\{i,\rho\},\\
 3W\wp'_\tau(\mu)-2g_2=0&\quad\Longleftrightarrow\quad
       \tau\in\operatorname{SL}_2(\Z)\rho.
 \end{aligned}
\]
On the coarse $t$-plane the function $L$ has no critical point
away from $t=1$. The point $t=0$ is a noncritical $C^1$ point,
with $L_t(0)=7/384$. The point $t=1$ is the unique global
maximum, of value $\log\sqrt2$, and has the fractional-power
local maximum described in C.3. In particular there is no
saddle in $0<t<1$. On the root charts, $h=0$ over $t=1$
is a nondegenerate maximum, while $\delta=0$ over $t=0$
is a degenerate critical point with nonzero harmonic cubic
leading term and three alternating ascending and descending
sectors. These are all critical points on the root charts.

For every angle $\theta$, $r\mapsto L(1+re^{i\theta})$ is
strictly decreasing for $r>0$, from $\log\sqrt2$ to $-\infty$.
For each $c<\log\sqrt2$, the level set $\{L=c\}$ is one
Jordan curve, strictly star-shaped about $1$. Its bounded
complement is $\{L>c\}$; its unbounded complement is
$\{L<c\}$. The curve is real analytic except possibly at
$t=0$ when $c=\frac14\log3$, where it is $C^1$ with nonzero
tangent. The closed superlevel sets form nested topological
discs. At the maximum the level set is the single point $1$,
and above it the level set is empty. The corresponding
statements for $D_{\rm part}$ follow by its increasing affine
change from $L$.
\end{proposition}


% ALL PRINTED PROOF MATERIAL MASKED































































\subsection{An independent interval certificate}\hypertarget{tag-034N}{}\label{push:section:E.4}
\label{sec:e4}

We write $\eta_D$ for the Dedekind eta function and set
\[
 h=z-(1-\tau)/2,\qquad
 \mathcal N^2(z,\tau)=
 \left|\frac{\vartheta_3(h,\tau)}{\eta_D(\tau)}\right|^2
 \exp\left(-\frac{2\pi(\Im h)^2}{\Im\tau}\right).
\]
The half-period theta identity identifies this with the norm used in
\cref{fam:src:C.1.1}, \cref{fam:src:C.2.3}. The coefficient of $\Im c$ in $D_c$ is one. Consequently the
identity from the period calculation above which we use is
\[
 D_{\rm part}=\frac6\pi
 \log\left(\frac{\mathcal N}{2\sqrt3}\right).
\]
It follows directly by taking the imaginary part of the theta formula
for $\beta^{\rm part}$; its Gaussian factor removes
$6(\Im\mu)^2/\Im\tau$. Thus no numerical normalization is inferred
from the new computation. The analytic formula and the first
certificate are recorded in the period calculation above, \cref{fam:src:C.1.1}, \cref{fam:src:C.2.3}. The result below supplies independent
arithmetic evidence for that same conclusion, with the independent certificate below.

\begin{theorem}[independent certification of every box]\hypertarget{tag-033L}{}\label{push:src:E.4.1}
\StatementRecord{push:src:E.4.1}{proved}{033L}{\citet[Sections 2--3]{Johansson2016}; \citet{PythonFlint2026}}



For all 19701 rational rectangles in the supplied certificate, and for
each $\tau$ in the rectangle, both zeros of $\wp_\tau$ satisfy
$\mathcal N^2<2$. The new largest certified upper bound is less than
\[
 1.999877571767410525.
\]
The 19701 rectangles, 36 local rectangles and 174 excluded rectangles
have disjoint interiors and cover exactly
$[0,1/2]\times[17/20,2]$. Every local rectangle is contained in
$|\tau-i|<3/20000$, and every excluded rectangle satisfies
$|\tau|<1$ throughout.
\end{theorem}



% ALL PRINTED PROOF MATERIAL MASKED














































































































































\begin{theorem}[quantitative local maximum]\hypertarget{tag-033M}{}\label{push:src:E.4.2}
\StatementRecord{push:src:E.4.2}{proved}{033M}{\SeriesRef{ari:C.1.1}, \SeriesRef{ari:C.8.15}}

If $|\tau-i|\leq3/20000$ and $\wp_\tau(z)=0$, then
$\mathcal N(z,\tau)\leq\sqrt2$, with equality precisely at
$\tau=i$. The explicit numerical second proof is retained in the compendium;
the global norm theorem gives the following short proof.
\end{theorem}

% ALL PRINTED PROOF MATERIAL MASKED






























\clearpage\appendix

\appendix
\section{Credited proofs of parameter equivalence and deformation results}

\begin{remark}[Source comparison for the credited re-derivations]\hypertarget{tag-02YS}{}\label{night:P5.7}
\StatementRecord{night:P5.7}{proved}{02YS}{\SeriesRef{night:P5.3}, \SeriesRef{night:P5.4}, \SeriesRef{night:P5.5}, \SeriesRef{night:P5.6}; \cite[Theorems 1.3--1.4, p.~3; Theorem 6.12, pp.~81--82; Proposition 6.16, p.~85; Theorem 7.14, p.~95]{DuEtAl2026v1}; \cite[title, abstract, contents and introduction, pp.~1--3]{DuEtAl2026v2}}
For the fixed standard $(3,4,\infty)$ data, our re-derived tangent
cohomology $(1,1,1,0)$, smooth one-dimensional Kuranishi germ,
automorphism group $\mathbb C^*$, and equivalence
$X_c\simeq X_{c'}\Longleftrightarrow c'-c\in2\mathbb Z$
agree with the corresponding assertions of
\cite{DuEtAl2026v1}. Version~2, entitled
\emph{Oka Geometry of Complex Six-Spheres}, omits those
tangent-cohomology, miniversality, full-automorphism and exact-period
assertions. The tangent and deformation calculations in this appendix
were re-derived after reading version~1. Our proof of the exact
$2\mathbb Z$ equivalence was in the public working edition by
8 September 2026; version~1, dated 23 September 2026, states that
equivalence as its Theorem~1.4.
\end{remark}
% ALL PRINTED PROOF MATERIAL MASKED


































\subsection{The canonical divisor and the exact anticanonical ring}
\hypertarget{tag-03CU}{}\label{wb:sec:canonical-ring-recovery}

The analytic construction admits a further explicit calculation which does
not use its claimed identification with a sphere. Retain the actual proper
map $f:X\to B=\mathbb P^1$, the cusp $p_0$, the finite points $p_1,p_2$,
the reduced finite fibres $S_1,S_2$, and
$f^*(p_1)=3S_1$, $f^*(p_2)=4S_2$. All line-bundle isomorphisms below
refer to this constructed analytic space. No replacement period matrix or
new compactification is introduced.

\begin{theorem}[Canonical bundle and recovery of the original fibration]
\hypertarget{tag-03DM}{}\label{wb:thm:canonical-ring-recovery}
\StatementRecord{wb:thm:canonical-ring-recovery}{proved}{03DM}{the definitions and constructions stated in the proof}

The evaluation map
\[
 f^*(f_*\omega_{X/B})\longrightarrow\omega_{X/B}
\]
has zero divisor exactly $2S_2$. Consequently
\[
 \omega_X\simeq f^*\mathcal O_B(-p_2)\otimes\mathcal O_X(2S_2)
       \simeq\mathcal O_X(-2S_2),\qquad
 \omega_X^{-2}\simeq f^*\mathcal O_B(p_2).
 \tag{CR1}
\]
For every integer $k$ the actual direct image of the vertical divisor is
\[
 f_*\mathcal O_X(kS_2)
       =\mathcal O_B\bigl(\lfloor k/4\rfloor p_2\bigr)
 \tag{CR2}
\]
as subsheaves of the meromorphic-function sheaf, using pullback to identify
their sections. In particular, for each integer $m\geq0$,
\[
 h^0(X,\omega_X^{-m})=\lfloor m/2\rfloor+1,
 \qquad h^0(X,\omega_X^m)=0\quad(m>0).
 \tag{CR3}
\]
There is a graded algebra isomorphism
\[
 \mathbb C[U,V]\xrightarrow{\sim}
       \bigoplus_{m\geq0}H^0(X,\omega_X^{-m}),
       \qquad\deg U=1,\quad\deg V=2.
 \tag{CR4}
\]
The degree-two complete linear system has no base points and its map is
the original $f$, followed by a specified projective-coordinate map of $B$.
The degree-one system consists of the single divisor $2S_2$.
\end{theorem}



% ALL PRINTED PROOF MATERIAL MASKED



















































In the already specified additive Picard coordinate \eqref{inv:eq:C-factor12}, with
$L_{12}=f^*\mathcal O_B(1)$, the identities retain the original constants:
\[
 \mathcal O_X(S_1)=L_4,\qquad
 \mathcal O_X(S_2)=L_3,\qquad
 \omega_X=L_{-6},\qquad \omega_X^{-2}=L_{12}.
\]
This last expression uses the previously proved Picard identification;
the divisor, direct-image and ring proofs above do not depend on it.


\subsection{The retained period constant is an effective deformation parameter}
\hypertarget{tag-03CV}{}\label{wb:sec:retained-period-deformation}

Keep the functions, lattice marking, base coordinate and local fillings of
\CompendiumRef{I}{wb:sec:candidate-geometry}. In particular, fix the actual cusp-regular
solution $\beta^{\rm part}$, not a replacement with a cusp pole, and write
\[
 \beta_c=\beta^{\rm part}+c,\qquad
 \mathcal P=\{c\in\mathbb C:\operatorname{Im}c<-M\},
 \quad M=\max_{B\setminus\{p_0\}}D_{\rm part}.
 \tag{PD1}
\]
The maximum and the exact admissible domain were proved in
\CompendiumRef{I}{wb:prop:period-exact-admissible-constants}. The following calculation
uses the analytic construction, not its integral cohomology or sphere
recognition. Its purpose is to determine whether the freedom in $c$
actually changes the constructed complex manifold.

\begin{theorem}[Exact period deformation]
\hypertarget{tag-03DP}{}\label{wb:thm:retained-period-deformation}
\StatementRecord{wb:thm:retained-period-deformation}{proved}{03DP}{\SeriesRef{wb:thm:canonical-ring-recovery}; \cite[Section~67, p.~15]{McMullen1996Complex}}

For every open disc $Q$ with $\overline Q\subset\mathcal P$, the retained
construction gives a proper holomorphic submersion
\[
 \pi:\mathcal X_Q\longrightarrow Q,
 \qquad F:\mathcal X_Q\longrightarrow B\times Q,
 \qquad \pi=\operatorname{pr}_Q\circ F,
 \tag{PD2}
\]
whose fibre over $c$ is the analytically constructed threefold $X_c$ and
whose restriction there is the original $f_c:X_c\to B$.
The three distinguished values remain the original
$t(p_1)=0$, $t(p_2)=1$, $t(p_0)=\infty$.

Fix a regular base point and one of its original lifts $z$, and fix
$c_*\in Q$. With all functions below evaluated at that $z$, set
\[
 T=\operatorname{Im}\tau>0,\quad m=\operatorname{Im}\mu,\quad
 D=\operatorname{Im}\beta_{c_*}-\frac{6m^2}{T}<0,\quad a=\frac mT,
 \qquad P=\begin{pmatrix}0&0\\-a&1\end{pmatrix}.
 \tag{PD3}
\]
For $\delta=c-c_*$, the exact marked real-linear map from the universal
cover of the fibre at $c_*$ to that at $c$ is
\[
 w=A_\delta\zeta+B_\delta\overline\zeta,\qquad
 A_\delta=I+\frac{\delta}{2iD}P,\qquad
 B_\delta=-\frac{\delta}{2iD}P.
 \tag{PD4}
\]
It maps the original period lattice onto the new period lattice, and
therefore descends to a real-analytic marked diffeomorphism of the two
tori. In the explicit convention that the transported $(0,1)$ tangent
space is spanned by
$\partial_{\overline\zeta_k}+\sum_j\nu^j_k\partial_{\zeta_j}$,
its exact coefficient matrix and its first derivative are
\[
 \nu_\delta=\frac{\delta}{2iD+\delta}P,\qquad
 \dot\nu_z=\frac{1}{2iD}\,
 \partial_{\zeta_2}\otimes
 (d\overline\zeta_2-a\,d\overline\zeta_1).
 \tag{PD5}
\]
The class of $\dot\nu_z$ in
$H^{0,1}(f_{c_*}^{-1}(t(z)),T^{1,0})$ is nonzero.
With the tangent-exact-sequence convention
$\operatorname{KS}(\partial_c)=[\bar\partial h]$ for a smooth
$(1,0)$ lift $h$ of $\partial_c$, the actual fibre class is
$\operatorname{KS}_z(\partial_c)=-[\dot\nu_z]$.
Moreover the first-order deformation of the entire $X_{c_*}$ obtained
from (PD2) is nontrivial: its Kodaira--Spencer class in
$H^1(X_{c_*},T_{X_{c_*}})$ is nonzero. In particular
$h^1(X_{c_*},T_{X_{c_*}})\geq1$.
\end{theorem}



% TARGET PROOF MASKED: wb:thm:retained-period-deformation

























































































































































































































































































































The theorem establishes a genuinely varying analytic family with its
original marked fibration and its intrinsic degree-two canonical pencil.
It does not assert that every two distinct parameters give nonisomorphic
unmarked manifolds, or that this family exhausts their deformation space.
The global integral topology and the CDP argument are not inputs to this
deformation calculation.


The global identifications within this retained family are computed in
\cref{wb:pshift:thm:classification}: they are precisely $c'-c\in2\mathbb Z$,
including the finite-filling obstruction and every cusp overlap.


\subsection{Exact global identifications of the retained period parameter}
\hypertarget{tag-03CR}{}\label{wb:pshift:sec:classification}

Keep the original base $B$, coordinate $t$, values
$t(p_1)=0,t(p_2)=1,t(p_0)=\infty$, and the original functions
$\tau,\mu,\beta^{\mathrm{part}}$. Write
\[
 \beta_c=\beta^{\mathrm{part}}+c,\qquad
 \mathcal P=\{c\in\mathbb C:\operatorname{Im}c<-M\},\qquad
 \Pi_c(z)=\begin{pmatrix}6\mu(z)&\tau(z)&1&0\\
                      \beta_c(z)&\mu(z)&0&1\end{pmatrix}.
 \tag{PS1}\hypertarget{tag-03CN}{}\label{wb:pshift:eq:periods}
\]
All assertions below concern the analytic pieces and their actual gluing.
No integral cohomology, sphere recognition, or CDP assertion is an input.
In particular, the global parameter classification is not inferred from
an equality between two fibre lattices.

Put $e=(0,1)^{\mathsf t}\in\mathbb C^2$ and retain the ordered lattice
\[
 \Lambda=\mathbb Z\langle\widehat\gamma,\widehat u,
                        \widehat w,\widehat\delta\rangle,
 \qquad E=\widehat\delta\otimes\gamma.
 \tag{PS2}\hypertarget{tag-03CO}{}\label{wb:pshift:eq:rank-one-map}
\]
Thus $E$ is the matrix with its only nonzero entry, equal to one, in
row four and column one; $E^2=0$ and $\Pi_cE=e\gamma$.
The finite data, with every sign retained, are
\[
\begin{gathered}
 A_1=\begin{pmatrix}1&0&0&0\\6&0&1&0\\-6&-1&-1&0\\-2&1&0&1\end{pmatrix},
 \qquad
 A_2=\begin{pmatrix}1&0&0&0\\0&0&-1&0\\-6&1&0&0\\3&0&1&1\end{pmatrix},\\
 v_1=(1,2,-4,0)^{\mathsf t},\quad
 v_2=(-1,-3,3,0)^{\mathsf t},\quad
 (m_1,m_2)=(3,4),\quad
 (\epsilon_1,\epsilon_2)=(1,-1),\quad \epsilon_j=\gamma(v_j).
\end{gathered}
\tag{PS3}\hypertarget{tag-03C8}{}\label{wb:pshift:eq:finite-data}
\]
The covariance is $R_g\Pi_c=\Pi_c(gz)A_g$, with
\[
 R_{g_1}=\begin{pmatrix}-1/\tau&0\\(1-\mu)/\tau&1\end{pmatrix},\qquad
 R_{g_2}=\begin{pmatrix}1/\tau&0\\-\mu/\tau&1\end{pmatrix},\qquad
 R_{g_0}=I_2.
 \tag{PS4}\hypertarget{tag-03C2}{}\label{wb:pshift:eq:covariance}
\]
In particular $R_ge=e$ for every $g$.

We give the full period-shift proof of Theorem~\ref{fam:st:A.5.1}.


\subsubsection{The middle family and its complete rational centralizer}

For an integer $n$, put $S_n=I+nE$. Since $S_n^{-1}=S_{-n}$,
\[
 \Pi_{c+n}=\Pi_c S_n,\qquad
 \Pi_{c+n}\Lambda=\Pi_c\Lambda.
 \tag{PS6}\hypertarget{tag-03CF}{}\label{wb:pshift:eq:integer-shear}
\]
Both original matrices fix $\widehat\delta$ and preserve $\gamma$.
Consequently $A_jE=EA_j=E$, and the same commutation holds for every
monodromy word. The identity of $\mathbb C^2$ therefore identifies
the two full middle torus families with their actual $R_g$-action.
Its source-to-target lattice map is $S_{-n}$, since
$\Pi_{c+n}S_{-n}=\Pi_c$.

\begin{lemma}[All middle-family linear parts]
\hypertarget{tag-03CQ}{}\label{wb:pshift:lem:centralizer}
\StatementRecord{wb:pshift:lem:centralizer}{proved}{03CQ}{the definitions and constructions stated in the proof}

The simultaneous rational centralizer of $A_1,A_2$ is exactly
\[
 \{aI+bE:a,b\in\mathbb Q\}.
 \tag{PS7}\hypertarget{tag-03BZ}{}\label{wb:pshift:eq:centralizer}
\]
A biholomorphism between the middle families for $c,c'$ over the
identity of $B^\circ$ forces $c'-c\in\mathbb Z$. Its complex-linear
part on every fibre is either $I_2$ or $-I_2$.
\end{lemma}



% ALL PRINTED PROOF MATERIAL MASKED












































































\subsubsection{The exact finite obstruction, including normalization}

The two cyclic norm matrices are, by multiplication of (PS3),
\[
 N_1=I+A_1+A_1^2=
 \begin{pmatrix}3&0&0&0\\6&0&0&0\\-12&0&0&0\\0&2&1&3\end{pmatrix},
\]
\[
 N_2=I+A_2+A_2^2+A_2^3=
 \begin{pmatrix}4&0&0&0\\12&0&0&0\\-12&0&0&0\\0&2&2&4\end{pmatrix}.
 \tag{PS9}\hypertarget{tag-03CK}{}\label{wb:pshift:eq:norms}
\]
For example these formulas give
$N_1\widehat w=\widehat\delta$,
$N_2\widehat u=N_2\widehat w=2\widehat\delta$,
$N_1\widehat\delta=3\widehat\delta$, and
$N_2\widehat\delta=4\widehat\delta$.

We first justify why a local isomorphism must satisfy the exact affine
equations on the original finite cover. The actual ramified base change
$t_j=u^{m_j}$ has local equation
\[
 s_j^{m_j}-u^{m_j}
   =\prod_{a=0}^{m_j-1}(s_j-e^{2\pi ia/m_j}u)=0.
\]
Its normalization separates all these smooth branches, each with
coordinates $(z_1,z_2,u)$. The original map
$y\mapsto(q_j(y),s_j(y))$ identifies the original torus family $J'_j$
with this normalization, as proved by the explicit branches in
\CompendiumRef{I}{wb:ffcert:thm:normalization}. A biholomorphism over the identity of
$D_j$ base-changes over the identity of the $u$-disc and lifts uniquely
to the normalization: on local rings it carries integral elements in
the total fraction ring to integral elements, and its inverse does
the same. This lift commutes with the original deck action because
both composites lift the same map $(x,u)\mapsto(F(x),\zeta_j u)$.
In particular one may not replace the generator by an arbitrary
coprime power.

Let $c'=c+n$, where $n$ is an integer, and suppose the linear part of
the lifted map is $aI_2$, $a\in\{1,-1\}$. The lifted map on the torus
family is $[z,\zeta]\mapsto[z,a\zeta+B_j(z)]$. Its translating
section has a holomorphic lift $B_j:\Delta_j\to\mathbb C^2$:
$\Delta_j\times\mathbb C^2\to J'_j$ is a covering map and the disc is
simply connected. Local inverses of this covering show that the lift
is holomorphic. The deck-equivariance identity holds modulo an
actual lattice vector, so, after choosing this lift, it is
\[
 B_j(g_jz)-R_{g_j}(z)B_j(z)
  =\Pi_c(g_jz)\left(
       \frac{(1-a)v_j+n\epsilon_j\widehat\delta}{m_j}
          +\ell_j\right),\qquad \ell_j\in\Lambda.
 \tag{PS10}\hypertarget{tag-03C9}{}\label{wb:pshift:eq:finite-equation}
\]
The integer column $\ell_j$ is constant: it is locally constant by
discreteness of the marked lattice and the disc is connected.
Iterating $m_j$ times, with $A_jv_j=v_j$ and
$A_j\widehat\delta=\widehat\delta$, yields the necessary equality
\[
 N_j\ell_j=-(1-a)v_j-n\epsilon_j\widehat\delta.
 \tag{PS11}\hypertarget{tag-03CB}{}\label{wb:pshift:eq:finite-norm-obstruction}
\]
Here the left side of the iterated equation is zero because the
base point and the original vector coordinate return after $m_j$
steps; real invertibility of $\Pi_c$ makes the lattice equality exact.

If $a=-1$, applying $\gamma$ to (PS11) gives
\[
                   m_j\gamma(\ell_j)=-2\epsilon_j.
\]
This is impossible already at $j=1$, since $3$ does not divide $2$.
Thus a global isomorphism cannot have negative linear part.

For $a=1$, the first row of (PS11) gives $\gamma(\ell_j)=0$.
Equation (PS9) then proves that the order-three filling imposes no
restriction on $n$, whereas the order-four filling imposes precisely
$n\in2\mathbb Z$. More intrinsically, the finite obstruction is the
class of $-n\epsilon_j\widehat\delta$ in
\[
 \Lambda^{A_j}/N_j\Lambda.
 \tag{PS12}\hypertarget{tag-03C7}{}\label{wb:pshift:eq:finite-class}
\]
In the original fixed bases this quotient is
$\mathbb Z/3$ at $j=1$, with $\widehat\delta$ zero, and
$\mathbb Z/4\oplus\mathbb Z/2$ at $j=2$, with
$\widehat\delta$ generating the second summand. Indeed (PS9) gives
\[
 N_1\Lambda=3\mathbb Z v_1\oplus\mathbb Z\widehat\delta,
 \quad
 N_2\Lambda=4\mathbb Z\varepsilon'\oplus
                      2\mathbb Z\widehat\delta,
 \qquad \varepsilon'=(1,3,-3,0)^{\mathsf t}.
\]
This proves the exact parity obstruction rather than merely failure
of the untranslated identity.

\subsubsection{Local maps for every even shift}

Let $n\in2\mathbb Z$. Define the following original rational columns,
integral cocycles, and holomorphic translation vectors:
\[
\begin{aligned}
 r_1&=-\frac n3(\widehat u+\widehat w),&
 \ell_1&=-n\widehat w,&
 b_1(z)&=\Pi_c(z)r_1=-\frac n3(\tau(z)+1,\mu(z))^{\mathsf t},\\
 r_2&= \frac n4(\widehat u+\widehat w),&
 \ell_2&=\frac n2\widehat u,&
 b_2(z)&=\Pi_c(z)r_2=\frac n4(\tau(z)+1,\mu(z))^{\mathsf t}.
\end{aligned}
\tag{PS13}\hypertarget{tag-03CH}{}\label{wb:pshift:eq:local-corrections}
\]
Direct multiplication of the complete matrices in (PS3) gives
\[
 (I-A_j)r_j=\frac{n\epsilon_j}{m_j}\widehat\delta+\ell_j.
 \tag{PS14}\hypertarget{tag-03CI}{}\label{wb:pshift:eq:local-matrix-identity}
\]
For $j=1$ the four entries are $(0,0,-n,n/3)$, and for $j=2$
they are $(0,n/2,0,-n/4)$. Pairing with the full $\Pi_c(g_jz)$
shows that $b_j$ satisfies (PS10) for $a=1$.
Translation by $b_j$ is therefore already an isomorphism of the
two finite local fillings. The remaining issue is its exact overlap
with the fixed middle identification.

Keep the original linearizing coordinate, radius, and phase
$s_j(g_jz)=e^{-2\pi i/m_j}s_j(z)$. On a logarithm branch put
\[
 w_j(z)=\frac{\log s_j(z)}{2\pi i},\qquad
 a_j(z)=b_j(z)+n\epsilon_j w_j(z)e.
 \tag{PS15}\hypertarget{tag-03CA}{}\label{wb:pshift:eq:finite-middle-sections}
\]
Changing the logarithm adds $n\epsilon_j\widehat\delta$ under the
period map, so $a_j$ is branch independent as a torus section.
For the specified negative generator, $w_j(g_jz)=w_j(z)-1/m_j$.
Equations (PS4) and (PS14) consequently give the exact equality
\[
 a_j(g_jz)=R_{g_j}(z)a_j(z)+\Pi_c(g_jz)\ell_j.
 \tag{PS16}\hypertarget{tag-03CG}{}\label{wb:pshift:eq:local-affine-cocycle}
\]
This is why the logarithmic term in (PS15) cannot be omitted.
In fact the original finite overlap maps are translations by
$\sigma_{j,c}=w_j\Pi_c v_j$, and
\[
                  \sigma_{j,c+n}-\sigma_{j,c}
                     =n\epsilon_jw_je.
 \tag{PS17}\hypertarget{tag-03BX}{}\label{wb:pshift:eq:actual-overlap-difference}
\]
Thus $a_j$ is exactly the middle translation induced by the local
translation $b_j$, with the original gluing maps retained.

At the cusp retain $s=\tau-h$, $e^{2\pi is}=t_c$, $s(g_0z)=s(z)-1$,
and the complete period matrix
\[
 Z_c=sB_0+C_c(t_c),\quad
 B_0=\begin{pmatrix}0&1\\-1&0\end{pmatrix},\quad
 C_c=\begin{pmatrix}6\mu&h\\ b^{\mathrm{part}}-h+c&\mu\end{pmatrix}.
 \tag{PS18}\hypertarget{tag-03C5}{}\label{wb:pshift:eq:cusp-periods}
\]
For $\lambda=(\lambda_1,\lambda_2)^{\mathsf t}\in\mathbb Z^2$,
$C_{c+n}\lambda-C_c\lambda=(0,n\lambda_1)^{\mathsf t}$. Therefore
\[
 \exp(2\pi iC_{c+n}\lambda)t_c^{B_0\lambda}
       =\exp(2\pi iC_c\lambda)t_c^{B_0\lambda}.
 \tag{PS19}\hypertarget{tag-03C3}{}\label{wb:pshift:eq:cusp-actions-identical}
\]
The entire cusp actions, including their nonconstant holomorphic units,
are identical for an integer shift.

The affine cocycle convention in (PS16) has multiplication law
\[
                    \ell_{gh}=\ell_g+A_g\ell_h.
 \tag{PS20}\hypertarget{tag-03C0}{}\label{wb:pshift:eq:cocycle-law}
\]
Since $g_0=(g_1g_2)^{-1}$, it gives
\[
\begin{aligned}
 \ell_1+A_1\ell_2
   &=-\frac{3n}{2}\widehat w+\frac n2\widehat\delta,\\
 \ell_0=-A_{g_0}(\ell_1+A_1\ell_2)
   &=\frac{3n}{2}\widehat w-\frac n2\widehat\delta.
\end{aligned}
\tag{PS21}\hypertarget{tag-03C4}{}\label{wb:pshift:eq:cusp-cocycle}
\]
Here $A_{g_0}$ fixes both toric lattice vectors, as is also immediate
from $A_{g_0}=\left(\begin{smallmatrix}I&0\\B_0&I\end{smallmatrix}\right)$.
Put
\[
 q_n=(-3n/2,n/2)^{\mathsf t}\in\mathbb Z^2,
                  \qquad a_0(z)=s(z)q_n.
 \tag{PS22}\hypertarget{tag-03C6}{}\label{wb:pshift:eq:cusp-translation}
\]
Then $a_0(g_0z)-a_0(z)=-q_n=\Pi_c(g_0z)\ell_0$, exactly as
required by (PS16). Exponentiating gives the map
$(x,t_c)\mapsto(t_c^{q_n}x,t_c)$ on the dense cusp torus.
It extends holomorphically and invertibly to the original infinite fan:
the integral map
\[
 (y,y_3)\longmapsto(y+y_3q_n,y_3)
\]
translates each height-one triangle by the integer vector $q_n$ and
therefore permutes the actual cones, with inverse given by $-q_n$.
It commutes with the cusp quotient action in (PS19). This proves
extension across every cusp stratum, not only across its dense torus.

\subsubsection{The global additive correction and all overlap identities}

We give the global correction explicitly as an additive gluing
construction and prove its solvability. Let $\widetilde B^\circ$ be
the universal cover of the three-punctured base. Its fundamental group
is free on the original two negative peripheral generators. Pull back
the original $\Pi_c,R$ to this cover. Assign $\ell_1,\ell_2$ from
(PS13) to its free generators and extend by (PS20). Freeness makes
this an actual integral cocycle; its third peripheral value is
precisely (PS21). The maps
\[
 (z,y)\longmapsto
       (gz,R_g(z)y+\Pi_c(gz)\ell_g)
 \tag{PS23}\hypertarget{tag-03BY}{}\label{wb:pshift:eq:affine-bundle-action}
\]
form an affine action. To check the action law, covariance gives
$R_g(hz)\Pi_c(hz)=\Pi_c(ghz)A_g$; substituting this equality
in the composition gives exactly (PS20).
The quotient is an affine holomorphic bundle over $B^\circ$, since
the deck action on the base is free and properly discontinuous and
one may use evenly covered base opens as affine charts.

On a finite punctured disc, (PS16) says that $a_j$ is a section of
this affine bundle. On the cusp punctured disc, (PS21)--(PS22) say
the same for $a_0$. Define a sheaf $\mathscr H_n$ on $B$ as follows.
Away from the special points it is the sheaf of sections of (PS23).
Near $p_j$, a section is required to have the form $a_j+\eta_j$,
where $\eta_j:\Delta_j\to\mathbb C^2$ is holomorphic and obeys
$\eta_j(g_jz)=R_{g_j}(z)\eta_j(z)$. Near the cusp it is required
to have the form $a_0+\eta_0$, with $\eta_0(t_c)$ holomorphic on
the full cusp disc. These are compatible definitions on the punctured
overlaps by the equal affine cocycles just proved. They also prove
local nonemptiness, including at all three special points.

The differences of these local sections form a vector sheaf
$\mathscr V$. Its ordinary sections have covariance $\eta(gz)=R_g\eta(z)$;
at the two finite points and the cusp its extensions are precisely
the holomorphic extensions specified above. Projection to the first
coordinate gives an exact sequence on the original compact base
\[
 0\longrightarrow\mathcal O_B e\longrightarrow\mathscr V
        \longrightarrow\mathcal L\longrightarrow0,
                  \qquad \mathcal L\simeq\mathcal O_B(-p_0).
 \tag{PS24}\hypertarget{tag-03CP}{}\label{wb:pshift:eq:vector-sheaf}
\]
Here $\mathcal L$ is exactly the homogeneous period sheaf from
\CompendiumRef{I}{wb:prop:period-input-audit}, not a newly chosen extension. Indeed
the first diagonal entries in (PS4) are its factors $-1/\tau$ and
$1/\tau$, and at the cusp both extensions require ordinary
holomorphicity. The original function
\[
 \Theta=\frac{(E_4\circ\tau)^2(E_6\circ\tau)^{1/2}}
                 {\Delta_{\mathrm{mod}}\circ\tau}
\]
has those exact factors, orders two and one in $s_1,s_2$, and a
simple pole in $t_c$. Dividing by $\Theta$ identifies the ordinary
and finite stalks with $\mathcal O_B$, and the cusp stalk with
$\mathfrak m_{p_0}$. This recovers the displayed $\mathcal L$
with its original cusp condition. These transformation and order
calculations are the ones used to construct the retained $\mu$;
no topological conclusion about $X_c$ is involved.

For completeness, exactness on the finite stalks in (PS24) retains
the off-diagonal entry of $R_g$. Given an equivariant first component
$u(z)$, begin with the vector $(u(z),0)^{\mathsf t}$ and average it
under the finite group action on vector-valued holomorphic functions.
The resulting vector is equivariant, and its first component is
still $u$. Thus projection is surjective there. Its kernel is
exactly the invariant second-component functions, which are
holomorphic functions of $s_j^{m_j}=t_j$, multiplied by $e$.
The ordinary and cusp stalks are immediate. This also proves that
$\mathscr V$ is locally free of rank two: lift a local frame of
$\mathcal L$ and join it to $e$.

Both $H^1(B,\mathcal O_B)$ and
$H^1(B,\mathcal O_B(-p_0))$ vanish. One can compute these two
groups without a degree approximation: on the two standard affine
charts of $\mathbb P^1$, split a holomorphic Laurent series on
their overlap into its nonnegative and negative powers. For
$\mathcal O_B$ these are the two coboundary contributions, and
for $\mathcal O_B(-p_0)$ its transition factor shifts the second
contribution by one, still leaving no missing exponent. The positive
part converges on the full first affine chart and the transformed
negative part on the full second affine chart. This proves the two
vanishings, equivalently the usual additive Cousin calculation on
these charts. The exact sequence (PS24) now gives
\[
                             H^1(B,\mathscr V)=0.
 \tag{PS25}\hypertarget{tag-03CJ}{}\label{wb:pshift:eq:no-analytic-obstruction}
\]

Here is how to choose the correction entering the maps. Choose
local sections $h_\alpha$ of $\mathscr H_n$, taking the displayed
$a_j,a_0$ on sufficiently small special discs and affine charts
elsewhere. Their differences $d_{\alpha\beta}=h_\beta-h_\alpha$
are a Cech $1$-cocycle with values in $\mathscr V$. Project it to
$\mathcal L$ and solve that additive cocycle, using its just proved
$H^1$ vanishing. Lift the resulting local $\mathcal L$ sections to
$\mathscr V$ by the local lifting construction in (PS24). Subtracting
their coboundary from $d$ leaves a scalar $\mathcal O_Be$ cocycle.
Solve that scalar additive cocycle by its $H^1$ vanishing. This
produces sections $v_\alpha$ with
$d_{\alpha\beta}=v_\beta-v_\alpha$. The local sections
$h_\alpha-v_\alpha$ agree and define a global section
$a\in\mathscr H_n(B)$. This is a two-stage construction of the
required correction, not an unproved additional gluing condition.
It gives holomorphic vectors $\eta_j,\eta_0$ with
\[
 a=a_j+\eta_j\quad(j=1,2),\qquad a=a_0+\eta_0
                              \quad\hbox{at the cusp}.
 \tag{PS26}\hypertarget{tag-03CE}{}\label{wb:pshift:eq:global-correction}
\]
Any choices made in the additive solutions give maps with the same
asserted properties. Their differences as sections of this affine
bundle lie in $H^0(B,\mathscr V)=\mathbb C e$: this follows from
(PS24), $H^0(B,\mathcal L)=0$, and $H^0(B,\mathcal O_B)=\mathbb C$.

Projecting $a$ modulo $\Pi_c\Lambda$ gives a holomorphic section of
the middle torus family because every $\ell_g$ in (PS23) is integral.
The promised three kinds of local maps are therefore
\[
\begin{aligned}
 F_{\mathrm{mid}}[z,\zeta]&=[z,\zeta+a(z)],\\
 F_j[z,\zeta]&=[z,\zeta+b_j(z)+\eta_j(z)]\quad(j=1,2),\\
 F_0[x,t_c]&=[\exp(2\pi i\eta_0(t_c))t_c^{q_n}x,t_c].
\end{aligned}
\tag{PS27}\hypertarget{tag-03CD}{}\label{wb:pshift:eq:full-isomorphism}
\]
The first descends by (PS23) and equality of the integer-shift
lattices. The finite maps descend by (PS10), (PS14), and the exact
covariance of $\eta_j$. The cusp map descends by (PS19); its
holomorphic-unit factor extends by torus multiplication on every
original chart, and its integer factor extends by the fan map in
(PS22). Each map is a biholomorphism, with inverse given by the
negative translation or the inverse torus multiplier and fan map.

On a finite overlap, write the original gluing as
$h_{j,c}(\zeta)=\zeta+\sigma_{j,c}$. Then (PS17), (PS26), and
(PS27) give the equality of actual torus maps
\[
                     h_{j,c+n}F_j=F_{\mathrm{mid}}h_{j,c}.
 \tag{PS28}\hypertarget{tag-03CC}{}\label{wb:pshift:eq:finite-overlap-identity}
\]
Both sides add
$\sigma_{j,c}+a=\sigma_{j,c+n}+b_j+\eta_j$.
On the cusp overlap the original map is exponentiation of $\zeta$;
(PS22) and (PS26) give
$\exp(2\pi ia)=t_c^{q_n}\exp(2\pi i\eta_0)$, exactly the
third line of (PS27). Thus that overlap also agrees, including
all nonconstant period units. The maps and their inverses glue to
a global biholomorphism $X_c\to X_{c+n}$ over $B$.

The choices can be made coherently for every parameter and every even
shift. Both initial cocycle columns in (PS13) have $\gamma=0$.
Because every $A_g$ preserves $\gamma$, (PS20) shows that
$\gamma(\ell_g)=0$ for every word. Therefore the affine term
$\Pi_c(gz)\ell_g$ in (PS23) is independent of $c$. The displayed
$a_j,a_0$, and the vector sheaf $\mathscr V$ are also independent
of $c$. Perform the additive construction once for $n=2$, keeping
one resulting section $a^{(2)}$. For $n=2k$ choose
$a^{(2k)}=k a^{(2)}$ and multiply its local correcting vectors by
$k$. All their affine cocycles and local formulas are then exactly
those for $n=2k$. The resulting maps satisfy the full group law
\[
 F_{2h,c+2k}\circ F_{2k,c}=F_{2(h+k),c},\qquad
                 F_{0,c}=\mathrm{id}_{X_c}.
 \tag{PS28a}\hypertarget{tag-03C1}{}\label{wb:pshift:eq:coherent-shift-action}
\]
This follows by adding the actual translation vectors in the first
two lines of (PS27), and by multiplying their units and adding the
integer exponents in its third line. The formulas are jointly
holomorphic in $c$, since their corrections are independent of it.

\subsubsection{Intrinsic necessity and the resulting parameter space}

The analytically computed complete linear system
$|\omega_{X_c}^{-2}|$ is the original map $f_c$, by
\cref{wb:thm:canonical-ring-recovery}. A biholomorphism preserves its
canonical bundle, the complete space of sections, and its evaluation
map. Consequently it induces $\alpha\in\operatorname{PGL}_2(\mathbb C)$
with $f_{c'}F=\alpha f_c$. The three special fibres have distinct
intrinsic types: the cusp is reduced and nonnormal, while the other
two have multiplicities three and four. All remaining fibres are
smooth reduced tori. Thus $\alpha$ fixes $p_0,p_1,p_2$ individually.
In the original coordinate, a fractional-linear map fixing
$\infty,0,1$ has zero lower-left and upper-right entries and equal
diagonal entries; hence $\alpha=\mathrm{id}_B$.

Now \cref{wb:pshift:lem:centralizer} forces $c'-c=n\in\mathbb Z$
and linear part $\pm I_2$. Equations (PS10)--(PS11), obtained on
the original normalized finite covers, exclude $-I_2$ and then
exclude odd $n$ for $+I_2$. Conversely (PS27)--(PS28) construct
the global map for every even $n$. This proves
\cref{wb:pshift:thm:classification} without converting a failed
particular map into a nonisomorphism assertion.

As a complex parameter space for the isomorphism classes represented
by this retained family, the quotient $\mathcal P/(2\mathbb Z)$ has
the explicit coordinate
\[
 w=\exp(-\pi ic),\qquad 0<|w|<\exp(-\pi M).
 \tag{PS29}\hypertarget{tag-03CL}{}\label{wb:pshift:eq:parameter-disc}
\]
Indeed $|w|=\exp(\pi\operatorname{Im}c)$, and two complex numbers
have the same value of this exponential exactly when they differ
by an even integer. Every point in the displayed punctured disc
has a logarithm giving a $c\in\mathcal P$. The exponential has
nonzero derivative, so these local inverse logarithms give the
claimed quotient complex structure. The integer-period quotient
seen by the middle family instead has coordinate $w^2$; it is
covered two-to-one by (PS29). Its two lifts are distinguished by
the exact order-two summand in (PS12). This statement classifies
the parameters of the retained family; it does not assert that
the family exhausts all deformations or all complex threefolds
with the same singular-fibre types.



\subsection{Re-derivation of Theorem~\ref{fam:src:A.5.1}}\hypertarget{tag-000C}{}\label{app:rederive:fam-src-A.5.1}
% ALL PRINTED PROOF MATERIAL MASKED








\subsection{Re-derivation of Theorem~\ref{fam:src:A.5.2}}\hypertarget{tag-000D}{}\label{app:rederive:fam-src-A.5.2}
% ALL PRINTED PROOF MATERIAL MASKED

















































































































































































































































































































































\subsection{Re-derivation of Proposition~\ref{fam:src:B.4.1}}\hypertarget{tag-000E}{}\label{app:rederive:fam-src-B.4.1}
% ALL PRINTED PROOF MATERIAL MASKED





















\section{Tangent cohomology and deformations}

We retain the standard filling data $(\ell _0,\ell _1,\ell _2)=(0,1,-1)$
and an arbitrary admissible $c$. All sheaves in this section are analytic,
and all cohomology groups have coefficients in $\C$. The integral twists
below mean the actual line bundles
\[
 M_n=\OO_X\bigl(n(S_1-S_2)\bigr),\qquad n\in\Z.
\]
We set $\mathfrak g=g_\gamma$. Thus, in the positively normalized
exponential coordinate of \ref{inv:C3.2},
$M_n=\exp(n\mathfrak g)$. The computations use the root base to retain the
nonreduced fibres. Passing to its coarse space is a separate, exact
invariant operation; it is never identified with taking the dual of an
invariant lattice.

\begin{proposition}[The root base and its explicit cohomology]\hypertarget{tag-01SM}{}\label{inv:D1.1}
\StatementRecord{inv:D1.1}{proved}{01SM}{the definitions and constructions stated in the proof}

Let $\rho:\mathcal B\to B$ be the analytic root curve with orders $3$ and
$4$ at $p_1,p_2$. Write $R_1,R_2$ for its tautological root line bundles.
Its Picard group is infinite cyclic, with generator
$R=R_1\otimes R_2^{-1}$, and
\[
 R_1=R^4,\quad R_2=R^3,\quad \rho^*\OO_B(1)=R^{12},
 \quad K_{\mathcal B}=R^{-7}.
\]
We abbreviate $R^n$ to $\OO_{\mathcal B}(n)$. Choose sections
$x\in H^0(R^3)$ and $y\in H^0(R^4)$ vanishing at $p_2$ and $p_1$,
respectively. For every integer $n$, put
\[
 r_1(n)\equiv n\pmod3,\quad 0\le r_1(n)<3,
 \qquad r_2(n)\equiv3n\pmod4,\quad0\le r_2(n)<4,
\]
\[
 d(n)=\frac{n-4r_1(n)-3r_2(n)}{12}.
\]
Then
\[
 \rho_*\OO_{\mathcal B}(n)=\OO_B(d(n)),\qquad
 H^i(\mathcal B,\OO(n))=H^i(B,\OO(d(n))).                 \tag{D.1}
\]
Equivalently, $h^0(\OO(n))$ is the number of pairs
$(a,b)\in\Z_{\ge0}^2$ with $3a+4b=n$, and
$h^1(\OO(n))=h^0(\OO(-n-7))$.
\end{proposition}



The root-curve Picard group and graded cohomology are
\citet[Lemma~6.1]{DuEtAl2026v1}. We re-derive them and retain the explicit
coarse invariant-descent function $d(n)$ in the same calculation.

% ALL PRINTED PROOF MATERIAL MASKED






































\begin{proposition}[Vertical and logarithmic tangent bundles]\hypertarget{tag-01SN}{}\label{inv:D1.2}
\StatementRecord{inv:D1.2}{proved}{01SN}{\SeriesRef{inv:D1.1}}

There is a canonical representable proper flat holomorphic map
$\bar f:X\to\mathcal B$ with $f=\rho\bar f$. If
\[
 \mathcal T=T_X(-\log W),\qquad
 \mathcal V=\ker\bigl(\mathcal T\to
                   \bar f^*T_{\mathcal B}(-p_0)\bigr),
\]
then, for rank-two vector bundles $E,G$ on $\mathcal B$,
\[
 \mathcal V=\bar f^*E,\quad E=\OO\oplus\OO(-1),\qquad
 G=R^1\bar f_*\OO_X=\OO\oplus\OO(-1),                 \tag{D.2}
\]
\[
 R^0\bar f_*\OO_X=\OO,\qquad R^2\bar f_*\OO_X=\OO(-1),
\]
and the following is an exact sequence of locally free sheaves on $X$:
\[
 0\longrightarrow\bar f^*E\longrightarrow\mathcal T
 \longrightarrow\bar f^*\OO(-5)\longrightarrow0.       \tag{D.3}
\]
The trivial summands in $E,G$ are generated by the vertical field $\xi$
and the constant integral cocycle $\gamma$, respectively.
\end{proposition}



The root lift, vertical splitting, logarithmic quotient and
structure direct images are re-derived after reading
\citet[Lemmas~6.2 and 6.7; Propositions~6.6 and 6.10]{DuEtAl2026v1}.
The complete root and ordinary-sheaf arguments are retained below.

% ALL PRINTED PROOF MATERIAL MASKED













































































































% ALL PRINTED PROOF MATERIAL MASKED



























































































































% ALL PRINTED PROOF MATERIAL MASKED






































% ALL PRINTED PROOF MATERIAL MASKED





























\section{The integral parameter comparison}

\subsection{Credited re-derivations: local charts and coherent direct images}

For the standard marked data, the vertical splitting and structure
direct images are re-derived from
\cite[Proposition~6.6, pp.~77--78; Lemma~6.7, p.~78;
Proposition~6.10, pp.~79--80]{DuEtAl2026v1}.
The tangent dimensions and Kuranishi germ are re-derived from its
Theorem~6.12, pp.~81--82, and Proposition~6.16, p.~85.
The full automorphism and period comparison, including the exact
parameter transport, is given in \cref{night:P5.7}.
The proofs below retain their stated ranges of affine twists.
Write $B=\mathbb P^1$, let $S_1,S_2$ be the reduced fibres of
multiplicities $3,4$, let $W$ be the cusp fibre, and let $D$ be its reduced
double locus.  Thus $D$ consists of three rational curves joining the two
triple points.  Set
\[
 L=f^*\mathcal O_B(-1)\otimes\mathcal O_X(2S_1+S_2),\quad
 F=\mathcal O_X(S_1+S_2),\quad
 E=\ker(df),\quad J=\mathcal I_D\otimes F.
\]
Let $\widehat B$ denote the analytic root orbifold with orders $3,4$ at
$p_1,p_2$, let $\pi:\widehat B\to B$ be its coarse map, and let
$h:X\to\widehat B$ be defined by the canonical sections of $S_1,S_2$.
In particular $f=\pi h$.  This notation means descent on the ordinary
root discs by the specified finite cyclic actions.  Taking invariants is
exact over $\mathbb C$, by averaging.  Denote by $\widehat L,\widehat F$
the root line bundles whose pullbacks are $L,F$.

\begin{theorem}\hypertarget{tag-02YM}{}\label{night:P5.1}\label{stmt:P5.1}
\StatementRecord{night:P5.1}{proved}{02YM}{the definitions and constructions stated in the proof}

The sheaf $E$ is locally free of rank two and is the pullback of the
translation Lie bundle $V$ on $\widehat B$.  There is an exact sequence
\begin{equation}\hypertarget{tag-01OK}{}\label{eq:P5tangent}
 0\longrightarrow E\longrightarrow T_X\longrightarrow J\longrightarrow0.
\end{equation}
\end{theorem}
% ALL PRINTED PROOF MATERIAL MASKED






























\begin{theorem}\hypertarget{tag-02YN}{}\label{night:P5.2}\label{stmt:P5.2}
\StatementRecord{night:P5.2}{proved}{02YN}{the definitions and constructions stated in the proof}

On the root orbifold,
\[
 V\simeq\mathcal O\oplus\widehat L,
 \qquad R^qh_*\mathcal O_X\simeq
 \bigwedge^q(\mathcal O\oplus\widehat L)\quad(0\leq q\leq2).
\]
On $B$, the direct images used below are
\[
\begin{array}{c|ccc}
 & f_* & R^1f_* &R^2f_*\\ \hline
 \mathcal O_X&\mathcal O&\mathcal O\oplus\mathcal O(-1)&\mathcal O(-1)\\
 L&\mathcal O(-1)&\mathcal O(-1)^{\oplus2}&\mathcal O(-1)\\
 F&\mathcal O&\mathcal O^{\oplus2}&\mathcal O\\
 E&\mathcal O\oplus\mathcal O(-1)&\mathcal O\oplus\mathcal O(-1)^{\oplus3}
       &\mathcal O(-1)^{\oplus2}\\
 J&\mathcal O(-1)&\mathcal O(-1)^{\oplus2}&\mathcal O .
\end{array}
\]
\end{theorem}
% ALL PRINTED PROOF MATERIAL MASKED








































































































































\begin{theorem}\hypertarget{tag-02YO}{}\label{night:P5.3}\label{stmt:P5.3}
\StatementRecord{night:P5.3}{proved}{02YO}{the definitions and constructions stated in the proof}

For every admissible $c$,
\[
 \bigl(h^0(T_{X_c}),h^1(T_{X_c}),h^2(T_{X_c}),h^3(T_{X_c})\bigr)
 =(1,1,1,0).
\]
\end{theorem}
% ALL PRINTED PROOF MATERIAL MASKED




















\begin{theorem}\hypertarget{tag-02YP}{}\label{night:P5.4}\label{stmt:P5.4}
\StatementRecord{night:P5.4}{proved}{02YP}{the definitions and constructions stated in the proof}

The germ of the $c$-parameter family at each admissible member is a
miniversal deformation germ.  The Kuranishi germ is a smooth disc.
\end{theorem}
% ALL PRINTED PROOF MATERIAL MASKED




































\subsection{All automorphisms and the integral parameter period}

\begin{theorem}\hypertarget{tag-02YQ}{}\label{night:P5.5}\label{stmt:P5.5}
\StatementRecord{night:P5.5}{proved}{02YQ}{the definitions and constructions stated in the proof}

For every admissible twist of the $(3,4)$ construction,
\[
 \operatorname{Aut}(X_c)=\mathbb C/\mathbb Z\simeq\mathbb C^*.
\]
The parameter $a\bmod\mathbb Z$ acts by translation by $a\widehat\delta$.
\end{theorem}
% ALL PRINTED PROOF MATERIAL MASKED
















































































































\begin{theorem}\hypertarget{tag-02YR}{}\label{night:P5.6}\label{stmt:P5.6}
\StatementRecord{night:P5.6}{proved}{02YR}{the definitions and constructions stated in the proof}

For fixed admissible twist data, $X_c\simeq X_{c'}$ precisely when
$c-c'\in2\mathbb Z$.  This fixes the integral period of $c$; the
construction still permits an additive choice of $\beta^{\rm part}$.
\end{theorem}
% ALL PRINTED PROOF MATERIAL MASKED









































































\clearpage
\section{Global coefficient representatives}
\hypertarget{tag-03L3}{}\label{inv:section:coefficient-atlas}

\subsection{A literal coefficient atlas for the two generators}

The Green representative of \cref{inv:C3.1} and the harmonic
normalization of a lift of \cref{inv:A3.2} give operator descriptions.
The following construction specifies a holomorphic \v Cech
cocycle and every smooth coefficient in its Dolbeault representative.
It uses the original chart maps through all three special fibres.
\paragraph{Relation to the source constructions.}
The period matrix, two affine finite fillings, logarithmic collar maps
and untranslated cusp fan are the original construction
\cite[Sections~3--6]{AC2026}.  The same local pieces in the two
versions of \cite{DuEtAl2026v1,DuEtAl2026v2} are identified by the
parameter translation $a=c+\delta$ of \cref{night:P5.7}; in particular
$\Omega_a=\Pi_c$, with the ordered lattice, vectors and fan retained.
The scalar cocycle below is precisely $f=-\eta_{\rm DGKX}$ on every
displayed generating arrow in
\cite[Section~6.4, scalar cocycle and Lemma~6.14 with proof, p.~84; conventions and generating arrows, pp.~82--83]{DuEtAl2026v1}.
The minus sign is the stated line-coordinate convention.  The
exponential quotient comparison with Engel's elliptic-surface
presentation is the map already given in
\cref{fam:src:B.1.1}.

Conversion of a holomorphic cocycle to a Dolbeault form by smooth
weights is the usual \v Cech--Dolbeault construction
\cite[Theorem~3.1 and proof, p.~10]{AbateBracciSuwaTovena2010Atiyah}.
On two charts its formula applied to $(0,0,c_{01}/12)$ is
$-\bar\partial w_0\wedge c_{01}/12
=\bar\partial(w_1c_{01})/12$, exactly the convention used here;
replace $c_{01}/12$ by $f_{01}$ for $G$.
The proof supplies the actual local theta and Cartier lines, their
annular holomorphic connections, the resulting one-form cocycle,
and a specified locally finite coordinate weight system with every
derivative evaluated.  Thus the general cohomological conversion and
the scalar input are credited together with the original family data.

\begin{proposition}[A specified global coefficient atlas]
\hypertarget{tag-01SL}{}\label{inv:C3.8}
\StatementRecord{inv:C3.8}{proved}{01SL}{\SeriesRef{inv:A2.1}, \SeriesRef{inv:A3.2}, \SeriesRef{inv:C3.1}, \SeriesRef{inv:C3.2}, \SeriesRef{inv:C3.7}}
Retain the standard marked family of \cref{wb:thm:retained-period-deformation},
its original vectors $v_1=(1,2,-4,0)^t$, $v_2=(-1,-3,3,0)^t$,
the untranslated original cusp fan, and any
$c\in\mathcal P=\{\operatorname{Im}c<3\log(6)/\pi\}$.
There are specified smooth forms
\[
 G\in\mathcal A^{0,1}(X),\qquad Q\in\mathcal A^{1,1}(X),
 \qquad \bar\partial G=\bar\partial Q=0,
\]
whose classes satisfy
\[
 [G]=g_\gamma,\qquad
 i_b^*[Q]_{\bar\partial}=\eta,\qquad
 \exp(12[G])=f^*\mathcal O_{\mathbb P^1}(1).
\]
Here $i_b$ is any regular fibre inclusion and
$\eta=u\wedge w+6\gamma\wedge\delta$ in the original integral marking.
The formulas \eqref{inv:eq:atlas-final-forms} and
\eqref{inv:eq:atlas-final-coefficients}, with the elementary recursion
\eqref{inv:eq:atlas-weight-derivative}, give their coefficients on a
specified countable refinement of the original atlas.  Every sum is
locally finite, also at the multiple fibres, cusp double curves and
both cusp triple points.  The form $Q$ is a complex representative;
no real $\bar\partial$-closed representative can have this fibre class.
\end{proposition}

% ALL PRINTED PROOF MATERIAL MASKED






















































































































































































































































































































































































































































































































\paragraph{Exact algebra checks.}
The programs and their same-stem \path{.output.json} files are in
\path{code/current_audits/audit_paper_i/open_atlas/}.
The program \path{exact_regular_atlas.py} records 25 exact checks of
the marking, lattice factors, curvature coefficients and regular
transition identities; \path{cusp_and_finite_exact.py} records 48
exact checks of the Cartier exponents, extending units, finite phases
and connection identities.  The separate first-jet implementation
\path{finite_collar_jet_verifier.py} records 106 checks of both affine
actions through their first twelve iterates, including the full base
component of the collar connection and the logarithmic gauge terms.
The universal identities in the proof establish descent on the stated
domains; the first-jet evaluations corroborate those identities at the
recorded exact inputs.

\Uses{the original marked family and full chart maps in
\cref{fam:section:A,fam:src:A.1.2,wb:thm:retained-period-deformation};
\CompendiumRef{I}{wb:lem:toric-chart-control};
\CompendiumRef{I}{wb:thm:candidate-cusp-geometry};
the exact determinant divisor in \cref{inv:A2.1};
the fibre class in \cref{inv:A3.2};
the Picard and primitive normalization in \cref{inv:C3.1,inv:C3.2}}

\subsection{Comparison with the Green and harmonic representatives}
\hypertarget{tag-03LI}{}\label{inv:C3.8:operator-comparison}

The real restriction map
\[
r_b:H_A^{1,1}(X;\R)\longrightarrow H_A^{1,1}(F_b;\R)
\]
has image $\R\eta$ by \ref{inv:A3.2} and \ref{inv:C3.7}.  Its source has real dimension
two and its target has real dimension four.  The map is induced by pullback of $\partial\bar\partial$-closed forms followed by the quotient by $\partial$- and $\bar\partial$-exact forms on $F_b$. Thus it specifies the class of each restricted representative through that exact quotient map.

The formula \eqref{inv:eq:C-G-Green} supplies a global $G$ after a metric
choice.  For the standard primitive marking, $G_\gamma=G_0/12$
represents $g_\gamma$ by \ref{inv:C3.2}.  There is a precise operator
construction of $Q_H$ as well.  Fix a Hermitian metric on $X$ and a
regular fibre $F$, identify $\eta$ with its real invariant form on
$F$, and choose an orthonormal basis $e_1,e_2$ of the Dolbeault-harmonic
$(1,1)$ forms for $\Delta_{\bar\partial}$.  Define
\[
\ell(e)=-\frac1{12}\int_F i_F^*e\wedge\eta,\qquad
Q_H=\frac{\overline{\ell(e_1)}e_1+
            \overline{\ell(e_2)}e_2}
          {|\ell(e_1)|^2+|\ell(e_2)|^2}.
\]
The denominator is positive by \ref{inv:A3.2} and
$\int_F\eta^2=-12$.  The expression is independent of the
orthonormal basis, is $\bar\partial$-closed, and has fibre
Dolbeault class $\eta$.  Indeed the image of fibre restriction on
$H^{1,1}_{\bar\partial}(X)$ is $\C\eta$, with kernel
$\C d_1g$; the displayed functional measures exactly the coefficient
of $\eta$, and the formula makes that coefficient one.

% ALL PRINTED PROOF MATERIAL MASKED

































The exact comparison on cohomology is
\[
 [G]=[G_0]/12=g_\gamma,\qquad
 [Q_H]-[Q]\in\C d_1g.
\]
The second inclusion follows from the restriction kernel in
\cref{inv:A3.2}.  Thus the two formulas choose lifts of the same fibre
class; their difference is the specified one-dimensional Leray
ambiguity.  The harmonic normalization remains basis independent for
the fixed metric and regular fibre.

\Uses{\cref{inv:A3.2,inv:C3.1,inv:C3.2,inv:C3.7};
the Green formula \eqref{inv:eq:C-G-Green}}


\subsection{Changes of the specified coefficient gauge}

\paragraph{The representative and its fibre restriction.}
\hypertarget{tag-03CT}{}\label{wb:rem:bcpr-literal-gauge}
The formulas \eqref{inv:eq:atlas-final-forms} and
\eqref{inv:eq:atlas-final-coefficients} choose actual global smooth
forms through all three special fibres.  Their restriction statement
uses the pullback map followed by the Dolbeault quotient on a regular
fibre.  In particular, its translation-invariant value is the class
represented in \CompendiumRef{I}{wb:eq:bcpr-q-fibre-expanded}.
The specified coefficient form on that fibre can be related to the
invariant representative by a Dolbeault coboundary.  The cocycle and
weight formulas below describe the changes produced by these choices.

\Uses{\cref{inv:A3.2,inv:C3.8};
\CompendiumRef{I}{wb:eq:bcpr-q-fibre-expanded}}

\paragraph{The exact changes between coefficient gauges.}
\hypertarget{tag-03CS}{}\label{wb:rem:bcpr-global-gauge-boundary}
Use the same holomorphic cocycles $c_{ab},f_{ab}$ on a common
refinement, and two locally finite weight systems $w_n,w'_n$ with
sum one.  In a chart $a$ set
\[
 K_a=\sum_n\sum_u(w'_n-w_n)c_{a,n},\qquad
 U_a=\sum_n\sum_u(w'_n-w_n)f_{a,n}.
\]
The cocycle identities and $\sum(w'_n-w_n)=0$ give
$K_a-K_b=0$ and $U_a-U_b=0$.  They therefore define a global
smooth $(1,0)$ form $K$ and a global smooth function $U$, with
\[
 Q'-Q=\tfrac1{12}\bar\partial K,\qquad
 G'-G=\bar\partial U.
\]
For a holomorphic change of cocycles
$c'_{ab}=c_{ab}+r_a-r_b$ and
$f'_{ab}=f_{ab}+u_a-u_b$, where $r_a$ is a holomorphic one-form
and $u_a$ a holomorphic function, the same weights give
\[
 B'_a-B_a=r_a-\sum_n\sum_u w_n r_n,\qquad
 H'_a-H_a=u_a-\sum_n\sum_u w_n u_n.
\]
The two sums on the right are globally defined smooth forms of the
indicated degrees.  Consequently
\[
 Q'-Q=-\tfrac1{12}\bar\partial\!\sum_n\sum_u w_n r_n,
 \qquad
 G'-G=-\bar\partial\!\sum_n\sum_u w_n u_n.
\]
These formulas identify the actual Dolbeault coboundaries under the
specified gauge changes and preserve all class and product
calculations.  A change of coherent-Leray lift of the fibre class is
measured separately by $\C d_1g$, as in the operator comparison above.

\Uses{the cocycles and coefficients in \cref{inv:C3.8};
the restriction kernel in \cref{inv:A3.2};
the class and product calculations in \cref{inv:C3.3}}

\clearpage
\part{Results not yet established}
\section{Periods, compact subgroups, and the transcendence question}

We use a Weierstrass equation over $\overline\Q$ and its invariant
differential. Its period lattice is $\Lambda=\Z\omega_1+\Z\omega_2$,
with $\Im(\omega_2/\omega_1)>0$. The quasi-periods are defined by
$\zeta(z+\omega_j)=\zeta(z)+\eta_j$. We extend
$\eta:\Lambda\to\C$ real-linearly to $\C$. The discriminant in
this section is the discriminant of this equation, equivalently of
the lattice in the specified differential. The quantity $N$ is
independent of changing this algebraic equation and differential.

\begin{proposition}[sigma expression and the third-kind period]\hypertarget{tag-0017}{}\label{ari:C.7.1}
\StatementRecord{ari:C.7.1}{proved}{0017}{\SeriesRef{ari:eq:cp-periods}; \citet[Definitions 2.3--2.4 and
Section 6]{bertolin2022}}

Let $P\in E(\overline\Q)$, with $P\ne O$ and $2P\ne O$, and
let $u$ be an elliptic logarithm of $P$. Then
\begin{equation}\hypertarget{tag-003K}{}\label{ari:eq:cp-sigma}
 2\log N(P)=\frac16\log|\Delta|
       +\Re\{\Lambda_1+u(\zeta(u)-\eta(u))\},
 \qquad \Lambda_1=2\log\sigma(u)-u\zeta(u).
\end{equation}
The choices of logarithms change the expression inside the real
part by purely imaginary integral multiples of $2\pi$.
There is an extension $G_P$ of $E$ by $\mathbb G_m$, defined
over $\overline\Q$ and corresponding, under the principal
polarization, to $P$, for which $\Lambda_1$, modulo a logarithm
of an algebraic number, is the negative of a relative period
of $[\Z\to G_P]$. Changing an algebraic differential of the
third kind adds an algebraic multiple of $u$ and algebraic
endpoint terms.
\end{proposition}


% ALL PRINTED PROOF MATERIAL MASKED




























































\begin{proposition}[the compact-subgroup criterion and the real formula]\hypertarget{tag-001B}{}\label{ari:C.7.2}
\StatementRecord{ari:C.7.2}{proved}{001B}{\SeriesRef{ari:C.7.1}, \SeriesRef{ari:eq:cp-periods}, \SeriesRef{ari:eq:cp-v}, \SeriesRef{ari:eq:cp-chi}; \citet[slides 5--8]{bertrand2017}}

For $P$ as in Statement \ref{ari:C.7.1}, $N(P)$ is algebraic if and only
if the fibre of $G_P$ above $P$ contains an algebraic point of
the maximal compact subgroup $G_P(\C)^c$.
For the family in the stated conventions, if $t>1$ is real, choose
$\tau=iT$, $T>1$, and the zero $\mu=1/2-iy$, $0<y<T/2$.
With $x=2y/T$, one has
\[
 N(P)^4=3^{3/4}(t-1)^{-1/4}
       \left|\frac{\vartheta_1(x,i/T)}{\eta(i/T)}\right|.
\]
\end{proposition}


% ALL PRINTED PROOF MATERIAL MASKED













































\subsection{Motivic groups from the point and extension data}

\citet[Construction~3.4, p.~12, Lemma~3.5, p.~13, and
Proposition~4.12, pp.~19--20]{bertolin2022} compute the motivic
group from its abelian point and extension data and the central
Weil-pairing bracket. For the $1$-motives used here, let $B$ be the
smallest abelian subvariety generated by these data, and let $Z(1)$
be the central torus obtained from that bracket and the remaining
algebraic lift. Their dimension formula is
\begin{equation}\hypertarget{tag-003D}{}\label{ari:eq:cp-MTdim}
 \dim G_{\mathrm{mot}}(M)
 =\dim G_{\mathrm{mot}}(A)+2\dim B+\dim Z(1).
\end{equation}
The bracket is described in their Section~6.1, pp.~26--27, and
their corrected Table~2, pp.~28--29, gives unipotent dimension
three in the dependent nondeficient case. The following calculation
determines these groups through the displayed polarization maps;
the arithmetic question for the metric norm is recorded in the
proved implication with its explicitly open hypothesis in \cref{ari:C.7.4}.

\begin{proposition}[the motivic groups needed on the diagonal]\hypertarget{tag-001C}{}\label{ari:C.7.3}
\StatementRecord{ari:C.7.3}{proved}{001C}{\SeriesRef{ari:eq:cp-MTdim}}

Let $P$ be nontorsion and let $R$ be any algebraic lift above
$P$ in $G_P$. Put $M_R=[\Z\to G_P]$, $1\mapsto R$.
Its connected motivic group is, after a choice of Levi subgroup,
the semidirect product of $G_{\mathrm{mot}}(E)$ and the
three-dimensional Heisenberg group attached to $H_1(E)$ and
its principal polarization. Its dimension is $5$ for CM and
$7$ without CM.

For a non-real embedding put
$\mathcal M=M_R\oplus\overline{M_R}$.
If $E$ and $\bar E$ are not isogenous, its dimension is $13$.
If they are isogenous, choose an isogeny
$\phi:\bar E\to E$ of degree $d$ and write
$Q=\phi(\bar P)$. Let $r$ be the dimension over
$F=\End(E)\otimes\Q$ of the span of $P,Q$, modulo torsion.
Then
\[
\begin{array}{c|c|c|c}
\text{case}&\dim B&\dim Z&\dim G_{\mathrm{mot}}(\mathcal M)\\ \hline
r=2,\ E\text{ CM}&2&2&8\\
r=2,\ E\text{ not CM}&2&2&10\\
r=1,\ E\text{ CM}&1&1+e&5+e\\
r=1,\ E\text{ not CM}&1&1+e&7+e.
\end{array}
\]
Here $e$ is either zero or one: it is zero precisely when the
algebraic lift in the one-dimensional quotient of the two
central tori by the bracket torus is torsion. This criterion is
independent of the choices made modulo isogeny. In the real
case we use the single motive $M_R$; no extra quotient torus
is necessary.
\end{proposition}


% ALL PRINTED PROOF MATERIAL MASKED



























































% Compendium-only ari:C.7.5: apparatus: verification/diagnostic record retained with source computation




\begin{proposition}[The period-conjecture implication for the local norm]\hypertarget{tag-001D}{}\label{ari:C.7.4}
\StatementRecord{ari:C.7.4}{proved(H)}{001D}{\SeriesRef{ari:C.7.2}, \SeriesRef{ari:C.7.3}, \SeriesRef{ari:eq:cp-v}, \SeriesRef{ari:eq:cp-positive-centre}; \cite[Construction~3.4, p.~12, Lemma~3.5, p.~13, Conjecture~4.7, p.~18, Proposition~4.12 and Remark~4.13, pp.~19--20]{bertolin2022}}
Let $E/\overline\Q$ be an elliptic curve and let
$P\in E(\overline\Q)$ be nontorsion. Use the local norm $N(P)$,
the algebraic extension $G_P$, and its algebraic fibre coordinate
from \cref{ari:C.7.1}, \cref{ari:C.7.2}.
The hypothesis $H$ is the Relative Semi-abelian Conjecture
\cite[Conjecture~4.7, p.~18]{bertolin2022}, restricted to
$G=G_P\times G_{\bar P}$ and the algebraic lifts used below.
Explicitly, for those lifts and their associated $B_Q,Z(1)$, it gives
\[
 \operatorname{trdeg}_{\overline\Q(\Omega_G)}
 \overline\Q\bigl(\Omega_G,R,\bar R,
                  \widetilde{\log}_G(R,\bar R)\bigr)
 \geq 2\dim B_Q+\dim Z(1),
\]
where $\Omega_G$ is the full de Rham--Betti period matrix,
including the second-kind periods. Under $H$, $N(P)$ is
transcendental. This applies to every algebraic parameter
$t\ne0,1$ with nontorsion $P_t$, including the conductor-$361$
CM point identified below. Hypothesis $H$ remains open for
these points. Grothendieck's period conjecture for the algebraic
$1$-motives in this proof implies $H$
\cite[Proposition~4.12, pp.~19--20]{bertolin2022}.
\end{proposition}

% ALL PRINTED PROOF MATERIAL MASKED

























































































\subsection{CM theta normalization and further unconditional deductions}

The holomorphic theta normalization proposed for the CM interpolation approach differs from the
norm occurring in the stated problem. We make that distinction explicit,
including the finite-character ambiguity in the multiplication law.

\begin{proposition}[the exact CM multiplication law]\hypertarget{tag-001F}{}\label{ari:C.7.6}
\StatementRecord{ari:C.7.6}{proved}{001F}{the definitions and constructions stated in the proof}

Let $E/\overline{\Q}$ have complex multiplication by an order
$\mathcal O$, and use an algebraic invariant differential. Write
$A=\operatorname{area}(\Lambda)$ and
\[
 \eta(z)=s_2z+\frac{\pi}{A}\bar z,\qquad
 \theta(z)=\sigma(z)e^{-s_2z^2/2}.
\]
Then $s_2\in\overline{\Q}$. For every nonzero
$\alpha\in\mathcal O$, with $d=\alpha\bar\alpha$, there is
$F_\alpha\in\overline{\Q}(E)$ such that
\begin{equation}\hypertarget{tag-003Q}{}\label{ari:eq:ct-CMmult}
 \theta(\alpha z)^2=F_\alpha(\exp_E z)\theta(z)^{2d}.
\end{equation}
Its divisor and leading Laurent coefficient are
\[
 \operatorname{div}(F_\alpha)=2[\alpha]^*(O)-2d(O),\qquad
 F_\alpha(\exp_Ez)=\alpha^2z^{2-2d}(1+O(z^2)).
\]
Thus the unsquared quotient is an algebraic function on $E$, but
it need not be a rational function on $E$. For a nontorsion
$P=\exp_Eu\in E(\overline{\Q})$, algebraicity of $\theta(u)$
would imply algebraicity of $\theta(\alpha u)$ for all
$\alpha\in\mathcal O$. The additive group $\mathcal O u$ has
rank two, not more than four.
\end{proposition}


% ALL PRINTED PROOF MATERIAL MASKED


















































\begin{cited}[with derivation: two unconditional normalized values]\hypertarget{tag-001G}{}\label{ari:C.7.7}
\StatementRecord{ari:C.7.7}{proved}{001G}{\SeriesRef{ari:C.7.6}, \SeriesRef{ari:C.7.1}, \SeriesRef{ari:eq:ct-Gaussian}, \SeriesRef{ari:C.7.2}, \SeriesRef{ari:C.7.9}; \citet[Theorem~1 and Corollary~1]{waldschmidt1979sigma}; \citet[Theorem~6.1]{huber2022}; \citet[final slide, comment (c)]{bertrand2017}; \citet[p.~15-03]{waldschmidt1988theta}; \citet[Corollary~2]{reyssat1982weierstrass}; \citet[p.~15-05]{waldschmidt1988theta}; \citet[Theorem 3.1.2]{bost2009}; \citet[Theorem 1.4 and Corollary 1.6]{nguyen2023}; \citet[Theorems 1.3 and 2.3]{xie2024algebraicity}}

Let $E/\overline{\Q}$ have an algebraic invariant differential,
and let $u$ be an elliptic logarithm of an algebraic nontorsion
point $P$. Then
\[
 \sigma(u)\exp\left(-\frac12u\zeta(u)\right)
\]
is transcendental. If $E$ has CM and $s_2$ is normalized as in
\ref{ari:C.7.6}, then $\zeta(u)-s_2u$ is also transcendental. The first
assertion concerns the displayed holomorphic normalization;
it does not assert transcendence of either $N(P)$ or
$\theta(u)=\sigma(u)e^{-s_2u^2/2}$.
\end{cited}
The expanded source derivation is retained in the proof compendium, Paper~I, label \nolinkurl{ari:C.7.7}.




The stated normalized values follow from \citet[Theorem~1 and Corollary~1]{waldschmidt1979sigma} and the analytic subgroup theorem in \citet[Theorem~6.1]{huber2022}. A complete re-derivation is given in \CompendiumRef{I}{app:rederive:ari-C.7.7}.


\paragraph{The exact dependence of the theta expression on its lifts.}
Put
$g(u,v)=\log(\sigma(u+v)/(\sigma(u)\sigma(v)))$ on a branch where
all three sigma values are nonzero. Replacing $v$ by $v+\omega$,
$\omega\in\Lambda$, changes $g(u,v)-s_2uv$, modulo $2\pi i\Z$, by
\[
 (\eta(\omega)-s_2\omega)u=\frac{\pi\bar\omega u}{A}.
\]
The map $v\mapsto\exp_Ev$ sends these two lifts to the same
elliptic point. Thus $N\exp_Ev=f(\exp_Eu)$ for a
Rosati-antisymmetric $f$ retains its value under this lift change,
while the displayed theta expression acquires the stated term.
A use of the pointwise criterion in
\citet[final slide, comment~(c)]{bertrand2017} must therefore retain
its compatible lift convention. The Gaussian-norm question here
remains the precise question stated above.

A subsequent source check resolves the normalization of the
one-variable holomorphic assertion itself:
\citet[p.~15-03]{waldschmidt1988theta} explicitly states that
$\sigma(u)e^{-s_2u^2/2}$ is transcendental when the invariants
are algebraic, $E$ has CM, and $\exp_Eu$ is an algebraic
nontorsion point.  That source defines $s_2$ by the regularized
lattice sum and attributes the result to Bertrand's 1982
paper \emph{Valeurs de fonctions th\^eta et hauteurs p-adiques}.
We have checked this complete normalized statement in
Waldschmidt's published exposition; the full original 1982
proof has not been retrieved.  This source clarification
does not validate the unrestricted-lift inference withdrawn
above, and the theorem concerns the holomorphic value rather
than the metric norm.  \ref{ari:C.7.7} retains the two assertions whose
proofs have been fully checked here.

\paragraph{The archimedean term still requiring control.}
The exact comparison is
\begin{equation}\hypertarget{tag-003R}{}\label{ari:eq:ct-Gaussian}
 N(P)=|\Delta|^{1/12}|\theta(u)|
                         \exp\left(-\frac{\pi|u|^2}{2A}\right).
\end{equation}
It follows from \ref{ari:C.7.1} and the real-linear decomposition of
$\eta$. The normalization proved transcendental in \ref{ari:C.7.7} instead
gives the exact identity
\[
 N(P)=|\Delta|^{1/12}
 \left|\sigma(u)e^{-u\zeta(u)/2}\right|
 \exp\left(\frac12\Re\{u(\zeta(u)-\eta(u))\}\right).
\]
Neither formula excludes algebraicity of $N(P)$: transcendence
of one complex factor does not control its absolute value or
its product with the remaining exponential.
The generic finite-order interpolation approach
supplies only the rank-two set $\mathcal O u$; it does not
meet the proposed rank-greater-than-four criterion. The
derivative approach would need the algebraicity of
$\theta'(u)/\theta(u)=\zeta(u)-s_2u$, which \ref{ari:C.7.7} excludes.
Replacing the norm by the holomorphic theta value would
therefore change the problem precisely by the factor in
\eqref{ari:eq:ct-Gaussian}.

The related independence theorem of
\citet[Corollary~2]{reyssat1982weierstrass} says that two of
$e^{i\pi u/\omega_1}$, $s(u)$ and $s'(u)$ are algebraically
independent, where $s(z)=\sigma(z)e^{-\eta_1z^2/(2\omega_1)}$.
For real $u$, hypothetical algebraicity of $N(P)$ makes
$s(u)$ algebraic, so the conclusion instead forces independence
of $e^{i\pi u/\omega_1}$ and
$\zeta(u)-\eta_1u/\omega_1$; it yields no contradiction.
The periodic-functional-equation theorem reported in
\citet[p.~15-05]{waldschmidt1988theta} assumes algebraic
$q=e^{i\pi\tau}$, whereas this $q$ is transcendental for
imaginary quadratic $\tau$ by Gelfond--Schneider.  These
specific checked applications therefore do not decide the
metric norm; their source scopes and the larger CM
independence tuples are recorded separately.

We also tested an additive-extension route.
\citet[Theorem 3.1.2]{bost2009} proves that an algebraic point
of the universal vector extension of an abelian variety which
lies in its maximal compact subgroup at one complex embedding
is torsion. The theorem's hypotheses require an algebraic
point of the vector extension, equivalently an algebraic line
bundle with an algebraic connection having unitary monodromy.
\ref{ari:C.7.2} supplies only an algebraic point in one fibre of a
multiplicative extension. It does not supply that algebraic
connection. Differentiating a metric value at one algebraic
point cannot establish algebraicity of the connection there
or globally. This is the exact missing implication in that
attempt.

Another route uses the Latt\`es map $f_E$ defined by
$x(2Q)=f_E(x(Q))$. Its canonical height is
$\widehat h_{f_E}(x(P))=2\widehat h_O(P)$. For rational
parameters C.6 therefore makes algebraicity of $N(P)$
equivalent to algebraicity of
$\exp(\widehat h_{f_E}(x(P)))$.
\citet[Theorem 1.4 and Corollary 1.6]{nguyen2023} settle the
corresponding multiplicative-height problem for nonintegrable
polynomials using B\"ottcher coordinates in the basin of
infinity. They do not apply to $f_E$. Indeed an elliptic
duplication Latt\`es map has no exceptional point: the inverse
images of any point contain the images of all successive
division points of an elliptic point, an infinite set. A
polynomial has the exceptional point $\infty$, so no conjugate
of $f_E$ is polynomial. The Julia set is the whole sphere
and the required basin of infinity is absent. This checks
the failed hypothesis, rather than treating a theorem for
polynomials as a theorem for all rational maps.

We also checked the arithmetic algebraicity criterion of
\citet[Theorems 1.3 and 2.3]{xie2024algebraicity}.
Besides sufficiently rapid height-normalized approximation by
points in a fixed number field, it requires the analytic germ
to be formally defined over that field. This hypothesis fails
for the compact leaf in \ref{ari:C.7.9}: its tangent has the proved
transcendental slope
$c=\zeta(u)-\eta_1u/\omega_1$.
A smooth formal germ defined over $\overline{\Q}$ has an
algebraic tangent; algebraic translations, coordinate changes
and isogenies preserve this obstruction. Assuming $N(P)$
algebraic supplies an algebraic point on the leaf and its
integer multiples, but not an algebraic formal germ. Attempting
to force convergence by successive division also loses the
fixed-number-field hypothesis: if $[2]^nQ_n=Q$ with all
$Q_n\in E(K)$ for one number field, then
$\widehat h_O(Q_n)=4^{-n}\widehat h_O(Q)>0$.
The $Q_n$ are distinct with bounded ample height, contrary to
Northcott finiteness. The full hypothesis audit and the opened
preprint are retained in the records.

\subsection{The precise nonarchimedean comparison}
\begin{cited}[with derivation: a nonzero ordinary height]\hypertarget{tag-001H}{}\label{ari:C.7.8}
\StatementRecord{ari:C.7.8}{proved}{001H}{\SeriesRef{ari:C.7.6}, \SeriesRef{ari:C.7.3}, \SeriesRef{ari:C.7.4}, \SeriesRef{ari:C.7.2}; \citet[Appendix A]{agboola2006}; \citet{agboola2006}; \citet[Theorem 2.1]{bertrand1995}; \citet[Theorem 3.1]{bertrand1995}; \citet{bost2009}}

Let $E/\Q$ have complex multiplication by the full ring of integers
of an imaginary quadratic field $K$ of class number one. Assume
that its functional-equation sign is $-1$, and let $p>3$ be a
prime of good ordinary reduction, $p\mathcal O_K=\mathfrak p
\mathfrak p^*$. For a nontorsion $P\in E(\Q)$ the
$\mathfrak p$-adic height and the cyclotomic $p$-adic height
of $P$ over $K$, in the normalization of
\citet[Appendix A]{agboola2006}, are nonzero.
In particular this applies to
\[
 E:y^2+y=x^3-38x+90,\qquad P=(0,9),\qquad p=5.
\]
\end{cited}
The expanded source derivation is retained in the proof compendium, Paper~I, label \nolinkurl{ari:C.7.8}.




The nonvanishing theorem is reproduced in Rubin's Appendix~A, Theorem~A.1 and Lemma~A.2 of \citet{agboola2006}; the exact hypotheses and application are retained. A complete re-derivation is given in \CompendiumRef{I}{app:rederive:ari-C.7.8}.


The comparison explains the missing step in an attempted
archimedean adaptation. The ordinary $p$-adic sigma construction
uses a holomorphic quadratic correction. In the CM normalization
the corresponding complex holomorphic function is $\theta$
of \ref{ari:C.7.6}, whereas the real metric also contains
$\exp(-\pi|u|^2/(2A))$. Theorem A.1 supplies no identity
between the ordinary $p$-adic height and that real factor.
Transferring the conclusion would require precisely such a
new argument; replacing the real factor by the holomorphic
CM correction does not preserve the norm being asked about.

There is also a definite limitation on the maximal-compact
approach. \citet[Theorem 2.1]{bertrand1995} constructs
nontrivial multiplicative extensions of polarized abelian
varieties with nontorsion algebraic points in their maximal
real compact subgroup. Its hypotheses include an endomorphism
not fixed by the Rosati involution and a nontorsion rational
point over a real number field. The construction uses a base
point related to the extension class by a Rosati-antisymmetric
endomorphism. It cannot be replaced by a theorem saying that
every algebraic compact point of every multiplicative
extension is torsion. For our diagonal class, an equality
$NP=f(P)$ with $N>0$ and $f+f^\dagger=0$ would force
$P$ torsion: $N-f$ has no zero eigenvalue, since
$(N-f)^\dagger(N-f)=N^2-f^2$ is positive. In the elliptic
case this is simply a nonzero isogeny. The diagonal positivity
used in \ref{ari:C.7.3}--\ref{ari:C.7.4} is therefore essential, but these
counterexamples neither prove nor refute the desired
unconditional diagonal statement. The additive result in
\citet[Theorem 3.1]{bertrand1995}, and its higher-dimensional
extension in \citet{bost2009}, requires a point in a vector
extension, a hypothesis not furnished by \ref{ari:C.7.2}.


\subsection{The fixed semiabelian extension and the diagonal subgroup attempt}

\begin{proposition}[the endomorphism and slope obstructions]\hypertarget{tag-001I}{}\label{ari:C.7.9}
\StatementRecord{ari:C.7.9}{proved}{001I}{\SeriesRef{ari:C.7.1}, \SeriesRef{ari:eq:cd-compatibility}, \SeriesRef{ari:C.6.R1}, \SeriesRef{ari:C.7.2}; \citep[Theorem 6.1]{huber2022}}

Let $E/\overline\Q$ be an elliptic curve, let
$q\in E^\vee(\overline\Q)$ be nontorsion, and let
\[
 1\longrightarrow\mathbb G_m\longrightarrow G_q
       \longrightarrow E\longrightarrow0
\]
be the extension of class $q$. Then
$\operatorname{End}_{\overline\Q}(G_q)=\Z$, including when
$E$ has complex multiplication. In particular a noninteger
CM endomorphism of $E$ does not lift to an endomorphism of this
fixed extension.

For
\[
 E:y^2+y=x^3-38x+90,\qquad P=(0,9),\qquad Q=2P=(4,-2),
\]
one has
\[
 N(Q)=19^{1/4}N(P)^4.
\]
Use the differential $dx/(2y+1)$, choose the positive real
period $\omega_1$, and choose the real logarithm
$u\in(0,\omega_1/2)$ of $Q$. Put
\begin{equation}\hypertarget{tag-0033}{}\label{ari:eq:cd-slope}
 c=\zeta(u)-\frac{\eta_1}{\omega_1}u,
 \qquad \eta_1=\eta(\omega_1).
\end{equation}
Then $c$ is a real transcendental number. The actual point
\[
 R_N=(Q,3\,19^{-1/2}N(Q)^2)\in G_Q(\C)
\]
is nontorsion, belongs to the maximal compact subgroup, and has
the real logarithm $(u,cu)$. Its algebraicity is equivalent to
the algebraicity of $N(P)$. In the algebraic logarithm coordinates
of \ref{ari:C.7.1}, the line $\ell_c=\C(1,c)$ has transcendental
slope. The proof establishes this slope statement from a period
of $G_Q$ without deciding the algebraicity of $R_N$.
\end{proposition}


% ALL PRINTED PROOF MATERIAL MASKED

















































































































\paragraph{What functional transcendence gives in this construction.}
The functional theorem can be applied fully, but its conclusion
does not contradict this hypothetical algebraic compact point.
Indeed consider
\[
 f:\C\longrightarrow\operatorname{Lie}(G_Q),\qquad
 f(z)=(z,cz).
\]
No proper connected algebraic subgroup of $G_Q$ projects
onto $E$: as in the proof above, such a subgroup would split
the nontorsion extension after an isogeny. Every proper
connected subgroup is therefore contained in the torus,
whereas $\exp_{G_Q}f$ has nonconstant elliptic projection.
Thus $f$ is $G_Q$-nondegenerate in Definition 1.2 of
\citet{noguchi2024}. Their Theorem 1.3 gives
\[
 \dim_{\C}\overline{\{(\exp_{G_Q}f(z),f(z)):z\in\C\}}^{\rm Zar}
       \ge \dim G_Q+1=3.
\]
The graph is contained in $G_Q\times\ell_c$, of dimension
three, so equality holds. This is a statement over $\C$
about a family of meromorphic functions. It remains valid
whether or not the single value $\exp_{G_Q}f(u)$ is
algebraic. Under the norm-algebraicity hypothesis the integer
values $\exp_{G_Q}f(nu)=[n]R$ are algebraic, but they do
not give an algebraic relation holding on the entire graph.
There is no contradiction with its full dimension.
One would need an additional arithmetic specialization
argument, or a way to obtain a further functional relation
from that single metric value. Neither is supplied by the
checked functional theorem. The endomorphism calculation
above also prevents obtaining extra CM multiplication
directions inside this fixed $G_Q$ by lifting the
noninteger endomorphisms of $E$.


\subsection{A CM-stable enlargement and its Gaussian eigencoordinate}

The fixed-extension calculation in \ref{ari:C.7.9} does not prevent
enlarging the torus. We test that alternative explicitly.
Retain its curve, its point $Q=2P$, and its real logarithm
$u$. Write $\beta=\sqrt{-19}=i\sqrt{19}$, $d=19$,
$D=\zeta(u)-s_2u$, and $k=\pi/A$. All endomorphisms and
extensions in what follows are defined over $\overline\Q$.

\begin{proposition}[the enlarged CM action does not isolate the Gaussian]\hypertarget{tag-0018}{}\label{ari:C.7.10}
\StatementRecord{ari:C.7.10}{proved}{0018}{\SeriesRef{ari:C.7.9}, \SeriesRef{ari:C.7.6}, \SeriesRef{ari:eq:ce-zeta}}

Let
\[
 \mathcal G=G_Q\mathbin{\times_E}G_{\beta Q},\qquad
 1\longrightarrow\mathbb G_m^2\longrightarrow\mathcal G
       \longrightarrow E\longrightarrow0.
\]
There is an endomorphism $B$ of $\mathcal G$ inducing
$\beta$ on $E$ and
$(t_1,t_2)\mapsto(t_2^{-1},t_1^d)$ on the torus.
It satisfies $B^2=[-d]$. Moreover
\[
 \operatorname{Hom}(\mathcal G,\mathbb G_m)=0,
 \qquad
 \operatorname{End}(\mathcal G)\otimes\Q=\Q(\beta).
\]
Thus $\mathcal G$ has no algebraic $\mathbb G_m$ quotient,
including after isogeny.

There is an algebraic compact point in the second extension
above $Q$ whose logarithm is
\[
 (u,v_2),\qquad v_2=-\beta uD+au,
 \qquad
 a=\frac1{2\beta}
       \frac{d}{du}\log F_\beta(Q)\in\overline\Q.
\]
Here the derivative means differentiation of the rational
function $F_\beta(\exp_E z)$ with respect to the algebraic
invariant differential, evaluated at $z=u$.
Pair this point with the actual compact point $R_N$ of
\ref{ari:C.7.9}. The resulting compact point of $\mathcal G(\C)$
has logarithm
\[
 U=(u,v_1,v_2),\qquad v_1=uD-ku^2.
\]
Its algebraicity is equivalent to the algebraicity of $N(P)$.
In these algebraic logarithm coordinates the differential of
$B$ and its $-\beta$ eigenprojection are
\[
 dB=\begin{pmatrix}
 \beta&0&0\\ a&0&-1\\ \beta a&d&0
 \end{pmatrix},\qquad
 \Pi_-:=\frac12\left(I-\frac{dB}{\beta}\right),
\]
\[
 \Pi_-U=-\frac{\pi u^2}{2A}(0,1,\beta).
\]
The first torus coordinate of $\exp(\Pi_-U)$ is exactly
the Gaussian $\exp(-\pi u^2/(2A))$. The projector is not
the differential of an algebraic quasi-endomorphism. No
sequence of rational polynomials in $B$, even allowing
inverses of isogenies, can converge to it as complex linear
operators. At least one of the two torus coordinates of
$\exp(\Pi_-U)$ is transcendental; the proof does not decide the
algebraicity of its first coordinate alone.
\end{proposition}


% ALL PRINTED PROOF MATERIAL MASKED








































































































































































































































































\subsection{Complexifying the Gaussian: values, jets, and the rank criterion}

The antiholomorphic Gaussian can be made holomorphic by introducing a
second variable. This gives more algebraic values under the hypothesis
in C.7(c), but the derivative and growth hypotheses of the resulting
transcendence arguments must still be checked. We do this on the
conductor-$361$ model, where all constants and conjugations are fixed.

\begin{proposition}[the complexified route and its exact obstructions]\hypertarget{tag-0019}{}\label{ari:C.7.11}
\StatementRecord{ari:C.7.11}{proved}{0019}{\SeriesRef{ari:C.7.9}, \SeriesRef{ari:C.7.1}, \SeriesRef{ari:C.7.6}, \SeriesRef{ari:eq:cc-fixedfield}, \SeriesRef{ari:C.7.7}; \citep[Theorem 6.1]{huber2022}}

Use the curve, point and real logarithm of \ref{ari:C.7.9}:
\[
 E:y^2+y=x^3-38x+90,\qquad Q=(4,-2)=\exp_Eu,
 \qquad 0<u<\omega_1/2.
\]
Put $\beta=i\sqrt{19}$, $k=\pi/A$, $s_2=-2$, and
\[
 \theta(z)=\sigma(z)e^{-s_2z^2/2},\qquad
 D(z)=\zeta(z)-s_2z,\qquad
 \mathcal F(z,w)=\sqrt{19}\,\theta(z)\theta(w)e^{-kzw},
 \qquad H=\mathcal F^2.
\]
Then $\mathcal F(z,\bar z)=N(z)^2$, and $\mathcal F$ is an
entire function of two complex variables of order at most two.
It is periodic under $(z,w)\mapsto(z+\omega,w+\bar\omega)$,
$\omega\in\Lambda$. The subgroup
\[
 \Gamma=\{(\alpha u+\omega,\bar\alpha u+\bar\omega):
                   \alpha\in\Z[\beta],\ \omega\in\Lambda\}
 \subset\C^2
\]
has rank four over $\Z$ and lies in the real plane
$\{(z,\bar z):z\in\C\}$. There is one number field $K$ such
that, with $H_0=H(u,u)=N(Q)^4$, every value of $H$ on
$\Gamma$ belongs to $K[H_0]$. More precisely, the nonzero values
are algebraic multiples of the powers $H_0^{\alpha\bar\alpha}$,
as specified in \eqref{ari:eq:cc-fixedfield}.

For every nonzero $\alpha\in\Z[\beta]$ there is an algebraic
number $a_\alpha$ such that
\[
 D(\alpha u)=\bar\alpha D(u)+a_\alpha.
\]
At $p_\alpha=(\alpha u,\bar\alpha u)$ the algebraic constant
direction $V_\alpha=\alpha\partial_z-\bar\alpha\partial_w$
satisfies
\begin{align*}
 V_\alpha\log\mathcal F(p_\alpha)
    &=\alpha a_\alpha-\bar\alpha\bar a_\alpha
                         \in\overline{\Q},\\
 V_\alpha^2\log\mathcal F(p_\alpha)
    &=-\alpha^2\{\wp(\alpha u)+s_2\}
      -\bar\alpha^2\{\wp(\bar\alpha u)+s_2\}
      +2\alpha\bar\alpha k.
\end{align*}
The number $k$ is transcendental. For every nonzero $\alpha$,
the exact normalized second derivative satisfies
\[
 \overline\Q\!\left(\frac{V_\alpha^2H(p_\alpha)}{H(p_\alpha)}\right)
      =\overline\Q(k).
\] In particular,
\[
 (\partial_z-\partial_w)H(u,u)=0,\qquad
 \frac{(\partial_z-\partial_w)^2H(u,u)}{H(u,u)}=4k-8.
\]
Thus the normalized first derivative is algebraic and the
normalized second derivative generates $\overline\Q(k)$. Moreover the rank-four set does not satisfy the
explicit derivative-free finite-order criterion audited below.
These are obstructions to the stated applications of those
criteria; they do not decide the algebraicity of $N(Q)$.
\end{proposition}



% ALL PRINTED PROOF MATERIAL MASKED









































































































































































\subsection{Division values and a variable-degree interpolation test}


\begin{proposition}[division ratios, scalar models and translated functions]\hypertarget{tag-001A}{}\label{ari:C.7.12}
\StatementRecord{ari:C.7.12}{proved}{001A}{\SeriesRef{ari:C.4.2}, \SeriesRef{ari:C.7.1}, \SeriesRef{ari:C.7.6}, \SeriesRef{ari:C.7.9}, \SeriesRef{ari:C.7.11}}

Let $E:y^2+y=x^3-38x+90$, $Q=(4,-2)$ and $N_0=N(Q)$.
Use the positive real logarithm $u$ and period $\omega_1$ of
\ref{ari:C.7.9}. Put
\[
 \delta=19^{1/4},\qquad b=\eta_1/\omega_1,\qquad
 g(z)=\delta\sigma(z)e^{-bz^2/2}.
\]
Thus $|g(z)|=N(\exp_Ez)$ on the real axis. Let $\mathcal D$
be the CM division hull of $Q$. For $R\in\mathcal D$, the
equations $[n]R=[\alpha]Q$, $n>0$, determine a unique
$\kappa_R=\alpha/n\in\operatorname{End}(E)\otimes\Q$.
Set $w_R=\kappa_R\bar\kappa_R$. For $R\ne O$,
\[
 A_R:=\frac{N(R)}{N_0^{w_R}}\in\overline\Q_{>0}.
\]
In particular, for $[n]R=[m]Q$ with $m\ne0$,
\begin{equation}\hypertarget{tag-003S}{}\label{ari:eq:division-ratio-model}
 A_R^{n^2}=\delta^{n^2-m^2}
                  \frac{|\psi_m(Q)|}{|\psi_n(R)|}.
\end{equation}
For each positive algebraic number $\lambda$, define
\[
 N_\lambda(R)=A_R\lambda^{w_R}
       =N(R)\left(\frac{\lambda}{N_0}\right)^{w_R}.
\]
All these values are algebraic. For $[n]R=[m]Q$, $|m|\le n$,
$R\ne O$, there are constants $D,H$ depending on $\lambda,E,Q$
such that
\[
 h(N_\lambda(R))\le H,\qquad
 [\Q(R,N_\lambda(R)):\Q]\le
 \begin{cases}Dn^4,&R\in E(\R),\\2Dn^6,&R\in E(\C).
 \end{cases}
\]
The transformation is exact for every positive real parameter:
$N_\lambda(R)/N_\mu(R)=(\lambda/\mu)^{w_R}$ and
$N_{N_0}(R)=N(R)$.

For any finite collection of real division-hull logarithms $a$,
the functions $g(z)$ and $g(z+a)/g(a)$, with $\exp_Ea\ne O$,
have function-field transcendence degree at most three over
$\overline\Q$ after adjoining the elliptic coordinates.
The three functions $g(z),g(z+v)/g(v),g(z-v)/g(v)$, $v=u/q$,
attain three. The normalized holomorphic extensions
$F_a(z)/N(a)^2$ of $N(\exp_E(z+a))^2$ for fixed complex
division-hull shifts belong to the same type of field.
The unnormalized translates have the same upper bound, attained by
the displayed triple, over the algebraic closure $L$ of $\overline\Q(N_0)$ in $\C$.

For the scalar models on the real division grids
$z=(mu+k\omega_1)/n$, $|m|\le Cn$, $k$ modulo $n$, with
fixed $C$, every fixed finite family of these values has a common
degree envelope $O(n^5)$ and bounded heights. For the actual values the same construction
gives a relative degree envelope over a fixed finite extension
of $\Q(N_0)$. The algebraicity of $N_0$ remains the open question
of \ref{ari:C.7.4}. The explicit interpolation and tangent
calculations below record what these bounds establish.
\end{proposition}

% ALL PRINTED PROOF MATERIAL MASKED






































































































































































































































































































































































































\section{Open arithmetic questions for the normalized norm}\hypertarget{tag-003T}{}\label{ari:open:norm}



Without hypothesis $H$, the arithmetic nature of the metric norm remains open at the
nontorsion points of \cref{ari:C.7.4}. An explicit small instance is
\[
 E:y^2=x^3-140x+784,\qquad P=(0,28),\qquad j=-3375,
 \qquad N(P)=2e^{-\widehat h(P)}.
\]
Here $P$ is a primitive generator modulo the rational
$2$-torsion and
\[
 \widehat h(P)=
 0.46736567718147128000752020126284678363562034370331\ldots.
\]
The exact unresolved assertion is that $2e^{-\widehat h(P)}$
is transcendental. The corresponding example of conductor
$361$, $y^2+y=x^3-38x+90$, $P=(0,9)$, has
$N=e^{-\widehat h(P)}$ and no finite local contribution.
The relative-period proof isolates its obstruction: an
unconditional theorem excluding
\eqref{ari:eq:cp-relation} for the positive diagonal central
direction. A theorem about linear $1$-period relations
does not supply that exclusion. The attempted approaches and the exact CM theta comparison are printed in the preceding period and CM-normalization sections;
their source checks and computational records are retained in the audit files.
These arguments do not resolve the stated exclusion.

\clearpage



\clearpage
\begin{thebibliography}{99}
\providecommand{\natexlab}[1]{#1}
\providecommand{\doi}[1]{\url{https://doi.org/#1}}
\bibitem[{Claude} and {L. Alp{\"o}ge}(2026)]{AC2026}
{Claude} and {L. Alp{\"o}ge}.

\newblock The (3,4,infinity) modular family of 2-tori, completed at its three
  special points, is a complex structure on {$S^6$}.

\newblock 108-page manuscript, 2026.

\newblock URL \url{https://alpo.ge/s6.pdf}.
\newblock \emph{Verified}.

\bibitem[Abramovich et~al.(2002)Abramovich, Karu, Matsuki, and
  W{\l}odarczyk]{AKMW02AB}
Dan Abramovich, Kalle Karu, Kenji Matsuki, and Jaros{\l}aw W{\l}odarczyk.

\newblock Torification and factorization of birational maps.

\newblock \emph{Journal of the American Mathematical Society}, 15\penalty0
  (3):\penalty0 531--572, 2002.

\newblock \doi{10.1090/S0894-0347-02-00396-X}.

\newblock URL \url{https://personal.math.ubc.ca/~karu/papers/semist.pdf}.
\newblock \emph{Verified}.

\bibitem[Abasheva and D{\'e}ev(2024)]{AbashevaDeev2024AB}
Anna Abasheva and Rodion D{\'e}ev.

\newblock Complex surfaces with many algebraic structures.

\newblock \emph{International Mathematics Research Notices}, 2024\penalty0
  (9):\penalty0 7379--7400, 2024.

\newblock \doi{10.1093/imrn/rnad190}.

\newblock URL \url{https://arxiv.org/pdf/2303.10764}.
\newblock \emph{Verified}.

\bibitem[Abate et~al.(2010)Abate, Bracci, Suwa, and Tovena]{AbateBracciSuwaTovena2010Atiyah}
Marco Abate, Filippo Bracci, Tatsuo Suwa, and Francesca Tovena.

\newblock \emph{Localization of Atiyah classes}.

\newblock Author manuscript dated 7 May 2010.

\newblock URL \url{https://pagine.dm.unipi.it/abate/articoli/artric/files/LocalAtiClasses.pdf}.

\newblock Verified at the cited scope: Theorem 3.1 and proof, p.~10;
  Remark 1.6, p.~3; formulas (1.4), (1.7), pp.~3--4; Section~4.1, p.~12.
\newblock \emph{Verified}.

\bibitem[Angella(2018)]{Angella2018}
Daniele Angella.

\newblock Hodge numbers of a hypothetical complex structure on {$S^6$}.

\newblock \emph{Differential Geometry and its Applications}, 57:\penalty0
  105--120, 2018.

\newblock \doi{10.1016/j.difgeo.2017.10.012}.

\newblock URL \url{https://arxiv.org/abs/1705.10518}.
\newblock \emph{Verified}.

\bibitem[Atiyah and Singer(1968)]{AtiyahSinger1968IIIE}
Michael~F. Atiyah and Isadore~M. Singer.

\newblock The index of elliptic operators: Iii.

\newblock \emph{Annals of Mathematics}, 87\penalty0 (3):\penalty0 546--604,
  1968.

\newblock \doi{10.2307/1970717}.

\newblock URL
  \url{https://math.uchicago.edu/~shmuel/tom-readings/AS%20III.pdf}.

\newblock Additional verified scope on 19 November 2026: Index Theorem (2.12), p.~556; associated-symbol Proposition (2.17), p.~559, with its structure-group setup on pp.~557--558; the Dolbeault-symbol derivation and Riemann--Roch Theorem (4.3), pp.~563--565. The actual inherited Spin-c symbol and its line twist are transported through their explicit graded characters.
\newblock \emph{Verified}.

\bibitem[Bierstone and Milman(1995)]{BM95AB}
Edward Bierstone and Pierre~D. Milman.

\newblock Canonical desingularization in characteristic zero by blowing up the
  maximum strata of a local invariant, 1995.

\newblock URL \url{https://arxiv.org/pdf/alg-geom/9508005}.
\newblock \emph{Verified}.

\bibitem[Bierstone and Milman(1997)]{BM97}
Edward Bierstone and Pierre~D. Milman.

\newblock Canonical desingularization in characteristic zero by blowing up the
  maximum strata of a local invariant.

\newblock \emph{Inventiones Mathematicae}, 128:\penalty0 207--302, 1997.

\newblock \doi{10.1007/s002220050141}.

\newblock URL \url{https://arxiv.org/pdf/alg-geom/9508005}.
\newblock \emph{Verified}.

\bibitem[B{\u a}nic{\u a} and St{\u a}n{\u a}{\c s}il{\u a}(1976)]{BS76}
Constantin B{\u a}nic{\u a} and Octavian St{\u a}n{\u a}{\c s}il{\u a}.

\newblock \emph{Algebraic Methods in the Global Theory of Complex Spaces}.

\newblock Editura Academiei and John Wiley \& Sons, Bucharest and London, 1976.

\newblock URL \url{https://djvu.online/file/2hQiOnz8WuoQk}.
\newblock \emph{Verified}.

\bibitem[Biliotti and Ghigi(2021)]{BiliottiGhigi2021Root}
Leonardo Biliotti and Alessandro Ghigi.

\newblock Meromorphic limits of automorphisms.

\newblock \emph{Transformation Groups}, 26:\penalty0 1147--1168, 2021.

\newblock \doi{10.1007/s00031-020-09551-x}.

\newblock URL \url{https://arxiv.org/abs/1901.10724}.
\newblock \emph{Verified}.

\bibitem[Bourdon et~al.(2017)Bourdon, Clark, and
  Stankewicz]{BourdonClarkStankewicz2017}
Abbey Bourdon, Pete~L. Clark, and James Stankewicz.

\newblock Torsion points on {CM} elliptic curves over real number fields.

\newblock \emph{Transactions of the American Mathematical Society},
  369\penalty0 (12):\penalty0 8457--8496, 2017.

\newblock \doi{10.1090/tran/6905}.

\newblock URL \url{https://arxiv.org/pdf/1411.2742}.
\newblock \emph{Verified}.

\bibitem[Buchdahl(1999)]{Buchdahl1999}
Nicholas Buchdahl.

\newblock On compact kähler surfaces.

\newblock \emph{Annales de l'Institut Fourier}, 49\penalty0 (1):\penalty0
  287--302, 1999.

\newblock \doi{10.5802/aif.1674}.

\newblock URL \url{https://www.numdam.org/article/AIF_1999__49_1_287_0.pdf}.
\newblock \emph{Verified}.

\bibitem[Cossec et~al.(2024)Cossec, Dolgachev, and Liedtke]{CDL24}
Fran\c{c}ois Cossec, Igor Dolgachev, and Christian Liedtke.

\newblock Enriques surfaces {I}, 2024.

\newblock URL
  \url{https://sites.lsa.umich.edu/idolga/wp-content/uploads/sites/1334/2024/08/EnriquesOne.pdf}.
\newblock \emph{Verified}.

\bibitem[Cartan(1952)]{Cartan1952B_AB}
Henri Cartan.

\newblock Faisceaux analytiques sur les vari{\'e}t{\'e}s de {Stein}:
  d{\'e}monstration des th{\'e}or{\`e}mes fondamentaux.

\newblock \emph{S{\'e}minaire Henri Cartan}, 4:\penalty0 1--15, 1952.

\newblock URL \url{https://www.numdam.org/article/SHC_1951-1952__4__A19_0.pdf}.
\newblock \emph{Verified}.

\bibitem[Belabas et~al.(2019)Belabas, Bernardi, and
  Perrin-Riou]{CgzBelabas2019}
Karim Belabas, Dominique Bernardi, and Bernadette Perrin-Riou.

\newblock La constante de {Manin} et le degr\'e modulaire d'une courbe
  elliptique.

\newblock \emph{Publications math\'ematiques de Besan\c con}, 2019/2:\penalty0
  81--103, 2019.

\newblock \doi{10.5802/pmb.37}.

\newblock URL \url{https://www.numdam.org/item/10.5802/pmb.37.pdf}.
\newblock \emph{Verified}.

\bibitem[Cai et~al.(2014)Cai, Shu, and Tian]{CgzCai2014}
Li~Cai, Jie Shu, and Ye~Tian.

\newblock Explicit {Gross--Zagier} and {Waldspurger} formulae.

\newblock \emph{Algebra \& Number Theory}, 8\penalty0 (10):\penalty0
  2523--2572, 2014.

\newblock \doi{10.2140/ant.2014.8.2523}.

\newblock URL \url{https://msp.org/ant/2014/8-10/ant-v8-n10-p05-p.pdf}.
\newblock \emph{Verified}.

\bibitem[Stein(2007)]{CgzStein2007}
William Stein.

\newblock \emph{Modular Forms: A Computational Approach}, volume~79 of
  \emph{Graduate Studies in Mathematics}.

\newblock American Mathematical Society, 2007.

\newblock URL \url{https://wstein.org/books/modform/stein-modform.pdf}.
\newblock \emph{Verified}.

\bibitem[Yuan et~al.(2011)Yuan, Zhang, and Zhang]{CgzYZZ2011}
Xinyi Yuan, Shou-Wu Zhang, and Wei Zhang.

\newblock Gross--zagier formula on shimura curves.

\newblock URL
  \url{https://www.researchgate.net/profile/Xinyi-Yuan-11/publication/267551160_Gross-Zagier_Formula_On_Shimura_Curves/links/548e842c0cf225bf66a5ff13/Gross-Zagier-Formula-On-Shimura-Curves.pdf?origin=publication_detail}.
\newblock \emph{Verified}.

\bibitem[Chen(2026)]{Chen2026HodgeRoot}
Kuan-Wen Chen.

\newblock On the {Hodge} structures of global smoothings of normal crossing
  varieties, 2026.

\newblock URL \url{https://arxiv.org/abs/2404.19229v2}.
\newblock \emph{Verified}.

\bibitem[Bertolini et~al.(2012)Bertolini, Darmon, and Prasanna]{CrBDP2012}
Massimo Bertolini, Henri Darmon, and Kartik Prasanna.

\newblock {$p$}-adic {Rankin} {L}-series and rational points on {CM} elliptic
  curves.

\newblock \emph{Pacific Journal of Mathematics}, 260\penalty0 (2):\penalty0
  261--303, 2012.

\newblock \doi{10.2140/pjm.2012.260.261}.

\newblock URL
  \url{https://www.math.mcgill.ca/darmon/pub/Articles/Research/52.BDP2/bdp-pjm.pdf}.
\newblock \emph{Verified}.

\bibitem[Castella(2025)]{CrCastella2025}
Francesc Castella.

\newblock Tamagawa number conjecture for {CM} modular forms and
  {R}ankin--{S}elberg convolutions.

\newblock \emph{Proceedings of the London Mathematical Society}, 131\penalty0
  (4):\penalty0 e70089, 2025.

\newblock \doi{10.1112/plms.70089}.

\newblock URL \url{https://web.math.ucsb.edu/~castella/Katz.pdf}.
\newblock \emph{Verified}.

\bibitem[Cremona and Siksek(2006)]{CrCremonaSiksek2006}
John~E. Cremona and Samir Siksek.

\newblock Computing a lower bound for the canonical height on elliptic curves
  over {$\mathbb Q$}.

\newblock In \emph{Algorithmic Number Theory}, volume 4076 of \emph{Lecture
  Notes in Computer Science}, pages 275--286. Springer, 2006.

\newblock URL \url{https://johncremona.github.io/papers/HeightLB.pdf}.
\newblock \emph{Verified}.

\bibitem[Darmon(2009)]{CrDarmon2009}
Henri Darmon.

\newblock Cycles on modular varieties and rational points on elliptic curves,
  2009.

\newblock URL
  \url{https://www.math.mcgill.ca/darmon/pub/Articles/Expository/13.Oberwolfach/paper.pdf}.
\newblock \emph{Verified}.

\bibitem[Gross(1991)]{CrGross1991}
Benedict~H. Gross.

\newblock Kolyvagin's work on modular elliptic curves.

\newblock In \emph{L-functions and Arithmetic}, volume 153 of \emph{London
  Mathematical Society Lecture Note Series}, pages 235--256. Cambridge
  University Press, 1991.

\newblock URL
  \url{https://wstein.org/papers/bib/gross-kolyvagins_work_on_modular_elliptic_curves.pdf}.
\newblock \emph{Verified}.

\bibitem[Gross and Zagier(1986)]{CrGrossZagier1986}
Benedict~H. Gross and Don~B. Zagier.

\newblock Heegner points and derivatives of {L}-series.

\newblock \emph{Inventiones Mathematicae}, 84\penalty0 (2):\penalty0 225--320,
  1986.

\newblock \doi{10.1007/BF01388809}.

\newblock URL
  \url{https://archive.mpim-bonn.mpg.de/id/eprint/3578/1/preprint_1985_30.pdf}.
\newblock \emph{Verified}.

\bibitem[Jetchev et~al.(2008)Jetchev, Lauter, and Stein]{CrJetchev2008}
Dimitar Jetchev, Kristin Lauter, and William Stein.

\newblock Explicit {H}eegner points: {K}olyvagin's conjecture and non-trivial
  elements in the {S}hafarevich--{T}ate group.

\newblock \emph{Journal of Number Theory}, 128:\penalty0 3394--3410, 2008.

\newblock URL \url{https://arxiv.org/pdf/0707.0032}.
\newblock \emph{Verified}.

\bibitem[Rubin(1999)]{CrRubin1999}
Karl Rubin.

\newblock Elliptic curves with complex multiplication and the conjecture of
  {B}irch and {S}winnerton-{D}yer.

\newblock In \emph{Arithmetic Theory of Elliptic Curves}, volume 1716 of
  \emph{Lecture Notes in Mathematics}, pages 167--234. Springer, 1999.

\newblock URL \url{https://swc-math.github.io/aws/1999/99RubinCM.pdf}.
\newblock \emph{Verified}.

\bibitem[{National Institute of Standards and Technology}(2026)]{DLMFTheta2026}
{National Institute of Standards and Technology}.

\newblock {NIST Digital Library of Mathematical Functions}, chapters 20 and 23,
  2026.

\newblock URL \url{https://dlmf.nist.gov/}.
\newblock \emph{Verified}.

\bibitem[de~Cataldo et~al.(2016)de~Cataldo, Migliorini, and
  Musta\c{t}\u{a}]{DMM16}
Mark Andrea~A. de~Cataldo, Luca Migliorini, and Mircea Musta\c{t}\u{a}.

\newblock Combinatorics and topology of proper toric maps, 2016.

\newblock URL \url{https://arxiv.org/pdf/1407.3497}.
\newblock \emph{Verified}.

\bibitem[Demailly(2012)]{Demailly2012E}
Jean-Pierre Demailly.

\newblock Complex analytic and differential geometry, 2012.

\newblock URL
  \url{https://www-fourier.univ-grenoble-alpes.fr/~demailly/manuscripts/agbook.pdf}.
\newblock \emph{Verified}.

\bibitem[Dloussky(2005)]{Dloussky2005AB}
Georges Dloussky.

\newblock On surfaces of class {VII}$_0^+$ with numerically anticanonical
  divisor, 2005.

\newblock URL \url{https://arxiv.org/pdf/math/0406387}.
\newblock \emph{Verified}.

\bibitem[Dloussky and Teleman(2012)]{DlousskyTeleman2012AB}
Georges Dloussky and Andrei Teleman.

\newblock Infinite bubbling in non-{K}{\"a}hlerian geometry.

\newblock \emph{Mathematische Annalen}, 353\penalty0 (4):\penalty0 1283--1314,
  2012.

\newblock \doi{10.1007/s00208-011-0713-9}.

\newblock URL \url{https://arxiv.org/pdf/1012.5247}.
\newblock \emph{Verified}.

\bibitem[Dloussky and Teleman(2020)]{DlousskyTeleman2020D}
Georges Dloussky and Andrei Teleman.

\newblock Smooth deformations of singular contractions of class {VII} surfaces.

\newblock \emph{Mathematische Zeitschrift}, 296:\penalty0 1521--1537, 2020.

\newblock \doi{10.1007/s00209-020-02481-0}.

\newblock URL \url{https://arxiv.org/pdf/1803.07631}.
\newblock \emph{Verified}.

\bibitem[Doi(2009)]{Doi2009Root}
Mamoru Doi.

\newblock Gluing construction of compact complex surfaces with trivial
  canonical bundle.

\newblock \emph{Journal of the Mathematical Society of Japan}, 61\penalty0
  (3):\penalty0 853--884, 2009.

\newblock \doi{10.2969/jmsj/06130853}.

\newblock URL
  \url{https://www.jstage.jst.go.jp/article/jmath/61/3/61_3_853/_pdf/-char/en}.
\newblock \emph{Verified}.

\bibitem[Du et~al.(2026{\natexlab{a}})Du, Guo, Kusakabe, and Xie]{DuEtAl2026v1}
Yun-Heng Du, Bin Guo, Yuta Kusakabe, and Song-Yan Xie.

\newblock Oka geometry and deformations of torus-fibred complex six-spheres.

\newblock arXiv:2610.00250v1, superseded version, 2026{\natexlab{a}}.

\newblock URL \url{https://arxiv.org/abs/2610.00250v1}.
\newblock \emph{Verified}.

\bibitem[Du et~al.(2026{\natexlab{b}})Du, Guo, Kusakabe, and Xie]{DuEtAl2026v2}
Yun-Heng Du, Bin Guo, Yuta Kusakabe, and Song-Yan Xie.

\newblock Oka geometry of complex six-spheres.

\newblock arXiv:2610.00250v2, 2026{\natexlab{b}}.

\newblock URL \url{https://arxiv.org/abs/2610.00250v2}.
\newblock \emph{Verified}.

\bibitem[Du et~al.(2026{\natexlab{c}})Du, Guo, Kusakabe, and
  Xie]{DuGuoKusakabeXie2026v1}
Yun-Heng Du, Bin Guo, Yuta Kusakabe, and Song-Yan Xie.

\newblock Oka geometry and deformations of torus-fibred complex six-spheres,
  2026{\natexlab{c}}.

\newblock URL \url{https://arxiv.org/abs/2610.00250v1}.
\newblock \emph{Verified}.

\bibitem[Duke and Imamo{\u g}lu(2008)]{DukeImamoglu2008}
William Duke and {\"O}zlem Imamo{\u g}lu.

\newblock The zeros of the {Weierstrass} {$\wp$}-function and hypergeometric
  series.

\newblock \emph{Mathematische Annalen}, 340\penalty0 (4):\penalty0 897--905,
  2008.

\newblock \doi{10.1007/s00208-007-0174-3}.

\newblock URL \url{https://www.math.ucla.edu/~wdduke/preprints/}.
\newblock \emph{Verified}.

\bibitem[Engel et~al.(2025)Engel, de~Gaay~Fortman, and Schreieder]{EGFS25}
Philip Engel, Olivier de~Gaay~Fortman, and Stefan Schreieder.

\newblock Combinatorics and {H}odge theory of degenerations of abelian
  varieties: A survey of the {M}umford construction, 2025.

\newblock URL \url{https://arxiv.org/html/2507.15695v3}.
\newblock \emph{Verified}.

\bibitem[Enokizono and Hashizume(2024)]{EH2024E}
Makoto Enokizono and Kenta Hashizume.

\newblock Semistable reduction for complex analytic spaces, 2024.

\newblock URL \url{https://arxiv.org/pdf/2404.05124}.
\newblock \emph{Verified}.

\bibitem[Eichler and Zagier(1982)]{EichlerZagier1982}
Martin Eichler and Don Zagier.

\newblock On the zeros of the {Weierstrass} {$\wp$}-function.

\newblock \emph{Mathematische Annalen}, 258\penalty0 (4):\penalty0 399--407,
  1982.

\newblock \doi{10.1007/BF01453974}.

\newblock URL
  \url{https://people.mpim-bonn.mpg.de/zagier/files/doi/10.1007/BF01453974/fulltext.pdf}.
\newblock \emph{Verified}.

\bibitem[Engel(2026)]{Engel2026}
Philip Engel.

\newblock Complex structures on {$S^6$}.

\newblock arXiv:2609.38442v1, 2026.

\newblock URL \url{https://arxiv.org/abs/2609.38442}.
\newblock \emph{Verified}.

\bibitem[Erban and Van~Gorder(2021)]{ErbanVanGorder2021}
Radek Erban and Robert~A. Van~Gorder.

\newblock On the sum of positive divisors functions.

\newblock \emph{Research in Number Theory}, 7\penalty0 (25), 2021.

\newblock \doi{10.1007/s40993-021-00240-6}.

\newblock URL \url{https://people.maths.ox.ac.uk/erban/papers/paperRNT.pdf}.

\newblock Primary PDF opened and page1 visually read on 19 November 2026.
  Equations(1.1)--(1.4), p.1, state the Eisenstein series normalizations and
  Ramanujan differential system. In the receiving notation x=exp(2pi i tau),
  P=E2, Q=E4 and R=E6;\allowbreak{} the explicit derivative conversion is printed in PaperI.
  This is a bounded source verification of those equations.
\newblock \emph{Verified}.

\bibitem[Fantechi and G{\"o}ttsche(2010)]{FG10}
Barbara Fantechi and Lothar G{\"o}ttsche.

\newblock Riemann--roch theorems and elliptic genus for virtually smooth
  schemes.

\newblock \emph{Geometry \& Topology}, 14\penalty0 (1):\penalty0 83--115, 2010.

\newblock \doi{10.2140/gt.2010.14.83}.

\newblock URL \url{https://msp.org/gt/2010/14-1/gt-v14-n1-p02-s.pdf}.
\newblock \emph{Verified}.

\bibitem[Friedman(2016)]{Friedman2016AB}
Robert Friedman.

\newblock On the geometry of anticanonical pairs, 2016.

\newblock URL \url{https://arxiv.org/pdf/1502.02560}.
\newblock \emph{Verified}.

\bibitem[Fujiki(2009)]{Fujiki2009D}
Akira Fujiki.

\newblock Automorphisms of parabolic {Inoue} surfaces, 2009.

\newblock URL \url{https://arxiv.org/pdf/0903.5374}.
\newblock \emph{Verified}.

\bibitem[Fujiki and Pontecorvo(2010)]{FujikiPontecorvo2010AB}
Akira Fujiki and Massimiliano Pontecorvo.

\newblock Anti-self-dual bihermitian structures on {Inoue} surfaces.

\newblock \emph{Journal of Differential Geometry}, 85\penalty0 (1):\penalty0
  15--71, 2010.

\newblock \doi{10.4310/jdg/1284557925}.

\newblock URL
  \url{https://archive.ymsc.tsinghua.edu.cn/pacm_download/21/834-JDG-85-1-A2-fujiki.pdf}.
\newblock \emph{Verified}.

\bibitem[Fujimoto(1988)]{Fujimoto1988Root}
Yoshio Fujimoto.

\newblock Logarithmic transformations on elliptic fiber spaces.

\newblock \emph{Journal of Mathematics of Kyoto University}, 28\penalty0
  (1):\penalty0 91--110, 1988.

\newblock URL
  \url{https://scispace.com/pdf/logarithmic-transformations-on-elliptic-fiber-spaces-1uinrvp8jj.pdf}.
\newblock \emph{Verified}.

\bibitem[Fujimoto(1990)]{Fujimoto1990AB}
Yoshio Fujimoto.

\newblock Generalized logarithmic transformations---elliptic bundle case.

\newblock In K.~Konno, M.-H. Saito, and S.~Usui, editors, \emph{Proceedings of
  the Meeting and the Workshop ``Algebraic Geometry and Hodge Theory'', Vol.
  I}, volume~16 of \emph{Hokkaido University Technical Report Series in
  Mathematics}, pages 81--102. Hokkaido University, 1990.

\newblock \doi{10.14943/5135}.

\newblock URL
  \url{https://eprints.lib.hokudai.ac.jp/repo/huscap/all/5450/16.pdf}.
\newblock \emph{Verified}.

\bibitem[Gauduchon(1985)]{Gauduchon1985}
Paul Gauduchon.

\newblock Les m\'etriques standard d'une surface complexe compacte \`a premier
  nombre de {Betti} pair.

\newblock \emph{Ast\'erisque}, 126:\penalty0 129--135, 1985.

\newblock URL \url{https://numdam.org/item/AST_1985__126__129_0/}.
\newblock \emph{Verified}.

\bibitem[Graf and Kov{\'a}cs(2014)]{GrafKovacs2014Root}
Patrick Graf and S{\'a}ndor~J. Kov{\'a}cs.

\newblock Potentially {Du Bois} spaces.

\newblock \emph{Journal of Singularities}, 8:\penalty0 117--134, 2014.

\newblock \doi{10.5427/jsing.2014.8i}.

\newblock URL
  \url{https://sites.math.washington.edu/~kovacs/current/papers/Graf_Kovacs__Potentially_Du_Bois_spaces.pdf}.
\newblock \emph{Verified}.

\bibitem[Grauert(1960{\natexlab{a}})]{Grauert1960D}
Hans Grauert.

\newblock Ein theorem der analytischen garbentheorie und die modulräume
  komplexer strukturen.

\newblock \emph{Publications Mathématiques de l'IHÉS}, 5:\penalty0 5--64,
  1960{\natexlab{a}}.

\newblock \doi{10.1007/BF02684746}.

\newblock URL \url{https://www.numdam.org/article/PMIHES_1960__5__5_0.pdf}.

\newblock Reopened in this continuation on19 November2026: Section7 Satz1,
  printed p60 (Remmert proper-image theorem), with the reduced-space and
  proper-holomorphic hypotheses matched to the compact meromorphic graph.
  Verified in this writing on 18 November 2026 at these loci only. Section6
  HauptsatzeI,II,IIap59;\allowbreak{}Section7 Satz3p62,Satz5p64. Scope: Coherence, analytic
  formal functions, upper semicontinuity and constant-dimension base change,
  with flatness hypothesis retained Section 7, Satz 5, pp. 63–64. Scope:
  Opened full Satz 5. Also read corrigendum 1963 pp. 35–36
  (https://www.numdam.org/item/10.1007/BF02684291.pdf);\allowbreak{} it does not alter this
  statement. Author/title and edition matched to the opened item. Verification
  applies only to its recorded locator and edition;\allowbreak{} aliasing does not transfer
  preprint pagination to a published edition.
\newblock \emph{Verified}.

\bibitem[Grauert(1963)]{Grauert1963ErratumD}
Hans Grauert.

\newblock Berichtigung: Ein theorem der analytischen garbentheorie und die
  modulräume komplexer strukturen.

\newblock \emph{Publications Mathématiques de l'IHÉS}, 16:\penalty0 35--36,
  1963.

\newblock URL \url{https://www.numdam.org/item/PMIHES_1963__16__35_0.pdf}.
\newblock \emph{Verified}.

\bibitem[Grauert(1960{\natexlab{b}})]{Grauert60}
Hans Grauert.

\newblock Ein theorem der analytischen garbentheorie und die modulräume
  komplexer strukturen.

\newblock \emph{Publications Mathématiques de l'IHÉS}, 5:\penalty0 5--64,
  1960{\natexlab{b}}.

\newblock \doi{10.1007/BF02684746}.

\newblock URL \url{https://www.numdam.org/article/PMIHES_1960__5__5_0.pdf}.
\newblock \emph{Verified}.

\bibitem[Hochschild and Serre(1953)]{HS53}
Gerhard Hochschild and Jean-Pierre Serre.

\newblock Cohomology of group extensions.

\newblock \emph{Transactions of the American Mathematical Society}, 74\penalty0
  (1):\penalty0 110--134, 1953.

\newblock \doi{10.1090/S0002-9947-1953-0052438-8}.

\newblock URL
  \url{https://www.sas.rochester.edu/mth/sites/doug-ravenel/otherpapers/hochschild-serre.pdf}.
\newblock \emph{Verified}.

\bibitem[Honda and Viaclovsky(2025)]{HondaViaclovsky2025D}
Nobuhiro Honda and Jeff Viaclovsky.

\newblock Some remarks on fibrations in complex geometry.

\newblock \emph{Journal of Mathematical Study}, 58\penalty0 (4):\penalty0
  552--574, 2025.

\newblock \doi{10.4208/jms.v58n4.25.11}.

\newblock URL \url{https://arxiv.org/abs/2504.14420}.
\newblock \emph{Verified}.

\bibitem[Illusie et~al.(2005)Illusie, Kato, and Nakayama]{IKN2005E}
Luc Illusie, Kazuya Kato, and Chikara Nakayama.

\newblock Quasi-unipotent logarithmic {Riemann--Hilbert} correspondences.

\newblock \emph{Journal of Mathematical Sciences, the University of Tokyo},
  12\penalty0 (1):\penalty0 1--66, 2005.

\newblock URL \url{https://www.ms.u-tokyo.ac.jp/journal/pdf/jms120101.pdf}.
\newblock \emph{Verified}.

\bibitem[Illusie et~al.(2007)Illusie, Kato, and Nakayama]{IKN2007ErratumE}
Luc Illusie, Kazuya Kato, and Chikara Nakayama.

\newblock Erratum to ``quasi-unipotent logarithmic {Riemann--Hilbert}
  correspondences''.

\newblock \emph{Journal of Mathematical Sciences, the University of Tokyo},
  14\penalty0 (1):\penalty0 113--116, 2007.

\newblock URL \url{https://www.ms.u-tokyo.ac.jp/journal/pdf/jms140105.pdf}.
\newblock \emph{Verified}.

\bibitem[Illman(1972)]{Illman1972}
S{\"o}ren Illman.

\newblock \emph{Equivariant Algebraic Topology}.

\newblock PhD thesis, Princeton University, 1972.

\newblock URL
  \url{https://ncatlab.org/nlab/files/Illman_EquivariantAlgebraicTopology1972.pdf}.
\newblock \emph{Verified}.

\bibitem[Istrati et~al.(2022)Istrati, Otiman, Pontecorvo, and
  Ruggiero]{IstratiEtAl2022Root}
Nicolina Istrati, Alexandra Otiman, Massimiliano Pontecorvo, and Matteo
  Ruggiero.

\newblock Toric {Kato} manifolds.

\newblock \emph{Journal de l'{\'E}cole polytechnique---Math{\'e}matiques},
  9:\penalty0 1347--1395, 2022.

\newblock \doi{10.5802/jep.208}.

\newblock URL \url{https://www.numdam.org/item/JEP_2022__9__1347_0.pdf}.
\newblock \emph{Verified}.

\bibitem[Johansson(2016)]{Johansson2016}
Fredrik Johansson.

\newblock Arb: Efficient arbitrary-precision midpoint-radius interval
  arithmetic, 2016.

\newblock URL \url{https://arxiv.org/abs/1611.02831}.
\newblock \emph{Verified}.

\bibitem[Kajiwara et~al.(2008{\natexlab{a}})Kajiwara, Kato, and
  Nakayama]{KKN2008PicardE}
Takeshi Kajiwara, Kazuya Kato, and Chikara Nakayama.

\newblock Analytic log {Picard} varieties.

\newblock \emph{Nagoya Mathematical Journal}, 191:\penalty0 149--180,
  2008{\natexlab{a}}.

\newblock \doi{10.1017/S0027763000025940}.

\newblock URL
  \url{https://www.cambridge.org/core/services/aop-cambridge-core/content/view/E8B336E3252DE1F448583CC5865DB1DA/S0027763000025940a.pdf/analytic_log_picard_varieties.pdf}.
\newblock \emph{Verified}.

\bibitem[Kajiwara et~al.(2008{\natexlab{b}})Kajiwara, Kato, and
  Nakayama]{KKN2008Root}
Takeshi Kajiwara, Kazuya Kato, and Chikara Nakayama.

\newblock Logarithmic abelian varieties, {Part I}: Complex analytic theory.

\newblock \emph{Journal of Mathematical Sciences, the University of Tokyo},
  15\penalty0 (1):\penalty0 69--193, 2008{\natexlab{b}}.

\newblock URL \url{https://www.ms.u-tokyo.ac.jp/journal/pdf/jms150104.pdf}.
\newblock \emph{Verified}.

\bibitem[Kamenova(n.d.)]{KamenovaToriRoot}
Ljudmila Kamenova.

\newblock Degenerations of 2-dimensional tori.

\newblock Author-hosted manuscript, n.d.

\newblock URL
  \url{https://www.math.stonybrook.edu/~kamenova/homepage_files/degen.pdf}.
\newblock \emph{Verified}.

\bibitem[Katzarkov et~al.(2025)Katzarkov, Lee, Lupercio, and
  Meersseman]{KatzarkovEtAl2025Root}
Ludmil Katzarkov, Kyoung-Seog Lee, Ernesto Lupercio, and Laurent Meersseman.

\newblock On {Hodge} polynomials for nonalgebraic complex manifolds.

\newblock \emph{Proceedings of the National Academy of Sciences}, 122\penalty0
  (15):\penalty0 e2415722122, 2025.

\newblock \doi{10.1073/pnas.2415722122}.

\newblock URL \url{https://www.pnas.org/doi/10.1073/pnas.2415722122}.
\newblock \emph{Verified}.

\bibitem[Knudson(2004)]{Knudson2004}
Kevin~P. Knudson.

\newblock Invariant chains and the homology of quotient spaces, 2004.

\newblock URL \url{https://arxiv.org/pdf/math/0409133}.
\newblock \emph{Verified}.

\bibitem[Levy(2008)]{Levy08}
Roni~N. Levy.

\newblock Riemann--roch theorem for higher bivariant {K}-functors.

\newblock \emph{Annales de l'Institut Fourier}, 58\penalty0 (2):\penalty0
  571--601, 2008.

\newblock \doi{10.5802/aif.2361}.

\newblock URL \url{https://www.numdam.org/article/AIF_2008__58_2_571_0.pdf}.
\newblock \emph{Verified}.

\bibitem[\L{}ojasiewicz(1964)]{Loj64}
Stanis{\l}aw \L{}ojasiewicz.

\newblock Triangulation of semi-analytic sets.

\newblock \emph{Annali della Scuola Normale Superiore di Pisa, Classe di
  Scienze}, 18\penalty0 (4):\penalty0 449--474, 1964.

\newblock URL \url{https://www.numdam.org/item/ASNSP_1964_3_18_4_449_0.pdf}.
\newblock \emph{Verified}.

\bibitem[Manetti(1991)]{Manetti1991D}
Marco Manetti.

\newblock Normal degenerations of the complex projective plane.

\newblock \emph{Journal f{\"u}r die reine und angewandte Mathematik},
  419:\penalty0 89--118, 1991.

\newblock \doi{10.1515/crll.1991.419.89}.

\newblock URL \url{https://gdz.sub.uni-goettingen.de/id/PPN243919689_0419}.
\newblock \emph{Verified}.

\bibitem[Matsuki(2000)]{Matsuki00ErrAB}
Kenji Matsuki.

\newblock Erratum to the paper ``a note on the factorization theorem of toric
  birational maps after {Morelli} and its toroidal extension'', 2000.

\newblock URL \url{https://arxiv.org/abs/math/0006038}.
\newblock \emph{Verified}.

\bibitem[McHugh(2018)]{McHugh2018}
Andrew McHugh.

\newblock Bott--chern--aeppli, dolbeault and fr{\"o}licher on compact complex
  3-folds.

\newblock arXiv:1708.03251v4, 2018.

\newblock URL \url{https://arxiv.org/abs/1708.03251v4}.
\newblock \emph{Verified}.

\bibitem[McMullen(1996)]{McMullen1996Complex}
Curtis~T. McMullen.

\newblock Complex manifolds: Math 241 course outline.

\newblock University of California, Berkeley, course notes, 1996.

\newblock URL
  \url{https://people.math.harvard.edu/~ctm/home/text/class/berkeley/241/96/course/course.pdf}.
\newblock \emph{Verified}.

\bibitem[Meersseman(2011)]{Meersseman2011Root}
Laurent Meersseman.

\newblock Foliated structure of the {Kuranishi} space and isomorphisms of
  deformation families of compact complex manifolds.

\newblock \emph{Annales scientifiques de l'{\'E}cole Normale Sup{\'e}rieure},
  44\penalty0 (3):\penalty0 495--525, 2011.

\newblock \doi{10.24033/asens.2148}.

\newblock URL \url{https://www.numdam.org/item/10.24033/asens.2148.pdf}.
\newblock \emph{Verified}.

\bibitem[Merkulov(1995)]{Merkulov1995D}
Sergey~A. Merkulov.

\newblock Moduli spaces of compact complex submanifolds of complex fibered
  manifolds.

\newblock \emph{Mathematical Proceedings of the Cambridge Philosophical
  Society}, 118\penalty0 (1):\penalty0 71--91, 1995.

\newblock \doi{10.1017/S0305004100073473}.

\newblock URL \url{https://doi.org/10.1017/S0305004100073473}.
\newblock \emph{Verified}.

\bibitem[Movasati(2024)]{Movasati24}
Hossein Movasati.

\newblock A course in complex geometry and holomorphic foliations.

\newblock Author version dated 13 May 2024, 2024.

\newblock URL \url{https://w3.impa.br/~hossein/myarticles/ComplexGeometry.pdf}.
\newblock \emph{Verified}.

\bibitem[Nakayama and Ogus(2010)]{NO10}
Chikara Nakayama and Arthur Ogus.

\newblock Relative rounding in toric and logarithmic geometry.

\newblock \emph{Geometry \& Topology}, 14\penalty0 (4):\penalty0 2189--2241,
  2010.

\newblock \doi{10.2140/gt.2010.14.2189}.

\newblock URL \url{https://msp.org/gt/2010/14-4/gt-v14-n4-p09-s.pdf}.
\newblock \emph{Verified}.

\bibitem[Nakayama(2002)]{Nakayama2002Root}
Noboru Nakayama.

\newblock Global structure of an elliptic fibration.

\newblock \emph{Publications of the Research Institute for Mathematical
  Sciences}, 38\penalty0 (3):\penalty0 451--649, 2002.

\newblock \doi{10.2977/PRIMS/1145476270}.

\newblock URL \url{https://ems.press/content/serial-article-files/40808?nt=1}.
\newblock \emph{Verified}.

\bibitem[Nishiguchi(1988)]{Nishiguchi1988IIRoot}
Kenji Nishiguchi.

\newblock Degeneration of {K3} surfaces, {II}.

\newblock Max-Planck-Institut f{\"u}r Mathematik preprint MPI/88-48, 1988.

\newblock URL
  \url{https://archive.mpim-bonn.mpg.de/id/eprint/3390/1/preprint_1988_48.pdf}.
\newblock \emph{Verified}.

\bibitem[Oda(1978)]{Oda1978D}
Tadao Oda.

\newblock \emph{Lectures on Torus Embeddings and Applications}.

\newblock Number~58 in Tata Institute of Fundamental Research Lectures on
  Mathematics and Physics. Springer-Verlag, Berlin, Heidelberg, New York, 1978.

\newblock URL
  \url{https://mathweb.tifr.res.in/Documents/Publications/Lectures/tifr58.pdf}.
\newblock \emph{Verified}.

\bibitem[Parusi{\'n}ski(1990)]{Parusinski1990E}
Adam Parusi{\'n}ski.

\newblock Subanalytic functions.

\newblock Technical Report MPI/90-58, Max-Planck-Institut f{\"u}r Mathematik,
  1990.

\newblock URL
  \url{https://archive.mpim-bonn.mpg.de/id/eprint/1622/1/preprint_1990_58.pdf}.
\newblock \emph{Verified}.

\bibitem[Pinkham(1974)]{Pinkham1974Root}
Henry~C. Pinkham.

\newblock \emph{Deformations of algebraic varieties with {$G_m$} action}.

\newblock Number~20 in Ast\'erisque. Soci\'et\'e math\'ematique de France,
  1974.

\newblock URL \url{https://www.numdam.org/item/AST_1974__20__1_0/}.
\newblock \emph{Verified}.

\bibitem[{python-flint developers}(2026)]{PythonFlint2026}
{python-flint developers}.

\newblock Python-flint documentation: Arb real ball arithmetic, 2026.

\newblock URL \url{https://python-flint.readthedocs.io/en/latest/arb.html}.
\newblock \emph{Verified}.

\bibitem[Schr{\"o}er and Siebert(2002{\natexlab{a}})]{SS00}
Stefan Schr{\"o}er and Bernd Siebert.

\newblock Irreducible degenerations of primary {K}odaira surfaces.

\newblock In \emph{Complex geometry (G{\"o}ttingen, 2000)}, pages 193--222.
  Springer, Berlin, 2002{\natexlab{a}}.

\newblock URL
  \url{https://www.math.uni-duesseldorf.de/~schroeer/publications_pdf/kodaira.pdf}.
\newblock \emph{Verified}.

\bibitem[Sz\'ekelyhidi et~al.(2017)Sz\'ekelyhidi, Tosatti, and
  Weinkove]{STW2017}
G\'abor Sz\'ekelyhidi, Valentino Tosatti, and Ben Weinkove.

\newblock Gauduchon metrics with prescribed volume form, 2017.

\newblock URL \url{https://arxiv.org/abs/1503.04491}.
\newblock \emph{Verified}.

\bibitem[Schnell(2019)]{Schnell19}
Christian Schnell.

\newblock Lecture 4: The {Appel--Humbert} theorem.

\newblock MAT 615, Stony Brook University, lecture of 6 February, 2019.

\newblock URL
  \url{https://www.math.stonybrook.edu/~cschnell/mat615/lectures/lecture4.pdf}.
\newblock \emph{Verified}.

\bibitem[Schr{\"o}er and Siebert(2002{\natexlab{b}})]{SchroeerSiebert2002Root}
Stefan Schr{\"o}er and Bernd Siebert.

\newblock Irreducible degenerations of primary {Kodaira} surfaces.

\newblock In \emph{Complex geometry ({G}{\"o}ttingen, 2000)}, pages 193--222.
  Springer, Berlin, 2002{\natexlab{b}}.

\newblock URL
  \url{https://www.math.uni-duesseldorf.de/~schroeer/publications_pdf/kodaira.pdf}.
\newblock \emph{Verified}.

\bibitem[Shin(2015)]{Shin2015}
Dong~Hwa Shin.

\newblock Generation of class fields by {Siegel--Ramachandra} invariants.

\newblock \emph{Journal of the Korean Mathematical Society}, 52\penalty0
  (5):\penalty0 907--928, 2015.

\newblock URL \url{https://arxiv.org/pdf/1009.2253}.
\newblock \emph{Verified}.

\bibitem[Shioda(1990)]{Shioda90}
Tetsuji Shioda.

\newblock On the {M}ordell--{W}eil lattices.

\newblock \emph{Commentarii Mathematici Universitatis Sancti Pauli},
  39\penalty0 (2):\penalty0 211--240, 1990.

\newblock \doi{10.14992/00009994}.

\newblock URL
  \url{https://www.researchgate.net/profile/Tetsuji-Shioda/publication/247043151_On_the_Mordell-Weil_Lattices/links/5704b7fe08ae13eb88b69184/On-the-Mordell-Weil-Lattices.pdf}.
\newblock \emph{Verified}.

\bibitem[Stelzig(2021)]{Stelzig2021}
Jonas Stelzig.

\newblock On the structure of double complexes.

\newblock \emph{Journal of the London Mathematical Society}, 104\penalty0
  (2):\penalty0 956--988, 2021.

\newblock \doi{10.1112/jlms.12453}.
\newblock \emph{Verified}.

\bibitem[Sutherland(2023)]{Sutherland2023CM}
Andrew~V. Sutherland.

\newblock The {Hilbert} class polynomial.

\newblock MIT 18.783 Elliptic Curves, Lecture 20, November 2023.

\newblock URL
  \url{https://math.mit.edu/classes/18.783/2023/LectureNotes20.pdf}.
\newblock \emph{Verified}.

\bibitem[Teleman(2010)]{Teleman2010AB}
Andrei Teleman.

\newblock Instantons and curves on class {VII} surfaces.

\newblock \emph{Annals of Mathematics}, 172\penalty0 (3):\penalty0 1749--1804,
  2010.

\newblock \doi{10.4007/annals.2010.172.1749}.

\newblock URL
  \url{https://annals.math.princeton.edu/wp-content/uploads/annals-v172-n3-p07-p.pdf}.
\newblock \emph{Verified}.

\bibitem[Viaclovsky(2026{\natexlab{a}})]{ViaclovskyLog2026}
Jeff Viaclovsky.

\newblock Logarithmic transformations on the six-sphere.

\newblock arXiv:2610.39102v1, 2026{\natexlab{a}}.

\newblock URL \url{https://arxiv.org/abs/2610.39102}.
\newblock \emph{Verified}.

\bibitem[Viaclovsky(2026{\natexlab{b}})]{ViaclovskyTwo2026}
Jeff Viaclovsky.

\newblock A two-parameter family of complex structures on {$S^6$}.

\newblock arXiv:2609.33785v1, 2026{\natexlab{b}}.

\newblock URL \url{https://arxiv.org/abs/2609.33785}.
\newblock \emph{Verified}.

\bibitem[Xiao(1991)]{Xiao91}
Gang Xiao.

\newblock Elliptic surfaces with fixed jacobian.

\newblock Technical Report MPI/91-30, Max-Planck-Institut f{\"u}r Mathematik,
  1991.

\newblock URL
  \url{https://archive.mpim-bonn.mpg.de/id/eprint/2953/1/preprint_1991_30.pdf}.
\newblock \emph{Verified}.

\bibitem[Yang(2026)]{Yang2026}
Xiufan Yang.

\newblock Compact complex threefolds fibred by two-dimensional complex tori.

\newblock arXiv:2609.29666v1, 2026.

\newblock URL \url{https://arxiv.org/abs/2609.29666}.
\newblock \emph{Verified}.

\bibitem[Adamczewski and Faverjon(2017)]{adamczewski2017mahler}
Boris Adamczewski and Colin Faverjon.

\newblock M\'ethode de {Mahler}: relations lin\'eaires, transcendance et
  applications aux nombres automatiques.

\newblock \emph{Proceedings of the London Mathematical Society}, 115:\penalty0
  55--90, 2017.

\newblock \doi{10.1112/plms.12038}.

\newblock URL
  \url{https://faverjon.perso.math.cnrs.fr/AdamczewskiFaverjon_mahler1.pdf}.
\newblock \emph{Verified}.

\bibitem[Agboola and Howard(2006)]{agboola2006}
Adebisi Agboola and Benjamin Howard.

\newblock Anticyclotomic {Iwasawa} theory of {CM} elliptic curves.

\newblock \emph{Annales de l'Institut Fourier}, 56\penalty0 (4):\penalty0
  1001--1048, 2006.

\newblock URL
  \url{https://web.math.ucsb.edu/~agboola/papers/bisi_aif_2006.pdf}.
\newblock \emph{Verified}.

\bibitem[Bertolin et~al.(2022)Bertolin, Philippon, Saha, and
  Saha]{bertolin2022}
Cristiana Bertolin, Patrice Philippon, Biswajyoti Saha, and Ekata Saha.

\newblock Semi-abelian analogues of {Schanuel} conjecture and applications.

\newblock \emph{Journal of Algebra}, 596:\penalty0 250--288, 2022.

\newblock URL \url{https://arxiv.org/pdf/2010.15170v2}.
\newblock \emph{Verified}.

\bibitem[Bertrand(1995)]{bertrand1995}
Daniel Bertrand.

\newblock Points rationnels sur les sous-groupes compacts des groupes
  alg{\'e}briques.

\newblock \emph{Experimental Mathematics}, 4\penalty0 (2):\penalty0 145--151,
  1995.

\newblock \doi{10.1080/10586458.1995.10504315}.

\newblock URL
  \url{https://www.maths.tcd.ie/EMIS/journals/EM/restricted/4/4.2/bertrand.ps}.
\newblock \emph{Verified}.

\bibitem[Bertrand(2017)]{bertrand2017}
Daniel Bertrand.

\newblock Transcendence problems related to heights.

\newblock Fields Institute, Heights and Applications to Unlikely Intersections,
  February 13--17, 2017.

\newblock URL
  \url{https://webusers.imj-prg.fr/~daniel.bertrand/Recherche/rpdf/DB_Fieldspost.pdf}.
\newblock \emph{Verified}.

\bibitem[Bost and K{\"u}nnemann(2009)]{bost2009}
Jean-Beno{\^i}t Bost and Klaus K{\"u}nnemann.

\newblock Hermitian vector bundles and extension groups on arithmetic schemes
  {II}. {The} arithmetic {Atiyah} extension.

\newblock \emph{Ast{\'e}risque}, 327:\penalty0 361--424, 2009.

\newblock URL \url{https://www.numdam.org/article/AST_2009__327__361_0.pdf}.
\newblock \emph{Verified}.

\bibitem[Call and Silverman(1993)]{callSilverman1993}
Gregory~S. Call and Joseph~H. Silverman.

\newblock Canonical heights on varieties with morphisms.

\newblock \emph{Compositio Mathematica}, 89\penalty0 (2):\penalty0 163--205,
  1993.

\newblock URL \url{https://www.numdam.org/item/CM_1993__89_2_163_0.pdf}.
\newblock \emph{Verified}.

\bibitem[DeMarco and Mavraki(2020)]{demarco2017}
Laura DeMarco and Niki~Myrto Mavraki.

\newblock Variation of canonical height and equidistribution.

\newblock \emph{American Journal of Mathematics}, 142\penalty0 (2):\penalty0
  443--473, 2020.

\newblock URL \url{https://arxiv.org/pdf/1701.07947}.
\newblock \emph{Verified}.

\bibitem[Einsiedler et~al.(2001)Einsiedler, Everest, and
  Ward]{einsiedler2001entropy}
Manfred Einsiedler, Graham Everest, and Thomas Ward.

\newblock Entropy and the canonical height.

\newblock \emph{Journal of Number Theory}, 91\penalty0 (2):\penalty0 256--273,
  2001.

\newblock \doi{10.1006/jnth.2001.2682}.

\newblock URL \url{https://people.math.ethz.ch/~einsiedl/height.pdf}.
\newblock \emph{Verified}.

\bibitem[Gasbarri(2010)]{gasbarri2010analytic}
Carlo Gasbarri.

\newblock Analytic subvarieties with many rational points.

\newblock \emph{Mathematische Annalen}, 346\penalty0 (1):\penalty0 199--243,
  2010.

\newblock \doi{10.1007/s00208-009-0394-9}.

\newblock URL \url{https://preprints.ihes.fr/storage/newscheiderlang.pdf}.
\newblock \emph{Verified}.

\bibitem[Gramain et~al.(1986)Gramain, Mignotte, and
  Waldschmidt]{gramain1986values}
Fran\c{c}ois Gramain, Maurice Mignotte, and Michel Waldschmidt.

\newblock Valeurs alg\'ebriques de fonctions analytiques.

\newblock \emph{Acta Arithmetica}, 47\penalty0 (2):\penalty0 97--121, 1986.

\newblock \doi{10.4064/aa-47-2-97-121}.

\newblock URL \url{https://matwbn.icm.edu.pl/ksiazki/aa/aa47/aa4721.pdf}.
\newblock \emph{Verified}.

\bibitem[Herblot(2012)]{herblot2012meromorphic}
Mathilde Herblot.

\newblock Algebraic points on meromorphic curves, 2012.

\newblock URL \url{https://webusers.imj-prg.fr/~mathilde.herblot/publi/SL.pdf}.
\newblock \emph{Verified}.

\bibitem[Huber and W{\"u}stholz(2023)]{huber2022}
Annette Huber and Gisbert W{\"u}stholz.

\newblock \emph{Transcendence and linear relations of 1-periods}, volume 284 of
  \emph{Memoirs of the American Mathematical Society}.

\newblock American Mathematical Society, 2023.

\newblock URL \url{https://arxiv.org/abs/1805.10104}.
\newblock \emph{Verified}.

\bibitem[Miller(2011)]{miller2011}
Robert~L. Miller.

\newblock Proving the {Birch} and {Swinnerton-Dyer} conjecture for specific
  elliptic curves of analytic rank zero and one.

\newblock \emph{LMS Journal of Computation and Mathematics}, 14:\penalty0
  327--350, 2011.

\newblock \doi{10.1112/S1461157011000180}.

\newblock URL \url{https://arxiv.org/pdf/1010.2431v3}.
\newblock \emph{Verified}.

\bibitem[Nguyen(2023)]{nguyen2023}
Khoa~D. Nguyen.

\newblock Transcendence of polynomial canonical heights.

\newblock \emph{Mathematische Annalen}, 387:\penalty0 1--15, 2023.

\newblock URL \url{https://arxiv.org/pdf/2112.14937}.
\newblock \emph{Verified}.

\bibitem[Noguchi(2024)]{noguchi2024}
Junjiro Noguchi.

\newblock Analytic {Ax--Schanuel} for semi-abelian varieties and {Nevanlinna}
  theory.

\newblock \emph{Journal of the Mathematical Society of Japan}, 76\penalty0
  (1):\penalty0 1--22, 2024.

\newblock \doi{10.2969/jmsj/89588958}.

\newblock URL \url{https://www.ms.u-tokyo.ac.jp/~noguchi/An-axsch.pdf}.
\newblock \emph{Verified}.

\bibitem[Reyssat(1982)]{reyssat1982weierstrass}
Eric Reyssat.

\newblock Propri\'et\'es d'ind\'ependance alg\'ebrique de nombres li\'es aux
  fonctions de weierstrass.

\newblock \emph{Acta Arithmetica}, 41\penalty0 (3):\penalty0 291--310, 1982.

\newblock \doi{10.4064/aa-41-3-291-310}.

\newblock URL \url{https://matwbn.icm.edu.pl/ksiazki/aa/aa41/aa4137.pdf}.
\newblock \emph{Verified}.

\bibitem[Roy and Waldschmidt(1993)]{roy1993periodrank}
Damien Roy and Michel Waldschmidt.

\newblock Autour du th\'eor\`eme du sous-groupe alg\'ebrique.

\newblock \emph{Canadian Mathematical Bulletin}, 36\penalty0 (3):\penalty0
  358--367, 1993.

\newblock \doi{10.4153/CMB-1993-049-8}.

\newblock URL
  \url{https://www.cambridge.org/core/journals/canadian-mathematical-bulletin/article/autour-du-theoreme-du-sousgroupe-algebrique/A1E35F3798690D084664ACD83D1557E8}.
\newblock \emph{Verified}.

\bibitem[Silverman(2006)]{silverman2006heights}
Joseph~H. Silverman.

\newblock An introduction to height functions.

\newblock MSRI workshop lecture slides, January 2006, 2006.

\newblock URL
  \url{https://slmath-production-assets.sfo3.digitaloceanspaces.com/attachments/workshops/301/HtSurveyMSRIJan06.pdf}.
\newblock \emph{Verified}.

\bibitem[Waldschmidt(1975)]{waldschmidt1975dimension}
Michel Waldschmidt.

\newblock Dimension alg{\'e}brique de sous-groupes analytiques de
  vari{\'e}t{\'e}s de groupe.

\newblock \emph{Annales de l'Institut Fourier}, 25\penalty0 (1):\penalty0
  23--33, 1975.

\newblock \doi{10.5802/aif.543}.

\newblock URL
  \url{https://webusers.imj-prg.fr/~michel.waldschmidt/articles/pdf/Fourier.pdf}.
\newblock \emph{Verified}.

\bibitem[Waldschmidt(1979)]{waldschmidt1979sigma}
Michel Waldschmidt.

\newblock Nombres transcendants et fonctions sigma de {Weierstrass}.

\newblock \emph{Comptes Rendus Mathematiques de l'Academie des Sciences,
  Canada}, 1\penalty0 (2):\penalty0 111--114, 1979.

\newblock URL
  \url{https://webusers.imj-prg.fr/~michel.waldschmidt/articles/pdf/GDL/SigmaNoteCanada1979-eng.pdf}.
\newblock \emph{Verified}.

\bibitem[Waldschmidt(1988{\natexlab{a}})]{waldschmidt1988periodrank}
Michel Waldschmidt.

\newblock On the transcendence methods of {Gelfond} and {Schneider} in several
  variables.

\newblock In Alan Baker, editor, \emph{New Advances in Transcendence Theory},
  chapter~24, pages 375--398. Cambridge University Press, 1988{\natexlab{a}}.

\newblock URL
  \url{https://webusers.imj-prg.fr/~michel.waldschmidt/articles/pdf/Durham1986TMGSSV.pdf}.
\newblock \emph{Verified}.

\bibitem[Waldschmidt(1988{\natexlab{b}})]{waldschmidt1988theta}
Michel Waldschmidt.

\newblock Fonctions th\^eta et nombres transcendants.

\newblock S\'eminaire de Th\'eorie des Nombres de Bordeaux, 1987--1988,
  expos\'e 15, pp. 1--7, 1988{\natexlab{b}}.

\newblock URL
  \url{https://webusers.imj-prg.fr/~michel.waldschmidt/articles/pdf/PPN320141322_0016__0_0.pdf}.
\newblock \emph{Verified}.

\bibitem[Waldschmidt(2009)]{waldschmidt2009auxiliary}
Michel Waldschmidt.

\newblock Auxiliary functions in transcendental number theory.

\newblock \emph{The Ramanujan Journal}, 20\penalty0 (3):\penalty0 341--373,
  2009.

\newblock \doi{10.1007/s11139-009-9204-y}.

\newblock URL
  \url{https://webusers.imj-prg.fr/~michel.waldschmidt/articles/pdf/SASTRAaftp.pdf}.
\newblock \emph{Verified}.

\bibitem[Xie(2024)]{xie2024algebraicity}
Junyi Xie.

\newblock Algebraicity criteria, invariant subvarieties and transcendence
  problems from arithmetic dynamics.

\newblock \emph{Peking Mathematical Journal}, 7:\penalty0 345--398, 2024.

\newblock \doi{10.1007/s42543-022-00059-9}.

\newblock URL \url{https://arxiv.org/pdf/2202.09032}.
\newblock \emph{Verified}.

\end{thebibliography}

\end{document}
